\documentclass[11pt,left=2cm,bottom=2cm,oneside]{book}
%\usepackage[text={7in,10in},centering]{geometry}
\usepackage{booktabs}
\usepackage{graphicx}
\usepackage{ifpdf}
\usepackage{xcolor}
\usepackage{titlesec}
\usepackage{fancyhdr}
\usepackage{lmodern}
\usepackage[version=latest]{pgf}
\usepackage{tikz}
\usepackage{hyperref}
%\let\svtikzpic\tikzpicture
%\def\tikzpicture{\noindent\svtikzpic\raggedright}
\hypersetup{colorlinks,linkcolor={black}, citecolor={black},urlcolor={darkblue}}
%\vspace{6cm}
\newcommand{\mquot}[1]{\begin{quotation}\textit{#1}\end{quotation}}
\begin{document}
\pagenumbering{}
\begin{titlepage}
\nopagebreak
\centering
{\LARGE A Dynamic Stateful Multicast Firewall}\\
\vspace*{2cm}
{\Large Author: Bewketu Tadilo}\\
\vspace*{2cm}
{\Large Date: Ramadan 1438}\\
\vspace{2cm}
{\hspace{-3cm}\Large stuff to:}\\
\vspace{6cm}
\includegraphics[height=5cm]{unsw-logo.jpg}
\end{titlepage}
\clearpage

 \section*{\centering Abstract}
\begin{quotation}
Security has been one major hindrance to the spread of multicasting and its use. Network managers are still rightly paranoid
enabling the technology when a whole array of internal resources
can be at risk. Efforts have been made to alleviate some of the
concern in the form modified service model namely source specific
multicasting. The attractiveness of the original flat service model
still remains at large and so do the multitude of routers and hosts
that only support this traditional version. The newer protocols
have been made backward compatible for this reason. This thesis aims to create a security layer that can sit in between these
two models. It draws upon experiences learned from the design
of stateful unicast firewalls and the newer protocols to track connections and associations to form its rules of a stateful multicast
firewall. It enables both coexistence and smoother transition of
these two service models.
\end{quotation}
\clearpage
\vspace*{8cm}
\subsubsection*{\underline{Originality Statement}}
I hereby declare that this submission is the work of my own. This work has not been submitted previously for any other degree and to the best of my knowledge, does not contain materials previously published or written by another person except where due acknowledgment is made.
\subsection*{Signed:\underline{~~~~~~~~~~~~~~~~~~}}
\clearpage
\vspace*{5pt}
\begin{minipage}{15cm}
\raggedright
\begin{quote}
\textit{
Dedicated to my European and Australian friends whom I met in Gondar,
Ethiopia. Thank you for taking an extra burden to sponsor my studies in Australia
and your collective support. May God bless you all!}\\
\raggedleft
\-\-Bewketu Tadilo, June 2009.
\end{quote}
\end{minipage}
\clearpage
\tableofcontents{}
%\Abbreviations{}
\listoffigures{}
\listoftables{}
\chapter{Introduction}
\pagenumbering{arabic}
Multicasting has been supported by IP networks for two decades since Steve
Deering did his PhD on it, and wrote the Internet Engineering Task Force (IETF)standard
at the end of 1989\cite{Deering}. But, the technology have been little used up until recently.
As can be inferred from its name, multicasting is a technique for transmitting data
from a single source to only a set of recipeints. The very idea brings concepts that
are not encountered in broadcasting, in which a message is sent to all computers on
the network, and unicasting in which a message is passed between a single source
and receiver. One of them is how we manage these set of recipients and is discussed
in section 2.2 on page 6.
Advances in Audiovisual technologies, that naturally demand high bandwidth,
have spurred the recent interest in multicasting in Small and Medium Enterprises
(SMEs) and consumer services. This is because we can significantly reduce traffic
load if we use multicast efficiently as we will see in section 1.1 on the next page. Nevertheless, it has remained elusive to many Network Managers, Software Developers,
System Admins and end users who might otherwise benefit from this technology.
Multicasting is heavily used in the stoke market where its model of updating all
interested users at the same time fits in perfectly. The defence departments of
countries are also big investors to the advancement the technology for better coordination in battlefields. These are areas among others in which multicasting can/is
being used
• Clustering
– load sharing
– database replication
– failover clusters
• IPTV and live Internet Audio (radio)
• Interactive gaming
• Software distribution.
It even comes in very handy, if a Network Manager wants to install one program
to set of computers on his/her network, including operating systems.
We focus this thesis with medium enterprises in our mind and from a perspective of the Network Manager for that Enterprise. We are on Layer 3 (at the network
layer) in the entire thesis.
Multicasting reduces to either unicasting and broadcasting depending on whether
one or all members of the network are participating in communication respectively.
The use of broadcasting in many cases has been discouraged in favour of multicasting as it is prone to wasting huge resource.
\section{WHY MULTICATING?}
Figure 1.1. Data delivery methods
\begin{flushleft}% Graphic for TeX using PGF
% Title: /home/bewketu/diag1.dia
% Creator: Dia v0.97
% CreationDate: Thu Aug  3 03:29:36 2017
% For: bewketu
% \usepackage{tikz}
% The following commands are not supported in PSTricks at present
% We define them conditionally, so when they are implemented,
% this pgf file will use them.
\ifx\du\undefined
  \newlength{\du}
\fi
\setlength{\du}{15\unitlength}
\raggedright\begin{tikzpicture}
\pgftransformxscale{1.000000}
\pgftransformyscale{-1.000000}
\definecolor{dialinecolor}{rgb}{0.000000, 0.000000, 0.000000}
\pgfsetstrokecolor{dialinecolor}
\definecolor{dialinecolor}{rgb}{1.000000, 1.000000, 1.000000}
\pgfsetfillcolor{dialinecolor}
\pgfsetlinewidth{0.100000\du}
\pgfsetdash{}{0pt}
\pgfsetdash{}{0pt}
\pgfsetbuttcap
\pgfsetmiterjoin
\pgfsetlinewidth{0.001000\du}
\pgfsetbuttcap
\pgfsetmiterjoin
\pgfsetdash{}{0pt}
\definecolor{dialinecolor}{rgb}{0.000000, 0.470588, 0.666667}
\pgfsetfillcolor{dialinecolor}
\pgfpathmoveto{\pgfpoint{-18.257654\du}{0.403708\du}}
\pgfpathlineto{\pgfpoint{-18.259158\du}{0.387917\du}}
\pgfpathlineto{\pgfpoint{-18.262416\du}{0.372125\du}}
\pgfpathlineto{\pgfpoint{-18.267931\du}{0.356584\du}}
\pgfpathlineto{\pgfpoint{-18.275701\du}{0.341294\du}}
\pgfpathlineto{\pgfpoint{-18.286229\du}{0.325753\du}}
\pgfpathlineto{\pgfpoint{-18.298511\du}{0.310714\du}}
\pgfpathlineto{\pgfpoint{-18.313300\du}{0.295925\du}}
\pgfpathlineto{\pgfpoint{-18.329844\du}{0.281387\du}}
\pgfpathlineto{\pgfpoint{-18.348894\du}{0.267099\du}}
\pgfpathlineto{\pgfpoint{-18.369949\du}{0.253313\du}}
\pgfpathlineto{\pgfpoint{-18.392759\du}{0.240028\du}}
\pgfpathlineto{\pgfpoint{-18.417574\du}{0.226492\du}}
\pgfpathlineto{\pgfpoint{-18.444395\du}{0.213709\du}}
\pgfpathlineto{\pgfpoint{-18.472970\du}{0.201677\du}}
\pgfpathlineto{\pgfpoint{-18.503551\du}{0.190147\du}}
\pgfpathlineto{\pgfpoint{-18.535635\du}{0.179118\du}}
\pgfpathlineto{\pgfpoint{-18.568973\du}{0.168339\du}}
\pgfpathlineto{\pgfpoint{-18.604065\du}{0.158564\du}}
\pgfpathlineto{\pgfpoint{-18.640410\du}{0.149289\du}}
\pgfpathlineto{\pgfpoint{-18.678260\du}{0.140516\du}}
\pgfpathlineto{\pgfpoint{-18.717363\du}{0.132495\du}}
\pgfpathlineto{\pgfpoint{-18.757719\du}{0.125727\du}}
\pgfpathlineto{\pgfpoint{-18.798827\du}{0.118709\du}}
\pgfpathlineto{\pgfpoint{-18.840436\du}{0.113194\du}}
\pgfpathlineto{\pgfpoint{-18.883299\du}{0.107930\du}}
\pgfpathlineto{\pgfpoint{-18.926914\du}{0.103920\du}}
\pgfpathlineto{\pgfpoint{-18.971030\du}{0.100160\du}}
\pgfpathlineto{\pgfpoint{-19.015647\du}{0.097653\du}}
\pgfpathlineto{\pgfpoint{-19.060264\du}{0.095147\du}}
\pgfpathlineto{\pgfpoint{-19.105383\du}{0.094395\du}}
\pgfpathlineto{\pgfpoint{-19.150502\du}{0.093643\du}}
\pgfpathlineto{\pgfpoint{-19.150502\du}{0.093643\du}}
\pgfpathlineto{\pgfpoint{-19.195871\du}{0.094395\du}}
\pgfpathlineto{\pgfpoint{-19.240990\du}{0.095147\du}}
\pgfpathlineto{\pgfpoint{-19.285607\du}{0.097653\du}}
\pgfpathlineto{\pgfpoint{-19.330224\du}{0.100160\du}}
\pgfpathlineto{\pgfpoint{-19.374340\du}{0.103920\du}}
\pgfpathlineto{\pgfpoint{-19.417955\du}{0.107930\du}}
\pgfpathlineto{\pgfpoint{-19.460818\du}{0.113194\du}}
\pgfpathlineto{\pgfpoint{-19.502678\du}{0.118709\du}}
\pgfpathlineto{\pgfpoint{-19.543786\du}{0.125727\du}}
\pgfpathlineto{\pgfpoint{-19.583891\du}{0.132495\du}}
\pgfpathlineto{\pgfpoint{-19.622994\du}{0.140516\du}}
\pgfpathlineto{\pgfpoint{-19.660844\du}{0.149289\du}}
\pgfpathlineto{\pgfpoint{-19.697189\du}{0.158564\du}}
\pgfpathlineto{\pgfpoint{-19.732281\du}{0.168339\du}}
\pgfpathlineto{\pgfpoint{-19.765870\du}{0.179118\du}}
\pgfpathlineto{\pgfpoint{-19.797703\du}{0.190147\du}}
\pgfpathlineto{\pgfpoint{-19.828284\du}{0.201677\du}}
\pgfpathlineto{\pgfpoint{-19.856859\du}{0.213709\du}}
\pgfpathlineto{\pgfpoint{-19.883679\du}{0.226492\du}}
\pgfpathlineto{\pgfpoint{-19.908495\du}{0.240028\du}}
\pgfpathlineto{\pgfpoint{-19.931305\du}{0.253313\du}}
\pgfpathlineto{\pgfpoint{-19.952360\du}{0.267099\du}}
\pgfpathlineto{\pgfpoint{-19.971410\du}{0.281387\du}}
\pgfpathlineto{\pgfpoint{-19.987954\du}{0.295925\du}}
\pgfpathlineto{\pgfpoint{-20.002743\du}{0.310714\du}}
\pgfpathlineto{\pgfpoint{-20.015025\du}{0.325753\du}}
\pgfpathlineto{\pgfpoint{-20.025553\du}{0.341294\du}}
\pgfpathlineto{\pgfpoint{-20.033574\du}{0.356584\du}}
\pgfpathlineto{\pgfpoint{-20.039339\du}{0.372125\du}}
\pgfpathlineto{\pgfpoint{-20.042347\du}{0.387917\du}}
\pgfpathlineto{\pgfpoint{-20.043600\du}{0.403708\du}}
\pgfpathlineto{\pgfpoint{-20.043600\du}{0.403708\du}}
\pgfpathlineto{\pgfpoint{-20.042347\du}{0.419249\du}}
\pgfpathlineto{\pgfpoint{-20.039339\du}{0.435041\du}}
\pgfpathlineto{\pgfpoint{-20.033574\du}{0.450581\du}}
\pgfpathlineto{\pgfpoint{-20.025553\du}{0.465872\du}}
\pgfpathlineto{\pgfpoint{-20.015025\du}{0.481412\du}}
\pgfpathlineto{\pgfpoint{-20.002743\du}{0.496452\du}}
\pgfpathlineto{\pgfpoint{-19.987954\du}{0.511241\du}}
\pgfpathlineto{\pgfpoint{-19.971410\du}{0.525779\du}}
\pgfpathlineto{\pgfpoint{-19.952360\du}{0.540067\du}}
\pgfpathlineto{\pgfpoint{-19.931305\du}{0.553853\du}}
\pgfpathlineto{\pgfpoint{-19.908495\du}{0.567639\du}}
\pgfpathlineto{\pgfpoint{-19.883679\du}{0.580673\du}}
\pgfpathlineto{\pgfpoint{-19.856859\du}{0.593457\du}}
\pgfpathlineto{\pgfpoint{-19.828284\du}{0.605489\du}}
\pgfpathlineto{\pgfpoint{-19.797703\du}{0.617520\du}}
\pgfpathlineto{\pgfpoint{-19.765870\du}{0.628549\du}}
\pgfpathlineto{\pgfpoint{-19.732281\du}{0.638826\du}}
\pgfpathlineto{\pgfpoint{-19.697189\du}{0.648853\du}}
\pgfpathlineto{\pgfpoint{-19.660844\du}{0.657876\du}}
\pgfpathlineto{\pgfpoint{-19.622994\du}{0.666650\du}}
\pgfpathlineto{\pgfpoint{-19.583891\du}{0.674671\du}}
\pgfpathlineto{\pgfpoint{-19.543786\du}{0.681940\du}}
\pgfpathlineto{\pgfpoint{-19.502678\du}{0.688457\du}}
\pgfpathlineto{\pgfpoint{-19.460818\du}{0.694473\du}}
\pgfpathlineto{\pgfpoint{-19.417955\du}{0.699235\du}}
\pgfpathlineto{\pgfpoint{-19.374340\du}{0.703747\du}}
\pgfpathlineto{\pgfpoint{-19.330224\du}{0.707006\du}}
\pgfpathlineto{\pgfpoint{-19.285607\du}{0.710014\du}}
\pgfpathlineto{\pgfpoint{-19.240990\du}{0.712019\du}}
\pgfpathlineto{\pgfpoint{-19.195871\du}{0.713272\du}}
\pgfpathlineto{\pgfpoint{-19.150502\du}{0.713523\du}}
\pgfpathlineto{\pgfpoint{-19.150502\du}{0.713523\du}}
\pgfpathlineto{\pgfpoint{-19.105383\du}{0.713272\du}}
\pgfpathlineto{\pgfpoint{-19.060264\du}{0.712019\du}}
\pgfpathlineto{\pgfpoint{-19.015647\du}{0.710014\du}}
\pgfpathlineto{\pgfpoint{-18.971030\du}{0.707006\du}}
\pgfpathlineto{\pgfpoint{-18.926914\du}{0.703747\du}}
\pgfpathlineto{\pgfpoint{-18.883299\du}{0.699235\du}}
\pgfpathlineto{\pgfpoint{-18.840436\du}{0.694473\du}}
\pgfpathlineto{\pgfpoint{-18.798827\du}{0.688457\du}}
\pgfpathlineto{\pgfpoint{-18.757719\du}{0.681940\du}}
\pgfpathlineto{\pgfpoint{-18.717363\du}{0.674671\du}}
\pgfpathlineto{\pgfpoint{-18.678260\du}{0.666650\du}}
\pgfpathlineto{\pgfpoint{-18.640410\du}{0.657876\du}}
\pgfpathlineto{\pgfpoint{-18.604065\du}{0.648853\du}}
\pgfpathlineto{\pgfpoint{-18.568973\du}{0.638826\du}}
\pgfpathlineto{\pgfpoint{-18.535635\du}{0.628549\du}}
\pgfpathlineto{\pgfpoint{-18.503551\du}{0.617520\du}}
\pgfpathlineto{\pgfpoint{-18.472970\du}{0.605489\du}}
\pgfpathlineto{\pgfpoint{-18.444395\du}{0.593457\du}}
\pgfpathlineto{\pgfpoint{-18.417574\du}{0.580673\du}}
\pgfpathlineto{\pgfpoint{-18.392759\du}{0.567639\du}}
\pgfpathlineto{\pgfpoint{-18.369949\du}{0.553853\du}}
\pgfpathlineto{\pgfpoint{-18.348894\du}{0.540067\du}}
\pgfpathlineto{\pgfpoint{-18.329844\du}{0.525779\du}}
\pgfpathlineto{\pgfpoint{-18.313300\du}{0.511241\du}}
\pgfpathlineto{\pgfpoint{-18.298511\du}{0.496452\du}}
\pgfpathlineto{\pgfpoint{-18.286229\du}{0.481412\du}}
\pgfpathlineto{\pgfpoint{-18.275701\du}{0.465872\du}}
\pgfpathlineto{\pgfpoint{-18.267931\du}{0.450581\du}}
\pgfpathlineto{\pgfpoint{-18.262416\du}{0.435041\du}}
\pgfpathlineto{\pgfpoint{-18.259158\du}{0.419249\du}}
\pgfpathlineto{\pgfpoint{-18.257654\du}{0.403708\du}}
\pgfusepath{fill}
\pgfsetbuttcap
\pgfsetmiterjoin
\pgfsetdash{}{0pt}
\definecolor{dialinecolor}{rgb}{0.666667, 0.901961, 1.000000}
\pgfsetstrokecolor{dialinecolor}
\pgfpathmoveto{\pgfpoint{-18.268182\du}{0.398695\du}}
\pgfpathlineto{\pgfpoint{-18.269435\du}{0.382903\du}}
\pgfpathlineto{\pgfpoint{-18.272443\du}{0.367864\du}}
\pgfpathlineto{\pgfpoint{-18.278208\du}{0.352323\du}}
\pgfpathlineto{\pgfpoint{-18.285978\du}{0.337284\du}}
\pgfpathlineto{\pgfpoint{-18.296255\du}{0.322244\du}}
\pgfpathlineto{\pgfpoint{-18.309290\du}{0.307455\du}}
\pgfpathlineto{\pgfpoint{-18.322825\du}{0.293168\du}}
\pgfpathlineto{\pgfpoint{-18.339870\du}{0.278629\du}}
\pgfpathlineto{\pgfpoint{-18.358670\du}{0.264592\du}}
\pgfpathlineto{\pgfpoint{-18.379474\du}{0.250555\du}}
\pgfpathlineto{\pgfpoint{-18.402535\du}{0.237521\du}}
\pgfpathlineto{\pgfpoint{-18.427100\du}{0.224487\du}}
\pgfpathlineto{\pgfpoint{-18.453669\du}{0.211954\du}}
\pgfpathlineto{\pgfpoint{-18.481994\du}{0.199922\du}}
\pgfpathlineto{\pgfpoint{-18.512073\du}{0.188392\du}}
\pgfpathlineto{\pgfpoint{-18.544157\du}{0.177864\du}}
\pgfpathlineto{\pgfpoint{-18.577746\du}{0.167838\du}}
\pgfpathlineto{\pgfpoint{-18.612337\du}{0.157561\du}}
\pgfpathlineto{\pgfpoint{-18.648431\du}{0.149039\du}}
\pgfpathlineto{\pgfpoint{-18.686030\du}{0.140015\du}}
\pgfpathlineto{\pgfpoint{-18.725133\du}{0.132244\du}}
\pgfpathlineto{\pgfpoint{-18.764737\du}{0.125226\du}}
\pgfpathlineto{\pgfpoint{-18.805344\du}{0.118709\du}}
\pgfpathlineto{\pgfpoint{-18.847204\du}{0.113194\du}}
\pgfpathlineto{\pgfpoint{-18.889566\du}{0.107930\du}}
\pgfpathlineto{\pgfpoint{-18.933180\du}{0.103920\du}}
\pgfpathlineto{\pgfpoint{-18.976795\du}{0.100160\du}}
\pgfpathlineto{\pgfpoint{-19.021162\du}{0.097653\du}}
\pgfpathlineto{\pgfpoint{-19.065779\du}{0.095899\du}}
\pgfpathlineto{\pgfpoint{-19.110647\du}{0.094395\du}}
\pgfpathlineto{\pgfpoint{-19.155515\du}{0.094395\du}}
\pgfpathlineto{\pgfpoint{-19.155515\du}{0.094395\du}}
\pgfpathlineto{\pgfpoint{-19.200884\du}{0.094395\du}}
\pgfpathlineto{\pgfpoint{-19.245501\du}{0.095899\du}}
\pgfpathlineto{\pgfpoint{-19.290369\du}{0.097653\du}}
\pgfpathlineto{\pgfpoint{-19.334235\du}{0.100160\du}}
\pgfpathlineto{\pgfpoint{-19.378100\du}{0.103920\du}}
\pgfpathlineto{\pgfpoint{-19.421464\du}{0.107930\du}}
\pgfpathlineto{\pgfpoint{-19.464076\du}{0.113194\du}}
\pgfpathlineto{\pgfpoint{-19.505936\du}{0.118709\du}}
\pgfpathlineto{\pgfpoint{-19.546794\du}{0.125226\du}}
\pgfpathlineto{\pgfpoint{-19.586398\du}{0.132244\du}}
\pgfpathlineto{\pgfpoint{-19.625501\du}{0.140015\du}}
\pgfpathlineto{\pgfpoint{-19.663099\du}{0.149039\du}}
\pgfpathlineto{\pgfpoint{-19.698944\du}{0.157561\du}}
\pgfpathlineto{\pgfpoint{-19.733785\du}{0.167838\du}}
\pgfpathlineto{\pgfpoint{-19.767374\du}{0.177864\du}}
\pgfpathlineto{\pgfpoint{-19.799207\du}{0.188392\du}}
\pgfpathlineto{\pgfpoint{-19.829036\du}{0.199922\du}}
\pgfpathlineto{\pgfpoint{-19.857360\du}{0.211954\du}}
\pgfpathlineto{\pgfpoint{-19.884181\du}{0.224487\du}}
\pgfpathlineto{\pgfpoint{-19.909247\du}{0.237521\du}}
\pgfpathlineto{\pgfpoint{-19.931555\du}{0.250555\du}}
\pgfpathlineto{\pgfpoint{-19.952611\du}{0.264592\du}}
\pgfpathlineto{\pgfpoint{-19.971410\du}{0.278629\du}}
\pgfpathlineto{\pgfpoint{-19.987954\du}{0.293168\du}}
\pgfpathlineto{\pgfpoint{-20.002492\du}{0.307455\du}}
\pgfpathlineto{\pgfpoint{-20.015025\du}{0.322244\du}}
\pgfpathlineto{\pgfpoint{-20.025051\du}{0.337284\du}}
\pgfpathlineto{\pgfpoint{-20.033072\du}{0.352323\du}}
\pgfpathlineto{\pgfpoint{-20.038837\du}{0.367864\du}}
\pgfpathlineto{\pgfpoint{-20.042096\du}{0.382903\du}}
\pgfpathlineto{\pgfpoint{-20.043349\du}{0.398695\du}}
\pgfpathlineto{\pgfpoint{-20.043349\du}{0.398695\du}}
\pgfpathlineto{\pgfpoint{-20.042096\du}{0.413985\du}}
\pgfpathlineto{\pgfpoint{-20.038837\du}{0.429777\du}}
\pgfpathlineto{\pgfpoint{-20.033072\du}{0.444566\du}}
\pgfpathlineto{\pgfpoint{-20.025051\du}{0.459856\du}}
\pgfpathlineto{\pgfpoint{-20.015025\du}{0.475146\du}}
\pgfpathlineto{\pgfpoint{-20.002492\du}{0.489684\du}}
\pgfpathlineto{\pgfpoint{-19.987954\du}{0.504222\du}}
\pgfpathlineto{\pgfpoint{-19.971410\du}{0.518761\du}}
\pgfpathlineto{\pgfpoint{-19.952611\du}{0.532547\du}}
\pgfpathlineto{\pgfpoint{-19.931555\du}{0.546584\du}}
\pgfpathlineto{\pgfpoint{-19.909247\du}{0.559869\du}}
\pgfpathlineto{\pgfpoint{-19.884181\du}{0.572402\du}}
\pgfpathlineto{\pgfpoint{-19.857360\du}{0.585185\du}}
\pgfpathlineto{\pgfpoint{-19.829036\du}{0.597468\du}}
\pgfpathlineto{\pgfpoint{-19.799207\du}{0.608497\du}}
\pgfpathlineto{\pgfpoint{-19.767374\du}{0.619275\du}}
\pgfpathlineto{\pgfpoint{-19.733785\du}{0.629552\du}}
\pgfpathlineto{\pgfpoint{-19.698944\du}{0.639578\du}}
\pgfpathlineto{\pgfpoint{-19.663099\du}{0.648351\du}}
\pgfpathlineto{\pgfpoint{-19.625501\du}{0.656874\du}}
\pgfpathlineto{\pgfpoint{-19.586398\du}{0.664895\du}}
\pgfpathlineto{\pgfpoint{-19.546794\du}{0.671913\du}}
\pgfpathlineto{\pgfpoint{-19.505936\du}{0.678180\du}}
\pgfpathlineto{\pgfpoint{-19.464076\du}{0.684196\du}}
\pgfpathlineto{\pgfpoint{-19.421464\du}{0.689209\du}}
\pgfpathlineto{\pgfpoint{-19.378100\du}{0.693219\du}}
\pgfpathlineto{\pgfpoint{-19.334235\du}{0.696729\du}}
\pgfpathlineto{\pgfpoint{-19.290369\du}{0.699235\du}}
\pgfpathlineto{\pgfpoint{-19.245501\du}{0.701491\du}}
\pgfpathlineto{\pgfpoint{-19.200884\du}{0.702744\du}}
\pgfpathlineto{\pgfpoint{-19.155515\du}{0.702995\du}}
\pgfpathlineto{\pgfpoint{-19.155515\du}{0.702995\du}}
\pgfpathlineto{\pgfpoint{-19.110647\du}{0.702744\du}}
\pgfpathlineto{\pgfpoint{-19.065779\du}{0.701491\du}}
\pgfpathlineto{\pgfpoint{-19.021162\du}{0.699235\du}}
\pgfpathlineto{\pgfpoint{-18.976795\du}{0.696729\du}}
\pgfpathlineto{\pgfpoint{-18.933180\du}{0.693219\du}}
\pgfpathlineto{\pgfpoint{-18.889566\du}{0.689209\du}}
\pgfpathlineto{\pgfpoint{-18.847204\du}{0.684196\du}}
\pgfpathlineto{\pgfpoint{-18.805344\du}{0.678180\du}}
\pgfpathlineto{\pgfpoint{-18.764737\du}{0.671913\du}}
\pgfpathlineto{\pgfpoint{-18.725133\du}{0.664895\du}}
\pgfpathlineto{\pgfpoint{-18.686030\du}{0.656874\du}}
\pgfpathlineto{\pgfpoint{-18.648431\du}{0.648351\du}}
\pgfpathlineto{\pgfpoint{-18.612337\du}{0.639578\du}}
\pgfpathlineto{\pgfpoint{-18.577746\du}{0.629552\du}}
\pgfpathlineto{\pgfpoint{-18.544157\du}{0.619275\du}}
\pgfpathlineto{\pgfpoint{-18.512073\du}{0.608497\du}}
\pgfpathlineto{\pgfpoint{-18.481994\du}{0.597468\du}}
\pgfpathlineto{\pgfpoint{-18.453669\du}{0.585185\du}}
\pgfpathlineto{\pgfpoint{-18.427100\du}{0.572402\du}}
\pgfpathlineto{\pgfpoint{-18.402535\du}{0.559869\du}}
\pgfpathlineto{\pgfpoint{-18.379474\du}{0.546584\du}}
\pgfpathlineto{\pgfpoint{-18.358670\du}{0.532547\du}}
\pgfpathlineto{\pgfpoint{-18.339870\du}{0.518761\du}}
\pgfpathlineto{\pgfpoint{-18.322825\du}{0.504222\du}}
\pgfpathlineto{\pgfpoint{-18.309290\du}{0.489684\du}}
\pgfpathlineto{\pgfpoint{-18.296255\du}{0.475146\du}}
\pgfpathlineto{\pgfpoint{-18.285978\du}{0.459856\du}}
\pgfpathlineto{\pgfpoint{-18.278208\du}{0.444566\du}}
\pgfpathlineto{\pgfpoint{-18.272443\du}{0.429777\du}}
\pgfpathlineto{\pgfpoint{-18.269435\du}{0.413985\du}}
\pgfpathlineto{\pgfpoint{-18.268182\du}{0.398695\du}}
\pgfusepath{stroke}
\pgfsetbuttcap
\pgfsetmiterjoin
\pgfsetdash{}{0pt}
\definecolor{dialinecolor}{rgb}{0.000000, 0.470588, 0.666667}
\pgfsetfillcolor{dialinecolor}
\pgfpathmoveto{\pgfpoint{-20.043349\du}{-0.025671\du}}
\pgfpathlineto{\pgfpoint{-20.043349\du}{0.409473\du}}
\pgfpathlineto{\pgfpoint{-18.268182\du}{0.409473\du}}
\pgfpathlineto{\pgfpoint{-18.268182\du}{-0.025671\du}}
\pgfpathlineto{\pgfpoint{-20.043349\du}{-0.025671\du}}
\pgfusepath{fill}
\pgfsetbuttcap
\pgfsetmiterjoin
\pgfsetdash{}{0pt}
\definecolor{dialinecolor}{rgb}{0.000000, 0.705882, 1.000000}
\pgfsetfillcolor{dialinecolor}
\pgfpathmoveto{\pgfpoint{-18.257654\du}{-0.030935\du}}
\pgfpathlineto{\pgfpoint{-18.259158\du}{-0.046726\du}}
\pgfpathlineto{\pgfpoint{-18.262416\du}{-0.062768\du}}
\pgfpathlineto{\pgfpoint{-18.267931\du}{-0.078059\du}}
\pgfpathlineto{\pgfpoint{-18.275701\du}{-0.093599\du}}
\pgfpathlineto{\pgfpoint{-18.286229\du}{-0.109140\du}}
\pgfpathlineto{\pgfpoint{-18.298511\du}{-0.124180\du}}
\pgfpathlineto{\pgfpoint{-18.313300\du}{-0.138969\du}}
\pgfpathlineto{\pgfpoint{-18.329844\du}{-0.153256\du}}
\pgfpathlineto{\pgfpoint{-18.348894\du}{-0.167795\du}}
\pgfpathlineto{\pgfpoint{-18.369949\du}{-0.181330\du}}
\pgfpathlineto{\pgfpoint{-18.392759\du}{-0.195116\du}}
\pgfpathlineto{\pgfpoint{-18.417574\du}{-0.208151\du}}
\pgfpathlineto{\pgfpoint{-18.444395\du}{-0.221185\du}}
\pgfpathlineto{\pgfpoint{-18.472970\du}{-0.232966\du}}
\pgfpathlineto{\pgfpoint{-18.503551\du}{-0.244747\du}}
\pgfpathlineto{\pgfpoint{-18.535635\du}{-0.256027\du}}
\pgfpathlineto{\pgfpoint{-18.568973\du}{-0.266554\du}}
\pgfpathlineto{\pgfpoint{-18.604065\du}{-0.276330\du}}
\pgfpathlineto{\pgfpoint{-18.640410\du}{-0.285604\du}}
\pgfpathlineto{\pgfpoint{-18.678260\du}{-0.294127\du}}
\pgfpathlineto{\pgfpoint{-18.717363\du}{-0.301897\du}}
\pgfpathlineto{\pgfpoint{-18.757719\du}{-0.309417\du}}
\pgfpathlineto{\pgfpoint{-18.798827\du}{-0.315934\du}}
\pgfpathlineto{\pgfpoint{-18.840436\du}{-0.321950\du}}
\pgfpathlineto{\pgfpoint{-18.883299\du}{-0.326963\du}}
\pgfpathlineto{\pgfpoint{-18.926914\du}{-0.331224\du}}
\pgfpathlineto{\pgfpoint{-18.971030\du}{-0.334734\du}}
\pgfpathlineto{\pgfpoint{-19.015647\du}{-0.337491\du}}
\pgfpathlineto{\pgfpoint{-19.060264\du}{-0.339496\du}}
\pgfpathlineto{\pgfpoint{-19.105383\du}{-0.340749\du}}
\pgfpathlineto{\pgfpoint{-19.150502\du}{-0.341000\du}}
\pgfpathlineto{\pgfpoint{-19.150502\du}{-0.341000\du}}
\pgfpathlineto{\pgfpoint{-19.195871\du}{-0.340749\du}}
\pgfpathlineto{\pgfpoint{-19.240990\du}{-0.339496\du}}
\pgfpathlineto{\pgfpoint{-19.285607\du}{-0.337491\du}}
\pgfpathlineto{\pgfpoint{-19.330224\du}{-0.334734\du}}
\pgfpathlineto{\pgfpoint{-19.374340\du}{-0.331224\du}}
\pgfpathlineto{\pgfpoint{-19.417955\du}{-0.326963\du}}
\pgfpathlineto{\pgfpoint{-19.460818\du}{-0.321950\du}}
\pgfpathlineto{\pgfpoint{-19.502678\du}{-0.315934\du}}
\pgfpathlineto{\pgfpoint{-19.543786\du}{-0.309417\du}}
\pgfpathlineto{\pgfpoint{-19.583891\du}{-0.301897\du}}
\pgfpathlineto{\pgfpoint{-19.622994\du}{-0.294127\du}}
\pgfpathlineto{\pgfpoint{-19.660844\du}{-0.285604\du}}
\pgfpathlineto{\pgfpoint{-19.697189\du}{-0.276330\du}}
\pgfpathlineto{\pgfpoint{-19.732281\du}{-0.266554\du}}
\pgfpathlineto{\pgfpoint{-19.765870\du}{-0.256027\du}}
\pgfpathlineto{\pgfpoint{-19.797703\du}{-0.244747\du}}
\pgfpathlineto{\pgfpoint{-19.828284\du}{-0.232966\du}}
\pgfpathlineto{\pgfpoint{-19.856859\du}{-0.221185\du}}
\pgfpathlineto{\pgfpoint{-19.883679\du}{-0.208151\du}}
\pgfpathlineto{\pgfpoint{-19.908495\du}{-0.195116\du}}
\pgfpathlineto{\pgfpoint{-19.931305\du}{-0.181330\du}}
\pgfpathlineto{\pgfpoint{-19.952360\du}{-0.167795\du}}
\pgfpathlineto{\pgfpoint{-19.971410\du}{-0.153256\du}}
\pgfpathlineto{\pgfpoint{-19.987954\du}{-0.138969\du}}
\pgfpathlineto{\pgfpoint{-20.002743\du}{-0.124180\du}}
\pgfpathlineto{\pgfpoint{-20.015025\du}{-0.109140\du}}
\pgfpathlineto{\pgfpoint{-20.025553\du}{-0.093599\du}}
\pgfpathlineto{\pgfpoint{-20.033574\du}{-0.078059\du}}
\pgfpathlineto{\pgfpoint{-20.039339\du}{-0.062768\du}}
\pgfpathlineto{\pgfpoint{-20.042347\du}{-0.046726\du}}
\pgfpathlineto{\pgfpoint{-20.043600\du}{-0.030935\du}}
\pgfpathlineto{\pgfpoint{-20.043600\du}{-0.030935\du}}
\pgfpathlineto{\pgfpoint{-20.042347\du}{-0.015394\du}}
\pgfpathlineto{\pgfpoint{-20.039339\du}{0.000398\du}}
\pgfpathlineto{\pgfpoint{-20.033574\du}{0.015939\du}}
\pgfpathlineto{\pgfpoint{-20.025553\du}{0.030978\du}}
\pgfpathlineto{\pgfpoint{-20.015025\du}{0.046519\du}}
\pgfpathlineto{\pgfpoint{-20.002743\du}{0.061559\du}}
\pgfpathlineto{\pgfpoint{-19.987954\du}{0.076347\du}}
\pgfpathlineto{\pgfpoint{-19.971410\du}{0.091136\du}}
\pgfpathlineto{\pgfpoint{-19.952360\du}{0.105424\du}}
\pgfpathlineto{\pgfpoint{-19.931305\du}{0.118959\du}}
\pgfpathlineto{\pgfpoint{-19.908495\du}{0.132495\du}}
\pgfpathlineto{\pgfpoint{-19.883679\du}{0.146031\du}}
\pgfpathlineto{\pgfpoint{-19.856859\du}{0.158564\du}}
\pgfpathlineto{\pgfpoint{-19.828284\du}{0.170595\du}}
\pgfpathlineto{\pgfpoint{-19.797703\du}{0.182376\du}}
\pgfpathlineto{\pgfpoint{-19.765870\du}{0.193405\du}}
\pgfpathlineto{\pgfpoint{-19.732281\du}{0.204184\du}}
\pgfpathlineto{\pgfpoint{-19.697189\du}{0.213709\du}}
\pgfpathlineto{\pgfpoint{-19.660844\du}{0.222983\du}}
\pgfpathlineto{\pgfpoint{-19.622994\du}{0.231756\du}}
\pgfpathlineto{\pgfpoint{-19.583891\du}{0.240028\du}}
\pgfpathlineto{\pgfpoint{-19.543786\du}{0.246796\du}}
\pgfpathlineto{\pgfpoint{-19.502678\du}{0.253563\du}}
\pgfpathlineto{\pgfpoint{-19.460818\du}{0.259579\du}}
\pgfpathlineto{\pgfpoint{-19.417955\du}{0.264342\du}}
\pgfpathlineto{\pgfpoint{-19.374340\du}{0.268603\du}}
\pgfpathlineto{\pgfpoint{-19.330224\du}{0.272112\du}}
\pgfpathlineto{\pgfpoint{-19.285607\du}{0.274869\du}}
\pgfpathlineto{\pgfpoint{-19.240990\du}{0.277125\du}}
\pgfpathlineto{\pgfpoint{-19.195871\du}{0.278379\du}}
\pgfpathlineto{\pgfpoint{-19.150502\du}{0.278629\du}}
\pgfpathlineto{\pgfpoint{-19.150502\du}{0.278629\du}}
\pgfpathlineto{\pgfpoint{-19.105383\du}{0.278379\du}}
\pgfpathlineto{\pgfpoint{-19.060264\du}{0.277125\du}}
\pgfpathlineto{\pgfpoint{-19.015647\du}{0.274869\du}}
\pgfpathlineto{\pgfpoint{-18.971030\du}{0.272112\du}}
\pgfpathlineto{\pgfpoint{-18.926914\du}{0.268603\du}}
\pgfpathlineto{\pgfpoint{-18.883299\du}{0.264342\du}}
\pgfpathlineto{\pgfpoint{-18.840436\du}{0.259579\du}}
\pgfpathlineto{\pgfpoint{-18.798827\du}{0.253563\du}}
\pgfpathlineto{\pgfpoint{-18.757719\du}{0.246796\du}}
\pgfpathlineto{\pgfpoint{-18.717363\du}{0.240028\du}}
\pgfpathlineto{\pgfpoint{-18.678260\du}{0.231756\du}}
\pgfpathlineto{\pgfpoint{-18.640410\du}{0.222983\du}}
\pgfpathlineto{\pgfpoint{-18.604065\du}{0.213709\du}}
\pgfpathlineto{\pgfpoint{-18.568973\du}{0.204184\du}}
\pgfpathlineto{\pgfpoint{-18.535635\du}{0.193405\du}}
\pgfpathlineto{\pgfpoint{-18.503551\du}{0.182376\du}}
\pgfpathlineto{\pgfpoint{-18.472970\du}{0.170595\du}}
\pgfpathlineto{\pgfpoint{-18.444395\du}{0.158564\du}}
\pgfpathlineto{\pgfpoint{-18.417574\du}{0.146031\du}}
\pgfpathlineto{\pgfpoint{-18.392759\du}{0.132495\du}}
\pgfpathlineto{\pgfpoint{-18.369949\du}{0.118959\du}}
\pgfpathlineto{\pgfpoint{-18.348894\du}{0.105424\du}}
\pgfpathlineto{\pgfpoint{-18.329844\du}{0.091136\du}}
\pgfpathlineto{\pgfpoint{-18.313300\du}{0.076347\du}}
\pgfpathlineto{\pgfpoint{-18.298511\du}{0.061559\du}}
\pgfpathlineto{\pgfpoint{-18.286229\du}{0.046519\du}}
\pgfpathlineto{\pgfpoint{-18.275701\du}{0.030978\du}}
\pgfpathlineto{\pgfpoint{-18.267931\du}{0.015939\du}}
\pgfpathlineto{\pgfpoint{-18.262416\du}{0.000398\du}}
\pgfpathlineto{\pgfpoint{-18.259158\du}{-0.015394\du}}
\pgfpathlineto{\pgfpoint{-18.257654\du}{-0.030935\du}}
\pgfusepath{fill}
\pgfsetbuttcap
\pgfsetmiterjoin
\pgfsetdash{}{0pt}
\definecolor{dialinecolor}{rgb}{0.666667, 0.901961, 1.000000}
\pgfsetstrokecolor{dialinecolor}
\pgfpathmoveto{\pgfpoint{-18.268182\du}{-0.036199\du}}
\pgfpathlineto{\pgfpoint{-18.269435\du}{-0.051990\du}}
\pgfpathlineto{\pgfpoint{-18.272443\du}{-0.067030\du}}
\pgfpathlineto{\pgfpoint{-18.278208\du}{-0.082320\du}}
\pgfpathlineto{\pgfpoint{-18.285978\du}{-0.097610\du}}
\pgfpathlineto{\pgfpoint{-18.296255\du}{-0.112399\du}}
\pgfpathlineto{\pgfpoint{-18.309290\du}{-0.127689\du}}
\pgfpathlineto{\pgfpoint{-18.322825\du}{-0.141726\du}}
\pgfpathlineto{\pgfpoint{-18.339870\du}{-0.156264\du}}
\pgfpathlineto{\pgfpoint{-18.358670\du}{-0.170050\du}}
\pgfpathlineto{\pgfpoint{-18.379474\du}{-0.184338\du}}
\pgfpathlineto{\pgfpoint{-18.402535\du}{-0.197122\du}}
\pgfpathlineto{\pgfpoint{-18.427100\du}{-0.210407\du}}
\pgfpathlineto{\pgfpoint{-18.453669\du}{-0.222940\du}}
\pgfpathlineto{\pgfpoint{-18.481994\du}{-0.234470\du}}
\pgfpathlineto{\pgfpoint{-18.512073\du}{-0.246502\du}}
\pgfpathlineto{\pgfpoint{-18.544157\du}{-0.256528\du}}
\pgfpathlineto{\pgfpoint{-18.577746\du}{-0.267306\du}}
\pgfpathlineto{\pgfpoint{-18.612337\du}{-0.277082\du}}
\pgfpathlineto{\pgfpoint{-18.648431\du}{-0.286106\du}}
\pgfpathlineto{\pgfpoint{-18.686030\du}{-0.294879\du}}
\pgfpathlineto{\pgfpoint{-18.725133\du}{-0.302649\du}}
\pgfpathlineto{\pgfpoint{-18.764737\du}{-0.309668\du}}
\pgfpathlineto{\pgfpoint{-18.805344\du}{-0.315934\du}}
\pgfpathlineto{\pgfpoint{-18.847204\du}{-0.321950\du}}
\pgfpathlineto{\pgfpoint{-18.889566\du}{-0.326963\du}}
\pgfpathlineto{\pgfpoint{-18.933180\du}{-0.330974\du}}
\pgfpathlineto{\pgfpoint{-18.976795\du}{-0.334734\du}}
\pgfpathlineto{\pgfpoint{-19.021162\du}{-0.337240\du}}
\pgfpathlineto{\pgfpoint{-19.065779\du}{-0.339245\du}}
\pgfpathlineto{\pgfpoint{-19.110647\du}{-0.340499\du}}
\pgfpathlineto{\pgfpoint{-19.155515\du}{-0.340749\du}}
\pgfpathlineto{\pgfpoint{-19.155515\du}{-0.340749\du}}
\pgfpathlineto{\pgfpoint{-19.200884\du}{-0.340499\du}}
\pgfpathlineto{\pgfpoint{-19.245501\du}{-0.339245\du}}
\pgfpathlineto{\pgfpoint{-19.290369\du}{-0.337240\du}}
\pgfpathlineto{\pgfpoint{-19.334235\du}{-0.334734\du}}
\pgfpathlineto{\pgfpoint{-19.378100\du}{-0.330974\du}}
\pgfpathlineto{\pgfpoint{-19.421464\du}{-0.326963\du}}
\pgfpathlineto{\pgfpoint{-19.464076\du}{-0.321950\du}}
\pgfpathlineto{\pgfpoint{-19.505936\du}{-0.315934\du}}
\pgfpathlineto{\pgfpoint{-19.546794\du}{-0.309668\du}}
\pgfpathlineto{\pgfpoint{-19.586398\du}{-0.302649\du}}
\pgfpathlineto{\pgfpoint{-19.625501\du}{-0.294879\du}}
\pgfpathlineto{\pgfpoint{-19.663099\du}{-0.286106\du}}
\pgfpathlineto{\pgfpoint{-19.698944\du}{-0.277082\du}}
\pgfpathlineto{\pgfpoint{-19.733785\du}{-0.267306\du}}
\pgfpathlineto{\pgfpoint{-19.767374\du}{-0.256528\du}}
\pgfpathlineto{\pgfpoint{-19.799207\du}{-0.246502\du}}
\pgfpathlineto{\pgfpoint{-19.829036\du}{-0.234470\du}}
\pgfpathlineto{\pgfpoint{-19.857360\du}{-0.222940\du}}
\pgfpathlineto{\pgfpoint{-19.884181\du}{-0.210407\du}}
\pgfpathlineto{\pgfpoint{-19.909247\du}{-0.197122\du}}
\pgfpathlineto{\pgfpoint{-19.931555\du}{-0.184338\du}}
\pgfpathlineto{\pgfpoint{-19.952611\du}{-0.170050\du}}
\pgfpathlineto{\pgfpoint{-19.971410\du}{-0.156264\du}}
\pgfpathlineto{\pgfpoint{-19.987954\du}{-0.141726\du}}
\pgfpathlineto{\pgfpoint{-20.002492\du}{-0.127689\du}}
\pgfpathlineto{\pgfpoint{-20.015025\du}{-0.112399\du}}
\pgfpathlineto{\pgfpoint{-20.025051\du}{-0.097610\du}}
\pgfpathlineto{\pgfpoint{-20.033072\du}{-0.082320\du}}
\pgfpathlineto{\pgfpoint{-20.038837\du}{-0.067030\du}}
\pgfpathlineto{\pgfpoint{-20.042096\du}{-0.051990\du}}
\pgfpathlineto{\pgfpoint{-20.043349\du}{-0.036199\du}}
\pgfpathlineto{\pgfpoint{-20.043349\du}{-0.036199\du}}
\pgfpathlineto{\pgfpoint{-20.042096\du}{-0.020658\du}}
\pgfpathlineto{\pgfpoint{-20.038837\du}{-0.005367\du}}
\pgfpathlineto{\pgfpoint{-20.033072\du}{0.009923\du}}
\pgfpathlineto{\pgfpoint{-20.025051\du}{0.024962\du}}
\pgfpathlineto{\pgfpoint{-20.015025\du}{0.040002\du}}
\pgfpathlineto{\pgfpoint{-20.002492\du}{0.055041\du}}
\pgfpathlineto{\pgfpoint{-19.987954\du}{0.069329\du}}
\pgfpathlineto{\pgfpoint{-19.971410\du}{0.083867\du}}
\pgfpathlineto{\pgfpoint{-19.952611\du}{0.097653\du}}
\pgfpathlineto{\pgfpoint{-19.931555\du}{0.111690\du}}
\pgfpathlineto{\pgfpoint{-19.909247\du}{0.124975\du}}
\pgfpathlineto{\pgfpoint{-19.884181\du}{0.137508\du}}
\pgfpathlineto{\pgfpoint{-19.857360\du}{0.150292\du}}
\pgfpathlineto{\pgfpoint{-19.829036\du}{0.162323\du}}
\pgfpathlineto{\pgfpoint{-19.799207\du}{0.173603\du}}
\pgfpathlineto{\pgfpoint{-19.767374\du}{0.184381\du}}
\pgfpathlineto{\pgfpoint{-19.733785\du}{0.194658\du}}
\pgfpathlineto{\pgfpoint{-19.698944\du}{0.204434\du}}
\pgfpathlineto{\pgfpoint{-19.663099\du}{0.213458\du}}
\pgfpathlineto{\pgfpoint{-19.625501\du}{0.222231\du}}
\pgfpathlineto{\pgfpoint{-19.586398\du}{0.230503\du}}
\pgfpathlineto{\pgfpoint{-19.546794\du}{0.237020\du}}
\pgfpathlineto{\pgfpoint{-19.505936\du}{0.243788\du}}
\pgfpathlineto{\pgfpoint{-19.464076\du}{0.249553\du}}
\pgfpathlineto{\pgfpoint{-19.421464\du}{0.254315\du}}
\pgfpathlineto{\pgfpoint{-19.378100\du}{0.258827\du}}
\pgfpathlineto{\pgfpoint{-19.334235\du}{0.262336\du}}
\pgfpathlineto{\pgfpoint{-19.290369\du}{0.264592\du}}
\pgfpathlineto{\pgfpoint{-19.245501\du}{0.266848\du}}
\pgfpathlineto{\pgfpoint{-19.200884\du}{0.268102\du}}
\pgfpathlineto{\pgfpoint{-19.155515\du}{0.268352\du}}
\pgfpathlineto{\pgfpoint{-19.155515\du}{0.268352\du}}
\pgfpathlineto{\pgfpoint{-19.110647\du}{0.268102\du}}
\pgfpathlineto{\pgfpoint{-19.065779\du}{0.266848\du}}
\pgfpathlineto{\pgfpoint{-19.021162\du}{0.264592\du}}
\pgfpathlineto{\pgfpoint{-18.976795\du}{0.262336\du}}
\pgfpathlineto{\pgfpoint{-18.933180\du}{0.258827\du}}
\pgfpathlineto{\pgfpoint{-18.889566\du}{0.254315\du}}
\pgfpathlineto{\pgfpoint{-18.847204\du}{0.249553\du}}
\pgfpathlineto{\pgfpoint{-18.805344\du}{0.243788\du}}
\pgfpathlineto{\pgfpoint{-18.764737\du}{0.237020\du}}
\pgfpathlineto{\pgfpoint{-18.725133\du}{0.230503\du}}
\pgfpathlineto{\pgfpoint{-18.686030\du}{0.222231\du}}
\pgfpathlineto{\pgfpoint{-18.648431\du}{0.213458\du}}
\pgfpathlineto{\pgfpoint{-18.612337\du}{0.204434\du}}
\pgfpathlineto{\pgfpoint{-18.577746\du}{0.194658\du}}
\pgfpathlineto{\pgfpoint{-18.544157\du}{0.184381\du}}
\pgfpathlineto{\pgfpoint{-18.512073\du}{0.173603\du}}
\pgfpathlineto{\pgfpoint{-18.481994\du}{0.162323\du}}
\pgfpathlineto{\pgfpoint{-18.453669\du}{0.150292\du}}
\pgfpathlineto{\pgfpoint{-18.427100\du}{0.137508\du}}
\pgfpathlineto{\pgfpoint{-18.402535\du}{0.124975\du}}
\pgfpathlineto{\pgfpoint{-18.379474\du}{0.111690\du}}
\pgfpathlineto{\pgfpoint{-18.358670\du}{0.097653\du}}
\pgfpathlineto{\pgfpoint{-18.339870\du}{0.083867\du}}
\pgfpathlineto{\pgfpoint{-18.322825\du}{0.069329\du}}
\pgfpathlineto{\pgfpoint{-18.309290\du}{0.055041\du}}
\pgfpathlineto{\pgfpoint{-18.296255\du}{0.040002\du}}
\pgfpathlineto{\pgfpoint{-18.285978\du}{0.024962\du}}
\pgfpathlineto{\pgfpoint{-18.278208\du}{0.009923\du}}
\pgfpathlineto{\pgfpoint{-18.272443\du}{-0.005367\du}}
\pgfpathlineto{\pgfpoint{-18.269435\du}{-0.020658\du}}
\pgfpathlineto{\pgfpoint{-18.268182\du}{-0.036199\du}}
\pgfusepath{stroke}
\pgfsetbuttcap
\pgfsetmiterjoin
\pgfsetdash{}{0pt}
\definecolor{dialinecolor}{rgb}{0.000000, 0.000000, 0.000000}
\pgfsetfillcolor{dialinecolor}
\pgfpathmoveto{\pgfpoint{-19.133958\du}{-0.101370\du}}
\pgfpathlineto{\pgfpoint{-19.004117\du}{-0.057755\du}}
\pgfpathlineto{\pgfpoint{-18.690041\du}{-0.188599\du}}
\pgfpathlineto{\pgfpoint{-18.549421\du}{-0.145235\du}}
\pgfpathlineto{\pgfpoint{-18.625371\du}{-0.253520\du}}
\pgfpathlineto{\pgfpoint{-18.993088\du}{-0.253520\du}}
\pgfpathlineto{\pgfpoint{-18.841940\du}{-0.221185\du}}
\pgfpathlineto{\pgfpoint{-19.133958\du}{-0.101370\du}}
\pgfusepath{fill}
\pgfsetbuttcap
\pgfsetmiterjoin
\pgfsetdash{}{0pt}
\definecolor{dialinecolor}{rgb}{0.000000, 0.000000, 0.000000}
\pgfsetfillcolor{dialinecolor}
\pgfpathmoveto{\pgfpoint{-19.188351\du}{0.018194\du}}
\pgfpathlineto{\pgfpoint{-19.317942\du}{-0.025671\du}}
\pgfpathlineto{\pgfpoint{-19.621239\du}{0.105424\du}}
\pgfpathlineto{\pgfpoint{-19.772638\du}{0.061559\du}}
\pgfpathlineto{\pgfpoint{-19.697189\du}{0.181123\du}}
\pgfpathlineto{\pgfpoint{-19.317942\du}{0.181123\du}}
\pgfpathlineto{\pgfpoint{-19.480620\du}{0.137508\du}}
\pgfpathlineto{\pgfpoint{-19.188351\du}{0.018194\du}}
\pgfusepath{fill}
\pgfsetbuttcap
\pgfsetmiterjoin
\pgfsetdash{}{0pt}
\definecolor{dialinecolor}{rgb}{0.000000, 0.000000, 0.000000}
\pgfsetfillcolor{dialinecolor}
\pgfpathmoveto{\pgfpoint{-19.740302\du}{-0.221185\du}}
\pgfpathlineto{\pgfpoint{-19.610712\du}{-0.264549\du}}
\pgfpathlineto{\pgfpoint{-19.296385\du}{-0.145235\du}}
\pgfpathlineto{\pgfpoint{-19.155515\du}{-0.177570\du}}
\pgfpathlineto{\pgfpoint{-19.231214\du}{-0.069286\du}}
\pgfpathlineto{\pgfpoint{-19.599683\du}{-0.069286\du}}
\pgfpathlineto{\pgfpoint{-19.448285\du}{-0.101370\du}}
\pgfpathlineto{\pgfpoint{-19.740302\du}{-0.221185\du}}
\pgfusepath{fill}
\pgfsetbuttcap
\pgfsetmiterjoin
\pgfsetdash{}{0pt}
\definecolor{dialinecolor}{rgb}{0.000000, 0.000000, 0.000000}
\pgfsetfillcolor{dialinecolor}
\pgfpathmoveto{\pgfpoint{-18.571228\du}{0.148287\du}}
\pgfpathlineto{\pgfpoint{-18.701321\du}{0.191901\du}}
\pgfpathlineto{\pgfpoint{-19.004117\du}{0.061559\du}}
\pgfpathlineto{\pgfpoint{-19.155515\du}{0.105424\du}}
\pgfpathlineto{\pgfpoint{-19.079816\du}{-0.003864\du}}
\pgfpathlineto{\pgfpoint{-18.701321\du}{-0.003864\du}}
\pgfpathlineto{\pgfpoint{-18.863497\du}{0.028722\du}}
\pgfpathlineto{\pgfpoint{-18.571228\du}{0.148287\du}}
\pgfusepath{fill}
\pgfsetbuttcap
\pgfsetmiterjoin
\pgfsetdash{}{0pt}
\definecolor{dialinecolor}{rgb}{1.000000, 1.000000, 1.000000}
\pgfsetfillcolor{dialinecolor}
\pgfpathmoveto{\pgfpoint{-19.123180\du}{-0.090592\du}}
\pgfpathlineto{\pgfpoint{-18.993088\du}{-0.047228\du}}
\pgfpathlineto{\pgfpoint{-18.679764\du}{-0.177570\du}}
\pgfpathlineto{\pgfpoint{-18.538643\du}{-0.133956\du}}
\pgfpathlineto{\pgfpoint{-18.614593\du}{-0.242992\du}}
\pgfpathlineto{\pgfpoint{-18.982560\du}{-0.242992\du}}
\pgfpathlineto{\pgfpoint{-18.830911\du}{-0.210657\du}}
\pgfpathlineto{\pgfpoint{-19.123180\du}{-0.090592\du}}
\pgfusepath{fill}
\pgfsetbuttcap
\pgfsetmiterjoin
\pgfsetdash{}{0pt}
\definecolor{dialinecolor}{rgb}{1.000000, 1.000000, 1.000000}
\pgfsetfillcolor{dialinecolor}
\pgfpathmoveto{\pgfpoint{-19.177322\du}{0.028722\du}}
\pgfpathlineto{\pgfpoint{-19.307665\du}{-0.014642\du}}
\pgfpathlineto{\pgfpoint{-19.610712\du}{0.115701\du}}
\pgfpathlineto{\pgfpoint{-19.762110\du}{0.072337\du}}
\pgfpathlineto{\pgfpoint{-19.686411\du}{0.191901\du}}
\pgfpathlineto{\pgfpoint{-19.307665\du}{0.191901\du}}
\pgfpathlineto{\pgfpoint{-19.469841\du}{0.148287\du}}
\pgfpathlineto{\pgfpoint{-19.177322\du}{0.028722\du}}
\pgfusepath{fill}
\pgfsetbuttcap
\pgfsetmiterjoin
\pgfsetdash{}{0pt}
\definecolor{dialinecolor}{rgb}{1.000000, 1.000000, 1.000000}
\pgfsetfillcolor{dialinecolor}
\pgfpathmoveto{\pgfpoint{-19.729524\du}{-0.210657\du}}
\pgfpathlineto{\pgfpoint{-19.599683\du}{-0.253520\du}}
\pgfpathlineto{\pgfpoint{-19.285607\du}{-0.133956\du}}
\pgfpathlineto{\pgfpoint{-19.144987\du}{-0.167043\du}}
\pgfpathlineto{\pgfpoint{-19.220436\du}{-0.057755\du}}
\pgfpathlineto{\pgfpoint{-19.588904\du}{-0.057755\du}}
\pgfpathlineto{\pgfpoint{-19.437005\du}{-0.090592\du}}
\pgfpathlineto{\pgfpoint{-19.729524\du}{-0.210657\du}}
\pgfusepath{fill}
\pgfsetbuttcap
\pgfsetmiterjoin
\pgfsetdash{}{0pt}
\definecolor{dialinecolor}{rgb}{1.000000, 1.000000, 1.000000}
\pgfsetfillcolor{dialinecolor}
\pgfpathmoveto{\pgfpoint{-18.560701\du}{0.159566\du}}
\pgfpathlineto{\pgfpoint{-18.690041\du}{0.202930\du}}
\pgfpathlineto{\pgfpoint{-18.993088\du}{0.072337\du}}
\pgfpathlineto{\pgfpoint{-19.144987\du}{0.115701\du}}
\pgfpathlineto{\pgfpoint{-19.069288\du}{0.007416\du}}
\pgfpathlineto{\pgfpoint{-18.690041\du}{0.007416\du}}
\pgfpathlineto{\pgfpoint{-18.852468\du}{0.039751\du}}
\pgfpathlineto{\pgfpoint{-18.560701\du}{0.159566\du}}
\pgfusepath{fill}
\pgfsetbuttcap
\pgfsetmiterjoin
\pgfsetdash{}{0pt}
\definecolor{dialinecolor}{rgb}{0.666667, 0.901961, 1.000000}
\pgfsetstrokecolor{dialinecolor}
\pgfpathmoveto{\pgfpoint{-20.043349\du}{-0.036199\du}}
\pgfpathlineto{\pgfpoint{-20.043349\du}{0.398194\du}}
\pgfusepath{stroke}
\pgfsetbuttcap
\pgfsetmiterjoin
\pgfsetdash{}{0pt}
\definecolor{dialinecolor}{rgb}{0.666667, 0.901961, 1.000000}
\pgfsetstrokecolor{dialinecolor}
\pgfpathmoveto{\pgfpoint{-18.268182\du}{-0.036199\du}}
\pgfpathlineto{\pgfpoint{-18.268182\du}{0.398194\du}}
\pgfusepath{stroke}
\pgfsetlinewidth{0.100000\du}
\pgfsetdash{}{0pt}
\pgfsetdash{}{0pt}
\pgfsetbuttcap
\pgfsetmiterjoin
\pgfsetlinewidth{0.001000\du}
\pgfsetbuttcap
\pgfsetmiterjoin
\pgfsetdash{}{0pt}
\definecolor{dialinecolor}{rgb}{0.000000, 0.470588, 0.666667}
\pgfsetfillcolor{dialinecolor}
\pgfpathmoveto{\pgfpoint{-22.153366\du}{3.385770\du}}
\pgfpathlineto{\pgfpoint{-22.154864\du}{3.368094\du}}
\pgfpathlineto{\pgfpoint{-22.158759\du}{3.350118\du}}
\pgfpathlineto{\pgfpoint{-22.165350\du}{3.332441\du}}
\pgfpathlineto{\pgfpoint{-22.174638\du}{3.315065\du}}
\pgfpathlineto{\pgfpoint{-22.186622\du}{3.297388\du}}
\pgfpathlineto{\pgfpoint{-22.201003\du}{3.280011\du}}
\pgfpathlineto{\pgfpoint{-22.218080\du}{3.263234\du}}
\pgfpathlineto{\pgfpoint{-22.237254\du}{3.246756\du}}
\pgfpathlineto{\pgfpoint{-22.259425\du}{3.230577\du}}
\pgfpathlineto{\pgfpoint{-22.283692\du}{3.214998\du}}
\pgfpathlineto{\pgfpoint{-22.310357\du}{3.199419\du}}
\pgfpathlineto{\pgfpoint{-22.339418\du}{3.184439\du}}
\pgfpathlineto{\pgfpoint{-22.370577\du}{3.170058\du}}
\pgfpathlineto{\pgfpoint{-22.403832\du}{3.155977\du}}
\pgfpathlineto{\pgfpoint{-22.439185\du}{3.142794\du}}
\pgfpathlineto{\pgfpoint{-22.476036\du}{3.130211\du}}
\pgfpathlineto{\pgfpoint{-22.514984\du}{3.118227\du}}
\pgfpathlineto{\pgfpoint{-22.556029\du}{3.107142\du}}
\pgfpathlineto{\pgfpoint{-22.598273\du}{3.096056\du}}
\pgfpathlineto{\pgfpoint{-22.642015\du}{3.086769\du}}
\pgfpathlineto{\pgfpoint{-22.687255\du}{3.077481\du}}
\pgfpathlineto{\pgfpoint{-22.733992\du}{3.069392\du}}
\pgfpathlineto{\pgfpoint{-22.781629\du}{3.061602\du}}
\pgfpathlineto{\pgfpoint{-22.830763\du}{3.055011\du}}
\pgfpathlineto{\pgfpoint{-22.880197\du}{3.049319\du}}
\pgfpathlineto{\pgfpoint{-22.931130\du}{3.044525\du}}
\pgfpathlineto{\pgfpoint{-22.982361\du}{3.040331\du}}
\pgfpathlineto{\pgfpoint{-23.033893\du}{3.037035\du}}
\pgfpathlineto{\pgfpoint{-23.086023\du}{3.034638\du}}
\pgfpathlineto{\pgfpoint{-23.138453\du}{3.033739\du}}
\pgfpathlineto{\pgfpoint{-23.191483\du}{3.032841\du}}
\pgfpathlineto{\pgfpoint{-23.191483\du}{3.032841\du}}
\pgfpathlineto{\pgfpoint{-23.243913\du}{3.033739\du}}
\pgfpathlineto{\pgfpoint{-23.296043\du}{3.034638\du}}
\pgfpathlineto{\pgfpoint{-23.348174\du}{3.037035\du}}
\pgfpathlineto{\pgfpoint{-23.399705\du}{3.040331\du}}
\pgfpathlineto{\pgfpoint{-23.450937\du}{3.044525\du}}
\pgfpathlineto{\pgfpoint{-23.501869\du}{3.049319\du}}
\pgfpathlineto{\pgfpoint{-23.551303\du}{3.055011\du}}
\pgfpathlineto{\pgfpoint{-23.600437\du}{3.061602\du}}
\pgfpathlineto{\pgfpoint{-23.648074\du}{3.069392\du}}
\pgfpathlineto{\pgfpoint{-23.694812\du}{3.077481\du}}
\pgfpathlineto{\pgfpoint{-23.740351\du}{3.086769\du}}
\pgfpathlineto{\pgfpoint{-23.783793\du}{3.096056\du}}
\pgfpathlineto{\pgfpoint{-23.826037\du}{3.107142\du}}
\pgfpathlineto{\pgfpoint{-23.867082\du}{3.118227\du}}
\pgfpathlineto{\pgfpoint{-23.906030\du}{3.130211\du}}
\pgfpathlineto{\pgfpoint{-23.942881\du}{3.142794\du}}
\pgfpathlineto{\pgfpoint{-23.978234\du}{3.155977\du}}
\pgfpathlineto{\pgfpoint{-24.011490\du}{3.170058\du}}
\pgfpathlineto{\pgfpoint{-24.042648\du}{3.184439\du}}
\pgfpathlineto{\pgfpoint{-24.072009\du}{3.199419\du}}
\pgfpathlineto{\pgfpoint{-24.098374\du}{3.214998\du}}
\pgfpathlineto{\pgfpoint{-24.122941\du}{3.230577\du}}
\pgfpathlineto{\pgfpoint{-24.144812\du}{3.246756\du}}
\pgfpathlineto{\pgfpoint{-24.163986\du}{3.263234\du}}
\pgfpathlineto{\pgfpoint{-24.181063\du}{3.280011\du}}
\pgfpathlineto{\pgfpoint{-24.195444\du}{3.297388\du}}
\pgfpathlineto{\pgfpoint{-24.207428\du}{3.315065\du}}
\pgfpathlineto{\pgfpoint{-24.216716\du}{3.332441\du}}
\pgfpathlineto{\pgfpoint{-24.223607\du}{3.350118\du}}
\pgfpathlineto{\pgfpoint{-24.227202\du}{3.368094\du}}
\pgfpathlineto{\pgfpoint{-24.228700\du}{3.385770\du}}
\pgfpathlineto{\pgfpoint{-24.228700\du}{3.385770\du}}
\pgfpathlineto{\pgfpoint{-24.227202\du}{3.404046\du}}
\pgfpathlineto{\pgfpoint{-24.223607\du}{3.421722\du}}
\pgfpathlineto{\pgfpoint{-24.216716\du}{3.439399\du}}
\pgfpathlineto{\pgfpoint{-24.207428\du}{3.457075\du}}
\pgfpathlineto{\pgfpoint{-24.195444\du}{3.474452\du}}
\pgfpathlineto{\pgfpoint{-24.181063\du}{3.491829\du}}
\pgfpathlineto{\pgfpoint{-24.163986\du}{3.508607\du}}
\pgfpathlineto{\pgfpoint{-24.144812\du}{3.525085\du}}
\pgfpathlineto{\pgfpoint{-24.122941\du}{3.541263\du}}
\pgfpathlineto{\pgfpoint{-24.098374\du}{3.557142\du}}
\pgfpathlineto{\pgfpoint{-24.072009\du}{3.572122\du}}
\pgfpathlineto{\pgfpoint{-24.042648\du}{3.587701\du}}
\pgfpathlineto{\pgfpoint{-24.011490\du}{3.602082\du}}
\pgfpathlineto{\pgfpoint{-23.978234\du}{3.615864\du}}
\pgfpathlineto{\pgfpoint{-23.942881\du}{3.628746\du}}
\pgfpathlineto{\pgfpoint{-23.906030\du}{3.641629\du}}
\pgfpathlineto{\pgfpoint{-23.867082\du}{3.653613\du}}
\pgfpathlineto{\pgfpoint{-23.826037\du}{3.664998\du}}
\pgfpathlineto{\pgfpoint{-23.783793\du}{3.675484\du}}
\pgfpathlineto{\pgfpoint{-23.740351\du}{3.685371\du}}
\pgfpathlineto{\pgfpoint{-23.694812\du}{3.694359\du}}
\pgfpathlineto{\pgfpoint{-23.648074\du}{3.702448\du}}
\pgfpathlineto{\pgfpoint{-23.600437\du}{3.710238\du}}
\pgfpathlineto{\pgfpoint{-23.551303\du}{3.716829\du}}
\pgfpathlineto{\pgfpoint{-23.501869\du}{3.722521\du}}
\pgfpathlineto{\pgfpoint{-23.450937\du}{3.727615\du}}
\pgfpathlineto{\pgfpoint{-23.399705\du}{3.731210\du}}
\pgfpathlineto{\pgfpoint{-23.348174\du}{3.734505\du}}
\pgfpathlineto{\pgfpoint{-23.296043\du}{3.736902\du}}
\pgfpathlineto{\pgfpoint{-23.243913\du}{3.738400\du}}
\pgfpathlineto{\pgfpoint{-23.191483\du}{3.738700\du}}
\pgfpathlineto{\pgfpoint{-23.191483\du}{3.738700\du}}
\pgfpathlineto{\pgfpoint{-23.138453\du}{3.738400\du}}
\pgfpathlineto{\pgfpoint{-23.086023\du}{3.736902\du}}
\pgfpathlineto{\pgfpoint{-23.033893\du}{3.734505\du}}
\pgfpathlineto{\pgfpoint{-22.982361\du}{3.731210\du}}
\pgfpathlineto{\pgfpoint{-22.931130\du}{3.727615\du}}
\pgfpathlineto{\pgfpoint{-22.880197\du}{3.722521\du}}
\pgfpathlineto{\pgfpoint{-22.830763\du}{3.716829\du}}
\pgfpathlineto{\pgfpoint{-22.781629\du}{3.710238\du}}
\pgfpathlineto{\pgfpoint{-22.733992\du}{3.702448\du}}
\pgfpathlineto{\pgfpoint{-22.687255\du}{3.694359\du}}
\pgfpathlineto{\pgfpoint{-22.642015\du}{3.685371\du}}
\pgfpathlineto{\pgfpoint{-22.598273\du}{3.675484\du}}
\pgfpathlineto{\pgfpoint{-22.556029\du}{3.664998\du}}
\pgfpathlineto{\pgfpoint{-22.514984\du}{3.653613\du}}
\pgfpathlineto{\pgfpoint{-22.476036\du}{3.641629\du}}
\pgfpathlineto{\pgfpoint{-22.439185\du}{3.628746\du}}
\pgfpathlineto{\pgfpoint{-22.403832\du}{3.615864\du}}
\pgfpathlineto{\pgfpoint{-22.370577\du}{3.602082\du}}
\pgfpathlineto{\pgfpoint{-22.339418\du}{3.587701\du}}
\pgfpathlineto{\pgfpoint{-22.310357\du}{3.572122\du}}
\pgfpathlineto{\pgfpoint{-22.283692\du}{3.557142\du}}
\pgfpathlineto{\pgfpoint{-22.259425\du}{3.541263\du}}
\pgfpathlineto{\pgfpoint{-22.237254\du}{3.525085\du}}
\pgfpathlineto{\pgfpoint{-22.218080\du}{3.508607\du}}
\pgfpathlineto{\pgfpoint{-22.201003\du}{3.491829\du}}
\pgfpathlineto{\pgfpoint{-22.186622\du}{3.474452\du}}
\pgfpathlineto{\pgfpoint{-22.174638\du}{3.457075\du}}
\pgfpathlineto{\pgfpoint{-22.165350\du}{3.439399\du}}
\pgfpathlineto{\pgfpoint{-22.158759\du}{3.421722\du}}
\pgfpathlineto{\pgfpoint{-22.154864\du}{3.404046\du}}
\pgfpathlineto{\pgfpoint{-22.153366\du}{3.385770\du}}
\pgfusepath{fill}
\pgfsetbuttcap
\pgfsetmiterjoin
\pgfsetdash{}{0pt}
\definecolor{dialinecolor}{rgb}{0.666667, 0.901961, 1.000000}
\pgfsetstrokecolor{dialinecolor}
\pgfpathmoveto{\pgfpoint{-22.165350\du}{3.380377\du}}
\pgfpathlineto{\pgfpoint{-22.166848\du}{3.362401\du}}
\pgfpathlineto{\pgfpoint{-22.170743\du}{3.345324\du}}
\pgfpathlineto{\pgfpoint{-22.177334\du}{3.327648\du}}
\pgfpathlineto{\pgfpoint{-22.186322\du}{3.310571\du}}
\pgfpathlineto{\pgfpoint{-22.198306\du}{3.293194\du}}
\pgfpathlineto{\pgfpoint{-22.212987\du}{3.276416\du}}
\pgfpathlineto{\pgfpoint{-22.229165\du}{3.259938\du}}
\pgfpathlineto{\pgfpoint{-22.248939\du}{3.243460\du}}
\pgfpathlineto{\pgfpoint{-22.270810\du}{3.227581\du}}
\pgfpathlineto{\pgfpoint{-22.295077\du}{3.211702\du}}
\pgfpathlineto{\pgfpoint{-22.321442\du}{3.197022\du}}
\pgfpathlineto{\pgfpoint{-22.350503\du}{3.182341\du}}
\pgfpathlineto{\pgfpoint{-22.381063\du}{3.167661\du}}
\pgfpathlineto{\pgfpoint{-22.414019\du}{3.154179\du}}
\pgfpathlineto{\pgfpoint{-22.449372\du}{3.140997\du}}
\pgfpathlineto{\pgfpoint{-22.486522\du}{3.129013\du}}
\pgfpathlineto{\pgfpoint{-22.525171\du}{3.117028\du}}
\pgfpathlineto{\pgfpoint{-22.565617\du}{3.105943\du}}
\pgfpathlineto{\pgfpoint{-22.607261\du}{3.095757\du}}
\pgfpathlineto{\pgfpoint{-22.651303\du}{3.085870\du}}
\pgfpathlineto{\pgfpoint{-22.696542\du}{3.077182\du}}
\pgfpathlineto{\pgfpoint{-22.742381\du}{3.069092\du}}
\pgfpathlineto{\pgfpoint{-22.790018\du}{3.061602\du}}
\pgfpathlineto{\pgfpoint{-22.838553\du}{3.055011\du}}
\pgfpathlineto{\pgfpoint{-22.887987\du}{3.049319\du}}
\pgfpathlineto{\pgfpoint{-22.938620\du}{3.044525\du}}
\pgfpathlineto{\pgfpoint{-22.989252\du}{3.040331\du}}
\pgfpathlineto{\pgfpoint{-23.040484\du}{3.037934\du}}
\pgfpathlineto{\pgfpoint{-23.092614\du}{3.035237\du}}
\pgfpathlineto{\pgfpoint{-23.144745\du}{3.033739\du}}
\pgfpathlineto{\pgfpoint{-23.197175\du}{3.033739\du}}
\pgfpathlineto{\pgfpoint{-23.197175\du}{3.033739\du}}
\pgfpathlineto{\pgfpoint{-23.249305\du}{3.033739\du}}
\pgfpathlineto{\pgfpoint{-23.301736\du}{3.035237\du}}
\pgfpathlineto{\pgfpoint{-23.353566\du}{3.037934\du}}
\pgfpathlineto{\pgfpoint{-23.404798\du}{3.040331\du}}
\pgfpathlineto{\pgfpoint{-23.455730\du}{3.044525\du}}
\pgfpathlineto{\pgfpoint{-23.506063\du}{3.049319\du}}
\pgfpathlineto{\pgfpoint{-23.555497\du}{3.055011\du}}
\pgfpathlineto{\pgfpoint{-23.604332\du}{3.061602\du}}
\pgfpathlineto{\pgfpoint{-23.651669\du}{3.069092\du}}
\pgfpathlineto{\pgfpoint{-23.697508\du}{3.077182\du}}
\pgfpathlineto{\pgfpoint{-23.743047\du}{3.085870\du}}
\pgfpathlineto{\pgfpoint{-23.786489\du}{3.095757\du}}
\pgfpathlineto{\pgfpoint{-23.828134\du}{3.105943\du}}
\pgfpathlineto{\pgfpoint{-23.868880\du}{3.117028\du}}
\pgfpathlineto{\pgfpoint{-23.908127\du}{3.129013\du}}
\pgfpathlineto{\pgfpoint{-23.945278\du}{3.140997\du}}
\pgfpathlineto{\pgfpoint{-23.980031\du}{3.154179\du}}
\pgfpathlineto{\pgfpoint{-24.013287\du}{3.167661\du}}
\pgfpathlineto{\pgfpoint{-24.044146\du}{3.182341\du}}
\pgfpathlineto{\pgfpoint{-24.072608\du}{3.197022\du}}
\pgfpathlineto{\pgfpoint{-24.098673\du}{3.211702\du}}
\pgfpathlineto{\pgfpoint{-24.123241\du}{3.227581\du}}
\pgfpathlineto{\pgfpoint{-24.145411\du}{3.243460\du}}
\pgfpathlineto{\pgfpoint{-24.164585\du}{3.259938\du}}
\pgfpathlineto{\pgfpoint{-24.181363\du}{3.276416\du}}
\pgfpathlineto{\pgfpoint{-24.196343\du}{3.293194\du}}
\pgfpathlineto{\pgfpoint{-24.207428\du}{3.310571\du}}
\pgfpathlineto{\pgfpoint{-24.217016\du}{3.327648\du}}
\pgfpathlineto{\pgfpoint{-24.223607\du}{3.345324\du}}
\pgfpathlineto{\pgfpoint{-24.227202\du}{3.362401\du}}
\pgfpathlineto{\pgfpoint{-24.228700\du}{3.380377\du}}
\pgfpathlineto{\pgfpoint{-24.228700\du}{3.380377\du}}
\pgfpathlineto{\pgfpoint{-24.227202\du}{3.398054\du}}
\pgfpathlineto{\pgfpoint{-24.223607\du}{3.415431\du}}
\pgfpathlineto{\pgfpoint{-24.217016\du}{3.432808\du}}
\pgfpathlineto{\pgfpoint{-24.207428\du}{3.449885\du}}
\pgfpathlineto{\pgfpoint{-24.196343\du}{3.467262\du}}
\pgfpathlineto{\pgfpoint{-24.181363\du}{3.484039\du}}
\pgfpathlineto{\pgfpoint{-24.164585\du}{3.500517\du}}
\pgfpathlineto{\pgfpoint{-24.145411\du}{3.516396\du}}
\pgfpathlineto{\pgfpoint{-24.123241\du}{3.532874\du}}
\pgfpathlineto{\pgfpoint{-24.098673\du}{3.548753\du}}
\pgfpathlineto{\pgfpoint{-24.072608\du}{3.563733\du}}
\pgfpathlineto{\pgfpoint{-24.044146\du}{3.578114\du}}
\pgfpathlineto{\pgfpoint{-24.013287\du}{3.592495\du}}
\pgfpathlineto{\pgfpoint{-23.980031\du}{3.606276\du}}
\pgfpathlineto{\pgfpoint{-23.945278\du}{3.619159\du}}
\pgfpathlineto{\pgfpoint{-23.908127\du}{3.631443\du}}
\pgfpathlineto{\pgfpoint{-23.868880\du}{3.642828\du}}
\pgfpathlineto{\pgfpoint{-23.828134\du}{3.654512\du}}
\pgfpathlineto{\pgfpoint{-23.786489\du}{3.664698\du}}
\pgfpathlineto{\pgfpoint{-23.743047\du}{3.674286\du}}
\pgfpathlineto{\pgfpoint{-23.697508\du}{3.683573\du}}
\pgfpathlineto{\pgfpoint{-23.651669\du}{3.691363\du}}
\pgfpathlineto{\pgfpoint{-23.604332\du}{3.698553\du}}
\pgfpathlineto{\pgfpoint{-23.555497\du}{3.705444\du}}
\pgfpathlineto{\pgfpoint{-23.506063\du}{3.710837\du}}
\pgfpathlineto{\pgfpoint{-23.455730\du}{3.715331\du}}
\pgfpathlineto{\pgfpoint{-23.404798\du}{3.719825\du}}
\pgfpathlineto{\pgfpoint{-23.353566\du}{3.722521\du}}
\pgfpathlineto{\pgfpoint{-23.301736\du}{3.724918\du}}
\pgfpathlineto{\pgfpoint{-23.249305\du}{3.726117\du}}
\pgfpathlineto{\pgfpoint{-23.197175\du}{3.726716\du}}
\pgfpathlineto{\pgfpoint{-23.197175\du}{3.726716\du}}
\pgfpathlineto{\pgfpoint{-23.144745\du}{3.726117\du}}
\pgfpathlineto{\pgfpoint{-23.092614\du}{3.724918\du}}
\pgfpathlineto{\pgfpoint{-23.040484\du}{3.722521\du}}
\pgfpathlineto{\pgfpoint{-22.989252\du}{3.719825\du}}
\pgfpathlineto{\pgfpoint{-22.938620\du}{3.715331\du}}
\pgfpathlineto{\pgfpoint{-22.887987\du}{3.710837\du}}
\pgfpathlineto{\pgfpoint{-22.838553\du}{3.705444\du}}
\pgfpathlineto{\pgfpoint{-22.790018\du}{3.698553\du}}
\pgfpathlineto{\pgfpoint{-22.742381\du}{3.691363\du}}
\pgfpathlineto{\pgfpoint{-22.696542\du}{3.683573\du}}
\pgfpathlineto{\pgfpoint{-22.651303\du}{3.674286\du}}
\pgfpathlineto{\pgfpoint{-22.607261\du}{3.664698\du}}
\pgfpathlineto{\pgfpoint{-22.565617\du}{3.654512\du}}
\pgfpathlineto{\pgfpoint{-22.525171\du}{3.642828\du}}
\pgfpathlineto{\pgfpoint{-22.486522\du}{3.631443\du}}
\pgfpathlineto{\pgfpoint{-22.449372\du}{3.619159\du}}
\pgfpathlineto{\pgfpoint{-22.414019\du}{3.606276\du}}
\pgfpathlineto{\pgfpoint{-22.381063\du}{3.592495\du}}
\pgfpathlineto{\pgfpoint{-22.350503\du}{3.578114\du}}
\pgfpathlineto{\pgfpoint{-22.321442\du}{3.563733\du}}
\pgfpathlineto{\pgfpoint{-22.295077\du}{3.548753\du}}
\pgfpathlineto{\pgfpoint{-22.270810\du}{3.532874\du}}
\pgfpathlineto{\pgfpoint{-22.248939\du}{3.516396\du}}
\pgfpathlineto{\pgfpoint{-22.229165\du}{3.500517\du}}
\pgfpathlineto{\pgfpoint{-22.212987\du}{3.484039\du}}
\pgfpathlineto{\pgfpoint{-22.198306\du}{3.467262\du}}
\pgfpathlineto{\pgfpoint{-22.186322\du}{3.449885\du}}
\pgfpathlineto{\pgfpoint{-22.177334\du}{3.432808\du}}
\pgfpathlineto{\pgfpoint{-22.170743\du}{3.415431\du}}
\pgfpathlineto{\pgfpoint{-22.166848\du}{3.398054\du}}
\pgfpathlineto{\pgfpoint{-22.165350\du}{3.380377\du}}
\pgfusepath{stroke}
\pgfsetbuttcap
\pgfsetmiterjoin
\pgfsetdash{}{0pt}
\definecolor{dialinecolor}{rgb}{0.000000, 0.470588, 0.666667}
\pgfsetfillcolor{dialinecolor}
\pgfpathmoveto{\pgfpoint{-24.228700\du}{2.897421\du}}
\pgfpathlineto{\pgfpoint{-24.228700\du}{3.392361\du}}
\pgfpathlineto{\pgfpoint{-22.165350\du}{3.392361\du}}
\pgfpathlineto{\pgfpoint{-22.165350\du}{2.897421\du}}
\pgfpathlineto{\pgfpoint{-24.228700\du}{2.897421\du}}
\pgfusepath{fill}
\pgfsetbuttcap
\pgfsetmiterjoin
\pgfsetdash{}{0pt}
\definecolor{dialinecolor}{rgb}{0.000000, 0.705882, 1.000000}
\pgfsetfillcolor{dialinecolor}
\pgfpathmoveto{\pgfpoint{-22.153366\du}{2.890530\du}}
\pgfpathlineto{\pgfpoint{-22.154864\du}{2.873154\du}}
\pgfpathlineto{\pgfpoint{-22.158759\du}{2.855178\du}}
\pgfpathlineto{\pgfpoint{-22.165350\du}{2.837501\du}}
\pgfpathlineto{\pgfpoint{-22.174638\du}{2.819825\du}}
\pgfpathlineto{\pgfpoint{-22.186622\du}{2.802148\du}}
\pgfpathlineto{\pgfpoint{-22.201003\du}{2.785071\du}}
\pgfpathlineto{\pgfpoint{-22.218080\du}{2.768293\du}}
\pgfpathlineto{\pgfpoint{-22.237254\du}{2.751516\du}}
\pgfpathlineto{\pgfpoint{-22.259425\du}{2.735637\du}}
\pgfpathlineto{\pgfpoint{-22.283692\du}{2.719758\du}}
\pgfpathlineto{\pgfpoint{-22.310357\du}{2.704478\du}}
\pgfpathlineto{\pgfpoint{-22.339418\du}{2.688899\du}}
\pgfpathlineto{\pgfpoint{-22.370577\du}{2.674818\du}}
\pgfpathlineto{\pgfpoint{-22.403832\du}{2.660737\du}}
\pgfpathlineto{\pgfpoint{-22.439185\du}{2.647554\du}}
\pgfpathlineto{\pgfpoint{-22.476036\du}{2.635271\du}}
\pgfpathlineto{\pgfpoint{-22.514984\du}{2.623287\du}}
\pgfpathlineto{\pgfpoint{-22.556029\du}{2.611902\du}}
\pgfpathlineto{\pgfpoint{-22.598273\du}{2.601416\du}}
\pgfpathlineto{\pgfpoint{-22.642015\du}{2.591529\du}}
\pgfpathlineto{\pgfpoint{-22.687255\du}{2.582241\du}}
\pgfpathlineto{\pgfpoint{-22.733992\du}{2.574452\du}}
\pgfpathlineto{\pgfpoint{-22.781629\du}{2.566362\du}}
\pgfpathlineto{\pgfpoint{-22.830763\du}{2.559771\du}}
\pgfpathlineto{\pgfpoint{-22.880197\du}{2.554079\du}}
\pgfpathlineto{\pgfpoint{-22.931130\du}{2.549285\du}}
\pgfpathlineto{\pgfpoint{-22.982361\du}{2.545690\du}}
\pgfpathlineto{\pgfpoint{-23.033893\du}{2.542394\du}}
\pgfpathlineto{\pgfpoint{-23.086023\du}{2.539998\du}}
\pgfpathlineto{\pgfpoint{-23.138453\du}{2.538200\du}}
\pgfpathlineto{\pgfpoint{-23.191483\du}{2.538200\du}}
\pgfpathlineto{\pgfpoint{-23.191483\du}{2.538200\du}}
\pgfpathlineto{\pgfpoint{-23.243913\du}{2.538200\du}}
\pgfpathlineto{\pgfpoint{-23.296043\du}{2.539998\du}}
\pgfpathlineto{\pgfpoint{-23.348174\du}{2.542394\du}}
\pgfpathlineto{\pgfpoint{-23.399705\du}{2.545690\du}}
\pgfpathlineto{\pgfpoint{-23.450937\du}{2.549285\du}}
\pgfpathlineto{\pgfpoint{-23.501869\du}{2.554079\du}}
\pgfpathlineto{\pgfpoint{-23.551303\du}{2.559771\du}}
\pgfpathlineto{\pgfpoint{-23.600437\du}{2.566362\du}}
\pgfpathlineto{\pgfpoint{-23.648074\du}{2.574452\du}}
\pgfpathlineto{\pgfpoint{-23.694812\du}{2.582241\du}}
\pgfpathlineto{\pgfpoint{-23.740351\du}{2.591529\du}}
\pgfpathlineto{\pgfpoint{-23.783793\du}{2.601416\du}}
\pgfpathlineto{\pgfpoint{-23.826037\du}{2.611902\du}}
\pgfpathlineto{\pgfpoint{-23.867082\du}{2.623287\du}}
\pgfpathlineto{\pgfpoint{-23.906030\du}{2.635271\du}}
\pgfpathlineto{\pgfpoint{-23.942881\du}{2.647554\du}}
\pgfpathlineto{\pgfpoint{-23.978234\du}{2.660737\du}}
\pgfpathlineto{\pgfpoint{-24.011490\du}{2.674818\du}}
\pgfpathlineto{\pgfpoint{-24.042648\du}{2.688899\du}}
\pgfpathlineto{\pgfpoint{-24.072009\du}{2.704478\du}}
\pgfpathlineto{\pgfpoint{-24.098374\du}{2.719758\du}}
\pgfpathlineto{\pgfpoint{-24.122941\du}{2.735637\du}}
\pgfpathlineto{\pgfpoint{-24.144812\du}{2.751516\du}}
\pgfpathlineto{\pgfpoint{-24.163986\du}{2.768293\du}}
\pgfpathlineto{\pgfpoint{-24.181063\du}{2.785071\du}}
\pgfpathlineto{\pgfpoint{-24.195444\du}{2.802148\du}}
\pgfpathlineto{\pgfpoint{-24.207428\du}{2.819825\du}}
\pgfpathlineto{\pgfpoint{-24.216716\du}{2.837501\du}}
\pgfpathlineto{\pgfpoint{-24.223607\du}{2.855178\du}}
\pgfpathlineto{\pgfpoint{-24.227202\du}{2.873154\du}}
\pgfpathlineto{\pgfpoint{-24.228700\du}{2.890530\du}}
\pgfpathlineto{\pgfpoint{-24.228700\du}{2.890530\du}}
\pgfpathlineto{\pgfpoint{-24.227202\du}{2.908806\du}}
\pgfpathlineto{\pgfpoint{-24.223607\du}{2.926782\du}}
\pgfpathlineto{\pgfpoint{-24.216716\du}{2.944458\du}}
\pgfpathlineto{\pgfpoint{-24.207428\du}{2.961835\du}}
\pgfpathlineto{\pgfpoint{-24.195444\du}{2.979512\du}}
\pgfpathlineto{\pgfpoint{-24.181063\du}{2.996289\du}}
\pgfpathlineto{\pgfpoint{-24.163986\du}{3.013666\du}}
\pgfpathlineto{\pgfpoint{-24.144812\du}{3.030144\du}}
\pgfpathlineto{\pgfpoint{-24.122941\du}{3.046323\du}}
\pgfpathlineto{\pgfpoint{-24.098374\du}{3.062202\du}}
\pgfpathlineto{\pgfpoint{-24.072009\du}{3.077481\du}}
\pgfpathlineto{\pgfpoint{-24.042648\du}{3.092761\du}}
\pgfpathlineto{\pgfpoint{-24.011490\du}{3.107142\du}}
\pgfpathlineto{\pgfpoint{-23.978234\du}{3.120324\du}}
\pgfpathlineto{\pgfpoint{-23.942881\du}{3.134106\du}}
\pgfpathlineto{\pgfpoint{-23.906030\du}{3.146689\du}}
\pgfpathlineto{\pgfpoint{-23.867082\du}{3.158673\du}}
\pgfpathlineto{\pgfpoint{-23.826037\du}{3.170058\du}}
\pgfpathlineto{\pgfpoint{-23.783793\du}{3.180544\du}}
\pgfpathlineto{\pgfpoint{-23.740351\du}{3.190131\du}}
\pgfpathlineto{\pgfpoint{-23.694812\du}{3.199419\du}}
\pgfpathlineto{\pgfpoint{-23.648074\du}{3.207508\du}}
\pgfpathlineto{\pgfpoint{-23.600437\du}{3.215297\du}}
\pgfpathlineto{\pgfpoint{-23.551303\du}{3.221889\du}}
\pgfpathlineto{\pgfpoint{-23.501869\du}{3.227581\du}}
\pgfpathlineto{\pgfpoint{-23.450937\du}{3.232375\du}}
\pgfpathlineto{\pgfpoint{-23.399705\du}{3.235970\du}}
\pgfpathlineto{\pgfpoint{-23.348174\du}{3.239266\du}}
\pgfpathlineto{\pgfpoint{-23.296043\du}{3.241962\du}}
\pgfpathlineto{\pgfpoint{-23.243913\du}{3.243460\du}}
\pgfpathlineto{\pgfpoint{-23.191483\du}{3.243460\du}}
\pgfpathlineto{\pgfpoint{-23.191483\du}{3.243460\du}}
\pgfpathlineto{\pgfpoint{-23.138453\du}{3.243460\du}}
\pgfpathlineto{\pgfpoint{-23.086023\du}{3.241962\du}}
\pgfpathlineto{\pgfpoint{-23.033893\du}{3.239266\du}}
\pgfpathlineto{\pgfpoint{-22.982361\du}{3.235970\du}}
\pgfpathlineto{\pgfpoint{-22.931130\du}{3.232375\du}}
\pgfpathlineto{\pgfpoint{-22.880197\du}{3.227581\du}}
\pgfpathlineto{\pgfpoint{-22.830763\du}{3.221889\du}}
\pgfpathlineto{\pgfpoint{-22.781629\du}{3.215297\du}}
\pgfpathlineto{\pgfpoint{-22.733992\du}{3.207508\du}}
\pgfpathlineto{\pgfpoint{-22.687255\du}{3.199419\du}}
\pgfpathlineto{\pgfpoint{-22.642015\du}{3.190131\du}}
\pgfpathlineto{\pgfpoint{-22.598273\du}{3.180544\du}}
\pgfpathlineto{\pgfpoint{-22.556029\du}{3.170058\du}}
\pgfpathlineto{\pgfpoint{-22.514984\du}{3.158673\du}}
\pgfpathlineto{\pgfpoint{-22.476036\du}{3.146689\du}}
\pgfpathlineto{\pgfpoint{-22.439185\du}{3.134106\du}}
\pgfpathlineto{\pgfpoint{-22.403832\du}{3.120324\du}}
\pgfpathlineto{\pgfpoint{-22.370577\du}{3.107142\du}}
\pgfpathlineto{\pgfpoint{-22.339418\du}{3.092761\du}}
\pgfpathlineto{\pgfpoint{-22.310357\du}{3.077481\du}}
\pgfpathlineto{\pgfpoint{-22.283692\du}{3.062202\du}}
\pgfpathlineto{\pgfpoint{-22.259425\du}{3.046323\du}}
\pgfpathlineto{\pgfpoint{-22.237254\du}{3.030144\du}}
\pgfpathlineto{\pgfpoint{-22.218080\du}{3.013666\du}}
\pgfpathlineto{\pgfpoint{-22.201003\du}{2.996289\du}}
\pgfpathlineto{\pgfpoint{-22.186622\du}{2.979512\du}}
\pgfpathlineto{\pgfpoint{-22.174638\du}{2.961835\du}}
\pgfpathlineto{\pgfpoint{-22.165350\du}{2.944458\du}}
\pgfpathlineto{\pgfpoint{-22.158759\du}{2.926782\du}}
\pgfpathlineto{\pgfpoint{-22.154864\du}{2.908806\du}}
\pgfpathlineto{\pgfpoint{-22.153366\du}{2.890530\du}}
\pgfusepath{fill}
\pgfsetbuttcap
\pgfsetmiterjoin
\pgfsetdash{}{0pt}
\definecolor{dialinecolor}{rgb}{0.666667, 0.901961, 1.000000}
\pgfsetstrokecolor{dialinecolor}
\pgfpathmoveto{\pgfpoint{-22.165350\du}{2.885437\du}}
\pgfpathlineto{\pgfpoint{-22.166848\du}{2.867461\du}}
\pgfpathlineto{\pgfpoint{-22.170743\du}{2.849785\du}}
\pgfpathlineto{\pgfpoint{-22.177334\du}{2.833007\du}}
\pgfpathlineto{\pgfpoint{-22.186322\du}{2.815331\du}}
\pgfpathlineto{\pgfpoint{-22.198306\du}{2.798253\du}}
\pgfpathlineto{\pgfpoint{-22.212987\du}{2.781476\du}}
\pgfpathlineto{\pgfpoint{-22.229165\du}{2.765297\du}}
\pgfpathlineto{\pgfpoint{-22.248939\du}{2.748520\du}}
\pgfpathlineto{\pgfpoint{-22.270810\du}{2.732940\du}}
\pgfpathlineto{\pgfpoint{-22.295077\du}{2.716762\du}}
\pgfpathlineto{\pgfpoint{-22.321442\du}{2.701782\du}}
\pgfpathlineto{\pgfpoint{-22.350503\du}{2.686802\du}}
\pgfpathlineto{\pgfpoint{-22.381063\du}{2.672721\du}}
\pgfpathlineto{\pgfpoint{-22.414019\du}{2.659538\du}}
\pgfpathlineto{\pgfpoint{-22.449372\du}{2.645757\du}}
\pgfpathlineto{\pgfpoint{-22.486522\du}{2.633773\du}}
\pgfpathlineto{\pgfpoint{-22.525171\du}{2.622088\du}}
\pgfpathlineto{\pgfpoint{-22.565617\du}{2.611003\du}}
\pgfpathlineto{\pgfpoint{-22.607261\du}{2.600817\du}}
\pgfpathlineto{\pgfpoint{-22.651303\du}{2.590630\du}}
\pgfpathlineto{\pgfpoint{-22.696542\du}{2.581942\du}}
\pgfpathlineto{\pgfpoint{-22.742381\du}{2.573852\du}}
\pgfpathlineto{\pgfpoint{-22.790018\du}{2.566362\du}}
\pgfpathlineto{\pgfpoint{-22.838553\du}{2.559771\du}}
\pgfpathlineto{\pgfpoint{-22.887987\du}{2.554079\du}}
\pgfpathlineto{\pgfpoint{-22.938620\du}{2.549884\du}}
\pgfpathlineto{\pgfpoint{-22.989252\du}{2.545690\du}}
\pgfpathlineto{\pgfpoint{-23.040484\du}{2.542694\du}}
\pgfpathlineto{\pgfpoint{-23.092614\du}{2.539998\du}}
\pgfpathlineto{\pgfpoint{-23.144745\du}{2.538799\du}}
\pgfpathlineto{\pgfpoint{-23.197175\du}{2.538200\du}}
\pgfpathlineto{\pgfpoint{-23.197175\du}{2.538200\du}}
\pgfpathlineto{\pgfpoint{-23.249305\du}{2.538799\du}}
\pgfpathlineto{\pgfpoint{-23.301736\du}{2.539998\du}}
\pgfpathlineto{\pgfpoint{-23.353566\du}{2.542694\du}}
\pgfpathlineto{\pgfpoint{-23.404798\du}{2.545690\du}}
\pgfpathlineto{\pgfpoint{-23.455730\du}{2.549884\du}}
\pgfpathlineto{\pgfpoint{-23.506063\du}{2.554079\du}}
\pgfpathlineto{\pgfpoint{-23.555497\du}{2.559771\du}}
\pgfpathlineto{\pgfpoint{-23.604332\du}{2.566362\du}}
\pgfpathlineto{\pgfpoint{-23.651669\du}{2.573852\du}}
\pgfpathlineto{\pgfpoint{-23.697508\du}{2.581942\du}}
\pgfpathlineto{\pgfpoint{-23.743047\du}{2.590630\du}}
\pgfpathlineto{\pgfpoint{-23.786489\du}{2.600817\du}}
\pgfpathlineto{\pgfpoint{-23.828134\du}{2.611003\du}}
\pgfpathlineto{\pgfpoint{-23.868880\du}{2.622088\du}}
\pgfpathlineto{\pgfpoint{-23.908127\du}{2.633773\du}}
\pgfpathlineto{\pgfpoint{-23.945278\du}{2.645757\du}}
\pgfpathlineto{\pgfpoint{-23.980031\du}{2.659538\du}}
\pgfpathlineto{\pgfpoint{-24.013287\du}{2.672721\du}}
\pgfpathlineto{\pgfpoint{-24.044146\du}{2.686802\du}}
\pgfpathlineto{\pgfpoint{-24.072608\du}{2.701782\du}}
\pgfpathlineto{\pgfpoint{-24.098673\du}{2.716762\du}}
\pgfpathlineto{\pgfpoint{-24.123241\du}{2.732940\du}}
\pgfpathlineto{\pgfpoint{-24.145411\du}{2.748520\du}}
\pgfpathlineto{\pgfpoint{-24.164585\du}{2.765297\du}}
\pgfpathlineto{\pgfpoint{-24.181363\du}{2.781476\du}}
\pgfpathlineto{\pgfpoint{-24.196343\du}{2.798253\du}}
\pgfpathlineto{\pgfpoint{-24.207428\du}{2.815331\du}}
\pgfpathlineto{\pgfpoint{-24.217016\du}{2.833007\du}}
\pgfpathlineto{\pgfpoint{-24.223607\du}{2.849785\du}}
\pgfpathlineto{\pgfpoint{-24.227202\du}{2.867461\du}}
\pgfpathlineto{\pgfpoint{-24.228700\du}{2.885437\du}}
\pgfpathlineto{\pgfpoint{-24.228700\du}{2.885437\du}}
\pgfpathlineto{\pgfpoint{-24.227202\du}{2.902814\du}}
\pgfpathlineto{\pgfpoint{-24.223607\du}{2.920490\du}}
\pgfpathlineto{\pgfpoint{-24.217016\du}{2.937268\du}}
\pgfpathlineto{\pgfpoint{-24.207428\du}{2.954645\du}}
\pgfpathlineto{\pgfpoint{-24.196343\du}{2.972022\du}}
\pgfpathlineto{\pgfpoint{-24.181363\du}{2.988799\du}}
\pgfpathlineto{\pgfpoint{-24.164585\du}{3.005577\du}}
\pgfpathlineto{\pgfpoint{-24.145411\du}{3.022055\du}}
\pgfpathlineto{\pgfpoint{-24.123241\du}{3.037934\du}}
\pgfpathlineto{\pgfpoint{-24.098673\du}{3.053513\du}}
\pgfpathlineto{\pgfpoint{-24.072608\du}{3.068793\du}}
\pgfpathlineto{\pgfpoint{-24.044146\du}{3.082874\du}}
\pgfpathlineto{\pgfpoint{-24.013287\du}{3.097554\du}}
\pgfpathlineto{\pgfpoint{-23.980031\du}{3.111036\du}}
\pgfpathlineto{\pgfpoint{-23.945278\du}{3.124219\du}}
\pgfpathlineto{\pgfpoint{-23.908127\du}{3.136503\du}}
\pgfpathlineto{\pgfpoint{-23.868880\du}{3.148187\du}}
\pgfpathlineto{\pgfpoint{-23.828134\du}{3.159272\du}}
\pgfpathlineto{\pgfpoint{-23.786489\du}{3.169459\du}}
\pgfpathlineto{\pgfpoint{-23.743047\du}{3.179345\du}}
\pgfpathlineto{\pgfpoint{-23.697508\du}{3.188633\du}}
\pgfpathlineto{\pgfpoint{-23.651669\du}{3.196123\du}}
\pgfpathlineto{\pgfpoint{-23.604332\du}{3.203913\du}}
\pgfpathlineto{\pgfpoint{-23.555497\du}{3.210204\du}}
\pgfpathlineto{\pgfpoint{-23.506063\du}{3.215897\du}}
\pgfpathlineto{\pgfpoint{-23.455730\du}{3.220990\du}}
\pgfpathlineto{\pgfpoint{-23.404798\du}{3.225184\du}}
\pgfpathlineto{\pgfpoint{-23.353566\du}{3.227581\du}}
\pgfpathlineto{\pgfpoint{-23.301736\du}{3.229978\du}}
\pgfpathlineto{\pgfpoint{-23.249305\du}{3.231476\du}}
\pgfpathlineto{\pgfpoint{-23.197175\du}{3.231776\du}}
\pgfpathlineto{\pgfpoint{-23.197175\du}{3.231776\du}}
\pgfpathlineto{\pgfpoint{-23.144745\du}{3.231476\du}}
\pgfpathlineto{\pgfpoint{-23.092614\du}{3.229978\du}}
\pgfpathlineto{\pgfpoint{-23.040484\du}{3.227581\du}}
\pgfpathlineto{\pgfpoint{-22.989252\du}{3.225184\du}}
\pgfpathlineto{\pgfpoint{-22.938620\du}{3.220990\du}}
\pgfpathlineto{\pgfpoint{-22.887987\du}{3.215897\du}}
\pgfpathlineto{\pgfpoint{-22.838553\du}{3.210204\du}}
\pgfpathlineto{\pgfpoint{-22.790018\du}{3.203913\du}}
\pgfpathlineto{\pgfpoint{-22.742381\du}{3.196123\du}}
\pgfpathlineto{\pgfpoint{-22.696542\du}{3.188633\du}}
\pgfpathlineto{\pgfpoint{-22.651303\du}{3.179345\du}}
\pgfpathlineto{\pgfpoint{-22.607261\du}{3.169459\du}}
\pgfpathlineto{\pgfpoint{-22.565617\du}{3.159272\du}}
\pgfpathlineto{\pgfpoint{-22.525171\du}{3.148187\du}}
\pgfpathlineto{\pgfpoint{-22.486522\du}{3.136503\du}}
\pgfpathlineto{\pgfpoint{-22.449372\du}{3.124219\du}}
\pgfpathlineto{\pgfpoint{-22.414019\du}{3.111036\du}}
\pgfpathlineto{\pgfpoint{-22.381063\du}{3.097554\du}}
\pgfpathlineto{\pgfpoint{-22.350503\du}{3.082874\du}}
\pgfpathlineto{\pgfpoint{-22.321442\du}{3.068793\du}}
\pgfpathlineto{\pgfpoint{-22.295077\du}{3.053513\du}}
\pgfpathlineto{\pgfpoint{-22.270810\du}{3.037934\du}}
\pgfpathlineto{\pgfpoint{-22.248939\du}{3.022055\du}}
\pgfpathlineto{\pgfpoint{-22.229165\du}{3.005577\du}}
\pgfpathlineto{\pgfpoint{-22.212987\du}{2.988799\du}}
\pgfpathlineto{\pgfpoint{-22.198306\du}{2.972022\du}}
\pgfpathlineto{\pgfpoint{-22.186322\du}{2.954645\du}}
\pgfpathlineto{\pgfpoint{-22.177334\du}{2.937268\du}}
\pgfpathlineto{\pgfpoint{-22.170743\du}{2.920490\du}}
\pgfpathlineto{\pgfpoint{-22.166848\du}{2.902814\du}}
\pgfpathlineto{\pgfpoint{-22.165350\du}{2.885437\du}}
\pgfusepath{stroke}
\pgfsetbuttcap
\pgfsetmiterjoin
\pgfsetdash{}{0pt}
\definecolor{dialinecolor}{rgb}{0.000000, 0.000000, 0.000000}
\pgfsetfillcolor{dialinecolor}
\pgfpathmoveto{\pgfpoint{-23.172308\du}{2.810537\du}}
\pgfpathlineto{\pgfpoint{-23.021010\du}{2.860271\du}}
\pgfpathlineto{\pgfpoint{-22.655797\du}{2.711369\du}}
\pgfpathlineto{\pgfpoint{-22.492514\du}{2.761103\du}}
\pgfpathlineto{\pgfpoint{-22.581196\du}{2.637368\du}}
\pgfpathlineto{\pgfpoint{-23.008127\du}{2.637368\du}}
\pgfpathlineto{\pgfpoint{-22.832561\du}{2.674818\du}}
\pgfpathlineto{\pgfpoint{-23.172308\du}{2.810537\du}}
\pgfusepath{fill}
\pgfsetbuttcap
\pgfsetmiterjoin
\pgfsetdash{}{0pt}
\definecolor{dialinecolor}{rgb}{0.000000, 0.000000, 0.000000}
\pgfsetfillcolor{dialinecolor}
\pgfpathmoveto{\pgfpoint{-23.235524\du}{2.947155\du}}
\pgfpathlineto{\pgfpoint{-23.385624\du}{2.897421\du}}
\pgfpathlineto{\pgfpoint{-23.738254\du}{3.046323\du}}
\pgfpathlineto{\pgfpoint{-23.914119\du}{2.996289\du}}
\pgfpathlineto{\pgfpoint{-23.826037\du}{3.132608\du}}
\pgfpathlineto{\pgfpoint{-23.385624\du}{3.132608\du}}
\pgfpathlineto{\pgfpoint{-23.574672\du}{3.082874\du}}
\pgfpathlineto{\pgfpoint{-23.235524\du}{2.947155\du}}
\pgfusepath{fill}
\pgfsetbuttcap
\pgfsetmiterjoin
\pgfsetdash{}{0pt}
\definecolor{dialinecolor}{rgb}{0.000000, 0.000000, 0.000000}
\pgfsetfillcolor{dialinecolor}
\pgfpathmoveto{\pgfpoint{-23.876669\du}{2.674818\du}}
\pgfpathlineto{\pgfpoint{-23.725670\du}{2.625084\du}}
\pgfpathlineto{\pgfpoint{-23.361056\du}{2.761103\du}}
\pgfpathlineto{\pgfpoint{-23.197175\du}{2.724252\du}}
\pgfpathlineto{\pgfpoint{-23.284958\du}{2.847687\du}}
\pgfpathlineto{\pgfpoint{-23.713087\du}{2.847687\du}}
\pgfpathlineto{\pgfpoint{-23.537222\du}{2.810537\du}}
\pgfpathlineto{\pgfpoint{-23.876669\du}{2.674818\du}}
\pgfusepath{fill}
\pgfsetbuttcap
\pgfsetmiterjoin
\pgfsetdash{}{0pt}
\definecolor{dialinecolor}{rgb}{0.000000, 0.000000, 0.000000}
\pgfsetfillcolor{dialinecolor}
\pgfpathmoveto{\pgfpoint{-22.517980\du}{3.095457\du}}
\pgfpathlineto{\pgfpoint{-22.668679\du}{3.144891\du}}
\pgfpathlineto{\pgfpoint{-23.021010\du}{2.996289\du}}
\pgfpathlineto{\pgfpoint{-23.197175\du}{3.046323\du}}
\pgfpathlineto{\pgfpoint{-23.109092\du}{2.922288\du}}
\pgfpathlineto{\pgfpoint{-22.668679\du}{2.922288\du}}
\pgfpathlineto{\pgfpoint{-22.857727\du}{2.958839\du}}
\pgfpathlineto{\pgfpoint{-22.517980\du}{3.095457\du}}
\pgfusepath{fill}
\pgfsetbuttcap
\pgfsetmiterjoin
\pgfsetdash{}{0pt}
\definecolor{dialinecolor}{rgb}{1.000000, 1.000000, 1.000000}
\pgfsetfillcolor{dialinecolor}
\pgfpathmoveto{\pgfpoint{-23.159725\du}{2.822821\du}}
\pgfpathlineto{\pgfpoint{-23.008127\du}{2.872255\du}}
\pgfpathlineto{\pgfpoint{-22.643513\du}{2.724252\du}}
\pgfpathlineto{\pgfpoint{-22.479931\du}{2.773986\du}}
\pgfpathlineto{\pgfpoint{-22.568014\du}{2.649651\du}}
\pgfpathlineto{\pgfpoint{-22.996143\du}{2.649651\du}}
\pgfpathlineto{\pgfpoint{-22.819678\du}{2.686802\du}}
\pgfpathlineto{\pgfpoint{-23.159725\du}{2.822821\du}}
\pgfusepath{fill}
\pgfsetbuttcap
\pgfsetmiterjoin
\pgfsetdash{}{0pt}
\definecolor{dialinecolor}{rgb}{1.000000, 1.000000, 1.000000}
\pgfsetfillcolor{dialinecolor}
\pgfpathmoveto{\pgfpoint{-23.222042\du}{2.958839\du}}
\pgfpathlineto{\pgfpoint{-23.373640\du}{2.909705\du}}
\pgfpathlineto{\pgfpoint{-23.725670\du}{3.058307\du}}
\pgfpathlineto{\pgfpoint{-23.901836\du}{3.008573\du}}
\pgfpathlineto{\pgfpoint{-23.813753\du}{3.144891\du}}
\pgfpathlineto{\pgfpoint{-23.373640\du}{3.144891\du}}
\pgfpathlineto{\pgfpoint{-23.561789\du}{3.095457\du}}
\pgfpathlineto{\pgfpoint{-23.222042\du}{2.958839\du}}
\pgfusepath{fill}
\pgfsetbuttcap
\pgfsetmiterjoin
\pgfsetdash{}{0pt}
\definecolor{dialinecolor}{rgb}{1.000000, 1.000000, 1.000000}
\pgfsetfillcolor{dialinecolor}
\pgfpathmoveto{\pgfpoint{-23.863786\du}{2.686802\du}}
\pgfpathlineto{\pgfpoint{-23.713087\du}{2.637368\du}}
\pgfpathlineto{\pgfpoint{-23.348174\du}{2.773986\du}}
\pgfpathlineto{\pgfpoint{-23.184592\du}{2.736236\du}}
\pgfpathlineto{\pgfpoint{-23.272674\du}{2.860271\du}}
\pgfpathlineto{\pgfpoint{-23.700504\du}{2.860271\du}}
\pgfpathlineto{\pgfpoint{-23.524039\du}{2.822821\du}}
\pgfpathlineto{\pgfpoint{-23.863786\du}{2.686802\du}}
\pgfusepath{fill}
\pgfsetbuttcap
\pgfsetmiterjoin
\pgfsetdash{}{0pt}
\definecolor{dialinecolor}{rgb}{1.000000, 1.000000, 1.000000}
\pgfsetfillcolor{dialinecolor}
\pgfpathmoveto{\pgfpoint{-22.505097\du}{3.107741\du}}
\pgfpathlineto{\pgfpoint{-22.655797\du}{3.157475\du}}
\pgfpathlineto{\pgfpoint{-23.008127\du}{3.008573\du}}
\pgfpathlineto{\pgfpoint{-23.184592\du}{3.058307\du}}
\pgfpathlineto{\pgfpoint{-23.096809\du}{2.934871\du}}
\pgfpathlineto{\pgfpoint{-22.655797\du}{2.934871\du}}
\pgfpathlineto{\pgfpoint{-22.844545\du}{2.971722\du}}
\pgfpathlineto{\pgfpoint{-22.505097\du}{3.107741\du}}
\pgfusepath{fill}
\pgfsetbuttcap
\pgfsetmiterjoin
\pgfsetdash{}{0pt}
\definecolor{dialinecolor}{rgb}{0.666667, 0.901961, 1.000000}
\pgfsetstrokecolor{dialinecolor}
\pgfpathmoveto{\pgfpoint{-24.228700\du}{2.885437\du}}
\pgfpathlineto{\pgfpoint{-24.228700\du}{3.379778\du}}
\pgfusepath{stroke}
\pgfsetbuttcap
\pgfsetmiterjoin
\pgfsetdash{}{0pt}
\definecolor{dialinecolor}{rgb}{0.666667, 0.901961, 1.000000}
\pgfsetstrokecolor{dialinecolor}
\pgfpathmoveto{\pgfpoint{-22.165350\du}{2.885437\du}}
\pgfpathlineto{\pgfpoint{-22.165350\du}{3.379778\du}}
\pgfusepath{stroke}
\pgfsetbuttcap
\pgfsetmiterjoin
\pgfsetdash{}{0pt}
\definecolor{dialinecolor}{rgb}{0.000000, 0.000000, 0.000000}
\pgfsetfillcolor{dialinecolor}
\pgfpathmoveto{\pgfpoint{-23.574672\du}{3.318360\du}}
\pgfpathlineto{\pgfpoint{-23.348174\du}{3.318360\du}}
\pgfpathlineto{\pgfpoint{-23.209758\du}{3.429812\du}}
\pgfpathlineto{\pgfpoint{-23.071642\du}{3.318360\du}}
\pgfpathlineto{\pgfpoint{-22.857727\du}{3.318360\du}}
\pgfpathlineto{\pgfpoint{-22.857727\du}{3.268926\du}}
\pgfpathlineto{\pgfpoint{-22.706429\du}{3.342628\du}}
\pgfpathlineto{\pgfpoint{-22.857727\du}{3.417228\du}}
\pgfpathlineto{\pgfpoint{-22.857727\du}{3.367794\du}}
\pgfpathlineto{\pgfpoint{-23.046176\du}{3.367794\du}}
\pgfpathlineto{\pgfpoint{-23.172308\du}{3.454379\du}}
\pgfpathlineto{\pgfpoint{-23.046176\du}{3.552947\du}}
\pgfpathlineto{\pgfpoint{-22.857727\du}{3.552947\du}}
\pgfpathlineto{\pgfpoint{-22.857727\du}{3.491829\du}}
\pgfpathlineto{\pgfpoint{-22.706429\du}{3.578114\du}}
\pgfpathlineto{\pgfpoint{-22.857727\du}{3.652714\du}}
\pgfpathlineto{\pgfpoint{-22.857727\du}{3.602681\du}}
\pgfpathlineto{\pgfpoint{-23.071642\du}{3.602681\du}}
\pgfpathlineto{\pgfpoint{-23.209758\du}{3.491829\du}}
\pgfpathlineto{\pgfpoint{-23.348174\du}{3.602681\du}}
\pgfpathlineto{\pgfpoint{-23.574672\du}{3.602681\du}}
\pgfpathlineto{\pgfpoint{-23.574672\du}{3.652714\du}}
\pgfpathlineto{\pgfpoint{-23.725670\du}{3.578114\du}}
\pgfpathlineto{\pgfpoint{-23.574672\du}{3.491829\du}}
\pgfpathlineto{\pgfpoint{-23.574672\du}{3.552947\du}}
\pgfpathlineto{\pgfpoint{-23.373640\du}{3.552947\du}}
\pgfpathlineto{\pgfpoint{-23.247508\du}{3.466662\du}}
\pgfpathlineto{\pgfpoint{-23.373640\du}{3.367794\du}}
\pgfpathlineto{\pgfpoint{-23.574672\du}{3.367794\du}}
\pgfpathlineto{\pgfpoint{-23.574672\du}{3.417228\du}}
\pgfpathlineto{\pgfpoint{-23.725670\du}{3.342628\du}}
\pgfpathlineto{\pgfpoint{-23.574672\du}{3.268926\du}}
\pgfpathlineto{\pgfpoint{-23.574672\du}{3.318360\du}}
\pgfusepath{fill}
\pgfsetbuttcap
\pgfsetmiterjoin
\pgfsetdash{}{0pt}
\definecolor{dialinecolor}{rgb}{1.000000, 1.000000, 1.000000}
\pgfsetfillcolor{dialinecolor}
\pgfpathmoveto{\pgfpoint{-23.561789\du}{3.330344\du}}
\pgfpathlineto{\pgfpoint{-23.335590\du}{3.330344\du}}
\pgfpathlineto{\pgfpoint{-23.197175\du}{3.441496\du}}
\pgfpathlineto{\pgfpoint{-23.059059\du}{3.330344\du}}
\pgfpathlineto{\pgfpoint{-22.844545\du}{3.330344\du}}
\pgfpathlineto{\pgfpoint{-22.844545\du}{3.280910\du}}
\pgfpathlineto{\pgfpoint{-22.694145\du}{3.355511\du}}
\pgfpathlineto{\pgfpoint{-22.844545\du}{3.429812\du}}
\pgfpathlineto{\pgfpoint{-22.844545\du}{3.380377\du}}
\pgfpathlineto{\pgfpoint{-23.033593\du}{3.380377\du}}
\pgfpathlineto{\pgfpoint{-23.159725\du}{3.466662\du}}
\pgfpathlineto{\pgfpoint{-23.033593\du}{3.565531\du}}
\pgfpathlineto{\pgfpoint{-22.844545\du}{3.565531\du}}
\pgfpathlineto{\pgfpoint{-22.844545\du}{3.503813\du}}
\pgfpathlineto{\pgfpoint{-22.694145\du}{3.590398\du}}
\pgfpathlineto{\pgfpoint{-22.844545\du}{3.664698\du}}
\pgfpathlineto{\pgfpoint{-22.844545\du}{3.614965\du}}
\pgfpathlineto{\pgfpoint{-23.059059\du}{3.614965\du}}
\pgfpathlineto{\pgfpoint{-23.197175\du}{3.503813\du}}
\pgfpathlineto{\pgfpoint{-23.335590\du}{3.614965\du}}
\pgfpathlineto{\pgfpoint{-23.561789\du}{3.614965\du}}
\pgfpathlineto{\pgfpoint{-23.561789\du}{3.664698\du}}
\pgfpathlineto{\pgfpoint{-23.713087\du}{3.590398\du}}
\pgfpathlineto{\pgfpoint{-23.561789\du}{3.503813\du}}
\pgfpathlineto{\pgfpoint{-23.561789\du}{3.565531\du}}
\pgfpathlineto{\pgfpoint{-23.361056\du}{3.565531\du}}
\pgfpathlineto{\pgfpoint{-23.235524\du}{3.478946\du}}
\pgfpathlineto{\pgfpoint{-23.361056\du}{3.380377\du}}
\pgfpathlineto{\pgfpoint{-23.561789\du}{3.380377\du}}
\pgfpathlineto{\pgfpoint{-23.561789\du}{3.429812\du}}
\pgfpathlineto{\pgfpoint{-23.713087\du}{3.355511\du}}
\pgfpathlineto{\pgfpoint{-23.561789\du}{3.280910\du}}
\pgfpathlineto{\pgfpoint{-23.561789\du}{3.330344\du}}
\pgfusepath{fill}
\pgfsetlinewidth{0.100000\du}
\pgfsetdash{}{0pt}
\pgfsetdash{}{0pt}
\pgfsetbuttcap
\pgfsetmiterjoin
\pgfsetlinewidth{0.001000\du}
\pgfsetbuttcap
\pgfsetmiterjoin
\pgfsetdash{}{0pt}
\definecolor{dialinecolor}{rgb}{0.000000, 0.470588, 0.666667}
\pgfsetstrokecolor{dialinecolor}
\pgfpathmoveto{\pgfpoint{-20.581183\du}{7.491454\du}}
\pgfpathlineto{\pgfpoint{-20.574615\du}{7.482497\du}}
\pgfpathlineto{\pgfpoint{-20.568345\du}{7.472944\du}}
\pgfpathlineto{\pgfpoint{-20.562971\du}{7.462793\du}}
\pgfpathlineto{\pgfpoint{-20.557896\du}{7.453538\du}}
\pgfpathlineto{\pgfpoint{-20.554015\du}{7.443387\du}}
\pgfpathlineto{\pgfpoint{-20.550731\du}{7.433236\du}}
\pgfpathlineto{\pgfpoint{-20.547745\du}{7.423086\du}}
\pgfpathlineto{\pgfpoint{-20.545954\du}{7.412636\du}}
\pgfpathlineto{\pgfpoint{-20.544760\du}{7.402485\du}}
\pgfpathlineto{\pgfpoint{-20.543864\du}{7.392036\du}}
\pgfpathlineto{\pgfpoint{-20.543864\du}{7.381587\du}}
\pgfpathlineto{\pgfpoint{-20.544760\du}{7.370839\du}}
\pgfpathlineto{\pgfpoint{-20.545954\du}{7.360688\du}}
\pgfpathlineto{\pgfpoint{-20.547745\du}{7.350239\du}}
\pgfpathlineto{\pgfpoint{-20.550432\du}{7.340387\du}}
\pgfpathlineto{\pgfpoint{-20.554015\du}{7.329937\du}}
\pgfpathlineto{\pgfpoint{-20.557896\du}{7.320085\du}}
\pgfpathlineto{\pgfpoint{-20.562971\du}{7.309934\du}}
\pgfpathlineto{\pgfpoint{-20.568345\du}{7.300381\du}}
\pgfpathlineto{\pgfpoint{-20.574615\du}{7.291424\du}}
\pgfpathlineto{\pgfpoint{-20.581183\du}{7.281870\du}}
\pgfpathlineto{\pgfpoint{-20.588647\du}{7.273212\du}}
\pgfpathlineto{\pgfpoint{-20.596111\du}{7.264554\du}}
\pgfpathlineto{\pgfpoint{-20.604470\du}{7.256195\du}}
\pgfpathlineto{\pgfpoint{-20.613128\du}{7.248134\du}}
\pgfpathlineto{\pgfpoint{-20.622383\du}{7.240372\du}}
\pgfpathlineto{\pgfpoint{-20.631937\du}{7.233206\du}}
\pgfpathlineto{\pgfpoint{-20.642685\du}{7.226638\du}}
\pgfpathlineto{\pgfpoint{-20.653134\du}{7.220070\du}}
\pgfpathlineto{\pgfpoint{-20.664479\du}{7.213801\du}}
\pgfpathlineto{\pgfpoint{-20.676421\du}{7.208128\du}}
\pgfpathlineto{\pgfpoint{-20.688363\du}{7.202754\du}}
\pgfpathlineto{\pgfpoint{-20.700604\du}{7.198574\du}}
\pgfpathlineto{\pgfpoint{-20.713143\du}{7.194395\du}}
\pgfpathlineto{\pgfpoint{-20.726279\du}{7.190215\du}}
\pgfpathlineto{\pgfpoint{-20.739117\du}{7.187229\du}}
\pgfpathlineto{\pgfpoint{-20.752253\du}{7.184542\du}}
\pgfpathlineto{\pgfpoint{-20.765987\du}{7.182154\du}}
\pgfpathlineto{\pgfpoint{-20.779123\du}{7.180661\du}}
\pgfpathlineto{\pgfpoint{-20.793155\du}{7.179467\du}}
\pgfpathlineto{\pgfpoint{-20.806888\du}{7.178571\du}}
\pgfpathlineto{\pgfpoint{-20.820920\du}{7.178571\du}}
\pgfpathlineto{\pgfpoint{-20.834355\du}{7.179169\du}}
\pgfpathlineto{\pgfpoint{-20.848089\du}{7.179766\du}}
\pgfpathlineto{\pgfpoint{-20.862121\du}{7.181258\du}}
\pgfpathlineto{\pgfpoint{-20.875257\du}{7.183647\du}}
\pgfpathlineto{\pgfpoint{-20.888990\du}{7.186035\du}}
\pgfpathlineto{\pgfpoint{-20.901529\du}{7.189021\du}}
\pgfpathlineto{\pgfpoint{-20.914666\du}{7.192006\du}}
\pgfpathlineto{\pgfpoint{-20.927205\du}{7.196186\du}}
\pgfpathlineto{\pgfpoint{-20.940043\du}{7.200963\du}}
\pgfpathlineto{\pgfpoint{-20.952283\du}{7.205441\du}}
\pgfpathlineto{\pgfpoint{-20.963628\du}{7.210815\du}}
\pgfpathlineto{\pgfpoint{-20.975272\du}{7.217085\du}}
\pgfpathlineto{\pgfpoint{-20.986318\du}{7.223056\du}}
\pgfpathlineto{\pgfpoint{-20.996768\du}{7.229922\du}}
\pgfpathlineto{\pgfpoint{-21.006918\du}{7.236491\du}}
\pgfpathlineto{\pgfpoint{-21.016472\du}{7.244253\du}}
\pgfpathlineto{\pgfpoint{-21.026026\du}{7.252314\du}}
\pgfpathlineto{\pgfpoint{-21.034385\du}{7.260375\du}}
\pgfpathlineto{\pgfpoint{-21.042446\du}{7.268734\du}}
\pgfpathlineto{\pgfpoint{-21.050208\du}{7.277691\du}}
\pgfusepath{stroke}
\pgfsetlinewidth{0.080000\du}
\pgfsetbuttcap
\pgfsetmiterjoin
\pgfsetdash{}{0pt}
\definecolor{dialinecolor}{rgb}{0.666667, 0.901961, 1.000000}
\pgfsetstrokecolor{dialinecolor}
\pgfpathmoveto{\pgfpoint{-20.590438\du}{7.483990\du}}
\pgfpathlineto{\pgfpoint{-20.585064\du}{7.475034\du}}
\pgfpathlineto{\pgfpoint{-20.578795\du}{7.466376\du}}
\pgfpathlineto{\pgfpoint{-20.573719\du}{7.457120\du}}
\pgfpathlineto{\pgfpoint{-20.569540\du}{7.447865\du}}
\pgfpathlineto{\pgfpoint{-20.565658\du}{7.438312\du}}
\pgfpathlineto{\pgfpoint{-20.562673\du}{7.428758\du}}
\pgfpathlineto{\pgfpoint{-20.559986\du}{7.418906\du}}
\pgfpathlineto{\pgfpoint{-20.557896\du}{7.409651\du}}
\pgfpathlineto{\pgfpoint{-20.557000\du}{7.399798\du}}
\pgfpathlineto{\pgfpoint{-20.556403\du}{7.389648\du}}
\pgfpathlineto{\pgfpoint{-20.556403\du}{7.379795\du}}
\pgfpathlineto{\pgfpoint{-20.557299\du}{7.369943\du}}
\pgfpathlineto{\pgfpoint{-20.558792\du}{7.359793\du}}
\pgfpathlineto{\pgfpoint{-20.560583\du}{7.350537\du}}
\pgfpathlineto{\pgfpoint{-20.563569\du}{7.340685\du}}
\pgfpathlineto{\pgfpoint{-20.567151\du}{7.331132\du}}
\pgfpathlineto{\pgfpoint{-20.571032\du}{7.321876\du}}
\pgfpathlineto{\pgfpoint{-20.575511\du}{7.312621\du}}
\pgfpathlineto{\pgfpoint{-20.580586\du}{7.303963\du}}
\pgfpathlineto{\pgfpoint{-20.586258\du}{7.294708\du}}
\pgfpathlineto{\pgfpoint{-20.593424\du}{7.286349\du}}
\pgfpathlineto{\pgfpoint{-20.599992\du}{7.277691\du}}
\pgfpathlineto{\pgfpoint{-20.607754\du}{7.269630\du}}
\pgfpathlineto{\pgfpoint{-20.615815\du}{7.261867\du}}
\pgfpathlineto{\pgfpoint{-20.624175\du}{7.255001\du}}
\pgfpathlineto{\pgfpoint{-20.632833\du}{7.247537\du}}
\pgfpathlineto{\pgfpoint{-20.642685\du}{7.240670\du}}
\pgfpathlineto{\pgfpoint{-20.653134\du}{7.233804\du}}
\pgfpathlineto{\pgfpoint{-20.663285\du}{7.228131\du}}
\pgfpathlineto{\pgfpoint{-20.673734\du}{7.222757\du}}
\pgfpathlineto{\pgfpoint{-20.685676\du}{7.217383\du}}
\pgfpathlineto{\pgfpoint{-20.696424\du}{7.212308\du}}
\pgfpathlineto{\pgfpoint{-20.708963\du}{7.208128\du}}
\pgfpathlineto{\pgfpoint{-20.720308\du}{7.204247\du}}
\pgfpathlineto{\pgfpoint{-20.732848\du}{7.200963\du}}
\pgfpathlineto{\pgfpoint{-20.745387\du}{7.197977\du}}
\pgfpathlineto{\pgfpoint{-20.757926\du}{7.195290\du}}
\pgfpathlineto{\pgfpoint{-20.771062\du}{7.193499\du}}
\pgfpathlineto{\pgfpoint{-20.783601\du}{7.192006\du}}
\pgfpathlineto{\pgfpoint{-20.796738\du}{7.190812\du}}
\pgfpathlineto{\pgfpoint{-20.810770\du}{7.190215\du}}
\pgfpathlineto{\pgfpoint{-20.823906\du}{7.190215\du}}
\pgfpathlineto{\pgfpoint{-20.836744\du}{7.190812\du}}
\pgfpathlineto{\pgfpoint{-20.849880\du}{7.191708\du}}
\pgfpathlineto{\pgfpoint{-20.862419\du}{7.193499\du}}
\pgfpathlineto{\pgfpoint{-20.875555\du}{7.195290\du}}
\pgfpathlineto{\pgfpoint{-20.888095\du}{7.197380\du}}
\pgfpathlineto{\pgfpoint{-20.900932\du}{7.200963\du}}
\pgfpathlineto{\pgfpoint{-20.913173\du}{7.203948\du}}
\pgfpathlineto{\pgfpoint{-20.925115\du}{7.207830\du}}
\pgfpathlineto{\pgfpoint{-20.937356\du}{7.212009\du}}
\pgfpathlineto{\pgfpoint{-20.948701\du}{7.217085\du}}
\pgfpathlineto{\pgfpoint{-20.959747\du}{7.221861\du}}
\pgfpathlineto{\pgfpoint{-20.970495\du}{7.227833\du}}
\pgfpathlineto{\pgfpoint{-20.980944\du}{7.233505\du}}
\pgfpathlineto{\pgfpoint{-20.991394\du}{7.240073\du}}
\pgfpathlineto{\pgfpoint{-21.000947\du}{7.246641\du}}
\pgfpathlineto{\pgfpoint{-21.009605\du}{7.253807\du}}
\pgfpathlineto{\pgfpoint{-21.019159\du}{7.260972\du}}
\pgfpathlineto{\pgfpoint{-21.026921\du}{7.268734\du}}
\pgfpathlineto{\pgfpoint{-21.034385\du}{7.276795\du}}
\pgfpathlineto{\pgfpoint{-21.041550\du}{7.285155\du}}
\pgfusepath{stroke}
\pgfsetbuttcap
\pgfsetmiterjoin
\pgfsetdash{}{0pt}
\definecolor{dialinecolor}{rgb}{0.000000, 0.705882, 1.000000}
\pgfsetfillcolor{dialinecolor}
\pgfpathmoveto{\pgfpoint{-22.507441\du}{7.166928\du}}
\pgfpathlineto{\pgfpoint{-22.304724\du}{6.976153\du}}
\pgfpathlineto{\pgfpoint{-20.793155\du}{6.976153\du}}
\pgfpathlineto{\pgfpoint{-20.995872\du}{7.166928\du}}
\pgfpathlineto{\pgfpoint{-22.507441\du}{7.166928\du}}
\pgfusepath{fill}
\pgfsetlinewidth{0.001000\du}
\pgfsetbuttcap
\pgfsetmiterjoin
\pgfsetdash{}{0pt}
\definecolor{dialinecolor}{rgb}{0.666667, 0.901961, 1.000000}
\pgfsetstrokecolor{dialinecolor}
\pgfpathmoveto{\pgfpoint{-22.507441\du}{7.166928\du}}
\pgfpathlineto{\pgfpoint{-22.304724\du}{6.976153\du}}
\pgfpathlineto{\pgfpoint{-20.793155\du}{6.976153\du}}
\pgfpathlineto{\pgfpoint{-20.995872\du}{7.166928\du}}
\pgfpathlineto{\pgfpoint{-22.507441\du}{7.166928\du}}
\pgfusepath{stroke}
\pgfsetbuttcap
\pgfsetmiterjoin
\pgfsetdash{}{0pt}
\definecolor{dialinecolor}{rgb}{0.000000, 0.588235, 0.831373}
\pgfsetfillcolor{dialinecolor}
\pgfpathmoveto{\pgfpoint{-22.507441\du}{7.166928\du}}
\pgfpathlineto{\pgfpoint{-20.995872\du}{7.166928\du}}
\pgfpathlineto{\pgfpoint{-20.995872\du}{7.452941\du}}
\pgfpathlineto{\pgfpoint{-22.507441\du}{7.452941\du}}
\pgfpathlineto{\pgfpoint{-22.507441\du}{7.166928\du}}
\pgfusepath{fill}
\pgfsetbuttcap
\pgfsetmiterjoin
\pgfsetdash{}{0pt}
\definecolor{dialinecolor}{rgb}{0.666667, 0.901961, 1.000000}
\pgfsetstrokecolor{dialinecolor}
\pgfpathmoveto{\pgfpoint{-22.507441\du}{7.166928\du}}
\pgfpathlineto{\pgfpoint{-20.995872\du}{7.166928\du}}
\pgfpathlineto{\pgfpoint{-20.995872\du}{7.452941\du}}
\pgfpathlineto{\pgfpoint{-22.507441\du}{7.452941\du}}
\pgfpathlineto{\pgfpoint{-22.507441\du}{7.166928\du}}
\pgfusepath{stroke}
\pgfsetbuttcap
\pgfsetmiterjoin
\pgfsetdash{}{0pt}
\definecolor{dialinecolor}{rgb}{0.000000, 0.352941, 0.501961}
\pgfsetfillcolor{dialinecolor}
\pgfpathmoveto{\pgfpoint{-20.995872\du}{7.452941\du}}
\pgfpathlineto{\pgfpoint{-20.793155\du}{7.250224\du}}
\pgfpathlineto{\pgfpoint{-20.793155\du}{6.976153\du}}
\pgfpathlineto{\pgfpoint{-20.995872\du}{7.166928\du}}
\pgfpathlineto{\pgfpoint{-20.995872\du}{7.452941\du}}
\pgfusepath{fill}
\pgfsetbuttcap
\pgfsetmiterjoin
\pgfsetdash{}{0pt}
\definecolor{dialinecolor}{rgb}{0.666667, 0.901961, 1.000000}
\pgfsetstrokecolor{dialinecolor}
\pgfpathmoveto{\pgfpoint{-20.995872\du}{7.452941\du}}
\pgfpathlineto{\pgfpoint{-20.793155\du}{7.250224\du}}
\pgfpathlineto{\pgfpoint{-20.793155\du}{6.976153\du}}
\pgfpathlineto{\pgfpoint{-20.995872\du}{7.166928\du}}
\pgfpathlineto{\pgfpoint{-20.995872\du}{7.452941\du}}
\pgfusepath{stroke}
\pgfsetbuttcap
\pgfsetmiterjoin
\pgfsetdash{}{0pt}
\definecolor{dialinecolor}{rgb}{0.000000, 0.000000, 0.000000}
\pgfsetfillcolor{dialinecolor}
\pgfpathmoveto{\pgfpoint{-22.459374\du}{7.130803\du}}
\pgfpathlineto{\pgfpoint{-22.304724\du}{6.976153\du}}
\pgfpathlineto{\pgfpoint{-20.829280\du}{6.976153\du}}
\pgfpathlineto{\pgfpoint{-20.983631\du}{7.130803\du}}
\pgfpathlineto{\pgfpoint{-22.459374\du}{7.130803\du}}
\pgfusepath{fill}
\pgfsetbuttcap
\pgfsetmiterjoin
\pgfsetdash{}{0pt}
\definecolor{dialinecolor}{rgb}{0.000000, 0.000000, 0.000000}
\pgfsetstrokecolor{dialinecolor}
\pgfpathmoveto{\pgfpoint{-22.459374\du}{7.130803\du}}
\pgfpathlineto{\pgfpoint{-22.304724\du}{6.976153\du}}
\pgfpathlineto{\pgfpoint{-20.829280\du}{6.976153\du}}
\pgfpathlineto{\pgfpoint{-20.983631\du}{7.130803\du}}
\pgfpathlineto{\pgfpoint{-22.459374\du}{7.130803\du}}
\pgfusepath{stroke}
\pgfsetbuttcap
\pgfsetmiterjoin
\pgfsetdash{}{0pt}
\definecolor{dialinecolor}{rgb}{0.000000, 0.705882, 1.000000}
\pgfsetfillcolor{dialinecolor}
\pgfpathmoveto{\pgfpoint{-22.507441\du}{5.904650\du}}
\pgfpathlineto{\pgfpoint{-22.340252\du}{5.750000\du}}
\pgfpathlineto{\pgfpoint{-20.829280\du}{5.750000\du}}
\pgfpathlineto{\pgfpoint{-20.995872\du}{5.904650\du}}
\pgfpathlineto{\pgfpoint{-22.507441\du}{5.904650\du}}
\pgfusepath{fill}
\pgfsetbuttcap
\pgfsetmiterjoin
\pgfsetdash{}{0pt}
\definecolor{dialinecolor}{rgb}{0.666667, 0.901961, 1.000000}
\pgfsetstrokecolor{dialinecolor}
\pgfpathmoveto{\pgfpoint{-22.507441\du}{5.904650\du}}
\pgfpathlineto{\pgfpoint{-22.340252\du}{5.750000\du}}
\pgfpathlineto{\pgfpoint{-20.829280\du}{5.750000\du}}
\pgfpathlineto{\pgfpoint{-20.995872\du}{5.904650\du}}
\pgfpathlineto{\pgfpoint{-22.507441\du}{5.904650\du}}
\pgfusepath{stroke}
\pgfsetbuttcap
\pgfsetmiterjoin
\pgfsetdash{}{0pt}
\definecolor{dialinecolor}{rgb}{0.000000, 0.588235, 0.831373}
\pgfsetfillcolor{dialinecolor}
\pgfpathmoveto{\pgfpoint{-22.507441\du}{5.904650\du}}
\pgfpathlineto{\pgfpoint{-20.983631\du}{5.904650\du}}
\pgfpathlineto{\pgfpoint{-20.983631\du}{7.106919\du}}
\pgfpathlineto{\pgfpoint{-22.507441\du}{7.106919\du}}
\pgfpathlineto{\pgfpoint{-22.507441\du}{5.904650\du}}
\pgfusepath{fill}
\pgfsetbuttcap
\pgfsetmiterjoin
\pgfsetdash{}{0pt}
\definecolor{dialinecolor}{rgb}{0.666667, 0.901961, 1.000000}
\pgfsetstrokecolor{dialinecolor}
\pgfpathmoveto{\pgfpoint{-22.507441\du}{5.904650\du}}
\pgfpathlineto{\pgfpoint{-20.983930\du}{5.904650\du}}
\pgfpathlineto{\pgfpoint{-20.983930\du}{7.106620\du}}
\pgfpathlineto{\pgfpoint{-22.507441\du}{7.106620\du}}
\pgfpathlineto{\pgfpoint{-22.507441\du}{5.904650\du}}
\pgfusepath{stroke}
\pgfsetbuttcap
\pgfsetmiterjoin
\pgfsetdash{}{0pt}
\definecolor{dialinecolor}{rgb}{1.000000, 1.000000, 1.000000}
\pgfsetfillcolor{dialinecolor}
\pgfpathmoveto{\pgfpoint{-22.376376\du}{6.059598\du}}
\pgfpathlineto{\pgfpoint{-21.114696\du}{6.059598\du}}
\pgfpathlineto{\pgfpoint{-21.114696\du}{6.988394\du}}
\pgfpathlineto{\pgfpoint{-22.376376\du}{6.988394\du}}
\pgfpathlineto{\pgfpoint{-22.376376\du}{6.059598\du}}
\pgfusepath{fill}
\pgfsetbuttcap
\pgfsetmiterjoin
\pgfsetdash{}{0pt}
\definecolor{dialinecolor}{rgb}{0.000000, 0.235294, 0.333333}
\pgfsetstrokecolor{dialinecolor}
\pgfpathmoveto{\pgfpoint{-22.376376\du}{6.059598\du}}
\pgfpathlineto{\pgfpoint{-21.114994\du}{6.059598\du}}
\pgfpathlineto{\pgfpoint{-21.114994\du}{6.988095\du}}
\pgfpathlineto{\pgfpoint{-22.376376\du}{6.988095\du}}
\pgfpathlineto{\pgfpoint{-22.376376\du}{6.059598\du}}
\pgfusepath{stroke}
\pgfsetbuttcap
\pgfsetmiterjoin
\pgfsetdash{}{0pt}
\definecolor{dialinecolor}{rgb}{0.000000, 0.352941, 0.501961}
\pgfsetfillcolor{dialinecolor}
\pgfpathmoveto{\pgfpoint{-20.995872\du}{7.095574\du}}
\pgfpathlineto{\pgfpoint{-20.829280\du}{6.928683\du}}
\pgfpathlineto{\pgfpoint{-20.829280\du}{5.750000\du}}
\pgfpathlineto{\pgfpoint{-20.995872\du}{5.904650\du}}
\pgfpathlineto{\pgfpoint{-20.995872\du}{7.095574\du}}
\pgfusepath{fill}
\pgfsetbuttcap
\pgfsetmiterjoin
\pgfsetdash{}{0pt}
\definecolor{dialinecolor}{rgb}{0.666667, 0.901961, 1.000000}
\pgfsetstrokecolor{dialinecolor}
\pgfpathmoveto{\pgfpoint{-20.995872\du}{7.095574\du}}
\pgfpathlineto{\pgfpoint{-20.829280\du}{6.928683\du}}
\pgfpathlineto{\pgfpoint{-20.829280\du}{5.750000\du}}
\pgfpathlineto{\pgfpoint{-20.995872\du}{5.904650\du}}
\pgfpathlineto{\pgfpoint{-20.995872\du}{7.095574\du}}
\pgfusepath{stroke}
\pgfsetbuttcap
\pgfsetmiterjoin
\pgfsetdash{}{0pt}
\definecolor{dialinecolor}{rgb}{0.000000, 0.705882, 1.000000}
\pgfsetfillcolor{dialinecolor}
\pgfpathmoveto{\pgfpoint{-22.804500\du}{7.690887\du}}
\pgfpathlineto{\pgfpoint{-22.566853\du}{7.392932\du}}
\pgfpathlineto{\pgfpoint{-20.924219\du}{7.392932\du}}
\pgfpathlineto{\pgfpoint{-21.162165\du}{7.690887\du}}
\pgfpathlineto{\pgfpoint{-22.804500\du}{7.690887\du}}
\pgfusepath{fill}
\pgfsetbuttcap
\pgfsetmiterjoin
\pgfsetdash{}{0pt}
\definecolor{dialinecolor}{rgb}{0.666667, 0.901961, 1.000000}
\pgfsetstrokecolor{dialinecolor}
\pgfpathmoveto{\pgfpoint{-22.804500\du}{7.690887\du}}
\pgfpathlineto{\pgfpoint{-22.566853\du}{7.392932\du}}
\pgfpathlineto{\pgfpoint{-20.924219\du}{7.392932\du}}
\pgfpathlineto{\pgfpoint{-21.162165\du}{7.690887\du}}
\pgfpathlineto{\pgfpoint{-22.804500\du}{7.690887\du}}
\pgfusepath{stroke}
\pgfsetbuttcap
\pgfsetmiterjoin
\pgfsetdash{}{0pt}
\definecolor{dialinecolor}{rgb}{0.000000, 0.352941, 0.501961}
\pgfsetfillcolor{dialinecolor}
\pgfpathmoveto{\pgfpoint{-21.162165\du}{7.750000\du}}
\pgfpathlineto{\pgfpoint{-20.924219\du}{7.500112\du}}
\pgfpathlineto{\pgfpoint{-20.924219\du}{7.392932\du}}
\pgfpathlineto{\pgfpoint{-21.162165\du}{7.690887\du}}
\pgfpathlineto{\pgfpoint{-21.162165\du}{7.750000\du}}
\pgfusepath{fill}
\pgfsetbuttcap
\pgfsetmiterjoin
\pgfsetdash{}{0pt}
\definecolor{dialinecolor}{rgb}{0.666667, 0.901961, 1.000000}
\pgfsetstrokecolor{dialinecolor}
\pgfpathmoveto{\pgfpoint{-21.162165\du}{7.750000\du}}
\pgfpathlineto{\pgfpoint{-20.924219\du}{7.500112\du}}
\pgfpathlineto{\pgfpoint{-20.924219\du}{7.392932\du}}
\pgfpathlineto{\pgfpoint{-21.162165\du}{7.690887\du}}
\pgfpathlineto{\pgfpoint{-21.162165\du}{7.750000\du}}
\pgfusepath{stroke}
\pgfsetbuttcap
\pgfsetmiterjoin
\pgfsetdash{}{0pt}
\definecolor{dialinecolor}{rgb}{0.000000, 0.588235, 0.831373}
\pgfsetfillcolor{dialinecolor}
\pgfpathmoveto{\pgfpoint{-22.804500\du}{7.690887\du}}
\pgfpathlineto{\pgfpoint{-21.162165\du}{7.690887\du}}
\pgfpathlineto{\pgfpoint{-21.162165\du}{7.750000\du}}
\pgfpathlineto{\pgfpoint{-22.804500\du}{7.750000\du}}
\pgfpathlineto{\pgfpoint{-22.804500\du}{7.690887\du}}
\pgfusepath{fill}
\pgfsetbuttcap
\pgfsetmiterjoin
\pgfsetdash{}{0pt}
\definecolor{dialinecolor}{rgb}{0.666667, 0.901961, 1.000000}
\pgfsetstrokecolor{dialinecolor}
\pgfpathmoveto{\pgfpoint{-22.804500\du}{7.690887\du}}
\pgfpathlineto{\pgfpoint{-21.162165\du}{7.690887\du}}
\pgfpathlineto{\pgfpoint{-21.162165\du}{7.750000\du}}
\pgfpathlineto{\pgfpoint{-22.804500\du}{7.750000\du}}
\pgfpathlineto{\pgfpoint{-22.804500\du}{7.690887\du}}
\pgfusepath{stroke}
\pgfsetbuttcap
\pgfsetmiterjoin
\pgfsetdash{}{0pt}
\definecolor{dialinecolor}{rgb}{0.000000, 0.705882, 1.000000}
\pgfsetfillcolor{dialinecolor}
\pgfpathmoveto{\pgfpoint{-20.829280\du}{7.654762\du}}
\pgfpathlineto{\pgfpoint{-20.710158\du}{7.476526\du}}
\pgfpathlineto{\pgfpoint{-20.495499\du}{7.476526\du}}
\pgfpathlineto{\pgfpoint{-20.614322\du}{7.654762\du}}
\pgfpathlineto{\pgfpoint{-20.829280\du}{7.654762\du}}
\pgfusepath{fill}
\pgfsetbuttcap
\pgfsetmiterjoin
\pgfsetdash{}{0pt}
\definecolor{dialinecolor}{rgb}{0.666667, 0.901961, 1.000000}
\pgfsetstrokecolor{dialinecolor}
\pgfpathmoveto{\pgfpoint{-20.829280\du}{7.654762\du}}
\pgfpathlineto{\pgfpoint{-20.710158\du}{7.476526\du}}
\pgfpathlineto{\pgfpoint{-20.495499\du}{7.476526\du}}
\pgfpathlineto{\pgfpoint{-20.614322\du}{7.654762\du}}
\pgfpathlineto{\pgfpoint{-20.829280\du}{7.654762\du}}
\pgfusepath{stroke}
\pgfsetbuttcap
\pgfsetmiterjoin
\pgfsetdash{}{0pt}
\definecolor{dialinecolor}{rgb}{0.000000, 0.352941, 0.501961}
\pgfsetfillcolor{dialinecolor}
\pgfpathmoveto{\pgfpoint{-20.614322\du}{7.714472\du}}
\pgfpathlineto{\pgfpoint{-20.495499\du}{7.583408\du}}
\pgfpathlineto{\pgfpoint{-20.495499\du}{7.476526\du}}
\pgfpathlineto{\pgfpoint{-20.614322\du}{7.654762\du}}
\pgfpathlineto{\pgfpoint{-20.614322\du}{7.714472\du}}
\pgfusepath{fill}
\pgfsetbuttcap
\pgfsetmiterjoin
\pgfsetdash{}{0pt}
\definecolor{dialinecolor}{rgb}{0.666667, 0.901961, 1.000000}
\pgfsetstrokecolor{dialinecolor}
\pgfpathmoveto{\pgfpoint{-20.614322\du}{7.714472\du}}
\pgfpathlineto{\pgfpoint{-20.495499\du}{7.583408\du}}
\pgfpathlineto{\pgfpoint{-20.495499\du}{7.476526\du}}
\pgfpathlineto{\pgfpoint{-20.614322\du}{7.654762\du}}
\pgfpathlineto{\pgfpoint{-20.614322\du}{7.714472\du}}
\pgfusepath{stroke}
\pgfsetbuttcap
\pgfsetmiterjoin
\pgfsetdash{}{0pt}
\definecolor{dialinecolor}{rgb}{0.000000, 0.588235, 0.831373}
\pgfsetfillcolor{dialinecolor}
\pgfpathmoveto{\pgfpoint{-20.829280\du}{7.654762\du}}
\pgfpathlineto{\pgfpoint{-20.614322\du}{7.654762\du}}
\pgfpathlineto{\pgfpoint{-20.614322\du}{7.714472\du}}
\pgfpathlineto{\pgfpoint{-20.829280\du}{7.714472\du}}
\pgfpathlineto{\pgfpoint{-20.829280\du}{7.654762\du}}
\pgfusepath{fill}
\pgfsetbuttcap
\pgfsetmiterjoin
\pgfsetdash{}{0pt}
\definecolor{dialinecolor}{rgb}{0.666667, 0.901961, 1.000000}
\pgfsetstrokecolor{dialinecolor}
\pgfpathmoveto{\pgfpoint{-20.829280\du}{7.654762\du}}
\pgfpathlineto{\pgfpoint{-20.614322\du}{7.654762\du}}
\pgfpathlineto{\pgfpoint{-20.614322\du}{7.714472\du}}
\pgfpathlineto{\pgfpoint{-20.829280\du}{7.714472\du}}
\pgfpathlineto{\pgfpoint{-20.829280\du}{7.654762\du}}
\pgfusepath{stroke}
\pgfsetbuttcap
\pgfsetmiterjoin
\pgfsetdash{}{0pt}
\definecolor{dialinecolor}{rgb}{0.811765, 0.054902, 0.188235}
\pgfsetstrokecolor{dialinecolor}
\pgfpathmoveto{\pgfpoint{-22.043192\du}{6.285901\du}}
\pgfpathlineto{\pgfpoint{-21.709411\du}{6.285901\du}}
\pgfpathlineto{\pgfpoint{-21.793006\du}{6.369197\du}}
\pgfpathlineto{\pgfpoint{-21.507291\du}{6.369197\du}}
\pgfusepath{stroke}
\pgfsetbuttcap
\pgfsetmiterjoin
\pgfsetdash{}{0pt}
\definecolor{dialinecolor}{rgb}{0.811765, 0.054902, 0.188235}
\pgfsetstrokecolor{dialinecolor}
\pgfpathmoveto{\pgfpoint{-22.079019\du}{6.285901\du}}
\pgfpathlineto{\pgfpoint{-21.840476\du}{6.523847\du}}
\pgfpathlineto{\pgfpoint{-21.959896\du}{6.523847\du}}
\pgfpathlineto{\pgfpoint{-21.757478\du}{6.726564\du}}
\pgfusepath{stroke}
\pgfsetbuttcap
\pgfsetmiterjoin
\pgfsetdash{}{0pt}
\definecolor{dialinecolor}{rgb}{0.811765, 0.054902, 0.188235}
\pgfsetstrokecolor{dialinecolor}
\pgfpathmoveto{\pgfpoint{-21.400708\du}{6.285901\du}}
\pgfpathlineto{\pgfpoint{-21.638057\du}{6.523847\du}}
\pgfpathlineto{\pgfpoint{-21.519532\du}{6.523847\du}}
\pgfpathlineto{\pgfpoint{-21.721652\du}{6.726564\du}}
\pgfusepath{stroke}
\pgfsetbuttcap
\pgfsetmiterjoin
\pgfsetdash{}{0pt}
\definecolor{dialinecolor}{rgb}{0.000000, 0.470588, 0.666667}
\pgfsetfillcolor{dialinecolor}
\pgfpathmoveto{\pgfpoint{-21.912427\du}{6.374869\du}}
\pgfpathlineto{\pgfpoint{-21.913024\du}{6.371585\du}}
\pgfpathlineto{\pgfpoint{-21.913919\du}{6.368301\du}}
\pgfpathlineto{\pgfpoint{-21.914815\du}{6.365017\du}}
\pgfpathlineto{\pgfpoint{-21.916606\du}{6.361733\du}}
\pgfpathlineto{\pgfpoint{-21.918995\du}{6.358748\du}}
\pgfpathlineto{\pgfpoint{-21.921682\du}{6.355464\du}}
\pgfpathlineto{\pgfpoint{-21.924667\du}{6.352179\du}}
\pgfpathlineto{\pgfpoint{-21.928250\du}{6.349194\du}}
\pgfpathlineto{\pgfpoint{-21.932430\du}{6.346208\du}}
\pgfpathlineto{\pgfpoint{-21.937206\du}{6.343223\du}}
\pgfpathlineto{\pgfpoint{-21.942282\du}{6.340536\du}}
\pgfpathlineto{\pgfpoint{-21.947954\du}{6.337550\du}}
\pgfpathlineto{\pgfpoint{-21.953627\du}{6.335162\du}}
\pgfpathlineto{\pgfpoint{-21.959896\du}{6.332475\du}}
\pgfpathlineto{\pgfpoint{-21.966464\du}{6.330087\du}}
\pgfpathlineto{\pgfpoint{-21.973630\du}{6.327698\du}}
\pgfpathlineto{\pgfpoint{-21.981094\du}{6.325608\du}}
\pgfpathlineto{\pgfpoint{-21.988557\du}{6.323220\du}}
\pgfpathlineto{\pgfpoint{-21.996917\du}{6.321727\du}}
\pgfpathlineto{\pgfpoint{-22.005276\du}{6.319637\du}}
\pgfpathlineto{\pgfpoint{-22.013337\du}{6.317846\du}}
\pgfpathlineto{\pgfpoint{-22.022592\du}{6.316652\du}}
\pgfpathlineto{\pgfpoint{-22.031250\du}{6.315159\du}}
\pgfpathlineto{\pgfpoint{-22.040804\du}{6.313965\du}}
\pgfpathlineto{\pgfpoint{-22.050059\du}{6.312771\du}}
\pgfpathlineto{\pgfpoint{-22.059314\du}{6.311875\du}}
\pgfpathlineto{\pgfpoint{-22.069166\du}{6.310979\du}}
\pgfpathlineto{\pgfpoint{-22.079019\du}{6.310382\du}}
\pgfpathlineto{\pgfpoint{-22.088871\du}{6.310084\du}}
\pgfpathlineto{\pgfpoint{-22.099022\du}{6.309785\du}}
\pgfpathlineto{\pgfpoint{-22.108874\du}{6.309785\du}}
\pgfpathlineto{\pgfpoint{-22.108874\du}{6.309785\du}}
\pgfpathlineto{\pgfpoint{-22.118726\du}{6.309785\du}}
\pgfpathlineto{\pgfpoint{-22.128578\du}{6.310084\du}}
\pgfpathlineto{\pgfpoint{-22.138729\du}{6.310382\du}}
\pgfpathlineto{\pgfpoint{-22.148581\du}{6.310979\du}}
\pgfpathlineto{\pgfpoint{-22.158135\du}{6.311875\du}}
\pgfpathlineto{\pgfpoint{-22.167689\du}{6.312771\du}}
\pgfpathlineto{\pgfpoint{-22.176944\du}{6.313965\du}}
\pgfpathlineto{\pgfpoint{-22.186199\du}{6.315159\du}}
\pgfpathlineto{\pgfpoint{-22.195155\du}{6.316652\du}}
\pgfpathlineto{\pgfpoint{-22.204112\du}{6.317846\du}}
\pgfpathlineto{\pgfpoint{-22.212471\du}{6.319637\du}}
\pgfpathlineto{\pgfpoint{-22.220831\du}{6.321727\du}}
\pgfpathlineto{\pgfpoint{-22.229190\du}{6.323220\du}}
\pgfpathlineto{\pgfpoint{-22.236356\du}{6.325608\du}}
\pgfpathlineto{\pgfpoint{-22.243819\du}{6.327698\du}}
\pgfpathlineto{\pgfpoint{-22.250985\du}{6.330087\du}}
\pgfpathlineto{\pgfpoint{-22.257851\du}{6.332475\du}}
\pgfpathlineto{\pgfpoint{-22.263822\du}{6.335162\du}}
\pgfpathlineto{\pgfpoint{-22.269793\du}{6.337550\du}}
\pgfpathlineto{\pgfpoint{-22.275466\du}{6.340536\du}}
\pgfpathlineto{\pgfpoint{-22.280541\du}{6.343223\du}}
\pgfpathlineto{\pgfpoint{-22.285318\du}{6.346208\du}}
\pgfpathlineto{\pgfpoint{-22.289498\du}{6.349194\du}}
\pgfpathlineto{\pgfpoint{-22.292782\du}{6.352179\du}}
\pgfpathlineto{\pgfpoint{-22.296066\du}{6.355464\du}}
\pgfpathlineto{\pgfpoint{-22.299051\du}{6.358748\du}}
\pgfpathlineto{\pgfpoint{-22.301141\du}{6.361733\du}}
\pgfpathlineto{\pgfpoint{-22.302634\du}{6.365017\du}}
\pgfpathlineto{\pgfpoint{-22.304425\du}{6.368301\du}}
\pgfpathlineto{\pgfpoint{-22.304724\du}{6.371585\du}}
\pgfpathlineto{\pgfpoint{-22.305022\du}{6.374869\du}}
\pgfpathlineto{\pgfpoint{-22.305022\du}{6.374869\du}}
\pgfpathlineto{\pgfpoint{-22.304724\du}{6.378452\du}}
\pgfpathlineto{\pgfpoint{-22.304425\du}{6.381438\du}}
\pgfpathlineto{\pgfpoint{-22.302634\du}{6.384722\du}}
\pgfpathlineto{\pgfpoint{-22.301141\du}{6.388006\du}}
\pgfpathlineto{\pgfpoint{-22.299051\du}{6.391290\du}}
\pgfpathlineto{\pgfpoint{-22.296066\du}{6.394574\du}}
\pgfpathlineto{\pgfpoint{-22.292782\du}{6.397559\du}}
\pgfpathlineto{\pgfpoint{-22.289498\du}{6.400545\du}}
\pgfpathlineto{\pgfpoint{-22.285318\du}{6.403530\du}}
\pgfpathlineto{\pgfpoint{-22.280541\du}{6.406814\du}}
\pgfpathlineto{\pgfpoint{-22.275466\du}{6.409501\du}}
\pgfpathlineto{\pgfpoint{-22.269793\du}{6.412188\du}}
\pgfpathlineto{\pgfpoint{-22.263822\du}{6.414875\du}}
\pgfpathlineto{\pgfpoint{-22.257851\du}{6.417562\du}}
\pgfpathlineto{\pgfpoint{-22.250985\du}{6.419951\du}}
\pgfpathlineto{\pgfpoint{-22.243819\du}{6.422339\du}}
\pgfpathlineto{\pgfpoint{-22.236356\du}{6.424130\du}}
\pgfpathlineto{\pgfpoint{-22.229190\du}{6.426220\du}}
\pgfpathlineto{\pgfpoint{-22.220831\du}{6.428310\du}}
\pgfpathlineto{\pgfpoint{-22.212471\du}{6.430400\du}}
\pgfpathlineto{\pgfpoint{-22.204112\du}{6.431893\du}}
\pgfpathlineto{\pgfpoint{-22.195155\du}{6.433386\du}}
\pgfpathlineto{\pgfpoint{-22.186199\du}{6.434878\du}}
\pgfpathlineto{\pgfpoint{-22.176944\du}{6.436073\du}}
\pgfpathlineto{\pgfpoint{-22.167689\du}{6.436968\du}}
\pgfpathlineto{\pgfpoint{-22.158135\du}{6.438162\du}}
\pgfpathlineto{\pgfpoint{-22.148581\du}{6.438760\du}}
\pgfpathlineto{\pgfpoint{-22.138729\du}{6.439357\du}}
\pgfpathlineto{\pgfpoint{-22.128578\du}{6.439954\du}}
\pgfpathlineto{\pgfpoint{-22.118726\du}{6.440252\du}}
\pgfpathlineto{\pgfpoint{-22.108874\du}{6.440252\du}}
\pgfpathlineto{\pgfpoint{-22.108874\du}{6.440252\du}}
\pgfpathlineto{\pgfpoint{-22.099022\du}{6.440252\du}}
\pgfpathlineto{\pgfpoint{-22.088871\du}{6.439954\du}}
\pgfpathlineto{\pgfpoint{-22.079019\du}{6.439357\du}}
\pgfpathlineto{\pgfpoint{-22.069166\du}{6.438760\du}}
\pgfpathlineto{\pgfpoint{-22.059314\du}{6.438162\du}}
\pgfpathlineto{\pgfpoint{-22.050059\du}{6.436968\du}}
\pgfpathlineto{\pgfpoint{-22.040804\du}{6.436073\du}}
\pgfpathlineto{\pgfpoint{-22.031250\du}{6.434878\du}}
\pgfpathlineto{\pgfpoint{-22.022592\du}{6.433386\du}}
\pgfpathlineto{\pgfpoint{-22.013337\du}{6.431893\du}}
\pgfpathlineto{\pgfpoint{-22.005276\du}{6.430400\du}}
\pgfpathlineto{\pgfpoint{-21.996917\du}{6.428310\du}}
\pgfpathlineto{\pgfpoint{-21.988557\du}{6.426220\du}}
\pgfpathlineto{\pgfpoint{-21.981094\du}{6.424130\du}}
\pgfpathlineto{\pgfpoint{-21.973630\du}{6.422339\du}}
\pgfpathlineto{\pgfpoint{-21.966464\du}{6.419951\du}}
\pgfpathlineto{\pgfpoint{-21.959896\du}{6.417562\du}}
\pgfpathlineto{\pgfpoint{-21.953627\du}{6.414875\du}}
\pgfpathlineto{\pgfpoint{-21.947954\du}{6.412188\du}}
\pgfpathlineto{\pgfpoint{-21.942282\du}{6.409501\du}}
\pgfpathlineto{\pgfpoint{-21.937206\du}{6.406814\du}}
\pgfpathlineto{\pgfpoint{-21.932430\du}{6.403530\du}}
\pgfpathlineto{\pgfpoint{-21.928250\du}{6.400545\du}}
\pgfpathlineto{\pgfpoint{-21.924667\du}{6.397559\du}}
\pgfpathlineto{\pgfpoint{-21.921682\du}{6.394574\du}}
\pgfpathlineto{\pgfpoint{-21.918995\du}{6.391290\du}}
\pgfpathlineto{\pgfpoint{-21.916606\du}{6.388006\du}}
\pgfpathlineto{\pgfpoint{-21.914815\du}{6.384722\du}}
\pgfpathlineto{\pgfpoint{-21.913919\du}{6.381438\du}}
\pgfpathlineto{\pgfpoint{-21.913024\du}{6.378452\du}}
\pgfpathlineto{\pgfpoint{-21.912427\du}{6.374869\du}}
\pgfusepath{fill}
\pgfsetbuttcap
\pgfsetmiterjoin
\pgfsetdash{}{0pt}
\definecolor{dialinecolor}{rgb}{0.666667, 0.901961, 1.000000}
\pgfsetstrokecolor{dialinecolor}
\pgfpathmoveto{\pgfpoint{-21.924667\du}{6.369197\du}}
\pgfpathlineto{\pgfpoint{-21.924667\du}{6.365614\du}}
\pgfpathlineto{\pgfpoint{-21.925563\du}{6.363226\du}}
\pgfpathlineto{\pgfpoint{-21.926757\du}{6.359942\du}}
\pgfpathlineto{\pgfpoint{-21.928250\du}{6.356956\du}}
\pgfpathlineto{\pgfpoint{-21.930638\du}{6.354269\du}}
\pgfpathlineto{\pgfpoint{-21.933027\du}{6.350687\du}}
\pgfpathlineto{\pgfpoint{-21.936311\du}{6.348298\du}}
\pgfpathlineto{\pgfpoint{-21.939595\du}{6.345611\du}}
\pgfpathlineto{\pgfpoint{-21.943775\du}{6.342924\du}}
\pgfpathlineto{\pgfpoint{-21.948253\du}{6.340536\du}}
\pgfpathlineto{\pgfpoint{-21.953030\du}{6.337550\du}}
\pgfpathlineto{\pgfpoint{-21.958702\du}{6.335461\du}}
\pgfpathlineto{\pgfpoint{-21.964076\du}{6.332774\du}}
\pgfpathlineto{\pgfpoint{-21.970346\du}{6.330087\du}}
\pgfpathlineto{\pgfpoint{-21.976317\du}{6.328295\du}}
\pgfpathlineto{\pgfpoint{-21.983780\du}{6.325907\du}}
\pgfpathlineto{\pgfpoint{-21.990647\du}{6.324414\du}}
\pgfpathlineto{\pgfpoint{-21.997812\du}{6.322324\du}}
\pgfpathlineto{\pgfpoint{-22.005575\du}{6.320533\du}}
\pgfpathlineto{\pgfpoint{-22.014233\du}{6.318742\du}}
\pgfpathlineto{\pgfpoint{-22.022592\du}{6.317547\du}}
\pgfpathlineto{\pgfpoint{-22.030952\du}{6.315159\du}}
\pgfpathlineto{\pgfpoint{-22.039908\du}{6.314860\du}}
\pgfpathlineto{\pgfpoint{-22.048566\du}{6.313368\du}}
\pgfpathlineto{\pgfpoint{-22.057523\du}{6.312472\du}}
\pgfpathlineto{\pgfpoint{-22.066778\du}{6.311576\du}}
\pgfpathlineto{\pgfpoint{-22.076332\du}{6.310979\du}}
\pgfpathlineto{\pgfpoint{-22.085885\du}{6.310681\du}}
\pgfpathlineto{\pgfpoint{-22.094842\du}{6.310084\du}}
\pgfpathlineto{\pgfpoint{-22.104993\du}{6.309785\du}}
\pgfpathlineto{\pgfpoint{-22.114546\du}{6.309785\du}}
\pgfpathlineto{\pgfpoint{-22.114546\du}{6.309785\du}}
\pgfpathlineto{\pgfpoint{-22.123801\du}{6.309785\du}}
\pgfpathlineto{\pgfpoint{-22.133654\du}{6.310084\du}}
\pgfpathlineto{\pgfpoint{-22.143207\du}{6.310681\du}}
\pgfpathlineto{\pgfpoint{-22.152761\du}{6.310979\du}}
\pgfpathlineto{\pgfpoint{-22.162016\du}{6.311576\du}}
\pgfpathlineto{\pgfpoint{-22.171868\du}{6.312472\du}}
\pgfpathlineto{\pgfpoint{-22.180526\du}{6.313368\du}}
\pgfpathlineto{\pgfpoint{-22.189483\du}{6.314860\du}}
\pgfpathlineto{\pgfpoint{-22.197842\du}{6.315159\du}}
\pgfpathlineto{\pgfpoint{-22.206799\du}{6.317547\du}}
\pgfpathlineto{\pgfpoint{-22.215158\du}{6.318742\du}}
\pgfpathlineto{\pgfpoint{-22.223219\du}{6.320533\du}}
\pgfpathlineto{\pgfpoint{-22.230982\du}{6.322324\du}}
\pgfpathlineto{\pgfpoint{-22.238445\du}{6.324414\du}}
\pgfpathlineto{\pgfpoint{-22.245312\du}{6.325907\du}}
\pgfpathlineto{\pgfpoint{-22.252776\du}{6.328295\du}}
\pgfpathlineto{\pgfpoint{-22.258747\du}{6.330087\du}}
\pgfpathlineto{\pgfpoint{-22.265016\du}{6.332774\du}}
\pgfpathlineto{\pgfpoint{-22.270689\du}{6.335461\du}}
\pgfpathlineto{\pgfpoint{-22.276063\du}{6.337550\du}}
\pgfpathlineto{\pgfpoint{-22.280840\du}{6.340536\du}}
\pgfpathlineto{\pgfpoint{-22.285318\du}{6.342924\du}}
\pgfpathlineto{\pgfpoint{-22.289498\du}{6.345611\du}}
\pgfpathlineto{\pgfpoint{-22.292782\du}{6.348298\du}}
\pgfpathlineto{\pgfpoint{-22.296364\du}{6.350687\du}}
\pgfpathlineto{\pgfpoint{-22.298156\du}{6.354269\du}}
\pgfpathlineto{\pgfpoint{-22.300544\du}{6.356956\du}}
\pgfpathlineto{\pgfpoint{-22.302335\du}{6.359942\du}}
\pgfpathlineto{\pgfpoint{-22.303828\du}{6.363226\du}}
\pgfpathlineto{\pgfpoint{-22.304425\du}{6.365614\du}}
\pgfpathlineto{\pgfpoint{-22.304724\du}{6.369197\du}}
\pgfpathlineto{\pgfpoint{-22.304724\du}{6.369197\du}}
\pgfpathlineto{\pgfpoint{-22.304425\du}{6.372182\du}}
\pgfpathlineto{\pgfpoint{-22.303828\du}{6.374869\du}}
\pgfpathlineto{\pgfpoint{-22.302335\du}{6.377855\du}}
\pgfpathlineto{\pgfpoint{-22.300544\du}{6.381438\du}}
\pgfpathlineto{\pgfpoint{-22.298156\du}{6.384124\du}}
\pgfpathlineto{\pgfpoint{-22.296364\du}{6.387110\du}}
\pgfpathlineto{\pgfpoint{-22.292782\du}{6.389797\du}}
\pgfpathlineto{\pgfpoint{-22.289498\du}{6.392484\du}}
\pgfpathlineto{\pgfpoint{-22.285318\du}{6.395171\du}}
\pgfpathlineto{\pgfpoint{-22.280840\du}{6.397858\du}}
\pgfpathlineto{\pgfpoint{-22.276063\du}{6.400545\du}}
\pgfpathlineto{\pgfpoint{-22.270689\du}{6.402933\du}}
\pgfpathlineto{\pgfpoint{-22.265016\du}{6.405620\du}}
\pgfpathlineto{\pgfpoint{-22.258747\du}{6.407710\du}}
\pgfpathlineto{\pgfpoint{-22.252776\du}{6.410099\du}}
\pgfpathlineto{\pgfpoint{-22.245312\du}{6.412487\du}}
\pgfpathlineto{\pgfpoint{-22.238445\du}{6.414278\du}}
\pgfpathlineto{\pgfpoint{-22.230982\du}{6.415771\du}}
\pgfpathlineto{\pgfpoint{-22.223219\du}{6.417861\du}}
\pgfpathlineto{\pgfpoint{-22.215158\du}{6.419652\du}}
\pgfpathlineto{\pgfpoint{-22.206799\du}{6.420846\du}}
\pgfpathlineto{\pgfpoint{-22.197842\du}{6.422638\du}}
\pgfpathlineto{\pgfpoint{-22.189483\du}{6.423533\du}}
\pgfpathlineto{\pgfpoint{-22.180526\du}{6.425026\du}}
\pgfpathlineto{\pgfpoint{-22.171868\du}{6.425922\du}}
\pgfpathlineto{\pgfpoint{-22.162016\du}{6.426220\du}}
\pgfpathlineto{\pgfpoint{-22.152761\du}{6.427415\du}}
\pgfpathlineto{\pgfpoint{-22.143207\du}{6.427713\du}}
\pgfpathlineto{\pgfpoint{-22.133654\du}{6.428012\du}}
\pgfpathlineto{\pgfpoint{-22.123801\du}{6.428609\du}}
\pgfpathlineto{\pgfpoint{-22.114546\du}{6.428609\du}}
\pgfpathlineto{\pgfpoint{-22.114546\du}{6.428609\du}}
\pgfpathlineto{\pgfpoint{-22.104993\du}{6.428609\du}}
\pgfpathlineto{\pgfpoint{-22.094842\du}{6.428012\du}}
\pgfpathlineto{\pgfpoint{-22.085885\du}{6.427713\du}}
\pgfpathlineto{\pgfpoint{-22.076332\du}{6.427415\du}}
\pgfpathlineto{\pgfpoint{-22.066778\du}{6.426220\du}}
\pgfpathlineto{\pgfpoint{-22.057523\du}{6.425922\du}}
\pgfpathlineto{\pgfpoint{-22.048566\du}{6.425026\du}}
\pgfpathlineto{\pgfpoint{-22.039908\du}{6.423533\du}}
\pgfpathlineto{\pgfpoint{-22.030952\du}{6.422638\du}}
\pgfpathlineto{\pgfpoint{-22.022592\du}{6.420846\du}}
\pgfpathlineto{\pgfpoint{-22.014233\du}{6.419652\du}}
\pgfpathlineto{\pgfpoint{-22.005575\du}{6.417861\du}}
\pgfpathlineto{\pgfpoint{-21.997812\du}{6.415771\du}}
\pgfpathlineto{\pgfpoint{-21.990647\du}{6.414278\du}}
\pgfpathlineto{\pgfpoint{-21.983780\du}{6.412487\du}}
\pgfpathlineto{\pgfpoint{-21.976317\du}{6.410099\du}}
\pgfpathlineto{\pgfpoint{-21.970346\du}{6.407710\du}}
\pgfpathlineto{\pgfpoint{-21.964076\du}{6.405620\du}}
\pgfpathlineto{\pgfpoint{-21.958702\du}{6.402933\du}}
\pgfpathlineto{\pgfpoint{-21.953030\du}{6.400545\du}}
\pgfpathlineto{\pgfpoint{-21.948253\du}{6.397858\du}}
\pgfpathlineto{\pgfpoint{-21.943775\du}{6.395171\du}}
\pgfpathlineto{\pgfpoint{-21.939595\du}{6.392484\du}}
\pgfpathlineto{\pgfpoint{-21.936311\du}{6.389797\du}}
\pgfpathlineto{\pgfpoint{-21.933027\du}{6.387110\du}}
\pgfpathlineto{\pgfpoint{-21.930638\du}{6.384124\du}}
\pgfpathlineto{\pgfpoint{-21.928250\du}{6.381438\du}}
\pgfpathlineto{\pgfpoint{-21.926757\du}{6.377855\du}}
\pgfpathlineto{\pgfpoint{-21.925563\du}{6.374869\du}}
\pgfpathlineto{\pgfpoint{-21.924667\du}{6.372182\du}}
\pgfpathlineto{\pgfpoint{-21.924667\du}{6.369197\du}}
\pgfusepath{stroke}
\pgfsetbuttcap
\pgfsetmiterjoin
\pgfsetdash{}{0pt}
\definecolor{dialinecolor}{rgb}{0.000000, 0.470588, 0.666667}
\pgfsetfillcolor{dialinecolor}
\pgfpathmoveto{\pgfpoint{-22.304724\du}{6.285901\du}}
\pgfpathlineto{\pgfpoint{-22.304724\du}{6.380542\du}}
\pgfpathlineto{\pgfpoint{-21.912128\du}{6.380542\du}}
\pgfpathlineto{\pgfpoint{-21.912128\du}{6.285901\du}}
\pgfpathlineto{\pgfpoint{-22.304724\du}{6.285901\du}}
\pgfusepath{fill}
\pgfsetbuttcap
\pgfsetmiterjoin
\pgfsetdash{}{0pt}
\definecolor{dialinecolor}{rgb}{0.000000, 0.705882, 1.000000}
\pgfsetfillcolor{dialinecolor}
\pgfpathmoveto{\pgfpoint{-21.912427\du}{6.279333\du}}
\pgfpathlineto{\pgfpoint{-21.913024\du}{6.276049\du}}
\pgfpathlineto{\pgfpoint{-21.913919\du}{6.272466\du}}
\pgfpathlineto{\pgfpoint{-21.914815\du}{6.269481\du}}
\pgfpathlineto{\pgfpoint{-21.916606\du}{6.266196\du}}
\pgfpathlineto{\pgfpoint{-21.918995\du}{6.262912\du}}
\pgfpathlineto{\pgfpoint{-21.921682\du}{6.259628\du}}
\pgfpathlineto{\pgfpoint{-21.924667\du}{6.256643\du}}
\pgfpathlineto{\pgfpoint{-21.928250\du}{6.253359\du}}
\pgfpathlineto{\pgfpoint{-21.932430\du}{6.250672\du}}
\pgfpathlineto{\pgfpoint{-21.937206\du}{6.247686\du}}
\pgfpathlineto{\pgfpoint{-21.942282\du}{6.244701\du}}
\pgfpathlineto{\pgfpoint{-21.947954\du}{6.242014\du}}
\pgfpathlineto{\pgfpoint{-21.953627\du}{6.238730\du}}
\pgfpathlineto{\pgfpoint{-21.959896\du}{6.236938\du}}
\pgfpathlineto{\pgfpoint{-21.966464\du}{6.234251\du}}
\pgfpathlineto{\pgfpoint{-21.973630\du}{6.231863\du}}
\pgfpathlineto{\pgfpoint{-21.981094\du}{6.229773\du}}
\pgfpathlineto{\pgfpoint{-21.988557\du}{6.227683\du}}
\pgfpathlineto{\pgfpoint{-21.996917\du}{6.225593\du}}
\pgfpathlineto{\pgfpoint{-22.005276\du}{6.223504\du}}
\pgfpathlineto{\pgfpoint{-22.013337\du}{6.222011\du}}
\pgfpathlineto{\pgfpoint{-22.022592\du}{6.220518\du}}
\pgfpathlineto{\pgfpoint{-22.031250\du}{6.219324\du}}
\pgfpathlineto{\pgfpoint{-22.040804\du}{6.217831\du}}
\pgfpathlineto{\pgfpoint{-22.050059\du}{6.216935\du}}
\pgfpathlineto{\pgfpoint{-22.059314\du}{6.215741\du}}
\pgfpathlineto{\pgfpoint{-22.069166\du}{6.215144\du}}
\pgfpathlineto{\pgfpoint{-22.079019\du}{6.214547\du}}
\pgfpathlineto{\pgfpoint{-22.088871\du}{6.214248\du}}
\pgfpathlineto{\pgfpoint{-22.099022\du}{6.213950\du}}
\pgfpathlineto{\pgfpoint{-22.108874\du}{6.213950\du}}
\pgfpathlineto{\pgfpoint{-22.108874\du}{6.213950\du}}
\pgfpathlineto{\pgfpoint{-22.118726\du}{6.213950\du}}
\pgfpathlineto{\pgfpoint{-22.128578\du}{6.214248\du}}
\pgfpathlineto{\pgfpoint{-22.138729\du}{6.214547\du}}
\pgfpathlineto{\pgfpoint{-22.148581\du}{6.215144\du}}
\pgfpathlineto{\pgfpoint{-22.158135\du}{6.215741\du}}
\pgfpathlineto{\pgfpoint{-22.167689\du}{6.216935\du}}
\pgfpathlineto{\pgfpoint{-22.176944\du}{6.217831\du}}
\pgfpathlineto{\pgfpoint{-22.186199\du}{6.219324\du}}
\pgfpathlineto{\pgfpoint{-22.195155\du}{6.220518\du}}
\pgfpathlineto{\pgfpoint{-22.204112\du}{6.222011\du}}
\pgfpathlineto{\pgfpoint{-22.212471\du}{6.223504\du}}
\pgfpathlineto{\pgfpoint{-22.220831\du}{6.225593\du}}
\pgfpathlineto{\pgfpoint{-22.229190\du}{6.227683\du}}
\pgfpathlineto{\pgfpoint{-22.236356\du}{6.229773\du}}
\pgfpathlineto{\pgfpoint{-22.243819\du}{6.231863\du}}
\pgfpathlineto{\pgfpoint{-22.250985\du}{6.234251\du}}
\pgfpathlineto{\pgfpoint{-22.257851\du}{6.236938\du}}
\pgfpathlineto{\pgfpoint{-22.263822\du}{6.238730\du}}
\pgfpathlineto{\pgfpoint{-22.269793\du}{6.242014\du}}
\pgfpathlineto{\pgfpoint{-22.275466\du}{6.244701\du}}
\pgfpathlineto{\pgfpoint{-22.280541\du}{6.247686\du}}
\pgfpathlineto{\pgfpoint{-22.285318\du}{6.250672\du}}
\pgfpathlineto{\pgfpoint{-22.289498\du}{6.253359\du}}
\pgfpathlineto{\pgfpoint{-22.292782\du}{6.256643\du}}
\pgfpathlineto{\pgfpoint{-22.296066\du}{6.259628\du}}
\pgfpathlineto{\pgfpoint{-22.299051\du}{6.262912\du}}
\pgfpathlineto{\pgfpoint{-22.301141\du}{6.266196\du}}
\pgfpathlineto{\pgfpoint{-22.302634\du}{6.269481\du}}
\pgfpathlineto{\pgfpoint{-22.304425\du}{6.272466\du}}
\pgfpathlineto{\pgfpoint{-22.304724\du}{6.276049\du}}
\pgfpathlineto{\pgfpoint{-22.305022\du}{6.279333\du}}
\pgfpathlineto{\pgfpoint{-22.305022\du}{6.279333\du}}
\pgfpathlineto{\pgfpoint{-22.304724\du}{6.282617\du}}
\pgfpathlineto{\pgfpoint{-22.304425\du}{6.286199\du}}
\pgfpathlineto{\pgfpoint{-22.302634\du}{6.289185\du}}
\pgfpathlineto{\pgfpoint{-22.301141\du}{6.292469\du}}
\pgfpathlineto{\pgfpoint{-22.299051\du}{6.296052\du}}
\pgfpathlineto{\pgfpoint{-22.296066\du}{6.299037\du}}
\pgfpathlineto{\pgfpoint{-22.292782\du}{6.302023\du}}
\pgfpathlineto{\pgfpoint{-22.289498\du}{6.305307\du}}
\pgfpathlineto{\pgfpoint{-22.285318\du}{6.307695\du}}
\pgfpathlineto{\pgfpoint{-22.280541\du}{6.310979\du}}
\pgfpathlineto{\pgfpoint{-22.275466\du}{6.314263\du}}
\pgfpathlineto{\pgfpoint{-22.269793\du}{6.316950\du}}
\pgfpathlineto{\pgfpoint{-22.263822\du}{6.319637\du}}
\pgfpathlineto{\pgfpoint{-22.257851\du}{6.321727\du}}
\pgfpathlineto{\pgfpoint{-22.250985\du}{6.324713\du}}
\pgfpathlineto{\pgfpoint{-22.243819\du}{6.327101\du}}
\pgfpathlineto{\pgfpoint{-22.236356\du}{6.328892\du}}
\pgfpathlineto{\pgfpoint{-22.229190\du}{6.330982\du}}
\pgfpathlineto{\pgfpoint{-22.220831\du}{6.333072\du}}
\pgfpathlineto{\pgfpoint{-22.212471\du}{6.335162\du}}
\pgfpathlineto{\pgfpoint{-22.204112\du}{6.336655\du}}
\pgfpathlineto{\pgfpoint{-22.195155\du}{6.338147\du}}
\pgfpathlineto{\pgfpoint{-22.186199\du}{6.339342\du}}
\pgfpathlineto{\pgfpoint{-22.176944\du}{6.340834\du}}
\pgfpathlineto{\pgfpoint{-22.167689\du}{6.341730\du}}
\pgfpathlineto{\pgfpoint{-22.158135\du}{6.342924\du}}
\pgfpathlineto{\pgfpoint{-22.148581\du}{6.343223\du}}
\pgfpathlineto{\pgfpoint{-22.138729\du}{6.344119\du}}
\pgfpathlineto{\pgfpoint{-22.128578\du}{6.344119\du}}
\pgfpathlineto{\pgfpoint{-22.118726\du}{6.345014\du}}
\pgfpathlineto{\pgfpoint{-22.108874\du}{6.345014\du}}
\pgfpathlineto{\pgfpoint{-22.108874\du}{6.345014\du}}
\pgfpathlineto{\pgfpoint{-22.099022\du}{6.345014\du}}
\pgfpathlineto{\pgfpoint{-22.088871\du}{6.344119\du}}
\pgfpathlineto{\pgfpoint{-22.079019\du}{6.344119\du}}
\pgfpathlineto{\pgfpoint{-22.069166\du}{6.343223\du}}
\pgfpathlineto{\pgfpoint{-22.059314\du}{6.342924\du}}
\pgfpathlineto{\pgfpoint{-22.050059\du}{6.341730\du}}
\pgfpathlineto{\pgfpoint{-22.040804\du}{6.340834\du}}
\pgfpathlineto{\pgfpoint{-22.031250\du}{6.339342\du}}
\pgfpathlineto{\pgfpoint{-22.022592\du}{6.338147\du}}
\pgfpathlineto{\pgfpoint{-22.013337\du}{6.336655\du}}
\pgfpathlineto{\pgfpoint{-22.005276\du}{6.335162\du}}
\pgfpathlineto{\pgfpoint{-21.996917\du}{6.333072\du}}
\pgfpathlineto{\pgfpoint{-21.988557\du}{6.330982\du}}
\pgfpathlineto{\pgfpoint{-21.981094\du}{6.328892\du}}
\pgfpathlineto{\pgfpoint{-21.973630\du}{6.327101\du}}
\pgfpathlineto{\pgfpoint{-21.966464\du}{6.324713\du}}
\pgfpathlineto{\pgfpoint{-21.959896\du}{6.321727\du}}
\pgfpathlineto{\pgfpoint{-21.953627\du}{6.319637\du}}
\pgfpathlineto{\pgfpoint{-21.947954\du}{6.316950\du}}
\pgfpathlineto{\pgfpoint{-21.942282\du}{6.314263\du}}
\pgfpathlineto{\pgfpoint{-21.937206\du}{6.310979\du}}
\pgfpathlineto{\pgfpoint{-21.932430\du}{6.307695\du}}
\pgfpathlineto{\pgfpoint{-21.928250\du}{6.305307\du}}
\pgfpathlineto{\pgfpoint{-21.924667\du}{6.302023\du}}
\pgfpathlineto{\pgfpoint{-21.921682\du}{6.299037\du}}
\pgfpathlineto{\pgfpoint{-21.918995\du}{6.296052\du}}
\pgfpathlineto{\pgfpoint{-21.916606\du}{6.292469\du}}
\pgfpathlineto{\pgfpoint{-21.914815\du}{6.289185\du}}
\pgfpathlineto{\pgfpoint{-21.913919\du}{6.286199\du}}
\pgfpathlineto{\pgfpoint{-21.913024\du}{6.282617\du}}
\pgfpathlineto{\pgfpoint{-21.912427\du}{6.279333\du}}
\pgfusepath{fill}
\pgfsetbuttcap
\pgfsetmiterjoin
\pgfsetdash{}{0pt}
\definecolor{dialinecolor}{rgb}{0.666667, 0.901961, 1.000000}
\pgfsetstrokecolor{dialinecolor}
\pgfpathmoveto{\pgfpoint{-21.924667\du}{6.273660\du}}
\pgfpathlineto{\pgfpoint{-21.924667\du}{6.270376\du}}
\pgfpathlineto{\pgfpoint{-21.925563\du}{6.267988\du}}
\pgfpathlineto{\pgfpoint{-21.926757\du}{6.264704\du}}
\pgfpathlineto{\pgfpoint{-21.928250\du}{6.261121\du}}
\pgfpathlineto{\pgfpoint{-21.930638\du}{6.258733\du}}
\pgfpathlineto{\pgfpoint{-21.933027\du}{6.255449\du}}
\pgfpathlineto{\pgfpoint{-21.936311\du}{6.253060\du}}
\pgfpathlineto{\pgfpoint{-21.939595\du}{6.250373\du}}
\pgfpathlineto{\pgfpoint{-21.943775\du}{6.247686\du}}
\pgfpathlineto{\pgfpoint{-21.948253\du}{6.244999\du}}
\pgfpathlineto{\pgfpoint{-21.953030\du}{6.242312\du}}
\pgfpathlineto{\pgfpoint{-21.958702\du}{6.239625\du}}
\pgfpathlineto{\pgfpoint{-21.964076\du}{6.237237\du}}
\pgfpathlineto{\pgfpoint{-21.970346\du}{6.235147\du}}
\pgfpathlineto{\pgfpoint{-21.976317\du}{6.232759\du}}
\pgfpathlineto{\pgfpoint{-21.983780\du}{6.230370\du}}
\pgfpathlineto{\pgfpoint{-21.990647\du}{6.228579\du}}
\pgfpathlineto{\pgfpoint{-21.997812\du}{6.227086\du}}
\pgfpathlineto{\pgfpoint{-22.005575\du}{6.224698\du}}
\pgfpathlineto{\pgfpoint{-22.014233\du}{6.223205\du}}
\pgfpathlineto{\pgfpoint{-22.022592\du}{6.222011\du}}
\pgfpathlineto{\pgfpoint{-22.030952\du}{6.220219\du}}
\pgfpathlineto{\pgfpoint{-22.039908\du}{6.219324\du}}
\pgfpathlineto{\pgfpoint{-22.048566\du}{6.217831\du}}
\pgfpathlineto{\pgfpoint{-22.057523\du}{6.216935\du}}
\pgfpathlineto{\pgfpoint{-22.066778\du}{6.216338\du}}
\pgfpathlineto{\pgfpoint{-22.076332\du}{6.215443\du}}
\pgfpathlineto{\pgfpoint{-22.085885\du}{6.215144\du}}
\pgfpathlineto{\pgfpoint{-22.094842\du}{6.214845\du}}
\pgfpathlineto{\pgfpoint{-22.104993\du}{6.214248\du}}
\pgfpathlineto{\pgfpoint{-22.114546\du}{6.214248\du}}
\pgfpathlineto{\pgfpoint{-22.114546\du}{6.214248\du}}
\pgfpathlineto{\pgfpoint{-22.123801\du}{6.214248\du}}
\pgfpathlineto{\pgfpoint{-22.133654\du}{6.214845\du}}
\pgfpathlineto{\pgfpoint{-22.143207\du}{6.215144\du}}
\pgfpathlineto{\pgfpoint{-22.152761\du}{6.215443\du}}
\pgfpathlineto{\pgfpoint{-22.162016\du}{6.216338\du}}
\pgfpathlineto{\pgfpoint{-22.171868\du}{6.216935\du}}
\pgfpathlineto{\pgfpoint{-22.180526\du}{6.217831\du}}
\pgfpathlineto{\pgfpoint{-22.189483\du}{6.219324\du}}
\pgfpathlineto{\pgfpoint{-22.197842\du}{6.220219\du}}
\pgfpathlineto{\pgfpoint{-22.206799\du}{6.222011\du}}
\pgfpathlineto{\pgfpoint{-22.215158\du}{6.223205\du}}
\pgfpathlineto{\pgfpoint{-22.223219\du}{6.224698\du}}
\pgfpathlineto{\pgfpoint{-22.230982\du}{6.227086\du}}
\pgfpathlineto{\pgfpoint{-22.238445\du}{6.228579\du}}
\pgfpathlineto{\pgfpoint{-22.245312\du}{6.230370\du}}
\pgfpathlineto{\pgfpoint{-22.252776\du}{6.232759\du}}
\pgfpathlineto{\pgfpoint{-22.258747\du}{6.235147\du}}
\pgfpathlineto{\pgfpoint{-22.265016\du}{6.237237\du}}
\pgfpathlineto{\pgfpoint{-22.270689\du}{6.239625\du}}
\pgfpathlineto{\pgfpoint{-22.276063\du}{6.242312\du}}
\pgfpathlineto{\pgfpoint{-22.280840\du}{6.244999\du}}
\pgfpathlineto{\pgfpoint{-22.285318\du}{6.247686\du}}
\pgfpathlineto{\pgfpoint{-22.289498\du}{6.250373\du}}
\pgfpathlineto{\pgfpoint{-22.292782\du}{6.253060\du}}
\pgfpathlineto{\pgfpoint{-22.296364\du}{6.255449\du}}
\pgfpathlineto{\pgfpoint{-22.298156\du}{6.258733\du}}
\pgfpathlineto{\pgfpoint{-22.300544\du}{6.261121\du}}
\pgfpathlineto{\pgfpoint{-22.302335\du}{6.264704\du}}
\pgfpathlineto{\pgfpoint{-22.303828\du}{6.267988\du}}
\pgfpathlineto{\pgfpoint{-22.304425\du}{6.270376\du}}
\pgfpathlineto{\pgfpoint{-22.304724\du}{6.273660\du}}
\pgfpathlineto{\pgfpoint{-22.304724\du}{6.273660\du}}
\pgfpathlineto{\pgfpoint{-22.304425\du}{6.276944\du}}
\pgfpathlineto{\pgfpoint{-22.303828\du}{6.279631\du}}
\pgfpathlineto{\pgfpoint{-22.302335\du}{6.282617\du}}
\pgfpathlineto{\pgfpoint{-22.300544\du}{6.285901\du}}
\pgfpathlineto{\pgfpoint{-22.298156\du}{6.288588\du}}
\pgfpathlineto{\pgfpoint{-22.296364\du}{6.291872\du}}
\pgfpathlineto{\pgfpoint{-22.292782\du}{6.294559\du}}
\pgfpathlineto{\pgfpoint{-22.289498\du}{6.297246\du}}
\pgfpathlineto{\pgfpoint{-22.285318\du}{6.299933\du}}
\pgfpathlineto{\pgfpoint{-22.280840\du}{6.302620\du}}
\pgfpathlineto{\pgfpoint{-22.276063\du}{6.305307\du}}
\pgfpathlineto{\pgfpoint{-22.270689\du}{6.307397\du}}
\pgfpathlineto{\pgfpoint{-22.265016\du}{6.310084\du}}
\pgfpathlineto{\pgfpoint{-22.258747\du}{6.312472\du}}
\pgfpathlineto{\pgfpoint{-22.252776\du}{6.314860\du}}
\pgfpathlineto{\pgfpoint{-22.245312\du}{6.316950\du}}
\pgfpathlineto{\pgfpoint{-22.238445\du}{6.318742\du}}
\pgfpathlineto{\pgfpoint{-22.230982\du}{6.320533\du}}
\pgfpathlineto{\pgfpoint{-22.223219\du}{6.322324\du}}
\pgfpathlineto{\pgfpoint{-22.215158\du}{6.324414\du}}
\pgfpathlineto{\pgfpoint{-22.206799\du}{6.325608\du}}
\pgfpathlineto{\pgfpoint{-22.197842\du}{6.327400\du}}
\pgfpathlineto{\pgfpoint{-22.189483\du}{6.328295\du}}
\pgfpathlineto{\pgfpoint{-22.180526\du}{6.329788\du}}
\pgfpathlineto{\pgfpoint{-22.171868\du}{6.330087\du}}
\pgfpathlineto{\pgfpoint{-22.162016\du}{6.331281\du}}
\pgfpathlineto{\pgfpoint{-22.152761\du}{6.331579\du}}
\pgfpathlineto{\pgfpoint{-22.143207\du}{6.332475\du}}
\pgfpathlineto{\pgfpoint{-22.133654\du}{6.332774\du}}
\pgfpathlineto{\pgfpoint{-22.123801\du}{6.333072\du}}
\pgfpathlineto{\pgfpoint{-22.114546\du}{6.333072\du}}
\pgfpathlineto{\pgfpoint{-22.114546\du}{6.333072\du}}
\pgfpathlineto{\pgfpoint{-22.104993\du}{6.333072\du}}
\pgfpathlineto{\pgfpoint{-22.094842\du}{6.332774\du}}
\pgfpathlineto{\pgfpoint{-22.085885\du}{6.332475\du}}
\pgfpathlineto{\pgfpoint{-22.076332\du}{6.331579\du}}
\pgfpathlineto{\pgfpoint{-22.066778\du}{6.331281\du}}
\pgfpathlineto{\pgfpoint{-22.057523\du}{6.330087\du}}
\pgfpathlineto{\pgfpoint{-22.048566\du}{6.329788\du}}
\pgfpathlineto{\pgfpoint{-22.039908\du}{6.328295\du}}
\pgfpathlineto{\pgfpoint{-22.030952\du}{6.327400\du}}
\pgfpathlineto{\pgfpoint{-22.022592\du}{6.325608\du}}
\pgfpathlineto{\pgfpoint{-22.014233\du}{6.324414\du}}
\pgfpathlineto{\pgfpoint{-22.005575\du}{6.322324\du}}
\pgfpathlineto{\pgfpoint{-21.997812\du}{6.320533\du}}
\pgfpathlineto{\pgfpoint{-21.990647\du}{6.318742\du}}
\pgfpathlineto{\pgfpoint{-21.983780\du}{6.316950\du}}
\pgfpathlineto{\pgfpoint{-21.976317\du}{6.314860\du}}
\pgfpathlineto{\pgfpoint{-21.970346\du}{6.312472\du}}
\pgfpathlineto{\pgfpoint{-21.964076\du}{6.310084\du}}
\pgfpathlineto{\pgfpoint{-21.958702\du}{6.307397\du}}
\pgfpathlineto{\pgfpoint{-21.953030\du}{6.305307\du}}
\pgfpathlineto{\pgfpoint{-21.948253\du}{6.302620\du}}
\pgfpathlineto{\pgfpoint{-21.943775\du}{6.299933\du}}
\pgfpathlineto{\pgfpoint{-21.939595\du}{6.297246\du}}
\pgfpathlineto{\pgfpoint{-21.936311\du}{6.294559\du}}
\pgfpathlineto{\pgfpoint{-21.933027\du}{6.291872\du}}
\pgfpathlineto{\pgfpoint{-21.930638\du}{6.288588\du}}
\pgfpathlineto{\pgfpoint{-21.928250\du}{6.285901\du}}
\pgfpathlineto{\pgfpoint{-21.926757\du}{6.282617\du}}
\pgfpathlineto{\pgfpoint{-21.925563\du}{6.279631\du}}
\pgfpathlineto{\pgfpoint{-21.924667\du}{6.276944\du}}
\pgfpathlineto{\pgfpoint{-21.924667\du}{6.273660\du}}
\pgfusepath{stroke}
\pgfsetbuttcap
\pgfsetmiterjoin
\pgfsetdash{}{0pt}
\definecolor{dialinecolor}{rgb}{0.000000, 0.000000, 0.000000}
\pgfsetfillcolor{dialinecolor}
\pgfpathmoveto{\pgfpoint{-22.102306\du}{6.262017\du}}
\pgfpathlineto{\pgfpoint{-22.079019\du}{6.273660\du}}
\pgfpathlineto{\pgfpoint{-22.007068\du}{6.250075\du}}
\pgfpathlineto{\pgfpoint{-21.971540\du}{6.262017\du}}
\pgfpathlineto{\pgfpoint{-21.995125\du}{6.238133\du}}
\pgfpathlineto{\pgfpoint{-22.079019\du}{6.238133\du}}
\pgfpathlineto{\pgfpoint{-22.043192\du}{6.238133\du}}
\pgfpathlineto{\pgfpoint{-22.102306\du}{6.262017\du}}
\pgfusepath{fill}
\pgfsetbuttcap
\pgfsetmiterjoin
\pgfsetdash{}{0pt}
\definecolor{dialinecolor}{rgb}{0.000000, 0.000000, 0.000000}
\pgfsetfillcolor{dialinecolor}
\pgfpathmoveto{\pgfpoint{-22.114546\du}{6.297544\du}}
\pgfpathlineto{\pgfpoint{-22.150373\du}{6.285901\du}}
\pgfpathlineto{\pgfpoint{-22.209486\du}{6.309785\du}}
\pgfpathlineto{\pgfpoint{-22.245312\du}{6.297544\du}}
\pgfpathlineto{\pgfpoint{-22.233071\du}{6.333072\du}}
\pgfpathlineto{\pgfpoint{-22.150373\du}{6.333072\du}}
\pgfpathlineto{\pgfpoint{-22.185900\du}{6.321727\du}}
\pgfpathlineto{\pgfpoint{-22.114546\du}{6.297544\du}}
\pgfusepath{fill}
\pgfsetbuttcap
\pgfsetmiterjoin
\pgfsetdash{}{0pt}
\definecolor{dialinecolor}{rgb}{0.000000, 0.000000, 0.000000}
\pgfsetfillcolor{dialinecolor}
\pgfpathmoveto{\pgfpoint{-22.233071\du}{6.238133\du}}
\pgfpathlineto{\pgfpoint{-22.209486\du}{6.225593\du}}
\pgfpathlineto{\pgfpoint{-22.137833\du}{6.262017\du}}
\pgfpathlineto{\pgfpoint{-22.114546\du}{6.250075\du}}
\pgfpathlineto{\pgfpoint{-22.126190\du}{6.273660\du}}
\pgfpathlineto{\pgfpoint{-22.209486\du}{6.273660\du}}
\pgfpathlineto{\pgfpoint{-22.173660\du}{6.262017\du}}
\pgfpathlineto{\pgfpoint{-22.233071\du}{6.238133\du}}
\pgfusepath{fill}
\pgfsetbuttcap
\pgfsetmiterjoin
\pgfsetdash{}{0pt}
\definecolor{dialinecolor}{rgb}{0.000000, 0.000000, 0.000000}
\pgfsetfillcolor{dialinecolor}
\pgfpathmoveto{\pgfpoint{-21.983780\du}{6.321727\du}}
\pgfpathlineto{\pgfpoint{-22.007068\du}{6.333072\du}}
\pgfpathlineto{\pgfpoint{-22.079019\du}{6.297544\du}}
\pgfpathlineto{\pgfpoint{-22.114546\du}{6.309785\du}}
\pgfpathlineto{\pgfpoint{-22.090662\du}{6.285901\du}}
\pgfpathlineto{\pgfpoint{-22.007068\du}{6.285901\du}}
\pgfpathlineto{\pgfpoint{-22.043192\du}{6.297544\du}}
\pgfpathlineto{\pgfpoint{-21.983780\du}{6.321727\du}}
\pgfusepath{fill}
\pgfsetbuttcap
\pgfsetmiterjoin
\pgfsetdash{}{0pt}
\definecolor{dialinecolor}{rgb}{1.000000, 1.000000, 1.000000}
\pgfsetfillcolor{dialinecolor}
\pgfpathmoveto{\pgfpoint{-22.102306\du}{6.273660\du}}
\pgfpathlineto{\pgfpoint{-22.079019\du}{6.273660\du}}
\pgfpathlineto{\pgfpoint{-22.007068\du}{6.250075\du}}
\pgfpathlineto{\pgfpoint{-21.971540\du}{6.262017\du}}
\pgfpathlineto{\pgfpoint{-21.995125\du}{6.238133\du}}
\pgfpathlineto{\pgfpoint{-22.066778\du}{6.238133\du}}
\pgfpathlineto{\pgfpoint{-22.043192\du}{6.238133\du}}
\pgfpathlineto{\pgfpoint{-22.102306\du}{6.273660\du}}
\pgfusepath{fill}
\pgfsetbuttcap
\pgfsetmiterjoin
\pgfsetdash{}{0pt}
\definecolor{dialinecolor}{rgb}{1.000000, 1.000000, 1.000000}
\pgfsetfillcolor{dialinecolor}
\pgfpathmoveto{\pgfpoint{-22.114546\du}{6.297544\du}}
\pgfpathlineto{\pgfpoint{-22.137833\du}{6.285901\du}}
\pgfpathlineto{\pgfpoint{-22.209486\du}{6.309785\du}}
\pgfpathlineto{\pgfpoint{-22.245312\du}{6.309785\du}}
\pgfpathlineto{\pgfpoint{-22.221726\du}{6.333072\du}}
\pgfpathlineto{\pgfpoint{-22.137833\du}{6.333072\du}}
\pgfpathlineto{\pgfpoint{-22.173660\du}{6.321727\du}}
\pgfpathlineto{\pgfpoint{-22.114546\du}{6.297544\du}}
\pgfusepath{fill}
\pgfsetbuttcap
\pgfsetmiterjoin
\pgfsetdash{}{0pt}
\definecolor{dialinecolor}{rgb}{1.000000, 1.000000, 1.000000}
\pgfsetfillcolor{dialinecolor}
\pgfpathmoveto{\pgfpoint{-22.233071\du}{6.238133\du}}
\pgfpathlineto{\pgfpoint{-22.209486\du}{6.238133\du}}
\pgfpathlineto{\pgfpoint{-22.137833\du}{6.262017\du}}
\pgfpathlineto{\pgfpoint{-22.102306\du}{6.250075\du}}
\pgfpathlineto{\pgfpoint{-22.126190\du}{6.273660\du}}
\pgfpathlineto{\pgfpoint{-22.209486\du}{6.273660\du}}
\pgfpathlineto{\pgfpoint{-22.173660\du}{6.273660\du}}
\pgfpathlineto{\pgfpoint{-22.233071\du}{6.238133\du}}
\pgfusepath{fill}
\pgfsetbuttcap
\pgfsetmiterjoin
\pgfsetdash{}{0pt}
\definecolor{dialinecolor}{rgb}{1.000000, 1.000000, 1.000000}
\pgfsetfillcolor{dialinecolor}
\pgfpathmoveto{\pgfpoint{-21.983780\du}{6.321727\du}}
\pgfpathlineto{\pgfpoint{-22.007068\du}{6.333072\du}}
\pgfpathlineto{\pgfpoint{-22.079019\du}{6.309785\du}}
\pgfpathlineto{\pgfpoint{-22.102306\du}{6.309785\du}}
\pgfpathlineto{\pgfpoint{-22.090662\du}{6.285901\du}}
\pgfpathlineto{\pgfpoint{-22.007068\du}{6.285901\du}}
\pgfpathlineto{\pgfpoint{-22.043192\du}{6.297544\du}}
\pgfpathlineto{\pgfpoint{-21.983780\du}{6.321727\du}}
\pgfusepath{fill}
\pgfsetbuttcap
\pgfsetmiterjoin
\pgfsetdash{}{0pt}
\definecolor{dialinecolor}{rgb}{0.666667, 0.901961, 1.000000}
\pgfsetstrokecolor{dialinecolor}
\pgfpathmoveto{\pgfpoint{-22.304724\du}{6.273660\du}}
\pgfpathlineto{\pgfpoint{-22.304724\du}{6.368600\du}}
\pgfusepath{stroke}
\pgfsetbuttcap
\pgfsetmiterjoin
\pgfsetdash{}{0pt}
\definecolor{dialinecolor}{rgb}{0.666667, 0.901961, 1.000000}
\pgfsetstrokecolor{dialinecolor}
\pgfpathmoveto{\pgfpoint{-21.924070\du}{6.273660\du}}
\pgfpathlineto{\pgfpoint{-21.924070\du}{6.368600\du}}
\pgfusepath{stroke}
\pgfsetbuttcap
\pgfsetmiterjoin
\pgfsetdash{}{0pt}
\definecolor{dialinecolor}{rgb}{0.000000, 0.470588, 0.666667}
\pgfsetfillcolor{dialinecolor}
\pgfpathmoveto{\pgfpoint{-21.198290\du}{6.374869\du}}
\pgfpathlineto{\pgfpoint{-21.198290\du}{6.371585\du}}
\pgfpathlineto{\pgfpoint{-21.199186\du}{6.368301\du}}
\pgfpathlineto{\pgfpoint{-21.200380\du}{6.365017\du}}
\pgfpathlineto{\pgfpoint{-21.201873\du}{6.361733\du}}
\pgfpathlineto{\pgfpoint{-21.204261\du}{6.358748\du}}
\pgfpathlineto{\pgfpoint{-21.206948\du}{6.355464\du}}
\pgfpathlineto{\pgfpoint{-21.209934\du}{6.352179\du}}
\pgfpathlineto{\pgfpoint{-21.214113\du}{6.349194\du}}
\pgfpathlineto{\pgfpoint{-21.217696\du}{6.346208\du}}
\pgfpathlineto{\pgfpoint{-21.222473\du}{6.343223\du}}
\pgfpathlineto{\pgfpoint{-21.227548\du}{6.340536\du}}
\pgfpathlineto{\pgfpoint{-21.232922\du}{6.337550\du}}
\pgfpathlineto{\pgfpoint{-21.239192\du}{6.335162\du}}
\pgfpathlineto{\pgfpoint{-21.245163\du}{6.332475\du}}
\pgfpathlineto{\pgfpoint{-21.252029\du}{6.330087\du}}
\pgfpathlineto{\pgfpoint{-21.258896\du}{6.327698\du}}
\pgfpathlineto{\pgfpoint{-21.266360\du}{6.325608\du}}
\pgfpathlineto{\pgfpoint{-21.273824\du}{6.323220\du}}
\pgfpathlineto{\pgfpoint{-21.281885\du}{6.321727\du}}
\pgfpathlineto{\pgfpoint{-21.290543\du}{6.319637\du}}
\pgfpathlineto{\pgfpoint{-21.298902\du}{6.317846\du}}
\pgfpathlineto{\pgfpoint{-21.307560\du}{6.316652\du}}
\pgfpathlineto{\pgfpoint{-21.316815\du}{6.315159\du}}
\pgfpathlineto{\pgfpoint{-21.326070\du}{6.313965\du}}
\pgfpathlineto{\pgfpoint{-21.335325\du}{6.312771\du}}
\pgfpathlineto{\pgfpoint{-21.344879\du}{6.311875\du}}
\pgfpathlineto{\pgfpoint{-21.355030\du}{6.310979\du}}
\pgfpathlineto{\pgfpoint{-21.364584\du}{6.310382\du}}
\pgfpathlineto{\pgfpoint{-21.374137\du}{6.310084\du}}
\pgfpathlineto{\pgfpoint{-21.384288\du}{6.309785\du}}
\pgfpathlineto{\pgfpoint{-21.394140\du}{6.309785\du}}
\pgfpathlineto{\pgfpoint{-21.394140\du}{6.309785\du}}
\pgfpathlineto{\pgfpoint{-21.403992\du}{6.309785\du}}
\pgfpathlineto{\pgfpoint{-21.414143\du}{6.310084\du}}
\pgfpathlineto{\pgfpoint{-21.423995\du}{6.310382\du}}
\pgfpathlineto{\pgfpoint{-21.433848\du}{6.310979\du}}
\pgfpathlineto{\pgfpoint{-21.443401\du}{6.311875\du}}
\pgfpathlineto{\pgfpoint{-21.452955\du}{6.312771\du}}
\pgfpathlineto{\pgfpoint{-21.462509\du}{6.313965\du}}
\pgfpathlineto{\pgfpoint{-21.471465\du}{6.315159\du}}
\pgfpathlineto{\pgfpoint{-21.480720\du}{6.316652\du}}
\pgfpathlineto{\pgfpoint{-21.489378\du}{6.317846\du}}
\pgfpathlineto{\pgfpoint{-21.498036\du}{6.319637\du}}
\pgfpathlineto{\pgfpoint{-21.506396\du}{6.321727\du}}
\pgfpathlineto{\pgfpoint{-21.514457\du}{6.323220\du}}
\pgfpathlineto{\pgfpoint{-21.521921\du}{6.325608\du}}
\pgfpathlineto{\pgfpoint{-21.529384\du}{6.327698\du}}
\pgfpathlineto{\pgfpoint{-21.536550\du}{6.330087\du}}
\pgfpathlineto{\pgfpoint{-21.543118\du}{6.332475\du}}
\pgfpathlineto{\pgfpoint{-21.549387\du}{6.335162\du}}
\pgfpathlineto{\pgfpoint{-21.555358\du}{6.337550\du}}
\pgfpathlineto{\pgfpoint{-21.560732\du}{6.340536\du}}
\pgfpathlineto{\pgfpoint{-21.565808\du}{6.343223\du}}
\pgfpathlineto{\pgfpoint{-21.570584\du}{6.346208\du}}
\pgfpathlineto{\pgfpoint{-21.574466\du}{6.349194\du}}
\pgfpathlineto{\pgfpoint{-21.578347\du}{6.352179\du}}
\pgfpathlineto{\pgfpoint{-21.581332\du}{6.355464\du}}
\pgfpathlineto{\pgfpoint{-21.584019\du}{6.358748\du}}
\pgfpathlineto{\pgfpoint{-21.586408\du}{6.361733\du}}
\pgfpathlineto{\pgfpoint{-21.588199\du}{6.365017\du}}
\pgfpathlineto{\pgfpoint{-21.589692\du}{6.368301\du}}
\pgfpathlineto{\pgfpoint{-21.590587\du}{6.371585\du}}
\pgfpathlineto{\pgfpoint{-21.590587\du}{6.374869\du}}
\pgfpathlineto{\pgfpoint{-21.590587\du}{6.374869\du}}
\pgfpathlineto{\pgfpoint{-21.590587\du}{6.378452\du}}
\pgfpathlineto{\pgfpoint{-21.589692\du}{6.381438\du}}
\pgfpathlineto{\pgfpoint{-21.588199\du}{6.384722\du}}
\pgfpathlineto{\pgfpoint{-21.586408\du}{6.388006\du}}
\pgfpathlineto{\pgfpoint{-21.584019\du}{6.391290\du}}
\pgfpathlineto{\pgfpoint{-21.581332\du}{6.394574\du}}
\pgfpathlineto{\pgfpoint{-21.578347\du}{6.397559\du}}
\pgfpathlineto{\pgfpoint{-21.574466\du}{6.400545\du}}
\pgfpathlineto{\pgfpoint{-21.570584\du}{6.403530\du}}
\pgfpathlineto{\pgfpoint{-21.565808\du}{6.406814\du}}
\pgfpathlineto{\pgfpoint{-21.560732\du}{6.409501\du}}
\pgfpathlineto{\pgfpoint{-21.555358\du}{6.412188\du}}
\pgfpathlineto{\pgfpoint{-21.549387\du}{6.414875\du}}
\pgfpathlineto{\pgfpoint{-21.543118\du}{6.417562\du}}
\pgfpathlineto{\pgfpoint{-21.536550\du}{6.419951\du}}
\pgfpathlineto{\pgfpoint{-21.529384\du}{6.422339\du}}
\pgfpathlineto{\pgfpoint{-21.521921\du}{6.424130\du}}
\pgfpathlineto{\pgfpoint{-21.514457\du}{6.426220\du}}
\pgfpathlineto{\pgfpoint{-21.506396\du}{6.428310\du}}
\pgfpathlineto{\pgfpoint{-21.498036\du}{6.430400\du}}
\pgfpathlineto{\pgfpoint{-21.489378\du}{6.431893\du}}
\pgfpathlineto{\pgfpoint{-21.480720\du}{6.433386\du}}
\pgfpathlineto{\pgfpoint{-21.471465\du}{6.434878\du}}
\pgfpathlineto{\pgfpoint{-21.462509\du}{6.436073\du}}
\pgfpathlineto{\pgfpoint{-21.452955\du}{6.436968\du}}
\pgfpathlineto{\pgfpoint{-21.443401\du}{6.438162\du}}
\pgfpathlineto{\pgfpoint{-21.433848\du}{6.438760\du}}
\pgfpathlineto{\pgfpoint{-21.423995\du}{6.439357\du}}
\pgfpathlineto{\pgfpoint{-21.414143\du}{6.439954\du}}
\pgfpathlineto{\pgfpoint{-21.403992\du}{6.440252\du}}
\pgfpathlineto{\pgfpoint{-21.394140\du}{6.440252\du}}
\pgfpathlineto{\pgfpoint{-21.394140\du}{6.440252\du}}
\pgfpathlineto{\pgfpoint{-21.384288\du}{6.440252\du}}
\pgfpathlineto{\pgfpoint{-21.374137\du}{6.439954\du}}
\pgfpathlineto{\pgfpoint{-21.364584\du}{6.439357\du}}
\pgfpathlineto{\pgfpoint{-21.355030\du}{6.438760\du}}
\pgfpathlineto{\pgfpoint{-21.344879\du}{6.438162\du}}
\pgfpathlineto{\pgfpoint{-21.335325\du}{6.436968\du}}
\pgfpathlineto{\pgfpoint{-21.326070\du}{6.436073\du}}
\pgfpathlineto{\pgfpoint{-21.316815\du}{6.434878\du}}
\pgfpathlineto{\pgfpoint{-21.307560\du}{6.433386\du}}
\pgfpathlineto{\pgfpoint{-21.298902\du}{6.431893\du}}
\pgfpathlineto{\pgfpoint{-21.290543\du}{6.430400\du}}
\pgfpathlineto{\pgfpoint{-21.281885\du}{6.428310\du}}
\pgfpathlineto{\pgfpoint{-21.273824\du}{6.426220\du}}
\pgfpathlineto{\pgfpoint{-21.266360\du}{6.424130\du}}
\pgfpathlineto{\pgfpoint{-21.258896\du}{6.422339\du}}
\pgfpathlineto{\pgfpoint{-21.252029\du}{6.419951\du}}
\pgfpathlineto{\pgfpoint{-21.245163\du}{6.417562\du}}
\pgfpathlineto{\pgfpoint{-21.239192\du}{6.414875\du}}
\pgfpathlineto{\pgfpoint{-21.232922\du}{6.412188\du}}
\pgfpathlineto{\pgfpoint{-21.227548\du}{6.409501\du}}
\pgfpathlineto{\pgfpoint{-21.222473\du}{6.406814\du}}
\pgfpathlineto{\pgfpoint{-21.217696\du}{6.403530\du}}
\pgfpathlineto{\pgfpoint{-21.214113\du}{6.400545\du}}
\pgfpathlineto{\pgfpoint{-21.209934\du}{6.397559\du}}
\pgfpathlineto{\pgfpoint{-21.206948\du}{6.394574\du}}
\pgfpathlineto{\pgfpoint{-21.204261\du}{6.391290\du}}
\pgfpathlineto{\pgfpoint{-21.201873\du}{6.388006\du}}
\pgfpathlineto{\pgfpoint{-21.200380\du}{6.384722\du}}
\pgfpathlineto{\pgfpoint{-21.199186\du}{6.381438\du}}
\pgfpathlineto{\pgfpoint{-21.198290\du}{6.378452\du}}
\pgfpathlineto{\pgfpoint{-21.198290\du}{6.374869\du}}
\pgfusepath{fill}
\pgfsetbuttcap
\pgfsetmiterjoin
\pgfsetdash{}{0pt}
\definecolor{dialinecolor}{rgb}{0.666667, 0.901961, 1.000000}
\pgfsetstrokecolor{dialinecolor}
\pgfpathmoveto{\pgfpoint{-21.209934\du}{6.369197\du}}
\pgfpathlineto{\pgfpoint{-21.209934\du}{6.365614\du}}
\pgfpathlineto{\pgfpoint{-21.210829\du}{6.363226\du}}
\pgfpathlineto{\pgfpoint{-21.212024\du}{6.359942\du}}
\pgfpathlineto{\pgfpoint{-21.214113\du}{6.356956\du}}
\pgfpathlineto{\pgfpoint{-21.216203\du}{6.354269\du}}
\pgfpathlineto{\pgfpoint{-21.218293\du}{6.350687\du}}
\pgfpathlineto{\pgfpoint{-21.221577\du}{6.348298\du}}
\pgfpathlineto{\pgfpoint{-21.224861\du}{6.345611\du}}
\pgfpathlineto{\pgfpoint{-21.229041\du}{6.342924\du}}
\pgfpathlineto{\pgfpoint{-21.233818\du}{6.340536\du}}
\pgfpathlineto{\pgfpoint{-21.238296\du}{6.337550\du}}
\pgfpathlineto{\pgfpoint{-21.243969\du}{6.335461\du}}
\pgfpathlineto{\pgfpoint{-21.249940\du}{6.332774\du}}
\pgfpathlineto{\pgfpoint{-21.255612\du}{6.330087\du}}
\pgfpathlineto{\pgfpoint{-21.262180\du}{6.328295\du}}
\pgfpathlineto{\pgfpoint{-21.269047\du}{6.325907\du}}
\pgfpathlineto{\pgfpoint{-21.275914\du}{6.324414\du}}
\pgfpathlineto{\pgfpoint{-21.283676\du}{6.322324\du}}
\pgfpathlineto{\pgfpoint{-21.291140\du}{6.320533\du}}
\pgfpathlineto{\pgfpoint{-21.299499\du}{6.318742\du}}
\pgfpathlineto{\pgfpoint{-21.308157\du}{6.317547\du}}
\pgfpathlineto{\pgfpoint{-21.316517\du}{6.315159\du}}
\pgfpathlineto{\pgfpoint{-21.325175\du}{6.314860\du}}
\pgfpathlineto{\pgfpoint{-21.334131\du}{6.313368\du}}
\pgfpathlineto{\pgfpoint{-21.342789\du}{6.312472\du}}
\pgfpathlineto{\pgfpoint{-21.352343\du}{6.311576\du}}
\pgfpathlineto{\pgfpoint{-21.361598\du}{6.310979\du}}
\pgfpathlineto{\pgfpoint{-21.371152\du}{6.310681\du}}
\pgfpathlineto{\pgfpoint{-21.380705\du}{6.310084\du}}
\pgfpathlineto{\pgfpoint{-21.390856\du}{6.309785\du}}
\pgfpathlineto{\pgfpoint{-21.400708\du}{6.309785\du}}
\pgfpathlineto{\pgfpoint{-21.400708\du}{6.309785\du}}
\pgfpathlineto{\pgfpoint{-21.409665\du}{6.309785\du}}
\pgfpathlineto{\pgfpoint{-21.419517\du}{6.310084\du}}
\pgfpathlineto{\pgfpoint{-21.429071\du}{6.310681\du}}
\pgfpathlineto{\pgfpoint{-21.438326\du}{6.310979\du}}
\pgfpathlineto{\pgfpoint{-21.447880\du}{6.311576\du}}
\pgfpathlineto{\pgfpoint{-21.457433\du}{6.312472\du}}
\pgfpathlineto{\pgfpoint{-21.466390\du}{6.313368\du}}
\pgfpathlineto{\pgfpoint{-21.475346\du}{6.314860\du}}
\pgfpathlineto{\pgfpoint{-21.484004\du}{6.315159\du}}
\pgfpathlineto{\pgfpoint{-21.492364\du}{6.317547\du}}
\pgfpathlineto{\pgfpoint{-21.501022\du}{6.318742\du}}
\pgfpathlineto{\pgfpoint{-21.509083\du}{6.320533\du}}
\pgfpathlineto{\pgfpoint{-21.516845\du}{6.322324\du}}
\pgfpathlineto{\pgfpoint{-21.524309\du}{6.324414\du}}
\pgfpathlineto{\pgfpoint{-21.530877\du}{6.325907\du}}
\pgfpathlineto{\pgfpoint{-21.538341\du}{6.328295\du}}
\pgfpathlineto{\pgfpoint{-21.544610\du}{6.330087\du}}
\pgfpathlineto{\pgfpoint{-21.550582\du}{6.332774\du}}
\pgfpathlineto{\pgfpoint{-21.556553\du}{6.335461\du}}
\pgfpathlineto{\pgfpoint{-21.561628\du}{6.337550\du}}
\pgfpathlineto{\pgfpoint{-21.566703\du}{6.340536\du}}
\pgfpathlineto{\pgfpoint{-21.571480\du}{6.342924\du}}
\pgfpathlineto{\pgfpoint{-21.575660\du}{6.345611\du}}
\pgfpathlineto{\pgfpoint{-21.578944\du}{6.348298\du}}
\pgfpathlineto{\pgfpoint{-21.582228\du}{6.350687\du}}
\pgfpathlineto{\pgfpoint{-21.584616\du}{6.354269\du}}
\pgfpathlineto{\pgfpoint{-21.586408\du}{6.356956\du}}
\pgfpathlineto{\pgfpoint{-21.588796\du}{6.359942\du}}
\pgfpathlineto{\pgfpoint{-21.589692\du}{6.363226\du}}
\pgfpathlineto{\pgfpoint{-21.590587\du}{6.365614\du}}
\pgfpathlineto{\pgfpoint{-21.590587\du}{6.369197\du}}
\pgfpathlineto{\pgfpoint{-21.590587\du}{6.369197\du}}
\pgfpathlineto{\pgfpoint{-21.590587\du}{6.372182\du}}
\pgfpathlineto{\pgfpoint{-21.589692\du}{6.374869\du}}
\pgfpathlineto{\pgfpoint{-21.588796\du}{6.377855\du}}
\pgfpathlineto{\pgfpoint{-21.586408\du}{6.381438\du}}
\pgfpathlineto{\pgfpoint{-21.584616\du}{6.384124\du}}
\pgfpathlineto{\pgfpoint{-21.582228\du}{6.387110\du}}
\pgfpathlineto{\pgfpoint{-21.578944\du}{6.389797\du}}
\pgfpathlineto{\pgfpoint{-21.575660\du}{6.392484\du}}
\pgfpathlineto{\pgfpoint{-21.571480\du}{6.395171\du}}
\pgfpathlineto{\pgfpoint{-21.566703\du}{6.397858\du}}
\pgfpathlineto{\pgfpoint{-21.561628\du}{6.400545\du}}
\pgfpathlineto{\pgfpoint{-21.556553\du}{6.402933\du}}
\pgfpathlineto{\pgfpoint{-21.550582\du}{6.405620\du}}
\pgfpathlineto{\pgfpoint{-21.544610\du}{6.407710\du}}
\pgfpathlineto{\pgfpoint{-21.538341\du}{6.410099\du}}
\pgfpathlineto{\pgfpoint{-21.530877\du}{6.412487\du}}
\pgfpathlineto{\pgfpoint{-21.524309\du}{6.414278\du}}
\pgfpathlineto{\pgfpoint{-21.516845\du}{6.415771\du}}
\pgfpathlineto{\pgfpoint{-21.509083\du}{6.417861\du}}
\pgfpathlineto{\pgfpoint{-21.501022\du}{6.419652\du}}
\pgfpathlineto{\pgfpoint{-21.492364\du}{6.420846\du}}
\pgfpathlineto{\pgfpoint{-21.484004\du}{6.422638\du}}
\pgfpathlineto{\pgfpoint{-21.475346\du}{6.423533\du}}
\pgfpathlineto{\pgfpoint{-21.466390\du}{6.425026\du}}
\pgfpathlineto{\pgfpoint{-21.457433\du}{6.425922\du}}
\pgfpathlineto{\pgfpoint{-21.447880\du}{6.426220\du}}
\pgfpathlineto{\pgfpoint{-21.438326\du}{6.427415\du}}
\pgfpathlineto{\pgfpoint{-21.429071\du}{6.427713\du}}
\pgfpathlineto{\pgfpoint{-21.419517\du}{6.428012\du}}
\pgfpathlineto{\pgfpoint{-21.409665\du}{6.428609\du}}
\pgfpathlineto{\pgfpoint{-21.400708\du}{6.428609\du}}
\pgfpathlineto{\pgfpoint{-21.400708\du}{6.428609\du}}
\pgfpathlineto{\pgfpoint{-21.390856\du}{6.428609\du}}
\pgfpathlineto{\pgfpoint{-21.380705\du}{6.428012\du}}
\pgfpathlineto{\pgfpoint{-21.371152\du}{6.427713\du}}
\pgfpathlineto{\pgfpoint{-21.361598\du}{6.427415\du}}
\pgfpathlineto{\pgfpoint{-21.352343\du}{6.426220\du}}
\pgfpathlineto{\pgfpoint{-21.342789\du}{6.425922\du}}
\pgfpathlineto{\pgfpoint{-21.334131\du}{6.425026\du}}
\pgfpathlineto{\pgfpoint{-21.325175\du}{6.423533\du}}
\pgfpathlineto{\pgfpoint{-21.316517\du}{6.422638\du}}
\pgfpathlineto{\pgfpoint{-21.308157\du}{6.420846\du}}
\pgfpathlineto{\pgfpoint{-21.299499\du}{6.419652\du}}
\pgfpathlineto{\pgfpoint{-21.291140\du}{6.417861\du}}
\pgfpathlineto{\pgfpoint{-21.283676\du}{6.415771\du}}
\pgfpathlineto{\pgfpoint{-21.275914\du}{6.414278\du}}
\pgfpathlineto{\pgfpoint{-21.269047\du}{6.412487\du}}
\pgfpathlineto{\pgfpoint{-21.262180\du}{6.410099\du}}
\pgfpathlineto{\pgfpoint{-21.255612\du}{6.407710\du}}
\pgfpathlineto{\pgfpoint{-21.249940\du}{6.405620\du}}
\pgfpathlineto{\pgfpoint{-21.243969\du}{6.402933\du}}
\pgfpathlineto{\pgfpoint{-21.238296\du}{6.400545\du}}
\pgfpathlineto{\pgfpoint{-21.233818\du}{6.397858\du}}
\pgfpathlineto{\pgfpoint{-21.229041\du}{6.395171\du}}
\pgfpathlineto{\pgfpoint{-21.224861\du}{6.392484\du}}
\pgfpathlineto{\pgfpoint{-21.221577\du}{6.389797\du}}
\pgfpathlineto{\pgfpoint{-21.218293\du}{6.387110\du}}
\pgfpathlineto{\pgfpoint{-21.216203\du}{6.384124\du}}
\pgfpathlineto{\pgfpoint{-21.214113\du}{6.381438\du}}
\pgfpathlineto{\pgfpoint{-21.212024\du}{6.377855\du}}
\pgfpathlineto{\pgfpoint{-21.210829\du}{6.374869\du}}
\pgfpathlineto{\pgfpoint{-21.209934\du}{6.372182\du}}
\pgfpathlineto{\pgfpoint{-21.209934\du}{6.369197\du}}
\pgfusepath{stroke}
\pgfsetbuttcap
\pgfsetmiterjoin
\pgfsetdash{}{0pt}
\definecolor{dialinecolor}{rgb}{0.000000, 0.470588, 0.666667}
\pgfsetfillcolor{dialinecolor}
\pgfpathmoveto{\pgfpoint{-21.590587\du}{6.285901\du}}
\pgfpathlineto{\pgfpoint{-21.590587\du}{6.380542\du}}
\pgfpathlineto{\pgfpoint{-21.197693\du}{6.380542\du}}
\pgfpathlineto{\pgfpoint{-21.197693\du}{6.285901\du}}
\pgfpathlineto{\pgfpoint{-21.590587\du}{6.285901\du}}
\pgfusepath{fill}
\pgfsetbuttcap
\pgfsetmiterjoin
\pgfsetdash{}{0pt}
\definecolor{dialinecolor}{rgb}{0.000000, 0.705882, 1.000000}
\pgfsetfillcolor{dialinecolor}
\pgfpathmoveto{\pgfpoint{-21.198290\du}{6.279333\du}}
\pgfpathlineto{\pgfpoint{-21.198290\du}{6.276049\du}}
\pgfpathlineto{\pgfpoint{-21.199186\du}{6.272466\du}}
\pgfpathlineto{\pgfpoint{-21.200380\du}{6.269481\du}}
\pgfpathlineto{\pgfpoint{-21.201873\du}{6.266196\du}}
\pgfpathlineto{\pgfpoint{-21.204261\du}{6.262912\du}}
\pgfpathlineto{\pgfpoint{-21.206948\du}{6.259628\du}}
\pgfpathlineto{\pgfpoint{-21.209934\du}{6.256643\du}}
\pgfpathlineto{\pgfpoint{-21.214113\du}{6.253359\du}}
\pgfpathlineto{\pgfpoint{-21.217696\du}{6.250672\du}}
\pgfpathlineto{\pgfpoint{-21.222473\du}{6.247686\du}}
\pgfpathlineto{\pgfpoint{-21.227548\du}{6.244701\du}}
\pgfpathlineto{\pgfpoint{-21.232922\du}{6.242014\du}}
\pgfpathlineto{\pgfpoint{-21.239192\du}{6.238730\du}}
\pgfpathlineto{\pgfpoint{-21.245163\du}{6.236938\du}}
\pgfpathlineto{\pgfpoint{-21.252029\du}{6.234251\du}}
\pgfpathlineto{\pgfpoint{-21.258896\du}{6.231863\du}}
\pgfpathlineto{\pgfpoint{-21.266360\du}{6.229773\du}}
\pgfpathlineto{\pgfpoint{-21.273824\du}{6.227683\du}}
\pgfpathlineto{\pgfpoint{-21.281885\du}{6.225593\du}}
\pgfpathlineto{\pgfpoint{-21.290543\du}{6.223504\du}}
\pgfpathlineto{\pgfpoint{-21.298902\du}{6.222011\du}}
\pgfpathlineto{\pgfpoint{-21.307560\du}{6.220518\du}}
\pgfpathlineto{\pgfpoint{-21.316815\du}{6.219324\du}}
\pgfpathlineto{\pgfpoint{-21.326070\du}{6.217831\du}}
\pgfpathlineto{\pgfpoint{-21.335325\du}{6.216935\du}}
\pgfpathlineto{\pgfpoint{-21.344879\du}{6.215741\du}}
\pgfpathlineto{\pgfpoint{-21.355030\du}{6.215144\du}}
\pgfpathlineto{\pgfpoint{-21.364584\du}{6.214547\du}}
\pgfpathlineto{\pgfpoint{-21.374137\du}{6.214248\du}}
\pgfpathlineto{\pgfpoint{-21.384288\du}{6.213950\du}}
\pgfpathlineto{\pgfpoint{-21.394140\du}{6.213950\du}}
\pgfpathlineto{\pgfpoint{-21.394140\du}{6.213950\du}}
\pgfpathlineto{\pgfpoint{-21.403992\du}{6.213950\du}}
\pgfpathlineto{\pgfpoint{-21.414143\du}{6.214248\du}}
\pgfpathlineto{\pgfpoint{-21.423995\du}{6.214547\du}}
\pgfpathlineto{\pgfpoint{-21.433848\du}{6.215144\du}}
\pgfpathlineto{\pgfpoint{-21.443401\du}{6.215741\du}}
\pgfpathlineto{\pgfpoint{-21.452955\du}{6.216935\du}}
\pgfpathlineto{\pgfpoint{-21.462509\du}{6.217831\du}}
\pgfpathlineto{\pgfpoint{-21.471465\du}{6.219324\du}}
\pgfpathlineto{\pgfpoint{-21.480720\du}{6.220518\du}}
\pgfpathlineto{\pgfpoint{-21.489378\du}{6.222011\du}}
\pgfpathlineto{\pgfpoint{-21.498036\du}{6.223504\du}}
\pgfpathlineto{\pgfpoint{-21.506396\du}{6.225593\du}}
\pgfpathlineto{\pgfpoint{-21.514457\du}{6.227683\du}}
\pgfpathlineto{\pgfpoint{-21.521921\du}{6.229773\du}}
\pgfpathlineto{\pgfpoint{-21.529384\du}{6.231863\du}}
\pgfpathlineto{\pgfpoint{-21.536550\du}{6.234251\du}}
\pgfpathlineto{\pgfpoint{-21.543118\du}{6.236938\du}}
\pgfpathlineto{\pgfpoint{-21.549387\du}{6.238730\du}}
\pgfpathlineto{\pgfpoint{-21.555358\du}{6.242014\du}}
\pgfpathlineto{\pgfpoint{-21.560732\du}{6.244701\du}}
\pgfpathlineto{\pgfpoint{-21.565808\du}{6.247686\du}}
\pgfpathlineto{\pgfpoint{-21.570584\du}{6.250672\du}}
\pgfpathlineto{\pgfpoint{-21.574466\du}{6.253359\du}}
\pgfpathlineto{\pgfpoint{-21.578347\du}{6.256643\du}}
\pgfpathlineto{\pgfpoint{-21.581332\du}{6.259628\du}}
\pgfpathlineto{\pgfpoint{-21.584019\du}{6.262912\du}}
\pgfpathlineto{\pgfpoint{-21.586408\du}{6.266196\du}}
\pgfpathlineto{\pgfpoint{-21.588199\du}{6.269481\du}}
\pgfpathlineto{\pgfpoint{-21.589692\du}{6.272466\du}}
\pgfpathlineto{\pgfpoint{-21.590587\du}{6.276049\du}}
\pgfpathlineto{\pgfpoint{-21.590587\du}{6.279333\du}}
\pgfpathlineto{\pgfpoint{-21.590587\du}{6.279333\du}}
\pgfpathlineto{\pgfpoint{-21.590587\du}{6.282617\du}}
\pgfpathlineto{\pgfpoint{-21.589692\du}{6.286199\du}}
\pgfpathlineto{\pgfpoint{-21.588199\du}{6.289185\du}}
\pgfpathlineto{\pgfpoint{-21.586408\du}{6.292469\du}}
\pgfpathlineto{\pgfpoint{-21.584019\du}{6.296052\du}}
\pgfpathlineto{\pgfpoint{-21.581332\du}{6.299037\du}}
\pgfpathlineto{\pgfpoint{-21.578347\du}{6.302023\du}}
\pgfpathlineto{\pgfpoint{-21.574466\du}{6.305307\du}}
\pgfpathlineto{\pgfpoint{-21.570584\du}{6.307695\du}}
\pgfpathlineto{\pgfpoint{-21.565808\du}{6.310979\du}}
\pgfpathlineto{\pgfpoint{-21.560732\du}{6.314263\du}}
\pgfpathlineto{\pgfpoint{-21.555358\du}{6.316950\du}}
\pgfpathlineto{\pgfpoint{-21.549387\du}{6.319637\du}}
\pgfpathlineto{\pgfpoint{-21.543118\du}{6.321727\du}}
\pgfpathlineto{\pgfpoint{-21.536550\du}{6.324713\du}}
\pgfpathlineto{\pgfpoint{-21.529384\du}{6.327101\du}}
\pgfpathlineto{\pgfpoint{-21.521921\du}{6.328892\du}}
\pgfpathlineto{\pgfpoint{-21.514457\du}{6.330982\du}}
\pgfpathlineto{\pgfpoint{-21.506396\du}{6.333072\du}}
\pgfpathlineto{\pgfpoint{-21.498036\du}{6.335162\du}}
\pgfpathlineto{\pgfpoint{-21.489378\du}{6.336655\du}}
\pgfpathlineto{\pgfpoint{-21.480720\du}{6.338147\du}}
\pgfpathlineto{\pgfpoint{-21.471465\du}{6.339342\du}}
\pgfpathlineto{\pgfpoint{-21.462509\du}{6.340834\du}}
\pgfpathlineto{\pgfpoint{-21.452955\du}{6.341730\du}}
\pgfpathlineto{\pgfpoint{-21.443401\du}{6.342924\du}}
\pgfpathlineto{\pgfpoint{-21.433848\du}{6.343223\du}}
\pgfpathlineto{\pgfpoint{-21.423995\du}{6.344119\du}}
\pgfpathlineto{\pgfpoint{-21.414143\du}{6.344119\du}}
\pgfpathlineto{\pgfpoint{-21.403992\du}{6.345014\du}}
\pgfpathlineto{\pgfpoint{-21.394140\du}{6.345014\du}}
\pgfpathlineto{\pgfpoint{-21.394140\du}{6.345014\du}}
\pgfpathlineto{\pgfpoint{-21.384288\du}{6.345014\du}}
\pgfpathlineto{\pgfpoint{-21.374137\du}{6.344119\du}}
\pgfpathlineto{\pgfpoint{-21.364584\du}{6.344119\du}}
\pgfpathlineto{\pgfpoint{-21.355030\du}{6.343223\du}}
\pgfpathlineto{\pgfpoint{-21.344879\du}{6.342924\du}}
\pgfpathlineto{\pgfpoint{-21.335325\du}{6.341730\du}}
\pgfpathlineto{\pgfpoint{-21.326070\du}{6.340834\du}}
\pgfpathlineto{\pgfpoint{-21.316815\du}{6.339342\du}}
\pgfpathlineto{\pgfpoint{-21.307560\du}{6.338147\du}}
\pgfpathlineto{\pgfpoint{-21.298902\du}{6.336655\du}}
\pgfpathlineto{\pgfpoint{-21.290543\du}{6.335162\du}}
\pgfpathlineto{\pgfpoint{-21.281885\du}{6.333072\du}}
\pgfpathlineto{\pgfpoint{-21.273824\du}{6.330982\du}}
\pgfpathlineto{\pgfpoint{-21.266360\du}{6.328892\du}}
\pgfpathlineto{\pgfpoint{-21.258896\du}{6.327101\du}}
\pgfpathlineto{\pgfpoint{-21.252029\du}{6.324713\du}}
\pgfpathlineto{\pgfpoint{-21.245163\du}{6.321727\du}}
\pgfpathlineto{\pgfpoint{-21.239192\du}{6.319637\du}}
\pgfpathlineto{\pgfpoint{-21.232922\du}{6.316950\du}}
\pgfpathlineto{\pgfpoint{-21.227548\du}{6.314263\du}}
\pgfpathlineto{\pgfpoint{-21.222473\du}{6.310979\du}}
\pgfpathlineto{\pgfpoint{-21.217696\du}{6.307695\du}}
\pgfpathlineto{\pgfpoint{-21.214113\du}{6.305307\du}}
\pgfpathlineto{\pgfpoint{-21.209934\du}{6.302023\du}}
\pgfpathlineto{\pgfpoint{-21.206948\du}{6.299037\du}}
\pgfpathlineto{\pgfpoint{-21.204261\du}{6.296052\du}}
\pgfpathlineto{\pgfpoint{-21.201873\du}{6.292469\du}}
\pgfpathlineto{\pgfpoint{-21.200380\du}{6.289185\du}}
\pgfpathlineto{\pgfpoint{-21.199186\du}{6.286199\du}}
\pgfpathlineto{\pgfpoint{-21.198290\du}{6.282617\du}}
\pgfpathlineto{\pgfpoint{-21.198290\du}{6.279333\du}}
\pgfusepath{fill}
\pgfsetbuttcap
\pgfsetmiterjoin
\pgfsetdash{}{0pt}
\definecolor{dialinecolor}{rgb}{0.666667, 0.901961, 1.000000}
\pgfsetstrokecolor{dialinecolor}
\pgfpathmoveto{\pgfpoint{-21.209934\du}{6.273660\du}}
\pgfpathlineto{\pgfpoint{-21.209934\du}{6.270376\du}}
\pgfpathlineto{\pgfpoint{-21.210829\du}{6.267988\du}}
\pgfpathlineto{\pgfpoint{-21.212024\du}{6.264704\du}}
\pgfpathlineto{\pgfpoint{-21.214113\du}{6.261121\du}}
\pgfpathlineto{\pgfpoint{-21.216203\du}{6.258733\du}}
\pgfpathlineto{\pgfpoint{-21.218293\du}{6.255449\du}}
\pgfpathlineto{\pgfpoint{-21.221577\du}{6.253060\du}}
\pgfpathlineto{\pgfpoint{-21.224861\du}{6.250373\du}}
\pgfpathlineto{\pgfpoint{-21.229041\du}{6.247686\du}}
\pgfpathlineto{\pgfpoint{-21.233818\du}{6.244999\du}}
\pgfpathlineto{\pgfpoint{-21.238296\du}{6.242312\du}}
\pgfpathlineto{\pgfpoint{-21.243969\du}{6.239625\du}}
\pgfpathlineto{\pgfpoint{-21.249940\du}{6.237237\du}}
\pgfpathlineto{\pgfpoint{-21.255612\du}{6.235147\du}}
\pgfpathlineto{\pgfpoint{-21.262180\du}{6.232759\du}}
\pgfpathlineto{\pgfpoint{-21.269047\du}{6.230370\du}}
\pgfpathlineto{\pgfpoint{-21.275914\du}{6.228579\du}}
\pgfpathlineto{\pgfpoint{-21.283676\du}{6.227086\du}}
\pgfpathlineto{\pgfpoint{-21.291140\du}{6.224698\du}}
\pgfpathlineto{\pgfpoint{-21.299499\du}{6.223205\du}}
\pgfpathlineto{\pgfpoint{-21.308157\du}{6.222011\du}}
\pgfpathlineto{\pgfpoint{-21.316517\du}{6.220219\du}}
\pgfpathlineto{\pgfpoint{-21.325175\du}{6.219324\du}}
\pgfpathlineto{\pgfpoint{-21.334131\du}{6.217831\du}}
\pgfpathlineto{\pgfpoint{-21.342789\du}{6.216935\du}}
\pgfpathlineto{\pgfpoint{-21.352343\du}{6.216338\du}}
\pgfpathlineto{\pgfpoint{-21.361598\du}{6.215443\du}}
\pgfpathlineto{\pgfpoint{-21.371152\du}{6.215144\du}}
\pgfpathlineto{\pgfpoint{-21.380705\du}{6.214845\du}}
\pgfpathlineto{\pgfpoint{-21.390856\du}{6.214248\du}}
\pgfpathlineto{\pgfpoint{-21.400708\du}{6.214248\du}}
\pgfpathlineto{\pgfpoint{-21.400708\du}{6.214248\du}}
\pgfpathlineto{\pgfpoint{-21.409665\du}{6.214248\du}}
\pgfpathlineto{\pgfpoint{-21.419517\du}{6.214845\du}}
\pgfpathlineto{\pgfpoint{-21.429071\du}{6.215144\du}}
\pgfpathlineto{\pgfpoint{-21.438326\du}{6.215443\du}}
\pgfpathlineto{\pgfpoint{-21.447880\du}{6.216338\du}}
\pgfpathlineto{\pgfpoint{-21.457433\du}{6.216935\du}}
\pgfpathlineto{\pgfpoint{-21.466390\du}{6.217831\du}}
\pgfpathlineto{\pgfpoint{-21.475346\du}{6.219324\du}}
\pgfpathlineto{\pgfpoint{-21.484004\du}{6.220219\du}}
\pgfpathlineto{\pgfpoint{-21.492364\du}{6.222011\du}}
\pgfpathlineto{\pgfpoint{-21.501022\du}{6.223205\du}}
\pgfpathlineto{\pgfpoint{-21.509083\du}{6.224698\du}}
\pgfpathlineto{\pgfpoint{-21.516845\du}{6.227086\du}}
\pgfpathlineto{\pgfpoint{-21.524309\du}{6.228579\du}}
\pgfpathlineto{\pgfpoint{-21.530877\du}{6.230370\du}}
\pgfpathlineto{\pgfpoint{-21.538341\du}{6.232759\du}}
\pgfpathlineto{\pgfpoint{-21.544610\du}{6.235147\du}}
\pgfpathlineto{\pgfpoint{-21.550582\du}{6.237237\du}}
\pgfpathlineto{\pgfpoint{-21.556553\du}{6.239625\du}}
\pgfpathlineto{\pgfpoint{-21.561628\du}{6.242312\du}}
\pgfpathlineto{\pgfpoint{-21.566703\du}{6.244999\du}}
\pgfpathlineto{\pgfpoint{-21.571480\du}{6.247686\du}}
\pgfpathlineto{\pgfpoint{-21.575660\du}{6.250373\du}}
\pgfpathlineto{\pgfpoint{-21.578944\du}{6.253060\du}}
\pgfpathlineto{\pgfpoint{-21.582228\du}{6.255449\du}}
\pgfpathlineto{\pgfpoint{-21.584616\du}{6.258733\du}}
\pgfpathlineto{\pgfpoint{-21.586408\du}{6.261121\du}}
\pgfpathlineto{\pgfpoint{-21.588796\du}{6.264704\du}}
\pgfpathlineto{\pgfpoint{-21.589692\du}{6.267988\du}}
\pgfpathlineto{\pgfpoint{-21.590587\du}{6.270376\du}}
\pgfpathlineto{\pgfpoint{-21.590587\du}{6.273660\du}}
\pgfpathlineto{\pgfpoint{-21.590587\du}{6.273660\du}}
\pgfpathlineto{\pgfpoint{-21.590587\du}{6.276944\du}}
\pgfpathlineto{\pgfpoint{-21.589692\du}{6.279631\du}}
\pgfpathlineto{\pgfpoint{-21.588796\du}{6.282617\du}}
\pgfpathlineto{\pgfpoint{-21.586408\du}{6.285901\du}}
\pgfpathlineto{\pgfpoint{-21.584616\du}{6.288588\du}}
\pgfpathlineto{\pgfpoint{-21.582228\du}{6.291872\du}}
\pgfpathlineto{\pgfpoint{-21.578944\du}{6.294559\du}}
\pgfpathlineto{\pgfpoint{-21.575660\du}{6.297246\du}}
\pgfpathlineto{\pgfpoint{-21.571480\du}{6.299933\du}}
\pgfpathlineto{\pgfpoint{-21.566703\du}{6.302620\du}}
\pgfpathlineto{\pgfpoint{-21.561628\du}{6.305307\du}}
\pgfpathlineto{\pgfpoint{-21.556553\du}{6.307397\du}}
\pgfpathlineto{\pgfpoint{-21.550582\du}{6.310084\du}}
\pgfpathlineto{\pgfpoint{-21.544610\du}{6.312472\du}}
\pgfpathlineto{\pgfpoint{-21.538341\du}{6.314860\du}}
\pgfpathlineto{\pgfpoint{-21.530877\du}{6.316950\du}}
\pgfpathlineto{\pgfpoint{-21.524309\du}{6.318742\du}}
\pgfpathlineto{\pgfpoint{-21.516845\du}{6.320533\du}}
\pgfpathlineto{\pgfpoint{-21.509083\du}{6.322324\du}}
\pgfpathlineto{\pgfpoint{-21.501022\du}{6.324414\du}}
\pgfpathlineto{\pgfpoint{-21.492364\du}{6.325608\du}}
\pgfpathlineto{\pgfpoint{-21.484004\du}{6.327400\du}}
\pgfpathlineto{\pgfpoint{-21.475346\du}{6.328295\du}}
\pgfpathlineto{\pgfpoint{-21.466390\du}{6.329788\du}}
\pgfpathlineto{\pgfpoint{-21.457433\du}{6.330087\du}}
\pgfpathlineto{\pgfpoint{-21.447880\du}{6.331281\du}}
\pgfpathlineto{\pgfpoint{-21.438326\du}{6.331579\du}}
\pgfpathlineto{\pgfpoint{-21.429071\du}{6.332475\du}}
\pgfpathlineto{\pgfpoint{-21.419517\du}{6.332774\du}}
\pgfpathlineto{\pgfpoint{-21.409665\du}{6.333072\du}}
\pgfpathlineto{\pgfpoint{-21.400708\du}{6.333072\du}}
\pgfpathlineto{\pgfpoint{-21.400708\du}{6.333072\du}}
\pgfpathlineto{\pgfpoint{-21.390856\du}{6.333072\du}}
\pgfpathlineto{\pgfpoint{-21.380705\du}{6.332774\du}}
\pgfpathlineto{\pgfpoint{-21.371152\du}{6.332475\du}}
\pgfpathlineto{\pgfpoint{-21.361598\du}{6.331579\du}}
\pgfpathlineto{\pgfpoint{-21.352343\du}{6.331281\du}}
\pgfpathlineto{\pgfpoint{-21.342789\du}{6.330087\du}}
\pgfpathlineto{\pgfpoint{-21.334131\du}{6.329788\du}}
\pgfpathlineto{\pgfpoint{-21.325175\du}{6.328295\du}}
\pgfpathlineto{\pgfpoint{-21.316517\du}{6.327400\du}}
\pgfpathlineto{\pgfpoint{-21.308157\du}{6.325608\du}}
\pgfpathlineto{\pgfpoint{-21.299499\du}{6.324414\du}}
\pgfpathlineto{\pgfpoint{-21.291140\du}{6.322324\du}}
\pgfpathlineto{\pgfpoint{-21.283676\du}{6.320533\du}}
\pgfpathlineto{\pgfpoint{-21.275914\du}{6.318742\du}}
\pgfpathlineto{\pgfpoint{-21.269047\du}{6.316950\du}}
\pgfpathlineto{\pgfpoint{-21.262180\du}{6.314860\du}}
\pgfpathlineto{\pgfpoint{-21.255612\du}{6.312472\du}}
\pgfpathlineto{\pgfpoint{-21.249940\du}{6.310084\du}}
\pgfpathlineto{\pgfpoint{-21.243969\du}{6.307397\du}}
\pgfpathlineto{\pgfpoint{-21.238296\du}{6.305307\du}}
\pgfpathlineto{\pgfpoint{-21.233818\du}{6.302620\du}}
\pgfpathlineto{\pgfpoint{-21.229041\du}{6.299933\du}}
\pgfpathlineto{\pgfpoint{-21.224861\du}{6.297246\du}}
\pgfpathlineto{\pgfpoint{-21.221577\du}{6.294559\du}}
\pgfpathlineto{\pgfpoint{-21.218293\du}{6.291872\du}}
\pgfpathlineto{\pgfpoint{-21.216203\du}{6.288588\du}}
\pgfpathlineto{\pgfpoint{-21.214113\du}{6.285901\du}}
\pgfpathlineto{\pgfpoint{-21.212024\du}{6.282617\du}}
\pgfpathlineto{\pgfpoint{-21.210829\du}{6.279631\du}}
\pgfpathlineto{\pgfpoint{-21.209934\du}{6.276944\du}}
\pgfpathlineto{\pgfpoint{-21.209934\du}{6.273660\du}}
\pgfusepath{stroke}
\pgfsetbuttcap
\pgfsetmiterjoin
\pgfsetdash{}{0pt}
\definecolor{dialinecolor}{rgb}{0.000000, 0.000000, 0.000000}
\pgfsetfillcolor{dialinecolor}
\pgfpathmoveto{\pgfpoint{-21.388468\du}{6.262017\du}}
\pgfpathlineto{\pgfpoint{-21.364584\du}{6.273660\du}}
\pgfpathlineto{\pgfpoint{-21.293230\du}{6.250075\du}}
\pgfpathlineto{\pgfpoint{-21.257403\du}{6.262017\du}}
\pgfpathlineto{\pgfpoint{-21.281288\du}{6.238133\du}}
\pgfpathlineto{\pgfpoint{-21.364584\du}{6.238133\du}}
\pgfpathlineto{\pgfpoint{-21.328757\du}{6.238133\du}}
\pgfpathlineto{\pgfpoint{-21.388468\du}{6.262017\du}}
\pgfusepath{fill}
\pgfsetbuttcap
\pgfsetmiterjoin
\pgfsetdash{}{0pt}
\definecolor{dialinecolor}{rgb}{0.000000, 0.000000, 0.000000}
\pgfsetfillcolor{dialinecolor}
\pgfpathmoveto{\pgfpoint{-21.400708\du}{6.297544\du}}
\pgfpathlineto{\pgfpoint{-21.435639\du}{6.285901\du}}
\pgfpathlineto{\pgfpoint{-21.495648\du}{6.309785\du}}
\pgfpathlineto{\pgfpoint{-21.530877\du}{6.297544\du}}
\pgfpathlineto{\pgfpoint{-21.519532\du}{6.333072\du}}
\pgfpathlineto{\pgfpoint{-21.435639\du}{6.333072\du}}
\pgfpathlineto{\pgfpoint{-21.471764\du}{6.321727\du}}
\pgfpathlineto{\pgfpoint{-21.400708\du}{6.297544\du}}
\pgfusepath{fill}
\pgfsetbuttcap
\pgfsetmiterjoin
\pgfsetdash{}{0pt}
\definecolor{dialinecolor}{rgb}{0.000000, 0.000000, 0.000000}
\pgfsetfillcolor{dialinecolor}
\pgfpathmoveto{\pgfpoint{-21.519532\du}{6.238133\du}}
\pgfpathlineto{\pgfpoint{-21.495648\du}{6.225593\du}}
\pgfpathlineto{\pgfpoint{-21.423995\du}{6.262017\du}}
\pgfpathlineto{\pgfpoint{-21.400708\du}{6.250075\du}}
\pgfpathlineto{\pgfpoint{-21.412352\du}{6.273660\du}}
\pgfpathlineto{\pgfpoint{-21.495648\du}{6.273660\du}}
\pgfpathlineto{\pgfpoint{-21.459822\du}{6.262017\du}}
\pgfpathlineto{\pgfpoint{-21.519532\du}{6.238133\du}}
\pgfusepath{fill}
\pgfsetbuttcap
\pgfsetmiterjoin
\pgfsetdash{}{0pt}
\definecolor{dialinecolor}{rgb}{0.000000, 0.000000, 0.000000}
\pgfsetfillcolor{dialinecolor}
\pgfpathmoveto{\pgfpoint{-21.269047\du}{6.321727\du}}
\pgfpathlineto{\pgfpoint{-21.293230\du}{6.333072\du}}
\pgfpathlineto{\pgfpoint{-21.364584\du}{6.297544\du}}
\pgfpathlineto{\pgfpoint{-21.400708\du}{6.309785\du}}
\pgfpathlineto{\pgfpoint{-21.376227\du}{6.285901\du}}
\pgfpathlineto{\pgfpoint{-21.293230\du}{6.285901\du}}
\pgfpathlineto{\pgfpoint{-21.328757\du}{6.297544\du}}
\pgfpathlineto{\pgfpoint{-21.269047\du}{6.321727\du}}
\pgfusepath{fill}
\pgfsetbuttcap
\pgfsetmiterjoin
\pgfsetdash{}{0pt}
\definecolor{dialinecolor}{rgb}{1.000000, 1.000000, 1.000000}
\pgfsetfillcolor{dialinecolor}
\pgfpathmoveto{\pgfpoint{-21.388468\du}{6.273660\du}}
\pgfpathlineto{\pgfpoint{-21.364584\du}{6.273660\du}}
\pgfpathlineto{\pgfpoint{-21.293230\du}{6.250075\du}}
\pgfpathlineto{\pgfpoint{-21.257403\du}{6.262017\du}}
\pgfpathlineto{\pgfpoint{-21.281288\du}{6.238133\du}}
\pgfpathlineto{\pgfpoint{-21.352343\du}{6.238133\du}}
\pgfpathlineto{\pgfpoint{-21.328757\du}{6.238133\du}}
\pgfpathlineto{\pgfpoint{-21.388468\du}{6.273660\du}}
\pgfusepath{fill}
\pgfsetbuttcap
\pgfsetmiterjoin
\pgfsetdash{}{0pt}
\definecolor{dialinecolor}{rgb}{1.000000, 1.000000, 1.000000}
\pgfsetfillcolor{dialinecolor}
\pgfpathmoveto{\pgfpoint{-21.400708\du}{6.297544\du}}
\pgfpathlineto{\pgfpoint{-21.423995\du}{6.285901\du}}
\pgfpathlineto{\pgfpoint{-21.495648\du}{6.309785\du}}
\pgfpathlineto{\pgfpoint{-21.530877\du}{6.309785\du}}
\pgfpathlineto{\pgfpoint{-21.507291\du}{6.333072\du}}
\pgfpathlineto{\pgfpoint{-21.423995\du}{6.333072\du}}
\pgfpathlineto{\pgfpoint{-21.459822\du}{6.321727\du}}
\pgfpathlineto{\pgfpoint{-21.400708\du}{6.297544\du}}
\pgfusepath{fill}
\pgfsetbuttcap
\pgfsetmiterjoin
\pgfsetdash{}{0pt}
\definecolor{dialinecolor}{rgb}{1.000000, 1.000000, 1.000000}
\pgfsetfillcolor{dialinecolor}
\pgfpathmoveto{\pgfpoint{-21.519532\du}{6.238133\du}}
\pgfpathlineto{\pgfpoint{-21.495648\du}{6.238133\du}}
\pgfpathlineto{\pgfpoint{-21.423995\du}{6.262017\du}}
\pgfpathlineto{\pgfpoint{-21.388468\du}{6.250075\du}}
\pgfpathlineto{\pgfpoint{-21.412352\du}{6.273660\du}}
\pgfpathlineto{\pgfpoint{-21.495648\du}{6.273660\du}}
\pgfpathlineto{\pgfpoint{-21.459822\du}{6.273660\du}}
\pgfpathlineto{\pgfpoint{-21.519532\du}{6.238133\du}}
\pgfusepath{fill}
\pgfsetbuttcap
\pgfsetmiterjoin
\pgfsetdash{}{0pt}
\definecolor{dialinecolor}{rgb}{1.000000, 1.000000, 1.000000}
\pgfsetfillcolor{dialinecolor}
\pgfpathmoveto{\pgfpoint{-21.269047\du}{6.321727\du}}
\pgfpathlineto{\pgfpoint{-21.293230\du}{6.333072\du}}
\pgfpathlineto{\pgfpoint{-21.364584\du}{6.309785\du}}
\pgfpathlineto{\pgfpoint{-21.388468\du}{6.309785\du}}
\pgfpathlineto{\pgfpoint{-21.376227\du}{6.285901\du}}
\pgfpathlineto{\pgfpoint{-21.293230\du}{6.285901\du}}
\pgfpathlineto{\pgfpoint{-21.328757\du}{6.297544\du}}
\pgfpathlineto{\pgfpoint{-21.269047\du}{6.321727\du}}
\pgfusepath{fill}
\pgfsetbuttcap
\pgfsetmiterjoin
\pgfsetdash{}{0pt}
\definecolor{dialinecolor}{rgb}{0.666667, 0.901961, 1.000000}
\pgfsetstrokecolor{dialinecolor}
\pgfpathmoveto{\pgfpoint{-21.590587\du}{6.273660\du}}
\pgfpathlineto{\pgfpoint{-21.590587\du}{6.368600\du}}
\pgfusepath{stroke}
\pgfsetbuttcap
\pgfsetmiterjoin
\pgfsetdash{}{0pt}
\definecolor{dialinecolor}{rgb}{0.666667, 0.901961, 1.000000}
\pgfsetstrokecolor{dialinecolor}
\pgfpathmoveto{\pgfpoint{-21.209934\du}{6.273660\du}}
\pgfpathlineto{\pgfpoint{-21.209934\du}{6.368600\du}}
\pgfusepath{stroke}
\pgfsetbuttcap
\pgfsetmiterjoin
\pgfsetdash{}{0pt}
\definecolor{dialinecolor}{rgb}{0.000000, 0.470588, 0.666667}
\pgfsetfillcolor{dialinecolor}
\pgfpathmoveto{\pgfpoint{-21.555358\du}{6.779706\du}}
\pgfpathlineto{\pgfpoint{-21.555657\du}{6.776123\du}}
\pgfpathlineto{\pgfpoint{-21.556553\du}{6.772839\du}}
\pgfpathlineto{\pgfpoint{-21.557747\du}{6.769555\du}}
\pgfpathlineto{\pgfpoint{-21.559538\du}{6.766570\du}}
\pgfpathlineto{\pgfpoint{-21.561628\du}{6.763286\du}}
\pgfpathlineto{\pgfpoint{-21.564912\du}{6.760001\du}}
\pgfpathlineto{\pgfpoint{-21.567599\du}{6.757016\du}}
\pgfpathlineto{\pgfpoint{-21.571480\du}{6.753732\du}}
\pgfpathlineto{\pgfpoint{-21.575660\du}{6.751045\du}}
\pgfpathlineto{\pgfpoint{-21.580138\du}{6.747761\du}}
\pgfpathlineto{\pgfpoint{-21.585512\du}{6.745074\du}}
\pgfpathlineto{\pgfpoint{-21.590587\du}{6.742387\du}}
\pgfpathlineto{\pgfpoint{-21.596559\du}{6.739700\du}}
\pgfpathlineto{\pgfpoint{-21.602530\du}{6.737013\du}}
\pgfpathlineto{\pgfpoint{-21.609695\du}{6.734625\du}}
\pgfpathlineto{\pgfpoint{-21.616562\du}{6.732236\du}}
\pgfpathlineto{\pgfpoint{-21.624025\du}{6.730445\du}}
\pgfpathlineto{\pgfpoint{-21.631489\du}{6.728056\du}}
\pgfpathlineto{\pgfpoint{-21.639550\du}{6.726265\du}}
\pgfpathlineto{\pgfpoint{-21.647909\du}{6.723877\du}}
\pgfpathlineto{\pgfpoint{-21.656269\du}{6.722384\du}}
\pgfpathlineto{\pgfpoint{-21.665225\du}{6.721190\du}}
\pgfpathlineto{\pgfpoint{-21.674182\du}{6.719697\du}}
\pgfpathlineto{\pgfpoint{-21.683437\du}{6.718503\du}}
\pgfpathlineto{\pgfpoint{-21.693289\du}{6.717309\du}}
\pgfpathlineto{\pgfpoint{-21.702544\du}{6.716413\du}}
\pgfpathlineto{\pgfpoint{-21.712397\du}{6.715816\du}}
\pgfpathlineto{\pgfpoint{-21.721950\du}{6.715219\du}}
\pgfpathlineto{\pgfpoint{-21.731803\du}{6.714622\du}}
\pgfpathlineto{\pgfpoint{-21.742252\du}{6.714323\du}}
\pgfpathlineto{\pgfpoint{-21.751507\du}{6.714323\du}}
\pgfpathlineto{\pgfpoint{-21.751507\du}{6.714323\du}}
\pgfpathlineto{\pgfpoint{-21.761359\du}{6.714323\du}}
\pgfpathlineto{\pgfpoint{-21.771510\du}{6.714622\du}}
\pgfpathlineto{\pgfpoint{-21.781362\du}{6.715219\du}}
\pgfpathlineto{\pgfpoint{-21.791214\du}{6.715816\du}}
\pgfpathlineto{\pgfpoint{-21.801067\du}{6.716413\du}}
\pgfpathlineto{\pgfpoint{-21.810322\du}{6.717309\du}}
\pgfpathlineto{\pgfpoint{-21.820174\du}{6.718503\du}}
\pgfpathlineto{\pgfpoint{-21.829429\du}{6.719697\du}}
\pgfpathlineto{\pgfpoint{-21.838386\du}{6.721190\du}}
\pgfpathlineto{\pgfpoint{-21.847044\du}{6.722384\du}}
\pgfpathlineto{\pgfpoint{-21.855403\du}{6.723877\du}}
\pgfpathlineto{\pgfpoint{-21.863763\du}{6.726265\du}}
\pgfpathlineto{\pgfpoint{-21.871823\du}{6.728056\du}}
\pgfpathlineto{\pgfpoint{-21.879287\du}{6.730445\du}}
\pgfpathlineto{\pgfpoint{-21.886751\du}{6.732236\du}}
\pgfpathlineto{\pgfpoint{-21.893916\du}{6.734625\du}}
\pgfpathlineto{\pgfpoint{-21.900783\du}{6.737013\du}}
\pgfpathlineto{\pgfpoint{-21.906754\du}{6.739700\du}}
\pgfpathlineto{\pgfpoint{-21.913024\du}{6.742387\du}}
\pgfpathlineto{\pgfpoint{-21.918099\du}{6.745074\du}}
\pgfpathlineto{\pgfpoint{-21.923174\du}{6.747761\du}}
\pgfpathlineto{\pgfpoint{-21.927951\du}{6.751045\du}}
\pgfpathlineto{\pgfpoint{-21.932131\du}{6.753732\du}}
\pgfpathlineto{\pgfpoint{-21.935714\du}{6.757016\du}}
\pgfpathlineto{\pgfpoint{-21.938699\du}{6.760001\du}}
\pgfpathlineto{\pgfpoint{-21.941386\du}{6.763286\du}}
\pgfpathlineto{\pgfpoint{-21.943775\du}{6.766570\du}}
\pgfpathlineto{\pgfpoint{-21.945566\du}{6.769555\du}}
\pgfpathlineto{\pgfpoint{-21.947059\du}{6.772839\du}}
\pgfpathlineto{\pgfpoint{-21.947954\du}{6.776123\du}}
\pgfpathlineto{\pgfpoint{-21.947954\du}{6.779706\du}}
\pgfpathlineto{\pgfpoint{-21.947954\du}{6.779706\du}}
\pgfpathlineto{\pgfpoint{-21.947954\du}{6.782990\du}}
\pgfpathlineto{\pgfpoint{-21.947059\du}{6.786274\du}}
\pgfpathlineto{\pgfpoint{-21.945566\du}{6.789558\du}}
\pgfpathlineto{\pgfpoint{-21.943775\du}{6.792842\du}}
\pgfpathlineto{\pgfpoint{-21.941386\du}{6.795828\du}}
\pgfpathlineto{\pgfpoint{-21.938699\du}{6.799112\du}}
\pgfpathlineto{\pgfpoint{-21.935714\du}{6.802097\du}}
\pgfpathlineto{\pgfpoint{-21.932131\du}{6.805381\du}}
\pgfpathlineto{\pgfpoint{-21.927951\du}{6.808367\du}}
\pgfpathlineto{\pgfpoint{-21.923174\du}{6.811352\du}}
\pgfpathlineto{\pgfpoint{-21.918099\du}{6.814039\du}}
\pgfpathlineto{\pgfpoint{-21.913024\du}{6.816726\du}}
\pgfpathlineto{\pgfpoint{-21.906754\du}{6.819413\du}}
\pgfpathlineto{\pgfpoint{-21.900783\du}{6.821802\du}}
\pgfpathlineto{\pgfpoint{-21.893916\du}{6.824489\du}}
\pgfpathlineto{\pgfpoint{-21.886751\du}{6.826877\du}}
\pgfpathlineto{\pgfpoint{-21.879287\du}{6.828967\du}}
\pgfpathlineto{\pgfpoint{-21.871823\du}{6.831355\du}}
\pgfpathlineto{\pgfpoint{-21.863763\du}{6.832848\du}}
\pgfpathlineto{\pgfpoint{-21.855403\du}{6.834938\du}}
\pgfpathlineto{\pgfpoint{-21.847044\du}{6.836729\du}}
\pgfpathlineto{\pgfpoint{-21.838386\du}{6.837924\du}}
\pgfpathlineto{\pgfpoint{-21.829429\du}{6.839416\du}}
\pgfpathlineto{\pgfpoint{-21.820174\du}{6.840611\du}}
\pgfpathlineto{\pgfpoint{-21.810322\du}{6.841805\du}}
\pgfpathlineto{\pgfpoint{-21.801067\du}{6.842700\du}}
\pgfpathlineto{\pgfpoint{-21.791214\du}{6.843298\du}}
\pgfpathlineto{\pgfpoint{-21.781362\du}{6.844193\du}}
\pgfpathlineto{\pgfpoint{-21.771510\du}{6.844193\du}}
\pgfpathlineto{\pgfpoint{-21.761359\du}{6.844790\du}}
\pgfpathlineto{\pgfpoint{-21.751507\du}{6.844790\du}}
\pgfpathlineto{\pgfpoint{-21.751507\du}{6.844790\du}}
\pgfpathlineto{\pgfpoint{-21.742252\du}{6.844790\du}}
\pgfpathlineto{\pgfpoint{-21.731803\du}{6.844193\du}}
\pgfpathlineto{\pgfpoint{-21.721950\du}{6.844193\du}}
\pgfpathlineto{\pgfpoint{-21.712397\du}{6.843298\du}}
\pgfpathlineto{\pgfpoint{-21.702544\du}{6.842700\du}}
\pgfpathlineto{\pgfpoint{-21.693289\du}{6.841805\du}}
\pgfpathlineto{\pgfpoint{-21.683437\du}{6.840611\du}}
\pgfpathlineto{\pgfpoint{-21.674182\du}{6.839416\du}}
\pgfpathlineto{\pgfpoint{-21.665225\du}{6.837924\du}}
\pgfpathlineto{\pgfpoint{-21.656269\du}{6.836729\du}}
\pgfpathlineto{\pgfpoint{-21.647909\du}{6.834938\du}}
\pgfpathlineto{\pgfpoint{-21.639550\du}{6.832848\du}}
\pgfpathlineto{\pgfpoint{-21.631489\du}{6.831355\du}}
\pgfpathlineto{\pgfpoint{-21.624025\du}{6.828967\du}}
\pgfpathlineto{\pgfpoint{-21.616562\du}{6.826877\du}}
\pgfpathlineto{\pgfpoint{-21.609695\du}{6.824489\du}}
\pgfpathlineto{\pgfpoint{-21.602530\du}{6.821802\du}}
\pgfpathlineto{\pgfpoint{-21.596559\du}{6.819413\du}}
\pgfpathlineto{\pgfpoint{-21.590587\du}{6.816726\du}}
\pgfpathlineto{\pgfpoint{-21.585512\du}{6.814039\du}}
\pgfpathlineto{\pgfpoint{-21.580138\du}{6.811352\du}}
\pgfpathlineto{\pgfpoint{-21.575660\du}{6.808367\du}}
\pgfpathlineto{\pgfpoint{-21.571480\du}{6.805381\du}}
\pgfpathlineto{\pgfpoint{-21.567599\du}{6.802097\du}}
\pgfpathlineto{\pgfpoint{-21.564912\du}{6.799112\du}}
\pgfpathlineto{\pgfpoint{-21.561628\du}{6.795828\du}}
\pgfpathlineto{\pgfpoint{-21.559538\du}{6.792842\du}}
\pgfpathlineto{\pgfpoint{-21.557747\du}{6.789558\du}}
\pgfpathlineto{\pgfpoint{-21.556553\du}{6.786274\du}}
\pgfpathlineto{\pgfpoint{-21.555657\du}{6.782990\du}}
\pgfpathlineto{\pgfpoint{-21.555358\du}{6.779706\du}}
\pgfusepath{fill}
\pgfsetbuttcap
\pgfsetmiterjoin
\pgfsetdash{}{0pt}
\definecolor{dialinecolor}{rgb}{0.666667, 0.901961, 1.000000}
\pgfsetstrokecolor{dialinecolor}
\pgfpathmoveto{\pgfpoint{-21.567300\du}{6.773436\du}}
\pgfpathlineto{\pgfpoint{-21.567599\du}{6.770451\du}}
\pgfpathlineto{\pgfpoint{-21.568196\du}{6.767764\du}}
\pgfpathlineto{\pgfpoint{-21.569390\du}{6.764778\du}}
\pgfpathlineto{\pgfpoint{-21.571480\du}{6.761196\du}}
\pgfpathlineto{\pgfpoint{-21.573570\du}{6.758807\du}}
\pgfpathlineto{\pgfpoint{-21.575660\du}{6.755523\du}}
\pgfpathlineto{\pgfpoint{-21.578944\du}{6.752836\du}}
\pgfpathlineto{\pgfpoint{-21.582228\du}{6.750149\du}}
\pgfpathlineto{\pgfpoint{-21.586408\du}{6.747164\du}}
\pgfpathlineto{\pgfpoint{-21.591185\du}{6.744775\du}}
\pgfpathlineto{\pgfpoint{-21.596260\du}{6.742088\du}}
\pgfpathlineto{\pgfpoint{-21.601335\du}{6.739700\du}}
\pgfpathlineto{\pgfpoint{-21.607008\du}{6.737312\du}}
\pgfpathlineto{\pgfpoint{-21.612979\du}{6.734923\du}}
\pgfpathlineto{\pgfpoint{-21.619547\du}{6.732236\du}}
\pgfpathlineto{\pgfpoint{-21.626414\du}{6.730743\du}}
\pgfpathlineto{\pgfpoint{-21.632982\du}{6.728654\du}}
\pgfpathlineto{\pgfpoint{-21.641043\du}{6.726862\du}}
\pgfpathlineto{\pgfpoint{-21.648507\du}{6.725369\du}}
\pgfpathlineto{\pgfpoint{-21.656269\du}{6.723280\du}}
\pgfpathlineto{\pgfpoint{-21.665225\du}{6.721787\du}}
\pgfpathlineto{\pgfpoint{-21.673585\du}{6.720294\du}}
\pgfpathlineto{\pgfpoint{-21.681944\du}{6.719100\du}}
\pgfpathlineto{\pgfpoint{-21.691200\du}{6.717906\du}}
\pgfpathlineto{\pgfpoint{-21.699858\du}{6.717010\du}}
\pgfpathlineto{\pgfpoint{-21.709411\du}{6.716114\du}}
\pgfpathlineto{\pgfpoint{-21.718965\du}{6.715816\du}}
\pgfpathlineto{\pgfpoint{-21.728220\du}{6.715219\du}}
\pgfpathlineto{\pgfpoint{-21.738072\du}{6.714622\du}}
\pgfpathlineto{\pgfpoint{-21.747924\du}{6.714323\du}}
\pgfpathlineto{\pgfpoint{-21.757478\du}{6.714323\du}}
\pgfpathlineto{\pgfpoint{-21.757478\du}{6.714323\du}}
\pgfpathlineto{\pgfpoint{-21.766733\du}{6.714323\du}}
\pgfpathlineto{\pgfpoint{-21.776585\du}{6.714622\du}}
\pgfpathlineto{\pgfpoint{-21.786139\du}{6.715219\du}}
\pgfpathlineto{\pgfpoint{-21.795693\du}{6.715816\du}}
\pgfpathlineto{\pgfpoint{-21.805246\du}{6.716114\du}}
\pgfpathlineto{\pgfpoint{-21.814501\du}{6.717010\du}}
\pgfpathlineto{\pgfpoint{-21.823458\du}{6.717906\du}}
\pgfpathlineto{\pgfpoint{-21.832713\du}{6.719100\du}}
\pgfpathlineto{\pgfpoint{-21.841073\du}{6.720294\du}}
\pgfpathlineto{\pgfpoint{-21.850029\du}{6.721787\du}}
\pgfpathlineto{\pgfpoint{-21.858389\du}{6.723280\du}}
\pgfpathlineto{\pgfpoint{-21.866151\du}{6.725369\du}}
\pgfpathlineto{\pgfpoint{-21.874212\du}{6.726862\du}}
\pgfpathlineto{\pgfpoint{-21.881676\du}{6.728654\du}}
\pgfpathlineto{\pgfpoint{-21.888244\du}{6.730743\du}}
\pgfpathlineto{\pgfpoint{-21.895708\du}{6.732236\du}}
\pgfpathlineto{\pgfpoint{-21.901679\du}{6.734923\du}}
\pgfpathlineto{\pgfpoint{-21.907650\du}{6.737312\du}}
\pgfpathlineto{\pgfpoint{-21.913919\du}{6.739700\du}}
\pgfpathlineto{\pgfpoint{-21.918995\du}{6.742088\du}}
\pgfpathlineto{\pgfpoint{-21.924070\du}{6.744775\du}}
\pgfpathlineto{\pgfpoint{-21.928250\du}{6.747164\du}}
\pgfpathlineto{\pgfpoint{-21.932430\du}{6.750149\du}}
\pgfpathlineto{\pgfpoint{-21.936311\du}{6.752836\du}}
\pgfpathlineto{\pgfpoint{-21.939595\du}{6.755523\du}}
\pgfpathlineto{\pgfpoint{-21.941386\du}{6.758807\du}}
\pgfpathlineto{\pgfpoint{-21.943775\du}{6.761196\du}}
\pgfpathlineto{\pgfpoint{-21.945566\du}{6.764778\du}}
\pgfpathlineto{\pgfpoint{-21.947059\du}{6.767764\du}}
\pgfpathlineto{\pgfpoint{-21.947357\du}{6.770451\du}}
\pgfpathlineto{\pgfpoint{-21.947954\du}{6.773436\du}}
\pgfpathlineto{\pgfpoint{-21.947954\du}{6.773436\du}}
\pgfpathlineto{\pgfpoint{-21.947357\du}{6.777019\du}}
\pgfpathlineto{\pgfpoint{-21.947059\du}{6.779706\du}}
\pgfpathlineto{\pgfpoint{-21.945566\du}{6.782691\du}}
\pgfpathlineto{\pgfpoint{-21.943775\du}{6.785677\du}}
\pgfpathlineto{\pgfpoint{-21.941386\du}{6.788364\du}}
\pgfpathlineto{\pgfpoint{-21.939595\du}{6.791648\du}}
\pgfpathlineto{\pgfpoint{-21.936311\du}{6.794335\du}}
\pgfpathlineto{\pgfpoint{-21.932430\du}{6.797320\du}}
\pgfpathlineto{\pgfpoint{-21.928250\du}{6.800007\du}}
\pgfpathlineto{\pgfpoint{-21.924070\du}{6.802694\du}}
\pgfpathlineto{\pgfpoint{-21.918995\du}{6.805381\du}}
\pgfpathlineto{\pgfpoint{-21.913919\du}{6.806874\du}}
\pgfpathlineto{\pgfpoint{-21.907650\du}{6.810158\du}}
\pgfpathlineto{\pgfpoint{-21.901679\du}{6.812248\du}}
\pgfpathlineto{\pgfpoint{-21.895708\du}{6.814338\du}}
\pgfpathlineto{\pgfpoint{-21.888244\du}{6.816726\du}}
\pgfpathlineto{\pgfpoint{-21.881676\du}{6.818816\du}}
\pgfpathlineto{\pgfpoint{-21.874212\du}{6.820608\du}}
\pgfpathlineto{\pgfpoint{-21.866151\du}{6.821802\du}}
\pgfpathlineto{\pgfpoint{-21.858389\du}{6.824190\du}}
\pgfpathlineto{\pgfpoint{-21.850029\du}{6.825683\du}}
\pgfpathlineto{\pgfpoint{-21.841073\du}{6.827176\du}}
\pgfpathlineto{\pgfpoint{-21.832713\du}{6.828370\du}}
\pgfpathlineto{\pgfpoint{-21.823458\du}{6.829266\du}}
\pgfpathlineto{\pgfpoint{-21.814501\du}{6.830161\du}}
\pgfpathlineto{\pgfpoint{-21.805246\du}{6.831355\du}}
\pgfpathlineto{\pgfpoint{-21.795693\du}{6.831654\du}}
\pgfpathlineto{\pgfpoint{-21.786139\du}{6.832251\du}}
\pgfpathlineto{\pgfpoint{-21.776585\du}{6.832550\du}}
\pgfpathlineto{\pgfpoint{-21.766733\du}{6.832848\du}}
\pgfpathlineto{\pgfpoint{-21.757478\du}{6.832848\du}}
\pgfpathlineto{\pgfpoint{-21.757478\du}{6.832848\du}}
\pgfpathlineto{\pgfpoint{-21.747924\du}{6.832848\du}}
\pgfpathlineto{\pgfpoint{-21.738072\du}{6.832550\du}}
\pgfpathlineto{\pgfpoint{-21.728220\du}{6.832251\du}}
\pgfpathlineto{\pgfpoint{-21.718965\du}{6.831654\du}}
\pgfpathlineto{\pgfpoint{-21.709411\du}{6.831355\du}}
\pgfpathlineto{\pgfpoint{-21.699858\du}{6.830161\du}}
\pgfpathlineto{\pgfpoint{-21.691200\du}{6.829266\du}}
\pgfpathlineto{\pgfpoint{-21.681944\du}{6.828370\du}}
\pgfpathlineto{\pgfpoint{-21.673585\du}{6.827176\du}}
\pgfpathlineto{\pgfpoint{-21.665225\du}{6.825683\du}}
\pgfpathlineto{\pgfpoint{-21.656269\du}{6.824190\du}}
\pgfpathlineto{\pgfpoint{-21.648507\du}{6.821802\du}}
\pgfpathlineto{\pgfpoint{-21.641043\du}{6.820608\du}}
\pgfpathlineto{\pgfpoint{-21.632982\du}{6.818816\du}}
\pgfpathlineto{\pgfpoint{-21.626414\du}{6.816726\du}}
\pgfpathlineto{\pgfpoint{-21.619547\du}{6.814338\du}}
\pgfpathlineto{\pgfpoint{-21.612979\du}{6.812248\du}}
\pgfpathlineto{\pgfpoint{-21.607008\du}{6.810158\du}}
\pgfpathlineto{\pgfpoint{-21.601335\du}{6.806874\du}}
\pgfpathlineto{\pgfpoint{-21.596260\du}{6.805381\du}}
\pgfpathlineto{\pgfpoint{-21.591185\du}{6.802694\du}}
\pgfpathlineto{\pgfpoint{-21.586408\du}{6.800007\du}}
\pgfpathlineto{\pgfpoint{-21.582228\du}{6.797320\du}}
\pgfpathlineto{\pgfpoint{-21.578944\du}{6.794335\du}}
\pgfpathlineto{\pgfpoint{-21.575660\du}{6.791648\du}}
\pgfpathlineto{\pgfpoint{-21.573570\du}{6.788364\du}}
\pgfpathlineto{\pgfpoint{-21.571480\du}{6.785677\du}}
\pgfpathlineto{\pgfpoint{-21.569390\du}{6.782691\du}}
\pgfpathlineto{\pgfpoint{-21.568196\du}{6.779706\du}}
\pgfpathlineto{\pgfpoint{-21.567599\du}{6.777019\du}}
\pgfpathlineto{\pgfpoint{-21.567300\du}{6.773436\du}}
\pgfusepath{stroke}
\pgfsetbuttcap
\pgfsetmiterjoin
\pgfsetdash{}{0pt}
\definecolor{dialinecolor}{rgb}{0.000000, 0.470588, 0.666667}
\pgfsetfillcolor{dialinecolor}
\pgfpathmoveto{\pgfpoint{-21.947954\du}{6.690439\du}}
\pgfpathlineto{\pgfpoint{-21.947954\du}{6.785677\du}}
\pgfpathlineto{\pgfpoint{-21.555060\du}{6.785677\du}}
\pgfpathlineto{\pgfpoint{-21.555060\du}{6.690439\du}}
\pgfpathlineto{\pgfpoint{-21.947954\du}{6.690439\du}}
\pgfusepath{fill}
\pgfsetbuttcap
\pgfsetmiterjoin
\pgfsetdash{}{0pt}
\definecolor{dialinecolor}{rgb}{0.000000, 0.705882, 1.000000}
\pgfsetfillcolor{dialinecolor}
\pgfpathmoveto{\pgfpoint{-21.555358\du}{6.684169\du}}
\pgfpathlineto{\pgfpoint{-21.555657\du}{6.680587\du}}
\pgfpathlineto{\pgfpoint{-21.556553\du}{6.677303\du}}
\pgfpathlineto{\pgfpoint{-21.557747\du}{6.674317\du}}
\pgfpathlineto{\pgfpoint{-21.559538\du}{6.670734\du}}
\pgfpathlineto{\pgfpoint{-21.561628\du}{6.667450\du}}
\pgfpathlineto{\pgfpoint{-21.564912\du}{6.664166\du}}
\pgfpathlineto{\pgfpoint{-21.567599\du}{6.661479\du}}
\pgfpathlineto{\pgfpoint{-21.571480\du}{6.657598\du}}
\pgfpathlineto{\pgfpoint{-21.575660\du}{6.655210\du}}
\pgfpathlineto{\pgfpoint{-21.580138\du}{6.652224\du}}
\pgfpathlineto{\pgfpoint{-21.585512\du}{6.649239\du}}
\pgfpathlineto{\pgfpoint{-21.590587\du}{6.646552\du}}
\pgfpathlineto{\pgfpoint{-21.596559\du}{6.643865\du}}
\pgfpathlineto{\pgfpoint{-21.602530\du}{6.641476\du}}
\pgfpathlineto{\pgfpoint{-21.609695\du}{6.638789\du}}
\pgfpathlineto{\pgfpoint{-21.616562\du}{6.636401\du}}
\pgfpathlineto{\pgfpoint{-21.624025\du}{6.634311\du}}
\pgfpathlineto{\pgfpoint{-21.631489\du}{6.632221\du}}
\pgfpathlineto{\pgfpoint{-21.639550\du}{6.630430\du}}
\pgfpathlineto{\pgfpoint{-21.647909\du}{6.628340\du}}
\pgfpathlineto{\pgfpoint{-21.656269\du}{6.626549\du}}
\pgfpathlineto{\pgfpoint{-21.665225\du}{6.625355\du}}
\pgfpathlineto{\pgfpoint{-21.674182\du}{6.623862\du}}
\pgfpathlineto{\pgfpoint{-21.683437\du}{6.622668\du}}
\pgfpathlineto{\pgfpoint{-21.693289\du}{6.621473\du}}
\pgfpathlineto{\pgfpoint{-21.702544\du}{6.620279\du}}
\pgfpathlineto{\pgfpoint{-21.712397\du}{6.619981\du}}
\pgfpathlineto{\pgfpoint{-21.721950\du}{6.619085\du}}
\pgfpathlineto{\pgfpoint{-21.731803\du}{6.618786\du}}
\pgfpathlineto{\pgfpoint{-21.742252\du}{6.618488\du}}
\pgfpathlineto{\pgfpoint{-21.751507\du}{6.618488\du}}
\pgfpathlineto{\pgfpoint{-21.751507\du}{6.618488\du}}
\pgfpathlineto{\pgfpoint{-21.761359\du}{6.618488\du}}
\pgfpathlineto{\pgfpoint{-21.771510\du}{6.618786\du}}
\pgfpathlineto{\pgfpoint{-21.781362\du}{6.619085\du}}
\pgfpathlineto{\pgfpoint{-21.791214\du}{6.619981\du}}
\pgfpathlineto{\pgfpoint{-21.801067\du}{6.620279\du}}
\pgfpathlineto{\pgfpoint{-21.810322\du}{6.621473\du}}
\pgfpathlineto{\pgfpoint{-21.820174\du}{6.622668\du}}
\pgfpathlineto{\pgfpoint{-21.829429\du}{6.623862\du}}
\pgfpathlineto{\pgfpoint{-21.838386\du}{6.625355\du}}
\pgfpathlineto{\pgfpoint{-21.847044\du}{6.626549\du}}
\pgfpathlineto{\pgfpoint{-21.855403\du}{6.628340\du}}
\pgfpathlineto{\pgfpoint{-21.863763\du}{6.630430\du}}
\pgfpathlineto{\pgfpoint{-21.871823\du}{6.632221\du}}
\pgfpathlineto{\pgfpoint{-21.879287\du}{6.634311\du}}
\pgfpathlineto{\pgfpoint{-21.886751\du}{6.636401\du}}
\pgfpathlineto{\pgfpoint{-21.893916\du}{6.638789\du}}
\pgfpathlineto{\pgfpoint{-21.900783\du}{6.641476\du}}
\pgfpathlineto{\pgfpoint{-21.906754\du}{6.643865\du}}
\pgfpathlineto{\pgfpoint{-21.913024\du}{6.646552\du}}
\pgfpathlineto{\pgfpoint{-21.918099\du}{6.649239\du}}
\pgfpathlineto{\pgfpoint{-21.923174\du}{6.652224\du}}
\pgfpathlineto{\pgfpoint{-21.927951\du}{6.655210\du}}
\pgfpathlineto{\pgfpoint{-21.932131\du}{6.657598\du}}
\pgfpathlineto{\pgfpoint{-21.935714\du}{6.661479\du}}
\pgfpathlineto{\pgfpoint{-21.938699\du}{6.664166\du}}
\pgfpathlineto{\pgfpoint{-21.941386\du}{6.667450\du}}
\pgfpathlineto{\pgfpoint{-21.943775\du}{6.670734\du}}
\pgfpathlineto{\pgfpoint{-21.945566\du}{6.674317\du}}
\pgfpathlineto{\pgfpoint{-21.947059\du}{6.677303\du}}
\pgfpathlineto{\pgfpoint{-21.947954\du}{6.680587\du}}
\pgfpathlineto{\pgfpoint{-21.947954\du}{6.684169\du}}
\pgfpathlineto{\pgfpoint{-21.947954\du}{6.684169\du}}
\pgfpathlineto{\pgfpoint{-21.947954\du}{6.687453\du}}
\pgfpathlineto{\pgfpoint{-21.947059\du}{6.690737\du}}
\pgfpathlineto{\pgfpoint{-21.945566\du}{6.693723\du}}
\pgfpathlineto{\pgfpoint{-21.943775\du}{6.697306\du}}
\pgfpathlineto{\pgfpoint{-21.941386\du}{6.700590\du}}
\pgfpathlineto{\pgfpoint{-21.938699\du}{6.703575\du}}
\pgfpathlineto{\pgfpoint{-21.935714\du}{6.706561\du}}
\pgfpathlineto{\pgfpoint{-21.932131\du}{6.709845\du}}
\pgfpathlineto{\pgfpoint{-21.927951\du}{6.712830\du}}
\pgfpathlineto{\pgfpoint{-21.923174\du}{6.715816\du}}
\pgfpathlineto{\pgfpoint{-21.918099\du}{6.718801\du}}
\pgfpathlineto{\pgfpoint{-21.913024\du}{6.721488\du}}
\pgfpathlineto{\pgfpoint{-21.906754\du}{6.723877\du}}
\pgfpathlineto{\pgfpoint{-21.900783\du}{6.726564\du}}
\pgfpathlineto{\pgfpoint{-21.893916\du}{6.729251\du}}
\pgfpathlineto{\pgfpoint{-21.886751\du}{6.731340\du}}
\pgfpathlineto{\pgfpoint{-21.879287\du}{6.733729\du}}
\pgfpathlineto{\pgfpoint{-21.871823\du}{6.735819\du}}
\pgfpathlineto{\pgfpoint{-21.863763\du}{6.737610\du}}
\pgfpathlineto{\pgfpoint{-21.855403\du}{6.739700\du}}
\pgfpathlineto{\pgfpoint{-21.847044\du}{6.741491\du}}
\pgfpathlineto{\pgfpoint{-21.838386\du}{6.742685\du}}
\pgfpathlineto{\pgfpoint{-21.829429\du}{6.744178\du}}
\pgfpathlineto{\pgfpoint{-21.820174\du}{6.745372\du}}
\pgfpathlineto{\pgfpoint{-21.810322\du}{6.746268\du}}
\pgfpathlineto{\pgfpoint{-21.801067\du}{6.747164\du}}
\pgfpathlineto{\pgfpoint{-21.791214\du}{6.748059\du}}
\pgfpathlineto{\pgfpoint{-21.781362\du}{6.748955\du}}
\pgfpathlineto{\pgfpoint{-21.771510\du}{6.749254\du}}
\pgfpathlineto{\pgfpoint{-21.761359\du}{6.749552\du}}
\pgfpathlineto{\pgfpoint{-21.751507\du}{6.749552\du}}
\pgfpathlineto{\pgfpoint{-21.751507\du}{6.749552\du}}
\pgfpathlineto{\pgfpoint{-21.742252\du}{6.749552\du}}
\pgfpathlineto{\pgfpoint{-21.731803\du}{6.749254\du}}
\pgfpathlineto{\pgfpoint{-21.721950\du}{6.748955\du}}
\pgfpathlineto{\pgfpoint{-21.712397\du}{6.748059\du}}
\pgfpathlineto{\pgfpoint{-21.702544\du}{6.747164\du}}
\pgfpathlineto{\pgfpoint{-21.693289\du}{6.746268\du}}
\pgfpathlineto{\pgfpoint{-21.683437\du}{6.745372\du}}
\pgfpathlineto{\pgfpoint{-21.674182\du}{6.744178\du}}
\pgfpathlineto{\pgfpoint{-21.665225\du}{6.742685\du}}
\pgfpathlineto{\pgfpoint{-21.656269\du}{6.741491\du}}
\pgfpathlineto{\pgfpoint{-21.647909\du}{6.739700\du}}
\pgfpathlineto{\pgfpoint{-21.639550\du}{6.737610\du}}
\pgfpathlineto{\pgfpoint{-21.631489\du}{6.735819\du}}
\pgfpathlineto{\pgfpoint{-21.624025\du}{6.733729\du}}
\pgfpathlineto{\pgfpoint{-21.616562\du}{6.731340\du}}
\pgfpathlineto{\pgfpoint{-21.609695\du}{6.729251\du}}
\pgfpathlineto{\pgfpoint{-21.602530\du}{6.726564\du}}
\pgfpathlineto{\pgfpoint{-21.596559\du}{6.723877\du}}
\pgfpathlineto{\pgfpoint{-21.590587\du}{6.721488\du}}
\pgfpathlineto{\pgfpoint{-21.585512\du}{6.718801\du}}
\pgfpathlineto{\pgfpoint{-21.580138\du}{6.715816\du}}
\pgfpathlineto{\pgfpoint{-21.575660\du}{6.712830\du}}
\pgfpathlineto{\pgfpoint{-21.571480\du}{6.709845\du}}
\pgfpathlineto{\pgfpoint{-21.567599\du}{6.706561\du}}
\pgfpathlineto{\pgfpoint{-21.564912\du}{6.703575\du}}
\pgfpathlineto{\pgfpoint{-21.561628\du}{6.700590\du}}
\pgfpathlineto{\pgfpoint{-21.559538\du}{6.697306\du}}
\pgfpathlineto{\pgfpoint{-21.557747\du}{6.693723\du}}
\pgfpathlineto{\pgfpoint{-21.556553\du}{6.690737\du}}
\pgfpathlineto{\pgfpoint{-21.555657\du}{6.687453\du}}
\pgfpathlineto{\pgfpoint{-21.555358\du}{6.684169\du}}
\pgfusepath{fill}
\pgfsetbuttcap
\pgfsetmiterjoin
\pgfsetdash{}{0pt}
\definecolor{dialinecolor}{rgb}{0.666667, 0.901961, 1.000000}
\pgfsetstrokecolor{dialinecolor}
\pgfpathmoveto{\pgfpoint{-21.567300\du}{6.678497\du}}
\pgfpathlineto{\pgfpoint{-21.567599\du}{6.675511\du}}
\pgfpathlineto{\pgfpoint{-21.568196\du}{6.672526\du}}
\pgfpathlineto{\pgfpoint{-21.569390\du}{6.669839\du}}
\pgfpathlineto{\pgfpoint{-21.571480\du}{6.666555\du}}
\pgfpathlineto{\pgfpoint{-21.573570\du}{6.663868\du}}
\pgfpathlineto{\pgfpoint{-21.575660\du}{6.660584\du}}
\pgfpathlineto{\pgfpoint{-21.578944\du}{6.657598\du}}
\pgfpathlineto{\pgfpoint{-21.582228\du}{6.655210\du}}
\pgfpathlineto{\pgfpoint{-21.586408\du}{6.652523\du}}
\pgfpathlineto{\pgfpoint{-21.591185\du}{6.649836\du}}
\pgfpathlineto{\pgfpoint{-21.596260\du}{6.647149\du}}
\pgfpathlineto{\pgfpoint{-21.601335\du}{6.644760\du}}
\pgfpathlineto{\pgfpoint{-21.607008\du}{6.641775\du}}
\pgfpathlineto{\pgfpoint{-21.612979\du}{6.639984\du}}
\pgfpathlineto{\pgfpoint{-21.619547\du}{6.637894\du}}
\pgfpathlineto{\pgfpoint{-21.626414\du}{6.635207\du}}
\pgfpathlineto{\pgfpoint{-21.632982\du}{6.633714\du}}
\pgfpathlineto{\pgfpoint{-21.641043\du}{6.631923\du}}
\pgfpathlineto{\pgfpoint{-21.648507\du}{6.629833\du}}
\pgfpathlineto{\pgfpoint{-21.656269\du}{6.628340\du}}
\pgfpathlineto{\pgfpoint{-21.665225\du}{6.626847\du}}
\pgfpathlineto{\pgfpoint{-21.673585\du}{6.625355\du}}
\pgfpathlineto{\pgfpoint{-21.681944\du}{6.624160\du}}
\pgfpathlineto{\pgfpoint{-21.691200\du}{6.622966\du}}
\pgfpathlineto{\pgfpoint{-21.699858\du}{6.621772\du}}
\pgfpathlineto{\pgfpoint{-21.709411\du}{6.621175\du}}
\pgfpathlineto{\pgfpoint{-21.718965\du}{6.620279\du}}
\pgfpathlineto{\pgfpoint{-21.728220\du}{6.620279\du}}
\pgfpathlineto{\pgfpoint{-21.738072\du}{6.619981\du}}
\pgfpathlineto{\pgfpoint{-21.747924\du}{6.619085\du}}
\pgfpathlineto{\pgfpoint{-21.757478\du}{6.619085\du}}
\pgfpathlineto{\pgfpoint{-21.757478\du}{6.619085\du}}
\pgfpathlineto{\pgfpoint{-21.766733\du}{6.619085\du}}
\pgfpathlineto{\pgfpoint{-21.776585\du}{6.619981\du}}
\pgfpathlineto{\pgfpoint{-21.786139\du}{6.620279\du}}
\pgfpathlineto{\pgfpoint{-21.795693\du}{6.620279\du}}
\pgfpathlineto{\pgfpoint{-21.805246\du}{6.621175\du}}
\pgfpathlineto{\pgfpoint{-21.814501\du}{6.621772\du}}
\pgfpathlineto{\pgfpoint{-21.823458\du}{6.622966\du}}
\pgfpathlineto{\pgfpoint{-21.832713\du}{6.624160\du}}
\pgfpathlineto{\pgfpoint{-21.841073\du}{6.625355\du}}
\pgfpathlineto{\pgfpoint{-21.850029\du}{6.626847\du}}
\pgfpathlineto{\pgfpoint{-21.858389\du}{6.628340\du}}
\pgfpathlineto{\pgfpoint{-21.866151\du}{6.629833\du}}
\pgfpathlineto{\pgfpoint{-21.874212\du}{6.631923\du}}
\pgfpathlineto{\pgfpoint{-21.881676\du}{6.633714\du}}
\pgfpathlineto{\pgfpoint{-21.888244\du}{6.635207\du}}
\pgfpathlineto{\pgfpoint{-21.895708\du}{6.637894\du}}
\pgfpathlineto{\pgfpoint{-21.901679\du}{6.639984\du}}
\pgfpathlineto{\pgfpoint{-21.907650\du}{6.641775\du}}
\pgfpathlineto{\pgfpoint{-21.913919\du}{6.644760\du}}
\pgfpathlineto{\pgfpoint{-21.918995\du}{6.647149\du}}
\pgfpathlineto{\pgfpoint{-21.924070\du}{6.649836\du}}
\pgfpathlineto{\pgfpoint{-21.928250\du}{6.652523\du}}
\pgfpathlineto{\pgfpoint{-21.932430\du}{6.655210\du}}
\pgfpathlineto{\pgfpoint{-21.936311\du}{6.657598\du}}
\pgfpathlineto{\pgfpoint{-21.939595\du}{6.660584\du}}
\pgfpathlineto{\pgfpoint{-21.941386\du}{6.663868\du}}
\pgfpathlineto{\pgfpoint{-21.943775\du}{6.666555\du}}
\pgfpathlineto{\pgfpoint{-21.945566\du}{6.669839\du}}
\pgfpathlineto{\pgfpoint{-21.947059\du}{6.672526\du}}
\pgfpathlineto{\pgfpoint{-21.947357\du}{6.675511\du}}
\pgfpathlineto{\pgfpoint{-21.947954\du}{6.678497\du}}
\pgfpathlineto{\pgfpoint{-21.947954\du}{6.678497\du}}
\pgfpathlineto{\pgfpoint{-21.947357\du}{6.682079\du}}
\pgfpathlineto{\pgfpoint{-21.947059\du}{6.684766\du}}
\pgfpathlineto{\pgfpoint{-21.945566\du}{6.687453\du}}
\pgfpathlineto{\pgfpoint{-21.943775\du}{6.690737\du}}
\pgfpathlineto{\pgfpoint{-21.941386\du}{6.693424\du}}
\pgfpathlineto{\pgfpoint{-21.939595\du}{6.697007\du}}
\pgfpathlineto{\pgfpoint{-21.936311\du}{6.699694\du}}
\pgfpathlineto{\pgfpoint{-21.932430\du}{6.702381\du}}
\pgfpathlineto{\pgfpoint{-21.928250\du}{6.705068\du}}
\pgfpathlineto{\pgfpoint{-21.924070\du}{6.707755\du}}
\pgfpathlineto{\pgfpoint{-21.918995\du}{6.710442\du}}
\pgfpathlineto{\pgfpoint{-21.913919\du}{6.712532\du}}
\pgfpathlineto{\pgfpoint{-21.907650\du}{6.715219\du}}
\pgfpathlineto{\pgfpoint{-21.901679\du}{6.717309\du}}
\pgfpathlineto{\pgfpoint{-21.895708\du}{6.719697\du}}
\pgfpathlineto{\pgfpoint{-21.888244\du}{6.721787\du}}
\pgfpathlineto{\pgfpoint{-21.881676\du}{6.723877\du}}
\pgfpathlineto{\pgfpoint{-21.874212\du}{6.725668\du}}
\pgfpathlineto{\pgfpoint{-21.866151\du}{6.727161\du}}
\pgfpathlineto{\pgfpoint{-21.858389\du}{6.729251\du}}
\pgfpathlineto{\pgfpoint{-21.850029\du}{6.730743\du}}
\pgfpathlineto{\pgfpoint{-21.841073\du}{6.732236\du}}
\pgfpathlineto{\pgfpoint{-21.832713\du}{6.733430\du}}
\pgfpathlineto{\pgfpoint{-21.823458\du}{6.734625\du}}
\pgfpathlineto{\pgfpoint{-21.814501\du}{6.735222\du}}
\pgfpathlineto{\pgfpoint{-21.805246\du}{6.736416\du}}
\pgfpathlineto{\pgfpoint{-21.795693\du}{6.736714\du}}
\pgfpathlineto{\pgfpoint{-21.786139\du}{6.737312\du}}
\pgfpathlineto{\pgfpoint{-21.776585\du}{6.737610\du}}
\pgfpathlineto{\pgfpoint{-21.766733\du}{6.738207\du}}
\pgfpathlineto{\pgfpoint{-21.757478\du}{6.738207\du}}
\pgfpathlineto{\pgfpoint{-21.757478\du}{6.738207\du}}
\pgfpathlineto{\pgfpoint{-21.747924\du}{6.738207\du}}
\pgfpathlineto{\pgfpoint{-21.738072\du}{6.737610\du}}
\pgfpathlineto{\pgfpoint{-21.728220\du}{6.737312\du}}
\pgfpathlineto{\pgfpoint{-21.718965\du}{6.736714\du}}
\pgfpathlineto{\pgfpoint{-21.709411\du}{6.736416\du}}
\pgfpathlineto{\pgfpoint{-21.699858\du}{6.735222\du}}
\pgfpathlineto{\pgfpoint{-21.691200\du}{6.734625\du}}
\pgfpathlineto{\pgfpoint{-21.681944\du}{6.733430\du}}
\pgfpathlineto{\pgfpoint{-21.673585\du}{6.732236\du}}
\pgfpathlineto{\pgfpoint{-21.665225\du}{6.730743\du}}
\pgfpathlineto{\pgfpoint{-21.656269\du}{6.729251\du}}
\pgfpathlineto{\pgfpoint{-21.648507\du}{6.727161\du}}
\pgfpathlineto{\pgfpoint{-21.641043\du}{6.725668\du}}
\pgfpathlineto{\pgfpoint{-21.632982\du}{6.723877\du}}
\pgfpathlineto{\pgfpoint{-21.626414\du}{6.721787\du}}
\pgfpathlineto{\pgfpoint{-21.619547\du}{6.719697\du}}
\pgfpathlineto{\pgfpoint{-21.612979\du}{6.717309\du}}
\pgfpathlineto{\pgfpoint{-21.607008\du}{6.715219\du}}
\pgfpathlineto{\pgfpoint{-21.601335\du}{6.712532\du}}
\pgfpathlineto{\pgfpoint{-21.596260\du}{6.710442\du}}
\pgfpathlineto{\pgfpoint{-21.591185\du}{6.707755\du}}
\pgfpathlineto{\pgfpoint{-21.586408\du}{6.705068\du}}
\pgfpathlineto{\pgfpoint{-21.582228\du}{6.702381\du}}
\pgfpathlineto{\pgfpoint{-21.578944\du}{6.699694\du}}
\pgfpathlineto{\pgfpoint{-21.575660\du}{6.697007\du}}
\pgfpathlineto{\pgfpoint{-21.573570\du}{6.693424\du}}
\pgfpathlineto{\pgfpoint{-21.571480\du}{6.690737\du}}
\pgfpathlineto{\pgfpoint{-21.569390\du}{6.687453\du}}
\pgfpathlineto{\pgfpoint{-21.568196\du}{6.684766\du}}
\pgfpathlineto{\pgfpoint{-21.567599\du}{6.682079\du}}
\pgfpathlineto{\pgfpoint{-21.567300\du}{6.678497\du}}
\pgfusepath{stroke}
\pgfsetbuttcap
\pgfsetmiterjoin
\pgfsetdash{}{0pt}
\definecolor{dialinecolor}{rgb}{0.000000, 0.000000, 0.000000}
\pgfsetfillcolor{dialinecolor}
\pgfpathmoveto{\pgfpoint{-21.745536\du}{6.666555\du}}
\pgfpathlineto{\pgfpoint{-21.721652\du}{6.678497\du}}
\pgfpathlineto{\pgfpoint{-21.650298\du}{6.654911\du}}
\pgfpathlineto{\pgfpoint{-21.614472\du}{6.666555\du}}
\pgfpathlineto{\pgfpoint{-21.638057\du}{6.642671\du}}
\pgfpathlineto{\pgfpoint{-21.721652\du}{6.642671\du}}
\pgfpathlineto{\pgfpoint{-21.685826\du}{6.642671\du}}
\pgfpathlineto{\pgfpoint{-21.745536\du}{6.666555\du}}
\pgfusepath{fill}
\pgfsetbuttcap
\pgfsetmiterjoin
\pgfsetdash{}{0pt}
\definecolor{dialinecolor}{rgb}{0.000000, 0.000000, 0.000000}
\pgfsetfillcolor{dialinecolor}
\pgfpathmoveto{\pgfpoint{-21.757478\du}{6.702381\du}}
\pgfpathlineto{\pgfpoint{-21.793006\du}{6.690439\du}}
\pgfpathlineto{\pgfpoint{-21.852716\du}{6.714323\du}}
\pgfpathlineto{\pgfpoint{-21.888244\du}{6.702381\du}}
\pgfpathlineto{\pgfpoint{-21.876003\du}{6.738207\du}}
\pgfpathlineto{\pgfpoint{-21.793006\du}{6.738207\du}}
\pgfpathlineto{\pgfpoint{-21.828533\du}{6.726564\du}}
\pgfpathlineto{\pgfpoint{-21.757478\du}{6.702381\du}}
\pgfusepath{fill}
\pgfsetbuttcap
\pgfsetmiterjoin
\pgfsetdash{}{0pt}
\definecolor{dialinecolor}{rgb}{0.000000, 0.000000, 0.000000}
\pgfsetfillcolor{dialinecolor}
\pgfpathmoveto{\pgfpoint{-21.876003\du}{6.642671\du}}
\pgfpathlineto{\pgfpoint{-21.852716\du}{6.631027\du}}
\pgfpathlineto{\pgfpoint{-21.781362\du}{6.666555\du}}
\pgfpathlineto{\pgfpoint{-21.757478\du}{6.654911\du}}
\pgfpathlineto{\pgfpoint{-21.768823\du}{6.678497\du}}
\pgfpathlineto{\pgfpoint{-21.852716\du}{6.678497\du}}
\pgfpathlineto{\pgfpoint{-21.816890\du}{6.666555\du}}
\pgfpathlineto{\pgfpoint{-21.876003\du}{6.642671\du}}
\pgfusepath{fill}
\pgfsetbuttcap
\pgfsetmiterjoin
\pgfsetdash{}{0pt}
\definecolor{dialinecolor}{rgb}{0.000000, 0.000000, 0.000000}
\pgfsetfillcolor{dialinecolor}
\pgfpathmoveto{\pgfpoint{-21.626414\du}{6.726564\du}}
\pgfpathlineto{\pgfpoint{-21.650298\du}{6.738207\du}}
\pgfpathlineto{\pgfpoint{-21.721652\du}{6.702381\du}}
\pgfpathlineto{\pgfpoint{-21.757478\du}{6.714323\du}}
\pgfpathlineto{\pgfpoint{-21.733295\du}{6.690439\du}}
\pgfpathlineto{\pgfpoint{-21.650298\du}{6.690439\du}}
\pgfpathlineto{\pgfpoint{-21.685826\du}{6.702381\du}}
\pgfpathlineto{\pgfpoint{-21.626414\du}{6.726564\du}}
\pgfusepath{fill}
\pgfsetbuttcap
\pgfsetmiterjoin
\pgfsetdash{}{0pt}
\definecolor{dialinecolor}{rgb}{1.000000, 1.000000, 1.000000}
\pgfsetfillcolor{dialinecolor}
\pgfpathmoveto{\pgfpoint{-21.745536\du}{6.678497\du}}
\pgfpathlineto{\pgfpoint{-21.721652\du}{6.678497\du}}
\pgfpathlineto{\pgfpoint{-21.650298\du}{6.654911\du}}
\pgfpathlineto{\pgfpoint{-21.614472\du}{6.666555\du}}
\pgfpathlineto{\pgfpoint{-21.638057\du}{6.642671\du}}
\pgfpathlineto{\pgfpoint{-21.709411\du}{6.642671\du}}
\pgfpathlineto{\pgfpoint{-21.685826\du}{6.642671\du}}
\pgfpathlineto{\pgfpoint{-21.745536\du}{6.678497\du}}
\pgfusepath{fill}
\pgfsetbuttcap
\pgfsetmiterjoin
\pgfsetdash{}{0pt}
\definecolor{dialinecolor}{rgb}{1.000000, 1.000000, 1.000000}
\pgfsetfillcolor{dialinecolor}
\pgfpathmoveto{\pgfpoint{-21.757478\du}{6.702381\du}}
\pgfpathlineto{\pgfpoint{-21.781362\du}{6.690439\du}}
\pgfpathlineto{\pgfpoint{-21.852716\du}{6.714323\du}}
\pgfpathlineto{\pgfpoint{-21.888244\du}{6.714323\du}}
\pgfpathlineto{\pgfpoint{-21.864360\du}{6.738207\du}}
\pgfpathlineto{\pgfpoint{-21.781362\du}{6.738207\du}}
\pgfpathlineto{\pgfpoint{-21.816890\du}{6.726564\du}}
\pgfpathlineto{\pgfpoint{-21.757478\du}{6.702381\du}}
\pgfusepath{fill}
\pgfsetbuttcap
\pgfsetmiterjoin
\pgfsetdash{}{0pt}
\definecolor{dialinecolor}{rgb}{1.000000, 1.000000, 1.000000}
\pgfsetfillcolor{dialinecolor}
\pgfpathmoveto{\pgfpoint{-21.876003\du}{6.642671\du}}
\pgfpathlineto{\pgfpoint{-21.852716\du}{6.642671\du}}
\pgfpathlineto{\pgfpoint{-21.781362\du}{6.666555\du}}
\pgfpathlineto{\pgfpoint{-21.745536\du}{6.654911\du}}
\pgfpathlineto{\pgfpoint{-21.768823\du}{6.678497\du}}
\pgfpathlineto{\pgfpoint{-21.852716\du}{6.678497\du}}
\pgfpathlineto{\pgfpoint{-21.816890\du}{6.678497\du}}
\pgfpathlineto{\pgfpoint{-21.876003\du}{6.642671\du}}
\pgfusepath{fill}
\pgfsetbuttcap
\pgfsetmiterjoin
\pgfsetdash{}{0pt}
\definecolor{dialinecolor}{rgb}{1.000000, 1.000000, 1.000000}
\pgfsetfillcolor{dialinecolor}
\pgfpathmoveto{\pgfpoint{-21.626414\du}{6.726564\du}}
\pgfpathlineto{\pgfpoint{-21.650298\du}{6.738207\du}}
\pgfpathlineto{\pgfpoint{-21.721652\du}{6.714323\du}}
\pgfpathlineto{\pgfpoint{-21.745536\du}{6.714323\du}}
\pgfpathlineto{\pgfpoint{-21.733295\du}{6.690439\du}}
\pgfpathlineto{\pgfpoint{-21.650298\du}{6.690439\du}}
\pgfpathlineto{\pgfpoint{-21.685826\du}{6.702381\du}}
\pgfpathlineto{\pgfpoint{-21.626414\du}{6.726564\du}}
\pgfusepath{fill}
\pgfsetbuttcap
\pgfsetmiterjoin
\pgfsetdash{}{0pt}
\definecolor{dialinecolor}{rgb}{0.666667, 0.901961, 1.000000}
\pgfsetstrokecolor{dialinecolor}
\pgfpathmoveto{\pgfpoint{-21.947954\du}{6.678497\du}}
\pgfpathlineto{\pgfpoint{-21.947954\du}{6.773436\du}}
\pgfusepath{stroke}
\pgfsetbuttcap
\pgfsetmiterjoin
\pgfsetdash{}{0pt}
\definecolor{dialinecolor}{rgb}{0.666667, 0.901961, 1.000000}
\pgfsetstrokecolor{dialinecolor}
\pgfpathmoveto{\pgfpoint{-21.567300\du}{6.678497\du}}
\pgfpathlineto{\pgfpoint{-21.567300\du}{6.773436\du}}
\pgfusepath{stroke}
\pgfsetlinewidth{0.100000\du}
\pgfsetdash{}{0pt}
\pgfsetdash{}{0pt}
\pgfsetbuttcap
\pgfsetmiterjoin
\pgfsetlinewidth{0.001000\du}
\pgfsetbuttcap
\pgfsetmiterjoin
\pgfsetdash{}{0pt}
\definecolor{dialinecolor}{rgb}{0.717647, 0.717647, 0.615686}
\pgfsetfillcolor{dialinecolor}
\pgfpathmoveto{\pgfpoint{-26.155200\du}{7.040680\du}}
\pgfpathlineto{\pgfpoint{-24.156774\du}{7.040680\du}}
\pgfpathlineto{\pgfpoint{-24.156774\du}{7.410013\du}}
\pgfpathlineto{\pgfpoint{-26.155200\du}{7.410013\du}}
\pgfpathlineto{\pgfpoint{-26.155200\du}{7.040680\du}}
\pgfusepath{fill}
\pgfsetbuttcap
\pgfsetmiterjoin
\pgfsetdash{}{0pt}
\definecolor{dialinecolor}{rgb}{0.286275, 0.286275, 0.211765}
\pgfsetstrokecolor{dialinecolor}
\pgfpathmoveto{\pgfpoint{-26.155200\du}{7.040680\du}}
\pgfpathlineto{\pgfpoint{-24.156774\du}{7.040680\du}}
\pgfpathlineto{\pgfpoint{-24.156774\du}{7.410013\du}}
\pgfpathlineto{\pgfpoint{-26.155200\du}{7.410013\du}}
\pgfpathlineto{\pgfpoint{-26.155200\du}{7.040680\du}}
\pgfusepath{stroke}
\pgfsetbuttcap
\pgfsetmiterjoin
\pgfsetdash{}{0pt}
\definecolor{dialinecolor}{rgb}{0.788235, 0.788235, 0.713726}
\pgfsetfillcolor{dialinecolor}
\pgfpathmoveto{\pgfpoint{-26.155200\du}{7.040680\du}}
\pgfpathlineto{\pgfpoint{-25.943298\du}{6.839798\du}}
\pgfpathlineto{\pgfpoint{-23.944873\du}{6.839798\du}}
\pgfpathlineto{\pgfpoint{-24.156774\du}{7.040680\du}}
\pgfpathlineto{\pgfpoint{-26.155200\du}{7.040680\du}}
\pgfusepath{fill}
\pgfsetbuttcap
\pgfsetmiterjoin
\pgfsetdash{}{0pt}
\definecolor{dialinecolor}{rgb}{0.286275, 0.286275, 0.211765}
\pgfsetstrokecolor{dialinecolor}
\pgfpathmoveto{\pgfpoint{-26.155200\du}{7.040680\du}}
\pgfpathlineto{\pgfpoint{-25.943298\du}{6.839798\du}}
\pgfpathlineto{\pgfpoint{-23.944873\du}{6.839798\du}}
\pgfpathlineto{\pgfpoint{-24.156774\du}{7.040680\du}}
\pgfpathlineto{\pgfpoint{-26.155200\du}{7.040680\du}}
\pgfusepath{stroke}
\pgfsetlinewidth{0.106000\du}
\pgfsetbuttcap
\pgfsetmiterjoin
\pgfsetdash{}{0pt}
\definecolor{dialinecolor}{rgb}{0.000000, 0.000000, 0.000000}
\pgfsetstrokecolor{dialinecolor}
\pgfpathmoveto{\pgfpoint{-24.268550\du}{7.208816\du}}
\pgfpathlineto{\pgfpoint{-24.748084\du}{7.208816\du}}
\pgfusepath{stroke}
\pgfsetlinewidth{0.001000\du}
\pgfsetbuttcap
\pgfsetmiterjoin
\pgfsetdash{}{0pt}
\definecolor{dialinecolor}{rgb}{0.478431, 0.478431, 0.352941}
\pgfsetfillcolor{dialinecolor}
\pgfpathmoveto{\pgfpoint{-24.156774\du}{7.410013\du}}
\pgfpathlineto{\pgfpoint{-23.944873\du}{7.197481\du}}
\pgfpathlineto{\pgfpoint{-23.944873\du}{6.839798\du}}
\pgfpathlineto{\pgfpoint{-24.156774\du}{7.040680\du}}
\pgfpathlineto{\pgfpoint{-24.156774\du}{7.410013\du}}
\pgfusepath{fill}
\pgfsetbuttcap
\pgfsetmiterjoin
\pgfsetdash{}{0pt}
\definecolor{dialinecolor}{rgb}{0.286275, 0.286275, 0.211765}
\pgfsetstrokecolor{dialinecolor}
\pgfpathmoveto{\pgfpoint{-24.156774\du}{7.410013\du}}
\pgfpathlineto{\pgfpoint{-23.944873\du}{7.197481\du}}
\pgfpathlineto{\pgfpoint{-23.944873\du}{6.839798\du}}
\pgfpathlineto{\pgfpoint{-24.156774\du}{7.040680\du}}
\pgfpathlineto{\pgfpoint{-24.156774\du}{7.410013\du}}
\pgfusepath{stroke}
\pgfsetbuttcap
\pgfsetmiterjoin
\pgfsetdash{}{0pt}
\definecolor{dialinecolor}{rgb}{0.788235, 0.788235, 0.713726}
\pgfsetfillcolor{dialinecolor}
\pgfpathmoveto{\pgfpoint{-26.143865\du}{7.644270\du}}
\pgfpathlineto{\pgfpoint{-25.920943\du}{7.364987\du}}
\pgfpathlineto{\pgfpoint{-24.380011\du}{7.364987\du}}
\pgfpathlineto{\pgfpoint{-24.602933\du}{7.644270\du}}
\pgfpathlineto{\pgfpoint{-26.143865\du}{7.644270\du}}
\pgfusepath{fill}
\pgfsetbuttcap
\pgfsetmiterjoin
\pgfsetdash{}{0pt}
\definecolor{dialinecolor}{rgb}{0.286275, 0.286275, 0.211765}
\pgfsetstrokecolor{dialinecolor}
\pgfpathmoveto{\pgfpoint{-26.143865\du}{7.644270\du}}
\pgfpathlineto{\pgfpoint{-25.920943\du}{7.364987\du}}
\pgfpathlineto{\pgfpoint{-24.380011\du}{7.364987\du}}
\pgfpathlineto{\pgfpoint{-24.602933\du}{7.644270\du}}
\pgfpathlineto{\pgfpoint{-26.143865\du}{7.644270\du}}
\pgfusepath{stroke}
\pgfsetbuttcap
\pgfsetmiterjoin
\pgfsetdash{}{0pt}
\definecolor{dialinecolor}{rgb}{0.478431, 0.478431, 0.352941}
\pgfsetfillcolor{dialinecolor}
\pgfpathmoveto{\pgfpoint{-24.602933\du}{7.700000\du}}
\pgfpathlineto{\pgfpoint{-24.380011\du}{7.465743\du}}
\pgfpathlineto{\pgfpoint{-24.380011\du}{7.364987\du}}
\pgfpathlineto{\pgfpoint{-24.602933\du}{7.644270\du}}
\pgfpathlineto{\pgfpoint{-24.602933\du}{7.700000\du}}
\pgfusepath{fill}
\pgfsetbuttcap
\pgfsetmiterjoin
\pgfsetdash{}{0pt}
\definecolor{dialinecolor}{rgb}{0.286275, 0.286275, 0.211765}
\pgfsetstrokecolor{dialinecolor}
\pgfpathmoveto{\pgfpoint{-24.602933\du}{7.700000\du}}
\pgfpathlineto{\pgfpoint{-24.380011\du}{7.465743\du}}
\pgfpathlineto{\pgfpoint{-24.380011\du}{7.364987\du}}
\pgfpathlineto{\pgfpoint{-24.602933\du}{7.644270\du}}
\pgfpathlineto{\pgfpoint{-24.602933\du}{7.700000\du}}
\pgfusepath{stroke}
\pgfsetbuttcap
\pgfsetmiterjoin
\pgfsetdash{}{0pt}
\definecolor{dialinecolor}{rgb}{0.717647, 0.717647, 0.615686}
\pgfsetfillcolor{dialinecolor}
\pgfpathmoveto{\pgfpoint{-26.143865\du}{7.644270\du}}
\pgfpathlineto{\pgfpoint{-24.602933\du}{7.644270\du}}
\pgfpathlineto{\pgfpoint{-24.602933\du}{7.700000\du}}
\pgfpathlineto{\pgfpoint{-26.143865\du}{7.700000\du}}
\pgfpathlineto{\pgfpoint{-26.143865\du}{7.644270\du}}
\pgfusepath{fill}
\pgfsetbuttcap
\pgfsetmiterjoin
\pgfsetdash{}{0pt}
\definecolor{dialinecolor}{rgb}{0.286275, 0.286275, 0.211765}
\pgfsetstrokecolor{dialinecolor}
\pgfpathmoveto{\pgfpoint{-26.143865\du}{7.644270\du}}
\pgfpathlineto{\pgfpoint{-24.602933\du}{7.644270\du}}
\pgfpathlineto{\pgfpoint{-24.602933\du}{7.700000\du}}
\pgfpathlineto{\pgfpoint{-26.143865\du}{7.700000\du}}
\pgfpathlineto{\pgfpoint{-26.143865\du}{7.644270\du}}
\pgfusepath{stroke}
\pgfsetbuttcap
\pgfsetmiterjoin
\pgfsetdash{}{0pt}
\definecolor{dialinecolor}{rgb}{0.000000, 0.000000, 0.000000}
\pgfsetfillcolor{dialinecolor}
\pgfpathmoveto{\pgfpoint{-25.853878\du}{6.996285\du}}
\pgfpathlineto{\pgfpoint{-25.686056\du}{6.839798\du}}
\pgfpathlineto{\pgfpoint{-24.268550\du}{6.839798\du}}
\pgfpathlineto{\pgfpoint{-24.424407\du}{6.996285\du}}
\pgfpathlineto{\pgfpoint{-25.853878\du}{6.996285\du}}
\pgfusepath{fill}
\pgfsetbuttcap
\pgfsetmiterjoin
\pgfsetdash{}{0pt}
\definecolor{dialinecolor}{rgb}{0.000000, 0.000000, 0.000000}
\pgfsetstrokecolor{dialinecolor}
\pgfpathmoveto{\pgfpoint{-25.853878\du}{6.996285\du}}
\pgfpathlineto{\pgfpoint{-25.686056\du}{6.839798\du}}
\pgfpathlineto{\pgfpoint{-24.268550\du}{6.839798\du}}
\pgfpathlineto{\pgfpoint{-24.424407\du}{6.996285\du}}
\pgfpathlineto{\pgfpoint{-25.853878\du}{6.996285\du}}
\pgfusepath{stroke}
\pgfsetbuttcap
\pgfsetmiterjoin
\pgfsetdash{}{0pt}
\definecolor{dialinecolor}{rgb}{0.788235, 0.788235, 0.713726}
\pgfsetfillcolor{dialinecolor}
\pgfpathmoveto{\pgfpoint{-25.865213\du}{5.845151\du}}
\pgfpathlineto{\pgfpoint{-25.708412\du}{5.700000\du}}
\pgfpathlineto{\pgfpoint{-24.290276\du}{5.700000\du}}
\pgfpathlineto{\pgfpoint{-24.447077\du}{5.845151\du}}
\pgfpathlineto{\pgfpoint{-25.865213\du}{5.845151\du}}
\pgfusepath{fill}
\pgfsetbuttcap
\pgfsetmiterjoin
\pgfsetdash{}{0pt}
\definecolor{dialinecolor}{rgb}{0.286275, 0.286275, 0.211765}
\pgfsetstrokecolor{dialinecolor}
\pgfpathmoveto{\pgfpoint{-25.865213\du}{5.845151\du}}
\pgfpathlineto{\pgfpoint{-25.708412\du}{5.700000\du}}
\pgfpathlineto{\pgfpoint{-24.290276\du}{5.700000\du}}
\pgfpathlineto{\pgfpoint{-24.447077\du}{5.845151\du}}
\pgfpathlineto{\pgfpoint{-25.865213\du}{5.845151\du}}
\pgfusepath{stroke}
\pgfsetbuttcap
\pgfsetmiterjoin
\pgfsetdash{}{0pt}
\definecolor{dialinecolor}{rgb}{0.717647, 0.717647, 0.615686}
\pgfsetfillcolor{dialinecolor}
\pgfpathmoveto{\pgfpoint{-25.865213\du}{5.845151\du}}
\pgfpathlineto{\pgfpoint{-24.435742\du}{5.845151\du}}
\pgfpathlineto{\pgfpoint{-24.435742\du}{6.973615\du}}
\pgfpathlineto{\pgfpoint{-25.865213\du}{6.973615\du}}
\pgfpathlineto{\pgfpoint{-25.865213\du}{5.845151\du}}
\pgfusepath{fill}
\pgfsetbuttcap
\pgfsetmiterjoin
\pgfsetdash{}{0pt}
\definecolor{dialinecolor}{rgb}{0.286275, 0.286275, 0.211765}
\pgfsetstrokecolor{dialinecolor}
\pgfpathmoveto{\pgfpoint{-25.865213\du}{5.845151\du}}
\pgfpathlineto{\pgfpoint{-24.436371\du}{5.845151\du}}
\pgfpathlineto{\pgfpoint{-24.436371\du}{6.973300\du}}
\pgfpathlineto{\pgfpoint{-25.865213\du}{6.973300\du}}
\pgfpathlineto{\pgfpoint{-25.865213\du}{5.845151\du}}
\pgfusepath{stroke}
\pgfsetbuttcap
\pgfsetmiterjoin
\pgfsetdash{}{0pt}
\definecolor{dialinecolor}{rgb}{1.000000, 1.000000, 1.000000}
\pgfsetfillcolor{dialinecolor}
\pgfpathmoveto{\pgfpoint{-25.742417\du}{5.989987\du}}
\pgfpathlineto{\pgfpoint{-24.558852\du}{5.989987\du}}
\pgfpathlineto{\pgfpoint{-24.558852\du}{6.861839\du}}
\pgfpathlineto{\pgfpoint{-25.742417\du}{6.861839\du}}
\pgfpathlineto{\pgfpoint{-25.742417\du}{5.989987\du}}
\pgfusepath{fill}
\pgfsetbuttcap
\pgfsetmiterjoin
\pgfsetdash{}{0pt}
\definecolor{dialinecolor}{rgb}{0.286275, 0.286275, 0.211765}
\pgfsetstrokecolor{dialinecolor}
\pgfpathmoveto{\pgfpoint{-25.742417\du}{5.989987\du}}
\pgfpathlineto{\pgfpoint{-24.558852\du}{5.989987\du}}
\pgfpathlineto{\pgfpoint{-24.558852\du}{6.861524\du}}
\pgfpathlineto{\pgfpoint{-25.742417\du}{6.861524\du}}
\pgfpathlineto{\pgfpoint{-25.742417\du}{5.989987\du}}
\pgfusepath{stroke}
\pgfsetbuttcap
\pgfsetmiterjoin
\pgfsetdash{}{0pt}
\definecolor{dialinecolor}{rgb}{0.478431, 0.478431, 0.352941}
\pgfsetfillcolor{dialinecolor}
\pgfpathmoveto{\pgfpoint{-24.447077\du}{6.962909\du}}
\pgfpathlineto{\pgfpoint{-24.290276\du}{6.806423\du}}
\pgfpathlineto{\pgfpoint{-24.290276\du}{5.700000\du}}
\pgfpathlineto{\pgfpoint{-24.447077\du}{5.845151\du}}
\pgfpathlineto{\pgfpoint{-24.447077\du}{6.962909\du}}
\pgfusepath{fill}
\pgfsetbuttcap
\pgfsetmiterjoin
\pgfsetdash{}{0pt}
\definecolor{dialinecolor}{rgb}{0.286275, 0.286275, 0.211765}
\pgfsetstrokecolor{dialinecolor}
\pgfpathmoveto{\pgfpoint{-24.447077\du}{6.962909\du}}
\pgfpathlineto{\pgfpoint{-24.290276\du}{6.806423\du}}
\pgfpathlineto{\pgfpoint{-24.290276\du}{5.700000\du}}
\pgfpathlineto{\pgfpoint{-24.447077\du}{5.845151\du}}
\pgfpathlineto{\pgfpoint{-24.447077\du}{6.962909\du}}
\pgfusepath{stroke}
\pgfsetlinewidth{0.100000\du}
\pgfsetdash{}{0pt}
\pgfsetdash{}{0pt}
\pgfsetbuttcap
\pgfsetmiterjoin
\pgfsetlinewidth{0.001000\du}
\pgfsetbuttcap
\pgfsetmiterjoin
\pgfsetdash{}{0pt}
\definecolor{dialinecolor}{rgb}{0.717647, 0.717647, 0.615686}
\pgfsetfillcolor{dialinecolor}
\pgfpathmoveto{\pgfpoint{-19.786900\du}{-2.780497\du}}
\pgfpathlineto{\pgfpoint{-19.786900\du}{-4.157964\du}}
\pgfpathlineto{\pgfpoint{-19.229882\du}{-4.157964\du}}
\pgfpathlineto{\pgfpoint{-19.229882\du}{-2.780497\du}}
\pgfpathlineto{\pgfpoint{-19.786900\du}{-2.780497\du}}
\pgfusepath{fill}
\pgfsetbuttcap
\pgfsetmiterjoin
\pgfsetdash{}{0pt}
\definecolor{dialinecolor}{rgb}{0.286275, 0.286275, 0.211765}
\pgfsetstrokecolor{dialinecolor}
\pgfpathmoveto{\pgfpoint{-19.786900\du}{-2.780497\du}}
\pgfpathlineto{\pgfpoint{-19.786900\du}{-4.157964\du}}
\pgfpathlineto{\pgfpoint{-19.229882\du}{-4.157964\du}}
\pgfpathlineto{\pgfpoint{-19.229882\du}{-2.780497\du}}
\pgfpathlineto{\pgfpoint{-19.786900\du}{-2.780497\du}}
\pgfusepath{stroke}
\pgfsetbuttcap
\pgfsetmiterjoin
\pgfsetdash{}{0pt}
\definecolor{dialinecolor}{rgb}{0.647059, 0.647059, 0.521569}
\pgfsetfillcolor{dialinecolor}
\pgfpathmoveto{\pgfpoint{-19.229882\du}{-2.780497\du}}
\pgfpathlineto{\pgfpoint{-19.229882\du}{-4.157964\du}}
\pgfpathlineto{\pgfpoint{-18.672420\du}{-4.157964\du}}
\pgfpathlineto{\pgfpoint{-18.672420\du}{-2.780497\du}}
\pgfpathlineto{\pgfpoint{-19.229882\du}{-2.780497\du}}
\pgfusepath{fill}
\pgfsetbuttcap
\pgfsetmiterjoin
\pgfsetdash{}{0pt}
\definecolor{dialinecolor}{rgb}{0.286275, 0.286275, 0.211765}
\pgfsetstrokecolor{dialinecolor}
\pgfpathmoveto{\pgfpoint{-19.229882\du}{-2.780497\du}}
\pgfpathlineto{\pgfpoint{-19.229882\du}{-4.157964\du}}
\pgfpathlineto{\pgfpoint{-18.672420\du}{-4.157964\du}}
\pgfpathlineto{\pgfpoint{-18.672420\du}{-2.780497\du}}
\pgfpathlineto{\pgfpoint{-19.229882\du}{-2.780497\du}}
\pgfusepath{stroke}
\pgfsetbuttcap
\pgfsetmiterjoin
\pgfsetdash{}{0pt}
\definecolor{dialinecolor}{rgb}{0.478431, 0.478431, 0.352941}
\pgfsetfillcolor{dialinecolor}
\pgfpathmoveto{\pgfpoint{-18.381493\du}{-3.072532\du}}
\pgfpathlineto{\pgfpoint{-18.381493\du}{-4.450000\du}}
\pgfpathlineto{\pgfpoint{-18.672420\du}{-4.157964\du}}
\pgfpathlineto{\pgfpoint{-18.672420\du}{-2.780497\du}}
\pgfpathlineto{\pgfpoint{-18.381493\du}{-3.072532\du}}
\pgfusepath{fill}
\pgfsetbuttcap
\pgfsetmiterjoin
\pgfsetdash{}{0pt}
\definecolor{dialinecolor}{rgb}{0.286275, 0.286275, 0.211765}
\pgfsetstrokecolor{dialinecolor}
\pgfpathmoveto{\pgfpoint{-18.381493\du}{-3.072532\du}}
\pgfpathlineto{\pgfpoint{-18.381493\du}{-4.450000\du}}
\pgfpathlineto{\pgfpoint{-18.672420\du}{-4.157964\du}}
\pgfpathlineto{\pgfpoint{-18.672420\du}{-2.780497\du}}
\pgfpathlineto{\pgfpoint{-18.381493\du}{-3.072532\du}}
\pgfusepath{stroke}
\pgfsetbuttcap
\pgfsetmiterjoin
\pgfsetdash{}{0pt}
\definecolor{dialinecolor}{rgb}{0.788235, 0.788235, 0.713726}
\pgfsetfillcolor{dialinecolor}
\pgfpathmoveto{\pgfpoint{-19.496195\du}{-4.450000\du}}
\pgfpathlineto{\pgfpoint{-19.786900\du}{-4.157964\du}}
\pgfpathlineto{\pgfpoint{-18.672420\du}{-4.157964\du}}
\pgfpathlineto{\pgfpoint{-18.381493\du}{-4.450000\du}}
\pgfpathlineto{\pgfpoint{-19.496195\du}{-4.450000\du}}
\pgfusepath{fill}
\pgfsetbuttcap
\pgfsetmiterjoin
\pgfsetdash{}{0pt}
\definecolor{dialinecolor}{rgb}{0.286275, 0.286275, 0.211765}
\pgfsetstrokecolor{dialinecolor}
\pgfpathmoveto{\pgfpoint{-19.496195\du}{-4.450000\du}}
\pgfpathlineto{\pgfpoint{-19.786900\du}{-4.157964\du}}
\pgfpathlineto{\pgfpoint{-18.672420\du}{-4.157964\du}}
\pgfpathlineto{\pgfpoint{-18.381493\du}{-4.450000\du}}
\pgfpathlineto{\pgfpoint{-19.496195\du}{-4.450000\du}}
\pgfusepath{stroke}
\pgfsetbuttcap
\pgfsetmiterjoin
\pgfsetdash{}{0pt}
\definecolor{dialinecolor}{rgb}{0.478431, 0.478431, 0.352941}
\pgfsetstrokecolor{dialinecolor}
\pgfpathmoveto{\pgfpoint{-18.971994\du}{-4.416739\du}}
\pgfpathlineto{\pgfpoint{-19.196398\du}{-4.191448\du}}
\pgfusepath{stroke}
\pgfsetbuttcap
\pgfsetmiterjoin
\pgfsetdash{}{0pt}
\definecolor{dialinecolor}{rgb}{0.717647, 0.717647, 0.615686}
\pgfsetfillcolor{dialinecolor}
\pgfpathmoveto{\pgfpoint{-19.504621\du}{-2.789145\du}}
\pgfpathlineto{\pgfpoint{-19.786900\du}{-2.789145\du}}
\pgfpathlineto{\pgfpoint{-19.786678\du}{-2.782936\du}}
\pgfpathlineto{\pgfpoint{-19.785570\du}{-2.776727\du}}
\pgfpathlineto{\pgfpoint{-19.783796\du}{-2.770297\du}}
\pgfpathlineto{\pgfpoint{-19.781356\du}{-2.764088\du}}
\pgfpathlineto{\pgfpoint{-19.777809\du}{-2.758101\du}}
\pgfpathlineto{\pgfpoint{-19.774039\du}{-2.751670\du}}
\pgfpathlineto{\pgfpoint{-19.769382\du}{-2.745683\du}}
\pgfpathlineto{\pgfpoint{-19.764060\du}{-2.740140\du}}
\pgfpathlineto{\pgfpoint{-19.758295\du}{-2.734374\du}}
\pgfpathlineto{\pgfpoint{-19.751643\du}{-2.728609\du}}
\pgfpathlineto{\pgfpoint{-19.744325\du}{-2.723065\du}}
\pgfpathlineto{\pgfpoint{-19.736343\du}{-2.717965\du}}
\pgfpathlineto{\pgfpoint{-19.727916\du}{-2.712865\du}}
\pgfpathlineto{\pgfpoint{-19.718825\du}{-2.707765\du}}
\pgfpathlineto{\pgfpoint{-19.709290\du}{-2.703109\du}}
\pgfpathlineto{\pgfpoint{-19.699312\du}{-2.698674\du}}
\pgfpathlineto{\pgfpoint{-19.688446\du}{-2.694239\du}}
\pgfpathlineto{\pgfpoint{-19.677359\du}{-2.690469\du}}
\pgfpathlineto{\pgfpoint{-19.666050\du}{-2.686700\du}}
\pgfpathlineto{\pgfpoint{-19.653854\du}{-2.683152\du}}
\pgfpathlineto{\pgfpoint{-19.641658\du}{-2.680047\du}}
\pgfpathlineto{\pgfpoint{-19.629019\du}{-2.676943\du}}
\pgfpathlineto{\pgfpoint{-19.615936\du}{-2.674504\du}}
\pgfpathlineto{\pgfpoint{-19.602632\du}{-2.672286\du}}
\pgfpathlineto{\pgfpoint{-19.589327\du}{-2.670069\du}}
\pgfpathlineto{\pgfpoint{-19.575357\du}{-2.668517\du}}
\pgfpathlineto{\pgfpoint{-19.561609\du}{-2.666743\du}}
\pgfpathlineto{\pgfpoint{-19.547418\du}{-2.665856\du}}
\pgfpathlineto{\pgfpoint{-19.533004\du}{-2.664969\du}}
\pgfpathlineto{\pgfpoint{-19.518813\du}{-2.664525\du}}
\pgfpathlineto{\pgfpoint{-19.504621\du}{-2.664303\du}}
\pgfpathlineto{\pgfpoint{-19.490430\du}{-2.664525\du}}
\pgfpathlineto{\pgfpoint{-19.476238\du}{-2.664969\du}}
\pgfpathlineto{\pgfpoint{-19.462046\du}{-2.665856\du}}
\pgfpathlineto{\pgfpoint{-19.447855\du}{-2.666743\du}}
\pgfpathlineto{\pgfpoint{-19.433885\du}{-2.668517\du}}
\pgfpathlineto{\pgfpoint{-19.420137\du}{-2.670069\du}}
\pgfpathlineto{\pgfpoint{-19.406832\du}{-2.672286\du}}
\pgfpathlineto{\pgfpoint{-19.393306\du}{-2.674504\du}}
\pgfpathlineto{\pgfpoint{-19.380223\du}{-2.676943\du}}
\pgfpathlineto{\pgfpoint{-19.367584\du}{-2.680047\du}}
\pgfpathlineto{\pgfpoint{-19.355166\du}{-2.683152\du}}
\pgfpathlineto{\pgfpoint{-19.343192\du}{-2.686700\du}}
\pgfpathlineto{\pgfpoint{-19.331662\du}{-2.690469\du}}
\pgfpathlineto{\pgfpoint{-19.320574\du}{-2.694239\du}}
\pgfpathlineto{\pgfpoint{-19.310153\du}{-2.698674\du}}
\pgfpathlineto{\pgfpoint{-19.299731\du}{-2.703109\du}}
\pgfpathlineto{\pgfpoint{-19.290639\du}{-2.707765\du}}
\pgfpathlineto{\pgfpoint{-19.281326\du}{-2.712865\du}}
\pgfpathlineto{\pgfpoint{-19.272678\du}{-2.717965\du}}
\pgfpathlineto{\pgfpoint{-19.264917\du}{-2.723065\du}}
\pgfpathlineto{\pgfpoint{-19.257599\du}{-2.728609\du}}
\pgfpathlineto{\pgfpoint{-19.250947\du}{-2.734374\du}}
\pgfpathlineto{\pgfpoint{-19.245404\du}{-2.740140\du}}
\pgfpathlineto{\pgfpoint{-19.240082\du}{-2.745683\du}}
\pgfpathlineto{\pgfpoint{-19.235203\du}{-2.751670\du}}
\pgfpathlineto{\pgfpoint{-19.231434\du}{-2.758101\du}}
\pgfpathlineto{\pgfpoint{-19.228108\du}{-2.764088\du}}
\pgfpathlineto{\pgfpoint{-19.225447\du}{-2.770297\du}}
\pgfpathlineto{\pgfpoint{-19.223894\du}{-2.776727\du}}
\pgfpathlineto{\pgfpoint{-19.222564\du}{-2.782936\du}}
\pgfpathlineto{\pgfpoint{-19.222121\du}{-2.789145\du}}
\pgfpathlineto{\pgfpoint{-19.504621\du}{-2.789145\du}}
\pgfusepath{fill}
\pgfsetbuttcap
\pgfsetmiterjoin
\pgfsetdash{}{0pt}
\definecolor{dialinecolor}{rgb}{0.286275, 0.286275, 0.211765}
\pgfsetstrokecolor{dialinecolor}
\pgfpathmoveto{\pgfpoint{-19.786900\du}{-2.789145\du}}
\pgfpathlineto{\pgfpoint{-19.786678\du}{-2.782936\du}}
\pgfpathlineto{\pgfpoint{-19.785570\du}{-2.776727\du}}
\pgfpathlineto{\pgfpoint{-19.783796\du}{-2.770297\du}}
\pgfpathlineto{\pgfpoint{-19.781356\du}{-2.764088\du}}
\pgfpathlineto{\pgfpoint{-19.777809\du}{-2.758101\du}}
\pgfpathlineto{\pgfpoint{-19.774039\du}{-2.751670\du}}
\pgfpathlineto{\pgfpoint{-19.769382\du}{-2.745683\du}}
\pgfpathlineto{\pgfpoint{-19.764060\du}{-2.740140\du}}
\pgfpathlineto{\pgfpoint{-19.758295\du}{-2.734374\du}}
\pgfpathlineto{\pgfpoint{-19.751643\du}{-2.728609\du}}
\pgfpathlineto{\pgfpoint{-19.744325\du}{-2.723065\du}}
\pgfpathlineto{\pgfpoint{-19.736343\du}{-2.717965\du}}
\pgfpathlineto{\pgfpoint{-19.727916\du}{-2.712865\du}}
\pgfpathlineto{\pgfpoint{-19.718825\du}{-2.707765\du}}
\pgfpathlineto{\pgfpoint{-19.709068\du}{-2.703109\du}}
\pgfpathlineto{\pgfpoint{-19.698868\du}{-2.698674\du}}
\pgfpathlineto{\pgfpoint{-19.688446\du}{-2.694239\du}}
\pgfpathlineto{\pgfpoint{-19.677359\du}{-2.690469\du}}
\pgfpathlineto{\pgfpoint{-19.665828\du}{-2.686700\du}}
\pgfpathlineto{\pgfpoint{-19.653854\du}{-2.683152\du}}
\pgfpathlineto{\pgfpoint{-19.641437\du}{-2.680047\du}}
\pgfpathlineto{\pgfpoint{-19.629019\du}{-2.676943\du}}
\pgfpathlineto{\pgfpoint{-19.615936\du}{-2.674504\du}}
\pgfpathlineto{\pgfpoint{-19.602632\du}{-2.672286\du}}
\pgfpathlineto{\pgfpoint{-19.589105\du}{-2.670069\du}}
\pgfpathlineto{\pgfpoint{-19.575135\du}{-2.668517\du}}
\pgfpathlineto{\pgfpoint{-19.561609\du}{-2.666743\du}}
\pgfpathlineto{\pgfpoint{-19.547418\du}{-2.665856\du}}
\pgfpathlineto{\pgfpoint{-19.533004\du}{-2.664969\du}}
\pgfpathlineto{\pgfpoint{-19.518813\du}{-2.664525\du}}
\pgfpathlineto{\pgfpoint{-19.504178\du}{-2.664303\du}}
\pgfpathlineto{\pgfpoint{-19.490208\du}{-2.664525\du}}
\pgfpathlineto{\pgfpoint{-19.475795\du}{-2.664969\du}}
\pgfpathlineto{\pgfpoint{-19.461825\du}{-2.665856\du}}
\pgfpathlineto{\pgfpoint{-19.447411\du}{-2.666743\du}}
\pgfpathlineto{\pgfpoint{-19.433885\du}{-2.668517\du}}
\pgfpathlineto{\pgfpoint{-19.420137\du}{-2.670069\du}}
\pgfpathlineto{\pgfpoint{-19.406389\du}{-2.672286\du}}
\pgfpathlineto{\pgfpoint{-19.393084\du}{-2.674504\du}}
\pgfpathlineto{\pgfpoint{-19.380002\du}{-2.676943\du}}
\pgfpathlineto{\pgfpoint{-19.367584\du}{-2.680047\du}}
\pgfpathlineto{\pgfpoint{-19.355166\du}{-2.683152\du}}
\pgfpathlineto{\pgfpoint{-19.343192\du}{-2.686700\du}}
\pgfpathlineto{\pgfpoint{-19.331662\du}{-2.690469\du}}
\pgfpathlineto{\pgfpoint{-19.320574\du}{-2.694239\du}}
\pgfpathlineto{\pgfpoint{-19.309487\du}{-2.698674\du}}
\pgfpathlineto{\pgfpoint{-19.299731\du}{-2.703109\du}}
\pgfpathlineto{\pgfpoint{-19.290196\du}{-2.707765\du}}
\pgfpathlineto{\pgfpoint{-19.281104\du}{-2.712865\du}}
\pgfpathlineto{\pgfpoint{-19.272456\du}{-2.717965\du}}
\pgfpathlineto{\pgfpoint{-19.264695\du}{-2.723065\du}}
\pgfpathlineto{\pgfpoint{-19.257378\du}{-2.728609\du}}
\pgfpathlineto{\pgfpoint{-19.250504\du}{-2.734374\du}}
\pgfpathlineto{\pgfpoint{-19.244738\du}{-2.740140\du}}
\pgfpathlineto{\pgfpoint{-19.239416\du}{-2.745683\du}}
\pgfpathlineto{\pgfpoint{-19.234982\du}{-2.751670\du}}
\pgfpathlineto{\pgfpoint{-19.231212\du}{-2.758101\du}}
\pgfpathlineto{\pgfpoint{-19.227664\du}{-2.764088\du}}
\pgfpathlineto{\pgfpoint{-19.225225\du}{-2.770297\du}}
\pgfpathlineto{\pgfpoint{-19.223673\du}{-2.776727\du}}
\pgfpathlineto{\pgfpoint{-19.222121\du}{-2.782936\du}}
\pgfpathlineto{\pgfpoint{-19.221899\du}{-2.789145\du}}
\pgfusepath{stroke}
\pgfsetbuttcap
\pgfsetmiterjoin
\pgfsetdash{}{0pt}
\definecolor{dialinecolor}{rgb}{0.576471, 0.576471, 0.423529}
\pgfsetfillcolor{dialinecolor}
\pgfpathmoveto{\pgfpoint{-19.395967\du}{-4.049310\du}}
\pgfpathlineto{\pgfpoint{-19.395967\du}{-3.982566\du}}
\pgfpathlineto{\pgfpoint{-19.395967\du}{-3.907395\du}}
\pgfpathlineto{\pgfpoint{-19.395967\du}{-3.815593\du}}
\pgfpathlineto{\pgfpoint{-19.395967\du}{-3.723792\du}}
\pgfpathlineto{\pgfpoint{-19.395967\du}{-3.515353\du}}
\pgfpathlineto{\pgfpoint{-19.395967\du}{-3.298267\du}}
\pgfpathlineto{\pgfpoint{-19.395967\du}{-3.197817\du}}
\pgfpathlineto{\pgfpoint{-19.395967\du}{-3.097811\du}}
\pgfpathlineto{\pgfpoint{-19.395967\du}{-3.014214\du}}
\pgfpathlineto{\pgfpoint{-19.395967\du}{-2.939043\du}}
\pgfpathlineto{\pgfpoint{-19.395967\du}{-2.872298\du}}
\pgfpathlineto{\pgfpoint{-19.395967\du}{-2.822406\du}}
\pgfpathlineto{\pgfpoint{-19.395967\du}{-2.797128\du}}
\pgfpathlineto{\pgfpoint{-19.395967\du}{-2.780497\du}}
\pgfpathlineto{\pgfpoint{-19.229882\du}{-2.780497\du}}
\pgfpathlineto{\pgfpoint{-19.229882\du}{-4.149760\du}}
\pgfpathlineto{\pgfpoint{-19.238086\du}{-4.149760\du}}
\pgfpathlineto{\pgfpoint{-19.229882\du}{-4.141112\du}}
\pgfpathlineto{\pgfpoint{-19.229882\du}{-4.132907\du}}
\pgfpathlineto{\pgfpoint{-19.254495\du}{-4.107851\du}}
\pgfpathlineto{\pgfpoint{-19.296626\du}{-4.082794\du}}
\pgfpathlineto{\pgfpoint{-19.346075\du}{-4.066163\du}}
\pgfpathlineto{\pgfpoint{-19.395967\du}{-4.049310\du}}
\pgfusepath{fill}
\pgfsetbuttcap
\pgfsetmiterjoin
\pgfsetdash{}{0pt}
\definecolor{dialinecolor}{rgb}{0.576471, 0.576471, 0.423529}
\pgfsetfillcolor{dialinecolor}
\pgfpathmoveto{\pgfpoint{-19.395967\du}{-2.680491\du}}
\pgfpathlineto{\pgfpoint{-19.337870\du}{-2.697121\du}}
\pgfpathlineto{\pgfpoint{-19.287978\du}{-2.721957\du}}
\pgfpathlineto{\pgfpoint{-19.263143\du}{-2.730605\du}}
\pgfpathlineto{\pgfpoint{-19.246291\du}{-2.747235\du}}
\pgfpathlineto{\pgfpoint{-19.238086\du}{-2.772292\du}}
\pgfpathlineto{\pgfpoint{-19.229882\du}{-2.788923\du}}
\pgfpathlineto{\pgfpoint{-19.246291\du}{-2.788923\du}}
\pgfpathlineto{\pgfpoint{-19.263143\du}{-2.788923\du}}
\pgfpathlineto{\pgfpoint{-19.287978\du}{-2.788923\du}}
\pgfpathlineto{\pgfpoint{-19.337870\du}{-2.788923\du}}
\pgfpathlineto{\pgfpoint{-19.395967\du}{-2.788923\du}}
\pgfpathlineto{\pgfpoint{-19.395967\du}{-2.755440\du}}
\pgfpathlineto{\pgfpoint{-19.395967\du}{-2.721957\du}}
\pgfpathlineto{\pgfpoint{-19.395967\du}{-2.705326\du}}
\pgfpathlineto{\pgfpoint{-19.395967\du}{-2.688695\du}}
\pgfpathlineto{\pgfpoint{-19.395967\du}{-2.680491\du}}
\pgfusepath{fill}
\pgfsetbuttcap
\pgfsetmiterjoin
\pgfsetdash{}{0pt}
\definecolor{dialinecolor}{rgb}{0.788235, 0.788235, 0.713726}
\pgfsetfillcolor{dialinecolor}
\pgfpathmoveto{\pgfpoint{-19.762065\du}{-4.174595\du}}
\pgfpathlineto{\pgfpoint{-19.786900\du}{-4.149760\du}}
\pgfpathlineto{\pgfpoint{-19.229882\du}{-4.149760\du}}
\pgfpathlineto{\pgfpoint{-19.229882\du}{-4.174595\du}}
\pgfpathlineto{\pgfpoint{-19.762065\du}{-4.174595\du}}
\pgfusepath{fill}
\pgfsetbuttcap
\pgfsetmiterjoin
\pgfsetdash{}{0pt}
\definecolor{dialinecolor}{rgb}{0.788235, 0.788235, 0.713726}
\pgfsetfillcolor{dialinecolor}
\pgfpathmoveto{\pgfpoint{-19.504621\du}{-4.158186\du}}
\pgfpathlineto{\pgfpoint{-19.786900\du}{-4.158186\du}}
\pgfpathlineto{\pgfpoint{-19.786678\du}{-4.152199\du}}
\pgfpathlineto{\pgfpoint{-19.785570\du}{-4.145547\du}}
\pgfpathlineto{\pgfpoint{-19.783796\du}{-4.139338\du}}
\pgfpathlineto{\pgfpoint{-19.781356\du}{-4.132907\du}}
\pgfpathlineto{\pgfpoint{-19.777809\du}{-4.126920\du}}
\pgfpathlineto{\pgfpoint{-19.774039\du}{-4.120712\du}}
\pgfpathlineto{\pgfpoint{-19.769382\du}{-4.114725\du}}
\pgfpathlineto{\pgfpoint{-19.764060\du}{-4.108959\du}}
\pgfpathlineto{\pgfpoint{-19.758295\du}{-4.103416\du}}
\pgfpathlineto{\pgfpoint{-19.751643\du}{-4.097650\du}}
\pgfpathlineto{\pgfpoint{-19.744325\du}{-4.092107\du}}
\pgfpathlineto{\pgfpoint{-19.736343\du}{-4.087007\du}}
\pgfpathlineto{\pgfpoint{-19.727916\du}{-4.081907\du}}
\pgfpathlineto{\pgfpoint{-19.718825\du}{-4.076806\du}}
\pgfpathlineto{\pgfpoint{-19.709290\du}{-4.072150\du}}
\pgfpathlineto{\pgfpoint{-19.699312\du}{-4.067715\du}}
\pgfpathlineto{\pgfpoint{-19.688446\du}{-4.063502\du}}
\pgfpathlineto{\pgfpoint{-19.677359\du}{-4.059511\du}}
\pgfpathlineto{\pgfpoint{-19.666050\du}{-4.055741\du}}
\pgfpathlineto{\pgfpoint{-19.653854\du}{-4.052415\du}}
\pgfpathlineto{\pgfpoint{-19.641658\du}{-4.049089\du}}
\pgfpathlineto{\pgfpoint{-19.629019\du}{-4.045984\du}}
\pgfpathlineto{\pgfpoint{-19.615936\du}{-4.043323\du}}
\pgfpathlineto{\pgfpoint{-19.602632\du}{-4.041328\du}}
\pgfpathlineto{\pgfpoint{-19.589327\du}{-4.038888\du}}
\pgfpathlineto{\pgfpoint{-19.575357\du}{-4.037336\du}}
\pgfpathlineto{\pgfpoint{-19.561609\du}{-4.035784\du}}
\pgfpathlineto{\pgfpoint{-19.547418\du}{-4.034897\du}}
\pgfpathlineto{\pgfpoint{-19.533004\du}{-4.034010\du}}
\pgfpathlineto{\pgfpoint{-19.518813\du}{-4.033567\du}}
\pgfpathlineto{\pgfpoint{-19.504621\du}{-4.033345\du}}
\pgfpathlineto{\pgfpoint{-19.490430\du}{-4.033567\du}}
\pgfpathlineto{\pgfpoint{-19.476238\du}{-4.034010\du}}
\pgfpathlineto{\pgfpoint{-19.462046\du}{-4.034897\du}}
\pgfpathlineto{\pgfpoint{-19.447855\du}{-4.035784\du}}
\pgfpathlineto{\pgfpoint{-19.433885\du}{-4.037336\du}}
\pgfpathlineto{\pgfpoint{-19.420137\du}{-4.038888\du}}
\pgfpathlineto{\pgfpoint{-19.406832\du}{-4.041328\du}}
\pgfpathlineto{\pgfpoint{-19.393306\du}{-4.043323\du}}
\pgfpathlineto{\pgfpoint{-19.380223\du}{-4.045984\du}}
\pgfpathlineto{\pgfpoint{-19.367584\du}{-4.049089\du}}
\pgfpathlineto{\pgfpoint{-19.355166\du}{-4.052415\du}}
\pgfpathlineto{\pgfpoint{-19.343192\du}{-4.055741\du}}
\pgfpathlineto{\pgfpoint{-19.331662\du}{-4.059511\du}}
\pgfpathlineto{\pgfpoint{-19.320574\du}{-4.063502\du}}
\pgfpathlineto{\pgfpoint{-19.310153\du}{-4.067715\du}}
\pgfpathlineto{\pgfpoint{-19.299731\du}{-4.072150\du}}
\pgfpathlineto{\pgfpoint{-19.290639\du}{-4.076806\du}}
\pgfpathlineto{\pgfpoint{-19.281326\du}{-4.081907\du}}
\pgfpathlineto{\pgfpoint{-19.272678\du}{-4.087007\du}}
\pgfpathlineto{\pgfpoint{-19.264917\du}{-4.092107\du}}
\pgfpathlineto{\pgfpoint{-19.257599\du}{-4.097650\du}}
\pgfpathlineto{\pgfpoint{-19.250947\du}{-4.103416\du}}
\pgfpathlineto{\pgfpoint{-19.245404\du}{-4.108959\du}}
\pgfpathlineto{\pgfpoint{-19.240082\du}{-4.114725\du}}
\pgfpathlineto{\pgfpoint{-19.235203\du}{-4.120712\du}}
\pgfpathlineto{\pgfpoint{-19.231434\du}{-4.126920\du}}
\pgfpathlineto{\pgfpoint{-19.228108\du}{-4.132907\du}}
\pgfpathlineto{\pgfpoint{-19.225447\du}{-4.139338\du}}
\pgfpathlineto{\pgfpoint{-19.223894\du}{-4.145547\du}}
\pgfpathlineto{\pgfpoint{-19.222564\du}{-4.152199\du}}
\pgfpathlineto{\pgfpoint{-19.222121\du}{-4.158186\du}}
\pgfpathlineto{\pgfpoint{-19.504621\du}{-4.158186\du}}
\pgfusepath{fill}
\pgfsetbuttcap
\pgfsetmiterjoin
\pgfsetdash{}{0pt}
\definecolor{dialinecolor}{rgb}{0.286275, 0.286275, 0.211765}
\pgfsetstrokecolor{dialinecolor}
\pgfpathmoveto{\pgfpoint{-19.786900\du}{-4.158186\du}}
\pgfpathlineto{\pgfpoint{-19.786678\du}{-4.152199\du}}
\pgfpathlineto{\pgfpoint{-19.785570\du}{-4.145547\du}}
\pgfpathlineto{\pgfpoint{-19.783796\du}{-4.139338\du}}
\pgfpathlineto{\pgfpoint{-19.781356\du}{-4.132907\du}}
\pgfpathlineto{\pgfpoint{-19.777809\du}{-4.126920\du}}
\pgfpathlineto{\pgfpoint{-19.774039\du}{-4.120712\du}}
\pgfpathlineto{\pgfpoint{-19.769382\du}{-4.114725\du}}
\pgfpathlineto{\pgfpoint{-19.764060\du}{-4.108959\du}}
\pgfpathlineto{\pgfpoint{-19.758295\du}{-4.103416\du}}
\pgfpathlineto{\pgfpoint{-19.751643\du}{-4.097650\du}}
\pgfpathlineto{\pgfpoint{-19.744325\du}{-4.092107\du}}
\pgfpathlineto{\pgfpoint{-19.736343\du}{-4.087007\du}}
\pgfpathlineto{\pgfpoint{-19.727916\du}{-4.081907\du}}
\pgfpathlineto{\pgfpoint{-19.718825\du}{-4.076806\du}}
\pgfpathlineto{\pgfpoint{-19.709068\du}{-4.072150\du}}
\pgfpathlineto{\pgfpoint{-19.698868\du}{-4.067715\du}}
\pgfpathlineto{\pgfpoint{-19.688446\du}{-4.063502\du}}
\pgfpathlineto{\pgfpoint{-19.677359\du}{-4.059511\du}}
\pgfpathlineto{\pgfpoint{-19.665828\du}{-4.055741\du}}
\pgfpathlineto{\pgfpoint{-19.653854\du}{-4.052415\du}}
\pgfpathlineto{\pgfpoint{-19.641437\du}{-4.049089\du}}
\pgfpathlineto{\pgfpoint{-19.629019\du}{-4.045984\du}}
\pgfpathlineto{\pgfpoint{-19.615936\du}{-4.043323\du}}
\pgfpathlineto{\pgfpoint{-19.602632\du}{-4.041328\du}}
\pgfpathlineto{\pgfpoint{-19.589105\du}{-4.038888\du}}
\pgfpathlineto{\pgfpoint{-19.575135\du}{-4.037336\du}}
\pgfpathlineto{\pgfpoint{-19.561609\du}{-4.035784\du}}
\pgfpathlineto{\pgfpoint{-19.547418\du}{-4.034897\du}}
\pgfpathlineto{\pgfpoint{-19.533004\du}{-4.034010\du}}
\pgfpathlineto{\pgfpoint{-19.518813\du}{-4.033567\du}}
\pgfpathlineto{\pgfpoint{-19.504178\du}{-4.033345\du}}
\pgfpathlineto{\pgfpoint{-19.490208\du}{-4.033567\du}}
\pgfpathlineto{\pgfpoint{-19.475795\du}{-4.034010\du}}
\pgfpathlineto{\pgfpoint{-19.461825\du}{-4.034897\du}}
\pgfpathlineto{\pgfpoint{-19.447411\du}{-4.035784\du}}
\pgfpathlineto{\pgfpoint{-19.433885\du}{-4.037336\du}}
\pgfpathlineto{\pgfpoint{-19.420137\du}{-4.038888\du}}
\pgfpathlineto{\pgfpoint{-19.406389\du}{-4.041328\du}}
\pgfpathlineto{\pgfpoint{-19.393084\du}{-4.043323\du}}
\pgfpathlineto{\pgfpoint{-19.380002\du}{-4.045984\du}}
\pgfpathlineto{\pgfpoint{-19.367584\du}{-4.049089\du}}
\pgfpathlineto{\pgfpoint{-19.355166\du}{-4.052415\du}}
\pgfpathlineto{\pgfpoint{-19.343192\du}{-4.055741\du}}
\pgfpathlineto{\pgfpoint{-19.331662\du}{-4.059511\du}}
\pgfpathlineto{\pgfpoint{-19.320574\du}{-4.063502\du}}
\pgfpathlineto{\pgfpoint{-19.309487\du}{-4.067715\du}}
\pgfpathlineto{\pgfpoint{-19.299731\du}{-4.072150\du}}
\pgfpathlineto{\pgfpoint{-19.290196\du}{-4.076806\du}}
\pgfpathlineto{\pgfpoint{-19.281104\du}{-4.081907\du}}
\pgfpathlineto{\pgfpoint{-19.272456\du}{-4.087007\du}}
\pgfpathlineto{\pgfpoint{-19.264695\du}{-4.092107\du}}
\pgfpathlineto{\pgfpoint{-19.257378\du}{-4.097650\du}}
\pgfpathlineto{\pgfpoint{-19.250504\du}{-4.103416\du}}
\pgfpathlineto{\pgfpoint{-19.244738\du}{-4.108959\du}}
\pgfpathlineto{\pgfpoint{-19.239416\du}{-4.114725\du}}
\pgfpathlineto{\pgfpoint{-19.234982\du}{-4.120712\du}}
\pgfpathlineto{\pgfpoint{-19.231212\du}{-4.126920\du}}
\pgfpathlineto{\pgfpoint{-19.227664\du}{-4.132907\du}}
\pgfpathlineto{\pgfpoint{-19.225225\du}{-4.139338\du}}
\pgfpathlineto{\pgfpoint{-19.223673\du}{-4.145547\du}}
\pgfpathlineto{\pgfpoint{-19.222121\du}{-4.152199\du}}
\pgfpathlineto{\pgfpoint{-19.221899\du}{-4.158186\du}}
\pgfusepath{stroke}
\pgfsetbuttcap
\pgfsetmiterjoin
\pgfsetdash{}{0pt}
\definecolor{dialinecolor}{rgb}{0.478431, 0.478431, 0.352941}
\pgfsetstrokecolor{dialinecolor}
\pgfpathmoveto{\pgfpoint{-19.737008\du}{-4.157964\du}}
\pgfpathlineto{\pgfpoint{-19.279774\du}{-4.157964\du}}
\pgfusepath{stroke}
\pgfsetbuttcap
\pgfsetmiterjoin
\pgfsetdash{}{0pt}
\definecolor{dialinecolor}{rgb}{0.576471, 0.576471, 0.423529}
\pgfsetfillcolor{dialinecolor}
\pgfpathmoveto{\pgfpoint{-19.221455\du}{-2.788923\du}}
\pgfpathlineto{\pgfpoint{-19.221455\du}{-4.149760\du}}
\pgfpathlineto{\pgfpoint{-19.121449\du}{-4.149760\du}}
\pgfpathlineto{\pgfpoint{-19.121449\du}{-2.788923\du}}
\pgfpathlineto{\pgfpoint{-19.221455\du}{-2.788923\du}}
\pgfusepath{fill}
\pgfsetbuttcap
\pgfsetmiterjoin
\pgfsetdash{}{0pt}
\definecolor{dialinecolor}{rgb}{0.576471, 0.576471, 0.423529}
\pgfsetfillcolor{dialinecolor}
\pgfpathmoveto{\pgfpoint{-19.138302\du}{-3.749070\du}}
\pgfpathlineto{\pgfpoint{-19.138302\du}{-4.066163\du}}
\pgfpathlineto{\pgfpoint{-18.763999\du}{-4.066163\du}}
\pgfpathlineto{\pgfpoint{-18.763999\du}{-3.749070\du}}
\pgfpathlineto{\pgfpoint{-19.138302\du}{-3.749070\du}}
\pgfusepath{fill}
\pgfsetbuttcap
\pgfsetmiterjoin
\pgfsetdash{}{0pt}
\definecolor{dialinecolor}{rgb}{0.286275, 0.286275, 0.211765}
\pgfsetstrokecolor{dialinecolor}
\pgfpathmoveto{\pgfpoint{-19.138302\du}{-3.749070\du}}
\pgfpathlineto{\pgfpoint{-19.138302\du}{-4.066163\du}}
\pgfpathlineto{\pgfpoint{-18.763999\du}{-4.066163\du}}
\pgfpathlineto{\pgfpoint{-18.763999\du}{-3.749070\du}}
\pgfpathlineto{\pgfpoint{-19.138302\du}{-3.749070\du}}
\pgfusepath{stroke}
\pgfsetbuttcap
\pgfsetmiterjoin
\pgfsetdash{}{0pt}
\definecolor{dialinecolor}{rgb}{0.478431, 0.478431, 0.352941}
\pgfsetfillcolor{dialinecolor}
\pgfpathmoveto{\pgfpoint{-19.130097\du}{-3.749070\du}}
\pgfpathlineto{\pgfpoint{-19.130097\du}{-4.057515\du}}
\pgfpathlineto{\pgfpoint{-19.121449\du}{-4.057515\du}}
\pgfpathlineto{\pgfpoint{-19.121449\du}{-3.749070\du}}
\pgfpathlineto{\pgfpoint{-19.130097\du}{-3.749070\du}}
\pgfusepath{fill}
\pgfsetbuttcap
\pgfsetmiterjoin
\pgfsetdash{}{0pt}
\definecolor{dialinecolor}{rgb}{0.831373, 0.000000, 0.039216}
\pgfsetfillcolor{dialinecolor}
\pgfpathmoveto{\pgfpoint{-19.778474\du}{-4.132907\du}}
\pgfpathlineto{\pgfpoint{-19.753639\du}{-4.107851\du}}
\pgfpathlineto{\pgfpoint{-19.753639\du}{-2.738809\du}}
\pgfpathlineto{\pgfpoint{-19.778474\du}{-2.763644\du}}
\pgfpathlineto{\pgfpoint{-19.778474\du}{-4.132907\du}}
\pgfusepath{fill}
\pgfsetbuttcap
\pgfsetmiterjoin
\pgfsetdash{}{0pt}
\definecolor{dialinecolor}{rgb}{0.286275, 0.286275, 0.211765}
\pgfsetstrokecolor{dialinecolor}
\pgfpathmoveto{\pgfpoint{-19.786900\du}{-4.158186\du}}
\pgfpathlineto{\pgfpoint{-19.786678\du}{-4.152199\du}}
\pgfpathlineto{\pgfpoint{-19.785570\du}{-4.145547\du}}
\pgfpathlineto{\pgfpoint{-19.783796\du}{-4.139338\du}}
\pgfpathlineto{\pgfpoint{-19.781356\du}{-4.132907\du}}
\pgfpathlineto{\pgfpoint{-19.777809\du}{-4.126920\du}}
\pgfpathlineto{\pgfpoint{-19.774039\du}{-4.120712\du}}
\pgfpathlineto{\pgfpoint{-19.769382\du}{-4.114725\du}}
\pgfpathlineto{\pgfpoint{-19.764060\du}{-4.108959\du}}
\pgfpathlineto{\pgfpoint{-19.758295\du}{-4.103416\du}}
\pgfpathlineto{\pgfpoint{-19.751643\du}{-4.097650\du}}
\pgfpathlineto{\pgfpoint{-19.744325\du}{-4.092107\du}}
\pgfpathlineto{\pgfpoint{-19.736343\du}{-4.087007\du}}
\pgfpathlineto{\pgfpoint{-19.727916\du}{-4.081907\du}}
\pgfpathlineto{\pgfpoint{-19.718825\du}{-4.076806\du}}
\pgfpathlineto{\pgfpoint{-19.709068\du}{-4.072150\du}}
\pgfpathlineto{\pgfpoint{-19.698868\du}{-4.067715\du}}
\pgfpathlineto{\pgfpoint{-19.688446\du}{-4.063502\du}}
\pgfpathlineto{\pgfpoint{-19.677359\du}{-4.059511\du}}
\pgfpathlineto{\pgfpoint{-19.665828\du}{-4.055741\du}}
\pgfpathlineto{\pgfpoint{-19.653854\du}{-4.052415\du}}
\pgfpathlineto{\pgfpoint{-19.641437\du}{-4.049089\du}}
\pgfpathlineto{\pgfpoint{-19.629019\du}{-4.045984\du}}
\pgfpathlineto{\pgfpoint{-19.615936\du}{-4.043323\du}}
\pgfpathlineto{\pgfpoint{-19.602632\du}{-4.041328\du}}
\pgfpathlineto{\pgfpoint{-19.589105\du}{-4.038888\du}}
\pgfpathlineto{\pgfpoint{-19.575135\du}{-4.037336\du}}
\pgfpathlineto{\pgfpoint{-19.561609\du}{-4.035784\du}}
\pgfpathlineto{\pgfpoint{-19.547418\du}{-4.034897\du}}
\pgfpathlineto{\pgfpoint{-19.533004\du}{-4.034010\du}}
\pgfpathlineto{\pgfpoint{-19.518813\du}{-4.033567\du}}
\pgfpathlineto{\pgfpoint{-19.504178\du}{-4.033345\du}}
\pgfpathlineto{\pgfpoint{-19.490208\du}{-4.033567\du}}
\pgfpathlineto{\pgfpoint{-19.475795\du}{-4.034010\du}}
\pgfpathlineto{\pgfpoint{-19.461825\du}{-4.034897\du}}
\pgfpathlineto{\pgfpoint{-19.447411\du}{-4.035784\du}}
\pgfpathlineto{\pgfpoint{-19.433885\du}{-4.037336\du}}
\pgfpathlineto{\pgfpoint{-19.420137\du}{-4.038888\du}}
\pgfpathlineto{\pgfpoint{-19.406389\du}{-4.041328\du}}
\pgfpathlineto{\pgfpoint{-19.393084\du}{-4.043323\du}}
\pgfpathlineto{\pgfpoint{-19.380002\du}{-4.045984\du}}
\pgfpathlineto{\pgfpoint{-19.367584\du}{-4.049089\du}}
\pgfpathlineto{\pgfpoint{-19.355166\du}{-4.052415\du}}
\pgfpathlineto{\pgfpoint{-19.343192\du}{-4.055741\du}}
\pgfpathlineto{\pgfpoint{-19.331662\du}{-4.059511\du}}
\pgfpathlineto{\pgfpoint{-19.320574\du}{-4.063502\du}}
\pgfpathlineto{\pgfpoint{-19.309487\du}{-4.067715\du}}
\pgfpathlineto{\pgfpoint{-19.299731\du}{-4.072150\du}}
\pgfpathlineto{\pgfpoint{-19.290196\du}{-4.076806\du}}
\pgfpathlineto{\pgfpoint{-19.281104\du}{-4.081907\du}}
\pgfpathlineto{\pgfpoint{-19.272456\du}{-4.087007\du}}
\pgfpathlineto{\pgfpoint{-19.264695\du}{-4.092107\du}}
\pgfpathlineto{\pgfpoint{-19.257378\du}{-4.097650\du}}
\pgfpathlineto{\pgfpoint{-19.250504\du}{-4.103416\du}}
\pgfpathlineto{\pgfpoint{-19.244738\du}{-4.108959\du}}
\pgfpathlineto{\pgfpoint{-19.239416\du}{-4.114725\du}}
\pgfpathlineto{\pgfpoint{-19.234982\du}{-4.120712\du}}
\pgfpathlineto{\pgfpoint{-19.231212\du}{-4.126920\du}}
\pgfpathlineto{\pgfpoint{-19.227664\du}{-4.132907\du}}
\pgfpathlineto{\pgfpoint{-19.225225\du}{-4.139338\du}}
\pgfpathlineto{\pgfpoint{-19.223673\du}{-4.145547\du}}
\pgfpathlineto{\pgfpoint{-19.222121\du}{-4.152199\du}}
\pgfpathlineto{\pgfpoint{-19.221899\du}{-4.158186\du}}
\pgfusepath{stroke}
\pgfsetbuttcap
\pgfsetmiterjoin
\pgfsetdash{}{0pt}
\definecolor{dialinecolor}{rgb}{0.286275, 0.286275, 0.211765}
\pgfsetstrokecolor{dialinecolor}
\pgfpathmoveto{\pgfpoint{-19.786900\du}{-2.785375\du}}
\pgfpathlineto{\pgfpoint{-19.786678\du}{-2.779388\du}}
\pgfpathlineto{\pgfpoint{-19.785791\du}{-2.773401\du}}
\pgfpathlineto{\pgfpoint{-19.783796\du}{-2.767192\du}}
\pgfpathlineto{\pgfpoint{-19.781356\du}{-2.761427\du}}
\pgfpathlineto{\pgfpoint{-19.778474\du}{-2.755440\du}}
\pgfpathlineto{\pgfpoint{-19.774704\du}{-2.749453\du}}
\pgfpathlineto{\pgfpoint{-19.770269\du}{-2.743909\du}}
\pgfpathlineto{\pgfpoint{-19.764726\du}{-2.738144\du}}
\pgfpathlineto{\pgfpoint{-19.759626\du}{-2.732822\du}}
\pgfpathlineto{\pgfpoint{-19.752973\du}{-2.727057\du}}
\pgfpathlineto{\pgfpoint{-19.746099\du}{-2.721957\du}}
\pgfpathlineto{\pgfpoint{-19.738117\du}{-2.717078\du}}
\pgfpathlineto{\pgfpoint{-19.729912\du}{-2.711978\du}}
\pgfpathlineto{\pgfpoint{-19.721042\du}{-2.707322\du}}
\pgfpathlineto{\pgfpoint{-19.711507\du}{-2.702887\du}}
\pgfpathlineto{\pgfpoint{-19.701751\du}{-2.698452\du}}
\pgfpathlineto{\pgfpoint{-19.691772\du}{-2.694239\du}}
\pgfpathlineto{\pgfpoint{-19.680685\du}{-2.690247\du}}
\pgfpathlineto{\pgfpoint{-19.669154\du}{-2.686700\du}}
\pgfpathlineto{\pgfpoint{-19.657624\du}{-2.683152\du}}
\pgfpathlineto{\pgfpoint{-19.645206\du}{-2.680047\du}}
\pgfpathlineto{\pgfpoint{-19.633010\du}{-2.677165\du}}
\pgfpathlineto{\pgfpoint{-19.620371\du}{-2.674725\du}}
\pgfpathlineto{\pgfpoint{-19.607066\du}{-2.672286\du}}
\pgfpathlineto{\pgfpoint{-19.593762\du}{-2.670291\du}}
\pgfpathlineto{\pgfpoint{-19.580236\du}{-2.668517\du}}
\pgfpathlineto{\pgfpoint{-19.566487\du}{-2.666743\du}}
\pgfpathlineto{\pgfpoint{-19.552739\du}{-2.665856\du}}
\pgfpathlineto{\pgfpoint{-19.538548\du}{-2.664969\du}}
\pgfpathlineto{\pgfpoint{-19.524356\du}{-2.664303\du}}
\pgfpathlineto{\pgfpoint{-19.509721\du}{-2.664082\du}}
\pgfpathlineto{\pgfpoint{-19.496195\du}{-2.664082\du}}
\pgfpathlineto{\pgfpoint{-19.482003\du}{-2.664303\du}}
\pgfpathlineto{\pgfpoint{-19.467812\du}{-2.664969\du}}
\pgfpathlineto{\pgfpoint{-19.453620\du}{-2.665856\du}}
\pgfpathlineto{\pgfpoint{-19.439650\du}{-2.667186\du}}
\pgfpathlineto{\pgfpoint{-19.426346\du}{-2.668738\du}}
\pgfpathlineto{\pgfpoint{-19.412820\du}{-2.670512\du}}
\pgfpathlineto{\pgfpoint{-19.399293\du}{-2.672730\du}}
\pgfpathlineto{\pgfpoint{-19.386210\du}{-2.674947\du}}
\pgfpathlineto{\pgfpoint{-19.373571\du}{-2.677830\du}}
\pgfpathlineto{\pgfpoint{-19.361375\du}{-2.680712\du}}
\pgfpathlineto{\pgfpoint{-19.348958\du}{-2.683817\du}}
\pgfpathlineto{\pgfpoint{-19.337649\du}{-2.687365\du}}
\pgfpathlineto{\pgfpoint{-19.326118\du}{-2.691134\du}}
\pgfpathlineto{\pgfpoint{-19.315474\du}{-2.694904\du}}
\pgfpathlineto{\pgfpoint{-19.305274\du}{-2.699117\du}}
\pgfpathlineto{\pgfpoint{-19.295517\du}{-2.703330\du}}
\pgfpathlineto{\pgfpoint{-19.286204\du}{-2.708209\du}}
\pgfpathlineto{\pgfpoint{-19.277556\du}{-2.712865\du}}
\pgfpathlineto{\pgfpoint{-19.269352\du}{-2.717965\du}}
\pgfpathlineto{\pgfpoint{-19.261591\du}{-2.723065\du}}
\pgfpathlineto{\pgfpoint{-19.254495\du}{-2.728387\du}}
\pgfpathlineto{\pgfpoint{-19.248286\du}{-2.733709\du}}
\pgfpathlineto{\pgfpoint{-19.242521\du}{-2.739474\du}}
\pgfpathlineto{\pgfpoint{-19.237864\du}{-2.745018\du}}
\pgfpathlineto{\pgfpoint{-19.233651\du}{-2.750783\du}}
\pgfpathlineto{\pgfpoint{-19.229882\du}{-2.756549\du}}
\pgfpathlineto{\pgfpoint{-19.226777\du}{-2.762536\du}}
\pgfpathlineto{\pgfpoint{-19.224560\du}{-2.768523\du}}
\pgfpathlineto{\pgfpoint{-19.222786\du}{-2.774510\du}}
\pgfpathlineto{\pgfpoint{-19.222121\du}{-2.780497\du}}
\pgfpathlineto{\pgfpoint{-19.221899\du}{-2.786706\du}}
\pgfusepath{stroke}
\pgfsetlinewidth{0.100000\du}
\pgfsetdash{}{0pt}
\pgfsetdash{}{0pt}
\pgfsetbuttcap
\pgfsetmiterjoin
\pgfsetlinewidth{0.001000\du}
\pgfsetbuttcap
\pgfsetmiterjoin
\pgfsetdash{}{0pt}
\definecolor{dialinecolor}{rgb}{0.000000, 0.470588, 0.666667}
\pgfsetstrokecolor{dialinecolor}
\pgfpathmoveto{\pgfpoint{-5.435683\du}{7.473354\du}}
\pgfpathlineto{\pgfpoint{-5.429115\du}{7.464397\du}}
\pgfpathlineto{\pgfpoint{-5.422845\du}{7.454844\du}}
\pgfpathlineto{\pgfpoint{-5.417471\du}{7.444693\du}}
\pgfpathlineto{\pgfpoint{-5.412396\du}{7.435438\du}}
\pgfpathlineto{\pgfpoint{-5.408515\du}{7.425287\du}}
\pgfpathlineto{\pgfpoint{-5.405231\du}{7.415136\du}}
\pgfpathlineto{\pgfpoint{-5.402245\du}{7.404986\du}}
\pgfpathlineto{\pgfpoint{-5.400454\du}{7.394536\du}}
\pgfpathlineto{\pgfpoint{-5.399260\du}{7.384385\du}}
\pgfpathlineto{\pgfpoint{-5.398364\du}{7.373936\du}}
\pgfpathlineto{\pgfpoint{-5.398364\du}{7.363487\du}}
\pgfpathlineto{\pgfpoint{-5.399260\du}{7.352739\du}}
\pgfpathlineto{\pgfpoint{-5.400454\du}{7.342588\du}}
\pgfpathlineto{\pgfpoint{-5.402245\du}{7.332139\du}}
\pgfpathlineto{\pgfpoint{-5.404932\du}{7.322287\du}}
\pgfpathlineto{\pgfpoint{-5.408515\du}{7.311837\du}}
\pgfpathlineto{\pgfpoint{-5.412396\du}{7.301985\du}}
\pgfpathlineto{\pgfpoint{-5.417471\du}{7.291834\du}}
\pgfpathlineto{\pgfpoint{-5.422845\du}{7.282281\du}}
\pgfpathlineto{\pgfpoint{-5.429115\du}{7.273324\du}}
\pgfpathlineto{\pgfpoint{-5.435683\du}{7.263770\du}}
\pgfpathlineto{\pgfpoint{-5.443147\du}{7.255112\du}}
\pgfpathlineto{\pgfpoint{-5.450611\du}{7.246454\du}}
\pgfpathlineto{\pgfpoint{-5.458970\du}{7.238095\du}}
\pgfpathlineto{\pgfpoint{-5.467628\du}{7.230034\du}}
\pgfpathlineto{\pgfpoint{-5.476883\du}{7.222272\du}}
\pgfpathlineto{\pgfpoint{-5.486437\du}{7.215106\du}}
\pgfpathlineto{\pgfpoint{-5.497185\du}{7.208538\du}}
\pgfpathlineto{\pgfpoint{-5.507634\du}{7.201970\du}}
\pgfpathlineto{\pgfpoint{-5.518979\du}{7.195701\du}}
\pgfpathlineto{\pgfpoint{-5.530921\du}{7.190028\du}}
\pgfpathlineto{\pgfpoint{-5.542863\du}{7.184654\du}}
\pgfpathlineto{\pgfpoint{-5.555104\du}{7.180474\du}}
\pgfpathlineto{\pgfpoint{-5.567643\du}{7.176295\du}}
\pgfpathlineto{\pgfpoint{-5.580779\du}{7.172115\du}}
\pgfpathlineto{\pgfpoint{-5.593617\du}{7.169129\du}}
\pgfpathlineto{\pgfpoint{-5.606753\du}{7.166442\du}}
\pgfpathlineto{\pgfpoint{-5.620487\du}{7.164054\du}}
\pgfpathlineto{\pgfpoint{-5.633623\du}{7.162561\du}}
\pgfpathlineto{\pgfpoint{-5.647655\du}{7.161367\du}}
\pgfpathlineto{\pgfpoint{-5.661388\du}{7.160471\du}}
\pgfpathlineto{\pgfpoint{-5.675420\du}{7.160471\du}}
\pgfpathlineto{\pgfpoint{-5.688855\du}{7.161069\du}}
\pgfpathlineto{\pgfpoint{-5.702589\du}{7.161666\du}}
\pgfpathlineto{\pgfpoint{-5.716621\du}{7.163158\du}}
\pgfpathlineto{\pgfpoint{-5.729757\du}{7.165547\du}}
\pgfpathlineto{\pgfpoint{-5.743490\du}{7.167935\du}}
\pgfpathlineto{\pgfpoint{-5.756029\du}{7.170921\du}}
\pgfpathlineto{\pgfpoint{-5.769166\du}{7.173906\du}}
\pgfpathlineto{\pgfpoint{-5.781705\du}{7.178086\du}}
\pgfpathlineto{\pgfpoint{-5.794543\du}{7.182863\du}}
\pgfpathlineto{\pgfpoint{-5.806783\du}{7.187341\du}}
\pgfpathlineto{\pgfpoint{-5.818128\du}{7.192715\du}}
\pgfpathlineto{\pgfpoint{-5.829772\du}{7.198985\du}}
\pgfpathlineto{\pgfpoint{-5.840818\du}{7.204956\du}}
\pgfpathlineto{\pgfpoint{-5.851268\du}{7.211822\du}}
\pgfpathlineto{\pgfpoint{-5.861418\du}{7.218391\du}}
\pgfpathlineto{\pgfpoint{-5.870972\du}{7.226153\du}}
\pgfpathlineto{\pgfpoint{-5.880526\du}{7.234214\du}}
\pgfpathlineto{\pgfpoint{-5.888885\du}{7.242275\du}}
\pgfpathlineto{\pgfpoint{-5.896946\du}{7.250634\du}}
\pgfpathlineto{\pgfpoint{-5.904708\du}{7.259591\du}}
\pgfusepath{stroke}
\pgfsetlinewidth{0.080000\du}
\pgfsetbuttcap
\pgfsetmiterjoin
\pgfsetdash{}{0pt}
\definecolor{dialinecolor}{rgb}{0.666667, 0.901961, 1.000000}
\pgfsetstrokecolor{dialinecolor}
\pgfpathmoveto{\pgfpoint{-5.444938\du}{7.465890\du}}
\pgfpathlineto{\pgfpoint{-5.439564\du}{7.456934\du}}
\pgfpathlineto{\pgfpoint{-5.433295\du}{7.448276\du}}
\pgfpathlineto{\pgfpoint{-5.428219\du}{7.439020\du}}
\pgfpathlineto{\pgfpoint{-5.424040\du}{7.429765\du}}
\pgfpathlineto{\pgfpoint{-5.420158\du}{7.420212\du}}
\pgfpathlineto{\pgfpoint{-5.417173\du}{7.410658\du}}
\pgfpathlineto{\pgfpoint{-5.414486\du}{7.400806\du}}
\pgfpathlineto{\pgfpoint{-5.412396\du}{7.391551\du}}
\pgfpathlineto{\pgfpoint{-5.411500\du}{7.381698\du}}
\pgfpathlineto{\pgfpoint{-5.410903\du}{7.371548\du}}
\pgfpathlineto{\pgfpoint{-5.410903\du}{7.361695\du}}
\pgfpathlineto{\pgfpoint{-5.411799\du}{7.351843\du}}
\pgfpathlineto{\pgfpoint{-5.413292\du}{7.341693\du}}
\pgfpathlineto{\pgfpoint{-5.415083\du}{7.332437\du}}
\pgfpathlineto{\pgfpoint{-5.418069\du}{7.322585\du}}
\pgfpathlineto{\pgfpoint{-5.421651\du}{7.313032\du}}
\pgfpathlineto{\pgfpoint{-5.425532\du}{7.303776\du}}
\pgfpathlineto{\pgfpoint{-5.430011\du}{7.294521\du}}
\pgfpathlineto{\pgfpoint{-5.435086\du}{7.285863\du}}
\pgfpathlineto{\pgfpoint{-5.440758\du}{7.276608\du}}
\pgfpathlineto{\pgfpoint{-5.447924\du}{7.268249\du}}
\pgfpathlineto{\pgfpoint{-5.454492\du}{7.259591\du}}
\pgfpathlineto{\pgfpoint{-5.462254\du}{7.251530\du}}
\pgfpathlineto{\pgfpoint{-5.470315\du}{7.243767\du}}
\pgfpathlineto{\pgfpoint{-5.478675\du}{7.236901\du}}
\pgfpathlineto{\pgfpoint{-5.487333\du}{7.229437\du}}
\pgfpathlineto{\pgfpoint{-5.497185\du}{7.222570\du}}
\pgfpathlineto{\pgfpoint{-5.507634\du}{7.215704\du}}
\pgfpathlineto{\pgfpoint{-5.517785\du}{7.210031\du}}
\pgfpathlineto{\pgfpoint{-5.528234\du}{7.204657\du}}
\pgfpathlineto{\pgfpoint{-5.540176\du}{7.199283\du}}
\pgfpathlineto{\pgfpoint{-5.550924\du}{7.194208\du}}
\pgfpathlineto{\pgfpoint{-5.563463\du}{7.190028\du}}
\pgfpathlineto{\pgfpoint{-5.574808\du}{7.186147\du}}
\pgfpathlineto{\pgfpoint{-5.587348\du}{7.182863\du}}
\pgfpathlineto{\pgfpoint{-5.599887\du}{7.179877\du}}
\pgfpathlineto{\pgfpoint{-5.612426\du}{7.177190\du}}
\pgfpathlineto{\pgfpoint{-5.625562\du}{7.175399\du}}
\pgfpathlineto{\pgfpoint{-5.638101\du}{7.173906\du}}
\pgfpathlineto{\pgfpoint{-5.651238\du}{7.172712\du}}
\pgfpathlineto{\pgfpoint{-5.665270\du}{7.172115\du}}
\pgfpathlineto{\pgfpoint{-5.678406\du}{7.172115\du}}
\pgfpathlineto{\pgfpoint{-5.691244\du}{7.172712\du}}
\pgfpathlineto{\pgfpoint{-5.704380\du}{7.173608\du}}
\pgfpathlineto{\pgfpoint{-5.716919\du}{7.175399\du}}
\pgfpathlineto{\pgfpoint{-5.730055\du}{7.177190\du}}
\pgfpathlineto{\pgfpoint{-5.742595\du}{7.179280\du}}
\pgfpathlineto{\pgfpoint{-5.755432\du}{7.182863\du}}
\pgfpathlineto{\pgfpoint{-5.767673\du}{7.185848\du}}
\pgfpathlineto{\pgfpoint{-5.779615\du}{7.189730\du}}
\pgfpathlineto{\pgfpoint{-5.791856\du}{7.193909\du}}
\pgfpathlineto{\pgfpoint{-5.803201\du}{7.198985\du}}
\pgfpathlineto{\pgfpoint{-5.814247\du}{7.203761\du}}
\pgfpathlineto{\pgfpoint{-5.824995\du}{7.209733\du}}
\pgfpathlineto{\pgfpoint{-5.835444\du}{7.215405\du}}
\pgfpathlineto{\pgfpoint{-5.845894\du}{7.221973\du}}
\pgfpathlineto{\pgfpoint{-5.855447\du}{7.228541\du}}
\pgfpathlineto{\pgfpoint{-5.864105\du}{7.235707\du}}
\pgfpathlineto{\pgfpoint{-5.873659\du}{7.242872\du}}
\pgfpathlineto{\pgfpoint{-5.881421\du}{7.250634\du}}
\pgfpathlineto{\pgfpoint{-5.888885\du}{7.258695\du}}
\pgfpathlineto{\pgfpoint{-5.896050\du}{7.267055\du}}
\pgfusepath{stroke}
\pgfsetbuttcap
\pgfsetmiterjoin
\pgfsetdash{}{0pt}
\definecolor{dialinecolor}{rgb}{0.000000, 0.705882, 1.000000}
\pgfsetfillcolor{dialinecolor}
\pgfpathmoveto{\pgfpoint{-7.361941\du}{7.148828\du}}
\pgfpathlineto{\pgfpoint{-7.159224\du}{6.958053\du}}
\pgfpathlineto{\pgfpoint{-5.647655\du}{6.958053\du}}
\pgfpathlineto{\pgfpoint{-5.850372\du}{7.148828\du}}
\pgfpathlineto{\pgfpoint{-7.361941\du}{7.148828\du}}
\pgfusepath{fill}
\pgfsetlinewidth{0.001000\du}
\pgfsetbuttcap
\pgfsetmiterjoin
\pgfsetdash{}{0pt}
\definecolor{dialinecolor}{rgb}{0.666667, 0.901961, 1.000000}
\pgfsetstrokecolor{dialinecolor}
\pgfpathmoveto{\pgfpoint{-7.361941\du}{7.148828\du}}
\pgfpathlineto{\pgfpoint{-7.159224\du}{6.958053\du}}
\pgfpathlineto{\pgfpoint{-5.647655\du}{6.958053\du}}
\pgfpathlineto{\pgfpoint{-5.850372\du}{7.148828\du}}
\pgfpathlineto{\pgfpoint{-7.361941\du}{7.148828\du}}
\pgfusepath{stroke}
\pgfsetbuttcap
\pgfsetmiterjoin
\pgfsetdash{}{0pt}
\definecolor{dialinecolor}{rgb}{0.000000, 0.588235, 0.831373}
\pgfsetfillcolor{dialinecolor}
\pgfpathmoveto{\pgfpoint{-7.361941\du}{7.148828\du}}
\pgfpathlineto{\pgfpoint{-5.850372\du}{7.148828\du}}
\pgfpathlineto{\pgfpoint{-5.850372\du}{7.434841\du}}
\pgfpathlineto{\pgfpoint{-7.361941\du}{7.434841\du}}
\pgfpathlineto{\pgfpoint{-7.361941\du}{7.148828\du}}
\pgfusepath{fill}
\pgfsetbuttcap
\pgfsetmiterjoin
\pgfsetdash{}{0pt}
\definecolor{dialinecolor}{rgb}{0.666667, 0.901961, 1.000000}
\pgfsetstrokecolor{dialinecolor}
\pgfpathmoveto{\pgfpoint{-7.361941\du}{7.148828\du}}
\pgfpathlineto{\pgfpoint{-5.850372\du}{7.148828\du}}
\pgfpathlineto{\pgfpoint{-5.850372\du}{7.434841\du}}
\pgfpathlineto{\pgfpoint{-7.361941\du}{7.434841\du}}
\pgfpathlineto{\pgfpoint{-7.361941\du}{7.148828\du}}
\pgfusepath{stroke}
\pgfsetbuttcap
\pgfsetmiterjoin
\pgfsetdash{}{0pt}
\definecolor{dialinecolor}{rgb}{0.000000, 0.352941, 0.501961}
\pgfsetfillcolor{dialinecolor}
\pgfpathmoveto{\pgfpoint{-5.850372\du}{7.434841\du}}
\pgfpathlineto{\pgfpoint{-5.647655\du}{7.232124\du}}
\pgfpathlineto{\pgfpoint{-5.647655\du}{6.958053\du}}
\pgfpathlineto{\pgfpoint{-5.850372\du}{7.148828\du}}
\pgfpathlineto{\pgfpoint{-5.850372\du}{7.434841\du}}
\pgfusepath{fill}
\pgfsetbuttcap
\pgfsetmiterjoin
\pgfsetdash{}{0pt}
\definecolor{dialinecolor}{rgb}{0.666667, 0.901961, 1.000000}
\pgfsetstrokecolor{dialinecolor}
\pgfpathmoveto{\pgfpoint{-5.850372\du}{7.434841\du}}
\pgfpathlineto{\pgfpoint{-5.647655\du}{7.232124\du}}
\pgfpathlineto{\pgfpoint{-5.647655\du}{6.958053\du}}
\pgfpathlineto{\pgfpoint{-5.850372\du}{7.148828\du}}
\pgfpathlineto{\pgfpoint{-5.850372\du}{7.434841\du}}
\pgfusepath{stroke}
\pgfsetbuttcap
\pgfsetmiterjoin
\pgfsetdash{}{0pt}
\definecolor{dialinecolor}{rgb}{0.000000, 0.000000, 0.000000}
\pgfsetfillcolor{dialinecolor}
\pgfpathmoveto{\pgfpoint{-7.313874\du}{7.112703\du}}
\pgfpathlineto{\pgfpoint{-7.159224\du}{6.958053\du}}
\pgfpathlineto{\pgfpoint{-5.683780\du}{6.958053\du}}
\pgfpathlineto{\pgfpoint{-5.838131\du}{7.112703\du}}
\pgfpathlineto{\pgfpoint{-7.313874\du}{7.112703\du}}
\pgfusepath{fill}
\pgfsetbuttcap
\pgfsetmiterjoin
\pgfsetdash{}{0pt}
\definecolor{dialinecolor}{rgb}{0.000000, 0.000000, 0.000000}
\pgfsetstrokecolor{dialinecolor}
\pgfpathmoveto{\pgfpoint{-7.313874\du}{7.112703\du}}
\pgfpathlineto{\pgfpoint{-7.159224\du}{6.958053\du}}
\pgfpathlineto{\pgfpoint{-5.683780\du}{6.958053\du}}
\pgfpathlineto{\pgfpoint{-5.838131\du}{7.112703\du}}
\pgfpathlineto{\pgfpoint{-7.313874\du}{7.112703\du}}
\pgfusepath{stroke}
\pgfsetbuttcap
\pgfsetmiterjoin
\pgfsetdash{}{0pt}
\definecolor{dialinecolor}{rgb}{0.000000, 0.705882, 1.000000}
\pgfsetfillcolor{dialinecolor}
\pgfpathmoveto{\pgfpoint{-7.361941\du}{5.886550\du}}
\pgfpathlineto{\pgfpoint{-7.194752\du}{5.731900\du}}
\pgfpathlineto{\pgfpoint{-5.683780\du}{5.731900\du}}
\pgfpathlineto{\pgfpoint{-5.850372\du}{5.886550\du}}
\pgfpathlineto{\pgfpoint{-7.361941\du}{5.886550\du}}
\pgfusepath{fill}
\pgfsetbuttcap
\pgfsetmiterjoin
\pgfsetdash{}{0pt}
\definecolor{dialinecolor}{rgb}{0.666667, 0.901961, 1.000000}
\pgfsetstrokecolor{dialinecolor}
\pgfpathmoveto{\pgfpoint{-7.361941\du}{5.886550\du}}
\pgfpathlineto{\pgfpoint{-7.194752\du}{5.731900\du}}
\pgfpathlineto{\pgfpoint{-5.683780\du}{5.731900\du}}
\pgfpathlineto{\pgfpoint{-5.850372\du}{5.886550\du}}
\pgfpathlineto{\pgfpoint{-7.361941\du}{5.886550\du}}
\pgfusepath{stroke}
\pgfsetbuttcap
\pgfsetmiterjoin
\pgfsetdash{}{0pt}
\definecolor{dialinecolor}{rgb}{0.000000, 0.588235, 0.831373}
\pgfsetfillcolor{dialinecolor}
\pgfpathmoveto{\pgfpoint{-7.361941\du}{5.886550\du}}
\pgfpathlineto{\pgfpoint{-5.838131\du}{5.886550\du}}
\pgfpathlineto{\pgfpoint{-5.838131\du}{7.088819\du}}
\pgfpathlineto{\pgfpoint{-7.361941\du}{7.088819\du}}
\pgfpathlineto{\pgfpoint{-7.361941\du}{5.886550\du}}
\pgfusepath{fill}
\pgfsetbuttcap
\pgfsetmiterjoin
\pgfsetdash{}{0pt}
\definecolor{dialinecolor}{rgb}{0.666667, 0.901961, 1.000000}
\pgfsetstrokecolor{dialinecolor}
\pgfpathmoveto{\pgfpoint{-7.361941\du}{5.886550\du}}
\pgfpathlineto{\pgfpoint{-5.838430\du}{5.886550\du}}
\pgfpathlineto{\pgfpoint{-5.838430\du}{7.088520\du}}
\pgfpathlineto{\pgfpoint{-7.361941\du}{7.088520\du}}
\pgfpathlineto{\pgfpoint{-7.361941\du}{5.886550\du}}
\pgfusepath{stroke}
\pgfsetbuttcap
\pgfsetmiterjoin
\pgfsetdash{}{0pt}
\definecolor{dialinecolor}{rgb}{1.000000, 1.000000, 1.000000}
\pgfsetfillcolor{dialinecolor}
\pgfpathmoveto{\pgfpoint{-7.230876\du}{6.041498\du}}
\pgfpathlineto{\pgfpoint{-5.969196\du}{6.041498\du}}
\pgfpathlineto{\pgfpoint{-5.969196\du}{6.970294\du}}
\pgfpathlineto{\pgfpoint{-7.230876\du}{6.970294\du}}
\pgfpathlineto{\pgfpoint{-7.230876\du}{6.041498\du}}
\pgfusepath{fill}
\pgfsetbuttcap
\pgfsetmiterjoin
\pgfsetdash{}{0pt}
\definecolor{dialinecolor}{rgb}{0.000000, 0.235294, 0.333333}
\pgfsetstrokecolor{dialinecolor}
\pgfpathmoveto{\pgfpoint{-7.230876\du}{6.041498\du}}
\pgfpathlineto{\pgfpoint{-5.969494\du}{6.041498\du}}
\pgfpathlineto{\pgfpoint{-5.969494\du}{6.969995\du}}
\pgfpathlineto{\pgfpoint{-7.230876\du}{6.969995\du}}
\pgfpathlineto{\pgfpoint{-7.230876\du}{6.041498\du}}
\pgfusepath{stroke}
\pgfsetbuttcap
\pgfsetmiterjoin
\pgfsetdash{}{0pt}
\definecolor{dialinecolor}{rgb}{0.000000, 0.352941, 0.501961}
\pgfsetfillcolor{dialinecolor}
\pgfpathmoveto{\pgfpoint{-5.850372\du}{7.077474\du}}
\pgfpathlineto{\pgfpoint{-5.683780\du}{6.910583\du}}
\pgfpathlineto{\pgfpoint{-5.683780\du}{5.731900\du}}
\pgfpathlineto{\pgfpoint{-5.850372\du}{5.886550\du}}
\pgfpathlineto{\pgfpoint{-5.850372\du}{7.077474\du}}
\pgfusepath{fill}
\pgfsetbuttcap
\pgfsetmiterjoin
\pgfsetdash{}{0pt}
\definecolor{dialinecolor}{rgb}{0.666667, 0.901961, 1.000000}
\pgfsetstrokecolor{dialinecolor}
\pgfpathmoveto{\pgfpoint{-5.850372\du}{7.077474\du}}
\pgfpathlineto{\pgfpoint{-5.683780\du}{6.910583\du}}
\pgfpathlineto{\pgfpoint{-5.683780\du}{5.731900\du}}
\pgfpathlineto{\pgfpoint{-5.850372\du}{5.886550\du}}
\pgfpathlineto{\pgfpoint{-5.850372\du}{7.077474\du}}
\pgfusepath{stroke}
\pgfsetbuttcap
\pgfsetmiterjoin
\pgfsetdash{}{0pt}
\definecolor{dialinecolor}{rgb}{0.000000, 0.705882, 1.000000}
\pgfsetfillcolor{dialinecolor}
\pgfpathmoveto{\pgfpoint{-7.659000\du}{7.672787\du}}
\pgfpathlineto{\pgfpoint{-7.421353\du}{7.374832\du}}
\pgfpathlineto{\pgfpoint{-5.778719\du}{7.374832\du}}
\pgfpathlineto{\pgfpoint{-6.016665\du}{7.672787\du}}
\pgfpathlineto{\pgfpoint{-7.659000\du}{7.672787\du}}
\pgfusepath{fill}
\pgfsetbuttcap
\pgfsetmiterjoin
\pgfsetdash{}{0pt}
\definecolor{dialinecolor}{rgb}{0.666667, 0.901961, 1.000000}
\pgfsetstrokecolor{dialinecolor}
\pgfpathmoveto{\pgfpoint{-7.659000\du}{7.672787\du}}
\pgfpathlineto{\pgfpoint{-7.421353\du}{7.374832\du}}
\pgfpathlineto{\pgfpoint{-5.778719\du}{7.374832\du}}
\pgfpathlineto{\pgfpoint{-6.016665\du}{7.672787\du}}
\pgfpathlineto{\pgfpoint{-7.659000\du}{7.672787\du}}
\pgfusepath{stroke}
\pgfsetbuttcap
\pgfsetmiterjoin
\pgfsetdash{}{0pt}
\definecolor{dialinecolor}{rgb}{0.000000, 0.352941, 0.501961}
\pgfsetfillcolor{dialinecolor}
\pgfpathmoveto{\pgfpoint{-6.016665\du}{7.731900\du}}
\pgfpathlineto{\pgfpoint{-5.778719\du}{7.482012\du}}
\pgfpathlineto{\pgfpoint{-5.778719\du}{7.374832\du}}
\pgfpathlineto{\pgfpoint{-6.016665\du}{7.672787\du}}
\pgfpathlineto{\pgfpoint{-6.016665\du}{7.731900\du}}
\pgfusepath{fill}
\pgfsetbuttcap
\pgfsetmiterjoin
\pgfsetdash{}{0pt}
\definecolor{dialinecolor}{rgb}{0.666667, 0.901961, 1.000000}
\pgfsetstrokecolor{dialinecolor}
\pgfpathmoveto{\pgfpoint{-6.016665\du}{7.731900\du}}
\pgfpathlineto{\pgfpoint{-5.778719\du}{7.482012\du}}
\pgfpathlineto{\pgfpoint{-5.778719\du}{7.374832\du}}
\pgfpathlineto{\pgfpoint{-6.016665\du}{7.672787\du}}
\pgfpathlineto{\pgfpoint{-6.016665\du}{7.731900\du}}
\pgfusepath{stroke}
\pgfsetbuttcap
\pgfsetmiterjoin
\pgfsetdash{}{0pt}
\definecolor{dialinecolor}{rgb}{0.000000, 0.588235, 0.831373}
\pgfsetfillcolor{dialinecolor}
\pgfpathmoveto{\pgfpoint{-7.659000\du}{7.672787\du}}
\pgfpathlineto{\pgfpoint{-6.016665\du}{7.672787\du}}
\pgfpathlineto{\pgfpoint{-6.016665\du}{7.731900\du}}
\pgfpathlineto{\pgfpoint{-7.659000\du}{7.731900\du}}
\pgfpathlineto{\pgfpoint{-7.659000\du}{7.672787\du}}
\pgfusepath{fill}
\pgfsetbuttcap
\pgfsetmiterjoin
\pgfsetdash{}{0pt}
\definecolor{dialinecolor}{rgb}{0.666667, 0.901961, 1.000000}
\pgfsetstrokecolor{dialinecolor}
\pgfpathmoveto{\pgfpoint{-7.659000\du}{7.672787\du}}
\pgfpathlineto{\pgfpoint{-6.016665\du}{7.672787\du}}
\pgfpathlineto{\pgfpoint{-6.016665\du}{7.731900\du}}
\pgfpathlineto{\pgfpoint{-7.659000\du}{7.731900\du}}
\pgfpathlineto{\pgfpoint{-7.659000\du}{7.672787\du}}
\pgfusepath{stroke}
\pgfsetbuttcap
\pgfsetmiterjoin
\pgfsetdash{}{0pt}
\definecolor{dialinecolor}{rgb}{0.000000, 0.705882, 1.000000}
\pgfsetfillcolor{dialinecolor}
\pgfpathmoveto{\pgfpoint{-5.683780\du}{7.636662\du}}
\pgfpathlineto{\pgfpoint{-5.564658\du}{7.458426\du}}
\pgfpathlineto{\pgfpoint{-5.349999\du}{7.458426\du}}
\pgfpathlineto{\pgfpoint{-5.468822\du}{7.636662\du}}
\pgfpathlineto{\pgfpoint{-5.683780\du}{7.636662\du}}
\pgfusepath{fill}
\pgfsetbuttcap
\pgfsetmiterjoin
\pgfsetdash{}{0pt}
\definecolor{dialinecolor}{rgb}{0.666667, 0.901961, 1.000000}
\pgfsetstrokecolor{dialinecolor}
\pgfpathmoveto{\pgfpoint{-5.683780\du}{7.636662\du}}
\pgfpathlineto{\pgfpoint{-5.564658\du}{7.458426\du}}
\pgfpathlineto{\pgfpoint{-5.349999\du}{7.458426\du}}
\pgfpathlineto{\pgfpoint{-5.468822\du}{7.636662\du}}
\pgfpathlineto{\pgfpoint{-5.683780\du}{7.636662\du}}
\pgfusepath{stroke}
\pgfsetbuttcap
\pgfsetmiterjoin
\pgfsetdash{}{0pt}
\definecolor{dialinecolor}{rgb}{0.000000, 0.352941, 0.501961}
\pgfsetfillcolor{dialinecolor}
\pgfpathmoveto{\pgfpoint{-5.468822\du}{7.696372\du}}
\pgfpathlineto{\pgfpoint{-5.349999\du}{7.565308\du}}
\pgfpathlineto{\pgfpoint{-5.349999\du}{7.458426\du}}
\pgfpathlineto{\pgfpoint{-5.468822\du}{7.636662\du}}
\pgfpathlineto{\pgfpoint{-5.468822\du}{7.696372\du}}
\pgfusepath{fill}
\pgfsetbuttcap
\pgfsetmiterjoin
\pgfsetdash{}{0pt}
\definecolor{dialinecolor}{rgb}{0.666667, 0.901961, 1.000000}
\pgfsetstrokecolor{dialinecolor}
\pgfpathmoveto{\pgfpoint{-5.468822\du}{7.696372\du}}
\pgfpathlineto{\pgfpoint{-5.349999\du}{7.565308\du}}
\pgfpathlineto{\pgfpoint{-5.349999\du}{7.458426\du}}
\pgfpathlineto{\pgfpoint{-5.468822\du}{7.636662\du}}
\pgfpathlineto{\pgfpoint{-5.468822\du}{7.696372\du}}
\pgfusepath{stroke}
\pgfsetbuttcap
\pgfsetmiterjoin
\pgfsetdash{}{0pt}
\definecolor{dialinecolor}{rgb}{0.000000, 0.588235, 0.831373}
\pgfsetfillcolor{dialinecolor}
\pgfpathmoveto{\pgfpoint{-5.683780\du}{7.636662\du}}
\pgfpathlineto{\pgfpoint{-5.468822\du}{7.636662\du}}
\pgfpathlineto{\pgfpoint{-5.468822\du}{7.696372\du}}
\pgfpathlineto{\pgfpoint{-5.683780\du}{7.696372\du}}
\pgfpathlineto{\pgfpoint{-5.683780\du}{7.636662\du}}
\pgfusepath{fill}
\pgfsetbuttcap
\pgfsetmiterjoin
\pgfsetdash{}{0pt}
\definecolor{dialinecolor}{rgb}{0.666667, 0.901961, 1.000000}
\pgfsetstrokecolor{dialinecolor}
\pgfpathmoveto{\pgfpoint{-5.683780\du}{7.636662\du}}
\pgfpathlineto{\pgfpoint{-5.468822\du}{7.636662\du}}
\pgfpathlineto{\pgfpoint{-5.468822\du}{7.696372\du}}
\pgfpathlineto{\pgfpoint{-5.683780\du}{7.696372\du}}
\pgfpathlineto{\pgfpoint{-5.683780\du}{7.636662\du}}
\pgfusepath{stroke}
\pgfsetbuttcap
\pgfsetmiterjoin
\pgfsetdash{}{0pt}
\definecolor{dialinecolor}{rgb}{0.811765, 0.054902, 0.188235}
\pgfsetstrokecolor{dialinecolor}
\pgfpathmoveto{\pgfpoint{-6.897692\du}{6.267801\du}}
\pgfpathlineto{\pgfpoint{-6.563911\du}{6.267801\du}}
\pgfpathlineto{\pgfpoint{-6.647506\du}{6.351097\du}}
\pgfpathlineto{\pgfpoint{-6.361791\du}{6.351097\du}}
\pgfusepath{stroke}
\pgfsetbuttcap
\pgfsetmiterjoin
\pgfsetdash{}{0pt}
\definecolor{dialinecolor}{rgb}{0.811765, 0.054902, 0.188235}
\pgfsetstrokecolor{dialinecolor}
\pgfpathmoveto{\pgfpoint{-6.933519\du}{6.267801\du}}
\pgfpathlineto{\pgfpoint{-6.694976\du}{6.505747\du}}
\pgfpathlineto{\pgfpoint{-6.814396\du}{6.505747\du}}
\pgfpathlineto{\pgfpoint{-6.611978\du}{6.708464\du}}
\pgfusepath{stroke}
\pgfsetbuttcap
\pgfsetmiterjoin
\pgfsetdash{}{0pt}
\definecolor{dialinecolor}{rgb}{0.811765, 0.054902, 0.188235}
\pgfsetstrokecolor{dialinecolor}
\pgfpathmoveto{\pgfpoint{-6.255208\du}{6.267801\du}}
\pgfpathlineto{\pgfpoint{-6.492557\du}{6.505747\du}}
\pgfpathlineto{\pgfpoint{-6.374032\du}{6.505747\du}}
\pgfpathlineto{\pgfpoint{-6.576152\du}{6.708464\du}}
\pgfusepath{stroke}
\pgfsetbuttcap
\pgfsetmiterjoin
\pgfsetdash{}{0pt}
\definecolor{dialinecolor}{rgb}{0.000000, 0.470588, 0.666667}
\pgfsetfillcolor{dialinecolor}
\pgfpathmoveto{\pgfpoint{-6.766927\du}{6.356769\du}}
\pgfpathlineto{\pgfpoint{-6.767524\du}{6.353485\du}}
\pgfpathlineto{\pgfpoint{-6.768419\du}{6.350201\du}}
\pgfpathlineto{\pgfpoint{-6.769315\du}{6.346917\du}}
\pgfpathlineto{\pgfpoint{-6.771106\du}{6.343633\du}}
\pgfpathlineto{\pgfpoint{-6.773495\du}{6.340648\du}}
\pgfpathlineto{\pgfpoint{-6.776182\du}{6.337364\du}}
\pgfpathlineto{\pgfpoint{-6.779167\du}{6.334079\du}}
\pgfpathlineto{\pgfpoint{-6.782750\du}{6.331094\du}}
\pgfpathlineto{\pgfpoint{-6.786930\du}{6.328108\du}}
\pgfpathlineto{\pgfpoint{-6.791706\du}{6.325123\du}}
\pgfpathlineto{\pgfpoint{-6.796782\du}{6.322436\du}}
\pgfpathlineto{\pgfpoint{-6.802454\du}{6.319450\du}}
\pgfpathlineto{\pgfpoint{-6.808127\du}{6.317062\du}}
\pgfpathlineto{\pgfpoint{-6.814396\du}{6.314375\du}}
\pgfpathlineto{\pgfpoint{-6.820964\du}{6.311987\du}}
\pgfpathlineto{\pgfpoint{-6.828130\du}{6.309598\du}}
\pgfpathlineto{\pgfpoint{-6.835594\du}{6.307508\du}}
\pgfpathlineto{\pgfpoint{-6.843057\du}{6.305120\du}}
\pgfpathlineto{\pgfpoint{-6.851417\du}{6.303627\du}}
\pgfpathlineto{\pgfpoint{-6.859776\du}{6.301537\du}}
\pgfpathlineto{\pgfpoint{-6.867837\du}{6.299746\du}}
\pgfpathlineto{\pgfpoint{-6.877092\du}{6.298552\du}}
\pgfpathlineto{\pgfpoint{-6.885750\du}{6.297059\du}}
\pgfpathlineto{\pgfpoint{-6.895304\du}{6.295865\du}}
\pgfpathlineto{\pgfpoint{-6.904559\du}{6.294671\du}}
\pgfpathlineto{\pgfpoint{-6.913814\du}{6.293775\du}}
\pgfpathlineto{\pgfpoint{-6.923666\du}{6.292879\du}}
\pgfpathlineto{\pgfpoint{-6.933519\du}{6.292282\du}}
\pgfpathlineto{\pgfpoint{-6.943371\du}{6.291984\du}}
\pgfpathlineto{\pgfpoint{-6.953522\du}{6.291685\du}}
\pgfpathlineto{\pgfpoint{-6.963374\du}{6.291685\du}}
\pgfpathlineto{\pgfpoint{-6.963374\du}{6.291685\du}}
\pgfpathlineto{\pgfpoint{-6.973226\du}{6.291685\du}}
\pgfpathlineto{\pgfpoint{-6.983078\du}{6.291984\du}}
\pgfpathlineto{\pgfpoint{-6.993229\du}{6.292282\du}}
\pgfpathlineto{\pgfpoint{-7.003081\du}{6.292879\du}}
\pgfpathlineto{\pgfpoint{-7.012635\du}{6.293775\du}}
\pgfpathlineto{\pgfpoint{-7.022189\du}{6.294671\du}}
\pgfpathlineto{\pgfpoint{-7.031444\du}{6.295865\du}}
\pgfpathlineto{\pgfpoint{-7.040699\du}{6.297059\du}}
\pgfpathlineto{\pgfpoint{-7.049655\du}{6.298552\du}}
\pgfpathlineto{\pgfpoint{-7.058612\du}{6.299746\du}}
\pgfpathlineto{\pgfpoint{-7.066971\du}{6.301537\du}}
\pgfpathlineto{\pgfpoint{-7.075331\du}{6.303627\du}}
\pgfpathlineto{\pgfpoint{-7.083690\du}{6.305120\du}}
\pgfpathlineto{\pgfpoint{-7.090856\du}{6.307508\du}}
\pgfpathlineto{\pgfpoint{-7.098319\du}{6.309598\du}}
\pgfpathlineto{\pgfpoint{-7.105485\du}{6.311987\du}}
\pgfpathlineto{\pgfpoint{-7.112351\du}{6.314375\du}}
\pgfpathlineto{\pgfpoint{-7.118322\du}{6.317062\du}}
\pgfpathlineto{\pgfpoint{-7.124293\du}{6.319450\du}}
\pgfpathlineto{\pgfpoint{-7.129966\du}{6.322436\du}}
\pgfpathlineto{\pgfpoint{-7.135041\du}{6.325123\du}}
\pgfpathlineto{\pgfpoint{-7.139818\du}{6.328108\du}}
\pgfpathlineto{\pgfpoint{-7.143998\du}{6.331094\du}}
\pgfpathlineto{\pgfpoint{-7.147282\du}{6.334079\du}}
\pgfpathlineto{\pgfpoint{-7.150566\du}{6.337364\du}}
\pgfpathlineto{\pgfpoint{-7.153551\du}{6.340648\du}}
\pgfpathlineto{\pgfpoint{-7.155641\du}{6.343633\du}}
\pgfpathlineto{\pgfpoint{-7.157134\du}{6.346917\du}}
\pgfpathlineto{\pgfpoint{-7.158925\du}{6.350201\du}}
\pgfpathlineto{\pgfpoint{-7.159224\du}{6.353485\du}}
\pgfpathlineto{\pgfpoint{-7.159522\du}{6.356769\du}}
\pgfpathlineto{\pgfpoint{-7.159522\du}{6.356769\du}}
\pgfpathlineto{\pgfpoint{-7.159224\du}{6.360352\du}}
\pgfpathlineto{\pgfpoint{-7.158925\du}{6.363338\du}}
\pgfpathlineto{\pgfpoint{-7.157134\du}{6.366622\du}}
\pgfpathlineto{\pgfpoint{-7.155641\du}{6.369906\du}}
\pgfpathlineto{\pgfpoint{-7.153551\du}{6.373190\du}}
\pgfpathlineto{\pgfpoint{-7.150566\du}{6.376474\du}}
\pgfpathlineto{\pgfpoint{-7.147282\du}{6.379459\du}}
\pgfpathlineto{\pgfpoint{-7.143998\du}{6.382445\du}}
\pgfpathlineto{\pgfpoint{-7.139818\du}{6.385430\du}}
\pgfpathlineto{\pgfpoint{-7.135041\du}{6.388714\du}}
\pgfpathlineto{\pgfpoint{-7.129966\du}{6.391401\du}}
\pgfpathlineto{\pgfpoint{-7.124293\du}{6.394088\du}}
\pgfpathlineto{\pgfpoint{-7.118322\du}{6.396775\du}}
\pgfpathlineto{\pgfpoint{-7.112351\du}{6.399462\du}}
\pgfpathlineto{\pgfpoint{-7.105485\du}{6.401851\du}}
\pgfpathlineto{\pgfpoint{-7.098319\du}{6.404239\du}}
\pgfpathlineto{\pgfpoint{-7.090856\du}{6.406030\du}}
\pgfpathlineto{\pgfpoint{-7.083690\du}{6.408120\du}}
\pgfpathlineto{\pgfpoint{-7.075331\du}{6.410210\du}}
\pgfpathlineto{\pgfpoint{-7.066971\du}{6.412300\du}}
\pgfpathlineto{\pgfpoint{-7.058612\du}{6.413793\du}}
\pgfpathlineto{\pgfpoint{-7.049655\du}{6.415286\du}}
\pgfpathlineto{\pgfpoint{-7.040699\du}{6.416778\du}}
\pgfpathlineto{\pgfpoint{-7.031444\du}{6.417973\du}}
\pgfpathlineto{\pgfpoint{-7.022189\du}{6.418868\du}}
\pgfpathlineto{\pgfpoint{-7.012635\du}{6.420062\du}}
\pgfpathlineto{\pgfpoint{-7.003081\du}{6.420660\du}}
\pgfpathlineto{\pgfpoint{-6.993229\du}{6.421257\du}}
\pgfpathlineto{\pgfpoint{-6.983078\du}{6.421854\du}}
\pgfpathlineto{\pgfpoint{-6.973226\du}{6.422152\du}}
\pgfpathlineto{\pgfpoint{-6.963374\du}{6.422152\du}}
\pgfpathlineto{\pgfpoint{-6.963374\du}{6.422152\du}}
\pgfpathlineto{\pgfpoint{-6.953522\du}{6.422152\du}}
\pgfpathlineto{\pgfpoint{-6.943371\du}{6.421854\du}}
\pgfpathlineto{\pgfpoint{-6.933519\du}{6.421257\du}}
\pgfpathlineto{\pgfpoint{-6.923666\du}{6.420660\du}}
\pgfpathlineto{\pgfpoint{-6.913814\du}{6.420062\du}}
\pgfpathlineto{\pgfpoint{-6.904559\du}{6.418868\du}}
\pgfpathlineto{\pgfpoint{-6.895304\du}{6.417973\du}}
\pgfpathlineto{\pgfpoint{-6.885750\du}{6.416778\du}}
\pgfpathlineto{\pgfpoint{-6.877092\du}{6.415286\du}}
\pgfpathlineto{\pgfpoint{-6.867837\du}{6.413793\du}}
\pgfpathlineto{\pgfpoint{-6.859776\du}{6.412300\du}}
\pgfpathlineto{\pgfpoint{-6.851417\du}{6.410210\du}}
\pgfpathlineto{\pgfpoint{-6.843057\du}{6.408120\du}}
\pgfpathlineto{\pgfpoint{-6.835594\du}{6.406030\du}}
\pgfpathlineto{\pgfpoint{-6.828130\du}{6.404239\du}}
\pgfpathlineto{\pgfpoint{-6.820964\du}{6.401851\du}}
\pgfpathlineto{\pgfpoint{-6.814396\du}{6.399462\du}}
\pgfpathlineto{\pgfpoint{-6.808127\du}{6.396775\du}}
\pgfpathlineto{\pgfpoint{-6.802454\du}{6.394088\du}}
\pgfpathlineto{\pgfpoint{-6.796782\du}{6.391401\du}}
\pgfpathlineto{\pgfpoint{-6.791706\du}{6.388714\du}}
\pgfpathlineto{\pgfpoint{-6.786930\du}{6.385430\du}}
\pgfpathlineto{\pgfpoint{-6.782750\du}{6.382445\du}}
\pgfpathlineto{\pgfpoint{-6.779167\du}{6.379459\du}}
\pgfpathlineto{\pgfpoint{-6.776182\du}{6.376474\du}}
\pgfpathlineto{\pgfpoint{-6.773495\du}{6.373190\du}}
\pgfpathlineto{\pgfpoint{-6.771106\du}{6.369906\du}}
\pgfpathlineto{\pgfpoint{-6.769315\du}{6.366622\du}}
\pgfpathlineto{\pgfpoint{-6.768419\du}{6.363338\du}}
\pgfpathlineto{\pgfpoint{-6.767524\du}{6.360352\du}}
\pgfpathlineto{\pgfpoint{-6.766927\du}{6.356769\du}}
\pgfusepath{fill}
\pgfsetbuttcap
\pgfsetmiterjoin
\pgfsetdash{}{0pt}
\definecolor{dialinecolor}{rgb}{0.666667, 0.901961, 1.000000}
\pgfsetstrokecolor{dialinecolor}
\pgfpathmoveto{\pgfpoint{-6.779167\du}{6.351097\du}}
\pgfpathlineto{\pgfpoint{-6.779167\du}{6.347514\du}}
\pgfpathlineto{\pgfpoint{-6.780063\du}{6.345126\du}}
\pgfpathlineto{\pgfpoint{-6.781257\du}{6.341842\du}}
\pgfpathlineto{\pgfpoint{-6.782750\du}{6.338856\du}}
\pgfpathlineto{\pgfpoint{-6.785138\du}{6.336169\du}}
\pgfpathlineto{\pgfpoint{-6.787527\du}{6.332587\du}}
\pgfpathlineto{\pgfpoint{-6.790811\du}{6.330198\du}}
\pgfpathlineto{\pgfpoint{-6.794095\du}{6.327511\du}}
\pgfpathlineto{\pgfpoint{-6.798275\du}{6.324824\du}}
\pgfpathlineto{\pgfpoint{-6.802753\du}{6.322436\du}}
\pgfpathlineto{\pgfpoint{-6.807530\du}{6.319450\du}}
\pgfpathlineto{\pgfpoint{-6.813202\du}{6.317361\du}}
\pgfpathlineto{\pgfpoint{-6.818576\du}{6.314674\du}}
\pgfpathlineto{\pgfpoint{-6.824846\du}{6.311987\du}}
\pgfpathlineto{\pgfpoint{-6.830817\du}{6.310195\du}}
\pgfpathlineto{\pgfpoint{-6.838280\du}{6.307807\du}}
\pgfpathlineto{\pgfpoint{-6.845147\du}{6.306314\du}}
\pgfpathlineto{\pgfpoint{-6.852312\du}{6.304224\du}}
\pgfpathlineto{\pgfpoint{-6.860075\du}{6.302433\du}}
\pgfpathlineto{\pgfpoint{-6.868733\du}{6.300642\du}}
\pgfpathlineto{\pgfpoint{-6.877092\du}{6.299447\du}}
\pgfpathlineto{\pgfpoint{-6.885452\du}{6.297059\du}}
\pgfpathlineto{\pgfpoint{-6.894408\du}{6.296760\du}}
\pgfpathlineto{\pgfpoint{-6.903066\du}{6.295268\du}}
\pgfpathlineto{\pgfpoint{-6.912023\du}{6.294372\du}}
\pgfpathlineto{\pgfpoint{-6.921278\du}{6.293476\du}}
\pgfpathlineto{\pgfpoint{-6.930832\du}{6.292879\du}}
\pgfpathlineto{\pgfpoint{-6.940385\du}{6.292581\du}}
\pgfpathlineto{\pgfpoint{-6.949342\du}{6.291984\du}}
\pgfpathlineto{\pgfpoint{-6.959493\du}{6.291685\du}}
\pgfpathlineto{\pgfpoint{-6.969046\du}{6.291685\du}}
\pgfpathlineto{\pgfpoint{-6.969046\du}{6.291685\du}}
\pgfpathlineto{\pgfpoint{-6.978301\du}{6.291685\du}}
\pgfpathlineto{\pgfpoint{-6.988154\du}{6.291984\du}}
\pgfpathlineto{\pgfpoint{-6.997707\du}{6.292581\du}}
\pgfpathlineto{\pgfpoint{-7.007261\du}{6.292879\du}}
\pgfpathlineto{\pgfpoint{-7.016516\du}{6.293476\du}}
\pgfpathlineto{\pgfpoint{-7.026368\du}{6.294372\du}}
\pgfpathlineto{\pgfpoint{-7.035026\du}{6.295268\du}}
\pgfpathlineto{\pgfpoint{-7.043983\du}{6.296760\du}}
\pgfpathlineto{\pgfpoint{-7.052342\du}{6.297059\du}}
\pgfpathlineto{\pgfpoint{-7.061299\du}{6.299447\du}}
\pgfpathlineto{\pgfpoint{-7.069658\du}{6.300642\du}}
\pgfpathlineto{\pgfpoint{-7.077719\du}{6.302433\du}}
\pgfpathlineto{\pgfpoint{-7.085482\du}{6.304224\du}}
\pgfpathlineto{\pgfpoint{-7.092945\du}{6.306314\du}}
\pgfpathlineto{\pgfpoint{-7.099812\du}{6.307807\du}}
\pgfpathlineto{\pgfpoint{-7.107276\du}{6.310195\du}}
\pgfpathlineto{\pgfpoint{-7.113247\du}{6.311987\du}}
\pgfpathlineto{\pgfpoint{-7.119516\du}{6.314674\du}}
\pgfpathlineto{\pgfpoint{-7.125189\du}{6.317361\du}}
\pgfpathlineto{\pgfpoint{-7.130563\du}{6.319450\du}}
\pgfpathlineto{\pgfpoint{-7.135340\du}{6.322436\du}}
\pgfpathlineto{\pgfpoint{-7.139818\du}{6.324824\du}}
\pgfpathlineto{\pgfpoint{-7.143998\du}{6.327511\du}}
\pgfpathlineto{\pgfpoint{-7.147282\du}{6.330198\du}}
\pgfpathlineto{\pgfpoint{-7.150864\du}{6.332587\du}}
\pgfpathlineto{\pgfpoint{-7.152656\du}{6.336169\du}}
\pgfpathlineto{\pgfpoint{-7.155044\du}{6.338856\du}}
\pgfpathlineto{\pgfpoint{-7.156835\du}{6.341842\du}}
\pgfpathlineto{\pgfpoint{-7.158328\du}{6.345126\du}}
\pgfpathlineto{\pgfpoint{-7.158925\du}{6.347514\du}}
\pgfpathlineto{\pgfpoint{-7.159224\du}{6.351097\du}}
\pgfpathlineto{\pgfpoint{-7.159224\du}{6.351097\du}}
\pgfpathlineto{\pgfpoint{-7.158925\du}{6.354082\du}}
\pgfpathlineto{\pgfpoint{-7.158328\du}{6.356769\du}}
\pgfpathlineto{\pgfpoint{-7.156835\du}{6.359755\du}}
\pgfpathlineto{\pgfpoint{-7.155044\du}{6.363338\du}}
\pgfpathlineto{\pgfpoint{-7.152656\du}{6.366024\du}}
\pgfpathlineto{\pgfpoint{-7.150864\du}{6.369010\du}}
\pgfpathlineto{\pgfpoint{-7.147282\du}{6.371697\du}}
\pgfpathlineto{\pgfpoint{-7.143998\du}{6.374384\du}}
\pgfpathlineto{\pgfpoint{-7.139818\du}{6.377071\du}}
\pgfpathlineto{\pgfpoint{-7.135340\du}{6.379758\du}}
\pgfpathlineto{\pgfpoint{-7.130563\du}{6.382445\du}}
\pgfpathlineto{\pgfpoint{-7.125189\du}{6.384833\du}}
\pgfpathlineto{\pgfpoint{-7.119516\du}{6.387520\du}}
\pgfpathlineto{\pgfpoint{-7.113247\du}{6.389610\du}}
\pgfpathlineto{\pgfpoint{-7.107276\du}{6.391999\du}}
\pgfpathlineto{\pgfpoint{-7.099812\du}{6.394387\du}}
\pgfpathlineto{\pgfpoint{-7.092945\du}{6.396178\du}}
\pgfpathlineto{\pgfpoint{-7.085482\du}{6.397671\du}}
\pgfpathlineto{\pgfpoint{-7.077719\du}{6.399761\du}}
\pgfpathlineto{\pgfpoint{-7.069658\du}{6.401552\du}}
\pgfpathlineto{\pgfpoint{-7.061299\du}{6.402746\du}}
\pgfpathlineto{\pgfpoint{-7.052342\du}{6.404538\du}}
\pgfpathlineto{\pgfpoint{-7.043983\du}{6.405433\du}}
\pgfpathlineto{\pgfpoint{-7.035026\du}{6.406926\du}}
\pgfpathlineto{\pgfpoint{-7.026368\du}{6.407822\du}}
\pgfpathlineto{\pgfpoint{-7.016516\du}{6.408120\du}}
\pgfpathlineto{\pgfpoint{-7.007261\du}{6.409315\du}}
\pgfpathlineto{\pgfpoint{-6.997707\du}{6.409613\du}}
\pgfpathlineto{\pgfpoint{-6.988154\du}{6.409912\du}}
\pgfpathlineto{\pgfpoint{-6.978301\du}{6.410509\du}}
\pgfpathlineto{\pgfpoint{-6.969046\du}{6.410509\du}}
\pgfpathlineto{\pgfpoint{-6.969046\du}{6.410509\du}}
\pgfpathlineto{\pgfpoint{-6.959493\du}{6.410509\du}}
\pgfpathlineto{\pgfpoint{-6.949342\du}{6.409912\du}}
\pgfpathlineto{\pgfpoint{-6.940385\du}{6.409613\du}}
\pgfpathlineto{\pgfpoint{-6.930832\du}{6.409315\du}}
\pgfpathlineto{\pgfpoint{-6.921278\du}{6.408120\du}}
\pgfpathlineto{\pgfpoint{-6.912023\du}{6.407822\du}}
\pgfpathlineto{\pgfpoint{-6.903066\du}{6.406926\du}}
\pgfpathlineto{\pgfpoint{-6.894408\du}{6.405433\du}}
\pgfpathlineto{\pgfpoint{-6.885452\du}{6.404538\du}}
\pgfpathlineto{\pgfpoint{-6.877092\du}{6.402746\du}}
\pgfpathlineto{\pgfpoint{-6.868733\du}{6.401552\du}}
\pgfpathlineto{\pgfpoint{-6.860075\du}{6.399761\du}}
\pgfpathlineto{\pgfpoint{-6.852312\du}{6.397671\du}}
\pgfpathlineto{\pgfpoint{-6.845147\du}{6.396178\du}}
\pgfpathlineto{\pgfpoint{-6.838280\du}{6.394387\du}}
\pgfpathlineto{\pgfpoint{-6.830817\du}{6.391999\du}}
\pgfpathlineto{\pgfpoint{-6.824846\du}{6.389610\du}}
\pgfpathlineto{\pgfpoint{-6.818576\du}{6.387520\du}}
\pgfpathlineto{\pgfpoint{-6.813202\du}{6.384833\du}}
\pgfpathlineto{\pgfpoint{-6.807530\du}{6.382445\du}}
\pgfpathlineto{\pgfpoint{-6.802753\du}{6.379758\du}}
\pgfpathlineto{\pgfpoint{-6.798275\du}{6.377071\du}}
\pgfpathlineto{\pgfpoint{-6.794095\du}{6.374384\du}}
\pgfpathlineto{\pgfpoint{-6.790811\du}{6.371697\du}}
\pgfpathlineto{\pgfpoint{-6.787527\du}{6.369010\du}}
\pgfpathlineto{\pgfpoint{-6.785138\du}{6.366024\du}}
\pgfpathlineto{\pgfpoint{-6.782750\du}{6.363338\du}}
\pgfpathlineto{\pgfpoint{-6.781257\du}{6.359755\du}}
\pgfpathlineto{\pgfpoint{-6.780063\du}{6.356769\du}}
\pgfpathlineto{\pgfpoint{-6.779167\du}{6.354082\du}}
\pgfpathlineto{\pgfpoint{-6.779167\du}{6.351097\du}}
\pgfusepath{stroke}
\pgfsetbuttcap
\pgfsetmiterjoin
\pgfsetdash{}{0pt}
\definecolor{dialinecolor}{rgb}{0.000000, 0.470588, 0.666667}
\pgfsetfillcolor{dialinecolor}
\pgfpathmoveto{\pgfpoint{-7.159224\du}{6.267801\du}}
\pgfpathlineto{\pgfpoint{-7.159224\du}{6.362442\du}}
\pgfpathlineto{\pgfpoint{-6.766628\du}{6.362442\du}}
\pgfpathlineto{\pgfpoint{-6.766628\du}{6.267801\du}}
\pgfpathlineto{\pgfpoint{-7.159224\du}{6.267801\du}}
\pgfusepath{fill}
\pgfsetbuttcap
\pgfsetmiterjoin
\pgfsetdash{}{0pt}
\definecolor{dialinecolor}{rgb}{0.000000, 0.705882, 1.000000}
\pgfsetfillcolor{dialinecolor}
\pgfpathmoveto{\pgfpoint{-6.766927\du}{6.261233\du}}
\pgfpathlineto{\pgfpoint{-6.767524\du}{6.257949\du}}
\pgfpathlineto{\pgfpoint{-6.768419\du}{6.254366\du}}
\pgfpathlineto{\pgfpoint{-6.769315\du}{6.251381\du}}
\pgfpathlineto{\pgfpoint{-6.771106\du}{6.248096\du}}
\pgfpathlineto{\pgfpoint{-6.773495\du}{6.244812\du}}
\pgfpathlineto{\pgfpoint{-6.776182\du}{6.241528\du}}
\pgfpathlineto{\pgfpoint{-6.779167\du}{6.238543\du}}
\pgfpathlineto{\pgfpoint{-6.782750\du}{6.235259\du}}
\pgfpathlineto{\pgfpoint{-6.786930\du}{6.232572\du}}
\pgfpathlineto{\pgfpoint{-6.791706\du}{6.229586\du}}
\pgfpathlineto{\pgfpoint{-6.796782\du}{6.226601\du}}
\pgfpathlineto{\pgfpoint{-6.802454\du}{6.223914\du}}
\pgfpathlineto{\pgfpoint{-6.808127\du}{6.220630\du}}
\pgfpathlineto{\pgfpoint{-6.814396\du}{6.218838\du}}
\pgfpathlineto{\pgfpoint{-6.820964\du}{6.216151\du}}
\pgfpathlineto{\pgfpoint{-6.828130\du}{6.213763\du}}
\pgfpathlineto{\pgfpoint{-6.835594\du}{6.211673\du}}
\pgfpathlineto{\pgfpoint{-6.843057\du}{6.209583\du}}
\pgfpathlineto{\pgfpoint{-6.851417\du}{6.207493\du}}
\pgfpathlineto{\pgfpoint{-6.859776\du}{6.205404\du}}
\pgfpathlineto{\pgfpoint{-6.867837\du}{6.203911\du}}
\pgfpathlineto{\pgfpoint{-6.877092\du}{6.202418\du}}
\pgfpathlineto{\pgfpoint{-6.885750\du}{6.201224\du}}
\pgfpathlineto{\pgfpoint{-6.895304\du}{6.199731\du}}
\pgfpathlineto{\pgfpoint{-6.904559\du}{6.198835\du}}
\pgfpathlineto{\pgfpoint{-6.913814\du}{6.197641\du}}
\pgfpathlineto{\pgfpoint{-6.923666\du}{6.197044\du}}
\pgfpathlineto{\pgfpoint{-6.933519\du}{6.196447\du}}
\pgfpathlineto{\pgfpoint{-6.943371\du}{6.196148\du}}
\pgfpathlineto{\pgfpoint{-6.953522\du}{6.195850\du}}
\pgfpathlineto{\pgfpoint{-6.963374\du}{6.195850\du}}
\pgfpathlineto{\pgfpoint{-6.963374\du}{6.195850\du}}
\pgfpathlineto{\pgfpoint{-6.973226\du}{6.195850\du}}
\pgfpathlineto{\pgfpoint{-6.983078\du}{6.196148\du}}
\pgfpathlineto{\pgfpoint{-6.993229\du}{6.196447\du}}
\pgfpathlineto{\pgfpoint{-7.003081\du}{6.197044\du}}
\pgfpathlineto{\pgfpoint{-7.012635\du}{6.197641\du}}
\pgfpathlineto{\pgfpoint{-7.022189\du}{6.198835\du}}
\pgfpathlineto{\pgfpoint{-7.031444\du}{6.199731\du}}
\pgfpathlineto{\pgfpoint{-7.040699\du}{6.201224\du}}
\pgfpathlineto{\pgfpoint{-7.049655\du}{6.202418\du}}
\pgfpathlineto{\pgfpoint{-7.058612\du}{6.203911\du}}
\pgfpathlineto{\pgfpoint{-7.066971\du}{6.205404\du}}
\pgfpathlineto{\pgfpoint{-7.075331\du}{6.207493\du}}
\pgfpathlineto{\pgfpoint{-7.083690\du}{6.209583\du}}
\pgfpathlineto{\pgfpoint{-7.090856\du}{6.211673\du}}
\pgfpathlineto{\pgfpoint{-7.098319\du}{6.213763\du}}
\pgfpathlineto{\pgfpoint{-7.105485\du}{6.216151\du}}
\pgfpathlineto{\pgfpoint{-7.112351\du}{6.218838\du}}
\pgfpathlineto{\pgfpoint{-7.118322\du}{6.220630\du}}
\pgfpathlineto{\pgfpoint{-7.124293\du}{6.223914\du}}
\pgfpathlineto{\pgfpoint{-7.129966\du}{6.226601\du}}
\pgfpathlineto{\pgfpoint{-7.135041\du}{6.229586\du}}
\pgfpathlineto{\pgfpoint{-7.139818\du}{6.232572\du}}
\pgfpathlineto{\pgfpoint{-7.143998\du}{6.235259\du}}
\pgfpathlineto{\pgfpoint{-7.147282\du}{6.238543\du}}
\pgfpathlineto{\pgfpoint{-7.150566\du}{6.241528\du}}
\pgfpathlineto{\pgfpoint{-7.153551\du}{6.244812\du}}
\pgfpathlineto{\pgfpoint{-7.155641\du}{6.248096\du}}
\pgfpathlineto{\pgfpoint{-7.157134\du}{6.251381\du}}
\pgfpathlineto{\pgfpoint{-7.158925\du}{6.254366\du}}
\pgfpathlineto{\pgfpoint{-7.159224\du}{6.257949\du}}
\pgfpathlineto{\pgfpoint{-7.159522\du}{6.261233\du}}
\pgfpathlineto{\pgfpoint{-7.159522\du}{6.261233\du}}
\pgfpathlineto{\pgfpoint{-7.159224\du}{6.264517\du}}
\pgfpathlineto{\pgfpoint{-7.158925\du}{6.268099\du}}
\pgfpathlineto{\pgfpoint{-7.157134\du}{6.271085\du}}
\pgfpathlineto{\pgfpoint{-7.155641\du}{6.274369\du}}
\pgfpathlineto{\pgfpoint{-7.153551\du}{6.277952\du}}
\pgfpathlineto{\pgfpoint{-7.150566\du}{6.280937\du}}
\pgfpathlineto{\pgfpoint{-7.147282\du}{6.283923\du}}
\pgfpathlineto{\pgfpoint{-7.143998\du}{6.287207\du}}
\pgfpathlineto{\pgfpoint{-7.139818\du}{6.289595\du}}
\pgfpathlineto{\pgfpoint{-7.135041\du}{6.292879\du}}
\pgfpathlineto{\pgfpoint{-7.129966\du}{6.296163\du}}
\pgfpathlineto{\pgfpoint{-7.124293\du}{6.298850\du}}
\pgfpathlineto{\pgfpoint{-7.118322\du}{6.301537\du}}
\pgfpathlineto{\pgfpoint{-7.112351\du}{6.303627\du}}
\pgfpathlineto{\pgfpoint{-7.105485\du}{6.306613\du}}
\pgfpathlineto{\pgfpoint{-7.098319\du}{6.309001\du}}
\pgfpathlineto{\pgfpoint{-7.090856\du}{6.310792\du}}
\pgfpathlineto{\pgfpoint{-7.083690\du}{6.312882\du}}
\pgfpathlineto{\pgfpoint{-7.075331\du}{6.314972\du}}
\pgfpathlineto{\pgfpoint{-7.066971\du}{6.317062\du}}
\pgfpathlineto{\pgfpoint{-7.058612\du}{6.318555\du}}
\pgfpathlineto{\pgfpoint{-7.049655\du}{6.320047\du}}
\pgfpathlineto{\pgfpoint{-7.040699\du}{6.321242\du}}
\pgfpathlineto{\pgfpoint{-7.031444\du}{6.322734\du}}
\pgfpathlineto{\pgfpoint{-7.022189\du}{6.323630\du}}
\pgfpathlineto{\pgfpoint{-7.012635\du}{6.324824\du}}
\pgfpathlineto{\pgfpoint{-7.003081\du}{6.325123\du}}
\pgfpathlineto{\pgfpoint{-6.993229\du}{6.326019\du}}
\pgfpathlineto{\pgfpoint{-6.983078\du}{6.326019\du}}
\pgfpathlineto{\pgfpoint{-6.973226\du}{6.326914\du}}
\pgfpathlineto{\pgfpoint{-6.963374\du}{6.326914\du}}
\pgfpathlineto{\pgfpoint{-6.963374\du}{6.326914\du}}
\pgfpathlineto{\pgfpoint{-6.953522\du}{6.326914\du}}
\pgfpathlineto{\pgfpoint{-6.943371\du}{6.326019\du}}
\pgfpathlineto{\pgfpoint{-6.933519\du}{6.326019\du}}
\pgfpathlineto{\pgfpoint{-6.923666\du}{6.325123\du}}
\pgfpathlineto{\pgfpoint{-6.913814\du}{6.324824\du}}
\pgfpathlineto{\pgfpoint{-6.904559\du}{6.323630\du}}
\pgfpathlineto{\pgfpoint{-6.895304\du}{6.322734\du}}
\pgfpathlineto{\pgfpoint{-6.885750\du}{6.321242\du}}
\pgfpathlineto{\pgfpoint{-6.877092\du}{6.320047\du}}
\pgfpathlineto{\pgfpoint{-6.867837\du}{6.318555\du}}
\pgfpathlineto{\pgfpoint{-6.859776\du}{6.317062\du}}
\pgfpathlineto{\pgfpoint{-6.851417\du}{6.314972\du}}
\pgfpathlineto{\pgfpoint{-6.843057\du}{6.312882\du}}
\pgfpathlineto{\pgfpoint{-6.835594\du}{6.310792\du}}
\pgfpathlineto{\pgfpoint{-6.828130\du}{6.309001\du}}
\pgfpathlineto{\pgfpoint{-6.820964\du}{6.306613\du}}
\pgfpathlineto{\pgfpoint{-6.814396\du}{6.303627\du}}
\pgfpathlineto{\pgfpoint{-6.808127\du}{6.301537\du}}
\pgfpathlineto{\pgfpoint{-6.802454\du}{6.298850\du}}
\pgfpathlineto{\pgfpoint{-6.796782\du}{6.296163\du}}
\pgfpathlineto{\pgfpoint{-6.791706\du}{6.292879\du}}
\pgfpathlineto{\pgfpoint{-6.786930\du}{6.289595\du}}
\pgfpathlineto{\pgfpoint{-6.782750\du}{6.287207\du}}
\pgfpathlineto{\pgfpoint{-6.779167\du}{6.283923\du}}
\pgfpathlineto{\pgfpoint{-6.776182\du}{6.280937\du}}
\pgfpathlineto{\pgfpoint{-6.773495\du}{6.277952\du}}
\pgfpathlineto{\pgfpoint{-6.771106\du}{6.274369\du}}
\pgfpathlineto{\pgfpoint{-6.769315\du}{6.271085\du}}
\pgfpathlineto{\pgfpoint{-6.768419\du}{6.268099\du}}
\pgfpathlineto{\pgfpoint{-6.767524\du}{6.264517\du}}
\pgfpathlineto{\pgfpoint{-6.766927\du}{6.261233\du}}
\pgfusepath{fill}
\pgfsetbuttcap
\pgfsetmiterjoin
\pgfsetdash{}{0pt}
\definecolor{dialinecolor}{rgb}{0.666667, 0.901961, 1.000000}
\pgfsetstrokecolor{dialinecolor}
\pgfpathmoveto{\pgfpoint{-6.779167\du}{6.255560\du}}
\pgfpathlineto{\pgfpoint{-6.779167\du}{6.252276\du}}
\pgfpathlineto{\pgfpoint{-6.780063\du}{6.249888\du}}
\pgfpathlineto{\pgfpoint{-6.781257\du}{6.246604\du}}
\pgfpathlineto{\pgfpoint{-6.782750\du}{6.243021\du}}
\pgfpathlineto{\pgfpoint{-6.785138\du}{6.240633\du}}
\pgfpathlineto{\pgfpoint{-6.787527\du}{6.237349\du}}
\pgfpathlineto{\pgfpoint{-6.790811\du}{6.234960\du}}
\pgfpathlineto{\pgfpoint{-6.794095\du}{6.232273\du}}
\pgfpathlineto{\pgfpoint{-6.798275\du}{6.229586\du}}
\pgfpathlineto{\pgfpoint{-6.802753\du}{6.226899\du}}
\pgfpathlineto{\pgfpoint{-6.807530\du}{6.224212\du}}
\pgfpathlineto{\pgfpoint{-6.813202\du}{6.221525\du}}
\pgfpathlineto{\pgfpoint{-6.818576\du}{6.219137\du}}
\pgfpathlineto{\pgfpoint{-6.824846\du}{6.217047\du}}
\pgfpathlineto{\pgfpoint{-6.830817\du}{6.214659\du}}
\pgfpathlineto{\pgfpoint{-6.838280\du}{6.212270\du}}
\pgfpathlineto{\pgfpoint{-6.845147\du}{6.210479\du}}
\pgfpathlineto{\pgfpoint{-6.852312\du}{6.208986\du}}
\pgfpathlineto{\pgfpoint{-6.860075\du}{6.206598\du}}
\pgfpathlineto{\pgfpoint{-6.868733\du}{6.205105\du}}
\pgfpathlineto{\pgfpoint{-6.877092\du}{6.203911\du}}
\pgfpathlineto{\pgfpoint{-6.885452\du}{6.202119\du}}
\pgfpathlineto{\pgfpoint{-6.894408\du}{6.201224\du}}
\pgfpathlineto{\pgfpoint{-6.903066\du}{6.199731\du}}
\pgfpathlineto{\pgfpoint{-6.912023\du}{6.198835\du}}
\pgfpathlineto{\pgfpoint{-6.921278\du}{6.198238\du}}
\pgfpathlineto{\pgfpoint{-6.930832\du}{6.197343\du}}
\pgfpathlineto{\pgfpoint{-6.940385\du}{6.197044\du}}
\pgfpathlineto{\pgfpoint{-6.949342\du}{6.196745\du}}
\pgfpathlineto{\pgfpoint{-6.959493\du}{6.196148\du}}
\pgfpathlineto{\pgfpoint{-6.969046\du}{6.196148\du}}
\pgfpathlineto{\pgfpoint{-6.969046\du}{6.196148\du}}
\pgfpathlineto{\pgfpoint{-6.978301\du}{6.196148\du}}
\pgfpathlineto{\pgfpoint{-6.988154\du}{6.196745\du}}
\pgfpathlineto{\pgfpoint{-6.997707\du}{6.197044\du}}
\pgfpathlineto{\pgfpoint{-7.007261\du}{6.197343\du}}
\pgfpathlineto{\pgfpoint{-7.016516\du}{6.198238\du}}
\pgfpathlineto{\pgfpoint{-7.026368\du}{6.198835\du}}
\pgfpathlineto{\pgfpoint{-7.035026\du}{6.199731\du}}
\pgfpathlineto{\pgfpoint{-7.043983\du}{6.201224\du}}
\pgfpathlineto{\pgfpoint{-7.052342\du}{6.202119\du}}
\pgfpathlineto{\pgfpoint{-7.061299\du}{6.203911\du}}
\pgfpathlineto{\pgfpoint{-7.069658\du}{6.205105\du}}
\pgfpathlineto{\pgfpoint{-7.077719\du}{6.206598\du}}
\pgfpathlineto{\pgfpoint{-7.085482\du}{6.208986\du}}
\pgfpathlineto{\pgfpoint{-7.092945\du}{6.210479\du}}
\pgfpathlineto{\pgfpoint{-7.099812\du}{6.212270\du}}
\pgfpathlineto{\pgfpoint{-7.107276\du}{6.214659\du}}
\pgfpathlineto{\pgfpoint{-7.113247\du}{6.217047\du}}
\pgfpathlineto{\pgfpoint{-7.119516\du}{6.219137\du}}
\pgfpathlineto{\pgfpoint{-7.125189\du}{6.221525\du}}
\pgfpathlineto{\pgfpoint{-7.130563\du}{6.224212\du}}
\pgfpathlineto{\pgfpoint{-7.135340\du}{6.226899\du}}
\pgfpathlineto{\pgfpoint{-7.139818\du}{6.229586\du}}
\pgfpathlineto{\pgfpoint{-7.143998\du}{6.232273\du}}
\pgfpathlineto{\pgfpoint{-7.147282\du}{6.234960\du}}
\pgfpathlineto{\pgfpoint{-7.150864\du}{6.237349\du}}
\pgfpathlineto{\pgfpoint{-7.152656\du}{6.240633\du}}
\pgfpathlineto{\pgfpoint{-7.155044\du}{6.243021\du}}
\pgfpathlineto{\pgfpoint{-7.156835\du}{6.246604\du}}
\pgfpathlineto{\pgfpoint{-7.158328\du}{6.249888\du}}
\pgfpathlineto{\pgfpoint{-7.158925\du}{6.252276\du}}
\pgfpathlineto{\pgfpoint{-7.159224\du}{6.255560\du}}
\pgfpathlineto{\pgfpoint{-7.159224\du}{6.255560\du}}
\pgfpathlineto{\pgfpoint{-7.158925\du}{6.258844\du}}
\pgfpathlineto{\pgfpoint{-7.158328\du}{6.261531\du}}
\pgfpathlineto{\pgfpoint{-7.156835\du}{6.264517\du}}
\pgfpathlineto{\pgfpoint{-7.155044\du}{6.267801\du}}
\pgfpathlineto{\pgfpoint{-7.152656\du}{6.270488\du}}
\pgfpathlineto{\pgfpoint{-7.150864\du}{6.273772\du}}
\pgfpathlineto{\pgfpoint{-7.147282\du}{6.276459\du}}
\pgfpathlineto{\pgfpoint{-7.143998\du}{6.279146\du}}
\pgfpathlineto{\pgfpoint{-7.139818\du}{6.281833\du}}
\pgfpathlineto{\pgfpoint{-7.135340\du}{6.284520\du}}
\pgfpathlineto{\pgfpoint{-7.130563\du}{6.287207\du}}
\pgfpathlineto{\pgfpoint{-7.125189\du}{6.289297\du}}
\pgfpathlineto{\pgfpoint{-7.119516\du}{6.291984\du}}
\pgfpathlineto{\pgfpoint{-7.113247\du}{6.294372\du}}
\pgfpathlineto{\pgfpoint{-7.107276\du}{6.296760\du}}
\pgfpathlineto{\pgfpoint{-7.099812\du}{6.298850\du}}
\pgfpathlineto{\pgfpoint{-7.092945\du}{6.300642\du}}
\pgfpathlineto{\pgfpoint{-7.085482\du}{6.302433\du}}
\pgfpathlineto{\pgfpoint{-7.077719\du}{6.304224\du}}
\pgfpathlineto{\pgfpoint{-7.069658\du}{6.306314\du}}
\pgfpathlineto{\pgfpoint{-7.061299\du}{6.307508\du}}
\pgfpathlineto{\pgfpoint{-7.052342\du}{6.309300\du}}
\pgfpathlineto{\pgfpoint{-7.043983\du}{6.310195\du}}
\pgfpathlineto{\pgfpoint{-7.035026\du}{6.311688\du}}
\pgfpathlineto{\pgfpoint{-7.026368\du}{6.311987\du}}
\pgfpathlineto{\pgfpoint{-7.016516\du}{6.313181\du}}
\pgfpathlineto{\pgfpoint{-7.007261\du}{6.313479\du}}
\pgfpathlineto{\pgfpoint{-6.997707\du}{6.314375\du}}
\pgfpathlineto{\pgfpoint{-6.988154\du}{6.314674\du}}
\pgfpathlineto{\pgfpoint{-6.978301\du}{6.314972\du}}
\pgfpathlineto{\pgfpoint{-6.969046\du}{6.314972\du}}
\pgfpathlineto{\pgfpoint{-6.969046\du}{6.314972\du}}
\pgfpathlineto{\pgfpoint{-6.959493\du}{6.314972\du}}
\pgfpathlineto{\pgfpoint{-6.949342\du}{6.314674\du}}
\pgfpathlineto{\pgfpoint{-6.940385\du}{6.314375\du}}
\pgfpathlineto{\pgfpoint{-6.930832\du}{6.313479\du}}
\pgfpathlineto{\pgfpoint{-6.921278\du}{6.313181\du}}
\pgfpathlineto{\pgfpoint{-6.912023\du}{6.311987\du}}
\pgfpathlineto{\pgfpoint{-6.903066\du}{6.311688\du}}
\pgfpathlineto{\pgfpoint{-6.894408\du}{6.310195\du}}
\pgfpathlineto{\pgfpoint{-6.885452\du}{6.309300\du}}
\pgfpathlineto{\pgfpoint{-6.877092\du}{6.307508\du}}
\pgfpathlineto{\pgfpoint{-6.868733\du}{6.306314\du}}
\pgfpathlineto{\pgfpoint{-6.860075\du}{6.304224\du}}
\pgfpathlineto{\pgfpoint{-6.852312\du}{6.302433\du}}
\pgfpathlineto{\pgfpoint{-6.845147\du}{6.300642\du}}
\pgfpathlineto{\pgfpoint{-6.838280\du}{6.298850\du}}
\pgfpathlineto{\pgfpoint{-6.830817\du}{6.296760\du}}
\pgfpathlineto{\pgfpoint{-6.824846\du}{6.294372\du}}
\pgfpathlineto{\pgfpoint{-6.818576\du}{6.291984\du}}
\pgfpathlineto{\pgfpoint{-6.813202\du}{6.289297\du}}
\pgfpathlineto{\pgfpoint{-6.807530\du}{6.287207\du}}
\pgfpathlineto{\pgfpoint{-6.802753\du}{6.284520\du}}
\pgfpathlineto{\pgfpoint{-6.798275\du}{6.281833\du}}
\pgfpathlineto{\pgfpoint{-6.794095\du}{6.279146\du}}
\pgfpathlineto{\pgfpoint{-6.790811\du}{6.276459\du}}
\pgfpathlineto{\pgfpoint{-6.787527\du}{6.273772\du}}
\pgfpathlineto{\pgfpoint{-6.785138\du}{6.270488\du}}
\pgfpathlineto{\pgfpoint{-6.782750\du}{6.267801\du}}
\pgfpathlineto{\pgfpoint{-6.781257\du}{6.264517\du}}
\pgfpathlineto{\pgfpoint{-6.780063\du}{6.261531\du}}
\pgfpathlineto{\pgfpoint{-6.779167\du}{6.258844\du}}
\pgfpathlineto{\pgfpoint{-6.779167\du}{6.255560\du}}
\pgfusepath{stroke}
\pgfsetbuttcap
\pgfsetmiterjoin
\pgfsetdash{}{0pt}
\definecolor{dialinecolor}{rgb}{0.000000, 0.000000, 0.000000}
\pgfsetfillcolor{dialinecolor}
\pgfpathmoveto{\pgfpoint{-6.956806\du}{6.243917\du}}
\pgfpathlineto{\pgfpoint{-6.933519\du}{6.255560\du}}
\pgfpathlineto{\pgfpoint{-6.861568\du}{6.231975\du}}
\pgfpathlineto{\pgfpoint{-6.826040\du}{6.243917\du}}
\pgfpathlineto{\pgfpoint{-6.849625\du}{6.220033\du}}
\pgfpathlineto{\pgfpoint{-6.933519\du}{6.220033\du}}
\pgfpathlineto{\pgfpoint{-6.897692\du}{6.220033\du}}
\pgfpathlineto{\pgfpoint{-6.956806\du}{6.243917\du}}
\pgfusepath{fill}
\pgfsetbuttcap
\pgfsetmiterjoin
\pgfsetdash{}{0pt}
\definecolor{dialinecolor}{rgb}{0.000000, 0.000000, 0.000000}
\pgfsetfillcolor{dialinecolor}
\pgfpathmoveto{\pgfpoint{-6.969046\du}{6.279444\du}}
\pgfpathlineto{\pgfpoint{-7.004873\du}{6.267801\du}}
\pgfpathlineto{\pgfpoint{-7.063986\du}{6.291685\du}}
\pgfpathlineto{\pgfpoint{-7.099812\du}{6.279444\du}}
\pgfpathlineto{\pgfpoint{-7.087571\du}{6.314972\du}}
\pgfpathlineto{\pgfpoint{-7.004873\du}{6.314972\du}}
\pgfpathlineto{\pgfpoint{-7.040400\du}{6.303627\du}}
\pgfpathlineto{\pgfpoint{-6.969046\du}{6.279444\du}}
\pgfusepath{fill}
\pgfsetbuttcap
\pgfsetmiterjoin
\pgfsetdash{}{0pt}
\definecolor{dialinecolor}{rgb}{0.000000, 0.000000, 0.000000}
\pgfsetfillcolor{dialinecolor}
\pgfpathmoveto{\pgfpoint{-7.087571\du}{6.220033\du}}
\pgfpathlineto{\pgfpoint{-7.063986\du}{6.207493\du}}
\pgfpathlineto{\pgfpoint{-6.992333\du}{6.243917\du}}
\pgfpathlineto{\pgfpoint{-6.969046\du}{6.231975\du}}
\pgfpathlineto{\pgfpoint{-6.980690\du}{6.255560\du}}
\pgfpathlineto{\pgfpoint{-7.063986\du}{6.255560\du}}
\pgfpathlineto{\pgfpoint{-7.028160\du}{6.243917\du}}
\pgfpathlineto{\pgfpoint{-7.087571\du}{6.220033\du}}
\pgfusepath{fill}
\pgfsetbuttcap
\pgfsetmiterjoin
\pgfsetdash{}{0pt}
\definecolor{dialinecolor}{rgb}{0.000000, 0.000000, 0.000000}
\pgfsetfillcolor{dialinecolor}
\pgfpathmoveto{\pgfpoint{-6.838280\du}{6.303627\du}}
\pgfpathlineto{\pgfpoint{-6.861568\du}{6.314972\du}}
\pgfpathlineto{\pgfpoint{-6.933519\du}{6.279444\du}}
\pgfpathlineto{\pgfpoint{-6.969046\du}{6.291685\du}}
\pgfpathlineto{\pgfpoint{-6.945162\du}{6.267801\du}}
\pgfpathlineto{\pgfpoint{-6.861568\du}{6.267801\du}}
\pgfpathlineto{\pgfpoint{-6.897692\du}{6.279444\du}}
\pgfpathlineto{\pgfpoint{-6.838280\du}{6.303627\du}}
\pgfusepath{fill}
\pgfsetbuttcap
\pgfsetmiterjoin
\pgfsetdash{}{0pt}
\definecolor{dialinecolor}{rgb}{1.000000, 1.000000, 1.000000}
\pgfsetfillcolor{dialinecolor}
\pgfpathmoveto{\pgfpoint{-6.956806\du}{6.255560\du}}
\pgfpathlineto{\pgfpoint{-6.933519\du}{6.255560\du}}
\pgfpathlineto{\pgfpoint{-6.861568\du}{6.231975\du}}
\pgfpathlineto{\pgfpoint{-6.826040\du}{6.243917\du}}
\pgfpathlineto{\pgfpoint{-6.849625\du}{6.220033\du}}
\pgfpathlineto{\pgfpoint{-6.921278\du}{6.220033\du}}
\pgfpathlineto{\pgfpoint{-6.897692\du}{6.220033\du}}
\pgfpathlineto{\pgfpoint{-6.956806\du}{6.255560\du}}
\pgfusepath{fill}
\pgfsetbuttcap
\pgfsetmiterjoin
\pgfsetdash{}{0pt}
\definecolor{dialinecolor}{rgb}{1.000000, 1.000000, 1.000000}
\pgfsetfillcolor{dialinecolor}
\pgfpathmoveto{\pgfpoint{-6.969046\du}{6.279444\du}}
\pgfpathlineto{\pgfpoint{-6.992333\du}{6.267801\du}}
\pgfpathlineto{\pgfpoint{-7.063986\du}{6.291685\du}}
\pgfpathlineto{\pgfpoint{-7.099812\du}{6.291685\du}}
\pgfpathlineto{\pgfpoint{-7.076226\du}{6.314972\du}}
\pgfpathlineto{\pgfpoint{-6.992333\du}{6.314972\du}}
\pgfpathlineto{\pgfpoint{-7.028160\du}{6.303627\du}}
\pgfpathlineto{\pgfpoint{-6.969046\du}{6.279444\du}}
\pgfusepath{fill}
\pgfsetbuttcap
\pgfsetmiterjoin
\pgfsetdash{}{0pt}
\definecolor{dialinecolor}{rgb}{1.000000, 1.000000, 1.000000}
\pgfsetfillcolor{dialinecolor}
\pgfpathmoveto{\pgfpoint{-7.087571\du}{6.220033\du}}
\pgfpathlineto{\pgfpoint{-7.063986\du}{6.220033\du}}
\pgfpathlineto{\pgfpoint{-6.992333\du}{6.243917\du}}
\pgfpathlineto{\pgfpoint{-6.956806\du}{6.231975\du}}
\pgfpathlineto{\pgfpoint{-6.980690\du}{6.255560\du}}
\pgfpathlineto{\pgfpoint{-7.063986\du}{6.255560\du}}
\pgfpathlineto{\pgfpoint{-7.028160\du}{6.255560\du}}
\pgfpathlineto{\pgfpoint{-7.087571\du}{6.220033\du}}
\pgfusepath{fill}
\pgfsetbuttcap
\pgfsetmiterjoin
\pgfsetdash{}{0pt}
\definecolor{dialinecolor}{rgb}{1.000000, 1.000000, 1.000000}
\pgfsetfillcolor{dialinecolor}
\pgfpathmoveto{\pgfpoint{-6.838280\du}{6.303627\du}}
\pgfpathlineto{\pgfpoint{-6.861568\du}{6.314972\du}}
\pgfpathlineto{\pgfpoint{-6.933519\du}{6.291685\du}}
\pgfpathlineto{\pgfpoint{-6.956806\du}{6.291685\du}}
\pgfpathlineto{\pgfpoint{-6.945162\du}{6.267801\du}}
\pgfpathlineto{\pgfpoint{-6.861568\du}{6.267801\du}}
\pgfpathlineto{\pgfpoint{-6.897692\du}{6.279444\du}}
\pgfpathlineto{\pgfpoint{-6.838280\du}{6.303627\du}}
\pgfusepath{fill}
\pgfsetbuttcap
\pgfsetmiterjoin
\pgfsetdash{}{0pt}
\definecolor{dialinecolor}{rgb}{0.666667, 0.901961, 1.000000}
\pgfsetstrokecolor{dialinecolor}
\pgfpathmoveto{\pgfpoint{-7.159224\du}{6.255560\du}}
\pgfpathlineto{\pgfpoint{-7.159224\du}{6.350500\du}}
\pgfusepath{stroke}
\pgfsetbuttcap
\pgfsetmiterjoin
\pgfsetdash{}{0pt}
\definecolor{dialinecolor}{rgb}{0.666667, 0.901961, 1.000000}
\pgfsetstrokecolor{dialinecolor}
\pgfpathmoveto{\pgfpoint{-6.778570\du}{6.255560\du}}
\pgfpathlineto{\pgfpoint{-6.778570\du}{6.350500\du}}
\pgfusepath{stroke}
\pgfsetbuttcap
\pgfsetmiterjoin
\pgfsetdash{}{0pt}
\definecolor{dialinecolor}{rgb}{0.000000, 0.470588, 0.666667}
\pgfsetfillcolor{dialinecolor}
\pgfpathmoveto{\pgfpoint{-6.052790\du}{6.356769\du}}
\pgfpathlineto{\pgfpoint{-6.052790\du}{6.353485\du}}
\pgfpathlineto{\pgfpoint{-6.053686\du}{6.350201\du}}
\pgfpathlineto{\pgfpoint{-6.054880\du}{6.346917\du}}
\pgfpathlineto{\pgfpoint{-6.056373\du}{6.343633\du}}
\pgfpathlineto{\pgfpoint{-6.058761\du}{6.340648\du}}
\pgfpathlineto{\pgfpoint{-6.061448\du}{6.337364\du}}
\pgfpathlineto{\pgfpoint{-6.064434\du}{6.334079\du}}
\pgfpathlineto{\pgfpoint{-6.068613\du}{6.331094\du}}
\pgfpathlineto{\pgfpoint{-6.072196\du}{6.328108\du}}
\pgfpathlineto{\pgfpoint{-6.076973\du}{6.325123\du}}
\pgfpathlineto{\pgfpoint{-6.082048\du}{6.322436\du}}
\pgfpathlineto{\pgfpoint{-6.087422\du}{6.319450\du}}
\pgfpathlineto{\pgfpoint{-6.093692\du}{6.317062\du}}
\pgfpathlineto{\pgfpoint{-6.099663\du}{6.314375\du}}
\pgfpathlineto{\pgfpoint{-6.106529\du}{6.311987\du}}
\pgfpathlineto{\pgfpoint{-6.113396\du}{6.309598\du}}
\pgfpathlineto{\pgfpoint{-6.120860\du}{6.307508\du}}
\pgfpathlineto{\pgfpoint{-6.128324\du}{6.305120\du}}
\pgfpathlineto{\pgfpoint{-6.136385\du}{6.303627\du}}
\pgfpathlineto{\pgfpoint{-6.145043\du}{6.301537\du}}
\pgfpathlineto{\pgfpoint{-6.153402\du}{6.299746\du}}
\pgfpathlineto{\pgfpoint{-6.162060\du}{6.298552\du}}
\pgfpathlineto{\pgfpoint{-6.171315\du}{6.297059\du}}
\pgfpathlineto{\pgfpoint{-6.180570\du}{6.295865\du}}
\pgfpathlineto{\pgfpoint{-6.189825\du}{6.294671\du}}
\pgfpathlineto{\pgfpoint{-6.199379\du}{6.293775\du}}
\pgfpathlineto{\pgfpoint{-6.209530\du}{6.292879\du}}
\pgfpathlineto{\pgfpoint{-6.219084\du}{6.292282\du}}
\pgfpathlineto{\pgfpoint{-6.228637\du}{6.291984\du}}
\pgfpathlineto{\pgfpoint{-6.238788\du}{6.291685\du}}
\pgfpathlineto{\pgfpoint{-6.248640\du}{6.291685\du}}
\pgfpathlineto{\pgfpoint{-6.248640\du}{6.291685\du}}
\pgfpathlineto{\pgfpoint{-6.258492\du}{6.291685\du}}
\pgfpathlineto{\pgfpoint{-6.268643\du}{6.291984\du}}
\pgfpathlineto{\pgfpoint{-6.278495\du}{6.292282\du}}
\pgfpathlineto{\pgfpoint{-6.288348\du}{6.292879\du}}
\pgfpathlineto{\pgfpoint{-6.297901\du}{6.293775\du}}
\pgfpathlineto{\pgfpoint{-6.307455\du}{6.294671\du}}
\pgfpathlineto{\pgfpoint{-6.317009\du}{6.295865\du}}
\pgfpathlineto{\pgfpoint{-6.325965\du}{6.297059\du}}
\pgfpathlineto{\pgfpoint{-6.335220\du}{6.298552\du}}
\pgfpathlineto{\pgfpoint{-6.343878\du}{6.299746\du}}
\pgfpathlineto{\pgfpoint{-6.352536\du}{6.301537\du}}
\pgfpathlineto{\pgfpoint{-6.360896\du}{6.303627\du}}
\pgfpathlineto{\pgfpoint{-6.368957\du}{6.305120\du}}
\pgfpathlineto{\pgfpoint{-6.376421\du}{6.307508\du}}
\pgfpathlineto{\pgfpoint{-6.383884\du}{6.309598\du}}
\pgfpathlineto{\pgfpoint{-6.391050\du}{6.311987\du}}
\pgfpathlineto{\pgfpoint{-6.397618\du}{6.314375\du}}
\pgfpathlineto{\pgfpoint{-6.403887\du}{6.317062\du}}
\pgfpathlineto{\pgfpoint{-6.409858\du}{6.319450\du}}
\pgfpathlineto{\pgfpoint{-6.415232\du}{6.322436\du}}
\pgfpathlineto{\pgfpoint{-6.420308\du}{6.325123\du}}
\pgfpathlineto{\pgfpoint{-6.425084\du}{6.328108\du}}
\pgfpathlineto{\pgfpoint{-6.428966\du}{6.331094\du}}
\pgfpathlineto{\pgfpoint{-6.432847\du}{6.334079\du}}
\pgfpathlineto{\pgfpoint{-6.435832\du}{6.337364\du}}
\pgfpathlineto{\pgfpoint{-6.438519\du}{6.340648\du}}
\pgfpathlineto{\pgfpoint{-6.440908\du}{6.343633\du}}
\pgfpathlineto{\pgfpoint{-6.442699\du}{6.346917\du}}
\pgfpathlineto{\pgfpoint{-6.444192\du}{6.350201\du}}
\pgfpathlineto{\pgfpoint{-6.445087\du}{6.353485\du}}
\pgfpathlineto{\pgfpoint{-6.445087\du}{6.356769\du}}
\pgfpathlineto{\pgfpoint{-6.445087\du}{6.356769\du}}
\pgfpathlineto{\pgfpoint{-6.445087\du}{6.360352\du}}
\pgfpathlineto{\pgfpoint{-6.444192\du}{6.363338\du}}
\pgfpathlineto{\pgfpoint{-6.442699\du}{6.366622\du}}
\pgfpathlineto{\pgfpoint{-6.440908\du}{6.369906\du}}
\pgfpathlineto{\pgfpoint{-6.438519\du}{6.373190\du}}
\pgfpathlineto{\pgfpoint{-6.435832\du}{6.376474\du}}
\pgfpathlineto{\pgfpoint{-6.432847\du}{6.379459\du}}
\pgfpathlineto{\pgfpoint{-6.428966\du}{6.382445\du}}
\pgfpathlineto{\pgfpoint{-6.425084\du}{6.385430\du}}
\pgfpathlineto{\pgfpoint{-6.420308\du}{6.388714\du}}
\pgfpathlineto{\pgfpoint{-6.415232\du}{6.391401\du}}
\pgfpathlineto{\pgfpoint{-6.409858\du}{6.394088\du}}
\pgfpathlineto{\pgfpoint{-6.403887\du}{6.396775\du}}
\pgfpathlineto{\pgfpoint{-6.397618\du}{6.399462\du}}
\pgfpathlineto{\pgfpoint{-6.391050\du}{6.401851\du}}
\pgfpathlineto{\pgfpoint{-6.383884\du}{6.404239\du}}
\pgfpathlineto{\pgfpoint{-6.376421\du}{6.406030\du}}
\pgfpathlineto{\pgfpoint{-6.368957\du}{6.408120\du}}
\pgfpathlineto{\pgfpoint{-6.360896\du}{6.410210\du}}
\pgfpathlineto{\pgfpoint{-6.352536\du}{6.412300\du}}
\pgfpathlineto{\pgfpoint{-6.343878\du}{6.413793\du}}
\pgfpathlineto{\pgfpoint{-6.335220\du}{6.415286\du}}
\pgfpathlineto{\pgfpoint{-6.325965\du}{6.416778\du}}
\pgfpathlineto{\pgfpoint{-6.317009\du}{6.417973\du}}
\pgfpathlineto{\pgfpoint{-6.307455\du}{6.418868\du}}
\pgfpathlineto{\pgfpoint{-6.297901\du}{6.420062\du}}
\pgfpathlineto{\pgfpoint{-6.288348\du}{6.420660\du}}
\pgfpathlineto{\pgfpoint{-6.278495\du}{6.421257\du}}
\pgfpathlineto{\pgfpoint{-6.268643\du}{6.421854\du}}
\pgfpathlineto{\pgfpoint{-6.258492\du}{6.422152\du}}
\pgfpathlineto{\pgfpoint{-6.248640\du}{6.422152\du}}
\pgfpathlineto{\pgfpoint{-6.248640\du}{6.422152\du}}
\pgfpathlineto{\pgfpoint{-6.238788\du}{6.422152\du}}
\pgfpathlineto{\pgfpoint{-6.228637\du}{6.421854\du}}
\pgfpathlineto{\pgfpoint{-6.219084\du}{6.421257\du}}
\pgfpathlineto{\pgfpoint{-6.209530\du}{6.420660\du}}
\pgfpathlineto{\pgfpoint{-6.199379\du}{6.420062\du}}
\pgfpathlineto{\pgfpoint{-6.189825\du}{6.418868\du}}
\pgfpathlineto{\pgfpoint{-6.180570\du}{6.417973\du}}
\pgfpathlineto{\pgfpoint{-6.171315\du}{6.416778\du}}
\pgfpathlineto{\pgfpoint{-6.162060\du}{6.415286\du}}
\pgfpathlineto{\pgfpoint{-6.153402\du}{6.413793\du}}
\pgfpathlineto{\pgfpoint{-6.145043\du}{6.412300\du}}
\pgfpathlineto{\pgfpoint{-6.136385\du}{6.410210\du}}
\pgfpathlineto{\pgfpoint{-6.128324\du}{6.408120\du}}
\pgfpathlineto{\pgfpoint{-6.120860\du}{6.406030\du}}
\pgfpathlineto{\pgfpoint{-6.113396\du}{6.404239\du}}
\pgfpathlineto{\pgfpoint{-6.106529\du}{6.401851\du}}
\pgfpathlineto{\pgfpoint{-6.099663\du}{6.399462\du}}
\pgfpathlineto{\pgfpoint{-6.093692\du}{6.396775\du}}
\pgfpathlineto{\pgfpoint{-6.087422\du}{6.394088\du}}
\pgfpathlineto{\pgfpoint{-6.082048\du}{6.391401\du}}
\pgfpathlineto{\pgfpoint{-6.076973\du}{6.388714\du}}
\pgfpathlineto{\pgfpoint{-6.072196\du}{6.385430\du}}
\pgfpathlineto{\pgfpoint{-6.068613\du}{6.382445\du}}
\pgfpathlineto{\pgfpoint{-6.064434\du}{6.379459\du}}
\pgfpathlineto{\pgfpoint{-6.061448\du}{6.376474\du}}
\pgfpathlineto{\pgfpoint{-6.058761\du}{6.373190\du}}
\pgfpathlineto{\pgfpoint{-6.056373\du}{6.369906\du}}
\pgfpathlineto{\pgfpoint{-6.054880\du}{6.366622\du}}
\pgfpathlineto{\pgfpoint{-6.053686\du}{6.363338\du}}
\pgfpathlineto{\pgfpoint{-6.052790\du}{6.360352\du}}
\pgfpathlineto{\pgfpoint{-6.052790\du}{6.356769\du}}
\pgfusepath{fill}
\pgfsetbuttcap
\pgfsetmiterjoin
\pgfsetdash{}{0pt}
\definecolor{dialinecolor}{rgb}{0.666667, 0.901961, 1.000000}
\pgfsetstrokecolor{dialinecolor}
\pgfpathmoveto{\pgfpoint{-6.064434\du}{6.351097\du}}
\pgfpathlineto{\pgfpoint{-6.064434\du}{6.347514\du}}
\pgfpathlineto{\pgfpoint{-6.065329\du}{6.345126\du}}
\pgfpathlineto{\pgfpoint{-6.066524\du}{6.341842\du}}
\pgfpathlineto{\pgfpoint{-6.068613\du}{6.338856\du}}
\pgfpathlineto{\pgfpoint{-6.070703\du}{6.336169\du}}
\pgfpathlineto{\pgfpoint{-6.072793\du}{6.332587\du}}
\pgfpathlineto{\pgfpoint{-6.076077\du}{6.330198\du}}
\pgfpathlineto{\pgfpoint{-6.079361\du}{6.327511\du}}
\pgfpathlineto{\pgfpoint{-6.083541\du}{6.324824\du}}
\pgfpathlineto{\pgfpoint{-6.088318\du}{6.322436\du}}
\pgfpathlineto{\pgfpoint{-6.092796\du}{6.319450\du}}
\pgfpathlineto{\pgfpoint{-6.098469\du}{6.317361\du}}
\pgfpathlineto{\pgfpoint{-6.104440\du}{6.314674\du}}
\pgfpathlineto{\pgfpoint{-6.110112\du}{6.311987\du}}
\pgfpathlineto{\pgfpoint{-6.116680\du}{6.310195\du}}
\pgfpathlineto{\pgfpoint{-6.123547\du}{6.307807\du}}
\pgfpathlineto{\pgfpoint{-6.130414\du}{6.306314\du}}
\pgfpathlineto{\pgfpoint{-6.138176\du}{6.304224\du}}
\pgfpathlineto{\pgfpoint{-6.145640\du}{6.302433\du}}
\pgfpathlineto{\pgfpoint{-6.153999\du}{6.300642\du}}
\pgfpathlineto{\pgfpoint{-6.162657\du}{6.299447\du}}
\pgfpathlineto{\pgfpoint{-6.171017\du}{6.297059\du}}
\pgfpathlineto{\pgfpoint{-6.179675\du}{6.296760\du}}
\pgfpathlineto{\pgfpoint{-6.188631\du}{6.295268\du}}
\pgfpathlineto{\pgfpoint{-6.197289\du}{6.294372\du}}
\pgfpathlineto{\pgfpoint{-6.206843\du}{6.293476\du}}
\pgfpathlineto{\pgfpoint{-6.216098\du}{6.292879\du}}
\pgfpathlineto{\pgfpoint{-6.225652\du}{6.292581\du}}
\pgfpathlineto{\pgfpoint{-6.235205\du}{6.291984\du}}
\pgfpathlineto{\pgfpoint{-6.245356\du}{6.291685\du}}
\pgfpathlineto{\pgfpoint{-6.255208\du}{6.291685\du}}
\pgfpathlineto{\pgfpoint{-6.255208\du}{6.291685\du}}
\pgfpathlineto{\pgfpoint{-6.264165\du}{6.291685\du}}
\pgfpathlineto{\pgfpoint{-6.274017\du}{6.291984\du}}
\pgfpathlineto{\pgfpoint{-6.283571\du}{6.292581\du}}
\pgfpathlineto{\pgfpoint{-6.292826\du}{6.292879\du}}
\pgfpathlineto{\pgfpoint{-6.302380\du}{6.293476\du}}
\pgfpathlineto{\pgfpoint{-6.311933\du}{6.294372\du}}
\pgfpathlineto{\pgfpoint{-6.320890\du}{6.295268\du}}
\pgfpathlineto{\pgfpoint{-6.329846\du}{6.296760\du}}
\pgfpathlineto{\pgfpoint{-6.338504\du}{6.297059\du}}
\pgfpathlineto{\pgfpoint{-6.346864\du}{6.299447\du}}
\pgfpathlineto{\pgfpoint{-6.355522\du}{6.300642\du}}
\pgfpathlineto{\pgfpoint{-6.363583\du}{6.302433\du}}
\pgfpathlineto{\pgfpoint{-6.371345\du}{6.304224\du}}
\pgfpathlineto{\pgfpoint{-6.378809\du}{6.306314\du}}
\pgfpathlineto{\pgfpoint{-6.385377\du}{6.307807\du}}
\pgfpathlineto{\pgfpoint{-6.392841\du}{6.310195\du}}
\pgfpathlineto{\pgfpoint{-6.399110\du}{6.311987\du}}
\pgfpathlineto{\pgfpoint{-6.405082\du}{6.314674\du}}
\pgfpathlineto{\pgfpoint{-6.411053\du}{6.317361\du}}
\pgfpathlineto{\pgfpoint{-6.416128\du}{6.319450\du}}
\pgfpathlineto{\pgfpoint{-6.421203\du}{6.322436\du}}
\pgfpathlineto{\pgfpoint{-6.425980\du}{6.324824\du}}
\pgfpathlineto{\pgfpoint{-6.430160\du}{6.327511\du}}
\pgfpathlineto{\pgfpoint{-6.433444\du}{6.330198\du}}
\pgfpathlineto{\pgfpoint{-6.436728\du}{6.332587\du}}
\pgfpathlineto{\pgfpoint{-6.439116\du}{6.336169\du}}
\pgfpathlineto{\pgfpoint{-6.440908\du}{6.338856\du}}
\pgfpathlineto{\pgfpoint{-6.443296\du}{6.341842\du}}
\pgfpathlineto{\pgfpoint{-6.444192\du}{6.345126\du}}
\pgfpathlineto{\pgfpoint{-6.445087\du}{6.347514\du}}
\pgfpathlineto{\pgfpoint{-6.445087\du}{6.351097\du}}
\pgfpathlineto{\pgfpoint{-6.445087\du}{6.351097\du}}
\pgfpathlineto{\pgfpoint{-6.445087\du}{6.354082\du}}
\pgfpathlineto{\pgfpoint{-6.444192\du}{6.356769\du}}
\pgfpathlineto{\pgfpoint{-6.443296\du}{6.359755\du}}
\pgfpathlineto{\pgfpoint{-6.440908\du}{6.363338\du}}
\pgfpathlineto{\pgfpoint{-6.439116\du}{6.366024\du}}
\pgfpathlineto{\pgfpoint{-6.436728\du}{6.369010\du}}
\pgfpathlineto{\pgfpoint{-6.433444\du}{6.371697\du}}
\pgfpathlineto{\pgfpoint{-6.430160\du}{6.374384\du}}
\pgfpathlineto{\pgfpoint{-6.425980\du}{6.377071\du}}
\pgfpathlineto{\pgfpoint{-6.421203\du}{6.379758\du}}
\pgfpathlineto{\pgfpoint{-6.416128\du}{6.382445\du}}
\pgfpathlineto{\pgfpoint{-6.411053\du}{6.384833\du}}
\pgfpathlineto{\pgfpoint{-6.405082\du}{6.387520\du}}
\pgfpathlineto{\pgfpoint{-6.399110\du}{6.389610\du}}
\pgfpathlineto{\pgfpoint{-6.392841\du}{6.391999\du}}
\pgfpathlineto{\pgfpoint{-6.385377\du}{6.394387\du}}
\pgfpathlineto{\pgfpoint{-6.378809\du}{6.396178\du}}
\pgfpathlineto{\pgfpoint{-6.371345\du}{6.397671\du}}
\pgfpathlineto{\pgfpoint{-6.363583\du}{6.399761\du}}
\pgfpathlineto{\pgfpoint{-6.355522\du}{6.401552\du}}
\pgfpathlineto{\pgfpoint{-6.346864\du}{6.402746\du}}
\pgfpathlineto{\pgfpoint{-6.338504\du}{6.404538\du}}
\pgfpathlineto{\pgfpoint{-6.329846\du}{6.405433\du}}
\pgfpathlineto{\pgfpoint{-6.320890\du}{6.406926\du}}
\pgfpathlineto{\pgfpoint{-6.311933\du}{6.407822\du}}
\pgfpathlineto{\pgfpoint{-6.302380\du}{6.408120\du}}
\pgfpathlineto{\pgfpoint{-6.292826\du}{6.409315\du}}
\pgfpathlineto{\pgfpoint{-6.283571\du}{6.409613\du}}
\pgfpathlineto{\pgfpoint{-6.274017\du}{6.409912\du}}
\pgfpathlineto{\pgfpoint{-6.264165\du}{6.410509\du}}
\pgfpathlineto{\pgfpoint{-6.255208\du}{6.410509\du}}
\pgfpathlineto{\pgfpoint{-6.255208\du}{6.410509\du}}
\pgfpathlineto{\pgfpoint{-6.245356\du}{6.410509\du}}
\pgfpathlineto{\pgfpoint{-6.235205\du}{6.409912\du}}
\pgfpathlineto{\pgfpoint{-6.225652\du}{6.409613\du}}
\pgfpathlineto{\pgfpoint{-6.216098\du}{6.409315\du}}
\pgfpathlineto{\pgfpoint{-6.206843\du}{6.408120\du}}
\pgfpathlineto{\pgfpoint{-6.197289\du}{6.407822\du}}
\pgfpathlineto{\pgfpoint{-6.188631\du}{6.406926\du}}
\pgfpathlineto{\pgfpoint{-6.179675\du}{6.405433\du}}
\pgfpathlineto{\pgfpoint{-6.171017\du}{6.404538\du}}
\pgfpathlineto{\pgfpoint{-6.162657\du}{6.402746\du}}
\pgfpathlineto{\pgfpoint{-6.153999\du}{6.401552\du}}
\pgfpathlineto{\pgfpoint{-6.145640\du}{6.399761\du}}
\pgfpathlineto{\pgfpoint{-6.138176\du}{6.397671\du}}
\pgfpathlineto{\pgfpoint{-6.130414\du}{6.396178\du}}
\pgfpathlineto{\pgfpoint{-6.123547\du}{6.394387\du}}
\pgfpathlineto{\pgfpoint{-6.116680\du}{6.391999\du}}
\pgfpathlineto{\pgfpoint{-6.110112\du}{6.389610\du}}
\pgfpathlineto{\pgfpoint{-6.104440\du}{6.387520\du}}
\pgfpathlineto{\pgfpoint{-6.098469\du}{6.384833\du}}
\pgfpathlineto{\pgfpoint{-6.092796\du}{6.382445\du}}
\pgfpathlineto{\pgfpoint{-6.088318\du}{6.379758\du}}
\pgfpathlineto{\pgfpoint{-6.083541\du}{6.377071\du}}
\pgfpathlineto{\pgfpoint{-6.079361\du}{6.374384\du}}
\pgfpathlineto{\pgfpoint{-6.076077\du}{6.371697\du}}
\pgfpathlineto{\pgfpoint{-6.072793\du}{6.369010\du}}
\pgfpathlineto{\pgfpoint{-6.070703\du}{6.366024\du}}
\pgfpathlineto{\pgfpoint{-6.068613\du}{6.363338\du}}
\pgfpathlineto{\pgfpoint{-6.066524\du}{6.359755\du}}
\pgfpathlineto{\pgfpoint{-6.065329\du}{6.356769\du}}
\pgfpathlineto{\pgfpoint{-6.064434\du}{6.354082\du}}
\pgfpathlineto{\pgfpoint{-6.064434\du}{6.351097\du}}
\pgfusepath{stroke}
\pgfsetbuttcap
\pgfsetmiterjoin
\pgfsetdash{}{0pt}
\definecolor{dialinecolor}{rgb}{0.000000, 0.470588, 0.666667}
\pgfsetfillcolor{dialinecolor}
\pgfpathmoveto{\pgfpoint{-6.445087\du}{6.267801\du}}
\pgfpathlineto{\pgfpoint{-6.445087\du}{6.362442\du}}
\pgfpathlineto{\pgfpoint{-6.052193\du}{6.362442\du}}
\pgfpathlineto{\pgfpoint{-6.052193\du}{6.267801\du}}
\pgfpathlineto{\pgfpoint{-6.445087\du}{6.267801\du}}
\pgfusepath{fill}
\pgfsetbuttcap
\pgfsetmiterjoin
\pgfsetdash{}{0pt}
\definecolor{dialinecolor}{rgb}{0.000000, 0.705882, 1.000000}
\pgfsetfillcolor{dialinecolor}
\pgfpathmoveto{\pgfpoint{-6.052790\du}{6.261233\du}}
\pgfpathlineto{\pgfpoint{-6.052790\du}{6.257949\du}}
\pgfpathlineto{\pgfpoint{-6.053686\du}{6.254366\du}}
\pgfpathlineto{\pgfpoint{-6.054880\du}{6.251381\du}}
\pgfpathlineto{\pgfpoint{-6.056373\du}{6.248096\du}}
\pgfpathlineto{\pgfpoint{-6.058761\du}{6.244812\du}}
\pgfpathlineto{\pgfpoint{-6.061448\du}{6.241528\du}}
\pgfpathlineto{\pgfpoint{-6.064434\du}{6.238543\du}}
\pgfpathlineto{\pgfpoint{-6.068613\du}{6.235259\du}}
\pgfpathlineto{\pgfpoint{-6.072196\du}{6.232572\du}}
\pgfpathlineto{\pgfpoint{-6.076973\du}{6.229586\du}}
\pgfpathlineto{\pgfpoint{-6.082048\du}{6.226601\du}}
\pgfpathlineto{\pgfpoint{-6.087422\du}{6.223914\du}}
\pgfpathlineto{\pgfpoint{-6.093692\du}{6.220630\du}}
\pgfpathlineto{\pgfpoint{-6.099663\du}{6.218838\du}}
\pgfpathlineto{\pgfpoint{-6.106529\du}{6.216151\du}}
\pgfpathlineto{\pgfpoint{-6.113396\du}{6.213763\du}}
\pgfpathlineto{\pgfpoint{-6.120860\du}{6.211673\du}}
\pgfpathlineto{\pgfpoint{-6.128324\du}{6.209583\du}}
\pgfpathlineto{\pgfpoint{-6.136385\du}{6.207493\du}}
\pgfpathlineto{\pgfpoint{-6.145043\du}{6.205404\du}}
\pgfpathlineto{\pgfpoint{-6.153402\du}{6.203911\du}}
\pgfpathlineto{\pgfpoint{-6.162060\du}{6.202418\du}}
\pgfpathlineto{\pgfpoint{-6.171315\du}{6.201224\du}}
\pgfpathlineto{\pgfpoint{-6.180570\du}{6.199731\du}}
\pgfpathlineto{\pgfpoint{-6.189825\du}{6.198835\du}}
\pgfpathlineto{\pgfpoint{-6.199379\du}{6.197641\du}}
\pgfpathlineto{\pgfpoint{-6.209530\du}{6.197044\du}}
\pgfpathlineto{\pgfpoint{-6.219084\du}{6.196447\du}}
\pgfpathlineto{\pgfpoint{-6.228637\du}{6.196148\du}}
\pgfpathlineto{\pgfpoint{-6.238788\du}{6.195850\du}}
\pgfpathlineto{\pgfpoint{-6.248640\du}{6.195850\du}}
\pgfpathlineto{\pgfpoint{-6.248640\du}{6.195850\du}}
\pgfpathlineto{\pgfpoint{-6.258492\du}{6.195850\du}}
\pgfpathlineto{\pgfpoint{-6.268643\du}{6.196148\du}}
\pgfpathlineto{\pgfpoint{-6.278495\du}{6.196447\du}}
\pgfpathlineto{\pgfpoint{-6.288348\du}{6.197044\du}}
\pgfpathlineto{\pgfpoint{-6.297901\du}{6.197641\du}}
\pgfpathlineto{\pgfpoint{-6.307455\du}{6.198835\du}}
\pgfpathlineto{\pgfpoint{-6.317009\du}{6.199731\du}}
\pgfpathlineto{\pgfpoint{-6.325965\du}{6.201224\du}}
\pgfpathlineto{\pgfpoint{-6.335220\du}{6.202418\du}}
\pgfpathlineto{\pgfpoint{-6.343878\du}{6.203911\du}}
\pgfpathlineto{\pgfpoint{-6.352536\du}{6.205404\du}}
\pgfpathlineto{\pgfpoint{-6.360896\du}{6.207493\du}}
\pgfpathlineto{\pgfpoint{-6.368957\du}{6.209583\du}}
\pgfpathlineto{\pgfpoint{-6.376421\du}{6.211673\du}}
\pgfpathlineto{\pgfpoint{-6.383884\du}{6.213763\du}}
\pgfpathlineto{\pgfpoint{-6.391050\du}{6.216151\du}}
\pgfpathlineto{\pgfpoint{-6.397618\du}{6.218838\du}}
\pgfpathlineto{\pgfpoint{-6.403887\du}{6.220630\du}}
\pgfpathlineto{\pgfpoint{-6.409858\du}{6.223914\du}}
\pgfpathlineto{\pgfpoint{-6.415232\du}{6.226601\du}}
\pgfpathlineto{\pgfpoint{-6.420308\du}{6.229586\du}}
\pgfpathlineto{\pgfpoint{-6.425084\du}{6.232572\du}}
\pgfpathlineto{\pgfpoint{-6.428966\du}{6.235259\du}}
\pgfpathlineto{\pgfpoint{-6.432847\du}{6.238543\du}}
\pgfpathlineto{\pgfpoint{-6.435832\du}{6.241528\du}}
\pgfpathlineto{\pgfpoint{-6.438519\du}{6.244812\du}}
\pgfpathlineto{\pgfpoint{-6.440908\du}{6.248096\du}}
\pgfpathlineto{\pgfpoint{-6.442699\du}{6.251381\du}}
\pgfpathlineto{\pgfpoint{-6.444192\du}{6.254366\du}}
\pgfpathlineto{\pgfpoint{-6.445087\du}{6.257949\du}}
\pgfpathlineto{\pgfpoint{-6.445087\du}{6.261233\du}}
\pgfpathlineto{\pgfpoint{-6.445087\du}{6.261233\du}}
\pgfpathlineto{\pgfpoint{-6.445087\du}{6.264517\du}}
\pgfpathlineto{\pgfpoint{-6.444192\du}{6.268099\du}}
\pgfpathlineto{\pgfpoint{-6.442699\du}{6.271085\du}}
\pgfpathlineto{\pgfpoint{-6.440908\du}{6.274369\du}}
\pgfpathlineto{\pgfpoint{-6.438519\du}{6.277952\du}}
\pgfpathlineto{\pgfpoint{-6.435832\du}{6.280937\du}}
\pgfpathlineto{\pgfpoint{-6.432847\du}{6.283923\du}}
\pgfpathlineto{\pgfpoint{-6.428966\du}{6.287207\du}}
\pgfpathlineto{\pgfpoint{-6.425084\du}{6.289595\du}}
\pgfpathlineto{\pgfpoint{-6.420308\du}{6.292879\du}}
\pgfpathlineto{\pgfpoint{-6.415232\du}{6.296163\du}}
\pgfpathlineto{\pgfpoint{-6.409858\du}{6.298850\du}}
\pgfpathlineto{\pgfpoint{-6.403887\du}{6.301537\du}}
\pgfpathlineto{\pgfpoint{-6.397618\du}{6.303627\du}}
\pgfpathlineto{\pgfpoint{-6.391050\du}{6.306613\du}}
\pgfpathlineto{\pgfpoint{-6.383884\du}{6.309001\du}}
\pgfpathlineto{\pgfpoint{-6.376421\du}{6.310792\du}}
\pgfpathlineto{\pgfpoint{-6.368957\du}{6.312882\du}}
\pgfpathlineto{\pgfpoint{-6.360896\du}{6.314972\du}}
\pgfpathlineto{\pgfpoint{-6.352536\du}{6.317062\du}}
\pgfpathlineto{\pgfpoint{-6.343878\du}{6.318555\du}}
\pgfpathlineto{\pgfpoint{-6.335220\du}{6.320047\du}}
\pgfpathlineto{\pgfpoint{-6.325965\du}{6.321242\du}}
\pgfpathlineto{\pgfpoint{-6.317009\du}{6.322734\du}}
\pgfpathlineto{\pgfpoint{-6.307455\du}{6.323630\du}}
\pgfpathlineto{\pgfpoint{-6.297901\du}{6.324824\du}}
\pgfpathlineto{\pgfpoint{-6.288348\du}{6.325123\du}}
\pgfpathlineto{\pgfpoint{-6.278495\du}{6.326019\du}}
\pgfpathlineto{\pgfpoint{-6.268643\du}{6.326019\du}}
\pgfpathlineto{\pgfpoint{-6.258492\du}{6.326914\du}}
\pgfpathlineto{\pgfpoint{-6.248640\du}{6.326914\du}}
\pgfpathlineto{\pgfpoint{-6.248640\du}{6.326914\du}}
\pgfpathlineto{\pgfpoint{-6.238788\du}{6.326914\du}}
\pgfpathlineto{\pgfpoint{-6.228637\du}{6.326019\du}}
\pgfpathlineto{\pgfpoint{-6.219084\du}{6.326019\du}}
\pgfpathlineto{\pgfpoint{-6.209530\du}{6.325123\du}}
\pgfpathlineto{\pgfpoint{-6.199379\du}{6.324824\du}}
\pgfpathlineto{\pgfpoint{-6.189825\du}{6.323630\du}}
\pgfpathlineto{\pgfpoint{-6.180570\du}{6.322734\du}}
\pgfpathlineto{\pgfpoint{-6.171315\du}{6.321242\du}}
\pgfpathlineto{\pgfpoint{-6.162060\du}{6.320047\du}}
\pgfpathlineto{\pgfpoint{-6.153402\du}{6.318555\du}}
\pgfpathlineto{\pgfpoint{-6.145043\du}{6.317062\du}}
\pgfpathlineto{\pgfpoint{-6.136385\du}{6.314972\du}}
\pgfpathlineto{\pgfpoint{-6.128324\du}{6.312882\du}}
\pgfpathlineto{\pgfpoint{-6.120860\du}{6.310792\du}}
\pgfpathlineto{\pgfpoint{-6.113396\du}{6.309001\du}}
\pgfpathlineto{\pgfpoint{-6.106529\du}{6.306613\du}}
\pgfpathlineto{\pgfpoint{-6.099663\du}{6.303627\du}}
\pgfpathlineto{\pgfpoint{-6.093692\du}{6.301537\du}}
\pgfpathlineto{\pgfpoint{-6.087422\du}{6.298850\du}}
\pgfpathlineto{\pgfpoint{-6.082048\du}{6.296163\du}}
\pgfpathlineto{\pgfpoint{-6.076973\du}{6.292879\du}}
\pgfpathlineto{\pgfpoint{-6.072196\du}{6.289595\du}}
\pgfpathlineto{\pgfpoint{-6.068613\du}{6.287207\du}}
\pgfpathlineto{\pgfpoint{-6.064434\du}{6.283923\du}}
\pgfpathlineto{\pgfpoint{-6.061448\du}{6.280937\du}}
\pgfpathlineto{\pgfpoint{-6.058761\du}{6.277952\du}}
\pgfpathlineto{\pgfpoint{-6.056373\du}{6.274369\du}}
\pgfpathlineto{\pgfpoint{-6.054880\du}{6.271085\du}}
\pgfpathlineto{\pgfpoint{-6.053686\du}{6.268099\du}}
\pgfpathlineto{\pgfpoint{-6.052790\du}{6.264517\du}}
\pgfpathlineto{\pgfpoint{-6.052790\du}{6.261233\du}}
\pgfusepath{fill}
\pgfsetbuttcap
\pgfsetmiterjoin
\pgfsetdash{}{0pt}
\definecolor{dialinecolor}{rgb}{0.666667, 0.901961, 1.000000}
\pgfsetstrokecolor{dialinecolor}
\pgfpathmoveto{\pgfpoint{-6.064434\du}{6.255560\du}}
\pgfpathlineto{\pgfpoint{-6.064434\du}{6.252276\du}}
\pgfpathlineto{\pgfpoint{-6.065329\du}{6.249888\du}}
\pgfpathlineto{\pgfpoint{-6.066524\du}{6.246604\du}}
\pgfpathlineto{\pgfpoint{-6.068613\du}{6.243021\du}}
\pgfpathlineto{\pgfpoint{-6.070703\du}{6.240633\du}}
\pgfpathlineto{\pgfpoint{-6.072793\du}{6.237349\du}}
\pgfpathlineto{\pgfpoint{-6.076077\du}{6.234960\du}}
\pgfpathlineto{\pgfpoint{-6.079361\du}{6.232273\du}}
\pgfpathlineto{\pgfpoint{-6.083541\du}{6.229586\du}}
\pgfpathlineto{\pgfpoint{-6.088318\du}{6.226899\du}}
\pgfpathlineto{\pgfpoint{-6.092796\du}{6.224212\du}}
\pgfpathlineto{\pgfpoint{-6.098469\du}{6.221525\du}}
\pgfpathlineto{\pgfpoint{-6.104440\du}{6.219137\du}}
\pgfpathlineto{\pgfpoint{-6.110112\du}{6.217047\du}}
\pgfpathlineto{\pgfpoint{-6.116680\du}{6.214659\du}}
\pgfpathlineto{\pgfpoint{-6.123547\du}{6.212270\du}}
\pgfpathlineto{\pgfpoint{-6.130414\du}{6.210479\du}}
\pgfpathlineto{\pgfpoint{-6.138176\du}{6.208986\du}}
\pgfpathlineto{\pgfpoint{-6.145640\du}{6.206598\du}}
\pgfpathlineto{\pgfpoint{-6.153999\du}{6.205105\du}}
\pgfpathlineto{\pgfpoint{-6.162657\du}{6.203911\du}}
\pgfpathlineto{\pgfpoint{-6.171017\du}{6.202119\du}}
\pgfpathlineto{\pgfpoint{-6.179675\du}{6.201224\du}}
\pgfpathlineto{\pgfpoint{-6.188631\du}{6.199731\du}}
\pgfpathlineto{\pgfpoint{-6.197289\du}{6.198835\du}}
\pgfpathlineto{\pgfpoint{-6.206843\du}{6.198238\du}}
\pgfpathlineto{\pgfpoint{-6.216098\du}{6.197343\du}}
\pgfpathlineto{\pgfpoint{-6.225652\du}{6.197044\du}}
\pgfpathlineto{\pgfpoint{-6.235205\du}{6.196745\du}}
\pgfpathlineto{\pgfpoint{-6.245356\du}{6.196148\du}}
\pgfpathlineto{\pgfpoint{-6.255208\du}{6.196148\du}}
\pgfpathlineto{\pgfpoint{-6.255208\du}{6.196148\du}}
\pgfpathlineto{\pgfpoint{-6.264165\du}{6.196148\du}}
\pgfpathlineto{\pgfpoint{-6.274017\du}{6.196745\du}}
\pgfpathlineto{\pgfpoint{-6.283571\du}{6.197044\du}}
\pgfpathlineto{\pgfpoint{-6.292826\du}{6.197343\du}}
\pgfpathlineto{\pgfpoint{-6.302380\du}{6.198238\du}}
\pgfpathlineto{\pgfpoint{-6.311933\du}{6.198835\du}}
\pgfpathlineto{\pgfpoint{-6.320890\du}{6.199731\du}}
\pgfpathlineto{\pgfpoint{-6.329846\du}{6.201224\du}}
\pgfpathlineto{\pgfpoint{-6.338504\du}{6.202119\du}}
\pgfpathlineto{\pgfpoint{-6.346864\du}{6.203911\du}}
\pgfpathlineto{\pgfpoint{-6.355522\du}{6.205105\du}}
\pgfpathlineto{\pgfpoint{-6.363583\du}{6.206598\du}}
\pgfpathlineto{\pgfpoint{-6.371345\du}{6.208986\du}}
\pgfpathlineto{\pgfpoint{-6.378809\du}{6.210479\du}}
\pgfpathlineto{\pgfpoint{-6.385377\du}{6.212270\du}}
\pgfpathlineto{\pgfpoint{-6.392841\du}{6.214659\du}}
\pgfpathlineto{\pgfpoint{-6.399110\du}{6.217047\du}}
\pgfpathlineto{\pgfpoint{-6.405082\du}{6.219137\du}}
\pgfpathlineto{\pgfpoint{-6.411053\du}{6.221525\du}}
\pgfpathlineto{\pgfpoint{-6.416128\du}{6.224212\du}}
\pgfpathlineto{\pgfpoint{-6.421203\du}{6.226899\du}}
\pgfpathlineto{\pgfpoint{-6.425980\du}{6.229586\du}}
\pgfpathlineto{\pgfpoint{-6.430160\du}{6.232273\du}}
\pgfpathlineto{\pgfpoint{-6.433444\du}{6.234960\du}}
\pgfpathlineto{\pgfpoint{-6.436728\du}{6.237349\du}}
\pgfpathlineto{\pgfpoint{-6.439116\du}{6.240633\du}}
\pgfpathlineto{\pgfpoint{-6.440908\du}{6.243021\du}}
\pgfpathlineto{\pgfpoint{-6.443296\du}{6.246604\du}}
\pgfpathlineto{\pgfpoint{-6.444192\du}{6.249888\du}}
\pgfpathlineto{\pgfpoint{-6.445087\du}{6.252276\du}}
\pgfpathlineto{\pgfpoint{-6.445087\du}{6.255560\du}}
\pgfpathlineto{\pgfpoint{-6.445087\du}{6.255560\du}}
\pgfpathlineto{\pgfpoint{-6.445087\du}{6.258844\du}}
\pgfpathlineto{\pgfpoint{-6.444192\du}{6.261531\du}}
\pgfpathlineto{\pgfpoint{-6.443296\du}{6.264517\du}}
\pgfpathlineto{\pgfpoint{-6.440908\du}{6.267801\du}}
\pgfpathlineto{\pgfpoint{-6.439116\du}{6.270488\du}}
\pgfpathlineto{\pgfpoint{-6.436728\du}{6.273772\du}}
\pgfpathlineto{\pgfpoint{-6.433444\du}{6.276459\du}}
\pgfpathlineto{\pgfpoint{-6.430160\du}{6.279146\du}}
\pgfpathlineto{\pgfpoint{-6.425980\du}{6.281833\du}}
\pgfpathlineto{\pgfpoint{-6.421203\du}{6.284520\du}}
\pgfpathlineto{\pgfpoint{-6.416128\du}{6.287207\du}}
\pgfpathlineto{\pgfpoint{-6.411053\du}{6.289297\du}}
\pgfpathlineto{\pgfpoint{-6.405082\du}{6.291984\du}}
\pgfpathlineto{\pgfpoint{-6.399110\du}{6.294372\du}}
\pgfpathlineto{\pgfpoint{-6.392841\du}{6.296760\du}}
\pgfpathlineto{\pgfpoint{-6.385377\du}{6.298850\du}}
\pgfpathlineto{\pgfpoint{-6.378809\du}{6.300642\du}}
\pgfpathlineto{\pgfpoint{-6.371345\du}{6.302433\du}}
\pgfpathlineto{\pgfpoint{-6.363583\du}{6.304224\du}}
\pgfpathlineto{\pgfpoint{-6.355522\du}{6.306314\du}}
\pgfpathlineto{\pgfpoint{-6.346864\du}{6.307508\du}}
\pgfpathlineto{\pgfpoint{-6.338504\du}{6.309300\du}}
\pgfpathlineto{\pgfpoint{-6.329846\du}{6.310195\du}}
\pgfpathlineto{\pgfpoint{-6.320890\du}{6.311688\du}}
\pgfpathlineto{\pgfpoint{-6.311933\du}{6.311987\du}}
\pgfpathlineto{\pgfpoint{-6.302380\du}{6.313181\du}}
\pgfpathlineto{\pgfpoint{-6.292826\du}{6.313479\du}}
\pgfpathlineto{\pgfpoint{-6.283571\du}{6.314375\du}}
\pgfpathlineto{\pgfpoint{-6.274017\du}{6.314674\du}}
\pgfpathlineto{\pgfpoint{-6.264165\du}{6.314972\du}}
\pgfpathlineto{\pgfpoint{-6.255208\du}{6.314972\du}}
\pgfpathlineto{\pgfpoint{-6.255208\du}{6.314972\du}}
\pgfpathlineto{\pgfpoint{-6.245356\du}{6.314972\du}}
\pgfpathlineto{\pgfpoint{-6.235205\du}{6.314674\du}}
\pgfpathlineto{\pgfpoint{-6.225652\du}{6.314375\du}}
\pgfpathlineto{\pgfpoint{-6.216098\du}{6.313479\du}}
\pgfpathlineto{\pgfpoint{-6.206843\du}{6.313181\du}}
\pgfpathlineto{\pgfpoint{-6.197289\du}{6.311987\du}}
\pgfpathlineto{\pgfpoint{-6.188631\du}{6.311688\du}}
\pgfpathlineto{\pgfpoint{-6.179675\du}{6.310195\du}}
\pgfpathlineto{\pgfpoint{-6.171017\du}{6.309300\du}}
\pgfpathlineto{\pgfpoint{-6.162657\du}{6.307508\du}}
\pgfpathlineto{\pgfpoint{-6.153999\du}{6.306314\du}}
\pgfpathlineto{\pgfpoint{-6.145640\du}{6.304224\du}}
\pgfpathlineto{\pgfpoint{-6.138176\du}{6.302433\du}}
\pgfpathlineto{\pgfpoint{-6.130414\du}{6.300642\du}}
\pgfpathlineto{\pgfpoint{-6.123547\du}{6.298850\du}}
\pgfpathlineto{\pgfpoint{-6.116680\du}{6.296760\du}}
\pgfpathlineto{\pgfpoint{-6.110112\du}{6.294372\du}}
\pgfpathlineto{\pgfpoint{-6.104440\du}{6.291984\du}}
\pgfpathlineto{\pgfpoint{-6.098469\du}{6.289297\du}}
\pgfpathlineto{\pgfpoint{-6.092796\du}{6.287207\du}}
\pgfpathlineto{\pgfpoint{-6.088318\du}{6.284520\du}}
\pgfpathlineto{\pgfpoint{-6.083541\du}{6.281833\du}}
\pgfpathlineto{\pgfpoint{-6.079361\du}{6.279146\du}}
\pgfpathlineto{\pgfpoint{-6.076077\du}{6.276459\du}}
\pgfpathlineto{\pgfpoint{-6.072793\du}{6.273772\du}}
\pgfpathlineto{\pgfpoint{-6.070703\du}{6.270488\du}}
\pgfpathlineto{\pgfpoint{-6.068613\du}{6.267801\du}}
\pgfpathlineto{\pgfpoint{-6.066524\du}{6.264517\du}}
\pgfpathlineto{\pgfpoint{-6.065329\du}{6.261531\du}}
\pgfpathlineto{\pgfpoint{-6.064434\du}{6.258844\du}}
\pgfpathlineto{\pgfpoint{-6.064434\du}{6.255560\du}}
\pgfusepath{stroke}
\pgfsetbuttcap
\pgfsetmiterjoin
\pgfsetdash{}{0pt}
\definecolor{dialinecolor}{rgb}{0.000000, 0.000000, 0.000000}
\pgfsetfillcolor{dialinecolor}
\pgfpathmoveto{\pgfpoint{-6.242968\du}{6.243917\du}}
\pgfpathlineto{\pgfpoint{-6.219084\du}{6.255560\du}}
\pgfpathlineto{\pgfpoint{-6.147730\du}{6.231975\du}}
\pgfpathlineto{\pgfpoint{-6.111903\du}{6.243917\du}}
\pgfpathlineto{\pgfpoint{-6.135788\du}{6.220033\du}}
\pgfpathlineto{\pgfpoint{-6.219084\du}{6.220033\du}}
\pgfpathlineto{\pgfpoint{-6.183257\du}{6.220033\du}}
\pgfpathlineto{\pgfpoint{-6.242968\du}{6.243917\du}}
\pgfusepath{fill}
\pgfsetbuttcap
\pgfsetmiterjoin
\pgfsetdash{}{0pt}
\definecolor{dialinecolor}{rgb}{0.000000, 0.000000, 0.000000}
\pgfsetfillcolor{dialinecolor}
\pgfpathmoveto{\pgfpoint{-6.255208\du}{6.279444\du}}
\pgfpathlineto{\pgfpoint{-6.290139\du}{6.267801\du}}
\pgfpathlineto{\pgfpoint{-6.350148\du}{6.291685\du}}
\pgfpathlineto{\pgfpoint{-6.385377\du}{6.279444\du}}
\pgfpathlineto{\pgfpoint{-6.374032\du}{6.314972\du}}
\pgfpathlineto{\pgfpoint{-6.290139\du}{6.314972\du}}
\pgfpathlineto{\pgfpoint{-6.326264\du}{6.303627\du}}
\pgfpathlineto{\pgfpoint{-6.255208\du}{6.279444\du}}
\pgfusepath{fill}
\pgfsetbuttcap
\pgfsetmiterjoin
\pgfsetdash{}{0pt}
\definecolor{dialinecolor}{rgb}{0.000000, 0.000000, 0.000000}
\pgfsetfillcolor{dialinecolor}
\pgfpathmoveto{\pgfpoint{-6.374032\du}{6.220033\du}}
\pgfpathlineto{\pgfpoint{-6.350148\du}{6.207493\du}}
\pgfpathlineto{\pgfpoint{-6.278495\du}{6.243917\du}}
\pgfpathlineto{\pgfpoint{-6.255208\du}{6.231975\du}}
\pgfpathlineto{\pgfpoint{-6.266852\du}{6.255560\du}}
\pgfpathlineto{\pgfpoint{-6.350148\du}{6.255560\du}}
\pgfpathlineto{\pgfpoint{-6.314322\du}{6.243917\du}}
\pgfpathlineto{\pgfpoint{-6.374032\du}{6.220033\du}}
\pgfusepath{fill}
\pgfsetbuttcap
\pgfsetmiterjoin
\pgfsetdash{}{0pt}
\definecolor{dialinecolor}{rgb}{0.000000, 0.000000, 0.000000}
\pgfsetfillcolor{dialinecolor}
\pgfpathmoveto{\pgfpoint{-6.123547\du}{6.303627\du}}
\pgfpathlineto{\pgfpoint{-6.147730\du}{6.314972\du}}
\pgfpathlineto{\pgfpoint{-6.219084\du}{6.279444\du}}
\pgfpathlineto{\pgfpoint{-6.255208\du}{6.291685\du}}
\pgfpathlineto{\pgfpoint{-6.230727\du}{6.267801\du}}
\pgfpathlineto{\pgfpoint{-6.147730\du}{6.267801\du}}
\pgfpathlineto{\pgfpoint{-6.183257\du}{6.279444\du}}
\pgfpathlineto{\pgfpoint{-6.123547\du}{6.303627\du}}
\pgfusepath{fill}
\pgfsetbuttcap
\pgfsetmiterjoin
\pgfsetdash{}{0pt}
\definecolor{dialinecolor}{rgb}{1.000000, 1.000000, 1.000000}
\pgfsetfillcolor{dialinecolor}
\pgfpathmoveto{\pgfpoint{-6.242968\du}{6.255560\du}}
\pgfpathlineto{\pgfpoint{-6.219084\du}{6.255560\du}}
\pgfpathlineto{\pgfpoint{-6.147730\du}{6.231975\du}}
\pgfpathlineto{\pgfpoint{-6.111903\du}{6.243917\du}}
\pgfpathlineto{\pgfpoint{-6.135788\du}{6.220033\du}}
\pgfpathlineto{\pgfpoint{-6.206843\du}{6.220033\du}}
\pgfpathlineto{\pgfpoint{-6.183257\du}{6.220033\du}}
\pgfpathlineto{\pgfpoint{-6.242968\du}{6.255560\du}}
\pgfusepath{fill}
\pgfsetbuttcap
\pgfsetmiterjoin
\pgfsetdash{}{0pt}
\definecolor{dialinecolor}{rgb}{1.000000, 1.000000, 1.000000}
\pgfsetfillcolor{dialinecolor}
\pgfpathmoveto{\pgfpoint{-6.255208\du}{6.279444\du}}
\pgfpathlineto{\pgfpoint{-6.278495\du}{6.267801\du}}
\pgfpathlineto{\pgfpoint{-6.350148\du}{6.291685\du}}
\pgfpathlineto{\pgfpoint{-6.385377\du}{6.291685\du}}
\pgfpathlineto{\pgfpoint{-6.361791\du}{6.314972\du}}
\pgfpathlineto{\pgfpoint{-6.278495\du}{6.314972\du}}
\pgfpathlineto{\pgfpoint{-6.314322\du}{6.303627\du}}
\pgfpathlineto{\pgfpoint{-6.255208\du}{6.279444\du}}
\pgfusepath{fill}
\pgfsetbuttcap
\pgfsetmiterjoin
\pgfsetdash{}{0pt}
\definecolor{dialinecolor}{rgb}{1.000000, 1.000000, 1.000000}
\pgfsetfillcolor{dialinecolor}
\pgfpathmoveto{\pgfpoint{-6.374032\du}{6.220033\du}}
\pgfpathlineto{\pgfpoint{-6.350148\du}{6.220033\du}}
\pgfpathlineto{\pgfpoint{-6.278495\du}{6.243917\du}}
\pgfpathlineto{\pgfpoint{-6.242968\du}{6.231975\du}}
\pgfpathlineto{\pgfpoint{-6.266852\du}{6.255560\du}}
\pgfpathlineto{\pgfpoint{-6.350148\du}{6.255560\du}}
\pgfpathlineto{\pgfpoint{-6.314322\du}{6.255560\du}}
\pgfpathlineto{\pgfpoint{-6.374032\du}{6.220033\du}}
\pgfusepath{fill}
\pgfsetbuttcap
\pgfsetmiterjoin
\pgfsetdash{}{0pt}
\definecolor{dialinecolor}{rgb}{1.000000, 1.000000, 1.000000}
\pgfsetfillcolor{dialinecolor}
\pgfpathmoveto{\pgfpoint{-6.123547\du}{6.303627\du}}
\pgfpathlineto{\pgfpoint{-6.147730\du}{6.314972\du}}
\pgfpathlineto{\pgfpoint{-6.219084\du}{6.291685\du}}
\pgfpathlineto{\pgfpoint{-6.242968\du}{6.291685\du}}
\pgfpathlineto{\pgfpoint{-6.230727\du}{6.267801\du}}
\pgfpathlineto{\pgfpoint{-6.147730\du}{6.267801\du}}
\pgfpathlineto{\pgfpoint{-6.183257\du}{6.279444\du}}
\pgfpathlineto{\pgfpoint{-6.123547\du}{6.303627\du}}
\pgfusepath{fill}
\pgfsetbuttcap
\pgfsetmiterjoin
\pgfsetdash{}{0pt}
\definecolor{dialinecolor}{rgb}{0.666667, 0.901961, 1.000000}
\pgfsetstrokecolor{dialinecolor}
\pgfpathmoveto{\pgfpoint{-6.445087\du}{6.255560\du}}
\pgfpathlineto{\pgfpoint{-6.445087\du}{6.350500\du}}
\pgfusepath{stroke}
\pgfsetbuttcap
\pgfsetmiterjoin
\pgfsetdash{}{0pt}
\definecolor{dialinecolor}{rgb}{0.666667, 0.901961, 1.000000}
\pgfsetstrokecolor{dialinecolor}
\pgfpathmoveto{\pgfpoint{-6.064434\du}{6.255560\du}}
\pgfpathlineto{\pgfpoint{-6.064434\du}{6.350500\du}}
\pgfusepath{stroke}
\pgfsetbuttcap
\pgfsetmiterjoin
\pgfsetdash{}{0pt}
\definecolor{dialinecolor}{rgb}{0.000000, 0.470588, 0.666667}
\pgfsetfillcolor{dialinecolor}
\pgfpathmoveto{\pgfpoint{-6.409858\du}{6.761606\du}}
\pgfpathlineto{\pgfpoint{-6.410157\du}{6.758023\du}}
\pgfpathlineto{\pgfpoint{-6.411053\du}{6.754739\du}}
\pgfpathlineto{\pgfpoint{-6.412247\du}{6.751455\du}}
\pgfpathlineto{\pgfpoint{-6.414038\du}{6.748470\du}}
\pgfpathlineto{\pgfpoint{-6.416128\du}{6.745186\du}}
\pgfpathlineto{\pgfpoint{-6.419412\du}{6.741901\du}}
\pgfpathlineto{\pgfpoint{-6.422099\du}{6.738916\du}}
\pgfpathlineto{\pgfpoint{-6.425980\du}{6.735632\du}}
\pgfpathlineto{\pgfpoint{-6.430160\du}{6.732945\du}}
\pgfpathlineto{\pgfpoint{-6.434638\du}{6.729661\du}}
\pgfpathlineto{\pgfpoint{-6.440012\du}{6.726974\du}}
\pgfpathlineto{\pgfpoint{-6.445087\du}{6.724287\du}}
\pgfpathlineto{\pgfpoint{-6.451059\du}{6.721600\du}}
\pgfpathlineto{\pgfpoint{-6.457030\du}{6.718913\du}}
\pgfpathlineto{\pgfpoint{-6.464195\du}{6.716525\du}}
\pgfpathlineto{\pgfpoint{-6.471062\du}{6.714136\du}}
\pgfpathlineto{\pgfpoint{-6.478525\du}{6.712345\du}}
\pgfpathlineto{\pgfpoint{-6.485989\du}{6.709956\du}}
\pgfpathlineto{\pgfpoint{-6.494050\du}{6.708165\du}}
\pgfpathlineto{\pgfpoint{-6.502409\du}{6.705777\du}}
\pgfpathlineto{\pgfpoint{-6.510769\du}{6.704284\du}}
\pgfpathlineto{\pgfpoint{-6.519725\du}{6.703090\du}}
\pgfpathlineto{\pgfpoint{-6.528682\du}{6.701597\du}}
\pgfpathlineto{\pgfpoint{-6.537937\du}{6.700403\du}}
\pgfpathlineto{\pgfpoint{-6.547789\du}{6.699209\du}}
\pgfpathlineto{\pgfpoint{-6.557044\du}{6.698313\du}}
\pgfpathlineto{\pgfpoint{-6.566897\du}{6.697716\du}}
\pgfpathlineto{\pgfpoint{-6.576450\du}{6.697119\du}}
\pgfpathlineto{\pgfpoint{-6.586303\du}{6.696522\du}}
\pgfpathlineto{\pgfpoint{-6.596752\du}{6.696223\du}}
\pgfpathlineto{\pgfpoint{-6.606007\du}{6.696223\du}}
\pgfpathlineto{\pgfpoint{-6.606007\du}{6.696223\du}}
\pgfpathlineto{\pgfpoint{-6.615859\du}{6.696223\du}}
\pgfpathlineto{\pgfpoint{-6.626010\du}{6.696522\du}}
\pgfpathlineto{\pgfpoint{-6.635862\du}{6.697119\du}}
\pgfpathlineto{\pgfpoint{-6.645714\du}{6.697716\du}}
\pgfpathlineto{\pgfpoint{-6.655567\du}{6.698313\du}}
\pgfpathlineto{\pgfpoint{-6.664822\du}{6.699209\du}}
\pgfpathlineto{\pgfpoint{-6.674674\du}{6.700403\du}}
\pgfpathlineto{\pgfpoint{-6.683929\du}{6.701597\du}}
\pgfpathlineto{\pgfpoint{-6.692886\du}{6.703090\du}}
\pgfpathlineto{\pgfpoint{-6.701544\du}{6.704284\du}}
\pgfpathlineto{\pgfpoint{-6.709903\du}{6.705777\du}}
\pgfpathlineto{\pgfpoint{-6.718263\du}{6.708165\du}}
\pgfpathlineto{\pgfpoint{-6.726323\du}{6.709956\du}}
\pgfpathlineto{\pgfpoint{-6.733787\du}{6.712345\du}}
\pgfpathlineto{\pgfpoint{-6.741251\du}{6.714136\du}}
\pgfpathlineto{\pgfpoint{-6.748416\du}{6.716525\du}}
\pgfpathlineto{\pgfpoint{-6.755283\du}{6.718913\du}}
\pgfpathlineto{\pgfpoint{-6.761254\du}{6.721600\du}}
\pgfpathlineto{\pgfpoint{-6.767524\du}{6.724287\du}}
\pgfpathlineto{\pgfpoint{-6.772599\du}{6.726974\du}}
\pgfpathlineto{\pgfpoint{-6.777674\du}{6.729661\du}}
\pgfpathlineto{\pgfpoint{-6.782451\du}{6.732945\du}}
\pgfpathlineto{\pgfpoint{-6.786631\du}{6.735632\du}}
\pgfpathlineto{\pgfpoint{-6.790214\du}{6.738916\du}}
\pgfpathlineto{\pgfpoint{-6.793199\du}{6.741901\du}}
\pgfpathlineto{\pgfpoint{-6.795886\du}{6.745186\du}}
\pgfpathlineto{\pgfpoint{-6.798275\du}{6.748470\du}}
\pgfpathlineto{\pgfpoint{-6.800066\du}{6.751455\du}}
\pgfpathlineto{\pgfpoint{-6.801559\du}{6.754739\du}}
\pgfpathlineto{\pgfpoint{-6.802454\du}{6.758023\du}}
\pgfpathlineto{\pgfpoint{-6.802454\du}{6.761606\du}}
\pgfpathlineto{\pgfpoint{-6.802454\du}{6.761606\du}}
\pgfpathlineto{\pgfpoint{-6.802454\du}{6.764890\du}}
\pgfpathlineto{\pgfpoint{-6.801559\du}{6.768174\du}}
\pgfpathlineto{\pgfpoint{-6.800066\du}{6.771458\du}}
\pgfpathlineto{\pgfpoint{-6.798275\du}{6.774742\du}}
\pgfpathlineto{\pgfpoint{-6.795886\du}{6.777728\du}}
\pgfpathlineto{\pgfpoint{-6.793199\du}{6.781012\du}}
\pgfpathlineto{\pgfpoint{-6.790214\du}{6.783997\du}}
\pgfpathlineto{\pgfpoint{-6.786631\du}{6.787281\du}}
\pgfpathlineto{\pgfpoint{-6.782451\du}{6.790267\du}}
\pgfpathlineto{\pgfpoint{-6.777674\du}{6.793252\du}}
\pgfpathlineto{\pgfpoint{-6.772599\du}{6.795939\du}}
\pgfpathlineto{\pgfpoint{-6.767524\du}{6.798626\du}}
\pgfpathlineto{\pgfpoint{-6.761254\du}{6.801313\du}}
\pgfpathlineto{\pgfpoint{-6.755283\du}{6.803702\du}}
\pgfpathlineto{\pgfpoint{-6.748416\du}{6.806389\du}}
\pgfpathlineto{\pgfpoint{-6.741251\du}{6.808777\du}}
\pgfpathlineto{\pgfpoint{-6.733787\du}{6.810867\du}}
\pgfpathlineto{\pgfpoint{-6.726323\du}{6.813255\du}}
\pgfpathlineto{\pgfpoint{-6.718263\du}{6.814748\du}}
\pgfpathlineto{\pgfpoint{-6.709903\du}{6.816838\du}}
\pgfpathlineto{\pgfpoint{-6.701544\du}{6.818629\du}}
\pgfpathlineto{\pgfpoint{-6.692886\du}{6.819824\du}}
\pgfpathlineto{\pgfpoint{-6.683929\du}{6.821316\du}}
\pgfpathlineto{\pgfpoint{-6.674674\du}{6.822511\du}}
\pgfpathlineto{\pgfpoint{-6.664822\du}{6.823705\du}}
\pgfpathlineto{\pgfpoint{-6.655567\du}{6.824600\du}}
\pgfpathlineto{\pgfpoint{-6.645714\du}{6.825198\du}}
\pgfpathlineto{\pgfpoint{-6.635862\du}{6.826093\du}}
\pgfpathlineto{\pgfpoint{-6.626010\du}{6.826093\du}}
\pgfpathlineto{\pgfpoint{-6.615859\du}{6.826690\du}}
\pgfpathlineto{\pgfpoint{-6.606007\du}{6.826690\du}}
\pgfpathlineto{\pgfpoint{-6.606007\du}{6.826690\du}}
\pgfpathlineto{\pgfpoint{-6.596752\du}{6.826690\du}}
\pgfpathlineto{\pgfpoint{-6.586303\du}{6.826093\du}}
\pgfpathlineto{\pgfpoint{-6.576450\du}{6.826093\du}}
\pgfpathlineto{\pgfpoint{-6.566897\du}{6.825198\du}}
\pgfpathlineto{\pgfpoint{-6.557044\du}{6.824600\du}}
\pgfpathlineto{\pgfpoint{-6.547789\du}{6.823705\du}}
\pgfpathlineto{\pgfpoint{-6.537937\du}{6.822511\du}}
\pgfpathlineto{\pgfpoint{-6.528682\du}{6.821316\du}}
\pgfpathlineto{\pgfpoint{-6.519725\du}{6.819824\du}}
\pgfpathlineto{\pgfpoint{-6.510769\du}{6.818629\du}}
\pgfpathlineto{\pgfpoint{-6.502409\du}{6.816838\du}}
\pgfpathlineto{\pgfpoint{-6.494050\du}{6.814748\du}}
\pgfpathlineto{\pgfpoint{-6.485989\du}{6.813255\du}}
\pgfpathlineto{\pgfpoint{-6.478525\du}{6.810867\du}}
\pgfpathlineto{\pgfpoint{-6.471062\du}{6.808777\du}}
\pgfpathlineto{\pgfpoint{-6.464195\du}{6.806389\du}}
\pgfpathlineto{\pgfpoint{-6.457030\du}{6.803702\du}}
\pgfpathlineto{\pgfpoint{-6.451059\du}{6.801313\du}}
\pgfpathlineto{\pgfpoint{-6.445087\du}{6.798626\du}}
\pgfpathlineto{\pgfpoint{-6.440012\du}{6.795939\du}}
\pgfpathlineto{\pgfpoint{-6.434638\du}{6.793252\du}}
\pgfpathlineto{\pgfpoint{-6.430160\du}{6.790267\du}}
\pgfpathlineto{\pgfpoint{-6.425980\du}{6.787281\du}}
\pgfpathlineto{\pgfpoint{-6.422099\du}{6.783997\du}}
\pgfpathlineto{\pgfpoint{-6.419412\du}{6.781012\du}}
\pgfpathlineto{\pgfpoint{-6.416128\du}{6.777728\du}}
\pgfpathlineto{\pgfpoint{-6.414038\du}{6.774742\du}}
\pgfpathlineto{\pgfpoint{-6.412247\du}{6.771458\du}}
\pgfpathlineto{\pgfpoint{-6.411053\du}{6.768174\du}}
\pgfpathlineto{\pgfpoint{-6.410157\du}{6.764890\du}}
\pgfpathlineto{\pgfpoint{-6.409858\du}{6.761606\du}}
\pgfusepath{fill}
\pgfsetbuttcap
\pgfsetmiterjoin
\pgfsetdash{}{0pt}
\definecolor{dialinecolor}{rgb}{0.666667, 0.901961, 1.000000}
\pgfsetstrokecolor{dialinecolor}
\pgfpathmoveto{\pgfpoint{-6.421800\du}{6.755336\du}}
\pgfpathlineto{\pgfpoint{-6.422099\du}{6.752351\du}}
\pgfpathlineto{\pgfpoint{-6.422696\du}{6.749664\du}}
\pgfpathlineto{\pgfpoint{-6.423890\du}{6.746678\du}}
\pgfpathlineto{\pgfpoint{-6.425980\du}{6.743096\du}}
\pgfpathlineto{\pgfpoint{-6.428070\du}{6.740707\du}}
\pgfpathlineto{\pgfpoint{-6.430160\du}{6.737423\du}}
\pgfpathlineto{\pgfpoint{-6.433444\du}{6.734736\du}}
\pgfpathlineto{\pgfpoint{-6.436728\du}{6.732049\du}}
\pgfpathlineto{\pgfpoint{-6.440908\du}{6.729064\du}}
\pgfpathlineto{\pgfpoint{-6.445685\du}{6.726675\du}}
\pgfpathlineto{\pgfpoint{-6.450760\du}{6.723988\du}}
\pgfpathlineto{\pgfpoint{-6.455835\du}{6.721600\du}}
\pgfpathlineto{\pgfpoint{-6.461508\du}{6.719212\du}}
\pgfpathlineto{\pgfpoint{-6.467479\du}{6.716823\du}}
\pgfpathlineto{\pgfpoint{-6.474047\du}{6.714136\du}}
\pgfpathlineto{\pgfpoint{-6.480914\du}{6.712643\du}}
\pgfpathlineto{\pgfpoint{-6.487482\du}{6.710554\du}}
\pgfpathlineto{\pgfpoint{-6.495543\du}{6.708762\du}}
\pgfpathlineto{\pgfpoint{-6.503007\du}{6.707269\du}}
\pgfpathlineto{\pgfpoint{-6.510769\du}{6.705180\du}}
\pgfpathlineto{\pgfpoint{-6.519725\du}{6.703687\du}}
\pgfpathlineto{\pgfpoint{-6.528085\du}{6.702194\du}}
\pgfpathlineto{\pgfpoint{-6.536444\du}{6.701000\du}}
\pgfpathlineto{\pgfpoint{-6.545700\du}{6.699806\du}}
\pgfpathlineto{\pgfpoint{-6.554358\du}{6.698910\du}}
\pgfpathlineto{\pgfpoint{-6.563911\du}{6.698014\du}}
\pgfpathlineto{\pgfpoint{-6.573465\du}{6.697716\du}}
\pgfpathlineto{\pgfpoint{-6.582720\du}{6.697119\du}}
\pgfpathlineto{\pgfpoint{-6.592572\du}{6.696522\du}}
\pgfpathlineto{\pgfpoint{-6.602424\du}{6.696223\du}}
\pgfpathlineto{\pgfpoint{-6.611978\du}{6.696223\du}}
\pgfpathlineto{\pgfpoint{-6.611978\du}{6.696223\du}}
\pgfpathlineto{\pgfpoint{-6.621233\du}{6.696223\du}}
\pgfpathlineto{\pgfpoint{-6.631085\du}{6.696522\du}}
\pgfpathlineto{\pgfpoint{-6.640639\du}{6.697119\du}}
\pgfpathlineto{\pgfpoint{-6.650193\du}{6.697716\du}}
\pgfpathlineto{\pgfpoint{-6.659746\du}{6.698014\du}}
\pgfpathlineto{\pgfpoint{-6.669001\du}{6.698910\du}}
\pgfpathlineto{\pgfpoint{-6.677958\du}{6.699806\du}}
\pgfpathlineto{\pgfpoint{-6.687213\du}{6.701000\du}}
\pgfpathlineto{\pgfpoint{-6.695573\du}{6.702194\du}}
\pgfpathlineto{\pgfpoint{-6.704529\du}{6.703687\du}}
\pgfpathlineto{\pgfpoint{-6.712889\du}{6.705180\du}}
\pgfpathlineto{\pgfpoint{-6.720651\du}{6.707269\du}}
\pgfpathlineto{\pgfpoint{-6.728712\du}{6.708762\du}}
\pgfpathlineto{\pgfpoint{-6.736176\du}{6.710554\du}}
\pgfpathlineto{\pgfpoint{-6.742744\du}{6.712643\du}}
\pgfpathlineto{\pgfpoint{-6.750208\du}{6.714136\du}}
\pgfpathlineto{\pgfpoint{-6.756179\du}{6.716823\du}}
\pgfpathlineto{\pgfpoint{-6.762150\du}{6.719212\du}}
\pgfpathlineto{\pgfpoint{-6.768419\du}{6.721600\du}}
\pgfpathlineto{\pgfpoint{-6.773495\du}{6.723988\du}}
\pgfpathlineto{\pgfpoint{-6.778570\du}{6.726675\du}}
\pgfpathlineto{\pgfpoint{-6.782750\du}{6.729064\du}}
\pgfpathlineto{\pgfpoint{-6.786930\du}{6.732049\du}}
\pgfpathlineto{\pgfpoint{-6.790811\du}{6.734736\du}}
\pgfpathlineto{\pgfpoint{-6.794095\du}{6.737423\du}}
\pgfpathlineto{\pgfpoint{-6.795886\du}{6.740707\du}}
\pgfpathlineto{\pgfpoint{-6.798275\du}{6.743096\du}}
\pgfpathlineto{\pgfpoint{-6.800066\du}{6.746678\du}}
\pgfpathlineto{\pgfpoint{-6.801559\du}{6.749664\du}}
\pgfpathlineto{\pgfpoint{-6.801857\du}{6.752351\du}}
\pgfpathlineto{\pgfpoint{-6.802454\du}{6.755336\du}}
\pgfpathlineto{\pgfpoint{-6.802454\du}{6.755336\du}}
\pgfpathlineto{\pgfpoint{-6.801857\du}{6.758919\du}}
\pgfpathlineto{\pgfpoint{-6.801559\du}{6.761606\du}}
\pgfpathlineto{\pgfpoint{-6.800066\du}{6.764591\du}}
\pgfpathlineto{\pgfpoint{-6.798275\du}{6.767577\du}}
\pgfpathlineto{\pgfpoint{-6.795886\du}{6.770264\du}}
\pgfpathlineto{\pgfpoint{-6.794095\du}{6.773548\du}}
\pgfpathlineto{\pgfpoint{-6.790811\du}{6.776235\du}}
\pgfpathlineto{\pgfpoint{-6.786930\du}{6.779220\du}}
\pgfpathlineto{\pgfpoint{-6.782750\du}{6.781907\du}}
\pgfpathlineto{\pgfpoint{-6.778570\du}{6.784594\du}}
\pgfpathlineto{\pgfpoint{-6.773495\du}{6.787281\du}}
\pgfpathlineto{\pgfpoint{-6.768419\du}{6.788774\du}}
\pgfpathlineto{\pgfpoint{-6.762150\du}{6.792058\du}}
\pgfpathlineto{\pgfpoint{-6.756179\du}{6.794148\du}}
\pgfpathlineto{\pgfpoint{-6.750208\du}{6.796238\du}}
\pgfpathlineto{\pgfpoint{-6.742744\du}{6.798626\du}}
\pgfpathlineto{\pgfpoint{-6.736176\du}{6.800716\du}}
\pgfpathlineto{\pgfpoint{-6.728712\du}{6.802508\du}}
\pgfpathlineto{\pgfpoint{-6.720651\du}{6.803702\du}}
\pgfpathlineto{\pgfpoint{-6.712889\du}{6.806090\du}}
\pgfpathlineto{\pgfpoint{-6.704529\du}{6.807583\du}}
\pgfpathlineto{\pgfpoint{-6.695573\du}{6.809076\du}}
\pgfpathlineto{\pgfpoint{-6.687213\du}{6.810270\du}}
\pgfpathlineto{\pgfpoint{-6.677958\du}{6.811166\du}}
\pgfpathlineto{\pgfpoint{-6.669001\du}{6.812061\du}}
\pgfpathlineto{\pgfpoint{-6.659746\du}{6.813255\du}}
\pgfpathlineto{\pgfpoint{-6.650193\du}{6.813554\du}}
\pgfpathlineto{\pgfpoint{-6.640639\du}{6.814151\du}}
\pgfpathlineto{\pgfpoint{-6.631085\du}{6.814450\du}}
\pgfpathlineto{\pgfpoint{-6.621233\du}{6.814748\du}}
\pgfpathlineto{\pgfpoint{-6.611978\du}{6.814748\du}}
\pgfpathlineto{\pgfpoint{-6.611978\du}{6.814748\du}}
\pgfpathlineto{\pgfpoint{-6.602424\du}{6.814748\du}}
\pgfpathlineto{\pgfpoint{-6.592572\du}{6.814450\du}}
\pgfpathlineto{\pgfpoint{-6.582720\du}{6.814151\du}}
\pgfpathlineto{\pgfpoint{-6.573465\du}{6.813554\du}}
\pgfpathlineto{\pgfpoint{-6.563911\du}{6.813255\du}}
\pgfpathlineto{\pgfpoint{-6.554358\du}{6.812061\du}}
\pgfpathlineto{\pgfpoint{-6.545700\du}{6.811166\du}}
\pgfpathlineto{\pgfpoint{-6.536444\du}{6.810270\du}}
\pgfpathlineto{\pgfpoint{-6.528085\du}{6.809076\du}}
\pgfpathlineto{\pgfpoint{-6.519725\du}{6.807583\du}}
\pgfpathlineto{\pgfpoint{-6.510769\du}{6.806090\du}}
\pgfpathlineto{\pgfpoint{-6.503007\du}{6.803702\du}}
\pgfpathlineto{\pgfpoint{-6.495543\du}{6.802508\du}}
\pgfpathlineto{\pgfpoint{-6.487482\du}{6.800716\du}}
\pgfpathlineto{\pgfpoint{-6.480914\du}{6.798626\du}}
\pgfpathlineto{\pgfpoint{-6.474047\du}{6.796238\du}}
\pgfpathlineto{\pgfpoint{-6.467479\du}{6.794148\du}}
\pgfpathlineto{\pgfpoint{-6.461508\du}{6.792058\du}}
\pgfpathlineto{\pgfpoint{-6.455835\du}{6.788774\du}}
\pgfpathlineto{\pgfpoint{-6.450760\du}{6.787281\du}}
\pgfpathlineto{\pgfpoint{-6.445685\du}{6.784594\du}}
\pgfpathlineto{\pgfpoint{-6.440908\du}{6.781907\du}}
\pgfpathlineto{\pgfpoint{-6.436728\du}{6.779220\du}}
\pgfpathlineto{\pgfpoint{-6.433444\du}{6.776235\du}}
\pgfpathlineto{\pgfpoint{-6.430160\du}{6.773548\du}}
\pgfpathlineto{\pgfpoint{-6.428070\du}{6.770264\du}}
\pgfpathlineto{\pgfpoint{-6.425980\du}{6.767577\du}}
\pgfpathlineto{\pgfpoint{-6.423890\du}{6.764591\du}}
\pgfpathlineto{\pgfpoint{-6.422696\du}{6.761606\du}}
\pgfpathlineto{\pgfpoint{-6.422099\du}{6.758919\du}}
\pgfpathlineto{\pgfpoint{-6.421800\du}{6.755336\du}}
\pgfusepath{stroke}
\pgfsetbuttcap
\pgfsetmiterjoin
\pgfsetdash{}{0pt}
\definecolor{dialinecolor}{rgb}{0.000000, 0.470588, 0.666667}
\pgfsetfillcolor{dialinecolor}
\pgfpathmoveto{\pgfpoint{-6.802454\du}{6.672339\du}}
\pgfpathlineto{\pgfpoint{-6.802454\du}{6.767577\du}}
\pgfpathlineto{\pgfpoint{-6.409560\du}{6.767577\du}}
\pgfpathlineto{\pgfpoint{-6.409560\du}{6.672339\du}}
\pgfpathlineto{\pgfpoint{-6.802454\du}{6.672339\du}}
\pgfusepath{fill}
\pgfsetbuttcap
\pgfsetmiterjoin
\pgfsetdash{}{0pt}
\definecolor{dialinecolor}{rgb}{0.000000, 0.705882, 1.000000}
\pgfsetfillcolor{dialinecolor}
\pgfpathmoveto{\pgfpoint{-6.409858\du}{6.666069\du}}
\pgfpathlineto{\pgfpoint{-6.410157\du}{6.662487\du}}
\pgfpathlineto{\pgfpoint{-6.411053\du}{6.659203\du}}
\pgfpathlineto{\pgfpoint{-6.412247\du}{6.656217\du}}
\pgfpathlineto{\pgfpoint{-6.414038\du}{6.652634\du}}
\pgfpathlineto{\pgfpoint{-6.416128\du}{6.649350\du}}
\pgfpathlineto{\pgfpoint{-6.419412\du}{6.646066\du}}
\pgfpathlineto{\pgfpoint{-6.422099\du}{6.643379\du}}
\pgfpathlineto{\pgfpoint{-6.425980\du}{6.639498\du}}
\pgfpathlineto{\pgfpoint{-6.430160\du}{6.637110\du}}
\pgfpathlineto{\pgfpoint{-6.434638\du}{6.634124\du}}
\pgfpathlineto{\pgfpoint{-6.440012\du}{6.631139\du}}
\pgfpathlineto{\pgfpoint{-6.445087\du}{6.628452\du}}
\pgfpathlineto{\pgfpoint{-6.451059\du}{6.625765\du}}
\pgfpathlineto{\pgfpoint{-6.457030\du}{6.623376\du}}
\pgfpathlineto{\pgfpoint{-6.464195\du}{6.620689\du}}
\pgfpathlineto{\pgfpoint{-6.471062\du}{6.618301\du}}
\pgfpathlineto{\pgfpoint{-6.478525\du}{6.616211\du}}
\pgfpathlineto{\pgfpoint{-6.485989\du}{6.614121\du}}
\pgfpathlineto{\pgfpoint{-6.494050\du}{6.612330\du}}
\pgfpathlineto{\pgfpoint{-6.502409\du}{6.610240\du}}
\pgfpathlineto{\pgfpoint{-6.510769\du}{6.608449\du}}
\pgfpathlineto{\pgfpoint{-6.519725\du}{6.607255\du}}
\pgfpathlineto{\pgfpoint{-6.528682\du}{6.605762\du}}
\pgfpathlineto{\pgfpoint{-6.537937\du}{6.604568\du}}
\pgfpathlineto{\pgfpoint{-6.547789\du}{6.603373\du}}
\pgfpathlineto{\pgfpoint{-6.557044\du}{6.602179\du}}
\pgfpathlineto{\pgfpoint{-6.566897\du}{6.601881\du}}
\pgfpathlineto{\pgfpoint{-6.576450\du}{6.600985\du}}
\pgfpathlineto{\pgfpoint{-6.586303\du}{6.600686\du}}
\pgfpathlineto{\pgfpoint{-6.596752\du}{6.600388\du}}
\pgfpathlineto{\pgfpoint{-6.606007\du}{6.600388\du}}
\pgfpathlineto{\pgfpoint{-6.606007\du}{6.600388\du}}
\pgfpathlineto{\pgfpoint{-6.615859\du}{6.600388\du}}
\pgfpathlineto{\pgfpoint{-6.626010\du}{6.600686\du}}
\pgfpathlineto{\pgfpoint{-6.635862\du}{6.600985\du}}
\pgfpathlineto{\pgfpoint{-6.645714\du}{6.601881\du}}
\pgfpathlineto{\pgfpoint{-6.655567\du}{6.602179\du}}
\pgfpathlineto{\pgfpoint{-6.664822\du}{6.603373\du}}
\pgfpathlineto{\pgfpoint{-6.674674\du}{6.604568\du}}
\pgfpathlineto{\pgfpoint{-6.683929\du}{6.605762\du}}
\pgfpathlineto{\pgfpoint{-6.692886\du}{6.607255\du}}
\pgfpathlineto{\pgfpoint{-6.701544\du}{6.608449\du}}
\pgfpathlineto{\pgfpoint{-6.709903\du}{6.610240\du}}
\pgfpathlineto{\pgfpoint{-6.718263\du}{6.612330\du}}
\pgfpathlineto{\pgfpoint{-6.726323\du}{6.614121\du}}
\pgfpathlineto{\pgfpoint{-6.733787\du}{6.616211\du}}
\pgfpathlineto{\pgfpoint{-6.741251\du}{6.618301\du}}
\pgfpathlineto{\pgfpoint{-6.748416\du}{6.620689\du}}
\pgfpathlineto{\pgfpoint{-6.755283\du}{6.623376\du}}
\pgfpathlineto{\pgfpoint{-6.761254\du}{6.625765\du}}
\pgfpathlineto{\pgfpoint{-6.767524\du}{6.628452\du}}
\pgfpathlineto{\pgfpoint{-6.772599\du}{6.631139\du}}
\pgfpathlineto{\pgfpoint{-6.777674\du}{6.634124\du}}
\pgfpathlineto{\pgfpoint{-6.782451\du}{6.637110\du}}
\pgfpathlineto{\pgfpoint{-6.786631\du}{6.639498\du}}
\pgfpathlineto{\pgfpoint{-6.790214\du}{6.643379\du}}
\pgfpathlineto{\pgfpoint{-6.793199\du}{6.646066\du}}
\pgfpathlineto{\pgfpoint{-6.795886\du}{6.649350\du}}
\pgfpathlineto{\pgfpoint{-6.798275\du}{6.652634\du}}
\pgfpathlineto{\pgfpoint{-6.800066\du}{6.656217\du}}
\pgfpathlineto{\pgfpoint{-6.801559\du}{6.659203\du}}
\pgfpathlineto{\pgfpoint{-6.802454\du}{6.662487\du}}
\pgfpathlineto{\pgfpoint{-6.802454\du}{6.666069\du}}
\pgfpathlineto{\pgfpoint{-6.802454\du}{6.666069\du}}
\pgfpathlineto{\pgfpoint{-6.802454\du}{6.669353\du}}
\pgfpathlineto{\pgfpoint{-6.801559\du}{6.672637\du}}
\pgfpathlineto{\pgfpoint{-6.800066\du}{6.675623\du}}
\pgfpathlineto{\pgfpoint{-6.798275\du}{6.679206\du}}
\pgfpathlineto{\pgfpoint{-6.795886\du}{6.682490\du}}
\pgfpathlineto{\pgfpoint{-6.793199\du}{6.685475\du}}
\pgfpathlineto{\pgfpoint{-6.790214\du}{6.688461\du}}
\pgfpathlineto{\pgfpoint{-6.786631\du}{6.691745\du}}
\pgfpathlineto{\pgfpoint{-6.782451\du}{6.694730\du}}
\pgfpathlineto{\pgfpoint{-6.777674\du}{6.697716\du}}
\pgfpathlineto{\pgfpoint{-6.772599\du}{6.700701\du}}
\pgfpathlineto{\pgfpoint{-6.767524\du}{6.703388\du}}
\pgfpathlineto{\pgfpoint{-6.761254\du}{6.705777\du}}
\pgfpathlineto{\pgfpoint{-6.755283\du}{6.708464\du}}
\pgfpathlineto{\pgfpoint{-6.748416\du}{6.711151\du}}
\pgfpathlineto{\pgfpoint{-6.741251\du}{6.713240\du}}
\pgfpathlineto{\pgfpoint{-6.733787\du}{6.715629\du}}
\pgfpathlineto{\pgfpoint{-6.726323\du}{6.717719\du}}
\pgfpathlineto{\pgfpoint{-6.718263\du}{6.719510\du}}
\pgfpathlineto{\pgfpoint{-6.709903\du}{6.721600\du}}
\pgfpathlineto{\pgfpoint{-6.701544\du}{6.723391\du}}
\pgfpathlineto{\pgfpoint{-6.692886\du}{6.724585\du}}
\pgfpathlineto{\pgfpoint{-6.683929\du}{6.726078\du}}
\pgfpathlineto{\pgfpoint{-6.674674\du}{6.727272\du}}
\pgfpathlineto{\pgfpoint{-6.664822\du}{6.728168\du}}
\pgfpathlineto{\pgfpoint{-6.655567\du}{6.729064\du}}
\pgfpathlineto{\pgfpoint{-6.645714\du}{6.729959\du}}
\pgfpathlineto{\pgfpoint{-6.635862\du}{6.730855\du}}
\pgfpathlineto{\pgfpoint{-6.626010\du}{6.731154\du}}
\pgfpathlineto{\pgfpoint{-6.615859\du}{6.731452\du}}
\pgfpathlineto{\pgfpoint{-6.606007\du}{6.731452\du}}
\pgfpathlineto{\pgfpoint{-6.606007\du}{6.731452\du}}
\pgfpathlineto{\pgfpoint{-6.596752\du}{6.731452\du}}
\pgfpathlineto{\pgfpoint{-6.586303\du}{6.731154\du}}
\pgfpathlineto{\pgfpoint{-6.576450\du}{6.730855\du}}
\pgfpathlineto{\pgfpoint{-6.566897\du}{6.729959\du}}
\pgfpathlineto{\pgfpoint{-6.557044\du}{6.729064\du}}
\pgfpathlineto{\pgfpoint{-6.547789\du}{6.728168\du}}
\pgfpathlineto{\pgfpoint{-6.537937\du}{6.727272\du}}
\pgfpathlineto{\pgfpoint{-6.528682\du}{6.726078\du}}
\pgfpathlineto{\pgfpoint{-6.519725\du}{6.724585\du}}
\pgfpathlineto{\pgfpoint{-6.510769\du}{6.723391\du}}
\pgfpathlineto{\pgfpoint{-6.502409\du}{6.721600\du}}
\pgfpathlineto{\pgfpoint{-6.494050\du}{6.719510\du}}
\pgfpathlineto{\pgfpoint{-6.485989\du}{6.717719\du}}
\pgfpathlineto{\pgfpoint{-6.478525\du}{6.715629\du}}
\pgfpathlineto{\pgfpoint{-6.471062\du}{6.713240\du}}
\pgfpathlineto{\pgfpoint{-6.464195\du}{6.711151\du}}
\pgfpathlineto{\pgfpoint{-6.457030\du}{6.708464\du}}
\pgfpathlineto{\pgfpoint{-6.451059\du}{6.705777\du}}
\pgfpathlineto{\pgfpoint{-6.445087\du}{6.703388\du}}
\pgfpathlineto{\pgfpoint{-6.440012\du}{6.700701\du}}
\pgfpathlineto{\pgfpoint{-6.434638\du}{6.697716\du}}
\pgfpathlineto{\pgfpoint{-6.430160\du}{6.694730\du}}
\pgfpathlineto{\pgfpoint{-6.425980\du}{6.691745\du}}
\pgfpathlineto{\pgfpoint{-6.422099\du}{6.688461\du}}
\pgfpathlineto{\pgfpoint{-6.419412\du}{6.685475\du}}
\pgfpathlineto{\pgfpoint{-6.416128\du}{6.682490\du}}
\pgfpathlineto{\pgfpoint{-6.414038\du}{6.679206\du}}
\pgfpathlineto{\pgfpoint{-6.412247\du}{6.675623\du}}
\pgfpathlineto{\pgfpoint{-6.411053\du}{6.672637\du}}
\pgfpathlineto{\pgfpoint{-6.410157\du}{6.669353\du}}
\pgfpathlineto{\pgfpoint{-6.409858\du}{6.666069\du}}
\pgfusepath{fill}
\pgfsetbuttcap
\pgfsetmiterjoin
\pgfsetdash{}{0pt}
\definecolor{dialinecolor}{rgb}{0.666667, 0.901961, 1.000000}
\pgfsetstrokecolor{dialinecolor}
\pgfpathmoveto{\pgfpoint{-6.421800\du}{6.660397\du}}
\pgfpathlineto{\pgfpoint{-6.422099\du}{6.657411\du}}
\pgfpathlineto{\pgfpoint{-6.422696\du}{6.654426\du}}
\pgfpathlineto{\pgfpoint{-6.423890\du}{6.651739\du}}
\pgfpathlineto{\pgfpoint{-6.425980\du}{6.648455\du}}
\pgfpathlineto{\pgfpoint{-6.428070\du}{6.645768\du}}
\pgfpathlineto{\pgfpoint{-6.430160\du}{6.642484\du}}
\pgfpathlineto{\pgfpoint{-6.433444\du}{6.639498\du}}
\pgfpathlineto{\pgfpoint{-6.436728\du}{6.637110\du}}
\pgfpathlineto{\pgfpoint{-6.440908\du}{6.634423\du}}
\pgfpathlineto{\pgfpoint{-6.445685\du}{6.631736\du}}
\pgfpathlineto{\pgfpoint{-6.450760\du}{6.629049\du}}
\pgfpathlineto{\pgfpoint{-6.455835\du}{6.626660\du}}
\pgfpathlineto{\pgfpoint{-6.461508\du}{6.623675\du}}
\pgfpathlineto{\pgfpoint{-6.467479\du}{6.621884\du}}
\pgfpathlineto{\pgfpoint{-6.474047\du}{6.619794\du}}
\pgfpathlineto{\pgfpoint{-6.480914\du}{6.617107\du}}
\pgfpathlineto{\pgfpoint{-6.487482\du}{6.615614\du}}
\pgfpathlineto{\pgfpoint{-6.495543\du}{6.613823\du}}
\pgfpathlineto{\pgfpoint{-6.503007\du}{6.611733\du}}
\pgfpathlineto{\pgfpoint{-6.510769\du}{6.610240\du}}
\pgfpathlineto{\pgfpoint{-6.519725\du}{6.608747\du}}
\pgfpathlineto{\pgfpoint{-6.528085\du}{6.607255\du}}
\pgfpathlineto{\pgfpoint{-6.536444\du}{6.606060\du}}
\pgfpathlineto{\pgfpoint{-6.545700\du}{6.604866\du}}
\pgfpathlineto{\pgfpoint{-6.554358\du}{6.603672\du}}
\pgfpathlineto{\pgfpoint{-6.563911\du}{6.603075\du}}
\pgfpathlineto{\pgfpoint{-6.573465\du}{6.602179\du}}
\pgfpathlineto{\pgfpoint{-6.582720\du}{6.602179\du}}
\pgfpathlineto{\pgfpoint{-6.592572\du}{6.601881\du}}
\pgfpathlineto{\pgfpoint{-6.602424\du}{6.600985\du}}
\pgfpathlineto{\pgfpoint{-6.611978\du}{6.600985\du}}
\pgfpathlineto{\pgfpoint{-6.611978\du}{6.600985\du}}
\pgfpathlineto{\pgfpoint{-6.621233\du}{6.600985\du}}
\pgfpathlineto{\pgfpoint{-6.631085\du}{6.601881\du}}
\pgfpathlineto{\pgfpoint{-6.640639\du}{6.602179\du}}
\pgfpathlineto{\pgfpoint{-6.650193\du}{6.602179\du}}
\pgfpathlineto{\pgfpoint{-6.659746\du}{6.603075\du}}
\pgfpathlineto{\pgfpoint{-6.669001\du}{6.603672\du}}
\pgfpathlineto{\pgfpoint{-6.677958\du}{6.604866\du}}
\pgfpathlineto{\pgfpoint{-6.687213\du}{6.606060\du}}
\pgfpathlineto{\pgfpoint{-6.695573\du}{6.607255\du}}
\pgfpathlineto{\pgfpoint{-6.704529\du}{6.608747\du}}
\pgfpathlineto{\pgfpoint{-6.712889\du}{6.610240\du}}
\pgfpathlineto{\pgfpoint{-6.720651\du}{6.611733\du}}
\pgfpathlineto{\pgfpoint{-6.728712\du}{6.613823\du}}
\pgfpathlineto{\pgfpoint{-6.736176\du}{6.615614\du}}
\pgfpathlineto{\pgfpoint{-6.742744\du}{6.617107\du}}
\pgfpathlineto{\pgfpoint{-6.750208\du}{6.619794\du}}
\pgfpathlineto{\pgfpoint{-6.756179\du}{6.621884\du}}
\pgfpathlineto{\pgfpoint{-6.762150\du}{6.623675\du}}
\pgfpathlineto{\pgfpoint{-6.768419\du}{6.626660\du}}
\pgfpathlineto{\pgfpoint{-6.773495\du}{6.629049\du}}
\pgfpathlineto{\pgfpoint{-6.778570\du}{6.631736\du}}
\pgfpathlineto{\pgfpoint{-6.782750\du}{6.634423\du}}
\pgfpathlineto{\pgfpoint{-6.786930\du}{6.637110\du}}
\pgfpathlineto{\pgfpoint{-6.790811\du}{6.639498\du}}
\pgfpathlineto{\pgfpoint{-6.794095\du}{6.642484\du}}
\pgfpathlineto{\pgfpoint{-6.795886\du}{6.645768\du}}
\pgfpathlineto{\pgfpoint{-6.798275\du}{6.648455\du}}
\pgfpathlineto{\pgfpoint{-6.800066\du}{6.651739\du}}
\pgfpathlineto{\pgfpoint{-6.801559\du}{6.654426\du}}
\pgfpathlineto{\pgfpoint{-6.801857\du}{6.657411\du}}
\pgfpathlineto{\pgfpoint{-6.802454\du}{6.660397\du}}
\pgfpathlineto{\pgfpoint{-6.802454\du}{6.660397\du}}
\pgfpathlineto{\pgfpoint{-6.801857\du}{6.663979\du}}
\pgfpathlineto{\pgfpoint{-6.801559\du}{6.666666\du}}
\pgfpathlineto{\pgfpoint{-6.800066\du}{6.669353\du}}
\pgfpathlineto{\pgfpoint{-6.798275\du}{6.672637\du}}
\pgfpathlineto{\pgfpoint{-6.795886\du}{6.675324\du}}
\pgfpathlineto{\pgfpoint{-6.794095\du}{6.678907\du}}
\pgfpathlineto{\pgfpoint{-6.790811\du}{6.681594\du}}
\pgfpathlineto{\pgfpoint{-6.786930\du}{6.684281\du}}
\pgfpathlineto{\pgfpoint{-6.782750\du}{6.686968\du}}
\pgfpathlineto{\pgfpoint{-6.778570\du}{6.689655\du}}
\pgfpathlineto{\pgfpoint{-6.773495\du}{6.692342\du}}
\pgfpathlineto{\pgfpoint{-6.768419\du}{6.694432\du}}
\pgfpathlineto{\pgfpoint{-6.762150\du}{6.697119\du}}
\pgfpathlineto{\pgfpoint{-6.756179\du}{6.699209\du}}
\pgfpathlineto{\pgfpoint{-6.750208\du}{6.701597\du}}
\pgfpathlineto{\pgfpoint{-6.742744\du}{6.703687\du}}
\pgfpathlineto{\pgfpoint{-6.736176\du}{6.705777\du}}
\pgfpathlineto{\pgfpoint{-6.728712\du}{6.707568\du}}
\pgfpathlineto{\pgfpoint{-6.720651\du}{6.709061\du}}
\pgfpathlineto{\pgfpoint{-6.712889\du}{6.711151\du}}
\pgfpathlineto{\pgfpoint{-6.704529\du}{6.712643\du}}
\pgfpathlineto{\pgfpoint{-6.695573\du}{6.714136\du}}
\pgfpathlineto{\pgfpoint{-6.687213\du}{6.715330\du}}
\pgfpathlineto{\pgfpoint{-6.677958\du}{6.716525\du}}
\pgfpathlineto{\pgfpoint{-6.669001\du}{6.717122\du}}
\pgfpathlineto{\pgfpoint{-6.659746\du}{6.718316\du}}
\pgfpathlineto{\pgfpoint{-6.650193\du}{6.718614\du}}
\pgfpathlineto{\pgfpoint{-6.640639\du}{6.719212\du}}
\pgfpathlineto{\pgfpoint{-6.631085\du}{6.719510\du}}
\pgfpathlineto{\pgfpoint{-6.621233\du}{6.720107\du}}
\pgfpathlineto{\pgfpoint{-6.611978\du}{6.720107\du}}
\pgfpathlineto{\pgfpoint{-6.611978\du}{6.720107\du}}
\pgfpathlineto{\pgfpoint{-6.602424\du}{6.720107\du}}
\pgfpathlineto{\pgfpoint{-6.592572\du}{6.719510\du}}
\pgfpathlineto{\pgfpoint{-6.582720\du}{6.719212\du}}
\pgfpathlineto{\pgfpoint{-6.573465\du}{6.718614\du}}
\pgfpathlineto{\pgfpoint{-6.563911\du}{6.718316\du}}
\pgfpathlineto{\pgfpoint{-6.554358\du}{6.717122\du}}
\pgfpathlineto{\pgfpoint{-6.545700\du}{6.716525\du}}
\pgfpathlineto{\pgfpoint{-6.536444\du}{6.715330\du}}
\pgfpathlineto{\pgfpoint{-6.528085\du}{6.714136\du}}
\pgfpathlineto{\pgfpoint{-6.519725\du}{6.712643\du}}
\pgfpathlineto{\pgfpoint{-6.510769\du}{6.711151\du}}
\pgfpathlineto{\pgfpoint{-6.503007\du}{6.709061\du}}
\pgfpathlineto{\pgfpoint{-6.495543\du}{6.707568\du}}
\pgfpathlineto{\pgfpoint{-6.487482\du}{6.705777\du}}
\pgfpathlineto{\pgfpoint{-6.480914\du}{6.703687\du}}
\pgfpathlineto{\pgfpoint{-6.474047\du}{6.701597\du}}
\pgfpathlineto{\pgfpoint{-6.467479\du}{6.699209\du}}
\pgfpathlineto{\pgfpoint{-6.461508\du}{6.697119\du}}
\pgfpathlineto{\pgfpoint{-6.455835\du}{6.694432\du}}
\pgfpathlineto{\pgfpoint{-6.450760\du}{6.692342\du}}
\pgfpathlineto{\pgfpoint{-6.445685\du}{6.689655\du}}
\pgfpathlineto{\pgfpoint{-6.440908\du}{6.686968\du}}
\pgfpathlineto{\pgfpoint{-6.436728\du}{6.684281\du}}
\pgfpathlineto{\pgfpoint{-6.433444\du}{6.681594\du}}
\pgfpathlineto{\pgfpoint{-6.430160\du}{6.678907\du}}
\pgfpathlineto{\pgfpoint{-6.428070\du}{6.675324\du}}
\pgfpathlineto{\pgfpoint{-6.425980\du}{6.672637\du}}
\pgfpathlineto{\pgfpoint{-6.423890\du}{6.669353\du}}
\pgfpathlineto{\pgfpoint{-6.422696\du}{6.666666\du}}
\pgfpathlineto{\pgfpoint{-6.422099\du}{6.663979\du}}
\pgfpathlineto{\pgfpoint{-6.421800\du}{6.660397\du}}
\pgfusepath{stroke}
\pgfsetbuttcap
\pgfsetmiterjoin
\pgfsetdash{}{0pt}
\definecolor{dialinecolor}{rgb}{0.000000, 0.000000, 0.000000}
\pgfsetfillcolor{dialinecolor}
\pgfpathmoveto{\pgfpoint{-6.600036\du}{6.648455\du}}
\pgfpathlineto{\pgfpoint{-6.576152\du}{6.660397\du}}
\pgfpathlineto{\pgfpoint{-6.504798\du}{6.636811\du}}
\pgfpathlineto{\pgfpoint{-6.468972\du}{6.648455\du}}
\pgfpathlineto{\pgfpoint{-6.492557\du}{6.624571\du}}
\pgfpathlineto{\pgfpoint{-6.576152\du}{6.624571\du}}
\pgfpathlineto{\pgfpoint{-6.540326\du}{6.624571\du}}
\pgfpathlineto{\pgfpoint{-6.600036\du}{6.648455\du}}
\pgfusepath{fill}
\pgfsetbuttcap
\pgfsetmiterjoin
\pgfsetdash{}{0pt}
\definecolor{dialinecolor}{rgb}{0.000000, 0.000000, 0.000000}
\pgfsetfillcolor{dialinecolor}
\pgfpathmoveto{\pgfpoint{-6.611978\du}{6.684281\du}}
\pgfpathlineto{\pgfpoint{-6.647506\du}{6.672339\du}}
\pgfpathlineto{\pgfpoint{-6.707216\du}{6.696223\du}}
\pgfpathlineto{\pgfpoint{-6.742744\du}{6.684281\du}}
\pgfpathlineto{\pgfpoint{-6.730503\du}{6.720107\du}}
\pgfpathlineto{\pgfpoint{-6.647506\du}{6.720107\du}}
\pgfpathlineto{\pgfpoint{-6.683033\du}{6.708464\du}}
\pgfpathlineto{\pgfpoint{-6.611978\du}{6.684281\du}}
\pgfusepath{fill}
\pgfsetbuttcap
\pgfsetmiterjoin
\pgfsetdash{}{0pt}
\definecolor{dialinecolor}{rgb}{0.000000, 0.000000, 0.000000}
\pgfsetfillcolor{dialinecolor}
\pgfpathmoveto{\pgfpoint{-6.730503\du}{6.624571\du}}
\pgfpathlineto{\pgfpoint{-6.707216\du}{6.612927\du}}
\pgfpathlineto{\pgfpoint{-6.635862\du}{6.648455\du}}
\pgfpathlineto{\pgfpoint{-6.611978\du}{6.636811\du}}
\pgfpathlineto{\pgfpoint{-6.623323\du}{6.660397\du}}
\pgfpathlineto{\pgfpoint{-6.707216\du}{6.660397\du}}
\pgfpathlineto{\pgfpoint{-6.671390\du}{6.648455\du}}
\pgfpathlineto{\pgfpoint{-6.730503\du}{6.624571\du}}
\pgfusepath{fill}
\pgfsetbuttcap
\pgfsetmiterjoin
\pgfsetdash{}{0pt}
\definecolor{dialinecolor}{rgb}{0.000000, 0.000000, 0.000000}
\pgfsetfillcolor{dialinecolor}
\pgfpathmoveto{\pgfpoint{-6.480914\du}{6.708464\du}}
\pgfpathlineto{\pgfpoint{-6.504798\du}{6.720107\du}}
\pgfpathlineto{\pgfpoint{-6.576152\du}{6.684281\du}}
\pgfpathlineto{\pgfpoint{-6.611978\du}{6.696223\du}}
\pgfpathlineto{\pgfpoint{-6.587795\du}{6.672339\du}}
\pgfpathlineto{\pgfpoint{-6.504798\du}{6.672339\du}}
\pgfpathlineto{\pgfpoint{-6.540326\du}{6.684281\du}}
\pgfpathlineto{\pgfpoint{-6.480914\du}{6.708464\du}}
\pgfusepath{fill}
\pgfsetbuttcap
\pgfsetmiterjoin
\pgfsetdash{}{0pt}
\definecolor{dialinecolor}{rgb}{1.000000, 1.000000, 1.000000}
\pgfsetfillcolor{dialinecolor}
\pgfpathmoveto{\pgfpoint{-6.600036\du}{6.660397\du}}
\pgfpathlineto{\pgfpoint{-6.576152\du}{6.660397\du}}
\pgfpathlineto{\pgfpoint{-6.504798\du}{6.636811\du}}
\pgfpathlineto{\pgfpoint{-6.468972\du}{6.648455\du}}
\pgfpathlineto{\pgfpoint{-6.492557\du}{6.624571\du}}
\pgfpathlineto{\pgfpoint{-6.563911\du}{6.624571\du}}
\pgfpathlineto{\pgfpoint{-6.540326\du}{6.624571\du}}
\pgfpathlineto{\pgfpoint{-6.600036\du}{6.660397\du}}
\pgfusepath{fill}
\pgfsetbuttcap
\pgfsetmiterjoin
\pgfsetdash{}{0pt}
\definecolor{dialinecolor}{rgb}{1.000000, 1.000000, 1.000000}
\pgfsetfillcolor{dialinecolor}
\pgfpathmoveto{\pgfpoint{-6.611978\du}{6.684281\du}}
\pgfpathlineto{\pgfpoint{-6.635862\du}{6.672339\du}}
\pgfpathlineto{\pgfpoint{-6.707216\du}{6.696223\du}}
\pgfpathlineto{\pgfpoint{-6.742744\du}{6.696223\du}}
\pgfpathlineto{\pgfpoint{-6.718860\du}{6.720107\du}}
\pgfpathlineto{\pgfpoint{-6.635862\du}{6.720107\du}}
\pgfpathlineto{\pgfpoint{-6.671390\du}{6.708464\du}}
\pgfpathlineto{\pgfpoint{-6.611978\du}{6.684281\du}}
\pgfusepath{fill}
\pgfsetbuttcap
\pgfsetmiterjoin
\pgfsetdash{}{0pt}
\definecolor{dialinecolor}{rgb}{1.000000, 1.000000, 1.000000}
\pgfsetfillcolor{dialinecolor}
\pgfpathmoveto{\pgfpoint{-6.730503\du}{6.624571\du}}
\pgfpathlineto{\pgfpoint{-6.707216\du}{6.624571\du}}
\pgfpathlineto{\pgfpoint{-6.635862\du}{6.648455\du}}
\pgfpathlineto{\pgfpoint{-6.600036\du}{6.636811\du}}
\pgfpathlineto{\pgfpoint{-6.623323\du}{6.660397\du}}
\pgfpathlineto{\pgfpoint{-6.707216\du}{6.660397\du}}
\pgfpathlineto{\pgfpoint{-6.671390\du}{6.660397\du}}
\pgfpathlineto{\pgfpoint{-6.730503\du}{6.624571\du}}
\pgfusepath{fill}
\pgfsetbuttcap
\pgfsetmiterjoin
\pgfsetdash{}{0pt}
\definecolor{dialinecolor}{rgb}{1.000000, 1.000000, 1.000000}
\pgfsetfillcolor{dialinecolor}
\pgfpathmoveto{\pgfpoint{-6.480914\du}{6.708464\du}}
\pgfpathlineto{\pgfpoint{-6.504798\du}{6.720107\du}}
\pgfpathlineto{\pgfpoint{-6.576152\du}{6.696223\du}}
\pgfpathlineto{\pgfpoint{-6.600036\du}{6.696223\du}}
\pgfpathlineto{\pgfpoint{-6.587795\du}{6.672339\du}}
\pgfpathlineto{\pgfpoint{-6.504798\du}{6.672339\du}}
\pgfpathlineto{\pgfpoint{-6.540326\du}{6.684281\du}}
\pgfpathlineto{\pgfpoint{-6.480914\du}{6.708464\du}}
\pgfusepath{fill}
\pgfsetbuttcap
\pgfsetmiterjoin
\pgfsetdash{}{0pt}
\definecolor{dialinecolor}{rgb}{0.666667, 0.901961, 1.000000}
\pgfsetstrokecolor{dialinecolor}
\pgfpathmoveto{\pgfpoint{-6.802454\du}{6.660397\du}}
\pgfpathlineto{\pgfpoint{-6.802454\du}{6.755336\du}}
\pgfusepath{stroke}
\pgfsetbuttcap
\pgfsetmiterjoin
\pgfsetdash{}{0pt}
\definecolor{dialinecolor}{rgb}{0.666667, 0.901961, 1.000000}
\pgfsetstrokecolor{dialinecolor}
\pgfpathmoveto{\pgfpoint{-6.421800\du}{6.660397\du}}
\pgfpathlineto{\pgfpoint{-6.421800\du}{6.755336\du}}
\pgfusepath{stroke}
\pgfsetlinewidth{0.100000\du}
\pgfsetdash{}{0pt}
\pgfsetdash{}{0pt}
\pgfsetbuttcap
\pgfsetmiterjoin
\pgfsetlinewidth{0.001000\du}
\pgfsetbuttcap
\pgfsetmiterjoin
\pgfsetdash{}{0pt}
\definecolor{dialinecolor}{rgb}{0.717647, 0.717647, 0.615686}
\pgfsetfillcolor{dialinecolor}
\pgfpathmoveto{\pgfpoint{-5.103800\du}{-2.919297\du}}
\pgfpathlineto{\pgfpoint{-5.103800\du}{-4.296764\du}}
\pgfpathlineto{\pgfpoint{-4.546782\du}{-4.296764\du}}
\pgfpathlineto{\pgfpoint{-4.546782\du}{-2.919297\du}}
\pgfpathlineto{\pgfpoint{-5.103800\du}{-2.919297\du}}
\pgfusepath{fill}
\pgfsetbuttcap
\pgfsetmiterjoin
\pgfsetdash{}{0pt}
\definecolor{dialinecolor}{rgb}{0.286275, 0.286275, 0.211765}
\pgfsetstrokecolor{dialinecolor}
\pgfpathmoveto{\pgfpoint{-5.103800\du}{-2.919297\du}}
\pgfpathlineto{\pgfpoint{-5.103800\du}{-4.296764\du}}
\pgfpathlineto{\pgfpoint{-4.546782\du}{-4.296764\du}}
\pgfpathlineto{\pgfpoint{-4.546782\du}{-2.919297\du}}
\pgfpathlineto{\pgfpoint{-5.103800\du}{-2.919297\du}}
\pgfusepath{stroke}
\pgfsetbuttcap
\pgfsetmiterjoin
\pgfsetdash{}{0pt}
\definecolor{dialinecolor}{rgb}{0.647059, 0.647059, 0.521569}
\pgfsetfillcolor{dialinecolor}
\pgfpathmoveto{\pgfpoint{-4.546782\du}{-2.919297\du}}
\pgfpathlineto{\pgfpoint{-4.546782\du}{-4.296764\du}}
\pgfpathlineto{\pgfpoint{-3.989320\du}{-4.296764\du}}
\pgfpathlineto{\pgfpoint{-3.989320\du}{-2.919297\du}}
\pgfpathlineto{\pgfpoint{-4.546782\du}{-2.919297\du}}
\pgfusepath{fill}
\pgfsetbuttcap
\pgfsetmiterjoin
\pgfsetdash{}{0pt}
\definecolor{dialinecolor}{rgb}{0.286275, 0.286275, 0.211765}
\pgfsetstrokecolor{dialinecolor}
\pgfpathmoveto{\pgfpoint{-4.546782\du}{-2.919297\du}}
\pgfpathlineto{\pgfpoint{-4.546782\du}{-4.296764\du}}
\pgfpathlineto{\pgfpoint{-3.989320\du}{-4.296764\du}}
\pgfpathlineto{\pgfpoint{-3.989320\du}{-2.919297\du}}
\pgfpathlineto{\pgfpoint{-4.546782\du}{-2.919297\du}}
\pgfusepath{stroke}
\pgfsetbuttcap
\pgfsetmiterjoin
\pgfsetdash{}{0pt}
\definecolor{dialinecolor}{rgb}{0.478431, 0.478431, 0.352941}
\pgfsetfillcolor{dialinecolor}
\pgfpathmoveto{\pgfpoint{-3.698393\du}{-3.211332\du}}
\pgfpathlineto{\pgfpoint{-3.698393\du}{-4.588800\du}}
\pgfpathlineto{\pgfpoint{-3.989320\du}{-4.296764\du}}
\pgfpathlineto{\pgfpoint{-3.989320\du}{-2.919297\du}}
\pgfpathlineto{\pgfpoint{-3.698393\du}{-3.211332\du}}
\pgfusepath{fill}
\pgfsetbuttcap
\pgfsetmiterjoin
\pgfsetdash{}{0pt}
\definecolor{dialinecolor}{rgb}{0.286275, 0.286275, 0.211765}
\pgfsetstrokecolor{dialinecolor}
\pgfpathmoveto{\pgfpoint{-3.698393\du}{-3.211332\du}}
\pgfpathlineto{\pgfpoint{-3.698393\du}{-4.588800\du}}
\pgfpathlineto{\pgfpoint{-3.989320\du}{-4.296764\du}}
\pgfpathlineto{\pgfpoint{-3.989320\du}{-2.919297\du}}
\pgfpathlineto{\pgfpoint{-3.698393\du}{-3.211332\du}}
\pgfusepath{stroke}
\pgfsetbuttcap
\pgfsetmiterjoin
\pgfsetdash{}{0pt}
\definecolor{dialinecolor}{rgb}{0.788235, 0.788235, 0.713726}
\pgfsetfillcolor{dialinecolor}
\pgfpathmoveto{\pgfpoint{-4.813095\du}{-4.588800\du}}
\pgfpathlineto{\pgfpoint{-5.103800\du}{-4.296764\du}}
\pgfpathlineto{\pgfpoint{-3.989320\du}{-4.296764\du}}
\pgfpathlineto{\pgfpoint{-3.698393\du}{-4.588800\du}}
\pgfpathlineto{\pgfpoint{-4.813095\du}{-4.588800\du}}
\pgfusepath{fill}
\pgfsetbuttcap
\pgfsetmiterjoin
\pgfsetdash{}{0pt}
\definecolor{dialinecolor}{rgb}{0.286275, 0.286275, 0.211765}
\pgfsetstrokecolor{dialinecolor}
\pgfpathmoveto{\pgfpoint{-4.813095\du}{-4.588800\du}}
\pgfpathlineto{\pgfpoint{-5.103800\du}{-4.296764\du}}
\pgfpathlineto{\pgfpoint{-3.989320\du}{-4.296764\du}}
\pgfpathlineto{\pgfpoint{-3.698393\du}{-4.588800\du}}
\pgfpathlineto{\pgfpoint{-4.813095\du}{-4.588800\du}}
\pgfusepath{stroke}
\pgfsetbuttcap
\pgfsetmiterjoin
\pgfsetdash{}{0pt}
\definecolor{dialinecolor}{rgb}{0.478431, 0.478431, 0.352941}
\pgfsetstrokecolor{dialinecolor}
\pgfpathmoveto{\pgfpoint{-4.288894\du}{-4.555539\du}}
\pgfpathlineto{\pgfpoint{-4.513298\du}{-4.330248\du}}
\pgfusepath{stroke}
\pgfsetbuttcap
\pgfsetmiterjoin
\pgfsetdash{}{0pt}
\definecolor{dialinecolor}{rgb}{0.717647, 0.717647, 0.615686}
\pgfsetfillcolor{dialinecolor}
\pgfpathmoveto{\pgfpoint{-4.821521\du}{-2.927945\du}}
\pgfpathlineto{\pgfpoint{-5.103800\du}{-2.927945\du}}
\pgfpathlineto{\pgfpoint{-5.103578\du}{-2.921736\du}}
\pgfpathlineto{\pgfpoint{-5.102470\du}{-2.915527\du}}
\pgfpathlineto{\pgfpoint{-5.100696\du}{-2.909097\du}}
\pgfpathlineto{\pgfpoint{-5.098256\du}{-2.902888\du}}
\pgfpathlineto{\pgfpoint{-5.094709\du}{-2.896901\du}}
\pgfpathlineto{\pgfpoint{-5.090939\du}{-2.890470\du}}
\pgfpathlineto{\pgfpoint{-5.086282\du}{-2.884483\du}}
\pgfpathlineto{\pgfpoint{-5.080960\du}{-2.878940\du}}
\pgfpathlineto{\pgfpoint{-5.075195\du}{-2.873174\du}}
\pgfpathlineto{\pgfpoint{-5.068543\du}{-2.867409\du}}
\pgfpathlineto{\pgfpoint{-5.061225\du}{-2.861865\du}}
\pgfpathlineto{\pgfpoint{-5.053243\du}{-2.856765\du}}
\pgfpathlineto{\pgfpoint{-5.044816\du}{-2.851665\du}}
\pgfpathlineto{\pgfpoint{-5.035725\du}{-2.846565\du}}
\pgfpathlineto{\pgfpoint{-5.026190\du}{-2.841909\du}}
\pgfpathlineto{\pgfpoint{-5.016212\du}{-2.837474\du}}
\pgfpathlineto{\pgfpoint{-5.005346\du}{-2.833039\du}}
\pgfpathlineto{\pgfpoint{-4.994259\du}{-2.829269\du}}
\pgfpathlineto{\pgfpoint{-4.982950\du}{-2.825500\du}}
\pgfpathlineto{\pgfpoint{-4.970754\du}{-2.821952\du}}
\pgfpathlineto{\pgfpoint{-4.958558\du}{-2.818847\du}}
\pgfpathlineto{\pgfpoint{-4.945919\du}{-2.815743\du}}
\pgfpathlineto{\pgfpoint{-4.932836\du}{-2.813304\du}}
\pgfpathlineto{\pgfpoint{-4.919532\du}{-2.811086\du}}
\pgfpathlineto{\pgfpoint{-4.906227\du}{-2.808869\du}}
\pgfpathlineto{\pgfpoint{-4.892257\du}{-2.807317\du}}
\pgfpathlineto{\pgfpoint{-4.878509\du}{-2.805543\du}}
\pgfpathlineto{\pgfpoint{-4.864318\du}{-2.804656\du}}
\pgfpathlineto{\pgfpoint{-4.849904\du}{-2.803769\du}}
\pgfpathlineto{\pgfpoint{-4.835713\du}{-2.803325\du}}
\pgfpathlineto{\pgfpoint{-4.821521\du}{-2.803103\du}}
\pgfpathlineto{\pgfpoint{-4.807330\du}{-2.803325\du}}
\pgfpathlineto{\pgfpoint{-4.793138\du}{-2.803769\du}}
\pgfpathlineto{\pgfpoint{-4.778946\du}{-2.804656\du}}
\pgfpathlineto{\pgfpoint{-4.764755\du}{-2.805543\du}}
\pgfpathlineto{\pgfpoint{-4.750785\du}{-2.807317\du}}
\pgfpathlineto{\pgfpoint{-4.737037\du}{-2.808869\du}}
\pgfpathlineto{\pgfpoint{-4.723732\du}{-2.811086\du}}
\pgfpathlineto{\pgfpoint{-4.710206\du}{-2.813304\du}}
\pgfpathlineto{\pgfpoint{-4.697123\du}{-2.815743\du}}
\pgfpathlineto{\pgfpoint{-4.684484\du}{-2.818847\du}}
\pgfpathlineto{\pgfpoint{-4.672066\du}{-2.821952\du}}
\pgfpathlineto{\pgfpoint{-4.660092\du}{-2.825500\du}}
\pgfpathlineto{\pgfpoint{-4.648562\du}{-2.829269\du}}
\pgfpathlineto{\pgfpoint{-4.637474\du}{-2.833039\du}}
\pgfpathlineto{\pgfpoint{-4.627053\du}{-2.837474\du}}
\pgfpathlineto{\pgfpoint{-4.616631\du}{-2.841909\du}}
\pgfpathlineto{\pgfpoint{-4.607539\du}{-2.846565\du}}
\pgfpathlineto{\pgfpoint{-4.598226\du}{-2.851665\du}}
\pgfpathlineto{\pgfpoint{-4.589578\du}{-2.856765\du}}
\pgfpathlineto{\pgfpoint{-4.581817\du}{-2.861865\du}}
\pgfpathlineto{\pgfpoint{-4.574499\du}{-2.867409\du}}
\pgfpathlineto{\pgfpoint{-4.567847\du}{-2.873174\du}}
\pgfpathlineto{\pgfpoint{-4.562304\du}{-2.878940\du}}
\pgfpathlineto{\pgfpoint{-4.556982\du}{-2.884483\du}}
\pgfpathlineto{\pgfpoint{-4.552103\du}{-2.890470\du}}
\pgfpathlineto{\pgfpoint{-4.548334\du}{-2.896901\du}}
\pgfpathlineto{\pgfpoint{-4.545008\du}{-2.902888\du}}
\pgfpathlineto{\pgfpoint{-4.542347\du}{-2.909097\du}}
\pgfpathlineto{\pgfpoint{-4.540794\du}{-2.915527\du}}
\pgfpathlineto{\pgfpoint{-4.539464\du}{-2.921736\du}}
\pgfpathlineto{\pgfpoint{-4.539021\du}{-2.927945\du}}
\pgfpathlineto{\pgfpoint{-4.821521\du}{-2.927945\du}}
\pgfusepath{fill}
\pgfsetbuttcap
\pgfsetmiterjoin
\pgfsetdash{}{0pt}
\definecolor{dialinecolor}{rgb}{0.286275, 0.286275, 0.211765}
\pgfsetstrokecolor{dialinecolor}
\pgfpathmoveto{\pgfpoint{-5.103800\du}{-2.927945\du}}
\pgfpathlineto{\pgfpoint{-5.103578\du}{-2.921736\du}}
\pgfpathlineto{\pgfpoint{-5.102470\du}{-2.915527\du}}
\pgfpathlineto{\pgfpoint{-5.100696\du}{-2.909097\du}}
\pgfpathlineto{\pgfpoint{-5.098256\du}{-2.902888\du}}
\pgfpathlineto{\pgfpoint{-5.094709\du}{-2.896901\du}}
\pgfpathlineto{\pgfpoint{-5.090939\du}{-2.890470\du}}
\pgfpathlineto{\pgfpoint{-5.086282\du}{-2.884483\du}}
\pgfpathlineto{\pgfpoint{-5.080960\du}{-2.878940\du}}
\pgfpathlineto{\pgfpoint{-5.075195\du}{-2.873174\du}}
\pgfpathlineto{\pgfpoint{-5.068543\du}{-2.867409\du}}
\pgfpathlineto{\pgfpoint{-5.061225\du}{-2.861865\du}}
\pgfpathlineto{\pgfpoint{-5.053243\du}{-2.856765\du}}
\pgfpathlineto{\pgfpoint{-5.044816\du}{-2.851665\du}}
\pgfpathlineto{\pgfpoint{-5.035725\du}{-2.846565\du}}
\pgfpathlineto{\pgfpoint{-5.025968\du}{-2.841909\du}}
\pgfpathlineto{\pgfpoint{-5.015768\du}{-2.837474\du}}
\pgfpathlineto{\pgfpoint{-5.005346\du}{-2.833039\du}}
\pgfpathlineto{\pgfpoint{-4.994259\du}{-2.829269\du}}
\pgfpathlineto{\pgfpoint{-4.982728\du}{-2.825500\du}}
\pgfpathlineto{\pgfpoint{-4.970754\du}{-2.821952\du}}
\pgfpathlineto{\pgfpoint{-4.958337\du}{-2.818847\du}}
\pgfpathlineto{\pgfpoint{-4.945919\du}{-2.815743\du}}
\pgfpathlineto{\pgfpoint{-4.932836\du}{-2.813304\du}}
\pgfpathlineto{\pgfpoint{-4.919532\du}{-2.811086\du}}
\pgfpathlineto{\pgfpoint{-4.906005\du}{-2.808869\du}}
\pgfpathlineto{\pgfpoint{-4.892035\du}{-2.807317\du}}
\pgfpathlineto{\pgfpoint{-4.878509\du}{-2.805543\du}}
\pgfpathlineto{\pgfpoint{-4.864318\du}{-2.804656\du}}
\pgfpathlineto{\pgfpoint{-4.849904\du}{-2.803769\du}}
\pgfpathlineto{\pgfpoint{-4.835713\du}{-2.803325\du}}
\pgfpathlineto{\pgfpoint{-4.821078\du}{-2.803103\du}}
\pgfpathlineto{\pgfpoint{-4.807108\du}{-2.803325\du}}
\pgfpathlineto{\pgfpoint{-4.792695\du}{-2.803769\du}}
\pgfpathlineto{\pgfpoint{-4.778725\du}{-2.804656\du}}
\pgfpathlineto{\pgfpoint{-4.764311\du}{-2.805543\du}}
\pgfpathlineto{\pgfpoint{-4.750785\du}{-2.807317\du}}
\pgfpathlineto{\pgfpoint{-4.737037\du}{-2.808869\du}}
\pgfpathlineto{\pgfpoint{-4.723289\du}{-2.811086\du}}
\pgfpathlineto{\pgfpoint{-4.709984\du}{-2.813304\du}}
\pgfpathlineto{\pgfpoint{-4.696902\du}{-2.815743\du}}
\pgfpathlineto{\pgfpoint{-4.684484\du}{-2.818847\du}}
\pgfpathlineto{\pgfpoint{-4.672066\du}{-2.821952\du}}
\pgfpathlineto{\pgfpoint{-4.660092\du}{-2.825500\du}}
\pgfpathlineto{\pgfpoint{-4.648562\du}{-2.829269\du}}
\pgfpathlineto{\pgfpoint{-4.637474\du}{-2.833039\du}}
\pgfpathlineto{\pgfpoint{-4.626387\du}{-2.837474\du}}
\pgfpathlineto{\pgfpoint{-4.616631\du}{-2.841909\du}}
\pgfpathlineto{\pgfpoint{-4.607096\du}{-2.846565\du}}
\pgfpathlineto{\pgfpoint{-4.598004\du}{-2.851665\du}}
\pgfpathlineto{\pgfpoint{-4.589356\du}{-2.856765\du}}
\pgfpathlineto{\pgfpoint{-4.581595\du}{-2.861865\du}}
\pgfpathlineto{\pgfpoint{-4.574278\du}{-2.867409\du}}
\pgfpathlineto{\pgfpoint{-4.567404\du}{-2.873174\du}}
\pgfpathlineto{\pgfpoint{-4.561638\du}{-2.878940\du}}
\pgfpathlineto{\pgfpoint{-4.556316\du}{-2.884483\du}}
\pgfpathlineto{\pgfpoint{-4.551882\du}{-2.890470\du}}
\pgfpathlineto{\pgfpoint{-4.548112\du}{-2.896901\du}}
\pgfpathlineto{\pgfpoint{-4.544564\du}{-2.902888\du}}
\pgfpathlineto{\pgfpoint{-4.542125\du}{-2.909097\du}}
\pgfpathlineto{\pgfpoint{-4.540573\du}{-2.915527\du}}
\pgfpathlineto{\pgfpoint{-4.539021\du}{-2.921736\du}}
\pgfpathlineto{\pgfpoint{-4.538799\du}{-2.927945\du}}
\pgfusepath{stroke}
\pgfsetbuttcap
\pgfsetmiterjoin
\pgfsetdash{}{0pt}
\definecolor{dialinecolor}{rgb}{0.576471, 0.576471, 0.423529}
\pgfsetfillcolor{dialinecolor}
\pgfpathmoveto{\pgfpoint{-4.712867\du}{-4.188110\du}}
\pgfpathlineto{\pgfpoint{-4.712867\du}{-4.121366\du}}
\pgfpathlineto{\pgfpoint{-4.712867\du}{-4.046195\du}}
\pgfpathlineto{\pgfpoint{-4.712867\du}{-3.954393\du}}
\pgfpathlineto{\pgfpoint{-4.712867\du}{-3.862592\du}}
\pgfpathlineto{\pgfpoint{-4.712867\du}{-3.654153\du}}
\pgfpathlineto{\pgfpoint{-4.712867\du}{-3.437067\du}}
\pgfpathlineto{\pgfpoint{-4.712867\du}{-3.336617\du}}
\pgfpathlineto{\pgfpoint{-4.712867\du}{-3.236611\du}}
\pgfpathlineto{\pgfpoint{-4.712867\du}{-3.153014\du}}
\pgfpathlineto{\pgfpoint{-4.712867\du}{-3.077843\du}}
\pgfpathlineto{\pgfpoint{-4.712867\du}{-3.011098\du}}
\pgfpathlineto{\pgfpoint{-4.712867\du}{-2.961206\du}}
\pgfpathlineto{\pgfpoint{-4.712867\du}{-2.935928\du}}
\pgfpathlineto{\pgfpoint{-4.712867\du}{-2.919297\du}}
\pgfpathlineto{\pgfpoint{-4.546782\du}{-2.919297\du}}
\pgfpathlineto{\pgfpoint{-4.546782\du}{-4.288560\du}}
\pgfpathlineto{\pgfpoint{-4.554986\du}{-4.288560\du}}
\pgfpathlineto{\pgfpoint{-4.546782\du}{-4.279912\du}}
\pgfpathlineto{\pgfpoint{-4.546782\du}{-4.271707\du}}
\pgfpathlineto{\pgfpoint{-4.571395\du}{-4.246651\du}}
\pgfpathlineto{\pgfpoint{-4.613526\du}{-4.221594\du}}
\pgfpathlineto{\pgfpoint{-4.662975\du}{-4.204963\du}}
\pgfpathlineto{\pgfpoint{-4.712867\du}{-4.188110\du}}
\pgfusepath{fill}
\pgfsetbuttcap
\pgfsetmiterjoin
\pgfsetdash{}{0pt}
\definecolor{dialinecolor}{rgb}{0.576471, 0.576471, 0.423529}
\pgfsetfillcolor{dialinecolor}
\pgfpathmoveto{\pgfpoint{-4.712867\du}{-2.819291\du}}
\pgfpathlineto{\pgfpoint{-4.654770\du}{-2.835921\du}}
\pgfpathlineto{\pgfpoint{-4.604878\du}{-2.860757\du}}
\pgfpathlineto{\pgfpoint{-4.580043\du}{-2.869405\du}}
\pgfpathlineto{\pgfpoint{-4.563191\du}{-2.886035\du}}
\pgfpathlineto{\pgfpoint{-4.554986\du}{-2.911092\du}}
\pgfpathlineto{\pgfpoint{-4.546782\du}{-2.927723\du}}
\pgfpathlineto{\pgfpoint{-4.563191\du}{-2.927723\du}}
\pgfpathlineto{\pgfpoint{-4.580043\du}{-2.927723\du}}
\pgfpathlineto{\pgfpoint{-4.604878\du}{-2.927723\du}}
\pgfpathlineto{\pgfpoint{-4.654770\du}{-2.927723\du}}
\pgfpathlineto{\pgfpoint{-4.712867\du}{-2.927723\du}}
\pgfpathlineto{\pgfpoint{-4.712867\du}{-2.894240\du}}
\pgfpathlineto{\pgfpoint{-4.712867\du}{-2.860757\du}}
\pgfpathlineto{\pgfpoint{-4.712867\du}{-2.844126\du}}
\pgfpathlineto{\pgfpoint{-4.712867\du}{-2.827495\du}}
\pgfpathlineto{\pgfpoint{-4.712867\du}{-2.819291\du}}
\pgfusepath{fill}
\pgfsetbuttcap
\pgfsetmiterjoin
\pgfsetdash{}{0pt}
\definecolor{dialinecolor}{rgb}{0.788235, 0.788235, 0.713726}
\pgfsetfillcolor{dialinecolor}
\pgfpathmoveto{\pgfpoint{-5.078965\du}{-4.313395\du}}
\pgfpathlineto{\pgfpoint{-5.103800\du}{-4.288560\du}}
\pgfpathlineto{\pgfpoint{-4.546782\du}{-4.288560\du}}
\pgfpathlineto{\pgfpoint{-4.546782\du}{-4.313395\du}}
\pgfpathlineto{\pgfpoint{-5.078965\du}{-4.313395\du}}
\pgfusepath{fill}
\pgfsetbuttcap
\pgfsetmiterjoin
\pgfsetdash{}{0pt}
\definecolor{dialinecolor}{rgb}{0.788235, 0.788235, 0.713726}
\pgfsetfillcolor{dialinecolor}
\pgfpathmoveto{\pgfpoint{-4.821521\du}{-4.296986\du}}
\pgfpathlineto{\pgfpoint{-5.103800\du}{-4.296986\du}}
\pgfpathlineto{\pgfpoint{-5.103578\du}{-4.290999\du}}
\pgfpathlineto{\pgfpoint{-5.102470\du}{-4.284347\du}}
\pgfpathlineto{\pgfpoint{-5.100696\du}{-4.278138\du}}
\pgfpathlineto{\pgfpoint{-5.098256\du}{-4.271707\du}}
\pgfpathlineto{\pgfpoint{-5.094709\du}{-4.265720\du}}
\pgfpathlineto{\pgfpoint{-5.090939\du}{-4.259512\du}}
\pgfpathlineto{\pgfpoint{-5.086282\du}{-4.253525\du}}
\pgfpathlineto{\pgfpoint{-5.080960\du}{-4.247759\du}}
\pgfpathlineto{\pgfpoint{-5.075195\du}{-4.242216\du}}
\pgfpathlineto{\pgfpoint{-5.068543\du}{-4.236450\du}}
\pgfpathlineto{\pgfpoint{-5.061225\du}{-4.230907\du}}
\pgfpathlineto{\pgfpoint{-5.053243\du}{-4.225807\du}}
\pgfpathlineto{\pgfpoint{-5.044816\du}{-4.220707\du}}
\pgfpathlineto{\pgfpoint{-5.035725\du}{-4.215606\du}}
\pgfpathlineto{\pgfpoint{-5.026190\du}{-4.210950\du}}
\pgfpathlineto{\pgfpoint{-5.016212\du}{-4.206515\du}}
\pgfpathlineto{\pgfpoint{-5.005346\du}{-4.202302\du}}
\pgfpathlineto{\pgfpoint{-4.994259\du}{-4.198311\du}}
\pgfpathlineto{\pgfpoint{-4.982950\du}{-4.194541\du}}
\pgfpathlineto{\pgfpoint{-4.970754\du}{-4.191215\du}}
\pgfpathlineto{\pgfpoint{-4.958558\du}{-4.187889\du}}
\pgfpathlineto{\pgfpoint{-4.945919\du}{-4.184784\du}}
\pgfpathlineto{\pgfpoint{-4.932836\du}{-4.182123\du}}
\pgfpathlineto{\pgfpoint{-4.919532\du}{-4.180128\du}}
\pgfpathlineto{\pgfpoint{-4.906227\du}{-4.177688\du}}
\pgfpathlineto{\pgfpoint{-4.892257\du}{-4.176136\du}}
\pgfpathlineto{\pgfpoint{-4.878509\du}{-4.174584\du}}
\pgfpathlineto{\pgfpoint{-4.864318\du}{-4.173697\du}}
\pgfpathlineto{\pgfpoint{-4.849904\du}{-4.172810\du}}
\pgfpathlineto{\pgfpoint{-4.835713\du}{-4.172367\du}}
\pgfpathlineto{\pgfpoint{-4.821521\du}{-4.172145\du}}
\pgfpathlineto{\pgfpoint{-4.807330\du}{-4.172367\du}}
\pgfpathlineto{\pgfpoint{-4.793138\du}{-4.172810\du}}
\pgfpathlineto{\pgfpoint{-4.778946\du}{-4.173697\du}}
\pgfpathlineto{\pgfpoint{-4.764755\du}{-4.174584\du}}
\pgfpathlineto{\pgfpoint{-4.750785\du}{-4.176136\du}}
\pgfpathlineto{\pgfpoint{-4.737037\du}{-4.177688\du}}
\pgfpathlineto{\pgfpoint{-4.723732\du}{-4.180128\du}}
\pgfpathlineto{\pgfpoint{-4.710206\du}{-4.182123\du}}
\pgfpathlineto{\pgfpoint{-4.697123\du}{-4.184784\du}}
\pgfpathlineto{\pgfpoint{-4.684484\du}{-4.187889\du}}
\pgfpathlineto{\pgfpoint{-4.672066\du}{-4.191215\du}}
\pgfpathlineto{\pgfpoint{-4.660092\du}{-4.194541\du}}
\pgfpathlineto{\pgfpoint{-4.648562\du}{-4.198311\du}}
\pgfpathlineto{\pgfpoint{-4.637474\du}{-4.202302\du}}
\pgfpathlineto{\pgfpoint{-4.627053\du}{-4.206515\du}}
\pgfpathlineto{\pgfpoint{-4.616631\du}{-4.210950\du}}
\pgfpathlineto{\pgfpoint{-4.607539\du}{-4.215606\du}}
\pgfpathlineto{\pgfpoint{-4.598226\du}{-4.220707\du}}
\pgfpathlineto{\pgfpoint{-4.589578\du}{-4.225807\du}}
\pgfpathlineto{\pgfpoint{-4.581817\du}{-4.230907\du}}
\pgfpathlineto{\pgfpoint{-4.574499\du}{-4.236450\du}}
\pgfpathlineto{\pgfpoint{-4.567847\du}{-4.242216\du}}
\pgfpathlineto{\pgfpoint{-4.562304\du}{-4.247759\du}}
\pgfpathlineto{\pgfpoint{-4.556982\du}{-4.253525\du}}
\pgfpathlineto{\pgfpoint{-4.552103\du}{-4.259512\du}}
\pgfpathlineto{\pgfpoint{-4.548334\du}{-4.265720\du}}
\pgfpathlineto{\pgfpoint{-4.545008\du}{-4.271707\du}}
\pgfpathlineto{\pgfpoint{-4.542347\du}{-4.278138\du}}
\pgfpathlineto{\pgfpoint{-4.540794\du}{-4.284347\du}}
\pgfpathlineto{\pgfpoint{-4.539464\du}{-4.290999\du}}
\pgfpathlineto{\pgfpoint{-4.539021\du}{-4.296986\du}}
\pgfpathlineto{\pgfpoint{-4.821521\du}{-4.296986\du}}
\pgfusepath{fill}
\pgfsetbuttcap
\pgfsetmiterjoin
\pgfsetdash{}{0pt}
\definecolor{dialinecolor}{rgb}{0.286275, 0.286275, 0.211765}
\pgfsetstrokecolor{dialinecolor}
\pgfpathmoveto{\pgfpoint{-5.103800\du}{-4.296986\du}}
\pgfpathlineto{\pgfpoint{-5.103578\du}{-4.290999\du}}
\pgfpathlineto{\pgfpoint{-5.102470\du}{-4.284347\du}}
\pgfpathlineto{\pgfpoint{-5.100696\du}{-4.278138\du}}
\pgfpathlineto{\pgfpoint{-5.098256\du}{-4.271707\du}}
\pgfpathlineto{\pgfpoint{-5.094709\du}{-4.265720\du}}
\pgfpathlineto{\pgfpoint{-5.090939\du}{-4.259512\du}}
\pgfpathlineto{\pgfpoint{-5.086282\du}{-4.253525\du}}
\pgfpathlineto{\pgfpoint{-5.080960\du}{-4.247759\du}}
\pgfpathlineto{\pgfpoint{-5.075195\du}{-4.242216\du}}
\pgfpathlineto{\pgfpoint{-5.068543\du}{-4.236450\du}}
\pgfpathlineto{\pgfpoint{-5.061225\du}{-4.230907\du}}
\pgfpathlineto{\pgfpoint{-5.053243\du}{-4.225807\du}}
\pgfpathlineto{\pgfpoint{-5.044816\du}{-4.220707\du}}
\pgfpathlineto{\pgfpoint{-5.035725\du}{-4.215606\du}}
\pgfpathlineto{\pgfpoint{-5.025968\du}{-4.210950\du}}
\pgfpathlineto{\pgfpoint{-5.015768\du}{-4.206515\du}}
\pgfpathlineto{\pgfpoint{-5.005346\du}{-4.202302\du}}
\pgfpathlineto{\pgfpoint{-4.994259\du}{-4.198311\du}}
\pgfpathlineto{\pgfpoint{-4.982728\du}{-4.194541\du}}
\pgfpathlineto{\pgfpoint{-4.970754\du}{-4.191215\du}}
\pgfpathlineto{\pgfpoint{-4.958337\du}{-4.187889\du}}
\pgfpathlineto{\pgfpoint{-4.945919\du}{-4.184784\du}}
\pgfpathlineto{\pgfpoint{-4.932836\du}{-4.182123\du}}
\pgfpathlineto{\pgfpoint{-4.919532\du}{-4.180128\du}}
\pgfpathlineto{\pgfpoint{-4.906005\du}{-4.177688\du}}
\pgfpathlineto{\pgfpoint{-4.892035\du}{-4.176136\du}}
\pgfpathlineto{\pgfpoint{-4.878509\du}{-4.174584\du}}
\pgfpathlineto{\pgfpoint{-4.864318\du}{-4.173697\du}}
\pgfpathlineto{\pgfpoint{-4.849904\du}{-4.172810\du}}
\pgfpathlineto{\pgfpoint{-4.835713\du}{-4.172367\du}}
\pgfpathlineto{\pgfpoint{-4.821078\du}{-4.172145\du}}
\pgfpathlineto{\pgfpoint{-4.807108\du}{-4.172367\du}}
\pgfpathlineto{\pgfpoint{-4.792695\du}{-4.172810\du}}
\pgfpathlineto{\pgfpoint{-4.778725\du}{-4.173697\du}}
\pgfpathlineto{\pgfpoint{-4.764311\du}{-4.174584\du}}
\pgfpathlineto{\pgfpoint{-4.750785\du}{-4.176136\du}}
\pgfpathlineto{\pgfpoint{-4.737037\du}{-4.177688\du}}
\pgfpathlineto{\pgfpoint{-4.723289\du}{-4.180128\du}}
\pgfpathlineto{\pgfpoint{-4.709984\du}{-4.182123\du}}
\pgfpathlineto{\pgfpoint{-4.696902\du}{-4.184784\du}}
\pgfpathlineto{\pgfpoint{-4.684484\du}{-4.187889\du}}
\pgfpathlineto{\pgfpoint{-4.672066\du}{-4.191215\du}}
\pgfpathlineto{\pgfpoint{-4.660092\du}{-4.194541\du}}
\pgfpathlineto{\pgfpoint{-4.648562\du}{-4.198311\du}}
\pgfpathlineto{\pgfpoint{-4.637474\du}{-4.202302\du}}
\pgfpathlineto{\pgfpoint{-4.626387\du}{-4.206515\du}}
\pgfpathlineto{\pgfpoint{-4.616631\du}{-4.210950\du}}
\pgfpathlineto{\pgfpoint{-4.607096\du}{-4.215606\du}}
\pgfpathlineto{\pgfpoint{-4.598004\du}{-4.220707\du}}
\pgfpathlineto{\pgfpoint{-4.589356\du}{-4.225807\du}}
\pgfpathlineto{\pgfpoint{-4.581595\du}{-4.230907\du}}
\pgfpathlineto{\pgfpoint{-4.574278\du}{-4.236450\du}}
\pgfpathlineto{\pgfpoint{-4.567404\du}{-4.242216\du}}
\pgfpathlineto{\pgfpoint{-4.561638\du}{-4.247759\du}}
\pgfpathlineto{\pgfpoint{-4.556316\du}{-4.253525\du}}
\pgfpathlineto{\pgfpoint{-4.551882\du}{-4.259512\du}}
\pgfpathlineto{\pgfpoint{-4.548112\du}{-4.265720\du}}
\pgfpathlineto{\pgfpoint{-4.544564\du}{-4.271707\du}}
\pgfpathlineto{\pgfpoint{-4.542125\du}{-4.278138\du}}
\pgfpathlineto{\pgfpoint{-4.540573\du}{-4.284347\du}}
\pgfpathlineto{\pgfpoint{-4.539021\du}{-4.290999\du}}
\pgfpathlineto{\pgfpoint{-4.538799\du}{-4.296986\du}}
\pgfusepath{stroke}
\pgfsetbuttcap
\pgfsetmiterjoin
\pgfsetdash{}{0pt}
\definecolor{dialinecolor}{rgb}{0.478431, 0.478431, 0.352941}
\pgfsetstrokecolor{dialinecolor}
\pgfpathmoveto{\pgfpoint{-5.053908\du}{-4.296764\du}}
\pgfpathlineto{\pgfpoint{-4.596674\du}{-4.296764\du}}
\pgfusepath{stroke}
\pgfsetbuttcap
\pgfsetmiterjoin
\pgfsetdash{}{0pt}
\definecolor{dialinecolor}{rgb}{0.576471, 0.576471, 0.423529}
\pgfsetfillcolor{dialinecolor}
\pgfpathmoveto{\pgfpoint{-4.538355\du}{-2.927723\du}}
\pgfpathlineto{\pgfpoint{-4.538355\du}{-4.288560\du}}
\pgfpathlineto{\pgfpoint{-4.438349\du}{-4.288560\du}}
\pgfpathlineto{\pgfpoint{-4.438349\du}{-2.927723\du}}
\pgfpathlineto{\pgfpoint{-4.538355\du}{-2.927723\du}}
\pgfusepath{fill}
\pgfsetbuttcap
\pgfsetmiterjoin
\pgfsetdash{}{0pt}
\definecolor{dialinecolor}{rgb}{0.576471, 0.576471, 0.423529}
\pgfsetfillcolor{dialinecolor}
\pgfpathmoveto{\pgfpoint{-4.455202\du}{-3.887870\du}}
\pgfpathlineto{\pgfpoint{-4.455202\du}{-4.204963\du}}
\pgfpathlineto{\pgfpoint{-4.080899\du}{-4.204963\du}}
\pgfpathlineto{\pgfpoint{-4.080899\du}{-3.887870\du}}
\pgfpathlineto{\pgfpoint{-4.455202\du}{-3.887870\du}}
\pgfusepath{fill}
\pgfsetbuttcap
\pgfsetmiterjoin
\pgfsetdash{}{0pt}
\definecolor{dialinecolor}{rgb}{0.286275, 0.286275, 0.211765}
\pgfsetstrokecolor{dialinecolor}
\pgfpathmoveto{\pgfpoint{-4.455202\du}{-3.887870\du}}
\pgfpathlineto{\pgfpoint{-4.455202\du}{-4.204963\du}}
\pgfpathlineto{\pgfpoint{-4.080899\du}{-4.204963\du}}
\pgfpathlineto{\pgfpoint{-4.080899\du}{-3.887870\du}}
\pgfpathlineto{\pgfpoint{-4.455202\du}{-3.887870\du}}
\pgfusepath{stroke}
\pgfsetbuttcap
\pgfsetmiterjoin
\pgfsetdash{}{0pt}
\definecolor{dialinecolor}{rgb}{0.478431, 0.478431, 0.352941}
\pgfsetfillcolor{dialinecolor}
\pgfpathmoveto{\pgfpoint{-4.446997\du}{-3.887870\du}}
\pgfpathlineto{\pgfpoint{-4.446997\du}{-4.196315\du}}
\pgfpathlineto{\pgfpoint{-4.438349\du}{-4.196315\du}}
\pgfpathlineto{\pgfpoint{-4.438349\du}{-3.887870\du}}
\pgfpathlineto{\pgfpoint{-4.446997\du}{-3.887870\du}}
\pgfusepath{fill}
\pgfsetbuttcap
\pgfsetmiterjoin
\pgfsetdash{}{0pt}
\definecolor{dialinecolor}{rgb}{0.831373, 0.000000, 0.039216}
\pgfsetfillcolor{dialinecolor}
\pgfpathmoveto{\pgfpoint{-5.095374\du}{-4.271707\du}}
\pgfpathlineto{\pgfpoint{-5.070539\du}{-4.246651\du}}
\pgfpathlineto{\pgfpoint{-5.070539\du}{-2.877609\du}}
\pgfpathlineto{\pgfpoint{-5.095374\du}{-2.902444\du}}
\pgfpathlineto{\pgfpoint{-5.095374\du}{-4.271707\du}}
\pgfusepath{fill}
\pgfsetbuttcap
\pgfsetmiterjoin
\pgfsetdash{}{0pt}
\definecolor{dialinecolor}{rgb}{0.286275, 0.286275, 0.211765}
\pgfsetstrokecolor{dialinecolor}
\pgfpathmoveto{\pgfpoint{-5.103800\du}{-4.296986\du}}
\pgfpathlineto{\pgfpoint{-5.103578\du}{-4.290999\du}}
\pgfpathlineto{\pgfpoint{-5.102470\du}{-4.284347\du}}
\pgfpathlineto{\pgfpoint{-5.100696\du}{-4.278138\du}}
\pgfpathlineto{\pgfpoint{-5.098256\du}{-4.271707\du}}
\pgfpathlineto{\pgfpoint{-5.094709\du}{-4.265720\du}}
\pgfpathlineto{\pgfpoint{-5.090939\du}{-4.259512\du}}
\pgfpathlineto{\pgfpoint{-5.086282\du}{-4.253525\du}}
\pgfpathlineto{\pgfpoint{-5.080960\du}{-4.247759\du}}
\pgfpathlineto{\pgfpoint{-5.075195\du}{-4.242216\du}}
\pgfpathlineto{\pgfpoint{-5.068543\du}{-4.236450\du}}
\pgfpathlineto{\pgfpoint{-5.061225\du}{-4.230907\du}}
\pgfpathlineto{\pgfpoint{-5.053243\du}{-4.225807\du}}
\pgfpathlineto{\pgfpoint{-5.044816\du}{-4.220707\du}}
\pgfpathlineto{\pgfpoint{-5.035725\du}{-4.215606\du}}
\pgfpathlineto{\pgfpoint{-5.025968\du}{-4.210950\du}}
\pgfpathlineto{\pgfpoint{-5.015768\du}{-4.206515\du}}
\pgfpathlineto{\pgfpoint{-5.005346\du}{-4.202302\du}}
\pgfpathlineto{\pgfpoint{-4.994259\du}{-4.198311\du}}
\pgfpathlineto{\pgfpoint{-4.982728\du}{-4.194541\du}}
\pgfpathlineto{\pgfpoint{-4.970754\du}{-4.191215\du}}
\pgfpathlineto{\pgfpoint{-4.958337\du}{-4.187889\du}}
\pgfpathlineto{\pgfpoint{-4.945919\du}{-4.184784\du}}
\pgfpathlineto{\pgfpoint{-4.932836\du}{-4.182123\du}}
\pgfpathlineto{\pgfpoint{-4.919532\du}{-4.180128\du}}
\pgfpathlineto{\pgfpoint{-4.906005\du}{-4.177688\du}}
\pgfpathlineto{\pgfpoint{-4.892035\du}{-4.176136\du}}
\pgfpathlineto{\pgfpoint{-4.878509\du}{-4.174584\du}}
\pgfpathlineto{\pgfpoint{-4.864318\du}{-4.173697\du}}
\pgfpathlineto{\pgfpoint{-4.849904\du}{-4.172810\du}}
\pgfpathlineto{\pgfpoint{-4.835713\du}{-4.172367\du}}
\pgfpathlineto{\pgfpoint{-4.821078\du}{-4.172145\du}}
\pgfpathlineto{\pgfpoint{-4.807108\du}{-4.172367\du}}
\pgfpathlineto{\pgfpoint{-4.792695\du}{-4.172810\du}}
\pgfpathlineto{\pgfpoint{-4.778725\du}{-4.173697\du}}
\pgfpathlineto{\pgfpoint{-4.764311\du}{-4.174584\du}}
\pgfpathlineto{\pgfpoint{-4.750785\du}{-4.176136\du}}
\pgfpathlineto{\pgfpoint{-4.737037\du}{-4.177688\du}}
\pgfpathlineto{\pgfpoint{-4.723289\du}{-4.180128\du}}
\pgfpathlineto{\pgfpoint{-4.709984\du}{-4.182123\du}}
\pgfpathlineto{\pgfpoint{-4.696902\du}{-4.184784\du}}
\pgfpathlineto{\pgfpoint{-4.684484\du}{-4.187889\du}}
\pgfpathlineto{\pgfpoint{-4.672066\du}{-4.191215\du}}
\pgfpathlineto{\pgfpoint{-4.660092\du}{-4.194541\du}}
\pgfpathlineto{\pgfpoint{-4.648562\du}{-4.198311\du}}
\pgfpathlineto{\pgfpoint{-4.637474\du}{-4.202302\du}}
\pgfpathlineto{\pgfpoint{-4.626387\du}{-4.206515\du}}
\pgfpathlineto{\pgfpoint{-4.616631\du}{-4.210950\du}}
\pgfpathlineto{\pgfpoint{-4.607096\du}{-4.215606\du}}
\pgfpathlineto{\pgfpoint{-4.598004\du}{-4.220707\du}}
\pgfpathlineto{\pgfpoint{-4.589356\du}{-4.225807\du}}
\pgfpathlineto{\pgfpoint{-4.581595\du}{-4.230907\du}}
\pgfpathlineto{\pgfpoint{-4.574278\du}{-4.236450\du}}
\pgfpathlineto{\pgfpoint{-4.567404\du}{-4.242216\du}}
\pgfpathlineto{\pgfpoint{-4.561638\du}{-4.247759\du}}
\pgfpathlineto{\pgfpoint{-4.556316\du}{-4.253525\du}}
\pgfpathlineto{\pgfpoint{-4.551882\du}{-4.259512\du}}
\pgfpathlineto{\pgfpoint{-4.548112\du}{-4.265720\du}}
\pgfpathlineto{\pgfpoint{-4.544564\du}{-4.271707\du}}
\pgfpathlineto{\pgfpoint{-4.542125\du}{-4.278138\du}}
\pgfpathlineto{\pgfpoint{-4.540573\du}{-4.284347\du}}
\pgfpathlineto{\pgfpoint{-4.539021\du}{-4.290999\du}}
\pgfpathlineto{\pgfpoint{-4.538799\du}{-4.296986\du}}
\pgfusepath{stroke}
\pgfsetbuttcap
\pgfsetmiterjoin
\pgfsetdash{}{0pt}
\definecolor{dialinecolor}{rgb}{0.286275, 0.286275, 0.211765}
\pgfsetstrokecolor{dialinecolor}
\pgfpathmoveto{\pgfpoint{-5.103800\du}{-2.924175\du}}
\pgfpathlineto{\pgfpoint{-5.103578\du}{-2.918188\du}}
\pgfpathlineto{\pgfpoint{-5.102691\du}{-2.912201\du}}
\pgfpathlineto{\pgfpoint{-5.100696\du}{-2.905992\du}}
\pgfpathlineto{\pgfpoint{-5.098256\du}{-2.900227\du}}
\pgfpathlineto{\pgfpoint{-5.095374\du}{-2.894240\du}}
\pgfpathlineto{\pgfpoint{-5.091604\du}{-2.888253\du}}
\pgfpathlineto{\pgfpoint{-5.087169\du}{-2.882709\du}}
\pgfpathlineto{\pgfpoint{-5.081626\du}{-2.876944\du}}
\pgfpathlineto{\pgfpoint{-5.076526\du}{-2.871622\du}}
\pgfpathlineto{\pgfpoint{-5.069873\du}{-2.865857\du}}
\pgfpathlineto{\pgfpoint{-5.062999\du}{-2.860757\du}}
\pgfpathlineto{\pgfpoint{-5.055017\du}{-2.855878\du}}
\pgfpathlineto{\pgfpoint{-5.046812\du}{-2.850778\du}}
\pgfpathlineto{\pgfpoint{-5.037942\du}{-2.846122\du}}
\pgfpathlineto{\pgfpoint{-5.028407\du}{-2.841687\du}}
\pgfpathlineto{\pgfpoint{-5.018651\du}{-2.837252\du}}
\pgfpathlineto{\pgfpoint{-5.008672\du}{-2.833039\du}}
\pgfpathlineto{\pgfpoint{-4.997585\du}{-2.829047\du}}
\pgfpathlineto{\pgfpoint{-4.986054\du}{-2.825500\du}}
\pgfpathlineto{\pgfpoint{-4.974524\du}{-2.821952\du}}
\pgfpathlineto{\pgfpoint{-4.962106\du}{-2.818847\du}}
\pgfpathlineto{\pgfpoint{-4.949910\du}{-2.815965\du}}
\pgfpathlineto{\pgfpoint{-4.937271\du}{-2.813525\du}}
\pgfpathlineto{\pgfpoint{-4.923966\du}{-2.811086\du}}
\pgfpathlineto{\pgfpoint{-4.910662\du}{-2.809091\du}}
\pgfpathlineto{\pgfpoint{-4.897136\du}{-2.807317\du}}
\pgfpathlineto{\pgfpoint{-4.883387\du}{-2.805543\du}}
\pgfpathlineto{\pgfpoint{-4.869639\du}{-2.804656\du}}
\pgfpathlineto{\pgfpoint{-4.855448\du}{-2.803769\du}}
\pgfpathlineto{\pgfpoint{-4.841256\du}{-2.803103\du}}
\pgfpathlineto{\pgfpoint{-4.826621\du}{-2.802882\du}}
\pgfpathlineto{\pgfpoint{-4.813095\du}{-2.802882\du}}
\pgfpathlineto{\pgfpoint{-4.798903\du}{-2.803103\du}}
\pgfpathlineto{\pgfpoint{-4.784712\du}{-2.803769\du}}
\pgfpathlineto{\pgfpoint{-4.770520\du}{-2.804656\du}}
\pgfpathlineto{\pgfpoint{-4.756550\du}{-2.805986\du}}
\pgfpathlineto{\pgfpoint{-4.743246\du}{-2.807538\du}}
\pgfpathlineto{\pgfpoint{-4.729720\du}{-2.809312\du}}
\pgfpathlineto{\pgfpoint{-4.716193\du}{-2.811530\du}}
\pgfpathlineto{\pgfpoint{-4.703110\du}{-2.813747\du}}
\pgfpathlineto{\pgfpoint{-4.690471\du}{-2.816630\du}}
\pgfpathlineto{\pgfpoint{-4.678275\du}{-2.819512\du}}
\pgfpathlineto{\pgfpoint{-4.665858\du}{-2.822617\du}}
\pgfpathlineto{\pgfpoint{-4.654549\du}{-2.826165\du}}
\pgfpathlineto{\pgfpoint{-4.643018\du}{-2.829934\du}}
\pgfpathlineto{\pgfpoint{-4.632374\du}{-2.833704\du}}
\pgfpathlineto{\pgfpoint{-4.622174\du}{-2.837917\du}}
\pgfpathlineto{\pgfpoint{-4.612417\du}{-2.842130\du}}
\pgfpathlineto{\pgfpoint{-4.603104\du}{-2.847009\du}}
\pgfpathlineto{\pgfpoint{-4.594456\du}{-2.851665\du}}
\pgfpathlineto{\pgfpoint{-4.586252\du}{-2.856765\du}}
\pgfpathlineto{\pgfpoint{-4.578491\du}{-2.861865\du}}
\pgfpathlineto{\pgfpoint{-4.571395\du}{-2.867187\du}}
\pgfpathlineto{\pgfpoint{-4.565186\du}{-2.872509\du}}
\pgfpathlineto{\pgfpoint{-4.559421\du}{-2.878274\du}}
\pgfpathlineto{\pgfpoint{-4.554764\du}{-2.883818\du}}
\pgfpathlineto{\pgfpoint{-4.550551\du}{-2.889583\du}}
\pgfpathlineto{\pgfpoint{-4.546782\du}{-2.895349\du}}
\pgfpathlineto{\pgfpoint{-4.543677\du}{-2.901336\du}}
\pgfpathlineto{\pgfpoint{-4.541460\du}{-2.907323\du}}
\pgfpathlineto{\pgfpoint{-4.539686\du}{-2.913310\du}}
\pgfpathlineto{\pgfpoint{-4.539021\du}{-2.919297\du}}
\pgfpathlineto{\pgfpoint{-4.538799\du}{-2.925506\du}}
\pgfusepath{stroke}
\pgfsetlinewidth{0.100000\du}
\pgfsetdash{}{0pt}
\pgfsetdash{}{0pt}
\pgfsetbuttcap
\pgfsetmiterjoin
\pgfsetlinewidth{0.001000\du}
\pgfsetbuttcap
\pgfsetmiterjoin
\pgfsetdash{}{0pt}
\definecolor{dialinecolor}{rgb}{0.717647, 0.717647, 0.615686}
\pgfsetfillcolor{dialinecolor}
\pgfpathmoveto{\pgfpoint{-18.206000\du}{11.584049\du}}
\pgfpathlineto{\pgfpoint{-18.206000\du}{10.212080\du}}
\pgfpathlineto{\pgfpoint{-17.651205\du}{10.212080\du}}
\pgfpathlineto{\pgfpoint{-17.651205\du}{11.584049\du}}
\pgfpathlineto{\pgfpoint{-18.206000\du}{11.584049\du}}
\pgfusepath{fill}
\pgfsetbuttcap
\pgfsetmiterjoin
\pgfsetdash{}{0pt}
\definecolor{dialinecolor}{rgb}{0.286275, 0.286275, 0.211765}
\pgfsetstrokecolor{dialinecolor}
\pgfpathmoveto{\pgfpoint{-18.206000\du}{11.584049\du}}
\pgfpathlineto{\pgfpoint{-18.206000\du}{10.212080\du}}
\pgfpathlineto{\pgfpoint{-17.651205\du}{10.212080\du}}
\pgfpathlineto{\pgfpoint{-17.651205\du}{11.584049\du}}
\pgfpathlineto{\pgfpoint{-18.206000\du}{11.584049\du}}
\pgfusepath{stroke}
\pgfsetbuttcap
\pgfsetmiterjoin
\pgfsetdash{}{0pt}
\definecolor{dialinecolor}{rgb}{0.647059, 0.647059, 0.521569}
\pgfsetfillcolor{dialinecolor}
\pgfpathmoveto{\pgfpoint{-17.651205\du}{11.584049\du}}
\pgfpathlineto{\pgfpoint{-17.651205\du}{10.212080\du}}
\pgfpathlineto{\pgfpoint{-17.095968\du}{10.212080\du}}
\pgfpathlineto{\pgfpoint{-17.095968\du}{11.584049\du}}
\pgfpathlineto{\pgfpoint{-17.651205\du}{11.584049\du}}
\pgfusepath{fill}
\pgfsetbuttcap
\pgfsetmiterjoin
\pgfsetdash{}{0pt}
\definecolor{dialinecolor}{rgb}{0.286275, 0.286275, 0.211765}
\pgfsetstrokecolor{dialinecolor}
\pgfpathmoveto{\pgfpoint{-17.651205\du}{11.584049\du}}
\pgfpathlineto{\pgfpoint{-17.651205\du}{10.212080\du}}
\pgfpathlineto{\pgfpoint{-17.095968\du}{10.212080\du}}
\pgfpathlineto{\pgfpoint{-17.095968\du}{11.584049\du}}
\pgfpathlineto{\pgfpoint{-17.651205\du}{11.584049\du}}
\pgfusepath{stroke}
\pgfsetbuttcap
\pgfsetmiterjoin
\pgfsetdash{}{0pt}
\definecolor{dialinecolor}{rgb}{0.478431, 0.478431, 0.352941}
\pgfsetfillcolor{dialinecolor}
\pgfpathmoveto{\pgfpoint{-16.806203\du}{11.293179\du}}
\pgfpathlineto{\pgfpoint{-16.806203\du}{9.921210\du}}
\pgfpathlineto{\pgfpoint{-17.095968\du}{10.212080\du}}
\pgfpathlineto{\pgfpoint{-17.095968\du}{11.584049\du}}
\pgfpathlineto{\pgfpoint{-16.806203\du}{11.293179\du}}
\pgfusepath{fill}
\pgfsetbuttcap
\pgfsetmiterjoin
\pgfsetdash{}{0pt}
\definecolor{dialinecolor}{rgb}{0.286275, 0.286275, 0.211765}
\pgfsetstrokecolor{dialinecolor}
\pgfpathmoveto{\pgfpoint{-16.806203\du}{11.293179\du}}
\pgfpathlineto{\pgfpoint{-16.806203\du}{9.921210\du}}
\pgfpathlineto{\pgfpoint{-17.095968\du}{10.212080\du}}
\pgfpathlineto{\pgfpoint{-17.095968\du}{11.584049\du}}
\pgfpathlineto{\pgfpoint{-16.806203\du}{11.293179\du}}
\pgfusepath{stroke}
\pgfsetbuttcap
\pgfsetmiterjoin
\pgfsetdash{}{0pt}
\definecolor{dialinecolor}{rgb}{0.788235, 0.788235, 0.713726}
\pgfsetfillcolor{dialinecolor}
\pgfpathmoveto{\pgfpoint{-17.916455\du}{9.921210\du}}
\pgfpathlineto{\pgfpoint{-18.206000\du}{10.212080\du}}
\pgfpathlineto{\pgfpoint{-17.095968\du}{10.212080\du}}
\pgfpathlineto{\pgfpoint{-16.806203\du}{9.921210\du}}
\pgfpathlineto{\pgfpoint{-17.916455\du}{9.921210\du}}
\pgfusepath{fill}
\pgfsetbuttcap
\pgfsetmiterjoin
\pgfsetdash{}{0pt}
\definecolor{dialinecolor}{rgb}{0.286275, 0.286275, 0.211765}
\pgfsetstrokecolor{dialinecolor}
\pgfpathmoveto{\pgfpoint{-17.916455\du}{9.921210\du}}
\pgfpathlineto{\pgfpoint{-18.206000\du}{10.212080\du}}
\pgfpathlineto{\pgfpoint{-17.095968\du}{10.212080\du}}
\pgfpathlineto{\pgfpoint{-16.806203\du}{9.921210\du}}
\pgfpathlineto{\pgfpoint{-17.916455\du}{9.921210\du}}
\pgfusepath{stroke}
\pgfsetbuttcap
\pgfsetmiterjoin
\pgfsetdash{}{0pt}
\definecolor{dialinecolor}{rgb}{0.478431, 0.478431, 0.352941}
\pgfsetstrokecolor{dialinecolor}
\pgfpathmoveto{\pgfpoint{-17.394347\du}{9.954339\du}}
\pgfpathlineto{\pgfpoint{-17.617855\du}{10.178730\du}}
\pgfusepath{stroke}
\pgfsetbuttcap
\pgfsetmiterjoin
\pgfsetdash{}{0pt}
\definecolor{dialinecolor}{rgb}{0.717647, 0.717647, 0.615686}
\pgfsetfillcolor{dialinecolor}
\pgfpathmoveto{\pgfpoint{-17.924848\du}{11.575436\du}}
\pgfpathlineto{\pgfpoint{-18.206000\du}{11.575436\du}}
\pgfpathlineto{\pgfpoint{-18.205779\du}{11.581620\du}}
\pgfpathlineto{\pgfpoint{-18.204675\du}{11.587804\du}}
\pgfpathlineto{\pgfpoint{-18.202908\du}{11.594209\du}}
\pgfpathlineto{\pgfpoint{-18.200479\du}{11.600393\du}}
\pgfpathlineto{\pgfpoint{-18.196945\du}{11.606356\du}}
\pgfpathlineto{\pgfpoint{-18.193190\du}{11.612761\du}}
\pgfpathlineto{\pgfpoint{-18.188552\du}{11.618724\du}}
\pgfpathlineto{\pgfpoint{-18.183252\du}{11.624245\du}}
\pgfpathlineto{\pgfpoint{-18.177509\du}{11.629988\du}}
\pgfpathlineto{\pgfpoint{-18.170884\du}{11.635730\du}}
\pgfpathlineto{\pgfpoint{-18.163595\du}{11.641251\du}}
\pgfpathlineto{\pgfpoint{-18.155644\du}{11.646331\du}}
\pgfpathlineto{\pgfpoint{-18.147252\du}{11.651411\du}}
\pgfpathlineto{\pgfpoint{-18.138197\du}{11.656491\du}}
\pgfpathlineto{\pgfpoint{-18.128700\du}{11.661129\du}}
\pgfpathlineto{\pgfpoint{-18.118761\du}{11.665546\du}}
\pgfpathlineto{\pgfpoint{-18.107939\du}{11.669963\du}}
\pgfpathlineto{\pgfpoint{-18.096896\du}{11.673718\du}}
\pgfpathlineto{\pgfpoint{-18.085632\du}{11.677472\du}}
\pgfpathlineto{\pgfpoint{-18.073485\du}{11.681006\du}}
\pgfpathlineto{\pgfpoint{-18.061338\du}{11.684098\du}}
\pgfpathlineto{\pgfpoint{-18.048749\du}{11.687190\du}}
\pgfpathlineto{\pgfpoint{-18.035719\du}{11.689619\du}}
\pgfpathlineto{\pgfpoint{-18.022467\du}{11.691828\du}}
\pgfpathlineto{\pgfpoint{-18.009216\du}{11.694036\du}}
\pgfpathlineto{\pgfpoint{-17.995302\du}{11.695582\du}}
\pgfpathlineto{\pgfpoint{-17.981608\du}{11.697349\du}}
\pgfpathlineto{\pgfpoint{-17.967473\du}{11.698233\du}}
\pgfpathlineto{\pgfpoint{-17.953118\du}{11.699116\du}}
\pgfpathlineto{\pgfpoint{-17.938983\du}{11.699558\du}}
\pgfpathlineto{\pgfpoint{-17.924848\du}{11.699779\du}}
\pgfpathlineto{\pgfpoint{-17.910713\du}{11.699558\du}}
\pgfpathlineto{\pgfpoint{-17.896578\du}{11.699116\du}}
\pgfpathlineto{\pgfpoint{-17.882443\du}{11.698233\du}}
\pgfpathlineto{\pgfpoint{-17.868308\du}{11.697349\du}}
\pgfpathlineto{\pgfpoint{-17.854394\du}{11.695582\du}}
\pgfpathlineto{\pgfpoint{-17.840701\du}{11.694036\du}}
\pgfpathlineto{\pgfpoint{-17.827450\du}{11.691828\du}}
\pgfpathlineto{\pgfpoint{-17.813977\du}{11.689619\du}}
\pgfpathlineto{\pgfpoint{-17.800947\du}{11.687190\du}}
\pgfpathlineto{\pgfpoint{-17.788358\du}{11.684098\du}}
\pgfpathlineto{\pgfpoint{-17.775990\du}{11.681006\du}}
\pgfpathlineto{\pgfpoint{-17.764063\du}{11.677472\du}}
\pgfpathlineto{\pgfpoint{-17.752579\du}{11.673718\du}}
\pgfpathlineto{\pgfpoint{-17.741536\du}{11.669963\du}}
\pgfpathlineto{\pgfpoint{-17.731155\du}{11.665546\du}}
\pgfpathlineto{\pgfpoint{-17.720775\du}{11.661129\du}}
\pgfpathlineto{\pgfpoint{-17.711720\du}{11.656491\du}}
\pgfpathlineto{\pgfpoint{-17.702444\du}{11.651411\du}}
\pgfpathlineto{\pgfpoint{-17.693830\du}{11.646331\du}}
\pgfpathlineto{\pgfpoint{-17.686100\du}{11.641251\du}}
\pgfpathlineto{\pgfpoint{-17.678812\du}{11.635730\du}}
\pgfpathlineto{\pgfpoint{-17.672186\du}{11.629988\du}}
\pgfpathlineto{\pgfpoint{-17.666665\du}{11.624245\du}}
\pgfpathlineto{\pgfpoint{-17.661364\du}{11.618724\du}}
\pgfpathlineto{\pgfpoint{-17.656506\du}{11.612761\du}}
\pgfpathlineto{\pgfpoint{-17.652751\du}{11.606356\du}}
\pgfpathlineto{\pgfpoint{-17.649438\du}{11.600393\du}}
\pgfpathlineto{\pgfpoint{-17.646788\du}{11.594209\du}}
\pgfpathlineto{\pgfpoint{-17.645242\du}{11.587804\du}}
\pgfpathlineto{\pgfpoint{-17.643917\du}{11.581620\du}}
\pgfpathlineto{\pgfpoint{-17.643475\du}{11.575436\du}}
\pgfpathlineto{\pgfpoint{-17.924848\du}{11.575436\du}}
\pgfusepath{fill}
\pgfsetbuttcap
\pgfsetmiterjoin
\pgfsetdash{}{0pt}
\definecolor{dialinecolor}{rgb}{0.286275, 0.286275, 0.211765}
\pgfsetstrokecolor{dialinecolor}
\pgfpathmoveto{\pgfpoint{-18.206000\du}{11.575436\du}}
\pgfpathlineto{\pgfpoint{-18.205779\du}{11.581620\du}}
\pgfpathlineto{\pgfpoint{-18.204675\du}{11.587804\du}}
\pgfpathlineto{\pgfpoint{-18.202908\du}{11.594209\du}}
\pgfpathlineto{\pgfpoint{-18.200479\du}{11.600393\du}}
\pgfpathlineto{\pgfpoint{-18.196945\du}{11.606356\du}}
\pgfpathlineto{\pgfpoint{-18.193190\du}{11.612761\du}}
\pgfpathlineto{\pgfpoint{-18.188552\du}{11.618724\du}}
\pgfpathlineto{\pgfpoint{-18.183252\du}{11.624245\du}}
\pgfpathlineto{\pgfpoint{-18.177509\du}{11.629988\du}}
\pgfpathlineto{\pgfpoint{-18.170884\du}{11.635730\du}}
\pgfpathlineto{\pgfpoint{-18.163595\du}{11.641251\du}}
\pgfpathlineto{\pgfpoint{-18.155644\du}{11.646331\du}}
\pgfpathlineto{\pgfpoint{-18.147252\du}{11.651411\du}}
\pgfpathlineto{\pgfpoint{-18.138197\du}{11.656491\du}}
\pgfpathlineto{\pgfpoint{-18.128479\du}{11.661129\du}}
\pgfpathlineto{\pgfpoint{-18.118319\du}{11.665546\du}}
\pgfpathlineto{\pgfpoint{-18.107939\du}{11.669963\du}}
\pgfpathlineto{\pgfpoint{-18.096896\du}{11.673718\du}}
\pgfpathlineto{\pgfpoint{-18.085412\du}{11.677472\du}}
\pgfpathlineto{\pgfpoint{-18.073485\du}{11.681006\du}}
\pgfpathlineto{\pgfpoint{-18.061117\du}{11.684098\du}}
\pgfpathlineto{\pgfpoint{-18.048749\du}{11.687190\du}}
\pgfpathlineto{\pgfpoint{-18.035719\du}{11.689619\du}}
\pgfpathlineto{\pgfpoint{-18.022467\du}{11.691828\du}}
\pgfpathlineto{\pgfpoint{-18.008995\du}{11.694036\du}}
\pgfpathlineto{\pgfpoint{-17.995081\du}{11.695582\du}}
\pgfpathlineto{\pgfpoint{-17.981608\du}{11.697349\du}}
\pgfpathlineto{\pgfpoint{-17.967473\du}{11.698233\du}}
\pgfpathlineto{\pgfpoint{-17.953118\du}{11.699116\du}}
\pgfpathlineto{\pgfpoint{-17.938983\du}{11.699558\du}}
\pgfpathlineto{\pgfpoint{-17.924406\du}{11.699779\du}}
\pgfpathlineto{\pgfpoint{-17.910492\du}{11.699558\du}}
\pgfpathlineto{\pgfpoint{-17.896136\du}{11.699116\du}}
\pgfpathlineto{\pgfpoint{-17.882222\du}{11.698233\du}}
\pgfpathlineto{\pgfpoint{-17.867867\du}{11.697349\du}}
\pgfpathlineto{\pgfpoint{-17.854394\du}{11.695582\du}}
\pgfpathlineto{\pgfpoint{-17.840701\du}{11.694036\du}}
\pgfpathlineto{\pgfpoint{-17.827008\du}{11.691828\du}}
\pgfpathlineto{\pgfpoint{-17.813756\du}{11.689619\du}}
\pgfpathlineto{\pgfpoint{-17.800726\du}{11.687190\du}}
\pgfpathlineto{\pgfpoint{-17.788358\du}{11.684098\du}}
\pgfpathlineto{\pgfpoint{-17.775990\du}{11.681006\du}}
\pgfpathlineto{\pgfpoint{-17.764063\du}{11.677472\du}}
\pgfpathlineto{\pgfpoint{-17.752579\du}{11.673718\du}}
\pgfpathlineto{\pgfpoint{-17.741536\du}{11.669963\du}}
\pgfpathlineto{\pgfpoint{-17.730493\du}{11.665546\du}}
\pgfpathlineto{\pgfpoint{-17.720775\du}{11.661129\du}}
\pgfpathlineto{\pgfpoint{-17.711278\du}{11.656491\du}}
\pgfpathlineto{\pgfpoint{-17.702223\du}{11.651411\du}}
\pgfpathlineto{\pgfpoint{-17.693610\du}{11.646331\du}}
\pgfpathlineto{\pgfpoint{-17.685880\du}{11.641251\du}}
\pgfpathlineto{\pgfpoint{-17.678591\du}{11.635730\du}}
\pgfpathlineto{\pgfpoint{-17.671745\du}{11.629988\du}}
\pgfpathlineto{\pgfpoint{-17.666002\du}{11.624245\du}}
\pgfpathlineto{\pgfpoint{-17.660702\du}{11.618724\du}}
\pgfpathlineto{\pgfpoint{-17.656285\du}{11.612761\du}}
\pgfpathlineto{\pgfpoint{-17.652530\du}{11.606356\du}}
\pgfpathlineto{\pgfpoint{-17.648996\du}{11.600393\du}}
\pgfpathlineto{\pgfpoint{-17.646567\du}{11.594209\du}}
\pgfpathlineto{\pgfpoint{-17.645021\du}{11.587804\du}}
\pgfpathlineto{\pgfpoint{-17.643475\du}{11.581620\du}}
\pgfpathlineto{\pgfpoint{-17.643254\du}{11.575436\du}}
\pgfusepath{stroke}
\pgfsetbuttcap
\pgfsetmiterjoin
\pgfsetdash{}{0pt}
\definecolor{dialinecolor}{rgb}{0.576471, 0.576471, 0.423529}
\pgfsetfillcolor{dialinecolor}
\pgfpathmoveto{\pgfpoint{-17.816627\du}{10.320300\du}}
\pgfpathlineto{\pgfpoint{-17.816627\du}{10.386778\du}}
\pgfpathlineto{\pgfpoint{-17.816627\du}{10.461649\du}}
\pgfpathlineto{\pgfpoint{-17.816627\du}{10.553084\du}}
\pgfpathlineto{\pgfpoint{-17.816627\du}{10.644520\du}}
\pgfpathlineto{\pgfpoint{-17.816627\du}{10.852126\du}}
\pgfpathlineto{\pgfpoint{-17.816627\du}{11.068346\du}}
\pgfpathlineto{\pgfpoint{-17.816627\du}{11.168395\du}}
\pgfpathlineto{\pgfpoint{-17.816627\du}{11.268002\du}}
\pgfpathlineto{\pgfpoint{-17.816627\du}{11.351265\du}}
\pgfpathlineto{\pgfpoint{-17.816627\du}{11.426136\du}}
\pgfpathlineto{\pgfpoint{-17.816627\du}{11.492614\du}}
\pgfpathlineto{\pgfpoint{-17.816627\du}{11.542307\du}}
\pgfpathlineto{\pgfpoint{-17.816627\du}{11.567485\du}}
\pgfpathlineto{\pgfpoint{-17.816627\du}{11.584049\du}}
\pgfpathlineto{\pgfpoint{-17.651205\du}{11.584049\du}}
\pgfpathlineto{\pgfpoint{-17.651205\du}{10.220252\du}}
\pgfpathlineto{\pgfpoint{-17.659377\du}{10.220252\du}}
\pgfpathlineto{\pgfpoint{-17.651205\du}{10.228865\du}}
\pgfpathlineto{\pgfpoint{-17.651205\du}{10.237037\du}}
\pgfpathlineto{\pgfpoint{-17.675720\du}{10.261994\du}}
\pgfpathlineto{\pgfpoint{-17.717683\du}{10.286951\du}}
\pgfpathlineto{\pgfpoint{-17.766934\du}{10.303515\du}}
\pgfpathlineto{\pgfpoint{-17.816627\du}{10.320300\du}}
\pgfusepath{fill}
\pgfsetbuttcap
\pgfsetmiterjoin
\pgfsetdash{}{0pt}
\definecolor{dialinecolor}{rgb}{0.576471, 0.576471, 0.423529}
\pgfsetfillcolor{dialinecolor}
\pgfpathmoveto{\pgfpoint{-17.816627\du}{11.683656\du}}
\pgfpathlineto{\pgfpoint{-17.758763\du}{11.667092\du}}
\pgfpathlineto{\pgfpoint{-17.709070\du}{11.642356\du}}
\pgfpathlineto{\pgfpoint{-17.684334\du}{11.633742\du}}
\pgfpathlineto{\pgfpoint{-17.667548\du}{11.617178\du}}
\pgfpathlineto{\pgfpoint{-17.659377\du}{11.592221\du}}
\pgfpathlineto{\pgfpoint{-17.651205\du}{11.575657\du}}
\pgfpathlineto{\pgfpoint{-17.667548\du}{11.575657\du}}
\pgfpathlineto{\pgfpoint{-17.684334\du}{11.575657\du}}
\pgfpathlineto{\pgfpoint{-17.709070\du}{11.575657\du}}
\pgfpathlineto{\pgfpoint{-17.758763\du}{11.575657\du}}
\pgfpathlineto{\pgfpoint{-17.816627\du}{11.575657\du}}
\pgfpathlineto{\pgfpoint{-17.816627\du}{11.609006\du}}
\pgfpathlineto{\pgfpoint{-17.816627\du}{11.642356\du}}
\pgfpathlineto{\pgfpoint{-17.816627\du}{11.658920\du}}
\pgfpathlineto{\pgfpoint{-17.816627\du}{11.675484\du}}
\pgfpathlineto{\pgfpoint{-17.816627\du}{11.683656\du}}
\pgfusepath{fill}
\pgfsetbuttcap
\pgfsetmiterjoin
\pgfsetdash{}{0pt}
\definecolor{dialinecolor}{rgb}{0.788235, 0.788235, 0.713726}
\pgfsetfillcolor{dialinecolor}
\pgfpathmoveto{\pgfpoint{-18.181264\du}{10.195516\du}}
\pgfpathlineto{\pgfpoint{-18.206000\du}{10.220252\du}}
\pgfpathlineto{\pgfpoint{-17.651205\du}{10.220252\du}}
\pgfpathlineto{\pgfpoint{-17.651205\du}{10.195516\du}}
\pgfpathlineto{\pgfpoint{-18.181264\du}{10.195516\du}}
\pgfusepath{fill}
\pgfsetbuttcap
\pgfsetmiterjoin
\pgfsetdash{}{0pt}
\definecolor{dialinecolor}{rgb}{0.788235, 0.788235, 0.713726}
\pgfsetfillcolor{dialinecolor}
\pgfpathmoveto{\pgfpoint{-17.924848\du}{10.211859\du}}
\pgfpathlineto{\pgfpoint{-18.206000\du}{10.211859\du}}
\pgfpathlineto{\pgfpoint{-18.205779\du}{10.217822\du}}
\pgfpathlineto{\pgfpoint{-18.204675\du}{10.224448\du}}
\pgfpathlineto{\pgfpoint{-18.202908\du}{10.230632\du}}
\pgfpathlineto{\pgfpoint{-18.200479\du}{10.237037\du}}
\pgfpathlineto{\pgfpoint{-18.196945\du}{10.243000\du}}
\pgfpathlineto{\pgfpoint{-18.193190\du}{10.249184\du}}
\pgfpathlineto{\pgfpoint{-18.188552\du}{10.255147\du}}
\pgfpathlineto{\pgfpoint{-18.183252\du}{10.260889\du}}
\pgfpathlineto{\pgfpoint{-18.177509\du}{10.266411\du}}
\pgfpathlineto{\pgfpoint{-18.170884\du}{10.272153\du}}
\pgfpathlineto{\pgfpoint{-18.163595\du}{10.277675\du}}
\pgfpathlineto{\pgfpoint{-18.155644\du}{10.282754\du}}
\pgfpathlineto{\pgfpoint{-18.147252\du}{10.287834\du}}
\pgfpathlineto{\pgfpoint{-18.138197\du}{10.292914\du}}
\pgfpathlineto{\pgfpoint{-18.128700\du}{10.297552\du}}
\pgfpathlineto{\pgfpoint{-18.118761\du}{10.301969\du}}
\pgfpathlineto{\pgfpoint{-18.107939\du}{10.306165\du}}
\pgfpathlineto{\pgfpoint{-18.096896\du}{10.310141\du}}
\pgfpathlineto{\pgfpoint{-18.085632\du}{10.313895\du}}
\pgfpathlineto{\pgfpoint{-18.073485\du}{10.317208\du}}
\pgfpathlineto{\pgfpoint{-18.061338\du}{10.320521\du}}
\pgfpathlineto{\pgfpoint{-18.048749\du}{10.323613\du}}
\pgfpathlineto{\pgfpoint{-18.035719\du}{10.326263\du}}
\pgfpathlineto{\pgfpoint{-18.022467\du}{10.328251\du}}
\pgfpathlineto{\pgfpoint{-18.009216\du}{10.330681\du}}
\pgfpathlineto{\pgfpoint{-17.995302\du}{10.332227\du}}
\pgfpathlineto{\pgfpoint{-17.981608\du}{10.333773\du}}
\pgfpathlineto{\pgfpoint{-17.967473\du}{10.334656\du}}
\pgfpathlineto{\pgfpoint{-17.953118\du}{10.335539\du}}
\pgfpathlineto{\pgfpoint{-17.938983\du}{10.335981\du}}
\pgfpathlineto{\pgfpoint{-17.924848\du}{10.336202\du}}
\pgfpathlineto{\pgfpoint{-17.910713\du}{10.335981\du}}
\pgfpathlineto{\pgfpoint{-17.896578\du}{10.335539\du}}
\pgfpathlineto{\pgfpoint{-17.882443\du}{10.334656\du}}
\pgfpathlineto{\pgfpoint{-17.868308\du}{10.333773\du}}
\pgfpathlineto{\pgfpoint{-17.854394\du}{10.332227\du}}
\pgfpathlineto{\pgfpoint{-17.840701\du}{10.330681\du}}
\pgfpathlineto{\pgfpoint{-17.827450\du}{10.328251\du}}
\pgfpathlineto{\pgfpoint{-17.813977\du}{10.326263\du}}
\pgfpathlineto{\pgfpoint{-17.800947\du}{10.323613\du}}
\pgfpathlineto{\pgfpoint{-17.788358\du}{10.320521\du}}
\pgfpathlineto{\pgfpoint{-17.775990\du}{10.317208\du}}
\pgfpathlineto{\pgfpoint{-17.764063\du}{10.313895\du}}
\pgfpathlineto{\pgfpoint{-17.752579\du}{10.310141\du}}
\pgfpathlineto{\pgfpoint{-17.741536\du}{10.306165\du}}
\pgfpathlineto{\pgfpoint{-17.731155\du}{10.301969\du}}
\pgfpathlineto{\pgfpoint{-17.720775\du}{10.297552\du}}
\pgfpathlineto{\pgfpoint{-17.711720\du}{10.292914\du}}
\pgfpathlineto{\pgfpoint{-17.702444\du}{10.287834\du}}
\pgfpathlineto{\pgfpoint{-17.693830\du}{10.282754\du}}
\pgfpathlineto{\pgfpoint{-17.686100\du}{10.277675\du}}
\pgfpathlineto{\pgfpoint{-17.678812\du}{10.272153\du}}
\pgfpathlineto{\pgfpoint{-17.672186\du}{10.266411\du}}
\pgfpathlineto{\pgfpoint{-17.666665\du}{10.260889\du}}
\pgfpathlineto{\pgfpoint{-17.661364\du}{10.255147\du}}
\pgfpathlineto{\pgfpoint{-17.656506\du}{10.249184\du}}
\pgfpathlineto{\pgfpoint{-17.652751\du}{10.243000\du}}
\pgfpathlineto{\pgfpoint{-17.649438\du}{10.237037\du}}
\pgfpathlineto{\pgfpoint{-17.646788\du}{10.230632\du}}
\pgfpathlineto{\pgfpoint{-17.645242\du}{10.224448\du}}
\pgfpathlineto{\pgfpoint{-17.643917\du}{10.217822\du}}
\pgfpathlineto{\pgfpoint{-17.643475\du}{10.211859\du}}
\pgfpathlineto{\pgfpoint{-17.924848\du}{10.211859\du}}
\pgfusepath{fill}
\pgfsetbuttcap
\pgfsetmiterjoin
\pgfsetdash{}{0pt}
\definecolor{dialinecolor}{rgb}{0.286275, 0.286275, 0.211765}
\pgfsetstrokecolor{dialinecolor}
\pgfpathmoveto{\pgfpoint{-18.206000\du}{10.211859\du}}
\pgfpathlineto{\pgfpoint{-18.205779\du}{10.217822\du}}
\pgfpathlineto{\pgfpoint{-18.204675\du}{10.224448\du}}
\pgfpathlineto{\pgfpoint{-18.202908\du}{10.230632\du}}
\pgfpathlineto{\pgfpoint{-18.200479\du}{10.237037\du}}
\pgfpathlineto{\pgfpoint{-18.196945\du}{10.243000\du}}
\pgfpathlineto{\pgfpoint{-18.193190\du}{10.249184\du}}
\pgfpathlineto{\pgfpoint{-18.188552\du}{10.255147\du}}
\pgfpathlineto{\pgfpoint{-18.183252\du}{10.260889\du}}
\pgfpathlineto{\pgfpoint{-18.177509\du}{10.266411\du}}
\pgfpathlineto{\pgfpoint{-18.170884\du}{10.272153\du}}
\pgfpathlineto{\pgfpoint{-18.163595\du}{10.277675\du}}
\pgfpathlineto{\pgfpoint{-18.155644\du}{10.282754\du}}
\pgfpathlineto{\pgfpoint{-18.147252\du}{10.287834\du}}
\pgfpathlineto{\pgfpoint{-18.138197\du}{10.292914\du}}
\pgfpathlineto{\pgfpoint{-18.128479\du}{10.297552\du}}
\pgfpathlineto{\pgfpoint{-18.118319\du}{10.301969\du}}
\pgfpathlineto{\pgfpoint{-18.107939\du}{10.306165\du}}
\pgfpathlineto{\pgfpoint{-18.096896\du}{10.310141\du}}
\pgfpathlineto{\pgfpoint{-18.085412\du}{10.313895\du}}
\pgfpathlineto{\pgfpoint{-18.073485\du}{10.317208\du}}
\pgfpathlineto{\pgfpoint{-18.061117\du}{10.320521\du}}
\pgfpathlineto{\pgfpoint{-18.048749\du}{10.323613\du}}
\pgfpathlineto{\pgfpoint{-18.035719\du}{10.326263\du}}
\pgfpathlineto{\pgfpoint{-18.022467\du}{10.328251\du}}
\pgfpathlineto{\pgfpoint{-18.008995\du}{10.330681\du}}
\pgfpathlineto{\pgfpoint{-17.995081\du}{10.332227\du}}
\pgfpathlineto{\pgfpoint{-17.981608\du}{10.333773\du}}
\pgfpathlineto{\pgfpoint{-17.967473\du}{10.334656\du}}
\pgfpathlineto{\pgfpoint{-17.953118\du}{10.335539\du}}
\pgfpathlineto{\pgfpoint{-17.938983\du}{10.335981\du}}
\pgfpathlineto{\pgfpoint{-17.924406\du}{10.336202\du}}
\pgfpathlineto{\pgfpoint{-17.910492\du}{10.335981\du}}
\pgfpathlineto{\pgfpoint{-17.896136\du}{10.335539\du}}
\pgfpathlineto{\pgfpoint{-17.882222\du}{10.334656\du}}
\pgfpathlineto{\pgfpoint{-17.867867\du}{10.333773\du}}
\pgfpathlineto{\pgfpoint{-17.854394\du}{10.332227\du}}
\pgfpathlineto{\pgfpoint{-17.840701\du}{10.330681\du}}
\pgfpathlineto{\pgfpoint{-17.827008\du}{10.328251\du}}
\pgfpathlineto{\pgfpoint{-17.813756\du}{10.326263\du}}
\pgfpathlineto{\pgfpoint{-17.800726\du}{10.323613\du}}
\pgfpathlineto{\pgfpoint{-17.788358\du}{10.320521\du}}
\pgfpathlineto{\pgfpoint{-17.775990\du}{10.317208\du}}
\pgfpathlineto{\pgfpoint{-17.764063\du}{10.313895\du}}
\pgfpathlineto{\pgfpoint{-17.752579\du}{10.310141\du}}
\pgfpathlineto{\pgfpoint{-17.741536\du}{10.306165\du}}
\pgfpathlineto{\pgfpoint{-17.730493\du}{10.301969\du}}
\pgfpathlineto{\pgfpoint{-17.720775\du}{10.297552\du}}
\pgfpathlineto{\pgfpoint{-17.711278\du}{10.292914\du}}
\pgfpathlineto{\pgfpoint{-17.702223\du}{10.287834\du}}
\pgfpathlineto{\pgfpoint{-17.693610\du}{10.282754\du}}
\pgfpathlineto{\pgfpoint{-17.685880\du}{10.277675\du}}
\pgfpathlineto{\pgfpoint{-17.678591\du}{10.272153\du}}
\pgfpathlineto{\pgfpoint{-17.671745\du}{10.266411\du}}
\pgfpathlineto{\pgfpoint{-17.666002\du}{10.260889\du}}
\pgfpathlineto{\pgfpoint{-17.660702\du}{10.255147\du}}
\pgfpathlineto{\pgfpoint{-17.656285\du}{10.249184\du}}
\pgfpathlineto{\pgfpoint{-17.652530\du}{10.243000\du}}
\pgfpathlineto{\pgfpoint{-17.648996\du}{10.237037\du}}
\pgfpathlineto{\pgfpoint{-17.646567\du}{10.230632\du}}
\pgfpathlineto{\pgfpoint{-17.645021\du}{10.224448\du}}
\pgfpathlineto{\pgfpoint{-17.643475\du}{10.217822\du}}
\pgfpathlineto{\pgfpoint{-17.643254\du}{10.211859\du}}
\pgfusepath{stroke}
\pgfsetbuttcap
\pgfsetmiterjoin
\pgfsetdash{}{0pt}
\definecolor{dialinecolor}{rgb}{0.478431, 0.478431, 0.352941}
\pgfsetstrokecolor{dialinecolor}
\pgfpathmoveto{\pgfpoint{-18.156307\du}{10.212080\du}}
\pgfpathlineto{\pgfpoint{-17.700898\du}{10.212080\du}}
\pgfusepath{stroke}
\pgfsetbuttcap
\pgfsetmiterjoin
\pgfsetdash{}{0pt}
\definecolor{dialinecolor}{rgb}{0.576471, 0.576471, 0.423529}
\pgfsetfillcolor{dialinecolor}
\pgfpathmoveto{\pgfpoint{-17.642812\du}{11.575657\du}}
\pgfpathlineto{\pgfpoint{-17.642812\du}{10.220252\du}}
\pgfpathlineto{\pgfpoint{-17.543205\du}{10.220252\du}}
\pgfpathlineto{\pgfpoint{-17.543205\du}{11.575657\du}}
\pgfpathlineto{\pgfpoint{-17.642812\du}{11.575657\du}}
\pgfusepath{fill}
\pgfsetbuttcap
\pgfsetmiterjoin
\pgfsetdash{}{0pt}
\definecolor{dialinecolor}{rgb}{0.576471, 0.576471, 0.423529}
\pgfsetfillcolor{dialinecolor}
\pgfpathmoveto{\pgfpoint{-17.559991\du}{10.619342\du}}
\pgfpathlineto{\pgfpoint{-17.559991\du}{10.303515\du}}
\pgfpathlineto{\pgfpoint{-17.187182\du}{10.303515\du}}
\pgfpathlineto{\pgfpoint{-17.187182\du}{10.619342\du}}
\pgfpathlineto{\pgfpoint{-17.559991\du}{10.619342\du}}
\pgfusepath{fill}
\pgfsetbuttcap
\pgfsetmiterjoin
\pgfsetdash{}{0pt}
\definecolor{dialinecolor}{rgb}{0.286275, 0.286275, 0.211765}
\pgfsetstrokecolor{dialinecolor}
\pgfpathmoveto{\pgfpoint{-17.559991\du}{10.619342\du}}
\pgfpathlineto{\pgfpoint{-17.559991\du}{10.303515\du}}
\pgfpathlineto{\pgfpoint{-17.187182\du}{10.303515\du}}
\pgfpathlineto{\pgfpoint{-17.187182\du}{10.619342\du}}
\pgfpathlineto{\pgfpoint{-17.559991\du}{10.619342\du}}
\pgfusepath{stroke}
\pgfsetbuttcap
\pgfsetmiterjoin
\pgfsetdash{}{0pt}
\definecolor{dialinecolor}{rgb}{0.478431, 0.478431, 0.352941}
\pgfsetfillcolor{dialinecolor}
\pgfpathmoveto{\pgfpoint{-17.551819\du}{10.619342\du}}
\pgfpathlineto{\pgfpoint{-17.551819\du}{10.312129\du}}
\pgfpathlineto{\pgfpoint{-17.543205\du}{10.312129\du}}
\pgfpathlineto{\pgfpoint{-17.543205\du}{10.619342\du}}
\pgfpathlineto{\pgfpoint{-17.551819\du}{10.619342\du}}
\pgfusepath{fill}
\pgfsetbuttcap
\pgfsetmiterjoin
\pgfsetdash{}{0pt}
\definecolor{dialinecolor}{rgb}{0.831373, 0.000000, 0.039216}
\pgfsetfillcolor{dialinecolor}
\pgfpathmoveto{\pgfpoint{-18.197607\du}{10.237037\du}}
\pgfpathlineto{\pgfpoint{-18.172871\du}{10.261994\du}}
\pgfpathlineto{\pgfpoint{-18.172871\du}{11.625571\du}}
\pgfpathlineto{\pgfpoint{-18.197607\du}{11.600834\du}}
\pgfpathlineto{\pgfpoint{-18.197607\du}{10.237037\du}}
\pgfusepath{fill}
\pgfsetbuttcap
\pgfsetmiterjoin
\pgfsetdash{}{0pt}
\definecolor{dialinecolor}{rgb}{0.286275, 0.286275, 0.211765}
\pgfsetstrokecolor{dialinecolor}
\pgfpathmoveto{\pgfpoint{-18.206000\du}{10.211859\du}}
\pgfpathlineto{\pgfpoint{-18.205779\du}{10.217822\du}}
\pgfpathlineto{\pgfpoint{-18.204675\du}{10.224448\du}}
\pgfpathlineto{\pgfpoint{-18.202908\du}{10.230632\du}}
\pgfpathlineto{\pgfpoint{-18.200479\du}{10.237037\du}}
\pgfpathlineto{\pgfpoint{-18.196945\du}{10.243000\du}}
\pgfpathlineto{\pgfpoint{-18.193190\du}{10.249184\du}}
\pgfpathlineto{\pgfpoint{-18.188552\du}{10.255147\du}}
\pgfpathlineto{\pgfpoint{-18.183252\du}{10.260889\du}}
\pgfpathlineto{\pgfpoint{-18.177509\du}{10.266411\du}}
\pgfpathlineto{\pgfpoint{-18.170884\du}{10.272153\du}}
\pgfpathlineto{\pgfpoint{-18.163595\du}{10.277675\du}}
\pgfpathlineto{\pgfpoint{-18.155644\du}{10.282754\du}}
\pgfpathlineto{\pgfpoint{-18.147252\du}{10.287834\du}}
\pgfpathlineto{\pgfpoint{-18.138197\du}{10.292914\du}}
\pgfpathlineto{\pgfpoint{-18.128479\du}{10.297552\du}}
\pgfpathlineto{\pgfpoint{-18.118319\du}{10.301969\du}}
\pgfpathlineto{\pgfpoint{-18.107939\du}{10.306165\du}}
\pgfpathlineto{\pgfpoint{-18.096896\du}{10.310141\du}}
\pgfpathlineto{\pgfpoint{-18.085412\du}{10.313895\du}}
\pgfpathlineto{\pgfpoint{-18.073485\du}{10.317208\du}}
\pgfpathlineto{\pgfpoint{-18.061117\du}{10.320521\du}}
\pgfpathlineto{\pgfpoint{-18.048749\du}{10.323613\du}}
\pgfpathlineto{\pgfpoint{-18.035719\du}{10.326263\du}}
\pgfpathlineto{\pgfpoint{-18.022467\du}{10.328251\du}}
\pgfpathlineto{\pgfpoint{-18.008995\du}{10.330681\du}}
\pgfpathlineto{\pgfpoint{-17.995081\du}{10.332227\du}}
\pgfpathlineto{\pgfpoint{-17.981608\du}{10.333773\du}}
\pgfpathlineto{\pgfpoint{-17.967473\du}{10.334656\du}}
\pgfpathlineto{\pgfpoint{-17.953118\du}{10.335539\du}}
\pgfpathlineto{\pgfpoint{-17.938983\du}{10.335981\du}}
\pgfpathlineto{\pgfpoint{-17.924406\du}{10.336202\du}}
\pgfpathlineto{\pgfpoint{-17.910492\du}{10.335981\du}}
\pgfpathlineto{\pgfpoint{-17.896136\du}{10.335539\du}}
\pgfpathlineto{\pgfpoint{-17.882222\du}{10.334656\du}}
\pgfpathlineto{\pgfpoint{-17.867867\du}{10.333773\du}}
\pgfpathlineto{\pgfpoint{-17.854394\du}{10.332227\du}}
\pgfpathlineto{\pgfpoint{-17.840701\du}{10.330681\du}}
\pgfpathlineto{\pgfpoint{-17.827008\du}{10.328251\du}}
\pgfpathlineto{\pgfpoint{-17.813756\du}{10.326263\du}}
\pgfpathlineto{\pgfpoint{-17.800726\du}{10.323613\du}}
\pgfpathlineto{\pgfpoint{-17.788358\du}{10.320521\du}}
\pgfpathlineto{\pgfpoint{-17.775990\du}{10.317208\du}}
\pgfpathlineto{\pgfpoint{-17.764063\du}{10.313895\du}}
\pgfpathlineto{\pgfpoint{-17.752579\du}{10.310141\du}}
\pgfpathlineto{\pgfpoint{-17.741536\du}{10.306165\du}}
\pgfpathlineto{\pgfpoint{-17.730493\du}{10.301969\du}}
\pgfpathlineto{\pgfpoint{-17.720775\du}{10.297552\du}}
\pgfpathlineto{\pgfpoint{-17.711278\du}{10.292914\du}}
\pgfpathlineto{\pgfpoint{-17.702223\du}{10.287834\du}}
\pgfpathlineto{\pgfpoint{-17.693610\du}{10.282754\du}}
\pgfpathlineto{\pgfpoint{-17.685880\du}{10.277675\du}}
\pgfpathlineto{\pgfpoint{-17.678591\du}{10.272153\du}}
\pgfpathlineto{\pgfpoint{-17.671745\du}{10.266411\du}}
\pgfpathlineto{\pgfpoint{-17.666002\du}{10.260889\du}}
\pgfpathlineto{\pgfpoint{-17.660702\du}{10.255147\du}}
\pgfpathlineto{\pgfpoint{-17.656285\du}{10.249184\du}}
\pgfpathlineto{\pgfpoint{-17.652530\du}{10.243000\du}}
\pgfpathlineto{\pgfpoint{-17.648996\du}{10.237037\du}}
\pgfpathlineto{\pgfpoint{-17.646567\du}{10.230632\du}}
\pgfpathlineto{\pgfpoint{-17.645021\du}{10.224448\du}}
\pgfpathlineto{\pgfpoint{-17.643475\du}{10.217822\du}}
\pgfpathlineto{\pgfpoint{-17.643254\du}{10.211859\du}}
\pgfusepath{stroke}
\pgfsetbuttcap
\pgfsetmiterjoin
\pgfsetdash{}{0pt}
\definecolor{dialinecolor}{rgb}{0.286275, 0.286275, 0.211765}
\pgfsetstrokecolor{dialinecolor}
\pgfpathmoveto{\pgfpoint{-18.206000\du}{11.579190\du}}
\pgfpathlineto{\pgfpoint{-18.205779\du}{11.585154\du}}
\pgfpathlineto{\pgfpoint{-18.204896\du}{11.591117\du}}
\pgfpathlineto{\pgfpoint{-18.202908\du}{11.597301\du}}
\pgfpathlineto{\pgfpoint{-18.200479\du}{11.603043\du}}
\pgfpathlineto{\pgfpoint{-18.197607\du}{11.609006\du}}
\pgfpathlineto{\pgfpoint{-18.193853\du}{11.614969\du}}
\pgfpathlineto{\pgfpoint{-18.189436\du}{11.620491\du}}
\pgfpathlineto{\pgfpoint{-18.183914\du}{11.626233\du}}
\pgfpathlineto{\pgfpoint{-18.178834\du}{11.631534\du}}
\pgfpathlineto{\pgfpoint{-18.172209\du}{11.637276\du}}
\pgfpathlineto{\pgfpoint{-18.165362\du}{11.642356\du}}
\pgfpathlineto{\pgfpoint{-18.157411\du}{11.647215\du}}
\pgfpathlineto{\pgfpoint{-18.149240\du}{11.652294\du}}
\pgfpathlineto{\pgfpoint{-18.140405\du}{11.656932\du}}
\pgfpathlineto{\pgfpoint{-18.130908\du}{11.661350\du}}
\pgfpathlineto{\pgfpoint{-18.121191\du}{11.665767\du}}
\pgfpathlineto{\pgfpoint{-18.111252\du}{11.669963\du}}
\pgfpathlineto{\pgfpoint{-18.100209\du}{11.673938\du}}
\pgfpathlineto{\pgfpoint{-18.088724\du}{11.677472\du}}
\pgfpathlineto{\pgfpoint{-18.077240\du}{11.681006\du}}
\pgfpathlineto{\pgfpoint{-18.064872\du}{11.684098\du}}
\pgfpathlineto{\pgfpoint{-18.052725\du}{11.686969\du}}
\pgfpathlineto{\pgfpoint{-18.040136\du}{11.689398\du}}
\pgfpathlineto{\pgfpoint{-18.026884\du}{11.691828\du}}
\pgfpathlineto{\pgfpoint{-18.013633\du}{11.693816\du}}
\pgfpathlineto{\pgfpoint{-18.000160\du}{11.695582\du}}
\pgfpathlineto{\pgfpoint{-17.986467\du}{11.697349\du}}
\pgfpathlineto{\pgfpoint{-17.972774\du}{11.698233\du}}
\pgfpathlineto{\pgfpoint{-17.958639\du}{11.699116\du}}
\pgfpathlineto{\pgfpoint{-17.944504\du}{11.699779\du}}
\pgfpathlineto{\pgfpoint{-17.929928\du}{11.700000\du}}
\pgfpathlineto{\pgfpoint{-17.916455\du}{11.700000\du}}
\pgfpathlineto{\pgfpoint{-17.902320\du}{11.699779\du}}
\pgfpathlineto{\pgfpoint{-17.888185\du}{11.699116\du}}
\pgfpathlineto{\pgfpoint{-17.874051\du}{11.698233\du}}
\pgfpathlineto{\pgfpoint{-17.860137\du}{11.696908\du}}
\pgfpathlineto{\pgfpoint{-17.846885\du}{11.695362\du}}
\pgfpathlineto{\pgfpoint{-17.833413\du}{11.693595\du}}
\pgfpathlineto{\pgfpoint{-17.819940\du}{11.691386\du}}
\pgfpathlineto{\pgfpoint{-17.806910\du}{11.689178\du}}
\pgfpathlineto{\pgfpoint{-17.794321\du}{11.686306\du}}
\pgfpathlineto{\pgfpoint{-17.782174\du}{11.683435\du}}
\pgfpathlineto{\pgfpoint{-17.769806\du}{11.680343\du}}
\pgfpathlineto{\pgfpoint{-17.758542\du}{11.676810\du}}
\pgfpathlineto{\pgfpoint{-17.747057\du}{11.673055\du}}
\pgfpathlineto{\pgfpoint{-17.736456\du}{11.669300\du}}
\pgfpathlineto{\pgfpoint{-17.726297\du}{11.665104\du}}
\pgfpathlineto{\pgfpoint{-17.716579\du}{11.660908\du}}
\pgfpathlineto{\pgfpoint{-17.707303\du}{11.656049\du}}
\pgfpathlineto{\pgfpoint{-17.698689\du}{11.651411\du}}
\pgfpathlineto{\pgfpoint{-17.690518\du}{11.646331\du}}
\pgfpathlineto{\pgfpoint{-17.682788\du}{11.641251\du}}
\pgfpathlineto{\pgfpoint{-17.675720\du}{11.635951\du}}
\pgfpathlineto{\pgfpoint{-17.669536\du}{11.630650\du}}
\pgfpathlineto{\pgfpoint{-17.663794\du}{11.624908\du}}
\pgfpathlineto{\pgfpoint{-17.659156\du}{11.619387\du}}
\pgfpathlineto{\pgfpoint{-17.654960\du}{11.613644\du}}
\pgfpathlineto{\pgfpoint{-17.651205\du}{11.607902\du}}
\pgfpathlineto{\pgfpoint{-17.648113\du}{11.601939\du}}
\pgfpathlineto{\pgfpoint{-17.645904\du}{11.595976\du}}
\pgfpathlineto{\pgfpoint{-17.644137\du}{11.590012\du}}
\pgfpathlineto{\pgfpoint{-17.643475\du}{11.584049\du}}
\pgfpathlineto{\pgfpoint{-17.643254\du}{11.577865\du}}
\pgfusepath{stroke}
\pgfsetlinewidth{0.100000\du}
\pgfsetdash{}{0pt}
\pgfsetdash{}{0pt}
\pgfsetbuttcap
{
\definecolor{dialinecolor}{rgb}{0.000000, 0.000000, 0.000000}
\pgfsetfillcolor{dialinecolor}
% was here!!!
\definecolor{dialinecolor}{rgb}{0.000000, 0.000000, 0.000000}
\pgfsetstrokecolor{dialinecolor}
\draw (-19.206613\du,-2.779780\du)--(-19.152256\du,-0.341000\du);
}
\pgfsetlinewidth{0.100000\du}
\pgfsetdash{}{0pt}
\pgfsetdash{}{0pt}
\pgfsetbuttcap
{
\definecolor{dialinecolor}{rgb}{0.000000, 0.000000, 0.000000}
\pgfsetfillcolor{dialinecolor}
% was here!!!
\definecolor{dialinecolor}{rgb}{0.000000, 0.000000, 0.000000}
\pgfsetstrokecolor{dialinecolor}
\draw (-19.759648\du,0.630507\du)--(-22.494322\du,2.629267\du);
}
\pgfsetlinewidth{0.100000\du}
\pgfsetdash{}{0pt}
\pgfsetdash{}{0pt}
\pgfsetbuttcap
{
\definecolor{dialinecolor}{rgb}{0.000000, 0.000000, 0.000000}
\pgfsetfillcolor{dialinecolor}
% was here!!!
\definecolor{dialinecolor}{rgb}{0.000000, 0.000000, 0.000000}
\pgfsetstrokecolor{dialinecolor}
\draw (-4.510683\du,-2.918716\du)--(-4.428314\du,-0.861800\du);
}
\pgfsetlinewidth{0.100000\du}
\pgfsetdash{}{0pt}
\pgfsetdash{}{0pt}
\pgfsetbuttcap
\pgfsetmiterjoin
\pgfsetlinewidth{0.001000\du}
\pgfsetbuttcap
\pgfsetmiterjoin
\pgfsetdash{}{0pt}
\definecolor{dialinecolor}{rgb}{0.000000, 0.470588, 0.666667}
\pgfsetfillcolor{dialinecolor}
\pgfpathmoveto{\pgfpoint{-3.513701\du}{-0.116140\du}}
\pgfpathlineto{\pgfpoint{-3.515019\du}{-0.131691\du}}
\pgfpathlineto{\pgfpoint{-3.518446\du}{-0.147506\du}}
\pgfpathlineto{\pgfpoint{-3.524244\du}{-0.163057\du}}
\pgfpathlineto{\pgfpoint{-3.532415\du}{-0.178345\du}}
\pgfpathlineto{\pgfpoint{-3.542958\du}{-0.193896\du}}
\pgfpathlineto{\pgfpoint{-3.555610\du}{-0.209183\du}}
\pgfpathlineto{\pgfpoint{-3.570634\du}{-0.223943\du}}
\pgfpathlineto{\pgfpoint{-3.587503\du}{-0.238440\du}}
\pgfpathlineto{\pgfpoint{-3.607008\du}{-0.252673\du}}
\pgfpathlineto{\pgfpoint{-3.628357\du}{-0.266379\du}}
\pgfpathlineto{\pgfpoint{-3.651816\du}{-0.280085\du}}
\pgfpathlineto{\pgfpoint{-3.677383\du}{-0.293264\du}}
\pgfpathlineto{\pgfpoint{-3.704795\du}{-0.305916\du}}
\pgfpathlineto{\pgfpoint{-3.734052\du}{-0.318304\du}}
\pgfpathlineto{\pgfpoint{-3.765154\du}{-0.329901\du}}
\pgfpathlineto{\pgfpoint{-3.797574\du}{-0.340972\du}}
\pgfpathlineto{\pgfpoint{-3.831839\du}{-0.351515\du}}
\pgfpathlineto{\pgfpoint{-3.867949\du}{-0.361267\du}}
\pgfpathlineto{\pgfpoint{-3.905113\du}{-0.371019\du}}
\pgfpathlineto{\pgfpoint{-3.943596\du}{-0.379190\du}}
\pgfpathlineto{\pgfpoint{-3.983396\du}{-0.387361\du}}
\pgfpathlineto{\pgfpoint{-4.024514\du}{-0.394478\du}}
\pgfpathlineto{\pgfpoint{-4.066423\du}{-0.401331\du}}
\pgfpathlineto{\pgfpoint{-4.109649\du}{-0.407129\du}}
\pgfpathlineto{\pgfpoint{-4.153139\du}{-0.412137\du}}
\pgfpathlineto{\pgfpoint{-4.197948\du}{-0.416355\du}}
\pgfpathlineto{\pgfpoint{-4.243019\du}{-0.420045\du}}
\pgfpathlineto{\pgfpoint{-4.288355\du}{-0.422944\du}}
\pgfpathlineto{\pgfpoint{-4.334217\du}{-0.425053\du}}
\pgfpathlineto{\pgfpoint{-4.380343\du}{-0.425843\du}}
\pgfpathlineto{\pgfpoint{-4.426996\du}{-0.426634\du}}
\pgfpathlineto{\pgfpoint{-4.426996\du}{-0.426634\du}}
\pgfpathlineto{\pgfpoint{-4.473122\du}{-0.425843\du}}
\pgfpathlineto{\pgfpoint{-4.518984\du}{-0.425053\du}}
\pgfpathlineto{\pgfpoint{-4.564847\du}{-0.422944\du}}
\pgfpathlineto{\pgfpoint{-4.610182\du}{-0.420045\du}}
\pgfpathlineto{\pgfpoint{-4.655254\du}{-0.416355\du}}
\pgfpathlineto{\pgfpoint{-4.700062\du}{-0.412137\du}}
\pgfpathlineto{\pgfpoint{-4.743552\du}{-0.407129\du}}
\pgfpathlineto{\pgfpoint{-4.786779\du}{-0.401331\du}}
\pgfpathlineto{\pgfpoint{-4.828688\du}{-0.394478\du}}
\pgfpathlineto{\pgfpoint{-4.869806\du}{-0.387361\du}}
\pgfpathlineto{\pgfpoint{-4.909869\du}{-0.379190\du}}
\pgfpathlineto{\pgfpoint{-4.948088\du}{-0.371019\du}}
\pgfpathlineto{\pgfpoint{-4.985252\du}{-0.361267\du}}
\pgfpathlineto{\pgfpoint{-5.021362\du}{-0.351515\du}}
\pgfpathlineto{\pgfpoint{-5.055627\du}{-0.340972\du}}
\pgfpathlineto{\pgfpoint{-5.088047\du}{-0.329901\du}}
\pgfpathlineto{\pgfpoint{-5.119150\du}{-0.318304\du}}
\pgfpathlineto{\pgfpoint{-5.148407\du}{-0.305916\du}}
\pgfpathlineto{\pgfpoint{-5.175819\du}{-0.293264\du}}
\pgfpathlineto{\pgfpoint{-5.201649\du}{-0.280085\du}}
\pgfpathlineto{\pgfpoint{-5.224844\du}{-0.266379\du}}
\pgfpathlineto{\pgfpoint{-5.246457\du}{-0.252673\du}}
\pgfpathlineto{\pgfpoint{-5.265698\du}{-0.238440\du}}
\pgfpathlineto{\pgfpoint{-5.282567\du}{-0.223943\du}}
\pgfpathlineto{\pgfpoint{-5.297591\du}{-0.209183\du}}
\pgfpathlineto{\pgfpoint{-5.310243\du}{-0.193896\du}}
\pgfpathlineto{\pgfpoint{-5.320786\du}{-0.178345\du}}
\pgfpathlineto{\pgfpoint{-5.328957\du}{-0.163057\du}}
\pgfpathlineto{\pgfpoint{-5.335019\du}{-0.147506\du}}
\pgfpathlineto{\pgfpoint{-5.338182\du}{-0.131691\du}}
\pgfpathlineto{\pgfpoint{-5.339500\du}{-0.116140\du}}
\pgfpathlineto{\pgfpoint{-5.339500\du}{-0.116140\du}}
\pgfpathlineto{\pgfpoint{-5.338182\du}{-0.100062\du}}
\pgfpathlineto{\pgfpoint{-5.335019\du}{-0.084511\du}}
\pgfpathlineto{\pgfpoint{-5.328957\du}{-0.068960\du}}
\pgfpathlineto{\pgfpoint{-5.320786\du}{-0.053409\du}}
\pgfpathlineto{\pgfpoint{-5.310243\du}{-0.038122\du}}
\pgfpathlineto{\pgfpoint{-5.297591\du}{-0.022834\du}}
\pgfpathlineto{\pgfpoint{-5.282567\du}{-0.008074\du}}
\pgfpathlineto{\pgfpoint{-5.265698\du}{0.006423\du}}
\pgfpathlineto{\pgfpoint{-5.246457\du}{0.020656\du}}
\pgfpathlineto{\pgfpoint{-5.224844\du}{0.034626\du}}
\pgfpathlineto{\pgfpoint{-5.201649\du}{0.047805\du}}
\pgfpathlineto{\pgfpoint{-5.175819\du}{0.061511\du}}
\pgfpathlineto{\pgfpoint{-5.148407\du}{0.074162\du}}
\pgfpathlineto{\pgfpoint{-5.119150\du}{0.086287\du}}
\pgfpathlineto{\pgfpoint{-5.088047\du}{0.097621\du}}
\pgfpathlineto{\pgfpoint{-5.055627\du}{0.108954\du}}
\pgfpathlineto{\pgfpoint{-5.021362\du}{0.119498\du}}
\pgfpathlineto{\pgfpoint{-4.985252\du}{0.129513\du}}
\pgfpathlineto{\pgfpoint{-4.948088\du}{0.138739\du}}
\pgfpathlineto{\pgfpoint{-4.909869\du}{0.147437\du}}
\pgfpathlineto{\pgfpoint{-4.869806\du}{0.155344\du}}
\pgfpathlineto{\pgfpoint{-4.828688\du}{0.162461\du}}
\pgfpathlineto{\pgfpoint{-4.786779\du}{0.169314\du}}
\pgfpathlineto{\pgfpoint{-4.743552\du}{0.175112\du}}
\pgfpathlineto{\pgfpoint{-4.700062\du}{0.180120\du}}
\pgfpathlineto{\pgfpoint{-4.655254\du}{0.184601\du}}
\pgfpathlineto{\pgfpoint{-4.610182\du}{0.187764\du}}
\pgfpathlineto{\pgfpoint{-4.564847\du}{0.190663\du}}
\pgfpathlineto{\pgfpoint{-4.518984\du}{0.192772\du}}
\pgfpathlineto{\pgfpoint{-4.473122\du}{0.194090\du}}
\pgfpathlineto{\pgfpoint{-4.426996\du}{0.194353\du}}
\pgfpathlineto{\pgfpoint{-4.426996\du}{0.194353\du}}
\pgfpathlineto{\pgfpoint{-4.380343\du}{0.194090\du}}
\pgfpathlineto{\pgfpoint{-4.334217\du}{0.192772\du}}
\pgfpathlineto{\pgfpoint{-4.288355\du}{0.190663\du}}
\pgfpathlineto{\pgfpoint{-4.243019\du}{0.187764\du}}
\pgfpathlineto{\pgfpoint{-4.197948\du}{0.184601\du}}
\pgfpathlineto{\pgfpoint{-4.153139\du}{0.180120\du}}
\pgfpathlineto{\pgfpoint{-4.109649\du}{0.175112\du}}
\pgfpathlineto{\pgfpoint{-4.066423\du}{0.169314\du}}
\pgfpathlineto{\pgfpoint{-4.024514\du}{0.162461\du}}
\pgfpathlineto{\pgfpoint{-3.983396\du}{0.155344\du}}
\pgfpathlineto{\pgfpoint{-3.943596\du}{0.147437\du}}
\pgfpathlineto{\pgfpoint{-3.905113\du}{0.138739\du}}
\pgfpathlineto{\pgfpoint{-3.867949\du}{0.129513\du}}
\pgfpathlineto{\pgfpoint{-3.831839\du}{0.119498\du}}
\pgfpathlineto{\pgfpoint{-3.797574\du}{0.108954\du}}
\pgfpathlineto{\pgfpoint{-3.765154\du}{0.097621\du}}
\pgfpathlineto{\pgfpoint{-3.734052\du}{0.086287\du}}
\pgfpathlineto{\pgfpoint{-3.704795\du}{0.074162\du}}
\pgfpathlineto{\pgfpoint{-3.677383\du}{0.061511\du}}
\pgfpathlineto{\pgfpoint{-3.651816\du}{0.047805\du}}
\pgfpathlineto{\pgfpoint{-3.628357\du}{0.034626\du}}
\pgfpathlineto{\pgfpoint{-3.607008\du}{0.020656\du}}
\pgfpathlineto{\pgfpoint{-3.587503\du}{0.006423\du}}
\pgfpathlineto{\pgfpoint{-3.570634\du}{-0.008074\du}}
\pgfpathlineto{\pgfpoint{-3.555610\du}{-0.022834\du}}
\pgfpathlineto{\pgfpoint{-3.542958\du}{-0.038122\du}}
\pgfpathlineto{\pgfpoint{-3.532415\du}{-0.053409\du}}
\pgfpathlineto{\pgfpoint{-3.524244\du}{-0.068960\du}}
\pgfpathlineto{\pgfpoint{-3.518446\du}{-0.084511\du}}
\pgfpathlineto{\pgfpoint{-3.515019\du}{-0.100062\du}}
\pgfpathlineto{\pgfpoint{-3.513701\du}{-0.116140\du}}
\pgfusepath{fill}
\pgfsetbuttcap
\pgfsetmiterjoin
\pgfsetdash{}{0pt}
\definecolor{dialinecolor}{rgb}{0.666667, 0.901961, 1.000000}
\pgfsetstrokecolor{dialinecolor}
\pgfpathmoveto{\pgfpoint{-3.524244\du}{-0.120885\du}}
\pgfpathlineto{\pgfpoint{-3.525562\du}{-0.136699\du}}
\pgfpathlineto{\pgfpoint{-3.528989\du}{-0.151723\du}}
\pgfpathlineto{\pgfpoint{-3.534788\du}{-0.167274\du}}
\pgfpathlineto{\pgfpoint{-3.542695\du}{-0.182298\du}}
\pgfpathlineto{\pgfpoint{-3.553238\du}{-0.197586\du}}
\pgfpathlineto{\pgfpoint{-3.566153\du}{-0.212346\du}}
\pgfpathlineto{\pgfpoint{-3.580386\du}{-0.226843\du}}
\pgfpathlineto{\pgfpoint{-3.597782\du}{-0.241339\du}}
\pgfpathlineto{\pgfpoint{-3.617024\du}{-0.255309\du}}
\pgfpathlineto{\pgfpoint{-3.638373\du}{-0.269279\du}}
\pgfpathlineto{\pgfpoint{-3.661568\du}{-0.282194\du}}
\pgfpathlineto{\pgfpoint{-3.687135\du}{-0.295109\du}}
\pgfpathlineto{\pgfpoint{-3.714020\du}{-0.308025\du}}
\pgfpathlineto{\pgfpoint{-3.743013\du}{-0.319885\du}}
\pgfpathlineto{\pgfpoint{-3.774116\du}{-0.331483\du}}
\pgfpathlineto{\pgfpoint{-3.806799\du}{-0.342026\du}}
\pgfpathlineto{\pgfpoint{-3.840801\du}{-0.352569\du}}
\pgfpathlineto{\pgfpoint{-3.876383\du}{-0.362321\du}}
\pgfpathlineto{\pgfpoint{-3.913021\du}{-0.371283\du}}
\pgfpathlineto{\pgfpoint{-3.951767\du}{-0.379981\du}}
\pgfpathlineto{\pgfpoint{-3.991567\du}{-0.387625\du}}
\pgfpathlineto{\pgfpoint{-4.031894\du}{-0.394741\du}}
\pgfpathlineto{\pgfpoint{-4.073803\du}{-0.401331\du}}
\pgfpathlineto{\pgfpoint{-4.116502\du}{-0.407129\du}}
\pgfpathlineto{\pgfpoint{-4.159992\du}{-0.412137\du}}
\pgfpathlineto{\pgfpoint{-4.204537\du}{-0.416355\du}}
\pgfpathlineto{\pgfpoint{-4.249082\du}{-0.420045\du}}
\pgfpathlineto{\pgfpoint{-4.294153\du}{-0.422153\du}}
\pgfpathlineto{\pgfpoint{-4.340016\du}{-0.424526\du}}
\pgfpathlineto{\pgfpoint{-4.385878\du}{-0.425843\du}}
\pgfpathlineto{\pgfpoint{-4.432004\du}{-0.425843\du}}
\pgfpathlineto{\pgfpoint{-4.432004\du}{-0.425843\du}}
\pgfpathlineto{\pgfpoint{-4.477866\du}{-0.425843\du}}
\pgfpathlineto{\pgfpoint{-4.523992\du}{-0.424526\du}}
\pgfpathlineto{\pgfpoint{-4.569591\du}{-0.422153\du}}
\pgfpathlineto{\pgfpoint{-4.614663\du}{-0.420045\du}}
\pgfpathlineto{\pgfpoint{-4.659471\du}{-0.416355\du}}
\pgfpathlineto{\pgfpoint{-4.703752\du}{-0.412137\du}}
\pgfpathlineto{\pgfpoint{-4.747242\du}{-0.407129\du}}
\pgfpathlineto{\pgfpoint{-4.790205\du}{-0.401331\du}}
\pgfpathlineto{\pgfpoint{-4.831850\du}{-0.394741\du}}
\pgfpathlineto{\pgfpoint{-4.872178\du}{-0.387625\du}}
\pgfpathlineto{\pgfpoint{-4.912242\du}{-0.379981\du}}
\pgfpathlineto{\pgfpoint{-4.950460\du}{-0.371283\du}}
\pgfpathlineto{\pgfpoint{-4.987097\du}{-0.362321\du}}
\pgfpathlineto{\pgfpoint{-5.022944\du}{-0.352569\du}}
\pgfpathlineto{\pgfpoint{-5.057472\du}{-0.342026\du}}
\pgfpathlineto{\pgfpoint{-5.090156\du}{-0.331483\du}}
\pgfpathlineto{\pgfpoint{-5.120731\du}{-0.319885\du}}
\pgfpathlineto{\pgfpoint{-5.149988\du}{-0.308025\du}}
\pgfpathlineto{\pgfpoint{-5.177137\du}{-0.295109\du}}
\pgfpathlineto{\pgfpoint{-5.202176\du}{-0.282194\du}}
\pgfpathlineto{\pgfpoint{-5.225108\du}{-0.269279\du}}
\pgfpathlineto{\pgfpoint{-5.246721\du}{-0.255309\du}}
\pgfpathlineto{\pgfpoint{-5.266226\du}{-0.241339\du}}
\pgfpathlineto{\pgfpoint{-5.283094\du}{-0.226843\du}}
\pgfpathlineto{\pgfpoint{-5.297855\du}{-0.212346\du}}
\pgfpathlineto{\pgfpoint{-5.311034\du}{-0.197586\du}}
\pgfpathlineto{\pgfpoint{-5.320786\du}{-0.182298\du}}
\pgfpathlineto{\pgfpoint{-5.329220\du}{-0.167274\du}}
\pgfpathlineto{\pgfpoint{-5.335019\du}{-0.151723\du}}
\pgfpathlineto{\pgfpoint{-5.338182\du}{-0.136699\du}}
\pgfpathlineto{\pgfpoint{-5.339500\du}{-0.120885\du}}
\pgfpathlineto{\pgfpoint{-5.339500\du}{-0.120885\du}}
\pgfpathlineto{\pgfpoint{-5.338182\du}{-0.105334\du}}
\pgfpathlineto{\pgfpoint{-5.335019\du}{-0.090046\du}}
\pgfpathlineto{\pgfpoint{-5.329220\du}{-0.074759\du}}
\pgfpathlineto{\pgfpoint{-5.320786\du}{-0.059735\du}}
\pgfpathlineto{\pgfpoint{-5.311034\du}{-0.044447\du}}
\pgfpathlineto{\pgfpoint{-5.297855\du}{-0.029687\du}}
\pgfpathlineto{\pgfpoint{-5.283094\du}{-0.015190\du}}
\pgfpathlineto{\pgfpoint{-5.266226\du}{-0.001221\du}}
\pgfpathlineto{\pgfpoint{-5.246721\du}{0.013276\du}}
\pgfpathlineto{\pgfpoint{-5.225108\du}{0.027246\du}}
\pgfpathlineto{\pgfpoint{-5.202176\du}{0.040424\du}}
\pgfpathlineto{\pgfpoint{-5.177137\du}{0.053076\du}}
\pgfpathlineto{\pgfpoint{-5.149988\du}{0.065728\du}}
\pgfpathlineto{\pgfpoint{-5.120731\du}{0.077852\du}}
\pgfpathlineto{\pgfpoint{-5.090156\du}{0.089186\du}}
\pgfpathlineto{\pgfpoint{-5.057472\du}{0.099993\du}}
\pgfpathlineto{\pgfpoint{-5.022944\du}{0.110009\du}}
\pgfpathlineto{\pgfpoint{-4.987097\du}{0.120288\du}}
\pgfpathlineto{\pgfpoint{-4.950460\du}{0.129250\du}}
\pgfpathlineto{\pgfpoint{-4.912242\du}{0.137684\du}}
\pgfpathlineto{\pgfpoint{-4.872178\du}{0.145855\du}}
\pgfpathlineto{\pgfpoint{-4.831850\du}{0.152708\du}}
\pgfpathlineto{\pgfpoint{-4.790205\du}{0.159034\du}}
\pgfpathlineto{\pgfpoint{-4.747242\du}{0.165096\du}}
\pgfpathlineto{\pgfpoint{-4.703752\du}{0.169841\du}}
\pgfpathlineto{\pgfpoint{-4.659471\du}{0.173794\du}}
\pgfpathlineto{\pgfpoint{-4.614663\du}{0.177748\du}}
\pgfpathlineto{\pgfpoint{-4.569591\du}{0.180120\du}}
\pgfpathlineto{\pgfpoint{-4.523992\du}{0.182229\du}}
\pgfpathlineto{\pgfpoint{-4.477866\du}{0.183283\du}}
\pgfpathlineto{\pgfpoint{-4.432004\du}{0.183810\du}}
\pgfpathlineto{\pgfpoint{-4.432004\du}{0.183810\du}}
\pgfpathlineto{\pgfpoint{-4.385878\du}{0.183283\du}}
\pgfpathlineto{\pgfpoint{-4.340016\du}{0.182229\du}}
\pgfpathlineto{\pgfpoint{-4.294153\du}{0.180120\du}}
\pgfpathlineto{\pgfpoint{-4.249082\du}{0.177748\du}}
\pgfpathlineto{\pgfpoint{-4.204537\du}{0.173794\du}}
\pgfpathlineto{\pgfpoint{-4.159992\du}{0.169841\du}}
\pgfpathlineto{\pgfpoint{-4.116502\du}{0.165096\du}}
\pgfpathlineto{\pgfpoint{-4.073803\du}{0.159034\du}}
\pgfpathlineto{\pgfpoint{-4.031894\du}{0.152708\du}}
\pgfpathlineto{\pgfpoint{-3.991567\du}{0.145855\du}}
\pgfpathlineto{\pgfpoint{-3.951767\du}{0.137684\du}}
\pgfpathlineto{\pgfpoint{-3.913021\du}{0.129250\du}}
\pgfpathlineto{\pgfpoint{-3.876383\du}{0.120288\du}}
\pgfpathlineto{\pgfpoint{-3.840801\du}{0.110009\du}}
\pgfpathlineto{\pgfpoint{-3.806799\du}{0.099993\du}}
\pgfpathlineto{\pgfpoint{-3.774116\du}{0.089186\du}}
\pgfpathlineto{\pgfpoint{-3.743013\du}{0.077852\du}}
\pgfpathlineto{\pgfpoint{-3.714020\du}{0.065728\du}}
\pgfpathlineto{\pgfpoint{-3.687135\du}{0.053076\du}}
\pgfpathlineto{\pgfpoint{-3.661568\du}{0.040424\du}}
\pgfpathlineto{\pgfpoint{-3.638373\du}{0.027246\du}}
\pgfpathlineto{\pgfpoint{-3.617024\du}{0.013276\du}}
\pgfpathlineto{\pgfpoint{-3.597782\du}{-0.001221\du}}
\pgfpathlineto{\pgfpoint{-3.580386\du}{-0.015190\du}}
\pgfpathlineto{\pgfpoint{-3.566153\du}{-0.029687\du}}
\pgfpathlineto{\pgfpoint{-3.553238\du}{-0.044447\du}}
\pgfpathlineto{\pgfpoint{-3.542695\du}{-0.059735\du}}
\pgfpathlineto{\pgfpoint{-3.534788\du}{-0.074759\du}}
\pgfpathlineto{\pgfpoint{-3.528989\du}{-0.090046\du}}
\pgfpathlineto{\pgfpoint{-3.525562\du}{-0.105334\du}}
\pgfpathlineto{\pgfpoint{-3.524244\du}{-0.120885\du}}
\pgfusepath{stroke}
\pgfsetbuttcap
\pgfsetmiterjoin
\pgfsetdash{}{0pt}
\definecolor{dialinecolor}{rgb}{0.000000, 0.470588, 0.666667}
\pgfsetfillcolor{dialinecolor}
\pgfpathmoveto{\pgfpoint{-5.339500\du}{-0.545771\du}}
\pgfpathlineto{\pgfpoint{-5.339500\du}{-0.110342\du}}
\pgfpathlineto{\pgfpoint{-3.524244\du}{-0.110342\du}}
\pgfpathlineto{\pgfpoint{-3.524244\du}{-0.545771\du}}
\pgfpathlineto{\pgfpoint{-5.339500\du}{-0.545771\du}}
\pgfusepath{fill}
\pgfsetbuttcap
\pgfsetmiterjoin
\pgfsetdash{}{0pt}
\definecolor{dialinecolor}{rgb}{0.000000, 0.705882, 1.000000}
\pgfsetfillcolor{dialinecolor}
\pgfpathmoveto{\pgfpoint{-3.513701\du}{-0.551833\du}}
\pgfpathlineto{\pgfpoint{-3.515019\du}{-0.567121\du}}
\pgfpathlineto{\pgfpoint{-3.518446\du}{-0.582935\du}}
\pgfpathlineto{\pgfpoint{-3.524244\du}{-0.598486\du}}
\pgfpathlineto{\pgfpoint{-3.532415\du}{-0.614038\du}}
\pgfpathlineto{\pgfpoint{-3.542958\du}{-0.629589\du}}
\pgfpathlineto{\pgfpoint{-3.555610\du}{-0.644612\du}}
\pgfpathlineto{\pgfpoint{-3.570634\du}{-0.659373\du}}
\pgfpathlineto{\pgfpoint{-3.587503\du}{-0.674133\du}}
\pgfpathlineto{\pgfpoint{-3.607008\du}{-0.688103\du}}
\pgfpathlineto{\pgfpoint{-3.628357\du}{-0.702072\du}}
\pgfpathlineto{\pgfpoint{-3.651816\du}{-0.715515\du}}
\pgfpathlineto{\pgfpoint{-3.677383\du}{-0.729221\du}}
\pgfpathlineto{\pgfpoint{-3.704795\du}{-0.741609\du}}
\pgfpathlineto{\pgfpoint{-3.734052\du}{-0.753997\du}}
\pgfpathlineto{\pgfpoint{-3.765154\du}{-0.765594\du}}
\pgfpathlineto{\pgfpoint{-3.797574\du}{-0.776401\du}}
\pgfpathlineto{\pgfpoint{-3.831839\du}{-0.786944\du}}
\pgfpathlineto{\pgfpoint{-3.867949\du}{-0.796960\du}}
\pgfpathlineto{\pgfpoint{-3.905113\du}{-0.806185\du}}
\pgfpathlineto{\pgfpoint{-3.943596\du}{-0.814883\du}}
\pgfpathlineto{\pgfpoint{-3.983396\du}{-0.823054\du}}
\pgfpathlineto{\pgfpoint{-4.024514\du}{-0.829907\du}}
\pgfpathlineto{\pgfpoint{-4.066423\du}{-0.837024\du}}
\pgfpathlineto{\pgfpoint{-4.109649\du}{-0.842822\du}}
\pgfpathlineto{\pgfpoint{-4.153139\du}{-0.847830\du}}
\pgfpathlineto{\pgfpoint{-4.197948\du}{-0.852048\du}}
\pgfpathlineto{\pgfpoint{-4.243019\du}{-0.855211\du}}
\pgfpathlineto{\pgfpoint{-4.288355\du}{-0.858110\du}}
\pgfpathlineto{\pgfpoint{-4.334217\du}{-0.860219\du}}
\pgfpathlineto{\pgfpoint{-4.380343\du}{-0.861800\du}}
\pgfpathlineto{\pgfpoint{-4.426996\du}{-0.861800\du}}
\pgfpathlineto{\pgfpoint{-4.426996\du}{-0.861800\du}}
\pgfpathlineto{\pgfpoint{-4.473122\du}{-0.861800\du}}
\pgfpathlineto{\pgfpoint{-4.518984\du}{-0.860219\du}}
\pgfpathlineto{\pgfpoint{-4.564847\du}{-0.858110\du}}
\pgfpathlineto{\pgfpoint{-4.610182\du}{-0.855211\du}}
\pgfpathlineto{\pgfpoint{-4.655254\du}{-0.852048\du}}
\pgfpathlineto{\pgfpoint{-4.700062\du}{-0.847830\du}}
\pgfpathlineto{\pgfpoint{-4.743552\du}{-0.842822\du}}
\pgfpathlineto{\pgfpoint{-4.786779\du}{-0.837024\du}}
\pgfpathlineto{\pgfpoint{-4.828688\du}{-0.829907\du}}
\pgfpathlineto{\pgfpoint{-4.869806\du}{-0.823054\du}}
\pgfpathlineto{\pgfpoint{-4.909869\du}{-0.814883\du}}
\pgfpathlineto{\pgfpoint{-4.948088\du}{-0.806185\du}}
\pgfpathlineto{\pgfpoint{-4.985252\du}{-0.796960\du}}
\pgfpathlineto{\pgfpoint{-5.021362\du}{-0.786944\du}}
\pgfpathlineto{\pgfpoint{-5.055627\du}{-0.776401\du}}
\pgfpathlineto{\pgfpoint{-5.088047\du}{-0.765594\du}}
\pgfpathlineto{\pgfpoint{-5.119150\du}{-0.753997\du}}
\pgfpathlineto{\pgfpoint{-5.148407\du}{-0.741609\du}}
\pgfpathlineto{\pgfpoint{-5.175819\du}{-0.729221\du}}
\pgfpathlineto{\pgfpoint{-5.201649\du}{-0.715515\du}}
\pgfpathlineto{\pgfpoint{-5.224844\du}{-0.702072\du}}
\pgfpathlineto{\pgfpoint{-5.246457\du}{-0.688103\du}}
\pgfpathlineto{\pgfpoint{-5.265698\du}{-0.674133\du}}
\pgfpathlineto{\pgfpoint{-5.282567\du}{-0.659373\du}}
\pgfpathlineto{\pgfpoint{-5.297591\du}{-0.644612\du}}
\pgfpathlineto{\pgfpoint{-5.310243\du}{-0.629589\du}}
\pgfpathlineto{\pgfpoint{-5.320786\du}{-0.614038\du}}
\pgfpathlineto{\pgfpoint{-5.328957\du}{-0.598486\du}}
\pgfpathlineto{\pgfpoint{-5.335019\du}{-0.582935\du}}
\pgfpathlineto{\pgfpoint{-5.338182\du}{-0.567121\du}}
\pgfpathlineto{\pgfpoint{-5.339500\du}{-0.551833\du}}
\pgfpathlineto{\pgfpoint{-5.339500\du}{-0.551833\du}}
\pgfpathlineto{\pgfpoint{-5.338182\du}{-0.535755\du}}
\pgfpathlineto{\pgfpoint{-5.335019\du}{-0.519940\du}}
\pgfpathlineto{\pgfpoint{-5.328957\du}{-0.504389\du}}
\pgfpathlineto{\pgfpoint{-5.320786\du}{-0.489102\du}}
\pgfpathlineto{\pgfpoint{-5.310243\du}{-0.473551\du}}
\pgfpathlineto{\pgfpoint{-5.297591\du}{-0.458791\du}}
\pgfpathlineto{\pgfpoint{-5.282567\du}{-0.443503\du}}
\pgfpathlineto{\pgfpoint{-5.265698\du}{-0.429006\du}}
\pgfpathlineto{\pgfpoint{-5.246457\du}{-0.414773\du}}
\pgfpathlineto{\pgfpoint{-5.224844\du}{-0.400804\du}}
\pgfpathlineto{\pgfpoint{-5.201649\du}{-0.387361\du}}
\pgfpathlineto{\pgfpoint{-5.175819\du}{-0.373919\du}}
\pgfpathlineto{\pgfpoint{-5.148407\du}{-0.361267\du}}
\pgfpathlineto{\pgfpoint{-5.119150\du}{-0.349670\du}}
\pgfpathlineto{\pgfpoint{-5.088047\du}{-0.337545\du}}
\pgfpathlineto{\pgfpoint{-5.055627\du}{-0.326475\du}}
\pgfpathlineto{\pgfpoint{-5.021362\du}{-0.315932\du}}
\pgfpathlineto{\pgfpoint{-4.985252\du}{-0.305916\du}}
\pgfpathlineto{\pgfpoint{-4.948088\du}{-0.296691\du}}
\pgfpathlineto{\pgfpoint{-4.909869\du}{-0.288256\du}}
\pgfpathlineto{\pgfpoint{-4.869806\du}{-0.280085\du}}
\pgfpathlineto{\pgfpoint{-4.828688\du}{-0.272969\du}}
\pgfpathlineto{\pgfpoint{-4.786779\du}{-0.266116\du}}
\pgfpathlineto{\pgfpoint{-4.743552\du}{-0.260317\du}}
\pgfpathlineto{\pgfpoint{-4.700062\du}{-0.255309\du}}
\pgfpathlineto{\pgfpoint{-4.655254\du}{-0.251092\du}}
\pgfpathlineto{\pgfpoint{-4.610182\du}{-0.247929\du}}
\pgfpathlineto{\pgfpoint{-4.564847\du}{-0.245030\du}}
\pgfpathlineto{\pgfpoint{-4.518984\du}{-0.242657\du}}
\pgfpathlineto{\pgfpoint{-4.473122\du}{-0.241339\du}}
\pgfpathlineto{\pgfpoint{-4.426996\du}{-0.241339\du}}
\pgfpathlineto{\pgfpoint{-4.426996\du}{-0.241339\du}}
\pgfpathlineto{\pgfpoint{-4.380343\du}{-0.241339\du}}
\pgfpathlineto{\pgfpoint{-4.334217\du}{-0.242657\du}}
\pgfpathlineto{\pgfpoint{-4.288355\du}{-0.245030\du}}
\pgfpathlineto{\pgfpoint{-4.243019\du}{-0.247929\du}}
\pgfpathlineto{\pgfpoint{-4.197948\du}{-0.251092\du}}
\pgfpathlineto{\pgfpoint{-4.153139\du}{-0.255309\du}}
\pgfpathlineto{\pgfpoint{-4.109649\du}{-0.260317\du}}
\pgfpathlineto{\pgfpoint{-4.066423\du}{-0.266116\du}}
\pgfpathlineto{\pgfpoint{-4.024514\du}{-0.272969\du}}
\pgfpathlineto{\pgfpoint{-3.983396\du}{-0.280085\du}}
\pgfpathlineto{\pgfpoint{-3.943596\du}{-0.288256\du}}
\pgfpathlineto{\pgfpoint{-3.905113\du}{-0.296691\du}}
\pgfpathlineto{\pgfpoint{-3.867949\du}{-0.305916\du}}
\pgfpathlineto{\pgfpoint{-3.831839\du}{-0.315932\du}}
\pgfpathlineto{\pgfpoint{-3.797574\du}{-0.326475\du}}
\pgfpathlineto{\pgfpoint{-3.765154\du}{-0.337545\du}}
\pgfpathlineto{\pgfpoint{-3.734052\du}{-0.349670\du}}
\pgfpathlineto{\pgfpoint{-3.704795\du}{-0.361267\du}}
\pgfpathlineto{\pgfpoint{-3.677383\du}{-0.373919\du}}
\pgfpathlineto{\pgfpoint{-3.651816\du}{-0.387361\du}}
\pgfpathlineto{\pgfpoint{-3.628357\du}{-0.400804\du}}
\pgfpathlineto{\pgfpoint{-3.607008\du}{-0.414773\du}}
\pgfpathlineto{\pgfpoint{-3.587503\du}{-0.429006\du}}
\pgfpathlineto{\pgfpoint{-3.570634\du}{-0.443503\du}}
\pgfpathlineto{\pgfpoint{-3.555610\du}{-0.458791\du}}
\pgfpathlineto{\pgfpoint{-3.542958\du}{-0.473551\du}}
\pgfpathlineto{\pgfpoint{-3.532415\du}{-0.489102\du}}
\pgfpathlineto{\pgfpoint{-3.524244\du}{-0.504389\du}}
\pgfpathlineto{\pgfpoint{-3.518446\du}{-0.519940\du}}
\pgfpathlineto{\pgfpoint{-3.515019\du}{-0.535755\du}}
\pgfpathlineto{\pgfpoint{-3.513701\du}{-0.551833\du}}
\pgfusepath{fill}
\pgfsetbuttcap
\pgfsetmiterjoin
\pgfsetdash{}{0pt}
\definecolor{dialinecolor}{rgb}{0.666667, 0.901961, 1.000000}
\pgfsetstrokecolor{dialinecolor}
\pgfpathmoveto{\pgfpoint{-3.524244\du}{-0.556314\du}}
\pgfpathlineto{\pgfpoint{-3.525562\du}{-0.572129\du}}
\pgfpathlineto{\pgfpoint{-3.528989\du}{-0.587680\du}}
\pgfpathlineto{\pgfpoint{-3.534788\du}{-0.602440\du}}
\pgfpathlineto{\pgfpoint{-3.542695\du}{-0.617991\du}}
\pgfpathlineto{\pgfpoint{-3.553238\du}{-0.633015\du}}
\pgfpathlineto{\pgfpoint{-3.566153\du}{-0.647775\du}}
\pgfpathlineto{\pgfpoint{-3.580386\du}{-0.662009\du}}
\pgfpathlineto{\pgfpoint{-3.597782\du}{-0.676769\du}}
\pgfpathlineto{\pgfpoint{-3.617024\du}{-0.690475\du}}
\pgfpathlineto{\pgfpoint{-3.638373\du}{-0.704708\du}}
\pgfpathlineto{\pgfpoint{-3.661568\du}{-0.717887\du}}
\pgfpathlineto{\pgfpoint{-3.687135\du}{-0.731066\du}}
\pgfpathlineto{\pgfpoint{-3.714020\du}{-0.743454\du}}
\pgfpathlineto{\pgfpoint{-3.743013\du}{-0.755051\du}}
\pgfpathlineto{\pgfpoint{-3.774116\du}{-0.767176\du}}
\pgfpathlineto{\pgfpoint{-3.806799\du}{-0.777719\du}}
\pgfpathlineto{\pgfpoint{-3.840801\du}{-0.787998\du}}
\pgfpathlineto{\pgfpoint{-3.876383\du}{-0.797751\du}}
\pgfpathlineto{\pgfpoint{-3.913021\du}{-0.806712\du}}
\pgfpathlineto{\pgfpoint{-3.951767\du}{-0.815674\du}}
\pgfpathlineto{\pgfpoint{-3.991567\du}{-0.823318\du}}
\pgfpathlineto{\pgfpoint{-4.031894\du}{-0.830434\du}}
\pgfpathlineto{\pgfpoint{-4.073803\du}{-0.837024\du}}
\pgfpathlineto{\pgfpoint{-4.116502\du}{-0.842822\du}}
\pgfpathlineto{\pgfpoint{-4.159992\du}{-0.847830\du}}
\pgfpathlineto{\pgfpoint{-4.204537\du}{-0.851520\du}}
\pgfpathlineto{\pgfpoint{-4.249082\du}{-0.855211\du}}
\pgfpathlineto{\pgfpoint{-4.294153\du}{-0.857846\du}}
\pgfpathlineto{\pgfpoint{-4.340016\du}{-0.860219\du}}
\pgfpathlineto{\pgfpoint{-4.385878\du}{-0.861273\du}}
\pgfpathlineto{\pgfpoint{-4.432004\du}{-0.861800\du}}
\pgfpathlineto{\pgfpoint{-4.432004\du}{-0.861800\du}}
\pgfpathlineto{\pgfpoint{-4.477866\du}{-0.861273\du}}
\pgfpathlineto{\pgfpoint{-4.523992\du}{-0.860219\du}}
\pgfpathlineto{\pgfpoint{-4.569591\du}{-0.857846\du}}
\pgfpathlineto{\pgfpoint{-4.614663\du}{-0.855211\du}}
\pgfpathlineto{\pgfpoint{-4.659471\du}{-0.851520\du}}
\pgfpathlineto{\pgfpoint{-4.703752\du}{-0.847830\du}}
\pgfpathlineto{\pgfpoint{-4.747242\du}{-0.842822\du}}
\pgfpathlineto{\pgfpoint{-4.790205\du}{-0.837024\du}}
\pgfpathlineto{\pgfpoint{-4.831850\du}{-0.830434\du}}
\pgfpathlineto{\pgfpoint{-4.872178\du}{-0.823318\du}}
\pgfpathlineto{\pgfpoint{-4.912242\du}{-0.815674\du}}
\pgfpathlineto{\pgfpoint{-4.950460\du}{-0.806712\du}}
\pgfpathlineto{\pgfpoint{-4.987097\du}{-0.797751\du}}
\pgfpathlineto{\pgfpoint{-5.022944\du}{-0.787998\du}}
\pgfpathlineto{\pgfpoint{-5.057472\du}{-0.777719\du}}
\pgfpathlineto{\pgfpoint{-5.090156\du}{-0.767176\du}}
\pgfpathlineto{\pgfpoint{-5.120731\du}{-0.755051\du}}
\pgfpathlineto{\pgfpoint{-5.149988\du}{-0.743454\du}}
\pgfpathlineto{\pgfpoint{-5.177137\du}{-0.731066\du}}
\pgfpathlineto{\pgfpoint{-5.202176\du}{-0.717887\du}}
\pgfpathlineto{\pgfpoint{-5.225108\du}{-0.704708\du}}
\pgfpathlineto{\pgfpoint{-5.246721\du}{-0.690475\du}}
\pgfpathlineto{\pgfpoint{-5.266226\du}{-0.676769\du}}
\pgfpathlineto{\pgfpoint{-5.283094\du}{-0.662009\du}}
\pgfpathlineto{\pgfpoint{-5.297855\du}{-0.647775\du}}
\pgfpathlineto{\pgfpoint{-5.311034\du}{-0.633015\du}}
\pgfpathlineto{\pgfpoint{-5.320786\du}{-0.617991\du}}
\pgfpathlineto{\pgfpoint{-5.329220\du}{-0.602440\du}}
\pgfpathlineto{\pgfpoint{-5.335019\du}{-0.587680\du}}
\pgfpathlineto{\pgfpoint{-5.338182\du}{-0.572129\du}}
\pgfpathlineto{\pgfpoint{-5.339500\du}{-0.556314\du}}
\pgfpathlineto{\pgfpoint{-5.339500\du}{-0.556314\du}}
\pgfpathlineto{\pgfpoint{-5.338182\du}{-0.541027\du}}
\pgfpathlineto{\pgfpoint{-5.335019\du}{-0.525476\du}}
\pgfpathlineto{\pgfpoint{-5.329220\du}{-0.510715\du}}
\pgfpathlineto{\pgfpoint{-5.320786\du}{-0.495428\du}}
\pgfpathlineto{\pgfpoint{-5.311034\du}{-0.480140\du}}
\pgfpathlineto{\pgfpoint{-5.297855\du}{-0.465380\du}}
\pgfpathlineto{\pgfpoint{-5.283094\du}{-0.450620\du}}
\pgfpathlineto{\pgfpoint{-5.266226\du}{-0.436123\du}}
\pgfpathlineto{\pgfpoint{-5.246721\du}{-0.422153\du}}
\pgfpathlineto{\pgfpoint{-5.225108\du}{-0.408447\du}}
\pgfpathlineto{\pgfpoint{-5.202176\du}{-0.395005\du}}
\pgfpathlineto{\pgfpoint{-5.177137\du}{-0.382617\du}}
\pgfpathlineto{\pgfpoint{-5.149988\du}{-0.369702\du}}
\pgfpathlineto{\pgfpoint{-5.120731\du}{-0.357841\du}}
\pgfpathlineto{\pgfpoint{-5.090156\du}{-0.346243\du}}
\pgfpathlineto{\pgfpoint{-5.057472\du}{-0.335437\du}}
\pgfpathlineto{\pgfpoint{-5.022944\du}{-0.325157\du}}
\pgfpathlineto{\pgfpoint{-4.987097\du}{-0.315405\du}}
\pgfpathlineto{\pgfpoint{-4.950460\du}{-0.306443\du}}
\pgfpathlineto{\pgfpoint{-4.912242\du}{-0.297745\du}}
\pgfpathlineto{\pgfpoint{-4.872178\du}{-0.289574\du}}
\pgfpathlineto{\pgfpoint{-4.831850\du}{-0.282985\du}}
\pgfpathlineto{\pgfpoint{-4.790205\du}{-0.276132\du}}
\pgfpathlineto{\pgfpoint{-4.747242\du}{-0.270597\du}}
\pgfpathlineto{\pgfpoint{-4.703752\du}{-0.265589\du}}
\pgfpathlineto{\pgfpoint{-4.659471\du}{-0.261108\du}}
\pgfpathlineto{\pgfpoint{-4.614663\du}{-0.257418\du}}
\pgfpathlineto{\pgfpoint{-4.569591\du}{-0.255309\du}}
\pgfpathlineto{\pgfpoint{-4.523992\du}{-0.253200\du}}
\pgfpathlineto{\pgfpoint{-4.477866\du}{-0.251883\du}}
\pgfpathlineto{\pgfpoint{-4.432004\du}{-0.251619\du}}
\pgfpathlineto{\pgfpoint{-4.432004\du}{-0.251619\du}}
\pgfpathlineto{\pgfpoint{-4.385878\du}{-0.251883\du}}
\pgfpathlineto{\pgfpoint{-4.340016\du}{-0.253200\du}}
\pgfpathlineto{\pgfpoint{-4.294153\du}{-0.255309\du}}
\pgfpathlineto{\pgfpoint{-4.249082\du}{-0.257418\du}}
\pgfpathlineto{\pgfpoint{-4.204537\du}{-0.261108\du}}
\pgfpathlineto{\pgfpoint{-4.159992\du}{-0.265589\du}}
\pgfpathlineto{\pgfpoint{-4.116502\du}{-0.270597\du}}
\pgfpathlineto{\pgfpoint{-4.073803\du}{-0.276132\du}}
\pgfpathlineto{\pgfpoint{-4.031894\du}{-0.282985\du}}
\pgfpathlineto{\pgfpoint{-3.991567\du}{-0.289574\du}}
\pgfpathlineto{\pgfpoint{-3.951767\du}{-0.297745\du}}
\pgfpathlineto{\pgfpoint{-3.913021\du}{-0.306443\du}}
\pgfpathlineto{\pgfpoint{-3.876383\du}{-0.315405\du}}
\pgfpathlineto{\pgfpoint{-3.840801\du}{-0.325157\du}}
\pgfpathlineto{\pgfpoint{-3.806799\du}{-0.335437\du}}
\pgfpathlineto{\pgfpoint{-3.774116\du}{-0.346243\du}}
\pgfpathlineto{\pgfpoint{-3.743013\du}{-0.357841\du}}
\pgfpathlineto{\pgfpoint{-3.714020\du}{-0.369702\du}}
\pgfpathlineto{\pgfpoint{-3.687135\du}{-0.382617\du}}
\pgfpathlineto{\pgfpoint{-3.661568\du}{-0.395005\du}}
\pgfpathlineto{\pgfpoint{-3.638373\du}{-0.408447\du}}
\pgfpathlineto{\pgfpoint{-3.617024\du}{-0.422153\du}}
\pgfpathlineto{\pgfpoint{-3.597782\du}{-0.436123\du}}
\pgfpathlineto{\pgfpoint{-3.580386\du}{-0.450620\du}}
\pgfpathlineto{\pgfpoint{-3.566153\du}{-0.465380\du}}
\pgfpathlineto{\pgfpoint{-3.553238\du}{-0.480140\du}}
\pgfpathlineto{\pgfpoint{-3.542695\du}{-0.495428\du}}
\pgfpathlineto{\pgfpoint{-3.534788\du}{-0.510715\du}}
\pgfpathlineto{\pgfpoint{-3.528989\du}{-0.525476\du}}
\pgfpathlineto{\pgfpoint{-3.525562\du}{-0.541027\du}}
\pgfpathlineto{\pgfpoint{-3.524244\du}{-0.556314\du}}
\pgfusepath{stroke}
\pgfsetbuttcap
\pgfsetmiterjoin
\pgfsetdash{}{0pt}
\definecolor{dialinecolor}{rgb}{0.000000, 0.000000, 0.000000}
\pgfsetfillcolor{dialinecolor}
\pgfpathmoveto{\pgfpoint{-4.410127\du}{-0.622208\du}}
\pgfpathlineto{\pgfpoint{-4.277021\du}{-0.578455\du}}
\pgfpathlineto{\pgfpoint{-3.955720\du}{-0.709452\du}}
\pgfpathlineto{\pgfpoint{-3.812071\du}{-0.665699\du}}
\pgfpathlineto{\pgfpoint{-3.890090\du}{-0.774556\du}}
\pgfpathlineto{\pgfpoint{-4.265687\du}{-0.774556\du}}
\pgfpathlineto{\pgfpoint{-4.111231\du}{-0.741609\du}}
\pgfpathlineto{\pgfpoint{-4.410127\du}{-0.622208\du}}
\pgfusepath{fill}
\pgfsetbuttcap
\pgfsetmiterjoin
\pgfsetdash{}{0pt}
\definecolor{dialinecolor}{rgb}{0.000000, 0.000000, 0.000000}
\pgfsetfillcolor{dialinecolor}
\pgfpathmoveto{\pgfpoint{-4.465742\du}{-0.502017\du}}
\pgfpathlineto{\pgfpoint{-4.597794\du}{-0.545771\du}}
\pgfpathlineto{\pgfpoint{-4.908024\du}{-0.414773\du}}
\pgfpathlineto{\pgfpoint{-5.062744\du}{-0.458791\du}}
\pgfpathlineto{\pgfpoint{-4.985252\du}{-0.338863\du}}
\pgfpathlineto{\pgfpoint{-4.597794\du}{-0.338863\du}}
\pgfpathlineto{\pgfpoint{-4.764111\du}{-0.382617\du}}
\pgfpathlineto{\pgfpoint{-4.465742\du}{-0.502017\du}}
\pgfusepath{fill}
\pgfsetbuttcap
\pgfsetmiterjoin
\pgfsetdash{}{0pt}
\definecolor{dialinecolor}{rgb}{0.000000, 0.000000, 0.000000}
\pgfsetfillcolor{dialinecolor}
\pgfpathmoveto{\pgfpoint{-5.029797\du}{-0.741609\du}}
\pgfpathlineto{\pgfpoint{-4.896954\du}{-0.785363\du}}
\pgfpathlineto{\pgfpoint{-4.576181\du}{-0.665699\du}}
\pgfpathlineto{\pgfpoint{-4.432004\du}{-0.698119\du}}
\pgfpathlineto{\pgfpoint{-4.509232\du}{-0.589525\du}}
\pgfpathlineto{\pgfpoint{-4.885884\du}{-0.589525\du}}
\pgfpathlineto{\pgfpoint{-4.731164\du}{-0.622208\du}}
\pgfpathlineto{\pgfpoint{-5.029797\du}{-0.741609\du}}
\pgfusepath{fill}
\pgfsetbuttcap
\pgfsetmiterjoin
\pgfsetdash{}{0pt}
\definecolor{dialinecolor}{rgb}{0.000000, 0.000000, 0.000000}
\pgfsetfillcolor{dialinecolor}
\pgfpathmoveto{\pgfpoint{-3.834475\du}{-0.371547\du}}
\pgfpathlineto{\pgfpoint{-3.967054\du}{-0.328056\du}}
\pgfpathlineto{\pgfpoint{-4.277021\du}{-0.458791\du}}
\pgfpathlineto{\pgfpoint{-4.432004\du}{-0.414773\du}}
\pgfpathlineto{\pgfpoint{-4.354512\du}{-0.523894\du}}
\pgfpathlineto{\pgfpoint{-3.967054\du}{-0.523894\du}}
\pgfpathlineto{\pgfpoint{-4.133371\du}{-0.491738\du}}
\pgfpathlineto{\pgfpoint{-3.834475\du}{-0.371547\du}}
\pgfusepath{fill}
\pgfsetbuttcap
\pgfsetmiterjoin
\pgfsetdash{}{0pt}
\definecolor{dialinecolor}{rgb}{1.000000, 1.000000, 1.000000}
\pgfsetfillcolor{dialinecolor}
\pgfpathmoveto{\pgfpoint{-4.399057\du}{-0.611402\du}}
\pgfpathlineto{\pgfpoint{-4.265687\du}{-0.567912\du}}
\pgfpathlineto{\pgfpoint{-3.944914\du}{-0.698119\du}}
\pgfpathlineto{\pgfpoint{-3.801000\du}{-0.654365\du}}
\pgfpathlineto{\pgfpoint{-3.878492\du}{-0.763749\du}}
\pgfpathlineto{\pgfpoint{-4.255144\du}{-0.763749\du}}
\pgfpathlineto{\pgfpoint{-4.099897\du}{-0.731066\du}}
\pgfpathlineto{\pgfpoint{-4.399057\du}{-0.611402\du}}
\pgfusepath{fill}
\pgfsetbuttcap
\pgfsetmiterjoin
\pgfsetdash{}{0pt}
\definecolor{dialinecolor}{rgb}{1.000000, 1.000000, 1.000000}
\pgfsetfillcolor{dialinecolor}
\pgfpathmoveto{\pgfpoint{-4.453881\du}{-0.491738\du}}
\pgfpathlineto{\pgfpoint{-4.587251\du}{-0.534964\du}}
\pgfpathlineto{\pgfpoint{-4.896954\du}{-0.404230\du}}
\pgfpathlineto{\pgfpoint{-5.051937\du}{-0.447984\du}}
\pgfpathlineto{\pgfpoint{-4.974446\du}{-0.328056\du}}
\pgfpathlineto{\pgfpoint{-4.587251\du}{-0.328056\du}}
\pgfpathlineto{\pgfpoint{-4.752777\du}{-0.371547\du}}
\pgfpathlineto{\pgfpoint{-4.453881\du}{-0.491738\du}}
\pgfusepath{fill}
\pgfsetbuttcap
\pgfsetmiterjoin
\pgfsetdash{}{0pt}
\definecolor{dialinecolor}{rgb}{1.000000, 1.000000, 1.000000}
\pgfsetfillcolor{dialinecolor}
\pgfpathmoveto{\pgfpoint{-5.018463\du}{-0.731066\du}}
\pgfpathlineto{\pgfpoint{-4.885884\du}{-0.774556\du}}
\pgfpathlineto{\pgfpoint{-4.564847\du}{-0.654365\du}}
\pgfpathlineto{\pgfpoint{-4.420934\du}{-0.687576\du}}
\pgfpathlineto{\pgfpoint{-4.498425\du}{-0.578455\du}}
\pgfpathlineto{\pgfpoint{-4.874814\du}{-0.578455\du}}
\pgfpathlineto{\pgfpoint{-4.719567\du}{-0.611402\du}}
\pgfpathlineto{\pgfpoint{-5.018463\du}{-0.731066\du}}
\pgfusepath{fill}
\pgfsetbuttcap
\pgfsetmiterjoin
\pgfsetdash{}{0pt}
\definecolor{dialinecolor}{rgb}{1.000000, 1.000000, 1.000000}
\pgfsetfillcolor{dialinecolor}
\pgfpathmoveto{\pgfpoint{-3.823141\du}{-0.360740\du}}
\pgfpathlineto{\pgfpoint{-3.955720\du}{-0.316986\du}}
\pgfpathlineto{\pgfpoint{-4.265687\du}{-0.447984\du}}
\pgfpathlineto{\pgfpoint{-4.420934\du}{-0.404230\du}}
\pgfpathlineto{\pgfpoint{-4.343706\du}{-0.512824\du}}
\pgfpathlineto{\pgfpoint{-3.955720\du}{-0.512824\du}}
\pgfpathlineto{\pgfpoint{-4.121774\du}{-0.480404\du}}
\pgfpathlineto{\pgfpoint{-3.823141\du}{-0.360740\du}}
\pgfusepath{fill}
\pgfsetbuttcap
\pgfsetmiterjoin
\pgfsetdash{}{0pt}
\definecolor{dialinecolor}{rgb}{0.666667, 0.901961, 1.000000}
\pgfsetstrokecolor{dialinecolor}
\pgfpathmoveto{\pgfpoint{-5.339500\du}{-0.556314\du}}
\pgfpathlineto{\pgfpoint{-5.339500\du}{-0.121412\du}}
\pgfusepath{stroke}
\pgfsetbuttcap
\pgfsetmiterjoin
\pgfsetdash{}{0pt}
\definecolor{dialinecolor}{rgb}{0.666667, 0.901961, 1.000000}
\pgfsetstrokecolor{dialinecolor}
\pgfpathmoveto{\pgfpoint{-3.524244\du}{-0.556314\du}}
\pgfpathlineto{\pgfpoint{-3.524244\du}{-0.121412\du}}
\pgfusepath{stroke}
\pgfsetbuttcap
\pgfsetmiterjoin
\pgfsetdash{}{0pt}
\definecolor{dialinecolor}{rgb}{0.000000, 0.000000, 0.000000}
\pgfsetfillcolor{dialinecolor}
\pgfpathmoveto{\pgfpoint{-4.764111\du}{-0.175445\du}}
\pgfpathlineto{\pgfpoint{-4.564847\du}{-0.175445\du}}
\pgfpathlineto{\pgfpoint{-4.443074\du}{-0.077395\du}}
\pgfpathlineto{\pgfpoint{-4.321565\du}{-0.175445\du}}
\pgfpathlineto{\pgfpoint{-4.133371\du}{-0.175445\du}}
\pgfpathlineto{\pgfpoint{-4.133371\du}{-0.218935\du}}
\pgfpathlineto{\pgfpoint{-4.000265\du}{-0.154095\du}}
\pgfpathlineto{\pgfpoint{-4.133371\du}{-0.088465\du}}
\pgfpathlineto{\pgfpoint{-4.133371\du}{-0.131955\du}}
\pgfpathlineto{\pgfpoint{-4.299161\du}{-0.131955\du}}
\pgfpathlineto{\pgfpoint{-4.410127\du}{-0.055781\du}}
\pgfpathlineto{\pgfpoint{-4.299161\du}{0.030936\du}}
\pgfpathlineto{\pgfpoint{-4.133371\du}{0.030936\du}}
\pgfpathlineto{\pgfpoint{-4.133371\du}{-0.022834\du}}
\pgfpathlineto{\pgfpoint{-4.000265\du}{0.053076\du}}
\pgfpathlineto{\pgfpoint{-4.133371\du}{0.118707\du}}
\pgfpathlineto{\pgfpoint{-4.133371\du}{0.074689\du}}
\pgfpathlineto{\pgfpoint{-4.321565\du}{0.074689\du}}
\pgfpathlineto{\pgfpoint{-4.443074\du}{-0.022834\du}}
\pgfpathlineto{\pgfpoint{-4.564847\du}{0.074689\du}}
\pgfpathlineto{\pgfpoint{-4.764111\du}{0.074689\du}}
\pgfpathlineto{\pgfpoint{-4.764111\du}{0.118707\du}}
\pgfpathlineto{\pgfpoint{-4.896954\du}{0.053076\du}}
\pgfpathlineto{\pgfpoint{-4.764111\du}{-0.022834\du}}
\pgfpathlineto{\pgfpoint{-4.764111\du}{0.030936\du}}
\pgfpathlineto{\pgfpoint{-4.587251\du}{0.030936\du}}
\pgfpathlineto{\pgfpoint{-4.476285\du}{-0.044975\du}}
\pgfpathlineto{\pgfpoint{-4.587251\du}{-0.131955\du}}
\pgfpathlineto{\pgfpoint{-4.764111\du}{-0.131955\du}}
\pgfpathlineto{\pgfpoint{-4.764111\du}{-0.088465\du}}
\pgfpathlineto{\pgfpoint{-4.896954\du}{-0.154095\du}}
\pgfpathlineto{\pgfpoint{-4.764111\du}{-0.218935\du}}
\pgfpathlineto{\pgfpoint{-4.764111\du}{-0.175445\du}}
\pgfusepath{fill}
\pgfsetbuttcap
\pgfsetmiterjoin
\pgfsetdash{}{0pt}
\definecolor{dialinecolor}{rgb}{1.000000, 1.000000, 1.000000}
\pgfsetfillcolor{dialinecolor}
\pgfpathmoveto{\pgfpoint{-4.752777\du}{-0.164902\du}}
\pgfpathlineto{\pgfpoint{-4.553777\du}{-0.164902\du}}
\pgfpathlineto{\pgfpoint{-4.432004\du}{-0.067115\du}}
\pgfpathlineto{\pgfpoint{-4.310495\du}{-0.164902\du}}
\pgfpathlineto{\pgfpoint{-4.121774\du}{-0.164902\du}}
\pgfpathlineto{\pgfpoint{-4.121774\du}{-0.208392\du}}
\pgfpathlineto{\pgfpoint{-3.989458\du}{-0.142762\du}}
\pgfpathlineto{\pgfpoint{-4.121774\du}{-0.077395\du}}
\pgfpathlineto{\pgfpoint{-4.121774\du}{-0.120885\du}}
\pgfpathlineto{\pgfpoint{-4.288091\du}{-0.120885\du}}
\pgfpathlineto{\pgfpoint{-4.399057\du}{-0.044975\du}}
\pgfpathlineto{\pgfpoint{-4.288091\du}{0.042006\du}}
\pgfpathlineto{\pgfpoint{-4.121774\du}{0.042006\du}}
\pgfpathlineto{\pgfpoint{-4.121774\du}{-0.012291\du}}
\pgfpathlineto{\pgfpoint{-3.989458\du}{0.063883\du}}
\pgfpathlineto{\pgfpoint{-4.121774\du}{0.129250\du}}
\pgfpathlineto{\pgfpoint{-4.121774\du}{0.085496\du}}
\pgfpathlineto{\pgfpoint{-4.310495\du}{0.085496\du}}
\pgfpathlineto{\pgfpoint{-4.432004\du}{-0.012291\du}}
\pgfpathlineto{\pgfpoint{-4.553777\du}{0.085496\du}}
\pgfpathlineto{\pgfpoint{-4.752777\du}{0.085496\du}}
\pgfpathlineto{\pgfpoint{-4.752777\du}{0.129250\du}}
\pgfpathlineto{\pgfpoint{-4.885884\du}{0.063883\du}}
\pgfpathlineto{\pgfpoint{-4.752777\du}{-0.012291\du}}
\pgfpathlineto{\pgfpoint{-4.752777\du}{0.042006\du}}
\pgfpathlineto{\pgfpoint{-4.576181\du}{0.042006\du}}
\pgfpathlineto{\pgfpoint{-4.465742\du}{-0.034168\du}}
\pgfpathlineto{\pgfpoint{-4.576181\du}{-0.120885\du}}
\pgfpathlineto{\pgfpoint{-4.752777\du}{-0.120885\du}}
\pgfpathlineto{\pgfpoint{-4.752777\du}{-0.077395\du}}
\pgfpathlineto{\pgfpoint{-4.885884\du}{-0.142762\du}}
\pgfpathlineto{\pgfpoint{-4.752777\du}{-0.208392\du}}
\pgfpathlineto{\pgfpoint{-4.752777\du}{-0.164902\du}}
\pgfusepath{fill}
\pgfsetlinewidth{0.100000\du}
\pgfsetdash{}{0pt}
\pgfsetdash{}{0pt}
\pgfsetbuttcap
\pgfsetmiterjoin
\pgfsetlinewidth{0.001000\du}
\pgfsetbuttcap
\pgfsetmiterjoin
\pgfsetdash{}{0pt}
\definecolor{dialinecolor}{rgb}{0.000000, 0.470588, 0.666667}
\pgfsetfillcolor{dialinecolor}
\pgfpathmoveto{\pgfpoint{-6.779166\du}{3.323470\du}}
\pgfpathlineto{\pgfpoint{-6.780664\du}{3.305794\du}}
\pgfpathlineto{\pgfpoint{-6.784559\du}{3.287818\du}}
\pgfpathlineto{\pgfpoint{-6.791150\du}{3.270141\du}}
\pgfpathlineto{\pgfpoint{-6.800438\du}{3.252765\du}}
\pgfpathlineto{\pgfpoint{-6.812422\du}{3.235088\du}}
\pgfpathlineto{\pgfpoint{-6.826803\du}{3.217711\du}}
\pgfpathlineto{\pgfpoint{-6.843880\du}{3.200934\du}}
\pgfpathlineto{\pgfpoint{-6.863054\du}{3.184456\du}}
\pgfpathlineto{\pgfpoint{-6.885225\du}{3.168277\du}}
\pgfpathlineto{\pgfpoint{-6.909492\du}{3.152698\du}}
\pgfpathlineto{\pgfpoint{-6.936157\du}{3.137119\du}}
\pgfpathlineto{\pgfpoint{-6.965218\du}{3.122139\du}}
\pgfpathlineto{\pgfpoint{-6.996377\du}{3.107758\du}}
\pgfpathlineto{\pgfpoint{-7.029632\du}{3.093677\du}}
\pgfpathlineto{\pgfpoint{-7.064985\du}{3.080494\du}}
\pgfpathlineto{\pgfpoint{-7.101836\du}{3.067911\du}}
\pgfpathlineto{\pgfpoint{-7.140784\du}{3.055927\du}}
\pgfpathlineto{\pgfpoint{-7.181829\du}{3.044842\du}}
\pgfpathlineto{\pgfpoint{-7.224073\du}{3.033756\du}}
\pgfpathlineto{\pgfpoint{-7.267815\du}{3.024469\du}}
\pgfpathlineto{\pgfpoint{-7.313055\du}{3.015181\du}}
\pgfpathlineto{\pgfpoint{-7.359792\du}{3.007092\du}}
\pgfpathlineto{\pgfpoint{-7.407429\du}{2.999302\du}}
\pgfpathlineto{\pgfpoint{-7.456563\du}{2.992711\du}}
\pgfpathlineto{\pgfpoint{-7.505997\du}{2.987019\du}}
\pgfpathlineto{\pgfpoint{-7.556930\du}{2.982225\du}}
\pgfpathlineto{\pgfpoint{-7.608161\du}{2.978031\du}}
\pgfpathlineto{\pgfpoint{-7.659693\du}{2.974735\du}}
\pgfpathlineto{\pgfpoint{-7.711823\du}{2.972338\du}}
\pgfpathlineto{\pgfpoint{-7.764253\du}{2.971439\du}}
\pgfpathlineto{\pgfpoint{-7.817283\du}{2.970541\du}}
\pgfpathlineto{\pgfpoint{-7.817283\du}{2.970541\du}}
\pgfpathlineto{\pgfpoint{-7.869713\du}{2.971439\du}}
\pgfpathlineto{\pgfpoint{-7.921843\du}{2.972338\du}}
\pgfpathlineto{\pgfpoint{-7.973974\du}{2.974735\du}}
\pgfpathlineto{\pgfpoint{-8.025505\du}{2.978031\du}}
\pgfpathlineto{\pgfpoint{-8.076737\du}{2.982225\du}}
\pgfpathlineto{\pgfpoint{-8.127669\du}{2.987019\du}}
\pgfpathlineto{\pgfpoint{-8.177103\du}{2.992711\du}}
\pgfpathlineto{\pgfpoint{-8.226237\du}{2.999302\du}}
\pgfpathlineto{\pgfpoint{-8.273874\du}{3.007092\du}}
\pgfpathlineto{\pgfpoint{-8.320612\du}{3.015181\du}}
\pgfpathlineto{\pgfpoint{-8.366151\du}{3.024469\du}}
\pgfpathlineto{\pgfpoint{-8.409593\du}{3.033756\du}}
\pgfpathlineto{\pgfpoint{-8.451837\du}{3.044842\du}}
\pgfpathlineto{\pgfpoint{-8.492882\du}{3.055927\du}}
\pgfpathlineto{\pgfpoint{-8.531830\du}{3.067911\du}}
\pgfpathlineto{\pgfpoint{-8.568681\du}{3.080494\du}}
\pgfpathlineto{\pgfpoint{-8.604034\du}{3.093677\du}}
\pgfpathlineto{\pgfpoint{-8.637290\du}{3.107758\du}}
\pgfpathlineto{\pgfpoint{-8.668448\du}{3.122139\du}}
\pgfpathlineto{\pgfpoint{-8.697809\du}{3.137119\du}}
\pgfpathlineto{\pgfpoint{-8.724174\du}{3.152698\du}}
\pgfpathlineto{\pgfpoint{-8.748741\du}{3.168277\du}}
\pgfpathlineto{\pgfpoint{-8.770612\du}{3.184456\du}}
\pgfpathlineto{\pgfpoint{-8.789786\du}{3.200934\du}}
\pgfpathlineto{\pgfpoint{-8.806863\du}{3.217711\du}}
\pgfpathlineto{\pgfpoint{-8.821244\du}{3.235088\du}}
\pgfpathlineto{\pgfpoint{-8.833228\du}{3.252765\du}}
\pgfpathlineto{\pgfpoint{-8.842516\du}{3.270141\du}}
\pgfpathlineto{\pgfpoint{-8.849407\du}{3.287818\du}}
\pgfpathlineto{\pgfpoint{-8.853002\du}{3.305794\du}}
\pgfpathlineto{\pgfpoint{-8.854500\du}{3.323470\du}}
\pgfpathlineto{\pgfpoint{-8.854500\du}{3.323470\du}}
\pgfpathlineto{\pgfpoint{-8.853002\du}{3.341746\du}}
\pgfpathlineto{\pgfpoint{-8.849407\du}{3.359422\du}}
\pgfpathlineto{\pgfpoint{-8.842516\du}{3.377099\du}}
\pgfpathlineto{\pgfpoint{-8.833228\du}{3.394775\du}}
\pgfpathlineto{\pgfpoint{-8.821244\du}{3.412152\du}}
\pgfpathlineto{\pgfpoint{-8.806863\du}{3.429529\du}}
\pgfpathlineto{\pgfpoint{-8.789786\du}{3.446307\du}}
\pgfpathlineto{\pgfpoint{-8.770612\du}{3.462785\du}}
\pgfpathlineto{\pgfpoint{-8.748741\du}{3.478963\du}}
\pgfpathlineto{\pgfpoint{-8.724174\du}{3.494842\du}}
\pgfpathlineto{\pgfpoint{-8.697809\du}{3.509822\du}}
\pgfpathlineto{\pgfpoint{-8.668448\du}{3.525401\du}}
\pgfpathlineto{\pgfpoint{-8.637290\du}{3.539782\du}}
\pgfpathlineto{\pgfpoint{-8.604034\du}{3.553564\du}}
\pgfpathlineto{\pgfpoint{-8.568681\du}{3.566446\du}}
\pgfpathlineto{\pgfpoint{-8.531830\du}{3.579329\du}}
\pgfpathlineto{\pgfpoint{-8.492882\du}{3.591313\du}}
\pgfpathlineto{\pgfpoint{-8.451837\du}{3.602698\du}}
\pgfpathlineto{\pgfpoint{-8.409593\du}{3.613184\du}}
\pgfpathlineto{\pgfpoint{-8.366151\du}{3.623071\du}}
\pgfpathlineto{\pgfpoint{-8.320612\du}{3.632059\du}}
\pgfpathlineto{\pgfpoint{-8.273874\du}{3.640148\du}}
\pgfpathlineto{\pgfpoint{-8.226237\du}{3.647938\du}}
\pgfpathlineto{\pgfpoint{-8.177103\du}{3.654529\du}}
\pgfpathlineto{\pgfpoint{-8.127669\du}{3.660221\du}}
\pgfpathlineto{\pgfpoint{-8.076737\du}{3.665315\du}}
\pgfpathlineto{\pgfpoint{-8.025505\du}{3.668910\du}}
\pgfpathlineto{\pgfpoint{-7.973974\du}{3.672205\du}}
\pgfpathlineto{\pgfpoint{-7.921843\du}{3.674602\du}}
\pgfpathlineto{\pgfpoint{-7.869713\du}{3.676100\du}}
\pgfpathlineto{\pgfpoint{-7.817283\du}{3.676400\du}}
\pgfpathlineto{\pgfpoint{-7.817283\du}{3.676400\du}}
\pgfpathlineto{\pgfpoint{-7.764253\du}{3.676100\du}}
\pgfpathlineto{\pgfpoint{-7.711823\du}{3.674602\du}}
\pgfpathlineto{\pgfpoint{-7.659693\du}{3.672205\du}}
\pgfpathlineto{\pgfpoint{-7.608161\du}{3.668910\du}}
\pgfpathlineto{\pgfpoint{-7.556930\du}{3.665315\du}}
\pgfpathlineto{\pgfpoint{-7.505997\du}{3.660221\du}}
\pgfpathlineto{\pgfpoint{-7.456563\du}{3.654529\du}}
\pgfpathlineto{\pgfpoint{-7.407429\du}{3.647938\du}}
\pgfpathlineto{\pgfpoint{-7.359792\du}{3.640148\du}}
\pgfpathlineto{\pgfpoint{-7.313055\du}{3.632059\du}}
\pgfpathlineto{\pgfpoint{-7.267815\du}{3.623071\du}}
\pgfpathlineto{\pgfpoint{-7.224073\du}{3.613184\du}}
\pgfpathlineto{\pgfpoint{-7.181829\du}{3.602698\du}}
\pgfpathlineto{\pgfpoint{-7.140784\du}{3.591313\du}}
\pgfpathlineto{\pgfpoint{-7.101836\du}{3.579329\du}}
\pgfpathlineto{\pgfpoint{-7.064985\du}{3.566446\du}}
\pgfpathlineto{\pgfpoint{-7.029632\du}{3.553564\du}}
\pgfpathlineto{\pgfpoint{-6.996377\du}{3.539782\du}}
\pgfpathlineto{\pgfpoint{-6.965218\du}{3.525401\du}}
\pgfpathlineto{\pgfpoint{-6.936157\du}{3.509822\du}}
\pgfpathlineto{\pgfpoint{-6.909492\du}{3.494842\du}}
\pgfpathlineto{\pgfpoint{-6.885225\du}{3.478963\du}}
\pgfpathlineto{\pgfpoint{-6.863054\du}{3.462785\du}}
\pgfpathlineto{\pgfpoint{-6.843880\du}{3.446307\du}}
\pgfpathlineto{\pgfpoint{-6.826803\du}{3.429529\du}}
\pgfpathlineto{\pgfpoint{-6.812422\du}{3.412152\du}}
\pgfpathlineto{\pgfpoint{-6.800438\du}{3.394775\du}}
\pgfpathlineto{\pgfpoint{-6.791150\du}{3.377099\du}}
\pgfpathlineto{\pgfpoint{-6.784559\du}{3.359422\du}}
\pgfpathlineto{\pgfpoint{-6.780664\du}{3.341746\du}}
\pgfpathlineto{\pgfpoint{-6.779166\du}{3.323470\du}}
\pgfusepath{fill}
\pgfsetbuttcap
\pgfsetmiterjoin
\pgfsetdash{}{0pt}
\definecolor{dialinecolor}{rgb}{0.666667, 0.901961, 1.000000}
\pgfsetstrokecolor{dialinecolor}
\pgfpathmoveto{\pgfpoint{-6.791150\du}{3.318077\du}}
\pgfpathlineto{\pgfpoint{-6.792648\du}{3.300101\du}}
\pgfpathlineto{\pgfpoint{-6.796543\du}{3.283024\du}}
\pgfpathlineto{\pgfpoint{-6.803134\du}{3.265348\du}}
\pgfpathlineto{\pgfpoint{-6.812122\du}{3.248271\du}}
\pgfpathlineto{\pgfpoint{-6.824106\du}{3.230894\du}}
\pgfpathlineto{\pgfpoint{-6.838787\du}{3.214116\du}}
\pgfpathlineto{\pgfpoint{-6.854965\du}{3.197638\du}}
\pgfpathlineto{\pgfpoint{-6.874739\du}{3.181160\du}}
\pgfpathlineto{\pgfpoint{-6.896610\du}{3.165281\du}}
\pgfpathlineto{\pgfpoint{-6.920877\du}{3.149402\du}}
\pgfpathlineto{\pgfpoint{-6.947242\du}{3.134722\du}}
\pgfpathlineto{\pgfpoint{-6.976303\du}{3.120041\du}}
\pgfpathlineto{\pgfpoint{-7.006863\du}{3.105361\du}}
\pgfpathlineto{\pgfpoint{-7.039819\du}{3.091879\du}}
\pgfpathlineto{\pgfpoint{-7.075172\du}{3.078697\du}}
\pgfpathlineto{\pgfpoint{-7.112322\du}{3.066713\du}}
\pgfpathlineto{\pgfpoint{-7.150971\du}{3.054728\du}}
\pgfpathlineto{\pgfpoint{-7.191417\du}{3.043643\du}}
\pgfpathlineto{\pgfpoint{-7.233061\du}{3.033457\du}}
\pgfpathlineto{\pgfpoint{-7.277103\du}{3.023570\du}}
\pgfpathlineto{\pgfpoint{-7.322342\du}{3.014882\du}}
\pgfpathlineto{\pgfpoint{-7.368181\du}{3.006792\du}}
\pgfpathlineto{\pgfpoint{-7.415818\du}{2.999302\du}}
\pgfpathlineto{\pgfpoint{-7.464353\du}{2.992711\du}}
\pgfpathlineto{\pgfpoint{-7.513787\du}{2.987019\du}}
\pgfpathlineto{\pgfpoint{-7.564420\du}{2.982225\du}}
\pgfpathlineto{\pgfpoint{-7.615052\du}{2.978031\du}}
\pgfpathlineto{\pgfpoint{-7.666284\du}{2.975634\du}}
\pgfpathlineto{\pgfpoint{-7.718414\du}{2.972937\du}}
\pgfpathlineto{\pgfpoint{-7.770545\du}{2.971439\du}}
\pgfpathlineto{\pgfpoint{-7.822975\du}{2.971439\du}}
\pgfpathlineto{\pgfpoint{-7.822975\du}{2.971439\du}}
\pgfpathlineto{\pgfpoint{-7.875105\du}{2.971439\du}}
\pgfpathlineto{\pgfpoint{-7.927536\du}{2.972937\du}}
\pgfpathlineto{\pgfpoint{-7.979366\du}{2.975634\du}}
\pgfpathlineto{\pgfpoint{-8.030598\du}{2.978031\du}}
\pgfpathlineto{\pgfpoint{-8.081530\du}{2.982225\du}}
\pgfpathlineto{\pgfpoint{-8.131863\du}{2.987019\du}}
\pgfpathlineto{\pgfpoint{-8.181297\du}{2.992711\du}}
\pgfpathlineto{\pgfpoint{-8.230132\du}{2.999302\du}}
\pgfpathlineto{\pgfpoint{-8.277469\du}{3.006792\du}}
\pgfpathlineto{\pgfpoint{-8.323308\du}{3.014882\du}}
\pgfpathlineto{\pgfpoint{-8.368847\du}{3.023570\du}}
\pgfpathlineto{\pgfpoint{-8.412289\du}{3.033457\du}}
\pgfpathlineto{\pgfpoint{-8.453934\du}{3.043643\du}}
\pgfpathlineto{\pgfpoint{-8.494680\du}{3.054728\du}}
\pgfpathlineto{\pgfpoint{-8.533927\du}{3.066713\du}}
\pgfpathlineto{\pgfpoint{-8.571078\du}{3.078697\du}}
\pgfpathlineto{\pgfpoint{-8.605831\du}{3.091879\du}}
\pgfpathlineto{\pgfpoint{-8.639087\du}{3.105361\du}}
\pgfpathlineto{\pgfpoint{-8.669946\du}{3.120041\du}}
\pgfpathlineto{\pgfpoint{-8.698408\du}{3.134722\du}}
\pgfpathlineto{\pgfpoint{-8.724473\du}{3.149402\du}}
\pgfpathlineto{\pgfpoint{-8.749041\du}{3.165281\du}}
\pgfpathlineto{\pgfpoint{-8.771211\du}{3.181160\du}}
\pgfpathlineto{\pgfpoint{-8.790385\du}{3.197638\du}}
\pgfpathlineto{\pgfpoint{-8.807163\du}{3.214116\du}}
\pgfpathlineto{\pgfpoint{-8.822143\du}{3.230894\du}}
\pgfpathlineto{\pgfpoint{-8.833228\du}{3.248271\du}}
\pgfpathlineto{\pgfpoint{-8.842816\du}{3.265348\du}}
\pgfpathlineto{\pgfpoint{-8.849407\du}{3.283024\du}}
\pgfpathlineto{\pgfpoint{-8.853002\du}{3.300101\du}}
\pgfpathlineto{\pgfpoint{-8.854500\du}{3.318077\du}}
\pgfpathlineto{\pgfpoint{-8.854500\du}{3.318077\du}}
\pgfpathlineto{\pgfpoint{-8.853002\du}{3.335754\du}}
\pgfpathlineto{\pgfpoint{-8.849407\du}{3.353131\du}}
\pgfpathlineto{\pgfpoint{-8.842816\du}{3.370508\du}}
\pgfpathlineto{\pgfpoint{-8.833228\du}{3.387585\du}}
\pgfpathlineto{\pgfpoint{-8.822143\du}{3.404962\du}}
\pgfpathlineto{\pgfpoint{-8.807163\du}{3.421739\du}}
\pgfpathlineto{\pgfpoint{-8.790385\du}{3.438217\du}}
\pgfpathlineto{\pgfpoint{-8.771211\du}{3.454096\du}}
\pgfpathlineto{\pgfpoint{-8.749041\du}{3.470574\du}}
\pgfpathlineto{\pgfpoint{-8.724473\du}{3.486453\du}}
\pgfpathlineto{\pgfpoint{-8.698408\du}{3.501433\du}}
\pgfpathlineto{\pgfpoint{-8.669946\du}{3.515814\du}}
\pgfpathlineto{\pgfpoint{-8.639087\du}{3.530195\du}}
\pgfpathlineto{\pgfpoint{-8.605831\du}{3.543976\du}}
\pgfpathlineto{\pgfpoint{-8.571078\du}{3.556859\du}}
\pgfpathlineto{\pgfpoint{-8.533927\du}{3.569143\du}}
\pgfpathlineto{\pgfpoint{-8.494680\du}{3.580528\du}}
\pgfpathlineto{\pgfpoint{-8.453934\du}{3.592212\du}}
\pgfpathlineto{\pgfpoint{-8.412289\du}{3.602398\du}}
\pgfpathlineto{\pgfpoint{-8.368847\du}{3.611986\du}}
\pgfpathlineto{\pgfpoint{-8.323308\du}{3.621273\du}}
\pgfpathlineto{\pgfpoint{-8.277469\du}{3.629063\du}}
\pgfpathlineto{\pgfpoint{-8.230132\du}{3.636253\du}}
\pgfpathlineto{\pgfpoint{-8.181297\du}{3.643144\du}}
\pgfpathlineto{\pgfpoint{-8.131863\du}{3.648537\du}}
\pgfpathlineto{\pgfpoint{-8.081530\du}{3.653031\du}}
\pgfpathlineto{\pgfpoint{-8.030598\du}{3.657525\du}}
\pgfpathlineto{\pgfpoint{-7.979366\du}{3.660221\du}}
\pgfpathlineto{\pgfpoint{-7.927536\du}{3.662618\du}}
\pgfpathlineto{\pgfpoint{-7.875105\du}{3.663817\du}}
\pgfpathlineto{\pgfpoint{-7.822975\du}{3.664416\du}}
\pgfpathlineto{\pgfpoint{-7.822975\du}{3.664416\du}}
\pgfpathlineto{\pgfpoint{-7.770545\du}{3.663817\du}}
\pgfpathlineto{\pgfpoint{-7.718414\du}{3.662618\du}}
\pgfpathlineto{\pgfpoint{-7.666284\du}{3.660221\du}}
\pgfpathlineto{\pgfpoint{-7.615052\du}{3.657525\du}}
\pgfpathlineto{\pgfpoint{-7.564420\du}{3.653031\du}}
\pgfpathlineto{\pgfpoint{-7.513787\du}{3.648537\du}}
\pgfpathlineto{\pgfpoint{-7.464353\du}{3.643144\du}}
\pgfpathlineto{\pgfpoint{-7.415818\du}{3.636253\du}}
\pgfpathlineto{\pgfpoint{-7.368181\du}{3.629063\du}}
\pgfpathlineto{\pgfpoint{-7.322342\du}{3.621273\du}}
\pgfpathlineto{\pgfpoint{-7.277103\du}{3.611986\du}}
\pgfpathlineto{\pgfpoint{-7.233061\du}{3.602398\du}}
\pgfpathlineto{\pgfpoint{-7.191417\du}{3.592212\du}}
\pgfpathlineto{\pgfpoint{-7.150971\du}{3.580528\du}}
\pgfpathlineto{\pgfpoint{-7.112322\du}{3.569143\du}}
\pgfpathlineto{\pgfpoint{-7.075172\du}{3.556859\du}}
\pgfpathlineto{\pgfpoint{-7.039819\du}{3.543976\du}}
\pgfpathlineto{\pgfpoint{-7.006863\du}{3.530195\du}}
\pgfpathlineto{\pgfpoint{-6.976303\du}{3.515814\du}}
\pgfpathlineto{\pgfpoint{-6.947242\du}{3.501433\du}}
\pgfpathlineto{\pgfpoint{-6.920877\du}{3.486453\du}}
\pgfpathlineto{\pgfpoint{-6.896610\du}{3.470574\du}}
\pgfpathlineto{\pgfpoint{-6.874739\du}{3.454096\du}}
\pgfpathlineto{\pgfpoint{-6.854965\du}{3.438217\du}}
\pgfpathlineto{\pgfpoint{-6.838787\du}{3.421739\du}}
\pgfpathlineto{\pgfpoint{-6.824106\du}{3.404962\du}}
\pgfpathlineto{\pgfpoint{-6.812122\du}{3.387585\du}}
\pgfpathlineto{\pgfpoint{-6.803134\du}{3.370508\du}}
\pgfpathlineto{\pgfpoint{-6.796543\du}{3.353131\du}}
\pgfpathlineto{\pgfpoint{-6.792648\du}{3.335754\du}}
\pgfpathlineto{\pgfpoint{-6.791150\du}{3.318077\du}}
\pgfusepath{stroke}
\pgfsetbuttcap
\pgfsetmiterjoin
\pgfsetdash{}{0pt}
\definecolor{dialinecolor}{rgb}{0.000000, 0.470588, 0.666667}
\pgfsetfillcolor{dialinecolor}
\pgfpathmoveto{\pgfpoint{-8.854500\du}{2.835121\du}}
\pgfpathlineto{\pgfpoint{-8.854500\du}{3.330061\du}}
\pgfpathlineto{\pgfpoint{-6.791150\du}{3.330061\du}}
\pgfpathlineto{\pgfpoint{-6.791150\du}{2.835121\du}}
\pgfpathlineto{\pgfpoint{-8.854500\du}{2.835121\du}}
\pgfusepath{fill}
\pgfsetbuttcap
\pgfsetmiterjoin
\pgfsetdash{}{0pt}
\definecolor{dialinecolor}{rgb}{0.000000, 0.705882, 1.000000}
\pgfsetfillcolor{dialinecolor}
\pgfpathmoveto{\pgfpoint{-6.779166\du}{2.828230\du}}
\pgfpathlineto{\pgfpoint{-6.780664\du}{2.810854\du}}
\pgfpathlineto{\pgfpoint{-6.784559\du}{2.792878\du}}
\pgfpathlineto{\pgfpoint{-6.791150\du}{2.775201\du}}
\pgfpathlineto{\pgfpoint{-6.800438\du}{2.757525\du}}
\pgfpathlineto{\pgfpoint{-6.812422\du}{2.739848\du}}
\pgfpathlineto{\pgfpoint{-6.826803\du}{2.722771\du}}
\pgfpathlineto{\pgfpoint{-6.843880\du}{2.705993\du}}
\pgfpathlineto{\pgfpoint{-6.863054\du}{2.689216\du}}
\pgfpathlineto{\pgfpoint{-6.885225\du}{2.673337\du}}
\pgfpathlineto{\pgfpoint{-6.909492\du}{2.657458\du}}
\pgfpathlineto{\pgfpoint{-6.936157\du}{2.642178\du}}
\pgfpathlineto{\pgfpoint{-6.965218\du}{2.626599\du}}
\pgfpathlineto{\pgfpoint{-6.996377\du}{2.612518\du}}
\pgfpathlineto{\pgfpoint{-7.029632\du}{2.598437\du}}
\pgfpathlineto{\pgfpoint{-7.064985\du}{2.585254\du}}
\pgfpathlineto{\pgfpoint{-7.101836\du}{2.572971\du}}
\pgfpathlineto{\pgfpoint{-7.140784\du}{2.560987\du}}
\pgfpathlineto{\pgfpoint{-7.181829\du}{2.549602\du}}
\pgfpathlineto{\pgfpoint{-7.224073\du}{2.539116\du}}
\pgfpathlineto{\pgfpoint{-7.267815\du}{2.529229\du}}
\pgfpathlineto{\pgfpoint{-7.313055\du}{2.519941\du}}
\pgfpathlineto{\pgfpoint{-7.359792\du}{2.512152\du}}
\pgfpathlineto{\pgfpoint{-7.407429\du}{2.504062\du}}
\pgfpathlineto{\pgfpoint{-7.456563\du}{2.497471\du}}
\pgfpathlineto{\pgfpoint{-7.505997\du}{2.491779\du}}
\pgfpathlineto{\pgfpoint{-7.556930\du}{2.486985\du}}
\pgfpathlineto{\pgfpoint{-7.608161\du}{2.483390\du}}
\pgfpathlineto{\pgfpoint{-7.659693\du}{2.480094\du}}
\pgfpathlineto{\pgfpoint{-7.711823\du}{2.477698\du}}
\pgfpathlineto{\pgfpoint{-7.764253\du}{2.475900\du}}
\pgfpathlineto{\pgfpoint{-7.817283\du}{2.475900\du}}
\pgfpathlineto{\pgfpoint{-7.817283\du}{2.475900\du}}
\pgfpathlineto{\pgfpoint{-7.869713\du}{2.475900\du}}
\pgfpathlineto{\pgfpoint{-7.921843\du}{2.477698\du}}
\pgfpathlineto{\pgfpoint{-7.973974\du}{2.480094\du}}
\pgfpathlineto{\pgfpoint{-8.025505\du}{2.483390\du}}
\pgfpathlineto{\pgfpoint{-8.076737\du}{2.486985\du}}
\pgfpathlineto{\pgfpoint{-8.127669\du}{2.491779\du}}
\pgfpathlineto{\pgfpoint{-8.177103\du}{2.497471\du}}
\pgfpathlineto{\pgfpoint{-8.226237\du}{2.504062\du}}
\pgfpathlineto{\pgfpoint{-8.273874\du}{2.512152\du}}
\pgfpathlineto{\pgfpoint{-8.320612\du}{2.519941\du}}
\pgfpathlineto{\pgfpoint{-8.366151\du}{2.529229\du}}
\pgfpathlineto{\pgfpoint{-8.409593\du}{2.539116\du}}
\pgfpathlineto{\pgfpoint{-8.451837\du}{2.549602\du}}
\pgfpathlineto{\pgfpoint{-8.492882\du}{2.560987\du}}
\pgfpathlineto{\pgfpoint{-8.531830\du}{2.572971\du}}
\pgfpathlineto{\pgfpoint{-8.568681\du}{2.585254\du}}
\pgfpathlineto{\pgfpoint{-8.604034\du}{2.598437\du}}
\pgfpathlineto{\pgfpoint{-8.637290\du}{2.612518\du}}
\pgfpathlineto{\pgfpoint{-8.668448\du}{2.626599\du}}
\pgfpathlineto{\pgfpoint{-8.697809\du}{2.642178\du}}
\pgfpathlineto{\pgfpoint{-8.724174\du}{2.657458\du}}
\pgfpathlineto{\pgfpoint{-8.748741\du}{2.673337\du}}
\pgfpathlineto{\pgfpoint{-8.770612\du}{2.689216\du}}
\pgfpathlineto{\pgfpoint{-8.789786\du}{2.705993\du}}
\pgfpathlineto{\pgfpoint{-8.806863\du}{2.722771\du}}
\pgfpathlineto{\pgfpoint{-8.821244\du}{2.739848\du}}
\pgfpathlineto{\pgfpoint{-8.833228\du}{2.757525\du}}
\pgfpathlineto{\pgfpoint{-8.842516\du}{2.775201\du}}
\pgfpathlineto{\pgfpoint{-8.849407\du}{2.792878\du}}
\pgfpathlineto{\pgfpoint{-8.853002\du}{2.810854\du}}
\pgfpathlineto{\pgfpoint{-8.854500\du}{2.828230\du}}
\pgfpathlineto{\pgfpoint{-8.854500\du}{2.828230\du}}
\pgfpathlineto{\pgfpoint{-8.853002\du}{2.846506\du}}
\pgfpathlineto{\pgfpoint{-8.849407\du}{2.864482\du}}
\pgfpathlineto{\pgfpoint{-8.842516\du}{2.882158\du}}
\pgfpathlineto{\pgfpoint{-8.833228\du}{2.899535\du}}
\pgfpathlineto{\pgfpoint{-8.821244\du}{2.917212\du}}
\pgfpathlineto{\pgfpoint{-8.806863\du}{2.933989\du}}
\pgfpathlineto{\pgfpoint{-8.789786\du}{2.951366\du}}
\pgfpathlineto{\pgfpoint{-8.770612\du}{2.967844\du}}
\pgfpathlineto{\pgfpoint{-8.748741\du}{2.984023\du}}
\pgfpathlineto{\pgfpoint{-8.724174\du}{2.999902\du}}
\pgfpathlineto{\pgfpoint{-8.697809\du}{3.015181\du}}
\pgfpathlineto{\pgfpoint{-8.668448\du}{3.030461\du}}
\pgfpathlineto{\pgfpoint{-8.637290\du}{3.044842\du}}
\pgfpathlineto{\pgfpoint{-8.604034\du}{3.058024\du}}
\pgfpathlineto{\pgfpoint{-8.568681\du}{3.071806\du}}
\pgfpathlineto{\pgfpoint{-8.531830\du}{3.084389\du}}
\pgfpathlineto{\pgfpoint{-8.492882\du}{3.096373\du}}
\pgfpathlineto{\pgfpoint{-8.451837\du}{3.107758\du}}
\pgfpathlineto{\pgfpoint{-8.409593\du}{3.118244\du}}
\pgfpathlineto{\pgfpoint{-8.366151\du}{3.127831\du}}
\pgfpathlineto{\pgfpoint{-8.320612\du}{3.137119\du}}
\pgfpathlineto{\pgfpoint{-8.273874\du}{3.145208\du}}
\pgfpathlineto{\pgfpoint{-8.226237\du}{3.152997\du}}
\pgfpathlineto{\pgfpoint{-8.177103\du}{3.159589\du}}
\pgfpathlineto{\pgfpoint{-8.127669\du}{3.165281\du}}
\pgfpathlineto{\pgfpoint{-8.076737\du}{3.170075\du}}
\pgfpathlineto{\pgfpoint{-8.025505\du}{3.173670\du}}
\pgfpathlineto{\pgfpoint{-7.973974\du}{3.176966\du}}
\pgfpathlineto{\pgfpoint{-7.921843\du}{3.179662\du}}
\pgfpathlineto{\pgfpoint{-7.869713\du}{3.181160\du}}
\pgfpathlineto{\pgfpoint{-7.817283\du}{3.181160\du}}
\pgfpathlineto{\pgfpoint{-7.817283\du}{3.181160\du}}
\pgfpathlineto{\pgfpoint{-7.764253\du}{3.181160\du}}
\pgfpathlineto{\pgfpoint{-7.711823\du}{3.179662\du}}
\pgfpathlineto{\pgfpoint{-7.659693\du}{3.176966\du}}
\pgfpathlineto{\pgfpoint{-7.608161\du}{3.173670\du}}
\pgfpathlineto{\pgfpoint{-7.556930\du}{3.170075\du}}
\pgfpathlineto{\pgfpoint{-7.505997\du}{3.165281\du}}
\pgfpathlineto{\pgfpoint{-7.456563\du}{3.159589\du}}
\pgfpathlineto{\pgfpoint{-7.407429\du}{3.152997\du}}
\pgfpathlineto{\pgfpoint{-7.359792\du}{3.145208\du}}
\pgfpathlineto{\pgfpoint{-7.313055\du}{3.137119\du}}
\pgfpathlineto{\pgfpoint{-7.267815\du}{3.127831\du}}
\pgfpathlineto{\pgfpoint{-7.224073\du}{3.118244\du}}
\pgfpathlineto{\pgfpoint{-7.181829\du}{3.107758\du}}
\pgfpathlineto{\pgfpoint{-7.140784\du}{3.096373\du}}
\pgfpathlineto{\pgfpoint{-7.101836\du}{3.084389\du}}
\pgfpathlineto{\pgfpoint{-7.064985\du}{3.071806\du}}
\pgfpathlineto{\pgfpoint{-7.029632\du}{3.058024\du}}
\pgfpathlineto{\pgfpoint{-6.996377\du}{3.044842\du}}
\pgfpathlineto{\pgfpoint{-6.965218\du}{3.030461\du}}
\pgfpathlineto{\pgfpoint{-6.936157\du}{3.015181\du}}
\pgfpathlineto{\pgfpoint{-6.909492\du}{2.999902\du}}
\pgfpathlineto{\pgfpoint{-6.885225\du}{2.984023\du}}
\pgfpathlineto{\pgfpoint{-6.863054\du}{2.967844\du}}
\pgfpathlineto{\pgfpoint{-6.843880\du}{2.951366\du}}
\pgfpathlineto{\pgfpoint{-6.826803\du}{2.933989\du}}
\pgfpathlineto{\pgfpoint{-6.812422\du}{2.917212\du}}
\pgfpathlineto{\pgfpoint{-6.800438\du}{2.899535\du}}
\pgfpathlineto{\pgfpoint{-6.791150\du}{2.882158\du}}
\pgfpathlineto{\pgfpoint{-6.784559\du}{2.864482\du}}
\pgfpathlineto{\pgfpoint{-6.780664\du}{2.846506\du}}
\pgfpathlineto{\pgfpoint{-6.779166\du}{2.828230\du}}
\pgfusepath{fill}
\pgfsetbuttcap
\pgfsetmiterjoin
\pgfsetdash{}{0pt}
\definecolor{dialinecolor}{rgb}{0.666667, 0.901961, 1.000000}
\pgfsetstrokecolor{dialinecolor}
\pgfpathmoveto{\pgfpoint{-6.791150\du}{2.823137\du}}
\pgfpathlineto{\pgfpoint{-6.792648\du}{2.805161\du}}
\pgfpathlineto{\pgfpoint{-6.796543\du}{2.787485\du}}
\pgfpathlineto{\pgfpoint{-6.803134\du}{2.770707\du}}
\pgfpathlineto{\pgfpoint{-6.812122\du}{2.753031\du}}
\pgfpathlineto{\pgfpoint{-6.824106\du}{2.735953\du}}
\pgfpathlineto{\pgfpoint{-6.838787\du}{2.719176\du}}
\pgfpathlineto{\pgfpoint{-6.854965\du}{2.702997\du}}
\pgfpathlineto{\pgfpoint{-6.874739\du}{2.686220\du}}
\pgfpathlineto{\pgfpoint{-6.896610\du}{2.670640\du}}
\pgfpathlineto{\pgfpoint{-6.920877\du}{2.654462\du}}
\pgfpathlineto{\pgfpoint{-6.947242\du}{2.639482\du}}
\pgfpathlineto{\pgfpoint{-6.976303\du}{2.624502\du}}
\pgfpathlineto{\pgfpoint{-7.006863\du}{2.610421\du}}
\pgfpathlineto{\pgfpoint{-7.039819\du}{2.597238\du}}
\pgfpathlineto{\pgfpoint{-7.075172\du}{2.583457\du}}
\pgfpathlineto{\pgfpoint{-7.112322\du}{2.571473\du}}
\pgfpathlineto{\pgfpoint{-7.150971\du}{2.559788\du}}
\pgfpathlineto{\pgfpoint{-7.191417\du}{2.548703\du}}
\pgfpathlineto{\pgfpoint{-7.233061\du}{2.538517\du}}
\pgfpathlineto{\pgfpoint{-7.277103\du}{2.528330\du}}
\pgfpathlineto{\pgfpoint{-7.322342\du}{2.519642\du}}
\pgfpathlineto{\pgfpoint{-7.368181\du}{2.511552\du}}
\pgfpathlineto{\pgfpoint{-7.415818\du}{2.504062\du}}
\pgfpathlineto{\pgfpoint{-7.464353\du}{2.497471\du}}
\pgfpathlineto{\pgfpoint{-7.513787\du}{2.491779\du}}
\pgfpathlineto{\pgfpoint{-7.564420\du}{2.487584\du}}
\pgfpathlineto{\pgfpoint{-7.615052\du}{2.483390\du}}
\pgfpathlineto{\pgfpoint{-7.666284\du}{2.480394\du}}
\pgfpathlineto{\pgfpoint{-7.718414\du}{2.477698\du}}
\pgfpathlineto{\pgfpoint{-7.770545\du}{2.476499\du}}
\pgfpathlineto{\pgfpoint{-7.822975\du}{2.475900\du}}
\pgfpathlineto{\pgfpoint{-7.822975\du}{2.475900\du}}
\pgfpathlineto{\pgfpoint{-7.875105\du}{2.476499\du}}
\pgfpathlineto{\pgfpoint{-7.927536\du}{2.477698\du}}
\pgfpathlineto{\pgfpoint{-7.979366\du}{2.480394\du}}
\pgfpathlineto{\pgfpoint{-8.030598\du}{2.483390\du}}
\pgfpathlineto{\pgfpoint{-8.081530\du}{2.487584\du}}
\pgfpathlineto{\pgfpoint{-8.131863\du}{2.491779\du}}
\pgfpathlineto{\pgfpoint{-8.181297\du}{2.497471\du}}
\pgfpathlineto{\pgfpoint{-8.230132\du}{2.504062\du}}
\pgfpathlineto{\pgfpoint{-8.277469\du}{2.511552\du}}
\pgfpathlineto{\pgfpoint{-8.323308\du}{2.519642\du}}
\pgfpathlineto{\pgfpoint{-8.368847\du}{2.528330\du}}
\pgfpathlineto{\pgfpoint{-8.412289\du}{2.538517\du}}
\pgfpathlineto{\pgfpoint{-8.453934\du}{2.548703\du}}
\pgfpathlineto{\pgfpoint{-8.494680\du}{2.559788\du}}
\pgfpathlineto{\pgfpoint{-8.533927\du}{2.571473\du}}
\pgfpathlineto{\pgfpoint{-8.571078\du}{2.583457\du}}
\pgfpathlineto{\pgfpoint{-8.605831\du}{2.597238\du}}
\pgfpathlineto{\pgfpoint{-8.639087\du}{2.610421\du}}
\pgfpathlineto{\pgfpoint{-8.669946\du}{2.624502\du}}
\pgfpathlineto{\pgfpoint{-8.698408\du}{2.639482\du}}
\pgfpathlineto{\pgfpoint{-8.724473\du}{2.654462\du}}
\pgfpathlineto{\pgfpoint{-8.749041\du}{2.670640\du}}
\pgfpathlineto{\pgfpoint{-8.771211\du}{2.686220\du}}
\pgfpathlineto{\pgfpoint{-8.790385\du}{2.702997\du}}
\pgfpathlineto{\pgfpoint{-8.807163\du}{2.719176\du}}
\pgfpathlineto{\pgfpoint{-8.822143\du}{2.735953\du}}
\pgfpathlineto{\pgfpoint{-8.833228\du}{2.753031\du}}
\pgfpathlineto{\pgfpoint{-8.842816\du}{2.770707\du}}
\pgfpathlineto{\pgfpoint{-8.849407\du}{2.787485\du}}
\pgfpathlineto{\pgfpoint{-8.853002\du}{2.805161\du}}
\pgfpathlineto{\pgfpoint{-8.854500\du}{2.823137\du}}
\pgfpathlineto{\pgfpoint{-8.854500\du}{2.823137\du}}
\pgfpathlineto{\pgfpoint{-8.853002\du}{2.840514\du}}
\pgfpathlineto{\pgfpoint{-8.849407\du}{2.858190\du}}
\pgfpathlineto{\pgfpoint{-8.842816\du}{2.874968\du}}
\pgfpathlineto{\pgfpoint{-8.833228\du}{2.892345\du}}
\pgfpathlineto{\pgfpoint{-8.822143\du}{2.909722\du}}
\pgfpathlineto{\pgfpoint{-8.807163\du}{2.926499\du}}
\pgfpathlineto{\pgfpoint{-8.790385\du}{2.943277\du}}
\pgfpathlineto{\pgfpoint{-8.771211\du}{2.959755\du}}
\pgfpathlineto{\pgfpoint{-8.749041\du}{2.975634\du}}
\pgfpathlineto{\pgfpoint{-8.724473\du}{2.991213\du}}
\pgfpathlineto{\pgfpoint{-8.698408\du}{3.006493\du}}
\pgfpathlineto{\pgfpoint{-8.669946\du}{3.020574\du}}
\pgfpathlineto{\pgfpoint{-8.639087\du}{3.035254\du}}
\pgfpathlineto{\pgfpoint{-8.605831\du}{3.048736\du}}
\pgfpathlineto{\pgfpoint{-8.571078\du}{3.061919\du}}
\pgfpathlineto{\pgfpoint{-8.533927\du}{3.074203\du}}
\pgfpathlineto{\pgfpoint{-8.494680\du}{3.085887\du}}
\pgfpathlineto{\pgfpoint{-8.453934\du}{3.096972\du}}
\pgfpathlineto{\pgfpoint{-8.412289\du}{3.107159\du}}
\pgfpathlineto{\pgfpoint{-8.368847\du}{3.117045\du}}
\pgfpathlineto{\pgfpoint{-8.323308\du}{3.126333\du}}
\pgfpathlineto{\pgfpoint{-8.277469\du}{3.133823\du}}
\pgfpathlineto{\pgfpoint{-8.230132\du}{3.141613\du}}
\pgfpathlineto{\pgfpoint{-8.181297\du}{3.147904\du}}
\pgfpathlineto{\pgfpoint{-8.131863\du}{3.153597\du}}
\pgfpathlineto{\pgfpoint{-8.081530\du}{3.158690\du}}
\pgfpathlineto{\pgfpoint{-8.030598\du}{3.162884\du}}
\pgfpathlineto{\pgfpoint{-7.979366\du}{3.165281\du}}
\pgfpathlineto{\pgfpoint{-7.927536\du}{3.167678\du}}
\pgfpathlineto{\pgfpoint{-7.875105\du}{3.169176\du}}
\pgfpathlineto{\pgfpoint{-7.822975\du}{3.169476\du}}
\pgfpathlineto{\pgfpoint{-7.822975\du}{3.169476\du}}
\pgfpathlineto{\pgfpoint{-7.770545\du}{3.169176\du}}
\pgfpathlineto{\pgfpoint{-7.718414\du}{3.167678\du}}
\pgfpathlineto{\pgfpoint{-7.666284\du}{3.165281\du}}
\pgfpathlineto{\pgfpoint{-7.615052\du}{3.162884\du}}
\pgfpathlineto{\pgfpoint{-7.564420\du}{3.158690\du}}
\pgfpathlineto{\pgfpoint{-7.513787\du}{3.153597\du}}
\pgfpathlineto{\pgfpoint{-7.464353\du}{3.147904\du}}
\pgfpathlineto{\pgfpoint{-7.415818\du}{3.141613\du}}
\pgfpathlineto{\pgfpoint{-7.368181\du}{3.133823\du}}
\pgfpathlineto{\pgfpoint{-7.322342\du}{3.126333\du}}
\pgfpathlineto{\pgfpoint{-7.277103\du}{3.117045\du}}
\pgfpathlineto{\pgfpoint{-7.233061\du}{3.107159\du}}
\pgfpathlineto{\pgfpoint{-7.191417\du}{3.096972\du}}
\pgfpathlineto{\pgfpoint{-7.150971\du}{3.085887\du}}
\pgfpathlineto{\pgfpoint{-7.112322\du}{3.074203\du}}
\pgfpathlineto{\pgfpoint{-7.075172\du}{3.061919\du}}
\pgfpathlineto{\pgfpoint{-7.039819\du}{3.048736\du}}
\pgfpathlineto{\pgfpoint{-7.006863\du}{3.035254\du}}
\pgfpathlineto{\pgfpoint{-6.976303\du}{3.020574\du}}
\pgfpathlineto{\pgfpoint{-6.947242\du}{3.006493\du}}
\pgfpathlineto{\pgfpoint{-6.920877\du}{2.991213\du}}
\pgfpathlineto{\pgfpoint{-6.896610\du}{2.975634\du}}
\pgfpathlineto{\pgfpoint{-6.874739\du}{2.959755\du}}
\pgfpathlineto{\pgfpoint{-6.854965\du}{2.943277\du}}
\pgfpathlineto{\pgfpoint{-6.838787\du}{2.926499\du}}
\pgfpathlineto{\pgfpoint{-6.824106\du}{2.909722\du}}
\pgfpathlineto{\pgfpoint{-6.812122\du}{2.892345\du}}
\pgfpathlineto{\pgfpoint{-6.803134\du}{2.874968\du}}
\pgfpathlineto{\pgfpoint{-6.796543\du}{2.858190\du}}
\pgfpathlineto{\pgfpoint{-6.792648\du}{2.840514\du}}
\pgfpathlineto{\pgfpoint{-6.791150\du}{2.823137\du}}
\pgfusepath{stroke}
\pgfsetbuttcap
\pgfsetmiterjoin
\pgfsetdash{}{0pt}
\definecolor{dialinecolor}{rgb}{0.000000, 0.000000, 0.000000}
\pgfsetfillcolor{dialinecolor}
\pgfpathmoveto{\pgfpoint{-7.798108\du}{2.748237\du}}
\pgfpathlineto{\pgfpoint{-7.646810\du}{2.797971\du}}
\pgfpathlineto{\pgfpoint{-7.281597\du}{2.649069\du}}
\pgfpathlineto{\pgfpoint{-7.118314\du}{2.698803\du}}
\pgfpathlineto{\pgfpoint{-7.206996\du}{2.575068\du}}
\pgfpathlineto{\pgfpoint{-7.633927\du}{2.575068\du}}
\pgfpathlineto{\pgfpoint{-7.458361\du}{2.612518\du}}
\pgfpathlineto{\pgfpoint{-7.798108\du}{2.748237\du}}
\pgfusepath{fill}
\pgfsetbuttcap
\pgfsetmiterjoin
\pgfsetdash{}{0pt}
\definecolor{dialinecolor}{rgb}{0.000000, 0.000000, 0.000000}
\pgfsetfillcolor{dialinecolor}
\pgfpathmoveto{\pgfpoint{-7.861324\du}{2.884855\du}}
\pgfpathlineto{\pgfpoint{-8.011424\du}{2.835121\du}}
\pgfpathlineto{\pgfpoint{-8.364054\du}{2.984023\du}}
\pgfpathlineto{\pgfpoint{-8.539919\du}{2.933989\du}}
\pgfpathlineto{\pgfpoint{-8.451837\du}{3.070308\du}}
\pgfpathlineto{\pgfpoint{-8.011424\du}{3.070308\du}}
\pgfpathlineto{\pgfpoint{-8.200472\du}{3.020574\du}}
\pgfpathlineto{\pgfpoint{-7.861324\du}{2.884855\du}}
\pgfusepath{fill}
\pgfsetbuttcap
\pgfsetmiterjoin
\pgfsetdash{}{0pt}
\definecolor{dialinecolor}{rgb}{0.000000, 0.000000, 0.000000}
\pgfsetfillcolor{dialinecolor}
\pgfpathmoveto{\pgfpoint{-8.502469\du}{2.612518\du}}
\pgfpathlineto{\pgfpoint{-8.351470\du}{2.562784\du}}
\pgfpathlineto{\pgfpoint{-7.986856\du}{2.698803\du}}
\pgfpathlineto{\pgfpoint{-7.822975\du}{2.661952\du}}
\pgfpathlineto{\pgfpoint{-7.910758\du}{2.785387\du}}
\pgfpathlineto{\pgfpoint{-8.338887\du}{2.785387\du}}
\pgfpathlineto{\pgfpoint{-8.163022\du}{2.748237\du}}
\pgfpathlineto{\pgfpoint{-8.502469\du}{2.612518\du}}
\pgfusepath{fill}
\pgfsetbuttcap
\pgfsetmiterjoin
\pgfsetdash{}{0pt}
\definecolor{dialinecolor}{rgb}{0.000000, 0.000000, 0.000000}
\pgfsetfillcolor{dialinecolor}
\pgfpathmoveto{\pgfpoint{-7.143780\du}{3.033157\du}}
\pgfpathlineto{\pgfpoint{-7.294479\du}{3.082591\du}}
\pgfpathlineto{\pgfpoint{-7.646810\du}{2.933989\du}}
\pgfpathlineto{\pgfpoint{-7.822975\du}{2.984023\du}}
\pgfpathlineto{\pgfpoint{-7.734892\du}{2.859988\du}}
\pgfpathlineto{\pgfpoint{-7.294479\du}{2.859988\du}}
\pgfpathlineto{\pgfpoint{-7.483527\du}{2.896539\du}}
\pgfpathlineto{\pgfpoint{-7.143780\du}{3.033157\du}}
\pgfusepath{fill}
\pgfsetbuttcap
\pgfsetmiterjoin
\pgfsetdash{}{0pt}
\definecolor{dialinecolor}{rgb}{1.000000, 1.000000, 1.000000}
\pgfsetfillcolor{dialinecolor}
\pgfpathmoveto{\pgfpoint{-7.785525\du}{2.760521\du}}
\pgfpathlineto{\pgfpoint{-7.633927\du}{2.809955\du}}
\pgfpathlineto{\pgfpoint{-7.269313\du}{2.661952\du}}
\pgfpathlineto{\pgfpoint{-7.105731\du}{2.711686\du}}
\pgfpathlineto{\pgfpoint{-7.193814\du}{2.587351\du}}
\pgfpathlineto{\pgfpoint{-7.621943\du}{2.587351\du}}
\pgfpathlineto{\pgfpoint{-7.445478\du}{2.624502\du}}
\pgfpathlineto{\pgfpoint{-7.785525\du}{2.760521\du}}
\pgfusepath{fill}
\pgfsetbuttcap
\pgfsetmiterjoin
\pgfsetdash{}{0pt}
\definecolor{dialinecolor}{rgb}{1.000000, 1.000000, 1.000000}
\pgfsetfillcolor{dialinecolor}
\pgfpathmoveto{\pgfpoint{-7.847842\du}{2.896539\du}}
\pgfpathlineto{\pgfpoint{-7.999440\du}{2.847405\du}}
\pgfpathlineto{\pgfpoint{-8.351470\du}{2.996007\du}}
\pgfpathlineto{\pgfpoint{-8.527636\du}{2.946273\du}}
\pgfpathlineto{\pgfpoint{-8.439553\du}{3.082591\du}}
\pgfpathlineto{\pgfpoint{-7.999440\du}{3.082591\du}}
\pgfpathlineto{\pgfpoint{-8.187589\du}{3.033157\du}}
\pgfpathlineto{\pgfpoint{-7.847842\du}{2.896539\du}}
\pgfusepath{fill}
\pgfsetbuttcap
\pgfsetmiterjoin
\pgfsetdash{}{0pt}
\definecolor{dialinecolor}{rgb}{1.000000, 1.000000, 1.000000}
\pgfsetfillcolor{dialinecolor}
\pgfpathmoveto{\pgfpoint{-8.489586\du}{2.624502\du}}
\pgfpathlineto{\pgfpoint{-8.338887\du}{2.575068\du}}
\pgfpathlineto{\pgfpoint{-7.973974\du}{2.711686\du}}
\pgfpathlineto{\pgfpoint{-7.810392\du}{2.673936\du}}
\pgfpathlineto{\pgfpoint{-7.898474\du}{2.797971\du}}
\pgfpathlineto{\pgfpoint{-8.326304\du}{2.797971\du}}
\pgfpathlineto{\pgfpoint{-8.149839\du}{2.760521\du}}
\pgfpathlineto{\pgfpoint{-8.489586\du}{2.624502\du}}
\pgfusepath{fill}
\pgfsetbuttcap
\pgfsetmiterjoin
\pgfsetdash{}{0pt}
\definecolor{dialinecolor}{rgb}{1.000000, 1.000000, 1.000000}
\pgfsetfillcolor{dialinecolor}
\pgfpathmoveto{\pgfpoint{-7.130897\du}{3.045441\du}}
\pgfpathlineto{\pgfpoint{-7.281597\du}{3.095175\du}}
\pgfpathlineto{\pgfpoint{-7.633927\du}{2.946273\du}}
\pgfpathlineto{\pgfpoint{-7.810392\du}{2.996007\du}}
\pgfpathlineto{\pgfpoint{-7.722609\du}{2.872571\du}}
\pgfpathlineto{\pgfpoint{-7.281597\du}{2.872571\du}}
\pgfpathlineto{\pgfpoint{-7.470345\du}{2.909422\du}}
\pgfpathlineto{\pgfpoint{-7.130897\du}{3.045441\du}}
\pgfusepath{fill}
\pgfsetbuttcap
\pgfsetmiterjoin
\pgfsetdash{}{0pt}
\definecolor{dialinecolor}{rgb}{0.666667, 0.901961, 1.000000}
\pgfsetstrokecolor{dialinecolor}
\pgfpathmoveto{\pgfpoint{-8.854500\du}{2.823137\du}}
\pgfpathlineto{\pgfpoint{-8.854500\du}{3.317478\du}}
\pgfusepath{stroke}
\pgfsetbuttcap
\pgfsetmiterjoin
\pgfsetdash{}{0pt}
\definecolor{dialinecolor}{rgb}{0.666667, 0.901961, 1.000000}
\pgfsetstrokecolor{dialinecolor}
\pgfpathmoveto{\pgfpoint{-6.791150\du}{2.823137\du}}
\pgfpathlineto{\pgfpoint{-6.791150\du}{3.317478\du}}
\pgfusepath{stroke}
\pgfsetbuttcap
\pgfsetmiterjoin
\pgfsetdash{}{0pt}
\definecolor{dialinecolor}{rgb}{0.000000, 0.000000, 0.000000}
\pgfsetfillcolor{dialinecolor}
\pgfpathmoveto{\pgfpoint{-8.200472\du}{3.256060\du}}
\pgfpathlineto{\pgfpoint{-7.973974\du}{3.256060\du}}
\pgfpathlineto{\pgfpoint{-7.835558\du}{3.367512\du}}
\pgfpathlineto{\pgfpoint{-7.697442\du}{3.256060\du}}
\pgfpathlineto{\pgfpoint{-7.483527\du}{3.256060\du}}
\pgfpathlineto{\pgfpoint{-7.483527\du}{3.206626\du}}
\pgfpathlineto{\pgfpoint{-7.332229\du}{3.280328\du}}
\pgfpathlineto{\pgfpoint{-7.483527\du}{3.354928\du}}
\pgfpathlineto{\pgfpoint{-7.483527\du}{3.305494\du}}
\pgfpathlineto{\pgfpoint{-7.671976\du}{3.305494\du}}
\pgfpathlineto{\pgfpoint{-7.798108\du}{3.392079\du}}
\pgfpathlineto{\pgfpoint{-7.671976\du}{3.490647\du}}
\pgfpathlineto{\pgfpoint{-7.483527\du}{3.490647\du}}
\pgfpathlineto{\pgfpoint{-7.483527\du}{3.429529\du}}
\pgfpathlineto{\pgfpoint{-7.332229\du}{3.515814\du}}
\pgfpathlineto{\pgfpoint{-7.483527\du}{3.590414\du}}
\pgfpathlineto{\pgfpoint{-7.483527\du}{3.540381\du}}
\pgfpathlineto{\pgfpoint{-7.697442\du}{3.540381\du}}
\pgfpathlineto{\pgfpoint{-7.835558\du}{3.429529\du}}
\pgfpathlineto{\pgfpoint{-7.973974\du}{3.540381\du}}
\pgfpathlineto{\pgfpoint{-8.200472\du}{3.540381\du}}
\pgfpathlineto{\pgfpoint{-8.200472\du}{3.590414\du}}
\pgfpathlineto{\pgfpoint{-8.351470\du}{3.515814\du}}
\pgfpathlineto{\pgfpoint{-8.200472\du}{3.429529\du}}
\pgfpathlineto{\pgfpoint{-8.200472\du}{3.490647\du}}
\pgfpathlineto{\pgfpoint{-7.999440\du}{3.490647\du}}
\pgfpathlineto{\pgfpoint{-7.873308\du}{3.404362\du}}
\pgfpathlineto{\pgfpoint{-7.999440\du}{3.305494\du}}
\pgfpathlineto{\pgfpoint{-8.200472\du}{3.305494\du}}
\pgfpathlineto{\pgfpoint{-8.200472\du}{3.354928\du}}
\pgfpathlineto{\pgfpoint{-8.351470\du}{3.280328\du}}
\pgfpathlineto{\pgfpoint{-8.200472\du}{3.206626\du}}
\pgfpathlineto{\pgfpoint{-8.200472\du}{3.256060\du}}
\pgfusepath{fill}
\pgfsetbuttcap
\pgfsetmiterjoin
\pgfsetdash{}{0pt}
\definecolor{dialinecolor}{rgb}{1.000000, 1.000000, 1.000000}
\pgfsetfillcolor{dialinecolor}
\pgfpathmoveto{\pgfpoint{-8.187589\du}{3.268044\du}}
\pgfpathlineto{\pgfpoint{-7.961390\du}{3.268044\du}}
\pgfpathlineto{\pgfpoint{-7.822975\du}{3.379196\du}}
\pgfpathlineto{\pgfpoint{-7.684859\du}{3.268044\du}}
\pgfpathlineto{\pgfpoint{-7.470345\du}{3.268044\du}}
\pgfpathlineto{\pgfpoint{-7.470345\du}{3.218610\du}}
\pgfpathlineto{\pgfpoint{-7.319945\du}{3.293211\du}}
\pgfpathlineto{\pgfpoint{-7.470345\du}{3.367512\du}}
\pgfpathlineto{\pgfpoint{-7.470345\du}{3.318077\du}}
\pgfpathlineto{\pgfpoint{-7.659393\du}{3.318077\du}}
\pgfpathlineto{\pgfpoint{-7.785525\du}{3.404362\du}}
\pgfpathlineto{\pgfpoint{-7.659393\du}{3.503231\du}}
\pgfpathlineto{\pgfpoint{-7.470345\du}{3.503231\du}}
\pgfpathlineto{\pgfpoint{-7.470345\du}{3.441513\du}}
\pgfpathlineto{\pgfpoint{-7.319945\du}{3.528098\du}}
\pgfpathlineto{\pgfpoint{-7.470345\du}{3.602398\du}}
\pgfpathlineto{\pgfpoint{-7.470345\du}{3.552665\du}}
\pgfpathlineto{\pgfpoint{-7.684859\du}{3.552665\du}}
\pgfpathlineto{\pgfpoint{-7.822975\du}{3.441513\du}}
\pgfpathlineto{\pgfpoint{-7.961390\du}{3.552665\du}}
\pgfpathlineto{\pgfpoint{-8.187589\du}{3.552665\du}}
\pgfpathlineto{\pgfpoint{-8.187589\du}{3.602398\du}}
\pgfpathlineto{\pgfpoint{-8.338887\du}{3.528098\du}}
\pgfpathlineto{\pgfpoint{-8.187589\du}{3.441513\du}}
\pgfpathlineto{\pgfpoint{-8.187589\du}{3.503231\du}}
\pgfpathlineto{\pgfpoint{-7.986856\du}{3.503231\du}}
\pgfpathlineto{\pgfpoint{-7.861324\du}{3.416646\du}}
\pgfpathlineto{\pgfpoint{-7.986856\du}{3.318077\du}}
\pgfpathlineto{\pgfpoint{-8.187589\du}{3.318077\du}}
\pgfpathlineto{\pgfpoint{-8.187589\du}{3.367512\du}}
\pgfpathlineto{\pgfpoint{-8.338887\du}{3.293211\du}}
\pgfpathlineto{\pgfpoint{-8.187589\du}{3.218610\du}}
\pgfpathlineto{\pgfpoint{-8.187589\du}{3.268044\du}}
\pgfusepath{fill}
\pgfsetlinewidth{0.100000\du}
\pgfsetdash{}{0pt}
\pgfsetdash{}{0pt}
\pgfsetbuttcap
\pgfsetmiterjoin
\pgfsetlinewidth{0.001000\du}
\pgfsetbuttcap
\pgfsetmiterjoin
\pgfsetdash{}{0pt}
\definecolor{dialinecolor}{rgb}{0.000000, 0.470588, 0.666667}
\pgfsetfillcolor{dialinecolor}
\pgfpathmoveto{\pgfpoint{-21.460118\du}{17.791822\du}}
\pgfpathlineto{\pgfpoint{-21.461406\du}{17.776627\du}}
\pgfpathlineto{\pgfpoint{-21.464754\du}{17.761173\du}}
\pgfpathlineto{\pgfpoint{-21.470420\du}{17.745978\du}}
\pgfpathlineto{\pgfpoint{-21.478405\du}{17.731039\du}}
\pgfpathlineto{\pgfpoint{-21.488707\du}{17.715844\du}}
\pgfpathlineto{\pgfpoint{-21.501069\du}{17.700906\du}}
\pgfpathlineto{\pgfpoint{-21.515750\du}{17.686483\du}}
\pgfpathlineto{\pgfpoint{-21.532234\du}{17.672317\du}}
\pgfpathlineto{\pgfpoint{-21.551293\du}{17.658409\du}}
\pgfpathlineto{\pgfpoint{-21.572155\du}{17.645016\du}}
\pgfpathlineto{\pgfpoint{-21.595077\du}{17.631623\du}}
\pgfpathlineto{\pgfpoint{-21.620060\du}{17.618746\du}}
\pgfpathlineto{\pgfpoint{-21.646845\du}{17.606383\du}}
\pgfpathlineto{\pgfpoint{-21.675434\du}{17.594278\du}}
\pgfpathlineto{\pgfpoint{-21.705825\du}{17.582946\du}}
\pgfpathlineto{\pgfpoint{-21.737505\du}{17.572128\du}}
\pgfpathlineto{\pgfpoint{-21.770987\du}{17.561826\du}}
\pgfpathlineto{\pgfpoint{-21.806272\du}{17.552296\du}}
\pgfpathlineto{\pgfpoint{-21.842587\du}{17.542767\du}}
\pgfpathlineto{\pgfpoint{-21.880190\du}{17.534783\du}}
\pgfpathlineto{\pgfpoint{-21.919081\du}{17.526799\du}}
\pgfpathlineto{\pgfpoint{-21.959259\du}{17.519845\du}}
\pgfpathlineto{\pgfpoint{-22.000211\du}{17.513148\du}}
\pgfpathlineto{\pgfpoint{-22.042449\du}{17.507482\du}}
\pgfpathlineto{\pgfpoint{-22.084946\du}{17.502588\du}}
\pgfpathlineto{\pgfpoint{-22.128730\du}{17.498468\du}}
\pgfpathlineto{\pgfpoint{-22.172772\du}{17.494862\du}}
\pgfpathlineto{\pgfpoint{-22.217072\du}{17.492029\du}}
\pgfpathlineto{\pgfpoint{-22.261886\du}{17.489968\du}}
\pgfpathlineto{\pgfpoint{-22.306958\du}{17.489196\du}}
\pgfpathlineto{\pgfpoint{-22.352545\du}{17.488423\du}}
\pgfpathlineto{\pgfpoint{-22.352545\du}{17.488423\du}}
\pgfpathlineto{\pgfpoint{-22.397618\du}{17.489196\du}}
\pgfpathlineto{\pgfpoint{-22.442432\du}{17.489968\du}}
\pgfpathlineto{\pgfpoint{-22.487247\du}{17.492029\du}}
\pgfpathlineto{\pgfpoint{-22.531546\du}{17.494862\du}}
\pgfpathlineto{\pgfpoint{-22.575588\du}{17.498468\du}}
\pgfpathlineto{\pgfpoint{-22.619372\du}{17.502588\du}}
\pgfpathlineto{\pgfpoint{-22.661869\du}{17.507482\du}}
\pgfpathlineto{\pgfpoint{-22.704108\du}{17.513148\du}}
\pgfpathlineto{\pgfpoint{-22.745059\du}{17.519845\du}}
\pgfpathlineto{\pgfpoint{-22.785237\du}{17.526799\du}}
\pgfpathlineto{\pgfpoint{-22.824386\du}{17.534783\du}}
\pgfpathlineto{\pgfpoint{-22.861731\du}{17.542767\du}}
\pgfpathlineto{\pgfpoint{-22.898046\du}{17.552296\du}}
\pgfpathlineto{\pgfpoint{-22.933331\du}{17.561826\du}}
\pgfpathlineto{\pgfpoint{-22.966814\du}{17.572128\du}}
\pgfpathlineto{\pgfpoint{-22.998493\du}{17.582946\du}}
\pgfpathlineto{\pgfpoint{-23.028884\du}{17.594278\du}}
\pgfpathlineto{\pgfpoint{-23.057473\du}{17.606383\du}}
\pgfpathlineto{\pgfpoint{-23.084258\du}{17.618746\du}}
\pgfpathlineto{\pgfpoint{-23.109499\du}{17.631623\du}}
\pgfpathlineto{\pgfpoint{-23.132164\du}{17.645016\du}}
\pgfpathlineto{\pgfpoint{-23.153283\du}{17.658409\du}}
\pgfpathlineto{\pgfpoint{-23.172085\du}{17.672317\du}}
\pgfpathlineto{\pgfpoint{-23.188568\du}{17.686483\du}}
\pgfpathlineto{\pgfpoint{-23.203249\du}{17.700906\du}}
\pgfpathlineto{\pgfpoint{-23.215611\du}{17.715844\du}}
\pgfpathlineto{\pgfpoint{-23.225914\du}{17.731039\du}}
\pgfpathlineto{\pgfpoint{-23.233898\du}{17.745978\du}}
\pgfpathlineto{\pgfpoint{-23.239822\du}{17.761173\du}}
\pgfpathlineto{\pgfpoint{-23.242912\du}{17.776627\du}}
\pgfpathlineto{\pgfpoint{-23.244200\du}{17.791822\du}}
\pgfpathlineto{\pgfpoint{-23.244200\du}{17.791822\du}}
\pgfpathlineto{\pgfpoint{-23.242912\du}{17.807533\du}}
\pgfpathlineto{\pgfpoint{-23.239822\du}{17.822729\du}}
\pgfpathlineto{\pgfpoint{-23.233898\du}{17.837925\du}}
\pgfpathlineto{\pgfpoint{-23.225914\du}{17.853120\du}}
\pgfpathlineto{\pgfpoint{-23.215611\du}{17.868059\du}}
\pgfpathlineto{\pgfpoint{-23.203249\du}{17.882997\du}}
\pgfpathlineto{\pgfpoint{-23.188568\du}{17.897420\du}}
\pgfpathlineto{\pgfpoint{-23.172085\du}{17.911585\du}}
\pgfpathlineto{\pgfpoint{-23.153283\du}{17.925493\du}}
\pgfpathlineto{\pgfpoint{-23.132164\du}{17.939144\du}}
\pgfpathlineto{\pgfpoint{-23.109499\du}{17.952021\du}}
\pgfpathlineto{\pgfpoint{-23.084258\du}{17.965414\du}}
\pgfpathlineto{\pgfpoint{-23.057473\du}{17.977777\du}}
\pgfpathlineto{\pgfpoint{-23.028884\du}{17.989624\du}}
\pgfpathlineto{\pgfpoint{-22.998493\du}{18.000699\du}}
\pgfpathlineto{\pgfpoint{-22.966814\du}{18.011774\du}}
\pgfpathlineto{\pgfpoint{-22.933331\du}{18.022076\du}}
\pgfpathlineto{\pgfpoint{-22.898046\du}{18.031863\du}}
\pgfpathlineto{\pgfpoint{-22.861731\du}{18.040878\du}}
\pgfpathlineto{\pgfpoint{-22.824386\du}{18.049377\du}}
\pgfpathlineto{\pgfpoint{-22.785237\du}{18.057104\du}}
\pgfpathlineto{\pgfpoint{-22.745059\du}{18.064058\du}}
\pgfpathlineto{\pgfpoint{-22.704108\du}{18.070754\du}}
\pgfpathlineto{\pgfpoint{-22.661869\du}{18.076420\du}}
\pgfpathlineto{\pgfpoint{-22.619372\du}{18.081314\du}}
\pgfpathlineto{\pgfpoint{-22.575588\du}{18.085692\du}}
\pgfpathlineto{\pgfpoint{-22.531546\du}{18.088783\du}}
\pgfpathlineto{\pgfpoint{-22.487247\du}{18.091616\du}}
\pgfpathlineto{\pgfpoint{-22.442432\du}{18.093677\du}}
\pgfpathlineto{\pgfpoint{-22.397618\du}{18.094964\du}}
\pgfpathlineto{\pgfpoint{-22.352545\du}{18.095222\du}}
\pgfpathlineto{\pgfpoint{-22.352545\du}{18.095222\du}}
\pgfpathlineto{\pgfpoint{-22.306958\du}{18.094964\du}}
\pgfpathlineto{\pgfpoint{-22.261886\du}{18.093677\du}}
\pgfpathlineto{\pgfpoint{-22.217072\du}{18.091616\du}}
\pgfpathlineto{\pgfpoint{-22.172772\du}{18.088783\du}}
\pgfpathlineto{\pgfpoint{-22.128730\du}{18.085692\du}}
\pgfpathlineto{\pgfpoint{-22.084946\du}{18.081314\du}}
\pgfpathlineto{\pgfpoint{-22.042449\du}{18.076420\du}}
\pgfpathlineto{\pgfpoint{-22.000211\du}{18.070754\du}}
\pgfpathlineto{\pgfpoint{-21.959259\du}{18.064058\du}}
\pgfpathlineto{\pgfpoint{-21.919081\du}{18.057104\du}}
\pgfpathlineto{\pgfpoint{-21.880190\du}{18.049377\du}}
\pgfpathlineto{\pgfpoint{-21.842587\du}{18.040878\du}}
\pgfpathlineto{\pgfpoint{-21.806272\du}{18.031863\du}}
\pgfpathlineto{\pgfpoint{-21.770987\du}{18.022076\du}}
\pgfpathlineto{\pgfpoint{-21.737505\du}{18.011774\du}}
\pgfpathlineto{\pgfpoint{-21.705825\du}{18.000699\du}}
\pgfpathlineto{\pgfpoint{-21.675434\du}{17.989624\du}}
\pgfpathlineto{\pgfpoint{-21.646845\du}{17.977777\du}}
\pgfpathlineto{\pgfpoint{-21.620060\du}{17.965414\du}}
\pgfpathlineto{\pgfpoint{-21.595077\du}{17.952021\du}}
\pgfpathlineto{\pgfpoint{-21.572155\du}{17.939144\du}}
\pgfpathlineto{\pgfpoint{-21.551293\du}{17.925493\du}}
\pgfpathlineto{\pgfpoint{-21.532234\du}{17.911585\du}}
\pgfpathlineto{\pgfpoint{-21.515750\du}{17.897420\du}}
\pgfpathlineto{\pgfpoint{-21.501069\du}{17.882997\du}}
\pgfpathlineto{\pgfpoint{-21.488707\du}{17.868059\du}}
\pgfpathlineto{\pgfpoint{-21.478405\du}{17.853120\du}}
\pgfpathlineto{\pgfpoint{-21.470420\du}{17.837925\du}}
\pgfpathlineto{\pgfpoint{-21.464754\du}{17.822729\du}}
\pgfpathlineto{\pgfpoint{-21.461406\du}{17.807533\du}}
\pgfpathlineto{\pgfpoint{-21.460118\du}{17.791822\du}}
\pgfusepath{fill}
\pgfsetbuttcap
\pgfsetmiterjoin
\pgfsetdash{}{0pt}
\definecolor{dialinecolor}{rgb}{0.666667, 0.901961, 1.000000}
\pgfsetstrokecolor{dialinecolor}
\pgfpathmoveto{\pgfpoint{-21.470420\du}{17.787186\du}}
\pgfpathlineto{\pgfpoint{-21.471708\du}{17.771733\du}}
\pgfpathlineto{\pgfpoint{-21.475056\du}{17.757053\du}}
\pgfpathlineto{\pgfpoint{-21.480723\du}{17.741857\du}}
\pgfpathlineto{\pgfpoint{-21.488449\du}{17.727176\du}}
\pgfpathlineto{\pgfpoint{-21.498751\du}{17.712238\du}}
\pgfpathlineto{\pgfpoint{-21.511372\du}{17.697815\du}}
\pgfpathlineto{\pgfpoint{-21.525280\du}{17.683649\du}}
\pgfpathlineto{\pgfpoint{-21.542278\du}{17.669484\du}}
\pgfpathlineto{\pgfpoint{-21.561080\du}{17.655834\du}}
\pgfpathlineto{\pgfpoint{-21.581942\du}{17.642183\du}}
\pgfpathlineto{\pgfpoint{-21.604606\du}{17.629563\du}}
\pgfpathlineto{\pgfpoint{-21.629589\du}{17.616943\du}}
\pgfpathlineto{\pgfpoint{-21.655860\du}{17.604323\du}}
\pgfpathlineto{\pgfpoint{-21.684191\du}{17.592733\du}}
\pgfpathlineto{\pgfpoint{-21.714582\du}{17.581400\du}}
\pgfpathlineto{\pgfpoint{-21.746519\du}{17.571098\du}}
\pgfpathlineto{\pgfpoint{-21.779744\du}{17.560796\du}}
\pgfpathlineto{\pgfpoint{-21.814514\du}{17.551266\du}}
\pgfpathlineto{\pgfpoint{-21.850314\du}{17.542509\du}}
\pgfpathlineto{\pgfpoint{-21.888174\du}{17.534010\du}}
\pgfpathlineto{\pgfpoint{-21.927065\du}{17.526541\du}}
\pgfpathlineto{\pgfpoint{-21.966471\du}{17.519587\du}}
\pgfpathlineto{\pgfpoint{-22.007422\du}{17.513148\du}}
\pgfpathlineto{\pgfpoint{-22.049146\du}{17.507482\du}}
\pgfpathlineto{\pgfpoint{-22.091642\du}{17.502588\du}}
\pgfpathlineto{\pgfpoint{-22.135169\du}{17.498468\du}}
\pgfpathlineto{\pgfpoint{-22.178696\du}{17.494862\du}}
\pgfpathlineto{\pgfpoint{-22.222738\du}{17.492801\du}}
\pgfpathlineto{\pgfpoint{-22.267552\du}{17.490483\du}}
\pgfpathlineto{\pgfpoint{-22.312367\du}{17.489196\du}}
\pgfpathlineto{\pgfpoint{-22.357439\du}{17.489196\du}}
\pgfpathlineto{\pgfpoint{-22.357439\du}{17.489196\du}}
\pgfpathlineto{\pgfpoint{-22.402253\du}{17.489196\du}}
\pgfpathlineto{\pgfpoint{-22.447326\du}{17.490483\du}}
\pgfpathlineto{\pgfpoint{-22.491883\du}{17.492801\du}}
\pgfpathlineto{\pgfpoint{-22.535924\du}{17.494862\du}}
\pgfpathlineto{\pgfpoint{-22.579709\du}{17.498468\du}}
\pgfpathlineto{\pgfpoint{-22.622978\du}{17.502588\du}}
\pgfpathlineto{\pgfpoint{-22.665474\du}{17.507482\du}}
\pgfpathlineto{\pgfpoint{-22.707456\du}{17.513148\du}}
\pgfpathlineto{\pgfpoint{-22.748150\du}{17.519587\du}}
\pgfpathlineto{\pgfpoint{-22.787555\du}{17.526541\du}}
\pgfpathlineto{\pgfpoint{-22.826704\du}{17.534010\du}}
\pgfpathlineto{\pgfpoint{-22.864049\du}{17.542509\du}}
\pgfpathlineto{\pgfpoint{-22.899849\du}{17.551266\du}}
\pgfpathlineto{\pgfpoint{-22.934877\du}{17.560796\du}}
\pgfpathlineto{\pgfpoint{-22.968616\du}{17.571098\du}}
\pgfpathlineto{\pgfpoint{-23.000553\du}{17.581400\du}}
\pgfpathlineto{\pgfpoint{-23.030430\du}{17.592733\du}}
\pgfpathlineto{\pgfpoint{-23.059018\du}{17.604323\du}}
\pgfpathlineto{\pgfpoint{-23.085546\du}{17.616943\du}}
\pgfpathlineto{\pgfpoint{-23.110014\du}{17.629563\du}}
\pgfpathlineto{\pgfpoint{-23.132421\du}{17.642183\du}}
\pgfpathlineto{\pgfpoint{-23.153541\du}{17.655834\du}}
\pgfpathlineto{\pgfpoint{-23.172600\du}{17.669484\du}}
\pgfpathlineto{\pgfpoint{-23.189083\du}{17.683649\du}}
\pgfpathlineto{\pgfpoint{-23.203506\du}{17.697815\du}}
\pgfpathlineto{\pgfpoint{-23.216384\du}{17.712238\du}}
\pgfpathlineto{\pgfpoint{-23.225914\du}{17.727176\du}}
\pgfpathlineto{\pgfpoint{-23.234155\du}{17.741857\du}}
\pgfpathlineto{\pgfpoint{-23.239822\du}{17.757053\du}}
\pgfpathlineto{\pgfpoint{-23.242912\du}{17.771733\du}}
\pgfpathlineto{\pgfpoint{-23.244200\du}{17.787186\du}}
\pgfpathlineto{\pgfpoint{-23.244200\du}{17.787186\du}}
\pgfpathlineto{\pgfpoint{-23.242912\du}{17.802382\du}}
\pgfpathlineto{\pgfpoint{-23.239822\du}{17.817320\du}}
\pgfpathlineto{\pgfpoint{-23.234155\du}{17.832259\du}}
\pgfpathlineto{\pgfpoint{-23.225914\du}{17.846939\du}}
\pgfpathlineto{\pgfpoint{-23.216384\du}{17.861877\du}}
\pgfpathlineto{\pgfpoint{-23.203506\du}{17.876300\du}}
\pgfpathlineto{\pgfpoint{-23.189083\du}{17.890466\du}}
\pgfpathlineto{\pgfpoint{-23.172600\du}{17.904116\du}}
\pgfpathlineto{\pgfpoint{-23.153541\du}{17.918282\du}}
\pgfpathlineto{\pgfpoint{-23.132421\du}{17.931932\du}}
\pgfpathlineto{\pgfpoint{-23.110014\du}{17.944810\du}}
\pgfpathlineto{\pgfpoint{-23.085546\du}{17.957173\du}}
\pgfpathlineto{\pgfpoint{-23.059018\du}{17.969535\du}}
\pgfpathlineto{\pgfpoint{-23.030430\du}{17.981383\du}}
\pgfpathlineto{\pgfpoint{-23.000553\du}{17.992458\du}}
\pgfpathlineto{\pgfpoint{-22.968616\du}{18.003017\du}}
\pgfpathlineto{\pgfpoint{-22.934877\du}{18.012804\du}}
\pgfpathlineto{\pgfpoint{-22.899849\du}{18.022849\du}}
\pgfpathlineto{\pgfpoint{-22.864049\du}{18.031606\du}}
\pgfpathlineto{\pgfpoint{-22.826704\du}{18.039848\du}}
\pgfpathlineto{\pgfpoint{-22.787555\du}{18.047832\du}}
\pgfpathlineto{\pgfpoint{-22.748150\du}{18.054528\du}}
\pgfpathlineto{\pgfpoint{-22.707456\du}{18.060710\du}}
\pgfpathlineto{\pgfpoint{-22.665474\du}{18.066633\du}}
\pgfpathlineto{\pgfpoint{-22.622978\du}{18.071269\du}}
\pgfpathlineto{\pgfpoint{-22.579709\du}{18.075133\du}}
\pgfpathlineto{\pgfpoint{-22.535924\du}{18.078996\du}}
\pgfpathlineto{\pgfpoint{-22.491883\du}{18.081314\du}}
\pgfpathlineto{\pgfpoint{-22.447326\du}{18.083374\du}}
\pgfpathlineto{\pgfpoint{-22.402253\du}{18.084405\du}}
\pgfpathlineto{\pgfpoint{-22.357439\du}{18.084920\du}}
\pgfpathlineto{\pgfpoint{-22.357439\du}{18.084920\du}}
\pgfpathlineto{\pgfpoint{-22.312367\du}{18.084405\du}}
\pgfpathlineto{\pgfpoint{-22.267552\du}{18.083374\du}}
\pgfpathlineto{\pgfpoint{-22.222738\du}{18.081314\du}}
\pgfpathlineto{\pgfpoint{-22.178696\du}{18.078996\du}}
\pgfpathlineto{\pgfpoint{-22.135169\du}{18.075133\du}}
\pgfpathlineto{\pgfpoint{-22.091642\du}{18.071269\du}}
\pgfpathlineto{\pgfpoint{-22.049146\du}{18.066633\du}}
\pgfpathlineto{\pgfpoint{-22.007422\du}{18.060710\du}}
\pgfpathlineto{\pgfpoint{-21.966471\du}{18.054528\du}}
\pgfpathlineto{\pgfpoint{-21.927065\du}{18.047832\du}}
\pgfpathlineto{\pgfpoint{-21.888174\du}{18.039848\du}}
\pgfpathlineto{\pgfpoint{-21.850314\du}{18.031606\du}}
\pgfpathlineto{\pgfpoint{-21.814514\du}{18.022849\du}}
\pgfpathlineto{\pgfpoint{-21.779744\du}{18.012804\du}}
\pgfpathlineto{\pgfpoint{-21.746519\du}{18.003017\du}}
\pgfpathlineto{\pgfpoint{-21.714582\du}{17.992458\du}}
\pgfpathlineto{\pgfpoint{-21.684191\du}{17.981383\du}}
\pgfpathlineto{\pgfpoint{-21.655860\du}{17.969535\du}}
\pgfpathlineto{\pgfpoint{-21.629589\du}{17.957173\du}}
\pgfpathlineto{\pgfpoint{-21.604606\du}{17.944810\du}}
\pgfpathlineto{\pgfpoint{-21.581942\du}{17.931932\du}}
\pgfpathlineto{\pgfpoint{-21.561080\du}{17.918282\du}}
\pgfpathlineto{\pgfpoint{-21.542278\du}{17.904116\du}}
\pgfpathlineto{\pgfpoint{-21.525280\du}{17.890466\du}}
\pgfpathlineto{\pgfpoint{-21.511372\du}{17.876300\du}}
\pgfpathlineto{\pgfpoint{-21.498751\du}{17.861877\du}}
\pgfpathlineto{\pgfpoint{-21.488449\du}{17.846939\du}}
\pgfpathlineto{\pgfpoint{-21.480723\du}{17.832259\du}}
\pgfpathlineto{\pgfpoint{-21.475056\du}{17.817320\du}}
\pgfpathlineto{\pgfpoint{-21.471708\du}{17.802382\du}}
\pgfpathlineto{\pgfpoint{-21.470420\du}{17.787186\du}}
\pgfusepath{stroke}
\pgfsetbuttcap
\pgfsetmiterjoin
\pgfsetdash{}{0pt}
\definecolor{dialinecolor}{rgb}{0.000000, 0.470588, 0.666667}
\pgfsetfillcolor{dialinecolor}
\pgfpathmoveto{\pgfpoint{-23.244200\du}{17.372008\du}}
\pgfpathlineto{\pgfpoint{-23.244200\du}{17.797489\du}}
\pgfpathlineto{\pgfpoint{-21.470420\du}{17.797489\du}}
\pgfpathlineto{\pgfpoint{-21.470420\du}{17.372008\du}}
\pgfpathlineto{\pgfpoint{-23.244200\du}{17.372008\du}}
\pgfusepath{fill}
\pgfsetbuttcap
\pgfsetmiterjoin
\pgfsetdash{}{0pt}
\definecolor{dialinecolor}{rgb}{0.000000, 0.705882, 1.000000}
\pgfsetfillcolor{dialinecolor}
\pgfpathmoveto{\pgfpoint{-21.460118\du}{17.366084\du}}
\pgfpathlineto{\pgfpoint{-21.461406\du}{17.351146\du}}
\pgfpathlineto{\pgfpoint{-21.464754\du}{17.335693\du}}
\pgfpathlineto{\pgfpoint{-21.470420\du}{17.320497\du}}
\pgfpathlineto{\pgfpoint{-21.478405\du}{17.305301\du}}
\pgfpathlineto{\pgfpoint{-21.488707\du}{17.290106\du}}
\pgfpathlineto{\pgfpoint{-21.501069\du}{17.275425\du}}
\pgfpathlineto{\pgfpoint{-21.515750\du}{17.261002\du}}
\pgfpathlineto{\pgfpoint{-21.532234\du}{17.246579\du}}
\pgfpathlineto{\pgfpoint{-21.551293\du}{17.232929\du}}
\pgfpathlineto{\pgfpoint{-21.572155\du}{17.219278\du}}
\pgfpathlineto{\pgfpoint{-21.595077\du}{17.206143\du}}
\pgfpathlineto{\pgfpoint{-21.620060\du}{17.192750\du}}
\pgfpathlineto{\pgfpoint{-21.646845\du}{17.180645\du}}
\pgfpathlineto{\pgfpoint{-21.675434\du}{17.168540\du}}
\pgfpathlineto{\pgfpoint{-21.705825\du}{17.157207\du}}
\pgfpathlineto{\pgfpoint{-21.737505\du}{17.146648\du}}
\pgfpathlineto{\pgfpoint{-21.770987\du}{17.136346\du}}
\pgfpathlineto{\pgfpoint{-21.806272\du}{17.126558\du}}
\pgfpathlineto{\pgfpoint{-21.842587\du}{17.117544\du}}
\pgfpathlineto{\pgfpoint{-21.880190\du}{17.109045\du}}
\pgfpathlineto{\pgfpoint{-21.919081\du}{17.101061\du}}
\pgfpathlineto{\pgfpoint{-21.959259\du}{17.094364\du}}
\pgfpathlineto{\pgfpoint{-22.000211\du}{17.087410\du}}
\pgfpathlineto{\pgfpoint{-22.042449\du}{17.081744\du}}
\pgfpathlineto{\pgfpoint{-22.084946\du}{17.076850\du}}
\pgfpathlineto{\pgfpoint{-22.128730\du}{17.072730\du}}
\pgfpathlineto{\pgfpoint{-22.172772\du}{17.069639\du}}
\pgfpathlineto{\pgfpoint{-22.217072\du}{17.066806\du}}
\pgfpathlineto{\pgfpoint{-22.261886\du}{17.064745\du}}
\pgfpathlineto{\pgfpoint{-22.306958\du}{17.063200\du}}
\pgfpathlineto{\pgfpoint{-22.352545\du}{17.063200\du}}
\pgfpathlineto{\pgfpoint{-22.352545\du}{17.063200\du}}
\pgfpathlineto{\pgfpoint{-22.397618\du}{17.063200\du}}
\pgfpathlineto{\pgfpoint{-22.442432\du}{17.064745\du}}
\pgfpathlineto{\pgfpoint{-22.487247\du}{17.066806\du}}
\pgfpathlineto{\pgfpoint{-22.531546\du}{17.069639\du}}
\pgfpathlineto{\pgfpoint{-22.575588\du}{17.072730\du}}
\pgfpathlineto{\pgfpoint{-22.619372\du}{17.076850\du}}
\pgfpathlineto{\pgfpoint{-22.661869\du}{17.081744\du}}
\pgfpathlineto{\pgfpoint{-22.704108\du}{17.087410\du}}
\pgfpathlineto{\pgfpoint{-22.745059\du}{17.094364\du}}
\pgfpathlineto{\pgfpoint{-22.785237\du}{17.101061\du}}
\pgfpathlineto{\pgfpoint{-22.824386\du}{17.109045\du}}
\pgfpathlineto{\pgfpoint{-22.861731\du}{17.117544\du}}
\pgfpathlineto{\pgfpoint{-22.898046\du}{17.126558\du}}
\pgfpathlineto{\pgfpoint{-22.933331\du}{17.136346\du}}
\pgfpathlineto{\pgfpoint{-22.966814\du}{17.146648\du}}
\pgfpathlineto{\pgfpoint{-22.998493\du}{17.157207\du}}
\pgfpathlineto{\pgfpoint{-23.028884\du}{17.168540\du}}
\pgfpathlineto{\pgfpoint{-23.057473\du}{17.180645\du}}
\pgfpathlineto{\pgfpoint{-23.084258\du}{17.192750\du}}
\pgfpathlineto{\pgfpoint{-23.109499\du}{17.206143\du}}
\pgfpathlineto{\pgfpoint{-23.132164\du}{17.219278\du}}
\pgfpathlineto{\pgfpoint{-23.153283\du}{17.232929\du}}
\pgfpathlineto{\pgfpoint{-23.172085\du}{17.246579\du}}
\pgfpathlineto{\pgfpoint{-23.188568\du}{17.261002\du}}
\pgfpathlineto{\pgfpoint{-23.203249\du}{17.275425\du}}
\pgfpathlineto{\pgfpoint{-23.215611\du}{17.290106\du}}
\pgfpathlineto{\pgfpoint{-23.225914\du}{17.305301\du}}
\pgfpathlineto{\pgfpoint{-23.233898\du}{17.320497\du}}
\pgfpathlineto{\pgfpoint{-23.239822\du}{17.335693\du}}
\pgfpathlineto{\pgfpoint{-23.242912\du}{17.351146\du}}
\pgfpathlineto{\pgfpoint{-23.244200\du}{17.366084\du}}
\pgfpathlineto{\pgfpoint{-23.244200\du}{17.366084\du}}
\pgfpathlineto{\pgfpoint{-23.242912\du}{17.381795\du}}
\pgfpathlineto{\pgfpoint{-23.239822\du}{17.397249\du}}
\pgfpathlineto{\pgfpoint{-23.233898\du}{17.412444\du}}
\pgfpathlineto{\pgfpoint{-23.225914\du}{17.427382\du}}
\pgfpathlineto{\pgfpoint{-23.215611\du}{17.442578\du}}
\pgfpathlineto{\pgfpoint{-23.203249\du}{17.457001\du}}
\pgfpathlineto{\pgfpoint{-23.188568\du}{17.471939\du}}
\pgfpathlineto{\pgfpoint{-23.172085\du}{17.486105\du}}
\pgfpathlineto{\pgfpoint{-23.153283\du}{17.500013\du}}
\pgfpathlineto{\pgfpoint{-23.132164\du}{17.513663\du}}
\pgfpathlineto{\pgfpoint{-23.109499\du}{17.526799\du}}
\pgfpathlineto{\pgfpoint{-23.084258\du}{17.539934\du}}
\pgfpathlineto{\pgfpoint{-23.057473\du}{17.552296\du}}
\pgfpathlineto{\pgfpoint{-23.028884\du}{17.563629\du}}
\pgfpathlineto{\pgfpoint{-22.998493\du}{17.575476\du}}
\pgfpathlineto{\pgfpoint{-22.966814\du}{17.586294\du}}
\pgfpathlineto{\pgfpoint{-22.933331\du}{17.596596\du}}
\pgfpathlineto{\pgfpoint{-22.898046\du}{17.606383\du}}
\pgfpathlineto{\pgfpoint{-22.861731\du}{17.615397\du}}
\pgfpathlineto{\pgfpoint{-22.824386\du}{17.623639\du}}
\pgfpathlineto{\pgfpoint{-22.785237\du}{17.631623\du}}
\pgfpathlineto{\pgfpoint{-22.745059\du}{17.638577\du}}
\pgfpathlineto{\pgfpoint{-22.704108\du}{17.645274\du}}
\pgfpathlineto{\pgfpoint{-22.661869\du}{17.650940\du}}
\pgfpathlineto{\pgfpoint{-22.619372\du}{17.655834\du}}
\pgfpathlineto{\pgfpoint{-22.575588\du}{17.659954\du}}
\pgfpathlineto{\pgfpoint{-22.531546\du}{17.663045\du}}
\pgfpathlineto{\pgfpoint{-22.487247\du}{17.665878\du}}
\pgfpathlineto{\pgfpoint{-22.442432\du}{17.668196\du}}
\pgfpathlineto{\pgfpoint{-22.397618\du}{17.669484\du}}
\pgfpathlineto{\pgfpoint{-22.352545\du}{17.669484\du}}
\pgfpathlineto{\pgfpoint{-22.352545\du}{17.669484\du}}
\pgfpathlineto{\pgfpoint{-22.306958\du}{17.669484\du}}
\pgfpathlineto{\pgfpoint{-22.261886\du}{17.668196\du}}
\pgfpathlineto{\pgfpoint{-22.217072\du}{17.665878\du}}
\pgfpathlineto{\pgfpoint{-22.172772\du}{17.663045\du}}
\pgfpathlineto{\pgfpoint{-22.128730\du}{17.659954\du}}
\pgfpathlineto{\pgfpoint{-22.084946\du}{17.655834\du}}
\pgfpathlineto{\pgfpoint{-22.042449\du}{17.650940\du}}
\pgfpathlineto{\pgfpoint{-22.000211\du}{17.645274\du}}
\pgfpathlineto{\pgfpoint{-21.959259\du}{17.638577\du}}
\pgfpathlineto{\pgfpoint{-21.919081\du}{17.631623\du}}
\pgfpathlineto{\pgfpoint{-21.880190\du}{17.623639\du}}
\pgfpathlineto{\pgfpoint{-21.842587\du}{17.615397\du}}
\pgfpathlineto{\pgfpoint{-21.806272\du}{17.606383\du}}
\pgfpathlineto{\pgfpoint{-21.770987\du}{17.596596\du}}
\pgfpathlineto{\pgfpoint{-21.737505\du}{17.586294\du}}
\pgfpathlineto{\pgfpoint{-21.705825\du}{17.575476\du}}
\pgfpathlineto{\pgfpoint{-21.675434\du}{17.563629\du}}
\pgfpathlineto{\pgfpoint{-21.646845\du}{17.552296\du}}
\pgfpathlineto{\pgfpoint{-21.620060\du}{17.539934\du}}
\pgfpathlineto{\pgfpoint{-21.595077\du}{17.526799\du}}
\pgfpathlineto{\pgfpoint{-21.572155\du}{17.513663\du}}
\pgfpathlineto{\pgfpoint{-21.551293\du}{17.500013\du}}
\pgfpathlineto{\pgfpoint{-21.532234\du}{17.486105\du}}
\pgfpathlineto{\pgfpoint{-21.515750\du}{17.471939\du}}
\pgfpathlineto{\pgfpoint{-21.501069\du}{17.457001\du}}
\pgfpathlineto{\pgfpoint{-21.488707\du}{17.442578\du}}
\pgfpathlineto{\pgfpoint{-21.478405\du}{17.427382\du}}
\pgfpathlineto{\pgfpoint{-21.470420\du}{17.412444\du}}
\pgfpathlineto{\pgfpoint{-21.464754\du}{17.397249\du}}
\pgfpathlineto{\pgfpoint{-21.461406\du}{17.381795\du}}
\pgfpathlineto{\pgfpoint{-21.460118\du}{17.366084\du}}
\pgfusepath{fill}
\pgfsetbuttcap
\pgfsetmiterjoin
\pgfsetdash{}{0pt}
\definecolor{dialinecolor}{rgb}{0.666667, 0.901961, 1.000000}
\pgfsetstrokecolor{dialinecolor}
\pgfpathmoveto{\pgfpoint{-21.470420\du}{17.361706\du}}
\pgfpathlineto{\pgfpoint{-21.471708\du}{17.346253\du}}
\pgfpathlineto{\pgfpoint{-21.475056\du}{17.331057\du}}
\pgfpathlineto{\pgfpoint{-21.480723\du}{17.316634\du}}
\pgfpathlineto{\pgfpoint{-21.488449\du}{17.301438\du}}
\pgfpathlineto{\pgfpoint{-21.498751\du}{17.286758\du}}
\pgfpathlineto{\pgfpoint{-21.511372\du}{17.272334\du}}
\pgfpathlineto{\pgfpoint{-21.525280\du}{17.258427\du}}
\pgfpathlineto{\pgfpoint{-21.542278\du}{17.244003\du}}
\pgfpathlineto{\pgfpoint{-21.561080\du}{17.230611\du}}
\pgfpathlineto{\pgfpoint{-21.581942\du}{17.216703\du}}
\pgfpathlineto{\pgfpoint{-21.604606\du}{17.203825\du}}
\pgfpathlineto{\pgfpoint{-21.629589\du}{17.190947\du}}
\pgfpathlineto{\pgfpoint{-21.655860\du}{17.178842\du}}
\pgfpathlineto{\pgfpoint{-21.684191\du}{17.167510\du}}
\pgfpathlineto{\pgfpoint{-21.714582\du}{17.155662\du}}
\pgfpathlineto{\pgfpoint{-21.746519\du}{17.145360\du}}
\pgfpathlineto{\pgfpoint{-21.779744\du}{17.135315\du}}
\pgfpathlineto{\pgfpoint{-21.814514\du}{17.125786\du}}
\pgfpathlineto{\pgfpoint{-21.850314\du}{17.117029\du}}
\pgfpathlineto{\pgfpoint{-21.888174\du}{17.108272\du}}
\pgfpathlineto{\pgfpoint{-21.927065\du}{17.100803\du}}
\pgfpathlineto{\pgfpoint{-21.966471\du}{17.093849\du}}
\pgfpathlineto{\pgfpoint{-22.007422\du}{17.087410\du}}
\pgfpathlineto{\pgfpoint{-22.049146\du}{17.081744\du}}
\pgfpathlineto{\pgfpoint{-22.091642\du}{17.076850\du}}
\pgfpathlineto{\pgfpoint{-22.135169\du}{17.073245\du}}
\pgfpathlineto{\pgfpoint{-22.178696\du}{17.069639\du}}
\pgfpathlineto{\pgfpoint{-22.222738\du}{17.067063\du}}
\pgfpathlineto{\pgfpoint{-22.267552\du}{17.064745\du}}
\pgfpathlineto{\pgfpoint{-22.312367\du}{17.063715\du}}
\pgfpathlineto{\pgfpoint{-22.357439\du}{17.063200\du}}
\pgfpathlineto{\pgfpoint{-22.357439\du}{17.063200\du}}
\pgfpathlineto{\pgfpoint{-22.402253\du}{17.063715\du}}
\pgfpathlineto{\pgfpoint{-22.447326\du}{17.064745\du}}
\pgfpathlineto{\pgfpoint{-22.491883\du}{17.067063\du}}
\pgfpathlineto{\pgfpoint{-22.535924\du}{17.069639\du}}
\pgfpathlineto{\pgfpoint{-22.579709\du}{17.073245\du}}
\pgfpathlineto{\pgfpoint{-22.622978\du}{17.076850\du}}
\pgfpathlineto{\pgfpoint{-22.665474\du}{17.081744\du}}
\pgfpathlineto{\pgfpoint{-22.707456\du}{17.087410\du}}
\pgfpathlineto{\pgfpoint{-22.748150\du}{17.093849\du}}
\pgfpathlineto{\pgfpoint{-22.787555\du}{17.100803\du}}
\pgfpathlineto{\pgfpoint{-22.826704\du}{17.108272\du}}
\pgfpathlineto{\pgfpoint{-22.864049\du}{17.117029\du}}
\pgfpathlineto{\pgfpoint{-22.899849\du}{17.125786\du}}
\pgfpathlineto{\pgfpoint{-22.934877\du}{17.135315\du}}
\pgfpathlineto{\pgfpoint{-22.968616\du}{17.145360\du}}
\pgfpathlineto{\pgfpoint{-23.000553\du}{17.155662\du}}
\pgfpathlineto{\pgfpoint{-23.030430\du}{17.167510\du}}
\pgfpathlineto{\pgfpoint{-23.059018\du}{17.178842\du}}
\pgfpathlineto{\pgfpoint{-23.085546\du}{17.190947\du}}
\pgfpathlineto{\pgfpoint{-23.110014\du}{17.203825\du}}
\pgfpathlineto{\pgfpoint{-23.132421\du}{17.216703\du}}
\pgfpathlineto{\pgfpoint{-23.153541\du}{17.230611\du}}
\pgfpathlineto{\pgfpoint{-23.172600\du}{17.244003\du}}
\pgfpathlineto{\pgfpoint{-23.189083\du}{17.258427\du}}
\pgfpathlineto{\pgfpoint{-23.203506\du}{17.272334\du}}
\pgfpathlineto{\pgfpoint{-23.216384\du}{17.286758\du}}
\pgfpathlineto{\pgfpoint{-23.225914\du}{17.301438\du}}
\pgfpathlineto{\pgfpoint{-23.234155\du}{17.316634\du}}
\pgfpathlineto{\pgfpoint{-23.239822\du}{17.331057\du}}
\pgfpathlineto{\pgfpoint{-23.242912\du}{17.346253\du}}
\pgfpathlineto{\pgfpoint{-23.244200\du}{17.361706\du}}
\pgfpathlineto{\pgfpoint{-23.244200\du}{17.361706\du}}
\pgfpathlineto{\pgfpoint{-23.242912\du}{17.376644\du}}
\pgfpathlineto{\pgfpoint{-23.239822\du}{17.391840\du}}
\pgfpathlineto{\pgfpoint{-23.234155\du}{17.406263\du}}
\pgfpathlineto{\pgfpoint{-23.225914\du}{17.421201\du}}
\pgfpathlineto{\pgfpoint{-23.216384\du}{17.436139\du}}
\pgfpathlineto{\pgfpoint{-23.203506\du}{17.450562\du}}
\pgfpathlineto{\pgfpoint{-23.189083\du}{17.464985\du}}
\pgfpathlineto{\pgfpoint{-23.172600\du}{17.479151\du}}
\pgfpathlineto{\pgfpoint{-23.153541\du}{17.492801\du}}
\pgfpathlineto{\pgfpoint{-23.132421\du}{17.506194\du}}
\pgfpathlineto{\pgfpoint{-23.110014\du}{17.519329\du}}
\pgfpathlineto{\pgfpoint{-23.085546\du}{17.531435\du}}
\pgfpathlineto{\pgfpoint{-23.059018\du}{17.544055\du}}
\pgfpathlineto{\pgfpoint{-23.030430\du}{17.555645\du}}
\pgfpathlineto{\pgfpoint{-23.000553\du}{17.566977\du}}
\pgfpathlineto{\pgfpoint{-22.968616\du}{17.577537\du}}
\pgfpathlineto{\pgfpoint{-22.934877\du}{17.587581\du}}
\pgfpathlineto{\pgfpoint{-22.899849\du}{17.597111\du}}
\pgfpathlineto{\pgfpoint{-22.864049\du}{17.605868\du}}
\pgfpathlineto{\pgfpoint{-22.826704\du}{17.614367\du}}
\pgfpathlineto{\pgfpoint{-22.787555\du}{17.622351\du}}
\pgfpathlineto{\pgfpoint{-22.748150\du}{17.628790\du}}
\pgfpathlineto{\pgfpoint{-22.707456\du}{17.635487\du}}
\pgfpathlineto{\pgfpoint{-22.665474\du}{17.640895\du}}
\pgfpathlineto{\pgfpoint{-22.622978\du}{17.645789\du}}
\pgfpathlineto{\pgfpoint{-22.579709\du}{17.650167\du}}
\pgfpathlineto{\pgfpoint{-22.535924\du}{17.653773\du}}
\pgfpathlineto{\pgfpoint{-22.491883\du}{17.655834\du}}
\pgfpathlineto{\pgfpoint{-22.447326\du}{17.657894\du}}
\pgfpathlineto{\pgfpoint{-22.402253\du}{17.659182\du}}
\pgfpathlineto{\pgfpoint{-22.357439\du}{17.659439\du}}
\pgfpathlineto{\pgfpoint{-22.357439\du}{17.659439\du}}
\pgfpathlineto{\pgfpoint{-22.312367\du}{17.659182\du}}
\pgfpathlineto{\pgfpoint{-22.267552\du}{17.657894\du}}
\pgfpathlineto{\pgfpoint{-22.222738\du}{17.655834\du}}
\pgfpathlineto{\pgfpoint{-22.178696\du}{17.653773\du}}
\pgfpathlineto{\pgfpoint{-22.135169\du}{17.650167\du}}
\pgfpathlineto{\pgfpoint{-22.091642\du}{17.645789\du}}
\pgfpathlineto{\pgfpoint{-22.049146\du}{17.640895\du}}
\pgfpathlineto{\pgfpoint{-22.007422\du}{17.635487\du}}
\pgfpathlineto{\pgfpoint{-21.966471\du}{17.628790\du}}
\pgfpathlineto{\pgfpoint{-21.927065\du}{17.622351\du}}
\pgfpathlineto{\pgfpoint{-21.888174\du}{17.614367\du}}
\pgfpathlineto{\pgfpoint{-21.850314\du}{17.605868\du}}
\pgfpathlineto{\pgfpoint{-21.814514\du}{17.597111\du}}
\pgfpathlineto{\pgfpoint{-21.779744\du}{17.587581\du}}
\pgfpathlineto{\pgfpoint{-21.746519\du}{17.577537\du}}
\pgfpathlineto{\pgfpoint{-21.714582\du}{17.566977\du}}
\pgfpathlineto{\pgfpoint{-21.684191\du}{17.555645\du}}
\pgfpathlineto{\pgfpoint{-21.655860\du}{17.544055\du}}
\pgfpathlineto{\pgfpoint{-21.629589\du}{17.531435\du}}
\pgfpathlineto{\pgfpoint{-21.604606\du}{17.519329\du}}
\pgfpathlineto{\pgfpoint{-21.581942\du}{17.506194\du}}
\pgfpathlineto{\pgfpoint{-21.561080\du}{17.492801\du}}
\pgfpathlineto{\pgfpoint{-21.542278\du}{17.479151\du}}
\pgfpathlineto{\pgfpoint{-21.525280\du}{17.464985\du}}
\pgfpathlineto{\pgfpoint{-21.511372\du}{17.450562\du}}
\pgfpathlineto{\pgfpoint{-21.498751\du}{17.436139\du}}
\pgfpathlineto{\pgfpoint{-21.488449\du}{17.421201\du}}
\pgfpathlineto{\pgfpoint{-21.480723\du}{17.406263\du}}
\pgfpathlineto{\pgfpoint{-21.475056\du}{17.391840\du}}
\pgfpathlineto{\pgfpoint{-21.471708\du}{17.376644\du}}
\pgfpathlineto{\pgfpoint{-21.470420\du}{17.361706\du}}
\pgfusepath{stroke}
\pgfsetbuttcap
\pgfsetmiterjoin
\pgfsetdash{}{0pt}
\definecolor{dialinecolor}{rgb}{0.000000, 0.000000, 0.000000}
\pgfsetfillcolor{dialinecolor}
\pgfpathmoveto{\pgfpoint{-22.336062\du}{17.297317\du}}
\pgfpathlineto{\pgfpoint{-22.205997\du}{17.340071\du}}
\pgfpathlineto{\pgfpoint{-21.892038\du}{17.212067\du}}
\pgfpathlineto{\pgfpoint{-21.751670\du}{17.254821\du}}
\pgfpathlineto{\pgfpoint{-21.827906\du}{17.148451\du}}
\pgfpathlineto{\pgfpoint{-22.194922\du}{17.148451\du}}
\pgfpathlineto{\pgfpoint{-22.043995\du}{17.180645\du}}
\pgfpathlineto{\pgfpoint{-22.336062\du}{17.297317\du}}
\pgfusepath{fill}
\pgfsetbuttcap
\pgfsetmiterjoin
\pgfsetdash{}{0pt}
\definecolor{dialinecolor}{rgb}{0.000000, 0.000000, 0.000000}
\pgfsetfillcolor{dialinecolor}
\pgfpathmoveto{\pgfpoint{-22.390406\du}{17.414762\du}}
\pgfpathlineto{\pgfpoint{-22.519441\du}{17.372008\du}}
\pgfpathlineto{\pgfpoint{-22.822583\du}{17.500013\du}}
\pgfpathlineto{\pgfpoint{-22.973768\du}{17.457001\du}}
\pgfpathlineto{\pgfpoint{-22.898046\du}{17.574189\du}}
\pgfpathlineto{\pgfpoint{-22.519441\du}{17.574189\du}}
\pgfpathlineto{\pgfpoint{-22.681958\du}{17.531435\du}}
\pgfpathlineto{\pgfpoint{-22.390406\du}{17.414762\du}}
\pgfusepath{fill}
\pgfsetbuttcap
\pgfsetmiterjoin
\pgfsetdash{}{0pt}
\definecolor{dialinecolor}{rgb}{0.000000, 0.000000, 0.000000}
\pgfsetfillcolor{dialinecolor}
\pgfpathmoveto{\pgfpoint{-22.941573\du}{17.180645\du}}
\pgfpathlineto{\pgfpoint{-22.811766\du}{17.137891\du}}
\pgfpathlineto{\pgfpoint{-22.498321\du}{17.254821\du}}
\pgfpathlineto{\pgfpoint{-22.357439\du}{17.223142\du}}
\pgfpathlineto{\pgfpoint{-22.432903\du}{17.329254\du}}
\pgfpathlineto{\pgfpoint{-22.800948\du}{17.329254\du}}
\pgfpathlineto{\pgfpoint{-22.649764\du}{17.297317\du}}
\pgfpathlineto{\pgfpoint{-22.941573\du}{17.180645\du}}
\pgfusepath{fill}
\pgfsetbuttcap
\pgfsetmiterjoin
\pgfsetdash{}{0pt}
\definecolor{dialinecolor}{rgb}{0.000000, 0.000000, 0.000000}
\pgfsetfillcolor{dialinecolor}
\pgfpathmoveto{\pgfpoint{-21.773562\du}{17.542252\du}}
\pgfpathlineto{\pgfpoint{-21.903112\du}{17.584748\du}}
\pgfpathlineto{\pgfpoint{-22.205997\du}{17.457001\du}}
\pgfpathlineto{\pgfpoint{-22.357439\du}{17.500013\du}}
\pgfpathlineto{\pgfpoint{-22.281718\du}{17.393385\du}}
\pgfpathlineto{\pgfpoint{-21.903112\du}{17.393385\du}}
\pgfpathlineto{\pgfpoint{-22.065629\du}{17.424807\du}}
\pgfpathlineto{\pgfpoint{-21.773562\du}{17.542252\du}}
\pgfusepath{fill}
\pgfsetbuttcap
\pgfsetmiterjoin
\pgfsetdash{}{0pt}
\definecolor{dialinecolor}{rgb}{1.000000, 1.000000, 1.000000}
\pgfsetfillcolor{dialinecolor}
\pgfpathmoveto{\pgfpoint{-22.325245\du}{17.307877\du}}
\pgfpathlineto{\pgfpoint{-22.194922\du}{17.350374\du}}
\pgfpathlineto{\pgfpoint{-21.881478\du}{17.223142\du}}
\pgfpathlineto{\pgfpoint{-21.740853\du}{17.265896\du}}
\pgfpathlineto{\pgfpoint{-21.816574\du}{17.159010\du}}
\pgfpathlineto{\pgfpoint{-22.184620\du}{17.159010\du}}
\pgfpathlineto{\pgfpoint{-22.032920\du}{17.190947\du}}
\pgfpathlineto{\pgfpoint{-22.325245\du}{17.307877\du}}
\pgfusepath{fill}
\pgfsetbuttcap
\pgfsetmiterjoin
\pgfsetdash{}{0pt}
\definecolor{dialinecolor}{rgb}{1.000000, 1.000000, 1.000000}
\pgfsetfillcolor{dialinecolor}
\pgfpathmoveto{\pgfpoint{-22.378816\du}{17.424807\du}}
\pgfpathlineto{\pgfpoint{-22.509139\du}{17.382568\du}}
\pgfpathlineto{\pgfpoint{-22.811766\du}{17.510315\du}}
\pgfpathlineto{\pgfpoint{-22.963208\du}{17.467561\du}}
\pgfpathlineto{\pgfpoint{-22.887487\du}{17.584748\du}}
\pgfpathlineto{\pgfpoint{-22.509139\du}{17.584748\du}}
\pgfpathlineto{\pgfpoint{-22.670883\du}{17.542252\du}}
\pgfpathlineto{\pgfpoint{-22.378816\du}{17.424807\du}}
\pgfusepath{fill}
\pgfsetbuttcap
\pgfsetmiterjoin
\pgfsetdash{}{0pt}
\definecolor{dialinecolor}{rgb}{1.000000, 1.000000, 1.000000}
\pgfsetfillcolor{dialinecolor}
\pgfpathmoveto{\pgfpoint{-22.930498\du}{17.190947\du}}
\pgfpathlineto{\pgfpoint{-22.800948\du}{17.148451\du}}
\pgfpathlineto{\pgfpoint{-22.487247\du}{17.265896\du}}
\pgfpathlineto{\pgfpoint{-22.346622\du}{17.233444\du}}
\pgfpathlineto{\pgfpoint{-22.422343\du}{17.340071\du}}
\pgfpathlineto{\pgfpoint{-22.790131\du}{17.340071\du}}
\pgfpathlineto{\pgfpoint{-22.638431\du}{17.307877\du}}
\pgfpathlineto{\pgfpoint{-22.930498\du}{17.190947\du}}
\pgfusepath{fill}
\pgfsetbuttcap
\pgfsetmiterjoin
\pgfsetdash{}{0pt}
\definecolor{dialinecolor}{rgb}{1.000000, 1.000000, 1.000000}
\pgfsetfillcolor{dialinecolor}
\pgfpathmoveto{\pgfpoint{-21.762487\du}{17.552812\du}}
\pgfpathlineto{\pgfpoint{-21.892038\du}{17.595566\du}}
\pgfpathlineto{\pgfpoint{-22.194922\du}{17.467561\du}}
\pgfpathlineto{\pgfpoint{-22.346622\du}{17.510315\du}}
\pgfpathlineto{\pgfpoint{-22.271158\du}{17.404203\du}}
\pgfpathlineto{\pgfpoint{-21.892038\du}{17.404203\du}}
\pgfpathlineto{\pgfpoint{-22.054297\du}{17.435882\du}}
\pgfpathlineto{\pgfpoint{-21.762487\du}{17.552812\du}}
\pgfusepath{fill}
\pgfsetbuttcap
\pgfsetmiterjoin
\pgfsetdash{}{0pt}
\definecolor{dialinecolor}{rgb}{0.666667, 0.901961, 1.000000}
\pgfsetstrokecolor{dialinecolor}
\pgfpathmoveto{\pgfpoint{-23.244200\du}{17.361706\du}}
\pgfpathlineto{\pgfpoint{-23.244200\du}{17.786671\du}}
\pgfusepath{stroke}
\pgfsetbuttcap
\pgfsetmiterjoin
\pgfsetdash{}{0pt}
\definecolor{dialinecolor}{rgb}{0.666667, 0.901961, 1.000000}
\pgfsetstrokecolor{dialinecolor}
\pgfpathmoveto{\pgfpoint{-21.470420\du}{17.361706\du}}
\pgfpathlineto{\pgfpoint{-21.470420\du}{17.786671\du}}
\pgfusepath{stroke}
\pgfsetbuttcap
\pgfsetmiterjoin
\pgfsetdash{}{0pt}
\definecolor{dialinecolor}{rgb}{0.000000, 0.000000, 0.000000}
\pgfsetfillcolor{dialinecolor}
\pgfpathmoveto{\pgfpoint{-22.681958\du}{17.733873\du}}
\pgfpathlineto{\pgfpoint{-22.487247\du}{17.733873\du}}
\pgfpathlineto{\pgfpoint{-22.368256\du}{17.829683\du}}
\pgfpathlineto{\pgfpoint{-22.249524\du}{17.733873\du}}
\pgfpathlineto{\pgfpoint{-22.065629\du}{17.733873\du}}
\pgfpathlineto{\pgfpoint{-22.065629\du}{17.691376\du}}
\pgfpathlineto{\pgfpoint{-21.935564\du}{17.754735\du}}
\pgfpathlineto{\pgfpoint{-22.065629\du}{17.818866\du}}
\pgfpathlineto{\pgfpoint{-22.065629\du}{17.776369\du}}
\pgfpathlineto{\pgfpoint{-22.227631\du}{17.776369\du}}
\pgfpathlineto{\pgfpoint{-22.336062\du}{17.850802\du}}
\pgfpathlineto{\pgfpoint{-22.227631\du}{17.935538\du}}
\pgfpathlineto{\pgfpoint{-22.065629\du}{17.935538\du}}
\pgfpathlineto{\pgfpoint{-22.065629\du}{17.882997\du}}
\pgfpathlineto{\pgfpoint{-21.935564\du}{17.957173\du}}
\pgfpathlineto{\pgfpoint{-22.065629\du}{18.021304\du}}
\pgfpathlineto{\pgfpoint{-22.065629\du}{17.978292\du}}
\pgfpathlineto{\pgfpoint{-22.249524\du}{17.978292\du}}
\pgfpathlineto{\pgfpoint{-22.368256\du}{17.882997\du}}
\pgfpathlineto{\pgfpoint{-22.487247\du}{17.978292\du}}
\pgfpathlineto{\pgfpoint{-22.681958\du}{17.978292\du}}
\pgfpathlineto{\pgfpoint{-22.681958\du}{18.021304\du}}
\pgfpathlineto{\pgfpoint{-22.811766\du}{17.957173\du}}
\pgfpathlineto{\pgfpoint{-22.681958\du}{17.882997\du}}
\pgfpathlineto{\pgfpoint{-22.681958\du}{17.935538\du}}
\pgfpathlineto{\pgfpoint{-22.509139\du}{17.935538\du}}
\pgfpathlineto{\pgfpoint{-22.400708\du}{17.861362\du}}
\pgfpathlineto{\pgfpoint{-22.509139\du}{17.776369\du}}
\pgfpathlineto{\pgfpoint{-22.681958\du}{17.776369\du}}
\pgfpathlineto{\pgfpoint{-22.681958\du}{17.818866\du}}
\pgfpathlineto{\pgfpoint{-22.811766\du}{17.754735\du}}
\pgfpathlineto{\pgfpoint{-22.681958\du}{17.691376\du}}
\pgfpathlineto{\pgfpoint{-22.681958\du}{17.733873\du}}
\pgfusepath{fill}
\pgfsetbuttcap
\pgfsetmiterjoin
\pgfsetdash{}{0pt}
\definecolor{dialinecolor}{rgb}{1.000000, 1.000000, 1.000000}
\pgfsetfillcolor{dialinecolor}
\pgfpathmoveto{\pgfpoint{-22.670883\du}{17.744175\du}}
\pgfpathlineto{\pgfpoint{-22.476429\du}{17.744175\du}}
\pgfpathlineto{\pgfpoint{-22.357439\du}{17.839728\du}}
\pgfpathlineto{\pgfpoint{-22.238706\du}{17.744175\du}}
\pgfpathlineto{\pgfpoint{-22.054297\du}{17.744175\du}}
\pgfpathlineto{\pgfpoint{-22.054297\du}{17.701678\du}}
\pgfpathlineto{\pgfpoint{-21.925005\du}{17.765809\du}}
\pgfpathlineto{\pgfpoint{-22.054297\du}{17.829683\du}}
\pgfpathlineto{\pgfpoint{-22.054297\du}{17.787186\du}}
\pgfpathlineto{\pgfpoint{-22.216814\du}{17.787186\du}}
\pgfpathlineto{\pgfpoint{-22.325245\du}{17.861362\du}}
\pgfpathlineto{\pgfpoint{-22.216814\du}{17.946355\du}}
\pgfpathlineto{\pgfpoint{-22.054297\du}{17.946355\du}}
\pgfpathlineto{\pgfpoint{-22.054297\du}{17.893299\du}}
\pgfpathlineto{\pgfpoint{-21.925005\du}{17.967732\du}}
\pgfpathlineto{\pgfpoint{-22.054297\du}{18.031606\du}}
\pgfpathlineto{\pgfpoint{-22.054297\du}{17.988852\du}}
\pgfpathlineto{\pgfpoint{-22.238706\du}{17.988852\du}}
\pgfpathlineto{\pgfpoint{-22.357439\du}{17.893299\du}}
\pgfpathlineto{\pgfpoint{-22.476429\du}{17.988852\du}}
\pgfpathlineto{\pgfpoint{-22.670883\du}{17.988852\du}}
\pgfpathlineto{\pgfpoint{-22.670883\du}{18.031606\du}}
\pgfpathlineto{\pgfpoint{-22.800948\du}{17.967732\du}}
\pgfpathlineto{\pgfpoint{-22.670883\du}{17.893299\du}}
\pgfpathlineto{\pgfpoint{-22.670883\du}{17.946355\du}}
\pgfpathlineto{\pgfpoint{-22.498321\du}{17.946355\du}}
\pgfpathlineto{\pgfpoint{-22.390406\du}{17.871922\du}}
\pgfpathlineto{\pgfpoint{-22.498321\du}{17.787186\du}}
\pgfpathlineto{\pgfpoint{-22.670883\du}{17.787186\du}}
\pgfpathlineto{\pgfpoint{-22.670883\du}{17.829683\du}}
\pgfpathlineto{\pgfpoint{-22.800948\du}{17.765809\du}}
\pgfpathlineto{\pgfpoint{-22.670883\du}{17.701678\du}}
\pgfpathlineto{\pgfpoint{-22.670883\du}{17.744175\du}}
\pgfusepath{fill}
\pgfsetlinewidth{0.100000\du}
\pgfsetdash{}{0pt}
\pgfsetdash{}{0pt}
\pgfsetbuttcap
\pgfsetmiterjoin
\pgfsetlinewidth{0.001000\du}
\pgfsetbuttcap
\pgfsetmiterjoin
\pgfsetdash{}{0pt}
\definecolor{dialinecolor}{rgb}{0.000000, 0.470588, 0.666667}
\pgfsetfillcolor{dialinecolor}
\pgfpathmoveto{\pgfpoint{-16.750016\du}{14.856661\du}}
\pgfpathlineto{\pgfpoint{-16.751776\du}{14.838184\du}}
\pgfpathlineto{\pgfpoint{-16.755589\du}{14.819707\du}}
\pgfpathlineto{\pgfpoint{-16.762041\du}{14.801523\du}}
\pgfpathlineto{\pgfpoint{-16.771133\du}{14.783632\du}}
\pgfpathlineto{\pgfpoint{-16.783451\du}{14.765449\du}}
\pgfpathlineto{\pgfpoint{-16.797822\du}{14.747851\du}}
\pgfpathlineto{\pgfpoint{-16.815126\du}{14.730547\du}}
\pgfpathlineto{\pgfpoint{-16.834484\du}{14.713536\du}}
\pgfpathlineto{\pgfpoint{-16.856774\du}{14.696819\du}}
\pgfpathlineto{\pgfpoint{-16.881410\du}{14.680688\du}}
\pgfpathlineto{\pgfpoint{-16.908099\du}{14.665144\du}}
\pgfpathlineto{\pgfpoint{-16.937135\du}{14.649306\du}}
\pgfpathlineto{\pgfpoint{-16.968517\du}{14.634348\du}}
\pgfpathlineto{\pgfpoint{-17.001952\du}{14.620271\du}}
\pgfpathlineto{\pgfpoint{-17.037733\du}{14.606779\du}}
\pgfpathlineto{\pgfpoint{-17.075274\du}{14.593875\du}}
\pgfpathlineto{\pgfpoint{-17.114281\du}{14.581263\du}}
\pgfpathlineto{\pgfpoint{-17.155342\du}{14.569825\du}}
\pgfpathlineto{\pgfpoint{-17.197869\du}{14.558973\du}}
\pgfpathlineto{\pgfpoint{-17.242155\du}{14.548708\du}}
\pgfpathlineto{\pgfpoint{-17.287908\du}{14.539323\du}}
\pgfpathlineto{\pgfpoint{-17.335128\du}{14.531404\du}}
\pgfpathlineto{\pgfpoint{-17.383227\du}{14.523192\du}}
\pgfpathlineto{\pgfpoint{-17.431913\du}{14.516740\du}}
\pgfpathlineto{\pgfpoint{-17.482066\du}{14.510580\du}}
\pgfpathlineto{\pgfpoint{-17.533098\du}{14.505888\du}}
\pgfpathlineto{\pgfpoint{-17.584717\du}{14.501489\du}}
\pgfpathlineto{\pgfpoint{-17.636922\du}{14.498556\du}}
\pgfpathlineto{\pgfpoint{-17.689128\du}{14.495623\du}}
\pgfpathlineto{\pgfpoint{-17.741920\du}{14.494743\du}}
\pgfpathlineto{\pgfpoint{-17.794712\du}{14.493863\du}}
\pgfpathlineto{\pgfpoint{-17.794712\du}{14.493863\du}}
\pgfpathlineto{\pgfpoint{-17.847797\du}{14.494743\du}}
\pgfpathlineto{\pgfpoint{-17.900589\du}{14.495623\du}}
\pgfpathlineto{\pgfpoint{-17.952794\du}{14.498556\du}}
\pgfpathlineto{\pgfpoint{-18.005000\du}{14.501489\du}}
\pgfpathlineto{\pgfpoint{-18.056619\du}{14.505888\du}}
\pgfpathlineto{\pgfpoint{-18.107651\du}{14.510580\du}}
\pgfpathlineto{\pgfpoint{-18.157803\du}{14.516740\du}}
\pgfpathlineto{\pgfpoint{-18.206782\du}{14.523192\du}}
\pgfpathlineto{\pgfpoint{-18.254882\du}{14.531404\du}}
\pgfpathlineto{\pgfpoint{-18.301808\du}{14.539323\du}}
\pgfpathlineto{\pgfpoint{-18.347561\du}{14.548708\du}}
\pgfpathlineto{\pgfpoint{-18.391848\du}{14.558973\du}}
\pgfpathlineto{\pgfpoint{-18.434375\du}{14.569825\du}}
\pgfpathlineto{\pgfpoint{-18.475435\du}{14.581263\du}}
\pgfpathlineto{\pgfpoint{-18.514736\du}{14.593875\du}}
\pgfpathlineto{\pgfpoint{-18.551984\du}{14.606779\du}}
\pgfpathlineto{\pgfpoint{-18.587765\du}{14.620271\du}}
\pgfpathlineto{\pgfpoint{-18.621200\du}{14.634348\du}}
\pgfpathlineto{\pgfpoint{-18.652582\du}{14.649306\du}}
\pgfpathlineto{\pgfpoint{-18.681617\du}{14.665144\du}}
\pgfpathlineto{\pgfpoint{-18.708307\du}{14.680688\du}}
\pgfpathlineto{\pgfpoint{-18.732943\du}{14.696819\du}}
\pgfpathlineto{\pgfpoint{-18.755233\du}{14.713536\du}}
\pgfpathlineto{\pgfpoint{-18.774590\du}{14.730547\du}}
\pgfpathlineto{\pgfpoint{-18.791894\du}{14.747851\du}}
\pgfpathlineto{\pgfpoint{-18.806265\du}{14.765449\du}}
\pgfpathlineto{\pgfpoint{-18.818583\du}{14.783632\du}}
\pgfpathlineto{\pgfpoint{-18.827968\du}{14.801523\du}}
\pgfpathlineto{\pgfpoint{-18.834714\du}{14.819707\du}}
\pgfpathlineto{\pgfpoint{-18.838234\du}{14.838184\du}}
\pgfpathlineto{\pgfpoint{-18.839700\du}{14.856661\du}}
\pgfpathlineto{\pgfpoint{-18.839700\du}{14.856661\du}}
\pgfpathlineto{\pgfpoint{-18.838234\du}{14.874845\du}}
\pgfpathlineto{\pgfpoint{-18.834714\du}{14.893323\du}}
\pgfpathlineto{\pgfpoint{-18.827968\du}{14.911506\du}}
\pgfpathlineto{\pgfpoint{-18.818583\du}{14.929397\du}}
\pgfpathlineto{\pgfpoint{-18.806265\du}{14.947581\du}}
\pgfpathlineto{\pgfpoint{-18.791894\du}{14.965178\du}}
\pgfpathlineto{\pgfpoint{-18.774590\du}{14.982482\du}}
\pgfpathlineto{\pgfpoint{-18.755233\du}{14.999493\du}}
\pgfpathlineto{\pgfpoint{-18.732943\du}{15.016211\du}}
\pgfpathlineto{\pgfpoint{-18.708307\du}{15.032341\du}}
\pgfpathlineto{\pgfpoint{-18.681617\du}{15.048472\du}}
\pgfpathlineto{\pgfpoint{-18.652582\du}{15.063723\du}}
\pgfpathlineto{\pgfpoint{-18.621200\du}{15.078681\du}}
\pgfpathlineto{\pgfpoint{-18.587765\du}{15.092759\du}}
\pgfpathlineto{\pgfpoint{-18.551984\du}{15.106837\du}}
\pgfpathlineto{\pgfpoint{-18.514736\du}{15.119742\du}}
\pgfpathlineto{\pgfpoint{-18.475435\du}{15.131766\du}}
\pgfpathlineto{\pgfpoint{-18.434375\du}{15.143498\du}}
\pgfpathlineto{\pgfpoint{-18.391848\du}{15.154056\du}}
\pgfpathlineto{\pgfpoint{-18.347561\du}{15.164322\du}}
\pgfpathlineto{\pgfpoint{-18.301808\du}{15.173707\du}}
\pgfpathlineto{\pgfpoint{-18.254882\du}{15.182212\du}}
\pgfpathlineto{\pgfpoint{-18.206782\du}{15.189838\du}}
\pgfpathlineto{\pgfpoint{-18.157803\du}{15.196877\du}}
\pgfpathlineto{\pgfpoint{-18.107651\du}{15.202449\du}}
\pgfpathlineto{\pgfpoint{-18.056619\du}{15.207728\du}}
\pgfpathlineto{\pgfpoint{-18.005000\du}{15.211541\du}}
\pgfpathlineto{\pgfpoint{-17.952794\du}{15.215061\du}}
\pgfpathlineto{\pgfpoint{-17.900589\du}{15.217407\du}}
\pgfpathlineto{\pgfpoint{-17.847797\du}{15.218873\du}}
\pgfpathlineto{\pgfpoint{-17.794712\du}{15.219167\du}}
\pgfpathlineto{\pgfpoint{-17.794712\du}{15.219167\du}}
\pgfpathlineto{\pgfpoint{-17.741920\du}{15.218873\du}}
\pgfpathlineto{\pgfpoint{-17.689128\du}{15.217407\du}}
\pgfpathlineto{\pgfpoint{-17.636922\du}{15.215061\du}}
\pgfpathlineto{\pgfpoint{-17.584717\du}{15.211541\du}}
\pgfpathlineto{\pgfpoint{-17.533098\du}{15.207728\du}}
\pgfpathlineto{\pgfpoint{-17.482066\du}{15.202449\du}}
\pgfpathlineto{\pgfpoint{-17.431913\du}{15.196877\du}}
\pgfpathlineto{\pgfpoint{-17.383227\du}{15.189838\du}}
\pgfpathlineto{\pgfpoint{-17.335128\du}{15.182212\du}}
\pgfpathlineto{\pgfpoint{-17.287908\du}{15.173707\du}}
\pgfpathlineto{\pgfpoint{-17.242155\du}{15.164322\du}}
\pgfpathlineto{\pgfpoint{-17.197869\du}{15.154056\du}}
\pgfpathlineto{\pgfpoint{-17.155342\du}{15.143498\du}}
\pgfpathlineto{\pgfpoint{-17.114281\du}{15.131766\du}}
\pgfpathlineto{\pgfpoint{-17.075274\du}{15.119742\du}}
\pgfpathlineto{\pgfpoint{-17.037733\du}{15.106837\du}}
\pgfpathlineto{\pgfpoint{-17.001952\du}{15.092759\du}}
\pgfpathlineto{\pgfpoint{-16.968517\du}{15.078681\du}}
\pgfpathlineto{\pgfpoint{-16.937135\du}{15.063723\du}}
\pgfpathlineto{\pgfpoint{-16.908099\du}{15.048472\du}}
\pgfpathlineto{\pgfpoint{-16.881410\du}{15.032341\du}}
\pgfpathlineto{\pgfpoint{-16.856774\du}{15.016211\du}}
\pgfpathlineto{\pgfpoint{-16.834484\du}{14.999493\du}}
\pgfpathlineto{\pgfpoint{-16.815126\du}{14.982482\du}}
\pgfpathlineto{\pgfpoint{-16.797822\du}{14.965178\du}}
\pgfpathlineto{\pgfpoint{-16.783451\du}{14.947581\du}}
\pgfpathlineto{\pgfpoint{-16.771133\du}{14.929397\du}}
\pgfpathlineto{\pgfpoint{-16.762041\du}{14.911506\du}}
\pgfpathlineto{\pgfpoint{-16.755589\du}{14.893323\du}}
\pgfpathlineto{\pgfpoint{-16.751776\du}{14.874845\du}}
\pgfpathlineto{\pgfpoint{-16.750016\du}{14.856661\du}}
\pgfusepath{fill}
\pgfsetbuttcap
\pgfsetmiterjoin
\pgfsetdash{}{0pt}
\definecolor{dialinecolor}{rgb}{0.666667, 0.901961, 1.000000}
\pgfsetstrokecolor{dialinecolor}
\pgfpathmoveto{\pgfpoint{-16.762334\du}{14.850796\du}}
\pgfpathlineto{\pgfpoint{-16.763801\du}{14.832318\du}}
\pgfpathlineto{\pgfpoint{-16.767320\du}{14.814721\du}}
\pgfpathlineto{\pgfpoint{-16.774066\du}{14.796537\du}}
\pgfpathlineto{\pgfpoint{-16.783158\du}{14.778940\du}}
\pgfpathlineto{\pgfpoint{-16.795183\du}{14.761343\du}}
\pgfpathlineto{\pgfpoint{-16.810434\du}{14.744038\du}}
\pgfpathlineto{\pgfpoint{-16.826271\du}{14.727321\du}}
\pgfpathlineto{\pgfpoint{-16.846215\du}{14.710310\du}}
\pgfpathlineto{\pgfpoint{-16.868212\du}{14.693886\du}}
\pgfpathlineto{\pgfpoint{-16.892555\du}{14.677462\du}}
\pgfpathlineto{\pgfpoint{-16.919537\du}{14.662211\du}}
\pgfpathlineto{\pgfpoint{-16.948280\du}{14.646960\du}}
\pgfpathlineto{\pgfpoint{-16.979368\du}{14.632295\du}}
\pgfpathlineto{\pgfpoint{-17.012510\du}{14.618218\du}}
\pgfpathlineto{\pgfpoint{-17.047705\du}{14.604726\du}}
\pgfpathlineto{\pgfpoint{-17.085246\du}{14.592408\du}}
\pgfpathlineto{\pgfpoint{-17.124546\du}{14.580677\du}}
\pgfpathlineto{\pgfpoint{-17.165020\du}{14.568652\du}}
\pgfpathlineto{\pgfpoint{-17.207254\du}{14.558680\du}}
\pgfpathlineto{\pgfpoint{-17.251247\du}{14.548121\du}}
\pgfpathlineto{\pgfpoint{-17.297000\du}{14.539030\du}}
\pgfpathlineto{\pgfpoint{-17.343340\du}{14.530817\du}}
\pgfpathlineto{\pgfpoint{-17.390853\du}{14.523192\du}}
\pgfpathlineto{\pgfpoint{-17.439832\du}{14.516740\du}}
\pgfpathlineto{\pgfpoint{-17.489398\du}{14.510580\du}}
\pgfpathlineto{\pgfpoint{-17.540430\du}{14.505888\du}}
\pgfpathlineto{\pgfpoint{-17.591462\du}{14.501489\du}}
\pgfpathlineto{\pgfpoint{-17.643374\du}{14.498556\du}}
\pgfpathlineto{\pgfpoint{-17.695580\du}{14.496503\du}}
\pgfpathlineto{\pgfpoint{-17.748079\du}{14.494743\du}}
\pgfpathlineto{\pgfpoint{-17.800577\du}{14.494743\du}}
\pgfpathlineto{\pgfpoint{-17.800577\du}{14.494743\du}}
\pgfpathlineto{\pgfpoint{-17.853663\du}{14.494743\du}}
\pgfpathlineto{\pgfpoint{-17.905868\du}{14.496503\du}}
\pgfpathlineto{\pgfpoint{-17.958367\du}{14.498556\du}}
\pgfpathlineto{\pgfpoint{-18.009692\du}{14.501489\du}}
\pgfpathlineto{\pgfpoint{-18.061018\du}{14.505888\du}}
\pgfpathlineto{\pgfpoint{-18.111757\du}{14.510580\du}}
\pgfpathlineto{\pgfpoint{-18.161616\du}{14.516740\du}}
\pgfpathlineto{\pgfpoint{-18.210595\du}{14.523192\du}}
\pgfpathlineto{\pgfpoint{-18.258401\du}{14.530817\du}}
\pgfpathlineto{\pgfpoint{-18.304741\du}{14.539030\du}}
\pgfpathlineto{\pgfpoint{-18.350494\du}{14.548121\du}}
\pgfpathlineto{\pgfpoint{-18.394487\du}{14.558680\du}}
\pgfpathlineto{\pgfpoint{-18.436428\du}{14.568652\du}}
\pgfpathlineto{\pgfpoint{-18.477195\du}{14.580677\du}}
\pgfpathlineto{\pgfpoint{-18.516496\du}{14.592408\du}}
\pgfpathlineto{\pgfpoint{-18.553743\du}{14.604726\du}}
\pgfpathlineto{\pgfpoint{-18.588645\du}{14.618218\du}}
\pgfpathlineto{\pgfpoint{-18.621786\du}{14.632295\du}}
\pgfpathlineto{\pgfpoint{-18.653168\du}{14.646960\du}}
\pgfpathlineto{\pgfpoint{-18.682497\du}{14.662211\du}}
\pgfpathlineto{\pgfpoint{-18.708600\du}{14.677462\du}}
\pgfpathlineto{\pgfpoint{-18.733236\du}{14.693886\du}}
\pgfpathlineto{\pgfpoint{-18.755233\du}{14.710310\du}}
\pgfpathlineto{\pgfpoint{-18.774590\du}{14.727321\du}}
\pgfpathlineto{\pgfpoint{-18.791601\du}{14.744038\du}}
\pgfpathlineto{\pgfpoint{-18.806265\du}{14.761343\du}}
\pgfpathlineto{\pgfpoint{-18.817997\du}{14.778940\du}}
\pgfpathlineto{\pgfpoint{-18.827382\du}{14.796537\du}}
\pgfpathlineto{\pgfpoint{-18.834128\du}{14.814721\du}}
\pgfpathlineto{\pgfpoint{-18.837940\du}{14.832318\du}}
\pgfpathlineto{\pgfpoint{-18.839407\du}{14.850796\du}}
\pgfpathlineto{\pgfpoint{-18.839407\du}{14.850796\du}}
\pgfpathlineto{\pgfpoint{-18.837940\du}{14.868686\du}}
\pgfpathlineto{\pgfpoint{-18.834128\du}{14.887163\du}}
\pgfpathlineto{\pgfpoint{-18.827382\du}{14.904468\du}}
\pgfpathlineto{\pgfpoint{-18.817997\du}{14.922358\du}}
\pgfpathlineto{\pgfpoint{-18.806265\du}{14.940249\du}}
\pgfpathlineto{\pgfpoint{-18.791601\du}{14.957260\du}}
\pgfpathlineto{\pgfpoint{-18.774590\du}{14.974270\du}}
\pgfpathlineto{\pgfpoint{-18.755233\du}{14.991281\du}}
\pgfpathlineto{\pgfpoint{-18.733236\du}{15.007412\du}}
\pgfpathlineto{\pgfpoint{-18.708600\du}{15.023836\du}}
\pgfpathlineto{\pgfpoint{-18.682497\du}{15.039380\du}}
\pgfpathlineto{\pgfpoint{-18.653168\du}{15.054045\du}}
\pgfpathlineto{\pgfpoint{-18.621786\du}{15.069003\du}}
\pgfpathlineto{\pgfpoint{-18.588645\du}{15.083374\du}}
\pgfpathlineto{\pgfpoint{-18.553743\du}{15.096278\du}}
\pgfpathlineto{\pgfpoint{-18.516496\du}{15.108890\du}}
\pgfpathlineto{\pgfpoint{-18.477195\du}{15.120915\du}}
\pgfpathlineto{\pgfpoint{-18.436428\du}{15.132646\du}}
\pgfpathlineto{\pgfpoint{-18.394487\du}{15.142911\du}}
\pgfpathlineto{\pgfpoint{-18.350494\du}{15.152883\du}}
\pgfpathlineto{\pgfpoint{-18.304741\du}{15.162268\du}}
\pgfpathlineto{\pgfpoint{-18.258401\du}{15.170481\du}}
\pgfpathlineto{\pgfpoint{-18.210595\du}{15.177813\du}}
\pgfpathlineto{\pgfpoint{-18.161616\du}{15.184852\du}}
\pgfpathlineto{\pgfpoint{-18.111757\du}{15.190718\du}}
\pgfpathlineto{\pgfpoint{-18.061018\du}{15.195410\du}}
\pgfpathlineto{\pgfpoint{-18.009692\du}{15.199516\du}}
\pgfpathlineto{\pgfpoint{-17.958367\du}{15.202449\du}}
\pgfpathlineto{\pgfpoint{-17.905868\du}{15.205089\du}}
\pgfpathlineto{\pgfpoint{-17.853663\du}{15.206555\du}}
\pgfpathlineto{\pgfpoint{-17.800577\du}{15.206848\du}}
\pgfpathlineto{\pgfpoint{-17.800577\du}{15.206848\du}}
\pgfpathlineto{\pgfpoint{-17.748079\du}{15.206555\du}}
\pgfpathlineto{\pgfpoint{-17.695580\du}{15.205089\du}}
\pgfpathlineto{\pgfpoint{-17.643374\du}{15.202449\du}}
\pgfpathlineto{\pgfpoint{-17.591462\du}{15.199516\du}}
\pgfpathlineto{\pgfpoint{-17.540430\du}{15.195410\du}}
\pgfpathlineto{\pgfpoint{-17.489398\du}{15.190718\du}}
\pgfpathlineto{\pgfpoint{-17.439832\du}{15.184852\du}}
\pgfpathlineto{\pgfpoint{-17.390853\du}{15.177813\du}}
\pgfpathlineto{\pgfpoint{-17.343340\du}{15.170481\du}}
\pgfpathlineto{\pgfpoint{-17.297000\du}{15.162268\du}}
\pgfpathlineto{\pgfpoint{-17.251247\du}{15.152883\du}}
\pgfpathlineto{\pgfpoint{-17.207254\du}{15.142911\du}}
\pgfpathlineto{\pgfpoint{-17.165020\du}{15.132646\du}}
\pgfpathlineto{\pgfpoint{-17.124546\du}{15.120915\du}}
\pgfpathlineto{\pgfpoint{-17.085246\du}{15.108890\du}}
\pgfpathlineto{\pgfpoint{-17.047705\du}{15.096278\du}}
\pgfpathlineto{\pgfpoint{-17.012510\du}{15.083374\du}}
\pgfpathlineto{\pgfpoint{-16.979368\du}{15.069003\du}}
\pgfpathlineto{\pgfpoint{-16.948280\du}{15.054045\du}}
\pgfpathlineto{\pgfpoint{-16.919537\du}{15.039380\du}}
\pgfpathlineto{\pgfpoint{-16.892555\du}{15.023836\du}}
\pgfpathlineto{\pgfpoint{-16.868212\du}{15.007412\du}}
\pgfpathlineto{\pgfpoint{-16.846215\du}{14.991281\du}}
\pgfpathlineto{\pgfpoint{-16.826271\du}{14.974270\du}}
\pgfpathlineto{\pgfpoint{-16.810434\du}{14.957260\du}}
\pgfpathlineto{\pgfpoint{-16.795183\du}{14.940249\du}}
\pgfpathlineto{\pgfpoint{-16.783158\du}{14.922358\du}}
\pgfpathlineto{\pgfpoint{-16.774066\du}{14.904468\du}}
\pgfpathlineto{\pgfpoint{-16.767320\du}{14.887163\du}}
\pgfpathlineto{\pgfpoint{-16.763801\du}{14.868686\du}}
\pgfpathlineto{\pgfpoint{-16.762334\du}{14.850796\du}}
\pgfusepath{stroke}
\pgfsetbuttcap
\pgfsetmiterjoin
\pgfsetdash{}{0pt}
\definecolor{dialinecolor}{rgb}{0.000000, 0.470588, 0.666667}
\pgfsetfillcolor{dialinecolor}
\pgfpathmoveto{\pgfpoint{-18.839407\du}{14.354257\du}}
\pgfpathlineto{\pgfpoint{-18.839407\du}{14.863407\du}}
\pgfpathlineto{\pgfpoint{-16.762334\du}{14.863407\du}}
\pgfpathlineto{\pgfpoint{-16.762334\du}{14.354257\du}}
\pgfpathlineto{\pgfpoint{-18.839407\du}{14.354257\du}}
\pgfusepath{fill}
\pgfsetbuttcap
\pgfsetmiterjoin
\pgfsetdash{}{0pt}
\definecolor{dialinecolor}{rgb}{0.000000, 0.705882, 1.000000}
\pgfsetfillcolor{dialinecolor}
\pgfpathmoveto{\pgfpoint{-16.750016\du}{14.348098\du}}
\pgfpathlineto{\pgfpoint{-16.751776\du}{14.329621\du}}
\pgfpathlineto{\pgfpoint{-16.755589\du}{14.310851\du}}
\pgfpathlineto{\pgfpoint{-16.762041\du}{14.292960\du}}
\pgfpathlineto{\pgfpoint{-16.771133\du}{14.274776\du}}
\pgfpathlineto{\pgfpoint{-16.783451\du}{14.256592\du}}
\pgfpathlineto{\pgfpoint{-16.797822\du}{14.238995\du}}
\pgfpathlineto{\pgfpoint{-16.815126\du}{14.221691\du}}
\pgfpathlineto{\pgfpoint{-16.834484\du}{14.204973\du}}
\pgfpathlineto{\pgfpoint{-16.856774\du}{14.187963\du}}
\pgfpathlineto{\pgfpoint{-16.881410\du}{14.172125\du}}
\pgfpathlineto{\pgfpoint{-16.908099\du}{14.155994\du}}
\pgfpathlineto{\pgfpoint{-16.937135\du}{14.140743\du}}
\pgfpathlineto{\pgfpoint{-16.968517\du}{14.125492\du}}
\pgfpathlineto{\pgfpoint{-17.001952\du}{14.111708\du}}
\pgfpathlineto{\pgfpoint{-17.037733\du}{14.097923\du}}
\pgfpathlineto{\pgfpoint{-17.075274\du}{14.084725\du}}
\pgfpathlineto{\pgfpoint{-17.114281\du}{14.072407\du}}
\pgfpathlineto{\pgfpoint{-17.155342\du}{14.060969\du}}
\pgfpathlineto{\pgfpoint{-17.197869\du}{14.050117\du}}
\pgfpathlineto{\pgfpoint{-17.242155\du}{14.040145\du}}
\pgfpathlineto{\pgfpoint{-17.287908\du}{14.031053\du}}
\pgfpathlineto{\pgfpoint{-17.335128\du}{14.022254\du}}
\pgfpathlineto{\pgfpoint{-17.383227\du}{14.014629\du}}
\pgfpathlineto{\pgfpoint{-17.431913\du}{14.007590\du}}
\pgfpathlineto{\pgfpoint{-17.482066\du}{14.001724\du}}
\pgfpathlineto{\pgfpoint{-17.533098\du}{13.996738\du}}
\pgfpathlineto{\pgfpoint{-17.584717\du}{13.992632\du}}
\pgfpathlineto{\pgfpoint{-17.636922\du}{13.989406\du}}
\pgfpathlineto{\pgfpoint{-17.689128\du}{13.987060\du}}
\pgfpathlineto{\pgfpoint{-17.741920\du}{13.985593\du}}
\pgfpathlineto{\pgfpoint{-17.794712\du}{13.985300\du}}
\pgfpathlineto{\pgfpoint{-17.794712\du}{13.985300\du}}
\pgfpathlineto{\pgfpoint{-17.847797\du}{13.985593\du}}
\pgfpathlineto{\pgfpoint{-17.900589\du}{13.987060\du}}
\pgfpathlineto{\pgfpoint{-17.952794\du}{13.989406\du}}
\pgfpathlineto{\pgfpoint{-18.005000\du}{13.992632\du}}
\pgfpathlineto{\pgfpoint{-18.056619\du}{13.996738\du}}
\pgfpathlineto{\pgfpoint{-18.107651\du}{14.001724\du}}
\pgfpathlineto{\pgfpoint{-18.157803\du}{14.007590\du}}
\pgfpathlineto{\pgfpoint{-18.206782\du}{14.014629\du}}
\pgfpathlineto{\pgfpoint{-18.254882\du}{14.022254\du}}
\pgfpathlineto{\pgfpoint{-18.301808\du}{14.031053\du}}
\pgfpathlineto{\pgfpoint{-18.347561\du}{14.040145\du}}
\pgfpathlineto{\pgfpoint{-18.391848\du}{14.050117\du}}
\pgfpathlineto{\pgfpoint{-18.434375\du}{14.060969\du}}
\pgfpathlineto{\pgfpoint{-18.475435\du}{14.072407\du}}
\pgfpathlineto{\pgfpoint{-18.514736\du}{14.084725\du}}
\pgfpathlineto{\pgfpoint{-18.551984\du}{14.097923\du}}
\pgfpathlineto{\pgfpoint{-18.587765\du}{14.111708\du}}
\pgfpathlineto{\pgfpoint{-18.621200\du}{14.125492\du}}
\pgfpathlineto{\pgfpoint{-18.652582\du}{14.140743\du}}
\pgfpathlineto{\pgfpoint{-18.681617\du}{14.155994\du}}
\pgfpathlineto{\pgfpoint{-18.708307\du}{14.172125\du}}
\pgfpathlineto{\pgfpoint{-18.732943\du}{14.187963\du}}
\pgfpathlineto{\pgfpoint{-18.755233\du}{14.204973\du}}
\pgfpathlineto{\pgfpoint{-18.774590\du}{14.221691\du}}
\pgfpathlineto{\pgfpoint{-18.791894\du}{14.238995\du}}
\pgfpathlineto{\pgfpoint{-18.806265\du}{14.256592\du}}
\pgfpathlineto{\pgfpoint{-18.818583\du}{14.274776\du}}
\pgfpathlineto{\pgfpoint{-18.827968\du}{14.292960\du}}
\pgfpathlineto{\pgfpoint{-18.834714\du}{14.310851\du}}
\pgfpathlineto{\pgfpoint{-18.838234\du}{14.329621\du}}
\pgfpathlineto{\pgfpoint{-18.839700\du}{14.348098\du}}
\pgfpathlineto{\pgfpoint{-18.839700\du}{14.348098\du}}
\pgfpathlineto{\pgfpoint{-18.838234\du}{14.366282\du}}
\pgfpathlineto{\pgfpoint{-18.834714\du}{14.384760\du}}
\pgfpathlineto{\pgfpoint{-18.827968\du}{14.402943\du}}
\pgfpathlineto{\pgfpoint{-18.818583\du}{14.420541\du}}
\pgfpathlineto{\pgfpoint{-18.806265\du}{14.438725\du}}
\pgfpathlineto{\pgfpoint{-18.791894\du}{14.456322\du}}
\pgfpathlineto{\pgfpoint{-18.774590\du}{14.473626\du}}
\pgfpathlineto{\pgfpoint{-18.755233\du}{14.490930\du}}
\pgfpathlineto{\pgfpoint{-18.732943\du}{14.507648\du}}
\pgfpathlineto{\pgfpoint{-18.708307\du}{14.523485\du}}
\pgfpathlineto{\pgfpoint{-18.681617\du}{14.539323\du}}
\pgfpathlineto{\pgfpoint{-18.652582\du}{14.555160\du}}
\pgfpathlineto{\pgfpoint{-18.621200\du}{14.569825\du}}
\pgfpathlineto{\pgfpoint{-18.587765\du}{14.583903\du}}
\pgfpathlineto{\pgfpoint{-18.551984\du}{14.597687\du}}
\pgfpathlineto{\pgfpoint{-18.514736\du}{14.610592\du}}
\pgfpathlineto{\pgfpoint{-18.475435\du}{14.623203\du}}
\pgfpathlineto{\pgfpoint{-18.434375\du}{14.634348\du}}
\pgfpathlineto{\pgfpoint{-18.391848\du}{14.645200\du}}
\pgfpathlineto{\pgfpoint{-18.347561\du}{14.655465\du}}
\pgfpathlineto{\pgfpoint{-18.301808\du}{14.665144\du}}
\pgfpathlineto{\pgfpoint{-18.254882\du}{14.673063\du}}
\pgfpathlineto{\pgfpoint{-18.206782\du}{14.680981\du}}
\pgfpathlineto{\pgfpoint{-18.157803\du}{14.688020\du}}
\pgfpathlineto{\pgfpoint{-18.107651\du}{14.693593\du}}
\pgfpathlineto{\pgfpoint{-18.056619\du}{14.698579\du}}
\pgfpathlineto{\pgfpoint{-18.005000\du}{14.702685\du}}
\pgfpathlineto{\pgfpoint{-17.952794\du}{14.705911\du}}
\pgfpathlineto{\pgfpoint{-17.900589\du}{14.708551\du}}
\pgfpathlineto{\pgfpoint{-17.847797\du}{14.710017\du}}
\pgfpathlineto{\pgfpoint{-17.794712\du}{14.710310\du}}
\pgfpathlineto{\pgfpoint{-17.794712\du}{14.710310\du}}
\pgfpathlineto{\pgfpoint{-17.741920\du}{14.710017\du}}
\pgfpathlineto{\pgfpoint{-17.689128\du}{14.708551\du}}
\pgfpathlineto{\pgfpoint{-17.636922\du}{14.705911\du}}
\pgfpathlineto{\pgfpoint{-17.584717\du}{14.702685\du}}
\pgfpathlineto{\pgfpoint{-17.533098\du}{14.698579\du}}
\pgfpathlineto{\pgfpoint{-17.482066\du}{14.693593\du}}
\pgfpathlineto{\pgfpoint{-17.431913\du}{14.688020\du}}
\pgfpathlineto{\pgfpoint{-17.383227\du}{14.680981\du}}
\pgfpathlineto{\pgfpoint{-17.335128\du}{14.673063\du}}
\pgfpathlineto{\pgfpoint{-17.287908\du}{14.665144\du}}
\pgfpathlineto{\pgfpoint{-17.242155\du}{14.655465\du}}
\pgfpathlineto{\pgfpoint{-17.197869\du}{14.645200\du}}
\pgfpathlineto{\pgfpoint{-17.155342\du}{14.634348\du}}
\pgfpathlineto{\pgfpoint{-17.114281\du}{14.623203\du}}
\pgfpathlineto{\pgfpoint{-17.075274\du}{14.610592\du}}
\pgfpathlineto{\pgfpoint{-17.037733\du}{14.597687\du}}
\pgfpathlineto{\pgfpoint{-17.001952\du}{14.583903\du}}
\pgfpathlineto{\pgfpoint{-16.968517\du}{14.569825\du}}
\pgfpathlineto{\pgfpoint{-16.937135\du}{14.555160\du}}
\pgfpathlineto{\pgfpoint{-16.908099\du}{14.539323\du}}
\pgfpathlineto{\pgfpoint{-16.881410\du}{14.523485\du}}
\pgfpathlineto{\pgfpoint{-16.856774\du}{14.507648\du}}
\pgfpathlineto{\pgfpoint{-16.834484\du}{14.490930\du}}
\pgfpathlineto{\pgfpoint{-16.815126\du}{14.473626\du}}
\pgfpathlineto{\pgfpoint{-16.797822\du}{14.456322\du}}
\pgfpathlineto{\pgfpoint{-16.783451\du}{14.438725\du}}
\pgfpathlineto{\pgfpoint{-16.771133\du}{14.420541\du}}
\pgfpathlineto{\pgfpoint{-16.762041\du}{14.402943\du}}
\pgfpathlineto{\pgfpoint{-16.755589\du}{14.384760\du}}
\pgfpathlineto{\pgfpoint{-16.751776\du}{14.366282\du}}
\pgfpathlineto{\pgfpoint{-16.750016\du}{14.348098\du}}
\pgfusepath{fill}
\pgfsetbuttcap
\pgfsetmiterjoin
\pgfsetdash{}{0pt}
\definecolor{dialinecolor}{rgb}{0.666667, 0.901961, 1.000000}
\pgfsetstrokecolor{dialinecolor}
\pgfpathmoveto{\pgfpoint{-16.762334\du}{14.341939\du}}
\pgfpathlineto{\pgfpoint{-16.763801\du}{14.323462\du}}
\pgfpathlineto{\pgfpoint{-16.767320\du}{14.305865\du}}
\pgfpathlineto{\pgfpoint{-16.774066\du}{14.287974\du}}
\pgfpathlineto{\pgfpoint{-16.783158\du}{14.270084\du}}
\pgfpathlineto{\pgfpoint{-16.795183\du}{14.252780\du}}
\pgfpathlineto{\pgfpoint{-16.810434\du}{14.234889\du}}
\pgfpathlineto{\pgfpoint{-16.826271\du}{14.218465\du}}
\pgfpathlineto{\pgfpoint{-16.846215\du}{14.201454\du}}
\pgfpathlineto{\pgfpoint{-16.868212\du}{14.185323\du}}
\pgfpathlineto{\pgfpoint{-16.892555\du}{14.168606\du}}
\pgfpathlineto{\pgfpoint{-16.919537\du}{14.153648\du}}
\pgfpathlineto{\pgfpoint{-16.948280\du}{14.138104\du}}
\pgfpathlineto{\pgfpoint{-16.979368\du}{14.123439\du}}
\pgfpathlineto{\pgfpoint{-17.012510\du}{14.109948\du}}
\pgfpathlineto{\pgfpoint{-17.047705\du}{14.095870\du}}
\pgfpathlineto{\pgfpoint{-17.085246\du}{14.084138\du}}
\pgfpathlineto{\pgfpoint{-17.124546\du}{14.071527\du}}
\pgfpathlineto{\pgfpoint{-17.165020\du}{14.060089\du}}
\pgfpathlineto{\pgfpoint{-17.207254\du}{14.049530\du}}
\pgfpathlineto{\pgfpoint{-17.251247\du}{14.039265\du}}
\pgfpathlineto{\pgfpoint{-17.297000\du}{14.030173\du}}
\pgfpathlineto{\pgfpoint{-17.343340\du}{14.021961\du}}
\pgfpathlineto{\pgfpoint{-17.390853\du}{14.014629\du}}
\pgfpathlineto{\pgfpoint{-17.439832\du}{14.007590\du}}
\pgfpathlineto{\pgfpoint{-17.489398\du}{14.001724\du}}
\pgfpathlineto{\pgfpoint{-17.540430\du}{13.997032\du}}
\pgfpathlineto{\pgfpoint{-17.591462\du}{13.992632\du}}
\pgfpathlineto{\pgfpoint{-17.643374\du}{13.989699\du}}
\pgfpathlineto{\pgfpoint{-17.695580\du}{13.987353\du}}
\pgfpathlineto{\pgfpoint{-17.748079\du}{13.985887\du}}
\pgfpathlineto{\pgfpoint{-17.800577\du}{13.985593\du}}
\pgfpathlineto{\pgfpoint{-17.800577\du}{13.985593\du}}
\pgfpathlineto{\pgfpoint{-17.853663\du}{13.985887\du}}
\pgfpathlineto{\pgfpoint{-17.905868\du}{13.987353\du}}
\pgfpathlineto{\pgfpoint{-17.958367\du}{13.989699\du}}
\pgfpathlineto{\pgfpoint{-18.009692\du}{13.992632\du}}
\pgfpathlineto{\pgfpoint{-18.061018\du}{13.997032\du}}
\pgfpathlineto{\pgfpoint{-18.111757\du}{14.001724\du}}
\pgfpathlineto{\pgfpoint{-18.161616\du}{14.007590\du}}
\pgfpathlineto{\pgfpoint{-18.210595\du}{14.014629\du}}
\pgfpathlineto{\pgfpoint{-18.258401\du}{14.021961\du}}
\pgfpathlineto{\pgfpoint{-18.304741\du}{14.030173\du}}
\pgfpathlineto{\pgfpoint{-18.350494\du}{14.039265\du}}
\pgfpathlineto{\pgfpoint{-18.394487\du}{14.049530\du}}
\pgfpathlineto{\pgfpoint{-18.436428\du}{14.060089\du}}
\pgfpathlineto{\pgfpoint{-18.477195\du}{14.071527\du}}
\pgfpathlineto{\pgfpoint{-18.516496\du}{14.084138\du}}
\pgfpathlineto{\pgfpoint{-18.553743\du}{14.095870\du}}
\pgfpathlineto{\pgfpoint{-18.588645\du}{14.109948\du}}
\pgfpathlineto{\pgfpoint{-18.621786\du}{14.123439\du}}
\pgfpathlineto{\pgfpoint{-18.653168\du}{14.138104\du}}
\pgfpathlineto{\pgfpoint{-18.682497\du}{14.153648\du}}
\pgfpathlineto{\pgfpoint{-18.708600\du}{14.168606\du}}
\pgfpathlineto{\pgfpoint{-18.733236\du}{14.185323\du}}
\pgfpathlineto{\pgfpoint{-18.755233\du}{14.201454\du}}
\pgfpathlineto{\pgfpoint{-18.774590\du}{14.218465\du}}
\pgfpathlineto{\pgfpoint{-18.791601\du}{14.234889\du}}
\pgfpathlineto{\pgfpoint{-18.806265\du}{14.252780\du}}
\pgfpathlineto{\pgfpoint{-18.817997\du}{14.270084\du}}
\pgfpathlineto{\pgfpoint{-18.827382\du}{14.287974\du}}
\pgfpathlineto{\pgfpoint{-18.834128\du}{14.305865\du}}
\pgfpathlineto{\pgfpoint{-18.837940\du}{14.323462\du}}
\pgfpathlineto{\pgfpoint{-18.839407\du}{14.341939\du}}
\pgfpathlineto{\pgfpoint{-18.839407\du}{14.341939\du}}
\pgfpathlineto{\pgfpoint{-18.837940\du}{14.360123\du}}
\pgfpathlineto{\pgfpoint{-18.834128\du}{14.378014\du}}
\pgfpathlineto{\pgfpoint{-18.827382\du}{14.395905\du}}
\pgfpathlineto{\pgfpoint{-18.817997\du}{14.413502\du}}
\pgfpathlineto{\pgfpoint{-18.806265\du}{14.431099\du}}
\pgfpathlineto{\pgfpoint{-18.791601\du}{14.448697\du}}
\pgfpathlineto{\pgfpoint{-18.774590\du}{14.465414\du}}
\pgfpathlineto{\pgfpoint{-18.755233\du}{14.482425\du}}
\pgfpathlineto{\pgfpoint{-18.733236\du}{14.498556\du}}
\pgfpathlineto{\pgfpoint{-18.708600\du}{14.514980\du}}
\pgfpathlineto{\pgfpoint{-18.682497\du}{14.530524\du}}
\pgfpathlineto{\pgfpoint{-18.653168\du}{14.545189\du}}
\pgfpathlineto{\pgfpoint{-18.621786\du}{14.560146\du}}
\pgfpathlineto{\pgfpoint{-18.588645\du}{14.574224\du}}
\pgfpathlineto{\pgfpoint{-18.553743\du}{14.587422\du}}
\pgfpathlineto{\pgfpoint{-18.516496\du}{14.600034\du}}
\pgfpathlineto{\pgfpoint{-18.477195\du}{14.612058\du}}
\pgfpathlineto{\pgfpoint{-18.436428\du}{14.623497\du}}
\pgfpathlineto{\pgfpoint{-18.394487\du}{14.634055\du}}
\pgfpathlineto{\pgfpoint{-18.350494\du}{14.644320\du}}
\pgfpathlineto{\pgfpoint{-18.304741\du}{14.653999\du}}
\pgfpathlineto{\pgfpoint{-18.258401\du}{14.661624\du}}
\pgfpathlineto{\pgfpoint{-18.210595\du}{14.669543\du}}
\pgfpathlineto{\pgfpoint{-18.161616\du}{14.676289\du}}
\pgfpathlineto{\pgfpoint{-18.111757\du}{14.681861\du}}
\pgfpathlineto{\pgfpoint{-18.061018\du}{14.687140\du}}
\pgfpathlineto{\pgfpoint{-18.009692\du}{14.691246\du}}
\pgfpathlineto{\pgfpoint{-17.958367\du}{14.693886\du}}
\pgfpathlineto{\pgfpoint{-17.905868\du}{14.696526\du}}
\pgfpathlineto{\pgfpoint{-17.853663\du}{14.697992\du}}
\pgfpathlineto{\pgfpoint{-17.800577\du}{14.698285\du}}
\pgfpathlineto{\pgfpoint{-17.800577\du}{14.698285\du}}
\pgfpathlineto{\pgfpoint{-17.748079\du}{14.697992\du}}
\pgfpathlineto{\pgfpoint{-17.695580\du}{14.696526\du}}
\pgfpathlineto{\pgfpoint{-17.643374\du}{14.693886\du}}
\pgfpathlineto{\pgfpoint{-17.591462\du}{14.691246\du}}
\pgfpathlineto{\pgfpoint{-17.540430\du}{14.687140\du}}
\pgfpathlineto{\pgfpoint{-17.489398\du}{14.681861\du}}
\pgfpathlineto{\pgfpoint{-17.439832\du}{14.676289\du}}
\pgfpathlineto{\pgfpoint{-17.390853\du}{14.669543\du}}
\pgfpathlineto{\pgfpoint{-17.343340\du}{14.661624\du}}
\pgfpathlineto{\pgfpoint{-17.297000\du}{14.653999\du}}
\pgfpathlineto{\pgfpoint{-17.251247\du}{14.644320\du}}
\pgfpathlineto{\pgfpoint{-17.207254\du}{14.634055\du}}
\pgfpathlineto{\pgfpoint{-17.165020\du}{14.623497\du}}
\pgfpathlineto{\pgfpoint{-17.124546\du}{14.612058\du}}
\pgfpathlineto{\pgfpoint{-17.085246\du}{14.600034\du}}
\pgfpathlineto{\pgfpoint{-17.047705\du}{14.587422\du}}
\pgfpathlineto{\pgfpoint{-17.012510\du}{14.574224\du}}
\pgfpathlineto{\pgfpoint{-16.979368\du}{14.560146\du}}
\pgfpathlineto{\pgfpoint{-16.948280\du}{14.545189\du}}
\pgfpathlineto{\pgfpoint{-16.919537\du}{14.530524\du}}
\pgfpathlineto{\pgfpoint{-16.892555\du}{14.514980\du}}
\pgfpathlineto{\pgfpoint{-16.868212\du}{14.498556\du}}
\pgfpathlineto{\pgfpoint{-16.846215\du}{14.482425\du}}
\pgfpathlineto{\pgfpoint{-16.826271\du}{14.465414\du}}
\pgfpathlineto{\pgfpoint{-16.810434\du}{14.448697\du}}
\pgfpathlineto{\pgfpoint{-16.795183\du}{14.431099\du}}
\pgfpathlineto{\pgfpoint{-16.783158\du}{14.413502\du}}
\pgfpathlineto{\pgfpoint{-16.774066\du}{14.395905\du}}
\pgfpathlineto{\pgfpoint{-16.767320\du}{14.378014\du}}
\pgfpathlineto{\pgfpoint{-16.763801\du}{14.360123\du}}
\pgfpathlineto{\pgfpoint{-16.762334\du}{14.341939\du}}
\pgfusepath{stroke}
\pgfsetbuttcap
\pgfsetmiterjoin
\pgfsetdash{}{0pt}
\definecolor{dialinecolor}{rgb}{0.000000, 0.000000, 0.000000}
\pgfsetfillcolor{dialinecolor}
\pgfpathmoveto{\pgfpoint{-17.775354\du}{14.265684\du}}
\pgfpathlineto{\pgfpoint{-17.623431\du}{14.316716\du}}
\pgfpathlineto{\pgfpoint{-17.255940\du}{14.163620\du}}
\pgfpathlineto{\pgfpoint{-17.091405\du}{14.214359\du}}
\pgfpathlineto{\pgfpoint{-17.180271\du}{14.087658\du}}
\pgfpathlineto{\pgfpoint{-17.610526\du}{14.087658\du}}
\pgfpathlineto{\pgfpoint{-17.433673\du}{14.125492\du}}
\pgfpathlineto{\pgfpoint{-17.775354\du}{14.265684\du}}
\pgfusepath{fill}
\pgfsetbuttcap
\pgfsetmiterjoin
\pgfsetdash{}{0pt}
\definecolor{dialinecolor}{rgb}{0.000000, 0.000000, 0.000000}
\pgfsetfillcolor{dialinecolor}
\pgfpathmoveto{\pgfpoint{-17.838998\du}{14.405583\du}}
\pgfpathlineto{\pgfpoint{-17.990629\du}{14.354257\du}}
\pgfpathlineto{\pgfpoint{-18.345508\du}{14.507648\du}}
\pgfpathlineto{\pgfpoint{-18.522655\du}{14.456322\du}}
\pgfpathlineto{\pgfpoint{-18.434375\du}{14.596221\du}}
\pgfpathlineto{\pgfpoint{-17.990629\du}{14.596221\du}}
\pgfpathlineto{\pgfpoint{-18.180973\du}{14.545189\du}}
\pgfpathlineto{\pgfpoint{-17.838998\du}{14.405583\du}}
\pgfusepath{fill}
\pgfsetbuttcap
\pgfsetmiterjoin
\pgfsetdash{}{0pt}
\definecolor{dialinecolor}{rgb}{0.000000, 0.000000, 0.000000}
\pgfsetfillcolor{dialinecolor}
\pgfpathmoveto{\pgfpoint{-18.484820\du}{14.125492\du}}
\pgfpathlineto{\pgfpoint{-18.333190\du}{14.074753\du}}
\pgfpathlineto{\pgfpoint{-17.965406\du}{14.214359\du}}
\pgfpathlineto{\pgfpoint{-17.800577\du}{14.176524\du}}
\pgfpathlineto{\pgfpoint{-17.889151\du}{14.303225\du}}
\pgfpathlineto{\pgfpoint{-18.320285\du}{14.303225\du}}
\pgfpathlineto{\pgfpoint{-18.143139\du}{14.265684\du}}
\pgfpathlineto{\pgfpoint{-18.484820\du}{14.125492\du}}
\pgfusepath{fill}
\pgfsetbuttcap
\pgfsetmiterjoin
\pgfsetdash{}{0pt}
\definecolor{dialinecolor}{rgb}{0.000000, 0.000000, 0.000000}
\pgfsetfillcolor{dialinecolor}
\pgfpathmoveto{\pgfpoint{-17.116921\du}{14.557800\du}}
\pgfpathlineto{\pgfpoint{-17.269138\du}{14.608832\du}}
\pgfpathlineto{\pgfpoint{-17.623431\du}{14.456322\du}}
\pgfpathlineto{\pgfpoint{-17.800577\du}{14.507648\du}}
\pgfpathlineto{\pgfpoint{-17.712004\du}{14.379774\du}}
\pgfpathlineto{\pgfpoint{-17.269138\du}{14.379774\du}}
\pgfpathlineto{\pgfpoint{-17.458896\du}{14.417901\du}}
\pgfpathlineto{\pgfpoint{-17.116921\du}{14.557800\du}}
\pgfusepath{fill}
\pgfsetbuttcap
\pgfsetmiterjoin
\pgfsetdash{}{0pt}
\definecolor{dialinecolor}{rgb}{1.000000, 1.000000, 1.000000}
\pgfsetfillcolor{dialinecolor}
\pgfpathmoveto{\pgfpoint{-17.762743\du}{14.278296\du}}
\pgfpathlineto{\pgfpoint{-17.610526\du}{14.329035\du}}
\pgfpathlineto{\pgfpoint{-17.243915\du}{14.176524\du}}
\pgfpathlineto{\pgfpoint{-17.078793\du}{14.227557\du}}
\pgfpathlineto{\pgfpoint{-17.167660\du}{14.099976\du}}
\pgfpathlineto{\pgfpoint{-17.598208\du}{14.099976\du}}
\pgfpathlineto{\pgfpoint{-17.420768\du}{14.137810\du}}
\pgfpathlineto{\pgfpoint{-17.762743\du}{14.278296\du}}
\pgfusepath{fill}
\pgfsetbuttcap
\pgfsetmiterjoin
\pgfsetdash{}{0pt}
\definecolor{dialinecolor}{rgb}{1.000000, 1.000000, 1.000000}
\pgfsetfillcolor{dialinecolor}
\pgfpathmoveto{\pgfpoint{-17.826093\du}{14.417901\du}}
\pgfpathlineto{\pgfpoint{-17.978604\du}{14.367162\du}}
\pgfpathlineto{\pgfpoint{-18.333190\du}{14.519672\du}}
\pgfpathlineto{\pgfpoint{-18.510337\du}{14.468933\du}}
\pgfpathlineto{\pgfpoint{-18.421763\du}{14.608832\du}}
\pgfpathlineto{\pgfpoint{-17.978604\du}{14.608832\du}}
\pgfpathlineto{\pgfpoint{-18.168362\du}{14.557800\du}}
\pgfpathlineto{\pgfpoint{-17.826093\du}{14.417901\du}}
\pgfusepath{fill}
\pgfsetbuttcap
\pgfsetmiterjoin
\pgfsetdash{}{0pt}
\definecolor{dialinecolor}{rgb}{1.000000, 1.000000, 1.000000}
\pgfsetfillcolor{dialinecolor}
\pgfpathmoveto{\pgfpoint{-18.472209\du}{14.137810\du}}
\pgfpathlineto{\pgfpoint{-18.320285\du}{14.087658\du}}
\pgfpathlineto{\pgfpoint{-17.952794\du}{14.227557\du}}
\pgfpathlineto{\pgfpoint{-17.788259\du}{14.188843\du}}
\pgfpathlineto{\pgfpoint{-17.876539\du}{14.316716\du}}
\pgfpathlineto{\pgfpoint{-18.307674\du}{14.316716\du}}
\pgfpathlineto{\pgfpoint{-18.129941\du}{14.278296\du}}
\pgfpathlineto{\pgfpoint{-18.472209\du}{14.137810\du}}
\pgfusepath{fill}
\pgfsetbuttcap
\pgfsetmiterjoin
\pgfsetdash{}{0pt}
\definecolor{dialinecolor}{rgb}{1.000000, 1.000000, 1.000000}
\pgfsetfillcolor{dialinecolor}
\pgfpathmoveto{\pgfpoint{-17.104603\du}{14.570998\du}}
\pgfpathlineto{\pgfpoint{-17.255940\du}{14.621737\du}}
\pgfpathlineto{\pgfpoint{-17.610526\du}{14.468933\du}}
\pgfpathlineto{\pgfpoint{-17.788259\du}{14.519672\du}}
\pgfpathlineto{\pgfpoint{-17.699686\du}{14.392972\du}}
\pgfpathlineto{\pgfpoint{-17.255940\du}{14.392972\du}}
\pgfpathlineto{\pgfpoint{-17.445991\du}{14.430806\du}}
\pgfpathlineto{\pgfpoint{-17.104603\du}{14.570998\du}}
\pgfusepath{fill}
\pgfsetbuttcap
\pgfsetmiterjoin
\pgfsetdash{}{0pt}
\definecolor{dialinecolor}{rgb}{0.666667, 0.901961, 1.000000}
\pgfsetstrokecolor{dialinecolor}
\pgfpathmoveto{\pgfpoint{-18.839407\du}{14.341939\du}}
\pgfpathlineto{\pgfpoint{-18.839407\du}{14.850209\du}}
\pgfusepath{stroke}
\pgfsetbuttcap
\pgfsetmiterjoin
\pgfsetdash{}{0pt}
\definecolor{dialinecolor}{rgb}{0.666667, 0.901961, 1.000000}
\pgfsetstrokecolor{dialinecolor}
\pgfpathmoveto{\pgfpoint{-16.762334\du}{14.341939\du}}
\pgfpathlineto{\pgfpoint{-16.762334\du}{14.850209\du}}
\pgfusepath{stroke}
\pgfsetlinewidth{0.100000\du}
\pgfsetdash{}{0pt}
\pgfsetdash{}{0pt}
\pgfsetbuttcap
{
\definecolor{dialinecolor}{rgb}{0.000000, 0.000000, 0.000000}
\pgfsetfillcolor{dialinecolor}
% was here!!!
\definecolor{dialinecolor}{rgb}{0.000000, 0.000000, 0.000000}
\pgfsetstrokecolor{dialinecolor}
\draw (-4.905142\du,0.147589\du)--(-7.271944\du,2.528103\du);
}
\pgfsetlinewidth{0.100000\du}
\pgfsetdash{}{0pt}
\pgfsetdash{}{0pt}
\pgfsetbuttcap
{
\definecolor{dialinecolor}{rgb}{0.000000, 0.000000, 0.000000}
\pgfsetfillcolor{dialinecolor}
% was here!!!
\definecolor{dialinecolor}{rgb}{0.000000, 0.000000, 0.000000}
\pgfsetstrokecolor{dialinecolor}
\draw (-17.769806\du,11.700000\du)--(-17.789954\du,13.986689\du);
}
\pgfsetlinewidth{0.100000\du}
\pgfsetdash{}{0pt}
\pgfsetdash{}{0pt}
\pgfsetbuttcap
{
\definecolor{dialinecolor}{rgb}{0.000000, 0.000000, 0.000000}
\pgfsetfillcolor{dialinecolor}
% was here!!!
\definecolor{dialinecolor}{rgb}{0.000000, 0.000000, 0.000000}
\pgfsetstrokecolor{dialinecolor}
\draw (-17.795005\du,15.219167\du)--(-21.612953\du,17.196445\du);
}
\pgfsetlinewidth{0.100000\du}
\pgfsetdash{}{0pt}
\pgfsetdash{}{0pt}
\pgfsetbuttcap
\pgfsetmiterjoin
\pgfsetlinewidth{0.001000\du}
\pgfsetbuttcap
\pgfsetmiterjoin
\pgfsetdash{}{0pt}
\definecolor{dialinecolor}{rgb}{0.000000, 0.470588, 0.666667}
\pgfsetstrokecolor{dialinecolor}
\pgfpathmoveto{\pgfpoint{-19.330683\du}{21.580654\du}}
\pgfpathlineto{\pgfpoint{-19.324115\du}{21.571697\du}}
\pgfpathlineto{\pgfpoint{-19.317845\du}{21.562144\du}}
\pgfpathlineto{\pgfpoint{-19.312471\du}{21.551993\du}}
\pgfpathlineto{\pgfpoint{-19.307396\du}{21.542738\du}}
\pgfpathlineto{\pgfpoint{-19.303515\du}{21.532587\du}}
\pgfpathlineto{\pgfpoint{-19.300231\du}{21.522436\du}}
\pgfpathlineto{\pgfpoint{-19.297245\du}{21.512286\du}}
\pgfpathlineto{\pgfpoint{-19.295454\du}{21.501836\du}}
\pgfpathlineto{\pgfpoint{-19.294260\du}{21.491685\du}}
\pgfpathlineto{\pgfpoint{-19.293364\du}{21.481236\du}}
\pgfpathlineto{\pgfpoint{-19.293364\du}{21.470787\du}}
\pgfpathlineto{\pgfpoint{-19.294260\du}{21.460039\du}}
\pgfpathlineto{\pgfpoint{-19.295454\du}{21.449888\du}}
\pgfpathlineto{\pgfpoint{-19.297245\du}{21.439439\du}}
\pgfpathlineto{\pgfpoint{-19.299932\du}{21.429587\du}}
\pgfpathlineto{\pgfpoint{-19.303515\du}{21.419137\du}}
\pgfpathlineto{\pgfpoint{-19.307396\du}{21.409285\du}}
\pgfpathlineto{\pgfpoint{-19.312471\du}{21.399134\du}}
\pgfpathlineto{\pgfpoint{-19.317845\du}{21.389581\du}}
\pgfpathlineto{\pgfpoint{-19.324115\du}{21.380624\du}}
\pgfpathlineto{\pgfpoint{-19.330683\du}{21.371070\du}}
\pgfpathlineto{\pgfpoint{-19.338147\du}{21.362412\du}}
\pgfpathlineto{\pgfpoint{-19.345611\du}{21.353754\du}}
\pgfpathlineto{\pgfpoint{-19.353970\du}{21.345395\du}}
\pgfpathlineto{\pgfpoint{-19.362628\du}{21.337334\du}}
\pgfpathlineto{\pgfpoint{-19.371883\du}{21.329572\du}}
\pgfpathlineto{\pgfpoint{-19.381437\du}{21.322406\du}}
\pgfpathlineto{\pgfpoint{-19.392185\du}{21.315838\du}}
\pgfpathlineto{\pgfpoint{-19.402634\du}{21.309270\du}}
\pgfpathlineto{\pgfpoint{-19.413979\du}{21.303001\du}}
\pgfpathlineto{\pgfpoint{-19.425921\du}{21.297328\du}}
\pgfpathlineto{\pgfpoint{-19.437863\du}{21.291954\du}}
\pgfpathlineto{\pgfpoint{-19.450104\du}{21.287774\du}}
\pgfpathlineto{\pgfpoint{-19.462643\du}{21.283595\du}}
\pgfpathlineto{\pgfpoint{-19.475779\du}{21.279415\du}}
\pgfpathlineto{\pgfpoint{-19.488617\du}{21.276429\du}}
\pgfpathlineto{\pgfpoint{-19.501753\du}{21.273742\du}}
\pgfpathlineto{\pgfpoint{-19.515487\du}{21.271354\du}}
\pgfpathlineto{\pgfpoint{-19.528623\du}{21.269861\du}}
\pgfpathlineto{\pgfpoint{-19.542655\du}{21.268667\du}}
\pgfpathlineto{\pgfpoint{-19.556388\du}{21.267771\du}}
\pgfpathlineto{\pgfpoint{-19.570420\du}{21.267771\du}}
\pgfpathlineto{\pgfpoint{-19.583855\du}{21.268369\du}}
\pgfpathlineto{\pgfpoint{-19.597589\du}{21.268966\du}}
\pgfpathlineto{\pgfpoint{-19.611621\du}{21.270458\du}}
\pgfpathlineto{\pgfpoint{-19.624757\du}{21.272847\du}}
\pgfpathlineto{\pgfpoint{-19.638490\du}{21.275235\du}}
\pgfpathlineto{\pgfpoint{-19.651029\du}{21.278221\du}}
\pgfpathlineto{\pgfpoint{-19.664166\du}{21.281206\du}}
\pgfpathlineto{\pgfpoint{-19.676705\du}{21.285386\du}}
\pgfpathlineto{\pgfpoint{-19.689543\du}{21.290163\du}}
\pgfpathlineto{\pgfpoint{-19.701783\du}{21.294641\du}}
\pgfpathlineto{\pgfpoint{-19.713128\du}{21.300015\du}}
\pgfpathlineto{\pgfpoint{-19.724772\du}{21.306285\du}}
\pgfpathlineto{\pgfpoint{-19.735818\du}{21.312256\du}}
\pgfpathlineto{\pgfpoint{-19.746268\du}{21.319122\du}}
\pgfpathlineto{\pgfpoint{-19.756418\du}{21.325691\du}}
\pgfpathlineto{\pgfpoint{-19.765972\du}{21.333453\du}}
\pgfpathlineto{\pgfpoint{-19.775526\du}{21.341514\du}}
\pgfpathlineto{\pgfpoint{-19.783885\du}{21.349575\du}}
\pgfpathlineto{\pgfpoint{-19.791946\du}{21.357934\du}}
\pgfpathlineto{\pgfpoint{-19.799708\du}{21.366891\du}}
\pgfusepath{stroke}
\pgfsetlinewidth{0.080000\du}
\pgfsetbuttcap
\pgfsetmiterjoin
\pgfsetdash{}{0pt}
\definecolor{dialinecolor}{rgb}{0.666667, 0.901961, 1.000000}
\pgfsetstrokecolor{dialinecolor}
\pgfpathmoveto{\pgfpoint{-19.339938\du}{21.573190\du}}
\pgfpathlineto{\pgfpoint{-19.334564\du}{21.564234\du}}
\pgfpathlineto{\pgfpoint{-19.328295\du}{21.555576\du}}
\pgfpathlineto{\pgfpoint{-19.323219\du}{21.546320\du}}
\pgfpathlineto{\pgfpoint{-19.319040\du}{21.537065\du}}
\pgfpathlineto{\pgfpoint{-19.315158\du}{21.527512\du}}
\pgfpathlineto{\pgfpoint{-19.312173\du}{21.517958\du}}
\pgfpathlineto{\pgfpoint{-19.309486\du}{21.508106\du}}
\pgfpathlineto{\pgfpoint{-19.307396\du}{21.498851\du}}
\pgfpathlineto{\pgfpoint{-19.306500\du}{21.488998\du}}
\pgfpathlineto{\pgfpoint{-19.305903\du}{21.478848\du}}
\pgfpathlineto{\pgfpoint{-19.305903\du}{21.468995\du}}
\pgfpathlineto{\pgfpoint{-19.306799\du}{21.459143\du}}
\pgfpathlineto{\pgfpoint{-19.308292\du}{21.448993\du}}
\pgfpathlineto{\pgfpoint{-19.310083\du}{21.439737\du}}
\pgfpathlineto{\pgfpoint{-19.313069\du}{21.429885\du}}
\pgfpathlineto{\pgfpoint{-19.316651\du}{21.420332\du}}
\pgfpathlineto{\pgfpoint{-19.320532\du}{21.411076\du}}
\pgfpathlineto{\pgfpoint{-19.325011\du}{21.401821\du}}
\pgfpathlineto{\pgfpoint{-19.330086\du}{21.393163\du}}
\pgfpathlineto{\pgfpoint{-19.335758\du}{21.383908\du}}
\pgfpathlineto{\pgfpoint{-19.342924\du}{21.375549\du}}
\pgfpathlineto{\pgfpoint{-19.349492\du}{21.366891\du}}
\pgfpathlineto{\pgfpoint{-19.357254\du}{21.358830\du}}
\pgfpathlineto{\pgfpoint{-19.365315\du}{21.351067\du}}
\pgfpathlineto{\pgfpoint{-19.373675\du}{21.344201\du}}
\pgfpathlineto{\pgfpoint{-19.382333\du}{21.336737\du}}
\pgfpathlineto{\pgfpoint{-19.392185\du}{21.329870\du}}
\pgfpathlineto{\pgfpoint{-19.402634\du}{21.323004\du}}
\pgfpathlineto{\pgfpoint{-19.412785\du}{21.317331\du}}
\pgfpathlineto{\pgfpoint{-19.423234\du}{21.311957\du}}
\pgfpathlineto{\pgfpoint{-19.435176\du}{21.306583\du}}
\pgfpathlineto{\pgfpoint{-19.445924\du}{21.301508\du}}
\pgfpathlineto{\pgfpoint{-19.458463\du}{21.297328\du}}
\pgfpathlineto{\pgfpoint{-19.469808\du}{21.293447\du}}
\pgfpathlineto{\pgfpoint{-19.482348\du}{21.290163\du}}
\pgfpathlineto{\pgfpoint{-19.494887\du}{21.287177\du}}
\pgfpathlineto{\pgfpoint{-19.507426\du}{21.284490\du}}
\pgfpathlineto{\pgfpoint{-19.520562\du}{21.282699\du}}
\pgfpathlineto{\pgfpoint{-19.533101\du}{21.281206\du}}
\pgfpathlineto{\pgfpoint{-19.546238\du}{21.280012\du}}
\pgfpathlineto{\pgfpoint{-19.560270\du}{21.279415\du}}
\pgfpathlineto{\pgfpoint{-19.573406\du}{21.279415\du}}
\pgfpathlineto{\pgfpoint{-19.586244\du}{21.280012\du}}
\pgfpathlineto{\pgfpoint{-19.599380\du}{21.280908\du}}
\pgfpathlineto{\pgfpoint{-19.611919\du}{21.282699\du}}
\pgfpathlineto{\pgfpoint{-19.625055\du}{21.284490\du}}
\pgfpathlineto{\pgfpoint{-19.637595\du}{21.286580\du}}
\pgfpathlineto{\pgfpoint{-19.650432\du}{21.290163\du}}
\pgfpathlineto{\pgfpoint{-19.662673\du}{21.293148\du}}
\pgfpathlineto{\pgfpoint{-19.674615\du}{21.297030\du}}
\pgfpathlineto{\pgfpoint{-19.686856\du}{21.301209\du}}
\pgfpathlineto{\pgfpoint{-19.698201\du}{21.306285\du}}
\pgfpathlineto{\pgfpoint{-19.709247\du}{21.311061\du}}
\pgfpathlineto{\pgfpoint{-19.719995\du}{21.317033\du}}
\pgfpathlineto{\pgfpoint{-19.730444\du}{21.322705\du}}
\pgfpathlineto{\pgfpoint{-19.740894\du}{21.329273\du}}
\pgfpathlineto{\pgfpoint{-19.750447\du}{21.335841\du}}
\pgfpathlineto{\pgfpoint{-19.759105\du}{21.343007\du}}
\pgfpathlineto{\pgfpoint{-19.768659\du}{21.350172\du}}
\pgfpathlineto{\pgfpoint{-19.776421\du}{21.357934\du}}
\pgfpathlineto{\pgfpoint{-19.783885\du}{21.365995\du}}
\pgfpathlineto{\pgfpoint{-19.791050\du}{21.374355\du}}
\pgfusepath{stroke}
\pgfsetbuttcap
\pgfsetmiterjoin
\pgfsetdash{}{0pt}
\definecolor{dialinecolor}{rgb}{0.000000, 0.705882, 1.000000}
\pgfsetfillcolor{dialinecolor}
\pgfpathmoveto{\pgfpoint{-21.256941\du}{21.256128\du}}
\pgfpathlineto{\pgfpoint{-21.054224\du}{21.065353\du}}
\pgfpathlineto{\pgfpoint{-19.542655\du}{21.065353\du}}
\pgfpathlineto{\pgfpoint{-19.745372\du}{21.256128\du}}
\pgfpathlineto{\pgfpoint{-21.256941\du}{21.256128\du}}
\pgfusepath{fill}
\pgfsetlinewidth{0.001000\du}
\pgfsetbuttcap
\pgfsetmiterjoin
\pgfsetdash{}{0pt}
\definecolor{dialinecolor}{rgb}{0.666667, 0.901961, 1.000000}
\pgfsetstrokecolor{dialinecolor}
\pgfpathmoveto{\pgfpoint{-21.256941\du}{21.256128\du}}
\pgfpathlineto{\pgfpoint{-21.054224\du}{21.065353\du}}
\pgfpathlineto{\pgfpoint{-19.542655\du}{21.065353\du}}
\pgfpathlineto{\pgfpoint{-19.745372\du}{21.256128\du}}
\pgfpathlineto{\pgfpoint{-21.256941\du}{21.256128\du}}
\pgfusepath{stroke}
\pgfsetbuttcap
\pgfsetmiterjoin
\pgfsetdash{}{0pt}
\definecolor{dialinecolor}{rgb}{0.000000, 0.588235, 0.831373}
\pgfsetfillcolor{dialinecolor}
\pgfpathmoveto{\pgfpoint{-21.256941\du}{21.256128\du}}
\pgfpathlineto{\pgfpoint{-19.745372\du}{21.256128\du}}
\pgfpathlineto{\pgfpoint{-19.745372\du}{21.542141\du}}
\pgfpathlineto{\pgfpoint{-21.256941\du}{21.542141\du}}
\pgfpathlineto{\pgfpoint{-21.256941\du}{21.256128\du}}
\pgfusepath{fill}
\pgfsetbuttcap
\pgfsetmiterjoin
\pgfsetdash{}{0pt}
\definecolor{dialinecolor}{rgb}{0.666667, 0.901961, 1.000000}
\pgfsetstrokecolor{dialinecolor}
\pgfpathmoveto{\pgfpoint{-21.256941\du}{21.256128\du}}
\pgfpathlineto{\pgfpoint{-19.745372\du}{21.256128\du}}
\pgfpathlineto{\pgfpoint{-19.745372\du}{21.542141\du}}
\pgfpathlineto{\pgfpoint{-21.256941\du}{21.542141\du}}
\pgfpathlineto{\pgfpoint{-21.256941\du}{21.256128\du}}
\pgfusepath{stroke}
\pgfsetbuttcap
\pgfsetmiterjoin
\pgfsetdash{}{0pt}
\definecolor{dialinecolor}{rgb}{0.000000, 0.352941, 0.501961}
\pgfsetfillcolor{dialinecolor}
\pgfpathmoveto{\pgfpoint{-19.745372\du}{21.542141\du}}
\pgfpathlineto{\pgfpoint{-19.542655\du}{21.339424\du}}
\pgfpathlineto{\pgfpoint{-19.542655\du}{21.065353\du}}
\pgfpathlineto{\pgfpoint{-19.745372\du}{21.256128\du}}
\pgfpathlineto{\pgfpoint{-19.745372\du}{21.542141\du}}
\pgfusepath{fill}
\pgfsetbuttcap
\pgfsetmiterjoin
\pgfsetdash{}{0pt}
\definecolor{dialinecolor}{rgb}{0.666667, 0.901961, 1.000000}
\pgfsetstrokecolor{dialinecolor}
\pgfpathmoveto{\pgfpoint{-19.745372\du}{21.542141\du}}
\pgfpathlineto{\pgfpoint{-19.542655\du}{21.339424\du}}
\pgfpathlineto{\pgfpoint{-19.542655\du}{21.065353\du}}
\pgfpathlineto{\pgfpoint{-19.745372\du}{21.256128\du}}
\pgfpathlineto{\pgfpoint{-19.745372\du}{21.542141\du}}
\pgfusepath{stroke}
\pgfsetbuttcap
\pgfsetmiterjoin
\pgfsetdash{}{0pt}
\definecolor{dialinecolor}{rgb}{0.000000, 0.000000, 0.000000}
\pgfsetfillcolor{dialinecolor}
\pgfpathmoveto{\pgfpoint{-21.208874\du}{21.220003\du}}
\pgfpathlineto{\pgfpoint{-21.054224\du}{21.065353\du}}
\pgfpathlineto{\pgfpoint{-19.578780\du}{21.065353\du}}
\pgfpathlineto{\pgfpoint{-19.733131\du}{21.220003\du}}
\pgfpathlineto{\pgfpoint{-21.208874\du}{21.220003\du}}
\pgfusepath{fill}
\pgfsetbuttcap
\pgfsetmiterjoin
\pgfsetdash{}{0pt}
\definecolor{dialinecolor}{rgb}{0.000000, 0.000000, 0.000000}
\pgfsetstrokecolor{dialinecolor}
\pgfpathmoveto{\pgfpoint{-21.208874\du}{21.220003\du}}
\pgfpathlineto{\pgfpoint{-21.054224\du}{21.065353\du}}
\pgfpathlineto{\pgfpoint{-19.578780\du}{21.065353\du}}
\pgfpathlineto{\pgfpoint{-19.733131\du}{21.220003\du}}
\pgfpathlineto{\pgfpoint{-21.208874\du}{21.220003\du}}
\pgfusepath{stroke}
\pgfsetbuttcap
\pgfsetmiterjoin
\pgfsetdash{}{0pt}
\definecolor{dialinecolor}{rgb}{0.000000, 0.705882, 1.000000}
\pgfsetfillcolor{dialinecolor}
\pgfpathmoveto{\pgfpoint{-21.256941\du}{19.993850\du}}
\pgfpathlineto{\pgfpoint{-21.089752\du}{19.839200\du}}
\pgfpathlineto{\pgfpoint{-19.578780\du}{19.839200\du}}
\pgfpathlineto{\pgfpoint{-19.745372\du}{19.993850\du}}
\pgfpathlineto{\pgfpoint{-21.256941\du}{19.993850\du}}
\pgfusepath{fill}
\pgfsetbuttcap
\pgfsetmiterjoin
\pgfsetdash{}{0pt}
\definecolor{dialinecolor}{rgb}{0.666667, 0.901961, 1.000000}
\pgfsetstrokecolor{dialinecolor}
\pgfpathmoveto{\pgfpoint{-21.256941\du}{19.993850\du}}
\pgfpathlineto{\pgfpoint{-21.089752\du}{19.839200\du}}
\pgfpathlineto{\pgfpoint{-19.578780\du}{19.839200\du}}
\pgfpathlineto{\pgfpoint{-19.745372\du}{19.993850\du}}
\pgfpathlineto{\pgfpoint{-21.256941\du}{19.993850\du}}
\pgfusepath{stroke}
\pgfsetbuttcap
\pgfsetmiterjoin
\pgfsetdash{}{0pt}
\definecolor{dialinecolor}{rgb}{0.000000, 0.588235, 0.831373}
\pgfsetfillcolor{dialinecolor}
\pgfpathmoveto{\pgfpoint{-21.256941\du}{19.993850\du}}
\pgfpathlineto{\pgfpoint{-19.733131\du}{19.993850\du}}
\pgfpathlineto{\pgfpoint{-19.733131\du}{21.196119\du}}
\pgfpathlineto{\pgfpoint{-21.256941\du}{21.196119\du}}
\pgfpathlineto{\pgfpoint{-21.256941\du}{19.993850\du}}
\pgfusepath{fill}
\pgfsetbuttcap
\pgfsetmiterjoin
\pgfsetdash{}{0pt}
\definecolor{dialinecolor}{rgb}{0.666667, 0.901961, 1.000000}
\pgfsetstrokecolor{dialinecolor}
\pgfpathmoveto{\pgfpoint{-21.256941\du}{19.993850\du}}
\pgfpathlineto{\pgfpoint{-19.733430\du}{19.993850\du}}
\pgfpathlineto{\pgfpoint{-19.733430\du}{21.195820\du}}
\pgfpathlineto{\pgfpoint{-21.256941\du}{21.195820\du}}
\pgfpathlineto{\pgfpoint{-21.256941\du}{19.993850\du}}
\pgfusepath{stroke}
\pgfsetbuttcap
\pgfsetmiterjoin
\pgfsetdash{}{0pt}
\definecolor{dialinecolor}{rgb}{1.000000, 1.000000, 1.000000}
\pgfsetfillcolor{dialinecolor}
\pgfpathmoveto{\pgfpoint{-21.125876\du}{20.148798\du}}
\pgfpathlineto{\pgfpoint{-19.864196\du}{20.148798\du}}
\pgfpathlineto{\pgfpoint{-19.864196\du}{21.077594\du}}
\pgfpathlineto{\pgfpoint{-21.125876\du}{21.077594\du}}
\pgfpathlineto{\pgfpoint{-21.125876\du}{20.148798\du}}
\pgfusepath{fill}
\pgfsetbuttcap
\pgfsetmiterjoin
\pgfsetdash{}{0pt}
\definecolor{dialinecolor}{rgb}{0.000000, 0.235294, 0.333333}
\pgfsetstrokecolor{dialinecolor}
\pgfpathmoveto{\pgfpoint{-21.125876\du}{20.148798\du}}
\pgfpathlineto{\pgfpoint{-19.864494\du}{20.148798\du}}
\pgfpathlineto{\pgfpoint{-19.864494\du}{21.077295\du}}
\pgfpathlineto{\pgfpoint{-21.125876\du}{21.077295\du}}
\pgfpathlineto{\pgfpoint{-21.125876\du}{20.148798\du}}
\pgfusepath{stroke}
\pgfsetbuttcap
\pgfsetmiterjoin
\pgfsetdash{}{0pt}
\definecolor{dialinecolor}{rgb}{0.000000, 0.352941, 0.501961}
\pgfsetfillcolor{dialinecolor}
\pgfpathmoveto{\pgfpoint{-19.745372\du}{21.184774\du}}
\pgfpathlineto{\pgfpoint{-19.578780\du}{21.017883\du}}
\pgfpathlineto{\pgfpoint{-19.578780\du}{19.839200\du}}
\pgfpathlineto{\pgfpoint{-19.745372\du}{19.993850\du}}
\pgfpathlineto{\pgfpoint{-19.745372\du}{21.184774\du}}
\pgfusepath{fill}
\pgfsetbuttcap
\pgfsetmiterjoin
\pgfsetdash{}{0pt}
\definecolor{dialinecolor}{rgb}{0.666667, 0.901961, 1.000000}
\pgfsetstrokecolor{dialinecolor}
\pgfpathmoveto{\pgfpoint{-19.745372\du}{21.184774\du}}
\pgfpathlineto{\pgfpoint{-19.578780\du}{21.017883\du}}
\pgfpathlineto{\pgfpoint{-19.578780\du}{19.839200\du}}
\pgfpathlineto{\pgfpoint{-19.745372\du}{19.993850\du}}
\pgfpathlineto{\pgfpoint{-19.745372\du}{21.184774\du}}
\pgfusepath{stroke}
\pgfsetbuttcap
\pgfsetmiterjoin
\pgfsetdash{}{0pt}
\definecolor{dialinecolor}{rgb}{0.000000, 0.705882, 1.000000}
\pgfsetfillcolor{dialinecolor}
\pgfpathmoveto{\pgfpoint{-21.554000\du}{21.780087\du}}
\pgfpathlineto{\pgfpoint{-21.316353\du}{21.482132\du}}
\pgfpathlineto{\pgfpoint{-19.673719\du}{21.482132\du}}
\pgfpathlineto{\pgfpoint{-19.911665\du}{21.780087\du}}
\pgfpathlineto{\pgfpoint{-21.554000\du}{21.780087\du}}
\pgfusepath{fill}
\pgfsetbuttcap
\pgfsetmiterjoin
\pgfsetdash{}{0pt}
\definecolor{dialinecolor}{rgb}{0.666667, 0.901961, 1.000000}
\pgfsetstrokecolor{dialinecolor}
\pgfpathmoveto{\pgfpoint{-21.554000\du}{21.780087\du}}
\pgfpathlineto{\pgfpoint{-21.316353\du}{21.482132\du}}
\pgfpathlineto{\pgfpoint{-19.673719\du}{21.482132\du}}
\pgfpathlineto{\pgfpoint{-19.911665\du}{21.780087\du}}
\pgfpathlineto{\pgfpoint{-21.554000\du}{21.780087\du}}
\pgfusepath{stroke}
\pgfsetbuttcap
\pgfsetmiterjoin
\pgfsetdash{}{0pt}
\definecolor{dialinecolor}{rgb}{0.000000, 0.352941, 0.501961}
\pgfsetfillcolor{dialinecolor}
\pgfpathmoveto{\pgfpoint{-19.911665\du}{21.839200\du}}
\pgfpathlineto{\pgfpoint{-19.673719\du}{21.589312\du}}
\pgfpathlineto{\pgfpoint{-19.673719\du}{21.482132\du}}
\pgfpathlineto{\pgfpoint{-19.911665\du}{21.780087\du}}
\pgfpathlineto{\pgfpoint{-19.911665\du}{21.839200\du}}
\pgfusepath{fill}
\pgfsetbuttcap
\pgfsetmiterjoin
\pgfsetdash{}{0pt}
\definecolor{dialinecolor}{rgb}{0.666667, 0.901961, 1.000000}
\pgfsetstrokecolor{dialinecolor}
\pgfpathmoveto{\pgfpoint{-19.911665\du}{21.839200\du}}
\pgfpathlineto{\pgfpoint{-19.673719\du}{21.589312\du}}
\pgfpathlineto{\pgfpoint{-19.673719\du}{21.482132\du}}
\pgfpathlineto{\pgfpoint{-19.911665\du}{21.780087\du}}
\pgfpathlineto{\pgfpoint{-19.911665\du}{21.839200\du}}
\pgfusepath{stroke}
\pgfsetbuttcap
\pgfsetmiterjoin
\pgfsetdash{}{0pt}
\definecolor{dialinecolor}{rgb}{0.000000, 0.588235, 0.831373}
\pgfsetfillcolor{dialinecolor}
\pgfpathmoveto{\pgfpoint{-21.554000\du}{21.780087\du}}
\pgfpathlineto{\pgfpoint{-19.911665\du}{21.780087\du}}
\pgfpathlineto{\pgfpoint{-19.911665\du}{21.839200\du}}
\pgfpathlineto{\pgfpoint{-21.554000\du}{21.839200\du}}
\pgfpathlineto{\pgfpoint{-21.554000\du}{21.780087\du}}
\pgfusepath{fill}
\pgfsetbuttcap
\pgfsetmiterjoin
\pgfsetdash{}{0pt}
\definecolor{dialinecolor}{rgb}{0.666667, 0.901961, 1.000000}
\pgfsetstrokecolor{dialinecolor}
\pgfpathmoveto{\pgfpoint{-21.554000\du}{21.780087\du}}
\pgfpathlineto{\pgfpoint{-19.911665\du}{21.780087\du}}
\pgfpathlineto{\pgfpoint{-19.911665\du}{21.839200\du}}
\pgfpathlineto{\pgfpoint{-21.554000\du}{21.839200\du}}
\pgfpathlineto{\pgfpoint{-21.554000\du}{21.780087\du}}
\pgfusepath{stroke}
\pgfsetbuttcap
\pgfsetmiterjoin
\pgfsetdash{}{0pt}
\definecolor{dialinecolor}{rgb}{0.000000, 0.705882, 1.000000}
\pgfsetfillcolor{dialinecolor}
\pgfpathmoveto{\pgfpoint{-19.578780\du}{21.743962\du}}
\pgfpathlineto{\pgfpoint{-19.459658\du}{21.565726\du}}
\pgfpathlineto{\pgfpoint{-19.244999\du}{21.565726\du}}
\pgfpathlineto{\pgfpoint{-19.363822\du}{21.743962\du}}
\pgfpathlineto{\pgfpoint{-19.578780\du}{21.743962\du}}
\pgfusepath{fill}
\pgfsetbuttcap
\pgfsetmiterjoin
\pgfsetdash{}{0pt}
\definecolor{dialinecolor}{rgb}{0.666667, 0.901961, 1.000000}
\pgfsetstrokecolor{dialinecolor}
\pgfpathmoveto{\pgfpoint{-19.578780\du}{21.743962\du}}
\pgfpathlineto{\pgfpoint{-19.459658\du}{21.565726\du}}
\pgfpathlineto{\pgfpoint{-19.244999\du}{21.565726\du}}
\pgfpathlineto{\pgfpoint{-19.363822\du}{21.743962\du}}
\pgfpathlineto{\pgfpoint{-19.578780\du}{21.743962\du}}
\pgfusepath{stroke}
\pgfsetbuttcap
\pgfsetmiterjoin
\pgfsetdash{}{0pt}
\definecolor{dialinecolor}{rgb}{0.000000, 0.352941, 0.501961}
\pgfsetfillcolor{dialinecolor}
\pgfpathmoveto{\pgfpoint{-19.363822\du}{21.803672\du}}
\pgfpathlineto{\pgfpoint{-19.244999\du}{21.672608\du}}
\pgfpathlineto{\pgfpoint{-19.244999\du}{21.565726\du}}
\pgfpathlineto{\pgfpoint{-19.363822\du}{21.743962\du}}
\pgfpathlineto{\pgfpoint{-19.363822\du}{21.803672\du}}
\pgfusepath{fill}
\pgfsetbuttcap
\pgfsetmiterjoin
\pgfsetdash{}{0pt}
\definecolor{dialinecolor}{rgb}{0.666667, 0.901961, 1.000000}
\pgfsetstrokecolor{dialinecolor}
\pgfpathmoveto{\pgfpoint{-19.363822\du}{21.803672\du}}
\pgfpathlineto{\pgfpoint{-19.244999\du}{21.672608\du}}
\pgfpathlineto{\pgfpoint{-19.244999\du}{21.565726\du}}
\pgfpathlineto{\pgfpoint{-19.363822\du}{21.743962\du}}
\pgfpathlineto{\pgfpoint{-19.363822\du}{21.803672\du}}
\pgfusepath{stroke}
\pgfsetbuttcap
\pgfsetmiterjoin
\pgfsetdash{}{0pt}
\definecolor{dialinecolor}{rgb}{0.000000, 0.588235, 0.831373}
\pgfsetfillcolor{dialinecolor}
\pgfpathmoveto{\pgfpoint{-19.578780\du}{21.743962\du}}
\pgfpathlineto{\pgfpoint{-19.363822\du}{21.743962\du}}
\pgfpathlineto{\pgfpoint{-19.363822\du}{21.803672\du}}
\pgfpathlineto{\pgfpoint{-19.578780\du}{21.803672\du}}
\pgfpathlineto{\pgfpoint{-19.578780\du}{21.743962\du}}
\pgfusepath{fill}
\pgfsetbuttcap
\pgfsetmiterjoin
\pgfsetdash{}{0pt}
\definecolor{dialinecolor}{rgb}{0.666667, 0.901961, 1.000000}
\pgfsetstrokecolor{dialinecolor}
\pgfpathmoveto{\pgfpoint{-19.578780\du}{21.743962\du}}
\pgfpathlineto{\pgfpoint{-19.363822\du}{21.743962\du}}
\pgfpathlineto{\pgfpoint{-19.363822\du}{21.803672\du}}
\pgfpathlineto{\pgfpoint{-19.578780\du}{21.803672\du}}
\pgfpathlineto{\pgfpoint{-19.578780\du}{21.743962\du}}
\pgfusepath{stroke}
\pgfsetbuttcap
\pgfsetmiterjoin
\pgfsetdash{}{0pt}
\definecolor{dialinecolor}{rgb}{0.811765, 0.054902, 0.188235}
\pgfsetstrokecolor{dialinecolor}
\pgfpathmoveto{\pgfpoint{-20.792692\du}{20.375101\du}}
\pgfpathlineto{\pgfpoint{-20.458911\du}{20.375101\du}}
\pgfpathlineto{\pgfpoint{-20.542506\du}{20.458397\du}}
\pgfpathlineto{\pgfpoint{-20.256791\du}{20.458397\du}}
\pgfusepath{stroke}
\pgfsetbuttcap
\pgfsetmiterjoin
\pgfsetdash{}{0pt}
\definecolor{dialinecolor}{rgb}{0.811765, 0.054902, 0.188235}
\pgfsetstrokecolor{dialinecolor}
\pgfpathmoveto{\pgfpoint{-20.828519\du}{20.375101\du}}
\pgfpathlineto{\pgfpoint{-20.589976\du}{20.613047\du}}
\pgfpathlineto{\pgfpoint{-20.709396\du}{20.613047\du}}
\pgfpathlineto{\pgfpoint{-20.506978\du}{20.815764\du}}
\pgfusepath{stroke}
\pgfsetbuttcap
\pgfsetmiterjoin
\pgfsetdash{}{0pt}
\definecolor{dialinecolor}{rgb}{0.811765, 0.054902, 0.188235}
\pgfsetstrokecolor{dialinecolor}
\pgfpathmoveto{\pgfpoint{-20.150208\du}{20.375101\du}}
\pgfpathlineto{\pgfpoint{-20.387557\du}{20.613047\du}}
\pgfpathlineto{\pgfpoint{-20.269032\du}{20.613047\du}}
\pgfpathlineto{\pgfpoint{-20.471152\du}{20.815764\du}}
\pgfusepath{stroke}
\pgfsetbuttcap
\pgfsetmiterjoin
\pgfsetdash{}{0pt}
\definecolor{dialinecolor}{rgb}{0.000000, 0.470588, 0.666667}
\pgfsetfillcolor{dialinecolor}
\pgfpathmoveto{\pgfpoint{-20.661927\du}{20.464069\du}}
\pgfpathlineto{\pgfpoint{-20.662524\du}{20.460785\du}}
\pgfpathlineto{\pgfpoint{-20.663419\du}{20.457501\du}}
\pgfpathlineto{\pgfpoint{-20.664315\du}{20.454217\du}}
\pgfpathlineto{\pgfpoint{-20.666106\du}{20.450933\du}}
\pgfpathlineto{\pgfpoint{-20.668495\du}{20.447948\du}}
\pgfpathlineto{\pgfpoint{-20.671182\du}{20.444664\du}}
\pgfpathlineto{\pgfpoint{-20.674167\du}{20.441379\du}}
\pgfpathlineto{\pgfpoint{-20.677750\du}{20.438394\du}}
\pgfpathlineto{\pgfpoint{-20.681930\du}{20.435408\du}}
\pgfpathlineto{\pgfpoint{-20.686706\du}{20.432423\du}}
\pgfpathlineto{\pgfpoint{-20.691782\du}{20.429736\du}}
\pgfpathlineto{\pgfpoint{-20.697454\du}{20.426750\du}}
\pgfpathlineto{\pgfpoint{-20.703127\du}{20.424362\du}}
\pgfpathlineto{\pgfpoint{-20.709396\du}{20.421675\du}}
\pgfpathlineto{\pgfpoint{-20.715964\du}{20.419287\du}}
\pgfpathlineto{\pgfpoint{-20.723130\du}{20.416898\du}}
\pgfpathlineto{\pgfpoint{-20.730594\du}{20.414808\du}}
\pgfpathlineto{\pgfpoint{-20.738057\du}{20.412420\du}}
\pgfpathlineto{\pgfpoint{-20.746417\du}{20.410927\du}}
\pgfpathlineto{\pgfpoint{-20.754776\du}{20.408837\du}}
\pgfpathlineto{\pgfpoint{-20.762837\du}{20.407046\du}}
\pgfpathlineto{\pgfpoint{-20.772092\du}{20.405852\du}}
\pgfpathlineto{\pgfpoint{-20.780750\du}{20.404359\du}}
\pgfpathlineto{\pgfpoint{-20.790304\du}{20.403165\du}}
\pgfpathlineto{\pgfpoint{-20.799559\du}{20.401971\du}}
\pgfpathlineto{\pgfpoint{-20.808814\du}{20.401075\du}}
\pgfpathlineto{\pgfpoint{-20.818666\du}{20.400179\du}}
\pgfpathlineto{\pgfpoint{-20.828519\du}{20.399582\du}}
\pgfpathlineto{\pgfpoint{-20.838371\du}{20.399284\du}}
\pgfpathlineto{\pgfpoint{-20.848522\du}{20.398985\du}}
\pgfpathlineto{\pgfpoint{-20.858374\du}{20.398985\du}}
\pgfpathlineto{\pgfpoint{-20.858374\du}{20.398985\du}}
\pgfpathlineto{\pgfpoint{-20.868226\du}{20.398985\du}}
\pgfpathlineto{\pgfpoint{-20.878078\du}{20.399284\du}}
\pgfpathlineto{\pgfpoint{-20.888229\du}{20.399582\du}}
\pgfpathlineto{\pgfpoint{-20.898081\du}{20.400179\du}}
\pgfpathlineto{\pgfpoint{-20.907635\du}{20.401075\du}}
\pgfpathlineto{\pgfpoint{-20.917189\du}{20.401971\du}}
\pgfpathlineto{\pgfpoint{-20.926444\du}{20.403165\du}}
\pgfpathlineto{\pgfpoint{-20.935699\du}{20.404359\du}}
\pgfpathlineto{\pgfpoint{-20.944655\du}{20.405852\du}}
\pgfpathlineto{\pgfpoint{-20.953612\du}{20.407046\du}}
\pgfpathlineto{\pgfpoint{-20.961971\du}{20.408837\du}}
\pgfpathlineto{\pgfpoint{-20.970331\du}{20.410927\du}}
\pgfpathlineto{\pgfpoint{-20.978690\du}{20.412420\du}}
\pgfpathlineto{\pgfpoint{-20.985856\du}{20.414808\du}}
\pgfpathlineto{\pgfpoint{-20.993319\du}{20.416898\du}}
\pgfpathlineto{\pgfpoint{-21.000485\du}{20.419287\du}}
\pgfpathlineto{\pgfpoint{-21.007351\du}{20.421675\du}}
\pgfpathlineto{\pgfpoint{-21.013322\du}{20.424362\du}}
\pgfpathlineto{\pgfpoint{-21.019293\du}{20.426750\du}}
\pgfpathlineto{\pgfpoint{-21.024966\du}{20.429736\du}}
\pgfpathlineto{\pgfpoint{-21.030041\du}{20.432423\du}}
\pgfpathlineto{\pgfpoint{-21.034818\du}{20.435408\du}}
\pgfpathlineto{\pgfpoint{-21.038998\du}{20.438394\du}}
\pgfpathlineto{\pgfpoint{-21.042282\du}{20.441379\du}}
\pgfpathlineto{\pgfpoint{-21.045566\du}{20.444664\du}}
\pgfpathlineto{\pgfpoint{-21.048551\du}{20.447948\du}}
\pgfpathlineto{\pgfpoint{-21.050641\du}{20.450933\du}}
\pgfpathlineto{\pgfpoint{-21.052134\du}{20.454217\du}}
\pgfpathlineto{\pgfpoint{-21.053925\du}{20.457501\du}}
\pgfpathlineto{\pgfpoint{-21.054224\du}{20.460785\du}}
\pgfpathlineto{\pgfpoint{-21.054522\du}{20.464069\du}}
\pgfpathlineto{\pgfpoint{-21.054522\du}{20.464069\du}}
\pgfpathlineto{\pgfpoint{-21.054224\du}{20.467652\du}}
\pgfpathlineto{\pgfpoint{-21.053925\du}{20.470638\du}}
\pgfpathlineto{\pgfpoint{-21.052134\du}{20.473922\du}}
\pgfpathlineto{\pgfpoint{-21.050641\du}{20.477206\du}}
\pgfpathlineto{\pgfpoint{-21.048551\du}{20.480490\du}}
\pgfpathlineto{\pgfpoint{-21.045566\du}{20.483774\du}}
\pgfpathlineto{\pgfpoint{-21.042282\du}{20.486759\du}}
\pgfpathlineto{\pgfpoint{-21.038998\du}{20.489745\du}}
\pgfpathlineto{\pgfpoint{-21.034818\du}{20.492730\du}}
\pgfpathlineto{\pgfpoint{-21.030041\du}{20.496014\du}}
\pgfpathlineto{\pgfpoint{-21.024966\du}{20.498701\du}}
\pgfpathlineto{\pgfpoint{-21.019293\du}{20.501388\du}}
\pgfpathlineto{\pgfpoint{-21.013322\du}{20.504075\du}}
\pgfpathlineto{\pgfpoint{-21.007351\du}{20.506762\du}}
\pgfpathlineto{\pgfpoint{-21.000485\du}{20.509151\du}}
\pgfpathlineto{\pgfpoint{-20.993319\du}{20.511539\du}}
\pgfpathlineto{\pgfpoint{-20.985856\du}{20.513330\du}}
\pgfpathlineto{\pgfpoint{-20.978690\du}{20.515420\du}}
\pgfpathlineto{\pgfpoint{-20.970331\du}{20.517510\du}}
\pgfpathlineto{\pgfpoint{-20.961971\du}{20.519600\du}}
\pgfpathlineto{\pgfpoint{-20.953612\du}{20.521093\du}}
\pgfpathlineto{\pgfpoint{-20.944655\du}{20.522586\du}}
\pgfpathlineto{\pgfpoint{-20.935699\du}{20.524078\du}}
\pgfpathlineto{\pgfpoint{-20.926444\du}{20.525273\du}}
\pgfpathlineto{\pgfpoint{-20.917189\du}{20.526168\du}}
\pgfpathlineto{\pgfpoint{-20.907635\du}{20.527362\du}}
\pgfpathlineto{\pgfpoint{-20.898081\du}{20.527960\du}}
\pgfpathlineto{\pgfpoint{-20.888229\du}{20.528557\du}}
\pgfpathlineto{\pgfpoint{-20.878078\du}{20.529154\du}}
\pgfpathlineto{\pgfpoint{-20.868226\du}{20.529452\du}}
\pgfpathlineto{\pgfpoint{-20.858374\du}{20.529452\du}}
\pgfpathlineto{\pgfpoint{-20.858374\du}{20.529452\du}}
\pgfpathlineto{\pgfpoint{-20.848522\du}{20.529452\du}}
\pgfpathlineto{\pgfpoint{-20.838371\du}{20.529154\du}}
\pgfpathlineto{\pgfpoint{-20.828519\du}{20.528557\du}}
\pgfpathlineto{\pgfpoint{-20.818666\du}{20.527960\du}}
\pgfpathlineto{\pgfpoint{-20.808814\du}{20.527362\du}}
\pgfpathlineto{\pgfpoint{-20.799559\du}{20.526168\du}}
\pgfpathlineto{\pgfpoint{-20.790304\du}{20.525273\du}}
\pgfpathlineto{\pgfpoint{-20.780750\du}{20.524078\du}}
\pgfpathlineto{\pgfpoint{-20.772092\du}{20.522586\du}}
\pgfpathlineto{\pgfpoint{-20.762837\du}{20.521093\du}}
\pgfpathlineto{\pgfpoint{-20.754776\du}{20.519600\du}}
\pgfpathlineto{\pgfpoint{-20.746417\du}{20.517510\du}}
\pgfpathlineto{\pgfpoint{-20.738057\du}{20.515420\du}}
\pgfpathlineto{\pgfpoint{-20.730594\du}{20.513330\du}}
\pgfpathlineto{\pgfpoint{-20.723130\du}{20.511539\du}}
\pgfpathlineto{\pgfpoint{-20.715964\du}{20.509151\du}}
\pgfpathlineto{\pgfpoint{-20.709396\du}{20.506762\du}}
\pgfpathlineto{\pgfpoint{-20.703127\du}{20.504075\du}}
\pgfpathlineto{\pgfpoint{-20.697454\du}{20.501388\du}}
\pgfpathlineto{\pgfpoint{-20.691782\du}{20.498701\du}}
\pgfpathlineto{\pgfpoint{-20.686706\du}{20.496014\du}}
\pgfpathlineto{\pgfpoint{-20.681930\du}{20.492730\du}}
\pgfpathlineto{\pgfpoint{-20.677750\du}{20.489745\du}}
\pgfpathlineto{\pgfpoint{-20.674167\du}{20.486759\du}}
\pgfpathlineto{\pgfpoint{-20.671182\du}{20.483774\du}}
\pgfpathlineto{\pgfpoint{-20.668495\du}{20.480490\du}}
\pgfpathlineto{\pgfpoint{-20.666106\du}{20.477206\du}}
\pgfpathlineto{\pgfpoint{-20.664315\du}{20.473922\du}}
\pgfpathlineto{\pgfpoint{-20.663419\du}{20.470638\du}}
\pgfpathlineto{\pgfpoint{-20.662524\du}{20.467652\du}}
\pgfpathlineto{\pgfpoint{-20.661927\du}{20.464069\du}}
\pgfusepath{fill}
\pgfsetbuttcap
\pgfsetmiterjoin
\pgfsetdash{}{0pt}
\definecolor{dialinecolor}{rgb}{0.666667, 0.901961, 1.000000}
\pgfsetstrokecolor{dialinecolor}
\pgfpathmoveto{\pgfpoint{-20.674167\du}{20.458397\du}}
\pgfpathlineto{\pgfpoint{-20.674167\du}{20.454814\du}}
\pgfpathlineto{\pgfpoint{-20.675063\du}{20.452426\du}}
\pgfpathlineto{\pgfpoint{-20.676257\du}{20.449142\du}}
\pgfpathlineto{\pgfpoint{-20.677750\du}{20.446156\du}}
\pgfpathlineto{\pgfpoint{-20.680138\du}{20.443469\du}}
\pgfpathlineto{\pgfpoint{-20.682527\du}{20.439887\du}}
\pgfpathlineto{\pgfpoint{-20.685811\du}{20.437498\du}}
\pgfpathlineto{\pgfpoint{-20.689095\du}{20.434811\du}}
\pgfpathlineto{\pgfpoint{-20.693275\du}{20.432124\du}}
\pgfpathlineto{\pgfpoint{-20.697753\du}{20.429736\du}}
\pgfpathlineto{\pgfpoint{-20.702530\du}{20.426750\du}}
\pgfpathlineto{\pgfpoint{-20.708202\du}{20.424661\du}}
\pgfpathlineto{\pgfpoint{-20.713576\du}{20.421974\du}}
\pgfpathlineto{\pgfpoint{-20.719846\du}{20.419287\du}}
\pgfpathlineto{\pgfpoint{-20.725817\du}{20.417495\du}}
\pgfpathlineto{\pgfpoint{-20.733280\du}{20.415107\du}}
\pgfpathlineto{\pgfpoint{-20.740147\du}{20.413614\du}}
\pgfpathlineto{\pgfpoint{-20.747312\du}{20.411524\du}}
\pgfpathlineto{\pgfpoint{-20.755075\du}{20.409733\du}}
\pgfpathlineto{\pgfpoint{-20.763733\du}{20.407942\du}}
\pgfpathlineto{\pgfpoint{-20.772092\du}{20.406747\du}}
\pgfpathlineto{\pgfpoint{-20.780452\du}{20.404359\du}}
\pgfpathlineto{\pgfpoint{-20.789408\du}{20.404060\du}}
\pgfpathlineto{\pgfpoint{-20.798066\du}{20.402568\du}}
\pgfpathlineto{\pgfpoint{-20.807023\du}{20.401672\du}}
\pgfpathlineto{\pgfpoint{-20.816278\du}{20.400776\du}}
\pgfpathlineto{\pgfpoint{-20.825832\du}{20.400179\du}}
\pgfpathlineto{\pgfpoint{-20.835385\du}{20.399881\du}}
\pgfpathlineto{\pgfpoint{-20.844342\du}{20.399284\du}}
\pgfpathlineto{\pgfpoint{-20.854493\du}{20.398985\du}}
\pgfpathlineto{\pgfpoint{-20.864046\du}{20.398985\du}}
\pgfpathlineto{\pgfpoint{-20.864046\du}{20.398985\du}}
\pgfpathlineto{\pgfpoint{-20.873301\du}{20.398985\du}}
\pgfpathlineto{\pgfpoint{-20.883154\du}{20.399284\du}}
\pgfpathlineto{\pgfpoint{-20.892707\du}{20.399881\du}}
\pgfpathlineto{\pgfpoint{-20.902261\du}{20.400179\du}}
\pgfpathlineto{\pgfpoint{-20.911516\du}{20.400776\du}}
\pgfpathlineto{\pgfpoint{-20.921368\du}{20.401672\du}}
\pgfpathlineto{\pgfpoint{-20.930026\du}{20.402568\du}}
\pgfpathlineto{\pgfpoint{-20.938983\du}{20.404060\du}}
\pgfpathlineto{\pgfpoint{-20.947342\du}{20.404359\du}}
\pgfpathlineto{\pgfpoint{-20.956299\du}{20.406747\du}}
\pgfpathlineto{\pgfpoint{-20.964658\du}{20.407942\du}}
\pgfpathlineto{\pgfpoint{-20.972719\du}{20.409733\du}}
\pgfpathlineto{\pgfpoint{-20.980482\du}{20.411524\du}}
\pgfpathlineto{\pgfpoint{-20.987945\du}{20.413614\du}}
\pgfpathlineto{\pgfpoint{-20.994812\du}{20.415107\du}}
\pgfpathlineto{\pgfpoint{-21.002276\du}{20.417495\du}}
\pgfpathlineto{\pgfpoint{-21.008247\du}{20.419287\du}}
\pgfpathlineto{\pgfpoint{-21.014516\du}{20.421974\du}}
\pgfpathlineto{\pgfpoint{-21.020189\du}{20.424661\du}}
\pgfpathlineto{\pgfpoint{-21.025563\du}{20.426750\du}}
\pgfpathlineto{\pgfpoint{-21.030340\du}{20.429736\du}}
\pgfpathlineto{\pgfpoint{-21.034818\du}{20.432124\du}}
\pgfpathlineto{\pgfpoint{-21.038998\du}{20.434811\du}}
\pgfpathlineto{\pgfpoint{-21.042282\du}{20.437498\du}}
\pgfpathlineto{\pgfpoint{-21.045864\du}{20.439887\du}}
\pgfpathlineto{\pgfpoint{-21.047656\du}{20.443469\du}}
\pgfpathlineto{\pgfpoint{-21.050044\du}{20.446156\du}}
\pgfpathlineto{\pgfpoint{-21.051835\du}{20.449142\du}}
\pgfpathlineto{\pgfpoint{-21.053328\du}{20.452426\du}}
\pgfpathlineto{\pgfpoint{-21.053925\du}{20.454814\du}}
\pgfpathlineto{\pgfpoint{-21.054224\du}{20.458397\du}}
\pgfpathlineto{\pgfpoint{-21.054224\du}{20.458397\du}}
\pgfpathlineto{\pgfpoint{-21.053925\du}{20.461382\du}}
\pgfpathlineto{\pgfpoint{-21.053328\du}{20.464069\du}}
\pgfpathlineto{\pgfpoint{-21.051835\du}{20.467055\du}}
\pgfpathlineto{\pgfpoint{-21.050044\du}{20.470638\du}}
\pgfpathlineto{\pgfpoint{-21.047656\du}{20.473324\du}}
\pgfpathlineto{\pgfpoint{-21.045864\du}{20.476310\du}}
\pgfpathlineto{\pgfpoint{-21.042282\du}{20.478997\du}}
\pgfpathlineto{\pgfpoint{-21.038998\du}{20.481684\du}}
\pgfpathlineto{\pgfpoint{-21.034818\du}{20.484371\du}}
\pgfpathlineto{\pgfpoint{-21.030340\du}{20.487058\du}}
\pgfpathlineto{\pgfpoint{-21.025563\du}{20.489745\du}}
\pgfpathlineto{\pgfpoint{-21.020189\du}{20.492133\du}}
\pgfpathlineto{\pgfpoint{-21.014516\du}{20.494820\du}}
\pgfpathlineto{\pgfpoint{-21.008247\du}{20.496910\du}}
\pgfpathlineto{\pgfpoint{-21.002276\du}{20.499299\du}}
\pgfpathlineto{\pgfpoint{-20.994812\du}{20.501687\du}}
\pgfpathlineto{\pgfpoint{-20.987945\du}{20.503478\du}}
\pgfpathlineto{\pgfpoint{-20.980482\du}{20.504971\du}}
\pgfpathlineto{\pgfpoint{-20.972719\du}{20.507061\du}}
\pgfpathlineto{\pgfpoint{-20.964658\du}{20.508852\du}}
\pgfpathlineto{\pgfpoint{-20.956299\du}{20.510046\du}}
\pgfpathlineto{\pgfpoint{-20.947342\du}{20.511838\du}}
\pgfpathlineto{\pgfpoint{-20.938983\du}{20.512733\du}}
\pgfpathlineto{\pgfpoint{-20.930026\du}{20.514226\du}}
\pgfpathlineto{\pgfpoint{-20.921368\du}{20.515122\du}}
\pgfpathlineto{\pgfpoint{-20.911516\du}{20.515420\du}}
\pgfpathlineto{\pgfpoint{-20.902261\du}{20.516615\du}}
\pgfpathlineto{\pgfpoint{-20.892707\du}{20.516913\du}}
\pgfpathlineto{\pgfpoint{-20.883154\du}{20.517212\du}}
\pgfpathlineto{\pgfpoint{-20.873301\du}{20.517809\du}}
\pgfpathlineto{\pgfpoint{-20.864046\du}{20.517809\du}}
\pgfpathlineto{\pgfpoint{-20.864046\du}{20.517809\du}}
\pgfpathlineto{\pgfpoint{-20.854493\du}{20.517809\du}}
\pgfpathlineto{\pgfpoint{-20.844342\du}{20.517212\du}}
\pgfpathlineto{\pgfpoint{-20.835385\du}{20.516913\du}}
\pgfpathlineto{\pgfpoint{-20.825832\du}{20.516615\du}}
\pgfpathlineto{\pgfpoint{-20.816278\du}{20.515420\du}}
\pgfpathlineto{\pgfpoint{-20.807023\du}{20.515122\du}}
\pgfpathlineto{\pgfpoint{-20.798066\du}{20.514226\du}}
\pgfpathlineto{\pgfpoint{-20.789408\du}{20.512733\du}}
\pgfpathlineto{\pgfpoint{-20.780452\du}{20.511838\du}}
\pgfpathlineto{\pgfpoint{-20.772092\du}{20.510046\du}}
\pgfpathlineto{\pgfpoint{-20.763733\du}{20.508852\du}}
\pgfpathlineto{\pgfpoint{-20.755075\du}{20.507061\du}}
\pgfpathlineto{\pgfpoint{-20.747312\du}{20.504971\du}}
\pgfpathlineto{\pgfpoint{-20.740147\du}{20.503478\du}}
\pgfpathlineto{\pgfpoint{-20.733280\du}{20.501687\du}}
\pgfpathlineto{\pgfpoint{-20.725817\du}{20.499299\du}}
\pgfpathlineto{\pgfpoint{-20.719846\du}{20.496910\du}}
\pgfpathlineto{\pgfpoint{-20.713576\du}{20.494820\du}}
\pgfpathlineto{\pgfpoint{-20.708202\du}{20.492133\du}}
\pgfpathlineto{\pgfpoint{-20.702530\du}{20.489745\du}}
\pgfpathlineto{\pgfpoint{-20.697753\du}{20.487058\du}}
\pgfpathlineto{\pgfpoint{-20.693275\du}{20.484371\du}}
\pgfpathlineto{\pgfpoint{-20.689095\du}{20.481684\du}}
\pgfpathlineto{\pgfpoint{-20.685811\du}{20.478997\du}}
\pgfpathlineto{\pgfpoint{-20.682527\du}{20.476310\du}}
\pgfpathlineto{\pgfpoint{-20.680138\du}{20.473324\du}}
\pgfpathlineto{\pgfpoint{-20.677750\du}{20.470638\du}}
\pgfpathlineto{\pgfpoint{-20.676257\du}{20.467055\du}}
\pgfpathlineto{\pgfpoint{-20.675063\du}{20.464069\du}}
\pgfpathlineto{\pgfpoint{-20.674167\du}{20.461382\du}}
\pgfpathlineto{\pgfpoint{-20.674167\du}{20.458397\du}}
\pgfusepath{stroke}
\pgfsetbuttcap
\pgfsetmiterjoin
\pgfsetdash{}{0pt}
\definecolor{dialinecolor}{rgb}{0.000000, 0.470588, 0.666667}
\pgfsetfillcolor{dialinecolor}
\pgfpathmoveto{\pgfpoint{-21.054224\du}{20.375101\du}}
\pgfpathlineto{\pgfpoint{-21.054224\du}{20.469742\du}}
\pgfpathlineto{\pgfpoint{-20.661628\du}{20.469742\du}}
\pgfpathlineto{\pgfpoint{-20.661628\du}{20.375101\du}}
\pgfpathlineto{\pgfpoint{-21.054224\du}{20.375101\du}}
\pgfusepath{fill}
\pgfsetbuttcap
\pgfsetmiterjoin
\pgfsetdash{}{0pt}
\definecolor{dialinecolor}{rgb}{0.000000, 0.705882, 1.000000}
\pgfsetfillcolor{dialinecolor}
\pgfpathmoveto{\pgfpoint{-20.661927\du}{20.368533\du}}
\pgfpathlineto{\pgfpoint{-20.662524\du}{20.365249\du}}
\pgfpathlineto{\pgfpoint{-20.663419\du}{20.361666\du}}
\pgfpathlineto{\pgfpoint{-20.664315\du}{20.358681\du}}
\pgfpathlineto{\pgfpoint{-20.666106\du}{20.355396\du}}
\pgfpathlineto{\pgfpoint{-20.668495\du}{20.352112\du}}
\pgfpathlineto{\pgfpoint{-20.671182\du}{20.348828\du}}
\pgfpathlineto{\pgfpoint{-20.674167\du}{20.345843\du}}
\pgfpathlineto{\pgfpoint{-20.677750\du}{20.342559\du}}
\pgfpathlineto{\pgfpoint{-20.681930\du}{20.339872\du}}
\pgfpathlineto{\pgfpoint{-20.686706\du}{20.336886\du}}
\pgfpathlineto{\pgfpoint{-20.691782\du}{20.333901\du}}
\pgfpathlineto{\pgfpoint{-20.697454\du}{20.331214\du}}
\pgfpathlineto{\pgfpoint{-20.703127\du}{20.327930\du}}
\pgfpathlineto{\pgfpoint{-20.709396\du}{20.326138\du}}
\pgfpathlineto{\pgfpoint{-20.715964\du}{20.323451\du}}
\pgfpathlineto{\pgfpoint{-20.723130\du}{20.321063\du}}
\pgfpathlineto{\pgfpoint{-20.730594\du}{20.318973\du}}
\pgfpathlineto{\pgfpoint{-20.738057\du}{20.316883\du}}
\pgfpathlineto{\pgfpoint{-20.746417\du}{20.314793\du}}
\pgfpathlineto{\pgfpoint{-20.754776\du}{20.312704\du}}
\pgfpathlineto{\pgfpoint{-20.762837\du}{20.311211\du}}
\pgfpathlineto{\pgfpoint{-20.772092\du}{20.309718\du}}
\pgfpathlineto{\pgfpoint{-20.780750\du}{20.308524\du}}
\pgfpathlineto{\pgfpoint{-20.790304\du}{20.307031\du}}
\pgfpathlineto{\pgfpoint{-20.799559\du}{20.306135\du}}
\pgfpathlineto{\pgfpoint{-20.808814\du}{20.304941\du}}
\pgfpathlineto{\pgfpoint{-20.818666\du}{20.304344\du}}
\pgfpathlineto{\pgfpoint{-20.828519\du}{20.303747\du}}
\pgfpathlineto{\pgfpoint{-20.838371\du}{20.303448\du}}
\pgfpathlineto{\pgfpoint{-20.848522\du}{20.303150\du}}
\pgfpathlineto{\pgfpoint{-20.858374\du}{20.303150\du}}
\pgfpathlineto{\pgfpoint{-20.858374\du}{20.303150\du}}
\pgfpathlineto{\pgfpoint{-20.868226\du}{20.303150\du}}
\pgfpathlineto{\pgfpoint{-20.878078\du}{20.303448\du}}
\pgfpathlineto{\pgfpoint{-20.888229\du}{20.303747\du}}
\pgfpathlineto{\pgfpoint{-20.898081\du}{20.304344\du}}
\pgfpathlineto{\pgfpoint{-20.907635\du}{20.304941\du}}
\pgfpathlineto{\pgfpoint{-20.917189\du}{20.306135\du}}
\pgfpathlineto{\pgfpoint{-20.926444\du}{20.307031\du}}
\pgfpathlineto{\pgfpoint{-20.935699\du}{20.308524\du}}
\pgfpathlineto{\pgfpoint{-20.944655\du}{20.309718\du}}
\pgfpathlineto{\pgfpoint{-20.953612\du}{20.311211\du}}
\pgfpathlineto{\pgfpoint{-20.961971\du}{20.312704\du}}
\pgfpathlineto{\pgfpoint{-20.970331\du}{20.314793\du}}
\pgfpathlineto{\pgfpoint{-20.978690\du}{20.316883\du}}
\pgfpathlineto{\pgfpoint{-20.985856\du}{20.318973\du}}
\pgfpathlineto{\pgfpoint{-20.993319\du}{20.321063\du}}
\pgfpathlineto{\pgfpoint{-21.000485\du}{20.323451\du}}
\pgfpathlineto{\pgfpoint{-21.007351\du}{20.326138\du}}
\pgfpathlineto{\pgfpoint{-21.013322\du}{20.327930\du}}
\pgfpathlineto{\pgfpoint{-21.019293\du}{20.331214\du}}
\pgfpathlineto{\pgfpoint{-21.024966\du}{20.333901\du}}
\pgfpathlineto{\pgfpoint{-21.030041\du}{20.336886\du}}
\pgfpathlineto{\pgfpoint{-21.034818\du}{20.339872\du}}
\pgfpathlineto{\pgfpoint{-21.038998\du}{20.342559\du}}
\pgfpathlineto{\pgfpoint{-21.042282\du}{20.345843\du}}
\pgfpathlineto{\pgfpoint{-21.045566\du}{20.348828\du}}
\pgfpathlineto{\pgfpoint{-21.048551\du}{20.352112\du}}
\pgfpathlineto{\pgfpoint{-21.050641\du}{20.355396\du}}
\pgfpathlineto{\pgfpoint{-21.052134\du}{20.358681\du}}
\pgfpathlineto{\pgfpoint{-21.053925\du}{20.361666\du}}
\pgfpathlineto{\pgfpoint{-21.054224\du}{20.365249\du}}
\pgfpathlineto{\pgfpoint{-21.054522\du}{20.368533\du}}
\pgfpathlineto{\pgfpoint{-21.054522\du}{20.368533\du}}
\pgfpathlineto{\pgfpoint{-21.054224\du}{20.371817\du}}
\pgfpathlineto{\pgfpoint{-21.053925\du}{20.375399\du}}
\pgfpathlineto{\pgfpoint{-21.052134\du}{20.378385\du}}
\pgfpathlineto{\pgfpoint{-21.050641\du}{20.381669\du}}
\pgfpathlineto{\pgfpoint{-21.048551\du}{20.385252\du}}
\pgfpathlineto{\pgfpoint{-21.045566\du}{20.388237\du}}
\pgfpathlineto{\pgfpoint{-21.042282\du}{20.391223\du}}
\pgfpathlineto{\pgfpoint{-21.038998\du}{20.394507\du}}
\pgfpathlineto{\pgfpoint{-21.034818\du}{20.396895\du}}
\pgfpathlineto{\pgfpoint{-21.030041\du}{20.400179\du}}
\pgfpathlineto{\pgfpoint{-21.024966\du}{20.403463\du}}
\pgfpathlineto{\pgfpoint{-21.019293\du}{20.406150\du}}
\pgfpathlineto{\pgfpoint{-21.013322\du}{20.408837\du}}
\pgfpathlineto{\pgfpoint{-21.007351\du}{20.410927\du}}
\pgfpathlineto{\pgfpoint{-21.000485\du}{20.413913\du}}
\pgfpathlineto{\pgfpoint{-20.993319\du}{20.416301\du}}
\pgfpathlineto{\pgfpoint{-20.985856\du}{20.418092\du}}
\pgfpathlineto{\pgfpoint{-20.978690\du}{20.420182\du}}
\pgfpathlineto{\pgfpoint{-20.970331\du}{20.422272\du}}
\pgfpathlineto{\pgfpoint{-20.961971\du}{20.424362\du}}
\pgfpathlineto{\pgfpoint{-20.953612\du}{20.425855\du}}
\pgfpathlineto{\pgfpoint{-20.944655\du}{20.427347\du}}
\pgfpathlineto{\pgfpoint{-20.935699\du}{20.428542\du}}
\pgfpathlineto{\pgfpoint{-20.926444\du}{20.430034\du}}
\pgfpathlineto{\pgfpoint{-20.917189\du}{20.430930\du}}
\pgfpathlineto{\pgfpoint{-20.907635\du}{20.432124\du}}
\pgfpathlineto{\pgfpoint{-20.898081\du}{20.432423\du}}
\pgfpathlineto{\pgfpoint{-20.888229\du}{20.433319\du}}
\pgfpathlineto{\pgfpoint{-20.878078\du}{20.433319\du}}
\pgfpathlineto{\pgfpoint{-20.868226\du}{20.434214\du}}
\pgfpathlineto{\pgfpoint{-20.858374\du}{20.434214\du}}
\pgfpathlineto{\pgfpoint{-20.858374\du}{20.434214\du}}
\pgfpathlineto{\pgfpoint{-20.848522\du}{20.434214\du}}
\pgfpathlineto{\pgfpoint{-20.838371\du}{20.433319\du}}
\pgfpathlineto{\pgfpoint{-20.828519\du}{20.433319\du}}
\pgfpathlineto{\pgfpoint{-20.818666\du}{20.432423\du}}
\pgfpathlineto{\pgfpoint{-20.808814\du}{20.432124\du}}
\pgfpathlineto{\pgfpoint{-20.799559\du}{20.430930\du}}
\pgfpathlineto{\pgfpoint{-20.790304\du}{20.430034\du}}
\pgfpathlineto{\pgfpoint{-20.780750\du}{20.428542\du}}
\pgfpathlineto{\pgfpoint{-20.772092\du}{20.427347\du}}
\pgfpathlineto{\pgfpoint{-20.762837\du}{20.425855\du}}
\pgfpathlineto{\pgfpoint{-20.754776\du}{20.424362\du}}
\pgfpathlineto{\pgfpoint{-20.746417\du}{20.422272\du}}
\pgfpathlineto{\pgfpoint{-20.738057\du}{20.420182\du}}
\pgfpathlineto{\pgfpoint{-20.730594\du}{20.418092\du}}
\pgfpathlineto{\pgfpoint{-20.723130\du}{20.416301\du}}
\pgfpathlineto{\pgfpoint{-20.715964\du}{20.413913\du}}
\pgfpathlineto{\pgfpoint{-20.709396\du}{20.410927\du}}
\pgfpathlineto{\pgfpoint{-20.703127\du}{20.408837\du}}
\pgfpathlineto{\pgfpoint{-20.697454\du}{20.406150\du}}
\pgfpathlineto{\pgfpoint{-20.691782\du}{20.403463\du}}
\pgfpathlineto{\pgfpoint{-20.686706\du}{20.400179\du}}
\pgfpathlineto{\pgfpoint{-20.681930\du}{20.396895\du}}
\pgfpathlineto{\pgfpoint{-20.677750\du}{20.394507\du}}
\pgfpathlineto{\pgfpoint{-20.674167\du}{20.391223\du}}
\pgfpathlineto{\pgfpoint{-20.671182\du}{20.388237\du}}
\pgfpathlineto{\pgfpoint{-20.668495\du}{20.385252\du}}
\pgfpathlineto{\pgfpoint{-20.666106\du}{20.381669\du}}
\pgfpathlineto{\pgfpoint{-20.664315\du}{20.378385\du}}
\pgfpathlineto{\pgfpoint{-20.663419\du}{20.375399\du}}
\pgfpathlineto{\pgfpoint{-20.662524\du}{20.371817\du}}
\pgfpathlineto{\pgfpoint{-20.661927\du}{20.368533\du}}
\pgfusepath{fill}
\pgfsetbuttcap
\pgfsetmiterjoin
\pgfsetdash{}{0pt}
\definecolor{dialinecolor}{rgb}{0.666667, 0.901961, 1.000000}
\pgfsetstrokecolor{dialinecolor}
\pgfpathmoveto{\pgfpoint{-20.674167\du}{20.362860\du}}
\pgfpathlineto{\pgfpoint{-20.674167\du}{20.359576\du}}
\pgfpathlineto{\pgfpoint{-20.675063\du}{20.357188\du}}
\pgfpathlineto{\pgfpoint{-20.676257\du}{20.353904\du}}
\pgfpathlineto{\pgfpoint{-20.677750\du}{20.350321\du}}
\pgfpathlineto{\pgfpoint{-20.680138\du}{20.347933\du}}
\pgfpathlineto{\pgfpoint{-20.682527\du}{20.344649\du}}
\pgfpathlineto{\pgfpoint{-20.685811\du}{20.342260\du}}
\pgfpathlineto{\pgfpoint{-20.689095\du}{20.339573\du}}
\pgfpathlineto{\pgfpoint{-20.693275\du}{20.336886\du}}
\pgfpathlineto{\pgfpoint{-20.697753\du}{20.334199\du}}
\pgfpathlineto{\pgfpoint{-20.702530\du}{20.331512\du}}
\pgfpathlineto{\pgfpoint{-20.708202\du}{20.328825\du}}
\pgfpathlineto{\pgfpoint{-20.713576\du}{20.326437\du}}
\pgfpathlineto{\pgfpoint{-20.719846\du}{20.324347\du}}
\pgfpathlineto{\pgfpoint{-20.725817\du}{20.321959\du}}
\pgfpathlineto{\pgfpoint{-20.733280\du}{20.319570\du}}
\pgfpathlineto{\pgfpoint{-20.740147\du}{20.317779\du}}
\pgfpathlineto{\pgfpoint{-20.747312\du}{20.316286\du}}
\pgfpathlineto{\pgfpoint{-20.755075\du}{20.313898\du}}
\pgfpathlineto{\pgfpoint{-20.763733\du}{20.312405\du}}
\pgfpathlineto{\pgfpoint{-20.772092\du}{20.311211\du}}
\pgfpathlineto{\pgfpoint{-20.780452\du}{20.309419\du}}
\pgfpathlineto{\pgfpoint{-20.789408\du}{20.308524\du}}
\pgfpathlineto{\pgfpoint{-20.798066\du}{20.307031\du}}
\pgfpathlineto{\pgfpoint{-20.807023\du}{20.306135\du}}
\pgfpathlineto{\pgfpoint{-20.816278\du}{20.305538\du}}
\pgfpathlineto{\pgfpoint{-20.825832\du}{20.304643\du}}
\pgfpathlineto{\pgfpoint{-20.835385\du}{20.304344\du}}
\pgfpathlineto{\pgfpoint{-20.844342\du}{20.304045\du}}
\pgfpathlineto{\pgfpoint{-20.854493\du}{20.303448\du}}
\pgfpathlineto{\pgfpoint{-20.864046\du}{20.303448\du}}
\pgfpathlineto{\pgfpoint{-20.864046\du}{20.303448\du}}
\pgfpathlineto{\pgfpoint{-20.873301\du}{20.303448\du}}
\pgfpathlineto{\pgfpoint{-20.883154\du}{20.304045\du}}
\pgfpathlineto{\pgfpoint{-20.892707\du}{20.304344\du}}
\pgfpathlineto{\pgfpoint{-20.902261\du}{20.304643\du}}
\pgfpathlineto{\pgfpoint{-20.911516\du}{20.305538\du}}
\pgfpathlineto{\pgfpoint{-20.921368\du}{20.306135\du}}
\pgfpathlineto{\pgfpoint{-20.930026\du}{20.307031\du}}
\pgfpathlineto{\pgfpoint{-20.938983\du}{20.308524\du}}
\pgfpathlineto{\pgfpoint{-20.947342\du}{20.309419\du}}
\pgfpathlineto{\pgfpoint{-20.956299\du}{20.311211\du}}
\pgfpathlineto{\pgfpoint{-20.964658\du}{20.312405\du}}
\pgfpathlineto{\pgfpoint{-20.972719\du}{20.313898\du}}
\pgfpathlineto{\pgfpoint{-20.980482\du}{20.316286\du}}
\pgfpathlineto{\pgfpoint{-20.987945\du}{20.317779\du}}
\pgfpathlineto{\pgfpoint{-20.994812\du}{20.319570\du}}
\pgfpathlineto{\pgfpoint{-21.002276\du}{20.321959\du}}
\pgfpathlineto{\pgfpoint{-21.008247\du}{20.324347\du}}
\pgfpathlineto{\pgfpoint{-21.014516\du}{20.326437\du}}
\pgfpathlineto{\pgfpoint{-21.020189\du}{20.328825\du}}
\pgfpathlineto{\pgfpoint{-21.025563\du}{20.331512\du}}
\pgfpathlineto{\pgfpoint{-21.030340\du}{20.334199\du}}
\pgfpathlineto{\pgfpoint{-21.034818\du}{20.336886\du}}
\pgfpathlineto{\pgfpoint{-21.038998\du}{20.339573\du}}
\pgfpathlineto{\pgfpoint{-21.042282\du}{20.342260\du}}
\pgfpathlineto{\pgfpoint{-21.045864\du}{20.344649\du}}
\pgfpathlineto{\pgfpoint{-21.047656\du}{20.347933\du}}
\pgfpathlineto{\pgfpoint{-21.050044\du}{20.350321\du}}
\pgfpathlineto{\pgfpoint{-21.051835\du}{20.353904\du}}
\pgfpathlineto{\pgfpoint{-21.053328\du}{20.357188\du}}
\pgfpathlineto{\pgfpoint{-21.053925\du}{20.359576\du}}
\pgfpathlineto{\pgfpoint{-21.054224\du}{20.362860\du}}
\pgfpathlineto{\pgfpoint{-21.054224\du}{20.362860\du}}
\pgfpathlineto{\pgfpoint{-21.053925\du}{20.366144\du}}
\pgfpathlineto{\pgfpoint{-21.053328\du}{20.368831\du}}
\pgfpathlineto{\pgfpoint{-21.051835\du}{20.371817\du}}
\pgfpathlineto{\pgfpoint{-21.050044\du}{20.375101\du}}
\pgfpathlineto{\pgfpoint{-21.047656\du}{20.377788\du}}
\pgfpathlineto{\pgfpoint{-21.045864\du}{20.381072\du}}
\pgfpathlineto{\pgfpoint{-21.042282\du}{20.383759\du}}
\pgfpathlineto{\pgfpoint{-21.038998\du}{20.386446\du}}
\pgfpathlineto{\pgfpoint{-21.034818\du}{20.389133\du}}
\pgfpathlineto{\pgfpoint{-21.030340\du}{20.391820\du}}
\pgfpathlineto{\pgfpoint{-21.025563\du}{20.394507\du}}
\pgfpathlineto{\pgfpoint{-21.020189\du}{20.396597\du}}
\pgfpathlineto{\pgfpoint{-21.014516\du}{20.399284\du}}
\pgfpathlineto{\pgfpoint{-21.008247\du}{20.401672\du}}
\pgfpathlineto{\pgfpoint{-21.002276\du}{20.404060\du}}
\pgfpathlineto{\pgfpoint{-20.994812\du}{20.406150\du}}
\pgfpathlineto{\pgfpoint{-20.987945\du}{20.407942\du}}
\pgfpathlineto{\pgfpoint{-20.980482\du}{20.409733\du}}
\pgfpathlineto{\pgfpoint{-20.972719\du}{20.411524\du}}
\pgfpathlineto{\pgfpoint{-20.964658\du}{20.413614\du}}
\pgfpathlineto{\pgfpoint{-20.956299\du}{20.414808\du}}
\pgfpathlineto{\pgfpoint{-20.947342\du}{20.416600\du}}
\pgfpathlineto{\pgfpoint{-20.938983\du}{20.417495\du}}
\pgfpathlineto{\pgfpoint{-20.930026\du}{20.418988\du}}
\pgfpathlineto{\pgfpoint{-20.921368\du}{20.419287\du}}
\pgfpathlineto{\pgfpoint{-20.911516\du}{20.420481\du}}
\pgfpathlineto{\pgfpoint{-20.902261\du}{20.420779\du}}
\pgfpathlineto{\pgfpoint{-20.892707\du}{20.421675\du}}
\pgfpathlineto{\pgfpoint{-20.883154\du}{20.421974\du}}
\pgfpathlineto{\pgfpoint{-20.873301\du}{20.422272\du}}
\pgfpathlineto{\pgfpoint{-20.864046\du}{20.422272\du}}
\pgfpathlineto{\pgfpoint{-20.864046\du}{20.422272\du}}
\pgfpathlineto{\pgfpoint{-20.854493\du}{20.422272\du}}
\pgfpathlineto{\pgfpoint{-20.844342\du}{20.421974\du}}
\pgfpathlineto{\pgfpoint{-20.835385\du}{20.421675\du}}
\pgfpathlineto{\pgfpoint{-20.825832\du}{20.420779\du}}
\pgfpathlineto{\pgfpoint{-20.816278\du}{20.420481\du}}
\pgfpathlineto{\pgfpoint{-20.807023\du}{20.419287\du}}
\pgfpathlineto{\pgfpoint{-20.798066\du}{20.418988\du}}
\pgfpathlineto{\pgfpoint{-20.789408\du}{20.417495\du}}
\pgfpathlineto{\pgfpoint{-20.780452\du}{20.416600\du}}
\pgfpathlineto{\pgfpoint{-20.772092\du}{20.414808\du}}
\pgfpathlineto{\pgfpoint{-20.763733\du}{20.413614\du}}
\pgfpathlineto{\pgfpoint{-20.755075\du}{20.411524\du}}
\pgfpathlineto{\pgfpoint{-20.747312\du}{20.409733\du}}
\pgfpathlineto{\pgfpoint{-20.740147\du}{20.407942\du}}
\pgfpathlineto{\pgfpoint{-20.733280\du}{20.406150\du}}
\pgfpathlineto{\pgfpoint{-20.725817\du}{20.404060\du}}
\pgfpathlineto{\pgfpoint{-20.719846\du}{20.401672\du}}
\pgfpathlineto{\pgfpoint{-20.713576\du}{20.399284\du}}
\pgfpathlineto{\pgfpoint{-20.708202\du}{20.396597\du}}
\pgfpathlineto{\pgfpoint{-20.702530\du}{20.394507\du}}
\pgfpathlineto{\pgfpoint{-20.697753\du}{20.391820\du}}
\pgfpathlineto{\pgfpoint{-20.693275\du}{20.389133\du}}
\pgfpathlineto{\pgfpoint{-20.689095\du}{20.386446\du}}
\pgfpathlineto{\pgfpoint{-20.685811\du}{20.383759\du}}
\pgfpathlineto{\pgfpoint{-20.682527\du}{20.381072\du}}
\pgfpathlineto{\pgfpoint{-20.680138\du}{20.377788\du}}
\pgfpathlineto{\pgfpoint{-20.677750\du}{20.375101\du}}
\pgfpathlineto{\pgfpoint{-20.676257\du}{20.371817\du}}
\pgfpathlineto{\pgfpoint{-20.675063\du}{20.368831\du}}
\pgfpathlineto{\pgfpoint{-20.674167\du}{20.366144\du}}
\pgfpathlineto{\pgfpoint{-20.674167\du}{20.362860\du}}
\pgfusepath{stroke}
\pgfsetbuttcap
\pgfsetmiterjoin
\pgfsetdash{}{0pt}
\definecolor{dialinecolor}{rgb}{0.000000, 0.000000, 0.000000}
\pgfsetfillcolor{dialinecolor}
\pgfpathmoveto{\pgfpoint{-20.851806\du}{20.351217\du}}
\pgfpathlineto{\pgfpoint{-20.828519\du}{20.362860\du}}
\pgfpathlineto{\pgfpoint{-20.756568\du}{20.339275\du}}
\pgfpathlineto{\pgfpoint{-20.721040\du}{20.351217\du}}
\pgfpathlineto{\pgfpoint{-20.744625\du}{20.327333\du}}
\pgfpathlineto{\pgfpoint{-20.828519\du}{20.327333\du}}
\pgfpathlineto{\pgfpoint{-20.792692\du}{20.327333\du}}
\pgfpathlineto{\pgfpoint{-20.851806\du}{20.351217\du}}
\pgfusepath{fill}
\pgfsetbuttcap
\pgfsetmiterjoin
\pgfsetdash{}{0pt}
\definecolor{dialinecolor}{rgb}{0.000000, 0.000000, 0.000000}
\pgfsetfillcolor{dialinecolor}
\pgfpathmoveto{\pgfpoint{-20.864046\du}{20.386744\du}}
\pgfpathlineto{\pgfpoint{-20.899873\du}{20.375101\du}}
\pgfpathlineto{\pgfpoint{-20.958986\du}{20.398985\du}}
\pgfpathlineto{\pgfpoint{-20.994812\du}{20.386744\du}}
\pgfpathlineto{\pgfpoint{-20.982571\du}{20.422272\du}}
\pgfpathlineto{\pgfpoint{-20.899873\du}{20.422272\du}}
\pgfpathlineto{\pgfpoint{-20.935400\du}{20.410927\du}}
\pgfpathlineto{\pgfpoint{-20.864046\du}{20.386744\du}}
\pgfusepath{fill}
\pgfsetbuttcap
\pgfsetmiterjoin
\pgfsetdash{}{0pt}
\definecolor{dialinecolor}{rgb}{0.000000, 0.000000, 0.000000}
\pgfsetfillcolor{dialinecolor}
\pgfpathmoveto{\pgfpoint{-20.982571\du}{20.327333\du}}
\pgfpathlineto{\pgfpoint{-20.958986\du}{20.314793\du}}
\pgfpathlineto{\pgfpoint{-20.887333\du}{20.351217\du}}
\pgfpathlineto{\pgfpoint{-20.864046\du}{20.339275\du}}
\pgfpathlineto{\pgfpoint{-20.875690\du}{20.362860\du}}
\pgfpathlineto{\pgfpoint{-20.958986\du}{20.362860\du}}
\pgfpathlineto{\pgfpoint{-20.923160\du}{20.351217\du}}
\pgfpathlineto{\pgfpoint{-20.982571\du}{20.327333\du}}
\pgfusepath{fill}
\pgfsetbuttcap
\pgfsetmiterjoin
\pgfsetdash{}{0pt}
\definecolor{dialinecolor}{rgb}{0.000000, 0.000000, 0.000000}
\pgfsetfillcolor{dialinecolor}
\pgfpathmoveto{\pgfpoint{-20.733280\du}{20.410927\du}}
\pgfpathlineto{\pgfpoint{-20.756568\du}{20.422272\du}}
\pgfpathlineto{\pgfpoint{-20.828519\du}{20.386744\du}}
\pgfpathlineto{\pgfpoint{-20.864046\du}{20.398985\du}}
\pgfpathlineto{\pgfpoint{-20.840162\du}{20.375101\du}}
\pgfpathlineto{\pgfpoint{-20.756568\du}{20.375101\du}}
\pgfpathlineto{\pgfpoint{-20.792692\du}{20.386744\du}}
\pgfpathlineto{\pgfpoint{-20.733280\du}{20.410927\du}}
\pgfusepath{fill}
\pgfsetbuttcap
\pgfsetmiterjoin
\pgfsetdash{}{0pt}
\definecolor{dialinecolor}{rgb}{1.000000, 1.000000, 1.000000}
\pgfsetfillcolor{dialinecolor}
\pgfpathmoveto{\pgfpoint{-20.851806\du}{20.362860\du}}
\pgfpathlineto{\pgfpoint{-20.828519\du}{20.362860\du}}
\pgfpathlineto{\pgfpoint{-20.756568\du}{20.339275\du}}
\pgfpathlineto{\pgfpoint{-20.721040\du}{20.351217\du}}
\pgfpathlineto{\pgfpoint{-20.744625\du}{20.327333\du}}
\pgfpathlineto{\pgfpoint{-20.816278\du}{20.327333\du}}
\pgfpathlineto{\pgfpoint{-20.792692\du}{20.327333\du}}
\pgfpathlineto{\pgfpoint{-20.851806\du}{20.362860\du}}
\pgfusepath{fill}
\pgfsetbuttcap
\pgfsetmiterjoin
\pgfsetdash{}{0pt}
\definecolor{dialinecolor}{rgb}{1.000000, 1.000000, 1.000000}
\pgfsetfillcolor{dialinecolor}
\pgfpathmoveto{\pgfpoint{-20.864046\du}{20.386744\du}}
\pgfpathlineto{\pgfpoint{-20.887333\du}{20.375101\du}}
\pgfpathlineto{\pgfpoint{-20.958986\du}{20.398985\du}}
\pgfpathlineto{\pgfpoint{-20.994812\du}{20.398985\du}}
\pgfpathlineto{\pgfpoint{-20.971226\du}{20.422272\du}}
\pgfpathlineto{\pgfpoint{-20.887333\du}{20.422272\du}}
\pgfpathlineto{\pgfpoint{-20.923160\du}{20.410927\du}}
\pgfpathlineto{\pgfpoint{-20.864046\du}{20.386744\du}}
\pgfusepath{fill}
\pgfsetbuttcap
\pgfsetmiterjoin
\pgfsetdash{}{0pt}
\definecolor{dialinecolor}{rgb}{1.000000, 1.000000, 1.000000}
\pgfsetfillcolor{dialinecolor}
\pgfpathmoveto{\pgfpoint{-20.982571\du}{20.327333\du}}
\pgfpathlineto{\pgfpoint{-20.958986\du}{20.327333\du}}
\pgfpathlineto{\pgfpoint{-20.887333\du}{20.351217\du}}
\pgfpathlineto{\pgfpoint{-20.851806\du}{20.339275\du}}
\pgfpathlineto{\pgfpoint{-20.875690\du}{20.362860\du}}
\pgfpathlineto{\pgfpoint{-20.958986\du}{20.362860\du}}
\pgfpathlineto{\pgfpoint{-20.923160\du}{20.362860\du}}
\pgfpathlineto{\pgfpoint{-20.982571\du}{20.327333\du}}
\pgfusepath{fill}
\pgfsetbuttcap
\pgfsetmiterjoin
\pgfsetdash{}{0pt}
\definecolor{dialinecolor}{rgb}{1.000000, 1.000000, 1.000000}
\pgfsetfillcolor{dialinecolor}
\pgfpathmoveto{\pgfpoint{-20.733280\du}{20.410927\du}}
\pgfpathlineto{\pgfpoint{-20.756568\du}{20.422272\du}}
\pgfpathlineto{\pgfpoint{-20.828519\du}{20.398985\du}}
\pgfpathlineto{\pgfpoint{-20.851806\du}{20.398985\du}}
\pgfpathlineto{\pgfpoint{-20.840162\du}{20.375101\du}}
\pgfpathlineto{\pgfpoint{-20.756568\du}{20.375101\du}}
\pgfpathlineto{\pgfpoint{-20.792692\du}{20.386744\du}}
\pgfpathlineto{\pgfpoint{-20.733280\du}{20.410927\du}}
\pgfusepath{fill}
\pgfsetbuttcap
\pgfsetmiterjoin
\pgfsetdash{}{0pt}
\definecolor{dialinecolor}{rgb}{0.666667, 0.901961, 1.000000}
\pgfsetstrokecolor{dialinecolor}
\pgfpathmoveto{\pgfpoint{-21.054224\du}{20.362860\du}}
\pgfpathlineto{\pgfpoint{-21.054224\du}{20.457800\du}}
\pgfusepath{stroke}
\pgfsetbuttcap
\pgfsetmiterjoin
\pgfsetdash{}{0pt}
\definecolor{dialinecolor}{rgb}{0.666667, 0.901961, 1.000000}
\pgfsetstrokecolor{dialinecolor}
\pgfpathmoveto{\pgfpoint{-20.673570\du}{20.362860\du}}
\pgfpathlineto{\pgfpoint{-20.673570\du}{20.457800\du}}
\pgfusepath{stroke}
\pgfsetbuttcap
\pgfsetmiterjoin
\pgfsetdash{}{0pt}
\definecolor{dialinecolor}{rgb}{0.000000, 0.470588, 0.666667}
\pgfsetfillcolor{dialinecolor}
\pgfpathmoveto{\pgfpoint{-19.947790\du}{20.464069\du}}
\pgfpathlineto{\pgfpoint{-19.947790\du}{20.460785\du}}
\pgfpathlineto{\pgfpoint{-19.948686\du}{20.457501\du}}
\pgfpathlineto{\pgfpoint{-19.949880\du}{20.454217\du}}
\pgfpathlineto{\pgfpoint{-19.951373\du}{20.450933\du}}
\pgfpathlineto{\pgfpoint{-19.953761\du}{20.447948\du}}
\pgfpathlineto{\pgfpoint{-19.956448\du}{20.444664\du}}
\pgfpathlineto{\pgfpoint{-19.959434\du}{20.441379\du}}
\pgfpathlineto{\pgfpoint{-19.963613\du}{20.438394\du}}
\pgfpathlineto{\pgfpoint{-19.967196\du}{20.435408\du}}
\pgfpathlineto{\pgfpoint{-19.971973\du}{20.432423\du}}
\pgfpathlineto{\pgfpoint{-19.977048\du}{20.429736\du}}
\pgfpathlineto{\pgfpoint{-19.982422\du}{20.426750\du}}
\pgfpathlineto{\pgfpoint{-19.988692\du}{20.424362\du}}
\pgfpathlineto{\pgfpoint{-19.994663\du}{20.421675\du}}
\pgfpathlineto{\pgfpoint{-20.001529\du}{20.419287\du}}
\pgfpathlineto{\pgfpoint{-20.008396\du}{20.416898\du}}
\pgfpathlineto{\pgfpoint{-20.015860\du}{20.414808\du}}
\pgfpathlineto{\pgfpoint{-20.023324\du}{20.412420\du}}
\pgfpathlineto{\pgfpoint{-20.031385\du}{20.410927\du}}
\pgfpathlineto{\pgfpoint{-20.040043\du}{20.408837\du}}
\pgfpathlineto{\pgfpoint{-20.048402\du}{20.407046\du}}
\pgfpathlineto{\pgfpoint{-20.057060\du}{20.405852\du}}
\pgfpathlineto{\pgfpoint{-20.066315\du}{20.404359\du}}
\pgfpathlineto{\pgfpoint{-20.075570\du}{20.403165\du}}
\pgfpathlineto{\pgfpoint{-20.084825\du}{20.401971\du}}
\pgfpathlineto{\pgfpoint{-20.094379\du}{20.401075\du}}
\pgfpathlineto{\pgfpoint{-20.104530\du}{20.400179\du}}
\pgfpathlineto{\pgfpoint{-20.114084\du}{20.399582\du}}
\pgfpathlineto{\pgfpoint{-20.123637\du}{20.399284\du}}
\pgfpathlineto{\pgfpoint{-20.133788\du}{20.398985\du}}
\pgfpathlineto{\pgfpoint{-20.143640\du}{20.398985\du}}
\pgfpathlineto{\pgfpoint{-20.143640\du}{20.398985\du}}
\pgfpathlineto{\pgfpoint{-20.153492\du}{20.398985\du}}
\pgfpathlineto{\pgfpoint{-20.163643\du}{20.399284\du}}
\pgfpathlineto{\pgfpoint{-20.173495\du}{20.399582\du}}
\pgfpathlineto{\pgfpoint{-20.183348\du}{20.400179\du}}
\pgfpathlineto{\pgfpoint{-20.192901\du}{20.401075\du}}
\pgfpathlineto{\pgfpoint{-20.202455\du}{20.401971\du}}
\pgfpathlineto{\pgfpoint{-20.212009\du}{20.403165\du}}
\pgfpathlineto{\pgfpoint{-20.220965\du}{20.404359\du}}
\pgfpathlineto{\pgfpoint{-20.230220\du}{20.405852\du}}
\pgfpathlineto{\pgfpoint{-20.238878\du}{20.407046\du}}
\pgfpathlineto{\pgfpoint{-20.247536\du}{20.408837\du}}
\pgfpathlineto{\pgfpoint{-20.255896\du}{20.410927\du}}
\pgfpathlineto{\pgfpoint{-20.263957\du}{20.412420\du}}
\pgfpathlineto{\pgfpoint{-20.271421\du}{20.414808\du}}
\pgfpathlineto{\pgfpoint{-20.278884\du}{20.416898\du}}
\pgfpathlineto{\pgfpoint{-20.286050\du}{20.419287\du}}
\pgfpathlineto{\pgfpoint{-20.292618\du}{20.421675\du}}
\pgfpathlineto{\pgfpoint{-20.298887\du}{20.424362\du}}
\pgfpathlineto{\pgfpoint{-20.304858\du}{20.426750\du}}
\pgfpathlineto{\pgfpoint{-20.310232\du}{20.429736\du}}
\pgfpathlineto{\pgfpoint{-20.315308\du}{20.432423\du}}
\pgfpathlineto{\pgfpoint{-20.320084\du}{20.435408\du}}
\pgfpathlineto{\pgfpoint{-20.323966\du}{20.438394\du}}
\pgfpathlineto{\pgfpoint{-20.327847\du}{20.441379\du}}
\pgfpathlineto{\pgfpoint{-20.330832\du}{20.444664\du}}
\pgfpathlineto{\pgfpoint{-20.333519\du}{20.447948\du}}
\pgfpathlineto{\pgfpoint{-20.335908\du}{20.450933\du}}
\pgfpathlineto{\pgfpoint{-20.337699\du}{20.454217\du}}
\pgfpathlineto{\pgfpoint{-20.339192\du}{20.457501\du}}
\pgfpathlineto{\pgfpoint{-20.340087\du}{20.460785\du}}
\pgfpathlineto{\pgfpoint{-20.340087\du}{20.464069\du}}
\pgfpathlineto{\pgfpoint{-20.340087\du}{20.464069\du}}
\pgfpathlineto{\pgfpoint{-20.340087\du}{20.467652\du}}
\pgfpathlineto{\pgfpoint{-20.339192\du}{20.470638\du}}
\pgfpathlineto{\pgfpoint{-20.337699\du}{20.473922\du}}
\pgfpathlineto{\pgfpoint{-20.335908\du}{20.477206\du}}
\pgfpathlineto{\pgfpoint{-20.333519\du}{20.480490\du}}
\pgfpathlineto{\pgfpoint{-20.330832\du}{20.483774\du}}
\pgfpathlineto{\pgfpoint{-20.327847\du}{20.486759\du}}
\pgfpathlineto{\pgfpoint{-20.323966\du}{20.489745\du}}
\pgfpathlineto{\pgfpoint{-20.320084\du}{20.492730\du}}
\pgfpathlineto{\pgfpoint{-20.315308\du}{20.496014\du}}
\pgfpathlineto{\pgfpoint{-20.310232\du}{20.498701\du}}
\pgfpathlineto{\pgfpoint{-20.304858\du}{20.501388\du}}
\pgfpathlineto{\pgfpoint{-20.298887\du}{20.504075\du}}
\pgfpathlineto{\pgfpoint{-20.292618\du}{20.506762\du}}
\pgfpathlineto{\pgfpoint{-20.286050\du}{20.509151\du}}
\pgfpathlineto{\pgfpoint{-20.278884\du}{20.511539\du}}
\pgfpathlineto{\pgfpoint{-20.271421\du}{20.513330\du}}
\pgfpathlineto{\pgfpoint{-20.263957\du}{20.515420\du}}
\pgfpathlineto{\pgfpoint{-20.255896\du}{20.517510\du}}
\pgfpathlineto{\pgfpoint{-20.247536\du}{20.519600\du}}
\pgfpathlineto{\pgfpoint{-20.238878\du}{20.521093\du}}
\pgfpathlineto{\pgfpoint{-20.230220\du}{20.522586\du}}
\pgfpathlineto{\pgfpoint{-20.220965\du}{20.524078\du}}
\pgfpathlineto{\pgfpoint{-20.212009\du}{20.525273\du}}
\pgfpathlineto{\pgfpoint{-20.202455\du}{20.526168\du}}
\pgfpathlineto{\pgfpoint{-20.192901\du}{20.527362\du}}
\pgfpathlineto{\pgfpoint{-20.183348\du}{20.527960\du}}
\pgfpathlineto{\pgfpoint{-20.173495\du}{20.528557\du}}
\pgfpathlineto{\pgfpoint{-20.163643\du}{20.529154\du}}
\pgfpathlineto{\pgfpoint{-20.153492\du}{20.529452\du}}
\pgfpathlineto{\pgfpoint{-20.143640\du}{20.529452\du}}
\pgfpathlineto{\pgfpoint{-20.143640\du}{20.529452\du}}
\pgfpathlineto{\pgfpoint{-20.133788\du}{20.529452\du}}
\pgfpathlineto{\pgfpoint{-20.123637\du}{20.529154\du}}
\pgfpathlineto{\pgfpoint{-20.114084\du}{20.528557\du}}
\pgfpathlineto{\pgfpoint{-20.104530\du}{20.527960\du}}
\pgfpathlineto{\pgfpoint{-20.094379\du}{20.527362\du}}
\pgfpathlineto{\pgfpoint{-20.084825\du}{20.526168\du}}
\pgfpathlineto{\pgfpoint{-20.075570\du}{20.525273\du}}
\pgfpathlineto{\pgfpoint{-20.066315\du}{20.524078\du}}
\pgfpathlineto{\pgfpoint{-20.057060\du}{20.522586\du}}
\pgfpathlineto{\pgfpoint{-20.048402\du}{20.521093\du}}
\pgfpathlineto{\pgfpoint{-20.040043\du}{20.519600\du}}
\pgfpathlineto{\pgfpoint{-20.031385\du}{20.517510\du}}
\pgfpathlineto{\pgfpoint{-20.023324\du}{20.515420\du}}
\pgfpathlineto{\pgfpoint{-20.015860\du}{20.513330\du}}
\pgfpathlineto{\pgfpoint{-20.008396\du}{20.511539\du}}
\pgfpathlineto{\pgfpoint{-20.001529\du}{20.509151\du}}
\pgfpathlineto{\pgfpoint{-19.994663\du}{20.506762\du}}
\pgfpathlineto{\pgfpoint{-19.988692\du}{20.504075\du}}
\pgfpathlineto{\pgfpoint{-19.982422\du}{20.501388\du}}
\pgfpathlineto{\pgfpoint{-19.977048\du}{20.498701\du}}
\pgfpathlineto{\pgfpoint{-19.971973\du}{20.496014\du}}
\pgfpathlineto{\pgfpoint{-19.967196\du}{20.492730\du}}
\pgfpathlineto{\pgfpoint{-19.963613\du}{20.489745\du}}
\pgfpathlineto{\pgfpoint{-19.959434\du}{20.486759\du}}
\pgfpathlineto{\pgfpoint{-19.956448\du}{20.483774\du}}
\pgfpathlineto{\pgfpoint{-19.953761\du}{20.480490\du}}
\pgfpathlineto{\pgfpoint{-19.951373\du}{20.477206\du}}
\pgfpathlineto{\pgfpoint{-19.949880\du}{20.473922\du}}
\pgfpathlineto{\pgfpoint{-19.948686\du}{20.470638\du}}
\pgfpathlineto{\pgfpoint{-19.947790\du}{20.467652\du}}
\pgfpathlineto{\pgfpoint{-19.947790\du}{20.464069\du}}
\pgfusepath{fill}
\pgfsetbuttcap
\pgfsetmiterjoin
\pgfsetdash{}{0pt}
\definecolor{dialinecolor}{rgb}{0.666667, 0.901961, 1.000000}
\pgfsetstrokecolor{dialinecolor}
\pgfpathmoveto{\pgfpoint{-19.959434\du}{20.458397\du}}
\pgfpathlineto{\pgfpoint{-19.959434\du}{20.454814\du}}
\pgfpathlineto{\pgfpoint{-19.960329\du}{20.452426\du}}
\pgfpathlineto{\pgfpoint{-19.961524\du}{20.449142\du}}
\pgfpathlineto{\pgfpoint{-19.963613\du}{20.446156\du}}
\pgfpathlineto{\pgfpoint{-19.965703\du}{20.443469\du}}
\pgfpathlineto{\pgfpoint{-19.967793\du}{20.439887\du}}
\pgfpathlineto{\pgfpoint{-19.971077\du}{20.437498\du}}
\pgfpathlineto{\pgfpoint{-19.974361\du}{20.434811\du}}
\pgfpathlineto{\pgfpoint{-19.978541\du}{20.432124\du}}
\pgfpathlineto{\pgfpoint{-19.983318\du}{20.429736\du}}
\pgfpathlineto{\pgfpoint{-19.987796\du}{20.426750\du}}
\pgfpathlineto{\pgfpoint{-19.993469\du}{20.424661\du}}
\pgfpathlineto{\pgfpoint{-19.999440\du}{20.421974\du}}
\pgfpathlineto{\pgfpoint{-20.005112\du}{20.419287\du}}
\pgfpathlineto{\pgfpoint{-20.011680\du}{20.417495\du}}
\pgfpathlineto{\pgfpoint{-20.018547\du}{20.415107\du}}
\pgfpathlineto{\pgfpoint{-20.025414\du}{20.413614\du}}
\pgfpathlineto{\pgfpoint{-20.033176\du}{20.411524\du}}
\pgfpathlineto{\pgfpoint{-20.040640\du}{20.409733\du}}
\pgfpathlineto{\pgfpoint{-20.048999\du}{20.407942\du}}
\pgfpathlineto{\pgfpoint{-20.057657\du}{20.406747\du}}
\pgfpathlineto{\pgfpoint{-20.066017\du}{20.404359\du}}
\pgfpathlineto{\pgfpoint{-20.074675\du}{20.404060\du}}
\pgfpathlineto{\pgfpoint{-20.083631\du}{20.402568\du}}
\pgfpathlineto{\pgfpoint{-20.092289\du}{20.401672\du}}
\pgfpathlineto{\pgfpoint{-20.101843\du}{20.400776\du}}
\pgfpathlineto{\pgfpoint{-20.111098\du}{20.400179\du}}
\pgfpathlineto{\pgfpoint{-20.120652\du}{20.399881\du}}
\pgfpathlineto{\pgfpoint{-20.130205\du}{20.399284\du}}
\pgfpathlineto{\pgfpoint{-20.140356\du}{20.398985\du}}
\pgfpathlineto{\pgfpoint{-20.150208\du}{20.398985\du}}
\pgfpathlineto{\pgfpoint{-20.150208\du}{20.398985\du}}
\pgfpathlineto{\pgfpoint{-20.159165\du}{20.398985\du}}
\pgfpathlineto{\pgfpoint{-20.169017\du}{20.399284\du}}
\pgfpathlineto{\pgfpoint{-20.178571\du}{20.399881\du}}
\pgfpathlineto{\pgfpoint{-20.187826\du}{20.400179\du}}
\pgfpathlineto{\pgfpoint{-20.197380\du}{20.400776\du}}
\pgfpathlineto{\pgfpoint{-20.206933\du}{20.401672\du}}
\pgfpathlineto{\pgfpoint{-20.215890\du}{20.402568\du}}
\pgfpathlineto{\pgfpoint{-20.224846\du}{20.404060\du}}
\pgfpathlineto{\pgfpoint{-20.233504\du}{20.404359\du}}
\pgfpathlineto{\pgfpoint{-20.241864\du}{20.406747\du}}
\pgfpathlineto{\pgfpoint{-20.250522\du}{20.407942\du}}
\pgfpathlineto{\pgfpoint{-20.258583\du}{20.409733\du}}
\pgfpathlineto{\pgfpoint{-20.266345\du}{20.411524\du}}
\pgfpathlineto{\pgfpoint{-20.273809\du}{20.413614\du}}
\pgfpathlineto{\pgfpoint{-20.280377\du}{20.415107\du}}
\pgfpathlineto{\pgfpoint{-20.287841\du}{20.417495\du}}
\pgfpathlineto{\pgfpoint{-20.294110\du}{20.419287\du}}
\pgfpathlineto{\pgfpoint{-20.300082\du}{20.421974\du}}
\pgfpathlineto{\pgfpoint{-20.306053\du}{20.424661\du}}
\pgfpathlineto{\pgfpoint{-20.311128\du}{20.426750\du}}
\pgfpathlineto{\pgfpoint{-20.316203\du}{20.429736\du}}
\pgfpathlineto{\pgfpoint{-20.320980\du}{20.432124\du}}
\pgfpathlineto{\pgfpoint{-20.325160\du}{20.434811\du}}
\pgfpathlineto{\pgfpoint{-20.328444\du}{20.437498\du}}
\pgfpathlineto{\pgfpoint{-20.331728\du}{20.439887\du}}
\pgfpathlineto{\pgfpoint{-20.334116\du}{20.443469\du}}
\pgfpathlineto{\pgfpoint{-20.335908\du}{20.446156\du}}
\pgfpathlineto{\pgfpoint{-20.338296\du}{20.449142\du}}
\pgfpathlineto{\pgfpoint{-20.339192\du}{20.452426\du}}
\pgfpathlineto{\pgfpoint{-20.340087\du}{20.454814\du}}
\pgfpathlineto{\pgfpoint{-20.340087\du}{20.458397\du}}
\pgfpathlineto{\pgfpoint{-20.340087\du}{20.458397\du}}
\pgfpathlineto{\pgfpoint{-20.340087\du}{20.461382\du}}
\pgfpathlineto{\pgfpoint{-20.339192\du}{20.464069\du}}
\pgfpathlineto{\pgfpoint{-20.338296\du}{20.467055\du}}
\pgfpathlineto{\pgfpoint{-20.335908\du}{20.470638\du}}
\pgfpathlineto{\pgfpoint{-20.334116\du}{20.473324\du}}
\pgfpathlineto{\pgfpoint{-20.331728\du}{20.476310\du}}
\pgfpathlineto{\pgfpoint{-20.328444\du}{20.478997\du}}
\pgfpathlineto{\pgfpoint{-20.325160\du}{20.481684\du}}
\pgfpathlineto{\pgfpoint{-20.320980\du}{20.484371\du}}
\pgfpathlineto{\pgfpoint{-20.316203\du}{20.487058\du}}
\pgfpathlineto{\pgfpoint{-20.311128\du}{20.489745\du}}
\pgfpathlineto{\pgfpoint{-20.306053\du}{20.492133\du}}
\pgfpathlineto{\pgfpoint{-20.300082\du}{20.494820\du}}
\pgfpathlineto{\pgfpoint{-20.294110\du}{20.496910\du}}
\pgfpathlineto{\pgfpoint{-20.287841\du}{20.499299\du}}
\pgfpathlineto{\pgfpoint{-20.280377\du}{20.501687\du}}
\pgfpathlineto{\pgfpoint{-20.273809\du}{20.503478\du}}
\pgfpathlineto{\pgfpoint{-20.266345\du}{20.504971\du}}
\pgfpathlineto{\pgfpoint{-20.258583\du}{20.507061\du}}
\pgfpathlineto{\pgfpoint{-20.250522\du}{20.508852\du}}
\pgfpathlineto{\pgfpoint{-20.241864\du}{20.510046\du}}
\pgfpathlineto{\pgfpoint{-20.233504\du}{20.511838\du}}
\pgfpathlineto{\pgfpoint{-20.224846\du}{20.512733\du}}
\pgfpathlineto{\pgfpoint{-20.215890\du}{20.514226\du}}
\pgfpathlineto{\pgfpoint{-20.206933\du}{20.515122\du}}
\pgfpathlineto{\pgfpoint{-20.197380\du}{20.515420\du}}
\pgfpathlineto{\pgfpoint{-20.187826\du}{20.516615\du}}
\pgfpathlineto{\pgfpoint{-20.178571\du}{20.516913\du}}
\pgfpathlineto{\pgfpoint{-20.169017\du}{20.517212\du}}
\pgfpathlineto{\pgfpoint{-20.159165\du}{20.517809\du}}
\pgfpathlineto{\pgfpoint{-20.150208\du}{20.517809\du}}
\pgfpathlineto{\pgfpoint{-20.150208\du}{20.517809\du}}
\pgfpathlineto{\pgfpoint{-20.140356\du}{20.517809\du}}
\pgfpathlineto{\pgfpoint{-20.130205\du}{20.517212\du}}
\pgfpathlineto{\pgfpoint{-20.120652\du}{20.516913\du}}
\pgfpathlineto{\pgfpoint{-20.111098\du}{20.516615\du}}
\pgfpathlineto{\pgfpoint{-20.101843\du}{20.515420\du}}
\pgfpathlineto{\pgfpoint{-20.092289\du}{20.515122\du}}
\pgfpathlineto{\pgfpoint{-20.083631\du}{20.514226\du}}
\pgfpathlineto{\pgfpoint{-20.074675\du}{20.512733\du}}
\pgfpathlineto{\pgfpoint{-20.066017\du}{20.511838\du}}
\pgfpathlineto{\pgfpoint{-20.057657\du}{20.510046\du}}
\pgfpathlineto{\pgfpoint{-20.048999\du}{20.508852\du}}
\pgfpathlineto{\pgfpoint{-20.040640\du}{20.507061\du}}
\pgfpathlineto{\pgfpoint{-20.033176\du}{20.504971\du}}
\pgfpathlineto{\pgfpoint{-20.025414\du}{20.503478\du}}
\pgfpathlineto{\pgfpoint{-20.018547\du}{20.501687\du}}
\pgfpathlineto{\pgfpoint{-20.011680\du}{20.499299\du}}
\pgfpathlineto{\pgfpoint{-20.005112\du}{20.496910\du}}
\pgfpathlineto{\pgfpoint{-19.999440\du}{20.494820\du}}
\pgfpathlineto{\pgfpoint{-19.993469\du}{20.492133\du}}
\pgfpathlineto{\pgfpoint{-19.987796\du}{20.489745\du}}
\pgfpathlineto{\pgfpoint{-19.983318\du}{20.487058\du}}
\pgfpathlineto{\pgfpoint{-19.978541\du}{20.484371\du}}
\pgfpathlineto{\pgfpoint{-19.974361\du}{20.481684\du}}
\pgfpathlineto{\pgfpoint{-19.971077\du}{20.478997\du}}
\pgfpathlineto{\pgfpoint{-19.967793\du}{20.476310\du}}
\pgfpathlineto{\pgfpoint{-19.965703\du}{20.473324\du}}
\pgfpathlineto{\pgfpoint{-19.963613\du}{20.470638\du}}
\pgfpathlineto{\pgfpoint{-19.961524\du}{20.467055\du}}
\pgfpathlineto{\pgfpoint{-19.960329\du}{20.464069\du}}
\pgfpathlineto{\pgfpoint{-19.959434\du}{20.461382\du}}
\pgfpathlineto{\pgfpoint{-19.959434\du}{20.458397\du}}
\pgfusepath{stroke}
\pgfsetbuttcap
\pgfsetmiterjoin
\pgfsetdash{}{0pt}
\definecolor{dialinecolor}{rgb}{0.000000, 0.470588, 0.666667}
\pgfsetfillcolor{dialinecolor}
\pgfpathmoveto{\pgfpoint{-20.340087\du}{20.375101\du}}
\pgfpathlineto{\pgfpoint{-20.340087\du}{20.469742\du}}
\pgfpathlineto{\pgfpoint{-19.947193\du}{20.469742\du}}
\pgfpathlineto{\pgfpoint{-19.947193\du}{20.375101\du}}
\pgfpathlineto{\pgfpoint{-20.340087\du}{20.375101\du}}
\pgfusepath{fill}
\pgfsetbuttcap
\pgfsetmiterjoin
\pgfsetdash{}{0pt}
\definecolor{dialinecolor}{rgb}{0.000000, 0.705882, 1.000000}
\pgfsetfillcolor{dialinecolor}
\pgfpathmoveto{\pgfpoint{-19.947790\du}{20.368533\du}}
\pgfpathlineto{\pgfpoint{-19.947790\du}{20.365249\du}}
\pgfpathlineto{\pgfpoint{-19.948686\du}{20.361666\du}}
\pgfpathlineto{\pgfpoint{-19.949880\du}{20.358681\du}}
\pgfpathlineto{\pgfpoint{-19.951373\du}{20.355396\du}}
\pgfpathlineto{\pgfpoint{-19.953761\du}{20.352112\du}}
\pgfpathlineto{\pgfpoint{-19.956448\du}{20.348828\du}}
\pgfpathlineto{\pgfpoint{-19.959434\du}{20.345843\du}}
\pgfpathlineto{\pgfpoint{-19.963613\du}{20.342559\du}}
\pgfpathlineto{\pgfpoint{-19.967196\du}{20.339872\du}}
\pgfpathlineto{\pgfpoint{-19.971973\du}{20.336886\du}}
\pgfpathlineto{\pgfpoint{-19.977048\du}{20.333901\du}}
\pgfpathlineto{\pgfpoint{-19.982422\du}{20.331214\du}}
\pgfpathlineto{\pgfpoint{-19.988692\du}{20.327930\du}}
\pgfpathlineto{\pgfpoint{-19.994663\du}{20.326138\du}}
\pgfpathlineto{\pgfpoint{-20.001529\du}{20.323451\du}}
\pgfpathlineto{\pgfpoint{-20.008396\du}{20.321063\du}}
\pgfpathlineto{\pgfpoint{-20.015860\du}{20.318973\du}}
\pgfpathlineto{\pgfpoint{-20.023324\du}{20.316883\du}}
\pgfpathlineto{\pgfpoint{-20.031385\du}{20.314793\du}}
\pgfpathlineto{\pgfpoint{-20.040043\du}{20.312704\du}}
\pgfpathlineto{\pgfpoint{-20.048402\du}{20.311211\du}}
\pgfpathlineto{\pgfpoint{-20.057060\du}{20.309718\du}}
\pgfpathlineto{\pgfpoint{-20.066315\du}{20.308524\du}}
\pgfpathlineto{\pgfpoint{-20.075570\du}{20.307031\du}}
\pgfpathlineto{\pgfpoint{-20.084825\du}{20.306135\du}}
\pgfpathlineto{\pgfpoint{-20.094379\du}{20.304941\du}}
\pgfpathlineto{\pgfpoint{-20.104530\du}{20.304344\du}}
\pgfpathlineto{\pgfpoint{-20.114084\du}{20.303747\du}}
\pgfpathlineto{\pgfpoint{-20.123637\du}{20.303448\du}}
\pgfpathlineto{\pgfpoint{-20.133788\du}{20.303150\du}}
\pgfpathlineto{\pgfpoint{-20.143640\du}{20.303150\du}}
\pgfpathlineto{\pgfpoint{-20.143640\du}{20.303150\du}}
\pgfpathlineto{\pgfpoint{-20.153492\du}{20.303150\du}}
\pgfpathlineto{\pgfpoint{-20.163643\du}{20.303448\du}}
\pgfpathlineto{\pgfpoint{-20.173495\du}{20.303747\du}}
\pgfpathlineto{\pgfpoint{-20.183348\du}{20.304344\du}}
\pgfpathlineto{\pgfpoint{-20.192901\du}{20.304941\du}}
\pgfpathlineto{\pgfpoint{-20.202455\du}{20.306135\du}}
\pgfpathlineto{\pgfpoint{-20.212009\du}{20.307031\du}}
\pgfpathlineto{\pgfpoint{-20.220965\du}{20.308524\du}}
\pgfpathlineto{\pgfpoint{-20.230220\du}{20.309718\du}}
\pgfpathlineto{\pgfpoint{-20.238878\du}{20.311211\du}}
\pgfpathlineto{\pgfpoint{-20.247536\du}{20.312704\du}}
\pgfpathlineto{\pgfpoint{-20.255896\du}{20.314793\du}}
\pgfpathlineto{\pgfpoint{-20.263957\du}{20.316883\du}}
\pgfpathlineto{\pgfpoint{-20.271421\du}{20.318973\du}}
\pgfpathlineto{\pgfpoint{-20.278884\du}{20.321063\du}}
\pgfpathlineto{\pgfpoint{-20.286050\du}{20.323451\du}}
\pgfpathlineto{\pgfpoint{-20.292618\du}{20.326138\du}}
\pgfpathlineto{\pgfpoint{-20.298887\du}{20.327930\du}}
\pgfpathlineto{\pgfpoint{-20.304858\du}{20.331214\du}}
\pgfpathlineto{\pgfpoint{-20.310232\du}{20.333901\du}}
\pgfpathlineto{\pgfpoint{-20.315308\du}{20.336886\du}}
\pgfpathlineto{\pgfpoint{-20.320084\du}{20.339872\du}}
\pgfpathlineto{\pgfpoint{-20.323966\du}{20.342559\du}}
\pgfpathlineto{\pgfpoint{-20.327847\du}{20.345843\du}}
\pgfpathlineto{\pgfpoint{-20.330832\du}{20.348828\du}}
\pgfpathlineto{\pgfpoint{-20.333519\du}{20.352112\du}}
\pgfpathlineto{\pgfpoint{-20.335908\du}{20.355396\du}}
\pgfpathlineto{\pgfpoint{-20.337699\du}{20.358681\du}}
\pgfpathlineto{\pgfpoint{-20.339192\du}{20.361666\du}}
\pgfpathlineto{\pgfpoint{-20.340087\du}{20.365249\du}}
\pgfpathlineto{\pgfpoint{-20.340087\du}{20.368533\du}}
\pgfpathlineto{\pgfpoint{-20.340087\du}{20.368533\du}}
\pgfpathlineto{\pgfpoint{-20.340087\du}{20.371817\du}}
\pgfpathlineto{\pgfpoint{-20.339192\du}{20.375399\du}}
\pgfpathlineto{\pgfpoint{-20.337699\du}{20.378385\du}}
\pgfpathlineto{\pgfpoint{-20.335908\du}{20.381669\du}}
\pgfpathlineto{\pgfpoint{-20.333519\du}{20.385252\du}}
\pgfpathlineto{\pgfpoint{-20.330832\du}{20.388237\du}}
\pgfpathlineto{\pgfpoint{-20.327847\du}{20.391223\du}}
\pgfpathlineto{\pgfpoint{-20.323966\du}{20.394507\du}}
\pgfpathlineto{\pgfpoint{-20.320084\du}{20.396895\du}}
\pgfpathlineto{\pgfpoint{-20.315308\du}{20.400179\du}}
\pgfpathlineto{\pgfpoint{-20.310232\du}{20.403463\du}}
\pgfpathlineto{\pgfpoint{-20.304858\du}{20.406150\du}}
\pgfpathlineto{\pgfpoint{-20.298887\du}{20.408837\du}}
\pgfpathlineto{\pgfpoint{-20.292618\du}{20.410927\du}}
\pgfpathlineto{\pgfpoint{-20.286050\du}{20.413913\du}}
\pgfpathlineto{\pgfpoint{-20.278884\du}{20.416301\du}}
\pgfpathlineto{\pgfpoint{-20.271421\du}{20.418092\du}}
\pgfpathlineto{\pgfpoint{-20.263957\du}{20.420182\du}}
\pgfpathlineto{\pgfpoint{-20.255896\du}{20.422272\du}}
\pgfpathlineto{\pgfpoint{-20.247536\du}{20.424362\du}}
\pgfpathlineto{\pgfpoint{-20.238878\du}{20.425855\du}}
\pgfpathlineto{\pgfpoint{-20.230220\du}{20.427347\du}}
\pgfpathlineto{\pgfpoint{-20.220965\du}{20.428542\du}}
\pgfpathlineto{\pgfpoint{-20.212009\du}{20.430034\du}}
\pgfpathlineto{\pgfpoint{-20.202455\du}{20.430930\du}}
\pgfpathlineto{\pgfpoint{-20.192901\du}{20.432124\du}}
\pgfpathlineto{\pgfpoint{-20.183348\du}{20.432423\du}}
\pgfpathlineto{\pgfpoint{-20.173495\du}{20.433319\du}}
\pgfpathlineto{\pgfpoint{-20.163643\du}{20.433319\du}}
\pgfpathlineto{\pgfpoint{-20.153492\du}{20.434214\du}}
\pgfpathlineto{\pgfpoint{-20.143640\du}{20.434214\du}}
\pgfpathlineto{\pgfpoint{-20.143640\du}{20.434214\du}}
\pgfpathlineto{\pgfpoint{-20.133788\du}{20.434214\du}}
\pgfpathlineto{\pgfpoint{-20.123637\du}{20.433319\du}}
\pgfpathlineto{\pgfpoint{-20.114084\du}{20.433319\du}}
\pgfpathlineto{\pgfpoint{-20.104530\du}{20.432423\du}}
\pgfpathlineto{\pgfpoint{-20.094379\du}{20.432124\du}}
\pgfpathlineto{\pgfpoint{-20.084825\du}{20.430930\du}}
\pgfpathlineto{\pgfpoint{-20.075570\du}{20.430034\du}}
\pgfpathlineto{\pgfpoint{-20.066315\du}{20.428542\du}}
\pgfpathlineto{\pgfpoint{-20.057060\du}{20.427347\du}}
\pgfpathlineto{\pgfpoint{-20.048402\du}{20.425855\du}}
\pgfpathlineto{\pgfpoint{-20.040043\du}{20.424362\du}}
\pgfpathlineto{\pgfpoint{-20.031385\du}{20.422272\du}}
\pgfpathlineto{\pgfpoint{-20.023324\du}{20.420182\du}}
\pgfpathlineto{\pgfpoint{-20.015860\du}{20.418092\du}}
\pgfpathlineto{\pgfpoint{-20.008396\du}{20.416301\du}}
\pgfpathlineto{\pgfpoint{-20.001529\du}{20.413913\du}}
\pgfpathlineto{\pgfpoint{-19.994663\du}{20.410927\du}}
\pgfpathlineto{\pgfpoint{-19.988692\du}{20.408837\du}}
\pgfpathlineto{\pgfpoint{-19.982422\du}{20.406150\du}}
\pgfpathlineto{\pgfpoint{-19.977048\du}{20.403463\du}}
\pgfpathlineto{\pgfpoint{-19.971973\du}{20.400179\du}}
\pgfpathlineto{\pgfpoint{-19.967196\du}{20.396895\du}}
\pgfpathlineto{\pgfpoint{-19.963613\du}{20.394507\du}}
\pgfpathlineto{\pgfpoint{-19.959434\du}{20.391223\du}}
\pgfpathlineto{\pgfpoint{-19.956448\du}{20.388237\du}}
\pgfpathlineto{\pgfpoint{-19.953761\du}{20.385252\du}}
\pgfpathlineto{\pgfpoint{-19.951373\du}{20.381669\du}}
\pgfpathlineto{\pgfpoint{-19.949880\du}{20.378385\du}}
\pgfpathlineto{\pgfpoint{-19.948686\du}{20.375399\du}}
\pgfpathlineto{\pgfpoint{-19.947790\du}{20.371817\du}}
\pgfpathlineto{\pgfpoint{-19.947790\du}{20.368533\du}}
\pgfusepath{fill}
\pgfsetbuttcap
\pgfsetmiterjoin
\pgfsetdash{}{0pt}
\definecolor{dialinecolor}{rgb}{0.666667, 0.901961, 1.000000}
\pgfsetstrokecolor{dialinecolor}
\pgfpathmoveto{\pgfpoint{-19.959434\du}{20.362860\du}}
\pgfpathlineto{\pgfpoint{-19.959434\du}{20.359576\du}}
\pgfpathlineto{\pgfpoint{-19.960329\du}{20.357188\du}}
\pgfpathlineto{\pgfpoint{-19.961524\du}{20.353904\du}}
\pgfpathlineto{\pgfpoint{-19.963613\du}{20.350321\du}}
\pgfpathlineto{\pgfpoint{-19.965703\du}{20.347933\du}}
\pgfpathlineto{\pgfpoint{-19.967793\du}{20.344649\du}}
\pgfpathlineto{\pgfpoint{-19.971077\du}{20.342260\du}}
\pgfpathlineto{\pgfpoint{-19.974361\du}{20.339573\du}}
\pgfpathlineto{\pgfpoint{-19.978541\du}{20.336886\du}}
\pgfpathlineto{\pgfpoint{-19.983318\du}{20.334199\du}}
\pgfpathlineto{\pgfpoint{-19.987796\du}{20.331512\du}}
\pgfpathlineto{\pgfpoint{-19.993469\du}{20.328825\du}}
\pgfpathlineto{\pgfpoint{-19.999440\du}{20.326437\du}}
\pgfpathlineto{\pgfpoint{-20.005112\du}{20.324347\du}}
\pgfpathlineto{\pgfpoint{-20.011680\du}{20.321959\du}}
\pgfpathlineto{\pgfpoint{-20.018547\du}{20.319570\du}}
\pgfpathlineto{\pgfpoint{-20.025414\du}{20.317779\du}}
\pgfpathlineto{\pgfpoint{-20.033176\du}{20.316286\du}}
\pgfpathlineto{\pgfpoint{-20.040640\du}{20.313898\du}}
\pgfpathlineto{\pgfpoint{-20.048999\du}{20.312405\du}}
\pgfpathlineto{\pgfpoint{-20.057657\du}{20.311211\du}}
\pgfpathlineto{\pgfpoint{-20.066017\du}{20.309419\du}}
\pgfpathlineto{\pgfpoint{-20.074675\du}{20.308524\du}}
\pgfpathlineto{\pgfpoint{-20.083631\du}{20.307031\du}}
\pgfpathlineto{\pgfpoint{-20.092289\du}{20.306135\du}}
\pgfpathlineto{\pgfpoint{-20.101843\du}{20.305538\du}}
\pgfpathlineto{\pgfpoint{-20.111098\du}{20.304643\du}}
\pgfpathlineto{\pgfpoint{-20.120652\du}{20.304344\du}}
\pgfpathlineto{\pgfpoint{-20.130205\du}{20.304045\du}}
\pgfpathlineto{\pgfpoint{-20.140356\du}{20.303448\du}}
\pgfpathlineto{\pgfpoint{-20.150208\du}{20.303448\du}}
\pgfpathlineto{\pgfpoint{-20.150208\du}{20.303448\du}}
\pgfpathlineto{\pgfpoint{-20.159165\du}{20.303448\du}}
\pgfpathlineto{\pgfpoint{-20.169017\du}{20.304045\du}}
\pgfpathlineto{\pgfpoint{-20.178571\du}{20.304344\du}}
\pgfpathlineto{\pgfpoint{-20.187826\du}{20.304643\du}}
\pgfpathlineto{\pgfpoint{-20.197380\du}{20.305538\du}}
\pgfpathlineto{\pgfpoint{-20.206933\du}{20.306135\du}}
\pgfpathlineto{\pgfpoint{-20.215890\du}{20.307031\du}}
\pgfpathlineto{\pgfpoint{-20.224846\du}{20.308524\du}}
\pgfpathlineto{\pgfpoint{-20.233504\du}{20.309419\du}}
\pgfpathlineto{\pgfpoint{-20.241864\du}{20.311211\du}}
\pgfpathlineto{\pgfpoint{-20.250522\du}{20.312405\du}}
\pgfpathlineto{\pgfpoint{-20.258583\du}{20.313898\du}}
\pgfpathlineto{\pgfpoint{-20.266345\du}{20.316286\du}}
\pgfpathlineto{\pgfpoint{-20.273809\du}{20.317779\du}}
\pgfpathlineto{\pgfpoint{-20.280377\du}{20.319570\du}}
\pgfpathlineto{\pgfpoint{-20.287841\du}{20.321959\du}}
\pgfpathlineto{\pgfpoint{-20.294110\du}{20.324347\du}}
\pgfpathlineto{\pgfpoint{-20.300082\du}{20.326437\du}}
\pgfpathlineto{\pgfpoint{-20.306053\du}{20.328825\du}}
\pgfpathlineto{\pgfpoint{-20.311128\du}{20.331512\du}}
\pgfpathlineto{\pgfpoint{-20.316203\du}{20.334199\du}}
\pgfpathlineto{\pgfpoint{-20.320980\du}{20.336886\du}}
\pgfpathlineto{\pgfpoint{-20.325160\du}{20.339573\du}}
\pgfpathlineto{\pgfpoint{-20.328444\du}{20.342260\du}}
\pgfpathlineto{\pgfpoint{-20.331728\du}{20.344649\du}}
\pgfpathlineto{\pgfpoint{-20.334116\du}{20.347933\du}}
\pgfpathlineto{\pgfpoint{-20.335908\du}{20.350321\du}}
\pgfpathlineto{\pgfpoint{-20.338296\du}{20.353904\du}}
\pgfpathlineto{\pgfpoint{-20.339192\du}{20.357188\du}}
\pgfpathlineto{\pgfpoint{-20.340087\du}{20.359576\du}}
\pgfpathlineto{\pgfpoint{-20.340087\du}{20.362860\du}}
\pgfpathlineto{\pgfpoint{-20.340087\du}{20.362860\du}}
\pgfpathlineto{\pgfpoint{-20.340087\du}{20.366144\du}}
\pgfpathlineto{\pgfpoint{-20.339192\du}{20.368831\du}}
\pgfpathlineto{\pgfpoint{-20.338296\du}{20.371817\du}}
\pgfpathlineto{\pgfpoint{-20.335908\du}{20.375101\du}}
\pgfpathlineto{\pgfpoint{-20.334116\du}{20.377788\du}}
\pgfpathlineto{\pgfpoint{-20.331728\du}{20.381072\du}}
\pgfpathlineto{\pgfpoint{-20.328444\du}{20.383759\du}}
\pgfpathlineto{\pgfpoint{-20.325160\du}{20.386446\du}}
\pgfpathlineto{\pgfpoint{-20.320980\du}{20.389133\du}}
\pgfpathlineto{\pgfpoint{-20.316203\du}{20.391820\du}}
\pgfpathlineto{\pgfpoint{-20.311128\du}{20.394507\du}}
\pgfpathlineto{\pgfpoint{-20.306053\du}{20.396597\du}}
\pgfpathlineto{\pgfpoint{-20.300082\du}{20.399284\du}}
\pgfpathlineto{\pgfpoint{-20.294110\du}{20.401672\du}}
\pgfpathlineto{\pgfpoint{-20.287841\du}{20.404060\du}}
\pgfpathlineto{\pgfpoint{-20.280377\du}{20.406150\du}}
\pgfpathlineto{\pgfpoint{-20.273809\du}{20.407942\du}}
\pgfpathlineto{\pgfpoint{-20.266345\du}{20.409733\du}}
\pgfpathlineto{\pgfpoint{-20.258583\du}{20.411524\du}}
\pgfpathlineto{\pgfpoint{-20.250522\du}{20.413614\du}}
\pgfpathlineto{\pgfpoint{-20.241864\du}{20.414808\du}}
\pgfpathlineto{\pgfpoint{-20.233504\du}{20.416600\du}}
\pgfpathlineto{\pgfpoint{-20.224846\du}{20.417495\du}}
\pgfpathlineto{\pgfpoint{-20.215890\du}{20.418988\du}}
\pgfpathlineto{\pgfpoint{-20.206933\du}{20.419287\du}}
\pgfpathlineto{\pgfpoint{-20.197380\du}{20.420481\du}}
\pgfpathlineto{\pgfpoint{-20.187826\du}{20.420779\du}}
\pgfpathlineto{\pgfpoint{-20.178571\du}{20.421675\du}}
\pgfpathlineto{\pgfpoint{-20.169017\du}{20.421974\du}}
\pgfpathlineto{\pgfpoint{-20.159165\du}{20.422272\du}}
\pgfpathlineto{\pgfpoint{-20.150208\du}{20.422272\du}}
\pgfpathlineto{\pgfpoint{-20.150208\du}{20.422272\du}}
\pgfpathlineto{\pgfpoint{-20.140356\du}{20.422272\du}}
\pgfpathlineto{\pgfpoint{-20.130205\du}{20.421974\du}}
\pgfpathlineto{\pgfpoint{-20.120652\du}{20.421675\du}}
\pgfpathlineto{\pgfpoint{-20.111098\du}{20.420779\du}}
\pgfpathlineto{\pgfpoint{-20.101843\du}{20.420481\du}}
\pgfpathlineto{\pgfpoint{-20.092289\du}{20.419287\du}}
\pgfpathlineto{\pgfpoint{-20.083631\du}{20.418988\du}}
\pgfpathlineto{\pgfpoint{-20.074675\du}{20.417495\du}}
\pgfpathlineto{\pgfpoint{-20.066017\du}{20.416600\du}}
\pgfpathlineto{\pgfpoint{-20.057657\du}{20.414808\du}}
\pgfpathlineto{\pgfpoint{-20.048999\du}{20.413614\du}}
\pgfpathlineto{\pgfpoint{-20.040640\du}{20.411524\du}}
\pgfpathlineto{\pgfpoint{-20.033176\du}{20.409733\du}}
\pgfpathlineto{\pgfpoint{-20.025414\du}{20.407942\du}}
\pgfpathlineto{\pgfpoint{-20.018547\du}{20.406150\du}}
\pgfpathlineto{\pgfpoint{-20.011680\du}{20.404060\du}}
\pgfpathlineto{\pgfpoint{-20.005112\du}{20.401672\du}}
\pgfpathlineto{\pgfpoint{-19.999440\du}{20.399284\du}}
\pgfpathlineto{\pgfpoint{-19.993469\du}{20.396597\du}}
\pgfpathlineto{\pgfpoint{-19.987796\du}{20.394507\du}}
\pgfpathlineto{\pgfpoint{-19.983318\du}{20.391820\du}}
\pgfpathlineto{\pgfpoint{-19.978541\du}{20.389133\du}}
\pgfpathlineto{\pgfpoint{-19.974361\du}{20.386446\du}}
\pgfpathlineto{\pgfpoint{-19.971077\du}{20.383759\du}}
\pgfpathlineto{\pgfpoint{-19.967793\du}{20.381072\du}}
\pgfpathlineto{\pgfpoint{-19.965703\du}{20.377788\du}}
\pgfpathlineto{\pgfpoint{-19.963613\du}{20.375101\du}}
\pgfpathlineto{\pgfpoint{-19.961524\du}{20.371817\du}}
\pgfpathlineto{\pgfpoint{-19.960329\du}{20.368831\du}}
\pgfpathlineto{\pgfpoint{-19.959434\du}{20.366144\du}}
\pgfpathlineto{\pgfpoint{-19.959434\du}{20.362860\du}}
\pgfusepath{stroke}
\pgfsetbuttcap
\pgfsetmiterjoin
\pgfsetdash{}{0pt}
\definecolor{dialinecolor}{rgb}{0.000000, 0.000000, 0.000000}
\pgfsetfillcolor{dialinecolor}
\pgfpathmoveto{\pgfpoint{-20.137968\du}{20.351217\du}}
\pgfpathlineto{\pgfpoint{-20.114084\du}{20.362860\du}}
\pgfpathlineto{\pgfpoint{-20.042730\du}{20.339275\du}}
\pgfpathlineto{\pgfpoint{-20.006903\du}{20.351217\du}}
\pgfpathlineto{\pgfpoint{-20.030788\du}{20.327333\du}}
\pgfpathlineto{\pgfpoint{-20.114084\du}{20.327333\du}}
\pgfpathlineto{\pgfpoint{-20.078257\du}{20.327333\du}}
\pgfpathlineto{\pgfpoint{-20.137968\du}{20.351217\du}}
\pgfusepath{fill}
\pgfsetbuttcap
\pgfsetmiterjoin
\pgfsetdash{}{0pt}
\definecolor{dialinecolor}{rgb}{0.000000, 0.000000, 0.000000}
\pgfsetfillcolor{dialinecolor}
\pgfpathmoveto{\pgfpoint{-20.150208\du}{20.386744\du}}
\pgfpathlineto{\pgfpoint{-20.185139\du}{20.375101\du}}
\pgfpathlineto{\pgfpoint{-20.245148\du}{20.398985\du}}
\pgfpathlineto{\pgfpoint{-20.280377\du}{20.386744\du}}
\pgfpathlineto{\pgfpoint{-20.269032\du}{20.422272\du}}
\pgfpathlineto{\pgfpoint{-20.185139\du}{20.422272\du}}
\pgfpathlineto{\pgfpoint{-20.221264\du}{20.410927\du}}
\pgfpathlineto{\pgfpoint{-20.150208\du}{20.386744\du}}
\pgfusepath{fill}
\pgfsetbuttcap
\pgfsetmiterjoin
\pgfsetdash{}{0pt}
\definecolor{dialinecolor}{rgb}{0.000000, 0.000000, 0.000000}
\pgfsetfillcolor{dialinecolor}
\pgfpathmoveto{\pgfpoint{-20.269032\du}{20.327333\du}}
\pgfpathlineto{\pgfpoint{-20.245148\du}{20.314793\du}}
\pgfpathlineto{\pgfpoint{-20.173495\du}{20.351217\du}}
\pgfpathlineto{\pgfpoint{-20.150208\du}{20.339275\du}}
\pgfpathlineto{\pgfpoint{-20.161852\du}{20.362860\du}}
\pgfpathlineto{\pgfpoint{-20.245148\du}{20.362860\du}}
\pgfpathlineto{\pgfpoint{-20.209322\du}{20.351217\du}}
\pgfpathlineto{\pgfpoint{-20.269032\du}{20.327333\du}}
\pgfusepath{fill}
\pgfsetbuttcap
\pgfsetmiterjoin
\pgfsetdash{}{0pt}
\definecolor{dialinecolor}{rgb}{0.000000, 0.000000, 0.000000}
\pgfsetfillcolor{dialinecolor}
\pgfpathmoveto{\pgfpoint{-20.018547\du}{20.410927\du}}
\pgfpathlineto{\pgfpoint{-20.042730\du}{20.422272\du}}
\pgfpathlineto{\pgfpoint{-20.114084\du}{20.386744\du}}
\pgfpathlineto{\pgfpoint{-20.150208\du}{20.398985\du}}
\pgfpathlineto{\pgfpoint{-20.125727\du}{20.375101\du}}
\pgfpathlineto{\pgfpoint{-20.042730\du}{20.375101\du}}
\pgfpathlineto{\pgfpoint{-20.078257\du}{20.386744\du}}
\pgfpathlineto{\pgfpoint{-20.018547\du}{20.410927\du}}
\pgfusepath{fill}
\pgfsetbuttcap
\pgfsetmiterjoin
\pgfsetdash{}{0pt}
\definecolor{dialinecolor}{rgb}{1.000000, 1.000000, 1.000000}
\pgfsetfillcolor{dialinecolor}
\pgfpathmoveto{\pgfpoint{-20.137968\du}{20.362860\du}}
\pgfpathlineto{\pgfpoint{-20.114084\du}{20.362860\du}}
\pgfpathlineto{\pgfpoint{-20.042730\du}{20.339275\du}}
\pgfpathlineto{\pgfpoint{-20.006903\du}{20.351217\du}}
\pgfpathlineto{\pgfpoint{-20.030788\du}{20.327333\du}}
\pgfpathlineto{\pgfpoint{-20.101843\du}{20.327333\du}}
\pgfpathlineto{\pgfpoint{-20.078257\du}{20.327333\du}}
\pgfpathlineto{\pgfpoint{-20.137968\du}{20.362860\du}}
\pgfusepath{fill}
\pgfsetbuttcap
\pgfsetmiterjoin
\pgfsetdash{}{0pt}
\definecolor{dialinecolor}{rgb}{1.000000, 1.000000, 1.000000}
\pgfsetfillcolor{dialinecolor}
\pgfpathmoveto{\pgfpoint{-20.150208\du}{20.386744\du}}
\pgfpathlineto{\pgfpoint{-20.173495\du}{20.375101\du}}
\pgfpathlineto{\pgfpoint{-20.245148\du}{20.398985\du}}
\pgfpathlineto{\pgfpoint{-20.280377\du}{20.398985\du}}
\pgfpathlineto{\pgfpoint{-20.256791\du}{20.422272\du}}
\pgfpathlineto{\pgfpoint{-20.173495\du}{20.422272\du}}
\pgfpathlineto{\pgfpoint{-20.209322\du}{20.410927\du}}
\pgfpathlineto{\pgfpoint{-20.150208\du}{20.386744\du}}
\pgfusepath{fill}
\pgfsetbuttcap
\pgfsetmiterjoin
\pgfsetdash{}{0pt}
\definecolor{dialinecolor}{rgb}{1.000000, 1.000000, 1.000000}
\pgfsetfillcolor{dialinecolor}
\pgfpathmoveto{\pgfpoint{-20.269032\du}{20.327333\du}}
\pgfpathlineto{\pgfpoint{-20.245148\du}{20.327333\du}}
\pgfpathlineto{\pgfpoint{-20.173495\du}{20.351217\du}}
\pgfpathlineto{\pgfpoint{-20.137968\du}{20.339275\du}}
\pgfpathlineto{\pgfpoint{-20.161852\du}{20.362860\du}}
\pgfpathlineto{\pgfpoint{-20.245148\du}{20.362860\du}}
\pgfpathlineto{\pgfpoint{-20.209322\du}{20.362860\du}}
\pgfpathlineto{\pgfpoint{-20.269032\du}{20.327333\du}}
\pgfusepath{fill}
\pgfsetbuttcap
\pgfsetmiterjoin
\pgfsetdash{}{0pt}
\definecolor{dialinecolor}{rgb}{1.000000, 1.000000, 1.000000}
\pgfsetfillcolor{dialinecolor}
\pgfpathmoveto{\pgfpoint{-20.018547\du}{20.410927\du}}
\pgfpathlineto{\pgfpoint{-20.042730\du}{20.422272\du}}
\pgfpathlineto{\pgfpoint{-20.114084\du}{20.398985\du}}
\pgfpathlineto{\pgfpoint{-20.137968\du}{20.398985\du}}
\pgfpathlineto{\pgfpoint{-20.125727\du}{20.375101\du}}
\pgfpathlineto{\pgfpoint{-20.042730\du}{20.375101\du}}
\pgfpathlineto{\pgfpoint{-20.078257\du}{20.386744\du}}
\pgfpathlineto{\pgfpoint{-20.018547\du}{20.410927\du}}
\pgfusepath{fill}
\pgfsetbuttcap
\pgfsetmiterjoin
\pgfsetdash{}{0pt}
\definecolor{dialinecolor}{rgb}{0.666667, 0.901961, 1.000000}
\pgfsetstrokecolor{dialinecolor}
\pgfpathmoveto{\pgfpoint{-20.340087\du}{20.362860\du}}
\pgfpathlineto{\pgfpoint{-20.340087\du}{20.457800\du}}
\pgfusepath{stroke}
\pgfsetbuttcap
\pgfsetmiterjoin
\pgfsetdash{}{0pt}
\definecolor{dialinecolor}{rgb}{0.666667, 0.901961, 1.000000}
\pgfsetstrokecolor{dialinecolor}
\pgfpathmoveto{\pgfpoint{-19.959434\du}{20.362860\du}}
\pgfpathlineto{\pgfpoint{-19.959434\du}{20.457800\du}}
\pgfusepath{stroke}
\pgfsetbuttcap
\pgfsetmiterjoin
\pgfsetdash{}{0pt}
\definecolor{dialinecolor}{rgb}{0.000000, 0.470588, 0.666667}
\pgfsetfillcolor{dialinecolor}
\pgfpathmoveto{\pgfpoint{-20.304858\du}{20.868906\du}}
\pgfpathlineto{\pgfpoint{-20.305157\du}{20.865323\du}}
\pgfpathlineto{\pgfpoint{-20.306053\du}{20.862039\du}}
\pgfpathlineto{\pgfpoint{-20.307247\du}{20.858755\du}}
\pgfpathlineto{\pgfpoint{-20.309038\du}{20.855770\du}}
\pgfpathlineto{\pgfpoint{-20.311128\du}{20.852486\du}}
\pgfpathlineto{\pgfpoint{-20.314412\du}{20.849201\du}}
\pgfpathlineto{\pgfpoint{-20.317099\du}{20.846216\du}}
\pgfpathlineto{\pgfpoint{-20.320980\du}{20.842932\du}}
\pgfpathlineto{\pgfpoint{-20.325160\du}{20.840245\du}}
\pgfpathlineto{\pgfpoint{-20.329638\du}{20.836961\du}}
\pgfpathlineto{\pgfpoint{-20.335012\du}{20.834274\du}}
\pgfpathlineto{\pgfpoint{-20.340087\du}{20.831587\du}}
\pgfpathlineto{\pgfpoint{-20.346059\du}{20.828900\du}}
\pgfpathlineto{\pgfpoint{-20.352030\du}{20.826213\du}}
\pgfpathlineto{\pgfpoint{-20.359195\du}{20.823825\du}}
\pgfpathlineto{\pgfpoint{-20.366062\du}{20.821436\du}}
\pgfpathlineto{\pgfpoint{-20.373525\du}{20.819645\du}}
\pgfpathlineto{\pgfpoint{-20.380989\du}{20.817256\du}}
\pgfpathlineto{\pgfpoint{-20.389050\du}{20.815465\du}}
\pgfpathlineto{\pgfpoint{-20.397409\du}{20.813077\du}}
\pgfpathlineto{\pgfpoint{-20.405769\du}{20.811584\du}}
\pgfpathlineto{\pgfpoint{-20.414725\du}{20.810390\du}}
\pgfpathlineto{\pgfpoint{-20.423682\du}{20.808897\du}}
\pgfpathlineto{\pgfpoint{-20.432937\du}{20.807703\du}}
\pgfpathlineto{\pgfpoint{-20.442789\du}{20.806509\du}}
\pgfpathlineto{\pgfpoint{-20.452044\du}{20.805613\du}}
\pgfpathlineto{\pgfpoint{-20.461897\du}{20.805016\du}}
\pgfpathlineto{\pgfpoint{-20.471450\du}{20.804419\du}}
\pgfpathlineto{\pgfpoint{-20.481303\du}{20.803822\du}}
\pgfpathlineto{\pgfpoint{-20.491752\du}{20.803523\du}}
\pgfpathlineto{\pgfpoint{-20.501007\du}{20.803523\du}}
\pgfpathlineto{\pgfpoint{-20.501007\du}{20.803523\du}}
\pgfpathlineto{\pgfpoint{-20.510859\du}{20.803523\du}}
\pgfpathlineto{\pgfpoint{-20.521010\du}{20.803822\du}}
\pgfpathlineto{\pgfpoint{-20.530862\du}{20.804419\du}}
\pgfpathlineto{\pgfpoint{-20.540714\du}{20.805016\du}}
\pgfpathlineto{\pgfpoint{-20.550567\du}{20.805613\du}}
\pgfpathlineto{\pgfpoint{-20.559822\du}{20.806509\du}}
\pgfpathlineto{\pgfpoint{-20.569674\du}{20.807703\du}}
\pgfpathlineto{\pgfpoint{-20.578929\du}{20.808897\du}}
\pgfpathlineto{\pgfpoint{-20.587886\du}{20.810390\du}}
\pgfpathlineto{\pgfpoint{-20.596544\du}{20.811584\du}}
\pgfpathlineto{\pgfpoint{-20.604903\du}{20.813077\du}}
\pgfpathlineto{\pgfpoint{-20.613263\du}{20.815465\du}}
\pgfpathlineto{\pgfpoint{-20.621323\du}{20.817256\du}}
\pgfpathlineto{\pgfpoint{-20.628787\du}{20.819645\du}}
\pgfpathlineto{\pgfpoint{-20.636251\du}{20.821436\du}}
\pgfpathlineto{\pgfpoint{-20.643416\du}{20.823825\du}}
\pgfpathlineto{\pgfpoint{-20.650283\du}{20.826213\du}}
\pgfpathlineto{\pgfpoint{-20.656254\du}{20.828900\du}}
\pgfpathlineto{\pgfpoint{-20.662524\du}{20.831587\du}}
\pgfpathlineto{\pgfpoint{-20.667599\du}{20.834274\du}}
\pgfpathlineto{\pgfpoint{-20.672674\du}{20.836961\du}}
\pgfpathlineto{\pgfpoint{-20.677451\du}{20.840245\du}}
\pgfpathlineto{\pgfpoint{-20.681631\du}{20.842932\du}}
\pgfpathlineto{\pgfpoint{-20.685214\du}{20.846216\du}}
\pgfpathlineto{\pgfpoint{-20.688199\du}{20.849201\du}}
\pgfpathlineto{\pgfpoint{-20.690886\du}{20.852486\du}}
\pgfpathlineto{\pgfpoint{-20.693275\du}{20.855770\du}}
\pgfpathlineto{\pgfpoint{-20.695066\du}{20.858755\du}}
\pgfpathlineto{\pgfpoint{-20.696559\du}{20.862039\du}}
\pgfpathlineto{\pgfpoint{-20.697454\du}{20.865323\du}}
\pgfpathlineto{\pgfpoint{-20.697454\du}{20.868906\du}}
\pgfpathlineto{\pgfpoint{-20.697454\du}{20.868906\du}}
\pgfpathlineto{\pgfpoint{-20.697454\du}{20.872190\du}}
\pgfpathlineto{\pgfpoint{-20.696559\du}{20.875474\du}}
\pgfpathlineto{\pgfpoint{-20.695066\du}{20.878758\du}}
\pgfpathlineto{\pgfpoint{-20.693275\du}{20.882042\du}}
\pgfpathlineto{\pgfpoint{-20.690886\du}{20.885028\du}}
\pgfpathlineto{\pgfpoint{-20.688199\du}{20.888312\du}}
\pgfpathlineto{\pgfpoint{-20.685214\du}{20.891297\du}}
\pgfpathlineto{\pgfpoint{-20.681631\du}{20.894581\du}}
\pgfpathlineto{\pgfpoint{-20.677451\du}{20.897567\du}}
\pgfpathlineto{\pgfpoint{-20.672674\du}{20.900552\du}}
\pgfpathlineto{\pgfpoint{-20.667599\du}{20.903239\du}}
\pgfpathlineto{\pgfpoint{-20.662524\du}{20.905926\du}}
\pgfpathlineto{\pgfpoint{-20.656254\du}{20.908613\du}}
\pgfpathlineto{\pgfpoint{-20.650283\du}{20.911002\du}}
\pgfpathlineto{\pgfpoint{-20.643416\du}{20.913689\du}}
\pgfpathlineto{\pgfpoint{-20.636251\du}{20.916077\du}}
\pgfpathlineto{\pgfpoint{-20.628787\du}{20.918167\du}}
\pgfpathlineto{\pgfpoint{-20.621323\du}{20.920555\du}}
\pgfpathlineto{\pgfpoint{-20.613263\du}{20.922048\du}}
\pgfpathlineto{\pgfpoint{-20.604903\du}{20.924138\du}}
\pgfpathlineto{\pgfpoint{-20.596544\du}{20.925929\du}}
\pgfpathlineto{\pgfpoint{-20.587886\du}{20.927124\du}}
\pgfpathlineto{\pgfpoint{-20.578929\du}{20.928616\du}}
\pgfpathlineto{\pgfpoint{-20.569674\du}{20.929811\du}}
\pgfpathlineto{\pgfpoint{-20.559822\du}{20.931005\du}}
\pgfpathlineto{\pgfpoint{-20.550567\du}{20.931900\du}}
\pgfpathlineto{\pgfpoint{-20.540714\du}{20.932498\du}}
\pgfpathlineto{\pgfpoint{-20.530862\du}{20.933393\du}}
\pgfpathlineto{\pgfpoint{-20.521010\du}{20.933393\du}}
\pgfpathlineto{\pgfpoint{-20.510859\du}{20.933990\du}}
\pgfpathlineto{\pgfpoint{-20.501007\du}{20.933990\du}}
\pgfpathlineto{\pgfpoint{-20.501007\du}{20.933990\du}}
\pgfpathlineto{\pgfpoint{-20.491752\du}{20.933990\du}}
\pgfpathlineto{\pgfpoint{-20.481303\du}{20.933393\du}}
\pgfpathlineto{\pgfpoint{-20.471450\du}{20.933393\du}}
\pgfpathlineto{\pgfpoint{-20.461897\du}{20.932498\du}}
\pgfpathlineto{\pgfpoint{-20.452044\du}{20.931900\du}}
\pgfpathlineto{\pgfpoint{-20.442789\du}{20.931005\du}}
\pgfpathlineto{\pgfpoint{-20.432937\du}{20.929811\du}}
\pgfpathlineto{\pgfpoint{-20.423682\du}{20.928616\du}}
\pgfpathlineto{\pgfpoint{-20.414725\du}{20.927124\du}}
\pgfpathlineto{\pgfpoint{-20.405769\du}{20.925929\du}}
\pgfpathlineto{\pgfpoint{-20.397409\du}{20.924138\du}}
\pgfpathlineto{\pgfpoint{-20.389050\du}{20.922048\du}}
\pgfpathlineto{\pgfpoint{-20.380989\du}{20.920555\du}}
\pgfpathlineto{\pgfpoint{-20.373525\du}{20.918167\du}}
\pgfpathlineto{\pgfpoint{-20.366062\du}{20.916077\du}}
\pgfpathlineto{\pgfpoint{-20.359195\du}{20.913689\du}}
\pgfpathlineto{\pgfpoint{-20.352030\du}{20.911002\du}}
\pgfpathlineto{\pgfpoint{-20.346059\du}{20.908613\du}}
\pgfpathlineto{\pgfpoint{-20.340087\du}{20.905926\du}}
\pgfpathlineto{\pgfpoint{-20.335012\du}{20.903239\du}}
\pgfpathlineto{\pgfpoint{-20.329638\du}{20.900552\du}}
\pgfpathlineto{\pgfpoint{-20.325160\du}{20.897567\du}}
\pgfpathlineto{\pgfpoint{-20.320980\du}{20.894581\du}}
\pgfpathlineto{\pgfpoint{-20.317099\du}{20.891297\du}}
\pgfpathlineto{\pgfpoint{-20.314412\du}{20.888312\du}}
\pgfpathlineto{\pgfpoint{-20.311128\du}{20.885028\du}}
\pgfpathlineto{\pgfpoint{-20.309038\du}{20.882042\du}}
\pgfpathlineto{\pgfpoint{-20.307247\du}{20.878758\du}}
\pgfpathlineto{\pgfpoint{-20.306053\du}{20.875474\du}}
\pgfpathlineto{\pgfpoint{-20.305157\du}{20.872190\du}}
\pgfpathlineto{\pgfpoint{-20.304858\du}{20.868906\du}}
\pgfusepath{fill}
\pgfsetbuttcap
\pgfsetmiterjoin
\pgfsetdash{}{0pt}
\definecolor{dialinecolor}{rgb}{0.666667, 0.901961, 1.000000}
\pgfsetstrokecolor{dialinecolor}
\pgfpathmoveto{\pgfpoint{-20.316800\du}{20.862636\du}}
\pgfpathlineto{\pgfpoint{-20.317099\du}{20.859651\du}}
\pgfpathlineto{\pgfpoint{-20.317696\du}{20.856964\du}}
\pgfpathlineto{\pgfpoint{-20.318890\du}{20.853978\du}}
\pgfpathlineto{\pgfpoint{-20.320980\du}{20.850396\du}}
\pgfpathlineto{\pgfpoint{-20.323070\du}{20.848007\du}}
\pgfpathlineto{\pgfpoint{-20.325160\du}{20.844723\du}}
\pgfpathlineto{\pgfpoint{-20.328444\du}{20.842036\du}}
\pgfpathlineto{\pgfpoint{-20.331728\du}{20.839349\du}}
\pgfpathlineto{\pgfpoint{-20.335908\du}{20.836364\du}}
\pgfpathlineto{\pgfpoint{-20.340685\du}{20.833975\du}}
\pgfpathlineto{\pgfpoint{-20.345760\du}{20.831288\du}}
\pgfpathlineto{\pgfpoint{-20.350835\du}{20.828900\du}}
\pgfpathlineto{\pgfpoint{-20.356508\du}{20.826512\du}}
\pgfpathlineto{\pgfpoint{-20.362479\du}{20.824123\du}}
\pgfpathlineto{\pgfpoint{-20.369047\du}{20.821436\du}}
\pgfpathlineto{\pgfpoint{-20.375914\du}{20.819943\du}}
\pgfpathlineto{\pgfpoint{-20.382482\du}{20.817854\du}}
\pgfpathlineto{\pgfpoint{-20.390543\du}{20.816062\du}}
\pgfpathlineto{\pgfpoint{-20.398007\du}{20.814569\du}}
\pgfpathlineto{\pgfpoint{-20.405769\du}{20.812480\du}}
\pgfpathlineto{\pgfpoint{-20.414725\du}{20.810987\du}}
\pgfpathlineto{\pgfpoint{-20.423085\du}{20.809494\du}}
\pgfpathlineto{\pgfpoint{-20.431444\du}{20.808300\du}}
\pgfpathlineto{\pgfpoint{-20.440700\du}{20.807106\du}}
\pgfpathlineto{\pgfpoint{-20.449358\du}{20.806210\du}}
\pgfpathlineto{\pgfpoint{-20.458911\du}{20.805314\du}}
\pgfpathlineto{\pgfpoint{-20.468465\du}{20.805016\du}}
\pgfpathlineto{\pgfpoint{-20.477720\du}{20.804419\du}}
\pgfpathlineto{\pgfpoint{-20.487572\du}{20.803822\du}}
\pgfpathlineto{\pgfpoint{-20.497424\du}{20.803523\du}}
\pgfpathlineto{\pgfpoint{-20.506978\du}{20.803523\du}}
\pgfpathlineto{\pgfpoint{-20.506978\du}{20.803523\du}}
\pgfpathlineto{\pgfpoint{-20.516233\du}{20.803523\du}}
\pgfpathlineto{\pgfpoint{-20.526085\du}{20.803822\du}}
\pgfpathlineto{\pgfpoint{-20.535639\du}{20.804419\du}}
\pgfpathlineto{\pgfpoint{-20.545193\du}{20.805016\du}}
\pgfpathlineto{\pgfpoint{-20.554746\du}{20.805314\du}}
\pgfpathlineto{\pgfpoint{-20.564001\du}{20.806210\du}}
\pgfpathlineto{\pgfpoint{-20.572958\du}{20.807106\du}}
\pgfpathlineto{\pgfpoint{-20.582213\du}{20.808300\du}}
\pgfpathlineto{\pgfpoint{-20.590573\du}{20.809494\du}}
\pgfpathlineto{\pgfpoint{-20.599529\du}{20.810987\du}}
\pgfpathlineto{\pgfpoint{-20.607889\du}{20.812480\du}}
\pgfpathlineto{\pgfpoint{-20.615651\du}{20.814569\du}}
\pgfpathlineto{\pgfpoint{-20.623712\du}{20.816062\du}}
\pgfpathlineto{\pgfpoint{-20.631176\du}{20.817854\du}}
\pgfpathlineto{\pgfpoint{-20.637744\du}{20.819943\du}}
\pgfpathlineto{\pgfpoint{-20.645208\du}{20.821436\du}}
\pgfpathlineto{\pgfpoint{-20.651179\du}{20.824123\du}}
\pgfpathlineto{\pgfpoint{-20.657150\du}{20.826512\du}}
\pgfpathlineto{\pgfpoint{-20.663419\du}{20.828900\du}}
\pgfpathlineto{\pgfpoint{-20.668495\du}{20.831288\du}}
\pgfpathlineto{\pgfpoint{-20.673570\du}{20.833975\du}}
\pgfpathlineto{\pgfpoint{-20.677750\du}{20.836364\du}}
\pgfpathlineto{\pgfpoint{-20.681930\du}{20.839349\du}}
\pgfpathlineto{\pgfpoint{-20.685811\du}{20.842036\du}}
\pgfpathlineto{\pgfpoint{-20.689095\du}{20.844723\du}}
\pgfpathlineto{\pgfpoint{-20.690886\du}{20.848007\du}}
\pgfpathlineto{\pgfpoint{-20.693275\du}{20.850396\du}}
\pgfpathlineto{\pgfpoint{-20.695066\du}{20.853978\du}}
\pgfpathlineto{\pgfpoint{-20.696559\du}{20.856964\du}}
\pgfpathlineto{\pgfpoint{-20.696857\du}{20.859651\du}}
\pgfpathlineto{\pgfpoint{-20.697454\du}{20.862636\du}}
\pgfpathlineto{\pgfpoint{-20.697454\du}{20.862636\du}}
\pgfpathlineto{\pgfpoint{-20.696857\du}{20.866219\du}}
\pgfpathlineto{\pgfpoint{-20.696559\du}{20.868906\du}}
\pgfpathlineto{\pgfpoint{-20.695066\du}{20.871891\du}}
\pgfpathlineto{\pgfpoint{-20.693275\du}{20.874877\du}}
\pgfpathlineto{\pgfpoint{-20.690886\du}{20.877564\du}}
\pgfpathlineto{\pgfpoint{-20.689095\du}{20.880848\du}}
\pgfpathlineto{\pgfpoint{-20.685811\du}{20.883535\du}}
\pgfpathlineto{\pgfpoint{-20.681930\du}{20.886520\du}}
\pgfpathlineto{\pgfpoint{-20.677750\du}{20.889207\du}}
\pgfpathlineto{\pgfpoint{-20.673570\du}{20.891894\du}}
\pgfpathlineto{\pgfpoint{-20.668495\du}{20.894581\du}}
\pgfpathlineto{\pgfpoint{-20.663419\du}{20.896074\du}}
\pgfpathlineto{\pgfpoint{-20.657150\du}{20.899358\du}}
\pgfpathlineto{\pgfpoint{-20.651179\du}{20.901448\du}}
\pgfpathlineto{\pgfpoint{-20.645208\du}{20.903538\du}}
\pgfpathlineto{\pgfpoint{-20.637744\du}{20.905926\du}}
\pgfpathlineto{\pgfpoint{-20.631176\du}{20.908016\du}}
\pgfpathlineto{\pgfpoint{-20.623712\du}{20.909808\du}}
\pgfpathlineto{\pgfpoint{-20.615651\du}{20.911002\du}}
\pgfpathlineto{\pgfpoint{-20.607889\du}{20.913390\du}}
\pgfpathlineto{\pgfpoint{-20.599529\du}{20.914883\du}}
\pgfpathlineto{\pgfpoint{-20.590573\du}{20.916376\du}}
\pgfpathlineto{\pgfpoint{-20.582213\du}{20.917570\du}}
\pgfpathlineto{\pgfpoint{-20.572958\du}{20.918466\du}}
\pgfpathlineto{\pgfpoint{-20.564001\du}{20.919361\du}}
\pgfpathlineto{\pgfpoint{-20.554746\du}{20.920555\du}}
\pgfpathlineto{\pgfpoint{-20.545193\du}{20.920854\du}}
\pgfpathlineto{\pgfpoint{-20.535639\du}{20.921451\du}}
\pgfpathlineto{\pgfpoint{-20.526085\du}{20.921750\du}}
\pgfpathlineto{\pgfpoint{-20.516233\du}{20.922048\du}}
\pgfpathlineto{\pgfpoint{-20.506978\du}{20.922048\du}}
\pgfpathlineto{\pgfpoint{-20.506978\du}{20.922048\du}}
\pgfpathlineto{\pgfpoint{-20.497424\du}{20.922048\du}}
\pgfpathlineto{\pgfpoint{-20.487572\du}{20.921750\du}}
\pgfpathlineto{\pgfpoint{-20.477720\du}{20.921451\du}}
\pgfpathlineto{\pgfpoint{-20.468465\du}{20.920854\du}}
\pgfpathlineto{\pgfpoint{-20.458911\du}{20.920555\du}}
\pgfpathlineto{\pgfpoint{-20.449358\du}{20.919361\du}}
\pgfpathlineto{\pgfpoint{-20.440700\du}{20.918466\du}}
\pgfpathlineto{\pgfpoint{-20.431444\du}{20.917570\du}}
\pgfpathlineto{\pgfpoint{-20.423085\du}{20.916376\du}}
\pgfpathlineto{\pgfpoint{-20.414725\du}{20.914883\du}}
\pgfpathlineto{\pgfpoint{-20.405769\du}{20.913390\du}}
\pgfpathlineto{\pgfpoint{-20.398007\du}{20.911002\du}}
\pgfpathlineto{\pgfpoint{-20.390543\du}{20.909808\du}}
\pgfpathlineto{\pgfpoint{-20.382482\du}{20.908016\du}}
\pgfpathlineto{\pgfpoint{-20.375914\du}{20.905926\du}}
\pgfpathlineto{\pgfpoint{-20.369047\du}{20.903538\du}}
\pgfpathlineto{\pgfpoint{-20.362479\du}{20.901448\du}}
\pgfpathlineto{\pgfpoint{-20.356508\du}{20.899358\du}}
\pgfpathlineto{\pgfpoint{-20.350835\du}{20.896074\du}}
\pgfpathlineto{\pgfpoint{-20.345760\du}{20.894581\du}}
\pgfpathlineto{\pgfpoint{-20.340685\du}{20.891894\du}}
\pgfpathlineto{\pgfpoint{-20.335908\du}{20.889207\du}}
\pgfpathlineto{\pgfpoint{-20.331728\du}{20.886520\du}}
\pgfpathlineto{\pgfpoint{-20.328444\du}{20.883535\du}}
\pgfpathlineto{\pgfpoint{-20.325160\du}{20.880848\du}}
\pgfpathlineto{\pgfpoint{-20.323070\du}{20.877564\du}}
\pgfpathlineto{\pgfpoint{-20.320980\du}{20.874877\du}}
\pgfpathlineto{\pgfpoint{-20.318890\du}{20.871891\du}}
\pgfpathlineto{\pgfpoint{-20.317696\du}{20.868906\du}}
\pgfpathlineto{\pgfpoint{-20.317099\du}{20.866219\du}}
\pgfpathlineto{\pgfpoint{-20.316800\du}{20.862636\du}}
\pgfusepath{stroke}
\pgfsetbuttcap
\pgfsetmiterjoin
\pgfsetdash{}{0pt}
\definecolor{dialinecolor}{rgb}{0.000000, 0.470588, 0.666667}
\pgfsetfillcolor{dialinecolor}
\pgfpathmoveto{\pgfpoint{-20.697454\du}{20.779639\du}}
\pgfpathlineto{\pgfpoint{-20.697454\du}{20.874877\du}}
\pgfpathlineto{\pgfpoint{-20.304560\du}{20.874877\du}}
\pgfpathlineto{\pgfpoint{-20.304560\du}{20.779639\du}}
\pgfpathlineto{\pgfpoint{-20.697454\du}{20.779639\du}}
\pgfusepath{fill}
\pgfsetbuttcap
\pgfsetmiterjoin
\pgfsetdash{}{0pt}
\definecolor{dialinecolor}{rgb}{0.000000, 0.705882, 1.000000}
\pgfsetfillcolor{dialinecolor}
\pgfpathmoveto{\pgfpoint{-20.304858\du}{20.773369\du}}
\pgfpathlineto{\pgfpoint{-20.305157\du}{20.769787\du}}
\pgfpathlineto{\pgfpoint{-20.306053\du}{20.766503\du}}
\pgfpathlineto{\pgfpoint{-20.307247\du}{20.763517\du}}
\pgfpathlineto{\pgfpoint{-20.309038\du}{20.759934\du}}
\pgfpathlineto{\pgfpoint{-20.311128\du}{20.756650\du}}
\pgfpathlineto{\pgfpoint{-20.314412\du}{20.753366\du}}
\pgfpathlineto{\pgfpoint{-20.317099\du}{20.750679\du}}
\pgfpathlineto{\pgfpoint{-20.320980\du}{20.746798\du}}
\pgfpathlineto{\pgfpoint{-20.325160\du}{20.744410\du}}
\pgfpathlineto{\pgfpoint{-20.329638\du}{20.741424\du}}
\pgfpathlineto{\pgfpoint{-20.335012\du}{20.738439\du}}
\pgfpathlineto{\pgfpoint{-20.340087\du}{20.735752\du}}
\pgfpathlineto{\pgfpoint{-20.346059\du}{20.733065\du}}
\pgfpathlineto{\pgfpoint{-20.352030\du}{20.730676\du}}
\pgfpathlineto{\pgfpoint{-20.359195\du}{20.727989\du}}
\pgfpathlineto{\pgfpoint{-20.366062\du}{20.725601\du}}
\pgfpathlineto{\pgfpoint{-20.373525\du}{20.723511\du}}
\pgfpathlineto{\pgfpoint{-20.380989\du}{20.721421\du}}
\pgfpathlineto{\pgfpoint{-20.389050\du}{20.719630\du}}
\pgfpathlineto{\pgfpoint{-20.397409\du}{20.717540\du}}
\pgfpathlineto{\pgfpoint{-20.405769\du}{20.715749\du}}
\pgfpathlineto{\pgfpoint{-20.414725\du}{20.714555\du}}
\pgfpathlineto{\pgfpoint{-20.423682\du}{20.713062\du}}
\pgfpathlineto{\pgfpoint{-20.432937\du}{20.711868\du}}
\pgfpathlineto{\pgfpoint{-20.442789\du}{20.710673\du}}
\pgfpathlineto{\pgfpoint{-20.452044\du}{20.709479\du}}
\pgfpathlineto{\pgfpoint{-20.461897\du}{20.709181\du}}
\pgfpathlineto{\pgfpoint{-20.471450\du}{20.708285\du}}
\pgfpathlineto{\pgfpoint{-20.481303\du}{20.707986\du}}
\pgfpathlineto{\pgfpoint{-20.491752\du}{20.707688\du}}
\pgfpathlineto{\pgfpoint{-20.501007\du}{20.707688\du}}
\pgfpathlineto{\pgfpoint{-20.501007\du}{20.707688\du}}
\pgfpathlineto{\pgfpoint{-20.510859\du}{20.707688\du}}
\pgfpathlineto{\pgfpoint{-20.521010\du}{20.707986\du}}
\pgfpathlineto{\pgfpoint{-20.530862\du}{20.708285\du}}
\pgfpathlineto{\pgfpoint{-20.540714\du}{20.709181\du}}
\pgfpathlineto{\pgfpoint{-20.550567\du}{20.709479\du}}
\pgfpathlineto{\pgfpoint{-20.559822\du}{20.710673\du}}
\pgfpathlineto{\pgfpoint{-20.569674\du}{20.711868\du}}
\pgfpathlineto{\pgfpoint{-20.578929\du}{20.713062\du}}
\pgfpathlineto{\pgfpoint{-20.587886\du}{20.714555\du}}
\pgfpathlineto{\pgfpoint{-20.596544\du}{20.715749\du}}
\pgfpathlineto{\pgfpoint{-20.604903\du}{20.717540\du}}
\pgfpathlineto{\pgfpoint{-20.613263\du}{20.719630\du}}
\pgfpathlineto{\pgfpoint{-20.621323\du}{20.721421\du}}
\pgfpathlineto{\pgfpoint{-20.628787\du}{20.723511\du}}
\pgfpathlineto{\pgfpoint{-20.636251\du}{20.725601\du}}
\pgfpathlineto{\pgfpoint{-20.643416\du}{20.727989\du}}
\pgfpathlineto{\pgfpoint{-20.650283\du}{20.730676\du}}
\pgfpathlineto{\pgfpoint{-20.656254\du}{20.733065\du}}
\pgfpathlineto{\pgfpoint{-20.662524\du}{20.735752\du}}
\pgfpathlineto{\pgfpoint{-20.667599\du}{20.738439\du}}
\pgfpathlineto{\pgfpoint{-20.672674\du}{20.741424\du}}
\pgfpathlineto{\pgfpoint{-20.677451\du}{20.744410\du}}
\pgfpathlineto{\pgfpoint{-20.681631\du}{20.746798\du}}
\pgfpathlineto{\pgfpoint{-20.685214\du}{20.750679\du}}
\pgfpathlineto{\pgfpoint{-20.688199\du}{20.753366\du}}
\pgfpathlineto{\pgfpoint{-20.690886\du}{20.756650\du}}
\pgfpathlineto{\pgfpoint{-20.693275\du}{20.759934\du}}
\pgfpathlineto{\pgfpoint{-20.695066\du}{20.763517\du}}
\pgfpathlineto{\pgfpoint{-20.696559\du}{20.766503\du}}
\pgfpathlineto{\pgfpoint{-20.697454\du}{20.769787\du}}
\pgfpathlineto{\pgfpoint{-20.697454\du}{20.773369\du}}
\pgfpathlineto{\pgfpoint{-20.697454\du}{20.773369\du}}
\pgfpathlineto{\pgfpoint{-20.697454\du}{20.776653\du}}
\pgfpathlineto{\pgfpoint{-20.696559\du}{20.779937\du}}
\pgfpathlineto{\pgfpoint{-20.695066\du}{20.782923\du}}
\pgfpathlineto{\pgfpoint{-20.693275\du}{20.786506\du}}
\pgfpathlineto{\pgfpoint{-20.690886\du}{20.789790\du}}
\pgfpathlineto{\pgfpoint{-20.688199\du}{20.792775\du}}
\pgfpathlineto{\pgfpoint{-20.685214\du}{20.795761\du}}
\pgfpathlineto{\pgfpoint{-20.681631\du}{20.799045\du}}
\pgfpathlineto{\pgfpoint{-20.677451\du}{20.802030\du}}
\pgfpathlineto{\pgfpoint{-20.672674\du}{20.805016\du}}
\pgfpathlineto{\pgfpoint{-20.667599\du}{20.808001\du}}
\pgfpathlineto{\pgfpoint{-20.662524\du}{20.810688\du}}
\pgfpathlineto{\pgfpoint{-20.656254\du}{20.813077\du}}
\pgfpathlineto{\pgfpoint{-20.650283\du}{20.815764\du}}
\pgfpathlineto{\pgfpoint{-20.643416\du}{20.818451\du}}
\pgfpathlineto{\pgfpoint{-20.636251\du}{20.820540\du}}
\pgfpathlineto{\pgfpoint{-20.628787\du}{20.822929\du}}
\pgfpathlineto{\pgfpoint{-20.621323\du}{20.825019\du}}
\pgfpathlineto{\pgfpoint{-20.613263\du}{20.826810\du}}
\pgfpathlineto{\pgfpoint{-20.604903\du}{20.828900\du}}
\pgfpathlineto{\pgfpoint{-20.596544\du}{20.830691\du}}
\pgfpathlineto{\pgfpoint{-20.587886\du}{20.831885\du}}
\pgfpathlineto{\pgfpoint{-20.578929\du}{20.833378\du}}
\pgfpathlineto{\pgfpoint{-20.569674\du}{20.834572\du}}
\pgfpathlineto{\pgfpoint{-20.559822\du}{20.835468\du}}
\pgfpathlineto{\pgfpoint{-20.550567\du}{20.836364\du}}
\pgfpathlineto{\pgfpoint{-20.540714\du}{20.837259\du}}
\pgfpathlineto{\pgfpoint{-20.530862\du}{20.838155\du}}
\pgfpathlineto{\pgfpoint{-20.521010\du}{20.838454\du}}
\pgfpathlineto{\pgfpoint{-20.510859\du}{20.838752\du}}
\pgfpathlineto{\pgfpoint{-20.501007\du}{20.838752\du}}
\pgfpathlineto{\pgfpoint{-20.501007\du}{20.838752\du}}
\pgfpathlineto{\pgfpoint{-20.491752\du}{20.838752\du}}
\pgfpathlineto{\pgfpoint{-20.481303\du}{20.838454\du}}
\pgfpathlineto{\pgfpoint{-20.471450\du}{20.838155\du}}
\pgfpathlineto{\pgfpoint{-20.461897\du}{20.837259\du}}
\pgfpathlineto{\pgfpoint{-20.452044\du}{20.836364\du}}
\pgfpathlineto{\pgfpoint{-20.442789\du}{20.835468\du}}
\pgfpathlineto{\pgfpoint{-20.432937\du}{20.834572\du}}
\pgfpathlineto{\pgfpoint{-20.423682\du}{20.833378\du}}
\pgfpathlineto{\pgfpoint{-20.414725\du}{20.831885\du}}
\pgfpathlineto{\pgfpoint{-20.405769\du}{20.830691\du}}
\pgfpathlineto{\pgfpoint{-20.397409\du}{20.828900\du}}
\pgfpathlineto{\pgfpoint{-20.389050\du}{20.826810\du}}
\pgfpathlineto{\pgfpoint{-20.380989\du}{20.825019\du}}
\pgfpathlineto{\pgfpoint{-20.373525\du}{20.822929\du}}
\pgfpathlineto{\pgfpoint{-20.366062\du}{20.820540\du}}
\pgfpathlineto{\pgfpoint{-20.359195\du}{20.818451\du}}
\pgfpathlineto{\pgfpoint{-20.352030\du}{20.815764\du}}
\pgfpathlineto{\pgfpoint{-20.346059\du}{20.813077\du}}
\pgfpathlineto{\pgfpoint{-20.340087\du}{20.810688\du}}
\pgfpathlineto{\pgfpoint{-20.335012\du}{20.808001\du}}
\pgfpathlineto{\pgfpoint{-20.329638\du}{20.805016\du}}
\pgfpathlineto{\pgfpoint{-20.325160\du}{20.802030\du}}
\pgfpathlineto{\pgfpoint{-20.320980\du}{20.799045\du}}
\pgfpathlineto{\pgfpoint{-20.317099\du}{20.795761\du}}
\pgfpathlineto{\pgfpoint{-20.314412\du}{20.792775\du}}
\pgfpathlineto{\pgfpoint{-20.311128\du}{20.789790\du}}
\pgfpathlineto{\pgfpoint{-20.309038\du}{20.786506\du}}
\pgfpathlineto{\pgfpoint{-20.307247\du}{20.782923\du}}
\pgfpathlineto{\pgfpoint{-20.306053\du}{20.779937\du}}
\pgfpathlineto{\pgfpoint{-20.305157\du}{20.776653\du}}
\pgfpathlineto{\pgfpoint{-20.304858\du}{20.773369\du}}
\pgfusepath{fill}
\pgfsetbuttcap
\pgfsetmiterjoin
\pgfsetdash{}{0pt}
\definecolor{dialinecolor}{rgb}{0.666667, 0.901961, 1.000000}
\pgfsetstrokecolor{dialinecolor}
\pgfpathmoveto{\pgfpoint{-20.316800\du}{20.767697\du}}
\pgfpathlineto{\pgfpoint{-20.317099\du}{20.764711\du}}
\pgfpathlineto{\pgfpoint{-20.317696\du}{20.761726\du}}
\pgfpathlineto{\pgfpoint{-20.318890\du}{20.759039\du}}
\pgfpathlineto{\pgfpoint{-20.320980\du}{20.755755\du}}
\pgfpathlineto{\pgfpoint{-20.323070\du}{20.753068\du}}
\pgfpathlineto{\pgfpoint{-20.325160\du}{20.749784\du}}
\pgfpathlineto{\pgfpoint{-20.328444\du}{20.746798\du}}
\pgfpathlineto{\pgfpoint{-20.331728\du}{20.744410\du}}
\pgfpathlineto{\pgfpoint{-20.335908\du}{20.741723\du}}
\pgfpathlineto{\pgfpoint{-20.340685\du}{20.739036\du}}
\pgfpathlineto{\pgfpoint{-20.345760\du}{20.736349\du}}
\pgfpathlineto{\pgfpoint{-20.350835\du}{20.733960\du}}
\pgfpathlineto{\pgfpoint{-20.356508\du}{20.730975\du}}
\pgfpathlineto{\pgfpoint{-20.362479\du}{20.729184\du}}
\pgfpathlineto{\pgfpoint{-20.369047\du}{20.727094\du}}
\pgfpathlineto{\pgfpoint{-20.375914\du}{20.724407\du}}
\pgfpathlineto{\pgfpoint{-20.382482\du}{20.722914\du}}
\pgfpathlineto{\pgfpoint{-20.390543\du}{20.721123\du}}
\pgfpathlineto{\pgfpoint{-20.398007\du}{20.719033\du}}
\pgfpathlineto{\pgfpoint{-20.405769\du}{20.717540\du}}
\pgfpathlineto{\pgfpoint{-20.414725\du}{20.716047\du}}
\pgfpathlineto{\pgfpoint{-20.423085\du}{20.714555\du}}
\pgfpathlineto{\pgfpoint{-20.431444\du}{20.713360\du}}
\pgfpathlineto{\pgfpoint{-20.440700\du}{20.712166\du}}
\pgfpathlineto{\pgfpoint{-20.449358\du}{20.710972\du}}
\pgfpathlineto{\pgfpoint{-20.458911\du}{20.710375\du}}
\pgfpathlineto{\pgfpoint{-20.468465\du}{20.709479\du}}
\pgfpathlineto{\pgfpoint{-20.477720\du}{20.709479\du}}
\pgfpathlineto{\pgfpoint{-20.487572\du}{20.709181\du}}
\pgfpathlineto{\pgfpoint{-20.497424\du}{20.708285\du}}
\pgfpathlineto{\pgfpoint{-20.506978\du}{20.708285\du}}
\pgfpathlineto{\pgfpoint{-20.506978\du}{20.708285\du}}
\pgfpathlineto{\pgfpoint{-20.516233\du}{20.708285\du}}
\pgfpathlineto{\pgfpoint{-20.526085\du}{20.709181\du}}
\pgfpathlineto{\pgfpoint{-20.535639\du}{20.709479\du}}
\pgfpathlineto{\pgfpoint{-20.545193\du}{20.709479\du}}
\pgfpathlineto{\pgfpoint{-20.554746\du}{20.710375\du}}
\pgfpathlineto{\pgfpoint{-20.564001\du}{20.710972\du}}
\pgfpathlineto{\pgfpoint{-20.572958\du}{20.712166\du}}
\pgfpathlineto{\pgfpoint{-20.582213\du}{20.713360\du}}
\pgfpathlineto{\pgfpoint{-20.590573\du}{20.714555\du}}
\pgfpathlineto{\pgfpoint{-20.599529\du}{20.716047\du}}
\pgfpathlineto{\pgfpoint{-20.607889\du}{20.717540\du}}
\pgfpathlineto{\pgfpoint{-20.615651\du}{20.719033\du}}
\pgfpathlineto{\pgfpoint{-20.623712\du}{20.721123\du}}
\pgfpathlineto{\pgfpoint{-20.631176\du}{20.722914\du}}
\pgfpathlineto{\pgfpoint{-20.637744\du}{20.724407\du}}
\pgfpathlineto{\pgfpoint{-20.645208\du}{20.727094\du}}
\pgfpathlineto{\pgfpoint{-20.651179\du}{20.729184\du}}
\pgfpathlineto{\pgfpoint{-20.657150\du}{20.730975\du}}
\pgfpathlineto{\pgfpoint{-20.663419\du}{20.733960\du}}
\pgfpathlineto{\pgfpoint{-20.668495\du}{20.736349\du}}
\pgfpathlineto{\pgfpoint{-20.673570\du}{20.739036\du}}
\pgfpathlineto{\pgfpoint{-20.677750\du}{20.741723\du}}
\pgfpathlineto{\pgfpoint{-20.681930\du}{20.744410\du}}
\pgfpathlineto{\pgfpoint{-20.685811\du}{20.746798\du}}
\pgfpathlineto{\pgfpoint{-20.689095\du}{20.749784\du}}
\pgfpathlineto{\pgfpoint{-20.690886\du}{20.753068\du}}
\pgfpathlineto{\pgfpoint{-20.693275\du}{20.755755\du}}
\pgfpathlineto{\pgfpoint{-20.695066\du}{20.759039\du}}
\pgfpathlineto{\pgfpoint{-20.696559\du}{20.761726\du}}
\pgfpathlineto{\pgfpoint{-20.696857\du}{20.764711\du}}
\pgfpathlineto{\pgfpoint{-20.697454\du}{20.767697\du}}
\pgfpathlineto{\pgfpoint{-20.697454\du}{20.767697\du}}
\pgfpathlineto{\pgfpoint{-20.696857\du}{20.771279\du}}
\pgfpathlineto{\pgfpoint{-20.696559\du}{20.773966\du}}
\pgfpathlineto{\pgfpoint{-20.695066\du}{20.776653\du}}
\pgfpathlineto{\pgfpoint{-20.693275\du}{20.779937\du}}
\pgfpathlineto{\pgfpoint{-20.690886\du}{20.782624\du}}
\pgfpathlineto{\pgfpoint{-20.689095\du}{20.786207\du}}
\pgfpathlineto{\pgfpoint{-20.685811\du}{20.788894\du}}
\pgfpathlineto{\pgfpoint{-20.681930\du}{20.791581\du}}
\pgfpathlineto{\pgfpoint{-20.677750\du}{20.794268\du}}
\pgfpathlineto{\pgfpoint{-20.673570\du}{20.796955\du}}
\pgfpathlineto{\pgfpoint{-20.668495\du}{20.799642\du}}
\pgfpathlineto{\pgfpoint{-20.663419\du}{20.801732\du}}
\pgfpathlineto{\pgfpoint{-20.657150\du}{20.804419\du}}
\pgfpathlineto{\pgfpoint{-20.651179\du}{20.806509\du}}
\pgfpathlineto{\pgfpoint{-20.645208\du}{20.808897\du}}
\pgfpathlineto{\pgfpoint{-20.637744\du}{20.810987\du}}
\pgfpathlineto{\pgfpoint{-20.631176\du}{20.813077\du}}
\pgfpathlineto{\pgfpoint{-20.623712\du}{20.814868\du}}
\pgfpathlineto{\pgfpoint{-20.615651\du}{20.816361\du}}
\pgfpathlineto{\pgfpoint{-20.607889\du}{20.818451\du}}
\pgfpathlineto{\pgfpoint{-20.599529\du}{20.819943\du}}
\pgfpathlineto{\pgfpoint{-20.590573\du}{20.821436\du}}
\pgfpathlineto{\pgfpoint{-20.582213\du}{20.822630\du}}
\pgfpathlineto{\pgfpoint{-20.572958\du}{20.823825\du}}
\pgfpathlineto{\pgfpoint{-20.564001\du}{20.824422\du}}
\pgfpathlineto{\pgfpoint{-20.554746\du}{20.825616\du}}
\pgfpathlineto{\pgfpoint{-20.545193\du}{20.825914\du}}
\pgfpathlineto{\pgfpoint{-20.535639\du}{20.826512\du}}
\pgfpathlineto{\pgfpoint{-20.526085\du}{20.826810\du}}
\pgfpathlineto{\pgfpoint{-20.516233\du}{20.827407\du}}
\pgfpathlineto{\pgfpoint{-20.506978\du}{20.827407\du}}
\pgfpathlineto{\pgfpoint{-20.506978\du}{20.827407\du}}
\pgfpathlineto{\pgfpoint{-20.497424\du}{20.827407\du}}
\pgfpathlineto{\pgfpoint{-20.487572\du}{20.826810\du}}
\pgfpathlineto{\pgfpoint{-20.477720\du}{20.826512\du}}
\pgfpathlineto{\pgfpoint{-20.468465\du}{20.825914\du}}
\pgfpathlineto{\pgfpoint{-20.458911\du}{20.825616\du}}
\pgfpathlineto{\pgfpoint{-20.449358\du}{20.824422\du}}
\pgfpathlineto{\pgfpoint{-20.440700\du}{20.823825\du}}
\pgfpathlineto{\pgfpoint{-20.431444\du}{20.822630\du}}
\pgfpathlineto{\pgfpoint{-20.423085\du}{20.821436\du}}
\pgfpathlineto{\pgfpoint{-20.414725\du}{20.819943\du}}
\pgfpathlineto{\pgfpoint{-20.405769\du}{20.818451\du}}
\pgfpathlineto{\pgfpoint{-20.398007\du}{20.816361\du}}
\pgfpathlineto{\pgfpoint{-20.390543\du}{20.814868\du}}
\pgfpathlineto{\pgfpoint{-20.382482\du}{20.813077\du}}
\pgfpathlineto{\pgfpoint{-20.375914\du}{20.810987\du}}
\pgfpathlineto{\pgfpoint{-20.369047\du}{20.808897\du}}
\pgfpathlineto{\pgfpoint{-20.362479\du}{20.806509\du}}
\pgfpathlineto{\pgfpoint{-20.356508\du}{20.804419\du}}
\pgfpathlineto{\pgfpoint{-20.350835\du}{20.801732\du}}
\pgfpathlineto{\pgfpoint{-20.345760\du}{20.799642\du}}
\pgfpathlineto{\pgfpoint{-20.340685\du}{20.796955\du}}
\pgfpathlineto{\pgfpoint{-20.335908\du}{20.794268\du}}
\pgfpathlineto{\pgfpoint{-20.331728\du}{20.791581\du}}
\pgfpathlineto{\pgfpoint{-20.328444\du}{20.788894\du}}
\pgfpathlineto{\pgfpoint{-20.325160\du}{20.786207\du}}
\pgfpathlineto{\pgfpoint{-20.323070\du}{20.782624\du}}
\pgfpathlineto{\pgfpoint{-20.320980\du}{20.779937\du}}
\pgfpathlineto{\pgfpoint{-20.318890\du}{20.776653\du}}
\pgfpathlineto{\pgfpoint{-20.317696\du}{20.773966\du}}
\pgfpathlineto{\pgfpoint{-20.317099\du}{20.771279\du}}
\pgfpathlineto{\pgfpoint{-20.316800\du}{20.767697\du}}
\pgfusepath{stroke}
\pgfsetbuttcap
\pgfsetmiterjoin
\pgfsetdash{}{0pt}
\definecolor{dialinecolor}{rgb}{0.000000, 0.000000, 0.000000}
\pgfsetfillcolor{dialinecolor}
\pgfpathmoveto{\pgfpoint{-20.495036\du}{20.755755\du}}
\pgfpathlineto{\pgfpoint{-20.471152\du}{20.767697\du}}
\pgfpathlineto{\pgfpoint{-20.399798\du}{20.744111\du}}
\pgfpathlineto{\pgfpoint{-20.363972\du}{20.755755\du}}
\pgfpathlineto{\pgfpoint{-20.387557\du}{20.731871\du}}
\pgfpathlineto{\pgfpoint{-20.471152\du}{20.731871\du}}
\pgfpathlineto{\pgfpoint{-20.435326\du}{20.731871\du}}
\pgfpathlineto{\pgfpoint{-20.495036\du}{20.755755\du}}
\pgfusepath{fill}
\pgfsetbuttcap
\pgfsetmiterjoin
\pgfsetdash{}{0pt}
\definecolor{dialinecolor}{rgb}{0.000000, 0.000000, 0.000000}
\pgfsetfillcolor{dialinecolor}
\pgfpathmoveto{\pgfpoint{-20.506978\du}{20.791581\du}}
\pgfpathlineto{\pgfpoint{-20.542506\du}{20.779639\du}}
\pgfpathlineto{\pgfpoint{-20.602216\du}{20.803523\du}}
\pgfpathlineto{\pgfpoint{-20.637744\du}{20.791581\du}}
\pgfpathlineto{\pgfpoint{-20.625503\du}{20.827407\du}}
\pgfpathlineto{\pgfpoint{-20.542506\du}{20.827407\du}}
\pgfpathlineto{\pgfpoint{-20.578033\du}{20.815764\du}}
\pgfpathlineto{\pgfpoint{-20.506978\du}{20.791581\du}}
\pgfusepath{fill}
\pgfsetbuttcap
\pgfsetmiterjoin
\pgfsetdash{}{0pt}
\definecolor{dialinecolor}{rgb}{0.000000, 0.000000, 0.000000}
\pgfsetfillcolor{dialinecolor}
\pgfpathmoveto{\pgfpoint{-20.625503\du}{20.731871\du}}
\pgfpathlineto{\pgfpoint{-20.602216\du}{20.720227\du}}
\pgfpathlineto{\pgfpoint{-20.530862\du}{20.755755\du}}
\pgfpathlineto{\pgfpoint{-20.506978\du}{20.744111\du}}
\pgfpathlineto{\pgfpoint{-20.518323\du}{20.767697\du}}
\pgfpathlineto{\pgfpoint{-20.602216\du}{20.767697\du}}
\pgfpathlineto{\pgfpoint{-20.566390\du}{20.755755\du}}
\pgfpathlineto{\pgfpoint{-20.625503\du}{20.731871\du}}
\pgfusepath{fill}
\pgfsetbuttcap
\pgfsetmiterjoin
\pgfsetdash{}{0pt}
\definecolor{dialinecolor}{rgb}{0.000000, 0.000000, 0.000000}
\pgfsetfillcolor{dialinecolor}
\pgfpathmoveto{\pgfpoint{-20.375914\du}{20.815764\du}}
\pgfpathlineto{\pgfpoint{-20.399798\du}{20.827407\du}}
\pgfpathlineto{\pgfpoint{-20.471152\du}{20.791581\du}}
\pgfpathlineto{\pgfpoint{-20.506978\du}{20.803523\du}}
\pgfpathlineto{\pgfpoint{-20.482795\du}{20.779639\du}}
\pgfpathlineto{\pgfpoint{-20.399798\du}{20.779639\du}}
\pgfpathlineto{\pgfpoint{-20.435326\du}{20.791581\du}}
\pgfpathlineto{\pgfpoint{-20.375914\du}{20.815764\du}}
\pgfusepath{fill}
\pgfsetbuttcap
\pgfsetmiterjoin
\pgfsetdash{}{0pt}
\definecolor{dialinecolor}{rgb}{1.000000, 1.000000, 1.000000}
\pgfsetfillcolor{dialinecolor}
\pgfpathmoveto{\pgfpoint{-20.495036\du}{20.767697\du}}
\pgfpathlineto{\pgfpoint{-20.471152\du}{20.767697\du}}
\pgfpathlineto{\pgfpoint{-20.399798\du}{20.744111\du}}
\pgfpathlineto{\pgfpoint{-20.363972\du}{20.755755\du}}
\pgfpathlineto{\pgfpoint{-20.387557\du}{20.731871\du}}
\pgfpathlineto{\pgfpoint{-20.458911\du}{20.731871\du}}
\pgfpathlineto{\pgfpoint{-20.435326\du}{20.731871\du}}
\pgfpathlineto{\pgfpoint{-20.495036\du}{20.767697\du}}
\pgfusepath{fill}
\pgfsetbuttcap
\pgfsetmiterjoin
\pgfsetdash{}{0pt}
\definecolor{dialinecolor}{rgb}{1.000000, 1.000000, 1.000000}
\pgfsetfillcolor{dialinecolor}
\pgfpathmoveto{\pgfpoint{-20.506978\du}{20.791581\du}}
\pgfpathlineto{\pgfpoint{-20.530862\du}{20.779639\du}}
\pgfpathlineto{\pgfpoint{-20.602216\du}{20.803523\du}}
\pgfpathlineto{\pgfpoint{-20.637744\du}{20.803523\du}}
\pgfpathlineto{\pgfpoint{-20.613860\du}{20.827407\du}}
\pgfpathlineto{\pgfpoint{-20.530862\du}{20.827407\du}}
\pgfpathlineto{\pgfpoint{-20.566390\du}{20.815764\du}}
\pgfpathlineto{\pgfpoint{-20.506978\du}{20.791581\du}}
\pgfusepath{fill}
\pgfsetbuttcap
\pgfsetmiterjoin
\pgfsetdash{}{0pt}
\definecolor{dialinecolor}{rgb}{1.000000, 1.000000, 1.000000}
\pgfsetfillcolor{dialinecolor}
\pgfpathmoveto{\pgfpoint{-20.625503\du}{20.731871\du}}
\pgfpathlineto{\pgfpoint{-20.602216\du}{20.731871\du}}
\pgfpathlineto{\pgfpoint{-20.530862\du}{20.755755\du}}
\pgfpathlineto{\pgfpoint{-20.495036\du}{20.744111\du}}
\pgfpathlineto{\pgfpoint{-20.518323\du}{20.767697\du}}
\pgfpathlineto{\pgfpoint{-20.602216\du}{20.767697\du}}
\pgfpathlineto{\pgfpoint{-20.566390\du}{20.767697\du}}
\pgfpathlineto{\pgfpoint{-20.625503\du}{20.731871\du}}
\pgfusepath{fill}
\pgfsetbuttcap
\pgfsetmiterjoin
\pgfsetdash{}{0pt}
\definecolor{dialinecolor}{rgb}{1.000000, 1.000000, 1.000000}
\pgfsetfillcolor{dialinecolor}
\pgfpathmoveto{\pgfpoint{-20.375914\du}{20.815764\du}}
\pgfpathlineto{\pgfpoint{-20.399798\du}{20.827407\du}}
\pgfpathlineto{\pgfpoint{-20.471152\du}{20.803523\du}}
\pgfpathlineto{\pgfpoint{-20.495036\du}{20.803523\du}}
\pgfpathlineto{\pgfpoint{-20.482795\du}{20.779639\du}}
\pgfpathlineto{\pgfpoint{-20.399798\du}{20.779639\du}}
\pgfpathlineto{\pgfpoint{-20.435326\du}{20.791581\du}}
\pgfpathlineto{\pgfpoint{-20.375914\du}{20.815764\du}}
\pgfusepath{fill}
\pgfsetbuttcap
\pgfsetmiterjoin
\pgfsetdash{}{0pt}
\definecolor{dialinecolor}{rgb}{0.666667, 0.901961, 1.000000}
\pgfsetstrokecolor{dialinecolor}
\pgfpathmoveto{\pgfpoint{-20.697454\du}{20.767697\du}}
\pgfpathlineto{\pgfpoint{-20.697454\du}{20.862636\du}}
\pgfusepath{stroke}
\pgfsetbuttcap
\pgfsetmiterjoin
\pgfsetdash{}{0pt}
\definecolor{dialinecolor}{rgb}{0.666667, 0.901961, 1.000000}
\pgfsetstrokecolor{dialinecolor}
\pgfpathmoveto{\pgfpoint{-20.316800\du}{20.767697\du}}
\pgfpathlineto{\pgfpoint{-20.316800\du}{20.862636\du}}
\pgfusepath{stroke}
\pgfsetlinewidth{0.100000\du}
\pgfsetdash{}{0pt}
\pgfsetdash{}{0pt}
\pgfsetbuttcap
\pgfsetmiterjoin
\pgfsetlinewidth{0.001000\du}
\pgfsetbuttcap
\pgfsetmiterjoin
\pgfsetdash{}{0pt}
\definecolor{dialinecolor}{rgb}{0.717647, 0.717647, 0.615686}
\pgfsetfillcolor{dialinecolor}
\pgfpathmoveto{\pgfpoint{-25.753700\du}{21.392280\du}}
\pgfpathlineto{\pgfpoint{-23.755274\du}{21.392280\du}}
\pgfpathlineto{\pgfpoint{-23.755274\du}{21.761613\du}}
\pgfpathlineto{\pgfpoint{-25.753700\du}{21.761613\du}}
\pgfpathlineto{\pgfpoint{-25.753700\du}{21.392280\du}}
\pgfusepath{fill}
\pgfsetbuttcap
\pgfsetmiterjoin
\pgfsetdash{}{0pt}
\definecolor{dialinecolor}{rgb}{0.286275, 0.286275, 0.211765}
\pgfsetstrokecolor{dialinecolor}
\pgfpathmoveto{\pgfpoint{-25.753700\du}{21.392280\du}}
\pgfpathlineto{\pgfpoint{-23.755274\du}{21.392280\du}}
\pgfpathlineto{\pgfpoint{-23.755274\du}{21.761613\du}}
\pgfpathlineto{\pgfpoint{-25.753700\du}{21.761613\du}}
\pgfpathlineto{\pgfpoint{-25.753700\du}{21.392280\du}}
\pgfusepath{stroke}
\pgfsetbuttcap
\pgfsetmiterjoin
\pgfsetdash{}{0pt}
\definecolor{dialinecolor}{rgb}{0.788235, 0.788235, 0.713726}
\pgfsetfillcolor{dialinecolor}
\pgfpathmoveto{\pgfpoint{-25.753700\du}{21.392280\du}}
\pgfpathlineto{\pgfpoint{-25.541798\du}{21.191398\du}}
\pgfpathlineto{\pgfpoint{-23.543373\du}{21.191398\du}}
\pgfpathlineto{\pgfpoint{-23.755274\du}{21.392280\du}}
\pgfpathlineto{\pgfpoint{-25.753700\du}{21.392280\du}}
\pgfusepath{fill}
\pgfsetbuttcap
\pgfsetmiterjoin
\pgfsetdash{}{0pt}
\definecolor{dialinecolor}{rgb}{0.286275, 0.286275, 0.211765}
\pgfsetstrokecolor{dialinecolor}
\pgfpathmoveto{\pgfpoint{-25.753700\du}{21.392280\du}}
\pgfpathlineto{\pgfpoint{-25.541798\du}{21.191398\du}}
\pgfpathlineto{\pgfpoint{-23.543373\du}{21.191398\du}}
\pgfpathlineto{\pgfpoint{-23.755274\du}{21.392280\du}}
\pgfpathlineto{\pgfpoint{-25.753700\du}{21.392280\du}}
\pgfusepath{stroke}
\pgfsetlinewidth{0.106000\du}
\pgfsetbuttcap
\pgfsetmiterjoin
\pgfsetdash{}{0pt}
\definecolor{dialinecolor}{rgb}{0.000000, 0.000000, 0.000000}
\pgfsetstrokecolor{dialinecolor}
\pgfpathmoveto{\pgfpoint{-23.867050\du}{21.560416\du}}
\pgfpathlineto{\pgfpoint{-24.346584\du}{21.560416\du}}
\pgfusepath{stroke}
\pgfsetlinewidth{0.001000\du}
\pgfsetbuttcap
\pgfsetmiterjoin
\pgfsetdash{}{0pt}
\definecolor{dialinecolor}{rgb}{0.478431, 0.478431, 0.352941}
\pgfsetfillcolor{dialinecolor}
\pgfpathmoveto{\pgfpoint{-23.755274\du}{21.761613\du}}
\pgfpathlineto{\pgfpoint{-23.543373\du}{21.549081\du}}
\pgfpathlineto{\pgfpoint{-23.543373\du}{21.191398\du}}
\pgfpathlineto{\pgfpoint{-23.755274\du}{21.392280\du}}
\pgfpathlineto{\pgfpoint{-23.755274\du}{21.761613\du}}
\pgfusepath{fill}
\pgfsetbuttcap
\pgfsetmiterjoin
\pgfsetdash{}{0pt}
\definecolor{dialinecolor}{rgb}{0.286275, 0.286275, 0.211765}
\pgfsetstrokecolor{dialinecolor}
\pgfpathmoveto{\pgfpoint{-23.755274\du}{21.761613\du}}
\pgfpathlineto{\pgfpoint{-23.543373\du}{21.549081\du}}
\pgfpathlineto{\pgfpoint{-23.543373\du}{21.191398\du}}
\pgfpathlineto{\pgfpoint{-23.755274\du}{21.392280\du}}
\pgfpathlineto{\pgfpoint{-23.755274\du}{21.761613\du}}
\pgfusepath{stroke}
\pgfsetbuttcap
\pgfsetmiterjoin
\pgfsetdash{}{0pt}
\definecolor{dialinecolor}{rgb}{0.788235, 0.788235, 0.713726}
\pgfsetfillcolor{dialinecolor}
\pgfpathmoveto{\pgfpoint{-25.742365\du}{21.995870\du}}
\pgfpathlineto{\pgfpoint{-25.519443\du}{21.716587\du}}
\pgfpathlineto{\pgfpoint{-23.978511\du}{21.716587\du}}
\pgfpathlineto{\pgfpoint{-24.201433\du}{21.995870\du}}
\pgfpathlineto{\pgfpoint{-25.742365\du}{21.995870\du}}
\pgfusepath{fill}
\pgfsetbuttcap
\pgfsetmiterjoin
\pgfsetdash{}{0pt}
\definecolor{dialinecolor}{rgb}{0.286275, 0.286275, 0.211765}
\pgfsetstrokecolor{dialinecolor}
\pgfpathmoveto{\pgfpoint{-25.742365\du}{21.995870\du}}
\pgfpathlineto{\pgfpoint{-25.519443\du}{21.716587\du}}
\pgfpathlineto{\pgfpoint{-23.978511\du}{21.716587\du}}
\pgfpathlineto{\pgfpoint{-24.201433\du}{21.995870\du}}
\pgfpathlineto{\pgfpoint{-25.742365\du}{21.995870\du}}
\pgfusepath{stroke}
\pgfsetbuttcap
\pgfsetmiterjoin
\pgfsetdash{}{0pt}
\definecolor{dialinecolor}{rgb}{0.478431, 0.478431, 0.352941}
\pgfsetfillcolor{dialinecolor}
\pgfpathmoveto{\pgfpoint{-24.201433\du}{22.051600\du}}
\pgfpathlineto{\pgfpoint{-23.978511\du}{21.817343\du}}
\pgfpathlineto{\pgfpoint{-23.978511\du}{21.716587\du}}
\pgfpathlineto{\pgfpoint{-24.201433\du}{21.995870\du}}
\pgfpathlineto{\pgfpoint{-24.201433\du}{22.051600\du}}
\pgfusepath{fill}
\pgfsetbuttcap
\pgfsetmiterjoin
\pgfsetdash{}{0pt}
\definecolor{dialinecolor}{rgb}{0.286275, 0.286275, 0.211765}
\pgfsetstrokecolor{dialinecolor}
\pgfpathmoveto{\pgfpoint{-24.201433\du}{22.051600\du}}
\pgfpathlineto{\pgfpoint{-23.978511\du}{21.817343\du}}
\pgfpathlineto{\pgfpoint{-23.978511\du}{21.716587\du}}
\pgfpathlineto{\pgfpoint{-24.201433\du}{21.995870\du}}
\pgfpathlineto{\pgfpoint{-24.201433\du}{22.051600\du}}
\pgfusepath{stroke}
\pgfsetbuttcap
\pgfsetmiterjoin
\pgfsetdash{}{0pt}
\definecolor{dialinecolor}{rgb}{0.717647, 0.717647, 0.615686}
\pgfsetfillcolor{dialinecolor}
\pgfpathmoveto{\pgfpoint{-25.742365\du}{21.995870\du}}
\pgfpathlineto{\pgfpoint{-24.201433\du}{21.995870\du}}
\pgfpathlineto{\pgfpoint{-24.201433\du}{22.051600\du}}
\pgfpathlineto{\pgfpoint{-25.742365\du}{22.051600\du}}
\pgfpathlineto{\pgfpoint{-25.742365\du}{21.995870\du}}
\pgfusepath{fill}
\pgfsetbuttcap
\pgfsetmiterjoin
\pgfsetdash{}{0pt}
\definecolor{dialinecolor}{rgb}{0.286275, 0.286275, 0.211765}
\pgfsetstrokecolor{dialinecolor}
\pgfpathmoveto{\pgfpoint{-25.742365\du}{21.995870\du}}
\pgfpathlineto{\pgfpoint{-24.201433\du}{21.995870\du}}
\pgfpathlineto{\pgfpoint{-24.201433\du}{22.051600\du}}
\pgfpathlineto{\pgfpoint{-25.742365\du}{22.051600\du}}
\pgfpathlineto{\pgfpoint{-25.742365\du}{21.995870\du}}
\pgfusepath{stroke}
\pgfsetbuttcap
\pgfsetmiterjoin
\pgfsetdash{}{0pt}
\definecolor{dialinecolor}{rgb}{0.000000, 0.000000, 0.000000}
\pgfsetfillcolor{dialinecolor}
\pgfpathmoveto{\pgfpoint{-25.452378\du}{21.347885\du}}
\pgfpathlineto{\pgfpoint{-25.284556\du}{21.191398\du}}
\pgfpathlineto{\pgfpoint{-23.867050\du}{21.191398\du}}
\pgfpathlineto{\pgfpoint{-24.022907\du}{21.347885\du}}
\pgfpathlineto{\pgfpoint{-25.452378\du}{21.347885\du}}
\pgfusepath{fill}
\pgfsetbuttcap
\pgfsetmiterjoin
\pgfsetdash{}{0pt}
\definecolor{dialinecolor}{rgb}{0.000000, 0.000000, 0.000000}
\pgfsetstrokecolor{dialinecolor}
\pgfpathmoveto{\pgfpoint{-25.452378\du}{21.347885\du}}
\pgfpathlineto{\pgfpoint{-25.284556\du}{21.191398\du}}
\pgfpathlineto{\pgfpoint{-23.867050\du}{21.191398\du}}
\pgfpathlineto{\pgfpoint{-24.022907\du}{21.347885\du}}
\pgfpathlineto{\pgfpoint{-25.452378\du}{21.347885\du}}
\pgfusepath{stroke}
\pgfsetbuttcap
\pgfsetmiterjoin
\pgfsetdash{}{0pt}
\definecolor{dialinecolor}{rgb}{0.788235, 0.788235, 0.713726}
\pgfsetfillcolor{dialinecolor}
\pgfpathmoveto{\pgfpoint{-25.463713\du}{20.196751\du}}
\pgfpathlineto{\pgfpoint{-25.306912\du}{20.051600\du}}
\pgfpathlineto{\pgfpoint{-23.888776\du}{20.051600\du}}
\pgfpathlineto{\pgfpoint{-24.045577\du}{20.196751\du}}
\pgfpathlineto{\pgfpoint{-25.463713\du}{20.196751\du}}
\pgfusepath{fill}
\pgfsetbuttcap
\pgfsetmiterjoin
\pgfsetdash{}{0pt}
\definecolor{dialinecolor}{rgb}{0.286275, 0.286275, 0.211765}
\pgfsetstrokecolor{dialinecolor}
\pgfpathmoveto{\pgfpoint{-25.463713\du}{20.196751\du}}
\pgfpathlineto{\pgfpoint{-25.306912\du}{20.051600\du}}
\pgfpathlineto{\pgfpoint{-23.888776\du}{20.051600\du}}
\pgfpathlineto{\pgfpoint{-24.045577\du}{20.196751\du}}
\pgfpathlineto{\pgfpoint{-25.463713\du}{20.196751\du}}
\pgfusepath{stroke}
\pgfsetbuttcap
\pgfsetmiterjoin
\pgfsetdash{}{0pt}
\definecolor{dialinecolor}{rgb}{0.717647, 0.717647, 0.615686}
\pgfsetfillcolor{dialinecolor}
\pgfpathmoveto{\pgfpoint{-25.463713\du}{20.196751\du}}
\pgfpathlineto{\pgfpoint{-24.034242\du}{20.196751\du}}
\pgfpathlineto{\pgfpoint{-24.034242\du}{21.325215\du}}
\pgfpathlineto{\pgfpoint{-25.463713\du}{21.325215\du}}
\pgfpathlineto{\pgfpoint{-25.463713\du}{20.196751\du}}
\pgfusepath{fill}
\pgfsetbuttcap
\pgfsetmiterjoin
\pgfsetdash{}{0pt}
\definecolor{dialinecolor}{rgb}{0.286275, 0.286275, 0.211765}
\pgfsetstrokecolor{dialinecolor}
\pgfpathmoveto{\pgfpoint{-25.463713\du}{20.196751\du}}
\pgfpathlineto{\pgfpoint{-24.034871\du}{20.196751\du}}
\pgfpathlineto{\pgfpoint{-24.034871\du}{21.324900\du}}
\pgfpathlineto{\pgfpoint{-25.463713\du}{21.324900\du}}
\pgfpathlineto{\pgfpoint{-25.463713\du}{20.196751\du}}
\pgfusepath{stroke}
\pgfsetbuttcap
\pgfsetmiterjoin
\pgfsetdash{}{0pt}
\definecolor{dialinecolor}{rgb}{1.000000, 1.000000, 1.000000}
\pgfsetfillcolor{dialinecolor}
\pgfpathmoveto{\pgfpoint{-25.340917\du}{20.341587\du}}
\pgfpathlineto{\pgfpoint{-24.157352\du}{20.341587\du}}
\pgfpathlineto{\pgfpoint{-24.157352\du}{21.213439\du}}
\pgfpathlineto{\pgfpoint{-25.340917\du}{21.213439\du}}
\pgfpathlineto{\pgfpoint{-25.340917\du}{20.341587\du}}
\pgfusepath{fill}
\pgfsetbuttcap
\pgfsetmiterjoin
\pgfsetdash{}{0pt}
\definecolor{dialinecolor}{rgb}{0.286275, 0.286275, 0.211765}
\pgfsetstrokecolor{dialinecolor}
\pgfpathmoveto{\pgfpoint{-25.340917\du}{20.341587\du}}
\pgfpathlineto{\pgfpoint{-24.157352\du}{20.341587\du}}
\pgfpathlineto{\pgfpoint{-24.157352\du}{21.213124\du}}
\pgfpathlineto{\pgfpoint{-25.340917\du}{21.213124\du}}
\pgfpathlineto{\pgfpoint{-25.340917\du}{20.341587\du}}
\pgfusepath{stroke}
\pgfsetbuttcap
\pgfsetmiterjoin
\pgfsetdash{}{0pt}
\definecolor{dialinecolor}{rgb}{0.478431, 0.478431, 0.352941}
\pgfsetfillcolor{dialinecolor}
\pgfpathmoveto{\pgfpoint{-24.045577\du}{21.314509\du}}
\pgfpathlineto{\pgfpoint{-23.888776\du}{21.158023\du}}
\pgfpathlineto{\pgfpoint{-23.888776\du}{20.051600\du}}
\pgfpathlineto{\pgfpoint{-24.045577\du}{20.196751\du}}
\pgfpathlineto{\pgfpoint{-24.045577\du}{21.314509\du}}
\pgfusepath{fill}
\pgfsetbuttcap
\pgfsetmiterjoin
\pgfsetdash{}{0pt}
\definecolor{dialinecolor}{rgb}{0.286275, 0.286275, 0.211765}
\pgfsetstrokecolor{dialinecolor}
\pgfpathmoveto{\pgfpoint{-24.045577\du}{21.314509\du}}
\pgfpathlineto{\pgfpoint{-23.888776\du}{21.158023\du}}
\pgfpathlineto{\pgfpoint{-23.888776\du}{20.051600\du}}
\pgfpathlineto{\pgfpoint{-24.045577\du}{20.196751\du}}
\pgfpathlineto{\pgfpoint{-24.045577\du}{21.314509\du}}
\pgfusepath{stroke}
\pgfsetlinewidth{0.100000\du}
\pgfsetdash{}{0pt}
\pgfsetdash{}{0pt}
\pgfsetbuttcap
\pgfsetmiterjoin
\pgfsetlinewidth{0.001000\du}
\pgfsetbuttcap
\pgfsetmiterjoin
\pgfsetdash{}{0pt}
\definecolor{dialinecolor}{rgb}{0.000000, 0.470588, 0.666667}
\pgfsetfillcolor{dialinecolor}
\pgfpathmoveto{\pgfpoint{-15.511355\du}{3.029107\du}}
\pgfpathlineto{\pgfpoint{-15.512648\du}{3.013851\du}}
\pgfpathlineto{\pgfpoint{-15.516010\du}{2.998337\du}}
\pgfpathlineto{\pgfpoint{-15.521698\du}{2.983081\du}}
\pgfpathlineto{\pgfpoint{-15.529714\du}{2.968083\du}}
\pgfpathlineto{\pgfpoint{-15.540057\du}{2.952828\du}}
\pgfpathlineto{\pgfpoint{-15.552469\du}{2.937830\du}}
\pgfpathlineto{\pgfpoint{-15.567207\du}{2.923350\du}}
\pgfpathlineto{\pgfpoint{-15.583756\du}{2.909128\du}}
\pgfpathlineto{\pgfpoint{-15.602891\du}{2.895165\du}}
\pgfpathlineto{\pgfpoint{-15.623835\du}{2.881720\du}}
\pgfpathlineto{\pgfpoint{-15.646848\du}{2.868274\du}}
\pgfpathlineto{\pgfpoint{-15.671930\du}{2.855345\du}}
\pgfpathlineto{\pgfpoint{-15.698822\du}{2.842933\du}}
\pgfpathlineto{\pgfpoint{-15.727523\du}{2.830780\du}}
\pgfpathlineto{\pgfpoint{-15.758035\du}{2.819403\du}}
\pgfpathlineto{\pgfpoint{-15.789840\du}{2.808543\du}}
\pgfpathlineto{\pgfpoint{-15.823454\du}{2.798200\du}}
\pgfpathlineto{\pgfpoint{-15.858879\du}{2.788633\du}}
\pgfpathlineto{\pgfpoint{-15.895338\du}{2.779066\du}}
\pgfpathlineto{\pgfpoint{-15.933090\du}{2.771050\du}}
\pgfpathlineto{\pgfpoint{-15.972135\du}{2.763034\du}}
\pgfpathlineto{\pgfpoint{-16.012472\du}{2.756052\du}}
\pgfpathlineto{\pgfpoint{-16.053586\du}{2.749329\du}}
\pgfpathlineto{\pgfpoint{-16.095992\du}{2.743641\du}}
\pgfpathlineto{\pgfpoint{-16.138657\du}{2.738728\du}}
\pgfpathlineto{\pgfpoint{-16.182614\du}{2.734591\du}}
\pgfpathlineto{\pgfpoint{-16.226831\du}{2.730971\du}}
\pgfpathlineto{\pgfpoint{-16.271305\du}{2.728126\du}}
\pgfpathlineto{\pgfpoint{-16.316297\du}{2.726058\du}}
\pgfpathlineto{\pgfpoint{-16.361548\du}{2.725282\du}}
\pgfpathlineto{\pgfpoint{-16.407315\du}{2.724506\du}}
\pgfpathlineto{\pgfpoint{-16.407315\du}{2.724506\du}}
\pgfpathlineto{\pgfpoint{-16.452566\du}{2.725282\du}}
\pgfpathlineto{\pgfpoint{-16.497558\du}{2.726058\du}}
\pgfpathlineto{\pgfpoint{-16.542550\du}{2.728126\du}}
\pgfpathlineto{\pgfpoint{-16.587025\du}{2.730971\du}}
\pgfpathlineto{\pgfpoint{-16.631241\du}{2.734591\du}}
\pgfpathlineto{\pgfpoint{-16.675199\du}{2.738728\du}}
\pgfpathlineto{\pgfpoint{-16.717863\du}{2.743641\du}}
\pgfpathlineto{\pgfpoint{-16.760270\du}{2.749329\du}}
\pgfpathlineto{\pgfpoint{-16.801383\du}{2.756052\du}}
\pgfpathlineto{\pgfpoint{-16.841720\du}{2.763034\du}}
\pgfpathlineto{\pgfpoint{-16.881024\du}{2.771050\du}}
\pgfpathlineto{\pgfpoint{-16.918517\du}{2.779066\du}}
\pgfpathlineto{\pgfpoint{-16.954976\du}{2.788633\du}}
\pgfpathlineto{\pgfpoint{-16.990401\du}{2.798200\du}}
\pgfpathlineto{\pgfpoint{-17.024015\du}{2.808543\du}}
\pgfpathlineto{\pgfpoint{-17.055820\du}{2.819403\du}}
\pgfpathlineto{\pgfpoint{-17.086332\du}{2.830780\du}}
\pgfpathlineto{\pgfpoint{-17.115034\du}{2.842933\du}}
\pgfpathlineto{\pgfpoint{-17.141925\du}{2.855345\du}}
\pgfpathlineto{\pgfpoint{-17.167266\du}{2.868274\du}}
\pgfpathlineto{\pgfpoint{-17.190020\du}{2.881720\du}}
\pgfpathlineto{\pgfpoint{-17.211223\du}{2.895165\du}}
\pgfpathlineto{\pgfpoint{-17.230099\du}{2.909128\du}}
\pgfpathlineto{\pgfpoint{-17.246648\du}{2.923350\du}}
\pgfpathlineto{\pgfpoint{-17.261387\du}{2.937830\du}}
\pgfpathlineto{\pgfpoint{-17.273798\du}{2.952828\du}}
\pgfpathlineto{\pgfpoint{-17.284141\du}{2.968083\du}}
\pgfpathlineto{\pgfpoint{-17.292157\du}{2.983081\du}}
\pgfpathlineto{\pgfpoint{-17.298104\du}{2.998337\du}}
\pgfpathlineto{\pgfpoint{-17.301207\du}{3.013851\du}}
\pgfpathlineto{\pgfpoint{-17.302500\du}{3.029107\du}}
\pgfpathlineto{\pgfpoint{-17.302500\du}{3.029107\du}}
\pgfpathlineto{\pgfpoint{-17.301207\du}{3.044880\du}}
\pgfpathlineto{\pgfpoint{-17.298104\du}{3.060136\du}}
\pgfpathlineto{\pgfpoint{-17.292157\du}{3.075392\du}}
\pgfpathlineto{\pgfpoint{-17.284141\du}{3.090648\du}}
\pgfpathlineto{\pgfpoint{-17.273798\du}{3.105645\du}}
\pgfpathlineto{\pgfpoint{-17.261387\du}{3.120642\du}}
\pgfpathlineto{\pgfpoint{-17.246648\du}{3.135122\du}}
\pgfpathlineto{\pgfpoint{-17.230099\du}{3.149344\du}}
\pgfpathlineto{\pgfpoint{-17.211223\du}{3.163307\du}}
\pgfpathlineto{\pgfpoint{-17.190020\du}{3.177012\du}}
\pgfpathlineto{\pgfpoint{-17.167266\du}{3.189940\du}}
\pgfpathlineto{\pgfpoint{-17.141925\du}{3.203386\du}}
\pgfpathlineto{\pgfpoint{-17.115034\du}{3.215798\du}}
\pgfpathlineto{\pgfpoint{-17.086332\du}{3.227692\du}}
\pgfpathlineto{\pgfpoint{-17.055820\du}{3.238811\du}}
\pgfpathlineto{\pgfpoint{-17.024015\du}{3.249929\du}}
\pgfpathlineto{\pgfpoint{-16.990401\du}{3.260272\du}}
\pgfpathlineto{\pgfpoint{-16.954976\du}{3.270098\du}}
\pgfpathlineto{\pgfpoint{-16.918517\du}{3.279148\du}}
\pgfpathlineto{\pgfpoint{-16.881024\du}{3.287681\du}}
\pgfpathlineto{\pgfpoint{-16.841720\du}{3.295439\du}}
\pgfpathlineto{\pgfpoint{-16.801383\du}{3.302420\du}}
\pgfpathlineto{\pgfpoint{-16.760270\du}{3.309143\du}}
\pgfpathlineto{\pgfpoint{-16.717863\du}{3.314832\du}}
\pgfpathlineto{\pgfpoint{-16.675199\du}{3.319745\du}}
\pgfpathlineto{\pgfpoint{-16.631241\du}{3.324140\du}}
\pgfpathlineto{\pgfpoint{-16.587025\du}{3.327243\du}}
\pgfpathlineto{\pgfpoint{-16.542550\du}{3.330088\du}}
\pgfpathlineto{\pgfpoint{-16.497558\du}{3.332156\du}}
\pgfpathlineto{\pgfpoint{-16.452566\du}{3.333449\du}}
\pgfpathlineto{\pgfpoint{-16.407315\du}{3.333708\du}}
\pgfpathlineto{\pgfpoint{-16.407315\du}{3.333708\du}}
\pgfpathlineto{\pgfpoint{-16.361548\du}{3.333449\du}}
\pgfpathlineto{\pgfpoint{-16.316297\du}{3.332156\du}}
\pgfpathlineto{\pgfpoint{-16.271305\du}{3.330088\du}}
\pgfpathlineto{\pgfpoint{-16.226831\du}{3.327243\du}}
\pgfpathlineto{\pgfpoint{-16.182614\du}{3.324140\du}}
\pgfpathlineto{\pgfpoint{-16.138657\du}{3.319745\du}}
\pgfpathlineto{\pgfpoint{-16.095992\du}{3.314832\du}}
\pgfpathlineto{\pgfpoint{-16.053586\du}{3.309143\du}}
\pgfpathlineto{\pgfpoint{-16.012472\du}{3.302420\du}}
\pgfpathlineto{\pgfpoint{-15.972135\du}{3.295439\du}}
\pgfpathlineto{\pgfpoint{-15.933090\du}{3.287681\du}}
\pgfpathlineto{\pgfpoint{-15.895338\du}{3.279148\du}}
\pgfpathlineto{\pgfpoint{-15.858879\du}{3.270098\du}}
\pgfpathlineto{\pgfpoint{-15.823454\du}{3.260272\du}}
\pgfpathlineto{\pgfpoint{-15.789840\du}{3.249929\du}}
\pgfpathlineto{\pgfpoint{-15.758035\du}{3.238811\du}}
\pgfpathlineto{\pgfpoint{-15.727523\du}{3.227692\du}}
\pgfpathlineto{\pgfpoint{-15.698822\du}{3.215798\du}}
\pgfpathlineto{\pgfpoint{-15.671930\du}{3.203386\du}}
\pgfpathlineto{\pgfpoint{-15.646848\du}{3.189940\du}}
\pgfpathlineto{\pgfpoint{-15.623835\du}{3.177012\du}}
\pgfpathlineto{\pgfpoint{-15.602891\du}{3.163307\du}}
\pgfpathlineto{\pgfpoint{-15.583756\du}{3.149344\du}}
\pgfpathlineto{\pgfpoint{-15.567207\du}{3.135122\du}}
\pgfpathlineto{\pgfpoint{-15.552469\du}{3.120642\du}}
\pgfpathlineto{\pgfpoint{-15.540057\du}{3.105645\du}}
\pgfpathlineto{\pgfpoint{-15.529714\du}{3.090648\du}}
\pgfpathlineto{\pgfpoint{-15.521698\du}{3.075392\du}}
\pgfpathlineto{\pgfpoint{-15.516010\du}{3.060136\du}}
\pgfpathlineto{\pgfpoint{-15.512648\du}{3.044880\du}}
\pgfpathlineto{\pgfpoint{-15.511355\du}{3.029107\du}}
\pgfusepath{fill}
\pgfsetbuttcap
\pgfsetmiterjoin
\pgfsetdash{}{0pt}
\definecolor{dialinecolor}{rgb}{0.666667, 0.901961, 1.000000}
\pgfsetstrokecolor{dialinecolor}
\pgfpathmoveto{\pgfpoint{-15.521698\du}{3.024453\du}}
\pgfpathlineto{\pgfpoint{-15.522991\du}{3.008938\du}}
\pgfpathlineto{\pgfpoint{-15.526353\du}{2.994199\du}}
\pgfpathlineto{\pgfpoint{-15.532041\du}{2.978944\du}}
\pgfpathlineto{\pgfpoint{-15.539798\du}{2.964205\du}}
\pgfpathlineto{\pgfpoint{-15.550141\du}{2.949207\du}}
\pgfpathlineto{\pgfpoint{-15.562811\du}{2.934727\du}}
\pgfpathlineto{\pgfpoint{-15.576775\du}{2.920506\du}}
\pgfpathlineto{\pgfpoint{-15.593840\du}{2.906284\du}}
\pgfpathlineto{\pgfpoint{-15.612716\du}{2.892580\du}}
\pgfpathlineto{\pgfpoint{-15.633661\du}{2.878875\du}}
\pgfpathlineto{\pgfpoint{-15.656415\du}{2.866205\du}}
\pgfpathlineto{\pgfpoint{-15.681497\du}{2.853535\du}}
\pgfpathlineto{\pgfpoint{-15.707872\du}{2.840865\du}}
\pgfpathlineto{\pgfpoint{-15.736315\du}{2.829229\du}}
\pgfpathlineto{\pgfpoint{-15.766827\du}{2.817852\du}}
\pgfpathlineto{\pgfpoint{-15.798890\du}{2.807509\du}}
\pgfpathlineto{\pgfpoint{-15.832246\du}{2.797166\du}}
\pgfpathlineto{\pgfpoint{-15.867154\du}{2.787598\du}}
\pgfpathlineto{\pgfpoint{-15.903095\du}{2.778807\du}}
\pgfpathlineto{\pgfpoint{-15.941106\du}{2.770274\du}}
\pgfpathlineto{\pgfpoint{-15.980151\du}{2.762775\du}}
\pgfpathlineto{\pgfpoint{-16.019712\du}{2.755794\du}}
\pgfpathlineto{\pgfpoint{-16.060826\du}{2.749329\du}}
\pgfpathlineto{\pgfpoint{-16.102715\du}{2.743641\du}}
\pgfpathlineto{\pgfpoint{-16.145380\du}{2.738728\du}}
\pgfpathlineto{\pgfpoint{-16.189079\du}{2.734591\du}}
\pgfpathlineto{\pgfpoint{-16.232778\du}{2.730971\du}}
\pgfpathlineto{\pgfpoint{-16.276994\du}{2.728902\du}}
\pgfpathlineto{\pgfpoint{-16.321986\du}{2.726575\du}}
\pgfpathlineto{\pgfpoint{-16.366978\du}{2.725282\du}}
\pgfpathlineto{\pgfpoint{-16.412228\du}{2.725282\du}}
\pgfpathlineto{\pgfpoint{-16.412228\du}{2.725282\du}}
\pgfpathlineto{\pgfpoint{-16.457220\du}{2.725282\du}}
\pgfpathlineto{\pgfpoint{-16.502471\du}{2.726575\du}}
\pgfpathlineto{\pgfpoint{-16.547204\du}{2.728902\du}}
\pgfpathlineto{\pgfpoint{-16.591420\du}{2.730971\du}}
\pgfpathlineto{\pgfpoint{-16.635378\du}{2.734591\du}}
\pgfpathlineto{\pgfpoint{-16.678819\du}{2.738728\du}}
\pgfpathlineto{\pgfpoint{-16.721483\du}{2.743641\du}}
\pgfpathlineto{\pgfpoint{-16.763631\du}{2.749329\du}}
\pgfpathlineto{\pgfpoint{-16.804486\du}{2.755794\du}}
\pgfpathlineto{\pgfpoint{-16.844048\du}{2.762775\du}}
\pgfpathlineto{\pgfpoint{-16.883351\du}{2.770274\du}}
\pgfpathlineto{\pgfpoint{-16.920844\du}{2.778807\du}}
\pgfpathlineto{\pgfpoint{-16.956786\du}{2.787598\du}}
\pgfpathlineto{\pgfpoint{-16.991952\du}{2.797166\du}}
\pgfpathlineto{\pgfpoint{-17.025825\du}{2.807509\du}}
\pgfpathlineto{\pgfpoint{-17.057889\du}{2.817852\du}}
\pgfpathlineto{\pgfpoint{-17.087883\du}{2.829229\du}}
\pgfpathlineto{\pgfpoint{-17.116585\du}{2.840865\du}}
\pgfpathlineto{\pgfpoint{-17.143218\du}{2.853535\du}}
\pgfpathlineto{\pgfpoint{-17.167783\du}{2.866205\du}}
\pgfpathlineto{\pgfpoint{-17.190279\du}{2.878875\du}}
\pgfpathlineto{\pgfpoint{-17.211482\du}{2.892580\du}}
\pgfpathlineto{\pgfpoint{-17.230616\du}{2.906284\du}}
\pgfpathlineto{\pgfpoint{-17.247165\du}{2.920506\du}}
\pgfpathlineto{\pgfpoint{-17.261645\du}{2.934727\du}}
\pgfpathlineto{\pgfpoint{-17.274574\du}{2.949207\du}}
\pgfpathlineto{\pgfpoint{-17.284141\du}{2.964205\du}}
\pgfpathlineto{\pgfpoint{-17.292416\du}{2.978944\du}}
\pgfpathlineto{\pgfpoint{-17.298104\du}{2.994199\du}}
\pgfpathlineto{\pgfpoint{-17.301207\du}{3.008938\du}}
\pgfpathlineto{\pgfpoint{-17.302500\du}{3.024453\du}}
\pgfpathlineto{\pgfpoint{-17.302500\du}{3.024453\du}}
\pgfpathlineto{\pgfpoint{-17.301207\du}{3.039709\du}}
\pgfpathlineto{\pgfpoint{-17.298104\du}{3.054706\du}}
\pgfpathlineto{\pgfpoint{-17.292416\du}{3.069703\du}}
\pgfpathlineto{\pgfpoint{-17.284141\du}{3.084442\du}}
\pgfpathlineto{\pgfpoint{-17.274574\du}{3.099439\du}}
\pgfpathlineto{\pgfpoint{-17.261645\du}{3.113919\du}}
\pgfpathlineto{\pgfpoint{-17.247165\du}{3.128141\du}}
\pgfpathlineto{\pgfpoint{-17.230616\du}{3.141845\du}}
\pgfpathlineto{\pgfpoint{-17.211482\du}{3.156067\du}}
\pgfpathlineto{\pgfpoint{-17.190279\du}{3.169771\du}}
\pgfpathlineto{\pgfpoint{-17.167783\du}{3.182700\du}}
\pgfpathlineto{\pgfpoint{-17.143218\du}{3.195112\du}}
\pgfpathlineto{\pgfpoint{-17.116585\du}{3.207523\du}}
\pgfpathlineto{\pgfpoint{-17.087883\du}{3.219418\du}}
\pgfpathlineto{\pgfpoint{-17.057889\du}{3.230536\du}}
\pgfpathlineto{\pgfpoint{-17.025825\du}{3.241138\du}}
\pgfpathlineto{\pgfpoint{-16.991952\du}{3.250964\du}}
\pgfpathlineto{\pgfpoint{-16.956786\du}{3.261048\du}}
\pgfpathlineto{\pgfpoint{-16.920844\du}{3.269840\du}}
\pgfpathlineto{\pgfpoint{-16.883351\du}{3.278114\du}}
\pgfpathlineto{\pgfpoint{-16.844048\du}{3.286130\du}}
\pgfpathlineto{\pgfpoint{-16.804486\du}{3.292853\du}}
\pgfpathlineto{\pgfpoint{-16.763631\du}{3.299059\du}}
\pgfpathlineto{\pgfpoint{-16.721483\du}{3.305006\du}}
\pgfpathlineto{\pgfpoint{-16.678819\du}{3.309660\du}}
\pgfpathlineto{\pgfpoint{-16.635378\du}{3.313539\du}}
\pgfpathlineto{\pgfpoint{-16.591420\du}{3.317417\du}}
\pgfpathlineto{\pgfpoint{-16.547204\du}{3.319745\du}}
\pgfpathlineto{\pgfpoint{-16.502471\du}{3.321813\du}}
\pgfpathlineto{\pgfpoint{-16.457220\du}{3.322847\du}}
\pgfpathlineto{\pgfpoint{-16.412228\du}{3.323365\du}}
\pgfpathlineto{\pgfpoint{-16.412228\du}{3.323365\du}}
\pgfpathlineto{\pgfpoint{-16.366978\du}{3.322847\du}}
\pgfpathlineto{\pgfpoint{-16.321986\du}{3.321813\du}}
\pgfpathlineto{\pgfpoint{-16.276994\du}{3.319745\du}}
\pgfpathlineto{\pgfpoint{-16.232778\du}{3.317417\du}}
\pgfpathlineto{\pgfpoint{-16.189079\du}{3.313539\du}}
\pgfpathlineto{\pgfpoint{-16.145380\du}{3.309660\du}}
\pgfpathlineto{\pgfpoint{-16.102715\du}{3.305006\du}}
\pgfpathlineto{\pgfpoint{-16.060826\du}{3.299059\du}}
\pgfpathlineto{\pgfpoint{-16.019712\du}{3.292853\du}}
\pgfpathlineto{\pgfpoint{-15.980151\du}{3.286130\du}}
\pgfpathlineto{\pgfpoint{-15.941106\du}{3.278114\du}}
\pgfpathlineto{\pgfpoint{-15.903095\du}{3.269840\du}}
\pgfpathlineto{\pgfpoint{-15.867154\du}{3.261048\du}}
\pgfpathlineto{\pgfpoint{-15.832246\du}{3.250964\du}}
\pgfpathlineto{\pgfpoint{-15.798890\du}{3.241138\du}}
\pgfpathlineto{\pgfpoint{-15.766827\du}{3.230536\du}}
\pgfpathlineto{\pgfpoint{-15.736315\du}{3.219418\du}}
\pgfpathlineto{\pgfpoint{-15.707872\du}{3.207523\du}}
\pgfpathlineto{\pgfpoint{-15.681497\du}{3.195112\du}}
\pgfpathlineto{\pgfpoint{-15.656415\du}{3.182700\du}}
\pgfpathlineto{\pgfpoint{-15.633661\du}{3.169771\du}}
\pgfpathlineto{\pgfpoint{-15.612716\du}{3.156067\du}}
\pgfpathlineto{\pgfpoint{-15.593840\du}{3.141845\du}}
\pgfpathlineto{\pgfpoint{-15.576775\du}{3.128141\du}}
\pgfpathlineto{\pgfpoint{-15.562811\du}{3.113919\du}}
\pgfpathlineto{\pgfpoint{-15.550141\du}{3.099439\du}}
\pgfpathlineto{\pgfpoint{-15.539798\du}{3.084442\du}}
\pgfpathlineto{\pgfpoint{-15.532041\du}{3.069703\du}}
\pgfpathlineto{\pgfpoint{-15.526353\du}{3.054706\du}}
\pgfpathlineto{\pgfpoint{-15.522991\du}{3.039709\du}}
\pgfpathlineto{\pgfpoint{-15.521698\du}{3.024453\du}}
\pgfusepath{stroke}
\pgfsetbuttcap
\pgfsetmiterjoin
\pgfsetdash{}{0pt}
\definecolor{dialinecolor}{rgb}{0.000000, 0.470588, 0.666667}
\pgfsetfillcolor{dialinecolor}
\pgfpathmoveto{\pgfpoint{-17.302500\du}{2.607631\du}}
\pgfpathlineto{\pgfpoint{-17.302500\du}{3.034796\du}}
\pgfpathlineto{\pgfpoint{-15.521698\du}{3.034796\du}}
\pgfpathlineto{\pgfpoint{-15.521698\du}{2.607631\du}}
\pgfpathlineto{\pgfpoint{-17.302500\du}{2.607631\du}}
\pgfusepath{fill}
\pgfsetbuttcap
\pgfsetmiterjoin
\pgfsetdash{}{0pt}
\definecolor{dialinecolor}{rgb}{0.000000, 0.705882, 1.000000}
\pgfsetfillcolor{dialinecolor}
\pgfpathmoveto{\pgfpoint{-15.511355\du}{2.601683\du}}
\pgfpathlineto{\pgfpoint{-15.512648\du}{2.586686\du}}
\pgfpathlineto{\pgfpoint{-15.516010\du}{2.571172\du}}
\pgfpathlineto{\pgfpoint{-15.521698\du}{2.555916\du}}
\pgfpathlineto{\pgfpoint{-15.529714\du}{2.540660\du}}
\pgfpathlineto{\pgfpoint{-15.540057\du}{2.525404\du}}
\pgfpathlineto{\pgfpoint{-15.552469\du}{2.510665\du}}
\pgfpathlineto{\pgfpoint{-15.567207\du}{2.496185\du}}
\pgfpathlineto{\pgfpoint{-15.583756\du}{2.481705\du}}
\pgfpathlineto{\pgfpoint{-15.602891\du}{2.468001\du}}
\pgfpathlineto{\pgfpoint{-15.623835\du}{2.454296\du}}
\pgfpathlineto{\pgfpoint{-15.646848\du}{2.441109\du}}
\pgfpathlineto{\pgfpoint{-15.671930\du}{2.427663\du}}
\pgfpathlineto{\pgfpoint{-15.698822\du}{2.415510\du}}
\pgfpathlineto{\pgfpoint{-15.727523\du}{2.403357\du}}
\pgfpathlineto{\pgfpoint{-15.758035\du}{2.391980\du}}
\pgfpathlineto{\pgfpoint{-15.789840\du}{2.381378\du}}
\pgfpathlineto{\pgfpoint{-15.823454\du}{2.371035\du}}
\pgfpathlineto{\pgfpoint{-15.858879\du}{2.361209\du}}
\pgfpathlineto{\pgfpoint{-15.895338\du}{2.352159\du}}
\pgfpathlineto{\pgfpoint{-15.933090\du}{2.343626\du}}
\pgfpathlineto{\pgfpoint{-15.972135\du}{2.335610\du}}
\pgfpathlineto{\pgfpoint{-16.012472\du}{2.328888\du}}
\pgfpathlineto{\pgfpoint{-16.053586\du}{2.321906\du}}
\pgfpathlineto{\pgfpoint{-16.095992\du}{2.316217\du}}
\pgfpathlineto{\pgfpoint{-16.138657\du}{2.311304\du}}
\pgfpathlineto{\pgfpoint{-16.182614\du}{2.307167\du}}
\pgfpathlineto{\pgfpoint{-16.226831\du}{2.304064\du}}
\pgfpathlineto{\pgfpoint{-16.271305\du}{2.301220\du}}
\pgfpathlineto{\pgfpoint{-16.316297\du}{2.299151\du}}
\pgfpathlineto{\pgfpoint{-16.361548\du}{2.297600\du}}
\pgfpathlineto{\pgfpoint{-16.407315\du}{2.297600\du}}
\pgfpathlineto{\pgfpoint{-16.407315\du}{2.297600\du}}
\pgfpathlineto{\pgfpoint{-16.452566\du}{2.297600\du}}
\pgfpathlineto{\pgfpoint{-16.497558\du}{2.299151\du}}
\pgfpathlineto{\pgfpoint{-16.542550\du}{2.301220\du}}
\pgfpathlineto{\pgfpoint{-16.587025\du}{2.304064\du}}
\pgfpathlineto{\pgfpoint{-16.631241\du}{2.307167\du}}
\pgfpathlineto{\pgfpoint{-16.675199\du}{2.311304\du}}
\pgfpathlineto{\pgfpoint{-16.717863\du}{2.316217\du}}
\pgfpathlineto{\pgfpoint{-16.760270\du}{2.321906\du}}
\pgfpathlineto{\pgfpoint{-16.801383\du}{2.328888\du}}
\pgfpathlineto{\pgfpoint{-16.841720\du}{2.335610\du}}
\pgfpathlineto{\pgfpoint{-16.881024\du}{2.343626\du}}
\pgfpathlineto{\pgfpoint{-16.918517\du}{2.352159\du}}
\pgfpathlineto{\pgfpoint{-16.954976\du}{2.361209\du}}
\pgfpathlineto{\pgfpoint{-16.990401\du}{2.371035\du}}
\pgfpathlineto{\pgfpoint{-17.024015\du}{2.381378\du}}
\pgfpathlineto{\pgfpoint{-17.055820\du}{2.391980\du}}
\pgfpathlineto{\pgfpoint{-17.086332\du}{2.403357\du}}
\pgfpathlineto{\pgfpoint{-17.115034\du}{2.415510\du}}
\pgfpathlineto{\pgfpoint{-17.141925\du}{2.427663\du}}
\pgfpathlineto{\pgfpoint{-17.167266\du}{2.441109\du}}
\pgfpathlineto{\pgfpoint{-17.190020\du}{2.454296\du}}
\pgfpathlineto{\pgfpoint{-17.211223\du}{2.468001\du}}
\pgfpathlineto{\pgfpoint{-17.230099\du}{2.481705\du}}
\pgfpathlineto{\pgfpoint{-17.246648\du}{2.496185\du}}
\pgfpathlineto{\pgfpoint{-17.261387\du}{2.510665\du}}
\pgfpathlineto{\pgfpoint{-17.273798\du}{2.525404\du}}
\pgfpathlineto{\pgfpoint{-17.284141\du}{2.540660\du}}
\pgfpathlineto{\pgfpoint{-17.292157\du}{2.555916\du}}
\pgfpathlineto{\pgfpoint{-17.298104\du}{2.571172\du}}
\pgfpathlineto{\pgfpoint{-17.301207\du}{2.586686\du}}
\pgfpathlineto{\pgfpoint{-17.302500\du}{2.601683\du}}
\pgfpathlineto{\pgfpoint{-17.302500\du}{2.601683\du}}
\pgfpathlineto{\pgfpoint{-17.301207\du}{2.617457\du}}
\pgfpathlineto{\pgfpoint{-17.298104\du}{2.632971\du}}
\pgfpathlineto{\pgfpoint{-17.292157\du}{2.648227\du}}
\pgfpathlineto{\pgfpoint{-17.284141\du}{2.663224\du}}
\pgfpathlineto{\pgfpoint{-17.273798\du}{2.678480\du}}
\pgfpathlineto{\pgfpoint{-17.261387\du}{2.692960\du}}
\pgfpathlineto{\pgfpoint{-17.246648\du}{2.707958\du}}
\pgfpathlineto{\pgfpoint{-17.230099\du}{2.722179\du}}
\pgfpathlineto{\pgfpoint{-17.211223\du}{2.736142\du}}
\pgfpathlineto{\pgfpoint{-17.190020\du}{2.749847\du}}
\pgfpathlineto{\pgfpoint{-17.167266\du}{2.763034\du}}
\pgfpathlineto{\pgfpoint{-17.141925\du}{2.776221\du}}
\pgfpathlineto{\pgfpoint{-17.115034\du}{2.788633\du}}
\pgfpathlineto{\pgfpoint{-17.086332\du}{2.800010\du}}
\pgfpathlineto{\pgfpoint{-17.055820\du}{2.811904\du}}
\pgfpathlineto{\pgfpoint{-17.024015\du}{2.822765\du}}
\pgfpathlineto{\pgfpoint{-16.990401\du}{2.833108\du}}
\pgfpathlineto{\pgfpoint{-16.954976\du}{2.842933\du}}
\pgfpathlineto{\pgfpoint{-16.918517\du}{2.851983\du}}
\pgfpathlineto{\pgfpoint{-16.881024\du}{2.860258\du}}
\pgfpathlineto{\pgfpoint{-16.841720\du}{2.868274\du}}
\pgfpathlineto{\pgfpoint{-16.801383\du}{2.875255\du}}
\pgfpathlineto{\pgfpoint{-16.760270\du}{2.881978\du}}
\pgfpathlineto{\pgfpoint{-16.717863\du}{2.887667\du}}
\pgfpathlineto{\pgfpoint{-16.675199\du}{2.892580\du}}
\pgfpathlineto{\pgfpoint{-16.631241\du}{2.896717\du}}
\pgfpathlineto{\pgfpoint{-16.587025\du}{2.899820\du}}
\pgfpathlineto{\pgfpoint{-16.542550\du}{2.902664\du}}
\pgfpathlineto{\pgfpoint{-16.497558\du}{2.904991\du}}
\pgfpathlineto{\pgfpoint{-16.452566\du}{2.906284\du}}
\pgfpathlineto{\pgfpoint{-16.407315\du}{2.906284\du}}
\pgfpathlineto{\pgfpoint{-16.407315\du}{2.906284\du}}
\pgfpathlineto{\pgfpoint{-16.361548\du}{2.906284\du}}
\pgfpathlineto{\pgfpoint{-16.316297\du}{2.904991\du}}
\pgfpathlineto{\pgfpoint{-16.271305\du}{2.902664\du}}
\pgfpathlineto{\pgfpoint{-16.226831\du}{2.899820\du}}
\pgfpathlineto{\pgfpoint{-16.182614\du}{2.896717\du}}
\pgfpathlineto{\pgfpoint{-16.138657\du}{2.892580\du}}
\pgfpathlineto{\pgfpoint{-16.095992\du}{2.887667\du}}
\pgfpathlineto{\pgfpoint{-16.053586\du}{2.881978\du}}
\pgfpathlineto{\pgfpoint{-16.012472\du}{2.875255\du}}
\pgfpathlineto{\pgfpoint{-15.972135\du}{2.868274\du}}
\pgfpathlineto{\pgfpoint{-15.933090\du}{2.860258\du}}
\pgfpathlineto{\pgfpoint{-15.895338\du}{2.851983\du}}
\pgfpathlineto{\pgfpoint{-15.858879\du}{2.842933\du}}
\pgfpathlineto{\pgfpoint{-15.823454\du}{2.833108\du}}
\pgfpathlineto{\pgfpoint{-15.789840\du}{2.822765\du}}
\pgfpathlineto{\pgfpoint{-15.758035\du}{2.811904\du}}
\pgfpathlineto{\pgfpoint{-15.727523\du}{2.800010\du}}
\pgfpathlineto{\pgfpoint{-15.698822\du}{2.788633\du}}
\pgfpathlineto{\pgfpoint{-15.671930\du}{2.776221\du}}
\pgfpathlineto{\pgfpoint{-15.646848\du}{2.763034\du}}
\pgfpathlineto{\pgfpoint{-15.623835\du}{2.749847\du}}
\pgfpathlineto{\pgfpoint{-15.602891\du}{2.736142\du}}
\pgfpathlineto{\pgfpoint{-15.583756\du}{2.722179\du}}
\pgfpathlineto{\pgfpoint{-15.567207\du}{2.707958\du}}
\pgfpathlineto{\pgfpoint{-15.552469\du}{2.692960\du}}
\pgfpathlineto{\pgfpoint{-15.540057\du}{2.678480\du}}
\pgfpathlineto{\pgfpoint{-15.529714\du}{2.663224\du}}
\pgfpathlineto{\pgfpoint{-15.521698\du}{2.648227\du}}
\pgfpathlineto{\pgfpoint{-15.516010\du}{2.632971\du}}
\pgfpathlineto{\pgfpoint{-15.512648\du}{2.617457\du}}
\pgfpathlineto{\pgfpoint{-15.511355\du}{2.601683\du}}
\pgfusepath{fill}
\pgfsetbuttcap
\pgfsetmiterjoin
\pgfsetdash{}{0pt}
\definecolor{dialinecolor}{rgb}{0.666667, 0.901961, 1.000000}
\pgfsetstrokecolor{dialinecolor}
\pgfpathmoveto{\pgfpoint{-15.521698\du}{2.597288\du}}
\pgfpathlineto{\pgfpoint{-15.522991\du}{2.581773\du}}
\pgfpathlineto{\pgfpoint{-15.526353\du}{2.566517\du}}
\pgfpathlineto{\pgfpoint{-15.532041\du}{2.552037\du}}
\pgfpathlineto{\pgfpoint{-15.539798\du}{2.536781\du}}
\pgfpathlineto{\pgfpoint{-15.550141\du}{2.522043\du}}
\pgfpathlineto{\pgfpoint{-15.562811\du}{2.507562\du}}
\pgfpathlineto{\pgfpoint{-15.576775\du}{2.493599\du}}
\pgfpathlineto{\pgfpoint{-15.593840\du}{2.479119\du}}
\pgfpathlineto{\pgfpoint{-15.612716\du}{2.465673\du}}
\pgfpathlineto{\pgfpoint{-15.633661\du}{2.451710\du}}
\pgfpathlineto{\pgfpoint{-15.656415\du}{2.438782\du}}
\pgfpathlineto{\pgfpoint{-15.681497\du}{2.425853\du}}
\pgfpathlineto{\pgfpoint{-15.707872\du}{2.413700\du}}
\pgfpathlineto{\pgfpoint{-15.736315\du}{2.402323\du}}
\pgfpathlineto{\pgfpoint{-15.766827\du}{2.390428\du}}
\pgfpathlineto{\pgfpoint{-15.798890\du}{2.380085\du}}
\pgfpathlineto{\pgfpoint{-15.832246\du}{2.370001\du}}
\pgfpathlineto{\pgfpoint{-15.867154\du}{2.360434\du}}
\pgfpathlineto{\pgfpoint{-15.903095\du}{2.351642\du}}
\pgfpathlineto{\pgfpoint{-15.941106\du}{2.342851\du}}
\pgfpathlineto{\pgfpoint{-15.980151\du}{2.335352\du}}
\pgfpathlineto{\pgfpoint{-16.019712\du}{2.328370\du}}
\pgfpathlineto{\pgfpoint{-16.060826\du}{2.321906\du}}
\pgfpathlineto{\pgfpoint{-16.102715\du}{2.316217\du}}
\pgfpathlineto{\pgfpoint{-16.145380\du}{2.311304\du}}
\pgfpathlineto{\pgfpoint{-16.189079\du}{2.307684\du}}
\pgfpathlineto{\pgfpoint{-16.232778\du}{2.304064\du}}
\pgfpathlineto{\pgfpoint{-16.276994\du}{2.301479\du}}
\pgfpathlineto{\pgfpoint{-16.321986\du}{2.299151\du}}
\pgfpathlineto{\pgfpoint{-16.366978\du}{2.298117\du}}
\pgfpathlineto{\pgfpoint{-16.412228\du}{2.297600\du}}
\pgfpathlineto{\pgfpoint{-16.412228\du}{2.297600\du}}
\pgfpathlineto{\pgfpoint{-16.457220\du}{2.298117\du}}
\pgfpathlineto{\pgfpoint{-16.502471\du}{2.299151\du}}
\pgfpathlineto{\pgfpoint{-16.547204\du}{2.301479\du}}
\pgfpathlineto{\pgfpoint{-16.591420\du}{2.304064\du}}
\pgfpathlineto{\pgfpoint{-16.635378\du}{2.307684\du}}
\pgfpathlineto{\pgfpoint{-16.678819\du}{2.311304\du}}
\pgfpathlineto{\pgfpoint{-16.721483\du}{2.316217\du}}
\pgfpathlineto{\pgfpoint{-16.763631\du}{2.321906\du}}
\pgfpathlineto{\pgfpoint{-16.804486\du}{2.328370\du}}
\pgfpathlineto{\pgfpoint{-16.844048\du}{2.335352\du}}
\pgfpathlineto{\pgfpoint{-16.883351\du}{2.342851\du}}
\pgfpathlineto{\pgfpoint{-16.920844\du}{2.351642\du}}
\pgfpathlineto{\pgfpoint{-16.956786\du}{2.360434\du}}
\pgfpathlineto{\pgfpoint{-16.991952\du}{2.370001\du}}
\pgfpathlineto{\pgfpoint{-17.025825\du}{2.380085\du}}
\pgfpathlineto{\pgfpoint{-17.057889\du}{2.390428\du}}
\pgfpathlineto{\pgfpoint{-17.087883\du}{2.402323\du}}
\pgfpathlineto{\pgfpoint{-17.116585\du}{2.413700\du}}
\pgfpathlineto{\pgfpoint{-17.143218\du}{2.425853\du}}
\pgfpathlineto{\pgfpoint{-17.167783\du}{2.438782\du}}
\pgfpathlineto{\pgfpoint{-17.190279\du}{2.451710\du}}
\pgfpathlineto{\pgfpoint{-17.211482\du}{2.465673\du}}
\pgfpathlineto{\pgfpoint{-17.230616\du}{2.479119\du}}
\pgfpathlineto{\pgfpoint{-17.247165\du}{2.493599\du}}
\pgfpathlineto{\pgfpoint{-17.261645\du}{2.507562\du}}
\pgfpathlineto{\pgfpoint{-17.274574\du}{2.522043\du}}
\pgfpathlineto{\pgfpoint{-17.284141\du}{2.536781\du}}
\pgfpathlineto{\pgfpoint{-17.292416\du}{2.552037\du}}
\pgfpathlineto{\pgfpoint{-17.298104\du}{2.566517\du}}
\pgfpathlineto{\pgfpoint{-17.301207\du}{2.581773\du}}
\pgfpathlineto{\pgfpoint{-17.302500\du}{2.597288\du}}
\pgfpathlineto{\pgfpoint{-17.302500\du}{2.597288\du}}
\pgfpathlineto{\pgfpoint{-17.301207\du}{2.612285\du}}
\pgfpathlineto{\pgfpoint{-17.298104\du}{2.627541\du}}
\pgfpathlineto{\pgfpoint{-17.292416\du}{2.642021\du}}
\pgfpathlineto{\pgfpoint{-17.284141\du}{2.657018\du}}
\pgfpathlineto{\pgfpoint{-17.274574\du}{2.672016\du}}
\pgfpathlineto{\pgfpoint{-17.261645\du}{2.686496\du}}
\pgfpathlineto{\pgfpoint{-17.247165\du}{2.700976\du}}
\pgfpathlineto{\pgfpoint{-17.230616\du}{2.715198\du}}
\pgfpathlineto{\pgfpoint{-17.211482\du}{2.728902\du}}
\pgfpathlineto{\pgfpoint{-17.190279\du}{2.742348\du}}
\pgfpathlineto{\pgfpoint{-17.167783\du}{2.755535\du}}
\pgfpathlineto{\pgfpoint{-17.143218\du}{2.767688\du}}
\pgfpathlineto{\pgfpoint{-17.116585\du}{2.780358\du}}
\pgfpathlineto{\pgfpoint{-17.087883\du}{2.791994\du}}
\pgfpathlineto{\pgfpoint{-17.057889\du}{2.803372\du}}
\pgfpathlineto{\pgfpoint{-17.025825\du}{2.813973\du}}
\pgfpathlineto{\pgfpoint{-16.991952\du}{2.824057\du}}
\pgfpathlineto{\pgfpoint{-16.956786\du}{2.833625\du}}
\pgfpathlineto{\pgfpoint{-16.920844\du}{2.842416\du}}
\pgfpathlineto{\pgfpoint{-16.883351\du}{2.850949\du}}
\pgfpathlineto{\pgfpoint{-16.844048\du}{2.858965\du}}
\pgfpathlineto{\pgfpoint{-16.804486\du}{2.865429\du}}
\pgfpathlineto{\pgfpoint{-16.763631\du}{2.872152\du}}
\pgfpathlineto{\pgfpoint{-16.721483\du}{2.877582\du}}
\pgfpathlineto{\pgfpoint{-16.678819\du}{2.882495\du}}
\pgfpathlineto{\pgfpoint{-16.635378\du}{2.886891\du}}
\pgfpathlineto{\pgfpoint{-16.591420\du}{2.890511\du}}
\pgfpathlineto{\pgfpoint{-16.547204\du}{2.892580\du}}
\pgfpathlineto{\pgfpoint{-16.502471\du}{2.894648\du}}
\pgfpathlineto{\pgfpoint{-16.457220\du}{2.895941\du}}
\pgfpathlineto{\pgfpoint{-16.412228\du}{2.896200\du}}
\pgfpathlineto{\pgfpoint{-16.412228\du}{2.896200\du}}
\pgfpathlineto{\pgfpoint{-16.366978\du}{2.895941\du}}
\pgfpathlineto{\pgfpoint{-16.321986\du}{2.894648\du}}
\pgfpathlineto{\pgfpoint{-16.276994\du}{2.892580\du}}
\pgfpathlineto{\pgfpoint{-16.232778\du}{2.890511\du}}
\pgfpathlineto{\pgfpoint{-16.189079\du}{2.886891\du}}
\pgfpathlineto{\pgfpoint{-16.145380\du}{2.882495\du}}
\pgfpathlineto{\pgfpoint{-16.102715\du}{2.877582\du}}
\pgfpathlineto{\pgfpoint{-16.060826\du}{2.872152\du}}
\pgfpathlineto{\pgfpoint{-16.019712\du}{2.865429\du}}
\pgfpathlineto{\pgfpoint{-15.980151\du}{2.858965\du}}
\pgfpathlineto{\pgfpoint{-15.941106\du}{2.850949\du}}
\pgfpathlineto{\pgfpoint{-15.903095\du}{2.842416\du}}
\pgfpathlineto{\pgfpoint{-15.867154\du}{2.833625\du}}
\pgfpathlineto{\pgfpoint{-15.832246\du}{2.824057\du}}
\pgfpathlineto{\pgfpoint{-15.798890\du}{2.813973\du}}
\pgfpathlineto{\pgfpoint{-15.766827\du}{2.803372\du}}
\pgfpathlineto{\pgfpoint{-15.736315\du}{2.791994\du}}
\pgfpathlineto{\pgfpoint{-15.707872\du}{2.780358\du}}
\pgfpathlineto{\pgfpoint{-15.681497\du}{2.767688\du}}
\pgfpathlineto{\pgfpoint{-15.656415\du}{2.755535\du}}
\pgfpathlineto{\pgfpoint{-15.633661\du}{2.742348\du}}
\pgfpathlineto{\pgfpoint{-15.612716\du}{2.728902\du}}
\pgfpathlineto{\pgfpoint{-15.593840\du}{2.715198\du}}
\pgfpathlineto{\pgfpoint{-15.576775\du}{2.700976\du}}
\pgfpathlineto{\pgfpoint{-15.562811\du}{2.686496\du}}
\pgfpathlineto{\pgfpoint{-15.550141\du}{2.672016\du}}
\pgfpathlineto{\pgfpoint{-15.539798\du}{2.657018\du}}
\pgfpathlineto{\pgfpoint{-15.532041\du}{2.642021\du}}
\pgfpathlineto{\pgfpoint{-15.526353\du}{2.627541\du}}
\pgfpathlineto{\pgfpoint{-15.522991\du}{2.612285\du}}
\pgfpathlineto{\pgfpoint{-15.521698\du}{2.597288\du}}
\pgfusepath{stroke}
\pgfsetbuttcap
\pgfsetmiterjoin
\pgfsetdash{}{0pt}
\definecolor{dialinecolor}{rgb}{0.000000, 0.000000, 0.000000}
\pgfsetfillcolor{dialinecolor}
\pgfpathmoveto{\pgfpoint{-16.390767\du}{2.532644\du}}
\pgfpathlineto{\pgfpoint{-16.260187\du}{2.575567\du}}
\pgfpathlineto{\pgfpoint{-15.944984\du}{2.447056\du}}
\pgfpathlineto{\pgfpoint{-15.804061\du}{2.489979\du}}
\pgfpathlineto{\pgfpoint{-15.880599\du}{2.383188\du}}
\pgfpathlineto{\pgfpoint{-16.249068\du}{2.383188\du}}
\pgfpathlineto{\pgfpoint{-16.097543\du}{2.415510\du}}
\pgfpathlineto{\pgfpoint{-16.390767\du}{2.532644\du}}
\pgfusepath{fill}
\pgfsetbuttcap
\pgfsetmiterjoin
\pgfsetdash{}{0pt}
\definecolor{dialinecolor}{rgb}{0.000000, 0.000000, 0.000000}
\pgfsetfillcolor{dialinecolor}
\pgfpathmoveto{\pgfpoint{-16.445326\du}{2.650554\du}}
\pgfpathlineto{\pgfpoint{-16.574872\du}{2.607631\du}}
\pgfpathlineto{\pgfpoint{-16.879214\du}{2.736142\du}}
\pgfpathlineto{\pgfpoint{-17.030997\du}{2.692960\du}}
\pgfpathlineto{\pgfpoint{-16.954976\du}{2.810612\du}}
\pgfpathlineto{\pgfpoint{-16.574872\du}{2.810612\du}}
\pgfpathlineto{\pgfpoint{-16.738032\du}{2.767688\du}}
\pgfpathlineto{\pgfpoint{-16.445326\du}{2.650554\du}}
\pgfusepath{fill}
\pgfsetbuttcap
\pgfsetmiterjoin
\pgfsetdash{}{0pt}
\definecolor{dialinecolor}{rgb}{0.000000, 0.000000, 0.000000}
\pgfsetfillcolor{dialinecolor}
\pgfpathmoveto{\pgfpoint{-16.998675\du}{2.415510\du}}
\pgfpathlineto{\pgfpoint{-16.868354\du}{2.372587\du}}
\pgfpathlineto{\pgfpoint{-16.553669\du}{2.489979\du}}
\pgfpathlineto{\pgfpoint{-16.412228\du}{2.458175\du}}
\pgfpathlineto{\pgfpoint{-16.487991\du}{2.564707\du}}
\pgfpathlineto{\pgfpoint{-16.857493\du}{2.564707\du}}
\pgfpathlineto{\pgfpoint{-16.705710\du}{2.532644\du}}
\pgfpathlineto{\pgfpoint{-16.998675\du}{2.415510\du}}
\pgfusepath{fill}
\pgfsetbuttcap
\pgfsetmiterjoin
\pgfsetdash{}{0pt}
\definecolor{dialinecolor}{rgb}{0.000000, 0.000000, 0.000000}
\pgfsetfillcolor{dialinecolor}
\pgfpathmoveto{\pgfpoint{-15.826040\du}{2.778548\du}}
\pgfpathlineto{\pgfpoint{-15.956103\du}{2.821213\du}}
\pgfpathlineto{\pgfpoint{-16.260187\du}{2.692960\du}}
\pgfpathlineto{\pgfpoint{-16.412228\du}{2.736142\du}}
\pgfpathlineto{\pgfpoint{-16.336207\du}{2.629092\du}}
\pgfpathlineto{\pgfpoint{-15.956103\du}{2.629092\du}}
\pgfpathlineto{\pgfpoint{-16.119264\du}{2.660638\du}}
\pgfpathlineto{\pgfpoint{-15.826040\du}{2.778548\du}}
\pgfusepath{fill}
\pgfsetbuttcap
\pgfsetmiterjoin
\pgfsetdash{}{0pt}
\definecolor{dialinecolor}{rgb}{1.000000, 1.000000, 1.000000}
\pgfsetfillcolor{dialinecolor}
\pgfpathmoveto{\pgfpoint{-16.379907\du}{2.543246\du}}
\pgfpathlineto{\pgfpoint{-16.249068\du}{2.585910\du}}
\pgfpathlineto{\pgfpoint{-15.934383\du}{2.458175\du}}
\pgfpathlineto{\pgfpoint{-15.793201\du}{2.501098\du}}
\pgfpathlineto{\pgfpoint{-15.869222\du}{2.393790\du}}
\pgfpathlineto{\pgfpoint{-16.238725\du}{2.393790\du}}
\pgfpathlineto{\pgfpoint{-16.086425\du}{2.425853\du}}
\pgfpathlineto{\pgfpoint{-16.379907\du}{2.543246\du}}
\pgfusepath{fill}
\pgfsetbuttcap
\pgfsetmiterjoin
\pgfsetdash{}{0pt}
\definecolor{dialinecolor}{rgb}{1.000000, 1.000000, 1.000000}
\pgfsetfillcolor{dialinecolor}
\pgfpathmoveto{\pgfpoint{-16.433690\du}{2.660638\du}}
\pgfpathlineto{\pgfpoint{-16.564529\du}{2.618232\du}}
\pgfpathlineto{\pgfpoint{-16.868354\du}{2.746485\du}}
\pgfpathlineto{\pgfpoint{-17.020395\du}{2.703562\du}}
\pgfpathlineto{\pgfpoint{-16.944374\du}{2.821213\du}}
\pgfpathlineto{\pgfpoint{-16.564529\du}{2.821213\du}}
\pgfpathlineto{\pgfpoint{-16.726913\du}{2.778548\du}}
\pgfpathlineto{\pgfpoint{-16.433690\du}{2.660638\du}}
\pgfusepath{fill}
\pgfsetbuttcap
\pgfsetmiterjoin
\pgfsetdash{}{0pt}
\definecolor{dialinecolor}{rgb}{1.000000, 1.000000, 1.000000}
\pgfsetfillcolor{dialinecolor}
\pgfpathmoveto{\pgfpoint{-16.987556\du}{2.425853\du}}
\pgfpathlineto{\pgfpoint{-16.857493\du}{2.383188\du}}
\pgfpathlineto{\pgfpoint{-16.542550\du}{2.501098\du}}
\pgfpathlineto{\pgfpoint{-16.401368\du}{2.468518\du}}
\pgfpathlineto{\pgfpoint{-16.477389\du}{2.575567\du}}
\pgfpathlineto{\pgfpoint{-16.846633\du}{2.575567\du}}
\pgfpathlineto{\pgfpoint{-16.694333\du}{2.543246\du}}
\pgfpathlineto{\pgfpoint{-16.987556\du}{2.425853\du}}
\pgfusepath{fill}
\pgfsetbuttcap
\pgfsetmiterjoin
\pgfsetdash{}{0pt}
\definecolor{dialinecolor}{rgb}{1.000000, 1.000000, 1.000000}
\pgfsetfillcolor{dialinecolor}
\pgfpathmoveto{\pgfpoint{-15.814922\du}{2.789150\du}}
\pgfpathlineto{\pgfpoint{-15.944984\du}{2.832073\du}}
\pgfpathlineto{\pgfpoint{-16.249068\du}{2.703562\du}}
\pgfpathlineto{\pgfpoint{-16.401368\du}{2.746485\du}}
\pgfpathlineto{\pgfpoint{-16.325606\du}{2.639952\du}}
\pgfpathlineto{\pgfpoint{-15.944984\du}{2.639952\du}}
\pgfpathlineto{\pgfpoint{-16.107886\du}{2.671757\du}}
\pgfpathlineto{\pgfpoint{-15.814922\du}{2.789150\du}}
\pgfusepath{fill}
\pgfsetbuttcap
\pgfsetmiterjoin
\pgfsetdash{}{0pt}
\definecolor{dialinecolor}{rgb}{0.666667, 0.901961, 1.000000}
\pgfsetstrokecolor{dialinecolor}
\pgfpathmoveto{\pgfpoint{-17.302500\du}{2.597288\du}}
\pgfpathlineto{\pgfpoint{-17.302500\du}{3.023935\du}}
\pgfusepath{stroke}
\pgfsetbuttcap
\pgfsetmiterjoin
\pgfsetdash{}{0pt}
\definecolor{dialinecolor}{rgb}{0.666667, 0.901961, 1.000000}
\pgfsetstrokecolor{dialinecolor}
\pgfpathmoveto{\pgfpoint{-15.521698\du}{2.597288\du}}
\pgfpathlineto{\pgfpoint{-15.521698\du}{3.023935\du}}
\pgfusepath{stroke}
\pgfsetbuttcap
\pgfsetmiterjoin
\pgfsetdash{}{0pt}
\definecolor{dialinecolor}{rgb}{0.000000, 0.000000, 0.000000}
\pgfsetfillcolor{dialinecolor}
\pgfpathmoveto{\pgfpoint{-16.738032\du}{2.970928\du}}
\pgfpathlineto{\pgfpoint{-16.542550\du}{2.970928\du}}
\pgfpathlineto{\pgfpoint{-16.423088\du}{3.067117\du}}
\pgfpathlineto{\pgfpoint{-16.303886\du}{2.970928\du}}
\pgfpathlineto{\pgfpoint{-16.119264\du}{2.970928\du}}
\pgfpathlineto{\pgfpoint{-16.119264\du}{2.928263\du}}
\pgfpathlineto{\pgfpoint{-15.988684\du}{2.991872\du}}
\pgfpathlineto{\pgfpoint{-16.119264\du}{3.056257\du}}
\pgfpathlineto{\pgfpoint{-16.119264\du}{3.013592\du}}
\pgfpathlineto{\pgfpoint{-16.281907\du}{3.013592\du}}
\pgfpathlineto{\pgfpoint{-16.390767\du}{3.088320\du}}
\pgfpathlineto{\pgfpoint{-16.281907\du}{3.173391\du}}
\pgfpathlineto{\pgfpoint{-16.119264\du}{3.173391\du}}
\pgfpathlineto{\pgfpoint{-16.119264\du}{3.120642\du}}
\pgfpathlineto{\pgfpoint{-15.988684\du}{3.195112\du}}
\pgfpathlineto{\pgfpoint{-16.119264\du}{3.259497\du}}
\pgfpathlineto{\pgfpoint{-16.119264\du}{3.216315\du}}
\pgfpathlineto{\pgfpoint{-16.303886\du}{3.216315\du}}
\pgfpathlineto{\pgfpoint{-16.423088\du}{3.120642\du}}
\pgfpathlineto{\pgfpoint{-16.542550\du}{3.216315\du}}
\pgfpathlineto{\pgfpoint{-16.738032\du}{3.216315\du}}
\pgfpathlineto{\pgfpoint{-16.738032\du}{3.259497\du}}
\pgfpathlineto{\pgfpoint{-16.868354\du}{3.195112\du}}
\pgfpathlineto{\pgfpoint{-16.738032\du}{3.120642\du}}
\pgfpathlineto{\pgfpoint{-16.738032\du}{3.173391\du}}
\pgfpathlineto{\pgfpoint{-16.564529\du}{3.173391\du}}
\pgfpathlineto{\pgfpoint{-16.455669\du}{3.098922\du}}
\pgfpathlineto{\pgfpoint{-16.564529\du}{3.013592\du}}
\pgfpathlineto{\pgfpoint{-16.738032\du}{3.013592\du}}
\pgfpathlineto{\pgfpoint{-16.738032\du}{3.056257\du}}
\pgfpathlineto{\pgfpoint{-16.868354\du}{2.991872\du}}
\pgfpathlineto{\pgfpoint{-16.738032\du}{2.928263\du}}
\pgfpathlineto{\pgfpoint{-16.738032\du}{2.970928\du}}
\pgfusepath{fill}
\pgfsetbuttcap
\pgfsetmiterjoin
\pgfsetdash{}{0pt}
\definecolor{dialinecolor}{rgb}{1.000000, 1.000000, 1.000000}
\pgfsetfillcolor{dialinecolor}
\pgfpathmoveto{\pgfpoint{-16.726913\du}{2.981271\du}}
\pgfpathlineto{\pgfpoint{-16.531690\du}{2.981271\du}}
\pgfpathlineto{\pgfpoint{-16.412228\du}{3.077202\du}}
\pgfpathlineto{\pgfpoint{-16.293026\du}{2.981271\du}}
\pgfpathlineto{\pgfpoint{-16.107886\du}{2.981271\du}}
\pgfpathlineto{\pgfpoint{-16.107886\du}{2.938606\du}}
\pgfpathlineto{\pgfpoint{-15.978082\du}{3.002991\du}}
\pgfpathlineto{\pgfpoint{-16.107886\du}{3.067117\du}}
\pgfpathlineto{\pgfpoint{-16.107886\du}{3.024453\du}}
\pgfpathlineto{\pgfpoint{-16.271047\du}{3.024453\du}}
\pgfpathlineto{\pgfpoint{-16.379907\du}{3.098922\du}}
\pgfpathlineto{\pgfpoint{-16.271047\du}{3.184252\du}}
\pgfpathlineto{\pgfpoint{-16.107886\du}{3.184252\du}}
\pgfpathlineto{\pgfpoint{-16.107886\du}{3.130985\du}}
\pgfpathlineto{\pgfpoint{-15.978082\du}{3.205713\du}}
\pgfpathlineto{\pgfpoint{-16.107886\du}{3.269840\du}}
\pgfpathlineto{\pgfpoint{-16.107886\du}{3.226916\du}}
\pgfpathlineto{\pgfpoint{-16.293026\du}{3.226916\du}}
\pgfpathlineto{\pgfpoint{-16.412228\du}{3.130985\du}}
\pgfpathlineto{\pgfpoint{-16.531690\du}{3.226916\du}}
\pgfpathlineto{\pgfpoint{-16.726913\du}{3.226916\du}}
\pgfpathlineto{\pgfpoint{-16.726913\du}{3.269840\du}}
\pgfpathlineto{\pgfpoint{-16.857493\du}{3.205713\du}}
\pgfpathlineto{\pgfpoint{-16.726913\du}{3.130985\du}}
\pgfpathlineto{\pgfpoint{-16.726913\du}{3.184252\du}}
\pgfpathlineto{\pgfpoint{-16.553669\du}{3.184252\du}}
\pgfpathlineto{\pgfpoint{-16.445326\du}{3.109524\du}}
\pgfpathlineto{\pgfpoint{-16.553669\du}{3.024453\du}}
\pgfpathlineto{\pgfpoint{-16.726913\du}{3.024453\du}}
\pgfpathlineto{\pgfpoint{-16.726913\du}{3.067117\du}}
\pgfpathlineto{\pgfpoint{-16.857493\du}{3.002991\du}}
\pgfpathlineto{\pgfpoint{-16.726913\du}{2.938606\du}}
\pgfpathlineto{\pgfpoint{-16.726913\du}{2.981271\du}}
\pgfusepath{fill}
\pgfsetlinewidth{0.100000\du}
\pgfsetdash{}{0pt}
\pgfsetdash{}{0pt}
\pgfsetbuttcap
{
\definecolor{dialinecolor}{rgb}{0.000000, 0.000000, 0.000000}
\pgfsetfillcolor{dialinecolor}
% was here!!!
\definecolor{dialinecolor}{rgb}{0.000000, 0.000000, 0.000000}
\pgfsetstrokecolor{dialinecolor}
\draw (-18.654652\du,0.661501\du)--(-16.897776\du,2.345437\du);
}
\pgfsetlinewidth{0.100000\du}
\pgfsetdash{}{0pt}
\pgfsetdash{}{0pt}
\pgfsetbuttcap
\pgfsetmiterjoin
\pgfsetlinewidth{0.001000\du}
\pgfsetbuttcap
\pgfsetmiterjoin
\pgfsetdash{}{0pt}
\definecolor{dialinecolor}{rgb}{0.000000, 0.470588, 0.666667}
\pgfsetfillcolor{dialinecolor}
\pgfpathmoveto{\pgfpoint{-0.529166\du}{2.945170\du}}
\pgfpathlineto{\pgfpoint{-0.530664\du}{2.927494\du}}
\pgfpathlineto{\pgfpoint{-0.534559\du}{2.909518\du}}
\pgfpathlineto{\pgfpoint{-0.541150\du}{2.891841\du}}
\pgfpathlineto{\pgfpoint{-0.550438\du}{2.874465\du}}
\pgfpathlineto{\pgfpoint{-0.562422\du}{2.856788\du}}
\pgfpathlineto{\pgfpoint{-0.576803\du}{2.839411\du}}
\pgfpathlineto{\pgfpoint{-0.593880\du}{2.822634\du}}
\pgfpathlineto{\pgfpoint{-0.613054\du}{2.806156\du}}
\pgfpathlineto{\pgfpoint{-0.635225\du}{2.789977\du}}
\pgfpathlineto{\pgfpoint{-0.659492\du}{2.774398\du}}
\pgfpathlineto{\pgfpoint{-0.686157\du}{2.758819\du}}
\pgfpathlineto{\pgfpoint{-0.715218\du}{2.743839\du}}
\pgfpathlineto{\pgfpoint{-0.746377\du}{2.729458\du}}
\pgfpathlineto{\pgfpoint{-0.779632\du}{2.715377\du}}
\pgfpathlineto{\pgfpoint{-0.814985\du}{2.702194\du}}
\pgfpathlineto{\pgfpoint{-0.851836\du}{2.689611\du}}
\pgfpathlineto{\pgfpoint{-0.890784\du}{2.677627\du}}
\pgfpathlineto{\pgfpoint{-0.931829\du}{2.666542\du}}
\pgfpathlineto{\pgfpoint{-0.974073\du}{2.655456\du}}
\pgfpathlineto{\pgfpoint{-1.017815\du}{2.646169\du}}
\pgfpathlineto{\pgfpoint{-1.063055\du}{2.636881\du}}
\pgfpathlineto{\pgfpoint{-1.109792\du}{2.628792\du}}
\pgfpathlineto{\pgfpoint{-1.157429\du}{2.621002\du}}
\pgfpathlineto{\pgfpoint{-1.206563\du}{2.614411\du}}
\pgfpathlineto{\pgfpoint{-1.255997\du}{2.608719\du}}
\pgfpathlineto{\pgfpoint{-1.306930\du}{2.603925\du}}
\pgfpathlineto{\pgfpoint{-1.358161\du}{2.599731\du}}
\pgfpathlineto{\pgfpoint{-1.409693\du}{2.596435\du}}
\pgfpathlineto{\pgfpoint{-1.461823\du}{2.594038\du}}
\pgfpathlineto{\pgfpoint{-1.514253\du}{2.593139\du}}
\pgfpathlineto{\pgfpoint{-1.567283\du}{2.592241\du}}
\pgfpathlineto{\pgfpoint{-1.567283\du}{2.592241\du}}
\pgfpathlineto{\pgfpoint{-1.619713\du}{2.593139\du}}
\pgfpathlineto{\pgfpoint{-1.671843\du}{2.594038\du}}
\pgfpathlineto{\pgfpoint{-1.723974\du}{2.596435\du}}
\pgfpathlineto{\pgfpoint{-1.775505\du}{2.599731\du}}
\pgfpathlineto{\pgfpoint{-1.826737\du}{2.603925\du}}
\pgfpathlineto{\pgfpoint{-1.877669\du}{2.608719\du}}
\pgfpathlineto{\pgfpoint{-1.927103\du}{2.614411\du}}
\pgfpathlineto{\pgfpoint{-1.976237\du}{2.621002\du}}
\pgfpathlineto{\pgfpoint{-2.023874\du}{2.628792\du}}
\pgfpathlineto{\pgfpoint{-2.070612\du}{2.636881\du}}
\pgfpathlineto{\pgfpoint{-2.116151\du}{2.646169\du}}
\pgfpathlineto{\pgfpoint{-2.159593\du}{2.655456\du}}
\pgfpathlineto{\pgfpoint{-2.201837\du}{2.666542\du}}
\pgfpathlineto{\pgfpoint{-2.242882\du}{2.677627\du}}
\pgfpathlineto{\pgfpoint{-2.281830\du}{2.689611\du}}
\pgfpathlineto{\pgfpoint{-2.318681\du}{2.702194\du}}
\pgfpathlineto{\pgfpoint{-2.354034\du}{2.715377\du}}
\pgfpathlineto{\pgfpoint{-2.387290\du}{2.729458\du}}
\pgfpathlineto{\pgfpoint{-2.418448\du}{2.743839\du}}
\pgfpathlineto{\pgfpoint{-2.447809\du}{2.758819\du}}
\pgfpathlineto{\pgfpoint{-2.474174\du}{2.774398\du}}
\pgfpathlineto{\pgfpoint{-2.498741\du}{2.789977\du}}
\pgfpathlineto{\pgfpoint{-2.520612\du}{2.806156\du}}
\pgfpathlineto{\pgfpoint{-2.539786\du}{2.822634\du}}
\pgfpathlineto{\pgfpoint{-2.556863\du}{2.839411\du}}
\pgfpathlineto{\pgfpoint{-2.571244\du}{2.856788\du}}
\pgfpathlineto{\pgfpoint{-2.583228\du}{2.874465\du}}
\pgfpathlineto{\pgfpoint{-2.592516\du}{2.891841\du}}
\pgfpathlineto{\pgfpoint{-2.599407\du}{2.909518\du}}
\pgfpathlineto{\pgfpoint{-2.603002\du}{2.927494\du}}
\pgfpathlineto{\pgfpoint{-2.604500\du}{2.945170\du}}
\pgfpathlineto{\pgfpoint{-2.604500\du}{2.945170\du}}
\pgfpathlineto{\pgfpoint{-2.603002\du}{2.963446\du}}
\pgfpathlineto{\pgfpoint{-2.599407\du}{2.981122\du}}
\pgfpathlineto{\pgfpoint{-2.592516\du}{2.998799\du}}
\pgfpathlineto{\pgfpoint{-2.583228\du}{3.016475\du}}
\pgfpathlineto{\pgfpoint{-2.571244\du}{3.033852\du}}
\pgfpathlineto{\pgfpoint{-2.556863\du}{3.051229\du}}
\pgfpathlineto{\pgfpoint{-2.539786\du}{3.068007\du}}
\pgfpathlineto{\pgfpoint{-2.520612\du}{3.084485\du}}
\pgfpathlineto{\pgfpoint{-2.498741\du}{3.100663\du}}
\pgfpathlineto{\pgfpoint{-2.474174\du}{3.116542\du}}
\pgfpathlineto{\pgfpoint{-2.447809\du}{3.131522\du}}
\pgfpathlineto{\pgfpoint{-2.418448\du}{3.147101\du}}
\pgfpathlineto{\pgfpoint{-2.387290\du}{3.161482\du}}
\pgfpathlineto{\pgfpoint{-2.354034\du}{3.175264\du}}
\pgfpathlineto{\pgfpoint{-2.318681\du}{3.188146\du}}
\pgfpathlineto{\pgfpoint{-2.281830\du}{3.201029\du}}
\pgfpathlineto{\pgfpoint{-2.242882\du}{3.213013\du}}
\pgfpathlineto{\pgfpoint{-2.201837\du}{3.224398\du}}
\pgfpathlineto{\pgfpoint{-2.159593\du}{3.234884\du}}
\pgfpathlineto{\pgfpoint{-2.116151\du}{3.244771\du}}
\pgfpathlineto{\pgfpoint{-2.070612\du}{3.253759\du}}
\pgfpathlineto{\pgfpoint{-2.023874\du}{3.261848\du}}
\pgfpathlineto{\pgfpoint{-1.976237\du}{3.269638\du}}
\pgfpathlineto{\pgfpoint{-1.927103\du}{3.276229\du}}
\pgfpathlineto{\pgfpoint{-1.877669\du}{3.281921\du}}
\pgfpathlineto{\pgfpoint{-1.826737\du}{3.287015\du}}
\pgfpathlineto{\pgfpoint{-1.775505\du}{3.290610\du}}
\pgfpathlineto{\pgfpoint{-1.723974\du}{3.293905\du}}
\pgfpathlineto{\pgfpoint{-1.671843\du}{3.296302\du}}
\pgfpathlineto{\pgfpoint{-1.619713\du}{3.297800\du}}
\pgfpathlineto{\pgfpoint{-1.567283\du}{3.298100\du}}
\pgfpathlineto{\pgfpoint{-1.567283\du}{3.298100\du}}
\pgfpathlineto{\pgfpoint{-1.514253\du}{3.297800\du}}
\pgfpathlineto{\pgfpoint{-1.461823\du}{3.296302\du}}
\pgfpathlineto{\pgfpoint{-1.409693\du}{3.293905\du}}
\pgfpathlineto{\pgfpoint{-1.358161\du}{3.290610\du}}
\pgfpathlineto{\pgfpoint{-1.306930\du}{3.287015\du}}
\pgfpathlineto{\pgfpoint{-1.255997\du}{3.281921\du}}
\pgfpathlineto{\pgfpoint{-1.206563\du}{3.276229\du}}
\pgfpathlineto{\pgfpoint{-1.157429\du}{3.269638\du}}
\pgfpathlineto{\pgfpoint{-1.109792\du}{3.261848\du}}
\pgfpathlineto{\pgfpoint{-1.063055\du}{3.253759\du}}
\pgfpathlineto{\pgfpoint{-1.017815\du}{3.244771\du}}
\pgfpathlineto{\pgfpoint{-0.974073\du}{3.234884\du}}
\pgfpathlineto{\pgfpoint{-0.931829\du}{3.224398\du}}
\pgfpathlineto{\pgfpoint{-0.890784\du}{3.213013\du}}
\pgfpathlineto{\pgfpoint{-0.851836\du}{3.201029\du}}
\pgfpathlineto{\pgfpoint{-0.814985\du}{3.188146\du}}
\pgfpathlineto{\pgfpoint{-0.779632\du}{3.175264\du}}
\pgfpathlineto{\pgfpoint{-0.746377\du}{3.161482\du}}
\pgfpathlineto{\pgfpoint{-0.715218\du}{3.147101\du}}
\pgfpathlineto{\pgfpoint{-0.686157\du}{3.131522\du}}
\pgfpathlineto{\pgfpoint{-0.659492\du}{3.116542\du}}
\pgfpathlineto{\pgfpoint{-0.635225\du}{3.100663\du}}
\pgfpathlineto{\pgfpoint{-0.613054\du}{3.084485\du}}
\pgfpathlineto{\pgfpoint{-0.593880\du}{3.068007\du}}
\pgfpathlineto{\pgfpoint{-0.576803\du}{3.051229\du}}
\pgfpathlineto{\pgfpoint{-0.562422\du}{3.033852\du}}
\pgfpathlineto{\pgfpoint{-0.550438\du}{3.016475\du}}
\pgfpathlineto{\pgfpoint{-0.541150\du}{2.998799\du}}
\pgfpathlineto{\pgfpoint{-0.534559\du}{2.981122\du}}
\pgfpathlineto{\pgfpoint{-0.530664\du}{2.963446\du}}
\pgfpathlineto{\pgfpoint{-0.529166\du}{2.945170\du}}
\pgfusepath{fill}
\pgfsetbuttcap
\pgfsetmiterjoin
\pgfsetdash{}{0pt}
\definecolor{dialinecolor}{rgb}{0.666667, 0.901961, 1.000000}
\pgfsetstrokecolor{dialinecolor}
\pgfpathmoveto{\pgfpoint{-0.541150\du}{2.939777\du}}
\pgfpathlineto{\pgfpoint{-0.542648\du}{2.921801\du}}
\pgfpathlineto{\pgfpoint{-0.546543\du}{2.904724\du}}
\pgfpathlineto{\pgfpoint{-0.553134\du}{2.887048\du}}
\pgfpathlineto{\pgfpoint{-0.562122\du}{2.869971\du}}
\pgfpathlineto{\pgfpoint{-0.574106\du}{2.852594\du}}
\pgfpathlineto{\pgfpoint{-0.588787\du}{2.835816\du}}
\pgfpathlineto{\pgfpoint{-0.604965\du}{2.819338\du}}
\pgfpathlineto{\pgfpoint{-0.624739\du}{2.802860\du}}
\pgfpathlineto{\pgfpoint{-0.646610\du}{2.786981\du}}
\pgfpathlineto{\pgfpoint{-0.670877\du}{2.771102\du}}
\pgfpathlineto{\pgfpoint{-0.697242\du}{2.756422\du}}
\pgfpathlineto{\pgfpoint{-0.726303\du}{2.741741\du}}
\pgfpathlineto{\pgfpoint{-0.756863\du}{2.727061\du}}
\pgfpathlineto{\pgfpoint{-0.789819\du}{2.713579\du}}
\pgfpathlineto{\pgfpoint{-0.825172\du}{2.700397\du}}
\pgfpathlineto{\pgfpoint{-0.862322\du}{2.688413\du}}
\pgfpathlineto{\pgfpoint{-0.900971\du}{2.676428\du}}
\pgfpathlineto{\pgfpoint{-0.941417\du}{2.665343\du}}
\pgfpathlineto{\pgfpoint{-0.983061\du}{2.655157\du}}
\pgfpathlineto{\pgfpoint{-1.027103\du}{2.645270\du}}
\pgfpathlineto{\pgfpoint{-1.072342\du}{2.636582\du}}
\pgfpathlineto{\pgfpoint{-1.118181\du}{2.628492\du}}
\pgfpathlineto{\pgfpoint{-1.165818\du}{2.621002\du}}
\pgfpathlineto{\pgfpoint{-1.214353\du}{2.614411\du}}
\pgfpathlineto{\pgfpoint{-1.263787\du}{2.608719\du}}
\pgfpathlineto{\pgfpoint{-1.314420\du}{2.603925\du}}
\pgfpathlineto{\pgfpoint{-1.365052\du}{2.599731\du}}
\pgfpathlineto{\pgfpoint{-1.416284\du}{2.597334\du}}
\pgfpathlineto{\pgfpoint{-1.468414\du}{2.594637\du}}
\pgfpathlineto{\pgfpoint{-1.520545\du}{2.593139\du}}
\pgfpathlineto{\pgfpoint{-1.572975\du}{2.593139\du}}
\pgfpathlineto{\pgfpoint{-1.572975\du}{2.593139\du}}
\pgfpathlineto{\pgfpoint{-1.625105\du}{2.593139\du}}
\pgfpathlineto{\pgfpoint{-1.677536\du}{2.594637\du}}
\pgfpathlineto{\pgfpoint{-1.729366\du}{2.597334\du}}
\pgfpathlineto{\pgfpoint{-1.780598\du}{2.599731\du}}
\pgfpathlineto{\pgfpoint{-1.831530\du}{2.603925\du}}
\pgfpathlineto{\pgfpoint{-1.881863\du}{2.608719\du}}
\pgfpathlineto{\pgfpoint{-1.931297\du}{2.614411\du}}
\pgfpathlineto{\pgfpoint{-1.980132\du}{2.621002\du}}
\pgfpathlineto{\pgfpoint{-2.027469\du}{2.628492\du}}
\pgfpathlineto{\pgfpoint{-2.073308\du}{2.636582\du}}
\pgfpathlineto{\pgfpoint{-2.118847\du}{2.645270\du}}
\pgfpathlineto{\pgfpoint{-2.162289\du}{2.655157\du}}
\pgfpathlineto{\pgfpoint{-2.203934\du}{2.665343\du}}
\pgfpathlineto{\pgfpoint{-2.244680\du}{2.676428\du}}
\pgfpathlineto{\pgfpoint{-2.283927\du}{2.688413\du}}
\pgfpathlineto{\pgfpoint{-2.321078\du}{2.700397\du}}
\pgfpathlineto{\pgfpoint{-2.355831\du}{2.713579\du}}
\pgfpathlineto{\pgfpoint{-2.389087\du}{2.727061\du}}
\pgfpathlineto{\pgfpoint{-2.419946\du}{2.741741\du}}
\pgfpathlineto{\pgfpoint{-2.448408\du}{2.756422\du}}
\pgfpathlineto{\pgfpoint{-2.474473\du}{2.771102\du}}
\pgfpathlineto{\pgfpoint{-2.499041\du}{2.786981\du}}
\pgfpathlineto{\pgfpoint{-2.521211\du}{2.802860\du}}
\pgfpathlineto{\pgfpoint{-2.540385\du}{2.819338\du}}
\pgfpathlineto{\pgfpoint{-2.557163\du}{2.835816\du}}
\pgfpathlineto{\pgfpoint{-2.572143\du}{2.852594\du}}
\pgfpathlineto{\pgfpoint{-2.583228\du}{2.869971\du}}
\pgfpathlineto{\pgfpoint{-2.592816\du}{2.887048\du}}
\pgfpathlineto{\pgfpoint{-2.599407\du}{2.904724\du}}
\pgfpathlineto{\pgfpoint{-2.603002\du}{2.921801\du}}
\pgfpathlineto{\pgfpoint{-2.604500\du}{2.939777\du}}
\pgfpathlineto{\pgfpoint{-2.604500\du}{2.939777\du}}
\pgfpathlineto{\pgfpoint{-2.603002\du}{2.957454\du}}
\pgfpathlineto{\pgfpoint{-2.599407\du}{2.974831\du}}
\pgfpathlineto{\pgfpoint{-2.592816\du}{2.992208\du}}
\pgfpathlineto{\pgfpoint{-2.583228\du}{3.009285\du}}
\pgfpathlineto{\pgfpoint{-2.572143\du}{3.026662\du}}
\pgfpathlineto{\pgfpoint{-2.557163\du}{3.043439\du}}
\pgfpathlineto{\pgfpoint{-2.540385\du}{3.059917\du}}
\pgfpathlineto{\pgfpoint{-2.521211\du}{3.075796\du}}
\pgfpathlineto{\pgfpoint{-2.499041\du}{3.092274\du}}
\pgfpathlineto{\pgfpoint{-2.474473\du}{3.108153\du}}
\pgfpathlineto{\pgfpoint{-2.448408\du}{3.123133\du}}
\pgfpathlineto{\pgfpoint{-2.419946\du}{3.137514\du}}
\pgfpathlineto{\pgfpoint{-2.389087\du}{3.151895\du}}
\pgfpathlineto{\pgfpoint{-2.355831\du}{3.165676\du}}
\pgfpathlineto{\pgfpoint{-2.321078\du}{3.178559\du}}
\pgfpathlineto{\pgfpoint{-2.283927\du}{3.190843\du}}
\pgfpathlineto{\pgfpoint{-2.244680\du}{3.202228\du}}
\pgfpathlineto{\pgfpoint{-2.203934\du}{3.213912\du}}
\pgfpathlineto{\pgfpoint{-2.162289\du}{3.224098\du}}
\pgfpathlineto{\pgfpoint{-2.118847\du}{3.233686\du}}
\pgfpathlineto{\pgfpoint{-2.073308\du}{3.242973\du}}
\pgfpathlineto{\pgfpoint{-2.027469\du}{3.250763\du}}
\pgfpathlineto{\pgfpoint{-1.980132\du}{3.257953\du}}
\pgfpathlineto{\pgfpoint{-1.931297\du}{3.264844\du}}
\pgfpathlineto{\pgfpoint{-1.881863\du}{3.270237\du}}
\pgfpathlineto{\pgfpoint{-1.831530\du}{3.274731\du}}
\pgfpathlineto{\pgfpoint{-1.780598\du}{3.279225\du}}
\pgfpathlineto{\pgfpoint{-1.729366\du}{3.281921\du}}
\pgfpathlineto{\pgfpoint{-1.677536\du}{3.284318\du}}
\pgfpathlineto{\pgfpoint{-1.625105\du}{3.285517\du}}
\pgfpathlineto{\pgfpoint{-1.572975\du}{3.286116\du}}
\pgfpathlineto{\pgfpoint{-1.572975\du}{3.286116\du}}
\pgfpathlineto{\pgfpoint{-1.520545\du}{3.285517\du}}
\pgfpathlineto{\pgfpoint{-1.468414\du}{3.284318\du}}
\pgfpathlineto{\pgfpoint{-1.416284\du}{3.281921\du}}
\pgfpathlineto{\pgfpoint{-1.365052\du}{3.279225\du}}
\pgfpathlineto{\pgfpoint{-1.314420\du}{3.274731\du}}
\pgfpathlineto{\pgfpoint{-1.263787\du}{3.270237\du}}
\pgfpathlineto{\pgfpoint{-1.214353\du}{3.264844\du}}
\pgfpathlineto{\pgfpoint{-1.165818\du}{3.257953\du}}
\pgfpathlineto{\pgfpoint{-1.118181\du}{3.250763\du}}
\pgfpathlineto{\pgfpoint{-1.072342\du}{3.242973\du}}
\pgfpathlineto{\pgfpoint{-1.027103\du}{3.233686\du}}
\pgfpathlineto{\pgfpoint{-0.983061\du}{3.224098\du}}
\pgfpathlineto{\pgfpoint{-0.941417\du}{3.213912\du}}
\pgfpathlineto{\pgfpoint{-0.900971\du}{3.202228\du}}
\pgfpathlineto{\pgfpoint{-0.862322\du}{3.190843\du}}
\pgfpathlineto{\pgfpoint{-0.825172\du}{3.178559\du}}
\pgfpathlineto{\pgfpoint{-0.789819\du}{3.165676\du}}
\pgfpathlineto{\pgfpoint{-0.756863\du}{3.151895\du}}
\pgfpathlineto{\pgfpoint{-0.726303\du}{3.137514\du}}
\pgfpathlineto{\pgfpoint{-0.697242\du}{3.123133\du}}
\pgfpathlineto{\pgfpoint{-0.670877\du}{3.108153\du}}
\pgfpathlineto{\pgfpoint{-0.646610\du}{3.092274\du}}
\pgfpathlineto{\pgfpoint{-0.624739\du}{3.075796\du}}
\pgfpathlineto{\pgfpoint{-0.604965\du}{3.059917\du}}
\pgfpathlineto{\pgfpoint{-0.588787\du}{3.043439\du}}
\pgfpathlineto{\pgfpoint{-0.574106\du}{3.026662\du}}
\pgfpathlineto{\pgfpoint{-0.562122\du}{3.009285\du}}
\pgfpathlineto{\pgfpoint{-0.553134\du}{2.992208\du}}
\pgfpathlineto{\pgfpoint{-0.546543\du}{2.974831\du}}
\pgfpathlineto{\pgfpoint{-0.542648\du}{2.957454\du}}
\pgfpathlineto{\pgfpoint{-0.541150\du}{2.939777\du}}
\pgfusepath{stroke}
\pgfsetbuttcap
\pgfsetmiterjoin
\pgfsetdash{}{0pt}
\definecolor{dialinecolor}{rgb}{0.000000, 0.470588, 0.666667}
\pgfsetfillcolor{dialinecolor}
\pgfpathmoveto{\pgfpoint{-2.604500\du}{2.456821\du}}
\pgfpathlineto{\pgfpoint{-2.604500\du}{2.951761\du}}
\pgfpathlineto{\pgfpoint{-0.541150\du}{2.951761\du}}
\pgfpathlineto{\pgfpoint{-0.541150\du}{2.456821\du}}
\pgfpathlineto{\pgfpoint{-2.604500\du}{2.456821\du}}
\pgfusepath{fill}
\pgfsetbuttcap
\pgfsetmiterjoin
\pgfsetdash{}{0pt}
\definecolor{dialinecolor}{rgb}{0.000000, 0.705882, 1.000000}
\pgfsetfillcolor{dialinecolor}
\pgfpathmoveto{\pgfpoint{-0.529166\du}{2.449930\du}}
\pgfpathlineto{\pgfpoint{-0.530664\du}{2.432554\du}}
\pgfpathlineto{\pgfpoint{-0.534559\du}{2.414578\du}}
\pgfpathlineto{\pgfpoint{-0.541150\du}{2.396901\du}}
\pgfpathlineto{\pgfpoint{-0.550438\du}{2.379225\du}}
\pgfpathlineto{\pgfpoint{-0.562422\du}{2.361548\du}}
\pgfpathlineto{\pgfpoint{-0.576803\du}{2.344471\du}}
\pgfpathlineto{\pgfpoint{-0.593880\du}{2.327693\du}}
\pgfpathlineto{\pgfpoint{-0.613054\du}{2.310916\du}}
\pgfpathlineto{\pgfpoint{-0.635225\du}{2.295037\du}}
\pgfpathlineto{\pgfpoint{-0.659492\du}{2.279158\du}}
\pgfpathlineto{\pgfpoint{-0.686157\du}{2.263878\du}}
\pgfpathlineto{\pgfpoint{-0.715218\du}{2.248299\du}}
\pgfpathlineto{\pgfpoint{-0.746377\du}{2.234218\du}}
\pgfpathlineto{\pgfpoint{-0.779632\du}{2.220137\du}}
\pgfpathlineto{\pgfpoint{-0.814985\du}{2.206954\du}}
\pgfpathlineto{\pgfpoint{-0.851836\du}{2.194671\du}}
\pgfpathlineto{\pgfpoint{-0.890784\du}{2.182687\du}}
\pgfpathlineto{\pgfpoint{-0.931829\du}{2.171302\du}}
\pgfpathlineto{\pgfpoint{-0.974073\du}{2.160816\du}}
\pgfpathlineto{\pgfpoint{-1.017815\du}{2.150929\du}}
\pgfpathlineto{\pgfpoint{-1.063055\du}{2.141641\du}}
\pgfpathlineto{\pgfpoint{-1.109792\du}{2.133852\du}}
\pgfpathlineto{\pgfpoint{-1.157429\du}{2.125762\du}}
\pgfpathlineto{\pgfpoint{-1.206563\du}{2.119171\du}}
\pgfpathlineto{\pgfpoint{-1.255997\du}{2.113479\du}}
\pgfpathlineto{\pgfpoint{-1.306930\du}{2.108685\du}}
\pgfpathlineto{\pgfpoint{-1.358161\du}{2.105090\du}}
\pgfpathlineto{\pgfpoint{-1.409693\du}{2.101794\du}}
\pgfpathlineto{\pgfpoint{-1.461823\du}{2.099398\du}}
\pgfpathlineto{\pgfpoint{-1.514253\du}{2.097600\du}}
\pgfpathlineto{\pgfpoint{-1.567283\du}{2.097600\du}}
\pgfpathlineto{\pgfpoint{-1.567283\du}{2.097600\du}}
\pgfpathlineto{\pgfpoint{-1.619713\du}{2.097600\du}}
\pgfpathlineto{\pgfpoint{-1.671843\du}{2.099398\du}}
\pgfpathlineto{\pgfpoint{-1.723974\du}{2.101794\du}}
\pgfpathlineto{\pgfpoint{-1.775505\du}{2.105090\du}}
\pgfpathlineto{\pgfpoint{-1.826737\du}{2.108685\du}}
\pgfpathlineto{\pgfpoint{-1.877669\du}{2.113479\du}}
\pgfpathlineto{\pgfpoint{-1.927103\du}{2.119171\du}}
\pgfpathlineto{\pgfpoint{-1.976237\du}{2.125762\du}}
\pgfpathlineto{\pgfpoint{-2.023874\du}{2.133852\du}}
\pgfpathlineto{\pgfpoint{-2.070612\du}{2.141641\du}}
\pgfpathlineto{\pgfpoint{-2.116151\du}{2.150929\du}}
\pgfpathlineto{\pgfpoint{-2.159593\du}{2.160816\du}}
\pgfpathlineto{\pgfpoint{-2.201837\du}{2.171302\du}}
\pgfpathlineto{\pgfpoint{-2.242882\du}{2.182687\du}}
\pgfpathlineto{\pgfpoint{-2.281830\du}{2.194671\du}}
\pgfpathlineto{\pgfpoint{-2.318681\du}{2.206954\du}}
\pgfpathlineto{\pgfpoint{-2.354034\du}{2.220137\du}}
\pgfpathlineto{\pgfpoint{-2.387290\du}{2.234218\du}}
\pgfpathlineto{\pgfpoint{-2.418448\du}{2.248299\du}}
\pgfpathlineto{\pgfpoint{-2.447809\du}{2.263878\du}}
\pgfpathlineto{\pgfpoint{-2.474174\du}{2.279158\du}}
\pgfpathlineto{\pgfpoint{-2.498741\du}{2.295037\du}}
\pgfpathlineto{\pgfpoint{-2.520612\du}{2.310916\du}}
\pgfpathlineto{\pgfpoint{-2.539786\du}{2.327693\du}}
\pgfpathlineto{\pgfpoint{-2.556863\du}{2.344471\du}}
\pgfpathlineto{\pgfpoint{-2.571244\du}{2.361548\du}}
\pgfpathlineto{\pgfpoint{-2.583228\du}{2.379225\du}}
\pgfpathlineto{\pgfpoint{-2.592516\du}{2.396901\du}}
\pgfpathlineto{\pgfpoint{-2.599407\du}{2.414578\du}}
\pgfpathlineto{\pgfpoint{-2.603002\du}{2.432554\du}}
\pgfpathlineto{\pgfpoint{-2.604500\du}{2.449930\du}}
\pgfpathlineto{\pgfpoint{-2.604500\du}{2.449930\du}}
\pgfpathlineto{\pgfpoint{-2.603002\du}{2.468206\du}}
\pgfpathlineto{\pgfpoint{-2.599407\du}{2.486182\du}}
\pgfpathlineto{\pgfpoint{-2.592516\du}{2.503858\du}}
\pgfpathlineto{\pgfpoint{-2.583228\du}{2.521235\du}}
\pgfpathlineto{\pgfpoint{-2.571244\du}{2.538912\du}}
\pgfpathlineto{\pgfpoint{-2.556863\du}{2.555689\du}}
\pgfpathlineto{\pgfpoint{-2.539786\du}{2.573066\du}}
\pgfpathlineto{\pgfpoint{-2.520612\du}{2.589544\du}}
\pgfpathlineto{\pgfpoint{-2.498741\du}{2.605723\du}}
\pgfpathlineto{\pgfpoint{-2.474174\du}{2.621602\du}}
\pgfpathlineto{\pgfpoint{-2.447809\du}{2.636881\du}}
\pgfpathlineto{\pgfpoint{-2.418448\du}{2.652161\du}}
\pgfpathlineto{\pgfpoint{-2.387290\du}{2.666542\du}}
\pgfpathlineto{\pgfpoint{-2.354034\du}{2.679724\du}}
\pgfpathlineto{\pgfpoint{-2.318681\du}{2.693506\du}}
\pgfpathlineto{\pgfpoint{-2.281830\du}{2.706089\du}}
\pgfpathlineto{\pgfpoint{-2.242882\du}{2.718073\du}}
\pgfpathlineto{\pgfpoint{-2.201837\du}{2.729458\du}}
\pgfpathlineto{\pgfpoint{-2.159593\du}{2.739944\du}}
\pgfpathlineto{\pgfpoint{-2.116151\du}{2.749531\du}}
\pgfpathlineto{\pgfpoint{-2.070612\du}{2.758819\du}}
\pgfpathlineto{\pgfpoint{-2.023874\du}{2.766908\du}}
\pgfpathlineto{\pgfpoint{-1.976237\du}{2.774697\du}}
\pgfpathlineto{\pgfpoint{-1.927103\du}{2.781289\du}}
\pgfpathlineto{\pgfpoint{-1.877669\du}{2.786981\du}}
\pgfpathlineto{\pgfpoint{-1.826737\du}{2.791775\du}}
\pgfpathlineto{\pgfpoint{-1.775505\du}{2.795370\du}}
\pgfpathlineto{\pgfpoint{-1.723974\du}{2.798666\du}}
\pgfpathlineto{\pgfpoint{-1.671843\du}{2.801362\du}}
\pgfpathlineto{\pgfpoint{-1.619713\du}{2.802860\du}}
\pgfpathlineto{\pgfpoint{-1.567283\du}{2.802860\du}}
\pgfpathlineto{\pgfpoint{-1.567283\du}{2.802860\du}}
\pgfpathlineto{\pgfpoint{-1.514253\du}{2.802860\du}}
\pgfpathlineto{\pgfpoint{-1.461823\du}{2.801362\du}}
\pgfpathlineto{\pgfpoint{-1.409693\du}{2.798666\du}}
\pgfpathlineto{\pgfpoint{-1.358161\du}{2.795370\du}}
\pgfpathlineto{\pgfpoint{-1.306930\du}{2.791775\du}}
\pgfpathlineto{\pgfpoint{-1.255997\du}{2.786981\du}}
\pgfpathlineto{\pgfpoint{-1.206563\du}{2.781289\du}}
\pgfpathlineto{\pgfpoint{-1.157429\du}{2.774697\du}}
\pgfpathlineto{\pgfpoint{-1.109792\du}{2.766908\du}}
\pgfpathlineto{\pgfpoint{-1.063055\du}{2.758819\du}}
\pgfpathlineto{\pgfpoint{-1.017815\du}{2.749531\du}}
\pgfpathlineto{\pgfpoint{-0.974073\du}{2.739944\du}}
\pgfpathlineto{\pgfpoint{-0.931829\du}{2.729458\du}}
\pgfpathlineto{\pgfpoint{-0.890784\du}{2.718073\du}}
\pgfpathlineto{\pgfpoint{-0.851836\du}{2.706089\du}}
\pgfpathlineto{\pgfpoint{-0.814985\du}{2.693506\du}}
\pgfpathlineto{\pgfpoint{-0.779632\du}{2.679724\du}}
\pgfpathlineto{\pgfpoint{-0.746377\du}{2.666542\du}}
\pgfpathlineto{\pgfpoint{-0.715218\du}{2.652161\du}}
\pgfpathlineto{\pgfpoint{-0.686157\du}{2.636881\du}}
\pgfpathlineto{\pgfpoint{-0.659492\du}{2.621602\du}}
\pgfpathlineto{\pgfpoint{-0.635225\du}{2.605723\du}}
\pgfpathlineto{\pgfpoint{-0.613054\du}{2.589544\du}}
\pgfpathlineto{\pgfpoint{-0.593880\du}{2.573066\du}}
\pgfpathlineto{\pgfpoint{-0.576803\du}{2.555689\du}}
\pgfpathlineto{\pgfpoint{-0.562422\du}{2.538912\du}}
\pgfpathlineto{\pgfpoint{-0.550438\du}{2.521235\du}}
\pgfpathlineto{\pgfpoint{-0.541150\du}{2.503858\du}}
\pgfpathlineto{\pgfpoint{-0.534559\du}{2.486182\du}}
\pgfpathlineto{\pgfpoint{-0.530664\du}{2.468206\du}}
\pgfpathlineto{\pgfpoint{-0.529166\du}{2.449930\du}}
\pgfusepath{fill}
\pgfsetbuttcap
\pgfsetmiterjoin
\pgfsetdash{}{0pt}
\definecolor{dialinecolor}{rgb}{0.666667, 0.901961, 1.000000}
\pgfsetstrokecolor{dialinecolor}
\pgfpathmoveto{\pgfpoint{-0.541150\du}{2.444837\du}}
\pgfpathlineto{\pgfpoint{-0.542648\du}{2.426861\du}}
\pgfpathlineto{\pgfpoint{-0.546543\du}{2.409185\du}}
\pgfpathlineto{\pgfpoint{-0.553134\du}{2.392407\du}}
\pgfpathlineto{\pgfpoint{-0.562122\du}{2.374731\du}}
\pgfpathlineto{\pgfpoint{-0.574106\du}{2.357653\du}}
\pgfpathlineto{\pgfpoint{-0.588787\du}{2.340876\du}}
\pgfpathlineto{\pgfpoint{-0.604965\du}{2.324697\du}}
\pgfpathlineto{\pgfpoint{-0.624739\du}{2.307920\du}}
\pgfpathlineto{\pgfpoint{-0.646610\du}{2.292340\du}}
\pgfpathlineto{\pgfpoint{-0.670877\du}{2.276162\du}}
\pgfpathlineto{\pgfpoint{-0.697242\du}{2.261182\du}}
\pgfpathlineto{\pgfpoint{-0.726303\du}{2.246202\du}}
\pgfpathlineto{\pgfpoint{-0.756863\du}{2.232121\du}}
\pgfpathlineto{\pgfpoint{-0.789819\du}{2.218938\du}}
\pgfpathlineto{\pgfpoint{-0.825172\du}{2.205157\du}}
\pgfpathlineto{\pgfpoint{-0.862322\du}{2.193173\du}}
\pgfpathlineto{\pgfpoint{-0.900971\du}{2.181488\du}}
\pgfpathlineto{\pgfpoint{-0.941417\du}{2.170403\du}}
\pgfpathlineto{\pgfpoint{-0.983061\du}{2.160217\du}}
\pgfpathlineto{\pgfpoint{-1.027103\du}{2.150030\du}}
\pgfpathlineto{\pgfpoint{-1.072342\du}{2.141342\du}}
\pgfpathlineto{\pgfpoint{-1.118181\du}{2.133252\du}}
\pgfpathlineto{\pgfpoint{-1.165818\du}{2.125762\du}}
\pgfpathlineto{\pgfpoint{-1.214353\du}{2.119171\du}}
\pgfpathlineto{\pgfpoint{-1.263787\du}{2.113479\du}}
\pgfpathlineto{\pgfpoint{-1.314420\du}{2.109284\du}}
\pgfpathlineto{\pgfpoint{-1.365052\du}{2.105090\du}}
\pgfpathlineto{\pgfpoint{-1.416284\du}{2.102094\du}}
\pgfpathlineto{\pgfpoint{-1.468414\du}{2.099398\du}}
\pgfpathlineto{\pgfpoint{-1.520545\du}{2.098199\du}}
\pgfpathlineto{\pgfpoint{-1.572975\du}{2.097600\du}}
\pgfpathlineto{\pgfpoint{-1.572975\du}{2.097600\du}}
\pgfpathlineto{\pgfpoint{-1.625105\du}{2.098199\du}}
\pgfpathlineto{\pgfpoint{-1.677536\du}{2.099398\du}}
\pgfpathlineto{\pgfpoint{-1.729366\du}{2.102094\du}}
\pgfpathlineto{\pgfpoint{-1.780598\du}{2.105090\du}}
\pgfpathlineto{\pgfpoint{-1.831530\du}{2.109284\du}}
\pgfpathlineto{\pgfpoint{-1.881863\du}{2.113479\du}}
\pgfpathlineto{\pgfpoint{-1.931297\du}{2.119171\du}}
\pgfpathlineto{\pgfpoint{-1.980132\du}{2.125762\du}}
\pgfpathlineto{\pgfpoint{-2.027469\du}{2.133252\du}}
\pgfpathlineto{\pgfpoint{-2.073308\du}{2.141342\du}}
\pgfpathlineto{\pgfpoint{-2.118847\du}{2.150030\du}}
\pgfpathlineto{\pgfpoint{-2.162289\du}{2.160217\du}}
\pgfpathlineto{\pgfpoint{-2.203934\du}{2.170403\du}}
\pgfpathlineto{\pgfpoint{-2.244680\du}{2.181488\du}}
\pgfpathlineto{\pgfpoint{-2.283927\du}{2.193173\du}}
\pgfpathlineto{\pgfpoint{-2.321078\du}{2.205157\du}}
\pgfpathlineto{\pgfpoint{-2.355831\du}{2.218938\du}}
\pgfpathlineto{\pgfpoint{-2.389087\du}{2.232121\du}}
\pgfpathlineto{\pgfpoint{-2.419946\du}{2.246202\du}}
\pgfpathlineto{\pgfpoint{-2.448408\du}{2.261182\du}}
\pgfpathlineto{\pgfpoint{-2.474473\du}{2.276162\du}}
\pgfpathlineto{\pgfpoint{-2.499041\du}{2.292340\du}}
\pgfpathlineto{\pgfpoint{-2.521211\du}{2.307920\du}}
\pgfpathlineto{\pgfpoint{-2.540385\du}{2.324697\du}}
\pgfpathlineto{\pgfpoint{-2.557163\du}{2.340876\du}}
\pgfpathlineto{\pgfpoint{-2.572143\du}{2.357653\du}}
\pgfpathlineto{\pgfpoint{-2.583228\du}{2.374731\du}}
\pgfpathlineto{\pgfpoint{-2.592816\du}{2.392407\du}}
\pgfpathlineto{\pgfpoint{-2.599407\du}{2.409185\du}}
\pgfpathlineto{\pgfpoint{-2.603002\du}{2.426861\du}}
\pgfpathlineto{\pgfpoint{-2.604500\du}{2.444837\du}}
\pgfpathlineto{\pgfpoint{-2.604500\du}{2.444837\du}}
\pgfpathlineto{\pgfpoint{-2.603002\du}{2.462214\du}}
\pgfpathlineto{\pgfpoint{-2.599407\du}{2.479890\du}}
\pgfpathlineto{\pgfpoint{-2.592816\du}{2.496668\du}}
\pgfpathlineto{\pgfpoint{-2.583228\du}{2.514045\du}}
\pgfpathlineto{\pgfpoint{-2.572143\du}{2.531422\du}}
\pgfpathlineto{\pgfpoint{-2.557163\du}{2.548199\du}}
\pgfpathlineto{\pgfpoint{-2.540385\du}{2.564977\du}}
\pgfpathlineto{\pgfpoint{-2.521211\du}{2.581455\du}}
\pgfpathlineto{\pgfpoint{-2.499041\du}{2.597334\du}}
\pgfpathlineto{\pgfpoint{-2.474473\du}{2.612913\du}}
\pgfpathlineto{\pgfpoint{-2.448408\du}{2.628193\du}}
\pgfpathlineto{\pgfpoint{-2.419946\du}{2.642274\du}}
\pgfpathlineto{\pgfpoint{-2.389087\du}{2.656954\du}}
\pgfpathlineto{\pgfpoint{-2.355831\du}{2.670436\du}}
\pgfpathlineto{\pgfpoint{-2.321078\du}{2.683619\du}}
\pgfpathlineto{\pgfpoint{-2.283927\du}{2.695903\du}}
\pgfpathlineto{\pgfpoint{-2.244680\du}{2.707587\du}}
\pgfpathlineto{\pgfpoint{-2.203934\du}{2.718672\du}}
\pgfpathlineto{\pgfpoint{-2.162289\du}{2.728859\du}}
\pgfpathlineto{\pgfpoint{-2.118847\du}{2.738745\du}}
\pgfpathlineto{\pgfpoint{-2.073308\du}{2.748033\du}}
\pgfpathlineto{\pgfpoint{-2.027469\du}{2.755523\du}}
\pgfpathlineto{\pgfpoint{-1.980132\du}{2.763313\du}}
\pgfpathlineto{\pgfpoint{-1.931297\du}{2.769604\du}}
\pgfpathlineto{\pgfpoint{-1.881863\du}{2.775297\du}}
\pgfpathlineto{\pgfpoint{-1.831530\du}{2.780390\du}}
\pgfpathlineto{\pgfpoint{-1.780598\du}{2.784584\du}}
\pgfpathlineto{\pgfpoint{-1.729366\du}{2.786981\du}}
\pgfpathlineto{\pgfpoint{-1.677536\du}{2.789378\du}}
\pgfpathlineto{\pgfpoint{-1.625105\du}{2.790876\du}}
\pgfpathlineto{\pgfpoint{-1.572975\du}{2.791176\du}}
\pgfpathlineto{\pgfpoint{-1.572975\du}{2.791176\du}}
\pgfpathlineto{\pgfpoint{-1.520545\du}{2.790876\du}}
\pgfpathlineto{\pgfpoint{-1.468414\du}{2.789378\du}}
\pgfpathlineto{\pgfpoint{-1.416284\du}{2.786981\du}}
\pgfpathlineto{\pgfpoint{-1.365052\du}{2.784584\du}}
\pgfpathlineto{\pgfpoint{-1.314420\du}{2.780390\du}}
\pgfpathlineto{\pgfpoint{-1.263787\du}{2.775297\du}}
\pgfpathlineto{\pgfpoint{-1.214353\du}{2.769604\du}}
\pgfpathlineto{\pgfpoint{-1.165818\du}{2.763313\du}}
\pgfpathlineto{\pgfpoint{-1.118181\du}{2.755523\du}}
\pgfpathlineto{\pgfpoint{-1.072342\du}{2.748033\du}}
\pgfpathlineto{\pgfpoint{-1.027103\du}{2.738745\du}}
\pgfpathlineto{\pgfpoint{-0.983061\du}{2.728859\du}}
\pgfpathlineto{\pgfpoint{-0.941417\du}{2.718672\du}}
\pgfpathlineto{\pgfpoint{-0.900971\du}{2.707587\du}}
\pgfpathlineto{\pgfpoint{-0.862322\du}{2.695903\du}}
\pgfpathlineto{\pgfpoint{-0.825172\du}{2.683619\du}}
\pgfpathlineto{\pgfpoint{-0.789819\du}{2.670436\du}}
\pgfpathlineto{\pgfpoint{-0.756863\du}{2.656954\du}}
\pgfpathlineto{\pgfpoint{-0.726303\du}{2.642274\du}}
\pgfpathlineto{\pgfpoint{-0.697242\du}{2.628193\du}}
\pgfpathlineto{\pgfpoint{-0.670877\du}{2.612913\du}}
\pgfpathlineto{\pgfpoint{-0.646610\du}{2.597334\du}}
\pgfpathlineto{\pgfpoint{-0.624739\du}{2.581455\du}}
\pgfpathlineto{\pgfpoint{-0.604965\du}{2.564977\du}}
\pgfpathlineto{\pgfpoint{-0.588787\du}{2.548199\du}}
\pgfpathlineto{\pgfpoint{-0.574106\du}{2.531422\du}}
\pgfpathlineto{\pgfpoint{-0.562122\du}{2.514045\du}}
\pgfpathlineto{\pgfpoint{-0.553134\du}{2.496668\du}}
\pgfpathlineto{\pgfpoint{-0.546543\du}{2.479890\du}}
\pgfpathlineto{\pgfpoint{-0.542648\du}{2.462214\du}}
\pgfpathlineto{\pgfpoint{-0.541150\du}{2.444837\du}}
\pgfusepath{stroke}
\pgfsetbuttcap
\pgfsetmiterjoin
\pgfsetdash{}{0pt}
\definecolor{dialinecolor}{rgb}{0.000000, 0.000000, 0.000000}
\pgfsetfillcolor{dialinecolor}
\pgfpathmoveto{\pgfpoint{-1.548108\du}{2.369937\du}}
\pgfpathlineto{\pgfpoint{-1.396810\du}{2.419671\du}}
\pgfpathlineto{\pgfpoint{-1.031597\du}{2.270769\du}}
\pgfpathlineto{\pgfpoint{-0.868314\du}{2.320503\du}}
\pgfpathlineto{\pgfpoint{-0.956996\du}{2.196768\du}}
\pgfpathlineto{\pgfpoint{-1.383927\du}{2.196768\du}}
\pgfpathlineto{\pgfpoint{-1.208361\du}{2.234218\du}}
\pgfpathlineto{\pgfpoint{-1.548108\du}{2.369937\du}}
\pgfusepath{fill}
\pgfsetbuttcap
\pgfsetmiterjoin
\pgfsetdash{}{0pt}
\definecolor{dialinecolor}{rgb}{0.000000, 0.000000, 0.000000}
\pgfsetfillcolor{dialinecolor}
\pgfpathmoveto{\pgfpoint{-1.611324\du}{2.506555\du}}
\pgfpathlineto{\pgfpoint{-1.761424\du}{2.456821\du}}
\pgfpathlineto{\pgfpoint{-2.114054\du}{2.605723\du}}
\pgfpathlineto{\pgfpoint{-2.289919\du}{2.555689\du}}
\pgfpathlineto{\pgfpoint{-2.201837\du}{2.692008\du}}
\pgfpathlineto{\pgfpoint{-1.761424\du}{2.692008\du}}
\pgfpathlineto{\pgfpoint{-1.950472\du}{2.642274\du}}
\pgfpathlineto{\pgfpoint{-1.611324\du}{2.506555\du}}
\pgfusepath{fill}
\pgfsetbuttcap
\pgfsetmiterjoin
\pgfsetdash{}{0pt}
\definecolor{dialinecolor}{rgb}{0.000000, 0.000000, 0.000000}
\pgfsetfillcolor{dialinecolor}
\pgfpathmoveto{\pgfpoint{-2.252469\du}{2.234218\du}}
\pgfpathlineto{\pgfpoint{-2.101470\du}{2.184484\du}}
\pgfpathlineto{\pgfpoint{-1.736856\du}{2.320503\du}}
\pgfpathlineto{\pgfpoint{-1.572975\du}{2.283652\du}}
\pgfpathlineto{\pgfpoint{-1.660758\du}{2.407087\du}}
\pgfpathlineto{\pgfpoint{-2.088887\du}{2.407087\du}}
\pgfpathlineto{\pgfpoint{-1.913022\du}{2.369937\du}}
\pgfpathlineto{\pgfpoint{-2.252469\du}{2.234218\du}}
\pgfusepath{fill}
\pgfsetbuttcap
\pgfsetmiterjoin
\pgfsetdash{}{0pt}
\definecolor{dialinecolor}{rgb}{0.000000, 0.000000, 0.000000}
\pgfsetfillcolor{dialinecolor}
\pgfpathmoveto{\pgfpoint{-0.893780\du}{2.654857\du}}
\pgfpathlineto{\pgfpoint{-1.044479\du}{2.704291\du}}
\pgfpathlineto{\pgfpoint{-1.396810\du}{2.555689\du}}
\pgfpathlineto{\pgfpoint{-1.572975\du}{2.605723\du}}
\pgfpathlineto{\pgfpoint{-1.484892\du}{2.481688\du}}
\pgfpathlineto{\pgfpoint{-1.044479\du}{2.481688\du}}
\pgfpathlineto{\pgfpoint{-1.233527\du}{2.518239\du}}
\pgfpathlineto{\pgfpoint{-0.893780\du}{2.654857\du}}
\pgfusepath{fill}
\pgfsetbuttcap
\pgfsetmiterjoin
\pgfsetdash{}{0pt}
\definecolor{dialinecolor}{rgb}{1.000000, 1.000000, 1.000000}
\pgfsetfillcolor{dialinecolor}
\pgfpathmoveto{\pgfpoint{-1.535525\du}{2.382221\du}}
\pgfpathlineto{\pgfpoint{-1.383927\du}{2.431655\du}}
\pgfpathlineto{\pgfpoint{-1.019313\du}{2.283652\du}}
\pgfpathlineto{\pgfpoint{-0.855731\du}{2.333386\du}}
\pgfpathlineto{\pgfpoint{-0.943814\du}{2.209051\du}}
\pgfpathlineto{\pgfpoint{-1.371943\du}{2.209051\du}}
\pgfpathlineto{\pgfpoint{-1.195478\du}{2.246202\du}}
\pgfpathlineto{\pgfpoint{-1.535525\du}{2.382221\du}}
\pgfusepath{fill}
\pgfsetbuttcap
\pgfsetmiterjoin
\pgfsetdash{}{0pt}
\definecolor{dialinecolor}{rgb}{1.000000, 1.000000, 1.000000}
\pgfsetfillcolor{dialinecolor}
\pgfpathmoveto{\pgfpoint{-1.597842\du}{2.518239\du}}
\pgfpathlineto{\pgfpoint{-1.749440\du}{2.469105\du}}
\pgfpathlineto{\pgfpoint{-2.101470\du}{2.617707\du}}
\pgfpathlineto{\pgfpoint{-2.277636\du}{2.567973\du}}
\pgfpathlineto{\pgfpoint{-2.189553\du}{2.704291\du}}
\pgfpathlineto{\pgfpoint{-1.749440\du}{2.704291\du}}
\pgfpathlineto{\pgfpoint{-1.937589\du}{2.654857\du}}
\pgfpathlineto{\pgfpoint{-1.597842\du}{2.518239\du}}
\pgfusepath{fill}
\pgfsetbuttcap
\pgfsetmiterjoin
\pgfsetdash{}{0pt}
\definecolor{dialinecolor}{rgb}{1.000000, 1.000000, 1.000000}
\pgfsetfillcolor{dialinecolor}
\pgfpathmoveto{\pgfpoint{-2.239586\du}{2.246202\du}}
\pgfpathlineto{\pgfpoint{-2.088887\du}{2.196768\du}}
\pgfpathlineto{\pgfpoint{-1.723974\du}{2.333386\du}}
\pgfpathlineto{\pgfpoint{-1.560392\du}{2.295636\du}}
\pgfpathlineto{\pgfpoint{-1.648474\du}{2.419671\du}}
\pgfpathlineto{\pgfpoint{-2.076304\du}{2.419671\du}}
\pgfpathlineto{\pgfpoint{-1.899839\du}{2.382221\du}}
\pgfpathlineto{\pgfpoint{-2.239586\du}{2.246202\du}}
\pgfusepath{fill}
\pgfsetbuttcap
\pgfsetmiterjoin
\pgfsetdash{}{0pt}
\definecolor{dialinecolor}{rgb}{1.000000, 1.000000, 1.000000}
\pgfsetfillcolor{dialinecolor}
\pgfpathmoveto{\pgfpoint{-0.880897\du}{2.667141\du}}
\pgfpathlineto{\pgfpoint{-1.031597\du}{2.716875\du}}
\pgfpathlineto{\pgfpoint{-1.383927\du}{2.567973\du}}
\pgfpathlineto{\pgfpoint{-1.560392\du}{2.617707\du}}
\pgfpathlineto{\pgfpoint{-1.472609\du}{2.494271\du}}
\pgfpathlineto{\pgfpoint{-1.031597\du}{2.494271\du}}
\pgfpathlineto{\pgfpoint{-1.220345\du}{2.531122\du}}
\pgfpathlineto{\pgfpoint{-0.880897\du}{2.667141\du}}
\pgfusepath{fill}
\pgfsetbuttcap
\pgfsetmiterjoin
\pgfsetdash{}{0pt}
\definecolor{dialinecolor}{rgb}{0.666667, 0.901961, 1.000000}
\pgfsetstrokecolor{dialinecolor}
\pgfpathmoveto{\pgfpoint{-2.604500\du}{2.444837\du}}
\pgfpathlineto{\pgfpoint{-2.604500\du}{2.939178\du}}
\pgfusepath{stroke}
\pgfsetbuttcap
\pgfsetmiterjoin
\pgfsetdash{}{0pt}
\definecolor{dialinecolor}{rgb}{0.666667, 0.901961, 1.000000}
\pgfsetstrokecolor{dialinecolor}
\pgfpathmoveto{\pgfpoint{-0.541150\du}{2.444837\du}}
\pgfpathlineto{\pgfpoint{-0.541150\du}{2.939178\du}}
\pgfusepath{stroke}
\pgfsetbuttcap
\pgfsetmiterjoin
\pgfsetdash{}{0pt}
\definecolor{dialinecolor}{rgb}{0.000000, 0.000000, 0.000000}
\pgfsetfillcolor{dialinecolor}
\pgfpathmoveto{\pgfpoint{-1.950472\du}{2.877760\du}}
\pgfpathlineto{\pgfpoint{-1.723974\du}{2.877760\du}}
\pgfpathlineto{\pgfpoint{-1.585558\du}{2.989212\du}}
\pgfpathlineto{\pgfpoint{-1.447442\du}{2.877760\du}}
\pgfpathlineto{\pgfpoint{-1.233527\du}{2.877760\du}}
\pgfpathlineto{\pgfpoint{-1.233527\du}{2.828326\du}}
\pgfpathlineto{\pgfpoint{-1.082229\du}{2.902028\du}}
\pgfpathlineto{\pgfpoint{-1.233527\du}{2.976628\du}}
\pgfpathlineto{\pgfpoint{-1.233527\du}{2.927194\du}}
\pgfpathlineto{\pgfpoint{-1.421976\du}{2.927194\du}}
\pgfpathlineto{\pgfpoint{-1.548108\du}{3.013779\du}}
\pgfpathlineto{\pgfpoint{-1.421976\du}{3.112347\du}}
\pgfpathlineto{\pgfpoint{-1.233527\du}{3.112347\du}}
\pgfpathlineto{\pgfpoint{-1.233527\du}{3.051229\du}}
\pgfpathlineto{\pgfpoint{-1.082229\du}{3.137514\du}}
\pgfpathlineto{\pgfpoint{-1.233527\du}{3.212114\du}}
\pgfpathlineto{\pgfpoint{-1.233527\du}{3.162081\du}}
\pgfpathlineto{\pgfpoint{-1.447442\du}{3.162081\du}}
\pgfpathlineto{\pgfpoint{-1.585558\du}{3.051229\du}}
\pgfpathlineto{\pgfpoint{-1.723974\du}{3.162081\du}}
\pgfpathlineto{\pgfpoint{-1.950472\du}{3.162081\du}}
\pgfpathlineto{\pgfpoint{-1.950472\du}{3.212114\du}}
\pgfpathlineto{\pgfpoint{-2.101470\du}{3.137514\du}}
\pgfpathlineto{\pgfpoint{-1.950472\du}{3.051229\du}}
\pgfpathlineto{\pgfpoint{-1.950472\du}{3.112347\du}}
\pgfpathlineto{\pgfpoint{-1.749440\du}{3.112347\du}}
\pgfpathlineto{\pgfpoint{-1.623308\du}{3.026062\du}}
\pgfpathlineto{\pgfpoint{-1.749440\du}{2.927194\du}}
\pgfpathlineto{\pgfpoint{-1.950472\du}{2.927194\du}}
\pgfpathlineto{\pgfpoint{-1.950472\du}{2.976628\du}}
\pgfpathlineto{\pgfpoint{-2.101470\du}{2.902028\du}}
\pgfpathlineto{\pgfpoint{-1.950472\du}{2.828326\du}}
\pgfpathlineto{\pgfpoint{-1.950472\du}{2.877760\du}}
\pgfusepath{fill}
\pgfsetbuttcap
\pgfsetmiterjoin
\pgfsetdash{}{0pt}
\definecolor{dialinecolor}{rgb}{1.000000, 1.000000, 1.000000}
\pgfsetfillcolor{dialinecolor}
\pgfpathmoveto{\pgfpoint{-1.937589\du}{2.889744\du}}
\pgfpathlineto{\pgfpoint{-1.711390\du}{2.889744\du}}
\pgfpathlineto{\pgfpoint{-1.572975\du}{3.000896\du}}
\pgfpathlineto{\pgfpoint{-1.434859\du}{2.889744\du}}
\pgfpathlineto{\pgfpoint{-1.220345\du}{2.889744\du}}
\pgfpathlineto{\pgfpoint{-1.220345\du}{2.840310\du}}
\pgfpathlineto{\pgfpoint{-1.069945\du}{2.914911\du}}
\pgfpathlineto{\pgfpoint{-1.220345\du}{2.989212\du}}
\pgfpathlineto{\pgfpoint{-1.220345\du}{2.939777\du}}
\pgfpathlineto{\pgfpoint{-1.409393\du}{2.939777\du}}
\pgfpathlineto{\pgfpoint{-1.535525\du}{3.026062\du}}
\pgfpathlineto{\pgfpoint{-1.409393\du}{3.124931\du}}
\pgfpathlineto{\pgfpoint{-1.220345\du}{3.124931\du}}
\pgfpathlineto{\pgfpoint{-1.220345\du}{3.063213\du}}
\pgfpathlineto{\pgfpoint{-1.069945\du}{3.149798\du}}
\pgfpathlineto{\pgfpoint{-1.220345\du}{3.224098\du}}
\pgfpathlineto{\pgfpoint{-1.220345\du}{3.174365\du}}
\pgfpathlineto{\pgfpoint{-1.434859\du}{3.174365\du}}
\pgfpathlineto{\pgfpoint{-1.572975\du}{3.063213\du}}
\pgfpathlineto{\pgfpoint{-1.711390\du}{3.174365\du}}
\pgfpathlineto{\pgfpoint{-1.937589\du}{3.174365\du}}
\pgfpathlineto{\pgfpoint{-1.937589\du}{3.224098\du}}
\pgfpathlineto{\pgfpoint{-2.088887\du}{3.149798\du}}
\pgfpathlineto{\pgfpoint{-1.937589\du}{3.063213\du}}
\pgfpathlineto{\pgfpoint{-1.937589\du}{3.124931\du}}
\pgfpathlineto{\pgfpoint{-1.736856\du}{3.124931\du}}
\pgfpathlineto{\pgfpoint{-1.611324\du}{3.038346\du}}
\pgfpathlineto{\pgfpoint{-1.736856\du}{2.939777\du}}
\pgfpathlineto{\pgfpoint{-1.937589\du}{2.939777\du}}
\pgfpathlineto{\pgfpoint{-1.937589\du}{2.989212\du}}
\pgfpathlineto{\pgfpoint{-2.088887\du}{2.914911\du}}
\pgfpathlineto{\pgfpoint{-1.937589\du}{2.840310\du}}
\pgfpathlineto{\pgfpoint{-1.937589\du}{2.889744\du}}
\pgfusepath{fill}
\pgfsetlinewidth{0.100000\du}
\pgfsetdash{}{0pt}
\pgfsetdash{}{0pt}
\pgfsetbuttcap
{
\definecolor{dialinecolor}{rgb}{0.000000, 0.000000, 0.000000}
\pgfsetfillcolor{dialinecolor}
% was here!!!
\definecolor{dialinecolor}{rgb}{0.000000, 0.000000, 0.000000}
\pgfsetstrokecolor{dialinecolor}
\draw (-3.968028\du,0.152677\du)--(-2.089571\du,2.144009\du);
}
\pgfsetlinewidth{0.100000\du}
\pgfsetdash{}{0pt}
\pgfsetdash{}{0pt}
\pgfsetbuttcap
{
\definecolor{dialinecolor}{rgb}{0.000000, 0.000000, 0.000000}
\pgfsetfillcolor{dialinecolor}
% was here!!!
\definecolor{dialinecolor}{rgb}{0.000000, 0.000000, 0.000000}
\pgfsetstrokecolor{dialinecolor}
\draw (-22.153366\du,3.457075\du)--(-21.584616\du,5.750000\du);
}
\pgfsetlinewidth{0.100000\du}
\pgfsetdash{}{0pt}
\pgfsetdash{}{0pt}
\pgfsetbuttcap
{
\definecolor{dialinecolor}{rgb}{0.000000, 0.000000, 0.000000}
\pgfsetfillcolor{dialinecolor}
% was here!!!
\definecolor{dialinecolor}{rgb}{0.000000, 0.000000, 0.000000}
\pgfsetstrokecolor{dialinecolor}
\draw (-23.478668\du,3.725295\du)--(-24.446723\du,5.700884\du);
}
\pgfsetlinewidth{0.100000\du}
\pgfsetdash{}{0pt}
\pgfsetdash{}{0pt}
\pgfsetbuttcap
\pgfsetmiterjoin
\pgfsetlinewidth{0.001000\du}
\pgfsetbuttcap
\pgfsetmiterjoin
\pgfsetdash{}{0pt}
\definecolor{dialinecolor}{rgb}{0.717647, 0.717647, 0.615686}
\pgfsetfillcolor{dialinecolor}
\pgfpathmoveto{\pgfpoint{-17.353700\du}{7.136880\du}}
\pgfpathlineto{\pgfpoint{-15.355274\du}{7.136880\du}}
\pgfpathlineto{\pgfpoint{-15.355274\du}{7.506213\du}}
\pgfpathlineto{\pgfpoint{-17.353700\du}{7.506213\du}}
\pgfpathlineto{\pgfpoint{-17.353700\du}{7.136880\du}}
\pgfusepath{fill}
\pgfsetbuttcap
\pgfsetmiterjoin
\pgfsetdash{}{0pt}
\definecolor{dialinecolor}{rgb}{0.286275, 0.286275, 0.211765}
\pgfsetstrokecolor{dialinecolor}
\pgfpathmoveto{\pgfpoint{-17.353700\du}{7.136880\du}}
\pgfpathlineto{\pgfpoint{-15.355274\du}{7.136880\du}}
\pgfpathlineto{\pgfpoint{-15.355274\du}{7.506213\du}}
\pgfpathlineto{\pgfpoint{-17.353700\du}{7.506213\du}}
\pgfpathlineto{\pgfpoint{-17.353700\du}{7.136880\du}}
\pgfusepath{stroke}
\pgfsetbuttcap
\pgfsetmiterjoin
\pgfsetdash{}{0pt}
\definecolor{dialinecolor}{rgb}{0.788235, 0.788235, 0.713726}
\pgfsetfillcolor{dialinecolor}
\pgfpathmoveto{\pgfpoint{-17.353700\du}{7.136880\du}}
\pgfpathlineto{\pgfpoint{-17.141798\du}{6.935998\du}}
\pgfpathlineto{\pgfpoint{-15.143373\du}{6.935998\du}}
\pgfpathlineto{\pgfpoint{-15.355274\du}{7.136880\du}}
\pgfpathlineto{\pgfpoint{-17.353700\du}{7.136880\du}}
\pgfusepath{fill}
\pgfsetbuttcap
\pgfsetmiterjoin
\pgfsetdash{}{0pt}
\definecolor{dialinecolor}{rgb}{0.286275, 0.286275, 0.211765}
\pgfsetstrokecolor{dialinecolor}
\pgfpathmoveto{\pgfpoint{-17.353700\du}{7.136880\du}}
\pgfpathlineto{\pgfpoint{-17.141798\du}{6.935998\du}}
\pgfpathlineto{\pgfpoint{-15.143373\du}{6.935998\du}}
\pgfpathlineto{\pgfpoint{-15.355274\du}{7.136880\du}}
\pgfpathlineto{\pgfpoint{-17.353700\du}{7.136880\du}}
\pgfusepath{stroke}
\pgfsetlinewidth{0.106000\du}
\pgfsetbuttcap
\pgfsetmiterjoin
\pgfsetdash{}{0pt}
\definecolor{dialinecolor}{rgb}{0.000000, 0.000000, 0.000000}
\pgfsetstrokecolor{dialinecolor}
\pgfpathmoveto{\pgfpoint{-15.467050\du}{7.305016\du}}
\pgfpathlineto{\pgfpoint{-15.946584\du}{7.305016\du}}
\pgfusepath{stroke}
\pgfsetlinewidth{0.001000\du}
\pgfsetbuttcap
\pgfsetmiterjoin
\pgfsetdash{}{0pt}
\definecolor{dialinecolor}{rgb}{0.478431, 0.478431, 0.352941}
\pgfsetfillcolor{dialinecolor}
\pgfpathmoveto{\pgfpoint{-15.355274\du}{7.506213\du}}
\pgfpathlineto{\pgfpoint{-15.143373\du}{7.293681\du}}
\pgfpathlineto{\pgfpoint{-15.143373\du}{6.935998\du}}
\pgfpathlineto{\pgfpoint{-15.355274\du}{7.136880\du}}
\pgfpathlineto{\pgfpoint{-15.355274\du}{7.506213\du}}
\pgfusepath{fill}
\pgfsetbuttcap
\pgfsetmiterjoin
\pgfsetdash{}{0pt}
\definecolor{dialinecolor}{rgb}{0.286275, 0.286275, 0.211765}
\pgfsetstrokecolor{dialinecolor}
\pgfpathmoveto{\pgfpoint{-15.355274\du}{7.506213\du}}
\pgfpathlineto{\pgfpoint{-15.143373\du}{7.293681\du}}
\pgfpathlineto{\pgfpoint{-15.143373\du}{6.935998\du}}
\pgfpathlineto{\pgfpoint{-15.355274\du}{7.136880\du}}
\pgfpathlineto{\pgfpoint{-15.355274\du}{7.506213\du}}
\pgfusepath{stroke}
\pgfsetbuttcap
\pgfsetmiterjoin
\pgfsetdash{}{0pt}
\definecolor{dialinecolor}{rgb}{0.788235, 0.788235, 0.713726}
\pgfsetfillcolor{dialinecolor}
\pgfpathmoveto{\pgfpoint{-17.342365\du}{7.740470\du}}
\pgfpathlineto{\pgfpoint{-17.119443\du}{7.461187\du}}
\pgfpathlineto{\pgfpoint{-15.578511\du}{7.461187\du}}
\pgfpathlineto{\pgfpoint{-15.801433\du}{7.740470\du}}
\pgfpathlineto{\pgfpoint{-17.342365\du}{7.740470\du}}
\pgfusepath{fill}
\pgfsetbuttcap
\pgfsetmiterjoin
\pgfsetdash{}{0pt}
\definecolor{dialinecolor}{rgb}{0.286275, 0.286275, 0.211765}
\pgfsetstrokecolor{dialinecolor}
\pgfpathmoveto{\pgfpoint{-17.342365\du}{7.740470\du}}
\pgfpathlineto{\pgfpoint{-17.119443\du}{7.461187\du}}
\pgfpathlineto{\pgfpoint{-15.578511\du}{7.461187\du}}
\pgfpathlineto{\pgfpoint{-15.801433\du}{7.740470\du}}
\pgfpathlineto{\pgfpoint{-17.342365\du}{7.740470\du}}
\pgfusepath{stroke}
\pgfsetbuttcap
\pgfsetmiterjoin
\pgfsetdash{}{0pt}
\definecolor{dialinecolor}{rgb}{0.478431, 0.478431, 0.352941}
\pgfsetfillcolor{dialinecolor}
\pgfpathmoveto{\pgfpoint{-15.801433\du}{7.796200\du}}
\pgfpathlineto{\pgfpoint{-15.578511\du}{7.561943\du}}
\pgfpathlineto{\pgfpoint{-15.578511\du}{7.461187\du}}
\pgfpathlineto{\pgfpoint{-15.801433\du}{7.740470\du}}
\pgfpathlineto{\pgfpoint{-15.801433\du}{7.796200\du}}
\pgfusepath{fill}
\pgfsetbuttcap
\pgfsetmiterjoin
\pgfsetdash{}{0pt}
\definecolor{dialinecolor}{rgb}{0.286275, 0.286275, 0.211765}
\pgfsetstrokecolor{dialinecolor}
\pgfpathmoveto{\pgfpoint{-15.801433\du}{7.796200\du}}
\pgfpathlineto{\pgfpoint{-15.578511\du}{7.561943\du}}
\pgfpathlineto{\pgfpoint{-15.578511\du}{7.461187\du}}
\pgfpathlineto{\pgfpoint{-15.801433\du}{7.740470\du}}
\pgfpathlineto{\pgfpoint{-15.801433\du}{7.796200\du}}
\pgfusepath{stroke}
\pgfsetbuttcap
\pgfsetmiterjoin
\pgfsetdash{}{0pt}
\definecolor{dialinecolor}{rgb}{0.717647, 0.717647, 0.615686}
\pgfsetfillcolor{dialinecolor}
\pgfpathmoveto{\pgfpoint{-17.342365\du}{7.740470\du}}
\pgfpathlineto{\pgfpoint{-15.801433\du}{7.740470\du}}
\pgfpathlineto{\pgfpoint{-15.801433\du}{7.796200\du}}
\pgfpathlineto{\pgfpoint{-17.342365\du}{7.796200\du}}
\pgfpathlineto{\pgfpoint{-17.342365\du}{7.740470\du}}
\pgfusepath{fill}
\pgfsetbuttcap
\pgfsetmiterjoin
\pgfsetdash{}{0pt}
\definecolor{dialinecolor}{rgb}{0.286275, 0.286275, 0.211765}
\pgfsetstrokecolor{dialinecolor}
\pgfpathmoveto{\pgfpoint{-17.342365\du}{7.740470\du}}
\pgfpathlineto{\pgfpoint{-15.801433\du}{7.740470\du}}
\pgfpathlineto{\pgfpoint{-15.801433\du}{7.796200\du}}
\pgfpathlineto{\pgfpoint{-17.342365\du}{7.796200\du}}
\pgfpathlineto{\pgfpoint{-17.342365\du}{7.740470\du}}
\pgfusepath{stroke}
\pgfsetbuttcap
\pgfsetmiterjoin
\pgfsetdash{}{0pt}
\definecolor{dialinecolor}{rgb}{0.000000, 0.000000, 0.000000}
\pgfsetfillcolor{dialinecolor}
\pgfpathmoveto{\pgfpoint{-17.052378\du}{7.092485\du}}
\pgfpathlineto{\pgfpoint{-16.884556\du}{6.935998\du}}
\pgfpathlineto{\pgfpoint{-15.467050\du}{6.935998\du}}
\pgfpathlineto{\pgfpoint{-15.622907\du}{7.092485\du}}
\pgfpathlineto{\pgfpoint{-17.052378\du}{7.092485\du}}
\pgfusepath{fill}
\pgfsetbuttcap
\pgfsetmiterjoin
\pgfsetdash{}{0pt}
\definecolor{dialinecolor}{rgb}{0.000000, 0.000000, 0.000000}
\pgfsetstrokecolor{dialinecolor}
\pgfpathmoveto{\pgfpoint{-17.052378\du}{7.092485\du}}
\pgfpathlineto{\pgfpoint{-16.884556\du}{6.935998\du}}
\pgfpathlineto{\pgfpoint{-15.467050\du}{6.935998\du}}
\pgfpathlineto{\pgfpoint{-15.622907\du}{7.092485\du}}
\pgfpathlineto{\pgfpoint{-17.052378\du}{7.092485\du}}
\pgfusepath{stroke}
\pgfsetbuttcap
\pgfsetmiterjoin
\pgfsetdash{}{0pt}
\definecolor{dialinecolor}{rgb}{0.788235, 0.788235, 0.713726}
\pgfsetfillcolor{dialinecolor}
\pgfpathmoveto{\pgfpoint{-17.063713\du}{5.941351\du}}
\pgfpathlineto{\pgfpoint{-16.906912\du}{5.796200\du}}
\pgfpathlineto{\pgfpoint{-15.488776\du}{5.796200\du}}
\pgfpathlineto{\pgfpoint{-15.645577\du}{5.941351\du}}
\pgfpathlineto{\pgfpoint{-17.063713\du}{5.941351\du}}
\pgfusepath{fill}
\pgfsetbuttcap
\pgfsetmiterjoin
\pgfsetdash{}{0pt}
\definecolor{dialinecolor}{rgb}{0.286275, 0.286275, 0.211765}
\pgfsetstrokecolor{dialinecolor}
\pgfpathmoveto{\pgfpoint{-17.063713\du}{5.941351\du}}
\pgfpathlineto{\pgfpoint{-16.906912\du}{5.796200\du}}
\pgfpathlineto{\pgfpoint{-15.488776\du}{5.796200\du}}
\pgfpathlineto{\pgfpoint{-15.645577\du}{5.941351\du}}
\pgfpathlineto{\pgfpoint{-17.063713\du}{5.941351\du}}
\pgfusepath{stroke}
\pgfsetbuttcap
\pgfsetmiterjoin
\pgfsetdash{}{0pt}
\definecolor{dialinecolor}{rgb}{0.717647, 0.717647, 0.615686}
\pgfsetfillcolor{dialinecolor}
\pgfpathmoveto{\pgfpoint{-17.063713\du}{5.941351\du}}
\pgfpathlineto{\pgfpoint{-15.634242\du}{5.941351\du}}
\pgfpathlineto{\pgfpoint{-15.634242\du}{7.069815\du}}
\pgfpathlineto{\pgfpoint{-17.063713\du}{7.069815\du}}
\pgfpathlineto{\pgfpoint{-17.063713\du}{5.941351\du}}
\pgfusepath{fill}
\pgfsetbuttcap
\pgfsetmiterjoin
\pgfsetdash{}{0pt}
\definecolor{dialinecolor}{rgb}{0.286275, 0.286275, 0.211765}
\pgfsetstrokecolor{dialinecolor}
\pgfpathmoveto{\pgfpoint{-17.063713\du}{5.941351\du}}
\pgfpathlineto{\pgfpoint{-15.634871\du}{5.941351\du}}
\pgfpathlineto{\pgfpoint{-15.634871\du}{7.069500\du}}
\pgfpathlineto{\pgfpoint{-17.063713\du}{7.069500\du}}
\pgfpathlineto{\pgfpoint{-17.063713\du}{5.941351\du}}
\pgfusepath{stroke}
\pgfsetbuttcap
\pgfsetmiterjoin
\pgfsetdash{}{0pt}
\definecolor{dialinecolor}{rgb}{1.000000, 1.000000, 1.000000}
\pgfsetfillcolor{dialinecolor}
\pgfpathmoveto{\pgfpoint{-16.940917\du}{6.086187\du}}
\pgfpathlineto{\pgfpoint{-15.757352\du}{6.086187\du}}
\pgfpathlineto{\pgfpoint{-15.757352\du}{6.958039\du}}
\pgfpathlineto{\pgfpoint{-16.940917\du}{6.958039\du}}
\pgfpathlineto{\pgfpoint{-16.940917\du}{6.086187\du}}
\pgfusepath{fill}
\pgfsetbuttcap
\pgfsetmiterjoin
\pgfsetdash{}{0pt}
\definecolor{dialinecolor}{rgb}{0.286275, 0.286275, 0.211765}
\pgfsetstrokecolor{dialinecolor}
\pgfpathmoveto{\pgfpoint{-16.940917\du}{6.086187\du}}
\pgfpathlineto{\pgfpoint{-15.757352\du}{6.086187\du}}
\pgfpathlineto{\pgfpoint{-15.757352\du}{6.957724\du}}
\pgfpathlineto{\pgfpoint{-16.940917\du}{6.957724\du}}
\pgfpathlineto{\pgfpoint{-16.940917\du}{6.086187\du}}
\pgfusepath{stroke}
\pgfsetbuttcap
\pgfsetmiterjoin
\pgfsetdash{}{0pt}
\definecolor{dialinecolor}{rgb}{0.478431, 0.478431, 0.352941}
\pgfsetfillcolor{dialinecolor}
\pgfpathmoveto{\pgfpoint{-15.645577\du}{7.059109\du}}
\pgfpathlineto{\pgfpoint{-15.488776\du}{6.902623\du}}
\pgfpathlineto{\pgfpoint{-15.488776\du}{5.796200\du}}
\pgfpathlineto{\pgfpoint{-15.645577\du}{5.941351\du}}
\pgfpathlineto{\pgfpoint{-15.645577\du}{7.059109\du}}
\pgfusepath{fill}
\pgfsetbuttcap
\pgfsetmiterjoin
\pgfsetdash{}{0pt}
\definecolor{dialinecolor}{rgb}{0.286275, 0.286275, 0.211765}
\pgfsetstrokecolor{dialinecolor}
\pgfpathmoveto{\pgfpoint{-15.645577\du}{7.059109\du}}
\pgfpathlineto{\pgfpoint{-15.488776\du}{6.902623\du}}
\pgfpathlineto{\pgfpoint{-15.488776\du}{5.796200\du}}
\pgfpathlineto{\pgfpoint{-15.645577\du}{5.941351\du}}
\pgfpathlineto{\pgfpoint{-15.645577\du}{7.059109\du}}
\pgfusepath{stroke}
\pgfsetlinewidth{0.100000\du}
\pgfsetdash{}{0pt}
\pgfsetdash{}{0pt}
\pgfsetbuttcap
\pgfsetmiterjoin
\pgfsetlinewidth{0.001000\du}
\pgfsetbuttcap
\pgfsetmiterjoin
\pgfsetdash{}{0pt}
\definecolor{dialinecolor}{rgb}{0.717647, 0.717647, 0.615686}
\pgfsetfillcolor{dialinecolor}
\pgfpathmoveto{\pgfpoint{-11.308700\du}{6.946880\du}}
\pgfpathlineto{\pgfpoint{-9.310274\du}{6.946880\du}}
\pgfpathlineto{\pgfpoint{-9.310274\du}{7.316213\du}}
\pgfpathlineto{\pgfpoint{-11.308700\du}{7.316213\du}}
\pgfpathlineto{\pgfpoint{-11.308700\du}{6.946880\du}}
\pgfusepath{fill}
\pgfsetbuttcap
\pgfsetmiterjoin
\pgfsetdash{}{0pt}
\definecolor{dialinecolor}{rgb}{0.286275, 0.286275, 0.211765}
\pgfsetstrokecolor{dialinecolor}
\pgfpathmoveto{\pgfpoint{-11.308700\du}{6.946880\du}}
\pgfpathlineto{\pgfpoint{-9.310274\du}{6.946880\du}}
\pgfpathlineto{\pgfpoint{-9.310274\du}{7.316213\du}}
\pgfpathlineto{\pgfpoint{-11.308700\du}{7.316213\du}}
\pgfpathlineto{\pgfpoint{-11.308700\du}{6.946880\du}}
\pgfusepath{stroke}
\pgfsetbuttcap
\pgfsetmiterjoin
\pgfsetdash{}{0pt}
\definecolor{dialinecolor}{rgb}{0.788235, 0.788235, 0.713726}
\pgfsetfillcolor{dialinecolor}
\pgfpathmoveto{\pgfpoint{-11.308700\du}{6.946880\du}}
\pgfpathlineto{\pgfpoint{-11.096798\du}{6.745998\du}}
\pgfpathlineto{\pgfpoint{-9.098373\du}{6.745998\du}}
\pgfpathlineto{\pgfpoint{-9.310274\du}{6.946880\du}}
\pgfpathlineto{\pgfpoint{-11.308700\du}{6.946880\du}}
\pgfusepath{fill}
\pgfsetbuttcap
\pgfsetmiterjoin
\pgfsetdash{}{0pt}
\definecolor{dialinecolor}{rgb}{0.286275, 0.286275, 0.211765}
\pgfsetstrokecolor{dialinecolor}
\pgfpathmoveto{\pgfpoint{-11.308700\du}{6.946880\du}}
\pgfpathlineto{\pgfpoint{-11.096798\du}{6.745998\du}}
\pgfpathlineto{\pgfpoint{-9.098373\du}{6.745998\du}}
\pgfpathlineto{\pgfpoint{-9.310274\du}{6.946880\du}}
\pgfpathlineto{\pgfpoint{-11.308700\du}{6.946880\du}}
\pgfusepath{stroke}
\pgfsetlinewidth{0.106000\du}
\pgfsetbuttcap
\pgfsetmiterjoin
\pgfsetdash{}{0pt}
\definecolor{dialinecolor}{rgb}{0.000000, 0.000000, 0.000000}
\pgfsetstrokecolor{dialinecolor}
\pgfpathmoveto{\pgfpoint{-9.422050\du}{7.115016\du}}
\pgfpathlineto{\pgfpoint{-9.901584\du}{7.115016\du}}
\pgfusepath{stroke}
\pgfsetlinewidth{0.001000\du}
\pgfsetbuttcap
\pgfsetmiterjoin
\pgfsetdash{}{0pt}
\definecolor{dialinecolor}{rgb}{0.478431, 0.478431, 0.352941}
\pgfsetfillcolor{dialinecolor}
\pgfpathmoveto{\pgfpoint{-9.310274\du}{7.316213\du}}
\pgfpathlineto{\pgfpoint{-9.098373\du}{7.103681\du}}
\pgfpathlineto{\pgfpoint{-9.098373\du}{6.745998\du}}
\pgfpathlineto{\pgfpoint{-9.310274\du}{6.946880\du}}
\pgfpathlineto{\pgfpoint{-9.310274\du}{7.316213\du}}
\pgfusepath{fill}
\pgfsetbuttcap
\pgfsetmiterjoin
\pgfsetdash{}{0pt}
\definecolor{dialinecolor}{rgb}{0.286275, 0.286275, 0.211765}
\pgfsetstrokecolor{dialinecolor}
\pgfpathmoveto{\pgfpoint{-9.310274\du}{7.316213\du}}
\pgfpathlineto{\pgfpoint{-9.098373\du}{7.103681\du}}
\pgfpathlineto{\pgfpoint{-9.098373\du}{6.745998\du}}
\pgfpathlineto{\pgfpoint{-9.310274\du}{6.946880\du}}
\pgfpathlineto{\pgfpoint{-9.310274\du}{7.316213\du}}
\pgfusepath{stroke}
\pgfsetbuttcap
\pgfsetmiterjoin
\pgfsetdash{}{0pt}
\definecolor{dialinecolor}{rgb}{0.788235, 0.788235, 0.713726}
\pgfsetfillcolor{dialinecolor}
\pgfpathmoveto{\pgfpoint{-11.297365\du}{7.550470\du}}
\pgfpathlineto{\pgfpoint{-11.074443\du}{7.271187\du}}
\pgfpathlineto{\pgfpoint{-9.533511\du}{7.271187\du}}
\pgfpathlineto{\pgfpoint{-9.756433\du}{7.550470\du}}
\pgfpathlineto{\pgfpoint{-11.297365\du}{7.550470\du}}
\pgfusepath{fill}
\pgfsetbuttcap
\pgfsetmiterjoin
\pgfsetdash{}{0pt}
\definecolor{dialinecolor}{rgb}{0.286275, 0.286275, 0.211765}
\pgfsetstrokecolor{dialinecolor}
\pgfpathmoveto{\pgfpoint{-11.297365\du}{7.550470\du}}
\pgfpathlineto{\pgfpoint{-11.074443\du}{7.271187\du}}
\pgfpathlineto{\pgfpoint{-9.533511\du}{7.271187\du}}
\pgfpathlineto{\pgfpoint{-9.756433\du}{7.550470\du}}
\pgfpathlineto{\pgfpoint{-11.297365\du}{7.550470\du}}
\pgfusepath{stroke}
\pgfsetbuttcap
\pgfsetmiterjoin
\pgfsetdash{}{0pt}
\definecolor{dialinecolor}{rgb}{0.478431, 0.478431, 0.352941}
\pgfsetfillcolor{dialinecolor}
\pgfpathmoveto{\pgfpoint{-9.756433\du}{7.606200\du}}
\pgfpathlineto{\pgfpoint{-9.533511\du}{7.371943\du}}
\pgfpathlineto{\pgfpoint{-9.533511\du}{7.271187\du}}
\pgfpathlineto{\pgfpoint{-9.756433\du}{7.550470\du}}
\pgfpathlineto{\pgfpoint{-9.756433\du}{7.606200\du}}
\pgfusepath{fill}
\pgfsetbuttcap
\pgfsetmiterjoin
\pgfsetdash{}{0pt}
\definecolor{dialinecolor}{rgb}{0.286275, 0.286275, 0.211765}
\pgfsetstrokecolor{dialinecolor}
\pgfpathmoveto{\pgfpoint{-9.756433\du}{7.606200\du}}
\pgfpathlineto{\pgfpoint{-9.533511\du}{7.371943\du}}
\pgfpathlineto{\pgfpoint{-9.533511\du}{7.271187\du}}
\pgfpathlineto{\pgfpoint{-9.756433\du}{7.550470\du}}
\pgfpathlineto{\pgfpoint{-9.756433\du}{7.606200\du}}
\pgfusepath{stroke}
\pgfsetbuttcap
\pgfsetmiterjoin
\pgfsetdash{}{0pt}
\definecolor{dialinecolor}{rgb}{0.717647, 0.717647, 0.615686}
\pgfsetfillcolor{dialinecolor}
\pgfpathmoveto{\pgfpoint{-11.297365\du}{7.550470\du}}
\pgfpathlineto{\pgfpoint{-9.756433\du}{7.550470\du}}
\pgfpathlineto{\pgfpoint{-9.756433\du}{7.606200\du}}
\pgfpathlineto{\pgfpoint{-11.297365\du}{7.606200\du}}
\pgfpathlineto{\pgfpoint{-11.297365\du}{7.550470\du}}
\pgfusepath{fill}
\pgfsetbuttcap
\pgfsetmiterjoin
\pgfsetdash{}{0pt}
\definecolor{dialinecolor}{rgb}{0.286275, 0.286275, 0.211765}
\pgfsetstrokecolor{dialinecolor}
\pgfpathmoveto{\pgfpoint{-11.297365\du}{7.550470\du}}
\pgfpathlineto{\pgfpoint{-9.756433\du}{7.550470\du}}
\pgfpathlineto{\pgfpoint{-9.756433\du}{7.606200\du}}
\pgfpathlineto{\pgfpoint{-11.297365\du}{7.606200\du}}
\pgfpathlineto{\pgfpoint{-11.297365\du}{7.550470\du}}
\pgfusepath{stroke}
\pgfsetbuttcap
\pgfsetmiterjoin
\pgfsetdash{}{0pt}
\definecolor{dialinecolor}{rgb}{0.000000, 0.000000, 0.000000}
\pgfsetfillcolor{dialinecolor}
\pgfpathmoveto{\pgfpoint{-11.007378\du}{6.902485\du}}
\pgfpathlineto{\pgfpoint{-10.839556\du}{6.745998\du}}
\pgfpathlineto{\pgfpoint{-9.422050\du}{6.745998\du}}
\pgfpathlineto{\pgfpoint{-9.577907\du}{6.902485\du}}
\pgfpathlineto{\pgfpoint{-11.007378\du}{6.902485\du}}
\pgfusepath{fill}
\pgfsetbuttcap
\pgfsetmiterjoin
\pgfsetdash{}{0pt}
\definecolor{dialinecolor}{rgb}{0.000000, 0.000000, 0.000000}
\pgfsetstrokecolor{dialinecolor}
\pgfpathmoveto{\pgfpoint{-11.007378\du}{6.902485\du}}
\pgfpathlineto{\pgfpoint{-10.839556\du}{6.745998\du}}
\pgfpathlineto{\pgfpoint{-9.422050\du}{6.745998\du}}
\pgfpathlineto{\pgfpoint{-9.577907\du}{6.902485\du}}
\pgfpathlineto{\pgfpoint{-11.007378\du}{6.902485\du}}
\pgfusepath{stroke}
\pgfsetbuttcap
\pgfsetmiterjoin
\pgfsetdash{}{0pt}
\definecolor{dialinecolor}{rgb}{0.788235, 0.788235, 0.713726}
\pgfsetfillcolor{dialinecolor}
\pgfpathmoveto{\pgfpoint{-11.018713\du}{5.751351\du}}
\pgfpathlineto{\pgfpoint{-10.861912\du}{5.606200\du}}
\pgfpathlineto{\pgfpoint{-9.443776\du}{5.606200\du}}
\pgfpathlineto{\pgfpoint{-9.600577\du}{5.751351\du}}
\pgfpathlineto{\pgfpoint{-11.018713\du}{5.751351\du}}
\pgfusepath{fill}
\pgfsetbuttcap
\pgfsetmiterjoin
\pgfsetdash{}{0pt}
\definecolor{dialinecolor}{rgb}{0.286275, 0.286275, 0.211765}
\pgfsetstrokecolor{dialinecolor}
\pgfpathmoveto{\pgfpoint{-11.018713\du}{5.751351\du}}
\pgfpathlineto{\pgfpoint{-10.861912\du}{5.606200\du}}
\pgfpathlineto{\pgfpoint{-9.443776\du}{5.606200\du}}
\pgfpathlineto{\pgfpoint{-9.600577\du}{5.751351\du}}
\pgfpathlineto{\pgfpoint{-11.018713\du}{5.751351\du}}
\pgfusepath{stroke}
\pgfsetbuttcap
\pgfsetmiterjoin
\pgfsetdash{}{0pt}
\definecolor{dialinecolor}{rgb}{0.717647, 0.717647, 0.615686}
\pgfsetfillcolor{dialinecolor}
\pgfpathmoveto{\pgfpoint{-11.018713\du}{5.751351\du}}
\pgfpathlineto{\pgfpoint{-9.589242\du}{5.751351\du}}
\pgfpathlineto{\pgfpoint{-9.589242\du}{6.879815\du}}
\pgfpathlineto{\pgfpoint{-11.018713\du}{6.879815\du}}
\pgfpathlineto{\pgfpoint{-11.018713\du}{5.751351\du}}
\pgfusepath{fill}
\pgfsetbuttcap
\pgfsetmiterjoin
\pgfsetdash{}{0pt}
\definecolor{dialinecolor}{rgb}{0.286275, 0.286275, 0.211765}
\pgfsetstrokecolor{dialinecolor}
\pgfpathmoveto{\pgfpoint{-11.018713\du}{5.751351\du}}
\pgfpathlineto{\pgfpoint{-9.589871\du}{5.751351\du}}
\pgfpathlineto{\pgfpoint{-9.589871\du}{6.879500\du}}
\pgfpathlineto{\pgfpoint{-11.018713\du}{6.879500\du}}
\pgfpathlineto{\pgfpoint{-11.018713\du}{5.751351\du}}
\pgfusepath{stroke}
\pgfsetbuttcap
\pgfsetmiterjoin
\pgfsetdash{}{0pt}
\definecolor{dialinecolor}{rgb}{1.000000, 1.000000, 1.000000}
\pgfsetfillcolor{dialinecolor}
\pgfpathmoveto{\pgfpoint{-10.895917\du}{5.896187\du}}
\pgfpathlineto{\pgfpoint{-9.712352\du}{5.896187\du}}
\pgfpathlineto{\pgfpoint{-9.712352\du}{6.768039\du}}
\pgfpathlineto{\pgfpoint{-10.895917\du}{6.768039\du}}
\pgfpathlineto{\pgfpoint{-10.895917\du}{5.896187\du}}
\pgfusepath{fill}
\pgfsetbuttcap
\pgfsetmiterjoin
\pgfsetdash{}{0pt}
\definecolor{dialinecolor}{rgb}{0.286275, 0.286275, 0.211765}
\pgfsetstrokecolor{dialinecolor}
\pgfpathmoveto{\pgfpoint{-10.895917\du}{5.896187\du}}
\pgfpathlineto{\pgfpoint{-9.712352\du}{5.896187\du}}
\pgfpathlineto{\pgfpoint{-9.712352\du}{6.767724\du}}
\pgfpathlineto{\pgfpoint{-10.895917\du}{6.767724\du}}
\pgfpathlineto{\pgfpoint{-10.895917\du}{5.896187\du}}
\pgfusepath{stroke}
\pgfsetbuttcap
\pgfsetmiterjoin
\pgfsetdash{}{0pt}
\definecolor{dialinecolor}{rgb}{0.478431, 0.478431, 0.352941}
\pgfsetfillcolor{dialinecolor}
\pgfpathmoveto{\pgfpoint{-9.600577\du}{6.869109\du}}
\pgfpathlineto{\pgfpoint{-9.443776\du}{6.712623\du}}
\pgfpathlineto{\pgfpoint{-9.443776\du}{5.606200\du}}
\pgfpathlineto{\pgfpoint{-9.600577\du}{5.751351\du}}
\pgfpathlineto{\pgfpoint{-9.600577\du}{6.869109\du}}
\pgfusepath{fill}
\pgfsetbuttcap
\pgfsetmiterjoin
\pgfsetdash{}{0pt}
\definecolor{dialinecolor}{rgb}{0.286275, 0.286275, 0.211765}
\pgfsetstrokecolor{dialinecolor}
\pgfpathmoveto{\pgfpoint{-9.600577\du}{6.869109\du}}
\pgfpathlineto{\pgfpoint{-9.443776\du}{6.712623\du}}
\pgfpathlineto{\pgfpoint{-9.443776\du}{5.606200\du}}
\pgfpathlineto{\pgfpoint{-9.600577\du}{5.751351\du}}
\pgfpathlineto{\pgfpoint{-9.600577\du}{6.869109\du}}
\pgfusepath{stroke}
\pgfsetlinewidth{0.100000\du}
\pgfsetdash{}{0pt}
\pgfsetdash{}{0pt}
\pgfsetbuttcap
\pgfsetmiterjoin
\pgfsetlinewidth{0.001000\du}
\pgfsetbuttcap
\pgfsetmiterjoin
\pgfsetdash{}{0pt}
\definecolor{dialinecolor}{rgb}{0.717647, 0.717647, 0.615686}
\pgfsetfillcolor{dialinecolor}
\pgfpathmoveto{\pgfpoint{-14.363700\du}{20.962280\du}}
\pgfpathlineto{\pgfpoint{-12.365274\du}{20.962280\du}}
\pgfpathlineto{\pgfpoint{-12.365274\du}{21.331613\du}}
\pgfpathlineto{\pgfpoint{-14.363700\du}{21.331613\du}}
\pgfpathlineto{\pgfpoint{-14.363700\du}{20.962280\du}}
\pgfusepath{fill}
\pgfsetbuttcap
\pgfsetmiterjoin
\pgfsetdash{}{0pt}
\definecolor{dialinecolor}{rgb}{0.286275, 0.286275, 0.211765}
\pgfsetstrokecolor{dialinecolor}
\pgfpathmoveto{\pgfpoint{-14.363700\du}{20.962280\du}}
\pgfpathlineto{\pgfpoint{-12.365274\du}{20.962280\du}}
\pgfpathlineto{\pgfpoint{-12.365274\du}{21.331613\du}}
\pgfpathlineto{\pgfpoint{-14.363700\du}{21.331613\du}}
\pgfpathlineto{\pgfpoint{-14.363700\du}{20.962280\du}}
\pgfusepath{stroke}
\pgfsetbuttcap
\pgfsetmiterjoin
\pgfsetdash{}{0pt}
\definecolor{dialinecolor}{rgb}{0.788235, 0.788235, 0.713726}
\pgfsetfillcolor{dialinecolor}
\pgfpathmoveto{\pgfpoint{-14.363700\du}{20.962280\du}}
\pgfpathlineto{\pgfpoint{-14.151798\du}{20.761398\du}}
\pgfpathlineto{\pgfpoint{-12.153373\du}{20.761398\du}}
\pgfpathlineto{\pgfpoint{-12.365274\du}{20.962280\du}}
\pgfpathlineto{\pgfpoint{-14.363700\du}{20.962280\du}}
\pgfusepath{fill}
\pgfsetbuttcap
\pgfsetmiterjoin
\pgfsetdash{}{0pt}
\definecolor{dialinecolor}{rgb}{0.286275, 0.286275, 0.211765}
\pgfsetstrokecolor{dialinecolor}
\pgfpathmoveto{\pgfpoint{-14.363700\du}{20.962280\du}}
\pgfpathlineto{\pgfpoint{-14.151798\du}{20.761398\du}}
\pgfpathlineto{\pgfpoint{-12.153373\du}{20.761398\du}}
\pgfpathlineto{\pgfpoint{-12.365274\du}{20.962280\du}}
\pgfpathlineto{\pgfpoint{-14.363700\du}{20.962280\du}}
\pgfusepath{stroke}
\pgfsetlinewidth{0.106000\du}
\pgfsetbuttcap
\pgfsetmiterjoin
\pgfsetdash{}{0pt}
\definecolor{dialinecolor}{rgb}{0.000000, 0.000000, 0.000000}
\pgfsetstrokecolor{dialinecolor}
\pgfpathmoveto{\pgfpoint{-12.477050\du}{21.130416\du}}
\pgfpathlineto{\pgfpoint{-12.956584\du}{21.130416\du}}
\pgfusepath{stroke}
\pgfsetlinewidth{0.001000\du}
\pgfsetbuttcap
\pgfsetmiterjoin
\pgfsetdash{}{0pt}
\definecolor{dialinecolor}{rgb}{0.478431, 0.478431, 0.352941}
\pgfsetfillcolor{dialinecolor}
\pgfpathmoveto{\pgfpoint{-12.365274\du}{21.331613\du}}
\pgfpathlineto{\pgfpoint{-12.153373\du}{21.119081\du}}
\pgfpathlineto{\pgfpoint{-12.153373\du}{20.761398\du}}
\pgfpathlineto{\pgfpoint{-12.365274\du}{20.962280\du}}
\pgfpathlineto{\pgfpoint{-12.365274\du}{21.331613\du}}
\pgfusepath{fill}
\pgfsetbuttcap
\pgfsetmiterjoin
\pgfsetdash{}{0pt}
\definecolor{dialinecolor}{rgb}{0.286275, 0.286275, 0.211765}
\pgfsetstrokecolor{dialinecolor}
\pgfpathmoveto{\pgfpoint{-12.365274\du}{21.331613\du}}
\pgfpathlineto{\pgfpoint{-12.153373\du}{21.119081\du}}
\pgfpathlineto{\pgfpoint{-12.153373\du}{20.761398\du}}
\pgfpathlineto{\pgfpoint{-12.365274\du}{20.962280\du}}
\pgfpathlineto{\pgfpoint{-12.365274\du}{21.331613\du}}
\pgfusepath{stroke}
\pgfsetbuttcap
\pgfsetmiterjoin
\pgfsetdash{}{0pt}
\definecolor{dialinecolor}{rgb}{0.788235, 0.788235, 0.713726}
\pgfsetfillcolor{dialinecolor}
\pgfpathmoveto{\pgfpoint{-14.352365\du}{21.565870\du}}
\pgfpathlineto{\pgfpoint{-14.129443\du}{21.286587\du}}
\pgfpathlineto{\pgfpoint{-12.588511\du}{21.286587\du}}
\pgfpathlineto{\pgfpoint{-12.811433\du}{21.565870\du}}
\pgfpathlineto{\pgfpoint{-14.352365\du}{21.565870\du}}
\pgfusepath{fill}
\pgfsetbuttcap
\pgfsetmiterjoin
\pgfsetdash{}{0pt}
\definecolor{dialinecolor}{rgb}{0.286275, 0.286275, 0.211765}
\pgfsetstrokecolor{dialinecolor}
\pgfpathmoveto{\pgfpoint{-14.352365\du}{21.565870\du}}
\pgfpathlineto{\pgfpoint{-14.129443\du}{21.286587\du}}
\pgfpathlineto{\pgfpoint{-12.588511\du}{21.286587\du}}
\pgfpathlineto{\pgfpoint{-12.811433\du}{21.565870\du}}
\pgfpathlineto{\pgfpoint{-14.352365\du}{21.565870\du}}
\pgfusepath{stroke}
\pgfsetbuttcap
\pgfsetmiterjoin
\pgfsetdash{}{0pt}
\definecolor{dialinecolor}{rgb}{0.478431, 0.478431, 0.352941}
\pgfsetfillcolor{dialinecolor}
\pgfpathmoveto{\pgfpoint{-12.811433\du}{21.621600\du}}
\pgfpathlineto{\pgfpoint{-12.588511\du}{21.387343\du}}
\pgfpathlineto{\pgfpoint{-12.588511\du}{21.286587\du}}
\pgfpathlineto{\pgfpoint{-12.811433\du}{21.565870\du}}
\pgfpathlineto{\pgfpoint{-12.811433\du}{21.621600\du}}
\pgfusepath{fill}
\pgfsetbuttcap
\pgfsetmiterjoin
\pgfsetdash{}{0pt}
\definecolor{dialinecolor}{rgb}{0.286275, 0.286275, 0.211765}
\pgfsetstrokecolor{dialinecolor}
\pgfpathmoveto{\pgfpoint{-12.811433\du}{21.621600\du}}
\pgfpathlineto{\pgfpoint{-12.588511\du}{21.387343\du}}
\pgfpathlineto{\pgfpoint{-12.588511\du}{21.286587\du}}
\pgfpathlineto{\pgfpoint{-12.811433\du}{21.565870\du}}
\pgfpathlineto{\pgfpoint{-12.811433\du}{21.621600\du}}
\pgfusepath{stroke}
\pgfsetbuttcap
\pgfsetmiterjoin
\pgfsetdash{}{0pt}
\definecolor{dialinecolor}{rgb}{0.717647, 0.717647, 0.615686}
\pgfsetfillcolor{dialinecolor}
\pgfpathmoveto{\pgfpoint{-14.352365\du}{21.565870\du}}
\pgfpathlineto{\pgfpoint{-12.811433\du}{21.565870\du}}
\pgfpathlineto{\pgfpoint{-12.811433\du}{21.621600\du}}
\pgfpathlineto{\pgfpoint{-14.352365\du}{21.621600\du}}
\pgfpathlineto{\pgfpoint{-14.352365\du}{21.565870\du}}
\pgfusepath{fill}
\pgfsetbuttcap
\pgfsetmiterjoin
\pgfsetdash{}{0pt}
\definecolor{dialinecolor}{rgb}{0.286275, 0.286275, 0.211765}
\pgfsetstrokecolor{dialinecolor}
\pgfpathmoveto{\pgfpoint{-14.352365\du}{21.565870\du}}
\pgfpathlineto{\pgfpoint{-12.811433\du}{21.565870\du}}
\pgfpathlineto{\pgfpoint{-12.811433\du}{21.621600\du}}
\pgfpathlineto{\pgfpoint{-14.352365\du}{21.621600\du}}
\pgfpathlineto{\pgfpoint{-14.352365\du}{21.565870\du}}
\pgfusepath{stroke}
\pgfsetbuttcap
\pgfsetmiterjoin
\pgfsetdash{}{0pt}
\definecolor{dialinecolor}{rgb}{0.000000, 0.000000, 0.000000}
\pgfsetfillcolor{dialinecolor}
\pgfpathmoveto{\pgfpoint{-14.062378\du}{20.917885\du}}
\pgfpathlineto{\pgfpoint{-13.894556\du}{20.761398\du}}
\pgfpathlineto{\pgfpoint{-12.477050\du}{20.761398\du}}
\pgfpathlineto{\pgfpoint{-12.632907\du}{20.917885\du}}
\pgfpathlineto{\pgfpoint{-14.062378\du}{20.917885\du}}
\pgfusepath{fill}
\pgfsetbuttcap
\pgfsetmiterjoin
\pgfsetdash{}{0pt}
\definecolor{dialinecolor}{rgb}{0.000000, 0.000000, 0.000000}
\pgfsetstrokecolor{dialinecolor}
\pgfpathmoveto{\pgfpoint{-14.062378\du}{20.917885\du}}
\pgfpathlineto{\pgfpoint{-13.894556\du}{20.761398\du}}
\pgfpathlineto{\pgfpoint{-12.477050\du}{20.761398\du}}
\pgfpathlineto{\pgfpoint{-12.632907\du}{20.917885\du}}
\pgfpathlineto{\pgfpoint{-14.062378\du}{20.917885\du}}
\pgfusepath{stroke}
\pgfsetbuttcap
\pgfsetmiterjoin
\pgfsetdash{}{0pt}
\definecolor{dialinecolor}{rgb}{0.788235, 0.788235, 0.713726}
\pgfsetfillcolor{dialinecolor}
\pgfpathmoveto{\pgfpoint{-14.073713\du}{19.766751\du}}
\pgfpathlineto{\pgfpoint{-13.916912\du}{19.621600\du}}
\pgfpathlineto{\pgfpoint{-12.498776\du}{19.621600\du}}
\pgfpathlineto{\pgfpoint{-12.655577\du}{19.766751\du}}
\pgfpathlineto{\pgfpoint{-14.073713\du}{19.766751\du}}
\pgfusepath{fill}
\pgfsetbuttcap
\pgfsetmiterjoin
\pgfsetdash{}{0pt}
\definecolor{dialinecolor}{rgb}{0.286275, 0.286275, 0.211765}
\pgfsetstrokecolor{dialinecolor}
\pgfpathmoveto{\pgfpoint{-14.073713\du}{19.766751\du}}
\pgfpathlineto{\pgfpoint{-13.916912\du}{19.621600\du}}
\pgfpathlineto{\pgfpoint{-12.498776\du}{19.621600\du}}
\pgfpathlineto{\pgfpoint{-12.655577\du}{19.766751\du}}
\pgfpathlineto{\pgfpoint{-14.073713\du}{19.766751\du}}
\pgfusepath{stroke}
\pgfsetbuttcap
\pgfsetmiterjoin
\pgfsetdash{}{0pt}
\definecolor{dialinecolor}{rgb}{0.717647, 0.717647, 0.615686}
\pgfsetfillcolor{dialinecolor}
\pgfpathmoveto{\pgfpoint{-14.073713\du}{19.766751\du}}
\pgfpathlineto{\pgfpoint{-12.644242\du}{19.766751\du}}
\pgfpathlineto{\pgfpoint{-12.644242\du}{20.895215\du}}
\pgfpathlineto{\pgfpoint{-14.073713\du}{20.895215\du}}
\pgfpathlineto{\pgfpoint{-14.073713\du}{19.766751\du}}
\pgfusepath{fill}
\pgfsetbuttcap
\pgfsetmiterjoin
\pgfsetdash{}{0pt}
\definecolor{dialinecolor}{rgb}{0.286275, 0.286275, 0.211765}
\pgfsetstrokecolor{dialinecolor}
\pgfpathmoveto{\pgfpoint{-14.073713\du}{19.766751\du}}
\pgfpathlineto{\pgfpoint{-12.644871\du}{19.766751\du}}
\pgfpathlineto{\pgfpoint{-12.644871\du}{20.894900\du}}
\pgfpathlineto{\pgfpoint{-14.073713\du}{20.894900\du}}
\pgfpathlineto{\pgfpoint{-14.073713\du}{19.766751\du}}
\pgfusepath{stroke}
\pgfsetbuttcap
\pgfsetmiterjoin
\pgfsetdash{}{0pt}
\definecolor{dialinecolor}{rgb}{1.000000, 1.000000, 1.000000}
\pgfsetfillcolor{dialinecolor}
\pgfpathmoveto{\pgfpoint{-13.950917\du}{19.911587\du}}
\pgfpathlineto{\pgfpoint{-12.767352\du}{19.911587\du}}
\pgfpathlineto{\pgfpoint{-12.767352\du}{20.783439\du}}
\pgfpathlineto{\pgfpoint{-13.950917\du}{20.783439\du}}
\pgfpathlineto{\pgfpoint{-13.950917\du}{19.911587\du}}
\pgfusepath{fill}
\pgfsetbuttcap
\pgfsetmiterjoin
\pgfsetdash{}{0pt}
\definecolor{dialinecolor}{rgb}{0.286275, 0.286275, 0.211765}
\pgfsetstrokecolor{dialinecolor}
\pgfpathmoveto{\pgfpoint{-13.950917\du}{19.911587\du}}
\pgfpathlineto{\pgfpoint{-12.767352\du}{19.911587\du}}
\pgfpathlineto{\pgfpoint{-12.767352\du}{20.783124\du}}
\pgfpathlineto{\pgfpoint{-13.950917\du}{20.783124\du}}
\pgfpathlineto{\pgfpoint{-13.950917\du}{19.911587\du}}
\pgfusepath{stroke}
\pgfsetbuttcap
\pgfsetmiterjoin
\pgfsetdash{}{0pt}
\definecolor{dialinecolor}{rgb}{0.478431, 0.478431, 0.352941}
\pgfsetfillcolor{dialinecolor}
\pgfpathmoveto{\pgfpoint{-12.655577\du}{20.884509\du}}
\pgfpathlineto{\pgfpoint{-12.498776\du}{20.728023\du}}
\pgfpathlineto{\pgfpoint{-12.498776\du}{19.621600\du}}
\pgfpathlineto{\pgfpoint{-12.655577\du}{19.766751\du}}
\pgfpathlineto{\pgfpoint{-12.655577\du}{20.884509\du}}
\pgfusepath{fill}
\pgfsetbuttcap
\pgfsetmiterjoin
\pgfsetdash{}{0pt}
\definecolor{dialinecolor}{rgb}{0.286275, 0.286275, 0.211765}
\pgfsetstrokecolor{dialinecolor}
\pgfpathmoveto{\pgfpoint{-12.655577\du}{20.884509\du}}
\pgfpathlineto{\pgfpoint{-12.498776\du}{20.728023\du}}
\pgfpathlineto{\pgfpoint{-12.498776\du}{19.621600\du}}
\pgfpathlineto{\pgfpoint{-12.655577\du}{19.766751\du}}
\pgfpathlineto{\pgfpoint{-12.655577\du}{20.884509\du}}
\pgfusepath{stroke}
\pgfsetlinewidth{0.100000\du}
\pgfsetdash{}{0pt}
\pgfsetdash{}{0pt}
\pgfsetbuttcap
\pgfsetmiterjoin
\pgfsetlinewidth{0.001000\du}
\pgfsetbuttcap
\pgfsetmiterjoin
\pgfsetdash{}{0pt}
\definecolor{dialinecolor}{rgb}{0.000000, 0.470588, 0.666667}
\pgfsetstrokecolor{dialinecolor}
\pgfpathmoveto{\pgfpoint{-17.980683\du}{7.553754\du}}
\pgfpathlineto{\pgfpoint{-17.974115\du}{7.544797\du}}
\pgfpathlineto{\pgfpoint{-17.967845\du}{7.535244\du}}
\pgfpathlineto{\pgfpoint{-17.962471\du}{7.525093\du}}
\pgfpathlineto{\pgfpoint{-17.957396\du}{7.515838\du}}
\pgfpathlineto{\pgfpoint{-17.953515\du}{7.505687\du}}
\pgfpathlineto{\pgfpoint{-17.950231\du}{7.495536\du}}
\pgfpathlineto{\pgfpoint{-17.947245\du}{7.485386\du}}
\pgfpathlineto{\pgfpoint{-17.945454\du}{7.474936\du}}
\pgfpathlineto{\pgfpoint{-17.944260\du}{7.464785\du}}
\pgfpathlineto{\pgfpoint{-17.943364\du}{7.454336\du}}
\pgfpathlineto{\pgfpoint{-17.943364\du}{7.443887\du}}
\pgfpathlineto{\pgfpoint{-17.944260\du}{7.433139\du}}
\pgfpathlineto{\pgfpoint{-17.945454\du}{7.422988\du}}
\pgfpathlineto{\pgfpoint{-17.947245\du}{7.412539\du}}
\pgfpathlineto{\pgfpoint{-17.949932\du}{7.402687\du}}
\pgfpathlineto{\pgfpoint{-17.953515\du}{7.392237\du}}
\pgfpathlineto{\pgfpoint{-17.957396\du}{7.382385\du}}
\pgfpathlineto{\pgfpoint{-17.962471\du}{7.372234\du}}
\pgfpathlineto{\pgfpoint{-17.967845\du}{7.362681\du}}
\pgfpathlineto{\pgfpoint{-17.974115\du}{7.353724\du}}
\pgfpathlineto{\pgfpoint{-17.980683\du}{7.344170\du}}
\pgfpathlineto{\pgfpoint{-17.988147\du}{7.335512\du}}
\pgfpathlineto{\pgfpoint{-17.995611\du}{7.326854\du}}
\pgfpathlineto{\pgfpoint{-18.003970\du}{7.318495\du}}
\pgfpathlineto{\pgfpoint{-18.012628\du}{7.310434\du}}
\pgfpathlineto{\pgfpoint{-18.021883\du}{7.302672\du}}
\pgfpathlineto{\pgfpoint{-18.031437\du}{7.295506\du}}
\pgfpathlineto{\pgfpoint{-18.042185\du}{7.288938\du}}
\pgfpathlineto{\pgfpoint{-18.052634\du}{7.282370\du}}
\pgfpathlineto{\pgfpoint{-18.063979\du}{7.276101\du}}
\pgfpathlineto{\pgfpoint{-18.075921\du}{7.270428\du}}
\pgfpathlineto{\pgfpoint{-18.087863\du}{7.265054\du}}
\pgfpathlineto{\pgfpoint{-18.100104\du}{7.260874\du}}
\pgfpathlineto{\pgfpoint{-18.112643\du}{7.256695\du}}
\pgfpathlineto{\pgfpoint{-18.125779\du}{7.252515\du}}
\pgfpathlineto{\pgfpoint{-18.138617\du}{7.249529\du}}
\pgfpathlineto{\pgfpoint{-18.151753\du}{7.246842\du}}
\pgfpathlineto{\pgfpoint{-18.165487\du}{7.244454\du}}
\pgfpathlineto{\pgfpoint{-18.178623\du}{7.242961\du}}
\pgfpathlineto{\pgfpoint{-18.192655\du}{7.241767\du}}
\pgfpathlineto{\pgfpoint{-18.206388\du}{7.240871\du}}
\pgfpathlineto{\pgfpoint{-18.220420\du}{7.240871\du}}
\pgfpathlineto{\pgfpoint{-18.233855\du}{7.241469\du}}
\pgfpathlineto{\pgfpoint{-18.247589\du}{7.242066\du}}
\pgfpathlineto{\pgfpoint{-18.261621\du}{7.243558\du}}
\pgfpathlineto{\pgfpoint{-18.274757\du}{7.245947\du}}
\pgfpathlineto{\pgfpoint{-18.288490\du}{7.248335\du}}
\pgfpathlineto{\pgfpoint{-18.301029\du}{7.251321\du}}
\pgfpathlineto{\pgfpoint{-18.314166\du}{7.254306\du}}
\pgfpathlineto{\pgfpoint{-18.326705\du}{7.258486\du}}
\pgfpathlineto{\pgfpoint{-18.339543\du}{7.263263\du}}
\pgfpathlineto{\pgfpoint{-18.351783\du}{7.267741\du}}
\pgfpathlineto{\pgfpoint{-18.363128\du}{7.273115\du}}
\pgfpathlineto{\pgfpoint{-18.374772\du}{7.279385\du}}
\pgfpathlineto{\pgfpoint{-18.385818\du}{7.285356\du}}
\pgfpathlineto{\pgfpoint{-18.396268\du}{7.292222\du}}
\pgfpathlineto{\pgfpoint{-18.406418\du}{7.298791\du}}
\pgfpathlineto{\pgfpoint{-18.415972\du}{7.306553\du}}
\pgfpathlineto{\pgfpoint{-18.425526\du}{7.314614\du}}
\pgfpathlineto{\pgfpoint{-18.433885\du}{7.322675\du}}
\pgfpathlineto{\pgfpoint{-18.441946\du}{7.331034\du}}
\pgfpathlineto{\pgfpoint{-18.449708\du}{7.339991\du}}
\pgfusepath{stroke}
\pgfsetlinewidth{0.080000\du}
\pgfsetbuttcap
\pgfsetmiterjoin
\pgfsetdash{}{0pt}
\definecolor{dialinecolor}{rgb}{0.666667, 0.901961, 1.000000}
\pgfsetstrokecolor{dialinecolor}
\pgfpathmoveto{\pgfpoint{-17.989938\du}{7.546290\du}}
\pgfpathlineto{\pgfpoint{-17.984564\du}{7.537334\du}}
\pgfpathlineto{\pgfpoint{-17.978295\du}{7.528676\du}}
\pgfpathlineto{\pgfpoint{-17.973219\du}{7.519420\du}}
\pgfpathlineto{\pgfpoint{-17.969040\du}{7.510165\du}}
\pgfpathlineto{\pgfpoint{-17.965158\du}{7.500612\du}}
\pgfpathlineto{\pgfpoint{-17.962173\du}{7.491058\du}}
\pgfpathlineto{\pgfpoint{-17.959486\du}{7.481206\du}}
\pgfpathlineto{\pgfpoint{-17.957396\du}{7.471951\du}}
\pgfpathlineto{\pgfpoint{-17.956500\du}{7.462098\du}}
\pgfpathlineto{\pgfpoint{-17.955903\du}{7.451948\du}}
\pgfpathlineto{\pgfpoint{-17.955903\du}{7.442095\du}}
\pgfpathlineto{\pgfpoint{-17.956799\du}{7.432243\du}}
\pgfpathlineto{\pgfpoint{-17.958292\du}{7.422093\du}}
\pgfpathlineto{\pgfpoint{-17.960083\du}{7.412837\du}}
\pgfpathlineto{\pgfpoint{-17.963069\du}{7.402985\du}}
\pgfpathlineto{\pgfpoint{-17.966651\du}{7.393432\du}}
\pgfpathlineto{\pgfpoint{-17.970532\du}{7.384176\du}}
\pgfpathlineto{\pgfpoint{-17.975011\du}{7.374921\du}}
\pgfpathlineto{\pgfpoint{-17.980086\du}{7.366263\du}}
\pgfpathlineto{\pgfpoint{-17.985758\du}{7.357008\du}}
\pgfpathlineto{\pgfpoint{-17.992924\du}{7.348649\du}}
\pgfpathlineto{\pgfpoint{-17.999492\du}{7.339991\du}}
\pgfpathlineto{\pgfpoint{-18.007254\du}{7.331930\du}}
\pgfpathlineto{\pgfpoint{-18.015315\du}{7.324167\du}}
\pgfpathlineto{\pgfpoint{-18.023675\du}{7.317301\du}}
\pgfpathlineto{\pgfpoint{-18.032333\du}{7.309837\du}}
\pgfpathlineto{\pgfpoint{-18.042185\du}{7.302970\du}}
\pgfpathlineto{\pgfpoint{-18.052634\du}{7.296104\du}}
\pgfpathlineto{\pgfpoint{-18.062785\du}{7.290431\du}}
\pgfpathlineto{\pgfpoint{-18.073234\du}{7.285057\du}}
\pgfpathlineto{\pgfpoint{-18.085176\du}{7.279683\du}}
\pgfpathlineto{\pgfpoint{-18.095924\du}{7.274608\du}}
\pgfpathlineto{\pgfpoint{-18.108463\du}{7.270428\du}}
\pgfpathlineto{\pgfpoint{-18.119808\du}{7.266547\du}}
\pgfpathlineto{\pgfpoint{-18.132348\du}{7.263263\du}}
\pgfpathlineto{\pgfpoint{-18.144887\du}{7.260277\du}}
\pgfpathlineto{\pgfpoint{-18.157426\du}{7.257590\du}}
\pgfpathlineto{\pgfpoint{-18.170562\du}{7.255799\du}}
\pgfpathlineto{\pgfpoint{-18.183101\du}{7.254306\du}}
\pgfpathlineto{\pgfpoint{-18.196238\du}{7.253112\du}}
\pgfpathlineto{\pgfpoint{-18.210270\du}{7.252515\du}}
\pgfpathlineto{\pgfpoint{-18.223406\du}{7.252515\du}}
\pgfpathlineto{\pgfpoint{-18.236244\du}{7.253112\du}}
\pgfpathlineto{\pgfpoint{-18.249380\du}{7.254008\du}}
\pgfpathlineto{\pgfpoint{-18.261919\du}{7.255799\du}}
\pgfpathlineto{\pgfpoint{-18.275055\du}{7.257590\du}}
\pgfpathlineto{\pgfpoint{-18.287595\du}{7.259680\du}}
\pgfpathlineto{\pgfpoint{-18.300432\du}{7.263263\du}}
\pgfpathlineto{\pgfpoint{-18.312673\du}{7.266248\du}}
\pgfpathlineto{\pgfpoint{-18.324615\du}{7.270130\du}}
\pgfpathlineto{\pgfpoint{-18.336856\du}{7.274309\du}}
\pgfpathlineto{\pgfpoint{-18.348201\du}{7.279385\du}}
\pgfpathlineto{\pgfpoint{-18.359247\du}{7.284161\du}}
\pgfpathlineto{\pgfpoint{-18.369995\du}{7.290133\du}}
\pgfpathlineto{\pgfpoint{-18.380444\du}{7.295805\du}}
\pgfpathlineto{\pgfpoint{-18.390894\du}{7.302373\du}}
\pgfpathlineto{\pgfpoint{-18.400447\du}{7.308941\du}}
\pgfpathlineto{\pgfpoint{-18.409105\du}{7.316107\du}}
\pgfpathlineto{\pgfpoint{-18.418659\du}{7.323272\du}}
\pgfpathlineto{\pgfpoint{-18.426421\du}{7.331034\du}}
\pgfpathlineto{\pgfpoint{-18.433885\du}{7.339095\du}}
\pgfpathlineto{\pgfpoint{-18.441050\du}{7.347455\du}}
\pgfusepath{stroke}
\pgfsetbuttcap
\pgfsetmiterjoin
\pgfsetdash{}{0pt}
\definecolor{dialinecolor}{rgb}{0.000000, 0.705882, 1.000000}
\pgfsetfillcolor{dialinecolor}
\pgfpathmoveto{\pgfpoint{-19.906941\du}{7.229228\du}}
\pgfpathlineto{\pgfpoint{-19.704224\du}{7.038453\du}}
\pgfpathlineto{\pgfpoint{-18.192655\du}{7.038453\du}}
\pgfpathlineto{\pgfpoint{-18.395372\du}{7.229228\du}}
\pgfpathlineto{\pgfpoint{-19.906941\du}{7.229228\du}}
\pgfusepath{fill}
\pgfsetlinewidth{0.001000\du}
\pgfsetbuttcap
\pgfsetmiterjoin
\pgfsetdash{}{0pt}
\definecolor{dialinecolor}{rgb}{0.666667, 0.901961, 1.000000}
\pgfsetstrokecolor{dialinecolor}
\pgfpathmoveto{\pgfpoint{-19.906941\du}{7.229228\du}}
\pgfpathlineto{\pgfpoint{-19.704224\du}{7.038453\du}}
\pgfpathlineto{\pgfpoint{-18.192655\du}{7.038453\du}}
\pgfpathlineto{\pgfpoint{-18.395372\du}{7.229228\du}}
\pgfpathlineto{\pgfpoint{-19.906941\du}{7.229228\du}}
\pgfusepath{stroke}
\pgfsetbuttcap
\pgfsetmiterjoin
\pgfsetdash{}{0pt}
\definecolor{dialinecolor}{rgb}{0.000000, 0.588235, 0.831373}
\pgfsetfillcolor{dialinecolor}
\pgfpathmoveto{\pgfpoint{-19.906941\du}{7.229228\du}}
\pgfpathlineto{\pgfpoint{-18.395372\du}{7.229228\du}}
\pgfpathlineto{\pgfpoint{-18.395372\du}{7.515241\du}}
\pgfpathlineto{\pgfpoint{-19.906941\du}{7.515241\du}}
\pgfpathlineto{\pgfpoint{-19.906941\du}{7.229228\du}}
\pgfusepath{fill}
\pgfsetbuttcap
\pgfsetmiterjoin
\pgfsetdash{}{0pt}
\definecolor{dialinecolor}{rgb}{0.666667, 0.901961, 1.000000}
\pgfsetstrokecolor{dialinecolor}
\pgfpathmoveto{\pgfpoint{-19.906941\du}{7.229228\du}}
\pgfpathlineto{\pgfpoint{-18.395372\du}{7.229228\du}}
\pgfpathlineto{\pgfpoint{-18.395372\du}{7.515241\du}}
\pgfpathlineto{\pgfpoint{-19.906941\du}{7.515241\du}}
\pgfpathlineto{\pgfpoint{-19.906941\du}{7.229228\du}}
\pgfusepath{stroke}
\pgfsetbuttcap
\pgfsetmiterjoin
\pgfsetdash{}{0pt}
\definecolor{dialinecolor}{rgb}{0.000000, 0.352941, 0.501961}
\pgfsetfillcolor{dialinecolor}
\pgfpathmoveto{\pgfpoint{-18.395372\du}{7.515241\du}}
\pgfpathlineto{\pgfpoint{-18.192655\du}{7.312524\du}}
\pgfpathlineto{\pgfpoint{-18.192655\du}{7.038453\du}}
\pgfpathlineto{\pgfpoint{-18.395372\du}{7.229228\du}}
\pgfpathlineto{\pgfpoint{-18.395372\du}{7.515241\du}}
\pgfusepath{fill}
\pgfsetbuttcap
\pgfsetmiterjoin
\pgfsetdash{}{0pt}
\definecolor{dialinecolor}{rgb}{0.666667, 0.901961, 1.000000}
\pgfsetstrokecolor{dialinecolor}
\pgfpathmoveto{\pgfpoint{-18.395372\du}{7.515241\du}}
\pgfpathlineto{\pgfpoint{-18.192655\du}{7.312524\du}}
\pgfpathlineto{\pgfpoint{-18.192655\du}{7.038453\du}}
\pgfpathlineto{\pgfpoint{-18.395372\du}{7.229228\du}}
\pgfpathlineto{\pgfpoint{-18.395372\du}{7.515241\du}}
\pgfusepath{stroke}
\pgfsetbuttcap
\pgfsetmiterjoin
\pgfsetdash{}{0pt}
\definecolor{dialinecolor}{rgb}{0.000000, 0.000000, 0.000000}
\pgfsetfillcolor{dialinecolor}
\pgfpathmoveto{\pgfpoint{-19.858874\du}{7.193103\du}}
\pgfpathlineto{\pgfpoint{-19.704224\du}{7.038453\du}}
\pgfpathlineto{\pgfpoint{-18.228780\du}{7.038453\du}}
\pgfpathlineto{\pgfpoint{-18.383131\du}{7.193103\du}}
\pgfpathlineto{\pgfpoint{-19.858874\du}{7.193103\du}}
\pgfusepath{fill}
\pgfsetbuttcap
\pgfsetmiterjoin
\pgfsetdash{}{0pt}
\definecolor{dialinecolor}{rgb}{0.000000, 0.000000, 0.000000}
\pgfsetstrokecolor{dialinecolor}
\pgfpathmoveto{\pgfpoint{-19.858874\du}{7.193103\du}}
\pgfpathlineto{\pgfpoint{-19.704224\du}{7.038453\du}}
\pgfpathlineto{\pgfpoint{-18.228780\du}{7.038453\du}}
\pgfpathlineto{\pgfpoint{-18.383131\du}{7.193103\du}}
\pgfpathlineto{\pgfpoint{-19.858874\du}{7.193103\du}}
\pgfusepath{stroke}
\pgfsetbuttcap
\pgfsetmiterjoin
\pgfsetdash{}{0pt}
\definecolor{dialinecolor}{rgb}{0.000000, 0.705882, 1.000000}
\pgfsetfillcolor{dialinecolor}
\pgfpathmoveto{\pgfpoint{-19.906941\du}{5.966950\du}}
\pgfpathlineto{\pgfpoint{-19.739752\du}{5.812300\du}}
\pgfpathlineto{\pgfpoint{-18.228780\du}{5.812300\du}}
\pgfpathlineto{\pgfpoint{-18.395372\du}{5.966950\du}}
\pgfpathlineto{\pgfpoint{-19.906941\du}{5.966950\du}}
\pgfusepath{fill}
\pgfsetbuttcap
\pgfsetmiterjoin
\pgfsetdash{}{0pt}
\definecolor{dialinecolor}{rgb}{0.666667, 0.901961, 1.000000}
\pgfsetstrokecolor{dialinecolor}
\pgfpathmoveto{\pgfpoint{-19.906941\du}{5.966950\du}}
\pgfpathlineto{\pgfpoint{-19.739752\du}{5.812300\du}}
\pgfpathlineto{\pgfpoint{-18.228780\du}{5.812300\du}}
\pgfpathlineto{\pgfpoint{-18.395372\du}{5.966950\du}}
\pgfpathlineto{\pgfpoint{-19.906941\du}{5.966950\du}}
\pgfusepath{stroke}
\pgfsetbuttcap
\pgfsetmiterjoin
\pgfsetdash{}{0pt}
\definecolor{dialinecolor}{rgb}{0.000000, 0.588235, 0.831373}
\pgfsetfillcolor{dialinecolor}
\pgfpathmoveto{\pgfpoint{-19.906941\du}{5.966950\du}}
\pgfpathlineto{\pgfpoint{-18.383131\du}{5.966950\du}}
\pgfpathlineto{\pgfpoint{-18.383131\du}{7.169219\du}}
\pgfpathlineto{\pgfpoint{-19.906941\du}{7.169219\du}}
\pgfpathlineto{\pgfpoint{-19.906941\du}{5.966950\du}}
\pgfusepath{fill}
\pgfsetbuttcap
\pgfsetmiterjoin
\pgfsetdash{}{0pt}
\definecolor{dialinecolor}{rgb}{0.666667, 0.901961, 1.000000}
\pgfsetstrokecolor{dialinecolor}
\pgfpathmoveto{\pgfpoint{-19.906941\du}{5.966950\du}}
\pgfpathlineto{\pgfpoint{-18.383430\du}{5.966950\du}}
\pgfpathlineto{\pgfpoint{-18.383430\du}{7.168920\du}}
\pgfpathlineto{\pgfpoint{-19.906941\du}{7.168920\du}}
\pgfpathlineto{\pgfpoint{-19.906941\du}{5.966950\du}}
\pgfusepath{stroke}
\pgfsetbuttcap
\pgfsetmiterjoin
\pgfsetdash{}{0pt}
\definecolor{dialinecolor}{rgb}{1.000000, 1.000000, 1.000000}
\pgfsetfillcolor{dialinecolor}
\pgfpathmoveto{\pgfpoint{-19.775876\du}{6.121898\du}}
\pgfpathlineto{\pgfpoint{-18.514196\du}{6.121898\du}}
\pgfpathlineto{\pgfpoint{-18.514196\du}{7.050694\du}}
\pgfpathlineto{\pgfpoint{-19.775876\du}{7.050694\du}}
\pgfpathlineto{\pgfpoint{-19.775876\du}{6.121898\du}}
\pgfusepath{fill}
\pgfsetbuttcap
\pgfsetmiterjoin
\pgfsetdash{}{0pt}
\definecolor{dialinecolor}{rgb}{0.000000, 0.235294, 0.333333}
\pgfsetstrokecolor{dialinecolor}
\pgfpathmoveto{\pgfpoint{-19.775876\du}{6.121898\du}}
\pgfpathlineto{\pgfpoint{-18.514494\du}{6.121898\du}}
\pgfpathlineto{\pgfpoint{-18.514494\du}{7.050395\du}}
\pgfpathlineto{\pgfpoint{-19.775876\du}{7.050395\du}}
\pgfpathlineto{\pgfpoint{-19.775876\du}{6.121898\du}}
\pgfusepath{stroke}
\pgfsetbuttcap
\pgfsetmiterjoin
\pgfsetdash{}{0pt}
\definecolor{dialinecolor}{rgb}{0.000000, 0.352941, 0.501961}
\pgfsetfillcolor{dialinecolor}
\pgfpathmoveto{\pgfpoint{-18.395372\du}{7.157874\du}}
\pgfpathlineto{\pgfpoint{-18.228780\du}{6.990983\du}}
\pgfpathlineto{\pgfpoint{-18.228780\du}{5.812300\du}}
\pgfpathlineto{\pgfpoint{-18.395372\du}{5.966950\du}}
\pgfpathlineto{\pgfpoint{-18.395372\du}{7.157874\du}}
\pgfusepath{fill}
\pgfsetbuttcap
\pgfsetmiterjoin
\pgfsetdash{}{0pt}
\definecolor{dialinecolor}{rgb}{0.666667, 0.901961, 1.000000}
\pgfsetstrokecolor{dialinecolor}
\pgfpathmoveto{\pgfpoint{-18.395372\du}{7.157874\du}}
\pgfpathlineto{\pgfpoint{-18.228780\du}{6.990983\du}}
\pgfpathlineto{\pgfpoint{-18.228780\du}{5.812300\du}}
\pgfpathlineto{\pgfpoint{-18.395372\du}{5.966950\du}}
\pgfpathlineto{\pgfpoint{-18.395372\du}{7.157874\du}}
\pgfusepath{stroke}
\pgfsetbuttcap
\pgfsetmiterjoin
\pgfsetdash{}{0pt}
\definecolor{dialinecolor}{rgb}{0.000000, 0.705882, 1.000000}
\pgfsetfillcolor{dialinecolor}
\pgfpathmoveto{\pgfpoint{-20.204000\du}{7.753187\du}}
\pgfpathlineto{\pgfpoint{-19.966353\du}{7.455232\du}}
\pgfpathlineto{\pgfpoint{-18.323719\du}{7.455232\du}}
\pgfpathlineto{\pgfpoint{-18.561665\du}{7.753187\du}}
\pgfpathlineto{\pgfpoint{-20.204000\du}{7.753187\du}}
\pgfusepath{fill}
\pgfsetbuttcap
\pgfsetmiterjoin
\pgfsetdash{}{0pt}
\definecolor{dialinecolor}{rgb}{0.666667, 0.901961, 1.000000}
\pgfsetstrokecolor{dialinecolor}
\pgfpathmoveto{\pgfpoint{-20.204000\du}{7.753187\du}}
\pgfpathlineto{\pgfpoint{-19.966353\du}{7.455232\du}}
\pgfpathlineto{\pgfpoint{-18.323719\du}{7.455232\du}}
\pgfpathlineto{\pgfpoint{-18.561665\du}{7.753187\du}}
\pgfpathlineto{\pgfpoint{-20.204000\du}{7.753187\du}}
\pgfusepath{stroke}
\pgfsetbuttcap
\pgfsetmiterjoin
\pgfsetdash{}{0pt}
\definecolor{dialinecolor}{rgb}{0.000000, 0.352941, 0.501961}
\pgfsetfillcolor{dialinecolor}
\pgfpathmoveto{\pgfpoint{-18.561665\du}{7.812300\du}}
\pgfpathlineto{\pgfpoint{-18.323719\du}{7.562412\du}}
\pgfpathlineto{\pgfpoint{-18.323719\du}{7.455232\du}}
\pgfpathlineto{\pgfpoint{-18.561665\du}{7.753187\du}}
\pgfpathlineto{\pgfpoint{-18.561665\du}{7.812300\du}}
\pgfusepath{fill}
\pgfsetbuttcap
\pgfsetmiterjoin
\pgfsetdash{}{0pt}
\definecolor{dialinecolor}{rgb}{0.666667, 0.901961, 1.000000}
\pgfsetstrokecolor{dialinecolor}
\pgfpathmoveto{\pgfpoint{-18.561665\du}{7.812300\du}}
\pgfpathlineto{\pgfpoint{-18.323719\du}{7.562412\du}}
\pgfpathlineto{\pgfpoint{-18.323719\du}{7.455232\du}}
\pgfpathlineto{\pgfpoint{-18.561665\du}{7.753187\du}}
\pgfpathlineto{\pgfpoint{-18.561665\du}{7.812300\du}}
\pgfusepath{stroke}
\pgfsetbuttcap
\pgfsetmiterjoin
\pgfsetdash{}{0pt}
\definecolor{dialinecolor}{rgb}{0.000000, 0.588235, 0.831373}
\pgfsetfillcolor{dialinecolor}
\pgfpathmoveto{\pgfpoint{-20.204000\du}{7.753187\du}}
\pgfpathlineto{\pgfpoint{-18.561665\du}{7.753187\du}}
\pgfpathlineto{\pgfpoint{-18.561665\du}{7.812300\du}}
\pgfpathlineto{\pgfpoint{-20.204000\du}{7.812300\du}}
\pgfpathlineto{\pgfpoint{-20.204000\du}{7.753187\du}}
\pgfusepath{fill}
\pgfsetbuttcap
\pgfsetmiterjoin
\pgfsetdash{}{0pt}
\definecolor{dialinecolor}{rgb}{0.666667, 0.901961, 1.000000}
\pgfsetstrokecolor{dialinecolor}
\pgfpathmoveto{\pgfpoint{-20.204000\du}{7.753187\du}}
\pgfpathlineto{\pgfpoint{-18.561665\du}{7.753187\du}}
\pgfpathlineto{\pgfpoint{-18.561665\du}{7.812300\du}}
\pgfpathlineto{\pgfpoint{-20.204000\du}{7.812300\du}}
\pgfpathlineto{\pgfpoint{-20.204000\du}{7.753187\du}}
\pgfusepath{stroke}
\pgfsetbuttcap
\pgfsetmiterjoin
\pgfsetdash{}{0pt}
\definecolor{dialinecolor}{rgb}{0.000000, 0.705882, 1.000000}
\pgfsetfillcolor{dialinecolor}
\pgfpathmoveto{\pgfpoint{-18.228780\du}{7.717062\du}}
\pgfpathlineto{\pgfpoint{-18.109658\du}{7.538826\du}}
\pgfpathlineto{\pgfpoint{-17.894999\du}{7.538826\du}}
\pgfpathlineto{\pgfpoint{-18.013822\du}{7.717062\du}}
\pgfpathlineto{\pgfpoint{-18.228780\du}{7.717062\du}}
\pgfusepath{fill}
\pgfsetbuttcap
\pgfsetmiterjoin
\pgfsetdash{}{0pt}
\definecolor{dialinecolor}{rgb}{0.666667, 0.901961, 1.000000}
\pgfsetstrokecolor{dialinecolor}
\pgfpathmoveto{\pgfpoint{-18.228780\du}{7.717062\du}}
\pgfpathlineto{\pgfpoint{-18.109658\du}{7.538826\du}}
\pgfpathlineto{\pgfpoint{-17.894999\du}{7.538826\du}}
\pgfpathlineto{\pgfpoint{-18.013822\du}{7.717062\du}}
\pgfpathlineto{\pgfpoint{-18.228780\du}{7.717062\du}}
\pgfusepath{stroke}
\pgfsetbuttcap
\pgfsetmiterjoin
\pgfsetdash{}{0pt}
\definecolor{dialinecolor}{rgb}{0.000000, 0.352941, 0.501961}
\pgfsetfillcolor{dialinecolor}
\pgfpathmoveto{\pgfpoint{-18.013822\du}{7.776772\du}}
\pgfpathlineto{\pgfpoint{-17.894999\du}{7.645708\du}}
\pgfpathlineto{\pgfpoint{-17.894999\du}{7.538826\du}}
\pgfpathlineto{\pgfpoint{-18.013822\du}{7.717062\du}}
\pgfpathlineto{\pgfpoint{-18.013822\du}{7.776772\du}}
\pgfusepath{fill}
\pgfsetbuttcap
\pgfsetmiterjoin
\pgfsetdash{}{0pt}
\definecolor{dialinecolor}{rgb}{0.666667, 0.901961, 1.000000}
\pgfsetstrokecolor{dialinecolor}
\pgfpathmoveto{\pgfpoint{-18.013822\du}{7.776772\du}}
\pgfpathlineto{\pgfpoint{-17.894999\du}{7.645708\du}}
\pgfpathlineto{\pgfpoint{-17.894999\du}{7.538826\du}}
\pgfpathlineto{\pgfpoint{-18.013822\du}{7.717062\du}}
\pgfpathlineto{\pgfpoint{-18.013822\du}{7.776772\du}}
\pgfusepath{stroke}
\pgfsetbuttcap
\pgfsetmiterjoin
\pgfsetdash{}{0pt}
\definecolor{dialinecolor}{rgb}{0.000000, 0.588235, 0.831373}
\pgfsetfillcolor{dialinecolor}
\pgfpathmoveto{\pgfpoint{-18.228780\du}{7.717062\du}}
\pgfpathlineto{\pgfpoint{-18.013822\du}{7.717062\du}}
\pgfpathlineto{\pgfpoint{-18.013822\du}{7.776772\du}}
\pgfpathlineto{\pgfpoint{-18.228780\du}{7.776772\du}}
\pgfpathlineto{\pgfpoint{-18.228780\du}{7.717062\du}}
\pgfusepath{fill}
\pgfsetbuttcap
\pgfsetmiterjoin
\pgfsetdash{}{0pt}
\definecolor{dialinecolor}{rgb}{0.666667, 0.901961, 1.000000}
\pgfsetstrokecolor{dialinecolor}
\pgfpathmoveto{\pgfpoint{-18.228780\du}{7.717062\du}}
\pgfpathlineto{\pgfpoint{-18.013822\du}{7.717062\du}}
\pgfpathlineto{\pgfpoint{-18.013822\du}{7.776772\du}}
\pgfpathlineto{\pgfpoint{-18.228780\du}{7.776772\du}}
\pgfpathlineto{\pgfpoint{-18.228780\du}{7.717062\du}}
\pgfusepath{stroke}
\pgfsetbuttcap
\pgfsetmiterjoin
\pgfsetdash{}{0pt}
\definecolor{dialinecolor}{rgb}{0.811765, 0.054902, 0.188235}
\pgfsetstrokecolor{dialinecolor}
\pgfpathmoveto{\pgfpoint{-19.442692\du}{6.348201\du}}
\pgfpathlineto{\pgfpoint{-19.108911\du}{6.348201\du}}
\pgfpathlineto{\pgfpoint{-19.192506\du}{6.431497\du}}
\pgfpathlineto{\pgfpoint{-18.906791\du}{6.431497\du}}
\pgfusepath{stroke}
\pgfsetbuttcap
\pgfsetmiterjoin
\pgfsetdash{}{0pt}
\definecolor{dialinecolor}{rgb}{0.811765, 0.054902, 0.188235}
\pgfsetstrokecolor{dialinecolor}
\pgfpathmoveto{\pgfpoint{-19.478519\du}{6.348201\du}}
\pgfpathlineto{\pgfpoint{-19.239976\du}{6.586147\du}}
\pgfpathlineto{\pgfpoint{-19.359396\du}{6.586147\du}}
\pgfpathlineto{\pgfpoint{-19.156978\du}{6.788864\du}}
\pgfusepath{stroke}
\pgfsetbuttcap
\pgfsetmiterjoin
\pgfsetdash{}{0pt}
\definecolor{dialinecolor}{rgb}{0.811765, 0.054902, 0.188235}
\pgfsetstrokecolor{dialinecolor}
\pgfpathmoveto{\pgfpoint{-18.800208\du}{6.348201\du}}
\pgfpathlineto{\pgfpoint{-19.037557\du}{6.586147\du}}
\pgfpathlineto{\pgfpoint{-18.919032\du}{6.586147\du}}
\pgfpathlineto{\pgfpoint{-19.121152\du}{6.788864\du}}
\pgfusepath{stroke}
\pgfsetbuttcap
\pgfsetmiterjoin
\pgfsetdash{}{0pt}
\definecolor{dialinecolor}{rgb}{0.000000, 0.470588, 0.666667}
\pgfsetfillcolor{dialinecolor}
\pgfpathmoveto{\pgfpoint{-19.311927\du}{6.437169\du}}
\pgfpathlineto{\pgfpoint{-19.312524\du}{6.433885\du}}
\pgfpathlineto{\pgfpoint{-19.313419\du}{6.430601\du}}
\pgfpathlineto{\pgfpoint{-19.314315\du}{6.427317\du}}
\pgfpathlineto{\pgfpoint{-19.316106\du}{6.424033\du}}
\pgfpathlineto{\pgfpoint{-19.318495\du}{6.421048\du}}
\pgfpathlineto{\pgfpoint{-19.321182\du}{6.417764\du}}
\pgfpathlineto{\pgfpoint{-19.324167\du}{6.414479\du}}
\pgfpathlineto{\pgfpoint{-19.327750\du}{6.411494\du}}
\pgfpathlineto{\pgfpoint{-19.331930\du}{6.408508\du}}
\pgfpathlineto{\pgfpoint{-19.336706\du}{6.405523\du}}
\pgfpathlineto{\pgfpoint{-19.341782\du}{6.402836\du}}
\pgfpathlineto{\pgfpoint{-19.347454\du}{6.399850\du}}
\pgfpathlineto{\pgfpoint{-19.353127\du}{6.397462\du}}
\pgfpathlineto{\pgfpoint{-19.359396\du}{6.394775\du}}
\pgfpathlineto{\pgfpoint{-19.365964\du}{6.392387\du}}
\pgfpathlineto{\pgfpoint{-19.373130\du}{6.389998\du}}
\pgfpathlineto{\pgfpoint{-19.380594\du}{6.387908\du}}
\pgfpathlineto{\pgfpoint{-19.388057\du}{6.385520\du}}
\pgfpathlineto{\pgfpoint{-19.396417\du}{6.384027\du}}
\pgfpathlineto{\pgfpoint{-19.404776\du}{6.381937\du}}
\pgfpathlineto{\pgfpoint{-19.412837\du}{6.380146\du}}
\pgfpathlineto{\pgfpoint{-19.422092\du}{6.378952\du}}
\pgfpathlineto{\pgfpoint{-19.430750\du}{6.377459\du}}
\pgfpathlineto{\pgfpoint{-19.440304\du}{6.376265\du}}
\pgfpathlineto{\pgfpoint{-19.449559\du}{6.375071\du}}
\pgfpathlineto{\pgfpoint{-19.458814\du}{6.374175\du}}
\pgfpathlineto{\pgfpoint{-19.468666\du}{6.373279\du}}
\pgfpathlineto{\pgfpoint{-19.478519\du}{6.372682\du}}
\pgfpathlineto{\pgfpoint{-19.488371\du}{6.372384\du}}
\pgfpathlineto{\pgfpoint{-19.498522\du}{6.372085\du}}
\pgfpathlineto{\pgfpoint{-19.508374\du}{6.372085\du}}
\pgfpathlineto{\pgfpoint{-19.508374\du}{6.372085\du}}
\pgfpathlineto{\pgfpoint{-19.518226\du}{6.372085\du}}
\pgfpathlineto{\pgfpoint{-19.528078\du}{6.372384\du}}
\pgfpathlineto{\pgfpoint{-19.538229\du}{6.372682\du}}
\pgfpathlineto{\pgfpoint{-19.548081\du}{6.373279\du}}
\pgfpathlineto{\pgfpoint{-19.557635\du}{6.374175\du}}
\pgfpathlineto{\pgfpoint{-19.567189\du}{6.375071\du}}
\pgfpathlineto{\pgfpoint{-19.576444\du}{6.376265\du}}
\pgfpathlineto{\pgfpoint{-19.585699\du}{6.377459\du}}
\pgfpathlineto{\pgfpoint{-19.594655\du}{6.378952\du}}
\pgfpathlineto{\pgfpoint{-19.603612\du}{6.380146\du}}
\pgfpathlineto{\pgfpoint{-19.611971\du}{6.381937\du}}
\pgfpathlineto{\pgfpoint{-19.620331\du}{6.384027\du}}
\pgfpathlineto{\pgfpoint{-19.628690\du}{6.385520\du}}
\pgfpathlineto{\pgfpoint{-19.635856\du}{6.387908\du}}
\pgfpathlineto{\pgfpoint{-19.643319\du}{6.389998\du}}
\pgfpathlineto{\pgfpoint{-19.650485\du}{6.392387\du}}
\pgfpathlineto{\pgfpoint{-19.657351\du}{6.394775\du}}
\pgfpathlineto{\pgfpoint{-19.663322\du}{6.397462\du}}
\pgfpathlineto{\pgfpoint{-19.669293\du}{6.399850\du}}
\pgfpathlineto{\pgfpoint{-19.674966\du}{6.402836\du}}
\pgfpathlineto{\pgfpoint{-19.680041\du}{6.405523\du}}
\pgfpathlineto{\pgfpoint{-19.684818\du}{6.408508\du}}
\pgfpathlineto{\pgfpoint{-19.688998\du}{6.411494\du}}
\pgfpathlineto{\pgfpoint{-19.692282\du}{6.414479\du}}
\pgfpathlineto{\pgfpoint{-19.695566\du}{6.417764\du}}
\pgfpathlineto{\pgfpoint{-19.698551\du}{6.421048\du}}
\pgfpathlineto{\pgfpoint{-19.700641\du}{6.424033\du}}
\pgfpathlineto{\pgfpoint{-19.702134\du}{6.427317\du}}
\pgfpathlineto{\pgfpoint{-19.703925\du}{6.430601\du}}
\pgfpathlineto{\pgfpoint{-19.704224\du}{6.433885\du}}
\pgfpathlineto{\pgfpoint{-19.704522\du}{6.437169\du}}
\pgfpathlineto{\pgfpoint{-19.704522\du}{6.437169\du}}
\pgfpathlineto{\pgfpoint{-19.704224\du}{6.440752\du}}
\pgfpathlineto{\pgfpoint{-19.703925\du}{6.443738\du}}
\pgfpathlineto{\pgfpoint{-19.702134\du}{6.447022\du}}
\pgfpathlineto{\pgfpoint{-19.700641\du}{6.450306\du}}
\pgfpathlineto{\pgfpoint{-19.698551\du}{6.453590\du}}
\pgfpathlineto{\pgfpoint{-19.695566\du}{6.456874\du}}
\pgfpathlineto{\pgfpoint{-19.692282\du}{6.459859\du}}
\pgfpathlineto{\pgfpoint{-19.688998\du}{6.462845\du}}
\pgfpathlineto{\pgfpoint{-19.684818\du}{6.465830\du}}
\pgfpathlineto{\pgfpoint{-19.680041\du}{6.469114\du}}
\pgfpathlineto{\pgfpoint{-19.674966\du}{6.471801\du}}
\pgfpathlineto{\pgfpoint{-19.669293\du}{6.474488\du}}
\pgfpathlineto{\pgfpoint{-19.663322\du}{6.477175\du}}
\pgfpathlineto{\pgfpoint{-19.657351\du}{6.479862\du}}
\pgfpathlineto{\pgfpoint{-19.650485\du}{6.482251\du}}
\pgfpathlineto{\pgfpoint{-19.643319\du}{6.484639\du}}
\pgfpathlineto{\pgfpoint{-19.635856\du}{6.486430\du}}
\pgfpathlineto{\pgfpoint{-19.628690\du}{6.488520\du}}
\pgfpathlineto{\pgfpoint{-19.620331\du}{6.490610\du}}
\pgfpathlineto{\pgfpoint{-19.611971\du}{6.492700\du}}
\pgfpathlineto{\pgfpoint{-19.603612\du}{6.494193\du}}
\pgfpathlineto{\pgfpoint{-19.594655\du}{6.495686\du}}
\pgfpathlineto{\pgfpoint{-19.585699\du}{6.497178\du}}
\pgfpathlineto{\pgfpoint{-19.576444\du}{6.498373\du}}
\pgfpathlineto{\pgfpoint{-19.567189\du}{6.499268\du}}
\pgfpathlineto{\pgfpoint{-19.557635\du}{6.500462\du}}
\pgfpathlineto{\pgfpoint{-19.548081\du}{6.501060\du}}
\pgfpathlineto{\pgfpoint{-19.538229\du}{6.501657\du}}
\pgfpathlineto{\pgfpoint{-19.528078\du}{6.502254\du}}
\pgfpathlineto{\pgfpoint{-19.518226\du}{6.502552\du}}
\pgfpathlineto{\pgfpoint{-19.508374\du}{6.502552\du}}
\pgfpathlineto{\pgfpoint{-19.508374\du}{6.502552\du}}
\pgfpathlineto{\pgfpoint{-19.498522\du}{6.502552\du}}
\pgfpathlineto{\pgfpoint{-19.488371\du}{6.502254\du}}
\pgfpathlineto{\pgfpoint{-19.478519\du}{6.501657\du}}
\pgfpathlineto{\pgfpoint{-19.468666\du}{6.501060\du}}
\pgfpathlineto{\pgfpoint{-19.458814\du}{6.500462\du}}
\pgfpathlineto{\pgfpoint{-19.449559\du}{6.499268\du}}
\pgfpathlineto{\pgfpoint{-19.440304\du}{6.498373\du}}
\pgfpathlineto{\pgfpoint{-19.430750\du}{6.497178\du}}
\pgfpathlineto{\pgfpoint{-19.422092\du}{6.495686\du}}
\pgfpathlineto{\pgfpoint{-19.412837\du}{6.494193\du}}
\pgfpathlineto{\pgfpoint{-19.404776\du}{6.492700\du}}
\pgfpathlineto{\pgfpoint{-19.396417\du}{6.490610\du}}
\pgfpathlineto{\pgfpoint{-19.388057\du}{6.488520\du}}
\pgfpathlineto{\pgfpoint{-19.380594\du}{6.486430\du}}
\pgfpathlineto{\pgfpoint{-19.373130\du}{6.484639\du}}
\pgfpathlineto{\pgfpoint{-19.365964\du}{6.482251\du}}
\pgfpathlineto{\pgfpoint{-19.359396\du}{6.479862\du}}
\pgfpathlineto{\pgfpoint{-19.353127\du}{6.477175\du}}
\pgfpathlineto{\pgfpoint{-19.347454\du}{6.474488\du}}
\pgfpathlineto{\pgfpoint{-19.341782\du}{6.471801\du}}
\pgfpathlineto{\pgfpoint{-19.336706\du}{6.469114\du}}
\pgfpathlineto{\pgfpoint{-19.331930\du}{6.465830\du}}
\pgfpathlineto{\pgfpoint{-19.327750\du}{6.462845\du}}
\pgfpathlineto{\pgfpoint{-19.324167\du}{6.459859\du}}
\pgfpathlineto{\pgfpoint{-19.321182\du}{6.456874\du}}
\pgfpathlineto{\pgfpoint{-19.318495\du}{6.453590\du}}
\pgfpathlineto{\pgfpoint{-19.316106\du}{6.450306\du}}
\pgfpathlineto{\pgfpoint{-19.314315\du}{6.447022\du}}
\pgfpathlineto{\pgfpoint{-19.313419\du}{6.443738\du}}
\pgfpathlineto{\pgfpoint{-19.312524\du}{6.440752\du}}
\pgfpathlineto{\pgfpoint{-19.311927\du}{6.437169\du}}
\pgfusepath{fill}
\pgfsetbuttcap
\pgfsetmiterjoin
\pgfsetdash{}{0pt}
\definecolor{dialinecolor}{rgb}{0.666667, 0.901961, 1.000000}
\pgfsetstrokecolor{dialinecolor}
\pgfpathmoveto{\pgfpoint{-19.324167\du}{6.431497\du}}
\pgfpathlineto{\pgfpoint{-19.324167\du}{6.427914\du}}
\pgfpathlineto{\pgfpoint{-19.325063\du}{6.425526\du}}
\pgfpathlineto{\pgfpoint{-19.326257\du}{6.422242\du}}
\pgfpathlineto{\pgfpoint{-19.327750\du}{6.419256\du}}
\pgfpathlineto{\pgfpoint{-19.330138\du}{6.416569\du}}
\pgfpathlineto{\pgfpoint{-19.332527\du}{6.412987\du}}
\pgfpathlineto{\pgfpoint{-19.335811\du}{6.410598\du}}
\pgfpathlineto{\pgfpoint{-19.339095\du}{6.407911\du}}
\pgfpathlineto{\pgfpoint{-19.343275\du}{6.405224\du}}
\pgfpathlineto{\pgfpoint{-19.347753\du}{6.402836\du}}
\pgfpathlineto{\pgfpoint{-19.352530\du}{6.399850\du}}
\pgfpathlineto{\pgfpoint{-19.358202\du}{6.397761\du}}
\pgfpathlineto{\pgfpoint{-19.363576\du}{6.395074\du}}
\pgfpathlineto{\pgfpoint{-19.369846\du}{6.392387\du}}
\pgfpathlineto{\pgfpoint{-19.375817\du}{6.390595\du}}
\pgfpathlineto{\pgfpoint{-19.383280\du}{6.388207\du}}
\pgfpathlineto{\pgfpoint{-19.390147\du}{6.386714\du}}
\pgfpathlineto{\pgfpoint{-19.397312\du}{6.384624\du}}
\pgfpathlineto{\pgfpoint{-19.405075\du}{6.382833\du}}
\pgfpathlineto{\pgfpoint{-19.413733\du}{6.381042\du}}
\pgfpathlineto{\pgfpoint{-19.422092\du}{6.379847\du}}
\pgfpathlineto{\pgfpoint{-19.430452\du}{6.377459\du}}
\pgfpathlineto{\pgfpoint{-19.439408\du}{6.377160\du}}
\pgfpathlineto{\pgfpoint{-19.448066\du}{6.375668\du}}
\pgfpathlineto{\pgfpoint{-19.457023\du}{6.374772\du}}
\pgfpathlineto{\pgfpoint{-19.466278\du}{6.373876\du}}
\pgfpathlineto{\pgfpoint{-19.475832\du}{6.373279\du}}
\pgfpathlineto{\pgfpoint{-19.485385\du}{6.372981\du}}
\pgfpathlineto{\pgfpoint{-19.494342\du}{6.372384\du}}
\pgfpathlineto{\pgfpoint{-19.504493\du}{6.372085\du}}
\pgfpathlineto{\pgfpoint{-19.514046\du}{6.372085\du}}
\pgfpathlineto{\pgfpoint{-19.514046\du}{6.372085\du}}
\pgfpathlineto{\pgfpoint{-19.523301\du}{6.372085\du}}
\pgfpathlineto{\pgfpoint{-19.533154\du}{6.372384\du}}
\pgfpathlineto{\pgfpoint{-19.542707\du}{6.372981\du}}
\pgfpathlineto{\pgfpoint{-19.552261\du}{6.373279\du}}
\pgfpathlineto{\pgfpoint{-19.561516\du}{6.373876\du}}
\pgfpathlineto{\pgfpoint{-19.571368\du}{6.374772\du}}
\pgfpathlineto{\pgfpoint{-19.580026\du}{6.375668\du}}
\pgfpathlineto{\pgfpoint{-19.588983\du}{6.377160\du}}
\pgfpathlineto{\pgfpoint{-19.597342\du}{6.377459\du}}
\pgfpathlineto{\pgfpoint{-19.606299\du}{6.379847\du}}
\pgfpathlineto{\pgfpoint{-19.614658\du}{6.381042\du}}
\pgfpathlineto{\pgfpoint{-19.622719\du}{6.382833\du}}
\pgfpathlineto{\pgfpoint{-19.630482\du}{6.384624\du}}
\pgfpathlineto{\pgfpoint{-19.637945\du}{6.386714\du}}
\pgfpathlineto{\pgfpoint{-19.644812\du}{6.388207\du}}
\pgfpathlineto{\pgfpoint{-19.652276\du}{6.390595\du}}
\pgfpathlineto{\pgfpoint{-19.658247\du}{6.392387\du}}
\pgfpathlineto{\pgfpoint{-19.664516\du}{6.395074\du}}
\pgfpathlineto{\pgfpoint{-19.670189\du}{6.397761\du}}
\pgfpathlineto{\pgfpoint{-19.675563\du}{6.399850\du}}
\pgfpathlineto{\pgfpoint{-19.680340\du}{6.402836\du}}
\pgfpathlineto{\pgfpoint{-19.684818\du}{6.405224\du}}
\pgfpathlineto{\pgfpoint{-19.688998\du}{6.407911\du}}
\pgfpathlineto{\pgfpoint{-19.692282\du}{6.410598\du}}
\pgfpathlineto{\pgfpoint{-19.695864\du}{6.412987\du}}
\pgfpathlineto{\pgfpoint{-19.697656\du}{6.416569\du}}
\pgfpathlineto{\pgfpoint{-19.700044\du}{6.419256\du}}
\pgfpathlineto{\pgfpoint{-19.701835\du}{6.422242\du}}
\pgfpathlineto{\pgfpoint{-19.703328\du}{6.425526\du}}
\pgfpathlineto{\pgfpoint{-19.703925\du}{6.427914\du}}
\pgfpathlineto{\pgfpoint{-19.704224\du}{6.431497\du}}
\pgfpathlineto{\pgfpoint{-19.704224\du}{6.431497\du}}
\pgfpathlineto{\pgfpoint{-19.703925\du}{6.434482\du}}
\pgfpathlineto{\pgfpoint{-19.703328\du}{6.437169\du}}
\pgfpathlineto{\pgfpoint{-19.701835\du}{6.440155\du}}
\pgfpathlineto{\pgfpoint{-19.700044\du}{6.443738\du}}
\pgfpathlineto{\pgfpoint{-19.697656\du}{6.446424\du}}
\pgfpathlineto{\pgfpoint{-19.695864\du}{6.449410\du}}
\pgfpathlineto{\pgfpoint{-19.692282\du}{6.452097\du}}
\pgfpathlineto{\pgfpoint{-19.688998\du}{6.454784\du}}
\pgfpathlineto{\pgfpoint{-19.684818\du}{6.457471\du}}
\pgfpathlineto{\pgfpoint{-19.680340\du}{6.460158\du}}
\pgfpathlineto{\pgfpoint{-19.675563\du}{6.462845\du}}
\pgfpathlineto{\pgfpoint{-19.670189\du}{6.465233\du}}
\pgfpathlineto{\pgfpoint{-19.664516\du}{6.467920\du}}
\pgfpathlineto{\pgfpoint{-19.658247\du}{6.470010\du}}
\pgfpathlineto{\pgfpoint{-19.652276\du}{6.472399\du}}
\pgfpathlineto{\pgfpoint{-19.644812\du}{6.474787\du}}
\pgfpathlineto{\pgfpoint{-19.637945\du}{6.476578\du}}
\pgfpathlineto{\pgfpoint{-19.630482\du}{6.478071\du}}
\pgfpathlineto{\pgfpoint{-19.622719\du}{6.480161\du}}
\pgfpathlineto{\pgfpoint{-19.614658\du}{6.481952\du}}
\pgfpathlineto{\pgfpoint{-19.606299\du}{6.483146\du}}
\pgfpathlineto{\pgfpoint{-19.597342\du}{6.484938\du}}
\pgfpathlineto{\pgfpoint{-19.588983\du}{6.485833\du}}
\pgfpathlineto{\pgfpoint{-19.580026\du}{6.487326\du}}
\pgfpathlineto{\pgfpoint{-19.571368\du}{6.488222\du}}
\pgfpathlineto{\pgfpoint{-19.561516\du}{6.488520\du}}
\pgfpathlineto{\pgfpoint{-19.552261\du}{6.489715\du}}
\pgfpathlineto{\pgfpoint{-19.542707\du}{6.490013\du}}
\pgfpathlineto{\pgfpoint{-19.533154\du}{6.490312\du}}
\pgfpathlineto{\pgfpoint{-19.523301\du}{6.490909\du}}
\pgfpathlineto{\pgfpoint{-19.514046\du}{6.490909\du}}
\pgfpathlineto{\pgfpoint{-19.514046\du}{6.490909\du}}
\pgfpathlineto{\pgfpoint{-19.504493\du}{6.490909\du}}
\pgfpathlineto{\pgfpoint{-19.494342\du}{6.490312\du}}
\pgfpathlineto{\pgfpoint{-19.485385\du}{6.490013\du}}
\pgfpathlineto{\pgfpoint{-19.475832\du}{6.489715\du}}
\pgfpathlineto{\pgfpoint{-19.466278\du}{6.488520\du}}
\pgfpathlineto{\pgfpoint{-19.457023\du}{6.488222\du}}
\pgfpathlineto{\pgfpoint{-19.448066\du}{6.487326\du}}
\pgfpathlineto{\pgfpoint{-19.439408\du}{6.485833\du}}
\pgfpathlineto{\pgfpoint{-19.430452\du}{6.484938\du}}
\pgfpathlineto{\pgfpoint{-19.422092\du}{6.483146\du}}
\pgfpathlineto{\pgfpoint{-19.413733\du}{6.481952\du}}
\pgfpathlineto{\pgfpoint{-19.405075\du}{6.480161\du}}
\pgfpathlineto{\pgfpoint{-19.397312\du}{6.478071\du}}
\pgfpathlineto{\pgfpoint{-19.390147\du}{6.476578\du}}
\pgfpathlineto{\pgfpoint{-19.383280\du}{6.474787\du}}
\pgfpathlineto{\pgfpoint{-19.375817\du}{6.472399\du}}
\pgfpathlineto{\pgfpoint{-19.369846\du}{6.470010\du}}
\pgfpathlineto{\pgfpoint{-19.363576\du}{6.467920\du}}
\pgfpathlineto{\pgfpoint{-19.358202\du}{6.465233\du}}
\pgfpathlineto{\pgfpoint{-19.352530\du}{6.462845\du}}
\pgfpathlineto{\pgfpoint{-19.347753\du}{6.460158\du}}
\pgfpathlineto{\pgfpoint{-19.343275\du}{6.457471\du}}
\pgfpathlineto{\pgfpoint{-19.339095\du}{6.454784\du}}
\pgfpathlineto{\pgfpoint{-19.335811\du}{6.452097\du}}
\pgfpathlineto{\pgfpoint{-19.332527\du}{6.449410\du}}
\pgfpathlineto{\pgfpoint{-19.330138\du}{6.446424\du}}
\pgfpathlineto{\pgfpoint{-19.327750\du}{6.443738\du}}
\pgfpathlineto{\pgfpoint{-19.326257\du}{6.440155\du}}
\pgfpathlineto{\pgfpoint{-19.325063\du}{6.437169\du}}
\pgfpathlineto{\pgfpoint{-19.324167\du}{6.434482\du}}
\pgfpathlineto{\pgfpoint{-19.324167\du}{6.431497\du}}
\pgfusepath{stroke}
\pgfsetbuttcap
\pgfsetmiterjoin
\pgfsetdash{}{0pt}
\definecolor{dialinecolor}{rgb}{0.000000, 0.470588, 0.666667}
\pgfsetfillcolor{dialinecolor}
\pgfpathmoveto{\pgfpoint{-19.704224\du}{6.348201\du}}
\pgfpathlineto{\pgfpoint{-19.704224\du}{6.442842\du}}
\pgfpathlineto{\pgfpoint{-19.311628\du}{6.442842\du}}
\pgfpathlineto{\pgfpoint{-19.311628\du}{6.348201\du}}
\pgfpathlineto{\pgfpoint{-19.704224\du}{6.348201\du}}
\pgfusepath{fill}
\pgfsetbuttcap
\pgfsetmiterjoin
\pgfsetdash{}{0pt}
\definecolor{dialinecolor}{rgb}{0.000000, 0.705882, 1.000000}
\pgfsetfillcolor{dialinecolor}
\pgfpathmoveto{\pgfpoint{-19.311927\du}{6.341633\du}}
\pgfpathlineto{\pgfpoint{-19.312524\du}{6.338349\du}}
\pgfpathlineto{\pgfpoint{-19.313419\du}{6.334766\du}}
\pgfpathlineto{\pgfpoint{-19.314315\du}{6.331781\du}}
\pgfpathlineto{\pgfpoint{-19.316106\du}{6.328496\du}}
\pgfpathlineto{\pgfpoint{-19.318495\du}{6.325212\du}}
\pgfpathlineto{\pgfpoint{-19.321182\du}{6.321928\du}}
\pgfpathlineto{\pgfpoint{-19.324167\du}{6.318943\du}}
\pgfpathlineto{\pgfpoint{-19.327750\du}{6.315659\du}}
\pgfpathlineto{\pgfpoint{-19.331930\du}{6.312972\du}}
\pgfpathlineto{\pgfpoint{-19.336706\du}{6.309986\du}}
\pgfpathlineto{\pgfpoint{-19.341782\du}{6.307001\du}}
\pgfpathlineto{\pgfpoint{-19.347454\du}{6.304314\du}}
\pgfpathlineto{\pgfpoint{-19.353127\du}{6.301030\du}}
\pgfpathlineto{\pgfpoint{-19.359396\du}{6.299238\du}}
\pgfpathlineto{\pgfpoint{-19.365964\du}{6.296551\du}}
\pgfpathlineto{\pgfpoint{-19.373130\du}{6.294163\du}}
\pgfpathlineto{\pgfpoint{-19.380594\du}{6.292073\du}}
\pgfpathlineto{\pgfpoint{-19.388057\du}{6.289983\du}}
\pgfpathlineto{\pgfpoint{-19.396417\du}{6.287893\du}}
\pgfpathlineto{\pgfpoint{-19.404776\du}{6.285804\du}}
\pgfpathlineto{\pgfpoint{-19.412837\du}{6.284311\du}}
\pgfpathlineto{\pgfpoint{-19.422092\du}{6.282818\du}}
\pgfpathlineto{\pgfpoint{-19.430750\du}{6.281624\du}}
\pgfpathlineto{\pgfpoint{-19.440304\du}{6.280131\du}}
\pgfpathlineto{\pgfpoint{-19.449559\du}{6.279235\du}}
\pgfpathlineto{\pgfpoint{-19.458814\du}{6.278041\du}}
\pgfpathlineto{\pgfpoint{-19.468666\du}{6.277444\du}}
\pgfpathlineto{\pgfpoint{-19.478519\du}{6.276847\du}}
\pgfpathlineto{\pgfpoint{-19.488371\du}{6.276548\du}}
\pgfpathlineto{\pgfpoint{-19.498522\du}{6.276250\du}}
\pgfpathlineto{\pgfpoint{-19.508374\du}{6.276250\du}}
\pgfpathlineto{\pgfpoint{-19.508374\du}{6.276250\du}}
\pgfpathlineto{\pgfpoint{-19.518226\du}{6.276250\du}}
\pgfpathlineto{\pgfpoint{-19.528078\du}{6.276548\du}}
\pgfpathlineto{\pgfpoint{-19.538229\du}{6.276847\du}}
\pgfpathlineto{\pgfpoint{-19.548081\du}{6.277444\du}}
\pgfpathlineto{\pgfpoint{-19.557635\du}{6.278041\du}}
\pgfpathlineto{\pgfpoint{-19.567189\du}{6.279235\du}}
\pgfpathlineto{\pgfpoint{-19.576444\du}{6.280131\du}}
\pgfpathlineto{\pgfpoint{-19.585699\du}{6.281624\du}}
\pgfpathlineto{\pgfpoint{-19.594655\du}{6.282818\du}}
\pgfpathlineto{\pgfpoint{-19.603612\du}{6.284311\du}}
\pgfpathlineto{\pgfpoint{-19.611971\du}{6.285804\du}}
\pgfpathlineto{\pgfpoint{-19.620331\du}{6.287893\du}}
\pgfpathlineto{\pgfpoint{-19.628690\du}{6.289983\du}}
\pgfpathlineto{\pgfpoint{-19.635856\du}{6.292073\du}}
\pgfpathlineto{\pgfpoint{-19.643319\du}{6.294163\du}}
\pgfpathlineto{\pgfpoint{-19.650485\du}{6.296551\du}}
\pgfpathlineto{\pgfpoint{-19.657351\du}{6.299238\du}}
\pgfpathlineto{\pgfpoint{-19.663322\du}{6.301030\du}}
\pgfpathlineto{\pgfpoint{-19.669293\du}{6.304314\du}}
\pgfpathlineto{\pgfpoint{-19.674966\du}{6.307001\du}}
\pgfpathlineto{\pgfpoint{-19.680041\du}{6.309986\du}}
\pgfpathlineto{\pgfpoint{-19.684818\du}{6.312972\du}}
\pgfpathlineto{\pgfpoint{-19.688998\du}{6.315659\du}}
\pgfpathlineto{\pgfpoint{-19.692282\du}{6.318943\du}}
\pgfpathlineto{\pgfpoint{-19.695566\du}{6.321928\du}}
\pgfpathlineto{\pgfpoint{-19.698551\du}{6.325212\du}}
\pgfpathlineto{\pgfpoint{-19.700641\du}{6.328496\du}}
\pgfpathlineto{\pgfpoint{-19.702134\du}{6.331781\du}}
\pgfpathlineto{\pgfpoint{-19.703925\du}{6.334766\du}}
\pgfpathlineto{\pgfpoint{-19.704224\du}{6.338349\du}}
\pgfpathlineto{\pgfpoint{-19.704522\du}{6.341633\du}}
\pgfpathlineto{\pgfpoint{-19.704522\du}{6.341633\du}}
\pgfpathlineto{\pgfpoint{-19.704224\du}{6.344917\du}}
\pgfpathlineto{\pgfpoint{-19.703925\du}{6.348499\du}}
\pgfpathlineto{\pgfpoint{-19.702134\du}{6.351485\du}}
\pgfpathlineto{\pgfpoint{-19.700641\du}{6.354769\du}}
\pgfpathlineto{\pgfpoint{-19.698551\du}{6.358352\du}}
\pgfpathlineto{\pgfpoint{-19.695566\du}{6.361337\du}}
\pgfpathlineto{\pgfpoint{-19.692282\du}{6.364323\du}}
\pgfpathlineto{\pgfpoint{-19.688998\du}{6.367607\du}}
\pgfpathlineto{\pgfpoint{-19.684818\du}{6.369995\du}}
\pgfpathlineto{\pgfpoint{-19.680041\du}{6.373279\du}}
\pgfpathlineto{\pgfpoint{-19.674966\du}{6.376563\du}}
\pgfpathlineto{\pgfpoint{-19.669293\du}{6.379250\du}}
\pgfpathlineto{\pgfpoint{-19.663322\du}{6.381937\du}}
\pgfpathlineto{\pgfpoint{-19.657351\du}{6.384027\du}}
\pgfpathlineto{\pgfpoint{-19.650485\du}{6.387013\du}}
\pgfpathlineto{\pgfpoint{-19.643319\du}{6.389401\du}}
\pgfpathlineto{\pgfpoint{-19.635856\du}{6.391192\du}}
\pgfpathlineto{\pgfpoint{-19.628690\du}{6.393282\du}}
\pgfpathlineto{\pgfpoint{-19.620331\du}{6.395372\du}}
\pgfpathlineto{\pgfpoint{-19.611971\du}{6.397462\du}}
\pgfpathlineto{\pgfpoint{-19.603612\du}{6.398955\du}}
\pgfpathlineto{\pgfpoint{-19.594655\du}{6.400447\du}}
\pgfpathlineto{\pgfpoint{-19.585699\du}{6.401642\du}}
\pgfpathlineto{\pgfpoint{-19.576444\du}{6.403134\du}}
\pgfpathlineto{\pgfpoint{-19.567189\du}{6.404030\du}}
\pgfpathlineto{\pgfpoint{-19.557635\du}{6.405224\du}}
\pgfpathlineto{\pgfpoint{-19.548081\du}{6.405523\du}}
\pgfpathlineto{\pgfpoint{-19.538229\du}{6.406419\du}}
\pgfpathlineto{\pgfpoint{-19.528078\du}{6.406419\du}}
\pgfpathlineto{\pgfpoint{-19.518226\du}{6.407314\du}}
\pgfpathlineto{\pgfpoint{-19.508374\du}{6.407314\du}}
\pgfpathlineto{\pgfpoint{-19.508374\du}{6.407314\du}}
\pgfpathlineto{\pgfpoint{-19.498522\du}{6.407314\du}}
\pgfpathlineto{\pgfpoint{-19.488371\du}{6.406419\du}}
\pgfpathlineto{\pgfpoint{-19.478519\du}{6.406419\du}}
\pgfpathlineto{\pgfpoint{-19.468666\du}{6.405523\du}}
\pgfpathlineto{\pgfpoint{-19.458814\du}{6.405224\du}}
\pgfpathlineto{\pgfpoint{-19.449559\du}{6.404030\du}}
\pgfpathlineto{\pgfpoint{-19.440304\du}{6.403134\du}}
\pgfpathlineto{\pgfpoint{-19.430750\du}{6.401642\du}}
\pgfpathlineto{\pgfpoint{-19.422092\du}{6.400447\du}}
\pgfpathlineto{\pgfpoint{-19.412837\du}{6.398955\du}}
\pgfpathlineto{\pgfpoint{-19.404776\du}{6.397462\du}}
\pgfpathlineto{\pgfpoint{-19.396417\du}{6.395372\du}}
\pgfpathlineto{\pgfpoint{-19.388057\du}{6.393282\du}}
\pgfpathlineto{\pgfpoint{-19.380594\du}{6.391192\du}}
\pgfpathlineto{\pgfpoint{-19.373130\du}{6.389401\du}}
\pgfpathlineto{\pgfpoint{-19.365964\du}{6.387013\du}}
\pgfpathlineto{\pgfpoint{-19.359396\du}{6.384027\du}}
\pgfpathlineto{\pgfpoint{-19.353127\du}{6.381937\du}}
\pgfpathlineto{\pgfpoint{-19.347454\du}{6.379250\du}}
\pgfpathlineto{\pgfpoint{-19.341782\du}{6.376563\du}}
\pgfpathlineto{\pgfpoint{-19.336706\du}{6.373279\du}}
\pgfpathlineto{\pgfpoint{-19.331930\du}{6.369995\du}}
\pgfpathlineto{\pgfpoint{-19.327750\du}{6.367607\du}}
\pgfpathlineto{\pgfpoint{-19.324167\du}{6.364323\du}}
\pgfpathlineto{\pgfpoint{-19.321182\du}{6.361337\du}}
\pgfpathlineto{\pgfpoint{-19.318495\du}{6.358352\du}}
\pgfpathlineto{\pgfpoint{-19.316106\du}{6.354769\du}}
\pgfpathlineto{\pgfpoint{-19.314315\du}{6.351485\du}}
\pgfpathlineto{\pgfpoint{-19.313419\du}{6.348499\du}}
\pgfpathlineto{\pgfpoint{-19.312524\du}{6.344917\du}}
\pgfpathlineto{\pgfpoint{-19.311927\du}{6.341633\du}}
\pgfusepath{fill}
\pgfsetbuttcap
\pgfsetmiterjoin
\pgfsetdash{}{0pt}
\definecolor{dialinecolor}{rgb}{0.666667, 0.901961, 1.000000}
\pgfsetstrokecolor{dialinecolor}
\pgfpathmoveto{\pgfpoint{-19.324167\du}{6.335960\du}}
\pgfpathlineto{\pgfpoint{-19.324167\du}{6.332676\du}}
\pgfpathlineto{\pgfpoint{-19.325063\du}{6.330288\du}}
\pgfpathlineto{\pgfpoint{-19.326257\du}{6.327004\du}}
\pgfpathlineto{\pgfpoint{-19.327750\du}{6.323421\du}}
\pgfpathlineto{\pgfpoint{-19.330138\du}{6.321033\du}}
\pgfpathlineto{\pgfpoint{-19.332527\du}{6.317749\du}}
\pgfpathlineto{\pgfpoint{-19.335811\du}{6.315360\du}}
\pgfpathlineto{\pgfpoint{-19.339095\du}{6.312673\du}}
\pgfpathlineto{\pgfpoint{-19.343275\du}{6.309986\du}}
\pgfpathlineto{\pgfpoint{-19.347753\du}{6.307299\du}}
\pgfpathlineto{\pgfpoint{-19.352530\du}{6.304612\du}}
\pgfpathlineto{\pgfpoint{-19.358202\du}{6.301925\du}}
\pgfpathlineto{\pgfpoint{-19.363576\du}{6.299537\du}}
\pgfpathlineto{\pgfpoint{-19.369846\du}{6.297447\du}}
\pgfpathlineto{\pgfpoint{-19.375817\du}{6.295059\du}}
\pgfpathlineto{\pgfpoint{-19.383280\du}{6.292670\du}}
\pgfpathlineto{\pgfpoint{-19.390147\du}{6.290879\du}}
\pgfpathlineto{\pgfpoint{-19.397312\du}{6.289386\du}}
\pgfpathlineto{\pgfpoint{-19.405075\du}{6.286998\du}}
\pgfpathlineto{\pgfpoint{-19.413733\du}{6.285505\du}}
\pgfpathlineto{\pgfpoint{-19.422092\du}{6.284311\du}}
\pgfpathlineto{\pgfpoint{-19.430452\du}{6.282519\du}}
\pgfpathlineto{\pgfpoint{-19.439408\du}{6.281624\du}}
\pgfpathlineto{\pgfpoint{-19.448066\du}{6.280131\du}}
\pgfpathlineto{\pgfpoint{-19.457023\du}{6.279235\du}}
\pgfpathlineto{\pgfpoint{-19.466278\du}{6.278638\du}}
\pgfpathlineto{\pgfpoint{-19.475832\du}{6.277743\du}}
\pgfpathlineto{\pgfpoint{-19.485385\du}{6.277444\du}}
\pgfpathlineto{\pgfpoint{-19.494342\du}{6.277145\du}}
\pgfpathlineto{\pgfpoint{-19.504493\du}{6.276548\du}}
\pgfpathlineto{\pgfpoint{-19.514046\du}{6.276548\du}}
\pgfpathlineto{\pgfpoint{-19.514046\du}{6.276548\du}}
\pgfpathlineto{\pgfpoint{-19.523301\du}{6.276548\du}}
\pgfpathlineto{\pgfpoint{-19.533154\du}{6.277145\du}}
\pgfpathlineto{\pgfpoint{-19.542707\du}{6.277444\du}}
\pgfpathlineto{\pgfpoint{-19.552261\du}{6.277743\du}}
\pgfpathlineto{\pgfpoint{-19.561516\du}{6.278638\du}}
\pgfpathlineto{\pgfpoint{-19.571368\du}{6.279235\du}}
\pgfpathlineto{\pgfpoint{-19.580026\du}{6.280131\du}}
\pgfpathlineto{\pgfpoint{-19.588983\du}{6.281624\du}}
\pgfpathlineto{\pgfpoint{-19.597342\du}{6.282519\du}}
\pgfpathlineto{\pgfpoint{-19.606299\du}{6.284311\du}}
\pgfpathlineto{\pgfpoint{-19.614658\du}{6.285505\du}}
\pgfpathlineto{\pgfpoint{-19.622719\du}{6.286998\du}}
\pgfpathlineto{\pgfpoint{-19.630482\du}{6.289386\du}}
\pgfpathlineto{\pgfpoint{-19.637945\du}{6.290879\du}}
\pgfpathlineto{\pgfpoint{-19.644812\du}{6.292670\du}}
\pgfpathlineto{\pgfpoint{-19.652276\du}{6.295059\du}}
\pgfpathlineto{\pgfpoint{-19.658247\du}{6.297447\du}}
\pgfpathlineto{\pgfpoint{-19.664516\du}{6.299537\du}}
\pgfpathlineto{\pgfpoint{-19.670189\du}{6.301925\du}}
\pgfpathlineto{\pgfpoint{-19.675563\du}{6.304612\du}}
\pgfpathlineto{\pgfpoint{-19.680340\du}{6.307299\du}}
\pgfpathlineto{\pgfpoint{-19.684818\du}{6.309986\du}}
\pgfpathlineto{\pgfpoint{-19.688998\du}{6.312673\du}}
\pgfpathlineto{\pgfpoint{-19.692282\du}{6.315360\du}}
\pgfpathlineto{\pgfpoint{-19.695864\du}{6.317749\du}}
\pgfpathlineto{\pgfpoint{-19.697656\du}{6.321033\du}}
\pgfpathlineto{\pgfpoint{-19.700044\du}{6.323421\du}}
\pgfpathlineto{\pgfpoint{-19.701835\du}{6.327004\du}}
\pgfpathlineto{\pgfpoint{-19.703328\du}{6.330288\du}}
\pgfpathlineto{\pgfpoint{-19.703925\du}{6.332676\du}}
\pgfpathlineto{\pgfpoint{-19.704224\du}{6.335960\du}}
\pgfpathlineto{\pgfpoint{-19.704224\du}{6.335960\du}}
\pgfpathlineto{\pgfpoint{-19.703925\du}{6.339244\du}}
\pgfpathlineto{\pgfpoint{-19.703328\du}{6.341931\du}}
\pgfpathlineto{\pgfpoint{-19.701835\du}{6.344917\du}}
\pgfpathlineto{\pgfpoint{-19.700044\du}{6.348201\du}}
\pgfpathlineto{\pgfpoint{-19.697656\du}{6.350888\du}}
\pgfpathlineto{\pgfpoint{-19.695864\du}{6.354172\du}}
\pgfpathlineto{\pgfpoint{-19.692282\du}{6.356859\du}}
\pgfpathlineto{\pgfpoint{-19.688998\du}{6.359546\du}}
\pgfpathlineto{\pgfpoint{-19.684818\du}{6.362233\du}}
\pgfpathlineto{\pgfpoint{-19.680340\du}{6.364920\du}}
\pgfpathlineto{\pgfpoint{-19.675563\du}{6.367607\du}}
\pgfpathlineto{\pgfpoint{-19.670189\du}{6.369697\du}}
\pgfpathlineto{\pgfpoint{-19.664516\du}{6.372384\du}}
\pgfpathlineto{\pgfpoint{-19.658247\du}{6.374772\du}}
\pgfpathlineto{\pgfpoint{-19.652276\du}{6.377160\du}}
\pgfpathlineto{\pgfpoint{-19.644812\du}{6.379250\du}}
\pgfpathlineto{\pgfpoint{-19.637945\du}{6.381042\du}}
\pgfpathlineto{\pgfpoint{-19.630482\du}{6.382833\du}}
\pgfpathlineto{\pgfpoint{-19.622719\du}{6.384624\du}}
\pgfpathlineto{\pgfpoint{-19.614658\du}{6.386714\du}}
\pgfpathlineto{\pgfpoint{-19.606299\du}{6.387908\du}}
\pgfpathlineto{\pgfpoint{-19.597342\du}{6.389700\du}}
\pgfpathlineto{\pgfpoint{-19.588983\du}{6.390595\du}}
\pgfpathlineto{\pgfpoint{-19.580026\du}{6.392088\du}}
\pgfpathlineto{\pgfpoint{-19.571368\du}{6.392387\du}}
\pgfpathlineto{\pgfpoint{-19.561516\du}{6.393581\du}}
\pgfpathlineto{\pgfpoint{-19.552261\du}{6.393879\du}}
\pgfpathlineto{\pgfpoint{-19.542707\du}{6.394775\du}}
\pgfpathlineto{\pgfpoint{-19.533154\du}{6.395074\du}}
\pgfpathlineto{\pgfpoint{-19.523301\du}{6.395372\du}}
\pgfpathlineto{\pgfpoint{-19.514046\du}{6.395372\du}}
\pgfpathlineto{\pgfpoint{-19.514046\du}{6.395372\du}}
\pgfpathlineto{\pgfpoint{-19.504493\du}{6.395372\du}}
\pgfpathlineto{\pgfpoint{-19.494342\du}{6.395074\du}}
\pgfpathlineto{\pgfpoint{-19.485385\du}{6.394775\du}}
\pgfpathlineto{\pgfpoint{-19.475832\du}{6.393879\du}}
\pgfpathlineto{\pgfpoint{-19.466278\du}{6.393581\du}}
\pgfpathlineto{\pgfpoint{-19.457023\du}{6.392387\du}}
\pgfpathlineto{\pgfpoint{-19.448066\du}{6.392088\du}}
\pgfpathlineto{\pgfpoint{-19.439408\du}{6.390595\du}}
\pgfpathlineto{\pgfpoint{-19.430452\du}{6.389700\du}}
\pgfpathlineto{\pgfpoint{-19.422092\du}{6.387908\du}}
\pgfpathlineto{\pgfpoint{-19.413733\du}{6.386714\du}}
\pgfpathlineto{\pgfpoint{-19.405075\du}{6.384624\du}}
\pgfpathlineto{\pgfpoint{-19.397312\du}{6.382833\du}}
\pgfpathlineto{\pgfpoint{-19.390147\du}{6.381042\du}}
\pgfpathlineto{\pgfpoint{-19.383280\du}{6.379250\du}}
\pgfpathlineto{\pgfpoint{-19.375817\du}{6.377160\du}}
\pgfpathlineto{\pgfpoint{-19.369846\du}{6.374772\du}}
\pgfpathlineto{\pgfpoint{-19.363576\du}{6.372384\du}}
\pgfpathlineto{\pgfpoint{-19.358202\du}{6.369697\du}}
\pgfpathlineto{\pgfpoint{-19.352530\du}{6.367607\du}}
\pgfpathlineto{\pgfpoint{-19.347753\du}{6.364920\du}}
\pgfpathlineto{\pgfpoint{-19.343275\du}{6.362233\du}}
\pgfpathlineto{\pgfpoint{-19.339095\du}{6.359546\du}}
\pgfpathlineto{\pgfpoint{-19.335811\du}{6.356859\du}}
\pgfpathlineto{\pgfpoint{-19.332527\du}{6.354172\du}}
\pgfpathlineto{\pgfpoint{-19.330138\du}{6.350888\du}}
\pgfpathlineto{\pgfpoint{-19.327750\du}{6.348201\du}}
\pgfpathlineto{\pgfpoint{-19.326257\du}{6.344917\du}}
\pgfpathlineto{\pgfpoint{-19.325063\du}{6.341931\du}}
\pgfpathlineto{\pgfpoint{-19.324167\du}{6.339244\du}}
\pgfpathlineto{\pgfpoint{-19.324167\du}{6.335960\du}}
\pgfusepath{stroke}
\pgfsetbuttcap
\pgfsetmiterjoin
\pgfsetdash{}{0pt}
\definecolor{dialinecolor}{rgb}{0.000000, 0.000000, 0.000000}
\pgfsetfillcolor{dialinecolor}
\pgfpathmoveto{\pgfpoint{-19.501806\du}{6.324317\du}}
\pgfpathlineto{\pgfpoint{-19.478519\du}{6.335960\du}}
\pgfpathlineto{\pgfpoint{-19.406568\du}{6.312375\du}}
\pgfpathlineto{\pgfpoint{-19.371040\du}{6.324317\du}}
\pgfpathlineto{\pgfpoint{-19.394625\du}{6.300433\du}}
\pgfpathlineto{\pgfpoint{-19.478519\du}{6.300433\du}}
\pgfpathlineto{\pgfpoint{-19.442692\du}{6.300433\du}}
\pgfpathlineto{\pgfpoint{-19.501806\du}{6.324317\du}}
\pgfusepath{fill}
\pgfsetbuttcap
\pgfsetmiterjoin
\pgfsetdash{}{0pt}
\definecolor{dialinecolor}{rgb}{0.000000, 0.000000, 0.000000}
\pgfsetfillcolor{dialinecolor}
\pgfpathmoveto{\pgfpoint{-19.514046\du}{6.359844\du}}
\pgfpathlineto{\pgfpoint{-19.549873\du}{6.348201\du}}
\pgfpathlineto{\pgfpoint{-19.608986\du}{6.372085\du}}
\pgfpathlineto{\pgfpoint{-19.644812\du}{6.359844\du}}
\pgfpathlineto{\pgfpoint{-19.632571\du}{6.395372\du}}
\pgfpathlineto{\pgfpoint{-19.549873\du}{6.395372\du}}
\pgfpathlineto{\pgfpoint{-19.585400\du}{6.384027\du}}
\pgfpathlineto{\pgfpoint{-19.514046\du}{6.359844\du}}
\pgfusepath{fill}
\pgfsetbuttcap
\pgfsetmiterjoin
\pgfsetdash{}{0pt}
\definecolor{dialinecolor}{rgb}{0.000000, 0.000000, 0.000000}
\pgfsetfillcolor{dialinecolor}
\pgfpathmoveto{\pgfpoint{-19.632571\du}{6.300433\du}}
\pgfpathlineto{\pgfpoint{-19.608986\du}{6.287893\du}}
\pgfpathlineto{\pgfpoint{-19.537333\du}{6.324317\du}}
\pgfpathlineto{\pgfpoint{-19.514046\du}{6.312375\du}}
\pgfpathlineto{\pgfpoint{-19.525690\du}{6.335960\du}}
\pgfpathlineto{\pgfpoint{-19.608986\du}{6.335960\du}}
\pgfpathlineto{\pgfpoint{-19.573160\du}{6.324317\du}}
\pgfpathlineto{\pgfpoint{-19.632571\du}{6.300433\du}}
\pgfusepath{fill}
\pgfsetbuttcap
\pgfsetmiterjoin
\pgfsetdash{}{0pt}
\definecolor{dialinecolor}{rgb}{0.000000, 0.000000, 0.000000}
\pgfsetfillcolor{dialinecolor}
\pgfpathmoveto{\pgfpoint{-19.383280\du}{6.384027\du}}
\pgfpathlineto{\pgfpoint{-19.406568\du}{6.395372\du}}
\pgfpathlineto{\pgfpoint{-19.478519\du}{6.359844\du}}
\pgfpathlineto{\pgfpoint{-19.514046\du}{6.372085\du}}
\pgfpathlineto{\pgfpoint{-19.490162\du}{6.348201\du}}
\pgfpathlineto{\pgfpoint{-19.406568\du}{6.348201\du}}
\pgfpathlineto{\pgfpoint{-19.442692\du}{6.359844\du}}
\pgfpathlineto{\pgfpoint{-19.383280\du}{6.384027\du}}
\pgfusepath{fill}
\pgfsetbuttcap
\pgfsetmiterjoin
\pgfsetdash{}{0pt}
\definecolor{dialinecolor}{rgb}{1.000000, 1.000000, 1.000000}
\pgfsetfillcolor{dialinecolor}
\pgfpathmoveto{\pgfpoint{-19.501806\du}{6.335960\du}}
\pgfpathlineto{\pgfpoint{-19.478519\du}{6.335960\du}}
\pgfpathlineto{\pgfpoint{-19.406568\du}{6.312375\du}}
\pgfpathlineto{\pgfpoint{-19.371040\du}{6.324317\du}}
\pgfpathlineto{\pgfpoint{-19.394625\du}{6.300433\du}}
\pgfpathlineto{\pgfpoint{-19.466278\du}{6.300433\du}}
\pgfpathlineto{\pgfpoint{-19.442692\du}{6.300433\du}}
\pgfpathlineto{\pgfpoint{-19.501806\du}{6.335960\du}}
\pgfusepath{fill}
\pgfsetbuttcap
\pgfsetmiterjoin
\pgfsetdash{}{0pt}
\definecolor{dialinecolor}{rgb}{1.000000, 1.000000, 1.000000}
\pgfsetfillcolor{dialinecolor}
\pgfpathmoveto{\pgfpoint{-19.514046\du}{6.359844\du}}
\pgfpathlineto{\pgfpoint{-19.537333\du}{6.348201\du}}
\pgfpathlineto{\pgfpoint{-19.608986\du}{6.372085\du}}
\pgfpathlineto{\pgfpoint{-19.644812\du}{6.372085\du}}
\pgfpathlineto{\pgfpoint{-19.621226\du}{6.395372\du}}
\pgfpathlineto{\pgfpoint{-19.537333\du}{6.395372\du}}
\pgfpathlineto{\pgfpoint{-19.573160\du}{6.384027\du}}
\pgfpathlineto{\pgfpoint{-19.514046\du}{6.359844\du}}
\pgfusepath{fill}
\pgfsetbuttcap
\pgfsetmiterjoin
\pgfsetdash{}{0pt}
\definecolor{dialinecolor}{rgb}{1.000000, 1.000000, 1.000000}
\pgfsetfillcolor{dialinecolor}
\pgfpathmoveto{\pgfpoint{-19.632571\du}{6.300433\du}}
\pgfpathlineto{\pgfpoint{-19.608986\du}{6.300433\du}}
\pgfpathlineto{\pgfpoint{-19.537333\du}{6.324317\du}}
\pgfpathlineto{\pgfpoint{-19.501806\du}{6.312375\du}}
\pgfpathlineto{\pgfpoint{-19.525690\du}{6.335960\du}}
\pgfpathlineto{\pgfpoint{-19.608986\du}{6.335960\du}}
\pgfpathlineto{\pgfpoint{-19.573160\du}{6.335960\du}}
\pgfpathlineto{\pgfpoint{-19.632571\du}{6.300433\du}}
\pgfusepath{fill}
\pgfsetbuttcap
\pgfsetmiterjoin
\pgfsetdash{}{0pt}
\definecolor{dialinecolor}{rgb}{1.000000, 1.000000, 1.000000}
\pgfsetfillcolor{dialinecolor}
\pgfpathmoveto{\pgfpoint{-19.383280\du}{6.384027\du}}
\pgfpathlineto{\pgfpoint{-19.406568\du}{6.395372\du}}
\pgfpathlineto{\pgfpoint{-19.478519\du}{6.372085\du}}
\pgfpathlineto{\pgfpoint{-19.501806\du}{6.372085\du}}
\pgfpathlineto{\pgfpoint{-19.490162\du}{6.348201\du}}
\pgfpathlineto{\pgfpoint{-19.406568\du}{6.348201\du}}
\pgfpathlineto{\pgfpoint{-19.442692\du}{6.359844\du}}
\pgfpathlineto{\pgfpoint{-19.383280\du}{6.384027\du}}
\pgfusepath{fill}
\pgfsetbuttcap
\pgfsetmiterjoin
\pgfsetdash{}{0pt}
\definecolor{dialinecolor}{rgb}{0.666667, 0.901961, 1.000000}
\pgfsetstrokecolor{dialinecolor}
\pgfpathmoveto{\pgfpoint{-19.704224\du}{6.335960\du}}
\pgfpathlineto{\pgfpoint{-19.704224\du}{6.430900\du}}
\pgfusepath{stroke}
\pgfsetbuttcap
\pgfsetmiterjoin
\pgfsetdash{}{0pt}
\definecolor{dialinecolor}{rgb}{0.666667, 0.901961, 1.000000}
\pgfsetstrokecolor{dialinecolor}
\pgfpathmoveto{\pgfpoint{-19.323570\du}{6.335960\du}}
\pgfpathlineto{\pgfpoint{-19.323570\du}{6.430900\du}}
\pgfusepath{stroke}
\pgfsetbuttcap
\pgfsetmiterjoin
\pgfsetdash{}{0pt}
\definecolor{dialinecolor}{rgb}{0.000000, 0.470588, 0.666667}
\pgfsetfillcolor{dialinecolor}
\pgfpathmoveto{\pgfpoint{-18.597790\du}{6.437169\du}}
\pgfpathlineto{\pgfpoint{-18.597790\du}{6.433885\du}}
\pgfpathlineto{\pgfpoint{-18.598686\du}{6.430601\du}}
\pgfpathlineto{\pgfpoint{-18.599880\du}{6.427317\du}}
\pgfpathlineto{\pgfpoint{-18.601373\du}{6.424033\du}}
\pgfpathlineto{\pgfpoint{-18.603761\du}{6.421048\du}}
\pgfpathlineto{\pgfpoint{-18.606448\du}{6.417764\du}}
\pgfpathlineto{\pgfpoint{-18.609434\du}{6.414479\du}}
\pgfpathlineto{\pgfpoint{-18.613613\du}{6.411494\du}}
\pgfpathlineto{\pgfpoint{-18.617196\du}{6.408508\du}}
\pgfpathlineto{\pgfpoint{-18.621973\du}{6.405523\du}}
\pgfpathlineto{\pgfpoint{-18.627048\du}{6.402836\du}}
\pgfpathlineto{\pgfpoint{-18.632422\du}{6.399850\du}}
\pgfpathlineto{\pgfpoint{-18.638692\du}{6.397462\du}}
\pgfpathlineto{\pgfpoint{-18.644663\du}{6.394775\du}}
\pgfpathlineto{\pgfpoint{-18.651529\du}{6.392387\du}}
\pgfpathlineto{\pgfpoint{-18.658396\du}{6.389998\du}}
\pgfpathlineto{\pgfpoint{-18.665860\du}{6.387908\du}}
\pgfpathlineto{\pgfpoint{-18.673324\du}{6.385520\du}}
\pgfpathlineto{\pgfpoint{-18.681385\du}{6.384027\du}}
\pgfpathlineto{\pgfpoint{-18.690043\du}{6.381937\du}}
\pgfpathlineto{\pgfpoint{-18.698402\du}{6.380146\du}}
\pgfpathlineto{\pgfpoint{-18.707060\du}{6.378952\du}}
\pgfpathlineto{\pgfpoint{-18.716315\du}{6.377459\du}}
\pgfpathlineto{\pgfpoint{-18.725570\du}{6.376265\du}}
\pgfpathlineto{\pgfpoint{-18.734825\du}{6.375071\du}}
\pgfpathlineto{\pgfpoint{-18.744379\du}{6.374175\du}}
\pgfpathlineto{\pgfpoint{-18.754530\du}{6.373279\du}}
\pgfpathlineto{\pgfpoint{-18.764084\du}{6.372682\du}}
\pgfpathlineto{\pgfpoint{-18.773637\du}{6.372384\du}}
\pgfpathlineto{\pgfpoint{-18.783788\du}{6.372085\du}}
\pgfpathlineto{\pgfpoint{-18.793640\du}{6.372085\du}}
\pgfpathlineto{\pgfpoint{-18.793640\du}{6.372085\du}}
\pgfpathlineto{\pgfpoint{-18.803492\du}{6.372085\du}}
\pgfpathlineto{\pgfpoint{-18.813643\du}{6.372384\du}}
\pgfpathlineto{\pgfpoint{-18.823495\du}{6.372682\du}}
\pgfpathlineto{\pgfpoint{-18.833348\du}{6.373279\du}}
\pgfpathlineto{\pgfpoint{-18.842901\du}{6.374175\du}}
\pgfpathlineto{\pgfpoint{-18.852455\du}{6.375071\du}}
\pgfpathlineto{\pgfpoint{-18.862009\du}{6.376265\du}}
\pgfpathlineto{\pgfpoint{-18.870965\du}{6.377459\du}}
\pgfpathlineto{\pgfpoint{-18.880220\du}{6.378952\du}}
\pgfpathlineto{\pgfpoint{-18.888878\du}{6.380146\du}}
\pgfpathlineto{\pgfpoint{-18.897536\du}{6.381937\du}}
\pgfpathlineto{\pgfpoint{-18.905896\du}{6.384027\du}}
\pgfpathlineto{\pgfpoint{-18.913957\du}{6.385520\du}}
\pgfpathlineto{\pgfpoint{-18.921421\du}{6.387908\du}}
\pgfpathlineto{\pgfpoint{-18.928884\du}{6.389998\du}}
\pgfpathlineto{\pgfpoint{-18.936050\du}{6.392387\du}}
\pgfpathlineto{\pgfpoint{-18.942618\du}{6.394775\du}}
\pgfpathlineto{\pgfpoint{-18.948887\du}{6.397462\du}}
\pgfpathlineto{\pgfpoint{-18.954858\du}{6.399850\du}}
\pgfpathlineto{\pgfpoint{-18.960232\du}{6.402836\du}}
\pgfpathlineto{\pgfpoint{-18.965308\du}{6.405523\du}}
\pgfpathlineto{\pgfpoint{-18.970084\du}{6.408508\du}}
\pgfpathlineto{\pgfpoint{-18.973966\du}{6.411494\du}}
\pgfpathlineto{\pgfpoint{-18.977847\du}{6.414479\du}}
\pgfpathlineto{\pgfpoint{-18.980832\du}{6.417764\du}}
\pgfpathlineto{\pgfpoint{-18.983519\du}{6.421048\du}}
\pgfpathlineto{\pgfpoint{-18.985908\du}{6.424033\du}}
\pgfpathlineto{\pgfpoint{-18.987699\du}{6.427317\du}}
\pgfpathlineto{\pgfpoint{-18.989192\du}{6.430601\du}}
\pgfpathlineto{\pgfpoint{-18.990087\du}{6.433885\du}}
\pgfpathlineto{\pgfpoint{-18.990087\du}{6.437169\du}}
\pgfpathlineto{\pgfpoint{-18.990087\du}{6.437169\du}}
\pgfpathlineto{\pgfpoint{-18.990087\du}{6.440752\du}}
\pgfpathlineto{\pgfpoint{-18.989192\du}{6.443738\du}}
\pgfpathlineto{\pgfpoint{-18.987699\du}{6.447022\du}}
\pgfpathlineto{\pgfpoint{-18.985908\du}{6.450306\du}}
\pgfpathlineto{\pgfpoint{-18.983519\du}{6.453590\du}}
\pgfpathlineto{\pgfpoint{-18.980832\du}{6.456874\du}}
\pgfpathlineto{\pgfpoint{-18.977847\du}{6.459859\du}}
\pgfpathlineto{\pgfpoint{-18.973966\du}{6.462845\du}}
\pgfpathlineto{\pgfpoint{-18.970084\du}{6.465830\du}}
\pgfpathlineto{\pgfpoint{-18.965308\du}{6.469114\du}}
\pgfpathlineto{\pgfpoint{-18.960232\du}{6.471801\du}}
\pgfpathlineto{\pgfpoint{-18.954858\du}{6.474488\du}}
\pgfpathlineto{\pgfpoint{-18.948887\du}{6.477175\du}}
\pgfpathlineto{\pgfpoint{-18.942618\du}{6.479862\du}}
\pgfpathlineto{\pgfpoint{-18.936050\du}{6.482251\du}}
\pgfpathlineto{\pgfpoint{-18.928884\du}{6.484639\du}}
\pgfpathlineto{\pgfpoint{-18.921421\du}{6.486430\du}}
\pgfpathlineto{\pgfpoint{-18.913957\du}{6.488520\du}}
\pgfpathlineto{\pgfpoint{-18.905896\du}{6.490610\du}}
\pgfpathlineto{\pgfpoint{-18.897536\du}{6.492700\du}}
\pgfpathlineto{\pgfpoint{-18.888878\du}{6.494193\du}}
\pgfpathlineto{\pgfpoint{-18.880220\du}{6.495686\du}}
\pgfpathlineto{\pgfpoint{-18.870965\du}{6.497178\du}}
\pgfpathlineto{\pgfpoint{-18.862009\du}{6.498373\du}}
\pgfpathlineto{\pgfpoint{-18.852455\du}{6.499268\du}}
\pgfpathlineto{\pgfpoint{-18.842901\du}{6.500462\du}}
\pgfpathlineto{\pgfpoint{-18.833348\du}{6.501060\du}}
\pgfpathlineto{\pgfpoint{-18.823495\du}{6.501657\du}}
\pgfpathlineto{\pgfpoint{-18.813643\du}{6.502254\du}}
\pgfpathlineto{\pgfpoint{-18.803492\du}{6.502552\du}}
\pgfpathlineto{\pgfpoint{-18.793640\du}{6.502552\du}}
\pgfpathlineto{\pgfpoint{-18.793640\du}{6.502552\du}}
\pgfpathlineto{\pgfpoint{-18.783788\du}{6.502552\du}}
\pgfpathlineto{\pgfpoint{-18.773637\du}{6.502254\du}}
\pgfpathlineto{\pgfpoint{-18.764084\du}{6.501657\du}}
\pgfpathlineto{\pgfpoint{-18.754530\du}{6.501060\du}}
\pgfpathlineto{\pgfpoint{-18.744379\du}{6.500462\du}}
\pgfpathlineto{\pgfpoint{-18.734825\du}{6.499268\du}}
\pgfpathlineto{\pgfpoint{-18.725570\du}{6.498373\du}}
\pgfpathlineto{\pgfpoint{-18.716315\du}{6.497178\du}}
\pgfpathlineto{\pgfpoint{-18.707060\du}{6.495686\du}}
\pgfpathlineto{\pgfpoint{-18.698402\du}{6.494193\du}}
\pgfpathlineto{\pgfpoint{-18.690043\du}{6.492700\du}}
\pgfpathlineto{\pgfpoint{-18.681385\du}{6.490610\du}}
\pgfpathlineto{\pgfpoint{-18.673324\du}{6.488520\du}}
\pgfpathlineto{\pgfpoint{-18.665860\du}{6.486430\du}}
\pgfpathlineto{\pgfpoint{-18.658396\du}{6.484639\du}}
\pgfpathlineto{\pgfpoint{-18.651529\du}{6.482251\du}}
\pgfpathlineto{\pgfpoint{-18.644663\du}{6.479862\du}}
\pgfpathlineto{\pgfpoint{-18.638692\du}{6.477175\du}}
\pgfpathlineto{\pgfpoint{-18.632422\du}{6.474488\du}}
\pgfpathlineto{\pgfpoint{-18.627048\du}{6.471801\du}}
\pgfpathlineto{\pgfpoint{-18.621973\du}{6.469114\du}}
\pgfpathlineto{\pgfpoint{-18.617196\du}{6.465830\du}}
\pgfpathlineto{\pgfpoint{-18.613613\du}{6.462845\du}}
\pgfpathlineto{\pgfpoint{-18.609434\du}{6.459859\du}}
\pgfpathlineto{\pgfpoint{-18.606448\du}{6.456874\du}}
\pgfpathlineto{\pgfpoint{-18.603761\du}{6.453590\du}}
\pgfpathlineto{\pgfpoint{-18.601373\du}{6.450306\du}}
\pgfpathlineto{\pgfpoint{-18.599880\du}{6.447022\du}}
\pgfpathlineto{\pgfpoint{-18.598686\du}{6.443738\du}}
\pgfpathlineto{\pgfpoint{-18.597790\du}{6.440752\du}}
\pgfpathlineto{\pgfpoint{-18.597790\du}{6.437169\du}}
\pgfusepath{fill}
\pgfsetbuttcap
\pgfsetmiterjoin
\pgfsetdash{}{0pt}
\definecolor{dialinecolor}{rgb}{0.666667, 0.901961, 1.000000}
\pgfsetstrokecolor{dialinecolor}
\pgfpathmoveto{\pgfpoint{-18.609434\du}{6.431497\du}}
\pgfpathlineto{\pgfpoint{-18.609434\du}{6.427914\du}}
\pgfpathlineto{\pgfpoint{-18.610329\du}{6.425526\du}}
\pgfpathlineto{\pgfpoint{-18.611524\du}{6.422242\du}}
\pgfpathlineto{\pgfpoint{-18.613613\du}{6.419256\du}}
\pgfpathlineto{\pgfpoint{-18.615703\du}{6.416569\du}}
\pgfpathlineto{\pgfpoint{-18.617793\du}{6.412987\du}}
\pgfpathlineto{\pgfpoint{-18.621077\du}{6.410598\du}}
\pgfpathlineto{\pgfpoint{-18.624361\du}{6.407911\du}}
\pgfpathlineto{\pgfpoint{-18.628541\du}{6.405224\du}}
\pgfpathlineto{\pgfpoint{-18.633318\du}{6.402836\du}}
\pgfpathlineto{\pgfpoint{-18.637796\du}{6.399850\du}}
\pgfpathlineto{\pgfpoint{-18.643469\du}{6.397761\du}}
\pgfpathlineto{\pgfpoint{-18.649440\du}{6.395074\du}}
\pgfpathlineto{\pgfpoint{-18.655112\du}{6.392387\du}}
\pgfpathlineto{\pgfpoint{-18.661680\du}{6.390595\du}}
\pgfpathlineto{\pgfpoint{-18.668547\du}{6.388207\du}}
\pgfpathlineto{\pgfpoint{-18.675414\du}{6.386714\du}}
\pgfpathlineto{\pgfpoint{-18.683176\du}{6.384624\du}}
\pgfpathlineto{\pgfpoint{-18.690640\du}{6.382833\du}}
\pgfpathlineto{\pgfpoint{-18.698999\du}{6.381042\du}}
\pgfpathlineto{\pgfpoint{-18.707657\du}{6.379847\du}}
\pgfpathlineto{\pgfpoint{-18.716017\du}{6.377459\du}}
\pgfpathlineto{\pgfpoint{-18.724675\du}{6.377160\du}}
\pgfpathlineto{\pgfpoint{-18.733631\du}{6.375668\du}}
\pgfpathlineto{\pgfpoint{-18.742289\du}{6.374772\du}}
\pgfpathlineto{\pgfpoint{-18.751843\du}{6.373876\du}}
\pgfpathlineto{\pgfpoint{-18.761098\du}{6.373279\du}}
\pgfpathlineto{\pgfpoint{-18.770652\du}{6.372981\du}}
\pgfpathlineto{\pgfpoint{-18.780205\du}{6.372384\du}}
\pgfpathlineto{\pgfpoint{-18.790356\du}{6.372085\du}}
\pgfpathlineto{\pgfpoint{-18.800208\du}{6.372085\du}}
\pgfpathlineto{\pgfpoint{-18.800208\du}{6.372085\du}}
\pgfpathlineto{\pgfpoint{-18.809165\du}{6.372085\du}}
\pgfpathlineto{\pgfpoint{-18.819017\du}{6.372384\du}}
\pgfpathlineto{\pgfpoint{-18.828571\du}{6.372981\du}}
\pgfpathlineto{\pgfpoint{-18.837826\du}{6.373279\du}}
\pgfpathlineto{\pgfpoint{-18.847380\du}{6.373876\du}}
\pgfpathlineto{\pgfpoint{-18.856933\du}{6.374772\du}}
\pgfpathlineto{\pgfpoint{-18.865890\du}{6.375668\du}}
\pgfpathlineto{\pgfpoint{-18.874846\du}{6.377160\du}}
\pgfpathlineto{\pgfpoint{-18.883504\du}{6.377459\du}}
\pgfpathlineto{\pgfpoint{-18.891864\du}{6.379847\du}}
\pgfpathlineto{\pgfpoint{-18.900522\du}{6.381042\du}}
\pgfpathlineto{\pgfpoint{-18.908583\du}{6.382833\du}}
\pgfpathlineto{\pgfpoint{-18.916345\du}{6.384624\du}}
\pgfpathlineto{\pgfpoint{-18.923809\du}{6.386714\du}}
\pgfpathlineto{\pgfpoint{-18.930377\du}{6.388207\du}}
\pgfpathlineto{\pgfpoint{-18.937841\du}{6.390595\du}}
\pgfpathlineto{\pgfpoint{-18.944110\du}{6.392387\du}}
\pgfpathlineto{\pgfpoint{-18.950082\du}{6.395074\du}}
\pgfpathlineto{\pgfpoint{-18.956053\du}{6.397761\du}}
\pgfpathlineto{\pgfpoint{-18.961128\du}{6.399850\du}}
\pgfpathlineto{\pgfpoint{-18.966203\du}{6.402836\du}}
\pgfpathlineto{\pgfpoint{-18.970980\du}{6.405224\du}}
\pgfpathlineto{\pgfpoint{-18.975160\du}{6.407911\du}}
\pgfpathlineto{\pgfpoint{-18.978444\du}{6.410598\du}}
\pgfpathlineto{\pgfpoint{-18.981728\du}{6.412987\du}}
\pgfpathlineto{\pgfpoint{-18.984116\du}{6.416569\du}}
\pgfpathlineto{\pgfpoint{-18.985908\du}{6.419256\du}}
\pgfpathlineto{\pgfpoint{-18.988296\du}{6.422242\du}}
\pgfpathlineto{\pgfpoint{-18.989192\du}{6.425526\du}}
\pgfpathlineto{\pgfpoint{-18.990087\du}{6.427914\du}}
\pgfpathlineto{\pgfpoint{-18.990087\du}{6.431497\du}}
\pgfpathlineto{\pgfpoint{-18.990087\du}{6.431497\du}}
\pgfpathlineto{\pgfpoint{-18.990087\du}{6.434482\du}}
\pgfpathlineto{\pgfpoint{-18.989192\du}{6.437169\du}}
\pgfpathlineto{\pgfpoint{-18.988296\du}{6.440155\du}}
\pgfpathlineto{\pgfpoint{-18.985908\du}{6.443738\du}}
\pgfpathlineto{\pgfpoint{-18.984116\du}{6.446424\du}}
\pgfpathlineto{\pgfpoint{-18.981728\du}{6.449410\du}}
\pgfpathlineto{\pgfpoint{-18.978444\du}{6.452097\du}}
\pgfpathlineto{\pgfpoint{-18.975160\du}{6.454784\du}}
\pgfpathlineto{\pgfpoint{-18.970980\du}{6.457471\du}}
\pgfpathlineto{\pgfpoint{-18.966203\du}{6.460158\du}}
\pgfpathlineto{\pgfpoint{-18.961128\du}{6.462845\du}}
\pgfpathlineto{\pgfpoint{-18.956053\du}{6.465233\du}}
\pgfpathlineto{\pgfpoint{-18.950082\du}{6.467920\du}}
\pgfpathlineto{\pgfpoint{-18.944110\du}{6.470010\du}}
\pgfpathlineto{\pgfpoint{-18.937841\du}{6.472399\du}}
\pgfpathlineto{\pgfpoint{-18.930377\du}{6.474787\du}}
\pgfpathlineto{\pgfpoint{-18.923809\du}{6.476578\du}}
\pgfpathlineto{\pgfpoint{-18.916345\du}{6.478071\du}}
\pgfpathlineto{\pgfpoint{-18.908583\du}{6.480161\du}}
\pgfpathlineto{\pgfpoint{-18.900522\du}{6.481952\du}}
\pgfpathlineto{\pgfpoint{-18.891864\du}{6.483146\du}}
\pgfpathlineto{\pgfpoint{-18.883504\du}{6.484938\du}}
\pgfpathlineto{\pgfpoint{-18.874846\du}{6.485833\du}}
\pgfpathlineto{\pgfpoint{-18.865890\du}{6.487326\du}}
\pgfpathlineto{\pgfpoint{-18.856933\du}{6.488222\du}}
\pgfpathlineto{\pgfpoint{-18.847380\du}{6.488520\du}}
\pgfpathlineto{\pgfpoint{-18.837826\du}{6.489715\du}}
\pgfpathlineto{\pgfpoint{-18.828571\du}{6.490013\du}}
\pgfpathlineto{\pgfpoint{-18.819017\du}{6.490312\du}}
\pgfpathlineto{\pgfpoint{-18.809165\du}{6.490909\du}}
\pgfpathlineto{\pgfpoint{-18.800208\du}{6.490909\du}}
\pgfpathlineto{\pgfpoint{-18.800208\du}{6.490909\du}}
\pgfpathlineto{\pgfpoint{-18.790356\du}{6.490909\du}}
\pgfpathlineto{\pgfpoint{-18.780205\du}{6.490312\du}}
\pgfpathlineto{\pgfpoint{-18.770652\du}{6.490013\du}}
\pgfpathlineto{\pgfpoint{-18.761098\du}{6.489715\du}}
\pgfpathlineto{\pgfpoint{-18.751843\du}{6.488520\du}}
\pgfpathlineto{\pgfpoint{-18.742289\du}{6.488222\du}}
\pgfpathlineto{\pgfpoint{-18.733631\du}{6.487326\du}}
\pgfpathlineto{\pgfpoint{-18.724675\du}{6.485833\du}}
\pgfpathlineto{\pgfpoint{-18.716017\du}{6.484938\du}}
\pgfpathlineto{\pgfpoint{-18.707657\du}{6.483146\du}}
\pgfpathlineto{\pgfpoint{-18.698999\du}{6.481952\du}}
\pgfpathlineto{\pgfpoint{-18.690640\du}{6.480161\du}}
\pgfpathlineto{\pgfpoint{-18.683176\du}{6.478071\du}}
\pgfpathlineto{\pgfpoint{-18.675414\du}{6.476578\du}}
\pgfpathlineto{\pgfpoint{-18.668547\du}{6.474787\du}}
\pgfpathlineto{\pgfpoint{-18.661680\du}{6.472399\du}}
\pgfpathlineto{\pgfpoint{-18.655112\du}{6.470010\du}}
\pgfpathlineto{\pgfpoint{-18.649440\du}{6.467920\du}}
\pgfpathlineto{\pgfpoint{-18.643469\du}{6.465233\du}}
\pgfpathlineto{\pgfpoint{-18.637796\du}{6.462845\du}}
\pgfpathlineto{\pgfpoint{-18.633318\du}{6.460158\du}}
\pgfpathlineto{\pgfpoint{-18.628541\du}{6.457471\du}}
\pgfpathlineto{\pgfpoint{-18.624361\du}{6.454784\du}}
\pgfpathlineto{\pgfpoint{-18.621077\du}{6.452097\du}}
\pgfpathlineto{\pgfpoint{-18.617793\du}{6.449410\du}}
\pgfpathlineto{\pgfpoint{-18.615703\du}{6.446424\du}}
\pgfpathlineto{\pgfpoint{-18.613613\du}{6.443738\du}}
\pgfpathlineto{\pgfpoint{-18.611524\du}{6.440155\du}}
\pgfpathlineto{\pgfpoint{-18.610329\du}{6.437169\du}}
\pgfpathlineto{\pgfpoint{-18.609434\du}{6.434482\du}}
\pgfpathlineto{\pgfpoint{-18.609434\du}{6.431497\du}}
\pgfusepath{stroke}
\pgfsetbuttcap
\pgfsetmiterjoin
\pgfsetdash{}{0pt}
\definecolor{dialinecolor}{rgb}{0.000000, 0.470588, 0.666667}
\pgfsetfillcolor{dialinecolor}
\pgfpathmoveto{\pgfpoint{-18.990087\du}{6.348201\du}}
\pgfpathlineto{\pgfpoint{-18.990087\du}{6.442842\du}}
\pgfpathlineto{\pgfpoint{-18.597193\du}{6.442842\du}}
\pgfpathlineto{\pgfpoint{-18.597193\du}{6.348201\du}}
\pgfpathlineto{\pgfpoint{-18.990087\du}{6.348201\du}}
\pgfusepath{fill}
\pgfsetbuttcap
\pgfsetmiterjoin
\pgfsetdash{}{0pt}
\definecolor{dialinecolor}{rgb}{0.000000, 0.705882, 1.000000}
\pgfsetfillcolor{dialinecolor}
\pgfpathmoveto{\pgfpoint{-18.597790\du}{6.341633\du}}
\pgfpathlineto{\pgfpoint{-18.597790\du}{6.338349\du}}
\pgfpathlineto{\pgfpoint{-18.598686\du}{6.334766\du}}
\pgfpathlineto{\pgfpoint{-18.599880\du}{6.331781\du}}
\pgfpathlineto{\pgfpoint{-18.601373\du}{6.328496\du}}
\pgfpathlineto{\pgfpoint{-18.603761\du}{6.325212\du}}
\pgfpathlineto{\pgfpoint{-18.606448\du}{6.321928\du}}
\pgfpathlineto{\pgfpoint{-18.609434\du}{6.318943\du}}
\pgfpathlineto{\pgfpoint{-18.613613\du}{6.315659\du}}
\pgfpathlineto{\pgfpoint{-18.617196\du}{6.312972\du}}
\pgfpathlineto{\pgfpoint{-18.621973\du}{6.309986\du}}
\pgfpathlineto{\pgfpoint{-18.627048\du}{6.307001\du}}
\pgfpathlineto{\pgfpoint{-18.632422\du}{6.304314\du}}
\pgfpathlineto{\pgfpoint{-18.638692\du}{6.301030\du}}
\pgfpathlineto{\pgfpoint{-18.644663\du}{6.299238\du}}
\pgfpathlineto{\pgfpoint{-18.651529\du}{6.296551\du}}
\pgfpathlineto{\pgfpoint{-18.658396\du}{6.294163\du}}
\pgfpathlineto{\pgfpoint{-18.665860\du}{6.292073\du}}
\pgfpathlineto{\pgfpoint{-18.673324\du}{6.289983\du}}
\pgfpathlineto{\pgfpoint{-18.681385\du}{6.287893\du}}
\pgfpathlineto{\pgfpoint{-18.690043\du}{6.285804\du}}
\pgfpathlineto{\pgfpoint{-18.698402\du}{6.284311\du}}
\pgfpathlineto{\pgfpoint{-18.707060\du}{6.282818\du}}
\pgfpathlineto{\pgfpoint{-18.716315\du}{6.281624\du}}
\pgfpathlineto{\pgfpoint{-18.725570\du}{6.280131\du}}
\pgfpathlineto{\pgfpoint{-18.734825\du}{6.279235\du}}
\pgfpathlineto{\pgfpoint{-18.744379\du}{6.278041\du}}
\pgfpathlineto{\pgfpoint{-18.754530\du}{6.277444\du}}
\pgfpathlineto{\pgfpoint{-18.764084\du}{6.276847\du}}
\pgfpathlineto{\pgfpoint{-18.773637\du}{6.276548\du}}
\pgfpathlineto{\pgfpoint{-18.783788\du}{6.276250\du}}
\pgfpathlineto{\pgfpoint{-18.793640\du}{6.276250\du}}
\pgfpathlineto{\pgfpoint{-18.793640\du}{6.276250\du}}
\pgfpathlineto{\pgfpoint{-18.803492\du}{6.276250\du}}
\pgfpathlineto{\pgfpoint{-18.813643\du}{6.276548\du}}
\pgfpathlineto{\pgfpoint{-18.823495\du}{6.276847\du}}
\pgfpathlineto{\pgfpoint{-18.833348\du}{6.277444\du}}
\pgfpathlineto{\pgfpoint{-18.842901\du}{6.278041\du}}
\pgfpathlineto{\pgfpoint{-18.852455\du}{6.279235\du}}
\pgfpathlineto{\pgfpoint{-18.862009\du}{6.280131\du}}
\pgfpathlineto{\pgfpoint{-18.870965\du}{6.281624\du}}
\pgfpathlineto{\pgfpoint{-18.880220\du}{6.282818\du}}
\pgfpathlineto{\pgfpoint{-18.888878\du}{6.284311\du}}
\pgfpathlineto{\pgfpoint{-18.897536\du}{6.285804\du}}
\pgfpathlineto{\pgfpoint{-18.905896\du}{6.287893\du}}
\pgfpathlineto{\pgfpoint{-18.913957\du}{6.289983\du}}
\pgfpathlineto{\pgfpoint{-18.921421\du}{6.292073\du}}
\pgfpathlineto{\pgfpoint{-18.928884\du}{6.294163\du}}
\pgfpathlineto{\pgfpoint{-18.936050\du}{6.296551\du}}
\pgfpathlineto{\pgfpoint{-18.942618\du}{6.299238\du}}
\pgfpathlineto{\pgfpoint{-18.948887\du}{6.301030\du}}
\pgfpathlineto{\pgfpoint{-18.954858\du}{6.304314\du}}
\pgfpathlineto{\pgfpoint{-18.960232\du}{6.307001\du}}
\pgfpathlineto{\pgfpoint{-18.965308\du}{6.309986\du}}
\pgfpathlineto{\pgfpoint{-18.970084\du}{6.312972\du}}
\pgfpathlineto{\pgfpoint{-18.973966\du}{6.315659\du}}
\pgfpathlineto{\pgfpoint{-18.977847\du}{6.318943\du}}
\pgfpathlineto{\pgfpoint{-18.980832\du}{6.321928\du}}
\pgfpathlineto{\pgfpoint{-18.983519\du}{6.325212\du}}
\pgfpathlineto{\pgfpoint{-18.985908\du}{6.328496\du}}
\pgfpathlineto{\pgfpoint{-18.987699\du}{6.331781\du}}
\pgfpathlineto{\pgfpoint{-18.989192\du}{6.334766\du}}
\pgfpathlineto{\pgfpoint{-18.990087\du}{6.338349\du}}
\pgfpathlineto{\pgfpoint{-18.990087\du}{6.341633\du}}
\pgfpathlineto{\pgfpoint{-18.990087\du}{6.341633\du}}
\pgfpathlineto{\pgfpoint{-18.990087\du}{6.344917\du}}
\pgfpathlineto{\pgfpoint{-18.989192\du}{6.348499\du}}
\pgfpathlineto{\pgfpoint{-18.987699\du}{6.351485\du}}
\pgfpathlineto{\pgfpoint{-18.985908\du}{6.354769\du}}
\pgfpathlineto{\pgfpoint{-18.983519\du}{6.358352\du}}
\pgfpathlineto{\pgfpoint{-18.980832\du}{6.361337\du}}
\pgfpathlineto{\pgfpoint{-18.977847\du}{6.364323\du}}
\pgfpathlineto{\pgfpoint{-18.973966\du}{6.367607\du}}
\pgfpathlineto{\pgfpoint{-18.970084\du}{6.369995\du}}
\pgfpathlineto{\pgfpoint{-18.965308\du}{6.373279\du}}
\pgfpathlineto{\pgfpoint{-18.960232\du}{6.376563\du}}
\pgfpathlineto{\pgfpoint{-18.954858\du}{6.379250\du}}
\pgfpathlineto{\pgfpoint{-18.948887\du}{6.381937\du}}
\pgfpathlineto{\pgfpoint{-18.942618\du}{6.384027\du}}
\pgfpathlineto{\pgfpoint{-18.936050\du}{6.387013\du}}
\pgfpathlineto{\pgfpoint{-18.928884\du}{6.389401\du}}
\pgfpathlineto{\pgfpoint{-18.921421\du}{6.391192\du}}
\pgfpathlineto{\pgfpoint{-18.913957\du}{6.393282\du}}
\pgfpathlineto{\pgfpoint{-18.905896\du}{6.395372\du}}
\pgfpathlineto{\pgfpoint{-18.897536\du}{6.397462\du}}
\pgfpathlineto{\pgfpoint{-18.888878\du}{6.398955\du}}
\pgfpathlineto{\pgfpoint{-18.880220\du}{6.400447\du}}
\pgfpathlineto{\pgfpoint{-18.870965\du}{6.401642\du}}
\pgfpathlineto{\pgfpoint{-18.862009\du}{6.403134\du}}
\pgfpathlineto{\pgfpoint{-18.852455\du}{6.404030\du}}
\pgfpathlineto{\pgfpoint{-18.842901\du}{6.405224\du}}
\pgfpathlineto{\pgfpoint{-18.833348\du}{6.405523\du}}
\pgfpathlineto{\pgfpoint{-18.823495\du}{6.406419\du}}
\pgfpathlineto{\pgfpoint{-18.813643\du}{6.406419\du}}
\pgfpathlineto{\pgfpoint{-18.803492\du}{6.407314\du}}
\pgfpathlineto{\pgfpoint{-18.793640\du}{6.407314\du}}
\pgfpathlineto{\pgfpoint{-18.793640\du}{6.407314\du}}
\pgfpathlineto{\pgfpoint{-18.783788\du}{6.407314\du}}
\pgfpathlineto{\pgfpoint{-18.773637\du}{6.406419\du}}
\pgfpathlineto{\pgfpoint{-18.764084\du}{6.406419\du}}
\pgfpathlineto{\pgfpoint{-18.754530\du}{6.405523\du}}
\pgfpathlineto{\pgfpoint{-18.744379\du}{6.405224\du}}
\pgfpathlineto{\pgfpoint{-18.734825\du}{6.404030\du}}
\pgfpathlineto{\pgfpoint{-18.725570\du}{6.403134\du}}
\pgfpathlineto{\pgfpoint{-18.716315\du}{6.401642\du}}
\pgfpathlineto{\pgfpoint{-18.707060\du}{6.400447\du}}
\pgfpathlineto{\pgfpoint{-18.698402\du}{6.398955\du}}
\pgfpathlineto{\pgfpoint{-18.690043\du}{6.397462\du}}
\pgfpathlineto{\pgfpoint{-18.681385\du}{6.395372\du}}
\pgfpathlineto{\pgfpoint{-18.673324\du}{6.393282\du}}
\pgfpathlineto{\pgfpoint{-18.665860\du}{6.391192\du}}
\pgfpathlineto{\pgfpoint{-18.658396\du}{6.389401\du}}
\pgfpathlineto{\pgfpoint{-18.651529\du}{6.387013\du}}
\pgfpathlineto{\pgfpoint{-18.644663\du}{6.384027\du}}
\pgfpathlineto{\pgfpoint{-18.638692\du}{6.381937\du}}
\pgfpathlineto{\pgfpoint{-18.632422\du}{6.379250\du}}
\pgfpathlineto{\pgfpoint{-18.627048\du}{6.376563\du}}
\pgfpathlineto{\pgfpoint{-18.621973\du}{6.373279\du}}
\pgfpathlineto{\pgfpoint{-18.617196\du}{6.369995\du}}
\pgfpathlineto{\pgfpoint{-18.613613\du}{6.367607\du}}
\pgfpathlineto{\pgfpoint{-18.609434\du}{6.364323\du}}
\pgfpathlineto{\pgfpoint{-18.606448\du}{6.361337\du}}
\pgfpathlineto{\pgfpoint{-18.603761\du}{6.358352\du}}
\pgfpathlineto{\pgfpoint{-18.601373\du}{6.354769\du}}
\pgfpathlineto{\pgfpoint{-18.599880\du}{6.351485\du}}
\pgfpathlineto{\pgfpoint{-18.598686\du}{6.348499\du}}
\pgfpathlineto{\pgfpoint{-18.597790\du}{6.344917\du}}
\pgfpathlineto{\pgfpoint{-18.597790\du}{6.341633\du}}
\pgfusepath{fill}
\pgfsetbuttcap
\pgfsetmiterjoin
\pgfsetdash{}{0pt}
\definecolor{dialinecolor}{rgb}{0.666667, 0.901961, 1.000000}
\pgfsetstrokecolor{dialinecolor}
\pgfpathmoveto{\pgfpoint{-18.609434\du}{6.335960\du}}
\pgfpathlineto{\pgfpoint{-18.609434\du}{6.332676\du}}
\pgfpathlineto{\pgfpoint{-18.610329\du}{6.330288\du}}
\pgfpathlineto{\pgfpoint{-18.611524\du}{6.327004\du}}
\pgfpathlineto{\pgfpoint{-18.613613\du}{6.323421\du}}
\pgfpathlineto{\pgfpoint{-18.615703\du}{6.321033\du}}
\pgfpathlineto{\pgfpoint{-18.617793\du}{6.317749\du}}
\pgfpathlineto{\pgfpoint{-18.621077\du}{6.315360\du}}
\pgfpathlineto{\pgfpoint{-18.624361\du}{6.312673\du}}
\pgfpathlineto{\pgfpoint{-18.628541\du}{6.309986\du}}
\pgfpathlineto{\pgfpoint{-18.633318\du}{6.307299\du}}
\pgfpathlineto{\pgfpoint{-18.637796\du}{6.304612\du}}
\pgfpathlineto{\pgfpoint{-18.643469\du}{6.301925\du}}
\pgfpathlineto{\pgfpoint{-18.649440\du}{6.299537\du}}
\pgfpathlineto{\pgfpoint{-18.655112\du}{6.297447\du}}
\pgfpathlineto{\pgfpoint{-18.661680\du}{6.295059\du}}
\pgfpathlineto{\pgfpoint{-18.668547\du}{6.292670\du}}
\pgfpathlineto{\pgfpoint{-18.675414\du}{6.290879\du}}
\pgfpathlineto{\pgfpoint{-18.683176\du}{6.289386\du}}
\pgfpathlineto{\pgfpoint{-18.690640\du}{6.286998\du}}
\pgfpathlineto{\pgfpoint{-18.698999\du}{6.285505\du}}
\pgfpathlineto{\pgfpoint{-18.707657\du}{6.284311\du}}
\pgfpathlineto{\pgfpoint{-18.716017\du}{6.282519\du}}
\pgfpathlineto{\pgfpoint{-18.724675\du}{6.281624\du}}
\pgfpathlineto{\pgfpoint{-18.733631\du}{6.280131\du}}
\pgfpathlineto{\pgfpoint{-18.742289\du}{6.279235\du}}
\pgfpathlineto{\pgfpoint{-18.751843\du}{6.278638\du}}
\pgfpathlineto{\pgfpoint{-18.761098\du}{6.277743\du}}
\pgfpathlineto{\pgfpoint{-18.770652\du}{6.277444\du}}
\pgfpathlineto{\pgfpoint{-18.780205\du}{6.277145\du}}
\pgfpathlineto{\pgfpoint{-18.790356\du}{6.276548\du}}
\pgfpathlineto{\pgfpoint{-18.800208\du}{6.276548\du}}
\pgfpathlineto{\pgfpoint{-18.800208\du}{6.276548\du}}
\pgfpathlineto{\pgfpoint{-18.809165\du}{6.276548\du}}
\pgfpathlineto{\pgfpoint{-18.819017\du}{6.277145\du}}
\pgfpathlineto{\pgfpoint{-18.828571\du}{6.277444\du}}
\pgfpathlineto{\pgfpoint{-18.837826\du}{6.277743\du}}
\pgfpathlineto{\pgfpoint{-18.847380\du}{6.278638\du}}
\pgfpathlineto{\pgfpoint{-18.856933\du}{6.279235\du}}
\pgfpathlineto{\pgfpoint{-18.865890\du}{6.280131\du}}
\pgfpathlineto{\pgfpoint{-18.874846\du}{6.281624\du}}
\pgfpathlineto{\pgfpoint{-18.883504\du}{6.282519\du}}
\pgfpathlineto{\pgfpoint{-18.891864\du}{6.284311\du}}
\pgfpathlineto{\pgfpoint{-18.900522\du}{6.285505\du}}
\pgfpathlineto{\pgfpoint{-18.908583\du}{6.286998\du}}
\pgfpathlineto{\pgfpoint{-18.916345\du}{6.289386\du}}
\pgfpathlineto{\pgfpoint{-18.923809\du}{6.290879\du}}
\pgfpathlineto{\pgfpoint{-18.930377\du}{6.292670\du}}
\pgfpathlineto{\pgfpoint{-18.937841\du}{6.295059\du}}
\pgfpathlineto{\pgfpoint{-18.944110\du}{6.297447\du}}
\pgfpathlineto{\pgfpoint{-18.950082\du}{6.299537\du}}
\pgfpathlineto{\pgfpoint{-18.956053\du}{6.301925\du}}
\pgfpathlineto{\pgfpoint{-18.961128\du}{6.304612\du}}
\pgfpathlineto{\pgfpoint{-18.966203\du}{6.307299\du}}
\pgfpathlineto{\pgfpoint{-18.970980\du}{6.309986\du}}
\pgfpathlineto{\pgfpoint{-18.975160\du}{6.312673\du}}
\pgfpathlineto{\pgfpoint{-18.978444\du}{6.315360\du}}
\pgfpathlineto{\pgfpoint{-18.981728\du}{6.317749\du}}
\pgfpathlineto{\pgfpoint{-18.984116\du}{6.321033\du}}
\pgfpathlineto{\pgfpoint{-18.985908\du}{6.323421\du}}
\pgfpathlineto{\pgfpoint{-18.988296\du}{6.327004\du}}
\pgfpathlineto{\pgfpoint{-18.989192\du}{6.330288\du}}
\pgfpathlineto{\pgfpoint{-18.990087\du}{6.332676\du}}
\pgfpathlineto{\pgfpoint{-18.990087\du}{6.335960\du}}
\pgfpathlineto{\pgfpoint{-18.990087\du}{6.335960\du}}
\pgfpathlineto{\pgfpoint{-18.990087\du}{6.339244\du}}
\pgfpathlineto{\pgfpoint{-18.989192\du}{6.341931\du}}
\pgfpathlineto{\pgfpoint{-18.988296\du}{6.344917\du}}
\pgfpathlineto{\pgfpoint{-18.985908\du}{6.348201\du}}
\pgfpathlineto{\pgfpoint{-18.984116\du}{6.350888\du}}
\pgfpathlineto{\pgfpoint{-18.981728\du}{6.354172\du}}
\pgfpathlineto{\pgfpoint{-18.978444\du}{6.356859\du}}
\pgfpathlineto{\pgfpoint{-18.975160\du}{6.359546\du}}
\pgfpathlineto{\pgfpoint{-18.970980\du}{6.362233\du}}
\pgfpathlineto{\pgfpoint{-18.966203\du}{6.364920\du}}
\pgfpathlineto{\pgfpoint{-18.961128\du}{6.367607\du}}
\pgfpathlineto{\pgfpoint{-18.956053\du}{6.369697\du}}
\pgfpathlineto{\pgfpoint{-18.950082\du}{6.372384\du}}
\pgfpathlineto{\pgfpoint{-18.944110\du}{6.374772\du}}
\pgfpathlineto{\pgfpoint{-18.937841\du}{6.377160\du}}
\pgfpathlineto{\pgfpoint{-18.930377\du}{6.379250\du}}
\pgfpathlineto{\pgfpoint{-18.923809\du}{6.381042\du}}
\pgfpathlineto{\pgfpoint{-18.916345\du}{6.382833\du}}
\pgfpathlineto{\pgfpoint{-18.908583\du}{6.384624\du}}
\pgfpathlineto{\pgfpoint{-18.900522\du}{6.386714\du}}
\pgfpathlineto{\pgfpoint{-18.891864\du}{6.387908\du}}
\pgfpathlineto{\pgfpoint{-18.883504\du}{6.389700\du}}
\pgfpathlineto{\pgfpoint{-18.874846\du}{6.390595\du}}
\pgfpathlineto{\pgfpoint{-18.865890\du}{6.392088\du}}
\pgfpathlineto{\pgfpoint{-18.856933\du}{6.392387\du}}
\pgfpathlineto{\pgfpoint{-18.847380\du}{6.393581\du}}
\pgfpathlineto{\pgfpoint{-18.837826\du}{6.393879\du}}
\pgfpathlineto{\pgfpoint{-18.828571\du}{6.394775\du}}
\pgfpathlineto{\pgfpoint{-18.819017\du}{6.395074\du}}
\pgfpathlineto{\pgfpoint{-18.809165\du}{6.395372\du}}
\pgfpathlineto{\pgfpoint{-18.800208\du}{6.395372\du}}
\pgfpathlineto{\pgfpoint{-18.800208\du}{6.395372\du}}
\pgfpathlineto{\pgfpoint{-18.790356\du}{6.395372\du}}
\pgfpathlineto{\pgfpoint{-18.780205\du}{6.395074\du}}
\pgfpathlineto{\pgfpoint{-18.770652\du}{6.394775\du}}
\pgfpathlineto{\pgfpoint{-18.761098\du}{6.393879\du}}
\pgfpathlineto{\pgfpoint{-18.751843\du}{6.393581\du}}
\pgfpathlineto{\pgfpoint{-18.742289\du}{6.392387\du}}
\pgfpathlineto{\pgfpoint{-18.733631\du}{6.392088\du}}
\pgfpathlineto{\pgfpoint{-18.724675\du}{6.390595\du}}
\pgfpathlineto{\pgfpoint{-18.716017\du}{6.389700\du}}
\pgfpathlineto{\pgfpoint{-18.707657\du}{6.387908\du}}
\pgfpathlineto{\pgfpoint{-18.698999\du}{6.386714\du}}
\pgfpathlineto{\pgfpoint{-18.690640\du}{6.384624\du}}
\pgfpathlineto{\pgfpoint{-18.683176\du}{6.382833\du}}
\pgfpathlineto{\pgfpoint{-18.675414\du}{6.381042\du}}
\pgfpathlineto{\pgfpoint{-18.668547\du}{6.379250\du}}
\pgfpathlineto{\pgfpoint{-18.661680\du}{6.377160\du}}
\pgfpathlineto{\pgfpoint{-18.655112\du}{6.374772\du}}
\pgfpathlineto{\pgfpoint{-18.649440\du}{6.372384\du}}
\pgfpathlineto{\pgfpoint{-18.643469\du}{6.369697\du}}
\pgfpathlineto{\pgfpoint{-18.637796\du}{6.367607\du}}
\pgfpathlineto{\pgfpoint{-18.633318\du}{6.364920\du}}
\pgfpathlineto{\pgfpoint{-18.628541\du}{6.362233\du}}
\pgfpathlineto{\pgfpoint{-18.624361\du}{6.359546\du}}
\pgfpathlineto{\pgfpoint{-18.621077\du}{6.356859\du}}
\pgfpathlineto{\pgfpoint{-18.617793\du}{6.354172\du}}
\pgfpathlineto{\pgfpoint{-18.615703\du}{6.350888\du}}
\pgfpathlineto{\pgfpoint{-18.613613\du}{6.348201\du}}
\pgfpathlineto{\pgfpoint{-18.611524\du}{6.344917\du}}
\pgfpathlineto{\pgfpoint{-18.610329\du}{6.341931\du}}
\pgfpathlineto{\pgfpoint{-18.609434\du}{6.339244\du}}
\pgfpathlineto{\pgfpoint{-18.609434\du}{6.335960\du}}
\pgfusepath{stroke}
\pgfsetbuttcap
\pgfsetmiterjoin
\pgfsetdash{}{0pt}
\definecolor{dialinecolor}{rgb}{0.000000, 0.000000, 0.000000}
\pgfsetfillcolor{dialinecolor}
\pgfpathmoveto{\pgfpoint{-18.787968\du}{6.324317\du}}
\pgfpathlineto{\pgfpoint{-18.764084\du}{6.335960\du}}
\pgfpathlineto{\pgfpoint{-18.692730\du}{6.312375\du}}
\pgfpathlineto{\pgfpoint{-18.656903\du}{6.324317\du}}
\pgfpathlineto{\pgfpoint{-18.680788\du}{6.300433\du}}
\pgfpathlineto{\pgfpoint{-18.764084\du}{6.300433\du}}
\pgfpathlineto{\pgfpoint{-18.728257\du}{6.300433\du}}
\pgfpathlineto{\pgfpoint{-18.787968\du}{6.324317\du}}
\pgfusepath{fill}
\pgfsetbuttcap
\pgfsetmiterjoin
\pgfsetdash{}{0pt}
\definecolor{dialinecolor}{rgb}{0.000000, 0.000000, 0.000000}
\pgfsetfillcolor{dialinecolor}
\pgfpathmoveto{\pgfpoint{-18.800208\du}{6.359844\du}}
\pgfpathlineto{\pgfpoint{-18.835139\du}{6.348201\du}}
\pgfpathlineto{\pgfpoint{-18.895148\du}{6.372085\du}}
\pgfpathlineto{\pgfpoint{-18.930377\du}{6.359844\du}}
\pgfpathlineto{\pgfpoint{-18.919032\du}{6.395372\du}}
\pgfpathlineto{\pgfpoint{-18.835139\du}{6.395372\du}}
\pgfpathlineto{\pgfpoint{-18.871264\du}{6.384027\du}}
\pgfpathlineto{\pgfpoint{-18.800208\du}{6.359844\du}}
\pgfusepath{fill}
\pgfsetbuttcap
\pgfsetmiterjoin
\pgfsetdash{}{0pt}
\definecolor{dialinecolor}{rgb}{0.000000, 0.000000, 0.000000}
\pgfsetfillcolor{dialinecolor}
\pgfpathmoveto{\pgfpoint{-18.919032\du}{6.300433\du}}
\pgfpathlineto{\pgfpoint{-18.895148\du}{6.287893\du}}
\pgfpathlineto{\pgfpoint{-18.823495\du}{6.324317\du}}
\pgfpathlineto{\pgfpoint{-18.800208\du}{6.312375\du}}
\pgfpathlineto{\pgfpoint{-18.811852\du}{6.335960\du}}
\pgfpathlineto{\pgfpoint{-18.895148\du}{6.335960\du}}
\pgfpathlineto{\pgfpoint{-18.859322\du}{6.324317\du}}
\pgfpathlineto{\pgfpoint{-18.919032\du}{6.300433\du}}
\pgfusepath{fill}
\pgfsetbuttcap
\pgfsetmiterjoin
\pgfsetdash{}{0pt}
\definecolor{dialinecolor}{rgb}{0.000000, 0.000000, 0.000000}
\pgfsetfillcolor{dialinecolor}
\pgfpathmoveto{\pgfpoint{-18.668547\du}{6.384027\du}}
\pgfpathlineto{\pgfpoint{-18.692730\du}{6.395372\du}}
\pgfpathlineto{\pgfpoint{-18.764084\du}{6.359844\du}}
\pgfpathlineto{\pgfpoint{-18.800208\du}{6.372085\du}}
\pgfpathlineto{\pgfpoint{-18.775727\du}{6.348201\du}}
\pgfpathlineto{\pgfpoint{-18.692730\du}{6.348201\du}}
\pgfpathlineto{\pgfpoint{-18.728257\du}{6.359844\du}}
\pgfpathlineto{\pgfpoint{-18.668547\du}{6.384027\du}}
\pgfusepath{fill}
\pgfsetbuttcap
\pgfsetmiterjoin
\pgfsetdash{}{0pt}
\definecolor{dialinecolor}{rgb}{1.000000, 1.000000, 1.000000}
\pgfsetfillcolor{dialinecolor}
\pgfpathmoveto{\pgfpoint{-18.787968\du}{6.335960\du}}
\pgfpathlineto{\pgfpoint{-18.764084\du}{6.335960\du}}
\pgfpathlineto{\pgfpoint{-18.692730\du}{6.312375\du}}
\pgfpathlineto{\pgfpoint{-18.656903\du}{6.324317\du}}
\pgfpathlineto{\pgfpoint{-18.680788\du}{6.300433\du}}
\pgfpathlineto{\pgfpoint{-18.751843\du}{6.300433\du}}
\pgfpathlineto{\pgfpoint{-18.728257\du}{6.300433\du}}
\pgfpathlineto{\pgfpoint{-18.787968\du}{6.335960\du}}
\pgfusepath{fill}
\pgfsetbuttcap
\pgfsetmiterjoin
\pgfsetdash{}{0pt}
\definecolor{dialinecolor}{rgb}{1.000000, 1.000000, 1.000000}
\pgfsetfillcolor{dialinecolor}
\pgfpathmoveto{\pgfpoint{-18.800208\du}{6.359844\du}}
\pgfpathlineto{\pgfpoint{-18.823495\du}{6.348201\du}}
\pgfpathlineto{\pgfpoint{-18.895148\du}{6.372085\du}}
\pgfpathlineto{\pgfpoint{-18.930377\du}{6.372085\du}}
\pgfpathlineto{\pgfpoint{-18.906791\du}{6.395372\du}}
\pgfpathlineto{\pgfpoint{-18.823495\du}{6.395372\du}}
\pgfpathlineto{\pgfpoint{-18.859322\du}{6.384027\du}}
\pgfpathlineto{\pgfpoint{-18.800208\du}{6.359844\du}}
\pgfusepath{fill}
\pgfsetbuttcap
\pgfsetmiterjoin
\pgfsetdash{}{0pt}
\definecolor{dialinecolor}{rgb}{1.000000, 1.000000, 1.000000}
\pgfsetfillcolor{dialinecolor}
\pgfpathmoveto{\pgfpoint{-18.919032\du}{6.300433\du}}
\pgfpathlineto{\pgfpoint{-18.895148\du}{6.300433\du}}
\pgfpathlineto{\pgfpoint{-18.823495\du}{6.324317\du}}
\pgfpathlineto{\pgfpoint{-18.787968\du}{6.312375\du}}
\pgfpathlineto{\pgfpoint{-18.811852\du}{6.335960\du}}
\pgfpathlineto{\pgfpoint{-18.895148\du}{6.335960\du}}
\pgfpathlineto{\pgfpoint{-18.859322\du}{6.335960\du}}
\pgfpathlineto{\pgfpoint{-18.919032\du}{6.300433\du}}
\pgfusepath{fill}
\pgfsetbuttcap
\pgfsetmiterjoin
\pgfsetdash{}{0pt}
\definecolor{dialinecolor}{rgb}{1.000000, 1.000000, 1.000000}
\pgfsetfillcolor{dialinecolor}
\pgfpathmoveto{\pgfpoint{-18.668547\du}{6.384027\du}}
\pgfpathlineto{\pgfpoint{-18.692730\du}{6.395372\du}}
\pgfpathlineto{\pgfpoint{-18.764084\du}{6.372085\du}}
\pgfpathlineto{\pgfpoint{-18.787968\du}{6.372085\du}}
\pgfpathlineto{\pgfpoint{-18.775727\du}{6.348201\du}}
\pgfpathlineto{\pgfpoint{-18.692730\du}{6.348201\du}}
\pgfpathlineto{\pgfpoint{-18.728257\du}{6.359844\du}}
\pgfpathlineto{\pgfpoint{-18.668547\du}{6.384027\du}}
\pgfusepath{fill}
\pgfsetbuttcap
\pgfsetmiterjoin
\pgfsetdash{}{0pt}
\definecolor{dialinecolor}{rgb}{0.666667, 0.901961, 1.000000}
\pgfsetstrokecolor{dialinecolor}
\pgfpathmoveto{\pgfpoint{-18.990087\du}{6.335960\du}}
\pgfpathlineto{\pgfpoint{-18.990087\du}{6.430900\du}}
\pgfusepath{stroke}
\pgfsetbuttcap
\pgfsetmiterjoin
\pgfsetdash{}{0pt}
\definecolor{dialinecolor}{rgb}{0.666667, 0.901961, 1.000000}
\pgfsetstrokecolor{dialinecolor}
\pgfpathmoveto{\pgfpoint{-18.609434\du}{6.335960\du}}
\pgfpathlineto{\pgfpoint{-18.609434\du}{6.430900\du}}
\pgfusepath{stroke}
\pgfsetbuttcap
\pgfsetmiterjoin
\pgfsetdash{}{0pt}
\definecolor{dialinecolor}{rgb}{0.000000, 0.470588, 0.666667}
\pgfsetfillcolor{dialinecolor}
\pgfpathmoveto{\pgfpoint{-18.954858\du}{6.842006\du}}
\pgfpathlineto{\pgfpoint{-18.955157\du}{6.838423\du}}
\pgfpathlineto{\pgfpoint{-18.956053\du}{6.835139\du}}
\pgfpathlineto{\pgfpoint{-18.957247\du}{6.831855\du}}
\pgfpathlineto{\pgfpoint{-18.959038\du}{6.828870\du}}
\pgfpathlineto{\pgfpoint{-18.961128\du}{6.825586\du}}
\pgfpathlineto{\pgfpoint{-18.964412\du}{6.822301\du}}
\pgfpathlineto{\pgfpoint{-18.967099\du}{6.819316\du}}
\pgfpathlineto{\pgfpoint{-18.970980\du}{6.816032\du}}
\pgfpathlineto{\pgfpoint{-18.975160\du}{6.813345\du}}
\pgfpathlineto{\pgfpoint{-18.979638\du}{6.810061\du}}
\pgfpathlineto{\pgfpoint{-18.985012\du}{6.807374\du}}
\pgfpathlineto{\pgfpoint{-18.990087\du}{6.804687\du}}
\pgfpathlineto{\pgfpoint{-18.996059\du}{6.802000\du}}
\pgfpathlineto{\pgfpoint{-19.002030\du}{6.799313\du}}
\pgfpathlineto{\pgfpoint{-19.009195\du}{6.796925\du}}
\pgfpathlineto{\pgfpoint{-19.016062\du}{6.794536\du}}
\pgfpathlineto{\pgfpoint{-19.023525\du}{6.792745\du}}
\pgfpathlineto{\pgfpoint{-19.030989\du}{6.790356\du}}
\pgfpathlineto{\pgfpoint{-19.039050\du}{6.788565\du}}
\pgfpathlineto{\pgfpoint{-19.047409\du}{6.786177\du}}
\pgfpathlineto{\pgfpoint{-19.055769\du}{6.784684\du}}
\pgfpathlineto{\pgfpoint{-19.064725\du}{6.783490\du}}
\pgfpathlineto{\pgfpoint{-19.073682\du}{6.781997\du}}
\pgfpathlineto{\pgfpoint{-19.082937\du}{6.780803\du}}
\pgfpathlineto{\pgfpoint{-19.092789\du}{6.779609\du}}
\pgfpathlineto{\pgfpoint{-19.102044\du}{6.778713\du}}
\pgfpathlineto{\pgfpoint{-19.111897\du}{6.778116\du}}
\pgfpathlineto{\pgfpoint{-19.121450\du}{6.777519\du}}
\pgfpathlineto{\pgfpoint{-19.131303\du}{6.776922\du}}
\pgfpathlineto{\pgfpoint{-19.141752\du}{6.776623\du}}
\pgfpathlineto{\pgfpoint{-19.151007\du}{6.776623\du}}
\pgfpathlineto{\pgfpoint{-19.151007\du}{6.776623\du}}
\pgfpathlineto{\pgfpoint{-19.160859\du}{6.776623\du}}
\pgfpathlineto{\pgfpoint{-19.171010\du}{6.776922\du}}
\pgfpathlineto{\pgfpoint{-19.180862\du}{6.777519\du}}
\pgfpathlineto{\pgfpoint{-19.190714\du}{6.778116\du}}
\pgfpathlineto{\pgfpoint{-19.200567\du}{6.778713\du}}
\pgfpathlineto{\pgfpoint{-19.209822\du}{6.779609\du}}
\pgfpathlineto{\pgfpoint{-19.219674\du}{6.780803\du}}
\pgfpathlineto{\pgfpoint{-19.228929\du}{6.781997\du}}
\pgfpathlineto{\pgfpoint{-19.237886\du}{6.783490\du}}
\pgfpathlineto{\pgfpoint{-19.246544\du}{6.784684\du}}
\pgfpathlineto{\pgfpoint{-19.254903\du}{6.786177\du}}
\pgfpathlineto{\pgfpoint{-19.263263\du}{6.788565\du}}
\pgfpathlineto{\pgfpoint{-19.271323\du}{6.790356\du}}
\pgfpathlineto{\pgfpoint{-19.278787\du}{6.792745\du}}
\pgfpathlineto{\pgfpoint{-19.286251\du}{6.794536\du}}
\pgfpathlineto{\pgfpoint{-19.293416\du}{6.796925\du}}
\pgfpathlineto{\pgfpoint{-19.300283\du}{6.799313\du}}
\pgfpathlineto{\pgfpoint{-19.306254\du}{6.802000\du}}
\pgfpathlineto{\pgfpoint{-19.312524\du}{6.804687\du}}
\pgfpathlineto{\pgfpoint{-19.317599\du}{6.807374\du}}
\pgfpathlineto{\pgfpoint{-19.322674\du}{6.810061\du}}
\pgfpathlineto{\pgfpoint{-19.327451\du}{6.813345\du}}
\pgfpathlineto{\pgfpoint{-19.331631\du}{6.816032\du}}
\pgfpathlineto{\pgfpoint{-19.335214\du}{6.819316\du}}
\pgfpathlineto{\pgfpoint{-19.338199\du}{6.822301\du}}
\pgfpathlineto{\pgfpoint{-19.340886\du}{6.825586\du}}
\pgfpathlineto{\pgfpoint{-19.343275\du}{6.828870\du}}
\pgfpathlineto{\pgfpoint{-19.345066\du}{6.831855\du}}
\pgfpathlineto{\pgfpoint{-19.346559\du}{6.835139\du}}
\pgfpathlineto{\pgfpoint{-19.347454\du}{6.838423\du}}
\pgfpathlineto{\pgfpoint{-19.347454\du}{6.842006\du}}
\pgfpathlineto{\pgfpoint{-19.347454\du}{6.842006\du}}
\pgfpathlineto{\pgfpoint{-19.347454\du}{6.845290\du}}
\pgfpathlineto{\pgfpoint{-19.346559\du}{6.848574\du}}
\pgfpathlineto{\pgfpoint{-19.345066\du}{6.851858\du}}
\pgfpathlineto{\pgfpoint{-19.343275\du}{6.855142\du}}
\pgfpathlineto{\pgfpoint{-19.340886\du}{6.858128\du}}
\pgfpathlineto{\pgfpoint{-19.338199\du}{6.861412\du}}
\pgfpathlineto{\pgfpoint{-19.335214\du}{6.864397\du}}
\pgfpathlineto{\pgfpoint{-19.331631\du}{6.867681\du}}
\pgfpathlineto{\pgfpoint{-19.327451\du}{6.870667\du}}
\pgfpathlineto{\pgfpoint{-19.322674\du}{6.873652\du}}
\pgfpathlineto{\pgfpoint{-19.317599\du}{6.876339\du}}
\pgfpathlineto{\pgfpoint{-19.312524\du}{6.879026\du}}
\pgfpathlineto{\pgfpoint{-19.306254\du}{6.881713\du}}
\pgfpathlineto{\pgfpoint{-19.300283\du}{6.884102\du}}
\pgfpathlineto{\pgfpoint{-19.293416\du}{6.886789\du}}
\pgfpathlineto{\pgfpoint{-19.286251\du}{6.889177\du}}
\pgfpathlineto{\pgfpoint{-19.278787\du}{6.891267\du}}
\pgfpathlineto{\pgfpoint{-19.271323\du}{6.893655\du}}
\pgfpathlineto{\pgfpoint{-19.263263\du}{6.895148\du}}
\pgfpathlineto{\pgfpoint{-19.254903\du}{6.897238\du}}
\pgfpathlineto{\pgfpoint{-19.246544\du}{6.899029\du}}
\pgfpathlineto{\pgfpoint{-19.237886\du}{6.900224\du}}
\pgfpathlineto{\pgfpoint{-19.228929\du}{6.901716\du}}
\pgfpathlineto{\pgfpoint{-19.219674\du}{6.902911\du}}
\pgfpathlineto{\pgfpoint{-19.209822\du}{6.904105\du}}
\pgfpathlineto{\pgfpoint{-19.200567\du}{6.905000\du}}
\pgfpathlineto{\pgfpoint{-19.190714\du}{6.905598\du}}
\pgfpathlineto{\pgfpoint{-19.180862\du}{6.906493\du}}
\pgfpathlineto{\pgfpoint{-19.171010\du}{6.906493\du}}
\pgfpathlineto{\pgfpoint{-19.160859\du}{6.907090\du}}
\pgfpathlineto{\pgfpoint{-19.151007\du}{6.907090\du}}
\pgfpathlineto{\pgfpoint{-19.151007\du}{6.907090\du}}
\pgfpathlineto{\pgfpoint{-19.141752\du}{6.907090\du}}
\pgfpathlineto{\pgfpoint{-19.131303\du}{6.906493\du}}
\pgfpathlineto{\pgfpoint{-19.121450\du}{6.906493\du}}
\pgfpathlineto{\pgfpoint{-19.111897\du}{6.905598\du}}
\pgfpathlineto{\pgfpoint{-19.102044\du}{6.905000\du}}
\pgfpathlineto{\pgfpoint{-19.092789\du}{6.904105\du}}
\pgfpathlineto{\pgfpoint{-19.082937\du}{6.902911\du}}
\pgfpathlineto{\pgfpoint{-19.073682\du}{6.901716\du}}
\pgfpathlineto{\pgfpoint{-19.064725\du}{6.900224\du}}
\pgfpathlineto{\pgfpoint{-19.055769\du}{6.899029\du}}
\pgfpathlineto{\pgfpoint{-19.047409\du}{6.897238\du}}
\pgfpathlineto{\pgfpoint{-19.039050\du}{6.895148\du}}
\pgfpathlineto{\pgfpoint{-19.030989\du}{6.893655\du}}
\pgfpathlineto{\pgfpoint{-19.023525\du}{6.891267\du}}
\pgfpathlineto{\pgfpoint{-19.016062\du}{6.889177\du}}
\pgfpathlineto{\pgfpoint{-19.009195\du}{6.886789\du}}
\pgfpathlineto{\pgfpoint{-19.002030\du}{6.884102\du}}
\pgfpathlineto{\pgfpoint{-18.996059\du}{6.881713\du}}
\pgfpathlineto{\pgfpoint{-18.990087\du}{6.879026\du}}
\pgfpathlineto{\pgfpoint{-18.985012\du}{6.876339\du}}
\pgfpathlineto{\pgfpoint{-18.979638\du}{6.873652\du}}
\pgfpathlineto{\pgfpoint{-18.975160\du}{6.870667\du}}
\pgfpathlineto{\pgfpoint{-18.970980\du}{6.867681\du}}
\pgfpathlineto{\pgfpoint{-18.967099\du}{6.864397\du}}
\pgfpathlineto{\pgfpoint{-18.964412\du}{6.861412\du}}
\pgfpathlineto{\pgfpoint{-18.961128\du}{6.858128\du}}
\pgfpathlineto{\pgfpoint{-18.959038\du}{6.855142\du}}
\pgfpathlineto{\pgfpoint{-18.957247\du}{6.851858\du}}
\pgfpathlineto{\pgfpoint{-18.956053\du}{6.848574\du}}
\pgfpathlineto{\pgfpoint{-18.955157\du}{6.845290\du}}
\pgfpathlineto{\pgfpoint{-18.954858\du}{6.842006\du}}
\pgfusepath{fill}
\pgfsetbuttcap
\pgfsetmiterjoin
\pgfsetdash{}{0pt}
\definecolor{dialinecolor}{rgb}{0.666667, 0.901961, 1.000000}
\pgfsetstrokecolor{dialinecolor}
\pgfpathmoveto{\pgfpoint{-18.966800\du}{6.835736\du}}
\pgfpathlineto{\pgfpoint{-18.967099\du}{6.832751\du}}
\pgfpathlineto{\pgfpoint{-18.967696\du}{6.830064\du}}
\pgfpathlineto{\pgfpoint{-18.968890\du}{6.827078\du}}
\pgfpathlineto{\pgfpoint{-18.970980\du}{6.823496\du}}
\pgfpathlineto{\pgfpoint{-18.973070\du}{6.821107\du}}
\pgfpathlineto{\pgfpoint{-18.975160\du}{6.817823\du}}
\pgfpathlineto{\pgfpoint{-18.978444\du}{6.815136\du}}
\pgfpathlineto{\pgfpoint{-18.981728\du}{6.812449\du}}
\pgfpathlineto{\pgfpoint{-18.985908\du}{6.809464\du}}
\pgfpathlineto{\pgfpoint{-18.990685\du}{6.807075\du}}
\pgfpathlineto{\pgfpoint{-18.995760\du}{6.804388\du}}
\pgfpathlineto{\pgfpoint{-19.000835\du}{6.802000\du}}
\pgfpathlineto{\pgfpoint{-19.006508\du}{6.799612\du}}
\pgfpathlineto{\pgfpoint{-19.012479\du}{6.797223\du}}
\pgfpathlineto{\pgfpoint{-19.019047\du}{6.794536\du}}
\pgfpathlineto{\pgfpoint{-19.025914\du}{6.793043\du}}
\pgfpathlineto{\pgfpoint{-19.032482\du}{6.790954\du}}
\pgfpathlineto{\pgfpoint{-19.040543\du}{6.789162\du}}
\pgfpathlineto{\pgfpoint{-19.048007\du}{6.787669\du}}
\pgfpathlineto{\pgfpoint{-19.055769\du}{6.785580\du}}
\pgfpathlineto{\pgfpoint{-19.064725\du}{6.784087\du}}
\pgfpathlineto{\pgfpoint{-19.073085\du}{6.782594\du}}
\pgfpathlineto{\pgfpoint{-19.081444\du}{6.781400\du}}
\pgfpathlineto{\pgfpoint{-19.090700\du}{6.780206\du}}
\pgfpathlineto{\pgfpoint{-19.099358\du}{6.779310\du}}
\pgfpathlineto{\pgfpoint{-19.108911\du}{6.778414\du}}
\pgfpathlineto{\pgfpoint{-19.118465\du}{6.778116\du}}
\pgfpathlineto{\pgfpoint{-19.127720\du}{6.777519\du}}
\pgfpathlineto{\pgfpoint{-19.137572\du}{6.776922\du}}
\pgfpathlineto{\pgfpoint{-19.147424\du}{6.776623\du}}
\pgfpathlineto{\pgfpoint{-19.156978\du}{6.776623\du}}
\pgfpathlineto{\pgfpoint{-19.156978\du}{6.776623\du}}
\pgfpathlineto{\pgfpoint{-19.166233\du}{6.776623\du}}
\pgfpathlineto{\pgfpoint{-19.176085\du}{6.776922\du}}
\pgfpathlineto{\pgfpoint{-19.185639\du}{6.777519\du}}
\pgfpathlineto{\pgfpoint{-19.195193\du}{6.778116\du}}
\pgfpathlineto{\pgfpoint{-19.204746\du}{6.778414\du}}
\pgfpathlineto{\pgfpoint{-19.214001\du}{6.779310\du}}
\pgfpathlineto{\pgfpoint{-19.222958\du}{6.780206\du}}
\pgfpathlineto{\pgfpoint{-19.232213\du}{6.781400\du}}
\pgfpathlineto{\pgfpoint{-19.240573\du}{6.782594\du}}
\pgfpathlineto{\pgfpoint{-19.249529\du}{6.784087\du}}
\pgfpathlineto{\pgfpoint{-19.257889\du}{6.785580\du}}
\pgfpathlineto{\pgfpoint{-19.265651\du}{6.787669\du}}
\pgfpathlineto{\pgfpoint{-19.273712\du}{6.789162\du}}
\pgfpathlineto{\pgfpoint{-19.281176\du}{6.790954\du}}
\pgfpathlineto{\pgfpoint{-19.287744\du}{6.793043\du}}
\pgfpathlineto{\pgfpoint{-19.295208\du}{6.794536\du}}
\pgfpathlineto{\pgfpoint{-19.301179\du}{6.797223\du}}
\pgfpathlineto{\pgfpoint{-19.307150\du}{6.799612\du}}
\pgfpathlineto{\pgfpoint{-19.313419\du}{6.802000\du}}
\pgfpathlineto{\pgfpoint{-19.318495\du}{6.804388\du}}
\pgfpathlineto{\pgfpoint{-19.323570\du}{6.807075\du}}
\pgfpathlineto{\pgfpoint{-19.327750\du}{6.809464\du}}
\pgfpathlineto{\pgfpoint{-19.331930\du}{6.812449\du}}
\pgfpathlineto{\pgfpoint{-19.335811\du}{6.815136\du}}
\pgfpathlineto{\pgfpoint{-19.339095\du}{6.817823\du}}
\pgfpathlineto{\pgfpoint{-19.340886\du}{6.821107\du}}
\pgfpathlineto{\pgfpoint{-19.343275\du}{6.823496\du}}
\pgfpathlineto{\pgfpoint{-19.345066\du}{6.827078\du}}
\pgfpathlineto{\pgfpoint{-19.346559\du}{6.830064\du}}
\pgfpathlineto{\pgfpoint{-19.346857\du}{6.832751\du}}
\pgfpathlineto{\pgfpoint{-19.347454\du}{6.835736\du}}
\pgfpathlineto{\pgfpoint{-19.347454\du}{6.835736\du}}
\pgfpathlineto{\pgfpoint{-19.346857\du}{6.839319\du}}
\pgfpathlineto{\pgfpoint{-19.346559\du}{6.842006\du}}
\pgfpathlineto{\pgfpoint{-19.345066\du}{6.844991\du}}
\pgfpathlineto{\pgfpoint{-19.343275\du}{6.847977\du}}
\pgfpathlineto{\pgfpoint{-19.340886\du}{6.850664\du}}
\pgfpathlineto{\pgfpoint{-19.339095\du}{6.853948\du}}
\pgfpathlineto{\pgfpoint{-19.335811\du}{6.856635\du}}
\pgfpathlineto{\pgfpoint{-19.331930\du}{6.859620\du}}
\pgfpathlineto{\pgfpoint{-19.327750\du}{6.862307\du}}
\pgfpathlineto{\pgfpoint{-19.323570\du}{6.864994\du}}
\pgfpathlineto{\pgfpoint{-19.318495\du}{6.867681\du}}
\pgfpathlineto{\pgfpoint{-19.313419\du}{6.869174\du}}
\pgfpathlineto{\pgfpoint{-19.307150\du}{6.872458\du}}
\pgfpathlineto{\pgfpoint{-19.301179\du}{6.874548\du}}
\pgfpathlineto{\pgfpoint{-19.295208\du}{6.876638\du}}
\pgfpathlineto{\pgfpoint{-19.287744\du}{6.879026\du}}
\pgfpathlineto{\pgfpoint{-19.281176\du}{6.881116\du}}
\pgfpathlineto{\pgfpoint{-19.273712\du}{6.882908\du}}
\pgfpathlineto{\pgfpoint{-19.265651\du}{6.884102\du}}
\pgfpathlineto{\pgfpoint{-19.257889\du}{6.886490\du}}
\pgfpathlineto{\pgfpoint{-19.249529\du}{6.887983\du}}
\pgfpathlineto{\pgfpoint{-19.240573\du}{6.889476\du}}
\pgfpathlineto{\pgfpoint{-19.232213\du}{6.890670\du}}
\pgfpathlineto{\pgfpoint{-19.222958\du}{6.891566\du}}
\pgfpathlineto{\pgfpoint{-19.214001\du}{6.892461\du}}
\pgfpathlineto{\pgfpoint{-19.204746\du}{6.893655\du}}
\pgfpathlineto{\pgfpoint{-19.195193\du}{6.893954\du}}
\pgfpathlineto{\pgfpoint{-19.185639\du}{6.894551\du}}
\pgfpathlineto{\pgfpoint{-19.176085\du}{6.894850\du}}
\pgfpathlineto{\pgfpoint{-19.166233\du}{6.895148\du}}
\pgfpathlineto{\pgfpoint{-19.156978\du}{6.895148\du}}
\pgfpathlineto{\pgfpoint{-19.156978\du}{6.895148\du}}
\pgfpathlineto{\pgfpoint{-19.147424\du}{6.895148\du}}
\pgfpathlineto{\pgfpoint{-19.137572\du}{6.894850\du}}
\pgfpathlineto{\pgfpoint{-19.127720\du}{6.894551\du}}
\pgfpathlineto{\pgfpoint{-19.118465\du}{6.893954\du}}
\pgfpathlineto{\pgfpoint{-19.108911\du}{6.893655\du}}
\pgfpathlineto{\pgfpoint{-19.099358\du}{6.892461\du}}
\pgfpathlineto{\pgfpoint{-19.090700\du}{6.891566\du}}
\pgfpathlineto{\pgfpoint{-19.081444\du}{6.890670\du}}
\pgfpathlineto{\pgfpoint{-19.073085\du}{6.889476\du}}
\pgfpathlineto{\pgfpoint{-19.064725\du}{6.887983\du}}
\pgfpathlineto{\pgfpoint{-19.055769\du}{6.886490\du}}
\pgfpathlineto{\pgfpoint{-19.048007\du}{6.884102\du}}
\pgfpathlineto{\pgfpoint{-19.040543\du}{6.882908\du}}
\pgfpathlineto{\pgfpoint{-19.032482\du}{6.881116\du}}
\pgfpathlineto{\pgfpoint{-19.025914\du}{6.879026\du}}
\pgfpathlineto{\pgfpoint{-19.019047\du}{6.876638\du}}
\pgfpathlineto{\pgfpoint{-19.012479\du}{6.874548\du}}
\pgfpathlineto{\pgfpoint{-19.006508\du}{6.872458\du}}
\pgfpathlineto{\pgfpoint{-19.000835\du}{6.869174\du}}
\pgfpathlineto{\pgfpoint{-18.995760\du}{6.867681\du}}
\pgfpathlineto{\pgfpoint{-18.990685\du}{6.864994\du}}
\pgfpathlineto{\pgfpoint{-18.985908\du}{6.862307\du}}
\pgfpathlineto{\pgfpoint{-18.981728\du}{6.859620\du}}
\pgfpathlineto{\pgfpoint{-18.978444\du}{6.856635\du}}
\pgfpathlineto{\pgfpoint{-18.975160\du}{6.853948\du}}
\pgfpathlineto{\pgfpoint{-18.973070\du}{6.850664\du}}
\pgfpathlineto{\pgfpoint{-18.970980\du}{6.847977\du}}
\pgfpathlineto{\pgfpoint{-18.968890\du}{6.844991\du}}
\pgfpathlineto{\pgfpoint{-18.967696\du}{6.842006\du}}
\pgfpathlineto{\pgfpoint{-18.967099\du}{6.839319\du}}
\pgfpathlineto{\pgfpoint{-18.966800\du}{6.835736\du}}
\pgfusepath{stroke}
\pgfsetbuttcap
\pgfsetmiterjoin
\pgfsetdash{}{0pt}
\definecolor{dialinecolor}{rgb}{0.000000, 0.470588, 0.666667}
\pgfsetfillcolor{dialinecolor}
\pgfpathmoveto{\pgfpoint{-19.347454\du}{6.752739\du}}
\pgfpathlineto{\pgfpoint{-19.347454\du}{6.847977\du}}
\pgfpathlineto{\pgfpoint{-18.954560\du}{6.847977\du}}
\pgfpathlineto{\pgfpoint{-18.954560\du}{6.752739\du}}
\pgfpathlineto{\pgfpoint{-19.347454\du}{6.752739\du}}
\pgfusepath{fill}
\pgfsetbuttcap
\pgfsetmiterjoin
\pgfsetdash{}{0pt}
\definecolor{dialinecolor}{rgb}{0.000000, 0.705882, 1.000000}
\pgfsetfillcolor{dialinecolor}
\pgfpathmoveto{\pgfpoint{-18.954858\du}{6.746469\du}}
\pgfpathlineto{\pgfpoint{-18.955157\du}{6.742887\du}}
\pgfpathlineto{\pgfpoint{-18.956053\du}{6.739603\du}}
\pgfpathlineto{\pgfpoint{-18.957247\du}{6.736617\du}}
\pgfpathlineto{\pgfpoint{-18.959038\du}{6.733034\du}}
\pgfpathlineto{\pgfpoint{-18.961128\du}{6.729750\du}}
\pgfpathlineto{\pgfpoint{-18.964412\du}{6.726466\du}}
\pgfpathlineto{\pgfpoint{-18.967099\du}{6.723779\du}}
\pgfpathlineto{\pgfpoint{-18.970980\du}{6.719898\du}}
\pgfpathlineto{\pgfpoint{-18.975160\du}{6.717510\du}}
\pgfpathlineto{\pgfpoint{-18.979638\du}{6.714524\du}}
\pgfpathlineto{\pgfpoint{-18.985012\du}{6.711539\du}}
\pgfpathlineto{\pgfpoint{-18.990087\du}{6.708852\du}}
\pgfpathlineto{\pgfpoint{-18.996059\du}{6.706165\du}}
\pgfpathlineto{\pgfpoint{-19.002030\du}{6.703776\du}}
\pgfpathlineto{\pgfpoint{-19.009195\du}{6.701089\du}}
\pgfpathlineto{\pgfpoint{-19.016062\du}{6.698701\du}}
\pgfpathlineto{\pgfpoint{-19.023525\du}{6.696611\du}}
\pgfpathlineto{\pgfpoint{-19.030989\du}{6.694521\du}}
\pgfpathlineto{\pgfpoint{-19.039050\du}{6.692730\du}}
\pgfpathlineto{\pgfpoint{-19.047409\du}{6.690640\du}}
\pgfpathlineto{\pgfpoint{-19.055769\du}{6.688849\du}}
\pgfpathlineto{\pgfpoint{-19.064725\du}{6.687655\du}}
\pgfpathlineto{\pgfpoint{-19.073682\du}{6.686162\du}}
\pgfpathlineto{\pgfpoint{-19.082937\du}{6.684968\du}}
\pgfpathlineto{\pgfpoint{-19.092789\du}{6.683773\du}}
\pgfpathlineto{\pgfpoint{-19.102044\du}{6.682579\du}}
\pgfpathlineto{\pgfpoint{-19.111897\du}{6.682281\du}}
\pgfpathlineto{\pgfpoint{-19.121450\du}{6.681385\du}}
\pgfpathlineto{\pgfpoint{-19.131303\du}{6.681086\du}}
\pgfpathlineto{\pgfpoint{-19.141752\du}{6.680788\du}}
\pgfpathlineto{\pgfpoint{-19.151007\du}{6.680788\du}}
\pgfpathlineto{\pgfpoint{-19.151007\du}{6.680788\du}}
\pgfpathlineto{\pgfpoint{-19.160859\du}{6.680788\du}}
\pgfpathlineto{\pgfpoint{-19.171010\du}{6.681086\du}}
\pgfpathlineto{\pgfpoint{-19.180862\du}{6.681385\du}}
\pgfpathlineto{\pgfpoint{-19.190714\du}{6.682281\du}}
\pgfpathlineto{\pgfpoint{-19.200567\du}{6.682579\du}}
\pgfpathlineto{\pgfpoint{-19.209822\du}{6.683773\du}}
\pgfpathlineto{\pgfpoint{-19.219674\du}{6.684968\du}}
\pgfpathlineto{\pgfpoint{-19.228929\du}{6.686162\du}}
\pgfpathlineto{\pgfpoint{-19.237886\du}{6.687655\du}}
\pgfpathlineto{\pgfpoint{-19.246544\du}{6.688849\du}}
\pgfpathlineto{\pgfpoint{-19.254903\du}{6.690640\du}}
\pgfpathlineto{\pgfpoint{-19.263263\du}{6.692730\du}}
\pgfpathlineto{\pgfpoint{-19.271323\du}{6.694521\du}}
\pgfpathlineto{\pgfpoint{-19.278787\du}{6.696611\du}}
\pgfpathlineto{\pgfpoint{-19.286251\du}{6.698701\du}}
\pgfpathlineto{\pgfpoint{-19.293416\du}{6.701089\du}}
\pgfpathlineto{\pgfpoint{-19.300283\du}{6.703776\du}}
\pgfpathlineto{\pgfpoint{-19.306254\du}{6.706165\du}}
\pgfpathlineto{\pgfpoint{-19.312524\du}{6.708852\du}}
\pgfpathlineto{\pgfpoint{-19.317599\du}{6.711539\du}}
\pgfpathlineto{\pgfpoint{-19.322674\du}{6.714524\du}}
\pgfpathlineto{\pgfpoint{-19.327451\du}{6.717510\du}}
\pgfpathlineto{\pgfpoint{-19.331631\du}{6.719898\du}}
\pgfpathlineto{\pgfpoint{-19.335214\du}{6.723779\du}}
\pgfpathlineto{\pgfpoint{-19.338199\du}{6.726466\du}}
\pgfpathlineto{\pgfpoint{-19.340886\du}{6.729750\du}}
\pgfpathlineto{\pgfpoint{-19.343275\du}{6.733034\du}}
\pgfpathlineto{\pgfpoint{-19.345066\du}{6.736617\du}}
\pgfpathlineto{\pgfpoint{-19.346559\du}{6.739603\du}}
\pgfpathlineto{\pgfpoint{-19.347454\du}{6.742887\du}}
\pgfpathlineto{\pgfpoint{-19.347454\du}{6.746469\du}}
\pgfpathlineto{\pgfpoint{-19.347454\du}{6.746469\du}}
\pgfpathlineto{\pgfpoint{-19.347454\du}{6.749753\du}}
\pgfpathlineto{\pgfpoint{-19.346559\du}{6.753037\du}}
\pgfpathlineto{\pgfpoint{-19.345066\du}{6.756023\du}}
\pgfpathlineto{\pgfpoint{-19.343275\du}{6.759606\du}}
\pgfpathlineto{\pgfpoint{-19.340886\du}{6.762890\du}}
\pgfpathlineto{\pgfpoint{-19.338199\du}{6.765875\du}}
\pgfpathlineto{\pgfpoint{-19.335214\du}{6.768861\du}}
\pgfpathlineto{\pgfpoint{-19.331631\du}{6.772145\du}}
\pgfpathlineto{\pgfpoint{-19.327451\du}{6.775130\du}}
\pgfpathlineto{\pgfpoint{-19.322674\du}{6.778116\du}}
\pgfpathlineto{\pgfpoint{-19.317599\du}{6.781101\du}}
\pgfpathlineto{\pgfpoint{-19.312524\du}{6.783788\du}}
\pgfpathlineto{\pgfpoint{-19.306254\du}{6.786177\du}}
\pgfpathlineto{\pgfpoint{-19.300283\du}{6.788864\du}}
\pgfpathlineto{\pgfpoint{-19.293416\du}{6.791551\du}}
\pgfpathlineto{\pgfpoint{-19.286251\du}{6.793640\du}}
\pgfpathlineto{\pgfpoint{-19.278787\du}{6.796029\du}}
\pgfpathlineto{\pgfpoint{-19.271323\du}{6.798119\du}}
\pgfpathlineto{\pgfpoint{-19.263263\du}{6.799910\du}}
\pgfpathlineto{\pgfpoint{-19.254903\du}{6.802000\du}}
\pgfpathlineto{\pgfpoint{-19.246544\du}{6.803791\du}}
\pgfpathlineto{\pgfpoint{-19.237886\du}{6.804985\du}}
\pgfpathlineto{\pgfpoint{-19.228929\du}{6.806478\du}}
\pgfpathlineto{\pgfpoint{-19.219674\du}{6.807672\du}}
\pgfpathlineto{\pgfpoint{-19.209822\du}{6.808568\du}}
\pgfpathlineto{\pgfpoint{-19.200567\du}{6.809464\du}}
\pgfpathlineto{\pgfpoint{-19.190714\du}{6.810359\du}}
\pgfpathlineto{\pgfpoint{-19.180862\du}{6.811255\du}}
\pgfpathlineto{\pgfpoint{-19.171010\du}{6.811554\du}}
\pgfpathlineto{\pgfpoint{-19.160859\du}{6.811852\du}}
\pgfpathlineto{\pgfpoint{-19.151007\du}{6.811852\du}}
\pgfpathlineto{\pgfpoint{-19.151007\du}{6.811852\du}}
\pgfpathlineto{\pgfpoint{-19.141752\du}{6.811852\du}}
\pgfpathlineto{\pgfpoint{-19.131303\du}{6.811554\du}}
\pgfpathlineto{\pgfpoint{-19.121450\du}{6.811255\du}}
\pgfpathlineto{\pgfpoint{-19.111897\du}{6.810359\du}}
\pgfpathlineto{\pgfpoint{-19.102044\du}{6.809464\du}}
\pgfpathlineto{\pgfpoint{-19.092789\du}{6.808568\du}}
\pgfpathlineto{\pgfpoint{-19.082937\du}{6.807672\du}}
\pgfpathlineto{\pgfpoint{-19.073682\du}{6.806478\du}}
\pgfpathlineto{\pgfpoint{-19.064725\du}{6.804985\du}}
\pgfpathlineto{\pgfpoint{-19.055769\du}{6.803791\du}}
\pgfpathlineto{\pgfpoint{-19.047409\du}{6.802000\du}}
\pgfpathlineto{\pgfpoint{-19.039050\du}{6.799910\du}}
\pgfpathlineto{\pgfpoint{-19.030989\du}{6.798119\du}}
\pgfpathlineto{\pgfpoint{-19.023525\du}{6.796029\du}}
\pgfpathlineto{\pgfpoint{-19.016062\du}{6.793640\du}}
\pgfpathlineto{\pgfpoint{-19.009195\du}{6.791551\du}}
\pgfpathlineto{\pgfpoint{-19.002030\du}{6.788864\du}}
\pgfpathlineto{\pgfpoint{-18.996059\du}{6.786177\du}}
\pgfpathlineto{\pgfpoint{-18.990087\du}{6.783788\du}}
\pgfpathlineto{\pgfpoint{-18.985012\du}{6.781101\du}}
\pgfpathlineto{\pgfpoint{-18.979638\du}{6.778116\du}}
\pgfpathlineto{\pgfpoint{-18.975160\du}{6.775130\du}}
\pgfpathlineto{\pgfpoint{-18.970980\du}{6.772145\du}}
\pgfpathlineto{\pgfpoint{-18.967099\du}{6.768861\du}}
\pgfpathlineto{\pgfpoint{-18.964412\du}{6.765875\du}}
\pgfpathlineto{\pgfpoint{-18.961128\du}{6.762890\du}}
\pgfpathlineto{\pgfpoint{-18.959038\du}{6.759606\du}}
\pgfpathlineto{\pgfpoint{-18.957247\du}{6.756023\du}}
\pgfpathlineto{\pgfpoint{-18.956053\du}{6.753037\du}}
\pgfpathlineto{\pgfpoint{-18.955157\du}{6.749753\du}}
\pgfpathlineto{\pgfpoint{-18.954858\du}{6.746469\du}}
\pgfusepath{fill}
\pgfsetbuttcap
\pgfsetmiterjoin
\pgfsetdash{}{0pt}
\definecolor{dialinecolor}{rgb}{0.666667, 0.901961, 1.000000}
\pgfsetstrokecolor{dialinecolor}
\pgfpathmoveto{\pgfpoint{-18.966800\du}{6.740797\du}}
\pgfpathlineto{\pgfpoint{-18.967099\du}{6.737811\du}}
\pgfpathlineto{\pgfpoint{-18.967696\du}{6.734826\du}}
\pgfpathlineto{\pgfpoint{-18.968890\du}{6.732139\du}}
\pgfpathlineto{\pgfpoint{-18.970980\du}{6.728855\du}}
\pgfpathlineto{\pgfpoint{-18.973070\du}{6.726168\du}}
\pgfpathlineto{\pgfpoint{-18.975160\du}{6.722884\du}}
\pgfpathlineto{\pgfpoint{-18.978444\du}{6.719898\du}}
\pgfpathlineto{\pgfpoint{-18.981728\du}{6.717510\du}}
\pgfpathlineto{\pgfpoint{-18.985908\du}{6.714823\du}}
\pgfpathlineto{\pgfpoint{-18.990685\du}{6.712136\du}}
\pgfpathlineto{\pgfpoint{-18.995760\du}{6.709449\du}}
\pgfpathlineto{\pgfpoint{-19.000835\du}{6.707060\du}}
\pgfpathlineto{\pgfpoint{-19.006508\du}{6.704075\du}}
\pgfpathlineto{\pgfpoint{-19.012479\du}{6.702284\du}}
\pgfpathlineto{\pgfpoint{-19.019047\du}{6.700194\du}}
\pgfpathlineto{\pgfpoint{-19.025914\du}{6.697507\du}}
\pgfpathlineto{\pgfpoint{-19.032482\du}{6.696014\du}}
\pgfpathlineto{\pgfpoint{-19.040543\du}{6.694223\du}}
\pgfpathlineto{\pgfpoint{-19.048007\du}{6.692133\du}}
\pgfpathlineto{\pgfpoint{-19.055769\du}{6.690640\du}}
\pgfpathlineto{\pgfpoint{-19.064725\du}{6.689147\du}}
\pgfpathlineto{\pgfpoint{-19.073085\du}{6.687655\du}}
\pgfpathlineto{\pgfpoint{-19.081444\du}{6.686460\du}}
\pgfpathlineto{\pgfpoint{-19.090700\du}{6.685266\du}}
\pgfpathlineto{\pgfpoint{-19.099358\du}{6.684072\du}}
\pgfpathlineto{\pgfpoint{-19.108911\du}{6.683475\du}}
\pgfpathlineto{\pgfpoint{-19.118465\du}{6.682579\du}}
\pgfpathlineto{\pgfpoint{-19.127720\du}{6.682579\du}}
\pgfpathlineto{\pgfpoint{-19.137572\du}{6.682281\du}}
\pgfpathlineto{\pgfpoint{-19.147424\du}{6.681385\du}}
\pgfpathlineto{\pgfpoint{-19.156978\du}{6.681385\du}}
\pgfpathlineto{\pgfpoint{-19.156978\du}{6.681385\du}}
\pgfpathlineto{\pgfpoint{-19.166233\du}{6.681385\du}}
\pgfpathlineto{\pgfpoint{-19.176085\du}{6.682281\du}}
\pgfpathlineto{\pgfpoint{-19.185639\du}{6.682579\du}}
\pgfpathlineto{\pgfpoint{-19.195193\du}{6.682579\du}}
\pgfpathlineto{\pgfpoint{-19.204746\du}{6.683475\du}}
\pgfpathlineto{\pgfpoint{-19.214001\du}{6.684072\du}}
\pgfpathlineto{\pgfpoint{-19.222958\du}{6.685266\du}}
\pgfpathlineto{\pgfpoint{-19.232213\du}{6.686460\du}}
\pgfpathlineto{\pgfpoint{-19.240573\du}{6.687655\du}}
\pgfpathlineto{\pgfpoint{-19.249529\du}{6.689147\du}}
\pgfpathlineto{\pgfpoint{-19.257889\du}{6.690640\du}}
\pgfpathlineto{\pgfpoint{-19.265651\du}{6.692133\du}}
\pgfpathlineto{\pgfpoint{-19.273712\du}{6.694223\du}}
\pgfpathlineto{\pgfpoint{-19.281176\du}{6.696014\du}}
\pgfpathlineto{\pgfpoint{-19.287744\du}{6.697507\du}}
\pgfpathlineto{\pgfpoint{-19.295208\du}{6.700194\du}}
\pgfpathlineto{\pgfpoint{-19.301179\du}{6.702284\du}}
\pgfpathlineto{\pgfpoint{-19.307150\du}{6.704075\du}}
\pgfpathlineto{\pgfpoint{-19.313419\du}{6.707060\du}}
\pgfpathlineto{\pgfpoint{-19.318495\du}{6.709449\du}}
\pgfpathlineto{\pgfpoint{-19.323570\du}{6.712136\du}}
\pgfpathlineto{\pgfpoint{-19.327750\du}{6.714823\du}}
\pgfpathlineto{\pgfpoint{-19.331930\du}{6.717510\du}}
\pgfpathlineto{\pgfpoint{-19.335811\du}{6.719898\du}}
\pgfpathlineto{\pgfpoint{-19.339095\du}{6.722884\du}}
\pgfpathlineto{\pgfpoint{-19.340886\du}{6.726168\du}}
\pgfpathlineto{\pgfpoint{-19.343275\du}{6.728855\du}}
\pgfpathlineto{\pgfpoint{-19.345066\du}{6.732139\du}}
\pgfpathlineto{\pgfpoint{-19.346559\du}{6.734826\du}}
\pgfpathlineto{\pgfpoint{-19.346857\du}{6.737811\du}}
\pgfpathlineto{\pgfpoint{-19.347454\du}{6.740797\du}}
\pgfpathlineto{\pgfpoint{-19.347454\du}{6.740797\du}}
\pgfpathlineto{\pgfpoint{-19.346857\du}{6.744379\du}}
\pgfpathlineto{\pgfpoint{-19.346559\du}{6.747066\du}}
\pgfpathlineto{\pgfpoint{-19.345066\du}{6.749753\du}}
\pgfpathlineto{\pgfpoint{-19.343275\du}{6.753037\du}}
\pgfpathlineto{\pgfpoint{-19.340886\du}{6.755724\du}}
\pgfpathlineto{\pgfpoint{-19.339095\du}{6.759307\du}}
\pgfpathlineto{\pgfpoint{-19.335811\du}{6.761994\du}}
\pgfpathlineto{\pgfpoint{-19.331930\du}{6.764681\du}}
\pgfpathlineto{\pgfpoint{-19.327750\du}{6.767368\du}}
\pgfpathlineto{\pgfpoint{-19.323570\du}{6.770055\du}}
\pgfpathlineto{\pgfpoint{-19.318495\du}{6.772742\du}}
\pgfpathlineto{\pgfpoint{-19.313419\du}{6.774832\du}}
\pgfpathlineto{\pgfpoint{-19.307150\du}{6.777519\du}}
\pgfpathlineto{\pgfpoint{-19.301179\du}{6.779609\du}}
\pgfpathlineto{\pgfpoint{-19.295208\du}{6.781997\du}}
\pgfpathlineto{\pgfpoint{-19.287744\du}{6.784087\du}}
\pgfpathlineto{\pgfpoint{-19.281176\du}{6.786177\du}}
\pgfpathlineto{\pgfpoint{-19.273712\du}{6.787968\du}}
\pgfpathlineto{\pgfpoint{-19.265651\du}{6.789461\du}}
\pgfpathlineto{\pgfpoint{-19.257889\du}{6.791551\du}}
\pgfpathlineto{\pgfpoint{-19.249529\du}{6.793043\du}}
\pgfpathlineto{\pgfpoint{-19.240573\du}{6.794536\du}}
\pgfpathlineto{\pgfpoint{-19.232213\du}{6.795730\du}}
\pgfpathlineto{\pgfpoint{-19.222958\du}{6.796925\du}}
\pgfpathlineto{\pgfpoint{-19.214001\du}{6.797522\du}}
\pgfpathlineto{\pgfpoint{-19.204746\du}{6.798716\du}}
\pgfpathlineto{\pgfpoint{-19.195193\du}{6.799014\du}}
\pgfpathlineto{\pgfpoint{-19.185639\du}{6.799612\du}}
\pgfpathlineto{\pgfpoint{-19.176085\du}{6.799910\du}}
\pgfpathlineto{\pgfpoint{-19.166233\du}{6.800507\du}}
\pgfpathlineto{\pgfpoint{-19.156978\du}{6.800507\du}}
\pgfpathlineto{\pgfpoint{-19.156978\du}{6.800507\du}}
\pgfpathlineto{\pgfpoint{-19.147424\du}{6.800507\du}}
\pgfpathlineto{\pgfpoint{-19.137572\du}{6.799910\du}}
\pgfpathlineto{\pgfpoint{-19.127720\du}{6.799612\du}}
\pgfpathlineto{\pgfpoint{-19.118465\du}{6.799014\du}}
\pgfpathlineto{\pgfpoint{-19.108911\du}{6.798716\du}}
\pgfpathlineto{\pgfpoint{-19.099358\du}{6.797522\du}}
\pgfpathlineto{\pgfpoint{-19.090700\du}{6.796925\du}}
\pgfpathlineto{\pgfpoint{-19.081444\du}{6.795730\du}}
\pgfpathlineto{\pgfpoint{-19.073085\du}{6.794536\du}}
\pgfpathlineto{\pgfpoint{-19.064725\du}{6.793043\du}}
\pgfpathlineto{\pgfpoint{-19.055769\du}{6.791551\du}}
\pgfpathlineto{\pgfpoint{-19.048007\du}{6.789461\du}}
\pgfpathlineto{\pgfpoint{-19.040543\du}{6.787968\du}}
\pgfpathlineto{\pgfpoint{-19.032482\du}{6.786177\du}}
\pgfpathlineto{\pgfpoint{-19.025914\du}{6.784087\du}}
\pgfpathlineto{\pgfpoint{-19.019047\du}{6.781997\du}}
\pgfpathlineto{\pgfpoint{-19.012479\du}{6.779609\du}}
\pgfpathlineto{\pgfpoint{-19.006508\du}{6.777519\du}}
\pgfpathlineto{\pgfpoint{-19.000835\du}{6.774832\du}}
\pgfpathlineto{\pgfpoint{-18.995760\du}{6.772742\du}}
\pgfpathlineto{\pgfpoint{-18.990685\du}{6.770055\du}}
\pgfpathlineto{\pgfpoint{-18.985908\du}{6.767368\du}}
\pgfpathlineto{\pgfpoint{-18.981728\du}{6.764681\du}}
\pgfpathlineto{\pgfpoint{-18.978444\du}{6.761994\du}}
\pgfpathlineto{\pgfpoint{-18.975160\du}{6.759307\du}}
\pgfpathlineto{\pgfpoint{-18.973070\du}{6.755724\du}}
\pgfpathlineto{\pgfpoint{-18.970980\du}{6.753037\du}}
\pgfpathlineto{\pgfpoint{-18.968890\du}{6.749753\du}}
\pgfpathlineto{\pgfpoint{-18.967696\du}{6.747066\du}}
\pgfpathlineto{\pgfpoint{-18.967099\du}{6.744379\du}}
\pgfpathlineto{\pgfpoint{-18.966800\du}{6.740797\du}}
\pgfusepath{stroke}
\pgfsetbuttcap
\pgfsetmiterjoin
\pgfsetdash{}{0pt}
\definecolor{dialinecolor}{rgb}{0.000000, 0.000000, 0.000000}
\pgfsetfillcolor{dialinecolor}
\pgfpathmoveto{\pgfpoint{-19.145036\du}{6.728855\du}}
\pgfpathlineto{\pgfpoint{-19.121152\du}{6.740797\du}}
\pgfpathlineto{\pgfpoint{-19.049798\du}{6.717211\du}}
\pgfpathlineto{\pgfpoint{-19.013972\du}{6.728855\du}}
\pgfpathlineto{\pgfpoint{-19.037557\du}{6.704971\du}}
\pgfpathlineto{\pgfpoint{-19.121152\du}{6.704971\du}}
\pgfpathlineto{\pgfpoint{-19.085326\du}{6.704971\du}}
\pgfpathlineto{\pgfpoint{-19.145036\du}{6.728855\du}}
\pgfusepath{fill}
\pgfsetbuttcap
\pgfsetmiterjoin
\pgfsetdash{}{0pt}
\definecolor{dialinecolor}{rgb}{0.000000, 0.000000, 0.000000}
\pgfsetfillcolor{dialinecolor}
\pgfpathmoveto{\pgfpoint{-19.156978\du}{6.764681\du}}
\pgfpathlineto{\pgfpoint{-19.192506\du}{6.752739\du}}
\pgfpathlineto{\pgfpoint{-19.252216\du}{6.776623\du}}
\pgfpathlineto{\pgfpoint{-19.287744\du}{6.764681\du}}
\pgfpathlineto{\pgfpoint{-19.275503\du}{6.800507\du}}
\pgfpathlineto{\pgfpoint{-19.192506\du}{6.800507\du}}
\pgfpathlineto{\pgfpoint{-19.228033\du}{6.788864\du}}
\pgfpathlineto{\pgfpoint{-19.156978\du}{6.764681\du}}
\pgfusepath{fill}
\pgfsetbuttcap
\pgfsetmiterjoin
\pgfsetdash{}{0pt}
\definecolor{dialinecolor}{rgb}{0.000000, 0.000000, 0.000000}
\pgfsetfillcolor{dialinecolor}
\pgfpathmoveto{\pgfpoint{-19.275503\du}{6.704971\du}}
\pgfpathlineto{\pgfpoint{-19.252216\du}{6.693327\du}}
\pgfpathlineto{\pgfpoint{-19.180862\du}{6.728855\du}}
\pgfpathlineto{\pgfpoint{-19.156978\du}{6.717211\du}}
\pgfpathlineto{\pgfpoint{-19.168323\du}{6.740797\du}}
\pgfpathlineto{\pgfpoint{-19.252216\du}{6.740797\du}}
\pgfpathlineto{\pgfpoint{-19.216390\du}{6.728855\du}}
\pgfpathlineto{\pgfpoint{-19.275503\du}{6.704971\du}}
\pgfusepath{fill}
\pgfsetbuttcap
\pgfsetmiterjoin
\pgfsetdash{}{0pt}
\definecolor{dialinecolor}{rgb}{0.000000, 0.000000, 0.000000}
\pgfsetfillcolor{dialinecolor}
\pgfpathmoveto{\pgfpoint{-19.025914\du}{6.788864\du}}
\pgfpathlineto{\pgfpoint{-19.049798\du}{6.800507\du}}
\pgfpathlineto{\pgfpoint{-19.121152\du}{6.764681\du}}
\pgfpathlineto{\pgfpoint{-19.156978\du}{6.776623\du}}
\pgfpathlineto{\pgfpoint{-19.132795\du}{6.752739\du}}
\pgfpathlineto{\pgfpoint{-19.049798\du}{6.752739\du}}
\pgfpathlineto{\pgfpoint{-19.085326\du}{6.764681\du}}
\pgfpathlineto{\pgfpoint{-19.025914\du}{6.788864\du}}
\pgfusepath{fill}
\pgfsetbuttcap
\pgfsetmiterjoin
\pgfsetdash{}{0pt}
\definecolor{dialinecolor}{rgb}{1.000000, 1.000000, 1.000000}
\pgfsetfillcolor{dialinecolor}
\pgfpathmoveto{\pgfpoint{-19.145036\du}{6.740797\du}}
\pgfpathlineto{\pgfpoint{-19.121152\du}{6.740797\du}}
\pgfpathlineto{\pgfpoint{-19.049798\du}{6.717211\du}}
\pgfpathlineto{\pgfpoint{-19.013972\du}{6.728855\du}}
\pgfpathlineto{\pgfpoint{-19.037557\du}{6.704971\du}}
\pgfpathlineto{\pgfpoint{-19.108911\du}{6.704971\du}}
\pgfpathlineto{\pgfpoint{-19.085326\du}{6.704971\du}}
\pgfpathlineto{\pgfpoint{-19.145036\du}{6.740797\du}}
\pgfusepath{fill}
\pgfsetbuttcap
\pgfsetmiterjoin
\pgfsetdash{}{0pt}
\definecolor{dialinecolor}{rgb}{1.000000, 1.000000, 1.000000}
\pgfsetfillcolor{dialinecolor}
\pgfpathmoveto{\pgfpoint{-19.156978\du}{6.764681\du}}
\pgfpathlineto{\pgfpoint{-19.180862\du}{6.752739\du}}
\pgfpathlineto{\pgfpoint{-19.252216\du}{6.776623\du}}
\pgfpathlineto{\pgfpoint{-19.287744\du}{6.776623\du}}
\pgfpathlineto{\pgfpoint{-19.263860\du}{6.800507\du}}
\pgfpathlineto{\pgfpoint{-19.180862\du}{6.800507\du}}
\pgfpathlineto{\pgfpoint{-19.216390\du}{6.788864\du}}
\pgfpathlineto{\pgfpoint{-19.156978\du}{6.764681\du}}
\pgfusepath{fill}
\pgfsetbuttcap
\pgfsetmiterjoin
\pgfsetdash{}{0pt}
\definecolor{dialinecolor}{rgb}{1.000000, 1.000000, 1.000000}
\pgfsetfillcolor{dialinecolor}
\pgfpathmoveto{\pgfpoint{-19.275503\du}{6.704971\du}}
\pgfpathlineto{\pgfpoint{-19.252216\du}{6.704971\du}}
\pgfpathlineto{\pgfpoint{-19.180862\du}{6.728855\du}}
\pgfpathlineto{\pgfpoint{-19.145036\du}{6.717211\du}}
\pgfpathlineto{\pgfpoint{-19.168323\du}{6.740797\du}}
\pgfpathlineto{\pgfpoint{-19.252216\du}{6.740797\du}}
\pgfpathlineto{\pgfpoint{-19.216390\du}{6.740797\du}}
\pgfpathlineto{\pgfpoint{-19.275503\du}{6.704971\du}}
\pgfusepath{fill}
\pgfsetbuttcap
\pgfsetmiterjoin
\pgfsetdash{}{0pt}
\definecolor{dialinecolor}{rgb}{1.000000, 1.000000, 1.000000}
\pgfsetfillcolor{dialinecolor}
\pgfpathmoveto{\pgfpoint{-19.025914\du}{6.788864\du}}
\pgfpathlineto{\pgfpoint{-19.049798\du}{6.800507\du}}
\pgfpathlineto{\pgfpoint{-19.121152\du}{6.776623\du}}
\pgfpathlineto{\pgfpoint{-19.145036\du}{6.776623\du}}
\pgfpathlineto{\pgfpoint{-19.132795\du}{6.752739\du}}
\pgfpathlineto{\pgfpoint{-19.049798\du}{6.752739\du}}
\pgfpathlineto{\pgfpoint{-19.085326\du}{6.764681\du}}
\pgfpathlineto{\pgfpoint{-19.025914\du}{6.788864\du}}
\pgfusepath{fill}
\pgfsetbuttcap
\pgfsetmiterjoin
\pgfsetdash{}{0pt}
\definecolor{dialinecolor}{rgb}{0.666667, 0.901961, 1.000000}
\pgfsetstrokecolor{dialinecolor}
\pgfpathmoveto{\pgfpoint{-19.347454\du}{6.740797\du}}
\pgfpathlineto{\pgfpoint{-19.347454\du}{6.835736\du}}
\pgfusepath{stroke}
\pgfsetbuttcap
\pgfsetmiterjoin
\pgfsetdash{}{0pt}
\definecolor{dialinecolor}{rgb}{0.666667, 0.901961, 1.000000}
\pgfsetstrokecolor{dialinecolor}
\pgfpathmoveto{\pgfpoint{-18.966800\du}{6.740797\du}}
\pgfpathlineto{\pgfpoint{-18.966800\du}{6.835736\du}}
\pgfusepath{stroke}
\pgfsetlinewidth{0.100000\du}
\pgfsetdash{}{0pt}
\pgfsetdash{}{0pt}
\pgfsetbuttcap
\pgfsetmiterjoin
\pgfsetlinewidth{0.001000\du}
\pgfsetbuttcap
\pgfsetmiterjoin
\pgfsetdash{}{0pt}
\definecolor{dialinecolor}{rgb}{0.000000, 0.470588, 0.666667}
\pgfsetstrokecolor{dialinecolor}
\pgfpathmoveto{\pgfpoint{-12.585683\du}{7.363754\du}}
\pgfpathlineto{\pgfpoint{-12.579115\du}{7.354797\du}}
\pgfpathlineto{\pgfpoint{-12.572845\du}{7.345244\du}}
\pgfpathlineto{\pgfpoint{-12.567471\du}{7.335093\du}}
\pgfpathlineto{\pgfpoint{-12.562396\du}{7.325838\du}}
\pgfpathlineto{\pgfpoint{-12.558515\du}{7.315687\du}}
\pgfpathlineto{\pgfpoint{-12.555231\du}{7.305536\du}}
\pgfpathlineto{\pgfpoint{-12.552245\du}{7.295386\du}}
\pgfpathlineto{\pgfpoint{-12.550454\du}{7.284936\du}}
\pgfpathlineto{\pgfpoint{-12.549260\du}{7.274785\du}}
\pgfpathlineto{\pgfpoint{-12.548364\du}{7.264336\du}}
\pgfpathlineto{\pgfpoint{-12.548364\du}{7.253887\du}}
\pgfpathlineto{\pgfpoint{-12.549260\du}{7.243139\du}}
\pgfpathlineto{\pgfpoint{-12.550454\du}{7.232988\du}}
\pgfpathlineto{\pgfpoint{-12.552245\du}{7.222539\du}}
\pgfpathlineto{\pgfpoint{-12.554932\du}{7.212687\du}}
\pgfpathlineto{\pgfpoint{-12.558515\du}{7.202237\du}}
\pgfpathlineto{\pgfpoint{-12.562396\du}{7.192385\du}}
\pgfpathlineto{\pgfpoint{-12.567471\du}{7.182234\du}}
\pgfpathlineto{\pgfpoint{-12.572845\du}{7.172681\du}}
\pgfpathlineto{\pgfpoint{-12.579115\du}{7.163724\du}}
\pgfpathlineto{\pgfpoint{-12.585683\du}{7.154170\du}}
\pgfpathlineto{\pgfpoint{-12.593147\du}{7.145512\du}}
\pgfpathlineto{\pgfpoint{-12.600611\du}{7.136854\du}}
\pgfpathlineto{\pgfpoint{-12.608970\du}{7.128495\du}}
\pgfpathlineto{\pgfpoint{-12.617628\du}{7.120434\du}}
\pgfpathlineto{\pgfpoint{-12.626883\du}{7.112672\du}}
\pgfpathlineto{\pgfpoint{-12.636437\du}{7.105506\du}}
\pgfpathlineto{\pgfpoint{-12.647185\du}{7.098938\du}}
\pgfpathlineto{\pgfpoint{-12.657634\du}{7.092370\du}}
\pgfpathlineto{\pgfpoint{-12.668979\du}{7.086101\du}}
\pgfpathlineto{\pgfpoint{-12.680921\du}{7.080428\du}}
\pgfpathlineto{\pgfpoint{-12.692863\du}{7.075054\du}}
\pgfpathlineto{\pgfpoint{-12.705104\du}{7.070874\du}}
\pgfpathlineto{\pgfpoint{-12.717643\du}{7.066695\du}}
\pgfpathlineto{\pgfpoint{-12.730779\du}{7.062515\du}}
\pgfpathlineto{\pgfpoint{-12.743617\du}{7.059529\du}}
\pgfpathlineto{\pgfpoint{-12.756753\du}{7.056842\du}}
\pgfpathlineto{\pgfpoint{-12.770487\du}{7.054454\du}}
\pgfpathlineto{\pgfpoint{-12.783623\du}{7.052961\du}}
\pgfpathlineto{\pgfpoint{-12.797655\du}{7.051767\du}}
\pgfpathlineto{\pgfpoint{-12.811388\du}{7.050871\du}}
\pgfpathlineto{\pgfpoint{-12.825420\du}{7.050871\du}}
\pgfpathlineto{\pgfpoint{-12.838855\du}{7.051469\du}}
\pgfpathlineto{\pgfpoint{-12.852589\du}{7.052066\du}}
\pgfpathlineto{\pgfpoint{-12.866621\du}{7.053558\du}}
\pgfpathlineto{\pgfpoint{-12.879757\du}{7.055947\du}}
\pgfpathlineto{\pgfpoint{-12.893490\du}{7.058335\du}}
\pgfpathlineto{\pgfpoint{-12.906029\du}{7.061321\du}}
\pgfpathlineto{\pgfpoint{-12.919166\du}{7.064306\du}}
\pgfpathlineto{\pgfpoint{-12.931705\du}{7.068486\du}}
\pgfpathlineto{\pgfpoint{-12.944543\du}{7.073263\du}}
\pgfpathlineto{\pgfpoint{-12.956783\du}{7.077741\du}}
\pgfpathlineto{\pgfpoint{-12.968128\du}{7.083115\du}}
\pgfpathlineto{\pgfpoint{-12.979772\du}{7.089385\du}}
\pgfpathlineto{\pgfpoint{-12.990818\du}{7.095356\du}}
\pgfpathlineto{\pgfpoint{-13.001268\du}{7.102222\du}}
\pgfpathlineto{\pgfpoint{-13.011418\du}{7.108791\du}}
\pgfpathlineto{\pgfpoint{-13.020972\du}{7.116553\du}}
\pgfpathlineto{\pgfpoint{-13.030526\du}{7.124614\du}}
\pgfpathlineto{\pgfpoint{-13.038885\du}{7.132675\du}}
\pgfpathlineto{\pgfpoint{-13.046946\du}{7.141034\du}}
\pgfpathlineto{\pgfpoint{-13.054708\du}{7.149991\du}}
\pgfusepath{stroke}
\pgfsetlinewidth{0.080000\du}
\pgfsetbuttcap
\pgfsetmiterjoin
\pgfsetdash{}{0pt}
\definecolor{dialinecolor}{rgb}{0.666667, 0.901961, 1.000000}
\pgfsetstrokecolor{dialinecolor}
\pgfpathmoveto{\pgfpoint{-12.594938\du}{7.356290\du}}
\pgfpathlineto{\pgfpoint{-12.589564\du}{7.347334\du}}
\pgfpathlineto{\pgfpoint{-12.583295\du}{7.338676\du}}
\pgfpathlineto{\pgfpoint{-12.578219\du}{7.329420\du}}
\pgfpathlineto{\pgfpoint{-12.574040\du}{7.320165\du}}
\pgfpathlineto{\pgfpoint{-12.570158\du}{7.310612\du}}
\pgfpathlineto{\pgfpoint{-12.567173\du}{7.301058\du}}
\pgfpathlineto{\pgfpoint{-12.564486\du}{7.291206\du}}
\pgfpathlineto{\pgfpoint{-12.562396\du}{7.281951\du}}
\pgfpathlineto{\pgfpoint{-12.561500\du}{7.272098\du}}
\pgfpathlineto{\pgfpoint{-12.560903\du}{7.261948\du}}
\pgfpathlineto{\pgfpoint{-12.560903\du}{7.252095\du}}
\pgfpathlineto{\pgfpoint{-12.561799\du}{7.242243\du}}
\pgfpathlineto{\pgfpoint{-12.563292\du}{7.232093\du}}
\pgfpathlineto{\pgfpoint{-12.565083\du}{7.222837\du}}
\pgfpathlineto{\pgfpoint{-12.568069\du}{7.212985\du}}
\pgfpathlineto{\pgfpoint{-12.571651\du}{7.203432\du}}
\pgfpathlineto{\pgfpoint{-12.575532\du}{7.194176\du}}
\pgfpathlineto{\pgfpoint{-12.580011\du}{7.184921\du}}
\pgfpathlineto{\pgfpoint{-12.585086\du}{7.176263\du}}
\pgfpathlineto{\pgfpoint{-12.590758\du}{7.167008\du}}
\pgfpathlineto{\pgfpoint{-12.597924\du}{7.158649\du}}
\pgfpathlineto{\pgfpoint{-12.604492\du}{7.149991\du}}
\pgfpathlineto{\pgfpoint{-12.612254\du}{7.141930\du}}
\pgfpathlineto{\pgfpoint{-12.620315\du}{7.134167\du}}
\pgfpathlineto{\pgfpoint{-12.628675\du}{7.127301\du}}
\pgfpathlineto{\pgfpoint{-12.637333\du}{7.119837\du}}
\pgfpathlineto{\pgfpoint{-12.647185\du}{7.112970\du}}
\pgfpathlineto{\pgfpoint{-12.657634\du}{7.106104\du}}
\pgfpathlineto{\pgfpoint{-12.667785\du}{7.100431\du}}
\pgfpathlineto{\pgfpoint{-12.678234\du}{7.095057\du}}
\pgfpathlineto{\pgfpoint{-12.690176\du}{7.089683\du}}
\pgfpathlineto{\pgfpoint{-12.700924\du}{7.084608\du}}
\pgfpathlineto{\pgfpoint{-12.713463\du}{7.080428\du}}
\pgfpathlineto{\pgfpoint{-12.724808\du}{7.076547\du}}
\pgfpathlineto{\pgfpoint{-12.737348\du}{7.073263\du}}
\pgfpathlineto{\pgfpoint{-12.749887\du}{7.070277\du}}
\pgfpathlineto{\pgfpoint{-12.762426\du}{7.067590\du}}
\pgfpathlineto{\pgfpoint{-12.775562\du}{7.065799\du}}
\pgfpathlineto{\pgfpoint{-12.788101\du}{7.064306\du}}
\pgfpathlineto{\pgfpoint{-12.801238\du}{7.063112\du}}
\pgfpathlineto{\pgfpoint{-12.815270\du}{7.062515\du}}
\pgfpathlineto{\pgfpoint{-12.828406\du}{7.062515\du}}
\pgfpathlineto{\pgfpoint{-12.841244\du}{7.063112\du}}
\pgfpathlineto{\pgfpoint{-12.854380\du}{7.064008\du}}
\pgfpathlineto{\pgfpoint{-12.866919\du}{7.065799\du}}
\pgfpathlineto{\pgfpoint{-12.880055\du}{7.067590\du}}
\pgfpathlineto{\pgfpoint{-12.892595\du}{7.069680\du}}
\pgfpathlineto{\pgfpoint{-12.905432\du}{7.073263\du}}
\pgfpathlineto{\pgfpoint{-12.917673\du}{7.076248\du}}
\pgfpathlineto{\pgfpoint{-12.929615\du}{7.080130\du}}
\pgfpathlineto{\pgfpoint{-12.941856\du}{7.084309\du}}
\pgfpathlineto{\pgfpoint{-12.953201\du}{7.089385\du}}
\pgfpathlineto{\pgfpoint{-12.964247\du}{7.094161\du}}
\pgfpathlineto{\pgfpoint{-12.974995\du}{7.100133\du}}
\pgfpathlineto{\pgfpoint{-12.985444\du}{7.105805\du}}
\pgfpathlineto{\pgfpoint{-12.995894\du}{7.112373\du}}
\pgfpathlineto{\pgfpoint{-13.005447\du}{7.118941\du}}
\pgfpathlineto{\pgfpoint{-13.014105\du}{7.126107\du}}
\pgfpathlineto{\pgfpoint{-13.023659\du}{7.133272\du}}
\pgfpathlineto{\pgfpoint{-13.031421\du}{7.141034\du}}
\pgfpathlineto{\pgfpoint{-13.038885\du}{7.149095\du}}
\pgfpathlineto{\pgfpoint{-13.046050\du}{7.157455\du}}
\pgfusepath{stroke}
\pgfsetbuttcap
\pgfsetmiterjoin
\pgfsetdash{}{0pt}
\definecolor{dialinecolor}{rgb}{0.000000, 0.705882, 1.000000}
\pgfsetfillcolor{dialinecolor}
\pgfpathmoveto{\pgfpoint{-14.511941\du}{7.039228\du}}
\pgfpathlineto{\pgfpoint{-14.309224\du}{6.848453\du}}
\pgfpathlineto{\pgfpoint{-12.797655\du}{6.848453\du}}
\pgfpathlineto{\pgfpoint{-13.000372\du}{7.039228\du}}
\pgfpathlineto{\pgfpoint{-14.511941\du}{7.039228\du}}
\pgfusepath{fill}
\pgfsetlinewidth{0.001000\du}
\pgfsetbuttcap
\pgfsetmiterjoin
\pgfsetdash{}{0pt}
\definecolor{dialinecolor}{rgb}{0.666667, 0.901961, 1.000000}
\pgfsetstrokecolor{dialinecolor}
\pgfpathmoveto{\pgfpoint{-14.511941\du}{7.039228\du}}
\pgfpathlineto{\pgfpoint{-14.309224\du}{6.848453\du}}
\pgfpathlineto{\pgfpoint{-12.797655\du}{6.848453\du}}
\pgfpathlineto{\pgfpoint{-13.000372\du}{7.039228\du}}
\pgfpathlineto{\pgfpoint{-14.511941\du}{7.039228\du}}
\pgfusepath{stroke}
\pgfsetbuttcap
\pgfsetmiterjoin
\pgfsetdash{}{0pt}
\definecolor{dialinecolor}{rgb}{0.000000, 0.588235, 0.831373}
\pgfsetfillcolor{dialinecolor}
\pgfpathmoveto{\pgfpoint{-14.511941\du}{7.039228\du}}
\pgfpathlineto{\pgfpoint{-13.000372\du}{7.039228\du}}
\pgfpathlineto{\pgfpoint{-13.000372\du}{7.325241\du}}
\pgfpathlineto{\pgfpoint{-14.511941\du}{7.325241\du}}
\pgfpathlineto{\pgfpoint{-14.511941\du}{7.039228\du}}
\pgfusepath{fill}
\pgfsetbuttcap
\pgfsetmiterjoin
\pgfsetdash{}{0pt}
\definecolor{dialinecolor}{rgb}{0.666667, 0.901961, 1.000000}
\pgfsetstrokecolor{dialinecolor}
\pgfpathmoveto{\pgfpoint{-14.511941\du}{7.039228\du}}
\pgfpathlineto{\pgfpoint{-13.000372\du}{7.039228\du}}
\pgfpathlineto{\pgfpoint{-13.000372\du}{7.325241\du}}
\pgfpathlineto{\pgfpoint{-14.511941\du}{7.325241\du}}
\pgfpathlineto{\pgfpoint{-14.511941\du}{7.039228\du}}
\pgfusepath{stroke}
\pgfsetbuttcap
\pgfsetmiterjoin
\pgfsetdash{}{0pt}
\definecolor{dialinecolor}{rgb}{0.000000, 0.352941, 0.501961}
\pgfsetfillcolor{dialinecolor}
\pgfpathmoveto{\pgfpoint{-13.000372\du}{7.325241\du}}
\pgfpathlineto{\pgfpoint{-12.797655\du}{7.122524\du}}
\pgfpathlineto{\pgfpoint{-12.797655\du}{6.848453\du}}
\pgfpathlineto{\pgfpoint{-13.000372\du}{7.039228\du}}
\pgfpathlineto{\pgfpoint{-13.000372\du}{7.325241\du}}
\pgfusepath{fill}
\pgfsetbuttcap
\pgfsetmiterjoin
\pgfsetdash{}{0pt}
\definecolor{dialinecolor}{rgb}{0.666667, 0.901961, 1.000000}
\pgfsetstrokecolor{dialinecolor}
\pgfpathmoveto{\pgfpoint{-13.000372\du}{7.325241\du}}
\pgfpathlineto{\pgfpoint{-12.797655\du}{7.122524\du}}
\pgfpathlineto{\pgfpoint{-12.797655\du}{6.848453\du}}
\pgfpathlineto{\pgfpoint{-13.000372\du}{7.039228\du}}
\pgfpathlineto{\pgfpoint{-13.000372\du}{7.325241\du}}
\pgfusepath{stroke}
\pgfsetbuttcap
\pgfsetmiterjoin
\pgfsetdash{}{0pt}
\definecolor{dialinecolor}{rgb}{0.000000, 0.000000, 0.000000}
\pgfsetfillcolor{dialinecolor}
\pgfpathmoveto{\pgfpoint{-14.463874\du}{7.003103\du}}
\pgfpathlineto{\pgfpoint{-14.309224\du}{6.848453\du}}
\pgfpathlineto{\pgfpoint{-12.833780\du}{6.848453\du}}
\pgfpathlineto{\pgfpoint{-12.988131\du}{7.003103\du}}
\pgfpathlineto{\pgfpoint{-14.463874\du}{7.003103\du}}
\pgfusepath{fill}
\pgfsetbuttcap
\pgfsetmiterjoin
\pgfsetdash{}{0pt}
\definecolor{dialinecolor}{rgb}{0.000000, 0.000000, 0.000000}
\pgfsetstrokecolor{dialinecolor}
\pgfpathmoveto{\pgfpoint{-14.463874\du}{7.003103\du}}
\pgfpathlineto{\pgfpoint{-14.309224\du}{6.848453\du}}
\pgfpathlineto{\pgfpoint{-12.833780\du}{6.848453\du}}
\pgfpathlineto{\pgfpoint{-12.988131\du}{7.003103\du}}
\pgfpathlineto{\pgfpoint{-14.463874\du}{7.003103\du}}
\pgfusepath{stroke}
\pgfsetbuttcap
\pgfsetmiterjoin
\pgfsetdash{}{0pt}
\definecolor{dialinecolor}{rgb}{0.000000, 0.705882, 1.000000}
\pgfsetfillcolor{dialinecolor}
\pgfpathmoveto{\pgfpoint{-14.511941\du}{5.776950\du}}
\pgfpathlineto{\pgfpoint{-14.344752\du}{5.622300\du}}
\pgfpathlineto{\pgfpoint{-12.833780\du}{5.622300\du}}
\pgfpathlineto{\pgfpoint{-13.000372\du}{5.776950\du}}
\pgfpathlineto{\pgfpoint{-14.511941\du}{5.776950\du}}
\pgfusepath{fill}
\pgfsetbuttcap
\pgfsetmiterjoin
\pgfsetdash{}{0pt}
\definecolor{dialinecolor}{rgb}{0.666667, 0.901961, 1.000000}
\pgfsetstrokecolor{dialinecolor}
\pgfpathmoveto{\pgfpoint{-14.511941\du}{5.776950\du}}
\pgfpathlineto{\pgfpoint{-14.344752\du}{5.622300\du}}
\pgfpathlineto{\pgfpoint{-12.833780\du}{5.622300\du}}
\pgfpathlineto{\pgfpoint{-13.000372\du}{5.776950\du}}
\pgfpathlineto{\pgfpoint{-14.511941\du}{5.776950\du}}
\pgfusepath{stroke}
\pgfsetbuttcap
\pgfsetmiterjoin
\pgfsetdash{}{0pt}
\definecolor{dialinecolor}{rgb}{0.000000, 0.588235, 0.831373}
\pgfsetfillcolor{dialinecolor}
\pgfpathmoveto{\pgfpoint{-14.511941\du}{5.776950\du}}
\pgfpathlineto{\pgfpoint{-12.988131\du}{5.776950\du}}
\pgfpathlineto{\pgfpoint{-12.988131\du}{6.979219\du}}
\pgfpathlineto{\pgfpoint{-14.511941\du}{6.979219\du}}
\pgfpathlineto{\pgfpoint{-14.511941\du}{5.776950\du}}
\pgfusepath{fill}
\pgfsetbuttcap
\pgfsetmiterjoin
\pgfsetdash{}{0pt}
\definecolor{dialinecolor}{rgb}{0.666667, 0.901961, 1.000000}
\pgfsetstrokecolor{dialinecolor}
\pgfpathmoveto{\pgfpoint{-14.511941\du}{5.776950\du}}
\pgfpathlineto{\pgfpoint{-12.988430\du}{5.776950\du}}
\pgfpathlineto{\pgfpoint{-12.988430\du}{6.978920\du}}
\pgfpathlineto{\pgfpoint{-14.511941\du}{6.978920\du}}
\pgfpathlineto{\pgfpoint{-14.511941\du}{5.776950\du}}
\pgfusepath{stroke}
\pgfsetbuttcap
\pgfsetmiterjoin
\pgfsetdash{}{0pt}
\definecolor{dialinecolor}{rgb}{1.000000, 1.000000, 1.000000}
\pgfsetfillcolor{dialinecolor}
\pgfpathmoveto{\pgfpoint{-14.380876\du}{5.931898\du}}
\pgfpathlineto{\pgfpoint{-13.119196\du}{5.931898\du}}
\pgfpathlineto{\pgfpoint{-13.119196\du}{6.860694\du}}
\pgfpathlineto{\pgfpoint{-14.380876\du}{6.860694\du}}
\pgfpathlineto{\pgfpoint{-14.380876\du}{5.931898\du}}
\pgfusepath{fill}
\pgfsetbuttcap
\pgfsetmiterjoin
\pgfsetdash{}{0pt}
\definecolor{dialinecolor}{rgb}{0.000000, 0.235294, 0.333333}
\pgfsetstrokecolor{dialinecolor}
\pgfpathmoveto{\pgfpoint{-14.380876\du}{5.931898\du}}
\pgfpathlineto{\pgfpoint{-13.119494\du}{5.931898\du}}
\pgfpathlineto{\pgfpoint{-13.119494\du}{6.860395\du}}
\pgfpathlineto{\pgfpoint{-14.380876\du}{6.860395\du}}
\pgfpathlineto{\pgfpoint{-14.380876\du}{5.931898\du}}
\pgfusepath{stroke}
\pgfsetbuttcap
\pgfsetmiterjoin
\pgfsetdash{}{0pt}
\definecolor{dialinecolor}{rgb}{0.000000, 0.352941, 0.501961}
\pgfsetfillcolor{dialinecolor}
\pgfpathmoveto{\pgfpoint{-13.000372\du}{6.967874\du}}
\pgfpathlineto{\pgfpoint{-12.833780\du}{6.800983\du}}
\pgfpathlineto{\pgfpoint{-12.833780\du}{5.622300\du}}
\pgfpathlineto{\pgfpoint{-13.000372\du}{5.776950\du}}
\pgfpathlineto{\pgfpoint{-13.000372\du}{6.967874\du}}
\pgfusepath{fill}
\pgfsetbuttcap
\pgfsetmiterjoin
\pgfsetdash{}{0pt}
\definecolor{dialinecolor}{rgb}{0.666667, 0.901961, 1.000000}
\pgfsetstrokecolor{dialinecolor}
\pgfpathmoveto{\pgfpoint{-13.000372\du}{6.967874\du}}
\pgfpathlineto{\pgfpoint{-12.833780\du}{6.800983\du}}
\pgfpathlineto{\pgfpoint{-12.833780\du}{5.622300\du}}
\pgfpathlineto{\pgfpoint{-13.000372\du}{5.776950\du}}
\pgfpathlineto{\pgfpoint{-13.000372\du}{6.967874\du}}
\pgfusepath{stroke}
\pgfsetbuttcap
\pgfsetmiterjoin
\pgfsetdash{}{0pt}
\definecolor{dialinecolor}{rgb}{0.000000, 0.705882, 1.000000}
\pgfsetfillcolor{dialinecolor}
\pgfpathmoveto{\pgfpoint{-14.809000\du}{7.563187\du}}
\pgfpathlineto{\pgfpoint{-14.571353\du}{7.265232\du}}
\pgfpathlineto{\pgfpoint{-12.928719\du}{7.265232\du}}
\pgfpathlineto{\pgfpoint{-13.166665\du}{7.563187\du}}
\pgfpathlineto{\pgfpoint{-14.809000\du}{7.563187\du}}
\pgfusepath{fill}
\pgfsetbuttcap
\pgfsetmiterjoin
\pgfsetdash{}{0pt}
\definecolor{dialinecolor}{rgb}{0.666667, 0.901961, 1.000000}
\pgfsetstrokecolor{dialinecolor}
\pgfpathmoveto{\pgfpoint{-14.809000\du}{7.563187\du}}
\pgfpathlineto{\pgfpoint{-14.571353\du}{7.265232\du}}
\pgfpathlineto{\pgfpoint{-12.928719\du}{7.265232\du}}
\pgfpathlineto{\pgfpoint{-13.166665\du}{7.563187\du}}
\pgfpathlineto{\pgfpoint{-14.809000\du}{7.563187\du}}
\pgfusepath{stroke}
\pgfsetbuttcap
\pgfsetmiterjoin
\pgfsetdash{}{0pt}
\definecolor{dialinecolor}{rgb}{0.000000, 0.352941, 0.501961}
\pgfsetfillcolor{dialinecolor}
\pgfpathmoveto{\pgfpoint{-13.166665\du}{7.622300\du}}
\pgfpathlineto{\pgfpoint{-12.928719\du}{7.372412\du}}
\pgfpathlineto{\pgfpoint{-12.928719\du}{7.265232\du}}
\pgfpathlineto{\pgfpoint{-13.166665\du}{7.563187\du}}
\pgfpathlineto{\pgfpoint{-13.166665\du}{7.622300\du}}
\pgfusepath{fill}
\pgfsetbuttcap
\pgfsetmiterjoin
\pgfsetdash{}{0pt}
\definecolor{dialinecolor}{rgb}{0.666667, 0.901961, 1.000000}
\pgfsetstrokecolor{dialinecolor}
\pgfpathmoveto{\pgfpoint{-13.166665\du}{7.622300\du}}
\pgfpathlineto{\pgfpoint{-12.928719\du}{7.372412\du}}
\pgfpathlineto{\pgfpoint{-12.928719\du}{7.265232\du}}
\pgfpathlineto{\pgfpoint{-13.166665\du}{7.563187\du}}
\pgfpathlineto{\pgfpoint{-13.166665\du}{7.622300\du}}
\pgfusepath{stroke}
\pgfsetbuttcap
\pgfsetmiterjoin
\pgfsetdash{}{0pt}
\definecolor{dialinecolor}{rgb}{0.000000, 0.588235, 0.831373}
\pgfsetfillcolor{dialinecolor}
\pgfpathmoveto{\pgfpoint{-14.809000\du}{7.563187\du}}
\pgfpathlineto{\pgfpoint{-13.166665\du}{7.563187\du}}
\pgfpathlineto{\pgfpoint{-13.166665\du}{7.622300\du}}
\pgfpathlineto{\pgfpoint{-14.809000\du}{7.622300\du}}
\pgfpathlineto{\pgfpoint{-14.809000\du}{7.563187\du}}
\pgfusepath{fill}
\pgfsetbuttcap
\pgfsetmiterjoin
\pgfsetdash{}{0pt}
\definecolor{dialinecolor}{rgb}{0.666667, 0.901961, 1.000000}
\pgfsetstrokecolor{dialinecolor}
\pgfpathmoveto{\pgfpoint{-14.809000\du}{7.563187\du}}
\pgfpathlineto{\pgfpoint{-13.166665\du}{7.563187\du}}
\pgfpathlineto{\pgfpoint{-13.166665\du}{7.622300\du}}
\pgfpathlineto{\pgfpoint{-14.809000\du}{7.622300\du}}
\pgfpathlineto{\pgfpoint{-14.809000\du}{7.563187\du}}
\pgfusepath{stroke}
\pgfsetbuttcap
\pgfsetmiterjoin
\pgfsetdash{}{0pt}
\definecolor{dialinecolor}{rgb}{0.000000, 0.705882, 1.000000}
\pgfsetfillcolor{dialinecolor}
\pgfpathmoveto{\pgfpoint{-12.833780\du}{7.527062\du}}
\pgfpathlineto{\pgfpoint{-12.714658\du}{7.348826\du}}
\pgfpathlineto{\pgfpoint{-12.499999\du}{7.348826\du}}
\pgfpathlineto{\pgfpoint{-12.618822\du}{7.527062\du}}
\pgfpathlineto{\pgfpoint{-12.833780\du}{7.527062\du}}
\pgfusepath{fill}
\pgfsetbuttcap
\pgfsetmiterjoin
\pgfsetdash{}{0pt}
\definecolor{dialinecolor}{rgb}{0.666667, 0.901961, 1.000000}
\pgfsetstrokecolor{dialinecolor}
\pgfpathmoveto{\pgfpoint{-12.833780\du}{7.527062\du}}
\pgfpathlineto{\pgfpoint{-12.714658\du}{7.348826\du}}
\pgfpathlineto{\pgfpoint{-12.499999\du}{7.348826\du}}
\pgfpathlineto{\pgfpoint{-12.618822\du}{7.527062\du}}
\pgfpathlineto{\pgfpoint{-12.833780\du}{7.527062\du}}
\pgfusepath{stroke}
\pgfsetbuttcap
\pgfsetmiterjoin
\pgfsetdash{}{0pt}
\definecolor{dialinecolor}{rgb}{0.000000, 0.352941, 0.501961}
\pgfsetfillcolor{dialinecolor}
\pgfpathmoveto{\pgfpoint{-12.618822\du}{7.586772\du}}
\pgfpathlineto{\pgfpoint{-12.499999\du}{7.455708\du}}
\pgfpathlineto{\pgfpoint{-12.499999\du}{7.348826\du}}
\pgfpathlineto{\pgfpoint{-12.618822\du}{7.527062\du}}
\pgfpathlineto{\pgfpoint{-12.618822\du}{7.586772\du}}
\pgfusepath{fill}
\pgfsetbuttcap
\pgfsetmiterjoin
\pgfsetdash{}{0pt}
\definecolor{dialinecolor}{rgb}{0.666667, 0.901961, 1.000000}
\pgfsetstrokecolor{dialinecolor}
\pgfpathmoveto{\pgfpoint{-12.618822\du}{7.586772\du}}
\pgfpathlineto{\pgfpoint{-12.499999\du}{7.455708\du}}
\pgfpathlineto{\pgfpoint{-12.499999\du}{7.348826\du}}
\pgfpathlineto{\pgfpoint{-12.618822\du}{7.527062\du}}
\pgfpathlineto{\pgfpoint{-12.618822\du}{7.586772\du}}
\pgfusepath{stroke}
\pgfsetbuttcap
\pgfsetmiterjoin
\pgfsetdash{}{0pt}
\definecolor{dialinecolor}{rgb}{0.000000, 0.588235, 0.831373}
\pgfsetfillcolor{dialinecolor}
\pgfpathmoveto{\pgfpoint{-12.833780\du}{7.527062\du}}
\pgfpathlineto{\pgfpoint{-12.618822\du}{7.527062\du}}
\pgfpathlineto{\pgfpoint{-12.618822\du}{7.586772\du}}
\pgfpathlineto{\pgfpoint{-12.833780\du}{7.586772\du}}
\pgfpathlineto{\pgfpoint{-12.833780\du}{7.527062\du}}
\pgfusepath{fill}
\pgfsetbuttcap
\pgfsetmiterjoin
\pgfsetdash{}{0pt}
\definecolor{dialinecolor}{rgb}{0.666667, 0.901961, 1.000000}
\pgfsetstrokecolor{dialinecolor}
\pgfpathmoveto{\pgfpoint{-12.833780\du}{7.527062\du}}
\pgfpathlineto{\pgfpoint{-12.618822\du}{7.527062\du}}
\pgfpathlineto{\pgfpoint{-12.618822\du}{7.586772\du}}
\pgfpathlineto{\pgfpoint{-12.833780\du}{7.586772\du}}
\pgfpathlineto{\pgfpoint{-12.833780\du}{7.527062\du}}
\pgfusepath{stroke}
\pgfsetbuttcap
\pgfsetmiterjoin
\pgfsetdash{}{0pt}
\definecolor{dialinecolor}{rgb}{0.811765, 0.054902, 0.188235}
\pgfsetstrokecolor{dialinecolor}
\pgfpathmoveto{\pgfpoint{-14.047692\du}{6.158201\du}}
\pgfpathlineto{\pgfpoint{-13.713911\du}{6.158201\du}}
\pgfpathlineto{\pgfpoint{-13.797506\du}{6.241497\du}}
\pgfpathlineto{\pgfpoint{-13.511791\du}{6.241497\du}}
\pgfusepath{stroke}
\pgfsetbuttcap
\pgfsetmiterjoin
\pgfsetdash{}{0pt}
\definecolor{dialinecolor}{rgb}{0.811765, 0.054902, 0.188235}
\pgfsetstrokecolor{dialinecolor}
\pgfpathmoveto{\pgfpoint{-14.083519\du}{6.158201\du}}
\pgfpathlineto{\pgfpoint{-13.844976\du}{6.396147\du}}
\pgfpathlineto{\pgfpoint{-13.964396\du}{6.396147\du}}
\pgfpathlineto{\pgfpoint{-13.761978\du}{6.598864\du}}
\pgfusepath{stroke}
\pgfsetbuttcap
\pgfsetmiterjoin
\pgfsetdash{}{0pt}
\definecolor{dialinecolor}{rgb}{0.811765, 0.054902, 0.188235}
\pgfsetstrokecolor{dialinecolor}
\pgfpathmoveto{\pgfpoint{-13.405208\du}{6.158201\du}}
\pgfpathlineto{\pgfpoint{-13.642557\du}{6.396147\du}}
\pgfpathlineto{\pgfpoint{-13.524032\du}{6.396147\du}}
\pgfpathlineto{\pgfpoint{-13.726152\du}{6.598864\du}}
\pgfusepath{stroke}
\pgfsetbuttcap
\pgfsetmiterjoin
\pgfsetdash{}{0pt}
\definecolor{dialinecolor}{rgb}{0.000000, 0.470588, 0.666667}
\pgfsetfillcolor{dialinecolor}
\pgfpathmoveto{\pgfpoint{-13.916927\du}{6.247169\du}}
\pgfpathlineto{\pgfpoint{-13.917524\du}{6.243885\du}}
\pgfpathlineto{\pgfpoint{-13.918419\du}{6.240601\du}}
\pgfpathlineto{\pgfpoint{-13.919315\du}{6.237317\du}}
\pgfpathlineto{\pgfpoint{-13.921106\du}{6.234033\du}}
\pgfpathlineto{\pgfpoint{-13.923495\du}{6.231048\du}}
\pgfpathlineto{\pgfpoint{-13.926182\du}{6.227764\du}}
\pgfpathlineto{\pgfpoint{-13.929167\du}{6.224479\du}}
\pgfpathlineto{\pgfpoint{-13.932750\du}{6.221494\du}}
\pgfpathlineto{\pgfpoint{-13.936930\du}{6.218508\du}}
\pgfpathlineto{\pgfpoint{-13.941706\du}{6.215523\du}}
\pgfpathlineto{\pgfpoint{-13.946782\du}{6.212836\du}}
\pgfpathlineto{\pgfpoint{-13.952454\du}{6.209850\du}}
\pgfpathlineto{\pgfpoint{-13.958127\du}{6.207462\du}}
\pgfpathlineto{\pgfpoint{-13.964396\du}{6.204775\du}}
\pgfpathlineto{\pgfpoint{-13.970964\du}{6.202387\du}}
\pgfpathlineto{\pgfpoint{-13.978130\du}{6.199998\du}}
\pgfpathlineto{\pgfpoint{-13.985594\du}{6.197908\du}}
\pgfpathlineto{\pgfpoint{-13.993057\du}{6.195520\du}}
\pgfpathlineto{\pgfpoint{-14.001417\du}{6.194027\du}}
\pgfpathlineto{\pgfpoint{-14.009776\du}{6.191937\du}}
\pgfpathlineto{\pgfpoint{-14.017837\du}{6.190146\du}}
\pgfpathlineto{\pgfpoint{-14.027092\du}{6.188952\du}}
\pgfpathlineto{\pgfpoint{-14.035750\du}{6.187459\du}}
\pgfpathlineto{\pgfpoint{-14.045304\du}{6.186265\du}}
\pgfpathlineto{\pgfpoint{-14.054559\du}{6.185071\du}}
\pgfpathlineto{\pgfpoint{-14.063814\du}{6.184175\du}}
\pgfpathlineto{\pgfpoint{-14.073666\du}{6.183279\du}}
\pgfpathlineto{\pgfpoint{-14.083519\du}{6.182682\du}}
\pgfpathlineto{\pgfpoint{-14.093371\du}{6.182384\du}}
\pgfpathlineto{\pgfpoint{-14.103522\du}{6.182085\du}}
\pgfpathlineto{\pgfpoint{-14.113374\du}{6.182085\du}}
\pgfpathlineto{\pgfpoint{-14.113374\du}{6.182085\du}}
\pgfpathlineto{\pgfpoint{-14.123226\du}{6.182085\du}}
\pgfpathlineto{\pgfpoint{-14.133078\du}{6.182384\du}}
\pgfpathlineto{\pgfpoint{-14.143229\du}{6.182682\du}}
\pgfpathlineto{\pgfpoint{-14.153081\du}{6.183279\du}}
\pgfpathlineto{\pgfpoint{-14.162635\du}{6.184175\du}}
\pgfpathlineto{\pgfpoint{-14.172189\du}{6.185071\du}}
\pgfpathlineto{\pgfpoint{-14.181444\du}{6.186265\du}}
\pgfpathlineto{\pgfpoint{-14.190699\du}{6.187459\du}}
\pgfpathlineto{\pgfpoint{-14.199655\du}{6.188952\du}}
\pgfpathlineto{\pgfpoint{-14.208612\du}{6.190146\du}}
\pgfpathlineto{\pgfpoint{-14.216971\du}{6.191937\du}}
\pgfpathlineto{\pgfpoint{-14.225331\du}{6.194027\du}}
\pgfpathlineto{\pgfpoint{-14.233690\du}{6.195520\du}}
\pgfpathlineto{\pgfpoint{-14.240856\du}{6.197908\du}}
\pgfpathlineto{\pgfpoint{-14.248319\du}{6.199998\du}}
\pgfpathlineto{\pgfpoint{-14.255485\du}{6.202387\du}}
\pgfpathlineto{\pgfpoint{-14.262351\du}{6.204775\du}}
\pgfpathlineto{\pgfpoint{-14.268322\du}{6.207462\du}}
\pgfpathlineto{\pgfpoint{-14.274293\du}{6.209850\du}}
\pgfpathlineto{\pgfpoint{-14.279966\du}{6.212836\du}}
\pgfpathlineto{\pgfpoint{-14.285041\du}{6.215523\du}}
\pgfpathlineto{\pgfpoint{-14.289818\du}{6.218508\du}}
\pgfpathlineto{\pgfpoint{-14.293998\du}{6.221494\du}}
\pgfpathlineto{\pgfpoint{-14.297282\du}{6.224479\du}}
\pgfpathlineto{\pgfpoint{-14.300566\du}{6.227764\du}}
\pgfpathlineto{\pgfpoint{-14.303551\du}{6.231048\du}}
\pgfpathlineto{\pgfpoint{-14.305641\du}{6.234033\du}}
\pgfpathlineto{\pgfpoint{-14.307134\du}{6.237317\du}}
\pgfpathlineto{\pgfpoint{-14.308925\du}{6.240601\du}}
\pgfpathlineto{\pgfpoint{-14.309224\du}{6.243885\du}}
\pgfpathlineto{\pgfpoint{-14.309522\du}{6.247169\du}}
\pgfpathlineto{\pgfpoint{-14.309522\du}{6.247169\du}}
\pgfpathlineto{\pgfpoint{-14.309224\du}{6.250752\du}}
\pgfpathlineto{\pgfpoint{-14.308925\du}{6.253738\du}}
\pgfpathlineto{\pgfpoint{-14.307134\du}{6.257022\du}}
\pgfpathlineto{\pgfpoint{-14.305641\du}{6.260306\du}}
\pgfpathlineto{\pgfpoint{-14.303551\du}{6.263590\du}}
\pgfpathlineto{\pgfpoint{-14.300566\du}{6.266874\du}}
\pgfpathlineto{\pgfpoint{-14.297282\du}{6.269859\du}}
\pgfpathlineto{\pgfpoint{-14.293998\du}{6.272845\du}}
\pgfpathlineto{\pgfpoint{-14.289818\du}{6.275830\du}}
\pgfpathlineto{\pgfpoint{-14.285041\du}{6.279114\du}}
\pgfpathlineto{\pgfpoint{-14.279966\du}{6.281801\du}}
\pgfpathlineto{\pgfpoint{-14.274293\du}{6.284488\du}}
\pgfpathlineto{\pgfpoint{-14.268322\du}{6.287175\du}}
\pgfpathlineto{\pgfpoint{-14.262351\du}{6.289862\du}}
\pgfpathlineto{\pgfpoint{-14.255485\du}{6.292251\du}}
\pgfpathlineto{\pgfpoint{-14.248319\du}{6.294639\du}}
\pgfpathlineto{\pgfpoint{-14.240856\du}{6.296430\du}}
\pgfpathlineto{\pgfpoint{-14.233690\du}{6.298520\du}}
\pgfpathlineto{\pgfpoint{-14.225331\du}{6.300610\du}}
\pgfpathlineto{\pgfpoint{-14.216971\du}{6.302700\du}}
\pgfpathlineto{\pgfpoint{-14.208612\du}{6.304193\du}}
\pgfpathlineto{\pgfpoint{-14.199655\du}{6.305686\du}}
\pgfpathlineto{\pgfpoint{-14.190699\du}{6.307178\du}}
\pgfpathlineto{\pgfpoint{-14.181444\du}{6.308373\du}}
\pgfpathlineto{\pgfpoint{-14.172189\du}{6.309268\du}}
\pgfpathlineto{\pgfpoint{-14.162635\du}{6.310462\du}}
\pgfpathlineto{\pgfpoint{-14.153081\du}{6.311060\du}}
\pgfpathlineto{\pgfpoint{-14.143229\du}{6.311657\du}}
\pgfpathlineto{\pgfpoint{-14.133078\du}{6.312254\du}}
\pgfpathlineto{\pgfpoint{-14.123226\du}{6.312552\du}}
\pgfpathlineto{\pgfpoint{-14.113374\du}{6.312552\du}}
\pgfpathlineto{\pgfpoint{-14.113374\du}{6.312552\du}}
\pgfpathlineto{\pgfpoint{-14.103522\du}{6.312552\du}}
\pgfpathlineto{\pgfpoint{-14.093371\du}{6.312254\du}}
\pgfpathlineto{\pgfpoint{-14.083519\du}{6.311657\du}}
\pgfpathlineto{\pgfpoint{-14.073666\du}{6.311060\du}}
\pgfpathlineto{\pgfpoint{-14.063814\du}{6.310462\du}}
\pgfpathlineto{\pgfpoint{-14.054559\du}{6.309268\du}}
\pgfpathlineto{\pgfpoint{-14.045304\du}{6.308373\du}}
\pgfpathlineto{\pgfpoint{-14.035750\du}{6.307178\du}}
\pgfpathlineto{\pgfpoint{-14.027092\du}{6.305686\du}}
\pgfpathlineto{\pgfpoint{-14.017837\du}{6.304193\du}}
\pgfpathlineto{\pgfpoint{-14.009776\du}{6.302700\du}}
\pgfpathlineto{\pgfpoint{-14.001417\du}{6.300610\du}}
\pgfpathlineto{\pgfpoint{-13.993057\du}{6.298520\du}}
\pgfpathlineto{\pgfpoint{-13.985594\du}{6.296430\du}}
\pgfpathlineto{\pgfpoint{-13.978130\du}{6.294639\du}}
\pgfpathlineto{\pgfpoint{-13.970964\du}{6.292251\du}}
\pgfpathlineto{\pgfpoint{-13.964396\du}{6.289862\du}}
\pgfpathlineto{\pgfpoint{-13.958127\du}{6.287175\du}}
\pgfpathlineto{\pgfpoint{-13.952454\du}{6.284488\du}}
\pgfpathlineto{\pgfpoint{-13.946782\du}{6.281801\du}}
\pgfpathlineto{\pgfpoint{-13.941706\du}{6.279114\du}}
\pgfpathlineto{\pgfpoint{-13.936930\du}{6.275830\du}}
\pgfpathlineto{\pgfpoint{-13.932750\du}{6.272845\du}}
\pgfpathlineto{\pgfpoint{-13.929167\du}{6.269859\du}}
\pgfpathlineto{\pgfpoint{-13.926182\du}{6.266874\du}}
\pgfpathlineto{\pgfpoint{-13.923495\du}{6.263590\du}}
\pgfpathlineto{\pgfpoint{-13.921106\du}{6.260306\du}}
\pgfpathlineto{\pgfpoint{-13.919315\du}{6.257022\du}}
\pgfpathlineto{\pgfpoint{-13.918419\du}{6.253738\du}}
\pgfpathlineto{\pgfpoint{-13.917524\du}{6.250752\du}}
\pgfpathlineto{\pgfpoint{-13.916927\du}{6.247169\du}}
\pgfusepath{fill}
\pgfsetbuttcap
\pgfsetmiterjoin
\pgfsetdash{}{0pt}
\definecolor{dialinecolor}{rgb}{0.666667, 0.901961, 1.000000}
\pgfsetstrokecolor{dialinecolor}
\pgfpathmoveto{\pgfpoint{-13.929167\du}{6.241497\du}}
\pgfpathlineto{\pgfpoint{-13.929167\du}{6.237914\du}}
\pgfpathlineto{\pgfpoint{-13.930063\du}{6.235526\du}}
\pgfpathlineto{\pgfpoint{-13.931257\du}{6.232242\du}}
\pgfpathlineto{\pgfpoint{-13.932750\du}{6.229256\du}}
\pgfpathlineto{\pgfpoint{-13.935138\du}{6.226569\du}}
\pgfpathlineto{\pgfpoint{-13.937527\du}{6.222987\du}}
\pgfpathlineto{\pgfpoint{-13.940811\du}{6.220598\du}}
\pgfpathlineto{\pgfpoint{-13.944095\du}{6.217911\du}}
\pgfpathlineto{\pgfpoint{-13.948275\du}{6.215224\du}}
\pgfpathlineto{\pgfpoint{-13.952753\du}{6.212836\du}}
\pgfpathlineto{\pgfpoint{-13.957530\du}{6.209850\du}}
\pgfpathlineto{\pgfpoint{-13.963202\du}{6.207761\du}}
\pgfpathlineto{\pgfpoint{-13.968576\du}{6.205074\du}}
\pgfpathlineto{\pgfpoint{-13.974846\du}{6.202387\du}}
\pgfpathlineto{\pgfpoint{-13.980817\du}{6.200595\du}}
\pgfpathlineto{\pgfpoint{-13.988280\du}{6.198207\du}}
\pgfpathlineto{\pgfpoint{-13.995147\du}{6.196714\du}}
\pgfpathlineto{\pgfpoint{-14.002312\du}{6.194624\du}}
\pgfpathlineto{\pgfpoint{-14.010075\du}{6.192833\du}}
\pgfpathlineto{\pgfpoint{-14.018733\du}{6.191042\du}}
\pgfpathlineto{\pgfpoint{-14.027092\du}{6.189847\du}}
\pgfpathlineto{\pgfpoint{-14.035452\du}{6.187459\du}}
\pgfpathlineto{\pgfpoint{-14.044408\du}{6.187160\du}}
\pgfpathlineto{\pgfpoint{-14.053066\du}{6.185668\du}}
\pgfpathlineto{\pgfpoint{-14.062023\du}{6.184772\du}}
\pgfpathlineto{\pgfpoint{-14.071278\du}{6.183876\du}}
\pgfpathlineto{\pgfpoint{-14.080832\du}{6.183279\du}}
\pgfpathlineto{\pgfpoint{-14.090385\du}{6.182981\du}}
\pgfpathlineto{\pgfpoint{-14.099342\du}{6.182384\du}}
\pgfpathlineto{\pgfpoint{-14.109493\du}{6.182085\du}}
\pgfpathlineto{\pgfpoint{-14.119046\du}{6.182085\du}}
\pgfpathlineto{\pgfpoint{-14.119046\du}{6.182085\du}}
\pgfpathlineto{\pgfpoint{-14.128301\du}{6.182085\du}}
\pgfpathlineto{\pgfpoint{-14.138154\du}{6.182384\du}}
\pgfpathlineto{\pgfpoint{-14.147707\du}{6.182981\du}}
\pgfpathlineto{\pgfpoint{-14.157261\du}{6.183279\du}}
\pgfpathlineto{\pgfpoint{-14.166516\du}{6.183876\du}}
\pgfpathlineto{\pgfpoint{-14.176368\du}{6.184772\du}}
\pgfpathlineto{\pgfpoint{-14.185026\du}{6.185668\du}}
\pgfpathlineto{\pgfpoint{-14.193983\du}{6.187160\du}}
\pgfpathlineto{\pgfpoint{-14.202342\du}{6.187459\du}}
\pgfpathlineto{\pgfpoint{-14.211299\du}{6.189847\du}}
\pgfpathlineto{\pgfpoint{-14.219658\du}{6.191042\du}}
\pgfpathlineto{\pgfpoint{-14.227719\du}{6.192833\du}}
\pgfpathlineto{\pgfpoint{-14.235482\du}{6.194624\du}}
\pgfpathlineto{\pgfpoint{-14.242945\du}{6.196714\du}}
\pgfpathlineto{\pgfpoint{-14.249812\du}{6.198207\du}}
\pgfpathlineto{\pgfpoint{-14.257276\du}{6.200595\du}}
\pgfpathlineto{\pgfpoint{-14.263247\du}{6.202387\du}}
\pgfpathlineto{\pgfpoint{-14.269516\du}{6.205074\du}}
\pgfpathlineto{\pgfpoint{-14.275189\du}{6.207761\du}}
\pgfpathlineto{\pgfpoint{-14.280563\du}{6.209850\du}}
\pgfpathlineto{\pgfpoint{-14.285340\du}{6.212836\du}}
\pgfpathlineto{\pgfpoint{-14.289818\du}{6.215224\du}}
\pgfpathlineto{\pgfpoint{-14.293998\du}{6.217911\du}}
\pgfpathlineto{\pgfpoint{-14.297282\du}{6.220598\du}}
\pgfpathlineto{\pgfpoint{-14.300864\du}{6.222987\du}}
\pgfpathlineto{\pgfpoint{-14.302656\du}{6.226569\du}}
\pgfpathlineto{\pgfpoint{-14.305044\du}{6.229256\du}}
\pgfpathlineto{\pgfpoint{-14.306835\du}{6.232242\du}}
\pgfpathlineto{\pgfpoint{-14.308328\du}{6.235526\du}}
\pgfpathlineto{\pgfpoint{-14.308925\du}{6.237914\du}}
\pgfpathlineto{\pgfpoint{-14.309224\du}{6.241497\du}}
\pgfpathlineto{\pgfpoint{-14.309224\du}{6.241497\du}}
\pgfpathlineto{\pgfpoint{-14.308925\du}{6.244482\du}}
\pgfpathlineto{\pgfpoint{-14.308328\du}{6.247169\du}}
\pgfpathlineto{\pgfpoint{-14.306835\du}{6.250155\du}}
\pgfpathlineto{\pgfpoint{-14.305044\du}{6.253738\du}}
\pgfpathlineto{\pgfpoint{-14.302656\du}{6.256424\du}}
\pgfpathlineto{\pgfpoint{-14.300864\du}{6.259410\du}}
\pgfpathlineto{\pgfpoint{-14.297282\du}{6.262097\du}}
\pgfpathlineto{\pgfpoint{-14.293998\du}{6.264784\du}}
\pgfpathlineto{\pgfpoint{-14.289818\du}{6.267471\du}}
\pgfpathlineto{\pgfpoint{-14.285340\du}{6.270158\du}}
\pgfpathlineto{\pgfpoint{-14.280563\du}{6.272845\du}}
\pgfpathlineto{\pgfpoint{-14.275189\du}{6.275233\du}}
\pgfpathlineto{\pgfpoint{-14.269516\du}{6.277920\du}}
\pgfpathlineto{\pgfpoint{-14.263247\du}{6.280010\du}}
\pgfpathlineto{\pgfpoint{-14.257276\du}{6.282399\du}}
\pgfpathlineto{\pgfpoint{-14.249812\du}{6.284787\du}}
\pgfpathlineto{\pgfpoint{-14.242945\du}{6.286578\du}}
\pgfpathlineto{\pgfpoint{-14.235482\du}{6.288071\du}}
\pgfpathlineto{\pgfpoint{-14.227719\du}{6.290161\du}}
\pgfpathlineto{\pgfpoint{-14.219658\du}{6.291952\du}}
\pgfpathlineto{\pgfpoint{-14.211299\du}{6.293146\du}}
\pgfpathlineto{\pgfpoint{-14.202342\du}{6.294938\du}}
\pgfpathlineto{\pgfpoint{-14.193983\du}{6.295833\du}}
\pgfpathlineto{\pgfpoint{-14.185026\du}{6.297326\du}}
\pgfpathlineto{\pgfpoint{-14.176368\du}{6.298222\du}}
\pgfpathlineto{\pgfpoint{-14.166516\du}{6.298520\du}}
\pgfpathlineto{\pgfpoint{-14.157261\du}{6.299715\du}}
\pgfpathlineto{\pgfpoint{-14.147707\du}{6.300013\du}}
\pgfpathlineto{\pgfpoint{-14.138154\du}{6.300312\du}}
\pgfpathlineto{\pgfpoint{-14.128301\du}{6.300909\du}}
\pgfpathlineto{\pgfpoint{-14.119046\du}{6.300909\du}}
\pgfpathlineto{\pgfpoint{-14.119046\du}{6.300909\du}}
\pgfpathlineto{\pgfpoint{-14.109493\du}{6.300909\du}}
\pgfpathlineto{\pgfpoint{-14.099342\du}{6.300312\du}}
\pgfpathlineto{\pgfpoint{-14.090385\du}{6.300013\du}}
\pgfpathlineto{\pgfpoint{-14.080832\du}{6.299715\du}}
\pgfpathlineto{\pgfpoint{-14.071278\du}{6.298520\du}}
\pgfpathlineto{\pgfpoint{-14.062023\du}{6.298222\du}}
\pgfpathlineto{\pgfpoint{-14.053066\du}{6.297326\du}}
\pgfpathlineto{\pgfpoint{-14.044408\du}{6.295833\du}}
\pgfpathlineto{\pgfpoint{-14.035452\du}{6.294938\du}}
\pgfpathlineto{\pgfpoint{-14.027092\du}{6.293146\du}}
\pgfpathlineto{\pgfpoint{-14.018733\du}{6.291952\du}}
\pgfpathlineto{\pgfpoint{-14.010075\du}{6.290161\du}}
\pgfpathlineto{\pgfpoint{-14.002312\du}{6.288071\du}}
\pgfpathlineto{\pgfpoint{-13.995147\du}{6.286578\du}}
\pgfpathlineto{\pgfpoint{-13.988280\du}{6.284787\du}}
\pgfpathlineto{\pgfpoint{-13.980817\du}{6.282399\du}}
\pgfpathlineto{\pgfpoint{-13.974846\du}{6.280010\du}}
\pgfpathlineto{\pgfpoint{-13.968576\du}{6.277920\du}}
\pgfpathlineto{\pgfpoint{-13.963202\du}{6.275233\du}}
\pgfpathlineto{\pgfpoint{-13.957530\du}{6.272845\du}}
\pgfpathlineto{\pgfpoint{-13.952753\du}{6.270158\du}}
\pgfpathlineto{\pgfpoint{-13.948275\du}{6.267471\du}}
\pgfpathlineto{\pgfpoint{-13.944095\du}{6.264784\du}}
\pgfpathlineto{\pgfpoint{-13.940811\du}{6.262097\du}}
\pgfpathlineto{\pgfpoint{-13.937527\du}{6.259410\du}}
\pgfpathlineto{\pgfpoint{-13.935138\du}{6.256424\du}}
\pgfpathlineto{\pgfpoint{-13.932750\du}{6.253738\du}}
\pgfpathlineto{\pgfpoint{-13.931257\du}{6.250155\du}}
\pgfpathlineto{\pgfpoint{-13.930063\du}{6.247169\du}}
\pgfpathlineto{\pgfpoint{-13.929167\du}{6.244482\du}}
\pgfpathlineto{\pgfpoint{-13.929167\du}{6.241497\du}}
\pgfusepath{stroke}
\pgfsetbuttcap
\pgfsetmiterjoin
\pgfsetdash{}{0pt}
\definecolor{dialinecolor}{rgb}{0.000000, 0.470588, 0.666667}
\pgfsetfillcolor{dialinecolor}
\pgfpathmoveto{\pgfpoint{-14.309224\du}{6.158201\du}}
\pgfpathlineto{\pgfpoint{-14.309224\du}{6.252842\du}}
\pgfpathlineto{\pgfpoint{-13.916628\du}{6.252842\du}}
\pgfpathlineto{\pgfpoint{-13.916628\du}{6.158201\du}}
\pgfpathlineto{\pgfpoint{-14.309224\du}{6.158201\du}}
\pgfusepath{fill}
\pgfsetbuttcap
\pgfsetmiterjoin
\pgfsetdash{}{0pt}
\definecolor{dialinecolor}{rgb}{0.000000, 0.705882, 1.000000}
\pgfsetfillcolor{dialinecolor}
\pgfpathmoveto{\pgfpoint{-13.916927\du}{6.151633\du}}
\pgfpathlineto{\pgfpoint{-13.917524\du}{6.148349\du}}
\pgfpathlineto{\pgfpoint{-13.918419\du}{6.144766\du}}
\pgfpathlineto{\pgfpoint{-13.919315\du}{6.141781\du}}
\pgfpathlineto{\pgfpoint{-13.921106\du}{6.138496\du}}
\pgfpathlineto{\pgfpoint{-13.923495\du}{6.135212\du}}
\pgfpathlineto{\pgfpoint{-13.926182\du}{6.131928\du}}
\pgfpathlineto{\pgfpoint{-13.929167\du}{6.128943\du}}
\pgfpathlineto{\pgfpoint{-13.932750\du}{6.125659\du}}
\pgfpathlineto{\pgfpoint{-13.936930\du}{6.122972\du}}
\pgfpathlineto{\pgfpoint{-13.941706\du}{6.119986\du}}
\pgfpathlineto{\pgfpoint{-13.946782\du}{6.117001\du}}
\pgfpathlineto{\pgfpoint{-13.952454\du}{6.114314\du}}
\pgfpathlineto{\pgfpoint{-13.958127\du}{6.111030\du}}
\pgfpathlineto{\pgfpoint{-13.964396\du}{6.109238\du}}
\pgfpathlineto{\pgfpoint{-13.970964\du}{6.106551\du}}
\pgfpathlineto{\pgfpoint{-13.978130\du}{6.104163\du}}
\pgfpathlineto{\pgfpoint{-13.985594\du}{6.102073\du}}
\pgfpathlineto{\pgfpoint{-13.993057\du}{6.099983\du}}
\pgfpathlineto{\pgfpoint{-14.001417\du}{6.097893\du}}
\pgfpathlineto{\pgfpoint{-14.009776\du}{6.095804\du}}
\pgfpathlineto{\pgfpoint{-14.017837\du}{6.094311\du}}
\pgfpathlineto{\pgfpoint{-14.027092\du}{6.092818\du}}
\pgfpathlineto{\pgfpoint{-14.035750\du}{6.091624\du}}
\pgfpathlineto{\pgfpoint{-14.045304\du}{6.090131\du}}
\pgfpathlineto{\pgfpoint{-14.054559\du}{6.089235\du}}
\pgfpathlineto{\pgfpoint{-14.063814\du}{6.088041\du}}
\pgfpathlineto{\pgfpoint{-14.073666\du}{6.087444\du}}
\pgfpathlineto{\pgfpoint{-14.083519\du}{6.086847\du}}
\pgfpathlineto{\pgfpoint{-14.093371\du}{6.086548\du}}
\pgfpathlineto{\pgfpoint{-14.103522\du}{6.086250\du}}
\pgfpathlineto{\pgfpoint{-14.113374\du}{6.086250\du}}
\pgfpathlineto{\pgfpoint{-14.113374\du}{6.086250\du}}
\pgfpathlineto{\pgfpoint{-14.123226\du}{6.086250\du}}
\pgfpathlineto{\pgfpoint{-14.133078\du}{6.086548\du}}
\pgfpathlineto{\pgfpoint{-14.143229\du}{6.086847\du}}
\pgfpathlineto{\pgfpoint{-14.153081\du}{6.087444\du}}
\pgfpathlineto{\pgfpoint{-14.162635\du}{6.088041\du}}
\pgfpathlineto{\pgfpoint{-14.172189\du}{6.089235\du}}
\pgfpathlineto{\pgfpoint{-14.181444\du}{6.090131\du}}
\pgfpathlineto{\pgfpoint{-14.190699\du}{6.091624\du}}
\pgfpathlineto{\pgfpoint{-14.199655\du}{6.092818\du}}
\pgfpathlineto{\pgfpoint{-14.208612\du}{6.094311\du}}
\pgfpathlineto{\pgfpoint{-14.216971\du}{6.095804\du}}
\pgfpathlineto{\pgfpoint{-14.225331\du}{6.097893\du}}
\pgfpathlineto{\pgfpoint{-14.233690\du}{6.099983\du}}
\pgfpathlineto{\pgfpoint{-14.240856\du}{6.102073\du}}
\pgfpathlineto{\pgfpoint{-14.248319\du}{6.104163\du}}
\pgfpathlineto{\pgfpoint{-14.255485\du}{6.106551\du}}
\pgfpathlineto{\pgfpoint{-14.262351\du}{6.109238\du}}
\pgfpathlineto{\pgfpoint{-14.268322\du}{6.111030\du}}
\pgfpathlineto{\pgfpoint{-14.274293\du}{6.114314\du}}
\pgfpathlineto{\pgfpoint{-14.279966\du}{6.117001\du}}
\pgfpathlineto{\pgfpoint{-14.285041\du}{6.119986\du}}
\pgfpathlineto{\pgfpoint{-14.289818\du}{6.122972\du}}
\pgfpathlineto{\pgfpoint{-14.293998\du}{6.125659\du}}
\pgfpathlineto{\pgfpoint{-14.297282\du}{6.128943\du}}
\pgfpathlineto{\pgfpoint{-14.300566\du}{6.131928\du}}
\pgfpathlineto{\pgfpoint{-14.303551\du}{6.135212\du}}
\pgfpathlineto{\pgfpoint{-14.305641\du}{6.138496\du}}
\pgfpathlineto{\pgfpoint{-14.307134\du}{6.141781\du}}
\pgfpathlineto{\pgfpoint{-14.308925\du}{6.144766\du}}
\pgfpathlineto{\pgfpoint{-14.309224\du}{6.148349\du}}
\pgfpathlineto{\pgfpoint{-14.309522\du}{6.151633\du}}
\pgfpathlineto{\pgfpoint{-14.309522\du}{6.151633\du}}
\pgfpathlineto{\pgfpoint{-14.309224\du}{6.154917\du}}
\pgfpathlineto{\pgfpoint{-14.308925\du}{6.158499\du}}
\pgfpathlineto{\pgfpoint{-14.307134\du}{6.161485\du}}
\pgfpathlineto{\pgfpoint{-14.305641\du}{6.164769\du}}
\pgfpathlineto{\pgfpoint{-14.303551\du}{6.168352\du}}
\pgfpathlineto{\pgfpoint{-14.300566\du}{6.171337\du}}
\pgfpathlineto{\pgfpoint{-14.297282\du}{6.174323\du}}
\pgfpathlineto{\pgfpoint{-14.293998\du}{6.177607\du}}
\pgfpathlineto{\pgfpoint{-14.289818\du}{6.179995\du}}
\pgfpathlineto{\pgfpoint{-14.285041\du}{6.183279\du}}
\pgfpathlineto{\pgfpoint{-14.279966\du}{6.186563\du}}
\pgfpathlineto{\pgfpoint{-14.274293\du}{6.189250\du}}
\pgfpathlineto{\pgfpoint{-14.268322\du}{6.191937\du}}
\pgfpathlineto{\pgfpoint{-14.262351\du}{6.194027\du}}
\pgfpathlineto{\pgfpoint{-14.255485\du}{6.197013\du}}
\pgfpathlineto{\pgfpoint{-14.248319\du}{6.199401\du}}
\pgfpathlineto{\pgfpoint{-14.240856\du}{6.201192\du}}
\pgfpathlineto{\pgfpoint{-14.233690\du}{6.203282\du}}
\pgfpathlineto{\pgfpoint{-14.225331\du}{6.205372\du}}
\pgfpathlineto{\pgfpoint{-14.216971\du}{6.207462\du}}
\pgfpathlineto{\pgfpoint{-14.208612\du}{6.208955\du}}
\pgfpathlineto{\pgfpoint{-14.199655\du}{6.210447\du}}
\pgfpathlineto{\pgfpoint{-14.190699\du}{6.211642\du}}
\pgfpathlineto{\pgfpoint{-14.181444\du}{6.213134\du}}
\pgfpathlineto{\pgfpoint{-14.172189\du}{6.214030\du}}
\pgfpathlineto{\pgfpoint{-14.162635\du}{6.215224\du}}
\pgfpathlineto{\pgfpoint{-14.153081\du}{6.215523\du}}
\pgfpathlineto{\pgfpoint{-14.143229\du}{6.216419\du}}
\pgfpathlineto{\pgfpoint{-14.133078\du}{6.216419\du}}
\pgfpathlineto{\pgfpoint{-14.123226\du}{6.217314\du}}
\pgfpathlineto{\pgfpoint{-14.113374\du}{6.217314\du}}
\pgfpathlineto{\pgfpoint{-14.113374\du}{6.217314\du}}
\pgfpathlineto{\pgfpoint{-14.103522\du}{6.217314\du}}
\pgfpathlineto{\pgfpoint{-14.093371\du}{6.216419\du}}
\pgfpathlineto{\pgfpoint{-14.083519\du}{6.216419\du}}
\pgfpathlineto{\pgfpoint{-14.073666\du}{6.215523\du}}
\pgfpathlineto{\pgfpoint{-14.063814\du}{6.215224\du}}
\pgfpathlineto{\pgfpoint{-14.054559\du}{6.214030\du}}
\pgfpathlineto{\pgfpoint{-14.045304\du}{6.213134\du}}
\pgfpathlineto{\pgfpoint{-14.035750\du}{6.211642\du}}
\pgfpathlineto{\pgfpoint{-14.027092\du}{6.210447\du}}
\pgfpathlineto{\pgfpoint{-14.017837\du}{6.208955\du}}
\pgfpathlineto{\pgfpoint{-14.009776\du}{6.207462\du}}
\pgfpathlineto{\pgfpoint{-14.001417\du}{6.205372\du}}
\pgfpathlineto{\pgfpoint{-13.993057\du}{6.203282\du}}
\pgfpathlineto{\pgfpoint{-13.985594\du}{6.201192\du}}
\pgfpathlineto{\pgfpoint{-13.978130\du}{6.199401\du}}
\pgfpathlineto{\pgfpoint{-13.970964\du}{6.197013\du}}
\pgfpathlineto{\pgfpoint{-13.964396\du}{6.194027\du}}
\pgfpathlineto{\pgfpoint{-13.958127\du}{6.191937\du}}
\pgfpathlineto{\pgfpoint{-13.952454\du}{6.189250\du}}
\pgfpathlineto{\pgfpoint{-13.946782\du}{6.186563\du}}
\pgfpathlineto{\pgfpoint{-13.941706\du}{6.183279\du}}
\pgfpathlineto{\pgfpoint{-13.936930\du}{6.179995\du}}
\pgfpathlineto{\pgfpoint{-13.932750\du}{6.177607\du}}
\pgfpathlineto{\pgfpoint{-13.929167\du}{6.174323\du}}
\pgfpathlineto{\pgfpoint{-13.926182\du}{6.171337\du}}
\pgfpathlineto{\pgfpoint{-13.923495\du}{6.168352\du}}
\pgfpathlineto{\pgfpoint{-13.921106\du}{6.164769\du}}
\pgfpathlineto{\pgfpoint{-13.919315\du}{6.161485\du}}
\pgfpathlineto{\pgfpoint{-13.918419\du}{6.158499\du}}
\pgfpathlineto{\pgfpoint{-13.917524\du}{6.154917\du}}
\pgfpathlineto{\pgfpoint{-13.916927\du}{6.151633\du}}
\pgfusepath{fill}
\pgfsetbuttcap
\pgfsetmiterjoin
\pgfsetdash{}{0pt}
\definecolor{dialinecolor}{rgb}{0.666667, 0.901961, 1.000000}
\pgfsetstrokecolor{dialinecolor}
\pgfpathmoveto{\pgfpoint{-13.929167\du}{6.145960\du}}
\pgfpathlineto{\pgfpoint{-13.929167\du}{6.142676\du}}
\pgfpathlineto{\pgfpoint{-13.930063\du}{6.140288\du}}
\pgfpathlineto{\pgfpoint{-13.931257\du}{6.137004\du}}
\pgfpathlineto{\pgfpoint{-13.932750\du}{6.133421\du}}
\pgfpathlineto{\pgfpoint{-13.935138\du}{6.131033\du}}
\pgfpathlineto{\pgfpoint{-13.937527\du}{6.127749\du}}
\pgfpathlineto{\pgfpoint{-13.940811\du}{6.125360\du}}
\pgfpathlineto{\pgfpoint{-13.944095\du}{6.122673\du}}
\pgfpathlineto{\pgfpoint{-13.948275\du}{6.119986\du}}
\pgfpathlineto{\pgfpoint{-13.952753\du}{6.117299\du}}
\pgfpathlineto{\pgfpoint{-13.957530\du}{6.114612\du}}
\pgfpathlineto{\pgfpoint{-13.963202\du}{6.111925\du}}
\pgfpathlineto{\pgfpoint{-13.968576\du}{6.109537\du}}
\pgfpathlineto{\pgfpoint{-13.974846\du}{6.107447\du}}
\pgfpathlineto{\pgfpoint{-13.980817\du}{6.105059\du}}
\pgfpathlineto{\pgfpoint{-13.988280\du}{6.102670\du}}
\pgfpathlineto{\pgfpoint{-13.995147\du}{6.100879\du}}
\pgfpathlineto{\pgfpoint{-14.002312\du}{6.099386\du}}
\pgfpathlineto{\pgfpoint{-14.010075\du}{6.096998\du}}
\pgfpathlineto{\pgfpoint{-14.018733\du}{6.095505\du}}
\pgfpathlineto{\pgfpoint{-14.027092\du}{6.094311\du}}
\pgfpathlineto{\pgfpoint{-14.035452\du}{6.092519\du}}
\pgfpathlineto{\pgfpoint{-14.044408\du}{6.091624\du}}
\pgfpathlineto{\pgfpoint{-14.053066\du}{6.090131\du}}
\pgfpathlineto{\pgfpoint{-14.062023\du}{6.089235\du}}
\pgfpathlineto{\pgfpoint{-14.071278\du}{6.088638\du}}
\pgfpathlineto{\pgfpoint{-14.080832\du}{6.087743\du}}
\pgfpathlineto{\pgfpoint{-14.090385\du}{6.087444\du}}
\pgfpathlineto{\pgfpoint{-14.099342\du}{6.087145\du}}
\pgfpathlineto{\pgfpoint{-14.109493\du}{6.086548\du}}
\pgfpathlineto{\pgfpoint{-14.119046\du}{6.086548\du}}
\pgfpathlineto{\pgfpoint{-14.119046\du}{6.086548\du}}
\pgfpathlineto{\pgfpoint{-14.128301\du}{6.086548\du}}
\pgfpathlineto{\pgfpoint{-14.138154\du}{6.087145\du}}
\pgfpathlineto{\pgfpoint{-14.147707\du}{6.087444\du}}
\pgfpathlineto{\pgfpoint{-14.157261\du}{6.087743\du}}
\pgfpathlineto{\pgfpoint{-14.166516\du}{6.088638\du}}
\pgfpathlineto{\pgfpoint{-14.176368\du}{6.089235\du}}
\pgfpathlineto{\pgfpoint{-14.185026\du}{6.090131\du}}
\pgfpathlineto{\pgfpoint{-14.193983\du}{6.091624\du}}
\pgfpathlineto{\pgfpoint{-14.202342\du}{6.092519\du}}
\pgfpathlineto{\pgfpoint{-14.211299\du}{6.094311\du}}
\pgfpathlineto{\pgfpoint{-14.219658\du}{6.095505\du}}
\pgfpathlineto{\pgfpoint{-14.227719\du}{6.096998\du}}
\pgfpathlineto{\pgfpoint{-14.235482\du}{6.099386\du}}
\pgfpathlineto{\pgfpoint{-14.242945\du}{6.100879\du}}
\pgfpathlineto{\pgfpoint{-14.249812\du}{6.102670\du}}
\pgfpathlineto{\pgfpoint{-14.257276\du}{6.105059\du}}
\pgfpathlineto{\pgfpoint{-14.263247\du}{6.107447\du}}
\pgfpathlineto{\pgfpoint{-14.269516\du}{6.109537\du}}
\pgfpathlineto{\pgfpoint{-14.275189\du}{6.111925\du}}
\pgfpathlineto{\pgfpoint{-14.280563\du}{6.114612\du}}
\pgfpathlineto{\pgfpoint{-14.285340\du}{6.117299\du}}
\pgfpathlineto{\pgfpoint{-14.289818\du}{6.119986\du}}
\pgfpathlineto{\pgfpoint{-14.293998\du}{6.122673\du}}
\pgfpathlineto{\pgfpoint{-14.297282\du}{6.125360\du}}
\pgfpathlineto{\pgfpoint{-14.300864\du}{6.127749\du}}
\pgfpathlineto{\pgfpoint{-14.302656\du}{6.131033\du}}
\pgfpathlineto{\pgfpoint{-14.305044\du}{6.133421\du}}
\pgfpathlineto{\pgfpoint{-14.306835\du}{6.137004\du}}
\pgfpathlineto{\pgfpoint{-14.308328\du}{6.140288\du}}
\pgfpathlineto{\pgfpoint{-14.308925\du}{6.142676\du}}
\pgfpathlineto{\pgfpoint{-14.309224\du}{6.145960\du}}
\pgfpathlineto{\pgfpoint{-14.309224\du}{6.145960\du}}
\pgfpathlineto{\pgfpoint{-14.308925\du}{6.149244\du}}
\pgfpathlineto{\pgfpoint{-14.308328\du}{6.151931\du}}
\pgfpathlineto{\pgfpoint{-14.306835\du}{6.154917\du}}
\pgfpathlineto{\pgfpoint{-14.305044\du}{6.158201\du}}
\pgfpathlineto{\pgfpoint{-14.302656\du}{6.160888\du}}
\pgfpathlineto{\pgfpoint{-14.300864\du}{6.164172\du}}
\pgfpathlineto{\pgfpoint{-14.297282\du}{6.166859\du}}
\pgfpathlineto{\pgfpoint{-14.293998\du}{6.169546\du}}
\pgfpathlineto{\pgfpoint{-14.289818\du}{6.172233\du}}
\pgfpathlineto{\pgfpoint{-14.285340\du}{6.174920\du}}
\pgfpathlineto{\pgfpoint{-14.280563\du}{6.177607\du}}
\pgfpathlineto{\pgfpoint{-14.275189\du}{6.179697\du}}
\pgfpathlineto{\pgfpoint{-14.269516\du}{6.182384\du}}
\pgfpathlineto{\pgfpoint{-14.263247\du}{6.184772\du}}
\pgfpathlineto{\pgfpoint{-14.257276\du}{6.187160\du}}
\pgfpathlineto{\pgfpoint{-14.249812\du}{6.189250\du}}
\pgfpathlineto{\pgfpoint{-14.242945\du}{6.191042\du}}
\pgfpathlineto{\pgfpoint{-14.235482\du}{6.192833\du}}
\pgfpathlineto{\pgfpoint{-14.227719\du}{6.194624\du}}
\pgfpathlineto{\pgfpoint{-14.219658\du}{6.196714\du}}
\pgfpathlineto{\pgfpoint{-14.211299\du}{6.197908\du}}
\pgfpathlineto{\pgfpoint{-14.202342\du}{6.199700\du}}
\pgfpathlineto{\pgfpoint{-14.193983\du}{6.200595\du}}
\pgfpathlineto{\pgfpoint{-14.185026\du}{6.202088\du}}
\pgfpathlineto{\pgfpoint{-14.176368\du}{6.202387\du}}
\pgfpathlineto{\pgfpoint{-14.166516\du}{6.203581\du}}
\pgfpathlineto{\pgfpoint{-14.157261\du}{6.203879\du}}
\pgfpathlineto{\pgfpoint{-14.147707\du}{6.204775\du}}
\pgfpathlineto{\pgfpoint{-14.138154\du}{6.205074\du}}
\pgfpathlineto{\pgfpoint{-14.128301\du}{6.205372\du}}
\pgfpathlineto{\pgfpoint{-14.119046\du}{6.205372\du}}
\pgfpathlineto{\pgfpoint{-14.119046\du}{6.205372\du}}
\pgfpathlineto{\pgfpoint{-14.109493\du}{6.205372\du}}
\pgfpathlineto{\pgfpoint{-14.099342\du}{6.205074\du}}
\pgfpathlineto{\pgfpoint{-14.090385\du}{6.204775\du}}
\pgfpathlineto{\pgfpoint{-14.080832\du}{6.203879\du}}
\pgfpathlineto{\pgfpoint{-14.071278\du}{6.203581\du}}
\pgfpathlineto{\pgfpoint{-14.062023\du}{6.202387\du}}
\pgfpathlineto{\pgfpoint{-14.053066\du}{6.202088\du}}
\pgfpathlineto{\pgfpoint{-14.044408\du}{6.200595\du}}
\pgfpathlineto{\pgfpoint{-14.035452\du}{6.199700\du}}
\pgfpathlineto{\pgfpoint{-14.027092\du}{6.197908\du}}
\pgfpathlineto{\pgfpoint{-14.018733\du}{6.196714\du}}
\pgfpathlineto{\pgfpoint{-14.010075\du}{6.194624\du}}
\pgfpathlineto{\pgfpoint{-14.002312\du}{6.192833\du}}
\pgfpathlineto{\pgfpoint{-13.995147\du}{6.191042\du}}
\pgfpathlineto{\pgfpoint{-13.988280\du}{6.189250\du}}
\pgfpathlineto{\pgfpoint{-13.980817\du}{6.187160\du}}
\pgfpathlineto{\pgfpoint{-13.974846\du}{6.184772\du}}
\pgfpathlineto{\pgfpoint{-13.968576\du}{6.182384\du}}
\pgfpathlineto{\pgfpoint{-13.963202\du}{6.179697\du}}
\pgfpathlineto{\pgfpoint{-13.957530\du}{6.177607\du}}
\pgfpathlineto{\pgfpoint{-13.952753\du}{6.174920\du}}
\pgfpathlineto{\pgfpoint{-13.948275\du}{6.172233\du}}
\pgfpathlineto{\pgfpoint{-13.944095\du}{6.169546\du}}
\pgfpathlineto{\pgfpoint{-13.940811\du}{6.166859\du}}
\pgfpathlineto{\pgfpoint{-13.937527\du}{6.164172\du}}
\pgfpathlineto{\pgfpoint{-13.935138\du}{6.160888\du}}
\pgfpathlineto{\pgfpoint{-13.932750\du}{6.158201\du}}
\pgfpathlineto{\pgfpoint{-13.931257\du}{6.154917\du}}
\pgfpathlineto{\pgfpoint{-13.930063\du}{6.151931\du}}
\pgfpathlineto{\pgfpoint{-13.929167\du}{6.149244\du}}
\pgfpathlineto{\pgfpoint{-13.929167\du}{6.145960\du}}
\pgfusepath{stroke}
\pgfsetbuttcap
\pgfsetmiterjoin
\pgfsetdash{}{0pt}
\definecolor{dialinecolor}{rgb}{0.000000, 0.000000, 0.000000}
\pgfsetfillcolor{dialinecolor}
\pgfpathmoveto{\pgfpoint{-14.106806\du}{6.134317\du}}
\pgfpathlineto{\pgfpoint{-14.083519\du}{6.145960\du}}
\pgfpathlineto{\pgfpoint{-14.011568\du}{6.122375\du}}
\pgfpathlineto{\pgfpoint{-13.976040\du}{6.134317\du}}
\pgfpathlineto{\pgfpoint{-13.999625\du}{6.110433\du}}
\pgfpathlineto{\pgfpoint{-14.083519\du}{6.110433\du}}
\pgfpathlineto{\pgfpoint{-14.047692\du}{6.110433\du}}
\pgfpathlineto{\pgfpoint{-14.106806\du}{6.134317\du}}
\pgfusepath{fill}
\pgfsetbuttcap
\pgfsetmiterjoin
\pgfsetdash{}{0pt}
\definecolor{dialinecolor}{rgb}{0.000000, 0.000000, 0.000000}
\pgfsetfillcolor{dialinecolor}
\pgfpathmoveto{\pgfpoint{-14.119046\du}{6.169844\du}}
\pgfpathlineto{\pgfpoint{-14.154873\du}{6.158201\du}}
\pgfpathlineto{\pgfpoint{-14.213986\du}{6.182085\du}}
\pgfpathlineto{\pgfpoint{-14.249812\du}{6.169844\du}}
\pgfpathlineto{\pgfpoint{-14.237571\du}{6.205372\du}}
\pgfpathlineto{\pgfpoint{-14.154873\du}{6.205372\du}}
\pgfpathlineto{\pgfpoint{-14.190400\du}{6.194027\du}}
\pgfpathlineto{\pgfpoint{-14.119046\du}{6.169844\du}}
\pgfusepath{fill}
\pgfsetbuttcap
\pgfsetmiterjoin
\pgfsetdash{}{0pt}
\definecolor{dialinecolor}{rgb}{0.000000, 0.000000, 0.000000}
\pgfsetfillcolor{dialinecolor}
\pgfpathmoveto{\pgfpoint{-14.237571\du}{6.110433\du}}
\pgfpathlineto{\pgfpoint{-14.213986\du}{6.097893\du}}
\pgfpathlineto{\pgfpoint{-14.142333\du}{6.134317\du}}
\pgfpathlineto{\pgfpoint{-14.119046\du}{6.122375\du}}
\pgfpathlineto{\pgfpoint{-14.130690\du}{6.145960\du}}
\pgfpathlineto{\pgfpoint{-14.213986\du}{6.145960\du}}
\pgfpathlineto{\pgfpoint{-14.178160\du}{6.134317\du}}
\pgfpathlineto{\pgfpoint{-14.237571\du}{6.110433\du}}
\pgfusepath{fill}
\pgfsetbuttcap
\pgfsetmiterjoin
\pgfsetdash{}{0pt}
\definecolor{dialinecolor}{rgb}{0.000000, 0.000000, 0.000000}
\pgfsetfillcolor{dialinecolor}
\pgfpathmoveto{\pgfpoint{-13.988280\du}{6.194027\du}}
\pgfpathlineto{\pgfpoint{-14.011568\du}{6.205372\du}}
\pgfpathlineto{\pgfpoint{-14.083519\du}{6.169844\du}}
\pgfpathlineto{\pgfpoint{-14.119046\du}{6.182085\du}}
\pgfpathlineto{\pgfpoint{-14.095162\du}{6.158201\du}}
\pgfpathlineto{\pgfpoint{-14.011568\du}{6.158201\du}}
\pgfpathlineto{\pgfpoint{-14.047692\du}{6.169844\du}}
\pgfpathlineto{\pgfpoint{-13.988280\du}{6.194027\du}}
\pgfusepath{fill}
\pgfsetbuttcap
\pgfsetmiterjoin
\pgfsetdash{}{0pt}
\definecolor{dialinecolor}{rgb}{1.000000, 1.000000, 1.000000}
\pgfsetfillcolor{dialinecolor}
\pgfpathmoveto{\pgfpoint{-14.106806\du}{6.145960\du}}
\pgfpathlineto{\pgfpoint{-14.083519\du}{6.145960\du}}
\pgfpathlineto{\pgfpoint{-14.011568\du}{6.122375\du}}
\pgfpathlineto{\pgfpoint{-13.976040\du}{6.134317\du}}
\pgfpathlineto{\pgfpoint{-13.999625\du}{6.110433\du}}
\pgfpathlineto{\pgfpoint{-14.071278\du}{6.110433\du}}
\pgfpathlineto{\pgfpoint{-14.047692\du}{6.110433\du}}
\pgfpathlineto{\pgfpoint{-14.106806\du}{6.145960\du}}
\pgfusepath{fill}
\pgfsetbuttcap
\pgfsetmiterjoin
\pgfsetdash{}{0pt}
\definecolor{dialinecolor}{rgb}{1.000000, 1.000000, 1.000000}
\pgfsetfillcolor{dialinecolor}
\pgfpathmoveto{\pgfpoint{-14.119046\du}{6.169844\du}}
\pgfpathlineto{\pgfpoint{-14.142333\du}{6.158201\du}}
\pgfpathlineto{\pgfpoint{-14.213986\du}{6.182085\du}}
\pgfpathlineto{\pgfpoint{-14.249812\du}{6.182085\du}}
\pgfpathlineto{\pgfpoint{-14.226226\du}{6.205372\du}}
\pgfpathlineto{\pgfpoint{-14.142333\du}{6.205372\du}}
\pgfpathlineto{\pgfpoint{-14.178160\du}{6.194027\du}}
\pgfpathlineto{\pgfpoint{-14.119046\du}{6.169844\du}}
\pgfusepath{fill}
\pgfsetbuttcap
\pgfsetmiterjoin
\pgfsetdash{}{0pt}
\definecolor{dialinecolor}{rgb}{1.000000, 1.000000, 1.000000}
\pgfsetfillcolor{dialinecolor}
\pgfpathmoveto{\pgfpoint{-14.237571\du}{6.110433\du}}
\pgfpathlineto{\pgfpoint{-14.213986\du}{6.110433\du}}
\pgfpathlineto{\pgfpoint{-14.142333\du}{6.134317\du}}
\pgfpathlineto{\pgfpoint{-14.106806\du}{6.122375\du}}
\pgfpathlineto{\pgfpoint{-14.130690\du}{6.145960\du}}
\pgfpathlineto{\pgfpoint{-14.213986\du}{6.145960\du}}
\pgfpathlineto{\pgfpoint{-14.178160\du}{6.145960\du}}
\pgfpathlineto{\pgfpoint{-14.237571\du}{6.110433\du}}
\pgfusepath{fill}
\pgfsetbuttcap
\pgfsetmiterjoin
\pgfsetdash{}{0pt}
\definecolor{dialinecolor}{rgb}{1.000000, 1.000000, 1.000000}
\pgfsetfillcolor{dialinecolor}
\pgfpathmoveto{\pgfpoint{-13.988280\du}{6.194027\du}}
\pgfpathlineto{\pgfpoint{-14.011568\du}{6.205372\du}}
\pgfpathlineto{\pgfpoint{-14.083519\du}{6.182085\du}}
\pgfpathlineto{\pgfpoint{-14.106806\du}{6.182085\du}}
\pgfpathlineto{\pgfpoint{-14.095162\du}{6.158201\du}}
\pgfpathlineto{\pgfpoint{-14.011568\du}{6.158201\du}}
\pgfpathlineto{\pgfpoint{-14.047692\du}{6.169844\du}}
\pgfpathlineto{\pgfpoint{-13.988280\du}{6.194027\du}}
\pgfusepath{fill}
\pgfsetbuttcap
\pgfsetmiterjoin
\pgfsetdash{}{0pt}
\definecolor{dialinecolor}{rgb}{0.666667, 0.901961, 1.000000}
\pgfsetstrokecolor{dialinecolor}
\pgfpathmoveto{\pgfpoint{-14.309224\du}{6.145960\du}}
\pgfpathlineto{\pgfpoint{-14.309224\du}{6.240900\du}}
\pgfusepath{stroke}
\pgfsetbuttcap
\pgfsetmiterjoin
\pgfsetdash{}{0pt}
\definecolor{dialinecolor}{rgb}{0.666667, 0.901961, 1.000000}
\pgfsetstrokecolor{dialinecolor}
\pgfpathmoveto{\pgfpoint{-13.928570\du}{6.145960\du}}
\pgfpathlineto{\pgfpoint{-13.928570\du}{6.240900\du}}
\pgfusepath{stroke}
\pgfsetbuttcap
\pgfsetmiterjoin
\pgfsetdash{}{0pt}
\definecolor{dialinecolor}{rgb}{0.000000, 0.470588, 0.666667}
\pgfsetfillcolor{dialinecolor}
\pgfpathmoveto{\pgfpoint{-13.202790\du}{6.247169\du}}
\pgfpathlineto{\pgfpoint{-13.202790\du}{6.243885\du}}
\pgfpathlineto{\pgfpoint{-13.203686\du}{6.240601\du}}
\pgfpathlineto{\pgfpoint{-13.204880\du}{6.237317\du}}
\pgfpathlineto{\pgfpoint{-13.206373\du}{6.234033\du}}
\pgfpathlineto{\pgfpoint{-13.208761\du}{6.231048\du}}
\pgfpathlineto{\pgfpoint{-13.211448\du}{6.227764\du}}
\pgfpathlineto{\pgfpoint{-13.214434\du}{6.224479\du}}
\pgfpathlineto{\pgfpoint{-13.218613\du}{6.221494\du}}
\pgfpathlineto{\pgfpoint{-13.222196\du}{6.218508\du}}
\pgfpathlineto{\pgfpoint{-13.226973\du}{6.215523\du}}
\pgfpathlineto{\pgfpoint{-13.232048\du}{6.212836\du}}
\pgfpathlineto{\pgfpoint{-13.237422\du}{6.209850\du}}
\pgfpathlineto{\pgfpoint{-13.243692\du}{6.207462\du}}
\pgfpathlineto{\pgfpoint{-13.249663\du}{6.204775\du}}
\pgfpathlineto{\pgfpoint{-13.256529\du}{6.202387\du}}
\pgfpathlineto{\pgfpoint{-13.263396\du}{6.199998\du}}
\pgfpathlineto{\pgfpoint{-13.270860\du}{6.197908\du}}
\pgfpathlineto{\pgfpoint{-13.278324\du}{6.195520\du}}
\pgfpathlineto{\pgfpoint{-13.286385\du}{6.194027\du}}
\pgfpathlineto{\pgfpoint{-13.295043\du}{6.191937\du}}
\pgfpathlineto{\pgfpoint{-13.303402\du}{6.190146\du}}
\pgfpathlineto{\pgfpoint{-13.312060\du}{6.188952\du}}
\pgfpathlineto{\pgfpoint{-13.321315\du}{6.187459\du}}
\pgfpathlineto{\pgfpoint{-13.330570\du}{6.186265\du}}
\pgfpathlineto{\pgfpoint{-13.339825\du}{6.185071\du}}
\pgfpathlineto{\pgfpoint{-13.349379\du}{6.184175\du}}
\pgfpathlineto{\pgfpoint{-13.359530\du}{6.183279\du}}
\pgfpathlineto{\pgfpoint{-13.369084\du}{6.182682\du}}
\pgfpathlineto{\pgfpoint{-13.378637\du}{6.182384\du}}
\pgfpathlineto{\pgfpoint{-13.388788\du}{6.182085\du}}
\pgfpathlineto{\pgfpoint{-13.398640\du}{6.182085\du}}
\pgfpathlineto{\pgfpoint{-13.398640\du}{6.182085\du}}
\pgfpathlineto{\pgfpoint{-13.408492\du}{6.182085\du}}
\pgfpathlineto{\pgfpoint{-13.418643\du}{6.182384\du}}
\pgfpathlineto{\pgfpoint{-13.428495\du}{6.182682\du}}
\pgfpathlineto{\pgfpoint{-13.438348\du}{6.183279\du}}
\pgfpathlineto{\pgfpoint{-13.447901\du}{6.184175\du}}
\pgfpathlineto{\pgfpoint{-13.457455\du}{6.185071\du}}
\pgfpathlineto{\pgfpoint{-13.467009\du}{6.186265\du}}
\pgfpathlineto{\pgfpoint{-13.475965\du}{6.187459\du}}
\pgfpathlineto{\pgfpoint{-13.485220\du}{6.188952\du}}
\pgfpathlineto{\pgfpoint{-13.493878\du}{6.190146\du}}
\pgfpathlineto{\pgfpoint{-13.502536\du}{6.191937\du}}
\pgfpathlineto{\pgfpoint{-13.510896\du}{6.194027\du}}
\pgfpathlineto{\pgfpoint{-13.518957\du}{6.195520\du}}
\pgfpathlineto{\pgfpoint{-13.526421\du}{6.197908\du}}
\pgfpathlineto{\pgfpoint{-13.533884\du}{6.199998\du}}
\pgfpathlineto{\pgfpoint{-13.541050\du}{6.202387\du}}
\pgfpathlineto{\pgfpoint{-13.547618\du}{6.204775\du}}
\pgfpathlineto{\pgfpoint{-13.553887\du}{6.207462\du}}
\pgfpathlineto{\pgfpoint{-13.559858\du}{6.209850\du}}
\pgfpathlineto{\pgfpoint{-13.565232\du}{6.212836\du}}
\pgfpathlineto{\pgfpoint{-13.570308\du}{6.215523\du}}
\pgfpathlineto{\pgfpoint{-13.575084\du}{6.218508\du}}
\pgfpathlineto{\pgfpoint{-13.578966\du}{6.221494\du}}
\pgfpathlineto{\pgfpoint{-13.582847\du}{6.224479\du}}
\pgfpathlineto{\pgfpoint{-13.585832\du}{6.227764\du}}
\pgfpathlineto{\pgfpoint{-13.588519\du}{6.231048\du}}
\pgfpathlineto{\pgfpoint{-13.590908\du}{6.234033\du}}
\pgfpathlineto{\pgfpoint{-13.592699\du}{6.237317\du}}
\pgfpathlineto{\pgfpoint{-13.594192\du}{6.240601\du}}
\pgfpathlineto{\pgfpoint{-13.595087\du}{6.243885\du}}
\pgfpathlineto{\pgfpoint{-13.595087\du}{6.247169\du}}
\pgfpathlineto{\pgfpoint{-13.595087\du}{6.247169\du}}
\pgfpathlineto{\pgfpoint{-13.595087\du}{6.250752\du}}
\pgfpathlineto{\pgfpoint{-13.594192\du}{6.253738\du}}
\pgfpathlineto{\pgfpoint{-13.592699\du}{6.257022\du}}
\pgfpathlineto{\pgfpoint{-13.590908\du}{6.260306\du}}
\pgfpathlineto{\pgfpoint{-13.588519\du}{6.263590\du}}
\pgfpathlineto{\pgfpoint{-13.585832\du}{6.266874\du}}
\pgfpathlineto{\pgfpoint{-13.582847\du}{6.269859\du}}
\pgfpathlineto{\pgfpoint{-13.578966\du}{6.272845\du}}
\pgfpathlineto{\pgfpoint{-13.575084\du}{6.275830\du}}
\pgfpathlineto{\pgfpoint{-13.570308\du}{6.279114\du}}
\pgfpathlineto{\pgfpoint{-13.565232\du}{6.281801\du}}
\pgfpathlineto{\pgfpoint{-13.559858\du}{6.284488\du}}
\pgfpathlineto{\pgfpoint{-13.553887\du}{6.287175\du}}
\pgfpathlineto{\pgfpoint{-13.547618\du}{6.289862\du}}
\pgfpathlineto{\pgfpoint{-13.541050\du}{6.292251\du}}
\pgfpathlineto{\pgfpoint{-13.533884\du}{6.294639\du}}
\pgfpathlineto{\pgfpoint{-13.526421\du}{6.296430\du}}
\pgfpathlineto{\pgfpoint{-13.518957\du}{6.298520\du}}
\pgfpathlineto{\pgfpoint{-13.510896\du}{6.300610\du}}
\pgfpathlineto{\pgfpoint{-13.502536\du}{6.302700\du}}
\pgfpathlineto{\pgfpoint{-13.493878\du}{6.304193\du}}
\pgfpathlineto{\pgfpoint{-13.485220\du}{6.305686\du}}
\pgfpathlineto{\pgfpoint{-13.475965\du}{6.307178\du}}
\pgfpathlineto{\pgfpoint{-13.467009\du}{6.308373\du}}
\pgfpathlineto{\pgfpoint{-13.457455\du}{6.309268\du}}
\pgfpathlineto{\pgfpoint{-13.447901\du}{6.310462\du}}
\pgfpathlineto{\pgfpoint{-13.438348\du}{6.311060\du}}
\pgfpathlineto{\pgfpoint{-13.428495\du}{6.311657\du}}
\pgfpathlineto{\pgfpoint{-13.418643\du}{6.312254\du}}
\pgfpathlineto{\pgfpoint{-13.408492\du}{6.312552\du}}
\pgfpathlineto{\pgfpoint{-13.398640\du}{6.312552\du}}
\pgfpathlineto{\pgfpoint{-13.398640\du}{6.312552\du}}
\pgfpathlineto{\pgfpoint{-13.388788\du}{6.312552\du}}
\pgfpathlineto{\pgfpoint{-13.378637\du}{6.312254\du}}
\pgfpathlineto{\pgfpoint{-13.369084\du}{6.311657\du}}
\pgfpathlineto{\pgfpoint{-13.359530\du}{6.311060\du}}
\pgfpathlineto{\pgfpoint{-13.349379\du}{6.310462\du}}
\pgfpathlineto{\pgfpoint{-13.339825\du}{6.309268\du}}
\pgfpathlineto{\pgfpoint{-13.330570\du}{6.308373\du}}
\pgfpathlineto{\pgfpoint{-13.321315\du}{6.307178\du}}
\pgfpathlineto{\pgfpoint{-13.312060\du}{6.305686\du}}
\pgfpathlineto{\pgfpoint{-13.303402\du}{6.304193\du}}
\pgfpathlineto{\pgfpoint{-13.295043\du}{6.302700\du}}
\pgfpathlineto{\pgfpoint{-13.286385\du}{6.300610\du}}
\pgfpathlineto{\pgfpoint{-13.278324\du}{6.298520\du}}
\pgfpathlineto{\pgfpoint{-13.270860\du}{6.296430\du}}
\pgfpathlineto{\pgfpoint{-13.263396\du}{6.294639\du}}
\pgfpathlineto{\pgfpoint{-13.256529\du}{6.292251\du}}
\pgfpathlineto{\pgfpoint{-13.249663\du}{6.289862\du}}
\pgfpathlineto{\pgfpoint{-13.243692\du}{6.287175\du}}
\pgfpathlineto{\pgfpoint{-13.237422\du}{6.284488\du}}
\pgfpathlineto{\pgfpoint{-13.232048\du}{6.281801\du}}
\pgfpathlineto{\pgfpoint{-13.226973\du}{6.279114\du}}
\pgfpathlineto{\pgfpoint{-13.222196\du}{6.275830\du}}
\pgfpathlineto{\pgfpoint{-13.218613\du}{6.272845\du}}
\pgfpathlineto{\pgfpoint{-13.214434\du}{6.269859\du}}
\pgfpathlineto{\pgfpoint{-13.211448\du}{6.266874\du}}
\pgfpathlineto{\pgfpoint{-13.208761\du}{6.263590\du}}
\pgfpathlineto{\pgfpoint{-13.206373\du}{6.260306\du}}
\pgfpathlineto{\pgfpoint{-13.204880\du}{6.257022\du}}
\pgfpathlineto{\pgfpoint{-13.203686\du}{6.253738\du}}
\pgfpathlineto{\pgfpoint{-13.202790\du}{6.250752\du}}
\pgfpathlineto{\pgfpoint{-13.202790\du}{6.247169\du}}
\pgfusepath{fill}
\pgfsetbuttcap
\pgfsetmiterjoin
\pgfsetdash{}{0pt}
\definecolor{dialinecolor}{rgb}{0.666667, 0.901961, 1.000000}
\pgfsetstrokecolor{dialinecolor}
\pgfpathmoveto{\pgfpoint{-13.214434\du}{6.241497\du}}
\pgfpathlineto{\pgfpoint{-13.214434\du}{6.237914\du}}
\pgfpathlineto{\pgfpoint{-13.215329\du}{6.235526\du}}
\pgfpathlineto{\pgfpoint{-13.216524\du}{6.232242\du}}
\pgfpathlineto{\pgfpoint{-13.218613\du}{6.229256\du}}
\pgfpathlineto{\pgfpoint{-13.220703\du}{6.226569\du}}
\pgfpathlineto{\pgfpoint{-13.222793\du}{6.222987\du}}
\pgfpathlineto{\pgfpoint{-13.226077\du}{6.220598\du}}
\pgfpathlineto{\pgfpoint{-13.229361\du}{6.217911\du}}
\pgfpathlineto{\pgfpoint{-13.233541\du}{6.215224\du}}
\pgfpathlineto{\pgfpoint{-13.238318\du}{6.212836\du}}
\pgfpathlineto{\pgfpoint{-13.242796\du}{6.209850\du}}
\pgfpathlineto{\pgfpoint{-13.248469\du}{6.207761\du}}
\pgfpathlineto{\pgfpoint{-13.254440\du}{6.205074\du}}
\pgfpathlineto{\pgfpoint{-13.260112\du}{6.202387\du}}
\pgfpathlineto{\pgfpoint{-13.266680\du}{6.200595\du}}
\pgfpathlineto{\pgfpoint{-13.273547\du}{6.198207\du}}
\pgfpathlineto{\pgfpoint{-13.280414\du}{6.196714\du}}
\pgfpathlineto{\pgfpoint{-13.288176\du}{6.194624\du}}
\pgfpathlineto{\pgfpoint{-13.295640\du}{6.192833\du}}
\pgfpathlineto{\pgfpoint{-13.303999\du}{6.191042\du}}
\pgfpathlineto{\pgfpoint{-13.312657\du}{6.189847\du}}
\pgfpathlineto{\pgfpoint{-13.321017\du}{6.187459\du}}
\pgfpathlineto{\pgfpoint{-13.329675\du}{6.187160\du}}
\pgfpathlineto{\pgfpoint{-13.338631\du}{6.185668\du}}
\pgfpathlineto{\pgfpoint{-13.347289\du}{6.184772\du}}
\pgfpathlineto{\pgfpoint{-13.356843\du}{6.183876\du}}
\pgfpathlineto{\pgfpoint{-13.366098\du}{6.183279\du}}
\pgfpathlineto{\pgfpoint{-13.375652\du}{6.182981\du}}
\pgfpathlineto{\pgfpoint{-13.385205\du}{6.182384\du}}
\pgfpathlineto{\pgfpoint{-13.395356\du}{6.182085\du}}
\pgfpathlineto{\pgfpoint{-13.405208\du}{6.182085\du}}
\pgfpathlineto{\pgfpoint{-13.405208\du}{6.182085\du}}
\pgfpathlineto{\pgfpoint{-13.414165\du}{6.182085\du}}
\pgfpathlineto{\pgfpoint{-13.424017\du}{6.182384\du}}
\pgfpathlineto{\pgfpoint{-13.433571\du}{6.182981\du}}
\pgfpathlineto{\pgfpoint{-13.442826\du}{6.183279\du}}
\pgfpathlineto{\pgfpoint{-13.452380\du}{6.183876\du}}
\pgfpathlineto{\pgfpoint{-13.461933\du}{6.184772\du}}
\pgfpathlineto{\pgfpoint{-13.470890\du}{6.185668\du}}
\pgfpathlineto{\pgfpoint{-13.479846\du}{6.187160\du}}
\pgfpathlineto{\pgfpoint{-13.488504\du}{6.187459\du}}
\pgfpathlineto{\pgfpoint{-13.496864\du}{6.189847\du}}
\pgfpathlineto{\pgfpoint{-13.505522\du}{6.191042\du}}
\pgfpathlineto{\pgfpoint{-13.513583\du}{6.192833\du}}
\pgfpathlineto{\pgfpoint{-13.521345\du}{6.194624\du}}
\pgfpathlineto{\pgfpoint{-13.528809\du}{6.196714\du}}
\pgfpathlineto{\pgfpoint{-13.535377\du}{6.198207\du}}
\pgfpathlineto{\pgfpoint{-13.542841\du}{6.200595\du}}
\pgfpathlineto{\pgfpoint{-13.549110\du}{6.202387\du}}
\pgfpathlineto{\pgfpoint{-13.555082\du}{6.205074\du}}
\pgfpathlineto{\pgfpoint{-13.561053\du}{6.207761\du}}
\pgfpathlineto{\pgfpoint{-13.566128\du}{6.209850\du}}
\pgfpathlineto{\pgfpoint{-13.571203\du}{6.212836\du}}
\pgfpathlineto{\pgfpoint{-13.575980\du}{6.215224\du}}
\pgfpathlineto{\pgfpoint{-13.580160\du}{6.217911\du}}
\pgfpathlineto{\pgfpoint{-13.583444\du}{6.220598\du}}
\pgfpathlineto{\pgfpoint{-13.586728\du}{6.222987\du}}
\pgfpathlineto{\pgfpoint{-13.589116\du}{6.226569\du}}
\pgfpathlineto{\pgfpoint{-13.590908\du}{6.229256\du}}
\pgfpathlineto{\pgfpoint{-13.593296\du}{6.232242\du}}
\pgfpathlineto{\pgfpoint{-13.594192\du}{6.235526\du}}
\pgfpathlineto{\pgfpoint{-13.595087\du}{6.237914\du}}
\pgfpathlineto{\pgfpoint{-13.595087\du}{6.241497\du}}
\pgfpathlineto{\pgfpoint{-13.595087\du}{6.241497\du}}
\pgfpathlineto{\pgfpoint{-13.595087\du}{6.244482\du}}
\pgfpathlineto{\pgfpoint{-13.594192\du}{6.247169\du}}
\pgfpathlineto{\pgfpoint{-13.593296\du}{6.250155\du}}
\pgfpathlineto{\pgfpoint{-13.590908\du}{6.253738\du}}
\pgfpathlineto{\pgfpoint{-13.589116\du}{6.256424\du}}
\pgfpathlineto{\pgfpoint{-13.586728\du}{6.259410\du}}
\pgfpathlineto{\pgfpoint{-13.583444\du}{6.262097\du}}
\pgfpathlineto{\pgfpoint{-13.580160\du}{6.264784\du}}
\pgfpathlineto{\pgfpoint{-13.575980\du}{6.267471\du}}
\pgfpathlineto{\pgfpoint{-13.571203\du}{6.270158\du}}
\pgfpathlineto{\pgfpoint{-13.566128\du}{6.272845\du}}
\pgfpathlineto{\pgfpoint{-13.561053\du}{6.275233\du}}
\pgfpathlineto{\pgfpoint{-13.555082\du}{6.277920\du}}
\pgfpathlineto{\pgfpoint{-13.549110\du}{6.280010\du}}
\pgfpathlineto{\pgfpoint{-13.542841\du}{6.282399\du}}
\pgfpathlineto{\pgfpoint{-13.535377\du}{6.284787\du}}
\pgfpathlineto{\pgfpoint{-13.528809\du}{6.286578\du}}
\pgfpathlineto{\pgfpoint{-13.521345\du}{6.288071\du}}
\pgfpathlineto{\pgfpoint{-13.513583\du}{6.290161\du}}
\pgfpathlineto{\pgfpoint{-13.505522\du}{6.291952\du}}
\pgfpathlineto{\pgfpoint{-13.496864\du}{6.293146\du}}
\pgfpathlineto{\pgfpoint{-13.488504\du}{6.294938\du}}
\pgfpathlineto{\pgfpoint{-13.479846\du}{6.295833\du}}
\pgfpathlineto{\pgfpoint{-13.470890\du}{6.297326\du}}
\pgfpathlineto{\pgfpoint{-13.461933\du}{6.298222\du}}
\pgfpathlineto{\pgfpoint{-13.452380\du}{6.298520\du}}
\pgfpathlineto{\pgfpoint{-13.442826\du}{6.299715\du}}
\pgfpathlineto{\pgfpoint{-13.433571\du}{6.300013\du}}
\pgfpathlineto{\pgfpoint{-13.424017\du}{6.300312\du}}
\pgfpathlineto{\pgfpoint{-13.414165\du}{6.300909\du}}
\pgfpathlineto{\pgfpoint{-13.405208\du}{6.300909\du}}
\pgfpathlineto{\pgfpoint{-13.405208\du}{6.300909\du}}
\pgfpathlineto{\pgfpoint{-13.395356\du}{6.300909\du}}
\pgfpathlineto{\pgfpoint{-13.385205\du}{6.300312\du}}
\pgfpathlineto{\pgfpoint{-13.375652\du}{6.300013\du}}
\pgfpathlineto{\pgfpoint{-13.366098\du}{6.299715\du}}
\pgfpathlineto{\pgfpoint{-13.356843\du}{6.298520\du}}
\pgfpathlineto{\pgfpoint{-13.347289\du}{6.298222\du}}
\pgfpathlineto{\pgfpoint{-13.338631\du}{6.297326\du}}
\pgfpathlineto{\pgfpoint{-13.329675\du}{6.295833\du}}
\pgfpathlineto{\pgfpoint{-13.321017\du}{6.294938\du}}
\pgfpathlineto{\pgfpoint{-13.312657\du}{6.293146\du}}
\pgfpathlineto{\pgfpoint{-13.303999\du}{6.291952\du}}
\pgfpathlineto{\pgfpoint{-13.295640\du}{6.290161\du}}
\pgfpathlineto{\pgfpoint{-13.288176\du}{6.288071\du}}
\pgfpathlineto{\pgfpoint{-13.280414\du}{6.286578\du}}
\pgfpathlineto{\pgfpoint{-13.273547\du}{6.284787\du}}
\pgfpathlineto{\pgfpoint{-13.266680\du}{6.282399\du}}
\pgfpathlineto{\pgfpoint{-13.260112\du}{6.280010\du}}
\pgfpathlineto{\pgfpoint{-13.254440\du}{6.277920\du}}
\pgfpathlineto{\pgfpoint{-13.248469\du}{6.275233\du}}
\pgfpathlineto{\pgfpoint{-13.242796\du}{6.272845\du}}
\pgfpathlineto{\pgfpoint{-13.238318\du}{6.270158\du}}
\pgfpathlineto{\pgfpoint{-13.233541\du}{6.267471\du}}
\pgfpathlineto{\pgfpoint{-13.229361\du}{6.264784\du}}
\pgfpathlineto{\pgfpoint{-13.226077\du}{6.262097\du}}
\pgfpathlineto{\pgfpoint{-13.222793\du}{6.259410\du}}
\pgfpathlineto{\pgfpoint{-13.220703\du}{6.256424\du}}
\pgfpathlineto{\pgfpoint{-13.218613\du}{6.253738\du}}
\pgfpathlineto{\pgfpoint{-13.216524\du}{6.250155\du}}
\pgfpathlineto{\pgfpoint{-13.215329\du}{6.247169\du}}
\pgfpathlineto{\pgfpoint{-13.214434\du}{6.244482\du}}
\pgfpathlineto{\pgfpoint{-13.214434\du}{6.241497\du}}
\pgfusepath{stroke}
\pgfsetbuttcap
\pgfsetmiterjoin
\pgfsetdash{}{0pt}
\definecolor{dialinecolor}{rgb}{0.000000, 0.470588, 0.666667}
\pgfsetfillcolor{dialinecolor}
\pgfpathmoveto{\pgfpoint{-13.595087\du}{6.158201\du}}
\pgfpathlineto{\pgfpoint{-13.595087\du}{6.252842\du}}
\pgfpathlineto{\pgfpoint{-13.202193\du}{6.252842\du}}
\pgfpathlineto{\pgfpoint{-13.202193\du}{6.158201\du}}
\pgfpathlineto{\pgfpoint{-13.595087\du}{6.158201\du}}
\pgfusepath{fill}
\pgfsetbuttcap
\pgfsetmiterjoin
\pgfsetdash{}{0pt}
\definecolor{dialinecolor}{rgb}{0.000000, 0.705882, 1.000000}
\pgfsetfillcolor{dialinecolor}
\pgfpathmoveto{\pgfpoint{-13.202790\du}{6.151633\du}}
\pgfpathlineto{\pgfpoint{-13.202790\du}{6.148349\du}}
\pgfpathlineto{\pgfpoint{-13.203686\du}{6.144766\du}}
\pgfpathlineto{\pgfpoint{-13.204880\du}{6.141781\du}}
\pgfpathlineto{\pgfpoint{-13.206373\du}{6.138496\du}}
\pgfpathlineto{\pgfpoint{-13.208761\du}{6.135212\du}}
\pgfpathlineto{\pgfpoint{-13.211448\du}{6.131928\du}}
\pgfpathlineto{\pgfpoint{-13.214434\du}{6.128943\du}}
\pgfpathlineto{\pgfpoint{-13.218613\du}{6.125659\du}}
\pgfpathlineto{\pgfpoint{-13.222196\du}{6.122972\du}}
\pgfpathlineto{\pgfpoint{-13.226973\du}{6.119986\du}}
\pgfpathlineto{\pgfpoint{-13.232048\du}{6.117001\du}}
\pgfpathlineto{\pgfpoint{-13.237422\du}{6.114314\du}}
\pgfpathlineto{\pgfpoint{-13.243692\du}{6.111030\du}}
\pgfpathlineto{\pgfpoint{-13.249663\du}{6.109238\du}}
\pgfpathlineto{\pgfpoint{-13.256529\du}{6.106551\du}}
\pgfpathlineto{\pgfpoint{-13.263396\du}{6.104163\du}}
\pgfpathlineto{\pgfpoint{-13.270860\du}{6.102073\du}}
\pgfpathlineto{\pgfpoint{-13.278324\du}{6.099983\du}}
\pgfpathlineto{\pgfpoint{-13.286385\du}{6.097893\du}}
\pgfpathlineto{\pgfpoint{-13.295043\du}{6.095804\du}}
\pgfpathlineto{\pgfpoint{-13.303402\du}{6.094311\du}}
\pgfpathlineto{\pgfpoint{-13.312060\du}{6.092818\du}}
\pgfpathlineto{\pgfpoint{-13.321315\du}{6.091624\du}}
\pgfpathlineto{\pgfpoint{-13.330570\du}{6.090131\du}}
\pgfpathlineto{\pgfpoint{-13.339825\du}{6.089235\du}}
\pgfpathlineto{\pgfpoint{-13.349379\du}{6.088041\du}}
\pgfpathlineto{\pgfpoint{-13.359530\du}{6.087444\du}}
\pgfpathlineto{\pgfpoint{-13.369084\du}{6.086847\du}}
\pgfpathlineto{\pgfpoint{-13.378637\du}{6.086548\du}}
\pgfpathlineto{\pgfpoint{-13.388788\du}{6.086250\du}}
\pgfpathlineto{\pgfpoint{-13.398640\du}{6.086250\du}}
\pgfpathlineto{\pgfpoint{-13.398640\du}{6.086250\du}}
\pgfpathlineto{\pgfpoint{-13.408492\du}{6.086250\du}}
\pgfpathlineto{\pgfpoint{-13.418643\du}{6.086548\du}}
\pgfpathlineto{\pgfpoint{-13.428495\du}{6.086847\du}}
\pgfpathlineto{\pgfpoint{-13.438348\du}{6.087444\du}}
\pgfpathlineto{\pgfpoint{-13.447901\du}{6.088041\du}}
\pgfpathlineto{\pgfpoint{-13.457455\du}{6.089235\du}}
\pgfpathlineto{\pgfpoint{-13.467009\du}{6.090131\du}}
\pgfpathlineto{\pgfpoint{-13.475965\du}{6.091624\du}}
\pgfpathlineto{\pgfpoint{-13.485220\du}{6.092818\du}}
\pgfpathlineto{\pgfpoint{-13.493878\du}{6.094311\du}}
\pgfpathlineto{\pgfpoint{-13.502536\du}{6.095804\du}}
\pgfpathlineto{\pgfpoint{-13.510896\du}{6.097893\du}}
\pgfpathlineto{\pgfpoint{-13.518957\du}{6.099983\du}}
\pgfpathlineto{\pgfpoint{-13.526421\du}{6.102073\du}}
\pgfpathlineto{\pgfpoint{-13.533884\du}{6.104163\du}}
\pgfpathlineto{\pgfpoint{-13.541050\du}{6.106551\du}}
\pgfpathlineto{\pgfpoint{-13.547618\du}{6.109238\du}}
\pgfpathlineto{\pgfpoint{-13.553887\du}{6.111030\du}}
\pgfpathlineto{\pgfpoint{-13.559858\du}{6.114314\du}}
\pgfpathlineto{\pgfpoint{-13.565232\du}{6.117001\du}}
\pgfpathlineto{\pgfpoint{-13.570308\du}{6.119986\du}}
\pgfpathlineto{\pgfpoint{-13.575084\du}{6.122972\du}}
\pgfpathlineto{\pgfpoint{-13.578966\du}{6.125659\du}}
\pgfpathlineto{\pgfpoint{-13.582847\du}{6.128943\du}}
\pgfpathlineto{\pgfpoint{-13.585832\du}{6.131928\du}}
\pgfpathlineto{\pgfpoint{-13.588519\du}{6.135212\du}}
\pgfpathlineto{\pgfpoint{-13.590908\du}{6.138496\du}}
\pgfpathlineto{\pgfpoint{-13.592699\du}{6.141781\du}}
\pgfpathlineto{\pgfpoint{-13.594192\du}{6.144766\du}}
\pgfpathlineto{\pgfpoint{-13.595087\du}{6.148349\du}}
\pgfpathlineto{\pgfpoint{-13.595087\du}{6.151633\du}}
\pgfpathlineto{\pgfpoint{-13.595087\du}{6.151633\du}}
\pgfpathlineto{\pgfpoint{-13.595087\du}{6.154917\du}}
\pgfpathlineto{\pgfpoint{-13.594192\du}{6.158499\du}}
\pgfpathlineto{\pgfpoint{-13.592699\du}{6.161485\du}}
\pgfpathlineto{\pgfpoint{-13.590908\du}{6.164769\du}}
\pgfpathlineto{\pgfpoint{-13.588519\du}{6.168352\du}}
\pgfpathlineto{\pgfpoint{-13.585832\du}{6.171337\du}}
\pgfpathlineto{\pgfpoint{-13.582847\du}{6.174323\du}}
\pgfpathlineto{\pgfpoint{-13.578966\du}{6.177607\du}}
\pgfpathlineto{\pgfpoint{-13.575084\du}{6.179995\du}}
\pgfpathlineto{\pgfpoint{-13.570308\du}{6.183279\du}}
\pgfpathlineto{\pgfpoint{-13.565232\du}{6.186563\du}}
\pgfpathlineto{\pgfpoint{-13.559858\du}{6.189250\du}}
\pgfpathlineto{\pgfpoint{-13.553887\du}{6.191937\du}}
\pgfpathlineto{\pgfpoint{-13.547618\du}{6.194027\du}}
\pgfpathlineto{\pgfpoint{-13.541050\du}{6.197013\du}}
\pgfpathlineto{\pgfpoint{-13.533884\du}{6.199401\du}}
\pgfpathlineto{\pgfpoint{-13.526421\du}{6.201192\du}}
\pgfpathlineto{\pgfpoint{-13.518957\du}{6.203282\du}}
\pgfpathlineto{\pgfpoint{-13.510896\du}{6.205372\du}}
\pgfpathlineto{\pgfpoint{-13.502536\du}{6.207462\du}}
\pgfpathlineto{\pgfpoint{-13.493878\du}{6.208955\du}}
\pgfpathlineto{\pgfpoint{-13.485220\du}{6.210447\du}}
\pgfpathlineto{\pgfpoint{-13.475965\du}{6.211642\du}}
\pgfpathlineto{\pgfpoint{-13.467009\du}{6.213134\du}}
\pgfpathlineto{\pgfpoint{-13.457455\du}{6.214030\du}}
\pgfpathlineto{\pgfpoint{-13.447901\du}{6.215224\du}}
\pgfpathlineto{\pgfpoint{-13.438348\du}{6.215523\du}}
\pgfpathlineto{\pgfpoint{-13.428495\du}{6.216419\du}}
\pgfpathlineto{\pgfpoint{-13.418643\du}{6.216419\du}}
\pgfpathlineto{\pgfpoint{-13.408492\du}{6.217314\du}}
\pgfpathlineto{\pgfpoint{-13.398640\du}{6.217314\du}}
\pgfpathlineto{\pgfpoint{-13.398640\du}{6.217314\du}}
\pgfpathlineto{\pgfpoint{-13.388788\du}{6.217314\du}}
\pgfpathlineto{\pgfpoint{-13.378637\du}{6.216419\du}}
\pgfpathlineto{\pgfpoint{-13.369084\du}{6.216419\du}}
\pgfpathlineto{\pgfpoint{-13.359530\du}{6.215523\du}}
\pgfpathlineto{\pgfpoint{-13.349379\du}{6.215224\du}}
\pgfpathlineto{\pgfpoint{-13.339825\du}{6.214030\du}}
\pgfpathlineto{\pgfpoint{-13.330570\du}{6.213134\du}}
\pgfpathlineto{\pgfpoint{-13.321315\du}{6.211642\du}}
\pgfpathlineto{\pgfpoint{-13.312060\du}{6.210447\du}}
\pgfpathlineto{\pgfpoint{-13.303402\du}{6.208955\du}}
\pgfpathlineto{\pgfpoint{-13.295043\du}{6.207462\du}}
\pgfpathlineto{\pgfpoint{-13.286385\du}{6.205372\du}}
\pgfpathlineto{\pgfpoint{-13.278324\du}{6.203282\du}}
\pgfpathlineto{\pgfpoint{-13.270860\du}{6.201192\du}}
\pgfpathlineto{\pgfpoint{-13.263396\du}{6.199401\du}}
\pgfpathlineto{\pgfpoint{-13.256529\du}{6.197013\du}}
\pgfpathlineto{\pgfpoint{-13.249663\du}{6.194027\du}}
\pgfpathlineto{\pgfpoint{-13.243692\du}{6.191937\du}}
\pgfpathlineto{\pgfpoint{-13.237422\du}{6.189250\du}}
\pgfpathlineto{\pgfpoint{-13.232048\du}{6.186563\du}}
\pgfpathlineto{\pgfpoint{-13.226973\du}{6.183279\du}}
\pgfpathlineto{\pgfpoint{-13.222196\du}{6.179995\du}}
\pgfpathlineto{\pgfpoint{-13.218613\du}{6.177607\du}}
\pgfpathlineto{\pgfpoint{-13.214434\du}{6.174323\du}}
\pgfpathlineto{\pgfpoint{-13.211448\du}{6.171337\du}}
\pgfpathlineto{\pgfpoint{-13.208761\du}{6.168352\du}}
\pgfpathlineto{\pgfpoint{-13.206373\du}{6.164769\du}}
\pgfpathlineto{\pgfpoint{-13.204880\du}{6.161485\du}}
\pgfpathlineto{\pgfpoint{-13.203686\du}{6.158499\du}}
\pgfpathlineto{\pgfpoint{-13.202790\du}{6.154917\du}}
\pgfpathlineto{\pgfpoint{-13.202790\du}{6.151633\du}}
\pgfusepath{fill}
\pgfsetbuttcap
\pgfsetmiterjoin
\pgfsetdash{}{0pt}
\definecolor{dialinecolor}{rgb}{0.666667, 0.901961, 1.000000}
\pgfsetstrokecolor{dialinecolor}
\pgfpathmoveto{\pgfpoint{-13.214434\du}{6.145960\du}}
\pgfpathlineto{\pgfpoint{-13.214434\du}{6.142676\du}}
\pgfpathlineto{\pgfpoint{-13.215329\du}{6.140288\du}}
\pgfpathlineto{\pgfpoint{-13.216524\du}{6.137004\du}}
\pgfpathlineto{\pgfpoint{-13.218613\du}{6.133421\du}}
\pgfpathlineto{\pgfpoint{-13.220703\du}{6.131033\du}}
\pgfpathlineto{\pgfpoint{-13.222793\du}{6.127749\du}}
\pgfpathlineto{\pgfpoint{-13.226077\du}{6.125360\du}}
\pgfpathlineto{\pgfpoint{-13.229361\du}{6.122673\du}}
\pgfpathlineto{\pgfpoint{-13.233541\du}{6.119986\du}}
\pgfpathlineto{\pgfpoint{-13.238318\du}{6.117299\du}}
\pgfpathlineto{\pgfpoint{-13.242796\du}{6.114612\du}}
\pgfpathlineto{\pgfpoint{-13.248469\du}{6.111925\du}}
\pgfpathlineto{\pgfpoint{-13.254440\du}{6.109537\du}}
\pgfpathlineto{\pgfpoint{-13.260112\du}{6.107447\du}}
\pgfpathlineto{\pgfpoint{-13.266680\du}{6.105059\du}}
\pgfpathlineto{\pgfpoint{-13.273547\du}{6.102670\du}}
\pgfpathlineto{\pgfpoint{-13.280414\du}{6.100879\du}}
\pgfpathlineto{\pgfpoint{-13.288176\du}{6.099386\du}}
\pgfpathlineto{\pgfpoint{-13.295640\du}{6.096998\du}}
\pgfpathlineto{\pgfpoint{-13.303999\du}{6.095505\du}}
\pgfpathlineto{\pgfpoint{-13.312657\du}{6.094311\du}}
\pgfpathlineto{\pgfpoint{-13.321017\du}{6.092519\du}}
\pgfpathlineto{\pgfpoint{-13.329675\du}{6.091624\du}}
\pgfpathlineto{\pgfpoint{-13.338631\du}{6.090131\du}}
\pgfpathlineto{\pgfpoint{-13.347289\du}{6.089235\du}}
\pgfpathlineto{\pgfpoint{-13.356843\du}{6.088638\du}}
\pgfpathlineto{\pgfpoint{-13.366098\du}{6.087743\du}}
\pgfpathlineto{\pgfpoint{-13.375652\du}{6.087444\du}}
\pgfpathlineto{\pgfpoint{-13.385205\du}{6.087145\du}}
\pgfpathlineto{\pgfpoint{-13.395356\du}{6.086548\du}}
\pgfpathlineto{\pgfpoint{-13.405208\du}{6.086548\du}}
\pgfpathlineto{\pgfpoint{-13.405208\du}{6.086548\du}}
\pgfpathlineto{\pgfpoint{-13.414165\du}{6.086548\du}}
\pgfpathlineto{\pgfpoint{-13.424017\du}{6.087145\du}}
\pgfpathlineto{\pgfpoint{-13.433571\du}{6.087444\du}}
\pgfpathlineto{\pgfpoint{-13.442826\du}{6.087743\du}}
\pgfpathlineto{\pgfpoint{-13.452380\du}{6.088638\du}}
\pgfpathlineto{\pgfpoint{-13.461933\du}{6.089235\du}}
\pgfpathlineto{\pgfpoint{-13.470890\du}{6.090131\du}}
\pgfpathlineto{\pgfpoint{-13.479846\du}{6.091624\du}}
\pgfpathlineto{\pgfpoint{-13.488504\du}{6.092519\du}}
\pgfpathlineto{\pgfpoint{-13.496864\du}{6.094311\du}}
\pgfpathlineto{\pgfpoint{-13.505522\du}{6.095505\du}}
\pgfpathlineto{\pgfpoint{-13.513583\du}{6.096998\du}}
\pgfpathlineto{\pgfpoint{-13.521345\du}{6.099386\du}}
\pgfpathlineto{\pgfpoint{-13.528809\du}{6.100879\du}}
\pgfpathlineto{\pgfpoint{-13.535377\du}{6.102670\du}}
\pgfpathlineto{\pgfpoint{-13.542841\du}{6.105059\du}}
\pgfpathlineto{\pgfpoint{-13.549110\du}{6.107447\du}}
\pgfpathlineto{\pgfpoint{-13.555082\du}{6.109537\du}}
\pgfpathlineto{\pgfpoint{-13.561053\du}{6.111925\du}}
\pgfpathlineto{\pgfpoint{-13.566128\du}{6.114612\du}}
\pgfpathlineto{\pgfpoint{-13.571203\du}{6.117299\du}}
\pgfpathlineto{\pgfpoint{-13.575980\du}{6.119986\du}}
\pgfpathlineto{\pgfpoint{-13.580160\du}{6.122673\du}}
\pgfpathlineto{\pgfpoint{-13.583444\du}{6.125360\du}}
\pgfpathlineto{\pgfpoint{-13.586728\du}{6.127749\du}}
\pgfpathlineto{\pgfpoint{-13.589116\du}{6.131033\du}}
\pgfpathlineto{\pgfpoint{-13.590908\du}{6.133421\du}}
\pgfpathlineto{\pgfpoint{-13.593296\du}{6.137004\du}}
\pgfpathlineto{\pgfpoint{-13.594192\du}{6.140288\du}}
\pgfpathlineto{\pgfpoint{-13.595087\du}{6.142676\du}}
\pgfpathlineto{\pgfpoint{-13.595087\du}{6.145960\du}}
\pgfpathlineto{\pgfpoint{-13.595087\du}{6.145960\du}}
\pgfpathlineto{\pgfpoint{-13.595087\du}{6.149244\du}}
\pgfpathlineto{\pgfpoint{-13.594192\du}{6.151931\du}}
\pgfpathlineto{\pgfpoint{-13.593296\du}{6.154917\du}}
\pgfpathlineto{\pgfpoint{-13.590908\du}{6.158201\du}}
\pgfpathlineto{\pgfpoint{-13.589116\du}{6.160888\du}}
\pgfpathlineto{\pgfpoint{-13.586728\du}{6.164172\du}}
\pgfpathlineto{\pgfpoint{-13.583444\du}{6.166859\du}}
\pgfpathlineto{\pgfpoint{-13.580160\du}{6.169546\du}}
\pgfpathlineto{\pgfpoint{-13.575980\du}{6.172233\du}}
\pgfpathlineto{\pgfpoint{-13.571203\du}{6.174920\du}}
\pgfpathlineto{\pgfpoint{-13.566128\du}{6.177607\du}}
\pgfpathlineto{\pgfpoint{-13.561053\du}{6.179697\du}}
\pgfpathlineto{\pgfpoint{-13.555082\du}{6.182384\du}}
\pgfpathlineto{\pgfpoint{-13.549110\du}{6.184772\du}}
\pgfpathlineto{\pgfpoint{-13.542841\du}{6.187160\du}}
\pgfpathlineto{\pgfpoint{-13.535377\du}{6.189250\du}}
\pgfpathlineto{\pgfpoint{-13.528809\du}{6.191042\du}}
\pgfpathlineto{\pgfpoint{-13.521345\du}{6.192833\du}}
\pgfpathlineto{\pgfpoint{-13.513583\du}{6.194624\du}}
\pgfpathlineto{\pgfpoint{-13.505522\du}{6.196714\du}}
\pgfpathlineto{\pgfpoint{-13.496864\du}{6.197908\du}}
\pgfpathlineto{\pgfpoint{-13.488504\du}{6.199700\du}}
\pgfpathlineto{\pgfpoint{-13.479846\du}{6.200595\du}}
\pgfpathlineto{\pgfpoint{-13.470890\du}{6.202088\du}}
\pgfpathlineto{\pgfpoint{-13.461933\du}{6.202387\du}}
\pgfpathlineto{\pgfpoint{-13.452380\du}{6.203581\du}}
\pgfpathlineto{\pgfpoint{-13.442826\du}{6.203879\du}}
\pgfpathlineto{\pgfpoint{-13.433571\du}{6.204775\du}}
\pgfpathlineto{\pgfpoint{-13.424017\du}{6.205074\du}}
\pgfpathlineto{\pgfpoint{-13.414165\du}{6.205372\du}}
\pgfpathlineto{\pgfpoint{-13.405208\du}{6.205372\du}}
\pgfpathlineto{\pgfpoint{-13.405208\du}{6.205372\du}}
\pgfpathlineto{\pgfpoint{-13.395356\du}{6.205372\du}}
\pgfpathlineto{\pgfpoint{-13.385205\du}{6.205074\du}}
\pgfpathlineto{\pgfpoint{-13.375652\du}{6.204775\du}}
\pgfpathlineto{\pgfpoint{-13.366098\du}{6.203879\du}}
\pgfpathlineto{\pgfpoint{-13.356843\du}{6.203581\du}}
\pgfpathlineto{\pgfpoint{-13.347289\du}{6.202387\du}}
\pgfpathlineto{\pgfpoint{-13.338631\du}{6.202088\du}}
\pgfpathlineto{\pgfpoint{-13.329675\du}{6.200595\du}}
\pgfpathlineto{\pgfpoint{-13.321017\du}{6.199700\du}}
\pgfpathlineto{\pgfpoint{-13.312657\du}{6.197908\du}}
\pgfpathlineto{\pgfpoint{-13.303999\du}{6.196714\du}}
\pgfpathlineto{\pgfpoint{-13.295640\du}{6.194624\du}}
\pgfpathlineto{\pgfpoint{-13.288176\du}{6.192833\du}}
\pgfpathlineto{\pgfpoint{-13.280414\du}{6.191042\du}}
\pgfpathlineto{\pgfpoint{-13.273547\du}{6.189250\du}}
\pgfpathlineto{\pgfpoint{-13.266680\du}{6.187160\du}}
\pgfpathlineto{\pgfpoint{-13.260112\du}{6.184772\du}}
\pgfpathlineto{\pgfpoint{-13.254440\du}{6.182384\du}}
\pgfpathlineto{\pgfpoint{-13.248469\du}{6.179697\du}}
\pgfpathlineto{\pgfpoint{-13.242796\du}{6.177607\du}}
\pgfpathlineto{\pgfpoint{-13.238318\du}{6.174920\du}}
\pgfpathlineto{\pgfpoint{-13.233541\du}{6.172233\du}}
\pgfpathlineto{\pgfpoint{-13.229361\du}{6.169546\du}}
\pgfpathlineto{\pgfpoint{-13.226077\du}{6.166859\du}}
\pgfpathlineto{\pgfpoint{-13.222793\du}{6.164172\du}}
\pgfpathlineto{\pgfpoint{-13.220703\du}{6.160888\du}}
\pgfpathlineto{\pgfpoint{-13.218613\du}{6.158201\du}}
\pgfpathlineto{\pgfpoint{-13.216524\du}{6.154917\du}}
\pgfpathlineto{\pgfpoint{-13.215329\du}{6.151931\du}}
\pgfpathlineto{\pgfpoint{-13.214434\du}{6.149244\du}}
\pgfpathlineto{\pgfpoint{-13.214434\du}{6.145960\du}}
\pgfusepath{stroke}
\pgfsetbuttcap
\pgfsetmiterjoin
\pgfsetdash{}{0pt}
\definecolor{dialinecolor}{rgb}{0.000000, 0.000000, 0.000000}
\pgfsetfillcolor{dialinecolor}
\pgfpathmoveto{\pgfpoint{-13.392968\du}{6.134317\du}}
\pgfpathlineto{\pgfpoint{-13.369084\du}{6.145960\du}}
\pgfpathlineto{\pgfpoint{-13.297730\du}{6.122375\du}}
\pgfpathlineto{\pgfpoint{-13.261903\du}{6.134317\du}}
\pgfpathlineto{\pgfpoint{-13.285788\du}{6.110433\du}}
\pgfpathlineto{\pgfpoint{-13.369084\du}{6.110433\du}}
\pgfpathlineto{\pgfpoint{-13.333257\du}{6.110433\du}}
\pgfpathlineto{\pgfpoint{-13.392968\du}{6.134317\du}}
\pgfusepath{fill}
\pgfsetbuttcap
\pgfsetmiterjoin
\pgfsetdash{}{0pt}
\definecolor{dialinecolor}{rgb}{0.000000, 0.000000, 0.000000}
\pgfsetfillcolor{dialinecolor}
\pgfpathmoveto{\pgfpoint{-13.405208\du}{6.169844\du}}
\pgfpathlineto{\pgfpoint{-13.440139\du}{6.158201\du}}
\pgfpathlineto{\pgfpoint{-13.500148\du}{6.182085\du}}
\pgfpathlineto{\pgfpoint{-13.535377\du}{6.169844\du}}
\pgfpathlineto{\pgfpoint{-13.524032\du}{6.205372\du}}
\pgfpathlineto{\pgfpoint{-13.440139\du}{6.205372\du}}
\pgfpathlineto{\pgfpoint{-13.476264\du}{6.194027\du}}
\pgfpathlineto{\pgfpoint{-13.405208\du}{6.169844\du}}
\pgfusepath{fill}
\pgfsetbuttcap
\pgfsetmiterjoin
\pgfsetdash{}{0pt}
\definecolor{dialinecolor}{rgb}{0.000000, 0.000000, 0.000000}
\pgfsetfillcolor{dialinecolor}
\pgfpathmoveto{\pgfpoint{-13.524032\du}{6.110433\du}}
\pgfpathlineto{\pgfpoint{-13.500148\du}{6.097893\du}}
\pgfpathlineto{\pgfpoint{-13.428495\du}{6.134317\du}}
\pgfpathlineto{\pgfpoint{-13.405208\du}{6.122375\du}}
\pgfpathlineto{\pgfpoint{-13.416852\du}{6.145960\du}}
\pgfpathlineto{\pgfpoint{-13.500148\du}{6.145960\du}}
\pgfpathlineto{\pgfpoint{-13.464322\du}{6.134317\du}}
\pgfpathlineto{\pgfpoint{-13.524032\du}{6.110433\du}}
\pgfusepath{fill}
\pgfsetbuttcap
\pgfsetmiterjoin
\pgfsetdash{}{0pt}
\definecolor{dialinecolor}{rgb}{0.000000, 0.000000, 0.000000}
\pgfsetfillcolor{dialinecolor}
\pgfpathmoveto{\pgfpoint{-13.273547\du}{6.194027\du}}
\pgfpathlineto{\pgfpoint{-13.297730\du}{6.205372\du}}
\pgfpathlineto{\pgfpoint{-13.369084\du}{6.169844\du}}
\pgfpathlineto{\pgfpoint{-13.405208\du}{6.182085\du}}
\pgfpathlineto{\pgfpoint{-13.380727\du}{6.158201\du}}
\pgfpathlineto{\pgfpoint{-13.297730\du}{6.158201\du}}
\pgfpathlineto{\pgfpoint{-13.333257\du}{6.169844\du}}
\pgfpathlineto{\pgfpoint{-13.273547\du}{6.194027\du}}
\pgfusepath{fill}
\pgfsetbuttcap
\pgfsetmiterjoin
\pgfsetdash{}{0pt}
\definecolor{dialinecolor}{rgb}{1.000000, 1.000000, 1.000000}
\pgfsetfillcolor{dialinecolor}
\pgfpathmoveto{\pgfpoint{-13.392968\du}{6.145960\du}}
\pgfpathlineto{\pgfpoint{-13.369084\du}{6.145960\du}}
\pgfpathlineto{\pgfpoint{-13.297730\du}{6.122375\du}}
\pgfpathlineto{\pgfpoint{-13.261903\du}{6.134317\du}}
\pgfpathlineto{\pgfpoint{-13.285788\du}{6.110433\du}}
\pgfpathlineto{\pgfpoint{-13.356843\du}{6.110433\du}}
\pgfpathlineto{\pgfpoint{-13.333257\du}{6.110433\du}}
\pgfpathlineto{\pgfpoint{-13.392968\du}{6.145960\du}}
\pgfusepath{fill}
\pgfsetbuttcap
\pgfsetmiterjoin
\pgfsetdash{}{0pt}
\definecolor{dialinecolor}{rgb}{1.000000, 1.000000, 1.000000}
\pgfsetfillcolor{dialinecolor}
\pgfpathmoveto{\pgfpoint{-13.405208\du}{6.169844\du}}
\pgfpathlineto{\pgfpoint{-13.428495\du}{6.158201\du}}
\pgfpathlineto{\pgfpoint{-13.500148\du}{6.182085\du}}
\pgfpathlineto{\pgfpoint{-13.535377\du}{6.182085\du}}
\pgfpathlineto{\pgfpoint{-13.511791\du}{6.205372\du}}
\pgfpathlineto{\pgfpoint{-13.428495\du}{6.205372\du}}
\pgfpathlineto{\pgfpoint{-13.464322\du}{6.194027\du}}
\pgfpathlineto{\pgfpoint{-13.405208\du}{6.169844\du}}
\pgfusepath{fill}
\pgfsetbuttcap
\pgfsetmiterjoin
\pgfsetdash{}{0pt}
\definecolor{dialinecolor}{rgb}{1.000000, 1.000000, 1.000000}
\pgfsetfillcolor{dialinecolor}
\pgfpathmoveto{\pgfpoint{-13.524032\du}{6.110433\du}}
\pgfpathlineto{\pgfpoint{-13.500148\du}{6.110433\du}}
\pgfpathlineto{\pgfpoint{-13.428495\du}{6.134317\du}}
\pgfpathlineto{\pgfpoint{-13.392968\du}{6.122375\du}}
\pgfpathlineto{\pgfpoint{-13.416852\du}{6.145960\du}}
\pgfpathlineto{\pgfpoint{-13.500148\du}{6.145960\du}}
\pgfpathlineto{\pgfpoint{-13.464322\du}{6.145960\du}}
\pgfpathlineto{\pgfpoint{-13.524032\du}{6.110433\du}}
\pgfusepath{fill}
\pgfsetbuttcap
\pgfsetmiterjoin
\pgfsetdash{}{0pt}
\definecolor{dialinecolor}{rgb}{1.000000, 1.000000, 1.000000}
\pgfsetfillcolor{dialinecolor}
\pgfpathmoveto{\pgfpoint{-13.273547\du}{6.194027\du}}
\pgfpathlineto{\pgfpoint{-13.297730\du}{6.205372\du}}
\pgfpathlineto{\pgfpoint{-13.369084\du}{6.182085\du}}
\pgfpathlineto{\pgfpoint{-13.392968\du}{6.182085\du}}
\pgfpathlineto{\pgfpoint{-13.380727\du}{6.158201\du}}
\pgfpathlineto{\pgfpoint{-13.297730\du}{6.158201\du}}
\pgfpathlineto{\pgfpoint{-13.333257\du}{6.169844\du}}
\pgfpathlineto{\pgfpoint{-13.273547\du}{6.194027\du}}
\pgfusepath{fill}
\pgfsetbuttcap
\pgfsetmiterjoin
\pgfsetdash{}{0pt}
\definecolor{dialinecolor}{rgb}{0.666667, 0.901961, 1.000000}
\pgfsetstrokecolor{dialinecolor}
\pgfpathmoveto{\pgfpoint{-13.595087\du}{6.145960\du}}
\pgfpathlineto{\pgfpoint{-13.595087\du}{6.240900\du}}
\pgfusepath{stroke}
\pgfsetbuttcap
\pgfsetmiterjoin
\pgfsetdash{}{0pt}
\definecolor{dialinecolor}{rgb}{0.666667, 0.901961, 1.000000}
\pgfsetstrokecolor{dialinecolor}
\pgfpathmoveto{\pgfpoint{-13.214434\du}{6.145960\du}}
\pgfpathlineto{\pgfpoint{-13.214434\du}{6.240900\du}}
\pgfusepath{stroke}
\pgfsetbuttcap
\pgfsetmiterjoin
\pgfsetdash{}{0pt}
\definecolor{dialinecolor}{rgb}{0.000000, 0.470588, 0.666667}
\pgfsetfillcolor{dialinecolor}
\pgfpathmoveto{\pgfpoint{-13.559858\du}{6.652006\du}}
\pgfpathlineto{\pgfpoint{-13.560157\du}{6.648423\du}}
\pgfpathlineto{\pgfpoint{-13.561053\du}{6.645139\du}}
\pgfpathlineto{\pgfpoint{-13.562247\du}{6.641855\du}}
\pgfpathlineto{\pgfpoint{-13.564038\du}{6.638870\du}}
\pgfpathlineto{\pgfpoint{-13.566128\du}{6.635586\du}}
\pgfpathlineto{\pgfpoint{-13.569412\du}{6.632301\du}}
\pgfpathlineto{\pgfpoint{-13.572099\du}{6.629316\du}}
\pgfpathlineto{\pgfpoint{-13.575980\du}{6.626032\du}}
\pgfpathlineto{\pgfpoint{-13.580160\du}{6.623345\du}}
\pgfpathlineto{\pgfpoint{-13.584638\du}{6.620061\du}}
\pgfpathlineto{\pgfpoint{-13.590012\du}{6.617374\du}}
\pgfpathlineto{\pgfpoint{-13.595087\du}{6.614687\du}}
\pgfpathlineto{\pgfpoint{-13.601059\du}{6.612000\du}}
\pgfpathlineto{\pgfpoint{-13.607030\du}{6.609313\du}}
\pgfpathlineto{\pgfpoint{-13.614195\du}{6.606925\du}}
\pgfpathlineto{\pgfpoint{-13.621062\du}{6.604536\du}}
\pgfpathlineto{\pgfpoint{-13.628525\du}{6.602745\du}}
\pgfpathlineto{\pgfpoint{-13.635989\du}{6.600356\du}}
\pgfpathlineto{\pgfpoint{-13.644050\du}{6.598565\du}}
\pgfpathlineto{\pgfpoint{-13.652409\du}{6.596177\du}}
\pgfpathlineto{\pgfpoint{-13.660769\du}{6.594684\du}}
\pgfpathlineto{\pgfpoint{-13.669725\du}{6.593490\du}}
\pgfpathlineto{\pgfpoint{-13.678682\du}{6.591997\du}}
\pgfpathlineto{\pgfpoint{-13.687937\du}{6.590803\du}}
\pgfpathlineto{\pgfpoint{-13.697789\du}{6.589609\du}}
\pgfpathlineto{\pgfpoint{-13.707044\du}{6.588713\du}}
\pgfpathlineto{\pgfpoint{-13.716897\du}{6.588116\du}}
\pgfpathlineto{\pgfpoint{-13.726450\du}{6.587519\du}}
\pgfpathlineto{\pgfpoint{-13.736303\du}{6.586922\du}}
\pgfpathlineto{\pgfpoint{-13.746752\du}{6.586623\du}}
\pgfpathlineto{\pgfpoint{-13.756007\du}{6.586623\du}}
\pgfpathlineto{\pgfpoint{-13.756007\du}{6.586623\du}}
\pgfpathlineto{\pgfpoint{-13.765859\du}{6.586623\du}}
\pgfpathlineto{\pgfpoint{-13.776010\du}{6.586922\du}}
\pgfpathlineto{\pgfpoint{-13.785862\du}{6.587519\du}}
\pgfpathlineto{\pgfpoint{-13.795714\du}{6.588116\du}}
\pgfpathlineto{\pgfpoint{-13.805567\du}{6.588713\du}}
\pgfpathlineto{\pgfpoint{-13.814822\du}{6.589609\du}}
\pgfpathlineto{\pgfpoint{-13.824674\du}{6.590803\du}}
\pgfpathlineto{\pgfpoint{-13.833929\du}{6.591997\du}}
\pgfpathlineto{\pgfpoint{-13.842886\du}{6.593490\du}}
\pgfpathlineto{\pgfpoint{-13.851544\du}{6.594684\du}}
\pgfpathlineto{\pgfpoint{-13.859903\du}{6.596177\du}}
\pgfpathlineto{\pgfpoint{-13.868263\du}{6.598565\du}}
\pgfpathlineto{\pgfpoint{-13.876323\du}{6.600356\du}}
\pgfpathlineto{\pgfpoint{-13.883787\du}{6.602745\du}}
\pgfpathlineto{\pgfpoint{-13.891251\du}{6.604536\du}}
\pgfpathlineto{\pgfpoint{-13.898416\du}{6.606925\du}}
\pgfpathlineto{\pgfpoint{-13.905283\du}{6.609313\du}}
\pgfpathlineto{\pgfpoint{-13.911254\du}{6.612000\du}}
\pgfpathlineto{\pgfpoint{-13.917524\du}{6.614687\du}}
\pgfpathlineto{\pgfpoint{-13.922599\du}{6.617374\du}}
\pgfpathlineto{\pgfpoint{-13.927674\du}{6.620061\du}}
\pgfpathlineto{\pgfpoint{-13.932451\du}{6.623345\du}}
\pgfpathlineto{\pgfpoint{-13.936631\du}{6.626032\du}}
\pgfpathlineto{\pgfpoint{-13.940214\du}{6.629316\du}}
\pgfpathlineto{\pgfpoint{-13.943199\du}{6.632301\du}}
\pgfpathlineto{\pgfpoint{-13.945886\du}{6.635586\du}}
\pgfpathlineto{\pgfpoint{-13.948275\du}{6.638870\du}}
\pgfpathlineto{\pgfpoint{-13.950066\du}{6.641855\du}}
\pgfpathlineto{\pgfpoint{-13.951559\du}{6.645139\du}}
\pgfpathlineto{\pgfpoint{-13.952454\du}{6.648423\du}}
\pgfpathlineto{\pgfpoint{-13.952454\du}{6.652006\du}}
\pgfpathlineto{\pgfpoint{-13.952454\du}{6.652006\du}}
\pgfpathlineto{\pgfpoint{-13.952454\du}{6.655290\du}}
\pgfpathlineto{\pgfpoint{-13.951559\du}{6.658574\du}}
\pgfpathlineto{\pgfpoint{-13.950066\du}{6.661858\du}}
\pgfpathlineto{\pgfpoint{-13.948275\du}{6.665142\du}}
\pgfpathlineto{\pgfpoint{-13.945886\du}{6.668128\du}}
\pgfpathlineto{\pgfpoint{-13.943199\du}{6.671412\du}}
\pgfpathlineto{\pgfpoint{-13.940214\du}{6.674397\du}}
\pgfpathlineto{\pgfpoint{-13.936631\du}{6.677681\du}}
\pgfpathlineto{\pgfpoint{-13.932451\du}{6.680667\du}}
\pgfpathlineto{\pgfpoint{-13.927674\du}{6.683652\du}}
\pgfpathlineto{\pgfpoint{-13.922599\du}{6.686339\du}}
\pgfpathlineto{\pgfpoint{-13.917524\du}{6.689026\du}}
\pgfpathlineto{\pgfpoint{-13.911254\du}{6.691713\du}}
\pgfpathlineto{\pgfpoint{-13.905283\du}{6.694102\du}}
\pgfpathlineto{\pgfpoint{-13.898416\du}{6.696789\du}}
\pgfpathlineto{\pgfpoint{-13.891251\du}{6.699177\du}}
\pgfpathlineto{\pgfpoint{-13.883787\du}{6.701267\du}}
\pgfpathlineto{\pgfpoint{-13.876323\du}{6.703655\du}}
\pgfpathlineto{\pgfpoint{-13.868263\du}{6.705148\du}}
\pgfpathlineto{\pgfpoint{-13.859903\du}{6.707238\du}}
\pgfpathlineto{\pgfpoint{-13.851544\du}{6.709029\du}}
\pgfpathlineto{\pgfpoint{-13.842886\du}{6.710224\du}}
\pgfpathlineto{\pgfpoint{-13.833929\du}{6.711716\du}}
\pgfpathlineto{\pgfpoint{-13.824674\du}{6.712911\du}}
\pgfpathlineto{\pgfpoint{-13.814822\du}{6.714105\du}}
\pgfpathlineto{\pgfpoint{-13.805567\du}{6.715000\du}}
\pgfpathlineto{\pgfpoint{-13.795714\du}{6.715598\du}}
\pgfpathlineto{\pgfpoint{-13.785862\du}{6.716493\du}}
\pgfpathlineto{\pgfpoint{-13.776010\du}{6.716493\du}}
\pgfpathlineto{\pgfpoint{-13.765859\du}{6.717090\du}}
\pgfpathlineto{\pgfpoint{-13.756007\du}{6.717090\du}}
\pgfpathlineto{\pgfpoint{-13.756007\du}{6.717090\du}}
\pgfpathlineto{\pgfpoint{-13.746752\du}{6.717090\du}}
\pgfpathlineto{\pgfpoint{-13.736303\du}{6.716493\du}}
\pgfpathlineto{\pgfpoint{-13.726450\du}{6.716493\du}}
\pgfpathlineto{\pgfpoint{-13.716897\du}{6.715598\du}}
\pgfpathlineto{\pgfpoint{-13.707044\du}{6.715000\du}}
\pgfpathlineto{\pgfpoint{-13.697789\du}{6.714105\du}}
\pgfpathlineto{\pgfpoint{-13.687937\du}{6.712911\du}}
\pgfpathlineto{\pgfpoint{-13.678682\du}{6.711716\du}}
\pgfpathlineto{\pgfpoint{-13.669725\du}{6.710224\du}}
\pgfpathlineto{\pgfpoint{-13.660769\du}{6.709029\du}}
\pgfpathlineto{\pgfpoint{-13.652409\du}{6.707238\du}}
\pgfpathlineto{\pgfpoint{-13.644050\du}{6.705148\du}}
\pgfpathlineto{\pgfpoint{-13.635989\du}{6.703655\du}}
\pgfpathlineto{\pgfpoint{-13.628525\du}{6.701267\du}}
\pgfpathlineto{\pgfpoint{-13.621062\du}{6.699177\du}}
\pgfpathlineto{\pgfpoint{-13.614195\du}{6.696789\du}}
\pgfpathlineto{\pgfpoint{-13.607030\du}{6.694102\du}}
\pgfpathlineto{\pgfpoint{-13.601059\du}{6.691713\du}}
\pgfpathlineto{\pgfpoint{-13.595087\du}{6.689026\du}}
\pgfpathlineto{\pgfpoint{-13.590012\du}{6.686339\du}}
\pgfpathlineto{\pgfpoint{-13.584638\du}{6.683652\du}}
\pgfpathlineto{\pgfpoint{-13.580160\du}{6.680667\du}}
\pgfpathlineto{\pgfpoint{-13.575980\du}{6.677681\du}}
\pgfpathlineto{\pgfpoint{-13.572099\du}{6.674397\du}}
\pgfpathlineto{\pgfpoint{-13.569412\du}{6.671412\du}}
\pgfpathlineto{\pgfpoint{-13.566128\du}{6.668128\du}}
\pgfpathlineto{\pgfpoint{-13.564038\du}{6.665142\du}}
\pgfpathlineto{\pgfpoint{-13.562247\du}{6.661858\du}}
\pgfpathlineto{\pgfpoint{-13.561053\du}{6.658574\du}}
\pgfpathlineto{\pgfpoint{-13.560157\du}{6.655290\du}}
\pgfpathlineto{\pgfpoint{-13.559858\du}{6.652006\du}}
\pgfusepath{fill}
\pgfsetbuttcap
\pgfsetmiterjoin
\pgfsetdash{}{0pt}
\definecolor{dialinecolor}{rgb}{0.666667, 0.901961, 1.000000}
\pgfsetstrokecolor{dialinecolor}
\pgfpathmoveto{\pgfpoint{-13.571800\du}{6.645736\du}}
\pgfpathlineto{\pgfpoint{-13.572099\du}{6.642751\du}}
\pgfpathlineto{\pgfpoint{-13.572696\du}{6.640064\du}}
\pgfpathlineto{\pgfpoint{-13.573890\du}{6.637078\du}}
\pgfpathlineto{\pgfpoint{-13.575980\du}{6.633496\du}}
\pgfpathlineto{\pgfpoint{-13.578070\du}{6.631107\du}}
\pgfpathlineto{\pgfpoint{-13.580160\du}{6.627823\du}}
\pgfpathlineto{\pgfpoint{-13.583444\du}{6.625136\du}}
\pgfpathlineto{\pgfpoint{-13.586728\du}{6.622449\du}}
\pgfpathlineto{\pgfpoint{-13.590908\du}{6.619464\du}}
\pgfpathlineto{\pgfpoint{-13.595685\du}{6.617075\du}}
\pgfpathlineto{\pgfpoint{-13.600760\du}{6.614388\du}}
\pgfpathlineto{\pgfpoint{-13.605835\du}{6.612000\du}}
\pgfpathlineto{\pgfpoint{-13.611508\du}{6.609612\du}}
\pgfpathlineto{\pgfpoint{-13.617479\du}{6.607223\du}}
\pgfpathlineto{\pgfpoint{-13.624047\du}{6.604536\du}}
\pgfpathlineto{\pgfpoint{-13.630914\du}{6.603043\du}}
\pgfpathlineto{\pgfpoint{-13.637482\du}{6.600954\du}}
\pgfpathlineto{\pgfpoint{-13.645543\du}{6.599162\du}}
\pgfpathlineto{\pgfpoint{-13.653007\du}{6.597669\du}}
\pgfpathlineto{\pgfpoint{-13.660769\du}{6.595580\du}}
\pgfpathlineto{\pgfpoint{-13.669725\du}{6.594087\du}}
\pgfpathlineto{\pgfpoint{-13.678085\du}{6.592594\du}}
\pgfpathlineto{\pgfpoint{-13.686444\du}{6.591400\du}}
\pgfpathlineto{\pgfpoint{-13.695700\du}{6.590206\du}}
\pgfpathlineto{\pgfpoint{-13.704358\du}{6.589310\du}}
\pgfpathlineto{\pgfpoint{-13.713911\du}{6.588414\du}}
\pgfpathlineto{\pgfpoint{-13.723465\du}{6.588116\du}}
\pgfpathlineto{\pgfpoint{-13.732720\du}{6.587519\du}}
\pgfpathlineto{\pgfpoint{-13.742572\du}{6.586922\du}}
\pgfpathlineto{\pgfpoint{-13.752424\du}{6.586623\du}}
\pgfpathlineto{\pgfpoint{-13.761978\du}{6.586623\du}}
\pgfpathlineto{\pgfpoint{-13.761978\du}{6.586623\du}}
\pgfpathlineto{\pgfpoint{-13.771233\du}{6.586623\du}}
\pgfpathlineto{\pgfpoint{-13.781085\du}{6.586922\du}}
\pgfpathlineto{\pgfpoint{-13.790639\du}{6.587519\du}}
\pgfpathlineto{\pgfpoint{-13.800193\du}{6.588116\du}}
\pgfpathlineto{\pgfpoint{-13.809746\du}{6.588414\du}}
\pgfpathlineto{\pgfpoint{-13.819001\du}{6.589310\du}}
\pgfpathlineto{\pgfpoint{-13.827958\du}{6.590206\du}}
\pgfpathlineto{\pgfpoint{-13.837213\du}{6.591400\du}}
\pgfpathlineto{\pgfpoint{-13.845573\du}{6.592594\du}}
\pgfpathlineto{\pgfpoint{-13.854529\du}{6.594087\du}}
\pgfpathlineto{\pgfpoint{-13.862889\du}{6.595580\du}}
\pgfpathlineto{\pgfpoint{-13.870651\du}{6.597669\du}}
\pgfpathlineto{\pgfpoint{-13.878712\du}{6.599162\du}}
\pgfpathlineto{\pgfpoint{-13.886176\du}{6.600954\du}}
\pgfpathlineto{\pgfpoint{-13.892744\du}{6.603043\du}}
\pgfpathlineto{\pgfpoint{-13.900208\du}{6.604536\du}}
\pgfpathlineto{\pgfpoint{-13.906179\du}{6.607223\du}}
\pgfpathlineto{\pgfpoint{-13.912150\du}{6.609612\du}}
\pgfpathlineto{\pgfpoint{-13.918419\du}{6.612000\du}}
\pgfpathlineto{\pgfpoint{-13.923495\du}{6.614388\du}}
\pgfpathlineto{\pgfpoint{-13.928570\du}{6.617075\du}}
\pgfpathlineto{\pgfpoint{-13.932750\du}{6.619464\du}}
\pgfpathlineto{\pgfpoint{-13.936930\du}{6.622449\du}}
\pgfpathlineto{\pgfpoint{-13.940811\du}{6.625136\du}}
\pgfpathlineto{\pgfpoint{-13.944095\du}{6.627823\du}}
\pgfpathlineto{\pgfpoint{-13.945886\du}{6.631107\du}}
\pgfpathlineto{\pgfpoint{-13.948275\du}{6.633496\du}}
\pgfpathlineto{\pgfpoint{-13.950066\du}{6.637078\du}}
\pgfpathlineto{\pgfpoint{-13.951559\du}{6.640064\du}}
\pgfpathlineto{\pgfpoint{-13.951857\du}{6.642751\du}}
\pgfpathlineto{\pgfpoint{-13.952454\du}{6.645736\du}}
\pgfpathlineto{\pgfpoint{-13.952454\du}{6.645736\du}}
\pgfpathlineto{\pgfpoint{-13.951857\du}{6.649319\du}}
\pgfpathlineto{\pgfpoint{-13.951559\du}{6.652006\du}}
\pgfpathlineto{\pgfpoint{-13.950066\du}{6.654991\du}}
\pgfpathlineto{\pgfpoint{-13.948275\du}{6.657977\du}}
\pgfpathlineto{\pgfpoint{-13.945886\du}{6.660664\du}}
\pgfpathlineto{\pgfpoint{-13.944095\du}{6.663948\du}}
\pgfpathlineto{\pgfpoint{-13.940811\du}{6.666635\du}}
\pgfpathlineto{\pgfpoint{-13.936930\du}{6.669620\du}}
\pgfpathlineto{\pgfpoint{-13.932750\du}{6.672307\du}}
\pgfpathlineto{\pgfpoint{-13.928570\du}{6.674994\du}}
\pgfpathlineto{\pgfpoint{-13.923495\du}{6.677681\du}}
\pgfpathlineto{\pgfpoint{-13.918419\du}{6.679174\du}}
\pgfpathlineto{\pgfpoint{-13.912150\du}{6.682458\du}}
\pgfpathlineto{\pgfpoint{-13.906179\du}{6.684548\du}}
\pgfpathlineto{\pgfpoint{-13.900208\du}{6.686638\du}}
\pgfpathlineto{\pgfpoint{-13.892744\du}{6.689026\du}}
\pgfpathlineto{\pgfpoint{-13.886176\du}{6.691116\du}}
\pgfpathlineto{\pgfpoint{-13.878712\du}{6.692908\du}}
\pgfpathlineto{\pgfpoint{-13.870651\du}{6.694102\du}}
\pgfpathlineto{\pgfpoint{-13.862889\du}{6.696490\du}}
\pgfpathlineto{\pgfpoint{-13.854529\du}{6.697983\du}}
\pgfpathlineto{\pgfpoint{-13.845573\du}{6.699476\du}}
\pgfpathlineto{\pgfpoint{-13.837213\du}{6.700670\du}}
\pgfpathlineto{\pgfpoint{-13.827958\du}{6.701566\du}}
\pgfpathlineto{\pgfpoint{-13.819001\du}{6.702461\du}}
\pgfpathlineto{\pgfpoint{-13.809746\du}{6.703655\du}}
\pgfpathlineto{\pgfpoint{-13.800193\du}{6.703954\du}}
\pgfpathlineto{\pgfpoint{-13.790639\du}{6.704551\du}}
\pgfpathlineto{\pgfpoint{-13.781085\du}{6.704850\du}}
\pgfpathlineto{\pgfpoint{-13.771233\du}{6.705148\du}}
\pgfpathlineto{\pgfpoint{-13.761978\du}{6.705148\du}}
\pgfpathlineto{\pgfpoint{-13.761978\du}{6.705148\du}}
\pgfpathlineto{\pgfpoint{-13.752424\du}{6.705148\du}}
\pgfpathlineto{\pgfpoint{-13.742572\du}{6.704850\du}}
\pgfpathlineto{\pgfpoint{-13.732720\du}{6.704551\du}}
\pgfpathlineto{\pgfpoint{-13.723465\du}{6.703954\du}}
\pgfpathlineto{\pgfpoint{-13.713911\du}{6.703655\du}}
\pgfpathlineto{\pgfpoint{-13.704358\du}{6.702461\du}}
\pgfpathlineto{\pgfpoint{-13.695700\du}{6.701566\du}}
\pgfpathlineto{\pgfpoint{-13.686444\du}{6.700670\du}}
\pgfpathlineto{\pgfpoint{-13.678085\du}{6.699476\du}}
\pgfpathlineto{\pgfpoint{-13.669725\du}{6.697983\du}}
\pgfpathlineto{\pgfpoint{-13.660769\du}{6.696490\du}}
\pgfpathlineto{\pgfpoint{-13.653007\du}{6.694102\du}}
\pgfpathlineto{\pgfpoint{-13.645543\du}{6.692908\du}}
\pgfpathlineto{\pgfpoint{-13.637482\du}{6.691116\du}}
\pgfpathlineto{\pgfpoint{-13.630914\du}{6.689026\du}}
\pgfpathlineto{\pgfpoint{-13.624047\du}{6.686638\du}}
\pgfpathlineto{\pgfpoint{-13.617479\du}{6.684548\du}}
\pgfpathlineto{\pgfpoint{-13.611508\du}{6.682458\du}}
\pgfpathlineto{\pgfpoint{-13.605835\du}{6.679174\du}}
\pgfpathlineto{\pgfpoint{-13.600760\du}{6.677681\du}}
\pgfpathlineto{\pgfpoint{-13.595685\du}{6.674994\du}}
\pgfpathlineto{\pgfpoint{-13.590908\du}{6.672307\du}}
\pgfpathlineto{\pgfpoint{-13.586728\du}{6.669620\du}}
\pgfpathlineto{\pgfpoint{-13.583444\du}{6.666635\du}}
\pgfpathlineto{\pgfpoint{-13.580160\du}{6.663948\du}}
\pgfpathlineto{\pgfpoint{-13.578070\du}{6.660664\du}}
\pgfpathlineto{\pgfpoint{-13.575980\du}{6.657977\du}}
\pgfpathlineto{\pgfpoint{-13.573890\du}{6.654991\du}}
\pgfpathlineto{\pgfpoint{-13.572696\du}{6.652006\du}}
\pgfpathlineto{\pgfpoint{-13.572099\du}{6.649319\du}}
\pgfpathlineto{\pgfpoint{-13.571800\du}{6.645736\du}}
\pgfusepath{stroke}
\pgfsetbuttcap
\pgfsetmiterjoin
\pgfsetdash{}{0pt}
\definecolor{dialinecolor}{rgb}{0.000000, 0.470588, 0.666667}
\pgfsetfillcolor{dialinecolor}
\pgfpathmoveto{\pgfpoint{-13.952454\du}{6.562739\du}}
\pgfpathlineto{\pgfpoint{-13.952454\du}{6.657977\du}}
\pgfpathlineto{\pgfpoint{-13.559560\du}{6.657977\du}}
\pgfpathlineto{\pgfpoint{-13.559560\du}{6.562739\du}}
\pgfpathlineto{\pgfpoint{-13.952454\du}{6.562739\du}}
\pgfusepath{fill}
\pgfsetbuttcap
\pgfsetmiterjoin
\pgfsetdash{}{0pt}
\definecolor{dialinecolor}{rgb}{0.000000, 0.705882, 1.000000}
\pgfsetfillcolor{dialinecolor}
\pgfpathmoveto{\pgfpoint{-13.559858\du}{6.556469\du}}
\pgfpathlineto{\pgfpoint{-13.560157\du}{6.552887\du}}
\pgfpathlineto{\pgfpoint{-13.561053\du}{6.549603\du}}
\pgfpathlineto{\pgfpoint{-13.562247\du}{6.546617\du}}
\pgfpathlineto{\pgfpoint{-13.564038\du}{6.543034\du}}
\pgfpathlineto{\pgfpoint{-13.566128\du}{6.539750\du}}
\pgfpathlineto{\pgfpoint{-13.569412\du}{6.536466\du}}
\pgfpathlineto{\pgfpoint{-13.572099\du}{6.533779\du}}
\pgfpathlineto{\pgfpoint{-13.575980\du}{6.529898\du}}
\pgfpathlineto{\pgfpoint{-13.580160\du}{6.527510\du}}
\pgfpathlineto{\pgfpoint{-13.584638\du}{6.524524\du}}
\pgfpathlineto{\pgfpoint{-13.590012\du}{6.521539\du}}
\pgfpathlineto{\pgfpoint{-13.595087\du}{6.518852\du}}
\pgfpathlineto{\pgfpoint{-13.601059\du}{6.516165\du}}
\pgfpathlineto{\pgfpoint{-13.607030\du}{6.513776\du}}
\pgfpathlineto{\pgfpoint{-13.614195\du}{6.511089\du}}
\pgfpathlineto{\pgfpoint{-13.621062\du}{6.508701\du}}
\pgfpathlineto{\pgfpoint{-13.628525\du}{6.506611\du}}
\pgfpathlineto{\pgfpoint{-13.635989\du}{6.504521\du}}
\pgfpathlineto{\pgfpoint{-13.644050\du}{6.502730\du}}
\pgfpathlineto{\pgfpoint{-13.652409\du}{6.500640\du}}
\pgfpathlineto{\pgfpoint{-13.660769\du}{6.498849\du}}
\pgfpathlineto{\pgfpoint{-13.669725\du}{6.497655\du}}
\pgfpathlineto{\pgfpoint{-13.678682\du}{6.496162\du}}
\pgfpathlineto{\pgfpoint{-13.687937\du}{6.494968\du}}
\pgfpathlineto{\pgfpoint{-13.697789\du}{6.493773\du}}
\pgfpathlineto{\pgfpoint{-13.707044\du}{6.492579\du}}
\pgfpathlineto{\pgfpoint{-13.716897\du}{6.492281\du}}
\pgfpathlineto{\pgfpoint{-13.726450\du}{6.491385\du}}
\pgfpathlineto{\pgfpoint{-13.736303\du}{6.491086\du}}
\pgfpathlineto{\pgfpoint{-13.746752\du}{6.490788\du}}
\pgfpathlineto{\pgfpoint{-13.756007\du}{6.490788\du}}
\pgfpathlineto{\pgfpoint{-13.756007\du}{6.490788\du}}
\pgfpathlineto{\pgfpoint{-13.765859\du}{6.490788\du}}
\pgfpathlineto{\pgfpoint{-13.776010\du}{6.491086\du}}
\pgfpathlineto{\pgfpoint{-13.785862\du}{6.491385\du}}
\pgfpathlineto{\pgfpoint{-13.795714\du}{6.492281\du}}
\pgfpathlineto{\pgfpoint{-13.805567\du}{6.492579\du}}
\pgfpathlineto{\pgfpoint{-13.814822\du}{6.493773\du}}
\pgfpathlineto{\pgfpoint{-13.824674\du}{6.494968\du}}
\pgfpathlineto{\pgfpoint{-13.833929\du}{6.496162\du}}
\pgfpathlineto{\pgfpoint{-13.842886\du}{6.497655\du}}
\pgfpathlineto{\pgfpoint{-13.851544\du}{6.498849\du}}
\pgfpathlineto{\pgfpoint{-13.859903\du}{6.500640\du}}
\pgfpathlineto{\pgfpoint{-13.868263\du}{6.502730\du}}
\pgfpathlineto{\pgfpoint{-13.876323\du}{6.504521\du}}
\pgfpathlineto{\pgfpoint{-13.883787\du}{6.506611\du}}
\pgfpathlineto{\pgfpoint{-13.891251\du}{6.508701\du}}
\pgfpathlineto{\pgfpoint{-13.898416\du}{6.511089\du}}
\pgfpathlineto{\pgfpoint{-13.905283\du}{6.513776\du}}
\pgfpathlineto{\pgfpoint{-13.911254\du}{6.516165\du}}
\pgfpathlineto{\pgfpoint{-13.917524\du}{6.518852\du}}
\pgfpathlineto{\pgfpoint{-13.922599\du}{6.521539\du}}
\pgfpathlineto{\pgfpoint{-13.927674\du}{6.524524\du}}
\pgfpathlineto{\pgfpoint{-13.932451\du}{6.527510\du}}
\pgfpathlineto{\pgfpoint{-13.936631\du}{6.529898\du}}
\pgfpathlineto{\pgfpoint{-13.940214\du}{6.533779\du}}
\pgfpathlineto{\pgfpoint{-13.943199\du}{6.536466\du}}
\pgfpathlineto{\pgfpoint{-13.945886\du}{6.539750\du}}
\pgfpathlineto{\pgfpoint{-13.948275\du}{6.543034\du}}
\pgfpathlineto{\pgfpoint{-13.950066\du}{6.546617\du}}
\pgfpathlineto{\pgfpoint{-13.951559\du}{6.549603\du}}
\pgfpathlineto{\pgfpoint{-13.952454\du}{6.552887\du}}
\pgfpathlineto{\pgfpoint{-13.952454\du}{6.556469\du}}
\pgfpathlineto{\pgfpoint{-13.952454\du}{6.556469\du}}
\pgfpathlineto{\pgfpoint{-13.952454\du}{6.559753\du}}
\pgfpathlineto{\pgfpoint{-13.951559\du}{6.563037\du}}
\pgfpathlineto{\pgfpoint{-13.950066\du}{6.566023\du}}
\pgfpathlineto{\pgfpoint{-13.948275\du}{6.569606\du}}
\pgfpathlineto{\pgfpoint{-13.945886\du}{6.572890\du}}
\pgfpathlineto{\pgfpoint{-13.943199\du}{6.575875\du}}
\pgfpathlineto{\pgfpoint{-13.940214\du}{6.578861\du}}
\pgfpathlineto{\pgfpoint{-13.936631\du}{6.582145\du}}
\pgfpathlineto{\pgfpoint{-13.932451\du}{6.585130\du}}
\pgfpathlineto{\pgfpoint{-13.927674\du}{6.588116\du}}
\pgfpathlineto{\pgfpoint{-13.922599\du}{6.591101\du}}
\pgfpathlineto{\pgfpoint{-13.917524\du}{6.593788\du}}
\pgfpathlineto{\pgfpoint{-13.911254\du}{6.596177\du}}
\pgfpathlineto{\pgfpoint{-13.905283\du}{6.598864\du}}
\pgfpathlineto{\pgfpoint{-13.898416\du}{6.601551\du}}
\pgfpathlineto{\pgfpoint{-13.891251\du}{6.603640\du}}
\pgfpathlineto{\pgfpoint{-13.883787\du}{6.606029\du}}
\pgfpathlineto{\pgfpoint{-13.876323\du}{6.608119\du}}
\pgfpathlineto{\pgfpoint{-13.868263\du}{6.609910\du}}
\pgfpathlineto{\pgfpoint{-13.859903\du}{6.612000\du}}
\pgfpathlineto{\pgfpoint{-13.851544\du}{6.613791\du}}
\pgfpathlineto{\pgfpoint{-13.842886\du}{6.614985\du}}
\pgfpathlineto{\pgfpoint{-13.833929\du}{6.616478\du}}
\pgfpathlineto{\pgfpoint{-13.824674\du}{6.617672\du}}
\pgfpathlineto{\pgfpoint{-13.814822\du}{6.618568\du}}
\pgfpathlineto{\pgfpoint{-13.805567\du}{6.619464\du}}
\pgfpathlineto{\pgfpoint{-13.795714\du}{6.620359\du}}
\pgfpathlineto{\pgfpoint{-13.785862\du}{6.621255\du}}
\pgfpathlineto{\pgfpoint{-13.776010\du}{6.621554\du}}
\pgfpathlineto{\pgfpoint{-13.765859\du}{6.621852\du}}
\pgfpathlineto{\pgfpoint{-13.756007\du}{6.621852\du}}
\pgfpathlineto{\pgfpoint{-13.756007\du}{6.621852\du}}
\pgfpathlineto{\pgfpoint{-13.746752\du}{6.621852\du}}
\pgfpathlineto{\pgfpoint{-13.736303\du}{6.621554\du}}
\pgfpathlineto{\pgfpoint{-13.726450\du}{6.621255\du}}
\pgfpathlineto{\pgfpoint{-13.716897\du}{6.620359\du}}
\pgfpathlineto{\pgfpoint{-13.707044\du}{6.619464\du}}
\pgfpathlineto{\pgfpoint{-13.697789\du}{6.618568\du}}
\pgfpathlineto{\pgfpoint{-13.687937\du}{6.617672\du}}
\pgfpathlineto{\pgfpoint{-13.678682\du}{6.616478\du}}
\pgfpathlineto{\pgfpoint{-13.669725\du}{6.614985\du}}
\pgfpathlineto{\pgfpoint{-13.660769\du}{6.613791\du}}
\pgfpathlineto{\pgfpoint{-13.652409\du}{6.612000\du}}
\pgfpathlineto{\pgfpoint{-13.644050\du}{6.609910\du}}
\pgfpathlineto{\pgfpoint{-13.635989\du}{6.608119\du}}
\pgfpathlineto{\pgfpoint{-13.628525\du}{6.606029\du}}
\pgfpathlineto{\pgfpoint{-13.621062\du}{6.603640\du}}
\pgfpathlineto{\pgfpoint{-13.614195\du}{6.601551\du}}
\pgfpathlineto{\pgfpoint{-13.607030\du}{6.598864\du}}
\pgfpathlineto{\pgfpoint{-13.601059\du}{6.596177\du}}
\pgfpathlineto{\pgfpoint{-13.595087\du}{6.593788\du}}
\pgfpathlineto{\pgfpoint{-13.590012\du}{6.591101\du}}
\pgfpathlineto{\pgfpoint{-13.584638\du}{6.588116\du}}
\pgfpathlineto{\pgfpoint{-13.580160\du}{6.585130\du}}
\pgfpathlineto{\pgfpoint{-13.575980\du}{6.582145\du}}
\pgfpathlineto{\pgfpoint{-13.572099\du}{6.578861\du}}
\pgfpathlineto{\pgfpoint{-13.569412\du}{6.575875\du}}
\pgfpathlineto{\pgfpoint{-13.566128\du}{6.572890\du}}
\pgfpathlineto{\pgfpoint{-13.564038\du}{6.569606\du}}
\pgfpathlineto{\pgfpoint{-13.562247\du}{6.566023\du}}
\pgfpathlineto{\pgfpoint{-13.561053\du}{6.563037\du}}
\pgfpathlineto{\pgfpoint{-13.560157\du}{6.559753\du}}
\pgfpathlineto{\pgfpoint{-13.559858\du}{6.556469\du}}
\pgfusepath{fill}
\pgfsetbuttcap
\pgfsetmiterjoin
\pgfsetdash{}{0pt}
\definecolor{dialinecolor}{rgb}{0.666667, 0.901961, 1.000000}
\pgfsetstrokecolor{dialinecolor}
\pgfpathmoveto{\pgfpoint{-13.571800\du}{6.550797\du}}
\pgfpathlineto{\pgfpoint{-13.572099\du}{6.547811\du}}
\pgfpathlineto{\pgfpoint{-13.572696\du}{6.544826\du}}
\pgfpathlineto{\pgfpoint{-13.573890\du}{6.542139\du}}
\pgfpathlineto{\pgfpoint{-13.575980\du}{6.538855\du}}
\pgfpathlineto{\pgfpoint{-13.578070\du}{6.536168\du}}
\pgfpathlineto{\pgfpoint{-13.580160\du}{6.532884\du}}
\pgfpathlineto{\pgfpoint{-13.583444\du}{6.529898\du}}
\pgfpathlineto{\pgfpoint{-13.586728\du}{6.527510\du}}
\pgfpathlineto{\pgfpoint{-13.590908\du}{6.524823\du}}
\pgfpathlineto{\pgfpoint{-13.595685\du}{6.522136\du}}
\pgfpathlineto{\pgfpoint{-13.600760\du}{6.519449\du}}
\pgfpathlineto{\pgfpoint{-13.605835\du}{6.517060\du}}
\pgfpathlineto{\pgfpoint{-13.611508\du}{6.514075\du}}
\pgfpathlineto{\pgfpoint{-13.617479\du}{6.512284\du}}
\pgfpathlineto{\pgfpoint{-13.624047\du}{6.510194\du}}
\pgfpathlineto{\pgfpoint{-13.630914\du}{6.507507\du}}
\pgfpathlineto{\pgfpoint{-13.637482\du}{6.506014\du}}
\pgfpathlineto{\pgfpoint{-13.645543\du}{6.504223\du}}
\pgfpathlineto{\pgfpoint{-13.653007\du}{6.502133\du}}
\pgfpathlineto{\pgfpoint{-13.660769\du}{6.500640\du}}
\pgfpathlineto{\pgfpoint{-13.669725\du}{6.499147\du}}
\pgfpathlineto{\pgfpoint{-13.678085\du}{6.497655\du}}
\pgfpathlineto{\pgfpoint{-13.686444\du}{6.496460\du}}
\pgfpathlineto{\pgfpoint{-13.695700\du}{6.495266\du}}
\pgfpathlineto{\pgfpoint{-13.704358\du}{6.494072\du}}
\pgfpathlineto{\pgfpoint{-13.713911\du}{6.493475\du}}
\pgfpathlineto{\pgfpoint{-13.723465\du}{6.492579\du}}
\pgfpathlineto{\pgfpoint{-13.732720\du}{6.492579\du}}
\pgfpathlineto{\pgfpoint{-13.742572\du}{6.492281\du}}
\pgfpathlineto{\pgfpoint{-13.752424\du}{6.491385\du}}
\pgfpathlineto{\pgfpoint{-13.761978\du}{6.491385\du}}
\pgfpathlineto{\pgfpoint{-13.761978\du}{6.491385\du}}
\pgfpathlineto{\pgfpoint{-13.771233\du}{6.491385\du}}
\pgfpathlineto{\pgfpoint{-13.781085\du}{6.492281\du}}
\pgfpathlineto{\pgfpoint{-13.790639\du}{6.492579\du}}
\pgfpathlineto{\pgfpoint{-13.800193\du}{6.492579\du}}
\pgfpathlineto{\pgfpoint{-13.809746\du}{6.493475\du}}
\pgfpathlineto{\pgfpoint{-13.819001\du}{6.494072\du}}
\pgfpathlineto{\pgfpoint{-13.827958\du}{6.495266\du}}
\pgfpathlineto{\pgfpoint{-13.837213\du}{6.496460\du}}
\pgfpathlineto{\pgfpoint{-13.845573\du}{6.497655\du}}
\pgfpathlineto{\pgfpoint{-13.854529\du}{6.499147\du}}
\pgfpathlineto{\pgfpoint{-13.862889\du}{6.500640\du}}
\pgfpathlineto{\pgfpoint{-13.870651\du}{6.502133\du}}
\pgfpathlineto{\pgfpoint{-13.878712\du}{6.504223\du}}
\pgfpathlineto{\pgfpoint{-13.886176\du}{6.506014\du}}
\pgfpathlineto{\pgfpoint{-13.892744\du}{6.507507\du}}
\pgfpathlineto{\pgfpoint{-13.900208\du}{6.510194\du}}
\pgfpathlineto{\pgfpoint{-13.906179\du}{6.512284\du}}
\pgfpathlineto{\pgfpoint{-13.912150\du}{6.514075\du}}
\pgfpathlineto{\pgfpoint{-13.918419\du}{6.517060\du}}
\pgfpathlineto{\pgfpoint{-13.923495\du}{6.519449\du}}
\pgfpathlineto{\pgfpoint{-13.928570\du}{6.522136\du}}
\pgfpathlineto{\pgfpoint{-13.932750\du}{6.524823\du}}
\pgfpathlineto{\pgfpoint{-13.936930\du}{6.527510\du}}
\pgfpathlineto{\pgfpoint{-13.940811\du}{6.529898\du}}
\pgfpathlineto{\pgfpoint{-13.944095\du}{6.532884\du}}
\pgfpathlineto{\pgfpoint{-13.945886\du}{6.536168\du}}
\pgfpathlineto{\pgfpoint{-13.948275\du}{6.538855\du}}
\pgfpathlineto{\pgfpoint{-13.950066\du}{6.542139\du}}
\pgfpathlineto{\pgfpoint{-13.951559\du}{6.544826\du}}
\pgfpathlineto{\pgfpoint{-13.951857\du}{6.547811\du}}
\pgfpathlineto{\pgfpoint{-13.952454\du}{6.550797\du}}
\pgfpathlineto{\pgfpoint{-13.952454\du}{6.550797\du}}
\pgfpathlineto{\pgfpoint{-13.951857\du}{6.554379\du}}
\pgfpathlineto{\pgfpoint{-13.951559\du}{6.557066\du}}
\pgfpathlineto{\pgfpoint{-13.950066\du}{6.559753\du}}
\pgfpathlineto{\pgfpoint{-13.948275\du}{6.563037\du}}
\pgfpathlineto{\pgfpoint{-13.945886\du}{6.565724\du}}
\pgfpathlineto{\pgfpoint{-13.944095\du}{6.569307\du}}
\pgfpathlineto{\pgfpoint{-13.940811\du}{6.571994\du}}
\pgfpathlineto{\pgfpoint{-13.936930\du}{6.574681\du}}
\pgfpathlineto{\pgfpoint{-13.932750\du}{6.577368\du}}
\pgfpathlineto{\pgfpoint{-13.928570\du}{6.580055\du}}
\pgfpathlineto{\pgfpoint{-13.923495\du}{6.582742\du}}
\pgfpathlineto{\pgfpoint{-13.918419\du}{6.584832\du}}
\pgfpathlineto{\pgfpoint{-13.912150\du}{6.587519\du}}
\pgfpathlineto{\pgfpoint{-13.906179\du}{6.589609\du}}
\pgfpathlineto{\pgfpoint{-13.900208\du}{6.591997\du}}
\pgfpathlineto{\pgfpoint{-13.892744\du}{6.594087\du}}
\pgfpathlineto{\pgfpoint{-13.886176\du}{6.596177\du}}
\pgfpathlineto{\pgfpoint{-13.878712\du}{6.597968\du}}
\pgfpathlineto{\pgfpoint{-13.870651\du}{6.599461\du}}
\pgfpathlineto{\pgfpoint{-13.862889\du}{6.601551\du}}
\pgfpathlineto{\pgfpoint{-13.854529\du}{6.603043\du}}
\pgfpathlineto{\pgfpoint{-13.845573\du}{6.604536\du}}
\pgfpathlineto{\pgfpoint{-13.837213\du}{6.605730\du}}
\pgfpathlineto{\pgfpoint{-13.827958\du}{6.606925\du}}
\pgfpathlineto{\pgfpoint{-13.819001\du}{6.607522\du}}
\pgfpathlineto{\pgfpoint{-13.809746\du}{6.608716\du}}
\pgfpathlineto{\pgfpoint{-13.800193\du}{6.609014\du}}
\pgfpathlineto{\pgfpoint{-13.790639\du}{6.609612\du}}
\pgfpathlineto{\pgfpoint{-13.781085\du}{6.609910\du}}
\pgfpathlineto{\pgfpoint{-13.771233\du}{6.610507\du}}
\pgfpathlineto{\pgfpoint{-13.761978\du}{6.610507\du}}
\pgfpathlineto{\pgfpoint{-13.761978\du}{6.610507\du}}
\pgfpathlineto{\pgfpoint{-13.752424\du}{6.610507\du}}
\pgfpathlineto{\pgfpoint{-13.742572\du}{6.609910\du}}
\pgfpathlineto{\pgfpoint{-13.732720\du}{6.609612\du}}
\pgfpathlineto{\pgfpoint{-13.723465\du}{6.609014\du}}
\pgfpathlineto{\pgfpoint{-13.713911\du}{6.608716\du}}
\pgfpathlineto{\pgfpoint{-13.704358\du}{6.607522\du}}
\pgfpathlineto{\pgfpoint{-13.695700\du}{6.606925\du}}
\pgfpathlineto{\pgfpoint{-13.686444\du}{6.605730\du}}
\pgfpathlineto{\pgfpoint{-13.678085\du}{6.604536\du}}
\pgfpathlineto{\pgfpoint{-13.669725\du}{6.603043\du}}
\pgfpathlineto{\pgfpoint{-13.660769\du}{6.601551\du}}
\pgfpathlineto{\pgfpoint{-13.653007\du}{6.599461\du}}
\pgfpathlineto{\pgfpoint{-13.645543\du}{6.597968\du}}
\pgfpathlineto{\pgfpoint{-13.637482\du}{6.596177\du}}
\pgfpathlineto{\pgfpoint{-13.630914\du}{6.594087\du}}
\pgfpathlineto{\pgfpoint{-13.624047\du}{6.591997\du}}
\pgfpathlineto{\pgfpoint{-13.617479\du}{6.589609\du}}
\pgfpathlineto{\pgfpoint{-13.611508\du}{6.587519\du}}
\pgfpathlineto{\pgfpoint{-13.605835\du}{6.584832\du}}
\pgfpathlineto{\pgfpoint{-13.600760\du}{6.582742\du}}
\pgfpathlineto{\pgfpoint{-13.595685\du}{6.580055\du}}
\pgfpathlineto{\pgfpoint{-13.590908\du}{6.577368\du}}
\pgfpathlineto{\pgfpoint{-13.586728\du}{6.574681\du}}
\pgfpathlineto{\pgfpoint{-13.583444\du}{6.571994\du}}
\pgfpathlineto{\pgfpoint{-13.580160\du}{6.569307\du}}
\pgfpathlineto{\pgfpoint{-13.578070\du}{6.565724\du}}
\pgfpathlineto{\pgfpoint{-13.575980\du}{6.563037\du}}
\pgfpathlineto{\pgfpoint{-13.573890\du}{6.559753\du}}
\pgfpathlineto{\pgfpoint{-13.572696\du}{6.557066\du}}
\pgfpathlineto{\pgfpoint{-13.572099\du}{6.554379\du}}
\pgfpathlineto{\pgfpoint{-13.571800\du}{6.550797\du}}
\pgfusepath{stroke}
\pgfsetbuttcap
\pgfsetmiterjoin
\pgfsetdash{}{0pt}
\definecolor{dialinecolor}{rgb}{0.000000, 0.000000, 0.000000}
\pgfsetfillcolor{dialinecolor}
\pgfpathmoveto{\pgfpoint{-13.750036\du}{6.538855\du}}
\pgfpathlineto{\pgfpoint{-13.726152\du}{6.550797\du}}
\pgfpathlineto{\pgfpoint{-13.654798\du}{6.527211\du}}
\pgfpathlineto{\pgfpoint{-13.618972\du}{6.538855\du}}
\pgfpathlineto{\pgfpoint{-13.642557\du}{6.514971\du}}
\pgfpathlineto{\pgfpoint{-13.726152\du}{6.514971\du}}
\pgfpathlineto{\pgfpoint{-13.690326\du}{6.514971\du}}
\pgfpathlineto{\pgfpoint{-13.750036\du}{6.538855\du}}
\pgfusepath{fill}
\pgfsetbuttcap
\pgfsetmiterjoin
\pgfsetdash{}{0pt}
\definecolor{dialinecolor}{rgb}{0.000000, 0.000000, 0.000000}
\pgfsetfillcolor{dialinecolor}
\pgfpathmoveto{\pgfpoint{-13.761978\du}{6.574681\du}}
\pgfpathlineto{\pgfpoint{-13.797506\du}{6.562739\du}}
\pgfpathlineto{\pgfpoint{-13.857216\du}{6.586623\du}}
\pgfpathlineto{\pgfpoint{-13.892744\du}{6.574681\du}}
\pgfpathlineto{\pgfpoint{-13.880503\du}{6.610507\du}}
\pgfpathlineto{\pgfpoint{-13.797506\du}{6.610507\du}}
\pgfpathlineto{\pgfpoint{-13.833033\du}{6.598864\du}}
\pgfpathlineto{\pgfpoint{-13.761978\du}{6.574681\du}}
\pgfusepath{fill}
\pgfsetbuttcap
\pgfsetmiterjoin
\pgfsetdash{}{0pt}
\definecolor{dialinecolor}{rgb}{0.000000, 0.000000, 0.000000}
\pgfsetfillcolor{dialinecolor}
\pgfpathmoveto{\pgfpoint{-13.880503\du}{6.514971\du}}
\pgfpathlineto{\pgfpoint{-13.857216\du}{6.503327\du}}
\pgfpathlineto{\pgfpoint{-13.785862\du}{6.538855\du}}
\pgfpathlineto{\pgfpoint{-13.761978\du}{6.527211\du}}
\pgfpathlineto{\pgfpoint{-13.773323\du}{6.550797\du}}
\pgfpathlineto{\pgfpoint{-13.857216\du}{6.550797\du}}
\pgfpathlineto{\pgfpoint{-13.821390\du}{6.538855\du}}
\pgfpathlineto{\pgfpoint{-13.880503\du}{6.514971\du}}
\pgfusepath{fill}
\pgfsetbuttcap
\pgfsetmiterjoin
\pgfsetdash{}{0pt}
\definecolor{dialinecolor}{rgb}{0.000000, 0.000000, 0.000000}
\pgfsetfillcolor{dialinecolor}
\pgfpathmoveto{\pgfpoint{-13.630914\du}{6.598864\du}}
\pgfpathlineto{\pgfpoint{-13.654798\du}{6.610507\du}}
\pgfpathlineto{\pgfpoint{-13.726152\du}{6.574681\du}}
\pgfpathlineto{\pgfpoint{-13.761978\du}{6.586623\du}}
\pgfpathlineto{\pgfpoint{-13.737795\du}{6.562739\du}}
\pgfpathlineto{\pgfpoint{-13.654798\du}{6.562739\du}}
\pgfpathlineto{\pgfpoint{-13.690326\du}{6.574681\du}}
\pgfpathlineto{\pgfpoint{-13.630914\du}{6.598864\du}}
\pgfusepath{fill}
\pgfsetbuttcap
\pgfsetmiterjoin
\pgfsetdash{}{0pt}
\definecolor{dialinecolor}{rgb}{1.000000, 1.000000, 1.000000}
\pgfsetfillcolor{dialinecolor}
\pgfpathmoveto{\pgfpoint{-13.750036\du}{6.550797\du}}
\pgfpathlineto{\pgfpoint{-13.726152\du}{6.550797\du}}
\pgfpathlineto{\pgfpoint{-13.654798\du}{6.527211\du}}
\pgfpathlineto{\pgfpoint{-13.618972\du}{6.538855\du}}
\pgfpathlineto{\pgfpoint{-13.642557\du}{6.514971\du}}
\pgfpathlineto{\pgfpoint{-13.713911\du}{6.514971\du}}
\pgfpathlineto{\pgfpoint{-13.690326\du}{6.514971\du}}
\pgfpathlineto{\pgfpoint{-13.750036\du}{6.550797\du}}
\pgfusepath{fill}
\pgfsetbuttcap
\pgfsetmiterjoin
\pgfsetdash{}{0pt}
\definecolor{dialinecolor}{rgb}{1.000000, 1.000000, 1.000000}
\pgfsetfillcolor{dialinecolor}
\pgfpathmoveto{\pgfpoint{-13.761978\du}{6.574681\du}}
\pgfpathlineto{\pgfpoint{-13.785862\du}{6.562739\du}}
\pgfpathlineto{\pgfpoint{-13.857216\du}{6.586623\du}}
\pgfpathlineto{\pgfpoint{-13.892744\du}{6.586623\du}}
\pgfpathlineto{\pgfpoint{-13.868860\du}{6.610507\du}}
\pgfpathlineto{\pgfpoint{-13.785862\du}{6.610507\du}}
\pgfpathlineto{\pgfpoint{-13.821390\du}{6.598864\du}}
\pgfpathlineto{\pgfpoint{-13.761978\du}{6.574681\du}}
\pgfusepath{fill}
\pgfsetbuttcap
\pgfsetmiterjoin
\pgfsetdash{}{0pt}
\definecolor{dialinecolor}{rgb}{1.000000, 1.000000, 1.000000}
\pgfsetfillcolor{dialinecolor}
\pgfpathmoveto{\pgfpoint{-13.880503\du}{6.514971\du}}
\pgfpathlineto{\pgfpoint{-13.857216\du}{6.514971\du}}
\pgfpathlineto{\pgfpoint{-13.785862\du}{6.538855\du}}
\pgfpathlineto{\pgfpoint{-13.750036\du}{6.527211\du}}
\pgfpathlineto{\pgfpoint{-13.773323\du}{6.550797\du}}
\pgfpathlineto{\pgfpoint{-13.857216\du}{6.550797\du}}
\pgfpathlineto{\pgfpoint{-13.821390\du}{6.550797\du}}
\pgfpathlineto{\pgfpoint{-13.880503\du}{6.514971\du}}
\pgfusepath{fill}
\pgfsetbuttcap
\pgfsetmiterjoin
\pgfsetdash{}{0pt}
\definecolor{dialinecolor}{rgb}{1.000000, 1.000000, 1.000000}
\pgfsetfillcolor{dialinecolor}
\pgfpathmoveto{\pgfpoint{-13.630914\du}{6.598864\du}}
\pgfpathlineto{\pgfpoint{-13.654798\du}{6.610507\du}}
\pgfpathlineto{\pgfpoint{-13.726152\du}{6.586623\du}}
\pgfpathlineto{\pgfpoint{-13.750036\du}{6.586623\du}}
\pgfpathlineto{\pgfpoint{-13.737795\du}{6.562739\du}}
\pgfpathlineto{\pgfpoint{-13.654798\du}{6.562739\du}}
\pgfpathlineto{\pgfpoint{-13.690326\du}{6.574681\du}}
\pgfpathlineto{\pgfpoint{-13.630914\du}{6.598864\du}}
\pgfusepath{fill}
\pgfsetbuttcap
\pgfsetmiterjoin
\pgfsetdash{}{0pt}
\definecolor{dialinecolor}{rgb}{0.666667, 0.901961, 1.000000}
\pgfsetstrokecolor{dialinecolor}
\pgfpathmoveto{\pgfpoint{-13.952454\du}{6.550797\du}}
\pgfpathlineto{\pgfpoint{-13.952454\du}{6.645736\du}}
\pgfusepath{stroke}
\pgfsetbuttcap
\pgfsetmiterjoin
\pgfsetdash{}{0pt}
\definecolor{dialinecolor}{rgb}{0.666667, 0.901961, 1.000000}
\pgfsetstrokecolor{dialinecolor}
\pgfpathmoveto{\pgfpoint{-13.571800\du}{6.550797\du}}
\pgfpathlineto{\pgfpoint{-13.571800\du}{6.645736\du}}
\pgfusepath{stroke}
\pgfsetlinewidth{0.100000\du}
\pgfsetdash{}{0pt}
\pgfsetdash{}{0pt}
\pgfsetbuttcap
\pgfsetmiterjoin
\pgfsetlinewidth{0.001000\du}
\pgfsetbuttcap
\pgfsetmiterjoin
\pgfsetdash{}{0pt}
\definecolor{dialinecolor}{rgb}{0.000000, 0.470588, 0.666667}
\pgfsetstrokecolor{dialinecolor}
\pgfpathmoveto{\pgfpoint{-2.140683\du}{7.493354\du}}
\pgfpathlineto{\pgfpoint{-2.134115\du}{7.484397\du}}
\pgfpathlineto{\pgfpoint{-2.127845\du}{7.474844\du}}
\pgfpathlineto{\pgfpoint{-2.122471\du}{7.464693\du}}
\pgfpathlineto{\pgfpoint{-2.117396\du}{7.455438\du}}
\pgfpathlineto{\pgfpoint{-2.113515\du}{7.445287\du}}
\pgfpathlineto{\pgfpoint{-2.110231\du}{7.435136\du}}
\pgfpathlineto{\pgfpoint{-2.107245\du}{7.424986\du}}
\pgfpathlineto{\pgfpoint{-2.105454\du}{7.414536\du}}
\pgfpathlineto{\pgfpoint{-2.104260\du}{7.404385\du}}
\pgfpathlineto{\pgfpoint{-2.103364\du}{7.393936\du}}
\pgfpathlineto{\pgfpoint{-2.103364\du}{7.383487\du}}
\pgfpathlineto{\pgfpoint{-2.104260\du}{7.372739\du}}
\pgfpathlineto{\pgfpoint{-2.105454\du}{7.362588\du}}
\pgfpathlineto{\pgfpoint{-2.107245\du}{7.352139\du}}
\pgfpathlineto{\pgfpoint{-2.109932\du}{7.342287\du}}
\pgfpathlineto{\pgfpoint{-2.113515\du}{7.331837\du}}
\pgfpathlineto{\pgfpoint{-2.117396\du}{7.321985\du}}
\pgfpathlineto{\pgfpoint{-2.122471\du}{7.311834\du}}
\pgfpathlineto{\pgfpoint{-2.127845\du}{7.302281\du}}
\pgfpathlineto{\pgfpoint{-2.134115\du}{7.293324\du}}
\pgfpathlineto{\pgfpoint{-2.140683\du}{7.283770\du}}
\pgfpathlineto{\pgfpoint{-2.148147\du}{7.275112\du}}
\pgfpathlineto{\pgfpoint{-2.155611\du}{7.266454\du}}
\pgfpathlineto{\pgfpoint{-2.163970\du}{7.258095\du}}
\pgfpathlineto{\pgfpoint{-2.172628\du}{7.250034\du}}
\pgfpathlineto{\pgfpoint{-2.181883\du}{7.242272\du}}
\pgfpathlineto{\pgfpoint{-2.191437\du}{7.235106\du}}
\pgfpathlineto{\pgfpoint{-2.202185\du}{7.228538\du}}
\pgfpathlineto{\pgfpoint{-2.212634\du}{7.221970\du}}
\pgfpathlineto{\pgfpoint{-2.223979\du}{7.215701\du}}
\pgfpathlineto{\pgfpoint{-2.235921\du}{7.210028\du}}
\pgfpathlineto{\pgfpoint{-2.247863\du}{7.204654\du}}
\pgfpathlineto{\pgfpoint{-2.260104\du}{7.200474\du}}
\pgfpathlineto{\pgfpoint{-2.272643\du}{7.196295\du}}
\pgfpathlineto{\pgfpoint{-2.285779\du}{7.192115\du}}
\pgfpathlineto{\pgfpoint{-2.298617\du}{7.189129\du}}
\pgfpathlineto{\pgfpoint{-2.311753\du}{7.186442\du}}
\pgfpathlineto{\pgfpoint{-2.325487\du}{7.184054\du}}
\pgfpathlineto{\pgfpoint{-2.338623\du}{7.182561\du}}
\pgfpathlineto{\pgfpoint{-2.352655\du}{7.181367\du}}
\pgfpathlineto{\pgfpoint{-2.366388\du}{7.180471\du}}
\pgfpathlineto{\pgfpoint{-2.380420\du}{7.180471\du}}
\pgfpathlineto{\pgfpoint{-2.393855\du}{7.181069\du}}
\pgfpathlineto{\pgfpoint{-2.407589\du}{7.181666\du}}
\pgfpathlineto{\pgfpoint{-2.421621\du}{7.183158\du}}
\pgfpathlineto{\pgfpoint{-2.434757\du}{7.185547\du}}
\pgfpathlineto{\pgfpoint{-2.448490\du}{7.187935\du}}
\pgfpathlineto{\pgfpoint{-2.461029\du}{7.190921\du}}
\pgfpathlineto{\pgfpoint{-2.474166\du}{7.193906\du}}
\pgfpathlineto{\pgfpoint{-2.486705\du}{7.198086\du}}
\pgfpathlineto{\pgfpoint{-2.499543\du}{7.202863\du}}
\pgfpathlineto{\pgfpoint{-2.511783\du}{7.207341\du}}
\pgfpathlineto{\pgfpoint{-2.523128\du}{7.212715\du}}
\pgfpathlineto{\pgfpoint{-2.534772\du}{7.218985\du}}
\pgfpathlineto{\pgfpoint{-2.545818\du}{7.224956\du}}
\pgfpathlineto{\pgfpoint{-2.556268\du}{7.231822\du}}
\pgfpathlineto{\pgfpoint{-2.566418\du}{7.238391\du}}
\pgfpathlineto{\pgfpoint{-2.575972\du}{7.246153\du}}
\pgfpathlineto{\pgfpoint{-2.585526\du}{7.254214\du}}
\pgfpathlineto{\pgfpoint{-2.593885\du}{7.262275\du}}
\pgfpathlineto{\pgfpoint{-2.601946\du}{7.270634\du}}
\pgfpathlineto{\pgfpoint{-2.609708\du}{7.279591\du}}
\pgfusepath{stroke}
\pgfsetlinewidth{0.080000\du}
\pgfsetbuttcap
\pgfsetmiterjoin
\pgfsetdash{}{0pt}
\definecolor{dialinecolor}{rgb}{0.666667, 0.901961, 1.000000}
\pgfsetstrokecolor{dialinecolor}
\pgfpathmoveto{\pgfpoint{-2.149938\du}{7.485890\du}}
\pgfpathlineto{\pgfpoint{-2.144564\du}{7.476934\du}}
\pgfpathlineto{\pgfpoint{-2.138295\du}{7.468276\du}}
\pgfpathlineto{\pgfpoint{-2.133219\du}{7.459020\du}}
\pgfpathlineto{\pgfpoint{-2.129040\du}{7.449765\du}}
\pgfpathlineto{\pgfpoint{-2.125158\du}{7.440212\du}}
\pgfpathlineto{\pgfpoint{-2.122173\du}{7.430658\du}}
\pgfpathlineto{\pgfpoint{-2.119486\du}{7.420806\du}}
\pgfpathlineto{\pgfpoint{-2.117396\du}{7.411551\du}}
\pgfpathlineto{\pgfpoint{-2.116500\du}{7.401698\du}}
\pgfpathlineto{\pgfpoint{-2.115903\du}{7.391548\du}}
\pgfpathlineto{\pgfpoint{-2.115903\du}{7.381695\du}}
\pgfpathlineto{\pgfpoint{-2.116799\du}{7.371843\du}}
\pgfpathlineto{\pgfpoint{-2.118292\du}{7.361693\du}}
\pgfpathlineto{\pgfpoint{-2.120083\du}{7.352437\du}}
\pgfpathlineto{\pgfpoint{-2.123069\du}{7.342585\du}}
\pgfpathlineto{\pgfpoint{-2.126651\du}{7.333032\du}}
\pgfpathlineto{\pgfpoint{-2.130532\du}{7.323776\du}}
\pgfpathlineto{\pgfpoint{-2.135011\du}{7.314521\du}}
\pgfpathlineto{\pgfpoint{-2.140086\du}{7.305863\du}}
\pgfpathlineto{\pgfpoint{-2.145758\du}{7.296608\du}}
\pgfpathlineto{\pgfpoint{-2.152924\du}{7.288249\du}}
\pgfpathlineto{\pgfpoint{-2.159492\du}{7.279591\du}}
\pgfpathlineto{\pgfpoint{-2.167254\du}{7.271530\du}}
\pgfpathlineto{\pgfpoint{-2.175315\du}{7.263767\du}}
\pgfpathlineto{\pgfpoint{-2.183675\du}{7.256901\du}}
\pgfpathlineto{\pgfpoint{-2.192333\du}{7.249437\du}}
\pgfpathlineto{\pgfpoint{-2.202185\du}{7.242570\du}}
\pgfpathlineto{\pgfpoint{-2.212634\du}{7.235704\du}}
\pgfpathlineto{\pgfpoint{-2.222785\du}{7.230031\du}}
\pgfpathlineto{\pgfpoint{-2.233234\du}{7.224657\du}}
\pgfpathlineto{\pgfpoint{-2.245176\du}{7.219283\du}}
\pgfpathlineto{\pgfpoint{-2.255924\du}{7.214208\du}}
\pgfpathlineto{\pgfpoint{-2.268463\du}{7.210028\du}}
\pgfpathlineto{\pgfpoint{-2.279808\du}{7.206147\du}}
\pgfpathlineto{\pgfpoint{-2.292348\du}{7.202863\du}}
\pgfpathlineto{\pgfpoint{-2.304887\du}{7.199877\du}}
\pgfpathlineto{\pgfpoint{-2.317426\du}{7.197190\du}}
\pgfpathlineto{\pgfpoint{-2.330562\du}{7.195399\du}}
\pgfpathlineto{\pgfpoint{-2.343101\du}{7.193906\du}}
\pgfpathlineto{\pgfpoint{-2.356238\du}{7.192712\du}}
\pgfpathlineto{\pgfpoint{-2.370270\du}{7.192115\du}}
\pgfpathlineto{\pgfpoint{-2.383406\du}{7.192115\du}}
\pgfpathlineto{\pgfpoint{-2.396244\du}{7.192712\du}}
\pgfpathlineto{\pgfpoint{-2.409380\du}{7.193608\du}}
\pgfpathlineto{\pgfpoint{-2.421919\du}{7.195399\du}}
\pgfpathlineto{\pgfpoint{-2.435055\du}{7.197190\du}}
\pgfpathlineto{\pgfpoint{-2.447595\du}{7.199280\du}}
\pgfpathlineto{\pgfpoint{-2.460432\du}{7.202863\du}}
\pgfpathlineto{\pgfpoint{-2.472673\du}{7.205848\du}}
\pgfpathlineto{\pgfpoint{-2.484615\du}{7.209730\du}}
\pgfpathlineto{\pgfpoint{-2.496856\du}{7.213909\du}}
\pgfpathlineto{\pgfpoint{-2.508201\du}{7.218985\du}}
\pgfpathlineto{\pgfpoint{-2.519247\du}{7.223761\du}}
\pgfpathlineto{\pgfpoint{-2.529995\du}{7.229733\du}}
\pgfpathlineto{\pgfpoint{-2.540444\du}{7.235405\du}}
\pgfpathlineto{\pgfpoint{-2.550894\du}{7.241973\du}}
\pgfpathlineto{\pgfpoint{-2.560447\du}{7.248541\du}}
\pgfpathlineto{\pgfpoint{-2.569105\du}{7.255707\du}}
\pgfpathlineto{\pgfpoint{-2.578659\du}{7.262872\du}}
\pgfpathlineto{\pgfpoint{-2.586421\du}{7.270634\du}}
\pgfpathlineto{\pgfpoint{-2.593885\du}{7.278695\du}}
\pgfpathlineto{\pgfpoint{-2.601050\du}{7.287055\du}}
\pgfusepath{stroke}
\pgfsetbuttcap
\pgfsetmiterjoin
\pgfsetdash{}{0pt}
\definecolor{dialinecolor}{rgb}{0.000000, 0.705882, 1.000000}
\pgfsetfillcolor{dialinecolor}
\pgfpathmoveto{\pgfpoint{-4.066941\du}{7.168828\du}}
\pgfpathlineto{\pgfpoint{-3.864224\du}{6.978053\du}}
\pgfpathlineto{\pgfpoint{-2.352655\du}{6.978053\du}}
\pgfpathlineto{\pgfpoint{-2.555372\du}{7.168828\du}}
\pgfpathlineto{\pgfpoint{-4.066941\du}{7.168828\du}}
\pgfusepath{fill}
\pgfsetlinewidth{0.001000\du}
\pgfsetbuttcap
\pgfsetmiterjoin
\pgfsetdash{}{0pt}
\definecolor{dialinecolor}{rgb}{0.666667, 0.901961, 1.000000}
\pgfsetstrokecolor{dialinecolor}
\pgfpathmoveto{\pgfpoint{-4.066941\du}{7.168828\du}}
\pgfpathlineto{\pgfpoint{-3.864224\du}{6.978053\du}}
\pgfpathlineto{\pgfpoint{-2.352655\du}{6.978053\du}}
\pgfpathlineto{\pgfpoint{-2.555372\du}{7.168828\du}}
\pgfpathlineto{\pgfpoint{-4.066941\du}{7.168828\du}}
\pgfusepath{stroke}
\pgfsetbuttcap
\pgfsetmiterjoin
\pgfsetdash{}{0pt}
\definecolor{dialinecolor}{rgb}{0.000000, 0.588235, 0.831373}
\pgfsetfillcolor{dialinecolor}
\pgfpathmoveto{\pgfpoint{-4.066941\du}{7.168828\du}}
\pgfpathlineto{\pgfpoint{-2.555372\du}{7.168828\du}}
\pgfpathlineto{\pgfpoint{-2.555372\du}{7.454841\du}}
\pgfpathlineto{\pgfpoint{-4.066941\du}{7.454841\du}}
\pgfpathlineto{\pgfpoint{-4.066941\du}{7.168828\du}}
\pgfusepath{fill}
\pgfsetbuttcap
\pgfsetmiterjoin
\pgfsetdash{}{0pt}
\definecolor{dialinecolor}{rgb}{0.666667, 0.901961, 1.000000}
\pgfsetstrokecolor{dialinecolor}
\pgfpathmoveto{\pgfpoint{-4.066941\du}{7.168828\du}}
\pgfpathlineto{\pgfpoint{-2.555372\du}{7.168828\du}}
\pgfpathlineto{\pgfpoint{-2.555372\du}{7.454841\du}}
\pgfpathlineto{\pgfpoint{-4.066941\du}{7.454841\du}}
\pgfpathlineto{\pgfpoint{-4.066941\du}{7.168828\du}}
\pgfusepath{stroke}
\pgfsetbuttcap
\pgfsetmiterjoin
\pgfsetdash{}{0pt}
\definecolor{dialinecolor}{rgb}{0.000000, 0.352941, 0.501961}
\pgfsetfillcolor{dialinecolor}
\pgfpathmoveto{\pgfpoint{-2.555372\du}{7.454841\du}}
\pgfpathlineto{\pgfpoint{-2.352655\du}{7.252124\du}}
\pgfpathlineto{\pgfpoint{-2.352655\du}{6.978053\du}}
\pgfpathlineto{\pgfpoint{-2.555372\du}{7.168828\du}}
\pgfpathlineto{\pgfpoint{-2.555372\du}{7.454841\du}}
\pgfusepath{fill}
\pgfsetbuttcap
\pgfsetmiterjoin
\pgfsetdash{}{0pt}
\definecolor{dialinecolor}{rgb}{0.666667, 0.901961, 1.000000}
\pgfsetstrokecolor{dialinecolor}
\pgfpathmoveto{\pgfpoint{-2.555372\du}{7.454841\du}}
\pgfpathlineto{\pgfpoint{-2.352655\du}{7.252124\du}}
\pgfpathlineto{\pgfpoint{-2.352655\du}{6.978053\du}}
\pgfpathlineto{\pgfpoint{-2.555372\du}{7.168828\du}}
\pgfpathlineto{\pgfpoint{-2.555372\du}{7.454841\du}}
\pgfusepath{stroke}
\pgfsetbuttcap
\pgfsetmiterjoin
\pgfsetdash{}{0pt}
\definecolor{dialinecolor}{rgb}{0.000000, 0.000000, 0.000000}
\pgfsetfillcolor{dialinecolor}
\pgfpathmoveto{\pgfpoint{-4.018874\du}{7.132703\du}}
\pgfpathlineto{\pgfpoint{-3.864224\du}{6.978053\du}}
\pgfpathlineto{\pgfpoint{-2.388780\du}{6.978053\du}}
\pgfpathlineto{\pgfpoint{-2.543131\du}{7.132703\du}}
\pgfpathlineto{\pgfpoint{-4.018874\du}{7.132703\du}}
\pgfusepath{fill}
\pgfsetbuttcap
\pgfsetmiterjoin
\pgfsetdash{}{0pt}
\definecolor{dialinecolor}{rgb}{0.000000, 0.000000, 0.000000}
\pgfsetstrokecolor{dialinecolor}
\pgfpathmoveto{\pgfpoint{-4.018874\du}{7.132703\du}}
\pgfpathlineto{\pgfpoint{-3.864224\du}{6.978053\du}}
\pgfpathlineto{\pgfpoint{-2.388780\du}{6.978053\du}}
\pgfpathlineto{\pgfpoint{-2.543131\du}{7.132703\du}}
\pgfpathlineto{\pgfpoint{-4.018874\du}{7.132703\du}}
\pgfusepath{stroke}
\pgfsetbuttcap
\pgfsetmiterjoin
\pgfsetdash{}{0pt}
\definecolor{dialinecolor}{rgb}{0.000000, 0.705882, 1.000000}
\pgfsetfillcolor{dialinecolor}
\pgfpathmoveto{\pgfpoint{-4.066941\du}{5.906550\du}}
\pgfpathlineto{\pgfpoint{-3.899752\du}{5.751900\du}}
\pgfpathlineto{\pgfpoint{-2.388780\du}{5.751900\du}}
\pgfpathlineto{\pgfpoint{-2.555372\du}{5.906550\du}}
\pgfpathlineto{\pgfpoint{-4.066941\du}{5.906550\du}}
\pgfusepath{fill}
\pgfsetbuttcap
\pgfsetmiterjoin
\pgfsetdash{}{0pt}
\definecolor{dialinecolor}{rgb}{0.666667, 0.901961, 1.000000}
\pgfsetstrokecolor{dialinecolor}
\pgfpathmoveto{\pgfpoint{-4.066941\du}{5.906550\du}}
\pgfpathlineto{\pgfpoint{-3.899752\du}{5.751900\du}}
\pgfpathlineto{\pgfpoint{-2.388780\du}{5.751900\du}}
\pgfpathlineto{\pgfpoint{-2.555372\du}{5.906550\du}}
\pgfpathlineto{\pgfpoint{-4.066941\du}{5.906550\du}}
\pgfusepath{stroke}
\pgfsetbuttcap
\pgfsetmiterjoin
\pgfsetdash{}{0pt}
\definecolor{dialinecolor}{rgb}{0.000000, 0.588235, 0.831373}
\pgfsetfillcolor{dialinecolor}
\pgfpathmoveto{\pgfpoint{-4.066941\du}{5.906550\du}}
\pgfpathlineto{\pgfpoint{-2.543131\du}{5.906550\du}}
\pgfpathlineto{\pgfpoint{-2.543131\du}{7.108819\du}}
\pgfpathlineto{\pgfpoint{-4.066941\du}{7.108819\du}}
\pgfpathlineto{\pgfpoint{-4.066941\du}{5.906550\du}}
\pgfusepath{fill}
\pgfsetbuttcap
\pgfsetmiterjoin
\pgfsetdash{}{0pt}
\definecolor{dialinecolor}{rgb}{0.666667, 0.901961, 1.000000}
\pgfsetstrokecolor{dialinecolor}
\pgfpathmoveto{\pgfpoint{-4.066941\du}{5.906550\du}}
\pgfpathlineto{\pgfpoint{-2.543430\du}{5.906550\du}}
\pgfpathlineto{\pgfpoint{-2.543430\du}{7.108520\du}}
\pgfpathlineto{\pgfpoint{-4.066941\du}{7.108520\du}}
\pgfpathlineto{\pgfpoint{-4.066941\du}{5.906550\du}}
\pgfusepath{stroke}
\pgfsetbuttcap
\pgfsetmiterjoin
\pgfsetdash{}{0pt}
\definecolor{dialinecolor}{rgb}{1.000000, 1.000000, 1.000000}
\pgfsetfillcolor{dialinecolor}
\pgfpathmoveto{\pgfpoint{-3.935876\du}{6.061498\du}}
\pgfpathlineto{\pgfpoint{-2.674196\du}{6.061498\du}}
\pgfpathlineto{\pgfpoint{-2.674196\du}{6.990294\du}}
\pgfpathlineto{\pgfpoint{-3.935876\du}{6.990294\du}}
\pgfpathlineto{\pgfpoint{-3.935876\du}{6.061498\du}}
\pgfusepath{fill}
\pgfsetbuttcap
\pgfsetmiterjoin
\pgfsetdash{}{0pt}
\definecolor{dialinecolor}{rgb}{0.000000, 0.235294, 0.333333}
\pgfsetstrokecolor{dialinecolor}
\pgfpathmoveto{\pgfpoint{-3.935876\du}{6.061498\du}}
\pgfpathlineto{\pgfpoint{-2.674494\du}{6.061498\du}}
\pgfpathlineto{\pgfpoint{-2.674494\du}{6.989995\du}}
\pgfpathlineto{\pgfpoint{-3.935876\du}{6.989995\du}}
\pgfpathlineto{\pgfpoint{-3.935876\du}{6.061498\du}}
\pgfusepath{stroke}
\pgfsetbuttcap
\pgfsetmiterjoin
\pgfsetdash{}{0pt}
\definecolor{dialinecolor}{rgb}{0.000000, 0.352941, 0.501961}
\pgfsetfillcolor{dialinecolor}
\pgfpathmoveto{\pgfpoint{-2.555372\du}{7.097474\du}}
\pgfpathlineto{\pgfpoint{-2.388780\du}{6.930583\du}}
\pgfpathlineto{\pgfpoint{-2.388780\du}{5.751900\du}}
\pgfpathlineto{\pgfpoint{-2.555372\du}{5.906550\du}}
\pgfpathlineto{\pgfpoint{-2.555372\du}{7.097474\du}}
\pgfusepath{fill}
\pgfsetbuttcap
\pgfsetmiterjoin
\pgfsetdash{}{0pt}
\definecolor{dialinecolor}{rgb}{0.666667, 0.901961, 1.000000}
\pgfsetstrokecolor{dialinecolor}
\pgfpathmoveto{\pgfpoint{-2.555372\du}{7.097474\du}}
\pgfpathlineto{\pgfpoint{-2.388780\du}{6.930583\du}}
\pgfpathlineto{\pgfpoint{-2.388780\du}{5.751900\du}}
\pgfpathlineto{\pgfpoint{-2.555372\du}{5.906550\du}}
\pgfpathlineto{\pgfpoint{-2.555372\du}{7.097474\du}}
\pgfusepath{stroke}
\pgfsetbuttcap
\pgfsetmiterjoin
\pgfsetdash{}{0pt}
\definecolor{dialinecolor}{rgb}{0.000000, 0.705882, 1.000000}
\pgfsetfillcolor{dialinecolor}
\pgfpathmoveto{\pgfpoint{-4.364000\du}{7.692787\du}}
\pgfpathlineto{\pgfpoint{-4.126353\du}{7.394832\du}}
\pgfpathlineto{\pgfpoint{-2.483719\du}{7.394832\du}}
\pgfpathlineto{\pgfpoint{-2.721665\du}{7.692787\du}}
\pgfpathlineto{\pgfpoint{-4.364000\du}{7.692787\du}}
\pgfusepath{fill}
\pgfsetbuttcap
\pgfsetmiterjoin
\pgfsetdash{}{0pt}
\definecolor{dialinecolor}{rgb}{0.666667, 0.901961, 1.000000}
\pgfsetstrokecolor{dialinecolor}
\pgfpathmoveto{\pgfpoint{-4.364000\du}{7.692787\du}}
\pgfpathlineto{\pgfpoint{-4.126353\du}{7.394832\du}}
\pgfpathlineto{\pgfpoint{-2.483719\du}{7.394832\du}}
\pgfpathlineto{\pgfpoint{-2.721665\du}{7.692787\du}}
\pgfpathlineto{\pgfpoint{-4.364000\du}{7.692787\du}}
\pgfusepath{stroke}
\pgfsetbuttcap
\pgfsetmiterjoin
\pgfsetdash{}{0pt}
\definecolor{dialinecolor}{rgb}{0.000000, 0.352941, 0.501961}
\pgfsetfillcolor{dialinecolor}
\pgfpathmoveto{\pgfpoint{-2.721665\du}{7.751900\du}}
\pgfpathlineto{\pgfpoint{-2.483719\du}{7.502012\du}}
\pgfpathlineto{\pgfpoint{-2.483719\du}{7.394832\du}}
\pgfpathlineto{\pgfpoint{-2.721665\du}{7.692787\du}}
\pgfpathlineto{\pgfpoint{-2.721665\du}{7.751900\du}}
\pgfusepath{fill}
\pgfsetbuttcap
\pgfsetmiterjoin
\pgfsetdash{}{0pt}
\definecolor{dialinecolor}{rgb}{0.666667, 0.901961, 1.000000}
\pgfsetstrokecolor{dialinecolor}
\pgfpathmoveto{\pgfpoint{-2.721665\du}{7.751900\du}}
\pgfpathlineto{\pgfpoint{-2.483719\du}{7.502012\du}}
\pgfpathlineto{\pgfpoint{-2.483719\du}{7.394832\du}}
\pgfpathlineto{\pgfpoint{-2.721665\du}{7.692787\du}}
\pgfpathlineto{\pgfpoint{-2.721665\du}{7.751900\du}}
\pgfusepath{stroke}
\pgfsetbuttcap
\pgfsetmiterjoin
\pgfsetdash{}{0pt}
\definecolor{dialinecolor}{rgb}{0.000000, 0.588235, 0.831373}
\pgfsetfillcolor{dialinecolor}
\pgfpathmoveto{\pgfpoint{-4.364000\du}{7.692787\du}}
\pgfpathlineto{\pgfpoint{-2.721665\du}{7.692787\du}}
\pgfpathlineto{\pgfpoint{-2.721665\du}{7.751900\du}}
\pgfpathlineto{\pgfpoint{-4.364000\du}{7.751900\du}}
\pgfpathlineto{\pgfpoint{-4.364000\du}{7.692787\du}}
\pgfusepath{fill}
\pgfsetbuttcap
\pgfsetmiterjoin
\pgfsetdash{}{0pt}
\definecolor{dialinecolor}{rgb}{0.666667, 0.901961, 1.000000}
\pgfsetstrokecolor{dialinecolor}
\pgfpathmoveto{\pgfpoint{-4.364000\du}{7.692787\du}}
\pgfpathlineto{\pgfpoint{-2.721665\du}{7.692787\du}}
\pgfpathlineto{\pgfpoint{-2.721665\du}{7.751900\du}}
\pgfpathlineto{\pgfpoint{-4.364000\du}{7.751900\du}}
\pgfpathlineto{\pgfpoint{-4.364000\du}{7.692787\du}}
\pgfusepath{stroke}
\pgfsetbuttcap
\pgfsetmiterjoin
\pgfsetdash{}{0pt}
\definecolor{dialinecolor}{rgb}{0.000000, 0.705882, 1.000000}
\pgfsetfillcolor{dialinecolor}
\pgfpathmoveto{\pgfpoint{-2.388780\du}{7.656662\du}}
\pgfpathlineto{\pgfpoint{-2.269658\du}{7.478426\du}}
\pgfpathlineto{\pgfpoint{-2.054999\du}{7.478426\du}}
\pgfpathlineto{\pgfpoint{-2.173822\du}{7.656662\du}}
\pgfpathlineto{\pgfpoint{-2.388780\du}{7.656662\du}}
\pgfusepath{fill}
\pgfsetbuttcap
\pgfsetmiterjoin
\pgfsetdash{}{0pt}
\definecolor{dialinecolor}{rgb}{0.666667, 0.901961, 1.000000}
\pgfsetstrokecolor{dialinecolor}
\pgfpathmoveto{\pgfpoint{-2.388780\du}{7.656662\du}}
\pgfpathlineto{\pgfpoint{-2.269658\du}{7.478426\du}}
\pgfpathlineto{\pgfpoint{-2.054999\du}{7.478426\du}}
\pgfpathlineto{\pgfpoint{-2.173822\du}{7.656662\du}}
\pgfpathlineto{\pgfpoint{-2.388780\du}{7.656662\du}}
\pgfusepath{stroke}
\pgfsetbuttcap
\pgfsetmiterjoin
\pgfsetdash{}{0pt}
\definecolor{dialinecolor}{rgb}{0.000000, 0.352941, 0.501961}
\pgfsetfillcolor{dialinecolor}
\pgfpathmoveto{\pgfpoint{-2.173822\du}{7.716372\du}}
\pgfpathlineto{\pgfpoint{-2.054999\du}{7.585308\du}}
\pgfpathlineto{\pgfpoint{-2.054999\du}{7.478426\du}}
\pgfpathlineto{\pgfpoint{-2.173822\du}{7.656662\du}}
\pgfpathlineto{\pgfpoint{-2.173822\du}{7.716372\du}}
\pgfusepath{fill}
\pgfsetbuttcap
\pgfsetmiterjoin
\pgfsetdash{}{0pt}
\definecolor{dialinecolor}{rgb}{0.666667, 0.901961, 1.000000}
\pgfsetstrokecolor{dialinecolor}
\pgfpathmoveto{\pgfpoint{-2.173822\du}{7.716372\du}}
\pgfpathlineto{\pgfpoint{-2.054999\du}{7.585308\du}}
\pgfpathlineto{\pgfpoint{-2.054999\du}{7.478426\du}}
\pgfpathlineto{\pgfpoint{-2.173822\du}{7.656662\du}}
\pgfpathlineto{\pgfpoint{-2.173822\du}{7.716372\du}}
\pgfusepath{stroke}
\pgfsetbuttcap
\pgfsetmiterjoin
\pgfsetdash{}{0pt}
\definecolor{dialinecolor}{rgb}{0.000000, 0.588235, 0.831373}
\pgfsetfillcolor{dialinecolor}
\pgfpathmoveto{\pgfpoint{-2.388780\du}{7.656662\du}}
\pgfpathlineto{\pgfpoint{-2.173822\du}{7.656662\du}}
\pgfpathlineto{\pgfpoint{-2.173822\du}{7.716372\du}}
\pgfpathlineto{\pgfpoint{-2.388780\du}{7.716372\du}}
\pgfpathlineto{\pgfpoint{-2.388780\du}{7.656662\du}}
\pgfusepath{fill}
\pgfsetbuttcap
\pgfsetmiterjoin
\pgfsetdash{}{0pt}
\definecolor{dialinecolor}{rgb}{0.666667, 0.901961, 1.000000}
\pgfsetstrokecolor{dialinecolor}
\pgfpathmoveto{\pgfpoint{-2.388780\du}{7.656662\du}}
\pgfpathlineto{\pgfpoint{-2.173822\du}{7.656662\du}}
\pgfpathlineto{\pgfpoint{-2.173822\du}{7.716372\du}}
\pgfpathlineto{\pgfpoint{-2.388780\du}{7.716372\du}}
\pgfpathlineto{\pgfpoint{-2.388780\du}{7.656662\du}}
\pgfusepath{stroke}
\pgfsetbuttcap
\pgfsetmiterjoin
\pgfsetdash{}{0pt}
\definecolor{dialinecolor}{rgb}{0.811765, 0.054902, 0.188235}
\pgfsetstrokecolor{dialinecolor}
\pgfpathmoveto{\pgfpoint{-3.602692\du}{6.287801\du}}
\pgfpathlineto{\pgfpoint{-3.268911\du}{6.287801\du}}
\pgfpathlineto{\pgfpoint{-3.352506\du}{6.371097\du}}
\pgfpathlineto{\pgfpoint{-3.066791\du}{6.371097\du}}
\pgfusepath{stroke}
\pgfsetbuttcap
\pgfsetmiterjoin
\pgfsetdash{}{0pt}
\definecolor{dialinecolor}{rgb}{0.811765, 0.054902, 0.188235}
\pgfsetstrokecolor{dialinecolor}
\pgfpathmoveto{\pgfpoint{-3.638519\du}{6.287801\du}}
\pgfpathlineto{\pgfpoint{-3.399976\du}{6.525747\du}}
\pgfpathlineto{\pgfpoint{-3.519396\du}{6.525747\du}}
\pgfpathlineto{\pgfpoint{-3.316978\du}{6.728464\du}}
\pgfusepath{stroke}
\pgfsetbuttcap
\pgfsetmiterjoin
\pgfsetdash{}{0pt}
\definecolor{dialinecolor}{rgb}{0.811765, 0.054902, 0.188235}
\pgfsetstrokecolor{dialinecolor}
\pgfpathmoveto{\pgfpoint{-2.960208\du}{6.287801\du}}
\pgfpathlineto{\pgfpoint{-3.197557\du}{6.525747\du}}
\pgfpathlineto{\pgfpoint{-3.079032\du}{6.525747\du}}
\pgfpathlineto{\pgfpoint{-3.281152\du}{6.728464\du}}
\pgfusepath{stroke}
\pgfsetbuttcap
\pgfsetmiterjoin
\pgfsetdash{}{0pt}
\definecolor{dialinecolor}{rgb}{0.000000, 0.470588, 0.666667}
\pgfsetfillcolor{dialinecolor}
\pgfpathmoveto{\pgfpoint{-3.471927\du}{6.376769\du}}
\pgfpathlineto{\pgfpoint{-3.472524\du}{6.373485\du}}
\pgfpathlineto{\pgfpoint{-3.473419\du}{6.370201\du}}
\pgfpathlineto{\pgfpoint{-3.474315\du}{6.366917\du}}
\pgfpathlineto{\pgfpoint{-3.476106\du}{6.363633\du}}
\pgfpathlineto{\pgfpoint{-3.478495\du}{6.360648\du}}
\pgfpathlineto{\pgfpoint{-3.481182\du}{6.357364\du}}
\pgfpathlineto{\pgfpoint{-3.484167\du}{6.354079\du}}
\pgfpathlineto{\pgfpoint{-3.487750\du}{6.351094\du}}
\pgfpathlineto{\pgfpoint{-3.491930\du}{6.348108\du}}
\pgfpathlineto{\pgfpoint{-3.496706\du}{6.345123\du}}
\pgfpathlineto{\pgfpoint{-3.501782\du}{6.342436\du}}
\pgfpathlineto{\pgfpoint{-3.507454\du}{6.339450\du}}
\pgfpathlineto{\pgfpoint{-3.513127\du}{6.337062\du}}
\pgfpathlineto{\pgfpoint{-3.519396\du}{6.334375\du}}
\pgfpathlineto{\pgfpoint{-3.525964\du}{6.331987\du}}
\pgfpathlineto{\pgfpoint{-3.533130\du}{6.329598\du}}
\pgfpathlineto{\pgfpoint{-3.540594\du}{6.327508\du}}
\pgfpathlineto{\pgfpoint{-3.548057\du}{6.325120\du}}
\pgfpathlineto{\pgfpoint{-3.556417\du}{6.323627\du}}
\pgfpathlineto{\pgfpoint{-3.564776\du}{6.321537\du}}
\pgfpathlineto{\pgfpoint{-3.572837\du}{6.319746\du}}
\pgfpathlineto{\pgfpoint{-3.582092\du}{6.318552\du}}
\pgfpathlineto{\pgfpoint{-3.590750\du}{6.317059\du}}
\pgfpathlineto{\pgfpoint{-3.600304\du}{6.315865\du}}
\pgfpathlineto{\pgfpoint{-3.609559\du}{6.314671\du}}
\pgfpathlineto{\pgfpoint{-3.618814\du}{6.313775\du}}
\pgfpathlineto{\pgfpoint{-3.628666\du}{6.312879\du}}
\pgfpathlineto{\pgfpoint{-3.638519\du}{6.312282\du}}
\pgfpathlineto{\pgfpoint{-3.648371\du}{6.311984\du}}
\pgfpathlineto{\pgfpoint{-3.658522\du}{6.311685\du}}
\pgfpathlineto{\pgfpoint{-3.668374\du}{6.311685\du}}
\pgfpathlineto{\pgfpoint{-3.668374\du}{6.311685\du}}
\pgfpathlineto{\pgfpoint{-3.678226\du}{6.311685\du}}
\pgfpathlineto{\pgfpoint{-3.688078\du}{6.311984\du}}
\pgfpathlineto{\pgfpoint{-3.698229\du}{6.312282\du}}
\pgfpathlineto{\pgfpoint{-3.708081\du}{6.312879\du}}
\pgfpathlineto{\pgfpoint{-3.717635\du}{6.313775\du}}
\pgfpathlineto{\pgfpoint{-3.727189\du}{6.314671\du}}
\pgfpathlineto{\pgfpoint{-3.736444\du}{6.315865\du}}
\pgfpathlineto{\pgfpoint{-3.745699\du}{6.317059\du}}
\pgfpathlineto{\pgfpoint{-3.754655\du}{6.318552\du}}
\pgfpathlineto{\pgfpoint{-3.763612\du}{6.319746\du}}
\pgfpathlineto{\pgfpoint{-3.771971\du}{6.321537\du}}
\pgfpathlineto{\pgfpoint{-3.780331\du}{6.323627\du}}
\pgfpathlineto{\pgfpoint{-3.788690\du}{6.325120\du}}
\pgfpathlineto{\pgfpoint{-3.795856\du}{6.327508\du}}
\pgfpathlineto{\pgfpoint{-3.803319\du}{6.329598\du}}
\pgfpathlineto{\pgfpoint{-3.810485\du}{6.331987\du}}
\pgfpathlineto{\pgfpoint{-3.817351\du}{6.334375\du}}
\pgfpathlineto{\pgfpoint{-3.823322\du}{6.337062\du}}
\pgfpathlineto{\pgfpoint{-3.829293\du}{6.339450\du}}
\pgfpathlineto{\pgfpoint{-3.834966\du}{6.342436\du}}
\pgfpathlineto{\pgfpoint{-3.840041\du}{6.345123\du}}
\pgfpathlineto{\pgfpoint{-3.844818\du}{6.348108\du}}
\pgfpathlineto{\pgfpoint{-3.848998\du}{6.351094\du}}
\pgfpathlineto{\pgfpoint{-3.852282\du}{6.354079\du}}
\pgfpathlineto{\pgfpoint{-3.855566\du}{6.357364\du}}
\pgfpathlineto{\pgfpoint{-3.858551\du}{6.360648\du}}
\pgfpathlineto{\pgfpoint{-3.860641\du}{6.363633\du}}
\pgfpathlineto{\pgfpoint{-3.862134\du}{6.366917\du}}
\pgfpathlineto{\pgfpoint{-3.863925\du}{6.370201\du}}
\pgfpathlineto{\pgfpoint{-3.864224\du}{6.373485\du}}
\pgfpathlineto{\pgfpoint{-3.864522\du}{6.376769\du}}
\pgfpathlineto{\pgfpoint{-3.864522\du}{6.376769\du}}
\pgfpathlineto{\pgfpoint{-3.864224\du}{6.380352\du}}
\pgfpathlineto{\pgfpoint{-3.863925\du}{6.383338\du}}
\pgfpathlineto{\pgfpoint{-3.862134\du}{6.386622\du}}
\pgfpathlineto{\pgfpoint{-3.860641\du}{6.389906\du}}
\pgfpathlineto{\pgfpoint{-3.858551\du}{6.393190\du}}
\pgfpathlineto{\pgfpoint{-3.855566\du}{6.396474\du}}
\pgfpathlineto{\pgfpoint{-3.852282\du}{6.399459\du}}
\pgfpathlineto{\pgfpoint{-3.848998\du}{6.402445\du}}
\pgfpathlineto{\pgfpoint{-3.844818\du}{6.405430\du}}
\pgfpathlineto{\pgfpoint{-3.840041\du}{6.408714\du}}
\pgfpathlineto{\pgfpoint{-3.834966\du}{6.411401\du}}
\pgfpathlineto{\pgfpoint{-3.829293\du}{6.414088\du}}
\pgfpathlineto{\pgfpoint{-3.823322\du}{6.416775\du}}
\pgfpathlineto{\pgfpoint{-3.817351\du}{6.419462\du}}
\pgfpathlineto{\pgfpoint{-3.810485\du}{6.421851\du}}
\pgfpathlineto{\pgfpoint{-3.803319\du}{6.424239\du}}
\pgfpathlineto{\pgfpoint{-3.795856\du}{6.426030\du}}
\pgfpathlineto{\pgfpoint{-3.788690\du}{6.428120\du}}
\pgfpathlineto{\pgfpoint{-3.780331\du}{6.430210\du}}
\pgfpathlineto{\pgfpoint{-3.771971\du}{6.432300\du}}
\pgfpathlineto{\pgfpoint{-3.763612\du}{6.433793\du}}
\pgfpathlineto{\pgfpoint{-3.754655\du}{6.435286\du}}
\pgfpathlineto{\pgfpoint{-3.745699\du}{6.436778\du}}
\pgfpathlineto{\pgfpoint{-3.736444\du}{6.437973\du}}
\pgfpathlineto{\pgfpoint{-3.727189\du}{6.438868\du}}
\pgfpathlineto{\pgfpoint{-3.717635\du}{6.440062\du}}
\pgfpathlineto{\pgfpoint{-3.708081\du}{6.440660\du}}
\pgfpathlineto{\pgfpoint{-3.698229\du}{6.441257\du}}
\pgfpathlineto{\pgfpoint{-3.688078\du}{6.441854\du}}
\pgfpathlineto{\pgfpoint{-3.678226\du}{6.442152\du}}
\pgfpathlineto{\pgfpoint{-3.668374\du}{6.442152\du}}
\pgfpathlineto{\pgfpoint{-3.668374\du}{6.442152\du}}
\pgfpathlineto{\pgfpoint{-3.658522\du}{6.442152\du}}
\pgfpathlineto{\pgfpoint{-3.648371\du}{6.441854\du}}
\pgfpathlineto{\pgfpoint{-3.638519\du}{6.441257\du}}
\pgfpathlineto{\pgfpoint{-3.628666\du}{6.440660\du}}
\pgfpathlineto{\pgfpoint{-3.618814\du}{6.440062\du}}
\pgfpathlineto{\pgfpoint{-3.609559\du}{6.438868\du}}
\pgfpathlineto{\pgfpoint{-3.600304\du}{6.437973\du}}
\pgfpathlineto{\pgfpoint{-3.590750\du}{6.436778\du}}
\pgfpathlineto{\pgfpoint{-3.582092\du}{6.435286\du}}
\pgfpathlineto{\pgfpoint{-3.572837\du}{6.433793\du}}
\pgfpathlineto{\pgfpoint{-3.564776\du}{6.432300\du}}
\pgfpathlineto{\pgfpoint{-3.556417\du}{6.430210\du}}
\pgfpathlineto{\pgfpoint{-3.548057\du}{6.428120\du}}
\pgfpathlineto{\pgfpoint{-3.540594\du}{6.426030\du}}
\pgfpathlineto{\pgfpoint{-3.533130\du}{6.424239\du}}
\pgfpathlineto{\pgfpoint{-3.525964\du}{6.421851\du}}
\pgfpathlineto{\pgfpoint{-3.519396\du}{6.419462\du}}
\pgfpathlineto{\pgfpoint{-3.513127\du}{6.416775\du}}
\pgfpathlineto{\pgfpoint{-3.507454\du}{6.414088\du}}
\pgfpathlineto{\pgfpoint{-3.501782\du}{6.411401\du}}
\pgfpathlineto{\pgfpoint{-3.496706\du}{6.408714\du}}
\pgfpathlineto{\pgfpoint{-3.491930\du}{6.405430\du}}
\pgfpathlineto{\pgfpoint{-3.487750\du}{6.402445\du}}
\pgfpathlineto{\pgfpoint{-3.484167\du}{6.399459\du}}
\pgfpathlineto{\pgfpoint{-3.481182\du}{6.396474\du}}
\pgfpathlineto{\pgfpoint{-3.478495\du}{6.393190\du}}
\pgfpathlineto{\pgfpoint{-3.476106\du}{6.389906\du}}
\pgfpathlineto{\pgfpoint{-3.474315\du}{6.386622\du}}
\pgfpathlineto{\pgfpoint{-3.473419\du}{6.383338\du}}
\pgfpathlineto{\pgfpoint{-3.472524\du}{6.380352\du}}
\pgfpathlineto{\pgfpoint{-3.471927\du}{6.376769\du}}
\pgfusepath{fill}
\pgfsetbuttcap
\pgfsetmiterjoin
\pgfsetdash{}{0pt}
\definecolor{dialinecolor}{rgb}{0.666667, 0.901961, 1.000000}
\pgfsetstrokecolor{dialinecolor}
\pgfpathmoveto{\pgfpoint{-3.484167\du}{6.371097\du}}
\pgfpathlineto{\pgfpoint{-3.484167\du}{6.367514\du}}
\pgfpathlineto{\pgfpoint{-3.485063\du}{6.365126\du}}
\pgfpathlineto{\pgfpoint{-3.486257\du}{6.361842\du}}
\pgfpathlineto{\pgfpoint{-3.487750\du}{6.358856\du}}
\pgfpathlineto{\pgfpoint{-3.490138\du}{6.356169\du}}
\pgfpathlineto{\pgfpoint{-3.492527\du}{6.352587\du}}
\pgfpathlineto{\pgfpoint{-3.495811\du}{6.350198\du}}
\pgfpathlineto{\pgfpoint{-3.499095\du}{6.347511\du}}
\pgfpathlineto{\pgfpoint{-3.503275\du}{6.344824\du}}
\pgfpathlineto{\pgfpoint{-3.507753\du}{6.342436\du}}
\pgfpathlineto{\pgfpoint{-3.512530\du}{6.339450\du}}
\pgfpathlineto{\pgfpoint{-3.518202\du}{6.337361\du}}
\pgfpathlineto{\pgfpoint{-3.523576\du}{6.334674\du}}
\pgfpathlineto{\pgfpoint{-3.529846\du}{6.331987\du}}
\pgfpathlineto{\pgfpoint{-3.535817\du}{6.330195\du}}
\pgfpathlineto{\pgfpoint{-3.543280\du}{6.327807\du}}
\pgfpathlineto{\pgfpoint{-3.550147\du}{6.326314\du}}
\pgfpathlineto{\pgfpoint{-3.557312\du}{6.324224\du}}
\pgfpathlineto{\pgfpoint{-3.565075\du}{6.322433\du}}
\pgfpathlineto{\pgfpoint{-3.573733\du}{6.320642\du}}
\pgfpathlineto{\pgfpoint{-3.582092\du}{6.319447\du}}
\pgfpathlineto{\pgfpoint{-3.590452\du}{6.317059\du}}
\pgfpathlineto{\pgfpoint{-3.599408\du}{6.316760\du}}
\pgfpathlineto{\pgfpoint{-3.608066\du}{6.315268\du}}
\pgfpathlineto{\pgfpoint{-3.617023\du}{6.314372\du}}
\pgfpathlineto{\pgfpoint{-3.626278\du}{6.313476\du}}
\pgfpathlineto{\pgfpoint{-3.635832\du}{6.312879\du}}
\pgfpathlineto{\pgfpoint{-3.645385\du}{6.312581\du}}
\pgfpathlineto{\pgfpoint{-3.654342\du}{6.311984\du}}
\pgfpathlineto{\pgfpoint{-3.664493\du}{6.311685\du}}
\pgfpathlineto{\pgfpoint{-3.674046\du}{6.311685\du}}
\pgfpathlineto{\pgfpoint{-3.674046\du}{6.311685\du}}
\pgfpathlineto{\pgfpoint{-3.683301\du}{6.311685\du}}
\pgfpathlineto{\pgfpoint{-3.693154\du}{6.311984\du}}
\pgfpathlineto{\pgfpoint{-3.702707\du}{6.312581\du}}
\pgfpathlineto{\pgfpoint{-3.712261\du}{6.312879\du}}
\pgfpathlineto{\pgfpoint{-3.721516\du}{6.313476\du}}
\pgfpathlineto{\pgfpoint{-3.731368\du}{6.314372\du}}
\pgfpathlineto{\pgfpoint{-3.740026\du}{6.315268\du}}
\pgfpathlineto{\pgfpoint{-3.748983\du}{6.316760\du}}
\pgfpathlineto{\pgfpoint{-3.757342\du}{6.317059\du}}
\pgfpathlineto{\pgfpoint{-3.766299\du}{6.319447\du}}
\pgfpathlineto{\pgfpoint{-3.774658\du}{6.320642\du}}
\pgfpathlineto{\pgfpoint{-3.782719\du}{6.322433\du}}
\pgfpathlineto{\pgfpoint{-3.790482\du}{6.324224\du}}
\pgfpathlineto{\pgfpoint{-3.797945\du}{6.326314\du}}
\pgfpathlineto{\pgfpoint{-3.804812\du}{6.327807\du}}
\pgfpathlineto{\pgfpoint{-3.812276\du}{6.330195\du}}
\pgfpathlineto{\pgfpoint{-3.818247\du}{6.331987\du}}
\pgfpathlineto{\pgfpoint{-3.824516\du}{6.334674\du}}
\pgfpathlineto{\pgfpoint{-3.830189\du}{6.337361\du}}
\pgfpathlineto{\pgfpoint{-3.835563\du}{6.339450\du}}
\pgfpathlineto{\pgfpoint{-3.840340\du}{6.342436\du}}
\pgfpathlineto{\pgfpoint{-3.844818\du}{6.344824\du}}
\pgfpathlineto{\pgfpoint{-3.848998\du}{6.347511\du}}
\pgfpathlineto{\pgfpoint{-3.852282\du}{6.350198\du}}
\pgfpathlineto{\pgfpoint{-3.855864\du}{6.352587\du}}
\pgfpathlineto{\pgfpoint{-3.857656\du}{6.356169\du}}
\pgfpathlineto{\pgfpoint{-3.860044\du}{6.358856\du}}
\pgfpathlineto{\pgfpoint{-3.861835\du}{6.361842\du}}
\pgfpathlineto{\pgfpoint{-3.863328\du}{6.365126\du}}
\pgfpathlineto{\pgfpoint{-3.863925\du}{6.367514\du}}
\pgfpathlineto{\pgfpoint{-3.864224\du}{6.371097\du}}
\pgfpathlineto{\pgfpoint{-3.864224\du}{6.371097\du}}
\pgfpathlineto{\pgfpoint{-3.863925\du}{6.374082\du}}
\pgfpathlineto{\pgfpoint{-3.863328\du}{6.376769\du}}
\pgfpathlineto{\pgfpoint{-3.861835\du}{6.379755\du}}
\pgfpathlineto{\pgfpoint{-3.860044\du}{6.383338\du}}
\pgfpathlineto{\pgfpoint{-3.857656\du}{6.386024\du}}
\pgfpathlineto{\pgfpoint{-3.855864\du}{6.389010\du}}
\pgfpathlineto{\pgfpoint{-3.852282\du}{6.391697\du}}
\pgfpathlineto{\pgfpoint{-3.848998\du}{6.394384\du}}
\pgfpathlineto{\pgfpoint{-3.844818\du}{6.397071\du}}
\pgfpathlineto{\pgfpoint{-3.840340\du}{6.399758\du}}
\pgfpathlineto{\pgfpoint{-3.835563\du}{6.402445\du}}
\pgfpathlineto{\pgfpoint{-3.830189\du}{6.404833\du}}
\pgfpathlineto{\pgfpoint{-3.824516\du}{6.407520\du}}
\pgfpathlineto{\pgfpoint{-3.818247\du}{6.409610\du}}
\pgfpathlineto{\pgfpoint{-3.812276\du}{6.411999\du}}
\pgfpathlineto{\pgfpoint{-3.804812\du}{6.414387\du}}
\pgfpathlineto{\pgfpoint{-3.797945\du}{6.416178\du}}
\pgfpathlineto{\pgfpoint{-3.790482\du}{6.417671\du}}
\pgfpathlineto{\pgfpoint{-3.782719\du}{6.419761\du}}
\pgfpathlineto{\pgfpoint{-3.774658\du}{6.421552\du}}
\pgfpathlineto{\pgfpoint{-3.766299\du}{6.422746\du}}
\pgfpathlineto{\pgfpoint{-3.757342\du}{6.424538\du}}
\pgfpathlineto{\pgfpoint{-3.748983\du}{6.425433\du}}
\pgfpathlineto{\pgfpoint{-3.740026\du}{6.426926\du}}
\pgfpathlineto{\pgfpoint{-3.731368\du}{6.427822\du}}
\pgfpathlineto{\pgfpoint{-3.721516\du}{6.428120\du}}
\pgfpathlineto{\pgfpoint{-3.712261\du}{6.429315\du}}
\pgfpathlineto{\pgfpoint{-3.702707\du}{6.429613\du}}
\pgfpathlineto{\pgfpoint{-3.693154\du}{6.429912\du}}
\pgfpathlineto{\pgfpoint{-3.683301\du}{6.430509\du}}
\pgfpathlineto{\pgfpoint{-3.674046\du}{6.430509\du}}
\pgfpathlineto{\pgfpoint{-3.674046\du}{6.430509\du}}
\pgfpathlineto{\pgfpoint{-3.664493\du}{6.430509\du}}
\pgfpathlineto{\pgfpoint{-3.654342\du}{6.429912\du}}
\pgfpathlineto{\pgfpoint{-3.645385\du}{6.429613\du}}
\pgfpathlineto{\pgfpoint{-3.635832\du}{6.429315\du}}
\pgfpathlineto{\pgfpoint{-3.626278\du}{6.428120\du}}
\pgfpathlineto{\pgfpoint{-3.617023\du}{6.427822\du}}
\pgfpathlineto{\pgfpoint{-3.608066\du}{6.426926\du}}
\pgfpathlineto{\pgfpoint{-3.599408\du}{6.425433\du}}
\pgfpathlineto{\pgfpoint{-3.590452\du}{6.424538\du}}
\pgfpathlineto{\pgfpoint{-3.582092\du}{6.422746\du}}
\pgfpathlineto{\pgfpoint{-3.573733\du}{6.421552\du}}
\pgfpathlineto{\pgfpoint{-3.565075\du}{6.419761\du}}
\pgfpathlineto{\pgfpoint{-3.557312\du}{6.417671\du}}
\pgfpathlineto{\pgfpoint{-3.550147\du}{6.416178\du}}
\pgfpathlineto{\pgfpoint{-3.543280\du}{6.414387\du}}
\pgfpathlineto{\pgfpoint{-3.535817\du}{6.411999\du}}
\pgfpathlineto{\pgfpoint{-3.529846\du}{6.409610\du}}
\pgfpathlineto{\pgfpoint{-3.523576\du}{6.407520\du}}
\pgfpathlineto{\pgfpoint{-3.518202\du}{6.404833\du}}
\pgfpathlineto{\pgfpoint{-3.512530\du}{6.402445\du}}
\pgfpathlineto{\pgfpoint{-3.507753\du}{6.399758\du}}
\pgfpathlineto{\pgfpoint{-3.503275\du}{6.397071\du}}
\pgfpathlineto{\pgfpoint{-3.499095\du}{6.394384\du}}
\pgfpathlineto{\pgfpoint{-3.495811\du}{6.391697\du}}
\pgfpathlineto{\pgfpoint{-3.492527\du}{6.389010\du}}
\pgfpathlineto{\pgfpoint{-3.490138\du}{6.386024\du}}
\pgfpathlineto{\pgfpoint{-3.487750\du}{6.383338\du}}
\pgfpathlineto{\pgfpoint{-3.486257\du}{6.379755\du}}
\pgfpathlineto{\pgfpoint{-3.485063\du}{6.376769\du}}
\pgfpathlineto{\pgfpoint{-3.484167\du}{6.374082\du}}
\pgfpathlineto{\pgfpoint{-3.484167\du}{6.371097\du}}
\pgfusepath{stroke}
\pgfsetbuttcap
\pgfsetmiterjoin
\pgfsetdash{}{0pt}
\definecolor{dialinecolor}{rgb}{0.000000, 0.470588, 0.666667}
\pgfsetfillcolor{dialinecolor}
\pgfpathmoveto{\pgfpoint{-3.864224\du}{6.287801\du}}
\pgfpathlineto{\pgfpoint{-3.864224\du}{6.382442\du}}
\pgfpathlineto{\pgfpoint{-3.471628\du}{6.382442\du}}
\pgfpathlineto{\pgfpoint{-3.471628\du}{6.287801\du}}
\pgfpathlineto{\pgfpoint{-3.864224\du}{6.287801\du}}
\pgfusepath{fill}
\pgfsetbuttcap
\pgfsetmiterjoin
\pgfsetdash{}{0pt}
\definecolor{dialinecolor}{rgb}{0.000000, 0.705882, 1.000000}
\pgfsetfillcolor{dialinecolor}
\pgfpathmoveto{\pgfpoint{-3.471927\du}{6.281233\du}}
\pgfpathlineto{\pgfpoint{-3.472524\du}{6.277949\du}}
\pgfpathlineto{\pgfpoint{-3.473419\du}{6.274366\du}}
\pgfpathlineto{\pgfpoint{-3.474315\du}{6.271381\du}}
\pgfpathlineto{\pgfpoint{-3.476106\du}{6.268096\du}}
\pgfpathlineto{\pgfpoint{-3.478495\du}{6.264812\du}}
\pgfpathlineto{\pgfpoint{-3.481182\du}{6.261528\du}}
\pgfpathlineto{\pgfpoint{-3.484167\du}{6.258543\du}}
\pgfpathlineto{\pgfpoint{-3.487750\du}{6.255259\du}}
\pgfpathlineto{\pgfpoint{-3.491930\du}{6.252572\du}}
\pgfpathlineto{\pgfpoint{-3.496706\du}{6.249586\du}}
\pgfpathlineto{\pgfpoint{-3.501782\du}{6.246601\du}}
\pgfpathlineto{\pgfpoint{-3.507454\du}{6.243914\du}}
\pgfpathlineto{\pgfpoint{-3.513127\du}{6.240630\du}}
\pgfpathlineto{\pgfpoint{-3.519396\du}{6.238838\du}}
\pgfpathlineto{\pgfpoint{-3.525964\du}{6.236151\du}}
\pgfpathlineto{\pgfpoint{-3.533130\du}{6.233763\du}}
\pgfpathlineto{\pgfpoint{-3.540594\du}{6.231673\du}}
\pgfpathlineto{\pgfpoint{-3.548057\du}{6.229583\du}}
\pgfpathlineto{\pgfpoint{-3.556417\du}{6.227493\du}}
\pgfpathlineto{\pgfpoint{-3.564776\du}{6.225404\du}}
\pgfpathlineto{\pgfpoint{-3.572837\du}{6.223911\du}}
\pgfpathlineto{\pgfpoint{-3.582092\du}{6.222418\du}}
\pgfpathlineto{\pgfpoint{-3.590750\du}{6.221224\du}}
\pgfpathlineto{\pgfpoint{-3.600304\du}{6.219731\du}}
\pgfpathlineto{\pgfpoint{-3.609559\du}{6.218835\du}}
\pgfpathlineto{\pgfpoint{-3.618814\du}{6.217641\du}}
\pgfpathlineto{\pgfpoint{-3.628666\du}{6.217044\du}}
\pgfpathlineto{\pgfpoint{-3.638519\du}{6.216447\du}}
\pgfpathlineto{\pgfpoint{-3.648371\du}{6.216148\du}}
\pgfpathlineto{\pgfpoint{-3.658522\du}{6.215850\du}}
\pgfpathlineto{\pgfpoint{-3.668374\du}{6.215850\du}}
\pgfpathlineto{\pgfpoint{-3.668374\du}{6.215850\du}}
\pgfpathlineto{\pgfpoint{-3.678226\du}{6.215850\du}}
\pgfpathlineto{\pgfpoint{-3.688078\du}{6.216148\du}}
\pgfpathlineto{\pgfpoint{-3.698229\du}{6.216447\du}}
\pgfpathlineto{\pgfpoint{-3.708081\du}{6.217044\du}}
\pgfpathlineto{\pgfpoint{-3.717635\du}{6.217641\du}}
\pgfpathlineto{\pgfpoint{-3.727189\du}{6.218835\du}}
\pgfpathlineto{\pgfpoint{-3.736444\du}{6.219731\du}}
\pgfpathlineto{\pgfpoint{-3.745699\du}{6.221224\du}}
\pgfpathlineto{\pgfpoint{-3.754655\du}{6.222418\du}}
\pgfpathlineto{\pgfpoint{-3.763612\du}{6.223911\du}}
\pgfpathlineto{\pgfpoint{-3.771971\du}{6.225404\du}}
\pgfpathlineto{\pgfpoint{-3.780331\du}{6.227493\du}}
\pgfpathlineto{\pgfpoint{-3.788690\du}{6.229583\du}}
\pgfpathlineto{\pgfpoint{-3.795856\du}{6.231673\du}}
\pgfpathlineto{\pgfpoint{-3.803319\du}{6.233763\du}}
\pgfpathlineto{\pgfpoint{-3.810485\du}{6.236151\du}}
\pgfpathlineto{\pgfpoint{-3.817351\du}{6.238838\du}}
\pgfpathlineto{\pgfpoint{-3.823322\du}{6.240630\du}}
\pgfpathlineto{\pgfpoint{-3.829293\du}{6.243914\du}}
\pgfpathlineto{\pgfpoint{-3.834966\du}{6.246601\du}}
\pgfpathlineto{\pgfpoint{-3.840041\du}{6.249586\du}}
\pgfpathlineto{\pgfpoint{-3.844818\du}{6.252572\du}}
\pgfpathlineto{\pgfpoint{-3.848998\du}{6.255259\du}}
\pgfpathlineto{\pgfpoint{-3.852282\du}{6.258543\du}}
\pgfpathlineto{\pgfpoint{-3.855566\du}{6.261528\du}}
\pgfpathlineto{\pgfpoint{-3.858551\du}{6.264812\du}}
\pgfpathlineto{\pgfpoint{-3.860641\du}{6.268096\du}}
\pgfpathlineto{\pgfpoint{-3.862134\du}{6.271381\du}}
\pgfpathlineto{\pgfpoint{-3.863925\du}{6.274366\du}}
\pgfpathlineto{\pgfpoint{-3.864224\du}{6.277949\du}}
\pgfpathlineto{\pgfpoint{-3.864522\du}{6.281233\du}}
\pgfpathlineto{\pgfpoint{-3.864522\du}{6.281233\du}}
\pgfpathlineto{\pgfpoint{-3.864224\du}{6.284517\du}}
\pgfpathlineto{\pgfpoint{-3.863925\du}{6.288099\du}}
\pgfpathlineto{\pgfpoint{-3.862134\du}{6.291085\du}}
\pgfpathlineto{\pgfpoint{-3.860641\du}{6.294369\du}}
\pgfpathlineto{\pgfpoint{-3.858551\du}{6.297952\du}}
\pgfpathlineto{\pgfpoint{-3.855566\du}{6.300937\du}}
\pgfpathlineto{\pgfpoint{-3.852282\du}{6.303923\du}}
\pgfpathlineto{\pgfpoint{-3.848998\du}{6.307207\du}}
\pgfpathlineto{\pgfpoint{-3.844818\du}{6.309595\du}}
\pgfpathlineto{\pgfpoint{-3.840041\du}{6.312879\du}}
\pgfpathlineto{\pgfpoint{-3.834966\du}{6.316163\du}}
\pgfpathlineto{\pgfpoint{-3.829293\du}{6.318850\du}}
\pgfpathlineto{\pgfpoint{-3.823322\du}{6.321537\du}}
\pgfpathlineto{\pgfpoint{-3.817351\du}{6.323627\du}}
\pgfpathlineto{\pgfpoint{-3.810485\du}{6.326613\du}}
\pgfpathlineto{\pgfpoint{-3.803319\du}{6.329001\du}}
\pgfpathlineto{\pgfpoint{-3.795856\du}{6.330792\du}}
\pgfpathlineto{\pgfpoint{-3.788690\du}{6.332882\du}}
\pgfpathlineto{\pgfpoint{-3.780331\du}{6.334972\du}}
\pgfpathlineto{\pgfpoint{-3.771971\du}{6.337062\du}}
\pgfpathlineto{\pgfpoint{-3.763612\du}{6.338555\du}}
\pgfpathlineto{\pgfpoint{-3.754655\du}{6.340047\du}}
\pgfpathlineto{\pgfpoint{-3.745699\du}{6.341242\du}}
\pgfpathlineto{\pgfpoint{-3.736444\du}{6.342734\du}}
\pgfpathlineto{\pgfpoint{-3.727189\du}{6.343630\du}}
\pgfpathlineto{\pgfpoint{-3.717635\du}{6.344824\du}}
\pgfpathlineto{\pgfpoint{-3.708081\du}{6.345123\du}}
\pgfpathlineto{\pgfpoint{-3.698229\du}{6.346019\du}}
\pgfpathlineto{\pgfpoint{-3.688078\du}{6.346019\du}}
\pgfpathlineto{\pgfpoint{-3.678226\du}{6.346914\du}}
\pgfpathlineto{\pgfpoint{-3.668374\du}{6.346914\du}}
\pgfpathlineto{\pgfpoint{-3.668374\du}{6.346914\du}}
\pgfpathlineto{\pgfpoint{-3.658522\du}{6.346914\du}}
\pgfpathlineto{\pgfpoint{-3.648371\du}{6.346019\du}}
\pgfpathlineto{\pgfpoint{-3.638519\du}{6.346019\du}}
\pgfpathlineto{\pgfpoint{-3.628666\du}{6.345123\du}}
\pgfpathlineto{\pgfpoint{-3.618814\du}{6.344824\du}}
\pgfpathlineto{\pgfpoint{-3.609559\du}{6.343630\du}}
\pgfpathlineto{\pgfpoint{-3.600304\du}{6.342734\du}}
\pgfpathlineto{\pgfpoint{-3.590750\du}{6.341242\du}}
\pgfpathlineto{\pgfpoint{-3.582092\du}{6.340047\du}}
\pgfpathlineto{\pgfpoint{-3.572837\du}{6.338555\du}}
\pgfpathlineto{\pgfpoint{-3.564776\du}{6.337062\du}}
\pgfpathlineto{\pgfpoint{-3.556417\du}{6.334972\du}}
\pgfpathlineto{\pgfpoint{-3.548057\du}{6.332882\du}}
\pgfpathlineto{\pgfpoint{-3.540594\du}{6.330792\du}}
\pgfpathlineto{\pgfpoint{-3.533130\du}{6.329001\du}}
\pgfpathlineto{\pgfpoint{-3.525964\du}{6.326613\du}}
\pgfpathlineto{\pgfpoint{-3.519396\du}{6.323627\du}}
\pgfpathlineto{\pgfpoint{-3.513127\du}{6.321537\du}}
\pgfpathlineto{\pgfpoint{-3.507454\du}{6.318850\du}}
\pgfpathlineto{\pgfpoint{-3.501782\du}{6.316163\du}}
\pgfpathlineto{\pgfpoint{-3.496706\du}{6.312879\du}}
\pgfpathlineto{\pgfpoint{-3.491930\du}{6.309595\du}}
\pgfpathlineto{\pgfpoint{-3.487750\du}{6.307207\du}}
\pgfpathlineto{\pgfpoint{-3.484167\du}{6.303923\du}}
\pgfpathlineto{\pgfpoint{-3.481182\du}{6.300937\du}}
\pgfpathlineto{\pgfpoint{-3.478495\du}{6.297952\du}}
\pgfpathlineto{\pgfpoint{-3.476106\du}{6.294369\du}}
\pgfpathlineto{\pgfpoint{-3.474315\du}{6.291085\du}}
\pgfpathlineto{\pgfpoint{-3.473419\du}{6.288099\du}}
\pgfpathlineto{\pgfpoint{-3.472524\du}{6.284517\du}}
\pgfpathlineto{\pgfpoint{-3.471927\du}{6.281233\du}}
\pgfusepath{fill}
\pgfsetbuttcap
\pgfsetmiterjoin
\pgfsetdash{}{0pt}
\definecolor{dialinecolor}{rgb}{0.666667, 0.901961, 1.000000}
\pgfsetstrokecolor{dialinecolor}
\pgfpathmoveto{\pgfpoint{-3.484167\du}{6.275560\du}}
\pgfpathlineto{\pgfpoint{-3.484167\du}{6.272276\du}}
\pgfpathlineto{\pgfpoint{-3.485063\du}{6.269888\du}}
\pgfpathlineto{\pgfpoint{-3.486257\du}{6.266604\du}}
\pgfpathlineto{\pgfpoint{-3.487750\du}{6.263021\du}}
\pgfpathlineto{\pgfpoint{-3.490138\du}{6.260633\du}}
\pgfpathlineto{\pgfpoint{-3.492527\du}{6.257349\du}}
\pgfpathlineto{\pgfpoint{-3.495811\du}{6.254960\du}}
\pgfpathlineto{\pgfpoint{-3.499095\du}{6.252273\du}}
\pgfpathlineto{\pgfpoint{-3.503275\du}{6.249586\du}}
\pgfpathlineto{\pgfpoint{-3.507753\du}{6.246899\du}}
\pgfpathlineto{\pgfpoint{-3.512530\du}{6.244212\du}}
\pgfpathlineto{\pgfpoint{-3.518202\du}{6.241525\du}}
\pgfpathlineto{\pgfpoint{-3.523576\du}{6.239137\du}}
\pgfpathlineto{\pgfpoint{-3.529846\du}{6.237047\du}}
\pgfpathlineto{\pgfpoint{-3.535817\du}{6.234659\du}}
\pgfpathlineto{\pgfpoint{-3.543280\du}{6.232270\du}}
\pgfpathlineto{\pgfpoint{-3.550147\du}{6.230479\du}}
\pgfpathlineto{\pgfpoint{-3.557312\du}{6.228986\du}}
\pgfpathlineto{\pgfpoint{-3.565075\du}{6.226598\du}}
\pgfpathlineto{\pgfpoint{-3.573733\du}{6.225105\du}}
\pgfpathlineto{\pgfpoint{-3.582092\du}{6.223911\du}}
\pgfpathlineto{\pgfpoint{-3.590452\du}{6.222119\du}}
\pgfpathlineto{\pgfpoint{-3.599408\du}{6.221224\du}}
\pgfpathlineto{\pgfpoint{-3.608066\du}{6.219731\du}}
\pgfpathlineto{\pgfpoint{-3.617023\du}{6.218835\du}}
\pgfpathlineto{\pgfpoint{-3.626278\du}{6.218238\du}}
\pgfpathlineto{\pgfpoint{-3.635832\du}{6.217343\du}}
\pgfpathlineto{\pgfpoint{-3.645385\du}{6.217044\du}}
\pgfpathlineto{\pgfpoint{-3.654342\du}{6.216745\du}}
\pgfpathlineto{\pgfpoint{-3.664493\du}{6.216148\du}}
\pgfpathlineto{\pgfpoint{-3.674046\du}{6.216148\du}}
\pgfpathlineto{\pgfpoint{-3.674046\du}{6.216148\du}}
\pgfpathlineto{\pgfpoint{-3.683301\du}{6.216148\du}}
\pgfpathlineto{\pgfpoint{-3.693154\du}{6.216745\du}}
\pgfpathlineto{\pgfpoint{-3.702707\du}{6.217044\du}}
\pgfpathlineto{\pgfpoint{-3.712261\du}{6.217343\du}}
\pgfpathlineto{\pgfpoint{-3.721516\du}{6.218238\du}}
\pgfpathlineto{\pgfpoint{-3.731368\du}{6.218835\du}}
\pgfpathlineto{\pgfpoint{-3.740026\du}{6.219731\du}}
\pgfpathlineto{\pgfpoint{-3.748983\du}{6.221224\du}}
\pgfpathlineto{\pgfpoint{-3.757342\du}{6.222119\du}}
\pgfpathlineto{\pgfpoint{-3.766299\du}{6.223911\du}}
\pgfpathlineto{\pgfpoint{-3.774658\du}{6.225105\du}}
\pgfpathlineto{\pgfpoint{-3.782719\du}{6.226598\du}}
\pgfpathlineto{\pgfpoint{-3.790482\du}{6.228986\du}}
\pgfpathlineto{\pgfpoint{-3.797945\du}{6.230479\du}}
\pgfpathlineto{\pgfpoint{-3.804812\du}{6.232270\du}}
\pgfpathlineto{\pgfpoint{-3.812276\du}{6.234659\du}}
\pgfpathlineto{\pgfpoint{-3.818247\du}{6.237047\du}}
\pgfpathlineto{\pgfpoint{-3.824516\du}{6.239137\du}}
\pgfpathlineto{\pgfpoint{-3.830189\du}{6.241525\du}}
\pgfpathlineto{\pgfpoint{-3.835563\du}{6.244212\du}}
\pgfpathlineto{\pgfpoint{-3.840340\du}{6.246899\du}}
\pgfpathlineto{\pgfpoint{-3.844818\du}{6.249586\du}}
\pgfpathlineto{\pgfpoint{-3.848998\du}{6.252273\du}}
\pgfpathlineto{\pgfpoint{-3.852282\du}{6.254960\du}}
\pgfpathlineto{\pgfpoint{-3.855864\du}{6.257349\du}}
\pgfpathlineto{\pgfpoint{-3.857656\du}{6.260633\du}}
\pgfpathlineto{\pgfpoint{-3.860044\du}{6.263021\du}}
\pgfpathlineto{\pgfpoint{-3.861835\du}{6.266604\du}}
\pgfpathlineto{\pgfpoint{-3.863328\du}{6.269888\du}}
\pgfpathlineto{\pgfpoint{-3.863925\du}{6.272276\du}}
\pgfpathlineto{\pgfpoint{-3.864224\du}{6.275560\du}}
\pgfpathlineto{\pgfpoint{-3.864224\du}{6.275560\du}}
\pgfpathlineto{\pgfpoint{-3.863925\du}{6.278844\du}}
\pgfpathlineto{\pgfpoint{-3.863328\du}{6.281531\du}}
\pgfpathlineto{\pgfpoint{-3.861835\du}{6.284517\du}}
\pgfpathlineto{\pgfpoint{-3.860044\du}{6.287801\du}}
\pgfpathlineto{\pgfpoint{-3.857656\du}{6.290488\du}}
\pgfpathlineto{\pgfpoint{-3.855864\du}{6.293772\du}}
\pgfpathlineto{\pgfpoint{-3.852282\du}{6.296459\du}}
\pgfpathlineto{\pgfpoint{-3.848998\du}{6.299146\du}}
\pgfpathlineto{\pgfpoint{-3.844818\du}{6.301833\du}}
\pgfpathlineto{\pgfpoint{-3.840340\du}{6.304520\du}}
\pgfpathlineto{\pgfpoint{-3.835563\du}{6.307207\du}}
\pgfpathlineto{\pgfpoint{-3.830189\du}{6.309297\du}}
\pgfpathlineto{\pgfpoint{-3.824516\du}{6.311984\du}}
\pgfpathlineto{\pgfpoint{-3.818247\du}{6.314372\du}}
\pgfpathlineto{\pgfpoint{-3.812276\du}{6.316760\du}}
\pgfpathlineto{\pgfpoint{-3.804812\du}{6.318850\du}}
\pgfpathlineto{\pgfpoint{-3.797945\du}{6.320642\du}}
\pgfpathlineto{\pgfpoint{-3.790482\du}{6.322433\du}}
\pgfpathlineto{\pgfpoint{-3.782719\du}{6.324224\du}}
\pgfpathlineto{\pgfpoint{-3.774658\du}{6.326314\du}}
\pgfpathlineto{\pgfpoint{-3.766299\du}{6.327508\du}}
\pgfpathlineto{\pgfpoint{-3.757342\du}{6.329300\du}}
\pgfpathlineto{\pgfpoint{-3.748983\du}{6.330195\du}}
\pgfpathlineto{\pgfpoint{-3.740026\du}{6.331688\du}}
\pgfpathlineto{\pgfpoint{-3.731368\du}{6.331987\du}}
\pgfpathlineto{\pgfpoint{-3.721516\du}{6.333181\du}}
\pgfpathlineto{\pgfpoint{-3.712261\du}{6.333479\du}}
\pgfpathlineto{\pgfpoint{-3.702707\du}{6.334375\du}}
\pgfpathlineto{\pgfpoint{-3.693154\du}{6.334674\du}}
\pgfpathlineto{\pgfpoint{-3.683301\du}{6.334972\du}}
\pgfpathlineto{\pgfpoint{-3.674046\du}{6.334972\du}}
\pgfpathlineto{\pgfpoint{-3.674046\du}{6.334972\du}}
\pgfpathlineto{\pgfpoint{-3.664493\du}{6.334972\du}}
\pgfpathlineto{\pgfpoint{-3.654342\du}{6.334674\du}}
\pgfpathlineto{\pgfpoint{-3.645385\du}{6.334375\du}}
\pgfpathlineto{\pgfpoint{-3.635832\du}{6.333479\du}}
\pgfpathlineto{\pgfpoint{-3.626278\du}{6.333181\du}}
\pgfpathlineto{\pgfpoint{-3.617023\du}{6.331987\du}}
\pgfpathlineto{\pgfpoint{-3.608066\du}{6.331688\du}}
\pgfpathlineto{\pgfpoint{-3.599408\du}{6.330195\du}}
\pgfpathlineto{\pgfpoint{-3.590452\du}{6.329300\du}}
\pgfpathlineto{\pgfpoint{-3.582092\du}{6.327508\du}}
\pgfpathlineto{\pgfpoint{-3.573733\du}{6.326314\du}}
\pgfpathlineto{\pgfpoint{-3.565075\du}{6.324224\du}}
\pgfpathlineto{\pgfpoint{-3.557312\du}{6.322433\du}}
\pgfpathlineto{\pgfpoint{-3.550147\du}{6.320642\du}}
\pgfpathlineto{\pgfpoint{-3.543280\du}{6.318850\du}}
\pgfpathlineto{\pgfpoint{-3.535817\du}{6.316760\du}}
\pgfpathlineto{\pgfpoint{-3.529846\du}{6.314372\du}}
\pgfpathlineto{\pgfpoint{-3.523576\du}{6.311984\du}}
\pgfpathlineto{\pgfpoint{-3.518202\du}{6.309297\du}}
\pgfpathlineto{\pgfpoint{-3.512530\du}{6.307207\du}}
\pgfpathlineto{\pgfpoint{-3.507753\du}{6.304520\du}}
\pgfpathlineto{\pgfpoint{-3.503275\du}{6.301833\du}}
\pgfpathlineto{\pgfpoint{-3.499095\du}{6.299146\du}}
\pgfpathlineto{\pgfpoint{-3.495811\du}{6.296459\du}}
\pgfpathlineto{\pgfpoint{-3.492527\du}{6.293772\du}}
\pgfpathlineto{\pgfpoint{-3.490138\du}{6.290488\du}}
\pgfpathlineto{\pgfpoint{-3.487750\du}{6.287801\du}}
\pgfpathlineto{\pgfpoint{-3.486257\du}{6.284517\du}}
\pgfpathlineto{\pgfpoint{-3.485063\du}{6.281531\du}}
\pgfpathlineto{\pgfpoint{-3.484167\du}{6.278844\du}}
\pgfpathlineto{\pgfpoint{-3.484167\du}{6.275560\du}}
\pgfusepath{stroke}
\pgfsetbuttcap
\pgfsetmiterjoin
\pgfsetdash{}{0pt}
\definecolor{dialinecolor}{rgb}{0.000000, 0.000000, 0.000000}
\pgfsetfillcolor{dialinecolor}
\pgfpathmoveto{\pgfpoint{-3.661806\du}{6.263917\du}}
\pgfpathlineto{\pgfpoint{-3.638519\du}{6.275560\du}}
\pgfpathlineto{\pgfpoint{-3.566568\du}{6.251975\du}}
\pgfpathlineto{\pgfpoint{-3.531040\du}{6.263917\du}}
\pgfpathlineto{\pgfpoint{-3.554625\du}{6.240033\du}}
\pgfpathlineto{\pgfpoint{-3.638519\du}{6.240033\du}}
\pgfpathlineto{\pgfpoint{-3.602692\du}{6.240033\du}}
\pgfpathlineto{\pgfpoint{-3.661806\du}{6.263917\du}}
\pgfusepath{fill}
\pgfsetbuttcap
\pgfsetmiterjoin
\pgfsetdash{}{0pt}
\definecolor{dialinecolor}{rgb}{0.000000, 0.000000, 0.000000}
\pgfsetfillcolor{dialinecolor}
\pgfpathmoveto{\pgfpoint{-3.674046\du}{6.299444\du}}
\pgfpathlineto{\pgfpoint{-3.709873\du}{6.287801\du}}
\pgfpathlineto{\pgfpoint{-3.768986\du}{6.311685\du}}
\pgfpathlineto{\pgfpoint{-3.804812\du}{6.299444\du}}
\pgfpathlineto{\pgfpoint{-3.792571\du}{6.334972\du}}
\pgfpathlineto{\pgfpoint{-3.709873\du}{6.334972\du}}
\pgfpathlineto{\pgfpoint{-3.745400\du}{6.323627\du}}
\pgfpathlineto{\pgfpoint{-3.674046\du}{6.299444\du}}
\pgfusepath{fill}
\pgfsetbuttcap
\pgfsetmiterjoin
\pgfsetdash{}{0pt}
\definecolor{dialinecolor}{rgb}{0.000000, 0.000000, 0.000000}
\pgfsetfillcolor{dialinecolor}
\pgfpathmoveto{\pgfpoint{-3.792571\du}{6.240033\du}}
\pgfpathlineto{\pgfpoint{-3.768986\du}{6.227493\du}}
\pgfpathlineto{\pgfpoint{-3.697333\du}{6.263917\du}}
\pgfpathlineto{\pgfpoint{-3.674046\du}{6.251975\du}}
\pgfpathlineto{\pgfpoint{-3.685690\du}{6.275560\du}}
\pgfpathlineto{\pgfpoint{-3.768986\du}{6.275560\du}}
\pgfpathlineto{\pgfpoint{-3.733160\du}{6.263917\du}}
\pgfpathlineto{\pgfpoint{-3.792571\du}{6.240033\du}}
\pgfusepath{fill}
\pgfsetbuttcap
\pgfsetmiterjoin
\pgfsetdash{}{0pt}
\definecolor{dialinecolor}{rgb}{0.000000, 0.000000, 0.000000}
\pgfsetfillcolor{dialinecolor}
\pgfpathmoveto{\pgfpoint{-3.543280\du}{6.323627\du}}
\pgfpathlineto{\pgfpoint{-3.566568\du}{6.334972\du}}
\pgfpathlineto{\pgfpoint{-3.638519\du}{6.299444\du}}
\pgfpathlineto{\pgfpoint{-3.674046\du}{6.311685\du}}
\pgfpathlineto{\pgfpoint{-3.650162\du}{6.287801\du}}
\pgfpathlineto{\pgfpoint{-3.566568\du}{6.287801\du}}
\pgfpathlineto{\pgfpoint{-3.602692\du}{6.299444\du}}
\pgfpathlineto{\pgfpoint{-3.543280\du}{6.323627\du}}
\pgfusepath{fill}
\pgfsetbuttcap
\pgfsetmiterjoin
\pgfsetdash{}{0pt}
\definecolor{dialinecolor}{rgb}{1.000000, 1.000000, 1.000000}
\pgfsetfillcolor{dialinecolor}
\pgfpathmoveto{\pgfpoint{-3.661806\du}{6.275560\du}}
\pgfpathlineto{\pgfpoint{-3.638519\du}{6.275560\du}}
\pgfpathlineto{\pgfpoint{-3.566568\du}{6.251975\du}}
\pgfpathlineto{\pgfpoint{-3.531040\du}{6.263917\du}}
\pgfpathlineto{\pgfpoint{-3.554625\du}{6.240033\du}}
\pgfpathlineto{\pgfpoint{-3.626278\du}{6.240033\du}}
\pgfpathlineto{\pgfpoint{-3.602692\du}{6.240033\du}}
\pgfpathlineto{\pgfpoint{-3.661806\du}{6.275560\du}}
\pgfusepath{fill}
\pgfsetbuttcap
\pgfsetmiterjoin
\pgfsetdash{}{0pt}
\definecolor{dialinecolor}{rgb}{1.000000, 1.000000, 1.000000}
\pgfsetfillcolor{dialinecolor}
\pgfpathmoveto{\pgfpoint{-3.674046\du}{6.299444\du}}
\pgfpathlineto{\pgfpoint{-3.697333\du}{6.287801\du}}
\pgfpathlineto{\pgfpoint{-3.768986\du}{6.311685\du}}
\pgfpathlineto{\pgfpoint{-3.804812\du}{6.311685\du}}
\pgfpathlineto{\pgfpoint{-3.781226\du}{6.334972\du}}
\pgfpathlineto{\pgfpoint{-3.697333\du}{6.334972\du}}
\pgfpathlineto{\pgfpoint{-3.733160\du}{6.323627\du}}
\pgfpathlineto{\pgfpoint{-3.674046\du}{6.299444\du}}
\pgfusepath{fill}
\pgfsetbuttcap
\pgfsetmiterjoin
\pgfsetdash{}{0pt}
\definecolor{dialinecolor}{rgb}{1.000000, 1.000000, 1.000000}
\pgfsetfillcolor{dialinecolor}
\pgfpathmoveto{\pgfpoint{-3.792571\du}{6.240033\du}}
\pgfpathlineto{\pgfpoint{-3.768986\du}{6.240033\du}}
\pgfpathlineto{\pgfpoint{-3.697333\du}{6.263917\du}}
\pgfpathlineto{\pgfpoint{-3.661806\du}{6.251975\du}}
\pgfpathlineto{\pgfpoint{-3.685690\du}{6.275560\du}}
\pgfpathlineto{\pgfpoint{-3.768986\du}{6.275560\du}}
\pgfpathlineto{\pgfpoint{-3.733160\du}{6.275560\du}}
\pgfpathlineto{\pgfpoint{-3.792571\du}{6.240033\du}}
\pgfusepath{fill}
\pgfsetbuttcap
\pgfsetmiterjoin
\pgfsetdash{}{0pt}
\definecolor{dialinecolor}{rgb}{1.000000, 1.000000, 1.000000}
\pgfsetfillcolor{dialinecolor}
\pgfpathmoveto{\pgfpoint{-3.543280\du}{6.323627\du}}
\pgfpathlineto{\pgfpoint{-3.566568\du}{6.334972\du}}
\pgfpathlineto{\pgfpoint{-3.638519\du}{6.311685\du}}
\pgfpathlineto{\pgfpoint{-3.661806\du}{6.311685\du}}
\pgfpathlineto{\pgfpoint{-3.650162\du}{6.287801\du}}
\pgfpathlineto{\pgfpoint{-3.566568\du}{6.287801\du}}
\pgfpathlineto{\pgfpoint{-3.602692\du}{6.299444\du}}
\pgfpathlineto{\pgfpoint{-3.543280\du}{6.323627\du}}
\pgfusepath{fill}
\pgfsetbuttcap
\pgfsetmiterjoin
\pgfsetdash{}{0pt}
\definecolor{dialinecolor}{rgb}{0.666667, 0.901961, 1.000000}
\pgfsetstrokecolor{dialinecolor}
\pgfpathmoveto{\pgfpoint{-3.864224\du}{6.275560\du}}
\pgfpathlineto{\pgfpoint{-3.864224\du}{6.370500\du}}
\pgfusepath{stroke}
\pgfsetbuttcap
\pgfsetmiterjoin
\pgfsetdash{}{0pt}
\definecolor{dialinecolor}{rgb}{0.666667, 0.901961, 1.000000}
\pgfsetstrokecolor{dialinecolor}
\pgfpathmoveto{\pgfpoint{-3.483570\du}{6.275560\du}}
\pgfpathlineto{\pgfpoint{-3.483570\du}{6.370500\du}}
\pgfusepath{stroke}
\pgfsetbuttcap
\pgfsetmiterjoin
\pgfsetdash{}{0pt}
\definecolor{dialinecolor}{rgb}{0.000000, 0.470588, 0.666667}
\pgfsetfillcolor{dialinecolor}
\pgfpathmoveto{\pgfpoint{-2.757790\du}{6.376769\du}}
\pgfpathlineto{\pgfpoint{-2.757790\du}{6.373485\du}}
\pgfpathlineto{\pgfpoint{-2.758686\du}{6.370201\du}}
\pgfpathlineto{\pgfpoint{-2.759880\du}{6.366917\du}}
\pgfpathlineto{\pgfpoint{-2.761373\du}{6.363633\du}}
\pgfpathlineto{\pgfpoint{-2.763761\du}{6.360648\du}}
\pgfpathlineto{\pgfpoint{-2.766448\du}{6.357364\du}}
\pgfpathlineto{\pgfpoint{-2.769434\du}{6.354079\du}}
\pgfpathlineto{\pgfpoint{-2.773613\du}{6.351094\du}}
\pgfpathlineto{\pgfpoint{-2.777196\du}{6.348108\du}}
\pgfpathlineto{\pgfpoint{-2.781973\du}{6.345123\du}}
\pgfpathlineto{\pgfpoint{-2.787048\du}{6.342436\du}}
\pgfpathlineto{\pgfpoint{-2.792422\du}{6.339450\du}}
\pgfpathlineto{\pgfpoint{-2.798692\du}{6.337062\du}}
\pgfpathlineto{\pgfpoint{-2.804663\du}{6.334375\du}}
\pgfpathlineto{\pgfpoint{-2.811529\du}{6.331987\du}}
\pgfpathlineto{\pgfpoint{-2.818396\du}{6.329598\du}}
\pgfpathlineto{\pgfpoint{-2.825860\du}{6.327508\du}}
\pgfpathlineto{\pgfpoint{-2.833324\du}{6.325120\du}}
\pgfpathlineto{\pgfpoint{-2.841385\du}{6.323627\du}}
\pgfpathlineto{\pgfpoint{-2.850043\du}{6.321537\du}}
\pgfpathlineto{\pgfpoint{-2.858402\du}{6.319746\du}}
\pgfpathlineto{\pgfpoint{-2.867060\du}{6.318552\du}}
\pgfpathlineto{\pgfpoint{-2.876315\du}{6.317059\du}}
\pgfpathlineto{\pgfpoint{-2.885570\du}{6.315865\du}}
\pgfpathlineto{\pgfpoint{-2.894825\du}{6.314671\du}}
\pgfpathlineto{\pgfpoint{-2.904379\du}{6.313775\du}}
\pgfpathlineto{\pgfpoint{-2.914530\du}{6.312879\du}}
\pgfpathlineto{\pgfpoint{-2.924084\du}{6.312282\du}}
\pgfpathlineto{\pgfpoint{-2.933637\du}{6.311984\du}}
\pgfpathlineto{\pgfpoint{-2.943788\du}{6.311685\du}}
\pgfpathlineto{\pgfpoint{-2.953640\du}{6.311685\du}}
\pgfpathlineto{\pgfpoint{-2.953640\du}{6.311685\du}}
\pgfpathlineto{\pgfpoint{-2.963492\du}{6.311685\du}}
\pgfpathlineto{\pgfpoint{-2.973643\du}{6.311984\du}}
\pgfpathlineto{\pgfpoint{-2.983495\du}{6.312282\du}}
\pgfpathlineto{\pgfpoint{-2.993348\du}{6.312879\du}}
\pgfpathlineto{\pgfpoint{-3.002901\du}{6.313775\du}}
\pgfpathlineto{\pgfpoint{-3.012455\du}{6.314671\du}}
\pgfpathlineto{\pgfpoint{-3.022009\du}{6.315865\du}}
\pgfpathlineto{\pgfpoint{-3.030965\du}{6.317059\du}}
\pgfpathlineto{\pgfpoint{-3.040220\du}{6.318552\du}}
\pgfpathlineto{\pgfpoint{-3.048878\du}{6.319746\du}}
\pgfpathlineto{\pgfpoint{-3.057536\du}{6.321537\du}}
\pgfpathlineto{\pgfpoint{-3.065896\du}{6.323627\du}}
\pgfpathlineto{\pgfpoint{-3.073957\du}{6.325120\du}}
\pgfpathlineto{\pgfpoint{-3.081421\du}{6.327508\du}}
\pgfpathlineto{\pgfpoint{-3.088884\du}{6.329598\du}}
\pgfpathlineto{\pgfpoint{-3.096050\du}{6.331987\du}}
\pgfpathlineto{\pgfpoint{-3.102618\du}{6.334375\du}}
\pgfpathlineto{\pgfpoint{-3.108887\du}{6.337062\du}}
\pgfpathlineto{\pgfpoint{-3.114858\du}{6.339450\du}}
\pgfpathlineto{\pgfpoint{-3.120232\du}{6.342436\du}}
\pgfpathlineto{\pgfpoint{-3.125308\du}{6.345123\du}}
\pgfpathlineto{\pgfpoint{-3.130084\du}{6.348108\du}}
\pgfpathlineto{\pgfpoint{-3.133966\du}{6.351094\du}}
\pgfpathlineto{\pgfpoint{-3.137847\du}{6.354079\du}}
\pgfpathlineto{\pgfpoint{-3.140832\du}{6.357364\du}}
\pgfpathlineto{\pgfpoint{-3.143519\du}{6.360648\du}}
\pgfpathlineto{\pgfpoint{-3.145908\du}{6.363633\du}}
\pgfpathlineto{\pgfpoint{-3.147699\du}{6.366917\du}}
\pgfpathlineto{\pgfpoint{-3.149192\du}{6.370201\du}}
\pgfpathlineto{\pgfpoint{-3.150087\du}{6.373485\du}}
\pgfpathlineto{\pgfpoint{-3.150087\du}{6.376769\du}}
\pgfpathlineto{\pgfpoint{-3.150087\du}{6.376769\du}}
\pgfpathlineto{\pgfpoint{-3.150087\du}{6.380352\du}}
\pgfpathlineto{\pgfpoint{-3.149192\du}{6.383338\du}}
\pgfpathlineto{\pgfpoint{-3.147699\du}{6.386622\du}}
\pgfpathlineto{\pgfpoint{-3.145908\du}{6.389906\du}}
\pgfpathlineto{\pgfpoint{-3.143519\du}{6.393190\du}}
\pgfpathlineto{\pgfpoint{-3.140832\du}{6.396474\du}}
\pgfpathlineto{\pgfpoint{-3.137847\du}{6.399459\du}}
\pgfpathlineto{\pgfpoint{-3.133966\du}{6.402445\du}}
\pgfpathlineto{\pgfpoint{-3.130084\du}{6.405430\du}}
\pgfpathlineto{\pgfpoint{-3.125308\du}{6.408714\du}}
\pgfpathlineto{\pgfpoint{-3.120232\du}{6.411401\du}}
\pgfpathlineto{\pgfpoint{-3.114858\du}{6.414088\du}}
\pgfpathlineto{\pgfpoint{-3.108887\du}{6.416775\du}}
\pgfpathlineto{\pgfpoint{-3.102618\du}{6.419462\du}}
\pgfpathlineto{\pgfpoint{-3.096050\du}{6.421851\du}}
\pgfpathlineto{\pgfpoint{-3.088884\du}{6.424239\du}}
\pgfpathlineto{\pgfpoint{-3.081421\du}{6.426030\du}}
\pgfpathlineto{\pgfpoint{-3.073957\du}{6.428120\du}}
\pgfpathlineto{\pgfpoint{-3.065896\du}{6.430210\du}}
\pgfpathlineto{\pgfpoint{-3.057536\du}{6.432300\du}}
\pgfpathlineto{\pgfpoint{-3.048878\du}{6.433793\du}}
\pgfpathlineto{\pgfpoint{-3.040220\du}{6.435286\du}}
\pgfpathlineto{\pgfpoint{-3.030965\du}{6.436778\du}}
\pgfpathlineto{\pgfpoint{-3.022009\du}{6.437973\du}}
\pgfpathlineto{\pgfpoint{-3.012455\du}{6.438868\du}}
\pgfpathlineto{\pgfpoint{-3.002901\du}{6.440062\du}}
\pgfpathlineto{\pgfpoint{-2.993348\du}{6.440660\du}}
\pgfpathlineto{\pgfpoint{-2.983495\du}{6.441257\du}}
\pgfpathlineto{\pgfpoint{-2.973643\du}{6.441854\du}}
\pgfpathlineto{\pgfpoint{-2.963492\du}{6.442152\du}}
\pgfpathlineto{\pgfpoint{-2.953640\du}{6.442152\du}}
\pgfpathlineto{\pgfpoint{-2.953640\du}{6.442152\du}}
\pgfpathlineto{\pgfpoint{-2.943788\du}{6.442152\du}}
\pgfpathlineto{\pgfpoint{-2.933637\du}{6.441854\du}}
\pgfpathlineto{\pgfpoint{-2.924084\du}{6.441257\du}}
\pgfpathlineto{\pgfpoint{-2.914530\du}{6.440660\du}}
\pgfpathlineto{\pgfpoint{-2.904379\du}{6.440062\du}}
\pgfpathlineto{\pgfpoint{-2.894825\du}{6.438868\du}}
\pgfpathlineto{\pgfpoint{-2.885570\du}{6.437973\du}}
\pgfpathlineto{\pgfpoint{-2.876315\du}{6.436778\du}}
\pgfpathlineto{\pgfpoint{-2.867060\du}{6.435286\du}}
\pgfpathlineto{\pgfpoint{-2.858402\du}{6.433793\du}}
\pgfpathlineto{\pgfpoint{-2.850043\du}{6.432300\du}}
\pgfpathlineto{\pgfpoint{-2.841385\du}{6.430210\du}}
\pgfpathlineto{\pgfpoint{-2.833324\du}{6.428120\du}}
\pgfpathlineto{\pgfpoint{-2.825860\du}{6.426030\du}}
\pgfpathlineto{\pgfpoint{-2.818396\du}{6.424239\du}}
\pgfpathlineto{\pgfpoint{-2.811529\du}{6.421851\du}}
\pgfpathlineto{\pgfpoint{-2.804663\du}{6.419462\du}}
\pgfpathlineto{\pgfpoint{-2.798692\du}{6.416775\du}}
\pgfpathlineto{\pgfpoint{-2.792422\du}{6.414088\du}}
\pgfpathlineto{\pgfpoint{-2.787048\du}{6.411401\du}}
\pgfpathlineto{\pgfpoint{-2.781973\du}{6.408714\du}}
\pgfpathlineto{\pgfpoint{-2.777196\du}{6.405430\du}}
\pgfpathlineto{\pgfpoint{-2.773613\du}{6.402445\du}}
\pgfpathlineto{\pgfpoint{-2.769434\du}{6.399459\du}}
\pgfpathlineto{\pgfpoint{-2.766448\du}{6.396474\du}}
\pgfpathlineto{\pgfpoint{-2.763761\du}{6.393190\du}}
\pgfpathlineto{\pgfpoint{-2.761373\du}{6.389906\du}}
\pgfpathlineto{\pgfpoint{-2.759880\du}{6.386622\du}}
\pgfpathlineto{\pgfpoint{-2.758686\du}{6.383338\du}}
\pgfpathlineto{\pgfpoint{-2.757790\du}{6.380352\du}}
\pgfpathlineto{\pgfpoint{-2.757790\du}{6.376769\du}}
\pgfusepath{fill}
\pgfsetbuttcap
\pgfsetmiterjoin
\pgfsetdash{}{0pt}
\definecolor{dialinecolor}{rgb}{0.666667, 0.901961, 1.000000}
\pgfsetstrokecolor{dialinecolor}
\pgfpathmoveto{\pgfpoint{-2.769434\du}{6.371097\du}}
\pgfpathlineto{\pgfpoint{-2.769434\du}{6.367514\du}}
\pgfpathlineto{\pgfpoint{-2.770329\du}{6.365126\du}}
\pgfpathlineto{\pgfpoint{-2.771524\du}{6.361842\du}}
\pgfpathlineto{\pgfpoint{-2.773613\du}{6.358856\du}}
\pgfpathlineto{\pgfpoint{-2.775703\du}{6.356169\du}}
\pgfpathlineto{\pgfpoint{-2.777793\du}{6.352587\du}}
\pgfpathlineto{\pgfpoint{-2.781077\du}{6.350198\du}}
\pgfpathlineto{\pgfpoint{-2.784361\du}{6.347511\du}}
\pgfpathlineto{\pgfpoint{-2.788541\du}{6.344824\du}}
\pgfpathlineto{\pgfpoint{-2.793318\du}{6.342436\du}}
\pgfpathlineto{\pgfpoint{-2.797796\du}{6.339450\du}}
\pgfpathlineto{\pgfpoint{-2.803469\du}{6.337361\du}}
\pgfpathlineto{\pgfpoint{-2.809440\du}{6.334674\du}}
\pgfpathlineto{\pgfpoint{-2.815112\du}{6.331987\du}}
\pgfpathlineto{\pgfpoint{-2.821680\du}{6.330195\du}}
\pgfpathlineto{\pgfpoint{-2.828547\du}{6.327807\du}}
\pgfpathlineto{\pgfpoint{-2.835414\du}{6.326314\du}}
\pgfpathlineto{\pgfpoint{-2.843176\du}{6.324224\du}}
\pgfpathlineto{\pgfpoint{-2.850640\du}{6.322433\du}}
\pgfpathlineto{\pgfpoint{-2.858999\du}{6.320642\du}}
\pgfpathlineto{\pgfpoint{-2.867657\du}{6.319447\du}}
\pgfpathlineto{\pgfpoint{-2.876017\du}{6.317059\du}}
\pgfpathlineto{\pgfpoint{-2.884675\du}{6.316760\du}}
\pgfpathlineto{\pgfpoint{-2.893631\du}{6.315268\du}}
\pgfpathlineto{\pgfpoint{-2.902289\du}{6.314372\du}}
\pgfpathlineto{\pgfpoint{-2.911843\du}{6.313476\du}}
\pgfpathlineto{\pgfpoint{-2.921098\du}{6.312879\du}}
\pgfpathlineto{\pgfpoint{-2.930652\du}{6.312581\du}}
\pgfpathlineto{\pgfpoint{-2.940205\du}{6.311984\du}}
\pgfpathlineto{\pgfpoint{-2.950356\du}{6.311685\du}}
\pgfpathlineto{\pgfpoint{-2.960208\du}{6.311685\du}}
\pgfpathlineto{\pgfpoint{-2.960208\du}{6.311685\du}}
\pgfpathlineto{\pgfpoint{-2.969165\du}{6.311685\du}}
\pgfpathlineto{\pgfpoint{-2.979017\du}{6.311984\du}}
\pgfpathlineto{\pgfpoint{-2.988571\du}{6.312581\du}}
\pgfpathlineto{\pgfpoint{-2.997826\du}{6.312879\du}}
\pgfpathlineto{\pgfpoint{-3.007380\du}{6.313476\du}}
\pgfpathlineto{\pgfpoint{-3.016933\du}{6.314372\du}}
\pgfpathlineto{\pgfpoint{-3.025890\du}{6.315268\du}}
\pgfpathlineto{\pgfpoint{-3.034846\du}{6.316760\du}}
\pgfpathlineto{\pgfpoint{-3.043504\du}{6.317059\du}}
\pgfpathlineto{\pgfpoint{-3.051864\du}{6.319447\du}}
\pgfpathlineto{\pgfpoint{-3.060522\du}{6.320642\du}}
\pgfpathlineto{\pgfpoint{-3.068583\du}{6.322433\du}}
\pgfpathlineto{\pgfpoint{-3.076345\du}{6.324224\du}}
\pgfpathlineto{\pgfpoint{-3.083809\du}{6.326314\du}}
\pgfpathlineto{\pgfpoint{-3.090377\du}{6.327807\du}}
\pgfpathlineto{\pgfpoint{-3.097841\du}{6.330195\du}}
\pgfpathlineto{\pgfpoint{-3.104110\du}{6.331987\du}}
\pgfpathlineto{\pgfpoint{-3.110082\du}{6.334674\du}}
\pgfpathlineto{\pgfpoint{-3.116053\du}{6.337361\du}}
\pgfpathlineto{\pgfpoint{-3.121128\du}{6.339450\du}}
\pgfpathlineto{\pgfpoint{-3.126203\du}{6.342436\du}}
\pgfpathlineto{\pgfpoint{-3.130980\du}{6.344824\du}}
\pgfpathlineto{\pgfpoint{-3.135160\du}{6.347511\du}}
\pgfpathlineto{\pgfpoint{-3.138444\du}{6.350198\du}}
\pgfpathlineto{\pgfpoint{-3.141728\du}{6.352587\du}}
\pgfpathlineto{\pgfpoint{-3.144116\du}{6.356169\du}}
\pgfpathlineto{\pgfpoint{-3.145908\du}{6.358856\du}}
\pgfpathlineto{\pgfpoint{-3.148296\du}{6.361842\du}}
\pgfpathlineto{\pgfpoint{-3.149192\du}{6.365126\du}}
\pgfpathlineto{\pgfpoint{-3.150087\du}{6.367514\du}}
\pgfpathlineto{\pgfpoint{-3.150087\du}{6.371097\du}}
\pgfpathlineto{\pgfpoint{-3.150087\du}{6.371097\du}}
\pgfpathlineto{\pgfpoint{-3.150087\du}{6.374082\du}}
\pgfpathlineto{\pgfpoint{-3.149192\du}{6.376769\du}}
\pgfpathlineto{\pgfpoint{-3.148296\du}{6.379755\du}}
\pgfpathlineto{\pgfpoint{-3.145908\du}{6.383338\du}}
\pgfpathlineto{\pgfpoint{-3.144116\du}{6.386024\du}}
\pgfpathlineto{\pgfpoint{-3.141728\du}{6.389010\du}}
\pgfpathlineto{\pgfpoint{-3.138444\du}{6.391697\du}}
\pgfpathlineto{\pgfpoint{-3.135160\du}{6.394384\du}}
\pgfpathlineto{\pgfpoint{-3.130980\du}{6.397071\du}}
\pgfpathlineto{\pgfpoint{-3.126203\du}{6.399758\du}}
\pgfpathlineto{\pgfpoint{-3.121128\du}{6.402445\du}}
\pgfpathlineto{\pgfpoint{-3.116053\du}{6.404833\du}}
\pgfpathlineto{\pgfpoint{-3.110082\du}{6.407520\du}}
\pgfpathlineto{\pgfpoint{-3.104110\du}{6.409610\du}}
\pgfpathlineto{\pgfpoint{-3.097841\du}{6.411999\du}}
\pgfpathlineto{\pgfpoint{-3.090377\du}{6.414387\du}}
\pgfpathlineto{\pgfpoint{-3.083809\du}{6.416178\du}}
\pgfpathlineto{\pgfpoint{-3.076345\du}{6.417671\du}}
\pgfpathlineto{\pgfpoint{-3.068583\du}{6.419761\du}}
\pgfpathlineto{\pgfpoint{-3.060522\du}{6.421552\du}}
\pgfpathlineto{\pgfpoint{-3.051864\du}{6.422746\du}}
\pgfpathlineto{\pgfpoint{-3.043504\du}{6.424538\du}}
\pgfpathlineto{\pgfpoint{-3.034846\du}{6.425433\du}}
\pgfpathlineto{\pgfpoint{-3.025890\du}{6.426926\du}}
\pgfpathlineto{\pgfpoint{-3.016933\du}{6.427822\du}}
\pgfpathlineto{\pgfpoint{-3.007380\du}{6.428120\du}}
\pgfpathlineto{\pgfpoint{-2.997826\du}{6.429315\du}}
\pgfpathlineto{\pgfpoint{-2.988571\du}{6.429613\du}}
\pgfpathlineto{\pgfpoint{-2.979017\du}{6.429912\du}}
\pgfpathlineto{\pgfpoint{-2.969165\du}{6.430509\du}}
\pgfpathlineto{\pgfpoint{-2.960208\du}{6.430509\du}}
\pgfpathlineto{\pgfpoint{-2.960208\du}{6.430509\du}}
\pgfpathlineto{\pgfpoint{-2.950356\du}{6.430509\du}}
\pgfpathlineto{\pgfpoint{-2.940205\du}{6.429912\du}}
\pgfpathlineto{\pgfpoint{-2.930652\du}{6.429613\du}}
\pgfpathlineto{\pgfpoint{-2.921098\du}{6.429315\du}}
\pgfpathlineto{\pgfpoint{-2.911843\du}{6.428120\du}}
\pgfpathlineto{\pgfpoint{-2.902289\du}{6.427822\du}}
\pgfpathlineto{\pgfpoint{-2.893631\du}{6.426926\du}}
\pgfpathlineto{\pgfpoint{-2.884675\du}{6.425433\du}}
\pgfpathlineto{\pgfpoint{-2.876017\du}{6.424538\du}}
\pgfpathlineto{\pgfpoint{-2.867657\du}{6.422746\du}}
\pgfpathlineto{\pgfpoint{-2.858999\du}{6.421552\du}}
\pgfpathlineto{\pgfpoint{-2.850640\du}{6.419761\du}}
\pgfpathlineto{\pgfpoint{-2.843176\du}{6.417671\du}}
\pgfpathlineto{\pgfpoint{-2.835414\du}{6.416178\du}}
\pgfpathlineto{\pgfpoint{-2.828547\du}{6.414387\du}}
\pgfpathlineto{\pgfpoint{-2.821680\du}{6.411999\du}}
\pgfpathlineto{\pgfpoint{-2.815112\du}{6.409610\du}}
\pgfpathlineto{\pgfpoint{-2.809440\du}{6.407520\du}}
\pgfpathlineto{\pgfpoint{-2.803469\du}{6.404833\du}}
\pgfpathlineto{\pgfpoint{-2.797796\du}{6.402445\du}}
\pgfpathlineto{\pgfpoint{-2.793318\du}{6.399758\du}}
\pgfpathlineto{\pgfpoint{-2.788541\du}{6.397071\du}}
\pgfpathlineto{\pgfpoint{-2.784361\du}{6.394384\du}}
\pgfpathlineto{\pgfpoint{-2.781077\du}{6.391697\du}}
\pgfpathlineto{\pgfpoint{-2.777793\du}{6.389010\du}}
\pgfpathlineto{\pgfpoint{-2.775703\du}{6.386024\du}}
\pgfpathlineto{\pgfpoint{-2.773613\du}{6.383338\du}}
\pgfpathlineto{\pgfpoint{-2.771524\du}{6.379755\du}}
\pgfpathlineto{\pgfpoint{-2.770329\du}{6.376769\du}}
\pgfpathlineto{\pgfpoint{-2.769434\du}{6.374082\du}}
\pgfpathlineto{\pgfpoint{-2.769434\du}{6.371097\du}}
\pgfusepath{stroke}
\pgfsetbuttcap
\pgfsetmiterjoin
\pgfsetdash{}{0pt}
\definecolor{dialinecolor}{rgb}{0.000000, 0.470588, 0.666667}
\pgfsetfillcolor{dialinecolor}
\pgfpathmoveto{\pgfpoint{-3.150087\du}{6.287801\du}}
\pgfpathlineto{\pgfpoint{-3.150087\du}{6.382442\du}}
\pgfpathlineto{\pgfpoint{-2.757193\du}{6.382442\du}}
\pgfpathlineto{\pgfpoint{-2.757193\du}{6.287801\du}}
\pgfpathlineto{\pgfpoint{-3.150087\du}{6.287801\du}}
\pgfusepath{fill}
\pgfsetbuttcap
\pgfsetmiterjoin
\pgfsetdash{}{0pt}
\definecolor{dialinecolor}{rgb}{0.000000, 0.705882, 1.000000}
\pgfsetfillcolor{dialinecolor}
\pgfpathmoveto{\pgfpoint{-2.757790\du}{6.281233\du}}
\pgfpathlineto{\pgfpoint{-2.757790\du}{6.277949\du}}
\pgfpathlineto{\pgfpoint{-2.758686\du}{6.274366\du}}
\pgfpathlineto{\pgfpoint{-2.759880\du}{6.271381\du}}
\pgfpathlineto{\pgfpoint{-2.761373\du}{6.268096\du}}
\pgfpathlineto{\pgfpoint{-2.763761\du}{6.264812\du}}
\pgfpathlineto{\pgfpoint{-2.766448\du}{6.261528\du}}
\pgfpathlineto{\pgfpoint{-2.769434\du}{6.258543\du}}
\pgfpathlineto{\pgfpoint{-2.773613\du}{6.255259\du}}
\pgfpathlineto{\pgfpoint{-2.777196\du}{6.252572\du}}
\pgfpathlineto{\pgfpoint{-2.781973\du}{6.249586\du}}
\pgfpathlineto{\pgfpoint{-2.787048\du}{6.246601\du}}
\pgfpathlineto{\pgfpoint{-2.792422\du}{6.243914\du}}
\pgfpathlineto{\pgfpoint{-2.798692\du}{6.240630\du}}
\pgfpathlineto{\pgfpoint{-2.804663\du}{6.238838\du}}
\pgfpathlineto{\pgfpoint{-2.811529\du}{6.236151\du}}
\pgfpathlineto{\pgfpoint{-2.818396\du}{6.233763\du}}
\pgfpathlineto{\pgfpoint{-2.825860\du}{6.231673\du}}
\pgfpathlineto{\pgfpoint{-2.833324\du}{6.229583\du}}
\pgfpathlineto{\pgfpoint{-2.841385\du}{6.227493\du}}
\pgfpathlineto{\pgfpoint{-2.850043\du}{6.225404\du}}
\pgfpathlineto{\pgfpoint{-2.858402\du}{6.223911\du}}
\pgfpathlineto{\pgfpoint{-2.867060\du}{6.222418\du}}
\pgfpathlineto{\pgfpoint{-2.876315\du}{6.221224\du}}
\pgfpathlineto{\pgfpoint{-2.885570\du}{6.219731\du}}
\pgfpathlineto{\pgfpoint{-2.894825\du}{6.218835\du}}
\pgfpathlineto{\pgfpoint{-2.904379\du}{6.217641\du}}
\pgfpathlineto{\pgfpoint{-2.914530\du}{6.217044\du}}
\pgfpathlineto{\pgfpoint{-2.924084\du}{6.216447\du}}
\pgfpathlineto{\pgfpoint{-2.933637\du}{6.216148\du}}
\pgfpathlineto{\pgfpoint{-2.943788\du}{6.215850\du}}
\pgfpathlineto{\pgfpoint{-2.953640\du}{6.215850\du}}
\pgfpathlineto{\pgfpoint{-2.953640\du}{6.215850\du}}
\pgfpathlineto{\pgfpoint{-2.963492\du}{6.215850\du}}
\pgfpathlineto{\pgfpoint{-2.973643\du}{6.216148\du}}
\pgfpathlineto{\pgfpoint{-2.983495\du}{6.216447\du}}
\pgfpathlineto{\pgfpoint{-2.993348\du}{6.217044\du}}
\pgfpathlineto{\pgfpoint{-3.002901\du}{6.217641\du}}
\pgfpathlineto{\pgfpoint{-3.012455\du}{6.218835\du}}
\pgfpathlineto{\pgfpoint{-3.022009\du}{6.219731\du}}
\pgfpathlineto{\pgfpoint{-3.030965\du}{6.221224\du}}
\pgfpathlineto{\pgfpoint{-3.040220\du}{6.222418\du}}
\pgfpathlineto{\pgfpoint{-3.048878\du}{6.223911\du}}
\pgfpathlineto{\pgfpoint{-3.057536\du}{6.225404\du}}
\pgfpathlineto{\pgfpoint{-3.065896\du}{6.227493\du}}
\pgfpathlineto{\pgfpoint{-3.073957\du}{6.229583\du}}
\pgfpathlineto{\pgfpoint{-3.081421\du}{6.231673\du}}
\pgfpathlineto{\pgfpoint{-3.088884\du}{6.233763\du}}
\pgfpathlineto{\pgfpoint{-3.096050\du}{6.236151\du}}
\pgfpathlineto{\pgfpoint{-3.102618\du}{6.238838\du}}
\pgfpathlineto{\pgfpoint{-3.108887\du}{6.240630\du}}
\pgfpathlineto{\pgfpoint{-3.114858\du}{6.243914\du}}
\pgfpathlineto{\pgfpoint{-3.120232\du}{6.246601\du}}
\pgfpathlineto{\pgfpoint{-3.125308\du}{6.249586\du}}
\pgfpathlineto{\pgfpoint{-3.130084\du}{6.252572\du}}
\pgfpathlineto{\pgfpoint{-3.133966\du}{6.255259\du}}
\pgfpathlineto{\pgfpoint{-3.137847\du}{6.258543\du}}
\pgfpathlineto{\pgfpoint{-3.140832\du}{6.261528\du}}
\pgfpathlineto{\pgfpoint{-3.143519\du}{6.264812\du}}
\pgfpathlineto{\pgfpoint{-3.145908\du}{6.268096\du}}
\pgfpathlineto{\pgfpoint{-3.147699\du}{6.271381\du}}
\pgfpathlineto{\pgfpoint{-3.149192\du}{6.274366\du}}
\pgfpathlineto{\pgfpoint{-3.150087\du}{6.277949\du}}
\pgfpathlineto{\pgfpoint{-3.150087\du}{6.281233\du}}
\pgfpathlineto{\pgfpoint{-3.150087\du}{6.281233\du}}
\pgfpathlineto{\pgfpoint{-3.150087\du}{6.284517\du}}
\pgfpathlineto{\pgfpoint{-3.149192\du}{6.288099\du}}
\pgfpathlineto{\pgfpoint{-3.147699\du}{6.291085\du}}
\pgfpathlineto{\pgfpoint{-3.145908\du}{6.294369\du}}
\pgfpathlineto{\pgfpoint{-3.143519\du}{6.297952\du}}
\pgfpathlineto{\pgfpoint{-3.140832\du}{6.300937\du}}
\pgfpathlineto{\pgfpoint{-3.137847\du}{6.303923\du}}
\pgfpathlineto{\pgfpoint{-3.133966\du}{6.307207\du}}
\pgfpathlineto{\pgfpoint{-3.130084\du}{6.309595\du}}
\pgfpathlineto{\pgfpoint{-3.125308\du}{6.312879\du}}
\pgfpathlineto{\pgfpoint{-3.120232\du}{6.316163\du}}
\pgfpathlineto{\pgfpoint{-3.114858\du}{6.318850\du}}
\pgfpathlineto{\pgfpoint{-3.108887\du}{6.321537\du}}
\pgfpathlineto{\pgfpoint{-3.102618\du}{6.323627\du}}
\pgfpathlineto{\pgfpoint{-3.096050\du}{6.326613\du}}
\pgfpathlineto{\pgfpoint{-3.088884\du}{6.329001\du}}
\pgfpathlineto{\pgfpoint{-3.081421\du}{6.330792\du}}
\pgfpathlineto{\pgfpoint{-3.073957\du}{6.332882\du}}
\pgfpathlineto{\pgfpoint{-3.065896\du}{6.334972\du}}
\pgfpathlineto{\pgfpoint{-3.057536\du}{6.337062\du}}
\pgfpathlineto{\pgfpoint{-3.048878\du}{6.338555\du}}
\pgfpathlineto{\pgfpoint{-3.040220\du}{6.340047\du}}
\pgfpathlineto{\pgfpoint{-3.030965\du}{6.341242\du}}
\pgfpathlineto{\pgfpoint{-3.022009\du}{6.342734\du}}
\pgfpathlineto{\pgfpoint{-3.012455\du}{6.343630\du}}
\pgfpathlineto{\pgfpoint{-3.002901\du}{6.344824\du}}
\pgfpathlineto{\pgfpoint{-2.993348\du}{6.345123\du}}
\pgfpathlineto{\pgfpoint{-2.983495\du}{6.346019\du}}
\pgfpathlineto{\pgfpoint{-2.973643\du}{6.346019\du}}
\pgfpathlineto{\pgfpoint{-2.963492\du}{6.346914\du}}
\pgfpathlineto{\pgfpoint{-2.953640\du}{6.346914\du}}
\pgfpathlineto{\pgfpoint{-2.953640\du}{6.346914\du}}
\pgfpathlineto{\pgfpoint{-2.943788\du}{6.346914\du}}
\pgfpathlineto{\pgfpoint{-2.933637\du}{6.346019\du}}
\pgfpathlineto{\pgfpoint{-2.924084\du}{6.346019\du}}
\pgfpathlineto{\pgfpoint{-2.914530\du}{6.345123\du}}
\pgfpathlineto{\pgfpoint{-2.904379\du}{6.344824\du}}
\pgfpathlineto{\pgfpoint{-2.894825\du}{6.343630\du}}
\pgfpathlineto{\pgfpoint{-2.885570\du}{6.342734\du}}
\pgfpathlineto{\pgfpoint{-2.876315\du}{6.341242\du}}
\pgfpathlineto{\pgfpoint{-2.867060\du}{6.340047\du}}
\pgfpathlineto{\pgfpoint{-2.858402\du}{6.338555\du}}
\pgfpathlineto{\pgfpoint{-2.850043\du}{6.337062\du}}
\pgfpathlineto{\pgfpoint{-2.841385\du}{6.334972\du}}
\pgfpathlineto{\pgfpoint{-2.833324\du}{6.332882\du}}
\pgfpathlineto{\pgfpoint{-2.825860\du}{6.330792\du}}
\pgfpathlineto{\pgfpoint{-2.818396\du}{6.329001\du}}
\pgfpathlineto{\pgfpoint{-2.811529\du}{6.326613\du}}
\pgfpathlineto{\pgfpoint{-2.804663\du}{6.323627\du}}
\pgfpathlineto{\pgfpoint{-2.798692\du}{6.321537\du}}
\pgfpathlineto{\pgfpoint{-2.792422\du}{6.318850\du}}
\pgfpathlineto{\pgfpoint{-2.787048\du}{6.316163\du}}
\pgfpathlineto{\pgfpoint{-2.781973\du}{6.312879\du}}
\pgfpathlineto{\pgfpoint{-2.777196\du}{6.309595\du}}
\pgfpathlineto{\pgfpoint{-2.773613\du}{6.307207\du}}
\pgfpathlineto{\pgfpoint{-2.769434\du}{6.303923\du}}
\pgfpathlineto{\pgfpoint{-2.766448\du}{6.300937\du}}
\pgfpathlineto{\pgfpoint{-2.763761\du}{6.297952\du}}
\pgfpathlineto{\pgfpoint{-2.761373\du}{6.294369\du}}
\pgfpathlineto{\pgfpoint{-2.759880\du}{6.291085\du}}
\pgfpathlineto{\pgfpoint{-2.758686\du}{6.288099\du}}
\pgfpathlineto{\pgfpoint{-2.757790\du}{6.284517\du}}
\pgfpathlineto{\pgfpoint{-2.757790\du}{6.281233\du}}
\pgfusepath{fill}
\pgfsetbuttcap
\pgfsetmiterjoin
\pgfsetdash{}{0pt}
\definecolor{dialinecolor}{rgb}{0.666667, 0.901961, 1.000000}
\pgfsetstrokecolor{dialinecolor}
\pgfpathmoveto{\pgfpoint{-2.769434\du}{6.275560\du}}
\pgfpathlineto{\pgfpoint{-2.769434\du}{6.272276\du}}
\pgfpathlineto{\pgfpoint{-2.770329\du}{6.269888\du}}
\pgfpathlineto{\pgfpoint{-2.771524\du}{6.266604\du}}
\pgfpathlineto{\pgfpoint{-2.773613\du}{6.263021\du}}
\pgfpathlineto{\pgfpoint{-2.775703\du}{6.260633\du}}
\pgfpathlineto{\pgfpoint{-2.777793\du}{6.257349\du}}
\pgfpathlineto{\pgfpoint{-2.781077\du}{6.254960\du}}
\pgfpathlineto{\pgfpoint{-2.784361\du}{6.252273\du}}
\pgfpathlineto{\pgfpoint{-2.788541\du}{6.249586\du}}
\pgfpathlineto{\pgfpoint{-2.793318\du}{6.246899\du}}
\pgfpathlineto{\pgfpoint{-2.797796\du}{6.244212\du}}
\pgfpathlineto{\pgfpoint{-2.803469\du}{6.241525\du}}
\pgfpathlineto{\pgfpoint{-2.809440\du}{6.239137\du}}
\pgfpathlineto{\pgfpoint{-2.815112\du}{6.237047\du}}
\pgfpathlineto{\pgfpoint{-2.821680\du}{6.234659\du}}
\pgfpathlineto{\pgfpoint{-2.828547\du}{6.232270\du}}
\pgfpathlineto{\pgfpoint{-2.835414\du}{6.230479\du}}
\pgfpathlineto{\pgfpoint{-2.843176\du}{6.228986\du}}
\pgfpathlineto{\pgfpoint{-2.850640\du}{6.226598\du}}
\pgfpathlineto{\pgfpoint{-2.858999\du}{6.225105\du}}
\pgfpathlineto{\pgfpoint{-2.867657\du}{6.223911\du}}
\pgfpathlineto{\pgfpoint{-2.876017\du}{6.222119\du}}
\pgfpathlineto{\pgfpoint{-2.884675\du}{6.221224\du}}
\pgfpathlineto{\pgfpoint{-2.893631\du}{6.219731\du}}
\pgfpathlineto{\pgfpoint{-2.902289\du}{6.218835\du}}
\pgfpathlineto{\pgfpoint{-2.911843\du}{6.218238\du}}
\pgfpathlineto{\pgfpoint{-2.921098\du}{6.217343\du}}
\pgfpathlineto{\pgfpoint{-2.930652\du}{6.217044\du}}
\pgfpathlineto{\pgfpoint{-2.940205\du}{6.216745\du}}
\pgfpathlineto{\pgfpoint{-2.950356\du}{6.216148\du}}
\pgfpathlineto{\pgfpoint{-2.960208\du}{6.216148\du}}
\pgfpathlineto{\pgfpoint{-2.960208\du}{6.216148\du}}
\pgfpathlineto{\pgfpoint{-2.969165\du}{6.216148\du}}
\pgfpathlineto{\pgfpoint{-2.979017\du}{6.216745\du}}
\pgfpathlineto{\pgfpoint{-2.988571\du}{6.217044\du}}
\pgfpathlineto{\pgfpoint{-2.997826\du}{6.217343\du}}
\pgfpathlineto{\pgfpoint{-3.007380\du}{6.218238\du}}
\pgfpathlineto{\pgfpoint{-3.016933\du}{6.218835\du}}
\pgfpathlineto{\pgfpoint{-3.025890\du}{6.219731\du}}
\pgfpathlineto{\pgfpoint{-3.034846\du}{6.221224\du}}
\pgfpathlineto{\pgfpoint{-3.043504\du}{6.222119\du}}
\pgfpathlineto{\pgfpoint{-3.051864\du}{6.223911\du}}
\pgfpathlineto{\pgfpoint{-3.060522\du}{6.225105\du}}
\pgfpathlineto{\pgfpoint{-3.068583\du}{6.226598\du}}
\pgfpathlineto{\pgfpoint{-3.076345\du}{6.228986\du}}
\pgfpathlineto{\pgfpoint{-3.083809\du}{6.230479\du}}
\pgfpathlineto{\pgfpoint{-3.090377\du}{6.232270\du}}
\pgfpathlineto{\pgfpoint{-3.097841\du}{6.234659\du}}
\pgfpathlineto{\pgfpoint{-3.104110\du}{6.237047\du}}
\pgfpathlineto{\pgfpoint{-3.110082\du}{6.239137\du}}
\pgfpathlineto{\pgfpoint{-3.116053\du}{6.241525\du}}
\pgfpathlineto{\pgfpoint{-3.121128\du}{6.244212\du}}
\pgfpathlineto{\pgfpoint{-3.126203\du}{6.246899\du}}
\pgfpathlineto{\pgfpoint{-3.130980\du}{6.249586\du}}
\pgfpathlineto{\pgfpoint{-3.135160\du}{6.252273\du}}
\pgfpathlineto{\pgfpoint{-3.138444\du}{6.254960\du}}
\pgfpathlineto{\pgfpoint{-3.141728\du}{6.257349\du}}
\pgfpathlineto{\pgfpoint{-3.144116\du}{6.260633\du}}
\pgfpathlineto{\pgfpoint{-3.145908\du}{6.263021\du}}
\pgfpathlineto{\pgfpoint{-3.148296\du}{6.266604\du}}
\pgfpathlineto{\pgfpoint{-3.149192\du}{6.269888\du}}
\pgfpathlineto{\pgfpoint{-3.150087\du}{6.272276\du}}
\pgfpathlineto{\pgfpoint{-3.150087\du}{6.275560\du}}
\pgfpathlineto{\pgfpoint{-3.150087\du}{6.275560\du}}
\pgfpathlineto{\pgfpoint{-3.150087\du}{6.278844\du}}
\pgfpathlineto{\pgfpoint{-3.149192\du}{6.281531\du}}
\pgfpathlineto{\pgfpoint{-3.148296\du}{6.284517\du}}
\pgfpathlineto{\pgfpoint{-3.145908\du}{6.287801\du}}
\pgfpathlineto{\pgfpoint{-3.144116\du}{6.290488\du}}
\pgfpathlineto{\pgfpoint{-3.141728\du}{6.293772\du}}
\pgfpathlineto{\pgfpoint{-3.138444\du}{6.296459\du}}
\pgfpathlineto{\pgfpoint{-3.135160\du}{6.299146\du}}
\pgfpathlineto{\pgfpoint{-3.130980\du}{6.301833\du}}
\pgfpathlineto{\pgfpoint{-3.126203\du}{6.304520\du}}
\pgfpathlineto{\pgfpoint{-3.121128\du}{6.307207\du}}
\pgfpathlineto{\pgfpoint{-3.116053\du}{6.309297\du}}
\pgfpathlineto{\pgfpoint{-3.110082\du}{6.311984\du}}
\pgfpathlineto{\pgfpoint{-3.104110\du}{6.314372\du}}
\pgfpathlineto{\pgfpoint{-3.097841\du}{6.316760\du}}
\pgfpathlineto{\pgfpoint{-3.090377\du}{6.318850\du}}
\pgfpathlineto{\pgfpoint{-3.083809\du}{6.320642\du}}
\pgfpathlineto{\pgfpoint{-3.076345\du}{6.322433\du}}
\pgfpathlineto{\pgfpoint{-3.068583\du}{6.324224\du}}
\pgfpathlineto{\pgfpoint{-3.060522\du}{6.326314\du}}
\pgfpathlineto{\pgfpoint{-3.051864\du}{6.327508\du}}
\pgfpathlineto{\pgfpoint{-3.043504\du}{6.329300\du}}
\pgfpathlineto{\pgfpoint{-3.034846\du}{6.330195\du}}
\pgfpathlineto{\pgfpoint{-3.025890\du}{6.331688\du}}
\pgfpathlineto{\pgfpoint{-3.016933\du}{6.331987\du}}
\pgfpathlineto{\pgfpoint{-3.007380\du}{6.333181\du}}
\pgfpathlineto{\pgfpoint{-2.997826\du}{6.333479\du}}
\pgfpathlineto{\pgfpoint{-2.988571\du}{6.334375\du}}
\pgfpathlineto{\pgfpoint{-2.979017\du}{6.334674\du}}
\pgfpathlineto{\pgfpoint{-2.969165\du}{6.334972\du}}
\pgfpathlineto{\pgfpoint{-2.960208\du}{6.334972\du}}
\pgfpathlineto{\pgfpoint{-2.960208\du}{6.334972\du}}
\pgfpathlineto{\pgfpoint{-2.950356\du}{6.334972\du}}
\pgfpathlineto{\pgfpoint{-2.940205\du}{6.334674\du}}
\pgfpathlineto{\pgfpoint{-2.930652\du}{6.334375\du}}
\pgfpathlineto{\pgfpoint{-2.921098\du}{6.333479\du}}
\pgfpathlineto{\pgfpoint{-2.911843\du}{6.333181\du}}
\pgfpathlineto{\pgfpoint{-2.902289\du}{6.331987\du}}
\pgfpathlineto{\pgfpoint{-2.893631\du}{6.331688\du}}
\pgfpathlineto{\pgfpoint{-2.884675\du}{6.330195\du}}
\pgfpathlineto{\pgfpoint{-2.876017\du}{6.329300\du}}
\pgfpathlineto{\pgfpoint{-2.867657\du}{6.327508\du}}
\pgfpathlineto{\pgfpoint{-2.858999\du}{6.326314\du}}
\pgfpathlineto{\pgfpoint{-2.850640\du}{6.324224\du}}
\pgfpathlineto{\pgfpoint{-2.843176\du}{6.322433\du}}
\pgfpathlineto{\pgfpoint{-2.835414\du}{6.320642\du}}
\pgfpathlineto{\pgfpoint{-2.828547\du}{6.318850\du}}
\pgfpathlineto{\pgfpoint{-2.821680\du}{6.316760\du}}
\pgfpathlineto{\pgfpoint{-2.815112\du}{6.314372\du}}
\pgfpathlineto{\pgfpoint{-2.809440\du}{6.311984\du}}
\pgfpathlineto{\pgfpoint{-2.803469\du}{6.309297\du}}
\pgfpathlineto{\pgfpoint{-2.797796\du}{6.307207\du}}
\pgfpathlineto{\pgfpoint{-2.793318\du}{6.304520\du}}
\pgfpathlineto{\pgfpoint{-2.788541\du}{6.301833\du}}
\pgfpathlineto{\pgfpoint{-2.784361\du}{6.299146\du}}
\pgfpathlineto{\pgfpoint{-2.781077\du}{6.296459\du}}
\pgfpathlineto{\pgfpoint{-2.777793\du}{6.293772\du}}
\pgfpathlineto{\pgfpoint{-2.775703\du}{6.290488\du}}
\pgfpathlineto{\pgfpoint{-2.773613\du}{6.287801\du}}
\pgfpathlineto{\pgfpoint{-2.771524\du}{6.284517\du}}
\pgfpathlineto{\pgfpoint{-2.770329\du}{6.281531\du}}
\pgfpathlineto{\pgfpoint{-2.769434\du}{6.278844\du}}
\pgfpathlineto{\pgfpoint{-2.769434\du}{6.275560\du}}
\pgfusepath{stroke}
\pgfsetbuttcap
\pgfsetmiterjoin
\pgfsetdash{}{0pt}
\definecolor{dialinecolor}{rgb}{0.000000, 0.000000, 0.000000}
\pgfsetfillcolor{dialinecolor}
\pgfpathmoveto{\pgfpoint{-2.947968\du}{6.263917\du}}
\pgfpathlineto{\pgfpoint{-2.924084\du}{6.275560\du}}
\pgfpathlineto{\pgfpoint{-2.852730\du}{6.251975\du}}
\pgfpathlineto{\pgfpoint{-2.816903\du}{6.263917\du}}
\pgfpathlineto{\pgfpoint{-2.840788\du}{6.240033\du}}
\pgfpathlineto{\pgfpoint{-2.924084\du}{6.240033\du}}
\pgfpathlineto{\pgfpoint{-2.888257\du}{6.240033\du}}
\pgfpathlineto{\pgfpoint{-2.947968\du}{6.263917\du}}
\pgfusepath{fill}
\pgfsetbuttcap
\pgfsetmiterjoin
\pgfsetdash{}{0pt}
\definecolor{dialinecolor}{rgb}{0.000000, 0.000000, 0.000000}
\pgfsetfillcolor{dialinecolor}
\pgfpathmoveto{\pgfpoint{-2.960208\du}{6.299444\du}}
\pgfpathlineto{\pgfpoint{-2.995139\du}{6.287801\du}}
\pgfpathlineto{\pgfpoint{-3.055148\du}{6.311685\du}}
\pgfpathlineto{\pgfpoint{-3.090377\du}{6.299444\du}}
\pgfpathlineto{\pgfpoint{-3.079032\du}{6.334972\du}}
\pgfpathlineto{\pgfpoint{-2.995139\du}{6.334972\du}}
\pgfpathlineto{\pgfpoint{-3.031264\du}{6.323627\du}}
\pgfpathlineto{\pgfpoint{-2.960208\du}{6.299444\du}}
\pgfusepath{fill}
\pgfsetbuttcap
\pgfsetmiterjoin
\pgfsetdash{}{0pt}
\definecolor{dialinecolor}{rgb}{0.000000, 0.000000, 0.000000}
\pgfsetfillcolor{dialinecolor}
\pgfpathmoveto{\pgfpoint{-3.079032\du}{6.240033\du}}
\pgfpathlineto{\pgfpoint{-3.055148\du}{6.227493\du}}
\pgfpathlineto{\pgfpoint{-2.983495\du}{6.263917\du}}
\pgfpathlineto{\pgfpoint{-2.960208\du}{6.251975\du}}
\pgfpathlineto{\pgfpoint{-2.971852\du}{6.275560\du}}
\pgfpathlineto{\pgfpoint{-3.055148\du}{6.275560\du}}
\pgfpathlineto{\pgfpoint{-3.019322\du}{6.263917\du}}
\pgfpathlineto{\pgfpoint{-3.079032\du}{6.240033\du}}
\pgfusepath{fill}
\pgfsetbuttcap
\pgfsetmiterjoin
\pgfsetdash{}{0pt}
\definecolor{dialinecolor}{rgb}{0.000000, 0.000000, 0.000000}
\pgfsetfillcolor{dialinecolor}
\pgfpathmoveto{\pgfpoint{-2.828547\du}{6.323627\du}}
\pgfpathlineto{\pgfpoint{-2.852730\du}{6.334972\du}}
\pgfpathlineto{\pgfpoint{-2.924084\du}{6.299444\du}}
\pgfpathlineto{\pgfpoint{-2.960208\du}{6.311685\du}}
\pgfpathlineto{\pgfpoint{-2.935727\du}{6.287801\du}}
\pgfpathlineto{\pgfpoint{-2.852730\du}{6.287801\du}}
\pgfpathlineto{\pgfpoint{-2.888257\du}{6.299444\du}}
\pgfpathlineto{\pgfpoint{-2.828547\du}{6.323627\du}}
\pgfusepath{fill}
\pgfsetbuttcap
\pgfsetmiterjoin
\pgfsetdash{}{0pt}
\definecolor{dialinecolor}{rgb}{1.000000, 1.000000, 1.000000}
\pgfsetfillcolor{dialinecolor}
\pgfpathmoveto{\pgfpoint{-2.947968\du}{6.275560\du}}
\pgfpathlineto{\pgfpoint{-2.924084\du}{6.275560\du}}
\pgfpathlineto{\pgfpoint{-2.852730\du}{6.251975\du}}
\pgfpathlineto{\pgfpoint{-2.816903\du}{6.263917\du}}
\pgfpathlineto{\pgfpoint{-2.840788\du}{6.240033\du}}
\pgfpathlineto{\pgfpoint{-2.911843\du}{6.240033\du}}
\pgfpathlineto{\pgfpoint{-2.888257\du}{6.240033\du}}
\pgfpathlineto{\pgfpoint{-2.947968\du}{6.275560\du}}
\pgfusepath{fill}
\pgfsetbuttcap
\pgfsetmiterjoin
\pgfsetdash{}{0pt}
\definecolor{dialinecolor}{rgb}{1.000000, 1.000000, 1.000000}
\pgfsetfillcolor{dialinecolor}
\pgfpathmoveto{\pgfpoint{-2.960208\du}{6.299444\du}}
\pgfpathlineto{\pgfpoint{-2.983495\du}{6.287801\du}}
\pgfpathlineto{\pgfpoint{-3.055148\du}{6.311685\du}}
\pgfpathlineto{\pgfpoint{-3.090377\du}{6.311685\du}}
\pgfpathlineto{\pgfpoint{-3.066791\du}{6.334972\du}}
\pgfpathlineto{\pgfpoint{-2.983495\du}{6.334972\du}}
\pgfpathlineto{\pgfpoint{-3.019322\du}{6.323627\du}}
\pgfpathlineto{\pgfpoint{-2.960208\du}{6.299444\du}}
\pgfusepath{fill}
\pgfsetbuttcap
\pgfsetmiterjoin
\pgfsetdash{}{0pt}
\definecolor{dialinecolor}{rgb}{1.000000, 1.000000, 1.000000}
\pgfsetfillcolor{dialinecolor}
\pgfpathmoveto{\pgfpoint{-3.079032\du}{6.240033\du}}
\pgfpathlineto{\pgfpoint{-3.055148\du}{6.240033\du}}
\pgfpathlineto{\pgfpoint{-2.983495\du}{6.263917\du}}
\pgfpathlineto{\pgfpoint{-2.947968\du}{6.251975\du}}
\pgfpathlineto{\pgfpoint{-2.971852\du}{6.275560\du}}
\pgfpathlineto{\pgfpoint{-3.055148\du}{6.275560\du}}
\pgfpathlineto{\pgfpoint{-3.019322\du}{6.275560\du}}
\pgfpathlineto{\pgfpoint{-3.079032\du}{6.240033\du}}
\pgfusepath{fill}
\pgfsetbuttcap
\pgfsetmiterjoin
\pgfsetdash{}{0pt}
\definecolor{dialinecolor}{rgb}{1.000000, 1.000000, 1.000000}
\pgfsetfillcolor{dialinecolor}
\pgfpathmoveto{\pgfpoint{-2.828547\du}{6.323627\du}}
\pgfpathlineto{\pgfpoint{-2.852730\du}{6.334972\du}}
\pgfpathlineto{\pgfpoint{-2.924084\du}{6.311685\du}}
\pgfpathlineto{\pgfpoint{-2.947968\du}{6.311685\du}}
\pgfpathlineto{\pgfpoint{-2.935727\du}{6.287801\du}}
\pgfpathlineto{\pgfpoint{-2.852730\du}{6.287801\du}}
\pgfpathlineto{\pgfpoint{-2.888257\du}{6.299444\du}}
\pgfpathlineto{\pgfpoint{-2.828547\du}{6.323627\du}}
\pgfusepath{fill}
\pgfsetbuttcap
\pgfsetmiterjoin
\pgfsetdash{}{0pt}
\definecolor{dialinecolor}{rgb}{0.666667, 0.901961, 1.000000}
\pgfsetstrokecolor{dialinecolor}
\pgfpathmoveto{\pgfpoint{-3.150087\du}{6.275560\du}}
\pgfpathlineto{\pgfpoint{-3.150087\du}{6.370500\du}}
\pgfusepath{stroke}
\pgfsetbuttcap
\pgfsetmiterjoin
\pgfsetdash{}{0pt}
\definecolor{dialinecolor}{rgb}{0.666667, 0.901961, 1.000000}
\pgfsetstrokecolor{dialinecolor}
\pgfpathmoveto{\pgfpoint{-2.769434\du}{6.275560\du}}
\pgfpathlineto{\pgfpoint{-2.769434\du}{6.370500\du}}
\pgfusepath{stroke}
\pgfsetbuttcap
\pgfsetmiterjoin
\pgfsetdash{}{0pt}
\definecolor{dialinecolor}{rgb}{0.000000, 0.470588, 0.666667}
\pgfsetfillcolor{dialinecolor}
\pgfpathmoveto{\pgfpoint{-3.114858\du}{6.781606\du}}
\pgfpathlineto{\pgfpoint{-3.115157\du}{6.778023\du}}
\pgfpathlineto{\pgfpoint{-3.116053\du}{6.774739\du}}
\pgfpathlineto{\pgfpoint{-3.117247\du}{6.771455\du}}
\pgfpathlineto{\pgfpoint{-3.119038\du}{6.768470\du}}
\pgfpathlineto{\pgfpoint{-3.121128\du}{6.765186\du}}
\pgfpathlineto{\pgfpoint{-3.124412\du}{6.761901\du}}
\pgfpathlineto{\pgfpoint{-3.127099\du}{6.758916\du}}
\pgfpathlineto{\pgfpoint{-3.130980\du}{6.755632\du}}
\pgfpathlineto{\pgfpoint{-3.135160\du}{6.752945\du}}
\pgfpathlineto{\pgfpoint{-3.139638\du}{6.749661\du}}
\pgfpathlineto{\pgfpoint{-3.145012\du}{6.746974\du}}
\pgfpathlineto{\pgfpoint{-3.150087\du}{6.744287\du}}
\pgfpathlineto{\pgfpoint{-3.156059\du}{6.741600\du}}
\pgfpathlineto{\pgfpoint{-3.162030\du}{6.738913\du}}
\pgfpathlineto{\pgfpoint{-3.169195\du}{6.736525\du}}
\pgfpathlineto{\pgfpoint{-3.176062\du}{6.734136\du}}
\pgfpathlineto{\pgfpoint{-3.183525\du}{6.732345\du}}
\pgfpathlineto{\pgfpoint{-3.190989\du}{6.729956\du}}
\pgfpathlineto{\pgfpoint{-3.199050\du}{6.728165\du}}
\pgfpathlineto{\pgfpoint{-3.207409\du}{6.725777\du}}
\pgfpathlineto{\pgfpoint{-3.215769\du}{6.724284\du}}
\pgfpathlineto{\pgfpoint{-3.224725\du}{6.723090\du}}
\pgfpathlineto{\pgfpoint{-3.233682\du}{6.721597\du}}
\pgfpathlineto{\pgfpoint{-3.242937\du}{6.720403\du}}
\pgfpathlineto{\pgfpoint{-3.252789\du}{6.719209\du}}
\pgfpathlineto{\pgfpoint{-3.262044\du}{6.718313\du}}
\pgfpathlineto{\pgfpoint{-3.271897\du}{6.717716\du}}
\pgfpathlineto{\pgfpoint{-3.281450\du}{6.717119\du}}
\pgfpathlineto{\pgfpoint{-3.291303\du}{6.716522\du}}
\pgfpathlineto{\pgfpoint{-3.301752\du}{6.716223\du}}
\pgfpathlineto{\pgfpoint{-3.311007\du}{6.716223\du}}
\pgfpathlineto{\pgfpoint{-3.311007\du}{6.716223\du}}
\pgfpathlineto{\pgfpoint{-3.320859\du}{6.716223\du}}
\pgfpathlineto{\pgfpoint{-3.331010\du}{6.716522\du}}
\pgfpathlineto{\pgfpoint{-3.340862\du}{6.717119\du}}
\pgfpathlineto{\pgfpoint{-3.350714\du}{6.717716\du}}
\pgfpathlineto{\pgfpoint{-3.360567\du}{6.718313\du}}
\pgfpathlineto{\pgfpoint{-3.369822\du}{6.719209\du}}
\pgfpathlineto{\pgfpoint{-3.379674\du}{6.720403\du}}
\pgfpathlineto{\pgfpoint{-3.388929\du}{6.721597\du}}
\pgfpathlineto{\pgfpoint{-3.397886\du}{6.723090\du}}
\pgfpathlineto{\pgfpoint{-3.406544\du}{6.724284\du}}
\pgfpathlineto{\pgfpoint{-3.414903\du}{6.725777\du}}
\pgfpathlineto{\pgfpoint{-3.423263\du}{6.728165\du}}
\pgfpathlineto{\pgfpoint{-3.431323\du}{6.729956\du}}
\pgfpathlineto{\pgfpoint{-3.438787\du}{6.732345\du}}
\pgfpathlineto{\pgfpoint{-3.446251\du}{6.734136\du}}
\pgfpathlineto{\pgfpoint{-3.453416\du}{6.736525\du}}
\pgfpathlineto{\pgfpoint{-3.460283\du}{6.738913\du}}
\pgfpathlineto{\pgfpoint{-3.466254\du}{6.741600\du}}
\pgfpathlineto{\pgfpoint{-3.472524\du}{6.744287\du}}
\pgfpathlineto{\pgfpoint{-3.477599\du}{6.746974\du}}
\pgfpathlineto{\pgfpoint{-3.482674\du}{6.749661\du}}
\pgfpathlineto{\pgfpoint{-3.487451\du}{6.752945\du}}
\pgfpathlineto{\pgfpoint{-3.491631\du}{6.755632\du}}
\pgfpathlineto{\pgfpoint{-3.495214\du}{6.758916\du}}
\pgfpathlineto{\pgfpoint{-3.498199\du}{6.761901\du}}
\pgfpathlineto{\pgfpoint{-3.500886\du}{6.765186\du}}
\pgfpathlineto{\pgfpoint{-3.503275\du}{6.768470\du}}
\pgfpathlineto{\pgfpoint{-3.505066\du}{6.771455\du}}
\pgfpathlineto{\pgfpoint{-3.506559\du}{6.774739\du}}
\pgfpathlineto{\pgfpoint{-3.507454\du}{6.778023\du}}
\pgfpathlineto{\pgfpoint{-3.507454\du}{6.781606\du}}
\pgfpathlineto{\pgfpoint{-3.507454\du}{6.781606\du}}
\pgfpathlineto{\pgfpoint{-3.507454\du}{6.784890\du}}
\pgfpathlineto{\pgfpoint{-3.506559\du}{6.788174\du}}
\pgfpathlineto{\pgfpoint{-3.505066\du}{6.791458\du}}
\pgfpathlineto{\pgfpoint{-3.503275\du}{6.794742\du}}
\pgfpathlineto{\pgfpoint{-3.500886\du}{6.797728\du}}
\pgfpathlineto{\pgfpoint{-3.498199\du}{6.801012\du}}
\pgfpathlineto{\pgfpoint{-3.495214\du}{6.803997\du}}
\pgfpathlineto{\pgfpoint{-3.491631\du}{6.807281\du}}
\pgfpathlineto{\pgfpoint{-3.487451\du}{6.810267\du}}
\pgfpathlineto{\pgfpoint{-3.482674\du}{6.813252\du}}
\pgfpathlineto{\pgfpoint{-3.477599\du}{6.815939\du}}
\pgfpathlineto{\pgfpoint{-3.472524\du}{6.818626\du}}
\pgfpathlineto{\pgfpoint{-3.466254\du}{6.821313\du}}
\pgfpathlineto{\pgfpoint{-3.460283\du}{6.823702\du}}
\pgfpathlineto{\pgfpoint{-3.453416\du}{6.826389\du}}
\pgfpathlineto{\pgfpoint{-3.446251\du}{6.828777\du}}
\pgfpathlineto{\pgfpoint{-3.438787\du}{6.830867\du}}
\pgfpathlineto{\pgfpoint{-3.431323\du}{6.833255\du}}
\pgfpathlineto{\pgfpoint{-3.423263\du}{6.834748\du}}
\pgfpathlineto{\pgfpoint{-3.414903\du}{6.836838\du}}
\pgfpathlineto{\pgfpoint{-3.406544\du}{6.838629\du}}
\pgfpathlineto{\pgfpoint{-3.397886\du}{6.839824\du}}
\pgfpathlineto{\pgfpoint{-3.388929\du}{6.841316\du}}
\pgfpathlineto{\pgfpoint{-3.379674\du}{6.842511\du}}
\pgfpathlineto{\pgfpoint{-3.369822\du}{6.843705\du}}
\pgfpathlineto{\pgfpoint{-3.360567\du}{6.844600\du}}
\pgfpathlineto{\pgfpoint{-3.350714\du}{6.845198\du}}
\pgfpathlineto{\pgfpoint{-3.340862\du}{6.846093\du}}
\pgfpathlineto{\pgfpoint{-3.331010\du}{6.846093\du}}
\pgfpathlineto{\pgfpoint{-3.320859\du}{6.846690\du}}
\pgfpathlineto{\pgfpoint{-3.311007\du}{6.846690\du}}
\pgfpathlineto{\pgfpoint{-3.311007\du}{6.846690\du}}
\pgfpathlineto{\pgfpoint{-3.301752\du}{6.846690\du}}
\pgfpathlineto{\pgfpoint{-3.291303\du}{6.846093\du}}
\pgfpathlineto{\pgfpoint{-3.281450\du}{6.846093\du}}
\pgfpathlineto{\pgfpoint{-3.271897\du}{6.845198\du}}
\pgfpathlineto{\pgfpoint{-3.262044\du}{6.844600\du}}
\pgfpathlineto{\pgfpoint{-3.252789\du}{6.843705\du}}
\pgfpathlineto{\pgfpoint{-3.242937\du}{6.842511\du}}
\pgfpathlineto{\pgfpoint{-3.233682\du}{6.841316\du}}
\pgfpathlineto{\pgfpoint{-3.224725\du}{6.839824\du}}
\pgfpathlineto{\pgfpoint{-3.215769\du}{6.838629\du}}
\pgfpathlineto{\pgfpoint{-3.207409\du}{6.836838\du}}
\pgfpathlineto{\pgfpoint{-3.199050\du}{6.834748\du}}
\pgfpathlineto{\pgfpoint{-3.190989\du}{6.833255\du}}
\pgfpathlineto{\pgfpoint{-3.183525\du}{6.830867\du}}
\pgfpathlineto{\pgfpoint{-3.176062\du}{6.828777\du}}
\pgfpathlineto{\pgfpoint{-3.169195\du}{6.826389\du}}
\pgfpathlineto{\pgfpoint{-3.162030\du}{6.823702\du}}
\pgfpathlineto{\pgfpoint{-3.156059\du}{6.821313\du}}
\pgfpathlineto{\pgfpoint{-3.150087\du}{6.818626\du}}
\pgfpathlineto{\pgfpoint{-3.145012\du}{6.815939\du}}
\pgfpathlineto{\pgfpoint{-3.139638\du}{6.813252\du}}
\pgfpathlineto{\pgfpoint{-3.135160\du}{6.810267\du}}
\pgfpathlineto{\pgfpoint{-3.130980\du}{6.807281\du}}
\pgfpathlineto{\pgfpoint{-3.127099\du}{6.803997\du}}
\pgfpathlineto{\pgfpoint{-3.124412\du}{6.801012\du}}
\pgfpathlineto{\pgfpoint{-3.121128\du}{6.797728\du}}
\pgfpathlineto{\pgfpoint{-3.119038\du}{6.794742\du}}
\pgfpathlineto{\pgfpoint{-3.117247\du}{6.791458\du}}
\pgfpathlineto{\pgfpoint{-3.116053\du}{6.788174\du}}
\pgfpathlineto{\pgfpoint{-3.115157\du}{6.784890\du}}
\pgfpathlineto{\pgfpoint{-3.114858\du}{6.781606\du}}
\pgfusepath{fill}
\pgfsetbuttcap
\pgfsetmiterjoin
\pgfsetdash{}{0pt}
\definecolor{dialinecolor}{rgb}{0.666667, 0.901961, 1.000000}
\pgfsetstrokecolor{dialinecolor}
\pgfpathmoveto{\pgfpoint{-3.126800\du}{6.775336\du}}
\pgfpathlineto{\pgfpoint{-3.127099\du}{6.772351\du}}
\pgfpathlineto{\pgfpoint{-3.127696\du}{6.769664\du}}
\pgfpathlineto{\pgfpoint{-3.128890\du}{6.766678\du}}
\pgfpathlineto{\pgfpoint{-3.130980\du}{6.763096\du}}
\pgfpathlineto{\pgfpoint{-3.133070\du}{6.760707\du}}
\pgfpathlineto{\pgfpoint{-3.135160\du}{6.757423\du}}
\pgfpathlineto{\pgfpoint{-3.138444\du}{6.754736\du}}
\pgfpathlineto{\pgfpoint{-3.141728\du}{6.752049\du}}
\pgfpathlineto{\pgfpoint{-3.145908\du}{6.749064\du}}
\pgfpathlineto{\pgfpoint{-3.150685\du}{6.746675\du}}
\pgfpathlineto{\pgfpoint{-3.155760\du}{6.743988\du}}
\pgfpathlineto{\pgfpoint{-3.160835\du}{6.741600\du}}
\pgfpathlineto{\pgfpoint{-3.166508\du}{6.739212\du}}
\pgfpathlineto{\pgfpoint{-3.172479\du}{6.736823\du}}
\pgfpathlineto{\pgfpoint{-3.179047\du}{6.734136\du}}
\pgfpathlineto{\pgfpoint{-3.185914\du}{6.732643\du}}
\pgfpathlineto{\pgfpoint{-3.192482\du}{6.730554\du}}
\pgfpathlineto{\pgfpoint{-3.200543\du}{6.728762\du}}
\pgfpathlineto{\pgfpoint{-3.208007\du}{6.727269\du}}
\pgfpathlineto{\pgfpoint{-3.215769\du}{6.725180\du}}
\pgfpathlineto{\pgfpoint{-3.224725\du}{6.723687\du}}
\pgfpathlineto{\pgfpoint{-3.233085\du}{6.722194\du}}
\pgfpathlineto{\pgfpoint{-3.241444\du}{6.721000\du}}
\pgfpathlineto{\pgfpoint{-3.250700\du}{6.719806\du}}
\pgfpathlineto{\pgfpoint{-3.259358\du}{6.718910\du}}
\pgfpathlineto{\pgfpoint{-3.268911\du}{6.718014\du}}
\pgfpathlineto{\pgfpoint{-3.278465\du}{6.717716\du}}
\pgfpathlineto{\pgfpoint{-3.287720\du}{6.717119\du}}
\pgfpathlineto{\pgfpoint{-3.297572\du}{6.716522\du}}
\pgfpathlineto{\pgfpoint{-3.307424\du}{6.716223\du}}
\pgfpathlineto{\pgfpoint{-3.316978\du}{6.716223\du}}
\pgfpathlineto{\pgfpoint{-3.316978\du}{6.716223\du}}
\pgfpathlineto{\pgfpoint{-3.326233\du}{6.716223\du}}
\pgfpathlineto{\pgfpoint{-3.336085\du}{6.716522\du}}
\pgfpathlineto{\pgfpoint{-3.345639\du}{6.717119\du}}
\pgfpathlineto{\pgfpoint{-3.355193\du}{6.717716\du}}
\pgfpathlineto{\pgfpoint{-3.364746\du}{6.718014\du}}
\pgfpathlineto{\pgfpoint{-3.374001\du}{6.718910\du}}
\pgfpathlineto{\pgfpoint{-3.382958\du}{6.719806\du}}
\pgfpathlineto{\pgfpoint{-3.392213\du}{6.721000\du}}
\pgfpathlineto{\pgfpoint{-3.400573\du}{6.722194\du}}
\pgfpathlineto{\pgfpoint{-3.409529\du}{6.723687\du}}
\pgfpathlineto{\pgfpoint{-3.417889\du}{6.725180\du}}
\pgfpathlineto{\pgfpoint{-3.425651\du}{6.727269\du}}
\pgfpathlineto{\pgfpoint{-3.433712\du}{6.728762\du}}
\pgfpathlineto{\pgfpoint{-3.441176\du}{6.730554\du}}
\pgfpathlineto{\pgfpoint{-3.447744\du}{6.732643\du}}
\pgfpathlineto{\pgfpoint{-3.455208\du}{6.734136\du}}
\pgfpathlineto{\pgfpoint{-3.461179\du}{6.736823\du}}
\pgfpathlineto{\pgfpoint{-3.467150\du}{6.739212\du}}
\pgfpathlineto{\pgfpoint{-3.473419\du}{6.741600\du}}
\pgfpathlineto{\pgfpoint{-3.478495\du}{6.743988\du}}
\pgfpathlineto{\pgfpoint{-3.483570\du}{6.746675\du}}
\pgfpathlineto{\pgfpoint{-3.487750\du}{6.749064\du}}
\pgfpathlineto{\pgfpoint{-3.491930\du}{6.752049\du}}
\pgfpathlineto{\pgfpoint{-3.495811\du}{6.754736\du}}
\pgfpathlineto{\pgfpoint{-3.499095\du}{6.757423\du}}
\pgfpathlineto{\pgfpoint{-3.500886\du}{6.760707\du}}
\pgfpathlineto{\pgfpoint{-3.503275\du}{6.763096\du}}
\pgfpathlineto{\pgfpoint{-3.505066\du}{6.766678\du}}
\pgfpathlineto{\pgfpoint{-3.506559\du}{6.769664\du}}
\pgfpathlineto{\pgfpoint{-3.506857\du}{6.772351\du}}
\pgfpathlineto{\pgfpoint{-3.507454\du}{6.775336\du}}
\pgfpathlineto{\pgfpoint{-3.507454\du}{6.775336\du}}
\pgfpathlineto{\pgfpoint{-3.506857\du}{6.778919\du}}
\pgfpathlineto{\pgfpoint{-3.506559\du}{6.781606\du}}
\pgfpathlineto{\pgfpoint{-3.505066\du}{6.784591\du}}
\pgfpathlineto{\pgfpoint{-3.503275\du}{6.787577\du}}
\pgfpathlineto{\pgfpoint{-3.500886\du}{6.790264\du}}
\pgfpathlineto{\pgfpoint{-3.499095\du}{6.793548\du}}
\pgfpathlineto{\pgfpoint{-3.495811\du}{6.796235\du}}
\pgfpathlineto{\pgfpoint{-3.491930\du}{6.799220\du}}
\pgfpathlineto{\pgfpoint{-3.487750\du}{6.801907\du}}
\pgfpathlineto{\pgfpoint{-3.483570\du}{6.804594\du}}
\pgfpathlineto{\pgfpoint{-3.478495\du}{6.807281\du}}
\pgfpathlineto{\pgfpoint{-3.473419\du}{6.808774\du}}
\pgfpathlineto{\pgfpoint{-3.467150\du}{6.812058\du}}
\pgfpathlineto{\pgfpoint{-3.461179\du}{6.814148\du}}
\pgfpathlineto{\pgfpoint{-3.455208\du}{6.816238\du}}
\pgfpathlineto{\pgfpoint{-3.447744\du}{6.818626\du}}
\pgfpathlineto{\pgfpoint{-3.441176\du}{6.820716\du}}
\pgfpathlineto{\pgfpoint{-3.433712\du}{6.822508\du}}
\pgfpathlineto{\pgfpoint{-3.425651\du}{6.823702\du}}
\pgfpathlineto{\pgfpoint{-3.417889\du}{6.826090\du}}
\pgfpathlineto{\pgfpoint{-3.409529\du}{6.827583\du}}
\pgfpathlineto{\pgfpoint{-3.400573\du}{6.829076\du}}
\pgfpathlineto{\pgfpoint{-3.392213\du}{6.830270\du}}
\pgfpathlineto{\pgfpoint{-3.382958\du}{6.831166\du}}
\pgfpathlineto{\pgfpoint{-3.374001\du}{6.832061\du}}
\pgfpathlineto{\pgfpoint{-3.364746\du}{6.833255\du}}
\pgfpathlineto{\pgfpoint{-3.355193\du}{6.833554\du}}
\pgfpathlineto{\pgfpoint{-3.345639\du}{6.834151\du}}
\pgfpathlineto{\pgfpoint{-3.336085\du}{6.834450\du}}
\pgfpathlineto{\pgfpoint{-3.326233\du}{6.834748\du}}
\pgfpathlineto{\pgfpoint{-3.316978\du}{6.834748\du}}
\pgfpathlineto{\pgfpoint{-3.316978\du}{6.834748\du}}
\pgfpathlineto{\pgfpoint{-3.307424\du}{6.834748\du}}
\pgfpathlineto{\pgfpoint{-3.297572\du}{6.834450\du}}
\pgfpathlineto{\pgfpoint{-3.287720\du}{6.834151\du}}
\pgfpathlineto{\pgfpoint{-3.278465\du}{6.833554\du}}
\pgfpathlineto{\pgfpoint{-3.268911\du}{6.833255\du}}
\pgfpathlineto{\pgfpoint{-3.259358\du}{6.832061\du}}
\pgfpathlineto{\pgfpoint{-3.250700\du}{6.831166\du}}
\pgfpathlineto{\pgfpoint{-3.241444\du}{6.830270\du}}
\pgfpathlineto{\pgfpoint{-3.233085\du}{6.829076\du}}
\pgfpathlineto{\pgfpoint{-3.224725\du}{6.827583\du}}
\pgfpathlineto{\pgfpoint{-3.215769\du}{6.826090\du}}
\pgfpathlineto{\pgfpoint{-3.208007\du}{6.823702\du}}
\pgfpathlineto{\pgfpoint{-3.200543\du}{6.822508\du}}
\pgfpathlineto{\pgfpoint{-3.192482\du}{6.820716\du}}
\pgfpathlineto{\pgfpoint{-3.185914\du}{6.818626\du}}
\pgfpathlineto{\pgfpoint{-3.179047\du}{6.816238\du}}
\pgfpathlineto{\pgfpoint{-3.172479\du}{6.814148\du}}
\pgfpathlineto{\pgfpoint{-3.166508\du}{6.812058\du}}
\pgfpathlineto{\pgfpoint{-3.160835\du}{6.808774\du}}
\pgfpathlineto{\pgfpoint{-3.155760\du}{6.807281\du}}
\pgfpathlineto{\pgfpoint{-3.150685\du}{6.804594\du}}
\pgfpathlineto{\pgfpoint{-3.145908\du}{6.801907\du}}
\pgfpathlineto{\pgfpoint{-3.141728\du}{6.799220\du}}
\pgfpathlineto{\pgfpoint{-3.138444\du}{6.796235\du}}
\pgfpathlineto{\pgfpoint{-3.135160\du}{6.793548\du}}
\pgfpathlineto{\pgfpoint{-3.133070\du}{6.790264\du}}
\pgfpathlineto{\pgfpoint{-3.130980\du}{6.787577\du}}
\pgfpathlineto{\pgfpoint{-3.128890\du}{6.784591\du}}
\pgfpathlineto{\pgfpoint{-3.127696\du}{6.781606\du}}
\pgfpathlineto{\pgfpoint{-3.127099\du}{6.778919\du}}
\pgfpathlineto{\pgfpoint{-3.126800\du}{6.775336\du}}
\pgfusepath{stroke}
\pgfsetbuttcap
\pgfsetmiterjoin
\pgfsetdash{}{0pt}
\definecolor{dialinecolor}{rgb}{0.000000, 0.470588, 0.666667}
\pgfsetfillcolor{dialinecolor}
\pgfpathmoveto{\pgfpoint{-3.507454\du}{6.692339\du}}
\pgfpathlineto{\pgfpoint{-3.507454\du}{6.787577\du}}
\pgfpathlineto{\pgfpoint{-3.114560\du}{6.787577\du}}
\pgfpathlineto{\pgfpoint{-3.114560\du}{6.692339\du}}
\pgfpathlineto{\pgfpoint{-3.507454\du}{6.692339\du}}
\pgfusepath{fill}
\pgfsetbuttcap
\pgfsetmiterjoin
\pgfsetdash{}{0pt}
\definecolor{dialinecolor}{rgb}{0.000000, 0.705882, 1.000000}
\pgfsetfillcolor{dialinecolor}
\pgfpathmoveto{\pgfpoint{-3.114858\du}{6.686069\du}}
\pgfpathlineto{\pgfpoint{-3.115157\du}{6.682487\du}}
\pgfpathlineto{\pgfpoint{-3.116053\du}{6.679203\du}}
\pgfpathlineto{\pgfpoint{-3.117247\du}{6.676217\du}}
\pgfpathlineto{\pgfpoint{-3.119038\du}{6.672634\du}}
\pgfpathlineto{\pgfpoint{-3.121128\du}{6.669350\du}}
\pgfpathlineto{\pgfpoint{-3.124412\du}{6.666066\du}}
\pgfpathlineto{\pgfpoint{-3.127099\du}{6.663379\du}}
\pgfpathlineto{\pgfpoint{-3.130980\du}{6.659498\du}}
\pgfpathlineto{\pgfpoint{-3.135160\du}{6.657110\du}}
\pgfpathlineto{\pgfpoint{-3.139638\du}{6.654124\du}}
\pgfpathlineto{\pgfpoint{-3.145012\du}{6.651139\du}}
\pgfpathlineto{\pgfpoint{-3.150087\du}{6.648452\du}}
\pgfpathlineto{\pgfpoint{-3.156059\du}{6.645765\du}}
\pgfpathlineto{\pgfpoint{-3.162030\du}{6.643376\du}}
\pgfpathlineto{\pgfpoint{-3.169195\du}{6.640689\du}}
\pgfpathlineto{\pgfpoint{-3.176062\du}{6.638301\du}}
\pgfpathlineto{\pgfpoint{-3.183525\du}{6.636211\du}}
\pgfpathlineto{\pgfpoint{-3.190989\du}{6.634121\du}}
\pgfpathlineto{\pgfpoint{-3.199050\du}{6.632330\du}}
\pgfpathlineto{\pgfpoint{-3.207409\du}{6.630240\du}}
\pgfpathlineto{\pgfpoint{-3.215769\du}{6.628449\du}}
\pgfpathlineto{\pgfpoint{-3.224725\du}{6.627255\du}}
\pgfpathlineto{\pgfpoint{-3.233682\du}{6.625762\du}}
\pgfpathlineto{\pgfpoint{-3.242937\du}{6.624568\du}}
\pgfpathlineto{\pgfpoint{-3.252789\du}{6.623373\du}}
\pgfpathlineto{\pgfpoint{-3.262044\du}{6.622179\du}}
\pgfpathlineto{\pgfpoint{-3.271897\du}{6.621881\du}}
\pgfpathlineto{\pgfpoint{-3.281450\du}{6.620985\du}}
\pgfpathlineto{\pgfpoint{-3.291303\du}{6.620686\du}}
\pgfpathlineto{\pgfpoint{-3.301752\du}{6.620388\du}}
\pgfpathlineto{\pgfpoint{-3.311007\du}{6.620388\du}}
\pgfpathlineto{\pgfpoint{-3.311007\du}{6.620388\du}}
\pgfpathlineto{\pgfpoint{-3.320859\du}{6.620388\du}}
\pgfpathlineto{\pgfpoint{-3.331010\du}{6.620686\du}}
\pgfpathlineto{\pgfpoint{-3.340862\du}{6.620985\du}}
\pgfpathlineto{\pgfpoint{-3.350714\du}{6.621881\du}}
\pgfpathlineto{\pgfpoint{-3.360567\du}{6.622179\du}}
\pgfpathlineto{\pgfpoint{-3.369822\du}{6.623373\du}}
\pgfpathlineto{\pgfpoint{-3.379674\du}{6.624568\du}}
\pgfpathlineto{\pgfpoint{-3.388929\du}{6.625762\du}}
\pgfpathlineto{\pgfpoint{-3.397886\du}{6.627255\du}}
\pgfpathlineto{\pgfpoint{-3.406544\du}{6.628449\du}}
\pgfpathlineto{\pgfpoint{-3.414903\du}{6.630240\du}}
\pgfpathlineto{\pgfpoint{-3.423263\du}{6.632330\du}}
\pgfpathlineto{\pgfpoint{-3.431323\du}{6.634121\du}}
\pgfpathlineto{\pgfpoint{-3.438787\du}{6.636211\du}}
\pgfpathlineto{\pgfpoint{-3.446251\du}{6.638301\du}}
\pgfpathlineto{\pgfpoint{-3.453416\du}{6.640689\du}}
\pgfpathlineto{\pgfpoint{-3.460283\du}{6.643376\du}}
\pgfpathlineto{\pgfpoint{-3.466254\du}{6.645765\du}}
\pgfpathlineto{\pgfpoint{-3.472524\du}{6.648452\du}}
\pgfpathlineto{\pgfpoint{-3.477599\du}{6.651139\du}}
\pgfpathlineto{\pgfpoint{-3.482674\du}{6.654124\du}}
\pgfpathlineto{\pgfpoint{-3.487451\du}{6.657110\du}}
\pgfpathlineto{\pgfpoint{-3.491631\du}{6.659498\du}}
\pgfpathlineto{\pgfpoint{-3.495214\du}{6.663379\du}}
\pgfpathlineto{\pgfpoint{-3.498199\du}{6.666066\du}}
\pgfpathlineto{\pgfpoint{-3.500886\du}{6.669350\du}}
\pgfpathlineto{\pgfpoint{-3.503275\du}{6.672634\du}}
\pgfpathlineto{\pgfpoint{-3.505066\du}{6.676217\du}}
\pgfpathlineto{\pgfpoint{-3.506559\du}{6.679203\du}}
\pgfpathlineto{\pgfpoint{-3.507454\du}{6.682487\du}}
\pgfpathlineto{\pgfpoint{-3.507454\du}{6.686069\du}}
\pgfpathlineto{\pgfpoint{-3.507454\du}{6.686069\du}}
\pgfpathlineto{\pgfpoint{-3.507454\du}{6.689353\du}}
\pgfpathlineto{\pgfpoint{-3.506559\du}{6.692637\du}}
\pgfpathlineto{\pgfpoint{-3.505066\du}{6.695623\du}}
\pgfpathlineto{\pgfpoint{-3.503275\du}{6.699206\du}}
\pgfpathlineto{\pgfpoint{-3.500886\du}{6.702490\du}}
\pgfpathlineto{\pgfpoint{-3.498199\du}{6.705475\du}}
\pgfpathlineto{\pgfpoint{-3.495214\du}{6.708461\du}}
\pgfpathlineto{\pgfpoint{-3.491631\du}{6.711745\du}}
\pgfpathlineto{\pgfpoint{-3.487451\du}{6.714730\du}}
\pgfpathlineto{\pgfpoint{-3.482674\du}{6.717716\du}}
\pgfpathlineto{\pgfpoint{-3.477599\du}{6.720701\du}}
\pgfpathlineto{\pgfpoint{-3.472524\du}{6.723388\du}}
\pgfpathlineto{\pgfpoint{-3.466254\du}{6.725777\du}}
\pgfpathlineto{\pgfpoint{-3.460283\du}{6.728464\du}}
\pgfpathlineto{\pgfpoint{-3.453416\du}{6.731151\du}}
\pgfpathlineto{\pgfpoint{-3.446251\du}{6.733240\du}}
\pgfpathlineto{\pgfpoint{-3.438787\du}{6.735629\du}}
\pgfpathlineto{\pgfpoint{-3.431323\du}{6.737719\du}}
\pgfpathlineto{\pgfpoint{-3.423263\du}{6.739510\du}}
\pgfpathlineto{\pgfpoint{-3.414903\du}{6.741600\du}}
\pgfpathlineto{\pgfpoint{-3.406544\du}{6.743391\du}}
\pgfpathlineto{\pgfpoint{-3.397886\du}{6.744585\du}}
\pgfpathlineto{\pgfpoint{-3.388929\du}{6.746078\du}}
\pgfpathlineto{\pgfpoint{-3.379674\du}{6.747272\du}}
\pgfpathlineto{\pgfpoint{-3.369822\du}{6.748168\du}}
\pgfpathlineto{\pgfpoint{-3.360567\du}{6.749064\du}}
\pgfpathlineto{\pgfpoint{-3.350714\du}{6.749959\du}}
\pgfpathlineto{\pgfpoint{-3.340862\du}{6.750855\du}}
\pgfpathlineto{\pgfpoint{-3.331010\du}{6.751154\du}}
\pgfpathlineto{\pgfpoint{-3.320859\du}{6.751452\du}}
\pgfpathlineto{\pgfpoint{-3.311007\du}{6.751452\du}}
\pgfpathlineto{\pgfpoint{-3.311007\du}{6.751452\du}}
\pgfpathlineto{\pgfpoint{-3.301752\du}{6.751452\du}}
\pgfpathlineto{\pgfpoint{-3.291303\du}{6.751154\du}}
\pgfpathlineto{\pgfpoint{-3.281450\du}{6.750855\du}}
\pgfpathlineto{\pgfpoint{-3.271897\du}{6.749959\du}}
\pgfpathlineto{\pgfpoint{-3.262044\du}{6.749064\du}}
\pgfpathlineto{\pgfpoint{-3.252789\du}{6.748168\du}}
\pgfpathlineto{\pgfpoint{-3.242937\du}{6.747272\du}}
\pgfpathlineto{\pgfpoint{-3.233682\du}{6.746078\du}}
\pgfpathlineto{\pgfpoint{-3.224725\du}{6.744585\du}}
\pgfpathlineto{\pgfpoint{-3.215769\du}{6.743391\du}}
\pgfpathlineto{\pgfpoint{-3.207409\du}{6.741600\du}}
\pgfpathlineto{\pgfpoint{-3.199050\du}{6.739510\du}}
\pgfpathlineto{\pgfpoint{-3.190989\du}{6.737719\du}}
\pgfpathlineto{\pgfpoint{-3.183525\du}{6.735629\du}}
\pgfpathlineto{\pgfpoint{-3.176062\du}{6.733240\du}}
\pgfpathlineto{\pgfpoint{-3.169195\du}{6.731151\du}}
\pgfpathlineto{\pgfpoint{-3.162030\du}{6.728464\du}}
\pgfpathlineto{\pgfpoint{-3.156059\du}{6.725777\du}}
\pgfpathlineto{\pgfpoint{-3.150087\du}{6.723388\du}}
\pgfpathlineto{\pgfpoint{-3.145012\du}{6.720701\du}}
\pgfpathlineto{\pgfpoint{-3.139638\du}{6.717716\du}}
\pgfpathlineto{\pgfpoint{-3.135160\du}{6.714730\du}}
\pgfpathlineto{\pgfpoint{-3.130980\du}{6.711745\du}}
\pgfpathlineto{\pgfpoint{-3.127099\du}{6.708461\du}}
\pgfpathlineto{\pgfpoint{-3.124412\du}{6.705475\du}}
\pgfpathlineto{\pgfpoint{-3.121128\du}{6.702490\du}}
\pgfpathlineto{\pgfpoint{-3.119038\du}{6.699206\du}}
\pgfpathlineto{\pgfpoint{-3.117247\du}{6.695623\du}}
\pgfpathlineto{\pgfpoint{-3.116053\du}{6.692637\du}}
\pgfpathlineto{\pgfpoint{-3.115157\du}{6.689353\du}}
\pgfpathlineto{\pgfpoint{-3.114858\du}{6.686069\du}}
\pgfusepath{fill}
\pgfsetbuttcap
\pgfsetmiterjoin
\pgfsetdash{}{0pt}
\definecolor{dialinecolor}{rgb}{0.666667, 0.901961, 1.000000}
\pgfsetstrokecolor{dialinecolor}
\pgfpathmoveto{\pgfpoint{-3.126800\du}{6.680397\du}}
\pgfpathlineto{\pgfpoint{-3.127099\du}{6.677411\du}}
\pgfpathlineto{\pgfpoint{-3.127696\du}{6.674426\du}}
\pgfpathlineto{\pgfpoint{-3.128890\du}{6.671739\du}}
\pgfpathlineto{\pgfpoint{-3.130980\du}{6.668455\du}}
\pgfpathlineto{\pgfpoint{-3.133070\du}{6.665768\du}}
\pgfpathlineto{\pgfpoint{-3.135160\du}{6.662484\du}}
\pgfpathlineto{\pgfpoint{-3.138444\du}{6.659498\du}}
\pgfpathlineto{\pgfpoint{-3.141728\du}{6.657110\du}}
\pgfpathlineto{\pgfpoint{-3.145908\du}{6.654423\du}}
\pgfpathlineto{\pgfpoint{-3.150685\du}{6.651736\du}}
\pgfpathlineto{\pgfpoint{-3.155760\du}{6.649049\du}}
\pgfpathlineto{\pgfpoint{-3.160835\du}{6.646660\du}}
\pgfpathlineto{\pgfpoint{-3.166508\du}{6.643675\du}}
\pgfpathlineto{\pgfpoint{-3.172479\du}{6.641884\du}}
\pgfpathlineto{\pgfpoint{-3.179047\du}{6.639794\du}}
\pgfpathlineto{\pgfpoint{-3.185914\du}{6.637107\du}}
\pgfpathlineto{\pgfpoint{-3.192482\du}{6.635614\du}}
\pgfpathlineto{\pgfpoint{-3.200543\du}{6.633823\du}}
\pgfpathlineto{\pgfpoint{-3.208007\du}{6.631733\du}}
\pgfpathlineto{\pgfpoint{-3.215769\du}{6.630240\du}}
\pgfpathlineto{\pgfpoint{-3.224725\du}{6.628747\du}}
\pgfpathlineto{\pgfpoint{-3.233085\du}{6.627255\du}}
\pgfpathlineto{\pgfpoint{-3.241444\du}{6.626060\du}}
\pgfpathlineto{\pgfpoint{-3.250700\du}{6.624866\du}}
\pgfpathlineto{\pgfpoint{-3.259358\du}{6.623672\du}}
\pgfpathlineto{\pgfpoint{-3.268911\du}{6.623075\du}}
\pgfpathlineto{\pgfpoint{-3.278465\du}{6.622179\du}}
\pgfpathlineto{\pgfpoint{-3.287720\du}{6.622179\du}}
\pgfpathlineto{\pgfpoint{-3.297572\du}{6.621881\du}}
\pgfpathlineto{\pgfpoint{-3.307424\du}{6.620985\du}}
\pgfpathlineto{\pgfpoint{-3.316978\du}{6.620985\du}}
\pgfpathlineto{\pgfpoint{-3.316978\du}{6.620985\du}}
\pgfpathlineto{\pgfpoint{-3.326233\du}{6.620985\du}}
\pgfpathlineto{\pgfpoint{-3.336085\du}{6.621881\du}}
\pgfpathlineto{\pgfpoint{-3.345639\du}{6.622179\du}}
\pgfpathlineto{\pgfpoint{-3.355193\du}{6.622179\du}}
\pgfpathlineto{\pgfpoint{-3.364746\du}{6.623075\du}}
\pgfpathlineto{\pgfpoint{-3.374001\du}{6.623672\du}}
\pgfpathlineto{\pgfpoint{-3.382958\du}{6.624866\du}}
\pgfpathlineto{\pgfpoint{-3.392213\du}{6.626060\du}}
\pgfpathlineto{\pgfpoint{-3.400573\du}{6.627255\du}}
\pgfpathlineto{\pgfpoint{-3.409529\du}{6.628747\du}}
\pgfpathlineto{\pgfpoint{-3.417889\du}{6.630240\du}}
\pgfpathlineto{\pgfpoint{-3.425651\du}{6.631733\du}}
\pgfpathlineto{\pgfpoint{-3.433712\du}{6.633823\du}}
\pgfpathlineto{\pgfpoint{-3.441176\du}{6.635614\du}}
\pgfpathlineto{\pgfpoint{-3.447744\du}{6.637107\du}}
\pgfpathlineto{\pgfpoint{-3.455208\du}{6.639794\du}}
\pgfpathlineto{\pgfpoint{-3.461179\du}{6.641884\du}}
\pgfpathlineto{\pgfpoint{-3.467150\du}{6.643675\du}}
\pgfpathlineto{\pgfpoint{-3.473419\du}{6.646660\du}}
\pgfpathlineto{\pgfpoint{-3.478495\du}{6.649049\du}}
\pgfpathlineto{\pgfpoint{-3.483570\du}{6.651736\du}}
\pgfpathlineto{\pgfpoint{-3.487750\du}{6.654423\du}}
\pgfpathlineto{\pgfpoint{-3.491930\du}{6.657110\du}}
\pgfpathlineto{\pgfpoint{-3.495811\du}{6.659498\du}}
\pgfpathlineto{\pgfpoint{-3.499095\du}{6.662484\du}}
\pgfpathlineto{\pgfpoint{-3.500886\du}{6.665768\du}}
\pgfpathlineto{\pgfpoint{-3.503275\du}{6.668455\du}}
\pgfpathlineto{\pgfpoint{-3.505066\du}{6.671739\du}}
\pgfpathlineto{\pgfpoint{-3.506559\du}{6.674426\du}}
\pgfpathlineto{\pgfpoint{-3.506857\du}{6.677411\du}}
\pgfpathlineto{\pgfpoint{-3.507454\du}{6.680397\du}}
\pgfpathlineto{\pgfpoint{-3.507454\du}{6.680397\du}}
\pgfpathlineto{\pgfpoint{-3.506857\du}{6.683979\du}}
\pgfpathlineto{\pgfpoint{-3.506559\du}{6.686666\du}}
\pgfpathlineto{\pgfpoint{-3.505066\du}{6.689353\du}}
\pgfpathlineto{\pgfpoint{-3.503275\du}{6.692637\du}}
\pgfpathlineto{\pgfpoint{-3.500886\du}{6.695324\du}}
\pgfpathlineto{\pgfpoint{-3.499095\du}{6.698907\du}}
\pgfpathlineto{\pgfpoint{-3.495811\du}{6.701594\du}}
\pgfpathlineto{\pgfpoint{-3.491930\du}{6.704281\du}}
\pgfpathlineto{\pgfpoint{-3.487750\du}{6.706968\du}}
\pgfpathlineto{\pgfpoint{-3.483570\du}{6.709655\du}}
\pgfpathlineto{\pgfpoint{-3.478495\du}{6.712342\du}}
\pgfpathlineto{\pgfpoint{-3.473419\du}{6.714432\du}}
\pgfpathlineto{\pgfpoint{-3.467150\du}{6.717119\du}}
\pgfpathlineto{\pgfpoint{-3.461179\du}{6.719209\du}}
\pgfpathlineto{\pgfpoint{-3.455208\du}{6.721597\du}}
\pgfpathlineto{\pgfpoint{-3.447744\du}{6.723687\du}}
\pgfpathlineto{\pgfpoint{-3.441176\du}{6.725777\du}}
\pgfpathlineto{\pgfpoint{-3.433712\du}{6.727568\du}}
\pgfpathlineto{\pgfpoint{-3.425651\du}{6.729061\du}}
\pgfpathlineto{\pgfpoint{-3.417889\du}{6.731151\du}}
\pgfpathlineto{\pgfpoint{-3.409529\du}{6.732643\du}}
\pgfpathlineto{\pgfpoint{-3.400573\du}{6.734136\du}}
\pgfpathlineto{\pgfpoint{-3.392213\du}{6.735330\du}}
\pgfpathlineto{\pgfpoint{-3.382958\du}{6.736525\du}}
\pgfpathlineto{\pgfpoint{-3.374001\du}{6.737122\du}}
\pgfpathlineto{\pgfpoint{-3.364746\du}{6.738316\du}}
\pgfpathlineto{\pgfpoint{-3.355193\du}{6.738614\du}}
\pgfpathlineto{\pgfpoint{-3.345639\du}{6.739212\du}}
\pgfpathlineto{\pgfpoint{-3.336085\du}{6.739510\du}}
\pgfpathlineto{\pgfpoint{-3.326233\du}{6.740107\du}}
\pgfpathlineto{\pgfpoint{-3.316978\du}{6.740107\du}}
\pgfpathlineto{\pgfpoint{-3.316978\du}{6.740107\du}}
\pgfpathlineto{\pgfpoint{-3.307424\du}{6.740107\du}}
\pgfpathlineto{\pgfpoint{-3.297572\du}{6.739510\du}}
\pgfpathlineto{\pgfpoint{-3.287720\du}{6.739212\du}}
\pgfpathlineto{\pgfpoint{-3.278465\du}{6.738614\du}}
\pgfpathlineto{\pgfpoint{-3.268911\du}{6.738316\du}}
\pgfpathlineto{\pgfpoint{-3.259358\du}{6.737122\du}}
\pgfpathlineto{\pgfpoint{-3.250700\du}{6.736525\du}}
\pgfpathlineto{\pgfpoint{-3.241444\du}{6.735330\du}}
\pgfpathlineto{\pgfpoint{-3.233085\du}{6.734136\du}}
\pgfpathlineto{\pgfpoint{-3.224725\du}{6.732643\du}}
\pgfpathlineto{\pgfpoint{-3.215769\du}{6.731151\du}}
\pgfpathlineto{\pgfpoint{-3.208007\du}{6.729061\du}}
\pgfpathlineto{\pgfpoint{-3.200543\du}{6.727568\du}}
\pgfpathlineto{\pgfpoint{-3.192482\du}{6.725777\du}}
\pgfpathlineto{\pgfpoint{-3.185914\du}{6.723687\du}}
\pgfpathlineto{\pgfpoint{-3.179047\du}{6.721597\du}}
\pgfpathlineto{\pgfpoint{-3.172479\du}{6.719209\du}}
\pgfpathlineto{\pgfpoint{-3.166508\du}{6.717119\du}}
\pgfpathlineto{\pgfpoint{-3.160835\du}{6.714432\du}}
\pgfpathlineto{\pgfpoint{-3.155760\du}{6.712342\du}}
\pgfpathlineto{\pgfpoint{-3.150685\du}{6.709655\du}}
\pgfpathlineto{\pgfpoint{-3.145908\du}{6.706968\du}}
\pgfpathlineto{\pgfpoint{-3.141728\du}{6.704281\du}}
\pgfpathlineto{\pgfpoint{-3.138444\du}{6.701594\du}}
\pgfpathlineto{\pgfpoint{-3.135160\du}{6.698907\du}}
\pgfpathlineto{\pgfpoint{-3.133070\du}{6.695324\du}}
\pgfpathlineto{\pgfpoint{-3.130980\du}{6.692637\du}}
\pgfpathlineto{\pgfpoint{-3.128890\du}{6.689353\du}}
\pgfpathlineto{\pgfpoint{-3.127696\du}{6.686666\du}}
\pgfpathlineto{\pgfpoint{-3.127099\du}{6.683979\du}}
\pgfpathlineto{\pgfpoint{-3.126800\du}{6.680397\du}}
\pgfusepath{stroke}
\pgfsetbuttcap
\pgfsetmiterjoin
\pgfsetdash{}{0pt}
\definecolor{dialinecolor}{rgb}{0.000000, 0.000000, 0.000000}
\pgfsetfillcolor{dialinecolor}
\pgfpathmoveto{\pgfpoint{-3.305036\du}{6.668455\du}}
\pgfpathlineto{\pgfpoint{-3.281152\du}{6.680397\du}}
\pgfpathlineto{\pgfpoint{-3.209798\du}{6.656811\du}}
\pgfpathlineto{\pgfpoint{-3.173972\du}{6.668455\du}}
\pgfpathlineto{\pgfpoint{-3.197557\du}{6.644571\du}}
\pgfpathlineto{\pgfpoint{-3.281152\du}{6.644571\du}}
\pgfpathlineto{\pgfpoint{-3.245326\du}{6.644571\du}}
\pgfpathlineto{\pgfpoint{-3.305036\du}{6.668455\du}}
\pgfusepath{fill}
\pgfsetbuttcap
\pgfsetmiterjoin
\pgfsetdash{}{0pt}
\definecolor{dialinecolor}{rgb}{0.000000, 0.000000, 0.000000}
\pgfsetfillcolor{dialinecolor}
\pgfpathmoveto{\pgfpoint{-3.316978\du}{6.704281\du}}
\pgfpathlineto{\pgfpoint{-3.352506\du}{6.692339\du}}
\pgfpathlineto{\pgfpoint{-3.412216\du}{6.716223\du}}
\pgfpathlineto{\pgfpoint{-3.447744\du}{6.704281\du}}
\pgfpathlineto{\pgfpoint{-3.435503\du}{6.740107\du}}
\pgfpathlineto{\pgfpoint{-3.352506\du}{6.740107\du}}
\pgfpathlineto{\pgfpoint{-3.388033\du}{6.728464\du}}
\pgfpathlineto{\pgfpoint{-3.316978\du}{6.704281\du}}
\pgfusepath{fill}
\pgfsetbuttcap
\pgfsetmiterjoin
\pgfsetdash{}{0pt}
\definecolor{dialinecolor}{rgb}{0.000000, 0.000000, 0.000000}
\pgfsetfillcolor{dialinecolor}
\pgfpathmoveto{\pgfpoint{-3.435503\du}{6.644571\du}}
\pgfpathlineto{\pgfpoint{-3.412216\du}{6.632927\du}}
\pgfpathlineto{\pgfpoint{-3.340862\du}{6.668455\du}}
\pgfpathlineto{\pgfpoint{-3.316978\du}{6.656811\du}}
\pgfpathlineto{\pgfpoint{-3.328323\du}{6.680397\du}}
\pgfpathlineto{\pgfpoint{-3.412216\du}{6.680397\du}}
\pgfpathlineto{\pgfpoint{-3.376390\du}{6.668455\du}}
\pgfpathlineto{\pgfpoint{-3.435503\du}{6.644571\du}}
\pgfusepath{fill}
\pgfsetbuttcap
\pgfsetmiterjoin
\pgfsetdash{}{0pt}
\definecolor{dialinecolor}{rgb}{0.000000, 0.000000, 0.000000}
\pgfsetfillcolor{dialinecolor}
\pgfpathmoveto{\pgfpoint{-3.185914\du}{6.728464\du}}
\pgfpathlineto{\pgfpoint{-3.209798\du}{6.740107\du}}
\pgfpathlineto{\pgfpoint{-3.281152\du}{6.704281\du}}
\pgfpathlineto{\pgfpoint{-3.316978\du}{6.716223\du}}
\pgfpathlineto{\pgfpoint{-3.292795\du}{6.692339\du}}
\pgfpathlineto{\pgfpoint{-3.209798\du}{6.692339\du}}
\pgfpathlineto{\pgfpoint{-3.245326\du}{6.704281\du}}
\pgfpathlineto{\pgfpoint{-3.185914\du}{6.728464\du}}
\pgfusepath{fill}
\pgfsetbuttcap
\pgfsetmiterjoin
\pgfsetdash{}{0pt}
\definecolor{dialinecolor}{rgb}{1.000000, 1.000000, 1.000000}
\pgfsetfillcolor{dialinecolor}
\pgfpathmoveto{\pgfpoint{-3.305036\du}{6.680397\du}}
\pgfpathlineto{\pgfpoint{-3.281152\du}{6.680397\du}}
\pgfpathlineto{\pgfpoint{-3.209798\du}{6.656811\du}}
\pgfpathlineto{\pgfpoint{-3.173972\du}{6.668455\du}}
\pgfpathlineto{\pgfpoint{-3.197557\du}{6.644571\du}}
\pgfpathlineto{\pgfpoint{-3.268911\du}{6.644571\du}}
\pgfpathlineto{\pgfpoint{-3.245326\du}{6.644571\du}}
\pgfpathlineto{\pgfpoint{-3.305036\du}{6.680397\du}}
\pgfusepath{fill}
\pgfsetbuttcap
\pgfsetmiterjoin
\pgfsetdash{}{0pt}
\definecolor{dialinecolor}{rgb}{1.000000, 1.000000, 1.000000}
\pgfsetfillcolor{dialinecolor}
\pgfpathmoveto{\pgfpoint{-3.316978\du}{6.704281\du}}
\pgfpathlineto{\pgfpoint{-3.340862\du}{6.692339\du}}
\pgfpathlineto{\pgfpoint{-3.412216\du}{6.716223\du}}
\pgfpathlineto{\pgfpoint{-3.447744\du}{6.716223\du}}
\pgfpathlineto{\pgfpoint{-3.423860\du}{6.740107\du}}
\pgfpathlineto{\pgfpoint{-3.340862\du}{6.740107\du}}
\pgfpathlineto{\pgfpoint{-3.376390\du}{6.728464\du}}
\pgfpathlineto{\pgfpoint{-3.316978\du}{6.704281\du}}
\pgfusepath{fill}
\pgfsetbuttcap
\pgfsetmiterjoin
\pgfsetdash{}{0pt}
\definecolor{dialinecolor}{rgb}{1.000000, 1.000000, 1.000000}
\pgfsetfillcolor{dialinecolor}
\pgfpathmoveto{\pgfpoint{-3.435503\du}{6.644571\du}}
\pgfpathlineto{\pgfpoint{-3.412216\du}{6.644571\du}}
\pgfpathlineto{\pgfpoint{-3.340862\du}{6.668455\du}}
\pgfpathlineto{\pgfpoint{-3.305036\du}{6.656811\du}}
\pgfpathlineto{\pgfpoint{-3.328323\du}{6.680397\du}}
\pgfpathlineto{\pgfpoint{-3.412216\du}{6.680397\du}}
\pgfpathlineto{\pgfpoint{-3.376390\du}{6.680397\du}}
\pgfpathlineto{\pgfpoint{-3.435503\du}{6.644571\du}}
\pgfusepath{fill}
\pgfsetbuttcap
\pgfsetmiterjoin
\pgfsetdash{}{0pt}
\definecolor{dialinecolor}{rgb}{1.000000, 1.000000, 1.000000}
\pgfsetfillcolor{dialinecolor}
\pgfpathmoveto{\pgfpoint{-3.185914\du}{6.728464\du}}
\pgfpathlineto{\pgfpoint{-3.209798\du}{6.740107\du}}
\pgfpathlineto{\pgfpoint{-3.281152\du}{6.716223\du}}
\pgfpathlineto{\pgfpoint{-3.305036\du}{6.716223\du}}
\pgfpathlineto{\pgfpoint{-3.292795\du}{6.692339\du}}
\pgfpathlineto{\pgfpoint{-3.209798\du}{6.692339\du}}
\pgfpathlineto{\pgfpoint{-3.245326\du}{6.704281\du}}
\pgfpathlineto{\pgfpoint{-3.185914\du}{6.728464\du}}
\pgfusepath{fill}
\pgfsetbuttcap
\pgfsetmiterjoin
\pgfsetdash{}{0pt}
\definecolor{dialinecolor}{rgb}{0.666667, 0.901961, 1.000000}
\pgfsetstrokecolor{dialinecolor}
\pgfpathmoveto{\pgfpoint{-3.507454\du}{6.680397\du}}
\pgfpathlineto{\pgfpoint{-3.507454\du}{6.775336\du}}
\pgfusepath{stroke}
\pgfsetbuttcap
\pgfsetmiterjoin
\pgfsetdash{}{0pt}
\definecolor{dialinecolor}{rgb}{0.666667, 0.901961, 1.000000}
\pgfsetstrokecolor{dialinecolor}
\pgfpathmoveto{\pgfpoint{-3.126800\du}{6.680397\du}}
\pgfpathlineto{\pgfpoint{-3.126800\du}{6.775336\du}}
\pgfusepath{stroke}
\pgfsetlinewidth{0.100000\du}
\pgfsetdash{}{0pt}
\pgfsetdash{}{0pt}
\pgfsetbuttcap
\pgfsetmiterjoin
\pgfsetlinewidth{0.001000\du}
\pgfsetbuttcap
\pgfsetmiterjoin
\pgfsetdash{}{0pt}
\definecolor{dialinecolor}{rgb}{0.000000, 0.470588, 0.666667}
\pgfsetstrokecolor{dialinecolor}
\pgfpathmoveto{\pgfpoint{3.654317\du}{7.303354\du}}
\pgfpathlineto{\pgfpoint{3.660885\du}{7.294397\du}}
\pgfpathlineto{\pgfpoint{3.667155\du}{7.284844\du}}
\pgfpathlineto{\pgfpoint{3.672529\du}{7.274693\du}}
\pgfpathlineto{\pgfpoint{3.677604\du}{7.265438\du}}
\pgfpathlineto{\pgfpoint{3.681485\du}{7.255287\du}}
\pgfpathlineto{\pgfpoint{3.684769\du}{7.245136\du}}
\pgfpathlineto{\pgfpoint{3.687755\du}{7.234986\du}}
\pgfpathlineto{\pgfpoint{3.689546\du}{7.224536\du}}
\pgfpathlineto{\pgfpoint{3.690740\du}{7.214385\du}}
\pgfpathlineto{\pgfpoint{3.691636\du}{7.203936\du}}
\pgfpathlineto{\pgfpoint{3.691636\du}{7.193487\du}}
\pgfpathlineto{\pgfpoint{3.690740\du}{7.182739\du}}
\pgfpathlineto{\pgfpoint{3.689546\du}{7.172588\du}}
\pgfpathlineto{\pgfpoint{3.687755\du}{7.162139\du}}
\pgfpathlineto{\pgfpoint{3.685068\du}{7.152287\du}}
\pgfpathlineto{\pgfpoint{3.681485\du}{7.141837\du}}
\pgfpathlineto{\pgfpoint{3.677604\du}{7.131985\du}}
\pgfpathlineto{\pgfpoint{3.672529\du}{7.121834\du}}
\pgfpathlineto{\pgfpoint{3.667155\du}{7.112281\du}}
\pgfpathlineto{\pgfpoint{3.660885\du}{7.103324\du}}
\pgfpathlineto{\pgfpoint{3.654317\du}{7.093770\du}}
\pgfpathlineto{\pgfpoint{3.646853\du}{7.085112\du}}
\pgfpathlineto{\pgfpoint{3.639389\du}{7.076454\du}}
\pgfpathlineto{\pgfpoint{3.631030\du}{7.068095\du}}
\pgfpathlineto{\pgfpoint{3.622372\du}{7.060034\du}}
\pgfpathlineto{\pgfpoint{3.613117\du}{7.052272\du}}
\pgfpathlineto{\pgfpoint{3.603563\du}{7.045106\du}}
\pgfpathlineto{\pgfpoint{3.592815\du}{7.038538\du}}
\pgfpathlineto{\pgfpoint{3.582366\du}{7.031970\du}}
\pgfpathlineto{\pgfpoint{3.571021\du}{7.025701\du}}
\pgfpathlineto{\pgfpoint{3.559079\du}{7.020028\du}}
\pgfpathlineto{\pgfpoint{3.547137\du}{7.014654\du}}
\pgfpathlineto{\pgfpoint{3.534896\du}{7.010474\du}}
\pgfpathlineto{\pgfpoint{3.522357\du}{7.006295\du}}
\pgfpathlineto{\pgfpoint{3.509221\du}{7.002115\du}}
\pgfpathlineto{\pgfpoint{3.496383\du}{6.999129\du}}
\pgfpathlineto{\pgfpoint{3.483247\du}{6.996442\du}}
\pgfpathlineto{\pgfpoint{3.469513\du}{6.994054\du}}
\pgfpathlineto{\pgfpoint{3.456377\du}{6.992561\du}}
\pgfpathlineto{\pgfpoint{3.442345\du}{6.991367\du}}
\pgfpathlineto{\pgfpoint{3.428612\du}{6.990471\du}}
\pgfpathlineto{\pgfpoint{3.414580\du}{6.990471\du}}
\pgfpathlineto{\pgfpoint{3.401145\du}{6.991069\du}}
\pgfpathlineto{\pgfpoint{3.387411\du}{6.991666\du}}
\pgfpathlineto{\pgfpoint{3.373379\du}{6.993158\du}}
\pgfpathlineto{\pgfpoint{3.360243\du}{6.995547\du}}
\pgfpathlineto{\pgfpoint{3.346510\du}{6.997935\du}}
\pgfpathlineto{\pgfpoint{3.333971\du}{7.000921\du}}
\pgfpathlineto{\pgfpoint{3.320834\du}{7.003906\du}}
\pgfpathlineto{\pgfpoint{3.308295\du}{7.008086\du}}
\pgfpathlineto{\pgfpoint{3.295457\du}{7.012863\du}}
\pgfpathlineto{\pgfpoint{3.283217\du}{7.017341\du}}
\pgfpathlineto{\pgfpoint{3.271872\du}{7.022715\du}}
\pgfpathlineto{\pgfpoint{3.260228\du}{7.028985\du}}
\pgfpathlineto{\pgfpoint{3.249182\du}{7.034956\du}}
\pgfpathlineto{\pgfpoint{3.238732\du}{7.041822\du}}
\pgfpathlineto{\pgfpoint{3.228582\du}{7.048391\du}}
\pgfpathlineto{\pgfpoint{3.219028\du}{7.056153\du}}
\pgfpathlineto{\pgfpoint{3.209474\du}{7.064214\du}}
\pgfpathlineto{\pgfpoint{3.201115\du}{7.072275\du}}
\pgfpathlineto{\pgfpoint{3.193054\du}{7.080634\du}}
\pgfpathlineto{\pgfpoint{3.185292\du}{7.089591\du}}
\pgfusepath{stroke}
\pgfsetlinewidth{0.080000\du}
\pgfsetbuttcap
\pgfsetmiterjoin
\pgfsetdash{}{0pt}
\definecolor{dialinecolor}{rgb}{0.666667, 0.901961, 1.000000}
\pgfsetstrokecolor{dialinecolor}
\pgfpathmoveto{\pgfpoint{3.645062\du}{7.295890\du}}
\pgfpathlineto{\pgfpoint{3.650436\du}{7.286934\du}}
\pgfpathlineto{\pgfpoint{3.656705\du}{7.278276\du}}
\pgfpathlineto{\pgfpoint{3.661781\du}{7.269020\du}}
\pgfpathlineto{\pgfpoint{3.665960\du}{7.259765\du}}
\pgfpathlineto{\pgfpoint{3.669842\du}{7.250212\du}}
\pgfpathlineto{\pgfpoint{3.672827\du}{7.240658\du}}
\pgfpathlineto{\pgfpoint{3.675514\du}{7.230806\du}}
\pgfpathlineto{\pgfpoint{3.677604\du}{7.221551\du}}
\pgfpathlineto{\pgfpoint{3.678500\du}{7.211698\du}}
\pgfpathlineto{\pgfpoint{3.679097\du}{7.201548\du}}
\pgfpathlineto{\pgfpoint{3.679097\du}{7.191695\du}}
\pgfpathlineto{\pgfpoint{3.678201\du}{7.181843\du}}
\pgfpathlineto{\pgfpoint{3.676708\du}{7.171693\du}}
\pgfpathlineto{\pgfpoint{3.674917\du}{7.162437\du}}
\pgfpathlineto{\pgfpoint{3.671931\du}{7.152585\du}}
\pgfpathlineto{\pgfpoint{3.668349\du}{7.143032\du}}
\pgfpathlineto{\pgfpoint{3.664468\du}{7.133776\du}}
\pgfpathlineto{\pgfpoint{3.659989\du}{7.124521\du}}
\pgfpathlineto{\pgfpoint{3.654914\du}{7.115863\du}}
\pgfpathlineto{\pgfpoint{3.649242\du}{7.106608\du}}
\pgfpathlineto{\pgfpoint{3.642076\du}{7.098249\du}}
\pgfpathlineto{\pgfpoint{3.635508\du}{7.089591\du}}
\pgfpathlineto{\pgfpoint{3.627746\du}{7.081530\du}}
\pgfpathlineto{\pgfpoint{3.619685\du}{7.073767\du}}
\pgfpathlineto{\pgfpoint{3.611325\du}{7.066901\du}}
\pgfpathlineto{\pgfpoint{3.602667\du}{7.059437\du}}
\pgfpathlineto{\pgfpoint{3.592815\du}{7.052570\du}}
\pgfpathlineto{\pgfpoint{3.582366\du}{7.045704\du}}
\pgfpathlineto{\pgfpoint{3.572215\du}{7.040031\du}}
\pgfpathlineto{\pgfpoint{3.561766\du}{7.034657\du}}
\pgfpathlineto{\pgfpoint{3.549824\du}{7.029283\du}}
\pgfpathlineto{\pgfpoint{3.539076\du}{7.024208\du}}
\pgfpathlineto{\pgfpoint{3.526537\du}{7.020028\du}}
\pgfpathlineto{\pgfpoint{3.515192\du}{7.016147\du}}
\pgfpathlineto{\pgfpoint{3.502652\du}{7.012863\du}}
\pgfpathlineto{\pgfpoint{3.490113\du}{7.009877\du}}
\pgfpathlineto{\pgfpoint{3.477574\du}{7.007190\du}}
\pgfpathlineto{\pgfpoint{3.464438\du}{7.005399\du}}
\pgfpathlineto{\pgfpoint{3.451899\du}{7.003906\du}}
\pgfpathlineto{\pgfpoint{3.438762\du}{7.002712\du}}
\pgfpathlineto{\pgfpoint{3.424730\du}{7.002115\du}}
\pgfpathlineto{\pgfpoint{3.411594\du}{7.002115\du}}
\pgfpathlineto{\pgfpoint{3.398756\du}{7.002712\du}}
\pgfpathlineto{\pgfpoint{3.385620\du}{7.003608\du}}
\pgfpathlineto{\pgfpoint{3.373081\du}{7.005399\du}}
\pgfpathlineto{\pgfpoint{3.359945\du}{7.007190\du}}
\pgfpathlineto{\pgfpoint{3.347405\du}{7.009280\du}}
\pgfpathlineto{\pgfpoint{3.334568\du}{7.012863\du}}
\pgfpathlineto{\pgfpoint{3.322327\du}{7.015848\du}}
\pgfpathlineto{\pgfpoint{3.310385\du}{7.019730\du}}
\pgfpathlineto{\pgfpoint{3.298144\du}{7.023909\du}}
\pgfpathlineto{\pgfpoint{3.286799\du}{7.028985\du}}
\pgfpathlineto{\pgfpoint{3.275753\du}{7.033761\du}}
\pgfpathlineto{\pgfpoint{3.265005\du}{7.039733\du}}
\pgfpathlineto{\pgfpoint{3.254556\du}{7.045405\du}}
\pgfpathlineto{\pgfpoint{3.244106\du}{7.051973\du}}
\pgfpathlineto{\pgfpoint{3.234553\du}{7.058541\du}}
\pgfpathlineto{\pgfpoint{3.225895\du}{7.065707\du}}
\pgfpathlineto{\pgfpoint{3.216341\du}{7.072872\du}}
\pgfpathlineto{\pgfpoint{3.208579\du}{7.080634\du}}
\pgfpathlineto{\pgfpoint{3.201115\du}{7.088695\du}}
\pgfpathlineto{\pgfpoint{3.193950\du}{7.097055\du}}
\pgfusepath{stroke}
\pgfsetbuttcap
\pgfsetmiterjoin
\pgfsetdash{}{0pt}
\definecolor{dialinecolor}{rgb}{0.000000, 0.705882, 1.000000}
\pgfsetfillcolor{dialinecolor}
\pgfpathmoveto{\pgfpoint{1.728059\du}{6.978828\du}}
\pgfpathlineto{\pgfpoint{1.930776\du}{6.788053\du}}
\pgfpathlineto{\pgfpoint{3.442345\du}{6.788053\du}}
\pgfpathlineto{\pgfpoint{3.239628\du}{6.978828\du}}
\pgfpathlineto{\pgfpoint{1.728059\du}{6.978828\du}}
\pgfusepath{fill}
\pgfsetlinewidth{0.001000\du}
\pgfsetbuttcap
\pgfsetmiterjoin
\pgfsetdash{}{0pt}
\definecolor{dialinecolor}{rgb}{0.666667, 0.901961, 1.000000}
\pgfsetstrokecolor{dialinecolor}
\pgfpathmoveto{\pgfpoint{1.728059\du}{6.978828\du}}
\pgfpathlineto{\pgfpoint{1.930776\du}{6.788053\du}}
\pgfpathlineto{\pgfpoint{3.442345\du}{6.788053\du}}
\pgfpathlineto{\pgfpoint{3.239628\du}{6.978828\du}}
\pgfpathlineto{\pgfpoint{1.728059\du}{6.978828\du}}
\pgfusepath{stroke}
\pgfsetbuttcap
\pgfsetmiterjoin
\pgfsetdash{}{0pt}
\definecolor{dialinecolor}{rgb}{0.000000, 0.588235, 0.831373}
\pgfsetfillcolor{dialinecolor}
\pgfpathmoveto{\pgfpoint{1.728059\du}{6.978828\du}}
\pgfpathlineto{\pgfpoint{3.239628\du}{6.978828\du}}
\pgfpathlineto{\pgfpoint{3.239628\du}{7.264841\du}}
\pgfpathlineto{\pgfpoint{1.728059\du}{7.264841\du}}
\pgfpathlineto{\pgfpoint{1.728059\du}{6.978828\du}}
\pgfusepath{fill}
\pgfsetbuttcap
\pgfsetmiterjoin
\pgfsetdash{}{0pt}
\definecolor{dialinecolor}{rgb}{0.666667, 0.901961, 1.000000}
\pgfsetstrokecolor{dialinecolor}
\pgfpathmoveto{\pgfpoint{1.728059\du}{6.978828\du}}
\pgfpathlineto{\pgfpoint{3.239628\du}{6.978828\du}}
\pgfpathlineto{\pgfpoint{3.239628\du}{7.264841\du}}
\pgfpathlineto{\pgfpoint{1.728059\du}{7.264841\du}}
\pgfpathlineto{\pgfpoint{1.728059\du}{6.978828\du}}
\pgfusepath{stroke}
\pgfsetbuttcap
\pgfsetmiterjoin
\pgfsetdash{}{0pt}
\definecolor{dialinecolor}{rgb}{0.000000, 0.352941, 0.501961}
\pgfsetfillcolor{dialinecolor}
\pgfpathmoveto{\pgfpoint{3.239628\du}{7.264841\du}}
\pgfpathlineto{\pgfpoint{3.442345\du}{7.062124\du}}
\pgfpathlineto{\pgfpoint{3.442345\du}{6.788053\du}}
\pgfpathlineto{\pgfpoint{3.239628\du}{6.978828\du}}
\pgfpathlineto{\pgfpoint{3.239628\du}{7.264841\du}}
\pgfusepath{fill}
\pgfsetbuttcap
\pgfsetmiterjoin
\pgfsetdash{}{0pt}
\definecolor{dialinecolor}{rgb}{0.666667, 0.901961, 1.000000}
\pgfsetstrokecolor{dialinecolor}
\pgfpathmoveto{\pgfpoint{3.239628\du}{7.264841\du}}
\pgfpathlineto{\pgfpoint{3.442345\du}{7.062124\du}}
\pgfpathlineto{\pgfpoint{3.442345\du}{6.788053\du}}
\pgfpathlineto{\pgfpoint{3.239628\du}{6.978828\du}}
\pgfpathlineto{\pgfpoint{3.239628\du}{7.264841\du}}
\pgfusepath{stroke}
\pgfsetbuttcap
\pgfsetmiterjoin
\pgfsetdash{}{0pt}
\definecolor{dialinecolor}{rgb}{0.000000, 0.000000, 0.000000}
\pgfsetfillcolor{dialinecolor}
\pgfpathmoveto{\pgfpoint{1.776126\du}{6.942703\du}}
\pgfpathlineto{\pgfpoint{1.930776\du}{6.788053\du}}
\pgfpathlineto{\pgfpoint{3.406220\du}{6.788053\du}}
\pgfpathlineto{\pgfpoint{3.251869\du}{6.942703\du}}
\pgfpathlineto{\pgfpoint{1.776126\du}{6.942703\du}}
\pgfusepath{fill}
\pgfsetbuttcap
\pgfsetmiterjoin
\pgfsetdash{}{0pt}
\definecolor{dialinecolor}{rgb}{0.000000, 0.000000, 0.000000}
\pgfsetstrokecolor{dialinecolor}
\pgfpathmoveto{\pgfpoint{1.776126\du}{6.942703\du}}
\pgfpathlineto{\pgfpoint{1.930776\du}{6.788053\du}}
\pgfpathlineto{\pgfpoint{3.406220\du}{6.788053\du}}
\pgfpathlineto{\pgfpoint{3.251869\du}{6.942703\du}}
\pgfpathlineto{\pgfpoint{1.776126\du}{6.942703\du}}
\pgfusepath{stroke}
\pgfsetbuttcap
\pgfsetmiterjoin
\pgfsetdash{}{0pt}
\definecolor{dialinecolor}{rgb}{0.000000, 0.705882, 1.000000}
\pgfsetfillcolor{dialinecolor}
\pgfpathmoveto{\pgfpoint{1.728059\du}{5.716550\du}}
\pgfpathlineto{\pgfpoint{1.895248\du}{5.561900\du}}
\pgfpathlineto{\pgfpoint{3.406220\du}{5.561900\du}}
\pgfpathlineto{\pgfpoint{3.239628\du}{5.716550\du}}
\pgfpathlineto{\pgfpoint{1.728059\du}{5.716550\du}}
\pgfusepath{fill}
\pgfsetbuttcap
\pgfsetmiterjoin
\pgfsetdash{}{0pt}
\definecolor{dialinecolor}{rgb}{0.666667, 0.901961, 1.000000}
\pgfsetstrokecolor{dialinecolor}
\pgfpathmoveto{\pgfpoint{1.728059\du}{5.716550\du}}
\pgfpathlineto{\pgfpoint{1.895248\du}{5.561900\du}}
\pgfpathlineto{\pgfpoint{3.406220\du}{5.561900\du}}
\pgfpathlineto{\pgfpoint{3.239628\du}{5.716550\du}}
\pgfpathlineto{\pgfpoint{1.728059\du}{5.716550\du}}
\pgfusepath{stroke}
\pgfsetbuttcap
\pgfsetmiterjoin
\pgfsetdash{}{0pt}
\definecolor{dialinecolor}{rgb}{0.000000, 0.588235, 0.831373}
\pgfsetfillcolor{dialinecolor}
\pgfpathmoveto{\pgfpoint{1.728059\du}{5.716550\du}}
\pgfpathlineto{\pgfpoint{3.251869\du}{5.716550\du}}
\pgfpathlineto{\pgfpoint{3.251869\du}{6.918819\du}}
\pgfpathlineto{\pgfpoint{1.728059\du}{6.918819\du}}
\pgfpathlineto{\pgfpoint{1.728059\du}{5.716550\du}}
\pgfusepath{fill}
\pgfsetbuttcap
\pgfsetmiterjoin
\pgfsetdash{}{0pt}
\definecolor{dialinecolor}{rgb}{0.666667, 0.901961, 1.000000}
\pgfsetstrokecolor{dialinecolor}
\pgfpathmoveto{\pgfpoint{1.728059\du}{5.716550\du}}
\pgfpathlineto{\pgfpoint{3.251570\du}{5.716550\du}}
\pgfpathlineto{\pgfpoint{3.251570\du}{6.918520\du}}
\pgfpathlineto{\pgfpoint{1.728059\du}{6.918520\du}}
\pgfpathlineto{\pgfpoint{1.728059\du}{5.716550\du}}
\pgfusepath{stroke}
\pgfsetbuttcap
\pgfsetmiterjoin
\pgfsetdash{}{0pt}
\definecolor{dialinecolor}{rgb}{1.000000, 1.000000, 1.000000}
\pgfsetfillcolor{dialinecolor}
\pgfpathmoveto{\pgfpoint{1.859124\du}{5.871498\du}}
\pgfpathlineto{\pgfpoint{3.120804\du}{5.871498\du}}
\pgfpathlineto{\pgfpoint{3.120804\du}{6.800294\du}}
\pgfpathlineto{\pgfpoint{1.859124\du}{6.800294\du}}
\pgfpathlineto{\pgfpoint{1.859124\du}{5.871498\du}}
\pgfusepath{fill}
\pgfsetbuttcap
\pgfsetmiterjoin
\pgfsetdash{}{0pt}
\definecolor{dialinecolor}{rgb}{0.000000, 0.235294, 0.333333}
\pgfsetstrokecolor{dialinecolor}
\pgfpathmoveto{\pgfpoint{1.859124\du}{5.871498\du}}
\pgfpathlineto{\pgfpoint{3.120506\du}{5.871498\du}}
\pgfpathlineto{\pgfpoint{3.120506\du}{6.799995\du}}
\pgfpathlineto{\pgfpoint{1.859124\du}{6.799995\du}}
\pgfpathlineto{\pgfpoint{1.859124\du}{5.871498\du}}
\pgfusepath{stroke}
\pgfsetbuttcap
\pgfsetmiterjoin
\pgfsetdash{}{0pt}
\definecolor{dialinecolor}{rgb}{0.000000, 0.352941, 0.501961}
\pgfsetfillcolor{dialinecolor}
\pgfpathmoveto{\pgfpoint{3.239628\du}{6.907474\du}}
\pgfpathlineto{\pgfpoint{3.406220\du}{6.740583\du}}
\pgfpathlineto{\pgfpoint{3.406220\du}{5.561900\du}}
\pgfpathlineto{\pgfpoint{3.239628\du}{5.716550\du}}
\pgfpathlineto{\pgfpoint{3.239628\du}{6.907474\du}}
\pgfusepath{fill}
\pgfsetbuttcap
\pgfsetmiterjoin
\pgfsetdash{}{0pt}
\definecolor{dialinecolor}{rgb}{0.666667, 0.901961, 1.000000}
\pgfsetstrokecolor{dialinecolor}
\pgfpathmoveto{\pgfpoint{3.239628\du}{6.907474\du}}
\pgfpathlineto{\pgfpoint{3.406220\du}{6.740583\du}}
\pgfpathlineto{\pgfpoint{3.406220\du}{5.561900\du}}
\pgfpathlineto{\pgfpoint{3.239628\du}{5.716550\du}}
\pgfpathlineto{\pgfpoint{3.239628\du}{6.907474\du}}
\pgfusepath{stroke}
\pgfsetbuttcap
\pgfsetmiterjoin
\pgfsetdash{}{0pt}
\definecolor{dialinecolor}{rgb}{0.000000, 0.705882, 1.000000}
\pgfsetfillcolor{dialinecolor}
\pgfpathmoveto{\pgfpoint{1.431000\du}{7.502787\du}}
\pgfpathlineto{\pgfpoint{1.668647\du}{7.204832\du}}
\pgfpathlineto{\pgfpoint{3.311281\du}{7.204832\du}}
\pgfpathlineto{\pgfpoint{3.073335\du}{7.502787\du}}
\pgfpathlineto{\pgfpoint{1.431000\du}{7.502787\du}}
\pgfusepath{fill}
\pgfsetbuttcap
\pgfsetmiterjoin
\pgfsetdash{}{0pt}
\definecolor{dialinecolor}{rgb}{0.666667, 0.901961, 1.000000}
\pgfsetstrokecolor{dialinecolor}
\pgfpathmoveto{\pgfpoint{1.431000\du}{7.502787\du}}
\pgfpathlineto{\pgfpoint{1.668647\du}{7.204832\du}}
\pgfpathlineto{\pgfpoint{3.311281\du}{7.204832\du}}
\pgfpathlineto{\pgfpoint{3.073335\du}{7.502787\du}}
\pgfpathlineto{\pgfpoint{1.431000\du}{7.502787\du}}
\pgfusepath{stroke}
\pgfsetbuttcap
\pgfsetmiterjoin
\pgfsetdash{}{0pt}
\definecolor{dialinecolor}{rgb}{0.000000, 0.352941, 0.501961}
\pgfsetfillcolor{dialinecolor}
\pgfpathmoveto{\pgfpoint{3.073335\du}{7.561900\du}}
\pgfpathlineto{\pgfpoint{3.311281\du}{7.312012\du}}
\pgfpathlineto{\pgfpoint{3.311281\du}{7.204832\du}}
\pgfpathlineto{\pgfpoint{3.073335\du}{7.502787\du}}
\pgfpathlineto{\pgfpoint{3.073335\du}{7.561900\du}}
\pgfusepath{fill}
\pgfsetbuttcap
\pgfsetmiterjoin
\pgfsetdash{}{0pt}
\definecolor{dialinecolor}{rgb}{0.666667, 0.901961, 1.000000}
\pgfsetstrokecolor{dialinecolor}
\pgfpathmoveto{\pgfpoint{3.073335\du}{7.561900\du}}
\pgfpathlineto{\pgfpoint{3.311281\du}{7.312012\du}}
\pgfpathlineto{\pgfpoint{3.311281\du}{7.204832\du}}
\pgfpathlineto{\pgfpoint{3.073335\du}{7.502787\du}}
\pgfpathlineto{\pgfpoint{3.073335\du}{7.561900\du}}
\pgfusepath{stroke}
\pgfsetbuttcap
\pgfsetmiterjoin
\pgfsetdash{}{0pt}
\definecolor{dialinecolor}{rgb}{0.000000, 0.588235, 0.831373}
\pgfsetfillcolor{dialinecolor}
\pgfpathmoveto{\pgfpoint{1.431000\du}{7.502787\du}}
\pgfpathlineto{\pgfpoint{3.073335\du}{7.502787\du}}
\pgfpathlineto{\pgfpoint{3.073335\du}{7.561900\du}}
\pgfpathlineto{\pgfpoint{1.431000\du}{7.561900\du}}
\pgfpathlineto{\pgfpoint{1.431000\du}{7.502787\du}}
\pgfusepath{fill}
\pgfsetbuttcap
\pgfsetmiterjoin
\pgfsetdash{}{0pt}
\definecolor{dialinecolor}{rgb}{0.666667, 0.901961, 1.000000}
\pgfsetstrokecolor{dialinecolor}
\pgfpathmoveto{\pgfpoint{1.431000\du}{7.502787\du}}
\pgfpathlineto{\pgfpoint{3.073335\du}{7.502787\du}}
\pgfpathlineto{\pgfpoint{3.073335\du}{7.561900\du}}
\pgfpathlineto{\pgfpoint{1.431000\du}{7.561900\du}}
\pgfpathlineto{\pgfpoint{1.431000\du}{7.502787\du}}
\pgfusepath{stroke}
\pgfsetbuttcap
\pgfsetmiterjoin
\pgfsetdash{}{0pt}
\definecolor{dialinecolor}{rgb}{0.000000, 0.705882, 1.000000}
\pgfsetfillcolor{dialinecolor}
\pgfpathmoveto{\pgfpoint{3.406220\du}{7.466662\du}}
\pgfpathlineto{\pgfpoint{3.525342\du}{7.288426\du}}
\pgfpathlineto{\pgfpoint{3.740001\du}{7.288426\du}}
\pgfpathlineto{\pgfpoint{3.621178\du}{7.466662\du}}
\pgfpathlineto{\pgfpoint{3.406220\du}{7.466662\du}}
\pgfusepath{fill}
\pgfsetbuttcap
\pgfsetmiterjoin
\pgfsetdash{}{0pt}
\definecolor{dialinecolor}{rgb}{0.666667, 0.901961, 1.000000}
\pgfsetstrokecolor{dialinecolor}
\pgfpathmoveto{\pgfpoint{3.406220\du}{7.466662\du}}
\pgfpathlineto{\pgfpoint{3.525342\du}{7.288426\du}}
\pgfpathlineto{\pgfpoint{3.740001\du}{7.288426\du}}
\pgfpathlineto{\pgfpoint{3.621178\du}{7.466662\du}}
\pgfpathlineto{\pgfpoint{3.406220\du}{7.466662\du}}
\pgfusepath{stroke}
\pgfsetbuttcap
\pgfsetmiterjoin
\pgfsetdash{}{0pt}
\definecolor{dialinecolor}{rgb}{0.000000, 0.352941, 0.501961}
\pgfsetfillcolor{dialinecolor}
\pgfpathmoveto{\pgfpoint{3.621178\du}{7.526372\du}}
\pgfpathlineto{\pgfpoint{3.740001\du}{7.395308\du}}
\pgfpathlineto{\pgfpoint{3.740001\du}{7.288426\du}}
\pgfpathlineto{\pgfpoint{3.621178\du}{7.466662\du}}
\pgfpathlineto{\pgfpoint{3.621178\du}{7.526372\du}}
\pgfusepath{fill}
\pgfsetbuttcap
\pgfsetmiterjoin
\pgfsetdash{}{0pt}
\definecolor{dialinecolor}{rgb}{0.666667, 0.901961, 1.000000}
\pgfsetstrokecolor{dialinecolor}
\pgfpathmoveto{\pgfpoint{3.621178\du}{7.526372\du}}
\pgfpathlineto{\pgfpoint{3.740001\du}{7.395308\du}}
\pgfpathlineto{\pgfpoint{3.740001\du}{7.288426\du}}
\pgfpathlineto{\pgfpoint{3.621178\du}{7.466662\du}}
\pgfpathlineto{\pgfpoint{3.621178\du}{7.526372\du}}
\pgfusepath{stroke}
\pgfsetbuttcap
\pgfsetmiterjoin
\pgfsetdash{}{0pt}
\definecolor{dialinecolor}{rgb}{0.000000, 0.588235, 0.831373}
\pgfsetfillcolor{dialinecolor}
\pgfpathmoveto{\pgfpoint{3.406220\du}{7.466662\du}}
\pgfpathlineto{\pgfpoint{3.621178\du}{7.466662\du}}
\pgfpathlineto{\pgfpoint{3.621178\du}{7.526372\du}}
\pgfpathlineto{\pgfpoint{3.406220\du}{7.526372\du}}
\pgfpathlineto{\pgfpoint{3.406220\du}{7.466662\du}}
\pgfusepath{fill}
\pgfsetbuttcap
\pgfsetmiterjoin
\pgfsetdash{}{0pt}
\definecolor{dialinecolor}{rgb}{0.666667, 0.901961, 1.000000}
\pgfsetstrokecolor{dialinecolor}
\pgfpathmoveto{\pgfpoint{3.406220\du}{7.466662\du}}
\pgfpathlineto{\pgfpoint{3.621178\du}{7.466662\du}}
\pgfpathlineto{\pgfpoint{3.621178\du}{7.526372\du}}
\pgfpathlineto{\pgfpoint{3.406220\du}{7.526372\du}}
\pgfpathlineto{\pgfpoint{3.406220\du}{7.466662\du}}
\pgfusepath{stroke}
\pgfsetbuttcap
\pgfsetmiterjoin
\pgfsetdash{}{0pt}
\definecolor{dialinecolor}{rgb}{0.811765, 0.054902, 0.188235}
\pgfsetstrokecolor{dialinecolor}
\pgfpathmoveto{\pgfpoint{2.192308\du}{6.097801\du}}
\pgfpathlineto{\pgfpoint{2.526089\du}{6.097801\du}}
\pgfpathlineto{\pgfpoint{2.442494\du}{6.181097\du}}
\pgfpathlineto{\pgfpoint{2.728209\du}{6.181097\du}}
\pgfusepath{stroke}
\pgfsetbuttcap
\pgfsetmiterjoin
\pgfsetdash{}{0pt}
\definecolor{dialinecolor}{rgb}{0.811765, 0.054902, 0.188235}
\pgfsetstrokecolor{dialinecolor}
\pgfpathmoveto{\pgfpoint{2.156481\du}{6.097801\du}}
\pgfpathlineto{\pgfpoint{2.395024\du}{6.335747\du}}
\pgfpathlineto{\pgfpoint{2.275604\du}{6.335747\du}}
\pgfpathlineto{\pgfpoint{2.478022\du}{6.538464\du}}
\pgfusepath{stroke}
\pgfsetbuttcap
\pgfsetmiterjoin
\pgfsetdash{}{0pt}
\definecolor{dialinecolor}{rgb}{0.811765, 0.054902, 0.188235}
\pgfsetstrokecolor{dialinecolor}
\pgfpathmoveto{\pgfpoint{2.834792\du}{6.097801\du}}
\pgfpathlineto{\pgfpoint{2.597443\du}{6.335747\du}}
\pgfpathlineto{\pgfpoint{2.715968\du}{6.335747\du}}
\pgfpathlineto{\pgfpoint{2.513848\du}{6.538464\du}}
\pgfusepath{stroke}
\pgfsetbuttcap
\pgfsetmiterjoin
\pgfsetdash{}{0pt}
\definecolor{dialinecolor}{rgb}{0.000000, 0.470588, 0.666667}
\pgfsetfillcolor{dialinecolor}
\pgfpathmoveto{\pgfpoint{2.323073\du}{6.186769\du}}
\pgfpathlineto{\pgfpoint{2.322476\du}{6.183485\du}}
\pgfpathlineto{\pgfpoint{2.321581\du}{6.180201\du}}
\pgfpathlineto{\pgfpoint{2.320685\du}{6.176917\du}}
\pgfpathlineto{\pgfpoint{2.318894\du}{6.173633\du}}
\pgfpathlineto{\pgfpoint{2.316505\du}{6.170648\du}}
\pgfpathlineto{\pgfpoint{2.313818\du}{6.167364\du}}
\pgfpathlineto{\pgfpoint{2.310833\du}{6.164079\du}}
\pgfpathlineto{\pgfpoint{2.307250\du}{6.161094\du}}
\pgfpathlineto{\pgfpoint{2.303070\du}{6.158108\du}}
\pgfpathlineto{\pgfpoint{2.298294\du}{6.155123\du}}
\pgfpathlineto{\pgfpoint{2.293218\du}{6.152436\du}}
\pgfpathlineto{\pgfpoint{2.287546\du}{6.149450\du}}
\pgfpathlineto{\pgfpoint{2.281873\du}{6.147062\du}}
\pgfpathlineto{\pgfpoint{2.275604\du}{6.144375\du}}
\pgfpathlineto{\pgfpoint{2.269036\du}{6.141987\du}}
\pgfpathlineto{\pgfpoint{2.261870\du}{6.139598\du}}
\pgfpathlineto{\pgfpoint{2.254406\du}{6.137508\du}}
\pgfpathlineto{\pgfpoint{2.246943\du}{6.135120\du}}
\pgfpathlineto{\pgfpoint{2.238583\du}{6.133627\du}}
\pgfpathlineto{\pgfpoint{2.230224\du}{6.131537\du}}
\pgfpathlineto{\pgfpoint{2.222163\du}{6.129746\du}}
\pgfpathlineto{\pgfpoint{2.212908\du}{6.128552\du}}
\pgfpathlineto{\pgfpoint{2.204250\du}{6.127059\du}}
\pgfpathlineto{\pgfpoint{2.194696\du}{6.125865\du}}
\pgfpathlineto{\pgfpoint{2.185441\du}{6.124671\du}}
\pgfpathlineto{\pgfpoint{2.176186\du}{6.123775\du}}
\pgfpathlineto{\pgfpoint{2.166334\du}{6.122879\du}}
\pgfpathlineto{\pgfpoint{2.156481\du}{6.122282\du}}
\pgfpathlineto{\pgfpoint{2.146629\du}{6.121984\du}}
\pgfpathlineto{\pgfpoint{2.136478\du}{6.121685\du}}
\pgfpathlineto{\pgfpoint{2.126626\du}{6.121685\du}}
\pgfpathlineto{\pgfpoint{2.126626\du}{6.121685\du}}
\pgfpathlineto{\pgfpoint{2.116774\du}{6.121685\du}}
\pgfpathlineto{\pgfpoint{2.106922\du}{6.121984\du}}
\pgfpathlineto{\pgfpoint{2.096771\du}{6.122282\du}}
\pgfpathlineto{\pgfpoint{2.086919\du}{6.122879\du}}
\pgfpathlineto{\pgfpoint{2.077365\du}{6.123775\du}}
\pgfpathlineto{\pgfpoint{2.067811\du}{6.124671\du}}
\pgfpathlineto{\pgfpoint{2.058556\du}{6.125865\du}}
\pgfpathlineto{\pgfpoint{2.049301\du}{6.127059\du}}
\pgfpathlineto{\pgfpoint{2.040345\du}{6.128552\du}}
\pgfpathlineto{\pgfpoint{2.031388\du}{6.129746\du}}
\pgfpathlineto{\pgfpoint{2.023029\du}{6.131537\du}}
\pgfpathlineto{\pgfpoint{2.014669\du}{6.133627\du}}
\pgfpathlineto{\pgfpoint{2.006310\du}{6.135120\du}}
\pgfpathlineto{\pgfpoint{1.999144\du}{6.137508\du}}
\pgfpathlineto{\pgfpoint{1.991681\du}{6.139598\du}}
\pgfpathlineto{\pgfpoint{1.984515\du}{6.141987\du}}
\pgfpathlineto{\pgfpoint{1.977649\du}{6.144375\du}}
\pgfpathlineto{\pgfpoint{1.971678\du}{6.147062\du}}
\pgfpathlineto{\pgfpoint{1.965707\du}{6.149450\du}}
\pgfpathlineto{\pgfpoint{1.960034\du}{6.152436\du}}
\pgfpathlineto{\pgfpoint{1.954959\du}{6.155123\du}}
\pgfpathlineto{\pgfpoint{1.950182\du}{6.158108\du}}
\pgfpathlineto{\pgfpoint{1.946002\du}{6.161094\du}}
\pgfpathlineto{\pgfpoint{1.942718\du}{6.164079\du}}
\pgfpathlineto{\pgfpoint{1.939434\du}{6.167364\du}}
\pgfpathlineto{\pgfpoint{1.936449\du}{6.170648\du}}
\pgfpathlineto{\pgfpoint{1.934359\du}{6.173633\du}}
\pgfpathlineto{\pgfpoint{1.932866\du}{6.176917\du}}
\pgfpathlineto{\pgfpoint{1.931075\du}{6.180201\du}}
\pgfpathlineto{\pgfpoint{1.930776\du}{6.183485\du}}
\pgfpathlineto{\pgfpoint{1.930478\du}{6.186769\du}}
\pgfpathlineto{\pgfpoint{1.930478\du}{6.186769\du}}
\pgfpathlineto{\pgfpoint{1.930776\du}{6.190352\du}}
\pgfpathlineto{\pgfpoint{1.931075\du}{6.193338\du}}
\pgfpathlineto{\pgfpoint{1.932866\du}{6.196622\du}}
\pgfpathlineto{\pgfpoint{1.934359\du}{6.199906\du}}
\pgfpathlineto{\pgfpoint{1.936449\du}{6.203190\du}}
\pgfpathlineto{\pgfpoint{1.939434\du}{6.206474\du}}
\pgfpathlineto{\pgfpoint{1.942718\du}{6.209459\du}}
\pgfpathlineto{\pgfpoint{1.946002\du}{6.212445\du}}
\pgfpathlineto{\pgfpoint{1.950182\du}{6.215430\du}}
\pgfpathlineto{\pgfpoint{1.954959\du}{6.218714\du}}
\pgfpathlineto{\pgfpoint{1.960034\du}{6.221401\du}}
\pgfpathlineto{\pgfpoint{1.965707\du}{6.224088\du}}
\pgfpathlineto{\pgfpoint{1.971678\du}{6.226775\du}}
\pgfpathlineto{\pgfpoint{1.977649\du}{6.229462\du}}
\pgfpathlineto{\pgfpoint{1.984515\du}{6.231851\du}}
\pgfpathlineto{\pgfpoint{1.991681\du}{6.234239\du}}
\pgfpathlineto{\pgfpoint{1.999144\du}{6.236030\du}}
\pgfpathlineto{\pgfpoint{2.006310\du}{6.238120\du}}
\pgfpathlineto{\pgfpoint{2.014669\du}{6.240210\du}}
\pgfpathlineto{\pgfpoint{2.023029\du}{6.242300\du}}
\pgfpathlineto{\pgfpoint{2.031388\du}{6.243793\du}}
\pgfpathlineto{\pgfpoint{2.040345\du}{6.245286\du}}
\pgfpathlineto{\pgfpoint{2.049301\du}{6.246778\du}}
\pgfpathlineto{\pgfpoint{2.058556\du}{6.247973\du}}
\pgfpathlineto{\pgfpoint{2.067811\du}{6.248868\du}}
\pgfpathlineto{\pgfpoint{2.077365\du}{6.250062\du}}
\pgfpathlineto{\pgfpoint{2.086919\du}{6.250660\du}}
\pgfpathlineto{\pgfpoint{2.096771\du}{6.251257\du}}
\pgfpathlineto{\pgfpoint{2.106922\du}{6.251854\du}}
\pgfpathlineto{\pgfpoint{2.116774\du}{6.252152\du}}
\pgfpathlineto{\pgfpoint{2.126626\du}{6.252152\du}}
\pgfpathlineto{\pgfpoint{2.126626\du}{6.252152\du}}
\pgfpathlineto{\pgfpoint{2.136478\du}{6.252152\du}}
\pgfpathlineto{\pgfpoint{2.146629\du}{6.251854\du}}
\pgfpathlineto{\pgfpoint{2.156481\du}{6.251257\du}}
\pgfpathlineto{\pgfpoint{2.166334\du}{6.250660\du}}
\pgfpathlineto{\pgfpoint{2.176186\du}{6.250062\du}}
\pgfpathlineto{\pgfpoint{2.185441\du}{6.248868\du}}
\pgfpathlineto{\pgfpoint{2.194696\du}{6.247973\du}}
\pgfpathlineto{\pgfpoint{2.204250\du}{6.246778\du}}
\pgfpathlineto{\pgfpoint{2.212908\du}{6.245286\du}}
\pgfpathlineto{\pgfpoint{2.222163\du}{6.243793\du}}
\pgfpathlineto{\pgfpoint{2.230224\du}{6.242300\du}}
\pgfpathlineto{\pgfpoint{2.238583\du}{6.240210\du}}
\pgfpathlineto{\pgfpoint{2.246943\du}{6.238120\du}}
\pgfpathlineto{\pgfpoint{2.254406\du}{6.236030\du}}
\pgfpathlineto{\pgfpoint{2.261870\du}{6.234239\du}}
\pgfpathlineto{\pgfpoint{2.269036\du}{6.231851\du}}
\pgfpathlineto{\pgfpoint{2.275604\du}{6.229462\du}}
\pgfpathlineto{\pgfpoint{2.281873\du}{6.226775\du}}
\pgfpathlineto{\pgfpoint{2.287546\du}{6.224088\du}}
\pgfpathlineto{\pgfpoint{2.293218\du}{6.221401\du}}
\pgfpathlineto{\pgfpoint{2.298294\du}{6.218714\du}}
\pgfpathlineto{\pgfpoint{2.303070\du}{6.215430\du}}
\pgfpathlineto{\pgfpoint{2.307250\du}{6.212445\du}}
\pgfpathlineto{\pgfpoint{2.310833\du}{6.209459\du}}
\pgfpathlineto{\pgfpoint{2.313818\du}{6.206474\du}}
\pgfpathlineto{\pgfpoint{2.316505\du}{6.203190\du}}
\pgfpathlineto{\pgfpoint{2.318894\du}{6.199906\du}}
\pgfpathlineto{\pgfpoint{2.320685\du}{6.196622\du}}
\pgfpathlineto{\pgfpoint{2.321581\du}{6.193338\du}}
\pgfpathlineto{\pgfpoint{2.322476\du}{6.190352\du}}
\pgfpathlineto{\pgfpoint{2.323073\du}{6.186769\du}}
\pgfusepath{fill}
\pgfsetbuttcap
\pgfsetmiterjoin
\pgfsetdash{}{0pt}
\definecolor{dialinecolor}{rgb}{0.666667, 0.901961, 1.000000}
\pgfsetstrokecolor{dialinecolor}
\pgfpathmoveto{\pgfpoint{2.310833\du}{6.181097\du}}
\pgfpathlineto{\pgfpoint{2.310833\du}{6.177514\du}}
\pgfpathlineto{\pgfpoint{2.309937\du}{6.175126\du}}
\pgfpathlineto{\pgfpoint{2.308743\du}{6.171842\du}}
\pgfpathlineto{\pgfpoint{2.307250\du}{6.168856\du}}
\pgfpathlineto{\pgfpoint{2.304862\du}{6.166169\du}}
\pgfpathlineto{\pgfpoint{2.302473\du}{6.162587\du}}
\pgfpathlineto{\pgfpoint{2.299189\du}{6.160198\du}}
\pgfpathlineto{\pgfpoint{2.295905\du}{6.157511\du}}
\pgfpathlineto{\pgfpoint{2.291725\du}{6.154824\du}}
\pgfpathlineto{\pgfpoint{2.287247\du}{6.152436\du}}
\pgfpathlineto{\pgfpoint{2.282470\du}{6.149450\du}}
\pgfpathlineto{\pgfpoint{2.276798\du}{6.147361\du}}
\pgfpathlineto{\pgfpoint{2.271424\du}{6.144674\du}}
\pgfpathlineto{\pgfpoint{2.265154\du}{6.141987\du}}
\pgfpathlineto{\pgfpoint{2.259183\du}{6.140195\du}}
\pgfpathlineto{\pgfpoint{2.251720\du}{6.137807\du}}
\pgfpathlineto{\pgfpoint{2.244853\du}{6.136314\du}}
\pgfpathlineto{\pgfpoint{2.237688\du}{6.134224\du}}
\pgfpathlineto{\pgfpoint{2.229925\du}{6.132433\du}}
\pgfpathlineto{\pgfpoint{2.221267\du}{6.130642\du}}
\pgfpathlineto{\pgfpoint{2.212908\du}{6.129447\du}}
\pgfpathlineto{\pgfpoint{2.204548\du}{6.127059\du}}
\pgfpathlineto{\pgfpoint{2.195592\du}{6.126760\du}}
\pgfpathlineto{\pgfpoint{2.186934\du}{6.125268\du}}
\pgfpathlineto{\pgfpoint{2.177977\du}{6.124372\du}}
\pgfpathlineto{\pgfpoint{2.168722\du}{6.123476\du}}
\pgfpathlineto{\pgfpoint{2.159168\du}{6.122879\du}}
\pgfpathlineto{\pgfpoint{2.149615\du}{6.122581\du}}
\pgfpathlineto{\pgfpoint{2.140658\du}{6.121984\du}}
\pgfpathlineto{\pgfpoint{2.130507\du}{6.121685\du}}
\pgfpathlineto{\pgfpoint{2.120954\du}{6.121685\du}}
\pgfpathlineto{\pgfpoint{2.120954\du}{6.121685\du}}
\pgfpathlineto{\pgfpoint{2.111699\du}{6.121685\du}}
\pgfpathlineto{\pgfpoint{2.101846\du}{6.121984\du}}
\pgfpathlineto{\pgfpoint{2.092293\du}{6.122581\du}}
\pgfpathlineto{\pgfpoint{2.082739\du}{6.122879\du}}
\pgfpathlineto{\pgfpoint{2.073484\du}{6.123476\du}}
\pgfpathlineto{\pgfpoint{2.063632\du}{6.124372\du}}
\pgfpathlineto{\pgfpoint{2.054974\du}{6.125268\du}}
\pgfpathlineto{\pgfpoint{2.046017\du}{6.126760\du}}
\pgfpathlineto{\pgfpoint{2.037658\du}{6.127059\du}}
\pgfpathlineto{\pgfpoint{2.028701\du}{6.129447\du}}
\pgfpathlineto{\pgfpoint{2.020342\du}{6.130642\du}}
\pgfpathlineto{\pgfpoint{2.012281\du}{6.132433\du}}
\pgfpathlineto{\pgfpoint{2.004518\du}{6.134224\du}}
\pgfpathlineto{\pgfpoint{1.997055\du}{6.136314\du}}
\pgfpathlineto{\pgfpoint{1.990188\du}{6.137807\du}}
\pgfpathlineto{\pgfpoint{1.982724\du}{6.140195\du}}
\pgfpathlineto{\pgfpoint{1.976753\du}{6.141987\du}}
\pgfpathlineto{\pgfpoint{1.970484\du}{6.144674\du}}
\pgfpathlineto{\pgfpoint{1.964811\du}{6.147361\du}}
\pgfpathlineto{\pgfpoint{1.959437\du}{6.149450\du}}
\pgfpathlineto{\pgfpoint{1.954660\du}{6.152436\du}}
\pgfpathlineto{\pgfpoint{1.950182\du}{6.154824\du}}
\pgfpathlineto{\pgfpoint{1.946002\du}{6.157511\du}}
\pgfpathlineto{\pgfpoint{1.942718\du}{6.160198\du}}
\pgfpathlineto{\pgfpoint{1.939136\du}{6.162587\du}}
\pgfpathlineto{\pgfpoint{1.937344\du}{6.166169\du}}
\pgfpathlineto{\pgfpoint{1.934956\du}{6.168856\du}}
\pgfpathlineto{\pgfpoint{1.933165\du}{6.171842\du}}
\pgfpathlineto{\pgfpoint{1.931672\du}{6.175126\du}}
\pgfpathlineto{\pgfpoint{1.931075\du}{6.177514\du}}
\pgfpathlineto{\pgfpoint{1.930776\du}{6.181097\du}}
\pgfpathlineto{\pgfpoint{1.930776\du}{6.181097\du}}
\pgfpathlineto{\pgfpoint{1.931075\du}{6.184082\du}}
\pgfpathlineto{\pgfpoint{1.931672\du}{6.186769\du}}
\pgfpathlineto{\pgfpoint{1.933165\du}{6.189755\du}}
\pgfpathlineto{\pgfpoint{1.934956\du}{6.193338\du}}
\pgfpathlineto{\pgfpoint{1.937344\du}{6.196024\du}}
\pgfpathlineto{\pgfpoint{1.939136\du}{6.199010\du}}
\pgfpathlineto{\pgfpoint{1.942718\du}{6.201697\du}}
\pgfpathlineto{\pgfpoint{1.946002\du}{6.204384\du}}
\pgfpathlineto{\pgfpoint{1.950182\du}{6.207071\du}}
\pgfpathlineto{\pgfpoint{1.954660\du}{6.209758\du}}
\pgfpathlineto{\pgfpoint{1.959437\du}{6.212445\du}}
\pgfpathlineto{\pgfpoint{1.964811\du}{6.214833\du}}
\pgfpathlineto{\pgfpoint{1.970484\du}{6.217520\du}}
\pgfpathlineto{\pgfpoint{1.976753\du}{6.219610\du}}
\pgfpathlineto{\pgfpoint{1.982724\du}{6.221999\du}}
\pgfpathlineto{\pgfpoint{1.990188\du}{6.224387\du}}
\pgfpathlineto{\pgfpoint{1.997055\du}{6.226178\du}}
\pgfpathlineto{\pgfpoint{2.004518\du}{6.227671\du}}
\pgfpathlineto{\pgfpoint{2.012281\du}{6.229761\du}}
\pgfpathlineto{\pgfpoint{2.020342\du}{6.231552\du}}
\pgfpathlineto{\pgfpoint{2.028701\du}{6.232746\du}}
\pgfpathlineto{\pgfpoint{2.037658\du}{6.234538\du}}
\pgfpathlineto{\pgfpoint{2.046017\du}{6.235433\du}}
\pgfpathlineto{\pgfpoint{2.054974\du}{6.236926\du}}
\pgfpathlineto{\pgfpoint{2.063632\du}{6.237822\du}}
\pgfpathlineto{\pgfpoint{2.073484\du}{6.238120\du}}
\pgfpathlineto{\pgfpoint{2.082739\du}{6.239315\du}}
\pgfpathlineto{\pgfpoint{2.092293\du}{6.239613\du}}
\pgfpathlineto{\pgfpoint{2.101846\du}{6.239912\du}}
\pgfpathlineto{\pgfpoint{2.111699\du}{6.240509\du}}
\pgfpathlineto{\pgfpoint{2.120954\du}{6.240509\du}}
\pgfpathlineto{\pgfpoint{2.120954\du}{6.240509\du}}
\pgfpathlineto{\pgfpoint{2.130507\du}{6.240509\du}}
\pgfpathlineto{\pgfpoint{2.140658\du}{6.239912\du}}
\pgfpathlineto{\pgfpoint{2.149615\du}{6.239613\du}}
\pgfpathlineto{\pgfpoint{2.159168\du}{6.239315\du}}
\pgfpathlineto{\pgfpoint{2.168722\du}{6.238120\du}}
\pgfpathlineto{\pgfpoint{2.177977\du}{6.237822\du}}
\pgfpathlineto{\pgfpoint{2.186934\du}{6.236926\du}}
\pgfpathlineto{\pgfpoint{2.195592\du}{6.235433\du}}
\pgfpathlineto{\pgfpoint{2.204548\du}{6.234538\du}}
\pgfpathlineto{\pgfpoint{2.212908\du}{6.232746\du}}
\pgfpathlineto{\pgfpoint{2.221267\du}{6.231552\du}}
\pgfpathlineto{\pgfpoint{2.229925\du}{6.229761\du}}
\pgfpathlineto{\pgfpoint{2.237688\du}{6.227671\du}}
\pgfpathlineto{\pgfpoint{2.244853\du}{6.226178\du}}
\pgfpathlineto{\pgfpoint{2.251720\du}{6.224387\du}}
\pgfpathlineto{\pgfpoint{2.259183\du}{6.221999\du}}
\pgfpathlineto{\pgfpoint{2.265154\du}{6.219610\du}}
\pgfpathlineto{\pgfpoint{2.271424\du}{6.217520\du}}
\pgfpathlineto{\pgfpoint{2.276798\du}{6.214833\du}}
\pgfpathlineto{\pgfpoint{2.282470\du}{6.212445\du}}
\pgfpathlineto{\pgfpoint{2.287247\du}{6.209758\du}}
\pgfpathlineto{\pgfpoint{2.291725\du}{6.207071\du}}
\pgfpathlineto{\pgfpoint{2.295905\du}{6.204384\du}}
\pgfpathlineto{\pgfpoint{2.299189\du}{6.201697\du}}
\pgfpathlineto{\pgfpoint{2.302473\du}{6.199010\du}}
\pgfpathlineto{\pgfpoint{2.304862\du}{6.196024\du}}
\pgfpathlineto{\pgfpoint{2.307250\du}{6.193338\du}}
\pgfpathlineto{\pgfpoint{2.308743\du}{6.189755\du}}
\pgfpathlineto{\pgfpoint{2.309937\du}{6.186769\du}}
\pgfpathlineto{\pgfpoint{2.310833\du}{6.184082\du}}
\pgfpathlineto{\pgfpoint{2.310833\du}{6.181097\du}}
\pgfusepath{stroke}
\pgfsetbuttcap
\pgfsetmiterjoin
\pgfsetdash{}{0pt}
\definecolor{dialinecolor}{rgb}{0.000000, 0.470588, 0.666667}
\pgfsetfillcolor{dialinecolor}
\pgfpathmoveto{\pgfpoint{1.930776\du}{6.097801\du}}
\pgfpathlineto{\pgfpoint{1.930776\du}{6.192442\du}}
\pgfpathlineto{\pgfpoint{2.323372\du}{6.192442\du}}
\pgfpathlineto{\pgfpoint{2.323372\du}{6.097801\du}}
\pgfpathlineto{\pgfpoint{1.930776\du}{6.097801\du}}
\pgfusepath{fill}
\pgfsetbuttcap
\pgfsetmiterjoin
\pgfsetdash{}{0pt}
\definecolor{dialinecolor}{rgb}{0.000000, 0.705882, 1.000000}
\pgfsetfillcolor{dialinecolor}
\pgfpathmoveto{\pgfpoint{2.323073\du}{6.091233\du}}
\pgfpathlineto{\pgfpoint{2.322476\du}{6.087949\du}}
\pgfpathlineto{\pgfpoint{2.321581\du}{6.084366\du}}
\pgfpathlineto{\pgfpoint{2.320685\du}{6.081381\du}}
\pgfpathlineto{\pgfpoint{2.318894\du}{6.078096\du}}
\pgfpathlineto{\pgfpoint{2.316505\du}{6.074812\du}}
\pgfpathlineto{\pgfpoint{2.313818\du}{6.071528\du}}
\pgfpathlineto{\pgfpoint{2.310833\du}{6.068543\du}}
\pgfpathlineto{\pgfpoint{2.307250\du}{6.065259\du}}
\pgfpathlineto{\pgfpoint{2.303070\du}{6.062572\du}}
\pgfpathlineto{\pgfpoint{2.298294\du}{6.059586\du}}
\pgfpathlineto{\pgfpoint{2.293218\du}{6.056601\du}}
\pgfpathlineto{\pgfpoint{2.287546\du}{6.053914\du}}
\pgfpathlineto{\pgfpoint{2.281873\du}{6.050630\du}}
\pgfpathlineto{\pgfpoint{2.275604\du}{6.048838\du}}
\pgfpathlineto{\pgfpoint{2.269036\du}{6.046151\du}}
\pgfpathlineto{\pgfpoint{2.261870\du}{6.043763\du}}
\pgfpathlineto{\pgfpoint{2.254406\du}{6.041673\du}}
\pgfpathlineto{\pgfpoint{2.246943\du}{6.039583\du}}
\pgfpathlineto{\pgfpoint{2.238583\du}{6.037493\du}}
\pgfpathlineto{\pgfpoint{2.230224\du}{6.035404\du}}
\pgfpathlineto{\pgfpoint{2.222163\du}{6.033911\du}}
\pgfpathlineto{\pgfpoint{2.212908\du}{6.032418\du}}
\pgfpathlineto{\pgfpoint{2.204250\du}{6.031224\du}}
\pgfpathlineto{\pgfpoint{2.194696\du}{6.029731\du}}
\pgfpathlineto{\pgfpoint{2.185441\du}{6.028835\du}}
\pgfpathlineto{\pgfpoint{2.176186\du}{6.027641\du}}
\pgfpathlineto{\pgfpoint{2.166334\du}{6.027044\du}}
\pgfpathlineto{\pgfpoint{2.156481\du}{6.026447\du}}
\pgfpathlineto{\pgfpoint{2.146629\du}{6.026148\du}}
\pgfpathlineto{\pgfpoint{2.136478\du}{6.025850\du}}
\pgfpathlineto{\pgfpoint{2.126626\du}{6.025850\du}}
\pgfpathlineto{\pgfpoint{2.126626\du}{6.025850\du}}
\pgfpathlineto{\pgfpoint{2.116774\du}{6.025850\du}}
\pgfpathlineto{\pgfpoint{2.106922\du}{6.026148\du}}
\pgfpathlineto{\pgfpoint{2.096771\du}{6.026447\du}}
\pgfpathlineto{\pgfpoint{2.086919\du}{6.027044\du}}
\pgfpathlineto{\pgfpoint{2.077365\du}{6.027641\du}}
\pgfpathlineto{\pgfpoint{2.067811\du}{6.028835\du}}
\pgfpathlineto{\pgfpoint{2.058556\du}{6.029731\du}}
\pgfpathlineto{\pgfpoint{2.049301\du}{6.031224\du}}
\pgfpathlineto{\pgfpoint{2.040345\du}{6.032418\du}}
\pgfpathlineto{\pgfpoint{2.031388\du}{6.033911\du}}
\pgfpathlineto{\pgfpoint{2.023029\du}{6.035404\du}}
\pgfpathlineto{\pgfpoint{2.014669\du}{6.037493\du}}
\pgfpathlineto{\pgfpoint{2.006310\du}{6.039583\du}}
\pgfpathlineto{\pgfpoint{1.999144\du}{6.041673\du}}
\pgfpathlineto{\pgfpoint{1.991681\du}{6.043763\du}}
\pgfpathlineto{\pgfpoint{1.984515\du}{6.046151\du}}
\pgfpathlineto{\pgfpoint{1.977649\du}{6.048838\du}}
\pgfpathlineto{\pgfpoint{1.971678\du}{6.050630\du}}
\pgfpathlineto{\pgfpoint{1.965707\du}{6.053914\du}}
\pgfpathlineto{\pgfpoint{1.960034\du}{6.056601\du}}
\pgfpathlineto{\pgfpoint{1.954959\du}{6.059586\du}}
\pgfpathlineto{\pgfpoint{1.950182\du}{6.062572\du}}
\pgfpathlineto{\pgfpoint{1.946002\du}{6.065259\du}}
\pgfpathlineto{\pgfpoint{1.942718\du}{6.068543\du}}
\pgfpathlineto{\pgfpoint{1.939434\du}{6.071528\du}}
\pgfpathlineto{\pgfpoint{1.936449\du}{6.074812\du}}
\pgfpathlineto{\pgfpoint{1.934359\du}{6.078096\du}}
\pgfpathlineto{\pgfpoint{1.932866\du}{6.081381\du}}
\pgfpathlineto{\pgfpoint{1.931075\du}{6.084366\du}}
\pgfpathlineto{\pgfpoint{1.930776\du}{6.087949\du}}
\pgfpathlineto{\pgfpoint{1.930478\du}{6.091233\du}}
\pgfpathlineto{\pgfpoint{1.930478\du}{6.091233\du}}
\pgfpathlineto{\pgfpoint{1.930776\du}{6.094517\du}}
\pgfpathlineto{\pgfpoint{1.931075\du}{6.098099\du}}
\pgfpathlineto{\pgfpoint{1.932866\du}{6.101085\du}}
\pgfpathlineto{\pgfpoint{1.934359\du}{6.104369\du}}
\pgfpathlineto{\pgfpoint{1.936449\du}{6.107952\du}}
\pgfpathlineto{\pgfpoint{1.939434\du}{6.110937\du}}
\pgfpathlineto{\pgfpoint{1.942718\du}{6.113923\du}}
\pgfpathlineto{\pgfpoint{1.946002\du}{6.117207\du}}
\pgfpathlineto{\pgfpoint{1.950182\du}{6.119595\du}}
\pgfpathlineto{\pgfpoint{1.954959\du}{6.122879\du}}
\pgfpathlineto{\pgfpoint{1.960034\du}{6.126163\du}}
\pgfpathlineto{\pgfpoint{1.965707\du}{6.128850\du}}
\pgfpathlineto{\pgfpoint{1.971678\du}{6.131537\du}}
\pgfpathlineto{\pgfpoint{1.977649\du}{6.133627\du}}
\pgfpathlineto{\pgfpoint{1.984515\du}{6.136613\du}}
\pgfpathlineto{\pgfpoint{1.991681\du}{6.139001\du}}
\pgfpathlineto{\pgfpoint{1.999144\du}{6.140792\du}}
\pgfpathlineto{\pgfpoint{2.006310\du}{6.142882\du}}
\pgfpathlineto{\pgfpoint{2.014669\du}{6.144972\du}}
\pgfpathlineto{\pgfpoint{2.023029\du}{6.147062\du}}
\pgfpathlineto{\pgfpoint{2.031388\du}{6.148555\du}}
\pgfpathlineto{\pgfpoint{2.040345\du}{6.150047\du}}
\pgfpathlineto{\pgfpoint{2.049301\du}{6.151242\du}}
\pgfpathlineto{\pgfpoint{2.058556\du}{6.152734\du}}
\pgfpathlineto{\pgfpoint{2.067811\du}{6.153630\du}}
\pgfpathlineto{\pgfpoint{2.077365\du}{6.154824\du}}
\pgfpathlineto{\pgfpoint{2.086919\du}{6.155123\du}}
\pgfpathlineto{\pgfpoint{2.096771\du}{6.156019\du}}
\pgfpathlineto{\pgfpoint{2.106922\du}{6.156019\du}}
\pgfpathlineto{\pgfpoint{2.116774\du}{6.156914\du}}
\pgfpathlineto{\pgfpoint{2.126626\du}{6.156914\du}}
\pgfpathlineto{\pgfpoint{2.126626\du}{6.156914\du}}
\pgfpathlineto{\pgfpoint{2.136478\du}{6.156914\du}}
\pgfpathlineto{\pgfpoint{2.146629\du}{6.156019\du}}
\pgfpathlineto{\pgfpoint{2.156481\du}{6.156019\du}}
\pgfpathlineto{\pgfpoint{2.166334\du}{6.155123\du}}
\pgfpathlineto{\pgfpoint{2.176186\du}{6.154824\du}}
\pgfpathlineto{\pgfpoint{2.185441\du}{6.153630\du}}
\pgfpathlineto{\pgfpoint{2.194696\du}{6.152734\du}}
\pgfpathlineto{\pgfpoint{2.204250\du}{6.151242\du}}
\pgfpathlineto{\pgfpoint{2.212908\du}{6.150047\du}}
\pgfpathlineto{\pgfpoint{2.222163\du}{6.148555\du}}
\pgfpathlineto{\pgfpoint{2.230224\du}{6.147062\du}}
\pgfpathlineto{\pgfpoint{2.238583\du}{6.144972\du}}
\pgfpathlineto{\pgfpoint{2.246943\du}{6.142882\du}}
\pgfpathlineto{\pgfpoint{2.254406\du}{6.140792\du}}
\pgfpathlineto{\pgfpoint{2.261870\du}{6.139001\du}}
\pgfpathlineto{\pgfpoint{2.269036\du}{6.136613\du}}
\pgfpathlineto{\pgfpoint{2.275604\du}{6.133627\du}}
\pgfpathlineto{\pgfpoint{2.281873\du}{6.131537\du}}
\pgfpathlineto{\pgfpoint{2.287546\du}{6.128850\du}}
\pgfpathlineto{\pgfpoint{2.293218\du}{6.126163\du}}
\pgfpathlineto{\pgfpoint{2.298294\du}{6.122879\du}}
\pgfpathlineto{\pgfpoint{2.303070\du}{6.119595\du}}
\pgfpathlineto{\pgfpoint{2.307250\du}{6.117207\du}}
\pgfpathlineto{\pgfpoint{2.310833\du}{6.113923\du}}
\pgfpathlineto{\pgfpoint{2.313818\du}{6.110937\du}}
\pgfpathlineto{\pgfpoint{2.316505\du}{6.107952\du}}
\pgfpathlineto{\pgfpoint{2.318894\du}{6.104369\du}}
\pgfpathlineto{\pgfpoint{2.320685\du}{6.101085\du}}
\pgfpathlineto{\pgfpoint{2.321581\du}{6.098099\du}}
\pgfpathlineto{\pgfpoint{2.322476\du}{6.094517\du}}
\pgfpathlineto{\pgfpoint{2.323073\du}{6.091233\du}}
\pgfusepath{fill}
\pgfsetbuttcap
\pgfsetmiterjoin
\pgfsetdash{}{0pt}
\definecolor{dialinecolor}{rgb}{0.666667, 0.901961, 1.000000}
\pgfsetstrokecolor{dialinecolor}
\pgfpathmoveto{\pgfpoint{2.310833\du}{6.085560\du}}
\pgfpathlineto{\pgfpoint{2.310833\du}{6.082276\du}}
\pgfpathlineto{\pgfpoint{2.309937\du}{6.079888\du}}
\pgfpathlineto{\pgfpoint{2.308743\du}{6.076604\du}}
\pgfpathlineto{\pgfpoint{2.307250\du}{6.073021\du}}
\pgfpathlineto{\pgfpoint{2.304862\du}{6.070633\du}}
\pgfpathlineto{\pgfpoint{2.302473\du}{6.067349\du}}
\pgfpathlineto{\pgfpoint{2.299189\du}{6.064960\du}}
\pgfpathlineto{\pgfpoint{2.295905\du}{6.062273\du}}
\pgfpathlineto{\pgfpoint{2.291725\du}{6.059586\du}}
\pgfpathlineto{\pgfpoint{2.287247\du}{6.056899\du}}
\pgfpathlineto{\pgfpoint{2.282470\du}{6.054212\du}}
\pgfpathlineto{\pgfpoint{2.276798\du}{6.051525\du}}
\pgfpathlineto{\pgfpoint{2.271424\du}{6.049137\du}}
\pgfpathlineto{\pgfpoint{2.265154\du}{6.047047\du}}
\pgfpathlineto{\pgfpoint{2.259183\du}{6.044659\du}}
\pgfpathlineto{\pgfpoint{2.251720\du}{6.042270\du}}
\pgfpathlineto{\pgfpoint{2.244853\du}{6.040479\du}}
\pgfpathlineto{\pgfpoint{2.237688\du}{6.038986\du}}
\pgfpathlineto{\pgfpoint{2.229925\du}{6.036598\du}}
\pgfpathlineto{\pgfpoint{2.221267\du}{6.035105\du}}
\pgfpathlineto{\pgfpoint{2.212908\du}{6.033911\du}}
\pgfpathlineto{\pgfpoint{2.204548\du}{6.032119\du}}
\pgfpathlineto{\pgfpoint{2.195592\du}{6.031224\du}}
\pgfpathlineto{\pgfpoint{2.186934\du}{6.029731\du}}
\pgfpathlineto{\pgfpoint{2.177977\du}{6.028835\du}}
\pgfpathlineto{\pgfpoint{2.168722\du}{6.028238\du}}
\pgfpathlineto{\pgfpoint{2.159168\du}{6.027343\du}}
\pgfpathlineto{\pgfpoint{2.149615\du}{6.027044\du}}
\pgfpathlineto{\pgfpoint{2.140658\du}{6.026745\du}}
\pgfpathlineto{\pgfpoint{2.130507\du}{6.026148\du}}
\pgfpathlineto{\pgfpoint{2.120954\du}{6.026148\du}}
\pgfpathlineto{\pgfpoint{2.120954\du}{6.026148\du}}
\pgfpathlineto{\pgfpoint{2.111699\du}{6.026148\du}}
\pgfpathlineto{\pgfpoint{2.101846\du}{6.026745\du}}
\pgfpathlineto{\pgfpoint{2.092293\du}{6.027044\du}}
\pgfpathlineto{\pgfpoint{2.082739\du}{6.027343\du}}
\pgfpathlineto{\pgfpoint{2.073484\du}{6.028238\du}}
\pgfpathlineto{\pgfpoint{2.063632\du}{6.028835\du}}
\pgfpathlineto{\pgfpoint{2.054974\du}{6.029731\du}}
\pgfpathlineto{\pgfpoint{2.046017\du}{6.031224\du}}
\pgfpathlineto{\pgfpoint{2.037658\du}{6.032119\du}}
\pgfpathlineto{\pgfpoint{2.028701\du}{6.033911\du}}
\pgfpathlineto{\pgfpoint{2.020342\du}{6.035105\du}}
\pgfpathlineto{\pgfpoint{2.012281\du}{6.036598\du}}
\pgfpathlineto{\pgfpoint{2.004518\du}{6.038986\du}}
\pgfpathlineto{\pgfpoint{1.997055\du}{6.040479\du}}
\pgfpathlineto{\pgfpoint{1.990188\du}{6.042270\du}}
\pgfpathlineto{\pgfpoint{1.982724\du}{6.044659\du}}
\pgfpathlineto{\pgfpoint{1.976753\du}{6.047047\du}}
\pgfpathlineto{\pgfpoint{1.970484\du}{6.049137\du}}
\pgfpathlineto{\pgfpoint{1.964811\du}{6.051525\du}}
\pgfpathlineto{\pgfpoint{1.959437\du}{6.054212\du}}
\pgfpathlineto{\pgfpoint{1.954660\du}{6.056899\du}}
\pgfpathlineto{\pgfpoint{1.950182\du}{6.059586\du}}
\pgfpathlineto{\pgfpoint{1.946002\du}{6.062273\du}}
\pgfpathlineto{\pgfpoint{1.942718\du}{6.064960\du}}
\pgfpathlineto{\pgfpoint{1.939136\du}{6.067349\du}}
\pgfpathlineto{\pgfpoint{1.937344\du}{6.070633\du}}
\pgfpathlineto{\pgfpoint{1.934956\du}{6.073021\du}}
\pgfpathlineto{\pgfpoint{1.933165\du}{6.076604\du}}
\pgfpathlineto{\pgfpoint{1.931672\du}{6.079888\du}}
\pgfpathlineto{\pgfpoint{1.931075\du}{6.082276\du}}
\pgfpathlineto{\pgfpoint{1.930776\du}{6.085560\du}}
\pgfpathlineto{\pgfpoint{1.930776\du}{6.085560\du}}
\pgfpathlineto{\pgfpoint{1.931075\du}{6.088844\du}}
\pgfpathlineto{\pgfpoint{1.931672\du}{6.091531\du}}
\pgfpathlineto{\pgfpoint{1.933165\du}{6.094517\du}}
\pgfpathlineto{\pgfpoint{1.934956\du}{6.097801\du}}
\pgfpathlineto{\pgfpoint{1.937344\du}{6.100488\du}}
\pgfpathlineto{\pgfpoint{1.939136\du}{6.103772\du}}
\pgfpathlineto{\pgfpoint{1.942718\du}{6.106459\du}}
\pgfpathlineto{\pgfpoint{1.946002\du}{6.109146\du}}
\pgfpathlineto{\pgfpoint{1.950182\du}{6.111833\du}}
\pgfpathlineto{\pgfpoint{1.954660\du}{6.114520\du}}
\pgfpathlineto{\pgfpoint{1.959437\du}{6.117207\du}}
\pgfpathlineto{\pgfpoint{1.964811\du}{6.119297\du}}
\pgfpathlineto{\pgfpoint{1.970484\du}{6.121984\du}}
\pgfpathlineto{\pgfpoint{1.976753\du}{6.124372\du}}
\pgfpathlineto{\pgfpoint{1.982724\du}{6.126760\du}}
\pgfpathlineto{\pgfpoint{1.990188\du}{6.128850\du}}
\pgfpathlineto{\pgfpoint{1.997055\du}{6.130642\du}}
\pgfpathlineto{\pgfpoint{2.004518\du}{6.132433\du}}
\pgfpathlineto{\pgfpoint{2.012281\du}{6.134224\du}}
\pgfpathlineto{\pgfpoint{2.020342\du}{6.136314\du}}
\pgfpathlineto{\pgfpoint{2.028701\du}{6.137508\du}}
\pgfpathlineto{\pgfpoint{2.037658\du}{6.139300\du}}
\pgfpathlineto{\pgfpoint{2.046017\du}{6.140195\du}}
\pgfpathlineto{\pgfpoint{2.054974\du}{6.141688\du}}
\pgfpathlineto{\pgfpoint{2.063632\du}{6.141987\du}}
\pgfpathlineto{\pgfpoint{2.073484\du}{6.143181\du}}
\pgfpathlineto{\pgfpoint{2.082739\du}{6.143479\du}}
\pgfpathlineto{\pgfpoint{2.092293\du}{6.144375\du}}
\pgfpathlineto{\pgfpoint{2.101846\du}{6.144674\du}}
\pgfpathlineto{\pgfpoint{2.111699\du}{6.144972\du}}
\pgfpathlineto{\pgfpoint{2.120954\du}{6.144972\du}}
\pgfpathlineto{\pgfpoint{2.120954\du}{6.144972\du}}
\pgfpathlineto{\pgfpoint{2.130507\du}{6.144972\du}}
\pgfpathlineto{\pgfpoint{2.140658\du}{6.144674\du}}
\pgfpathlineto{\pgfpoint{2.149615\du}{6.144375\du}}
\pgfpathlineto{\pgfpoint{2.159168\du}{6.143479\du}}
\pgfpathlineto{\pgfpoint{2.168722\du}{6.143181\du}}
\pgfpathlineto{\pgfpoint{2.177977\du}{6.141987\du}}
\pgfpathlineto{\pgfpoint{2.186934\du}{6.141688\du}}
\pgfpathlineto{\pgfpoint{2.195592\du}{6.140195\du}}
\pgfpathlineto{\pgfpoint{2.204548\du}{6.139300\du}}
\pgfpathlineto{\pgfpoint{2.212908\du}{6.137508\du}}
\pgfpathlineto{\pgfpoint{2.221267\du}{6.136314\du}}
\pgfpathlineto{\pgfpoint{2.229925\du}{6.134224\du}}
\pgfpathlineto{\pgfpoint{2.237688\du}{6.132433\du}}
\pgfpathlineto{\pgfpoint{2.244853\du}{6.130642\du}}
\pgfpathlineto{\pgfpoint{2.251720\du}{6.128850\du}}
\pgfpathlineto{\pgfpoint{2.259183\du}{6.126760\du}}
\pgfpathlineto{\pgfpoint{2.265154\du}{6.124372\du}}
\pgfpathlineto{\pgfpoint{2.271424\du}{6.121984\du}}
\pgfpathlineto{\pgfpoint{2.276798\du}{6.119297\du}}
\pgfpathlineto{\pgfpoint{2.282470\du}{6.117207\du}}
\pgfpathlineto{\pgfpoint{2.287247\du}{6.114520\du}}
\pgfpathlineto{\pgfpoint{2.291725\du}{6.111833\du}}
\pgfpathlineto{\pgfpoint{2.295905\du}{6.109146\du}}
\pgfpathlineto{\pgfpoint{2.299189\du}{6.106459\du}}
\pgfpathlineto{\pgfpoint{2.302473\du}{6.103772\du}}
\pgfpathlineto{\pgfpoint{2.304862\du}{6.100488\du}}
\pgfpathlineto{\pgfpoint{2.307250\du}{6.097801\du}}
\pgfpathlineto{\pgfpoint{2.308743\du}{6.094517\du}}
\pgfpathlineto{\pgfpoint{2.309937\du}{6.091531\du}}
\pgfpathlineto{\pgfpoint{2.310833\du}{6.088844\du}}
\pgfpathlineto{\pgfpoint{2.310833\du}{6.085560\du}}
\pgfusepath{stroke}
\pgfsetbuttcap
\pgfsetmiterjoin
\pgfsetdash{}{0pt}
\definecolor{dialinecolor}{rgb}{0.000000, 0.000000, 0.000000}
\pgfsetfillcolor{dialinecolor}
\pgfpathmoveto{\pgfpoint{2.133194\du}{6.073917\du}}
\pgfpathlineto{\pgfpoint{2.156481\du}{6.085560\du}}
\pgfpathlineto{\pgfpoint{2.228432\du}{6.061975\du}}
\pgfpathlineto{\pgfpoint{2.263960\du}{6.073917\du}}
\pgfpathlineto{\pgfpoint{2.240375\du}{6.050033\du}}
\pgfpathlineto{\pgfpoint{2.156481\du}{6.050033\du}}
\pgfpathlineto{\pgfpoint{2.192308\du}{6.050033\du}}
\pgfpathlineto{\pgfpoint{2.133194\du}{6.073917\du}}
\pgfusepath{fill}
\pgfsetbuttcap
\pgfsetmiterjoin
\pgfsetdash{}{0pt}
\definecolor{dialinecolor}{rgb}{0.000000, 0.000000, 0.000000}
\pgfsetfillcolor{dialinecolor}
\pgfpathmoveto{\pgfpoint{2.120954\du}{6.109444\du}}
\pgfpathlineto{\pgfpoint{2.085127\du}{6.097801\du}}
\pgfpathlineto{\pgfpoint{2.026014\du}{6.121685\du}}
\pgfpathlineto{\pgfpoint{1.990188\du}{6.109444\du}}
\pgfpathlineto{\pgfpoint{2.002429\du}{6.144972\du}}
\pgfpathlineto{\pgfpoint{2.085127\du}{6.144972\du}}
\pgfpathlineto{\pgfpoint{2.049600\du}{6.133627\du}}
\pgfpathlineto{\pgfpoint{2.120954\du}{6.109444\du}}
\pgfusepath{fill}
\pgfsetbuttcap
\pgfsetmiterjoin
\pgfsetdash{}{0pt}
\definecolor{dialinecolor}{rgb}{0.000000, 0.000000, 0.000000}
\pgfsetfillcolor{dialinecolor}
\pgfpathmoveto{\pgfpoint{2.002429\du}{6.050033\du}}
\pgfpathlineto{\pgfpoint{2.026014\du}{6.037493\du}}
\pgfpathlineto{\pgfpoint{2.097667\du}{6.073917\du}}
\pgfpathlineto{\pgfpoint{2.120954\du}{6.061975\du}}
\pgfpathlineto{\pgfpoint{2.109310\du}{6.085560\du}}
\pgfpathlineto{\pgfpoint{2.026014\du}{6.085560\du}}
\pgfpathlineto{\pgfpoint{2.061840\du}{6.073917\du}}
\pgfpathlineto{\pgfpoint{2.002429\du}{6.050033\du}}
\pgfusepath{fill}
\pgfsetbuttcap
\pgfsetmiterjoin
\pgfsetdash{}{0pt}
\definecolor{dialinecolor}{rgb}{0.000000, 0.000000, 0.000000}
\pgfsetfillcolor{dialinecolor}
\pgfpathmoveto{\pgfpoint{2.251720\du}{6.133627\du}}
\pgfpathlineto{\pgfpoint{2.228432\du}{6.144972\du}}
\pgfpathlineto{\pgfpoint{2.156481\du}{6.109444\du}}
\pgfpathlineto{\pgfpoint{2.120954\du}{6.121685\du}}
\pgfpathlineto{\pgfpoint{2.144838\du}{6.097801\du}}
\pgfpathlineto{\pgfpoint{2.228432\du}{6.097801\du}}
\pgfpathlineto{\pgfpoint{2.192308\du}{6.109444\du}}
\pgfpathlineto{\pgfpoint{2.251720\du}{6.133627\du}}
\pgfusepath{fill}
\pgfsetbuttcap
\pgfsetmiterjoin
\pgfsetdash{}{0pt}
\definecolor{dialinecolor}{rgb}{1.000000, 1.000000, 1.000000}
\pgfsetfillcolor{dialinecolor}
\pgfpathmoveto{\pgfpoint{2.133194\du}{6.085560\du}}
\pgfpathlineto{\pgfpoint{2.156481\du}{6.085560\du}}
\pgfpathlineto{\pgfpoint{2.228432\du}{6.061975\du}}
\pgfpathlineto{\pgfpoint{2.263960\du}{6.073917\du}}
\pgfpathlineto{\pgfpoint{2.240375\du}{6.050033\du}}
\pgfpathlineto{\pgfpoint{2.168722\du}{6.050033\du}}
\pgfpathlineto{\pgfpoint{2.192308\du}{6.050033\du}}
\pgfpathlineto{\pgfpoint{2.133194\du}{6.085560\du}}
\pgfusepath{fill}
\pgfsetbuttcap
\pgfsetmiterjoin
\pgfsetdash{}{0pt}
\definecolor{dialinecolor}{rgb}{1.000000, 1.000000, 1.000000}
\pgfsetfillcolor{dialinecolor}
\pgfpathmoveto{\pgfpoint{2.120954\du}{6.109444\du}}
\pgfpathlineto{\pgfpoint{2.097667\du}{6.097801\du}}
\pgfpathlineto{\pgfpoint{2.026014\du}{6.121685\du}}
\pgfpathlineto{\pgfpoint{1.990188\du}{6.121685\du}}
\pgfpathlineto{\pgfpoint{2.013774\du}{6.144972\du}}
\pgfpathlineto{\pgfpoint{2.097667\du}{6.144972\du}}
\pgfpathlineto{\pgfpoint{2.061840\du}{6.133627\du}}
\pgfpathlineto{\pgfpoint{2.120954\du}{6.109444\du}}
\pgfusepath{fill}
\pgfsetbuttcap
\pgfsetmiterjoin
\pgfsetdash{}{0pt}
\definecolor{dialinecolor}{rgb}{1.000000, 1.000000, 1.000000}
\pgfsetfillcolor{dialinecolor}
\pgfpathmoveto{\pgfpoint{2.002429\du}{6.050033\du}}
\pgfpathlineto{\pgfpoint{2.026014\du}{6.050033\du}}
\pgfpathlineto{\pgfpoint{2.097667\du}{6.073917\du}}
\pgfpathlineto{\pgfpoint{2.133194\du}{6.061975\du}}
\pgfpathlineto{\pgfpoint{2.109310\du}{6.085560\du}}
\pgfpathlineto{\pgfpoint{2.026014\du}{6.085560\du}}
\pgfpathlineto{\pgfpoint{2.061840\du}{6.085560\du}}
\pgfpathlineto{\pgfpoint{2.002429\du}{6.050033\du}}
\pgfusepath{fill}
\pgfsetbuttcap
\pgfsetmiterjoin
\pgfsetdash{}{0pt}
\definecolor{dialinecolor}{rgb}{1.000000, 1.000000, 1.000000}
\pgfsetfillcolor{dialinecolor}
\pgfpathmoveto{\pgfpoint{2.251720\du}{6.133627\du}}
\pgfpathlineto{\pgfpoint{2.228432\du}{6.144972\du}}
\pgfpathlineto{\pgfpoint{2.156481\du}{6.121685\du}}
\pgfpathlineto{\pgfpoint{2.133194\du}{6.121685\du}}
\pgfpathlineto{\pgfpoint{2.144838\du}{6.097801\du}}
\pgfpathlineto{\pgfpoint{2.228432\du}{6.097801\du}}
\pgfpathlineto{\pgfpoint{2.192308\du}{6.109444\du}}
\pgfpathlineto{\pgfpoint{2.251720\du}{6.133627\du}}
\pgfusepath{fill}
\pgfsetbuttcap
\pgfsetmiterjoin
\pgfsetdash{}{0pt}
\definecolor{dialinecolor}{rgb}{0.666667, 0.901961, 1.000000}
\pgfsetstrokecolor{dialinecolor}
\pgfpathmoveto{\pgfpoint{1.930776\du}{6.085560\du}}
\pgfpathlineto{\pgfpoint{1.930776\du}{6.180500\du}}
\pgfusepath{stroke}
\pgfsetbuttcap
\pgfsetmiterjoin
\pgfsetdash{}{0pt}
\definecolor{dialinecolor}{rgb}{0.666667, 0.901961, 1.000000}
\pgfsetstrokecolor{dialinecolor}
\pgfpathmoveto{\pgfpoint{2.311430\du}{6.085560\du}}
\pgfpathlineto{\pgfpoint{2.311430\du}{6.180500\du}}
\pgfusepath{stroke}
\pgfsetbuttcap
\pgfsetmiterjoin
\pgfsetdash{}{0pt}
\definecolor{dialinecolor}{rgb}{0.000000, 0.470588, 0.666667}
\pgfsetfillcolor{dialinecolor}
\pgfpathmoveto{\pgfpoint{3.037210\du}{6.186769\du}}
\pgfpathlineto{\pgfpoint{3.037210\du}{6.183485\du}}
\pgfpathlineto{\pgfpoint{3.036314\du}{6.180201\du}}
\pgfpathlineto{\pgfpoint{3.035120\du}{6.176917\du}}
\pgfpathlineto{\pgfpoint{3.033627\du}{6.173633\du}}
\pgfpathlineto{\pgfpoint{3.031239\du}{6.170648\du}}
\pgfpathlineto{\pgfpoint{3.028552\du}{6.167364\du}}
\pgfpathlineto{\pgfpoint{3.025566\du}{6.164079\du}}
\pgfpathlineto{\pgfpoint{3.021387\du}{6.161094\du}}
\pgfpathlineto{\pgfpoint{3.017804\du}{6.158108\du}}
\pgfpathlineto{\pgfpoint{3.013027\du}{6.155123\du}}
\pgfpathlineto{\pgfpoint{3.007952\du}{6.152436\du}}
\pgfpathlineto{\pgfpoint{3.002578\du}{6.149450\du}}
\pgfpathlineto{\pgfpoint{2.996308\du}{6.147062\du}}
\pgfpathlineto{\pgfpoint{2.990337\du}{6.144375\du}}
\pgfpathlineto{\pgfpoint{2.983471\du}{6.141987\du}}
\pgfpathlineto{\pgfpoint{2.976604\du}{6.139598\du}}
\pgfpathlineto{\pgfpoint{2.969140\du}{6.137508\du}}
\pgfpathlineto{\pgfpoint{2.961676\du}{6.135120\du}}
\pgfpathlineto{\pgfpoint{2.953615\du}{6.133627\du}}
\pgfpathlineto{\pgfpoint{2.944957\du}{6.131537\du}}
\pgfpathlineto{\pgfpoint{2.936598\du}{6.129746\du}}
\pgfpathlineto{\pgfpoint{2.927940\du}{6.128552\du}}
\pgfpathlineto{\pgfpoint{2.918685\du}{6.127059\du}}
\pgfpathlineto{\pgfpoint{2.909430\du}{6.125865\du}}
\pgfpathlineto{\pgfpoint{2.900175\du}{6.124671\du}}
\pgfpathlineto{\pgfpoint{2.890621\du}{6.123775\du}}
\pgfpathlineto{\pgfpoint{2.880470\du}{6.122879\du}}
\pgfpathlineto{\pgfpoint{2.870916\du}{6.122282\du}}
\pgfpathlineto{\pgfpoint{2.861363\du}{6.121984\du}}
\pgfpathlineto{\pgfpoint{2.851212\du}{6.121685\du}}
\pgfpathlineto{\pgfpoint{2.841360\du}{6.121685\du}}
\pgfpathlineto{\pgfpoint{2.841360\du}{6.121685\du}}
\pgfpathlineto{\pgfpoint{2.831508\du}{6.121685\du}}
\pgfpathlineto{\pgfpoint{2.821357\du}{6.121984\du}}
\pgfpathlineto{\pgfpoint{2.811505\du}{6.122282\du}}
\pgfpathlineto{\pgfpoint{2.801652\du}{6.122879\du}}
\pgfpathlineto{\pgfpoint{2.792099\du}{6.123775\du}}
\pgfpathlineto{\pgfpoint{2.782545\du}{6.124671\du}}
\pgfpathlineto{\pgfpoint{2.772991\du}{6.125865\du}}
\pgfpathlineto{\pgfpoint{2.764035\du}{6.127059\du}}
\pgfpathlineto{\pgfpoint{2.754780\du}{6.128552\du}}
\pgfpathlineto{\pgfpoint{2.746122\du}{6.129746\du}}
\pgfpathlineto{\pgfpoint{2.737464\du}{6.131537\du}}
\pgfpathlineto{\pgfpoint{2.729104\du}{6.133627\du}}
\pgfpathlineto{\pgfpoint{2.721043\du}{6.135120\du}}
\pgfpathlineto{\pgfpoint{2.713579\du}{6.137508\du}}
\pgfpathlineto{\pgfpoint{2.706116\du}{6.139598\du}}
\pgfpathlineto{\pgfpoint{2.698950\du}{6.141987\du}}
\pgfpathlineto{\pgfpoint{2.692382\du}{6.144375\du}}
\pgfpathlineto{\pgfpoint{2.686113\du}{6.147062\du}}
\pgfpathlineto{\pgfpoint{2.680142\du}{6.149450\du}}
\pgfpathlineto{\pgfpoint{2.674768\du}{6.152436\du}}
\pgfpathlineto{\pgfpoint{2.669692\du}{6.155123\du}}
\pgfpathlineto{\pgfpoint{2.664916\du}{6.158108\du}}
\pgfpathlineto{\pgfpoint{2.661034\du}{6.161094\du}}
\pgfpathlineto{\pgfpoint{2.657153\du}{6.164079\du}}
\pgfpathlineto{\pgfpoint{2.654168\du}{6.167364\du}}
\pgfpathlineto{\pgfpoint{2.651481\du}{6.170648\du}}
\pgfpathlineto{\pgfpoint{2.649092\du}{6.173633\du}}
\pgfpathlineto{\pgfpoint{2.647301\du}{6.176917\du}}
\pgfpathlineto{\pgfpoint{2.645808\du}{6.180201\du}}
\pgfpathlineto{\pgfpoint{2.644913\du}{6.183485\du}}
\pgfpathlineto{\pgfpoint{2.644913\du}{6.186769\du}}
\pgfpathlineto{\pgfpoint{2.644913\du}{6.186769\du}}
\pgfpathlineto{\pgfpoint{2.644913\du}{6.190352\du}}
\pgfpathlineto{\pgfpoint{2.645808\du}{6.193338\du}}
\pgfpathlineto{\pgfpoint{2.647301\du}{6.196622\du}}
\pgfpathlineto{\pgfpoint{2.649092\du}{6.199906\du}}
\pgfpathlineto{\pgfpoint{2.651481\du}{6.203190\du}}
\pgfpathlineto{\pgfpoint{2.654168\du}{6.206474\du}}
\pgfpathlineto{\pgfpoint{2.657153\du}{6.209459\du}}
\pgfpathlineto{\pgfpoint{2.661034\du}{6.212445\du}}
\pgfpathlineto{\pgfpoint{2.664916\du}{6.215430\du}}
\pgfpathlineto{\pgfpoint{2.669692\du}{6.218714\du}}
\pgfpathlineto{\pgfpoint{2.674768\du}{6.221401\du}}
\pgfpathlineto{\pgfpoint{2.680142\du}{6.224088\du}}
\pgfpathlineto{\pgfpoint{2.686113\du}{6.226775\du}}
\pgfpathlineto{\pgfpoint{2.692382\du}{6.229462\du}}
\pgfpathlineto{\pgfpoint{2.698950\du}{6.231851\du}}
\pgfpathlineto{\pgfpoint{2.706116\du}{6.234239\du}}
\pgfpathlineto{\pgfpoint{2.713579\du}{6.236030\du}}
\pgfpathlineto{\pgfpoint{2.721043\du}{6.238120\du}}
\pgfpathlineto{\pgfpoint{2.729104\du}{6.240210\du}}
\pgfpathlineto{\pgfpoint{2.737464\du}{6.242300\du}}
\pgfpathlineto{\pgfpoint{2.746122\du}{6.243793\du}}
\pgfpathlineto{\pgfpoint{2.754780\du}{6.245286\du}}
\pgfpathlineto{\pgfpoint{2.764035\du}{6.246778\du}}
\pgfpathlineto{\pgfpoint{2.772991\du}{6.247973\du}}
\pgfpathlineto{\pgfpoint{2.782545\du}{6.248868\du}}
\pgfpathlineto{\pgfpoint{2.792099\du}{6.250062\du}}
\pgfpathlineto{\pgfpoint{2.801652\du}{6.250660\du}}
\pgfpathlineto{\pgfpoint{2.811505\du}{6.251257\du}}
\pgfpathlineto{\pgfpoint{2.821357\du}{6.251854\du}}
\pgfpathlineto{\pgfpoint{2.831508\du}{6.252152\du}}
\pgfpathlineto{\pgfpoint{2.841360\du}{6.252152\du}}
\pgfpathlineto{\pgfpoint{2.841360\du}{6.252152\du}}
\pgfpathlineto{\pgfpoint{2.851212\du}{6.252152\du}}
\pgfpathlineto{\pgfpoint{2.861363\du}{6.251854\du}}
\pgfpathlineto{\pgfpoint{2.870916\du}{6.251257\du}}
\pgfpathlineto{\pgfpoint{2.880470\du}{6.250660\du}}
\pgfpathlineto{\pgfpoint{2.890621\du}{6.250062\du}}
\pgfpathlineto{\pgfpoint{2.900175\du}{6.248868\du}}
\pgfpathlineto{\pgfpoint{2.909430\du}{6.247973\du}}
\pgfpathlineto{\pgfpoint{2.918685\du}{6.246778\du}}
\pgfpathlineto{\pgfpoint{2.927940\du}{6.245286\du}}
\pgfpathlineto{\pgfpoint{2.936598\du}{6.243793\du}}
\pgfpathlineto{\pgfpoint{2.944957\du}{6.242300\du}}
\pgfpathlineto{\pgfpoint{2.953615\du}{6.240210\du}}
\pgfpathlineto{\pgfpoint{2.961676\du}{6.238120\du}}
\pgfpathlineto{\pgfpoint{2.969140\du}{6.236030\du}}
\pgfpathlineto{\pgfpoint{2.976604\du}{6.234239\du}}
\pgfpathlineto{\pgfpoint{2.983471\du}{6.231851\du}}
\pgfpathlineto{\pgfpoint{2.990337\du}{6.229462\du}}
\pgfpathlineto{\pgfpoint{2.996308\du}{6.226775\du}}
\pgfpathlineto{\pgfpoint{3.002578\du}{6.224088\du}}
\pgfpathlineto{\pgfpoint{3.007952\du}{6.221401\du}}
\pgfpathlineto{\pgfpoint{3.013027\du}{6.218714\du}}
\pgfpathlineto{\pgfpoint{3.017804\du}{6.215430\du}}
\pgfpathlineto{\pgfpoint{3.021387\du}{6.212445\du}}
\pgfpathlineto{\pgfpoint{3.025566\du}{6.209459\du}}
\pgfpathlineto{\pgfpoint{3.028552\du}{6.206474\du}}
\pgfpathlineto{\pgfpoint{3.031239\du}{6.203190\du}}
\pgfpathlineto{\pgfpoint{3.033627\du}{6.199906\du}}
\pgfpathlineto{\pgfpoint{3.035120\du}{6.196622\du}}
\pgfpathlineto{\pgfpoint{3.036314\du}{6.193338\du}}
\pgfpathlineto{\pgfpoint{3.037210\du}{6.190352\du}}
\pgfpathlineto{\pgfpoint{3.037210\du}{6.186769\du}}
\pgfusepath{fill}
\pgfsetbuttcap
\pgfsetmiterjoin
\pgfsetdash{}{0pt}
\definecolor{dialinecolor}{rgb}{0.666667, 0.901961, 1.000000}
\pgfsetstrokecolor{dialinecolor}
\pgfpathmoveto{\pgfpoint{3.025566\du}{6.181097\du}}
\pgfpathlineto{\pgfpoint{3.025566\du}{6.177514\du}}
\pgfpathlineto{\pgfpoint{3.024671\du}{6.175126\du}}
\pgfpathlineto{\pgfpoint{3.023476\du}{6.171842\du}}
\pgfpathlineto{\pgfpoint{3.021387\du}{6.168856\du}}
\pgfpathlineto{\pgfpoint{3.019297\du}{6.166169\du}}
\pgfpathlineto{\pgfpoint{3.017207\du}{6.162587\du}}
\pgfpathlineto{\pgfpoint{3.013923\du}{6.160198\du}}
\pgfpathlineto{\pgfpoint{3.010639\du}{6.157511\du}}
\pgfpathlineto{\pgfpoint{3.006459\du}{6.154824\du}}
\pgfpathlineto{\pgfpoint{3.001682\du}{6.152436\du}}
\pgfpathlineto{\pgfpoint{2.997204\du}{6.149450\du}}
\pgfpathlineto{\pgfpoint{2.991531\du}{6.147361\du}}
\pgfpathlineto{\pgfpoint{2.985560\du}{6.144674\du}}
\pgfpathlineto{\pgfpoint{2.979888\du}{6.141987\du}}
\pgfpathlineto{\pgfpoint{2.973320\du}{6.140195\du}}
\pgfpathlineto{\pgfpoint{2.966453\du}{6.137807\du}}
\pgfpathlineto{\pgfpoint{2.959586\du}{6.136314\du}}
\pgfpathlineto{\pgfpoint{2.951824\du}{6.134224\du}}
\pgfpathlineto{\pgfpoint{2.944360\du}{6.132433\du}}
\pgfpathlineto{\pgfpoint{2.936001\du}{6.130642\du}}
\pgfpathlineto{\pgfpoint{2.927343\du}{6.129447\du}}
\pgfpathlineto{\pgfpoint{2.918983\du}{6.127059\du}}
\pgfpathlineto{\pgfpoint{2.910325\du}{6.126760\du}}
\pgfpathlineto{\pgfpoint{2.901369\du}{6.125268\du}}
\pgfpathlineto{\pgfpoint{2.892711\du}{6.124372\du}}
\pgfpathlineto{\pgfpoint{2.883157\du}{6.123476\du}}
\pgfpathlineto{\pgfpoint{2.873902\du}{6.122879\du}}
\pgfpathlineto{\pgfpoint{2.864348\du}{6.122581\du}}
\pgfpathlineto{\pgfpoint{2.854795\du}{6.121984\du}}
\pgfpathlineto{\pgfpoint{2.844644\du}{6.121685\du}}
\pgfpathlineto{\pgfpoint{2.834792\du}{6.121685\du}}
\pgfpathlineto{\pgfpoint{2.834792\du}{6.121685\du}}
\pgfpathlineto{\pgfpoint{2.825835\du}{6.121685\du}}
\pgfpathlineto{\pgfpoint{2.815983\du}{6.121984\du}}
\pgfpathlineto{\pgfpoint{2.806429\du}{6.122581\du}}
\pgfpathlineto{\pgfpoint{2.797174\du}{6.122879\du}}
\pgfpathlineto{\pgfpoint{2.787620\du}{6.123476\du}}
\pgfpathlineto{\pgfpoint{2.778067\du}{6.124372\du}}
\pgfpathlineto{\pgfpoint{2.769110\du}{6.125268\du}}
\pgfpathlineto{\pgfpoint{2.760154\du}{6.126760\du}}
\pgfpathlineto{\pgfpoint{2.751496\du}{6.127059\du}}
\pgfpathlineto{\pgfpoint{2.743136\du}{6.129447\du}}
\pgfpathlineto{\pgfpoint{2.734478\du}{6.130642\du}}
\pgfpathlineto{\pgfpoint{2.726417\du}{6.132433\du}}
\pgfpathlineto{\pgfpoint{2.718655\du}{6.134224\du}}
\pgfpathlineto{\pgfpoint{2.711191\du}{6.136314\du}}
\pgfpathlineto{\pgfpoint{2.704623\du}{6.137807\du}}
\pgfpathlineto{\pgfpoint{2.697159\du}{6.140195\du}}
\pgfpathlineto{\pgfpoint{2.690890\du}{6.141987\du}}
\pgfpathlineto{\pgfpoint{2.684918\du}{6.144674\du}}
\pgfpathlineto{\pgfpoint{2.678947\du}{6.147361\du}}
\pgfpathlineto{\pgfpoint{2.673872\du}{6.149450\du}}
\pgfpathlineto{\pgfpoint{2.668797\du}{6.152436\du}}
\pgfpathlineto{\pgfpoint{2.664020\du}{6.154824\du}}
\pgfpathlineto{\pgfpoint{2.659840\du}{6.157511\du}}
\pgfpathlineto{\pgfpoint{2.656556\du}{6.160198\du}}
\pgfpathlineto{\pgfpoint{2.653272\du}{6.162587\du}}
\pgfpathlineto{\pgfpoint{2.650884\du}{6.166169\du}}
\pgfpathlineto{\pgfpoint{2.649092\du}{6.168856\du}}
\pgfpathlineto{\pgfpoint{2.646704\du}{6.171842\du}}
\pgfpathlineto{\pgfpoint{2.645808\du}{6.175126\du}}
\pgfpathlineto{\pgfpoint{2.644913\du}{6.177514\du}}
\pgfpathlineto{\pgfpoint{2.644913\du}{6.181097\du}}
\pgfpathlineto{\pgfpoint{2.644913\du}{6.181097\du}}
\pgfpathlineto{\pgfpoint{2.644913\du}{6.184082\du}}
\pgfpathlineto{\pgfpoint{2.645808\du}{6.186769\du}}
\pgfpathlineto{\pgfpoint{2.646704\du}{6.189755\du}}
\pgfpathlineto{\pgfpoint{2.649092\du}{6.193338\du}}
\pgfpathlineto{\pgfpoint{2.650884\du}{6.196024\du}}
\pgfpathlineto{\pgfpoint{2.653272\du}{6.199010\du}}
\pgfpathlineto{\pgfpoint{2.656556\du}{6.201697\du}}
\pgfpathlineto{\pgfpoint{2.659840\du}{6.204384\du}}
\pgfpathlineto{\pgfpoint{2.664020\du}{6.207071\du}}
\pgfpathlineto{\pgfpoint{2.668797\du}{6.209758\du}}
\pgfpathlineto{\pgfpoint{2.673872\du}{6.212445\du}}
\pgfpathlineto{\pgfpoint{2.678947\du}{6.214833\du}}
\pgfpathlineto{\pgfpoint{2.684918\du}{6.217520\du}}
\pgfpathlineto{\pgfpoint{2.690890\du}{6.219610\du}}
\pgfpathlineto{\pgfpoint{2.697159\du}{6.221999\du}}
\pgfpathlineto{\pgfpoint{2.704623\du}{6.224387\du}}
\pgfpathlineto{\pgfpoint{2.711191\du}{6.226178\du}}
\pgfpathlineto{\pgfpoint{2.718655\du}{6.227671\du}}
\pgfpathlineto{\pgfpoint{2.726417\du}{6.229761\du}}
\pgfpathlineto{\pgfpoint{2.734478\du}{6.231552\du}}
\pgfpathlineto{\pgfpoint{2.743136\du}{6.232746\du}}
\pgfpathlineto{\pgfpoint{2.751496\du}{6.234538\du}}
\pgfpathlineto{\pgfpoint{2.760154\du}{6.235433\du}}
\pgfpathlineto{\pgfpoint{2.769110\du}{6.236926\du}}
\pgfpathlineto{\pgfpoint{2.778067\du}{6.237822\du}}
\pgfpathlineto{\pgfpoint{2.787620\du}{6.238120\du}}
\pgfpathlineto{\pgfpoint{2.797174\du}{6.239315\du}}
\pgfpathlineto{\pgfpoint{2.806429\du}{6.239613\du}}
\pgfpathlineto{\pgfpoint{2.815983\du}{6.239912\du}}
\pgfpathlineto{\pgfpoint{2.825835\du}{6.240509\du}}
\pgfpathlineto{\pgfpoint{2.834792\du}{6.240509\du}}
\pgfpathlineto{\pgfpoint{2.834792\du}{6.240509\du}}
\pgfpathlineto{\pgfpoint{2.844644\du}{6.240509\du}}
\pgfpathlineto{\pgfpoint{2.854795\du}{6.239912\du}}
\pgfpathlineto{\pgfpoint{2.864348\du}{6.239613\du}}
\pgfpathlineto{\pgfpoint{2.873902\du}{6.239315\du}}
\pgfpathlineto{\pgfpoint{2.883157\du}{6.238120\du}}
\pgfpathlineto{\pgfpoint{2.892711\du}{6.237822\du}}
\pgfpathlineto{\pgfpoint{2.901369\du}{6.236926\du}}
\pgfpathlineto{\pgfpoint{2.910325\du}{6.235433\du}}
\pgfpathlineto{\pgfpoint{2.918983\du}{6.234538\du}}
\pgfpathlineto{\pgfpoint{2.927343\du}{6.232746\du}}
\pgfpathlineto{\pgfpoint{2.936001\du}{6.231552\du}}
\pgfpathlineto{\pgfpoint{2.944360\du}{6.229761\du}}
\pgfpathlineto{\pgfpoint{2.951824\du}{6.227671\du}}
\pgfpathlineto{\pgfpoint{2.959586\du}{6.226178\du}}
\pgfpathlineto{\pgfpoint{2.966453\du}{6.224387\du}}
\pgfpathlineto{\pgfpoint{2.973320\du}{6.221999\du}}
\pgfpathlineto{\pgfpoint{2.979888\du}{6.219610\du}}
\pgfpathlineto{\pgfpoint{2.985560\du}{6.217520\du}}
\pgfpathlineto{\pgfpoint{2.991531\du}{6.214833\du}}
\pgfpathlineto{\pgfpoint{2.997204\du}{6.212445\du}}
\pgfpathlineto{\pgfpoint{3.001682\du}{6.209758\du}}
\pgfpathlineto{\pgfpoint{3.006459\du}{6.207071\du}}
\pgfpathlineto{\pgfpoint{3.010639\du}{6.204384\du}}
\pgfpathlineto{\pgfpoint{3.013923\du}{6.201697\du}}
\pgfpathlineto{\pgfpoint{3.017207\du}{6.199010\du}}
\pgfpathlineto{\pgfpoint{3.019297\du}{6.196024\du}}
\pgfpathlineto{\pgfpoint{3.021387\du}{6.193338\du}}
\pgfpathlineto{\pgfpoint{3.023476\du}{6.189755\du}}
\pgfpathlineto{\pgfpoint{3.024671\du}{6.186769\du}}
\pgfpathlineto{\pgfpoint{3.025566\du}{6.184082\du}}
\pgfpathlineto{\pgfpoint{3.025566\du}{6.181097\du}}
\pgfusepath{stroke}
\pgfsetbuttcap
\pgfsetmiterjoin
\pgfsetdash{}{0pt}
\definecolor{dialinecolor}{rgb}{0.000000, 0.470588, 0.666667}
\pgfsetfillcolor{dialinecolor}
\pgfpathmoveto{\pgfpoint{2.644913\du}{6.097801\du}}
\pgfpathlineto{\pgfpoint{2.644913\du}{6.192442\du}}
\pgfpathlineto{\pgfpoint{3.037807\du}{6.192442\du}}
\pgfpathlineto{\pgfpoint{3.037807\du}{6.097801\du}}
\pgfpathlineto{\pgfpoint{2.644913\du}{6.097801\du}}
\pgfusepath{fill}
\pgfsetbuttcap
\pgfsetmiterjoin
\pgfsetdash{}{0pt}
\definecolor{dialinecolor}{rgb}{0.000000, 0.705882, 1.000000}
\pgfsetfillcolor{dialinecolor}
\pgfpathmoveto{\pgfpoint{3.037210\du}{6.091233\du}}
\pgfpathlineto{\pgfpoint{3.037210\du}{6.087949\du}}
\pgfpathlineto{\pgfpoint{3.036314\du}{6.084366\du}}
\pgfpathlineto{\pgfpoint{3.035120\du}{6.081381\du}}
\pgfpathlineto{\pgfpoint{3.033627\du}{6.078096\du}}
\pgfpathlineto{\pgfpoint{3.031239\du}{6.074812\du}}
\pgfpathlineto{\pgfpoint{3.028552\du}{6.071528\du}}
\pgfpathlineto{\pgfpoint{3.025566\du}{6.068543\du}}
\pgfpathlineto{\pgfpoint{3.021387\du}{6.065259\du}}
\pgfpathlineto{\pgfpoint{3.017804\du}{6.062572\du}}
\pgfpathlineto{\pgfpoint{3.013027\du}{6.059586\du}}
\pgfpathlineto{\pgfpoint{3.007952\du}{6.056601\du}}
\pgfpathlineto{\pgfpoint{3.002578\du}{6.053914\du}}
\pgfpathlineto{\pgfpoint{2.996308\du}{6.050630\du}}
\pgfpathlineto{\pgfpoint{2.990337\du}{6.048838\du}}
\pgfpathlineto{\pgfpoint{2.983471\du}{6.046151\du}}
\pgfpathlineto{\pgfpoint{2.976604\du}{6.043763\du}}
\pgfpathlineto{\pgfpoint{2.969140\du}{6.041673\du}}
\pgfpathlineto{\pgfpoint{2.961676\du}{6.039583\du}}
\pgfpathlineto{\pgfpoint{2.953615\du}{6.037493\du}}
\pgfpathlineto{\pgfpoint{2.944957\du}{6.035404\du}}
\pgfpathlineto{\pgfpoint{2.936598\du}{6.033911\du}}
\pgfpathlineto{\pgfpoint{2.927940\du}{6.032418\du}}
\pgfpathlineto{\pgfpoint{2.918685\du}{6.031224\du}}
\pgfpathlineto{\pgfpoint{2.909430\du}{6.029731\du}}
\pgfpathlineto{\pgfpoint{2.900175\du}{6.028835\du}}
\pgfpathlineto{\pgfpoint{2.890621\du}{6.027641\du}}
\pgfpathlineto{\pgfpoint{2.880470\du}{6.027044\du}}
\pgfpathlineto{\pgfpoint{2.870916\du}{6.026447\du}}
\pgfpathlineto{\pgfpoint{2.861363\du}{6.026148\du}}
\pgfpathlineto{\pgfpoint{2.851212\du}{6.025850\du}}
\pgfpathlineto{\pgfpoint{2.841360\du}{6.025850\du}}
\pgfpathlineto{\pgfpoint{2.841360\du}{6.025850\du}}
\pgfpathlineto{\pgfpoint{2.831508\du}{6.025850\du}}
\pgfpathlineto{\pgfpoint{2.821357\du}{6.026148\du}}
\pgfpathlineto{\pgfpoint{2.811505\du}{6.026447\du}}
\pgfpathlineto{\pgfpoint{2.801652\du}{6.027044\du}}
\pgfpathlineto{\pgfpoint{2.792099\du}{6.027641\du}}
\pgfpathlineto{\pgfpoint{2.782545\du}{6.028835\du}}
\pgfpathlineto{\pgfpoint{2.772991\du}{6.029731\du}}
\pgfpathlineto{\pgfpoint{2.764035\du}{6.031224\du}}
\pgfpathlineto{\pgfpoint{2.754780\du}{6.032418\du}}
\pgfpathlineto{\pgfpoint{2.746122\du}{6.033911\du}}
\pgfpathlineto{\pgfpoint{2.737464\du}{6.035404\du}}
\pgfpathlineto{\pgfpoint{2.729104\du}{6.037493\du}}
\pgfpathlineto{\pgfpoint{2.721043\du}{6.039583\du}}
\pgfpathlineto{\pgfpoint{2.713579\du}{6.041673\du}}
\pgfpathlineto{\pgfpoint{2.706116\du}{6.043763\du}}
\pgfpathlineto{\pgfpoint{2.698950\du}{6.046151\du}}
\pgfpathlineto{\pgfpoint{2.692382\du}{6.048838\du}}
\pgfpathlineto{\pgfpoint{2.686113\du}{6.050630\du}}
\pgfpathlineto{\pgfpoint{2.680142\du}{6.053914\du}}
\pgfpathlineto{\pgfpoint{2.674768\du}{6.056601\du}}
\pgfpathlineto{\pgfpoint{2.669692\du}{6.059586\du}}
\pgfpathlineto{\pgfpoint{2.664916\du}{6.062572\du}}
\pgfpathlineto{\pgfpoint{2.661034\du}{6.065259\du}}
\pgfpathlineto{\pgfpoint{2.657153\du}{6.068543\du}}
\pgfpathlineto{\pgfpoint{2.654168\du}{6.071528\du}}
\pgfpathlineto{\pgfpoint{2.651481\du}{6.074812\du}}
\pgfpathlineto{\pgfpoint{2.649092\du}{6.078096\du}}
\pgfpathlineto{\pgfpoint{2.647301\du}{6.081381\du}}
\pgfpathlineto{\pgfpoint{2.645808\du}{6.084366\du}}
\pgfpathlineto{\pgfpoint{2.644913\du}{6.087949\du}}
\pgfpathlineto{\pgfpoint{2.644913\du}{6.091233\du}}
\pgfpathlineto{\pgfpoint{2.644913\du}{6.091233\du}}
\pgfpathlineto{\pgfpoint{2.644913\du}{6.094517\du}}
\pgfpathlineto{\pgfpoint{2.645808\du}{6.098099\du}}
\pgfpathlineto{\pgfpoint{2.647301\du}{6.101085\du}}
\pgfpathlineto{\pgfpoint{2.649092\du}{6.104369\du}}
\pgfpathlineto{\pgfpoint{2.651481\du}{6.107952\du}}
\pgfpathlineto{\pgfpoint{2.654168\du}{6.110937\du}}
\pgfpathlineto{\pgfpoint{2.657153\du}{6.113923\du}}
\pgfpathlineto{\pgfpoint{2.661034\du}{6.117207\du}}
\pgfpathlineto{\pgfpoint{2.664916\du}{6.119595\du}}
\pgfpathlineto{\pgfpoint{2.669692\du}{6.122879\du}}
\pgfpathlineto{\pgfpoint{2.674768\du}{6.126163\du}}
\pgfpathlineto{\pgfpoint{2.680142\du}{6.128850\du}}
\pgfpathlineto{\pgfpoint{2.686113\du}{6.131537\du}}
\pgfpathlineto{\pgfpoint{2.692382\du}{6.133627\du}}
\pgfpathlineto{\pgfpoint{2.698950\du}{6.136613\du}}
\pgfpathlineto{\pgfpoint{2.706116\du}{6.139001\du}}
\pgfpathlineto{\pgfpoint{2.713579\du}{6.140792\du}}
\pgfpathlineto{\pgfpoint{2.721043\du}{6.142882\du}}
\pgfpathlineto{\pgfpoint{2.729104\du}{6.144972\du}}
\pgfpathlineto{\pgfpoint{2.737464\du}{6.147062\du}}
\pgfpathlineto{\pgfpoint{2.746122\du}{6.148555\du}}
\pgfpathlineto{\pgfpoint{2.754780\du}{6.150047\du}}
\pgfpathlineto{\pgfpoint{2.764035\du}{6.151242\du}}
\pgfpathlineto{\pgfpoint{2.772991\du}{6.152734\du}}
\pgfpathlineto{\pgfpoint{2.782545\du}{6.153630\du}}
\pgfpathlineto{\pgfpoint{2.792099\du}{6.154824\du}}
\pgfpathlineto{\pgfpoint{2.801652\du}{6.155123\du}}
\pgfpathlineto{\pgfpoint{2.811505\du}{6.156019\du}}
\pgfpathlineto{\pgfpoint{2.821357\du}{6.156019\du}}
\pgfpathlineto{\pgfpoint{2.831508\du}{6.156914\du}}
\pgfpathlineto{\pgfpoint{2.841360\du}{6.156914\du}}
\pgfpathlineto{\pgfpoint{2.841360\du}{6.156914\du}}
\pgfpathlineto{\pgfpoint{2.851212\du}{6.156914\du}}
\pgfpathlineto{\pgfpoint{2.861363\du}{6.156019\du}}
\pgfpathlineto{\pgfpoint{2.870916\du}{6.156019\du}}
\pgfpathlineto{\pgfpoint{2.880470\du}{6.155123\du}}
\pgfpathlineto{\pgfpoint{2.890621\du}{6.154824\du}}
\pgfpathlineto{\pgfpoint{2.900175\du}{6.153630\du}}
\pgfpathlineto{\pgfpoint{2.909430\du}{6.152734\du}}
\pgfpathlineto{\pgfpoint{2.918685\du}{6.151242\du}}
\pgfpathlineto{\pgfpoint{2.927940\du}{6.150047\du}}
\pgfpathlineto{\pgfpoint{2.936598\du}{6.148555\du}}
\pgfpathlineto{\pgfpoint{2.944957\du}{6.147062\du}}
\pgfpathlineto{\pgfpoint{2.953615\du}{6.144972\du}}
\pgfpathlineto{\pgfpoint{2.961676\du}{6.142882\du}}
\pgfpathlineto{\pgfpoint{2.969140\du}{6.140792\du}}
\pgfpathlineto{\pgfpoint{2.976604\du}{6.139001\du}}
\pgfpathlineto{\pgfpoint{2.983471\du}{6.136613\du}}
\pgfpathlineto{\pgfpoint{2.990337\du}{6.133627\du}}
\pgfpathlineto{\pgfpoint{2.996308\du}{6.131537\du}}
\pgfpathlineto{\pgfpoint{3.002578\du}{6.128850\du}}
\pgfpathlineto{\pgfpoint{3.007952\du}{6.126163\du}}
\pgfpathlineto{\pgfpoint{3.013027\du}{6.122879\du}}
\pgfpathlineto{\pgfpoint{3.017804\du}{6.119595\du}}
\pgfpathlineto{\pgfpoint{3.021387\du}{6.117207\du}}
\pgfpathlineto{\pgfpoint{3.025566\du}{6.113923\du}}
\pgfpathlineto{\pgfpoint{3.028552\du}{6.110937\du}}
\pgfpathlineto{\pgfpoint{3.031239\du}{6.107952\du}}
\pgfpathlineto{\pgfpoint{3.033627\du}{6.104369\du}}
\pgfpathlineto{\pgfpoint{3.035120\du}{6.101085\du}}
\pgfpathlineto{\pgfpoint{3.036314\du}{6.098099\du}}
\pgfpathlineto{\pgfpoint{3.037210\du}{6.094517\du}}
\pgfpathlineto{\pgfpoint{3.037210\du}{6.091233\du}}
\pgfusepath{fill}
\pgfsetbuttcap
\pgfsetmiterjoin
\pgfsetdash{}{0pt}
\definecolor{dialinecolor}{rgb}{0.666667, 0.901961, 1.000000}
\pgfsetstrokecolor{dialinecolor}
\pgfpathmoveto{\pgfpoint{3.025566\du}{6.085560\du}}
\pgfpathlineto{\pgfpoint{3.025566\du}{6.082276\du}}
\pgfpathlineto{\pgfpoint{3.024671\du}{6.079888\du}}
\pgfpathlineto{\pgfpoint{3.023476\du}{6.076604\du}}
\pgfpathlineto{\pgfpoint{3.021387\du}{6.073021\du}}
\pgfpathlineto{\pgfpoint{3.019297\du}{6.070633\du}}
\pgfpathlineto{\pgfpoint{3.017207\du}{6.067349\du}}
\pgfpathlineto{\pgfpoint{3.013923\du}{6.064960\du}}
\pgfpathlineto{\pgfpoint{3.010639\du}{6.062273\du}}
\pgfpathlineto{\pgfpoint{3.006459\du}{6.059586\du}}
\pgfpathlineto{\pgfpoint{3.001682\du}{6.056899\du}}
\pgfpathlineto{\pgfpoint{2.997204\du}{6.054212\du}}
\pgfpathlineto{\pgfpoint{2.991531\du}{6.051525\du}}
\pgfpathlineto{\pgfpoint{2.985560\du}{6.049137\du}}
\pgfpathlineto{\pgfpoint{2.979888\du}{6.047047\du}}
\pgfpathlineto{\pgfpoint{2.973320\du}{6.044659\du}}
\pgfpathlineto{\pgfpoint{2.966453\du}{6.042270\du}}
\pgfpathlineto{\pgfpoint{2.959586\du}{6.040479\du}}
\pgfpathlineto{\pgfpoint{2.951824\du}{6.038986\du}}
\pgfpathlineto{\pgfpoint{2.944360\du}{6.036598\du}}
\pgfpathlineto{\pgfpoint{2.936001\du}{6.035105\du}}
\pgfpathlineto{\pgfpoint{2.927343\du}{6.033911\du}}
\pgfpathlineto{\pgfpoint{2.918983\du}{6.032119\du}}
\pgfpathlineto{\pgfpoint{2.910325\du}{6.031224\du}}
\pgfpathlineto{\pgfpoint{2.901369\du}{6.029731\du}}
\pgfpathlineto{\pgfpoint{2.892711\du}{6.028835\du}}
\pgfpathlineto{\pgfpoint{2.883157\du}{6.028238\du}}
\pgfpathlineto{\pgfpoint{2.873902\du}{6.027343\du}}
\pgfpathlineto{\pgfpoint{2.864348\du}{6.027044\du}}
\pgfpathlineto{\pgfpoint{2.854795\du}{6.026745\du}}
\pgfpathlineto{\pgfpoint{2.844644\du}{6.026148\du}}
\pgfpathlineto{\pgfpoint{2.834792\du}{6.026148\du}}
\pgfpathlineto{\pgfpoint{2.834792\du}{6.026148\du}}
\pgfpathlineto{\pgfpoint{2.825835\du}{6.026148\du}}
\pgfpathlineto{\pgfpoint{2.815983\du}{6.026745\du}}
\pgfpathlineto{\pgfpoint{2.806429\du}{6.027044\du}}
\pgfpathlineto{\pgfpoint{2.797174\du}{6.027343\du}}
\pgfpathlineto{\pgfpoint{2.787620\du}{6.028238\du}}
\pgfpathlineto{\pgfpoint{2.778067\du}{6.028835\du}}
\pgfpathlineto{\pgfpoint{2.769110\du}{6.029731\du}}
\pgfpathlineto{\pgfpoint{2.760154\du}{6.031224\du}}
\pgfpathlineto{\pgfpoint{2.751496\du}{6.032119\du}}
\pgfpathlineto{\pgfpoint{2.743136\du}{6.033911\du}}
\pgfpathlineto{\pgfpoint{2.734478\du}{6.035105\du}}
\pgfpathlineto{\pgfpoint{2.726417\du}{6.036598\du}}
\pgfpathlineto{\pgfpoint{2.718655\du}{6.038986\du}}
\pgfpathlineto{\pgfpoint{2.711191\du}{6.040479\du}}
\pgfpathlineto{\pgfpoint{2.704623\du}{6.042270\du}}
\pgfpathlineto{\pgfpoint{2.697159\du}{6.044659\du}}
\pgfpathlineto{\pgfpoint{2.690890\du}{6.047047\du}}
\pgfpathlineto{\pgfpoint{2.684918\du}{6.049137\du}}
\pgfpathlineto{\pgfpoint{2.678947\du}{6.051525\du}}
\pgfpathlineto{\pgfpoint{2.673872\du}{6.054212\du}}
\pgfpathlineto{\pgfpoint{2.668797\du}{6.056899\du}}
\pgfpathlineto{\pgfpoint{2.664020\du}{6.059586\du}}
\pgfpathlineto{\pgfpoint{2.659840\du}{6.062273\du}}
\pgfpathlineto{\pgfpoint{2.656556\du}{6.064960\du}}
\pgfpathlineto{\pgfpoint{2.653272\du}{6.067349\du}}
\pgfpathlineto{\pgfpoint{2.650884\du}{6.070633\du}}
\pgfpathlineto{\pgfpoint{2.649092\du}{6.073021\du}}
\pgfpathlineto{\pgfpoint{2.646704\du}{6.076604\du}}
\pgfpathlineto{\pgfpoint{2.645808\du}{6.079888\du}}
\pgfpathlineto{\pgfpoint{2.644913\du}{6.082276\du}}
\pgfpathlineto{\pgfpoint{2.644913\du}{6.085560\du}}
\pgfpathlineto{\pgfpoint{2.644913\du}{6.085560\du}}
\pgfpathlineto{\pgfpoint{2.644913\du}{6.088844\du}}
\pgfpathlineto{\pgfpoint{2.645808\du}{6.091531\du}}
\pgfpathlineto{\pgfpoint{2.646704\du}{6.094517\du}}
\pgfpathlineto{\pgfpoint{2.649092\du}{6.097801\du}}
\pgfpathlineto{\pgfpoint{2.650884\du}{6.100488\du}}
\pgfpathlineto{\pgfpoint{2.653272\du}{6.103772\du}}
\pgfpathlineto{\pgfpoint{2.656556\du}{6.106459\du}}
\pgfpathlineto{\pgfpoint{2.659840\du}{6.109146\du}}
\pgfpathlineto{\pgfpoint{2.664020\du}{6.111833\du}}
\pgfpathlineto{\pgfpoint{2.668797\du}{6.114520\du}}
\pgfpathlineto{\pgfpoint{2.673872\du}{6.117207\du}}
\pgfpathlineto{\pgfpoint{2.678947\du}{6.119297\du}}
\pgfpathlineto{\pgfpoint{2.684918\du}{6.121984\du}}
\pgfpathlineto{\pgfpoint{2.690890\du}{6.124372\du}}
\pgfpathlineto{\pgfpoint{2.697159\du}{6.126760\du}}
\pgfpathlineto{\pgfpoint{2.704623\du}{6.128850\du}}
\pgfpathlineto{\pgfpoint{2.711191\du}{6.130642\du}}
\pgfpathlineto{\pgfpoint{2.718655\du}{6.132433\du}}
\pgfpathlineto{\pgfpoint{2.726417\du}{6.134224\du}}
\pgfpathlineto{\pgfpoint{2.734478\du}{6.136314\du}}
\pgfpathlineto{\pgfpoint{2.743136\du}{6.137508\du}}
\pgfpathlineto{\pgfpoint{2.751496\du}{6.139300\du}}
\pgfpathlineto{\pgfpoint{2.760154\du}{6.140195\du}}
\pgfpathlineto{\pgfpoint{2.769110\du}{6.141688\du}}
\pgfpathlineto{\pgfpoint{2.778067\du}{6.141987\du}}
\pgfpathlineto{\pgfpoint{2.787620\du}{6.143181\du}}
\pgfpathlineto{\pgfpoint{2.797174\du}{6.143479\du}}
\pgfpathlineto{\pgfpoint{2.806429\du}{6.144375\du}}
\pgfpathlineto{\pgfpoint{2.815983\du}{6.144674\du}}
\pgfpathlineto{\pgfpoint{2.825835\du}{6.144972\du}}
\pgfpathlineto{\pgfpoint{2.834792\du}{6.144972\du}}
\pgfpathlineto{\pgfpoint{2.834792\du}{6.144972\du}}
\pgfpathlineto{\pgfpoint{2.844644\du}{6.144972\du}}
\pgfpathlineto{\pgfpoint{2.854795\du}{6.144674\du}}
\pgfpathlineto{\pgfpoint{2.864348\du}{6.144375\du}}
\pgfpathlineto{\pgfpoint{2.873902\du}{6.143479\du}}
\pgfpathlineto{\pgfpoint{2.883157\du}{6.143181\du}}
\pgfpathlineto{\pgfpoint{2.892711\du}{6.141987\du}}
\pgfpathlineto{\pgfpoint{2.901369\du}{6.141688\du}}
\pgfpathlineto{\pgfpoint{2.910325\du}{6.140195\du}}
\pgfpathlineto{\pgfpoint{2.918983\du}{6.139300\du}}
\pgfpathlineto{\pgfpoint{2.927343\du}{6.137508\du}}
\pgfpathlineto{\pgfpoint{2.936001\du}{6.136314\du}}
\pgfpathlineto{\pgfpoint{2.944360\du}{6.134224\du}}
\pgfpathlineto{\pgfpoint{2.951824\du}{6.132433\du}}
\pgfpathlineto{\pgfpoint{2.959586\du}{6.130642\du}}
\pgfpathlineto{\pgfpoint{2.966453\du}{6.128850\du}}
\pgfpathlineto{\pgfpoint{2.973320\du}{6.126760\du}}
\pgfpathlineto{\pgfpoint{2.979888\du}{6.124372\du}}
\pgfpathlineto{\pgfpoint{2.985560\du}{6.121984\du}}
\pgfpathlineto{\pgfpoint{2.991531\du}{6.119297\du}}
\pgfpathlineto{\pgfpoint{2.997204\du}{6.117207\du}}
\pgfpathlineto{\pgfpoint{3.001682\du}{6.114520\du}}
\pgfpathlineto{\pgfpoint{3.006459\du}{6.111833\du}}
\pgfpathlineto{\pgfpoint{3.010639\du}{6.109146\du}}
\pgfpathlineto{\pgfpoint{3.013923\du}{6.106459\du}}
\pgfpathlineto{\pgfpoint{3.017207\du}{6.103772\du}}
\pgfpathlineto{\pgfpoint{3.019297\du}{6.100488\du}}
\pgfpathlineto{\pgfpoint{3.021387\du}{6.097801\du}}
\pgfpathlineto{\pgfpoint{3.023476\du}{6.094517\du}}
\pgfpathlineto{\pgfpoint{3.024671\du}{6.091531\du}}
\pgfpathlineto{\pgfpoint{3.025566\du}{6.088844\du}}
\pgfpathlineto{\pgfpoint{3.025566\du}{6.085560\du}}
\pgfusepath{stroke}
\pgfsetbuttcap
\pgfsetmiterjoin
\pgfsetdash{}{0pt}
\definecolor{dialinecolor}{rgb}{0.000000, 0.000000, 0.000000}
\pgfsetfillcolor{dialinecolor}
\pgfpathmoveto{\pgfpoint{2.847032\du}{6.073917\du}}
\pgfpathlineto{\pgfpoint{2.870916\du}{6.085560\du}}
\pgfpathlineto{\pgfpoint{2.942270\du}{6.061975\du}}
\pgfpathlineto{\pgfpoint{2.978097\du}{6.073917\du}}
\pgfpathlineto{\pgfpoint{2.954212\du}{6.050033\du}}
\pgfpathlineto{\pgfpoint{2.870916\du}{6.050033\du}}
\pgfpathlineto{\pgfpoint{2.906743\du}{6.050033\du}}
\pgfpathlineto{\pgfpoint{2.847032\du}{6.073917\du}}
\pgfusepath{fill}
\pgfsetbuttcap
\pgfsetmiterjoin
\pgfsetdash{}{0pt}
\definecolor{dialinecolor}{rgb}{0.000000, 0.000000, 0.000000}
\pgfsetfillcolor{dialinecolor}
\pgfpathmoveto{\pgfpoint{2.834792\du}{6.109444\du}}
\pgfpathlineto{\pgfpoint{2.799861\du}{6.097801\du}}
\pgfpathlineto{\pgfpoint{2.739852\du}{6.121685\du}}
\pgfpathlineto{\pgfpoint{2.704623\du}{6.109444\du}}
\pgfpathlineto{\pgfpoint{2.715968\du}{6.144972\du}}
\pgfpathlineto{\pgfpoint{2.799861\du}{6.144972\du}}
\pgfpathlineto{\pgfpoint{2.763736\du}{6.133627\du}}
\pgfpathlineto{\pgfpoint{2.834792\du}{6.109444\du}}
\pgfusepath{fill}
\pgfsetbuttcap
\pgfsetmiterjoin
\pgfsetdash{}{0pt}
\definecolor{dialinecolor}{rgb}{0.000000, 0.000000, 0.000000}
\pgfsetfillcolor{dialinecolor}
\pgfpathmoveto{\pgfpoint{2.715968\du}{6.050033\du}}
\pgfpathlineto{\pgfpoint{2.739852\du}{6.037493\du}}
\pgfpathlineto{\pgfpoint{2.811505\du}{6.073917\du}}
\pgfpathlineto{\pgfpoint{2.834792\du}{6.061975\du}}
\pgfpathlineto{\pgfpoint{2.823148\du}{6.085560\du}}
\pgfpathlineto{\pgfpoint{2.739852\du}{6.085560\du}}
\pgfpathlineto{\pgfpoint{2.775678\du}{6.073917\du}}
\pgfpathlineto{\pgfpoint{2.715968\du}{6.050033\du}}
\pgfusepath{fill}
\pgfsetbuttcap
\pgfsetmiterjoin
\pgfsetdash{}{0pt}
\definecolor{dialinecolor}{rgb}{0.000000, 0.000000, 0.000000}
\pgfsetfillcolor{dialinecolor}
\pgfpathmoveto{\pgfpoint{2.966453\du}{6.133627\du}}
\pgfpathlineto{\pgfpoint{2.942270\du}{6.144972\du}}
\pgfpathlineto{\pgfpoint{2.870916\du}{6.109444\du}}
\pgfpathlineto{\pgfpoint{2.834792\du}{6.121685\du}}
\pgfpathlineto{\pgfpoint{2.859273\du}{6.097801\du}}
\pgfpathlineto{\pgfpoint{2.942270\du}{6.097801\du}}
\pgfpathlineto{\pgfpoint{2.906743\du}{6.109444\du}}
\pgfpathlineto{\pgfpoint{2.966453\du}{6.133627\du}}
\pgfusepath{fill}
\pgfsetbuttcap
\pgfsetmiterjoin
\pgfsetdash{}{0pt}
\definecolor{dialinecolor}{rgb}{1.000000, 1.000000, 1.000000}
\pgfsetfillcolor{dialinecolor}
\pgfpathmoveto{\pgfpoint{2.847032\du}{6.085560\du}}
\pgfpathlineto{\pgfpoint{2.870916\du}{6.085560\du}}
\pgfpathlineto{\pgfpoint{2.942270\du}{6.061975\du}}
\pgfpathlineto{\pgfpoint{2.978097\du}{6.073917\du}}
\pgfpathlineto{\pgfpoint{2.954212\du}{6.050033\du}}
\pgfpathlineto{\pgfpoint{2.883157\du}{6.050033\du}}
\pgfpathlineto{\pgfpoint{2.906743\du}{6.050033\du}}
\pgfpathlineto{\pgfpoint{2.847032\du}{6.085560\du}}
\pgfusepath{fill}
\pgfsetbuttcap
\pgfsetmiterjoin
\pgfsetdash{}{0pt}
\definecolor{dialinecolor}{rgb}{1.000000, 1.000000, 1.000000}
\pgfsetfillcolor{dialinecolor}
\pgfpathmoveto{\pgfpoint{2.834792\du}{6.109444\du}}
\pgfpathlineto{\pgfpoint{2.811505\du}{6.097801\du}}
\pgfpathlineto{\pgfpoint{2.739852\du}{6.121685\du}}
\pgfpathlineto{\pgfpoint{2.704623\du}{6.121685\du}}
\pgfpathlineto{\pgfpoint{2.728209\du}{6.144972\du}}
\pgfpathlineto{\pgfpoint{2.811505\du}{6.144972\du}}
\pgfpathlineto{\pgfpoint{2.775678\du}{6.133627\du}}
\pgfpathlineto{\pgfpoint{2.834792\du}{6.109444\du}}
\pgfusepath{fill}
\pgfsetbuttcap
\pgfsetmiterjoin
\pgfsetdash{}{0pt}
\definecolor{dialinecolor}{rgb}{1.000000, 1.000000, 1.000000}
\pgfsetfillcolor{dialinecolor}
\pgfpathmoveto{\pgfpoint{2.715968\du}{6.050033\du}}
\pgfpathlineto{\pgfpoint{2.739852\du}{6.050033\du}}
\pgfpathlineto{\pgfpoint{2.811505\du}{6.073917\du}}
\pgfpathlineto{\pgfpoint{2.847032\du}{6.061975\du}}
\pgfpathlineto{\pgfpoint{2.823148\du}{6.085560\du}}
\pgfpathlineto{\pgfpoint{2.739852\du}{6.085560\du}}
\pgfpathlineto{\pgfpoint{2.775678\du}{6.085560\du}}
\pgfpathlineto{\pgfpoint{2.715968\du}{6.050033\du}}
\pgfusepath{fill}
\pgfsetbuttcap
\pgfsetmiterjoin
\pgfsetdash{}{0pt}
\definecolor{dialinecolor}{rgb}{1.000000, 1.000000, 1.000000}
\pgfsetfillcolor{dialinecolor}
\pgfpathmoveto{\pgfpoint{2.966453\du}{6.133627\du}}
\pgfpathlineto{\pgfpoint{2.942270\du}{6.144972\du}}
\pgfpathlineto{\pgfpoint{2.870916\du}{6.121685\du}}
\pgfpathlineto{\pgfpoint{2.847032\du}{6.121685\du}}
\pgfpathlineto{\pgfpoint{2.859273\du}{6.097801\du}}
\pgfpathlineto{\pgfpoint{2.942270\du}{6.097801\du}}
\pgfpathlineto{\pgfpoint{2.906743\du}{6.109444\du}}
\pgfpathlineto{\pgfpoint{2.966453\du}{6.133627\du}}
\pgfusepath{fill}
\pgfsetbuttcap
\pgfsetmiterjoin
\pgfsetdash{}{0pt}
\definecolor{dialinecolor}{rgb}{0.666667, 0.901961, 1.000000}
\pgfsetstrokecolor{dialinecolor}
\pgfpathmoveto{\pgfpoint{2.644913\du}{6.085560\du}}
\pgfpathlineto{\pgfpoint{2.644913\du}{6.180500\du}}
\pgfusepath{stroke}
\pgfsetbuttcap
\pgfsetmiterjoin
\pgfsetdash{}{0pt}
\definecolor{dialinecolor}{rgb}{0.666667, 0.901961, 1.000000}
\pgfsetstrokecolor{dialinecolor}
\pgfpathmoveto{\pgfpoint{3.025566\du}{6.085560\du}}
\pgfpathlineto{\pgfpoint{3.025566\du}{6.180500\du}}
\pgfusepath{stroke}
\pgfsetbuttcap
\pgfsetmiterjoin
\pgfsetdash{}{0pt}
\definecolor{dialinecolor}{rgb}{0.000000, 0.470588, 0.666667}
\pgfsetfillcolor{dialinecolor}
\pgfpathmoveto{\pgfpoint{2.680142\du}{6.591606\du}}
\pgfpathlineto{\pgfpoint{2.679843\du}{6.588023\du}}
\pgfpathlineto{\pgfpoint{2.678947\du}{6.584739\du}}
\pgfpathlineto{\pgfpoint{2.677753\du}{6.581455\du}}
\pgfpathlineto{\pgfpoint{2.675962\du}{6.578470\du}}
\pgfpathlineto{\pgfpoint{2.673872\du}{6.575186\du}}
\pgfpathlineto{\pgfpoint{2.670588\du}{6.571901\du}}
\pgfpathlineto{\pgfpoint{2.667901\du}{6.568916\du}}
\pgfpathlineto{\pgfpoint{2.664020\du}{6.565632\du}}
\pgfpathlineto{\pgfpoint{2.659840\du}{6.562945\du}}
\pgfpathlineto{\pgfpoint{2.655362\du}{6.559661\du}}
\pgfpathlineto{\pgfpoint{2.649988\du}{6.556974\du}}
\pgfpathlineto{\pgfpoint{2.644913\du}{6.554287\du}}
\pgfpathlineto{\pgfpoint{2.638941\du}{6.551600\du}}
\pgfpathlineto{\pgfpoint{2.632970\du}{6.548913\du}}
\pgfpathlineto{\pgfpoint{2.625805\du}{6.546525\du}}
\pgfpathlineto{\pgfpoint{2.618938\du}{6.544136\du}}
\pgfpathlineto{\pgfpoint{2.611475\du}{6.542345\du}}
\pgfpathlineto{\pgfpoint{2.604011\du}{6.539956\du}}
\pgfpathlineto{\pgfpoint{2.595950\du}{6.538165\du}}
\pgfpathlineto{\pgfpoint{2.587591\du}{6.535777\du}}
\pgfpathlineto{\pgfpoint{2.579231\du}{6.534284\du}}
\pgfpathlineto{\pgfpoint{2.570275\du}{6.533090\du}}
\pgfpathlineto{\pgfpoint{2.561318\du}{6.531597\du}}
\pgfpathlineto{\pgfpoint{2.552063\du}{6.530403\du}}
\pgfpathlineto{\pgfpoint{2.542211\du}{6.529209\du}}
\pgfpathlineto{\pgfpoint{2.532956\du}{6.528313\du}}
\pgfpathlineto{\pgfpoint{2.523103\du}{6.527716\du}}
\pgfpathlineto{\pgfpoint{2.513550\du}{6.527119\du}}
\pgfpathlineto{\pgfpoint{2.503697\du}{6.526522\du}}
\pgfpathlineto{\pgfpoint{2.493248\du}{6.526223\du}}
\pgfpathlineto{\pgfpoint{2.483993\du}{6.526223\du}}
\pgfpathlineto{\pgfpoint{2.483993\du}{6.526223\du}}
\pgfpathlineto{\pgfpoint{2.474141\du}{6.526223\du}}
\pgfpathlineto{\pgfpoint{2.463990\du}{6.526522\du}}
\pgfpathlineto{\pgfpoint{2.454138\du}{6.527119\du}}
\pgfpathlineto{\pgfpoint{2.444286\du}{6.527716\du}}
\pgfpathlineto{\pgfpoint{2.434433\du}{6.528313\du}}
\pgfpathlineto{\pgfpoint{2.425178\du}{6.529209\du}}
\pgfpathlineto{\pgfpoint{2.415326\du}{6.530403\du}}
\pgfpathlineto{\pgfpoint{2.406071\du}{6.531597\du}}
\pgfpathlineto{\pgfpoint{2.397114\du}{6.533090\du}}
\pgfpathlineto{\pgfpoint{2.388456\du}{6.534284\du}}
\pgfpathlineto{\pgfpoint{2.380097\du}{6.535777\du}}
\pgfpathlineto{\pgfpoint{2.371737\du}{6.538165\du}}
\pgfpathlineto{\pgfpoint{2.363677\du}{6.539956\du}}
\pgfpathlineto{\pgfpoint{2.356213\du}{6.542345\du}}
\pgfpathlineto{\pgfpoint{2.348749\du}{6.544136\du}}
\pgfpathlineto{\pgfpoint{2.341584\du}{6.546525\du}}
\pgfpathlineto{\pgfpoint{2.334717\du}{6.548913\du}}
\pgfpathlineto{\pgfpoint{2.328746\du}{6.551600\du}}
\pgfpathlineto{\pgfpoint{2.322476\du}{6.554287\du}}
\pgfpathlineto{\pgfpoint{2.317401\du}{6.556974\du}}
\pgfpathlineto{\pgfpoint{2.312326\du}{6.559661\du}}
\pgfpathlineto{\pgfpoint{2.307549\du}{6.562945\du}}
\pgfpathlineto{\pgfpoint{2.303369\du}{6.565632\du}}
\pgfpathlineto{\pgfpoint{2.299786\du}{6.568916\du}}
\pgfpathlineto{\pgfpoint{2.296801\du}{6.571901\du}}
\pgfpathlineto{\pgfpoint{2.294114\du}{6.575186\du}}
\pgfpathlineto{\pgfpoint{2.291725\du}{6.578470\du}}
\pgfpathlineto{\pgfpoint{2.289934\du}{6.581455\du}}
\pgfpathlineto{\pgfpoint{2.288441\du}{6.584739\du}}
\pgfpathlineto{\pgfpoint{2.287546\du}{6.588023\du}}
\pgfpathlineto{\pgfpoint{2.287546\du}{6.591606\du}}
\pgfpathlineto{\pgfpoint{2.287546\du}{6.591606\du}}
\pgfpathlineto{\pgfpoint{2.287546\du}{6.594890\du}}
\pgfpathlineto{\pgfpoint{2.288441\du}{6.598174\du}}
\pgfpathlineto{\pgfpoint{2.289934\du}{6.601458\du}}
\pgfpathlineto{\pgfpoint{2.291725\du}{6.604742\du}}
\pgfpathlineto{\pgfpoint{2.294114\du}{6.607728\du}}
\pgfpathlineto{\pgfpoint{2.296801\du}{6.611012\du}}
\pgfpathlineto{\pgfpoint{2.299786\du}{6.613997\du}}
\pgfpathlineto{\pgfpoint{2.303369\du}{6.617281\du}}
\pgfpathlineto{\pgfpoint{2.307549\du}{6.620267\du}}
\pgfpathlineto{\pgfpoint{2.312326\du}{6.623252\du}}
\pgfpathlineto{\pgfpoint{2.317401\du}{6.625939\du}}
\pgfpathlineto{\pgfpoint{2.322476\du}{6.628626\du}}
\pgfpathlineto{\pgfpoint{2.328746\du}{6.631313\du}}
\pgfpathlineto{\pgfpoint{2.334717\du}{6.633702\du}}
\pgfpathlineto{\pgfpoint{2.341584\du}{6.636389\du}}
\pgfpathlineto{\pgfpoint{2.348749\du}{6.638777\du}}
\pgfpathlineto{\pgfpoint{2.356213\du}{6.640867\du}}
\pgfpathlineto{\pgfpoint{2.363677\du}{6.643255\du}}
\pgfpathlineto{\pgfpoint{2.371737\du}{6.644748\du}}
\pgfpathlineto{\pgfpoint{2.380097\du}{6.646838\du}}
\pgfpathlineto{\pgfpoint{2.388456\du}{6.648629\du}}
\pgfpathlineto{\pgfpoint{2.397114\du}{6.649824\du}}
\pgfpathlineto{\pgfpoint{2.406071\du}{6.651316\du}}
\pgfpathlineto{\pgfpoint{2.415326\du}{6.652511\du}}
\pgfpathlineto{\pgfpoint{2.425178\du}{6.653705\du}}
\pgfpathlineto{\pgfpoint{2.434433\du}{6.654600\du}}
\pgfpathlineto{\pgfpoint{2.444286\du}{6.655198\du}}
\pgfpathlineto{\pgfpoint{2.454138\du}{6.656093\du}}
\pgfpathlineto{\pgfpoint{2.463990\du}{6.656093\du}}
\pgfpathlineto{\pgfpoint{2.474141\du}{6.656690\du}}
\pgfpathlineto{\pgfpoint{2.483993\du}{6.656690\du}}
\pgfpathlineto{\pgfpoint{2.483993\du}{6.656690\du}}
\pgfpathlineto{\pgfpoint{2.493248\du}{6.656690\du}}
\pgfpathlineto{\pgfpoint{2.503697\du}{6.656093\du}}
\pgfpathlineto{\pgfpoint{2.513550\du}{6.656093\du}}
\pgfpathlineto{\pgfpoint{2.523103\du}{6.655198\du}}
\pgfpathlineto{\pgfpoint{2.532956\du}{6.654600\du}}
\pgfpathlineto{\pgfpoint{2.542211\du}{6.653705\du}}
\pgfpathlineto{\pgfpoint{2.552063\du}{6.652511\du}}
\pgfpathlineto{\pgfpoint{2.561318\du}{6.651316\du}}
\pgfpathlineto{\pgfpoint{2.570275\du}{6.649824\du}}
\pgfpathlineto{\pgfpoint{2.579231\du}{6.648629\du}}
\pgfpathlineto{\pgfpoint{2.587591\du}{6.646838\du}}
\pgfpathlineto{\pgfpoint{2.595950\du}{6.644748\du}}
\pgfpathlineto{\pgfpoint{2.604011\du}{6.643255\du}}
\pgfpathlineto{\pgfpoint{2.611475\du}{6.640867\du}}
\pgfpathlineto{\pgfpoint{2.618938\du}{6.638777\du}}
\pgfpathlineto{\pgfpoint{2.625805\du}{6.636389\du}}
\pgfpathlineto{\pgfpoint{2.632970\du}{6.633702\du}}
\pgfpathlineto{\pgfpoint{2.638941\du}{6.631313\du}}
\pgfpathlineto{\pgfpoint{2.644913\du}{6.628626\du}}
\pgfpathlineto{\pgfpoint{2.649988\du}{6.625939\du}}
\pgfpathlineto{\pgfpoint{2.655362\du}{6.623252\du}}
\pgfpathlineto{\pgfpoint{2.659840\du}{6.620267\du}}
\pgfpathlineto{\pgfpoint{2.664020\du}{6.617281\du}}
\pgfpathlineto{\pgfpoint{2.667901\du}{6.613997\du}}
\pgfpathlineto{\pgfpoint{2.670588\du}{6.611012\du}}
\pgfpathlineto{\pgfpoint{2.673872\du}{6.607728\du}}
\pgfpathlineto{\pgfpoint{2.675962\du}{6.604742\du}}
\pgfpathlineto{\pgfpoint{2.677753\du}{6.601458\du}}
\pgfpathlineto{\pgfpoint{2.678947\du}{6.598174\du}}
\pgfpathlineto{\pgfpoint{2.679843\du}{6.594890\du}}
\pgfpathlineto{\pgfpoint{2.680142\du}{6.591606\du}}
\pgfusepath{fill}
\pgfsetbuttcap
\pgfsetmiterjoin
\pgfsetdash{}{0pt}
\definecolor{dialinecolor}{rgb}{0.666667, 0.901961, 1.000000}
\pgfsetstrokecolor{dialinecolor}
\pgfpathmoveto{\pgfpoint{2.668200\du}{6.585336\du}}
\pgfpathlineto{\pgfpoint{2.667901\du}{6.582351\du}}
\pgfpathlineto{\pgfpoint{2.667304\du}{6.579664\du}}
\pgfpathlineto{\pgfpoint{2.666110\du}{6.576678\du}}
\pgfpathlineto{\pgfpoint{2.664020\du}{6.573096\du}}
\pgfpathlineto{\pgfpoint{2.661930\du}{6.570707\du}}
\pgfpathlineto{\pgfpoint{2.659840\du}{6.567423\du}}
\pgfpathlineto{\pgfpoint{2.656556\du}{6.564736\du}}
\pgfpathlineto{\pgfpoint{2.653272\du}{6.562049\du}}
\pgfpathlineto{\pgfpoint{2.649092\du}{6.559064\du}}
\pgfpathlineto{\pgfpoint{2.644315\du}{6.556675\du}}
\pgfpathlineto{\pgfpoint{2.639240\du}{6.553988\du}}
\pgfpathlineto{\pgfpoint{2.634165\du}{6.551600\du}}
\pgfpathlineto{\pgfpoint{2.628492\du}{6.549212\du}}
\pgfpathlineto{\pgfpoint{2.622521\du}{6.546823\du}}
\pgfpathlineto{\pgfpoint{2.615953\du}{6.544136\du}}
\pgfpathlineto{\pgfpoint{2.609086\du}{6.542643\du}}
\pgfpathlineto{\pgfpoint{2.602518\du}{6.540554\du}}
\pgfpathlineto{\pgfpoint{2.594457\du}{6.538762\du}}
\pgfpathlineto{\pgfpoint{2.586993\du}{6.537269\du}}
\pgfpathlineto{\pgfpoint{2.579231\du}{6.535180\du}}
\pgfpathlineto{\pgfpoint{2.570275\du}{6.533687\du}}
\pgfpathlineto{\pgfpoint{2.561915\du}{6.532194\du}}
\pgfpathlineto{\pgfpoint{2.553556\du}{6.531000\du}}
\pgfpathlineto{\pgfpoint{2.544300\du}{6.529806\du}}
\pgfpathlineto{\pgfpoint{2.535642\du}{6.528910\du}}
\pgfpathlineto{\pgfpoint{2.526089\du}{6.528014\du}}
\pgfpathlineto{\pgfpoint{2.516535\du}{6.527716\du}}
\pgfpathlineto{\pgfpoint{2.507280\du}{6.527119\du}}
\pgfpathlineto{\pgfpoint{2.497428\du}{6.526522\du}}
\pgfpathlineto{\pgfpoint{2.487576\du}{6.526223\du}}
\pgfpathlineto{\pgfpoint{2.478022\du}{6.526223\du}}
\pgfpathlineto{\pgfpoint{2.478022\du}{6.526223\du}}
\pgfpathlineto{\pgfpoint{2.468767\du}{6.526223\du}}
\pgfpathlineto{\pgfpoint{2.458915\du}{6.526522\du}}
\pgfpathlineto{\pgfpoint{2.449361\du}{6.527119\du}}
\pgfpathlineto{\pgfpoint{2.439807\du}{6.527716\du}}
\pgfpathlineto{\pgfpoint{2.430254\du}{6.528014\du}}
\pgfpathlineto{\pgfpoint{2.420999\du}{6.528910\du}}
\pgfpathlineto{\pgfpoint{2.412042\du}{6.529806\du}}
\pgfpathlineto{\pgfpoint{2.402787\du}{6.531000\du}}
\pgfpathlineto{\pgfpoint{2.394427\du}{6.532194\du}}
\pgfpathlineto{\pgfpoint{2.385471\du}{6.533687\du}}
\pgfpathlineto{\pgfpoint{2.377111\du}{6.535180\du}}
\pgfpathlineto{\pgfpoint{2.369349\du}{6.537269\du}}
\pgfpathlineto{\pgfpoint{2.361288\du}{6.538762\du}}
\pgfpathlineto{\pgfpoint{2.353824\du}{6.540554\du}}
\pgfpathlineto{\pgfpoint{2.347256\du}{6.542643\du}}
\pgfpathlineto{\pgfpoint{2.339792\du}{6.544136\du}}
\pgfpathlineto{\pgfpoint{2.333821\du}{6.546823\du}}
\pgfpathlineto{\pgfpoint{2.327850\du}{6.549212\du}}
\pgfpathlineto{\pgfpoint{2.321581\du}{6.551600\du}}
\pgfpathlineto{\pgfpoint{2.316505\du}{6.553988\du}}
\pgfpathlineto{\pgfpoint{2.311430\du}{6.556675\du}}
\pgfpathlineto{\pgfpoint{2.307250\du}{6.559064\du}}
\pgfpathlineto{\pgfpoint{2.303070\du}{6.562049\du}}
\pgfpathlineto{\pgfpoint{2.299189\du}{6.564736\du}}
\pgfpathlineto{\pgfpoint{2.295905\du}{6.567423\du}}
\pgfpathlineto{\pgfpoint{2.294114\du}{6.570707\du}}
\pgfpathlineto{\pgfpoint{2.291725\du}{6.573096\du}}
\pgfpathlineto{\pgfpoint{2.289934\du}{6.576678\du}}
\pgfpathlineto{\pgfpoint{2.288441\du}{6.579664\du}}
\pgfpathlineto{\pgfpoint{2.288143\du}{6.582351\du}}
\pgfpathlineto{\pgfpoint{2.287546\du}{6.585336\du}}
\pgfpathlineto{\pgfpoint{2.287546\du}{6.585336\du}}
\pgfpathlineto{\pgfpoint{2.288143\du}{6.588919\du}}
\pgfpathlineto{\pgfpoint{2.288441\du}{6.591606\du}}
\pgfpathlineto{\pgfpoint{2.289934\du}{6.594591\du}}
\pgfpathlineto{\pgfpoint{2.291725\du}{6.597577\du}}
\pgfpathlineto{\pgfpoint{2.294114\du}{6.600264\du}}
\pgfpathlineto{\pgfpoint{2.295905\du}{6.603548\du}}
\pgfpathlineto{\pgfpoint{2.299189\du}{6.606235\du}}
\pgfpathlineto{\pgfpoint{2.303070\du}{6.609220\du}}
\pgfpathlineto{\pgfpoint{2.307250\du}{6.611907\du}}
\pgfpathlineto{\pgfpoint{2.311430\du}{6.614594\du}}
\pgfpathlineto{\pgfpoint{2.316505\du}{6.617281\du}}
\pgfpathlineto{\pgfpoint{2.321581\du}{6.618774\du}}
\pgfpathlineto{\pgfpoint{2.327850\du}{6.622058\du}}
\pgfpathlineto{\pgfpoint{2.333821\du}{6.624148\du}}
\pgfpathlineto{\pgfpoint{2.339792\du}{6.626238\du}}
\pgfpathlineto{\pgfpoint{2.347256\du}{6.628626\du}}
\pgfpathlineto{\pgfpoint{2.353824\du}{6.630716\du}}
\pgfpathlineto{\pgfpoint{2.361288\du}{6.632508\du}}
\pgfpathlineto{\pgfpoint{2.369349\du}{6.633702\du}}
\pgfpathlineto{\pgfpoint{2.377111\du}{6.636090\du}}
\pgfpathlineto{\pgfpoint{2.385471\du}{6.637583\du}}
\pgfpathlineto{\pgfpoint{2.394427\du}{6.639076\du}}
\pgfpathlineto{\pgfpoint{2.402787\du}{6.640270\du}}
\pgfpathlineto{\pgfpoint{2.412042\du}{6.641166\du}}
\pgfpathlineto{\pgfpoint{2.420999\du}{6.642061\du}}
\pgfpathlineto{\pgfpoint{2.430254\du}{6.643255\du}}
\pgfpathlineto{\pgfpoint{2.439807\du}{6.643554\du}}
\pgfpathlineto{\pgfpoint{2.449361\du}{6.644151\du}}
\pgfpathlineto{\pgfpoint{2.458915\du}{6.644450\du}}
\pgfpathlineto{\pgfpoint{2.468767\du}{6.644748\du}}
\pgfpathlineto{\pgfpoint{2.478022\du}{6.644748\du}}
\pgfpathlineto{\pgfpoint{2.478022\du}{6.644748\du}}
\pgfpathlineto{\pgfpoint{2.487576\du}{6.644748\du}}
\pgfpathlineto{\pgfpoint{2.497428\du}{6.644450\du}}
\pgfpathlineto{\pgfpoint{2.507280\du}{6.644151\du}}
\pgfpathlineto{\pgfpoint{2.516535\du}{6.643554\du}}
\pgfpathlineto{\pgfpoint{2.526089\du}{6.643255\du}}
\pgfpathlineto{\pgfpoint{2.535642\du}{6.642061\du}}
\pgfpathlineto{\pgfpoint{2.544300\du}{6.641166\du}}
\pgfpathlineto{\pgfpoint{2.553556\du}{6.640270\du}}
\pgfpathlineto{\pgfpoint{2.561915\du}{6.639076\du}}
\pgfpathlineto{\pgfpoint{2.570275\du}{6.637583\du}}
\pgfpathlineto{\pgfpoint{2.579231\du}{6.636090\du}}
\pgfpathlineto{\pgfpoint{2.586993\du}{6.633702\du}}
\pgfpathlineto{\pgfpoint{2.594457\du}{6.632508\du}}
\pgfpathlineto{\pgfpoint{2.602518\du}{6.630716\du}}
\pgfpathlineto{\pgfpoint{2.609086\du}{6.628626\du}}
\pgfpathlineto{\pgfpoint{2.615953\du}{6.626238\du}}
\pgfpathlineto{\pgfpoint{2.622521\du}{6.624148\du}}
\pgfpathlineto{\pgfpoint{2.628492\du}{6.622058\du}}
\pgfpathlineto{\pgfpoint{2.634165\du}{6.618774\du}}
\pgfpathlineto{\pgfpoint{2.639240\du}{6.617281\du}}
\pgfpathlineto{\pgfpoint{2.644315\du}{6.614594\du}}
\pgfpathlineto{\pgfpoint{2.649092\du}{6.611907\du}}
\pgfpathlineto{\pgfpoint{2.653272\du}{6.609220\du}}
\pgfpathlineto{\pgfpoint{2.656556\du}{6.606235\du}}
\pgfpathlineto{\pgfpoint{2.659840\du}{6.603548\du}}
\pgfpathlineto{\pgfpoint{2.661930\du}{6.600264\du}}
\pgfpathlineto{\pgfpoint{2.664020\du}{6.597577\du}}
\pgfpathlineto{\pgfpoint{2.666110\du}{6.594591\du}}
\pgfpathlineto{\pgfpoint{2.667304\du}{6.591606\du}}
\pgfpathlineto{\pgfpoint{2.667901\du}{6.588919\du}}
\pgfpathlineto{\pgfpoint{2.668200\du}{6.585336\du}}
\pgfusepath{stroke}
\pgfsetbuttcap
\pgfsetmiterjoin
\pgfsetdash{}{0pt}
\definecolor{dialinecolor}{rgb}{0.000000, 0.470588, 0.666667}
\pgfsetfillcolor{dialinecolor}
\pgfpathmoveto{\pgfpoint{2.287546\du}{6.502339\du}}
\pgfpathlineto{\pgfpoint{2.287546\du}{6.597577\du}}
\pgfpathlineto{\pgfpoint{2.680440\du}{6.597577\du}}
\pgfpathlineto{\pgfpoint{2.680440\du}{6.502339\du}}
\pgfpathlineto{\pgfpoint{2.287546\du}{6.502339\du}}
\pgfusepath{fill}
\pgfsetbuttcap
\pgfsetmiterjoin
\pgfsetdash{}{0pt}
\definecolor{dialinecolor}{rgb}{0.000000, 0.705882, 1.000000}
\pgfsetfillcolor{dialinecolor}
\pgfpathmoveto{\pgfpoint{2.680142\du}{6.496069\du}}
\pgfpathlineto{\pgfpoint{2.679843\du}{6.492487\du}}
\pgfpathlineto{\pgfpoint{2.678947\du}{6.489203\du}}
\pgfpathlineto{\pgfpoint{2.677753\du}{6.486217\du}}
\pgfpathlineto{\pgfpoint{2.675962\du}{6.482634\du}}
\pgfpathlineto{\pgfpoint{2.673872\du}{6.479350\du}}
\pgfpathlineto{\pgfpoint{2.670588\du}{6.476066\du}}
\pgfpathlineto{\pgfpoint{2.667901\du}{6.473379\du}}
\pgfpathlineto{\pgfpoint{2.664020\du}{6.469498\du}}
\pgfpathlineto{\pgfpoint{2.659840\du}{6.467110\du}}
\pgfpathlineto{\pgfpoint{2.655362\du}{6.464124\du}}
\pgfpathlineto{\pgfpoint{2.649988\du}{6.461139\du}}
\pgfpathlineto{\pgfpoint{2.644913\du}{6.458452\du}}
\pgfpathlineto{\pgfpoint{2.638941\du}{6.455765\du}}
\pgfpathlineto{\pgfpoint{2.632970\du}{6.453376\du}}
\pgfpathlineto{\pgfpoint{2.625805\du}{6.450689\du}}
\pgfpathlineto{\pgfpoint{2.618938\du}{6.448301\du}}
\pgfpathlineto{\pgfpoint{2.611475\du}{6.446211\du}}
\pgfpathlineto{\pgfpoint{2.604011\du}{6.444121\du}}
\pgfpathlineto{\pgfpoint{2.595950\du}{6.442330\du}}
\pgfpathlineto{\pgfpoint{2.587591\du}{6.440240\du}}
\pgfpathlineto{\pgfpoint{2.579231\du}{6.438449\du}}
\pgfpathlineto{\pgfpoint{2.570275\du}{6.437255\du}}
\pgfpathlineto{\pgfpoint{2.561318\du}{6.435762\du}}
\pgfpathlineto{\pgfpoint{2.552063\du}{6.434568\du}}
\pgfpathlineto{\pgfpoint{2.542211\du}{6.433373\du}}
\pgfpathlineto{\pgfpoint{2.532956\du}{6.432179\du}}
\pgfpathlineto{\pgfpoint{2.523103\du}{6.431881\du}}
\pgfpathlineto{\pgfpoint{2.513550\du}{6.430985\du}}
\pgfpathlineto{\pgfpoint{2.503697\du}{6.430686\du}}
\pgfpathlineto{\pgfpoint{2.493248\du}{6.430388\du}}
\pgfpathlineto{\pgfpoint{2.483993\du}{6.430388\du}}
\pgfpathlineto{\pgfpoint{2.483993\du}{6.430388\du}}
\pgfpathlineto{\pgfpoint{2.474141\du}{6.430388\du}}
\pgfpathlineto{\pgfpoint{2.463990\du}{6.430686\du}}
\pgfpathlineto{\pgfpoint{2.454138\du}{6.430985\du}}
\pgfpathlineto{\pgfpoint{2.444286\du}{6.431881\du}}
\pgfpathlineto{\pgfpoint{2.434433\du}{6.432179\du}}
\pgfpathlineto{\pgfpoint{2.425178\du}{6.433373\du}}
\pgfpathlineto{\pgfpoint{2.415326\du}{6.434568\du}}
\pgfpathlineto{\pgfpoint{2.406071\du}{6.435762\du}}
\pgfpathlineto{\pgfpoint{2.397114\du}{6.437255\du}}
\pgfpathlineto{\pgfpoint{2.388456\du}{6.438449\du}}
\pgfpathlineto{\pgfpoint{2.380097\du}{6.440240\du}}
\pgfpathlineto{\pgfpoint{2.371737\du}{6.442330\du}}
\pgfpathlineto{\pgfpoint{2.363677\du}{6.444121\du}}
\pgfpathlineto{\pgfpoint{2.356213\du}{6.446211\du}}
\pgfpathlineto{\pgfpoint{2.348749\du}{6.448301\du}}
\pgfpathlineto{\pgfpoint{2.341584\du}{6.450689\du}}
\pgfpathlineto{\pgfpoint{2.334717\du}{6.453376\du}}
\pgfpathlineto{\pgfpoint{2.328746\du}{6.455765\du}}
\pgfpathlineto{\pgfpoint{2.322476\du}{6.458452\du}}
\pgfpathlineto{\pgfpoint{2.317401\du}{6.461139\du}}
\pgfpathlineto{\pgfpoint{2.312326\du}{6.464124\du}}
\pgfpathlineto{\pgfpoint{2.307549\du}{6.467110\du}}
\pgfpathlineto{\pgfpoint{2.303369\du}{6.469498\du}}
\pgfpathlineto{\pgfpoint{2.299786\du}{6.473379\du}}
\pgfpathlineto{\pgfpoint{2.296801\du}{6.476066\du}}
\pgfpathlineto{\pgfpoint{2.294114\du}{6.479350\du}}
\pgfpathlineto{\pgfpoint{2.291725\du}{6.482634\du}}
\pgfpathlineto{\pgfpoint{2.289934\du}{6.486217\du}}
\pgfpathlineto{\pgfpoint{2.288441\du}{6.489203\du}}
\pgfpathlineto{\pgfpoint{2.287546\du}{6.492487\du}}
\pgfpathlineto{\pgfpoint{2.287546\du}{6.496069\du}}
\pgfpathlineto{\pgfpoint{2.287546\du}{6.496069\du}}
\pgfpathlineto{\pgfpoint{2.287546\du}{6.499353\du}}
\pgfpathlineto{\pgfpoint{2.288441\du}{6.502637\du}}
\pgfpathlineto{\pgfpoint{2.289934\du}{6.505623\du}}
\pgfpathlineto{\pgfpoint{2.291725\du}{6.509206\du}}
\pgfpathlineto{\pgfpoint{2.294114\du}{6.512490\du}}
\pgfpathlineto{\pgfpoint{2.296801\du}{6.515475\du}}
\pgfpathlineto{\pgfpoint{2.299786\du}{6.518461\du}}
\pgfpathlineto{\pgfpoint{2.303369\du}{6.521745\du}}
\pgfpathlineto{\pgfpoint{2.307549\du}{6.524730\du}}
\pgfpathlineto{\pgfpoint{2.312326\du}{6.527716\du}}
\pgfpathlineto{\pgfpoint{2.317401\du}{6.530701\du}}
\pgfpathlineto{\pgfpoint{2.322476\du}{6.533388\du}}
\pgfpathlineto{\pgfpoint{2.328746\du}{6.535777\du}}
\pgfpathlineto{\pgfpoint{2.334717\du}{6.538464\du}}
\pgfpathlineto{\pgfpoint{2.341584\du}{6.541151\du}}
\pgfpathlineto{\pgfpoint{2.348749\du}{6.543240\du}}
\pgfpathlineto{\pgfpoint{2.356213\du}{6.545629\du}}
\pgfpathlineto{\pgfpoint{2.363677\du}{6.547719\du}}
\pgfpathlineto{\pgfpoint{2.371737\du}{6.549510\du}}
\pgfpathlineto{\pgfpoint{2.380097\du}{6.551600\du}}
\pgfpathlineto{\pgfpoint{2.388456\du}{6.553391\du}}
\pgfpathlineto{\pgfpoint{2.397114\du}{6.554585\du}}
\pgfpathlineto{\pgfpoint{2.406071\du}{6.556078\du}}
\pgfpathlineto{\pgfpoint{2.415326\du}{6.557272\du}}
\pgfpathlineto{\pgfpoint{2.425178\du}{6.558168\du}}
\pgfpathlineto{\pgfpoint{2.434433\du}{6.559064\du}}
\pgfpathlineto{\pgfpoint{2.444286\du}{6.559959\du}}
\pgfpathlineto{\pgfpoint{2.454138\du}{6.560855\du}}
\pgfpathlineto{\pgfpoint{2.463990\du}{6.561154\du}}
\pgfpathlineto{\pgfpoint{2.474141\du}{6.561452\du}}
\pgfpathlineto{\pgfpoint{2.483993\du}{6.561452\du}}
\pgfpathlineto{\pgfpoint{2.483993\du}{6.561452\du}}
\pgfpathlineto{\pgfpoint{2.493248\du}{6.561452\du}}
\pgfpathlineto{\pgfpoint{2.503697\du}{6.561154\du}}
\pgfpathlineto{\pgfpoint{2.513550\du}{6.560855\du}}
\pgfpathlineto{\pgfpoint{2.523103\du}{6.559959\du}}
\pgfpathlineto{\pgfpoint{2.532956\du}{6.559064\du}}
\pgfpathlineto{\pgfpoint{2.542211\du}{6.558168\du}}
\pgfpathlineto{\pgfpoint{2.552063\du}{6.557272\du}}
\pgfpathlineto{\pgfpoint{2.561318\du}{6.556078\du}}
\pgfpathlineto{\pgfpoint{2.570275\du}{6.554585\du}}
\pgfpathlineto{\pgfpoint{2.579231\du}{6.553391\du}}
\pgfpathlineto{\pgfpoint{2.587591\du}{6.551600\du}}
\pgfpathlineto{\pgfpoint{2.595950\du}{6.549510\du}}
\pgfpathlineto{\pgfpoint{2.604011\du}{6.547719\du}}
\pgfpathlineto{\pgfpoint{2.611475\du}{6.545629\du}}
\pgfpathlineto{\pgfpoint{2.618938\du}{6.543240\du}}
\pgfpathlineto{\pgfpoint{2.625805\du}{6.541151\du}}
\pgfpathlineto{\pgfpoint{2.632970\du}{6.538464\du}}
\pgfpathlineto{\pgfpoint{2.638941\du}{6.535777\du}}
\pgfpathlineto{\pgfpoint{2.644913\du}{6.533388\du}}
\pgfpathlineto{\pgfpoint{2.649988\du}{6.530701\du}}
\pgfpathlineto{\pgfpoint{2.655362\du}{6.527716\du}}
\pgfpathlineto{\pgfpoint{2.659840\du}{6.524730\du}}
\pgfpathlineto{\pgfpoint{2.664020\du}{6.521745\du}}
\pgfpathlineto{\pgfpoint{2.667901\du}{6.518461\du}}
\pgfpathlineto{\pgfpoint{2.670588\du}{6.515475\du}}
\pgfpathlineto{\pgfpoint{2.673872\du}{6.512490\du}}
\pgfpathlineto{\pgfpoint{2.675962\du}{6.509206\du}}
\pgfpathlineto{\pgfpoint{2.677753\du}{6.505623\du}}
\pgfpathlineto{\pgfpoint{2.678947\du}{6.502637\du}}
\pgfpathlineto{\pgfpoint{2.679843\du}{6.499353\du}}
\pgfpathlineto{\pgfpoint{2.680142\du}{6.496069\du}}
\pgfusepath{fill}
\pgfsetbuttcap
\pgfsetmiterjoin
\pgfsetdash{}{0pt}
\definecolor{dialinecolor}{rgb}{0.666667, 0.901961, 1.000000}
\pgfsetstrokecolor{dialinecolor}
\pgfpathmoveto{\pgfpoint{2.668200\du}{6.490397\du}}
\pgfpathlineto{\pgfpoint{2.667901\du}{6.487411\du}}
\pgfpathlineto{\pgfpoint{2.667304\du}{6.484426\du}}
\pgfpathlineto{\pgfpoint{2.666110\du}{6.481739\du}}
\pgfpathlineto{\pgfpoint{2.664020\du}{6.478455\du}}
\pgfpathlineto{\pgfpoint{2.661930\du}{6.475768\du}}
\pgfpathlineto{\pgfpoint{2.659840\du}{6.472484\du}}
\pgfpathlineto{\pgfpoint{2.656556\du}{6.469498\du}}
\pgfpathlineto{\pgfpoint{2.653272\du}{6.467110\du}}
\pgfpathlineto{\pgfpoint{2.649092\du}{6.464423\du}}
\pgfpathlineto{\pgfpoint{2.644315\du}{6.461736\du}}
\pgfpathlineto{\pgfpoint{2.639240\du}{6.459049\du}}
\pgfpathlineto{\pgfpoint{2.634165\du}{6.456660\du}}
\pgfpathlineto{\pgfpoint{2.628492\du}{6.453675\du}}
\pgfpathlineto{\pgfpoint{2.622521\du}{6.451884\du}}
\pgfpathlineto{\pgfpoint{2.615953\du}{6.449794\du}}
\pgfpathlineto{\pgfpoint{2.609086\du}{6.447107\du}}
\pgfpathlineto{\pgfpoint{2.602518\du}{6.445614\du}}
\pgfpathlineto{\pgfpoint{2.594457\du}{6.443823\du}}
\pgfpathlineto{\pgfpoint{2.586993\du}{6.441733\du}}
\pgfpathlineto{\pgfpoint{2.579231\du}{6.440240\du}}
\pgfpathlineto{\pgfpoint{2.570275\du}{6.438747\du}}
\pgfpathlineto{\pgfpoint{2.561915\du}{6.437255\du}}
\pgfpathlineto{\pgfpoint{2.553556\du}{6.436060\du}}
\pgfpathlineto{\pgfpoint{2.544300\du}{6.434866\du}}
\pgfpathlineto{\pgfpoint{2.535642\du}{6.433672\du}}
\pgfpathlineto{\pgfpoint{2.526089\du}{6.433075\du}}
\pgfpathlineto{\pgfpoint{2.516535\du}{6.432179\du}}
\pgfpathlineto{\pgfpoint{2.507280\du}{6.432179\du}}
\pgfpathlineto{\pgfpoint{2.497428\du}{6.431881\du}}
\pgfpathlineto{\pgfpoint{2.487576\du}{6.430985\du}}
\pgfpathlineto{\pgfpoint{2.478022\du}{6.430985\du}}
\pgfpathlineto{\pgfpoint{2.478022\du}{6.430985\du}}
\pgfpathlineto{\pgfpoint{2.468767\du}{6.430985\du}}
\pgfpathlineto{\pgfpoint{2.458915\du}{6.431881\du}}
\pgfpathlineto{\pgfpoint{2.449361\du}{6.432179\du}}
\pgfpathlineto{\pgfpoint{2.439807\du}{6.432179\du}}
\pgfpathlineto{\pgfpoint{2.430254\du}{6.433075\du}}
\pgfpathlineto{\pgfpoint{2.420999\du}{6.433672\du}}
\pgfpathlineto{\pgfpoint{2.412042\du}{6.434866\du}}
\pgfpathlineto{\pgfpoint{2.402787\du}{6.436060\du}}
\pgfpathlineto{\pgfpoint{2.394427\du}{6.437255\du}}
\pgfpathlineto{\pgfpoint{2.385471\du}{6.438747\du}}
\pgfpathlineto{\pgfpoint{2.377111\du}{6.440240\du}}
\pgfpathlineto{\pgfpoint{2.369349\du}{6.441733\du}}
\pgfpathlineto{\pgfpoint{2.361288\du}{6.443823\du}}
\pgfpathlineto{\pgfpoint{2.353824\du}{6.445614\du}}
\pgfpathlineto{\pgfpoint{2.347256\du}{6.447107\du}}
\pgfpathlineto{\pgfpoint{2.339792\du}{6.449794\du}}
\pgfpathlineto{\pgfpoint{2.333821\du}{6.451884\du}}
\pgfpathlineto{\pgfpoint{2.327850\du}{6.453675\du}}
\pgfpathlineto{\pgfpoint{2.321581\du}{6.456660\du}}
\pgfpathlineto{\pgfpoint{2.316505\du}{6.459049\du}}
\pgfpathlineto{\pgfpoint{2.311430\du}{6.461736\du}}
\pgfpathlineto{\pgfpoint{2.307250\du}{6.464423\du}}
\pgfpathlineto{\pgfpoint{2.303070\du}{6.467110\du}}
\pgfpathlineto{\pgfpoint{2.299189\du}{6.469498\du}}
\pgfpathlineto{\pgfpoint{2.295905\du}{6.472484\du}}
\pgfpathlineto{\pgfpoint{2.294114\du}{6.475768\du}}
\pgfpathlineto{\pgfpoint{2.291725\du}{6.478455\du}}
\pgfpathlineto{\pgfpoint{2.289934\du}{6.481739\du}}
\pgfpathlineto{\pgfpoint{2.288441\du}{6.484426\du}}
\pgfpathlineto{\pgfpoint{2.288143\du}{6.487411\du}}
\pgfpathlineto{\pgfpoint{2.287546\du}{6.490397\du}}
\pgfpathlineto{\pgfpoint{2.287546\du}{6.490397\du}}
\pgfpathlineto{\pgfpoint{2.288143\du}{6.493979\du}}
\pgfpathlineto{\pgfpoint{2.288441\du}{6.496666\du}}
\pgfpathlineto{\pgfpoint{2.289934\du}{6.499353\du}}
\pgfpathlineto{\pgfpoint{2.291725\du}{6.502637\du}}
\pgfpathlineto{\pgfpoint{2.294114\du}{6.505324\du}}
\pgfpathlineto{\pgfpoint{2.295905\du}{6.508907\du}}
\pgfpathlineto{\pgfpoint{2.299189\du}{6.511594\du}}
\pgfpathlineto{\pgfpoint{2.303070\du}{6.514281\du}}
\pgfpathlineto{\pgfpoint{2.307250\du}{6.516968\du}}
\pgfpathlineto{\pgfpoint{2.311430\du}{6.519655\du}}
\pgfpathlineto{\pgfpoint{2.316505\du}{6.522342\du}}
\pgfpathlineto{\pgfpoint{2.321581\du}{6.524432\du}}
\pgfpathlineto{\pgfpoint{2.327850\du}{6.527119\du}}
\pgfpathlineto{\pgfpoint{2.333821\du}{6.529209\du}}
\pgfpathlineto{\pgfpoint{2.339792\du}{6.531597\du}}
\pgfpathlineto{\pgfpoint{2.347256\du}{6.533687\du}}
\pgfpathlineto{\pgfpoint{2.353824\du}{6.535777\du}}
\pgfpathlineto{\pgfpoint{2.361288\du}{6.537568\du}}
\pgfpathlineto{\pgfpoint{2.369349\du}{6.539061\du}}
\pgfpathlineto{\pgfpoint{2.377111\du}{6.541151\du}}
\pgfpathlineto{\pgfpoint{2.385471\du}{6.542643\du}}
\pgfpathlineto{\pgfpoint{2.394427\du}{6.544136\du}}
\pgfpathlineto{\pgfpoint{2.402787\du}{6.545330\du}}
\pgfpathlineto{\pgfpoint{2.412042\du}{6.546525\du}}
\pgfpathlineto{\pgfpoint{2.420999\du}{6.547122\du}}
\pgfpathlineto{\pgfpoint{2.430254\du}{6.548316\du}}
\pgfpathlineto{\pgfpoint{2.439807\du}{6.548614\du}}
\pgfpathlineto{\pgfpoint{2.449361\du}{6.549212\du}}
\pgfpathlineto{\pgfpoint{2.458915\du}{6.549510\du}}
\pgfpathlineto{\pgfpoint{2.468767\du}{6.550107\du}}
\pgfpathlineto{\pgfpoint{2.478022\du}{6.550107\du}}
\pgfpathlineto{\pgfpoint{2.478022\du}{6.550107\du}}
\pgfpathlineto{\pgfpoint{2.487576\du}{6.550107\du}}
\pgfpathlineto{\pgfpoint{2.497428\du}{6.549510\du}}
\pgfpathlineto{\pgfpoint{2.507280\du}{6.549212\du}}
\pgfpathlineto{\pgfpoint{2.516535\du}{6.548614\du}}
\pgfpathlineto{\pgfpoint{2.526089\du}{6.548316\du}}
\pgfpathlineto{\pgfpoint{2.535642\du}{6.547122\du}}
\pgfpathlineto{\pgfpoint{2.544300\du}{6.546525\du}}
\pgfpathlineto{\pgfpoint{2.553556\du}{6.545330\du}}
\pgfpathlineto{\pgfpoint{2.561915\du}{6.544136\du}}
\pgfpathlineto{\pgfpoint{2.570275\du}{6.542643\du}}
\pgfpathlineto{\pgfpoint{2.579231\du}{6.541151\du}}
\pgfpathlineto{\pgfpoint{2.586993\du}{6.539061\du}}
\pgfpathlineto{\pgfpoint{2.594457\du}{6.537568\du}}
\pgfpathlineto{\pgfpoint{2.602518\du}{6.535777\du}}
\pgfpathlineto{\pgfpoint{2.609086\du}{6.533687\du}}
\pgfpathlineto{\pgfpoint{2.615953\du}{6.531597\du}}
\pgfpathlineto{\pgfpoint{2.622521\du}{6.529209\du}}
\pgfpathlineto{\pgfpoint{2.628492\du}{6.527119\du}}
\pgfpathlineto{\pgfpoint{2.634165\du}{6.524432\du}}
\pgfpathlineto{\pgfpoint{2.639240\du}{6.522342\du}}
\pgfpathlineto{\pgfpoint{2.644315\du}{6.519655\du}}
\pgfpathlineto{\pgfpoint{2.649092\du}{6.516968\du}}
\pgfpathlineto{\pgfpoint{2.653272\du}{6.514281\du}}
\pgfpathlineto{\pgfpoint{2.656556\du}{6.511594\du}}
\pgfpathlineto{\pgfpoint{2.659840\du}{6.508907\du}}
\pgfpathlineto{\pgfpoint{2.661930\du}{6.505324\du}}
\pgfpathlineto{\pgfpoint{2.664020\du}{6.502637\du}}
\pgfpathlineto{\pgfpoint{2.666110\du}{6.499353\du}}
\pgfpathlineto{\pgfpoint{2.667304\du}{6.496666\du}}
\pgfpathlineto{\pgfpoint{2.667901\du}{6.493979\du}}
\pgfpathlineto{\pgfpoint{2.668200\du}{6.490397\du}}
\pgfusepath{stroke}
\pgfsetbuttcap
\pgfsetmiterjoin
\pgfsetdash{}{0pt}
\definecolor{dialinecolor}{rgb}{0.000000, 0.000000, 0.000000}
\pgfsetfillcolor{dialinecolor}
\pgfpathmoveto{\pgfpoint{2.489964\du}{6.478455\du}}
\pgfpathlineto{\pgfpoint{2.513848\du}{6.490397\du}}
\pgfpathlineto{\pgfpoint{2.585202\du}{6.466811\du}}
\pgfpathlineto{\pgfpoint{2.621028\du}{6.478455\du}}
\pgfpathlineto{\pgfpoint{2.597443\du}{6.454571\du}}
\pgfpathlineto{\pgfpoint{2.513848\du}{6.454571\du}}
\pgfpathlineto{\pgfpoint{2.549674\du}{6.454571\du}}
\pgfpathlineto{\pgfpoint{2.489964\du}{6.478455\du}}
\pgfusepath{fill}
\pgfsetbuttcap
\pgfsetmiterjoin
\pgfsetdash{}{0pt}
\definecolor{dialinecolor}{rgb}{0.000000, 0.000000, 0.000000}
\pgfsetfillcolor{dialinecolor}
\pgfpathmoveto{\pgfpoint{2.478022\du}{6.514281\du}}
\pgfpathlineto{\pgfpoint{2.442494\du}{6.502339\du}}
\pgfpathlineto{\pgfpoint{2.382784\du}{6.526223\du}}
\pgfpathlineto{\pgfpoint{2.347256\du}{6.514281\du}}
\pgfpathlineto{\pgfpoint{2.359497\du}{6.550107\du}}
\pgfpathlineto{\pgfpoint{2.442494\du}{6.550107\du}}
\pgfpathlineto{\pgfpoint{2.406967\du}{6.538464\du}}
\pgfpathlineto{\pgfpoint{2.478022\du}{6.514281\du}}
\pgfusepath{fill}
\pgfsetbuttcap
\pgfsetmiterjoin
\pgfsetdash{}{0pt}
\definecolor{dialinecolor}{rgb}{0.000000, 0.000000, 0.000000}
\pgfsetfillcolor{dialinecolor}
\pgfpathmoveto{\pgfpoint{2.359497\du}{6.454571\du}}
\pgfpathlineto{\pgfpoint{2.382784\du}{6.442927\du}}
\pgfpathlineto{\pgfpoint{2.454138\du}{6.478455\du}}
\pgfpathlineto{\pgfpoint{2.478022\du}{6.466811\du}}
\pgfpathlineto{\pgfpoint{2.466677\du}{6.490397\du}}
\pgfpathlineto{\pgfpoint{2.382784\du}{6.490397\du}}
\pgfpathlineto{\pgfpoint{2.418610\du}{6.478455\du}}
\pgfpathlineto{\pgfpoint{2.359497\du}{6.454571\du}}
\pgfusepath{fill}
\pgfsetbuttcap
\pgfsetmiterjoin
\pgfsetdash{}{0pt}
\definecolor{dialinecolor}{rgb}{0.000000, 0.000000, 0.000000}
\pgfsetfillcolor{dialinecolor}
\pgfpathmoveto{\pgfpoint{2.609086\du}{6.538464\du}}
\pgfpathlineto{\pgfpoint{2.585202\du}{6.550107\du}}
\pgfpathlineto{\pgfpoint{2.513848\du}{6.514281\du}}
\pgfpathlineto{\pgfpoint{2.478022\du}{6.526223\du}}
\pgfpathlineto{\pgfpoint{2.502205\du}{6.502339\du}}
\pgfpathlineto{\pgfpoint{2.585202\du}{6.502339\du}}
\pgfpathlineto{\pgfpoint{2.549674\du}{6.514281\du}}
\pgfpathlineto{\pgfpoint{2.609086\du}{6.538464\du}}
\pgfusepath{fill}
\pgfsetbuttcap
\pgfsetmiterjoin
\pgfsetdash{}{0pt}
\definecolor{dialinecolor}{rgb}{1.000000, 1.000000, 1.000000}
\pgfsetfillcolor{dialinecolor}
\pgfpathmoveto{\pgfpoint{2.489964\du}{6.490397\du}}
\pgfpathlineto{\pgfpoint{2.513848\du}{6.490397\du}}
\pgfpathlineto{\pgfpoint{2.585202\du}{6.466811\du}}
\pgfpathlineto{\pgfpoint{2.621028\du}{6.478455\du}}
\pgfpathlineto{\pgfpoint{2.597443\du}{6.454571\du}}
\pgfpathlineto{\pgfpoint{2.526089\du}{6.454571\du}}
\pgfpathlineto{\pgfpoint{2.549674\du}{6.454571\du}}
\pgfpathlineto{\pgfpoint{2.489964\du}{6.490397\du}}
\pgfusepath{fill}
\pgfsetbuttcap
\pgfsetmiterjoin
\pgfsetdash{}{0pt}
\definecolor{dialinecolor}{rgb}{1.000000, 1.000000, 1.000000}
\pgfsetfillcolor{dialinecolor}
\pgfpathmoveto{\pgfpoint{2.478022\du}{6.514281\du}}
\pgfpathlineto{\pgfpoint{2.454138\du}{6.502339\du}}
\pgfpathlineto{\pgfpoint{2.382784\du}{6.526223\du}}
\pgfpathlineto{\pgfpoint{2.347256\du}{6.526223\du}}
\pgfpathlineto{\pgfpoint{2.371140\du}{6.550107\du}}
\pgfpathlineto{\pgfpoint{2.454138\du}{6.550107\du}}
\pgfpathlineto{\pgfpoint{2.418610\du}{6.538464\du}}
\pgfpathlineto{\pgfpoint{2.478022\du}{6.514281\du}}
\pgfusepath{fill}
\pgfsetbuttcap
\pgfsetmiterjoin
\pgfsetdash{}{0pt}
\definecolor{dialinecolor}{rgb}{1.000000, 1.000000, 1.000000}
\pgfsetfillcolor{dialinecolor}
\pgfpathmoveto{\pgfpoint{2.359497\du}{6.454571\du}}
\pgfpathlineto{\pgfpoint{2.382784\du}{6.454571\du}}
\pgfpathlineto{\pgfpoint{2.454138\du}{6.478455\du}}
\pgfpathlineto{\pgfpoint{2.489964\du}{6.466811\du}}
\pgfpathlineto{\pgfpoint{2.466677\du}{6.490397\du}}
\pgfpathlineto{\pgfpoint{2.382784\du}{6.490397\du}}
\pgfpathlineto{\pgfpoint{2.418610\du}{6.490397\du}}
\pgfpathlineto{\pgfpoint{2.359497\du}{6.454571\du}}
\pgfusepath{fill}
\pgfsetbuttcap
\pgfsetmiterjoin
\pgfsetdash{}{0pt}
\definecolor{dialinecolor}{rgb}{1.000000, 1.000000, 1.000000}
\pgfsetfillcolor{dialinecolor}
\pgfpathmoveto{\pgfpoint{2.609086\du}{6.538464\du}}
\pgfpathlineto{\pgfpoint{2.585202\du}{6.550107\du}}
\pgfpathlineto{\pgfpoint{2.513848\du}{6.526223\du}}
\pgfpathlineto{\pgfpoint{2.489964\du}{6.526223\du}}
\pgfpathlineto{\pgfpoint{2.502205\du}{6.502339\du}}
\pgfpathlineto{\pgfpoint{2.585202\du}{6.502339\du}}
\pgfpathlineto{\pgfpoint{2.549674\du}{6.514281\du}}
\pgfpathlineto{\pgfpoint{2.609086\du}{6.538464\du}}
\pgfusepath{fill}
\pgfsetbuttcap
\pgfsetmiterjoin
\pgfsetdash{}{0pt}
\definecolor{dialinecolor}{rgb}{0.666667, 0.901961, 1.000000}
\pgfsetstrokecolor{dialinecolor}
\pgfpathmoveto{\pgfpoint{2.287546\du}{6.490397\du}}
\pgfpathlineto{\pgfpoint{2.287546\du}{6.585336\du}}
\pgfusepath{stroke}
\pgfsetbuttcap
\pgfsetmiterjoin
\pgfsetdash{}{0pt}
\definecolor{dialinecolor}{rgb}{0.666667, 0.901961, 1.000000}
\pgfsetstrokecolor{dialinecolor}
\pgfpathmoveto{\pgfpoint{2.668200\du}{6.490397\du}}
\pgfpathlineto{\pgfpoint{2.668200\du}{6.585336\du}}
\pgfusepath{stroke}
\pgfsetlinewidth{0.100000\du}
\pgfsetdash{}{0pt}
\pgfsetdash{}{0pt}
\pgfsetbuttcap
\pgfsetmiterjoin
\pgfsetlinewidth{0.001000\du}
\pgfsetbuttcap
\pgfsetmiterjoin
\pgfsetdash{}{0pt}
\definecolor{dialinecolor}{rgb}{0.717647, 0.717647, 0.615686}
\pgfsetfillcolor{dialinecolor}
\pgfpathmoveto{\pgfpoint{-1.353800\du}{7.161780\du}}
\pgfpathlineto{\pgfpoint{0.644626\du}{7.161780\du}}
\pgfpathlineto{\pgfpoint{0.644626\du}{7.531113\du}}
\pgfpathlineto{\pgfpoint{-1.353800\du}{7.531113\du}}
\pgfpathlineto{\pgfpoint{-1.353800\du}{7.161780\du}}
\pgfusepath{fill}
\pgfsetbuttcap
\pgfsetmiterjoin
\pgfsetdash{}{0pt}
\definecolor{dialinecolor}{rgb}{0.286275, 0.286275, 0.211765}
\pgfsetstrokecolor{dialinecolor}
\pgfpathmoveto{\pgfpoint{-1.353800\du}{7.161780\du}}
\pgfpathlineto{\pgfpoint{0.644626\du}{7.161780\du}}
\pgfpathlineto{\pgfpoint{0.644626\du}{7.531113\du}}
\pgfpathlineto{\pgfpoint{-1.353800\du}{7.531113\du}}
\pgfpathlineto{\pgfpoint{-1.353800\du}{7.161780\du}}
\pgfusepath{stroke}
\pgfsetbuttcap
\pgfsetmiterjoin
\pgfsetdash{}{0pt}
\definecolor{dialinecolor}{rgb}{0.788235, 0.788235, 0.713726}
\pgfsetfillcolor{dialinecolor}
\pgfpathmoveto{\pgfpoint{-1.353800\du}{7.161780\du}}
\pgfpathlineto{\pgfpoint{-1.141898\du}{6.960898\du}}
\pgfpathlineto{\pgfpoint{0.856527\du}{6.960898\du}}
\pgfpathlineto{\pgfpoint{0.644626\du}{7.161780\du}}
\pgfpathlineto{\pgfpoint{-1.353800\du}{7.161780\du}}
\pgfusepath{fill}
\pgfsetbuttcap
\pgfsetmiterjoin
\pgfsetdash{}{0pt}
\definecolor{dialinecolor}{rgb}{0.286275, 0.286275, 0.211765}
\pgfsetstrokecolor{dialinecolor}
\pgfpathmoveto{\pgfpoint{-1.353800\du}{7.161780\du}}
\pgfpathlineto{\pgfpoint{-1.141898\du}{6.960898\du}}
\pgfpathlineto{\pgfpoint{0.856527\du}{6.960898\du}}
\pgfpathlineto{\pgfpoint{0.644626\du}{7.161780\du}}
\pgfpathlineto{\pgfpoint{-1.353800\du}{7.161780\du}}
\pgfusepath{stroke}
\pgfsetlinewidth{0.106000\du}
\pgfsetbuttcap
\pgfsetmiterjoin
\pgfsetdash{}{0pt}
\definecolor{dialinecolor}{rgb}{0.000000, 0.000000, 0.000000}
\pgfsetstrokecolor{dialinecolor}
\pgfpathmoveto{\pgfpoint{0.532850\du}{7.329916\du}}
\pgfpathlineto{\pgfpoint{0.053316\du}{7.329916\du}}
\pgfusepath{stroke}
\pgfsetlinewidth{0.001000\du}
\pgfsetbuttcap
\pgfsetmiterjoin
\pgfsetdash{}{0pt}
\definecolor{dialinecolor}{rgb}{0.478431, 0.478431, 0.352941}
\pgfsetfillcolor{dialinecolor}
\pgfpathmoveto{\pgfpoint{0.644626\du}{7.531113\du}}
\pgfpathlineto{\pgfpoint{0.856527\du}{7.318581\du}}
\pgfpathlineto{\pgfpoint{0.856527\du}{6.960898\du}}
\pgfpathlineto{\pgfpoint{0.644626\du}{7.161780\du}}
\pgfpathlineto{\pgfpoint{0.644626\du}{7.531113\du}}
\pgfusepath{fill}
\pgfsetbuttcap
\pgfsetmiterjoin
\pgfsetdash{}{0pt}
\definecolor{dialinecolor}{rgb}{0.286275, 0.286275, 0.211765}
\pgfsetstrokecolor{dialinecolor}
\pgfpathmoveto{\pgfpoint{0.644626\du}{7.531113\du}}
\pgfpathlineto{\pgfpoint{0.856527\du}{7.318581\du}}
\pgfpathlineto{\pgfpoint{0.856527\du}{6.960898\du}}
\pgfpathlineto{\pgfpoint{0.644626\du}{7.161780\du}}
\pgfpathlineto{\pgfpoint{0.644626\du}{7.531113\du}}
\pgfusepath{stroke}
\pgfsetbuttcap
\pgfsetmiterjoin
\pgfsetdash{}{0pt}
\definecolor{dialinecolor}{rgb}{0.788235, 0.788235, 0.713726}
\pgfsetfillcolor{dialinecolor}
\pgfpathmoveto{\pgfpoint{-1.342465\du}{7.765370\du}}
\pgfpathlineto{\pgfpoint{-1.119543\du}{7.486087\du}}
\pgfpathlineto{\pgfpoint{0.421389\du}{7.486087\du}}
\pgfpathlineto{\pgfpoint{0.198467\du}{7.765370\du}}
\pgfpathlineto{\pgfpoint{-1.342465\du}{7.765370\du}}
\pgfusepath{fill}
\pgfsetbuttcap
\pgfsetmiterjoin
\pgfsetdash{}{0pt}
\definecolor{dialinecolor}{rgb}{0.286275, 0.286275, 0.211765}
\pgfsetstrokecolor{dialinecolor}
\pgfpathmoveto{\pgfpoint{-1.342465\du}{7.765370\du}}
\pgfpathlineto{\pgfpoint{-1.119543\du}{7.486087\du}}
\pgfpathlineto{\pgfpoint{0.421389\du}{7.486087\du}}
\pgfpathlineto{\pgfpoint{0.198467\du}{7.765370\du}}
\pgfpathlineto{\pgfpoint{-1.342465\du}{7.765370\du}}
\pgfusepath{stroke}
\pgfsetbuttcap
\pgfsetmiterjoin
\pgfsetdash{}{0pt}
\definecolor{dialinecolor}{rgb}{0.478431, 0.478431, 0.352941}
\pgfsetfillcolor{dialinecolor}
\pgfpathmoveto{\pgfpoint{0.198467\du}{7.821100\du}}
\pgfpathlineto{\pgfpoint{0.421389\du}{7.586843\du}}
\pgfpathlineto{\pgfpoint{0.421389\du}{7.486087\du}}
\pgfpathlineto{\pgfpoint{0.198467\du}{7.765370\du}}
\pgfpathlineto{\pgfpoint{0.198467\du}{7.821100\du}}
\pgfusepath{fill}
\pgfsetbuttcap
\pgfsetmiterjoin
\pgfsetdash{}{0pt}
\definecolor{dialinecolor}{rgb}{0.286275, 0.286275, 0.211765}
\pgfsetstrokecolor{dialinecolor}
\pgfpathmoveto{\pgfpoint{0.198467\du}{7.821100\du}}
\pgfpathlineto{\pgfpoint{0.421389\du}{7.586843\du}}
\pgfpathlineto{\pgfpoint{0.421389\du}{7.486087\du}}
\pgfpathlineto{\pgfpoint{0.198467\du}{7.765370\du}}
\pgfpathlineto{\pgfpoint{0.198467\du}{7.821100\du}}
\pgfusepath{stroke}
\pgfsetbuttcap
\pgfsetmiterjoin
\pgfsetdash{}{0pt}
\definecolor{dialinecolor}{rgb}{0.717647, 0.717647, 0.615686}
\pgfsetfillcolor{dialinecolor}
\pgfpathmoveto{\pgfpoint{-1.342465\du}{7.765370\du}}
\pgfpathlineto{\pgfpoint{0.198467\du}{7.765370\du}}
\pgfpathlineto{\pgfpoint{0.198467\du}{7.821100\du}}
\pgfpathlineto{\pgfpoint{-1.342465\du}{7.821100\du}}
\pgfpathlineto{\pgfpoint{-1.342465\du}{7.765370\du}}
\pgfusepath{fill}
\pgfsetbuttcap
\pgfsetmiterjoin
\pgfsetdash{}{0pt}
\definecolor{dialinecolor}{rgb}{0.286275, 0.286275, 0.211765}
\pgfsetstrokecolor{dialinecolor}
\pgfpathmoveto{\pgfpoint{-1.342465\du}{7.765370\du}}
\pgfpathlineto{\pgfpoint{0.198467\du}{7.765370\du}}
\pgfpathlineto{\pgfpoint{0.198467\du}{7.821100\du}}
\pgfpathlineto{\pgfpoint{-1.342465\du}{7.821100\du}}
\pgfpathlineto{\pgfpoint{-1.342465\du}{7.765370\du}}
\pgfusepath{stroke}
\pgfsetbuttcap
\pgfsetmiterjoin
\pgfsetdash{}{0pt}
\definecolor{dialinecolor}{rgb}{0.000000, 0.000000, 0.000000}
\pgfsetfillcolor{dialinecolor}
\pgfpathmoveto{\pgfpoint{-1.052478\du}{7.117385\du}}
\pgfpathlineto{\pgfpoint{-0.884656\du}{6.960898\du}}
\pgfpathlineto{\pgfpoint{0.532850\du}{6.960898\du}}
\pgfpathlineto{\pgfpoint{0.376993\du}{7.117385\du}}
\pgfpathlineto{\pgfpoint{-1.052478\du}{7.117385\du}}
\pgfusepath{fill}
\pgfsetbuttcap
\pgfsetmiterjoin
\pgfsetdash{}{0pt}
\definecolor{dialinecolor}{rgb}{0.000000, 0.000000, 0.000000}
\pgfsetstrokecolor{dialinecolor}
\pgfpathmoveto{\pgfpoint{-1.052478\du}{7.117385\du}}
\pgfpathlineto{\pgfpoint{-0.884656\du}{6.960898\du}}
\pgfpathlineto{\pgfpoint{0.532850\du}{6.960898\du}}
\pgfpathlineto{\pgfpoint{0.376993\du}{7.117385\du}}
\pgfpathlineto{\pgfpoint{-1.052478\du}{7.117385\du}}
\pgfusepath{stroke}
\pgfsetbuttcap
\pgfsetmiterjoin
\pgfsetdash{}{0pt}
\definecolor{dialinecolor}{rgb}{0.788235, 0.788235, 0.713726}
\pgfsetfillcolor{dialinecolor}
\pgfpathmoveto{\pgfpoint{-1.063813\du}{5.966251\du}}
\pgfpathlineto{\pgfpoint{-0.907012\du}{5.821100\du}}
\pgfpathlineto{\pgfpoint{0.511124\du}{5.821100\du}}
\pgfpathlineto{\pgfpoint{0.354323\du}{5.966251\du}}
\pgfpathlineto{\pgfpoint{-1.063813\du}{5.966251\du}}
\pgfusepath{fill}
\pgfsetbuttcap
\pgfsetmiterjoin
\pgfsetdash{}{0pt}
\definecolor{dialinecolor}{rgb}{0.286275, 0.286275, 0.211765}
\pgfsetstrokecolor{dialinecolor}
\pgfpathmoveto{\pgfpoint{-1.063813\du}{5.966251\du}}
\pgfpathlineto{\pgfpoint{-0.907012\du}{5.821100\du}}
\pgfpathlineto{\pgfpoint{0.511124\du}{5.821100\du}}
\pgfpathlineto{\pgfpoint{0.354323\du}{5.966251\du}}
\pgfpathlineto{\pgfpoint{-1.063813\du}{5.966251\du}}
\pgfusepath{stroke}
\pgfsetbuttcap
\pgfsetmiterjoin
\pgfsetdash{}{0pt}
\definecolor{dialinecolor}{rgb}{0.717647, 0.717647, 0.615686}
\pgfsetfillcolor{dialinecolor}
\pgfpathmoveto{\pgfpoint{-1.063813\du}{5.966251\du}}
\pgfpathlineto{\pgfpoint{0.365658\du}{5.966251\du}}
\pgfpathlineto{\pgfpoint{0.365658\du}{7.094715\du}}
\pgfpathlineto{\pgfpoint{-1.063813\du}{7.094715\du}}
\pgfpathlineto{\pgfpoint{-1.063813\du}{5.966251\du}}
\pgfusepath{fill}
\pgfsetbuttcap
\pgfsetmiterjoin
\pgfsetdash{}{0pt}
\definecolor{dialinecolor}{rgb}{0.286275, 0.286275, 0.211765}
\pgfsetstrokecolor{dialinecolor}
\pgfpathmoveto{\pgfpoint{-1.063813\du}{5.966251\du}}
\pgfpathlineto{\pgfpoint{0.365029\du}{5.966251\du}}
\pgfpathlineto{\pgfpoint{0.365029\du}{7.094400\du}}
\pgfpathlineto{\pgfpoint{-1.063813\du}{7.094400\du}}
\pgfpathlineto{\pgfpoint{-1.063813\du}{5.966251\du}}
\pgfusepath{stroke}
\pgfsetbuttcap
\pgfsetmiterjoin
\pgfsetdash{}{0pt}
\definecolor{dialinecolor}{rgb}{1.000000, 1.000000, 1.000000}
\pgfsetfillcolor{dialinecolor}
\pgfpathmoveto{\pgfpoint{-0.941017\du}{6.111087\du}}
\pgfpathlineto{\pgfpoint{0.242548\du}{6.111087\du}}
\pgfpathlineto{\pgfpoint{0.242548\du}{6.982939\du}}
\pgfpathlineto{\pgfpoint{-0.941017\du}{6.982939\du}}
\pgfpathlineto{\pgfpoint{-0.941017\du}{6.111087\du}}
\pgfusepath{fill}
\pgfsetbuttcap
\pgfsetmiterjoin
\pgfsetdash{}{0pt}
\definecolor{dialinecolor}{rgb}{0.286275, 0.286275, 0.211765}
\pgfsetstrokecolor{dialinecolor}
\pgfpathmoveto{\pgfpoint{-0.941017\du}{6.111087\du}}
\pgfpathlineto{\pgfpoint{0.242548\du}{6.111087\du}}
\pgfpathlineto{\pgfpoint{0.242548\du}{6.982624\du}}
\pgfpathlineto{\pgfpoint{-0.941017\du}{6.982624\du}}
\pgfpathlineto{\pgfpoint{-0.941017\du}{6.111087\du}}
\pgfusepath{stroke}
\pgfsetbuttcap
\pgfsetmiterjoin
\pgfsetdash{}{0pt}
\definecolor{dialinecolor}{rgb}{0.478431, 0.478431, 0.352941}
\pgfsetfillcolor{dialinecolor}
\pgfpathmoveto{\pgfpoint{0.354323\du}{7.084009\du}}
\pgfpathlineto{\pgfpoint{0.511124\du}{6.927523\du}}
\pgfpathlineto{\pgfpoint{0.511124\du}{5.821100\du}}
\pgfpathlineto{\pgfpoint{0.354323\du}{5.966251\du}}
\pgfpathlineto{\pgfpoint{0.354323\du}{7.084009\du}}
\pgfusepath{fill}
\pgfsetbuttcap
\pgfsetmiterjoin
\pgfsetdash{}{0pt}
\definecolor{dialinecolor}{rgb}{0.286275, 0.286275, 0.211765}
\pgfsetstrokecolor{dialinecolor}
\pgfpathmoveto{\pgfpoint{0.354323\du}{7.084009\du}}
\pgfpathlineto{\pgfpoint{0.511124\du}{6.927523\du}}
\pgfpathlineto{\pgfpoint{0.511124\du}{5.821100\du}}
\pgfpathlineto{\pgfpoint{0.354323\du}{5.966251\du}}
\pgfpathlineto{\pgfpoint{0.354323\du}{7.084009\du}}
\pgfusepath{stroke}
\pgfsetlinewidth{0.100000\du}
\pgfsetdash{}{0pt}
\pgfsetdash{}{0pt}
\pgfsetbuttcap
\pgfsetmiterjoin
\pgfsetlinewidth{0.001000\du}
\pgfsetbuttcap
\pgfsetmiterjoin
\pgfsetdash{}{0pt}
\definecolor{dialinecolor}{rgb}{0.000000, 0.470588, 0.666667}
\pgfsetstrokecolor{dialinecolor}
\pgfpathmoveto{\pgfpoint{-15.480683\du}{21.366254\du}}
\pgfpathlineto{\pgfpoint{-15.474115\du}{21.357297\du}}
\pgfpathlineto{\pgfpoint{-15.467845\du}{21.347744\du}}
\pgfpathlineto{\pgfpoint{-15.462471\du}{21.337593\du}}
\pgfpathlineto{\pgfpoint{-15.457396\du}{21.328338\du}}
\pgfpathlineto{\pgfpoint{-15.453515\du}{21.318187\du}}
\pgfpathlineto{\pgfpoint{-15.450231\du}{21.308036\du}}
\pgfpathlineto{\pgfpoint{-15.447245\du}{21.297886\du}}
\pgfpathlineto{\pgfpoint{-15.445454\du}{21.287436\du}}
\pgfpathlineto{\pgfpoint{-15.444260\du}{21.277285\du}}
\pgfpathlineto{\pgfpoint{-15.443364\du}{21.266836\du}}
\pgfpathlineto{\pgfpoint{-15.443364\du}{21.256387\du}}
\pgfpathlineto{\pgfpoint{-15.444260\du}{21.245639\du}}
\pgfpathlineto{\pgfpoint{-15.445454\du}{21.235488\du}}
\pgfpathlineto{\pgfpoint{-15.447245\du}{21.225039\du}}
\pgfpathlineto{\pgfpoint{-15.449932\du}{21.215187\du}}
\pgfpathlineto{\pgfpoint{-15.453515\du}{21.204737\du}}
\pgfpathlineto{\pgfpoint{-15.457396\du}{21.194885\du}}
\pgfpathlineto{\pgfpoint{-15.462471\du}{21.184734\du}}
\pgfpathlineto{\pgfpoint{-15.467845\du}{21.175181\du}}
\pgfpathlineto{\pgfpoint{-15.474115\du}{21.166224\du}}
\pgfpathlineto{\pgfpoint{-15.480683\du}{21.156670\du}}
\pgfpathlineto{\pgfpoint{-15.488147\du}{21.148012\du}}
\pgfpathlineto{\pgfpoint{-15.495611\du}{21.139354\du}}
\pgfpathlineto{\pgfpoint{-15.503970\du}{21.130995\du}}
\pgfpathlineto{\pgfpoint{-15.512628\du}{21.122934\du}}
\pgfpathlineto{\pgfpoint{-15.521883\du}{21.115172\du}}
\pgfpathlineto{\pgfpoint{-15.531437\du}{21.108006\du}}
\pgfpathlineto{\pgfpoint{-15.542185\du}{21.101438\du}}
\pgfpathlineto{\pgfpoint{-15.552634\du}{21.094870\du}}
\pgfpathlineto{\pgfpoint{-15.563979\du}{21.088601\du}}
\pgfpathlineto{\pgfpoint{-15.575921\du}{21.082928\du}}
\pgfpathlineto{\pgfpoint{-15.587863\du}{21.077554\du}}
\pgfpathlineto{\pgfpoint{-15.600104\du}{21.073374\du}}
\pgfpathlineto{\pgfpoint{-15.612643\du}{21.069195\du}}
\pgfpathlineto{\pgfpoint{-15.625779\du}{21.065015\du}}
\pgfpathlineto{\pgfpoint{-15.638617\du}{21.062029\du}}
\pgfpathlineto{\pgfpoint{-15.651753\du}{21.059342\du}}
\pgfpathlineto{\pgfpoint{-15.665487\du}{21.056954\du}}
\pgfpathlineto{\pgfpoint{-15.678623\du}{21.055461\du}}
\pgfpathlineto{\pgfpoint{-15.692655\du}{21.054267\du}}
\pgfpathlineto{\pgfpoint{-15.706388\du}{21.053371\du}}
\pgfpathlineto{\pgfpoint{-15.720420\du}{21.053371\du}}
\pgfpathlineto{\pgfpoint{-15.733855\du}{21.053969\du}}
\pgfpathlineto{\pgfpoint{-15.747589\du}{21.054566\du}}
\pgfpathlineto{\pgfpoint{-15.761621\du}{21.056058\du}}
\pgfpathlineto{\pgfpoint{-15.774757\du}{21.058447\du}}
\pgfpathlineto{\pgfpoint{-15.788490\du}{21.060835\du}}
\pgfpathlineto{\pgfpoint{-15.801029\du}{21.063821\du}}
\pgfpathlineto{\pgfpoint{-15.814166\du}{21.066806\du}}
\pgfpathlineto{\pgfpoint{-15.826705\du}{21.070986\du}}
\pgfpathlineto{\pgfpoint{-15.839543\du}{21.075763\du}}
\pgfpathlineto{\pgfpoint{-15.851783\du}{21.080241\du}}
\pgfpathlineto{\pgfpoint{-15.863128\du}{21.085615\du}}
\pgfpathlineto{\pgfpoint{-15.874772\du}{21.091885\du}}
\pgfpathlineto{\pgfpoint{-15.885818\du}{21.097856\du}}
\pgfpathlineto{\pgfpoint{-15.896268\du}{21.104722\du}}
\pgfpathlineto{\pgfpoint{-15.906418\du}{21.111291\du}}
\pgfpathlineto{\pgfpoint{-15.915972\du}{21.119053\du}}
\pgfpathlineto{\pgfpoint{-15.925526\du}{21.127114\du}}
\pgfpathlineto{\pgfpoint{-15.933885\du}{21.135175\du}}
\pgfpathlineto{\pgfpoint{-15.941946\du}{21.143534\du}}
\pgfpathlineto{\pgfpoint{-15.949708\du}{21.152491\du}}
\pgfusepath{stroke}
\pgfsetlinewidth{0.080000\du}
\pgfsetbuttcap
\pgfsetmiterjoin
\pgfsetdash{}{0pt}
\definecolor{dialinecolor}{rgb}{0.666667, 0.901961, 1.000000}
\pgfsetstrokecolor{dialinecolor}
\pgfpathmoveto{\pgfpoint{-15.489938\du}{21.358790\du}}
\pgfpathlineto{\pgfpoint{-15.484564\du}{21.349834\du}}
\pgfpathlineto{\pgfpoint{-15.478295\du}{21.341176\du}}
\pgfpathlineto{\pgfpoint{-15.473219\du}{21.331920\du}}
\pgfpathlineto{\pgfpoint{-15.469040\du}{21.322665\du}}
\pgfpathlineto{\pgfpoint{-15.465158\du}{21.313112\du}}
\pgfpathlineto{\pgfpoint{-15.462173\du}{21.303558\du}}
\pgfpathlineto{\pgfpoint{-15.459486\du}{21.293706\du}}
\pgfpathlineto{\pgfpoint{-15.457396\du}{21.284451\du}}
\pgfpathlineto{\pgfpoint{-15.456500\du}{21.274598\du}}
\pgfpathlineto{\pgfpoint{-15.455903\du}{21.264448\du}}
\pgfpathlineto{\pgfpoint{-15.455903\du}{21.254595\du}}
\pgfpathlineto{\pgfpoint{-15.456799\du}{21.244743\du}}
\pgfpathlineto{\pgfpoint{-15.458292\du}{21.234593\du}}
\pgfpathlineto{\pgfpoint{-15.460083\du}{21.225337\du}}
\pgfpathlineto{\pgfpoint{-15.463069\du}{21.215485\du}}
\pgfpathlineto{\pgfpoint{-15.466651\du}{21.205932\du}}
\pgfpathlineto{\pgfpoint{-15.470532\du}{21.196676\du}}
\pgfpathlineto{\pgfpoint{-15.475011\du}{21.187421\du}}
\pgfpathlineto{\pgfpoint{-15.480086\du}{21.178763\du}}
\pgfpathlineto{\pgfpoint{-15.485758\du}{21.169508\du}}
\pgfpathlineto{\pgfpoint{-15.492924\du}{21.161149\du}}
\pgfpathlineto{\pgfpoint{-15.499492\du}{21.152491\du}}
\pgfpathlineto{\pgfpoint{-15.507254\du}{21.144430\du}}
\pgfpathlineto{\pgfpoint{-15.515315\du}{21.136667\du}}
\pgfpathlineto{\pgfpoint{-15.523675\du}{21.129801\du}}
\pgfpathlineto{\pgfpoint{-15.532333\du}{21.122337\du}}
\pgfpathlineto{\pgfpoint{-15.542185\du}{21.115470\du}}
\pgfpathlineto{\pgfpoint{-15.552634\du}{21.108604\du}}
\pgfpathlineto{\pgfpoint{-15.562785\du}{21.102931\du}}
\pgfpathlineto{\pgfpoint{-15.573234\du}{21.097557\du}}
\pgfpathlineto{\pgfpoint{-15.585176\du}{21.092183\du}}
\pgfpathlineto{\pgfpoint{-15.595924\du}{21.087108\du}}
\pgfpathlineto{\pgfpoint{-15.608463\du}{21.082928\du}}
\pgfpathlineto{\pgfpoint{-15.619808\du}{21.079047\du}}
\pgfpathlineto{\pgfpoint{-15.632348\du}{21.075763\du}}
\pgfpathlineto{\pgfpoint{-15.644887\du}{21.072777\du}}
\pgfpathlineto{\pgfpoint{-15.657426\du}{21.070090\du}}
\pgfpathlineto{\pgfpoint{-15.670562\du}{21.068299\du}}
\pgfpathlineto{\pgfpoint{-15.683101\du}{21.066806\du}}
\pgfpathlineto{\pgfpoint{-15.696238\du}{21.065612\du}}
\pgfpathlineto{\pgfpoint{-15.710270\du}{21.065015\du}}
\pgfpathlineto{\pgfpoint{-15.723406\du}{21.065015\du}}
\pgfpathlineto{\pgfpoint{-15.736244\du}{21.065612\du}}
\pgfpathlineto{\pgfpoint{-15.749380\du}{21.066508\du}}
\pgfpathlineto{\pgfpoint{-15.761919\du}{21.068299\du}}
\pgfpathlineto{\pgfpoint{-15.775055\du}{21.070090\du}}
\pgfpathlineto{\pgfpoint{-15.787595\du}{21.072180\du}}
\pgfpathlineto{\pgfpoint{-15.800432\du}{21.075763\du}}
\pgfpathlineto{\pgfpoint{-15.812673\du}{21.078748\du}}
\pgfpathlineto{\pgfpoint{-15.824615\du}{21.082630\du}}
\pgfpathlineto{\pgfpoint{-15.836856\du}{21.086809\du}}
\pgfpathlineto{\pgfpoint{-15.848201\du}{21.091885\du}}
\pgfpathlineto{\pgfpoint{-15.859247\du}{21.096661\du}}
\pgfpathlineto{\pgfpoint{-15.869995\du}{21.102633\du}}
\pgfpathlineto{\pgfpoint{-15.880444\du}{21.108305\du}}
\pgfpathlineto{\pgfpoint{-15.890894\du}{21.114873\du}}
\pgfpathlineto{\pgfpoint{-15.900447\du}{21.121441\du}}
\pgfpathlineto{\pgfpoint{-15.909105\du}{21.128607\du}}
\pgfpathlineto{\pgfpoint{-15.918659\du}{21.135772\du}}
\pgfpathlineto{\pgfpoint{-15.926421\du}{21.143534\du}}
\pgfpathlineto{\pgfpoint{-15.933885\du}{21.151595\du}}
\pgfpathlineto{\pgfpoint{-15.941050\du}{21.159955\du}}
\pgfusepath{stroke}
\pgfsetbuttcap
\pgfsetmiterjoin
\pgfsetdash{}{0pt}
\definecolor{dialinecolor}{rgb}{0.000000, 0.705882, 1.000000}
\pgfsetfillcolor{dialinecolor}
\pgfpathmoveto{\pgfpoint{-17.406941\du}{21.041728\du}}
\pgfpathlineto{\pgfpoint{-17.204224\du}{20.850953\du}}
\pgfpathlineto{\pgfpoint{-15.692655\du}{20.850953\du}}
\pgfpathlineto{\pgfpoint{-15.895372\du}{21.041728\du}}
\pgfpathlineto{\pgfpoint{-17.406941\du}{21.041728\du}}
\pgfusepath{fill}
\pgfsetlinewidth{0.001000\du}
\pgfsetbuttcap
\pgfsetmiterjoin
\pgfsetdash{}{0pt}
\definecolor{dialinecolor}{rgb}{0.666667, 0.901961, 1.000000}
\pgfsetstrokecolor{dialinecolor}
\pgfpathmoveto{\pgfpoint{-17.406941\du}{21.041728\du}}
\pgfpathlineto{\pgfpoint{-17.204224\du}{20.850953\du}}
\pgfpathlineto{\pgfpoint{-15.692655\du}{20.850953\du}}
\pgfpathlineto{\pgfpoint{-15.895372\du}{21.041728\du}}
\pgfpathlineto{\pgfpoint{-17.406941\du}{21.041728\du}}
\pgfusepath{stroke}
\pgfsetbuttcap
\pgfsetmiterjoin
\pgfsetdash{}{0pt}
\definecolor{dialinecolor}{rgb}{0.000000, 0.588235, 0.831373}
\pgfsetfillcolor{dialinecolor}
\pgfpathmoveto{\pgfpoint{-17.406941\du}{21.041728\du}}
\pgfpathlineto{\pgfpoint{-15.895372\du}{21.041728\du}}
\pgfpathlineto{\pgfpoint{-15.895372\du}{21.327741\du}}
\pgfpathlineto{\pgfpoint{-17.406941\du}{21.327741\du}}
\pgfpathlineto{\pgfpoint{-17.406941\du}{21.041728\du}}
\pgfusepath{fill}
\pgfsetbuttcap
\pgfsetmiterjoin
\pgfsetdash{}{0pt}
\definecolor{dialinecolor}{rgb}{0.666667, 0.901961, 1.000000}
\pgfsetstrokecolor{dialinecolor}
\pgfpathmoveto{\pgfpoint{-17.406941\du}{21.041728\du}}
\pgfpathlineto{\pgfpoint{-15.895372\du}{21.041728\du}}
\pgfpathlineto{\pgfpoint{-15.895372\du}{21.327741\du}}
\pgfpathlineto{\pgfpoint{-17.406941\du}{21.327741\du}}
\pgfpathlineto{\pgfpoint{-17.406941\du}{21.041728\du}}
\pgfusepath{stroke}
\pgfsetbuttcap
\pgfsetmiterjoin
\pgfsetdash{}{0pt}
\definecolor{dialinecolor}{rgb}{0.000000, 0.352941, 0.501961}
\pgfsetfillcolor{dialinecolor}
\pgfpathmoveto{\pgfpoint{-15.895372\du}{21.327741\du}}
\pgfpathlineto{\pgfpoint{-15.692655\du}{21.125024\du}}
\pgfpathlineto{\pgfpoint{-15.692655\du}{20.850953\du}}
\pgfpathlineto{\pgfpoint{-15.895372\du}{21.041728\du}}
\pgfpathlineto{\pgfpoint{-15.895372\du}{21.327741\du}}
\pgfusepath{fill}
\pgfsetbuttcap
\pgfsetmiterjoin
\pgfsetdash{}{0pt}
\definecolor{dialinecolor}{rgb}{0.666667, 0.901961, 1.000000}
\pgfsetstrokecolor{dialinecolor}
\pgfpathmoveto{\pgfpoint{-15.895372\du}{21.327741\du}}
\pgfpathlineto{\pgfpoint{-15.692655\du}{21.125024\du}}
\pgfpathlineto{\pgfpoint{-15.692655\du}{20.850953\du}}
\pgfpathlineto{\pgfpoint{-15.895372\du}{21.041728\du}}
\pgfpathlineto{\pgfpoint{-15.895372\du}{21.327741\du}}
\pgfusepath{stroke}
\pgfsetbuttcap
\pgfsetmiterjoin
\pgfsetdash{}{0pt}
\definecolor{dialinecolor}{rgb}{0.000000, 0.000000, 0.000000}
\pgfsetfillcolor{dialinecolor}
\pgfpathmoveto{\pgfpoint{-17.358874\du}{21.005603\du}}
\pgfpathlineto{\pgfpoint{-17.204224\du}{20.850953\du}}
\pgfpathlineto{\pgfpoint{-15.728780\du}{20.850953\du}}
\pgfpathlineto{\pgfpoint{-15.883131\du}{21.005603\du}}
\pgfpathlineto{\pgfpoint{-17.358874\du}{21.005603\du}}
\pgfusepath{fill}
\pgfsetbuttcap
\pgfsetmiterjoin
\pgfsetdash{}{0pt}
\definecolor{dialinecolor}{rgb}{0.000000, 0.000000, 0.000000}
\pgfsetstrokecolor{dialinecolor}
\pgfpathmoveto{\pgfpoint{-17.358874\du}{21.005603\du}}
\pgfpathlineto{\pgfpoint{-17.204224\du}{20.850953\du}}
\pgfpathlineto{\pgfpoint{-15.728780\du}{20.850953\du}}
\pgfpathlineto{\pgfpoint{-15.883131\du}{21.005603\du}}
\pgfpathlineto{\pgfpoint{-17.358874\du}{21.005603\du}}
\pgfusepath{stroke}
\pgfsetbuttcap
\pgfsetmiterjoin
\pgfsetdash{}{0pt}
\definecolor{dialinecolor}{rgb}{0.000000, 0.705882, 1.000000}
\pgfsetfillcolor{dialinecolor}
\pgfpathmoveto{\pgfpoint{-17.406941\du}{19.779450\du}}
\pgfpathlineto{\pgfpoint{-17.239752\du}{19.624800\du}}
\pgfpathlineto{\pgfpoint{-15.728780\du}{19.624800\du}}
\pgfpathlineto{\pgfpoint{-15.895372\du}{19.779450\du}}
\pgfpathlineto{\pgfpoint{-17.406941\du}{19.779450\du}}
\pgfusepath{fill}
\pgfsetbuttcap
\pgfsetmiterjoin
\pgfsetdash{}{0pt}
\definecolor{dialinecolor}{rgb}{0.666667, 0.901961, 1.000000}
\pgfsetstrokecolor{dialinecolor}
\pgfpathmoveto{\pgfpoint{-17.406941\du}{19.779450\du}}
\pgfpathlineto{\pgfpoint{-17.239752\du}{19.624800\du}}
\pgfpathlineto{\pgfpoint{-15.728780\du}{19.624800\du}}
\pgfpathlineto{\pgfpoint{-15.895372\du}{19.779450\du}}
\pgfpathlineto{\pgfpoint{-17.406941\du}{19.779450\du}}
\pgfusepath{stroke}
\pgfsetbuttcap
\pgfsetmiterjoin
\pgfsetdash{}{0pt}
\definecolor{dialinecolor}{rgb}{0.000000, 0.588235, 0.831373}
\pgfsetfillcolor{dialinecolor}
\pgfpathmoveto{\pgfpoint{-17.406941\du}{19.779450\du}}
\pgfpathlineto{\pgfpoint{-15.883131\du}{19.779450\du}}
\pgfpathlineto{\pgfpoint{-15.883131\du}{20.981719\du}}
\pgfpathlineto{\pgfpoint{-17.406941\du}{20.981719\du}}
\pgfpathlineto{\pgfpoint{-17.406941\du}{19.779450\du}}
\pgfusepath{fill}
\pgfsetbuttcap
\pgfsetmiterjoin
\pgfsetdash{}{0pt}
\definecolor{dialinecolor}{rgb}{0.666667, 0.901961, 1.000000}
\pgfsetstrokecolor{dialinecolor}
\pgfpathmoveto{\pgfpoint{-17.406941\du}{19.779450\du}}
\pgfpathlineto{\pgfpoint{-15.883430\du}{19.779450\du}}
\pgfpathlineto{\pgfpoint{-15.883430\du}{20.981420\du}}
\pgfpathlineto{\pgfpoint{-17.406941\du}{20.981420\du}}
\pgfpathlineto{\pgfpoint{-17.406941\du}{19.779450\du}}
\pgfusepath{stroke}
\pgfsetbuttcap
\pgfsetmiterjoin
\pgfsetdash{}{0pt}
\definecolor{dialinecolor}{rgb}{1.000000, 1.000000, 1.000000}
\pgfsetfillcolor{dialinecolor}
\pgfpathmoveto{\pgfpoint{-17.275876\du}{19.934398\du}}
\pgfpathlineto{\pgfpoint{-16.014196\du}{19.934398\du}}
\pgfpathlineto{\pgfpoint{-16.014196\du}{20.863194\du}}
\pgfpathlineto{\pgfpoint{-17.275876\du}{20.863194\du}}
\pgfpathlineto{\pgfpoint{-17.275876\du}{19.934398\du}}
\pgfusepath{fill}
\pgfsetbuttcap
\pgfsetmiterjoin
\pgfsetdash{}{0pt}
\definecolor{dialinecolor}{rgb}{0.000000, 0.235294, 0.333333}
\pgfsetstrokecolor{dialinecolor}
\pgfpathmoveto{\pgfpoint{-17.275876\du}{19.934398\du}}
\pgfpathlineto{\pgfpoint{-16.014494\du}{19.934398\du}}
\pgfpathlineto{\pgfpoint{-16.014494\du}{20.862895\du}}
\pgfpathlineto{\pgfpoint{-17.275876\du}{20.862895\du}}
\pgfpathlineto{\pgfpoint{-17.275876\du}{19.934398\du}}
\pgfusepath{stroke}
\pgfsetbuttcap
\pgfsetmiterjoin
\pgfsetdash{}{0pt}
\definecolor{dialinecolor}{rgb}{0.000000, 0.352941, 0.501961}
\pgfsetfillcolor{dialinecolor}
\pgfpathmoveto{\pgfpoint{-15.895372\du}{20.970374\du}}
\pgfpathlineto{\pgfpoint{-15.728780\du}{20.803483\du}}
\pgfpathlineto{\pgfpoint{-15.728780\du}{19.624800\du}}
\pgfpathlineto{\pgfpoint{-15.895372\du}{19.779450\du}}
\pgfpathlineto{\pgfpoint{-15.895372\du}{20.970374\du}}
\pgfusepath{fill}
\pgfsetbuttcap
\pgfsetmiterjoin
\pgfsetdash{}{0pt}
\definecolor{dialinecolor}{rgb}{0.666667, 0.901961, 1.000000}
\pgfsetstrokecolor{dialinecolor}
\pgfpathmoveto{\pgfpoint{-15.895372\du}{20.970374\du}}
\pgfpathlineto{\pgfpoint{-15.728780\du}{20.803483\du}}
\pgfpathlineto{\pgfpoint{-15.728780\du}{19.624800\du}}
\pgfpathlineto{\pgfpoint{-15.895372\du}{19.779450\du}}
\pgfpathlineto{\pgfpoint{-15.895372\du}{20.970374\du}}
\pgfusepath{stroke}
\pgfsetbuttcap
\pgfsetmiterjoin
\pgfsetdash{}{0pt}
\definecolor{dialinecolor}{rgb}{0.000000, 0.705882, 1.000000}
\pgfsetfillcolor{dialinecolor}
\pgfpathmoveto{\pgfpoint{-17.704000\du}{21.565687\du}}
\pgfpathlineto{\pgfpoint{-17.466353\du}{21.267732\du}}
\pgfpathlineto{\pgfpoint{-15.823719\du}{21.267732\du}}
\pgfpathlineto{\pgfpoint{-16.061665\du}{21.565687\du}}
\pgfpathlineto{\pgfpoint{-17.704000\du}{21.565687\du}}
\pgfusepath{fill}
\pgfsetbuttcap
\pgfsetmiterjoin
\pgfsetdash{}{0pt}
\definecolor{dialinecolor}{rgb}{0.666667, 0.901961, 1.000000}
\pgfsetstrokecolor{dialinecolor}
\pgfpathmoveto{\pgfpoint{-17.704000\du}{21.565687\du}}
\pgfpathlineto{\pgfpoint{-17.466353\du}{21.267732\du}}
\pgfpathlineto{\pgfpoint{-15.823719\du}{21.267732\du}}
\pgfpathlineto{\pgfpoint{-16.061665\du}{21.565687\du}}
\pgfpathlineto{\pgfpoint{-17.704000\du}{21.565687\du}}
\pgfusepath{stroke}
\pgfsetbuttcap
\pgfsetmiterjoin
\pgfsetdash{}{0pt}
\definecolor{dialinecolor}{rgb}{0.000000, 0.352941, 0.501961}
\pgfsetfillcolor{dialinecolor}
\pgfpathmoveto{\pgfpoint{-16.061665\du}{21.624800\du}}
\pgfpathlineto{\pgfpoint{-15.823719\du}{21.374912\du}}
\pgfpathlineto{\pgfpoint{-15.823719\du}{21.267732\du}}
\pgfpathlineto{\pgfpoint{-16.061665\du}{21.565687\du}}
\pgfpathlineto{\pgfpoint{-16.061665\du}{21.624800\du}}
\pgfusepath{fill}
\pgfsetbuttcap
\pgfsetmiterjoin
\pgfsetdash{}{0pt}
\definecolor{dialinecolor}{rgb}{0.666667, 0.901961, 1.000000}
\pgfsetstrokecolor{dialinecolor}
\pgfpathmoveto{\pgfpoint{-16.061665\du}{21.624800\du}}
\pgfpathlineto{\pgfpoint{-15.823719\du}{21.374912\du}}
\pgfpathlineto{\pgfpoint{-15.823719\du}{21.267732\du}}
\pgfpathlineto{\pgfpoint{-16.061665\du}{21.565687\du}}
\pgfpathlineto{\pgfpoint{-16.061665\du}{21.624800\du}}
\pgfusepath{stroke}
\pgfsetbuttcap
\pgfsetmiterjoin
\pgfsetdash{}{0pt}
\definecolor{dialinecolor}{rgb}{0.000000, 0.588235, 0.831373}
\pgfsetfillcolor{dialinecolor}
\pgfpathmoveto{\pgfpoint{-17.704000\du}{21.565687\du}}
\pgfpathlineto{\pgfpoint{-16.061665\du}{21.565687\du}}
\pgfpathlineto{\pgfpoint{-16.061665\du}{21.624800\du}}
\pgfpathlineto{\pgfpoint{-17.704000\du}{21.624800\du}}
\pgfpathlineto{\pgfpoint{-17.704000\du}{21.565687\du}}
\pgfusepath{fill}
\pgfsetbuttcap
\pgfsetmiterjoin
\pgfsetdash{}{0pt}
\definecolor{dialinecolor}{rgb}{0.666667, 0.901961, 1.000000}
\pgfsetstrokecolor{dialinecolor}
\pgfpathmoveto{\pgfpoint{-17.704000\du}{21.565687\du}}
\pgfpathlineto{\pgfpoint{-16.061665\du}{21.565687\du}}
\pgfpathlineto{\pgfpoint{-16.061665\du}{21.624800\du}}
\pgfpathlineto{\pgfpoint{-17.704000\du}{21.624800\du}}
\pgfpathlineto{\pgfpoint{-17.704000\du}{21.565687\du}}
\pgfusepath{stroke}
\pgfsetbuttcap
\pgfsetmiterjoin
\pgfsetdash{}{0pt}
\definecolor{dialinecolor}{rgb}{0.000000, 0.705882, 1.000000}
\pgfsetfillcolor{dialinecolor}
\pgfpathmoveto{\pgfpoint{-15.728780\du}{21.529562\du}}
\pgfpathlineto{\pgfpoint{-15.609658\du}{21.351326\du}}
\pgfpathlineto{\pgfpoint{-15.394999\du}{21.351326\du}}
\pgfpathlineto{\pgfpoint{-15.513822\du}{21.529562\du}}
\pgfpathlineto{\pgfpoint{-15.728780\du}{21.529562\du}}
\pgfusepath{fill}
\pgfsetbuttcap
\pgfsetmiterjoin
\pgfsetdash{}{0pt}
\definecolor{dialinecolor}{rgb}{0.666667, 0.901961, 1.000000}
\pgfsetstrokecolor{dialinecolor}
\pgfpathmoveto{\pgfpoint{-15.728780\du}{21.529562\du}}
\pgfpathlineto{\pgfpoint{-15.609658\du}{21.351326\du}}
\pgfpathlineto{\pgfpoint{-15.394999\du}{21.351326\du}}
\pgfpathlineto{\pgfpoint{-15.513822\du}{21.529562\du}}
\pgfpathlineto{\pgfpoint{-15.728780\du}{21.529562\du}}
\pgfusepath{stroke}
\pgfsetbuttcap
\pgfsetmiterjoin
\pgfsetdash{}{0pt}
\definecolor{dialinecolor}{rgb}{0.000000, 0.352941, 0.501961}
\pgfsetfillcolor{dialinecolor}
\pgfpathmoveto{\pgfpoint{-15.513822\du}{21.589272\du}}
\pgfpathlineto{\pgfpoint{-15.394999\du}{21.458208\du}}
\pgfpathlineto{\pgfpoint{-15.394999\du}{21.351326\du}}
\pgfpathlineto{\pgfpoint{-15.513822\du}{21.529562\du}}
\pgfpathlineto{\pgfpoint{-15.513822\du}{21.589272\du}}
\pgfusepath{fill}
\pgfsetbuttcap
\pgfsetmiterjoin
\pgfsetdash{}{0pt}
\definecolor{dialinecolor}{rgb}{0.666667, 0.901961, 1.000000}
\pgfsetstrokecolor{dialinecolor}
\pgfpathmoveto{\pgfpoint{-15.513822\du}{21.589272\du}}
\pgfpathlineto{\pgfpoint{-15.394999\du}{21.458208\du}}
\pgfpathlineto{\pgfpoint{-15.394999\du}{21.351326\du}}
\pgfpathlineto{\pgfpoint{-15.513822\du}{21.529562\du}}
\pgfpathlineto{\pgfpoint{-15.513822\du}{21.589272\du}}
\pgfusepath{stroke}
\pgfsetbuttcap
\pgfsetmiterjoin
\pgfsetdash{}{0pt}
\definecolor{dialinecolor}{rgb}{0.000000, 0.588235, 0.831373}
\pgfsetfillcolor{dialinecolor}
\pgfpathmoveto{\pgfpoint{-15.728780\du}{21.529562\du}}
\pgfpathlineto{\pgfpoint{-15.513822\du}{21.529562\du}}
\pgfpathlineto{\pgfpoint{-15.513822\du}{21.589272\du}}
\pgfpathlineto{\pgfpoint{-15.728780\du}{21.589272\du}}
\pgfpathlineto{\pgfpoint{-15.728780\du}{21.529562\du}}
\pgfusepath{fill}
\pgfsetbuttcap
\pgfsetmiterjoin
\pgfsetdash{}{0pt}
\definecolor{dialinecolor}{rgb}{0.666667, 0.901961, 1.000000}
\pgfsetstrokecolor{dialinecolor}
\pgfpathmoveto{\pgfpoint{-15.728780\du}{21.529562\du}}
\pgfpathlineto{\pgfpoint{-15.513822\du}{21.529562\du}}
\pgfpathlineto{\pgfpoint{-15.513822\du}{21.589272\du}}
\pgfpathlineto{\pgfpoint{-15.728780\du}{21.589272\du}}
\pgfpathlineto{\pgfpoint{-15.728780\du}{21.529562\du}}
\pgfusepath{stroke}
\pgfsetbuttcap
\pgfsetmiterjoin
\pgfsetdash{}{0pt}
\definecolor{dialinecolor}{rgb}{0.811765, 0.054902, 0.188235}
\pgfsetstrokecolor{dialinecolor}
\pgfpathmoveto{\pgfpoint{-16.942692\du}{20.160701\du}}
\pgfpathlineto{\pgfpoint{-16.608911\du}{20.160701\du}}
\pgfpathlineto{\pgfpoint{-16.692506\du}{20.243997\du}}
\pgfpathlineto{\pgfpoint{-16.406791\du}{20.243997\du}}
\pgfusepath{stroke}
\pgfsetbuttcap
\pgfsetmiterjoin
\pgfsetdash{}{0pt}
\definecolor{dialinecolor}{rgb}{0.811765, 0.054902, 0.188235}
\pgfsetstrokecolor{dialinecolor}
\pgfpathmoveto{\pgfpoint{-16.978519\du}{20.160701\du}}
\pgfpathlineto{\pgfpoint{-16.739976\du}{20.398647\du}}
\pgfpathlineto{\pgfpoint{-16.859396\du}{20.398647\du}}
\pgfpathlineto{\pgfpoint{-16.656978\du}{20.601364\du}}
\pgfusepath{stroke}
\pgfsetbuttcap
\pgfsetmiterjoin
\pgfsetdash{}{0pt}
\definecolor{dialinecolor}{rgb}{0.811765, 0.054902, 0.188235}
\pgfsetstrokecolor{dialinecolor}
\pgfpathmoveto{\pgfpoint{-16.300208\du}{20.160701\du}}
\pgfpathlineto{\pgfpoint{-16.537557\du}{20.398647\du}}
\pgfpathlineto{\pgfpoint{-16.419032\du}{20.398647\du}}
\pgfpathlineto{\pgfpoint{-16.621152\du}{20.601364\du}}
\pgfusepath{stroke}
\pgfsetbuttcap
\pgfsetmiterjoin
\pgfsetdash{}{0pt}
\definecolor{dialinecolor}{rgb}{0.000000, 0.470588, 0.666667}
\pgfsetfillcolor{dialinecolor}
\pgfpathmoveto{\pgfpoint{-16.811927\du}{20.249669\du}}
\pgfpathlineto{\pgfpoint{-16.812524\du}{20.246385\du}}
\pgfpathlineto{\pgfpoint{-16.813419\du}{20.243101\du}}
\pgfpathlineto{\pgfpoint{-16.814315\du}{20.239817\du}}
\pgfpathlineto{\pgfpoint{-16.816106\du}{20.236533\du}}
\pgfpathlineto{\pgfpoint{-16.818495\du}{20.233548\du}}
\pgfpathlineto{\pgfpoint{-16.821182\du}{20.230264\du}}
\pgfpathlineto{\pgfpoint{-16.824167\du}{20.226979\du}}
\pgfpathlineto{\pgfpoint{-16.827750\du}{20.223994\du}}
\pgfpathlineto{\pgfpoint{-16.831930\du}{20.221008\du}}
\pgfpathlineto{\pgfpoint{-16.836706\du}{20.218023\du}}
\pgfpathlineto{\pgfpoint{-16.841782\du}{20.215336\du}}
\pgfpathlineto{\pgfpoint{-16.847454\du}{20.212350\du}}
\pgfpathlineto{\pgfpoint{-16.853127\du}{20.209962\du}}
\pgfpathlineto{\pgfpoint{-16.859396\du}{20.207275\du}}
\pgfpathlineto{\pgfpoint{-16.865964\du}{20.204887\du}}
\pgfpathlineto{\pgfpoint{-16.873130\du}{20.202498\du}}
\pgfpathlineto{\pgfpoint{-16.880594\du}{20.200408\du}}
\pgfpathlineto{\pgfpoint{-16.888057\du}{20.198020\du}}
\pgfpathlineto{\pgfpoint{-16.896417\du}{20.196527\du}}
\pgfpathlineto{\pgfpoint{-16.904776\du}{20.194437\du}}
\pgfpathlineto{\pgfpoint{-16.912837\du}{20.192646\du}}
\pgfpathlineto{\pgfpoint{-16.922092\du}{20.191452\du}}
\pgfpathlineto{\pgfpoint{-16.930750\du}{20.189959\du}}
\pgfpathlineto{\pgfpoint{-16.940304\du}{20.188765\du}}
\pgfpathlineto{\pgfpoint{-16.949559\du}{20.187571\du}}
\pgfpathlineto{\pgfpoint{-16.958814\du}{20.186675\du}}
\pgfpathlineto{\pgfpoint{-16.968666\du}{20.185779\du}}
\pgfpathlineto{\pgfpoint{-16.978519\du}{20.185182\du}}
\pgfpathlineto{\pgfpoint{-16.988371\du}{20.184884\du}}
\pgfpathlineto{\pgfpoint{-16.998522\du}{20.184585\du}}
\pgfpathlineto{\pgfpoint{-17.008374\du}{20.184585\du}}
\pgfpathlineto{\pgfpoint{-17.008374\du}{20.184585\du}}
\pgfpathlineto{\pgfpoint{-17.018226\du}{20.184585\du}}
\pgfpathlineto{\pgfpoint{-17.028078\du}{20.184884\du}}
\pgfpathlineto{\pgfpoint{-17.038229\du}{20.185182\du}}
\pgfpathlineto{\pgfpoint{-17.048081\du}{20.185779\du}}
\pgfpathlineto{\pgfpoint{-17.057635\du}{20.186675\du}}
\pgfpathlineto{\pgfpoint{-17.067189\du}{20.187571\du}}
\pgfpathlineto{\pgfpoint{-17.076444\du}{20.188765\du}}
\pgfpathlineto{\pgfpoint{-17.085699\du}{20.189959\du}}
\pgfpathlineto{\pgfpoint{-17.094655\du}{20.191452\du}}
\pgfpathlineto{\pgfpoint{-17.103612\du}{20.192646\du}}
\pgfpathlineto{\pgfpoint{-17.111971\du}{20.194437\du}}
\pgfpathlineto{\pgfpoint{-17.120331\du}{20.196527\du}}
\pgfpathlineto{\pgfpoint{-17.128690\du}{20.198020\du}}
\pgfpathlineto{\pgfpoint{-17.135856\du}{20.200408\du}}
\pgfpathlineto{\pgfpoint{-17.143319\du}{20.202498\du}}
\pgfpathlineto{\pgfpoint{-17.150485\du}{20.204887\du}}
\pgfpathlineto{\pgfpoint{-17.157351\du}{20.207275\du}}
\pgfpathlineto{\pgfpoint{-17.163322\du}{20.209962\du}}
\pgfpathlineto{\pgfpoint{-17.169293\du}{20.212350\du}}
\pgfpathlineto{\pgfpoint{-17.174966\du}{20.215336\du}}
\pgfpathlineto{\pgfpoint{-17.180041\du}{20.218023\du}}
\pgfpathlineto{\pgfpoint{-17.184818\du}{20.221008\du}}
\pgfpathlineto{\pgfpoint{-17.188998\du}{20.223994\du}}
\pgfpathlineto{\pgfpoint{-17.192282\du}{20.226979\du}}
\pgfpathlineto{\pgfpoint{-17.195566\du}{20.230264\du}}
\pgfpathlineto{\pgfpoint{-17.198551\du}{20.233548\du}}
\pgfpathlineto{\pgfpoint{-17.200641\du}{20.236533\du}}
\pgfpathlineto{\pgfpoint{-17.202134\du}{20.239817\du}}
\pgfpathlineto{\pgfpoint{-17.203925\du}{20.243101\du}}
\pgfpathlineto{\pgfpoint{-17.204224\du}{20.246385\du}}
\pgfpathlineto{\pgfpoint{-17.204522\du}{20.249669\du}}
\pgfpathlineto{\pgfpoint{-17.204522\du}{20.249669\du}}
\pgfpathlineto{\pgfpoint{-17.204224\du}{20.253252\du}}
\pgfpathlineto{\pgfpoint{-17.203925\du}{20.256238\du}}
\pgfpathlineto{\pgfpoint{-17.202134\du}{20.259522\du}}
\pgfpathlineto{\pgfpoint{-17.200641\du}{20.262806\du}}
\pgfpathlineto{\pgfpoint{-17.198551\du}{20.266090\du}}
\pgfpathlineto{\pgfpoint{-17.195566\du}{20.269374\du}}
\pgfpathlineto{\pgfpoint{-17.192282\du}{20.272359\du}}
\pgfpathlineto{\pgfpoint{-17.188998\du}{20.275345\du}}
\pgfpathlineto{\pgfpoint{-17.184818\du}{20.278330\du}}
\pgfpathlineto{\pgfpoint{-17.180041\du}{20.281614\du}}
\pgfpathlineto{\pgfpoint{-17.174966\du}{20.284301\du}}
\pgfpathlineto{\pgfpoint{-17.169293\du}{20.286988\du}}
\pgfpathlineto{\pgfpoint{-17.163322\du}{20.289675\du}}
\pgfpathlineto{\pgfpoint{-17.157351\du}{20.292362\du}}
\pgfpathlineto{\pgfpoint{-17.150485\du}{20.294751\du}}
\pgfpathlineto{\pgfpoint{-17.143319\du}{20.297139\du}}
\pgfpathlineto{\pgfpoint{-17.135856\du}{20.298930\du}}
\pgfpathlineto{\pgfpoint{-17.128690\du}{20.301020\du}}
\pgfpathlineto{\pgfpoint{-17.120331\du}{20.303110\du}}
\pgfpathlineto{\pgfpoint{-17.111971\du}{20.305200\du}}
\pgfpathlineto{\pgfpoint{-17.103612\du}{20.306693\du}}
\pgfpathlineto{\pgfpoint{-17.094655\du}{20.308186\du}}
\pgfpathlineto{\pgfpoint{-17.085699\du}{20.309678\du}}
\pgfpathlineto{\pgfpoint{-17.076444\du}{20.310873\du}}
\pgfpathlineto{\pgfpoint{-17.067189\du}{20.311768\du}}
\pgfpathlineto{\pgfpoint{-17.057635\du}{20.312962\du}}
\pgfpathlineto{\pgfpoint{-17.048081\du}{20.313560\du}}
\pgfpathlineto{\pgfpoint{-17.038229\du}{20.314157\du}}
\pgfpathlineto{\pgfpoint{-17.028078\du}{20.314754\du}}
\pgfpathlineto{\pgfpoint{-17.018226\du}{20.315052\du}}
\pgfpathlineto{\pgfpoint{-17.008374\du}{20.315052\du}}
\pgfpathlineto{\pgfpoint{-17.008374\du}{20.315052\du}}
\pgfpathlineto{\pgfpoint{-16.998522\du}{20.315052\du}}
\pgfpathlineto{\pgfpoint{-16.988371\du}{20.314754\du}}
\pgfpathlineto{\pgfpoint{-16.978519\du}{20.314157\du}}
\pgfpathlineto{\pgfpoint{-16.968666\du}{20.313560\du}}
\pgfpathlineto{\pgfpoint{-16.958814\du}{20.312962\du}}
\pgfpathlineto{\pgfpoint{-16.949559\du}{20.311768\du}}
\pgfpathlineto{\pgfpoint{-16.940304\du}{20.310873\du}}
\pgfpathlineto{\pgfpoint{-16.930750\du}{20.309678\du}}
\pgfpathlineto{\pgfpoint{-16.922092\du}{20.308186\du}}
\pgfpathlineto{\pgfpoint{-16.912837\du}{20.306693\du}}
\pgfpathlineto{\pgfpoint{-16.904776\du}{20.305200\du}}
\pgfpathlineto{\pgfpoint{-16.896417\du}{20.303110\du}}
\pgfpathlineto{\pgfpoint{-16.888057\du}{20.301020\du}}
\pgfpathlineto{\pgfpoint{-16.880594\du}{20.298930\du}}
\pgfpathlineto{\pgfpoint{-16.873130\du}{20.297139\du}}
\pgfpathlineto{\pgfpoint{-16.865964\du}{20.294751\du}}
\pgfpathlineto{\pgfpoint{-16.859396\du}{20.292362\du}}
\pgfpathlineto{\pgfpoint{-16.853127\du}{20.289675\du}}
\pgfpathlineto{\pgfpoint{-16.847454\du}{20.286988\du}}
\pgfpathlineto{\pgfpoint{-16.841782\du}{20.284301\du}}
\pgfpathlineto{\pgfpoint{-16.836706\du}{20.281614\du}}
\pgfpathlineto{\pgfpoint{-16.831930\du}{20.278330\du}}
\pgfpathlineto{\pgfpoint{-16.827750\du}{20.275345\du}}
\pgfpathlineto{\pgfpoint{-16.824167\du}{20.272359\du}}
\pgfpathlineto{\pgfpoint{-16.821182\du}{20.269374\du}}
\pgfpathlineto{\pgfpoint{-16.818495\du}{20.266090\du}}
\pgfpathlineto{\pgfpoint{-16.816106\du}{20.262806\du}}
\pgfpathlineto{\pgfpoint{-16.814315\du}{20.259522\du}}
\pgfpathlineto{\pgfpoint{-16.813419\du}{20.256238\du}}
\pgfpathlineto{\pgfpoint{-16.812524\du}{20.253252\du}}
\pgfpathlineto{\pgfpoint{-16.811927\du}{20.249669\du}}
\pgfusepath{fill}
\pgfsetbuttcap
\pgfsetmiterjoin
\pgfsetdash{}{0pt}
\definecolor{dialinecolor}{rgb}{0.666667, 0.901961, 1.000000}
\pgfsetstrokecolor{dialinecolor}
\pgfpathmoveto{\pgfpoint{-16.824167\du}{20.243997\du}}
\pgfpathlineto{\pgfpoint{-16.824167\du}{20.240414\du}}
\pgfpathlineto{\pgfpoint{-16.825063\du}{20.238026\du}}
\pgfpathlineto{\pgfpoint{-16.826257\du}{20.234742\du}}
\pgfpathlineto{\pgfpoint{-16.827750\du}{20.231756\du}}
\pgfpathlineto{\pgfpoint{-16.830138\du}{20.229069\du}}
\pgfpathlineto{\pgfpoint{-16.832527\du}{20.225487\du}}
\pgfpathlineto{\pgfpoint{-16.835811\du}{20.223098\du}}
\pgfpathlineto{\pgfpoint{-16.839095\du}{20.220411\du}}
\pgfpathlineto{\pgfpoint{-16.843275\du}{20.217724\du}}
\pgfpathlineto{\pgfpoint{-16.847753\du}{20.215336\du}}
\pgfpathlineto{\pgfpoint{-16.852530\du}{20.212350\du}}
\pgfpathlineto{\pgfpoint{-16.858202\du}{20.210261\du}}
\pgfpathlineto{\pgfpoint{-16.863576\du}{20.207574\du}}
\pgfpathlineto{\pgfpoint{-16.869846\du}{20.204887\du}}
\pgfpathlineto{\pgfpoint{-16.875817\du}{20.203095\du}}
\pgfpathlineto{\pgfpoint{-16.883280\du}{20.200707\du}}
\pgfpathlineto{\pgfpoint{-16.890147\du}{20.199214\du}}
\pgfpathlineto{\pgfpoint{-16.897312\du}{20.197124\du}}
\pgfpathlineto{\pgfpoint{-16.905075\du}{20.195333\du}}
\pgfpathlineto{\pgfpoint{-16.913733\du}{20.193542\du}}
\pgfpathlineto{\pgfpoint{-16.922092\du}{20.192347\du}}
\pgfpathlineto{\pgfpoint{-16.930452\du}{20.189959\du}}
\pgfpathlineto{\pgfpoint{-16.939408\du}{20.189660\du}}
\pgfpathlineto{\pgfpoint{-16.948066\du}{20.188168\du}}
\pgfpathlineto{\pgfpoint{-16.957023\du}{20.187272\du}}
\pgfpathlineto{\pgfpoint{-16.966278\du}{20.186376\du}}
\pgfpathlineto{\pgfpoint{-16.975832\du}{20.185779\du}}
\pgfpathlineto{\pgfpoint{-16.985385\du}{20.185481\du}}
\pgfpathlineto{\pgfpoint{-16.994342\du}{20.184884\du}}
\pgfpathlineto{\pgfpoint{-17.004493\du}{20.184585\du}}
\pgfpathlineto{\pgfpoint{-17.014046\du}{20.184585\du}}
\pgfpathlineto{\pgfpoint{-17.014046\du}{20.184585\du}}
\pgfpathlineto{\pgfpoint{-17.023301\du}{20.184585\du}}
\pgfpathlineto{\pgfpoint{-17.033154\du}{20.184884\du}}
\pgfpathlineto{\pgfpoint{-17.042707\du}{20.185481\du}}
\pgfpathlineto{\pgfpoint{-17.052261\du}{20.185779\du}}
\pgfpathlineto{\pgfpoint{-17.061516\du}{20.186376\du}}
\pgfpathlineto{\pgfpoint{-17.071368\du}{20.187272\du}}
\pgfpathlineto{\pgfpoint{-17.080026\du}{20.188168\du}}
\pgfpathlineto{\pgfpoint{-17.088983\du}{20.189660\du}}
\pgfpathlineto{\pgfpoint{-17.097342\du}{20.189959\du}}
\pgfpathlineto{\pgfpoint{-17.106299\du}{20.192347\du}}
\pgfpathlineto{\pgfpoint{-17.114658\du}{20.193542\du}}
\pgfpathlineto{\pgfpoint{-17.122719\du}{20.195333\du}}
\pgfpathlineto{\pgfpoint{-17.130482\du}{20.197124\du}}
\pgfpathlineto{\pgfpoint{-17.137945\du}{20.199214\du}}
\pgfpathlineto{\pgfpoint{-17.144812\du}{20.200707\du}}
\pgfpathlineto{\pgfpoint{-17.152276\du}{20.203095\du}}
\pgfpathlineto{\pgfpoint{-17.158247\du}{20.204887\du}}
\pgfpathlineto{\pgfpoint{-17.164516\du}{20.207574\du}}
\pgfpathlineto{\pgfpoint{-17.170189\du}{20.210261\du}}
\pgfpathlineto{\pgfpoint{-17.175563\du}{20.212350\du}}
\pgfpathlineto{\pgfpoint{-17.180340\du}{20.215336\du}}
\pgfpathlineto{\pgfpoint{-17.184818\du}{20.217724\du}}
\pgfpathlineto{\pgfpoint{-17.188998\du}{20.220411\du}}
\pgfpathlineto{\pgfpoint{-17.192282\du}{20.223098\du}}
\pgfpathlineto{\pgfpoint{-17.195864\du}{20.225487\du}}
\pgfpathlineto{\pgfpoint{-17.197656\du}{20.229069\du}}
\pgfpathlineto{\pgfpoint{-17.200044\du}{20.231756\du}}
\pgfpathlineto{\pgfpoint{-17.201835\du}{20.234742\du}}
\pgfpathlineto{\pgfpoint{-17.203328\du}{20.238026\du}}
\pgfpathlineto{\pgfpoint{-17.203925\du}{20.240414\du}}
\pgfpathlineto{\pgfpoint{-17.204224\du}{20.243997\du}}
\pgfpathlineto{\pgfpoint{-17.204224\du}{20.243997\du}}
\pgfpathlineto{\pgfpoint{-17.203925\du}{20.246982\du}}
\pgfpathlineto{\pgfpoint{-17.203328\du}{20.249669\du}}
\pgfpathlineto{\pgfpoint{-17.201835\du}{20.252655\du}}
\pgfpathlineto{\pgfpoint{-17.200044\du}{20.256238\du}}
\pgfpathlineto{\pgfpoint{-17.197656\du}{20.258924\du}}
\pgfpathlineto{\pgfpoint{-17.195864\du}{20.261910\du}}
\pgfpathlineto{\pgfpoint{-17.192282\du}{20.264597\du}}
\pgfpathlineto{\pgfpoint{-17.188998\du}{20.267284\du}}
\pgfpathlineto{\pgfpoint{-17.184818\du}{20.269971\du}}
\pgfpathlineto{\pgfpoint{-17.180340\du}{20.272658\du}}
\pgfpathlineto{\pgfpoint{-17.175563\du}{20.275345\du}}
\pgfpathlineto{\pgfpoint{-17.170189\du}{20.277733\du}}
\pgfpathlineto{\pgfpoint{-17.164516\du}{20.280420\du}}
\pgfpathlineto{\pgfpoint{-17.158247\du}{20.282510\du}}
\pgfpathlineto{\pgfpoint{-17.152276\du}{20.284899\du}}
\pgfpathlineto{\pgfpoint{-17.144812\du}{20.287287\du}}
\pgfpathlineto{\pgfpoint{-17.137945\du}{20.289078\du}}
\pgfpathlineto{\pgfpoint{-17.130482\du}{20.290571\du}}
\pgfpathlineto{\pgfpoint{-17.122719\du}{20.292661\du}}
\pgfpathlineto{\pgfpoint{-17.114658\du}{20.294452\du}}
\pgfpathlineto{\pgfpoint{-17.106299\du}{20.295646\du}}
\pgfpathlineto{\pgfpoint{-17.097342\du}{20.297438\du}}
\pgfpathlineto{\pgfpoint{-17.088983\du}{20.298333\du}}
\pgfpathlineto{\pgfpoint{-17.080026\du}{20.299826\du}}
\pgfpathlineto{\pgfpoint{-17.071368\du}{20.300722\du}}
\pgfpathlineto{\pgfpoint{-17.061516\du}{20.301020\du}}
\pgfpathlineto{\pgfpoint{-17.052261\du}{20.302215\du}}
\pgfpathlineto{\pgfpoint{-17.042707\du}{20.302513\du}}
\pgfpathlineto{\pgfpoint{-17.033154\du}{20.302812\du}}
\pgfpathlineto{\pgfpoint{-17.023301\du}{20.303409\du}}
\pgfpathlineto{\pgfpoint{-17.014046\du}{20.303409\du}}
\pgfpathlineto{\pgfpoint{-17.014046\du}{20.303409\du}}
\pgfpathlineto{\pgfpoint{-17.004493\du}{20.303409\du}}
\pgfpathlineto{\pgfpoint{-16.994342\du}{20.302812\du}}
\pgfpathlineto{\pgfpoint{-16.985385\du}{20.302513\du}}
\pgfpathlineto{\pgfpoint{-16.975832\du}{20.302215\du}}
\pgfpathlineto{\pgfpoint{-16.966278\du}{20.301020\du}}
\pgfpathlineto{\pgfpoint{-16.957023\du}{20.300722\du}}
\pgfpathlineto{\pgfpoint{-16.948066\du}{20.299826\du}}
\pgfpathlineto{\pgfpoint{-16.939408\du}{20.298333\du}}
\pgfpathlineto{\pgfpoint{-16.930452\du}{20.297438\du}}
\pgfpathlineto{\pgfpoint{-16.922092\du}{20.295646\du}}
\pgfpathlineto{\pgfpoint{-16.913733\du}{20.294452\du}}
\pgfpathlineto{\pgfpoint{-16.905075\du}{20.292661\du}}
\pgfpathlineto{\pgfpoint{-16.897312\du}{20.290571\du}}
\pgfpathlineto{\pgfpoint{-16.890147\du}{20.289078\du}}
\pgfpathlineto{\pgfpoint{-16.883280\du}{20.287287\du}}
\pgfpathlineto{\pgfpoint{-16.875817\du}{20.284899\du}}
\pgfpathlineto{\pgfpoint{-16.869846\du}{20.282510\du}}
\pgfpathlineto{\pgfpoint{-16.863576\du}{20.280420\du}}
\pgfpathlineto{\pgfpoint{-16.858202\du}{20.277733\du}}
\pgfpathlineto{\pgfpoint{-16.852530\du}{20.275345\du}}
\pgfpathlineto{\pgfpoint{-16.847753\du}{20.272658\du}}
\pgfpathlineto{\pgfpoint{-16.843275\du}{20.269971\du}}
\pgfpathlineto{\pgfpoint{-16.839095\du}{20.267284\du}}
\pgfpathlineto{\pgfpoint{-16.835811\du}{20.264597\du}}
\pgfpathlineto{\pgfpoint{-16.832527\du}{20.261910\du}}
\pgfpathlineto{\pgfpoint{-16.830138\du}{20.258924\du}}
\pgfpathlineto{\pgfpoint{-16.827750\du}{20.256238\du}}
\pgfpathlineto{\pgfpoint{-16.826257\du}{20.252655\du}}
\pgfpathlineto{\pgfpoint{-16.825063\du}{20.249669\du}}
\pgfpathlineto{\pgfpoint{-16.824167\du}{20.246982\du}}
\pgfpathlineto{\pgfpoint{-16.824167\du}{20.243997\du}}
\pgfusepath{stroke}
\pgfsetbuttcap
\pgfsetmiterjoin
\pgfsetdash{}{0pt}
\definecolor{dialinecolor}{rgb}{0.000000, 0.470588, 0.666667}
\pgfsetfillcolor{dialinecolor}
\pgfpathmoveto{\pgfpoint{-17.204224\du}{20.160701\du}}
\pgfpathlineto{\pgfpoint{-17.204224\du}{20.255342\du}}
\pgfpathlineto{\pgfpoint{-16.811628\du}{20.255342\du}}
\pgfpathlineto{\pgfpoint{-16.811628\du}{20.160701\du}}
\pgfpathlineto{\pgfpoint{-17.204224\du}{20.160701\du}}
\pgfusepath{fill}
\pgfsetbuttcap
\pgfsetmiterjoin
\pgfsetdash{}{0pt}
\definecolor{dialinecolor}{rgb}{0.000000, 0.705882, 1.000000}
\pgfsetfillcolor{dialinecolor}
\pgfpathmoveto{\pgfpoint{-16.811927\du}{20.154133\du}}
\pgfpathlineto{\pgfpoint{-16.812524\du}{20.150849\du}}
\pgfpathlineto{\pgfpoint{-16.813419\du}{20.147266\du}}
\pgfpathlineto{\pgfpoint{-16.814315\du}{20.144281\du}}
\pgfpathlineto{\pgfpoint{-16.816106\du}{20.140996\du}}
\pgfpathlineto{\pgfpoint{-16.818495\du}{20.137712\du}}
\pgfpathlineto{\pgfpoint{-16.821182\du}{20.134428\du}}
\pgfpathlineto{\pgfpoint{-16.824167\du}{20.131443\du}}
\pgfpathlineto{\pgfpoint{-16.827750\du}{20.128159\du}}
\pgfpathlineto{\pgfpoint{-16.831930\du}{20.125472\du}}
\pgfpathlineto{\pgfpoint{-16.836706\du}{20.122486\du}}
\pgfpathlineto{\pgfpoint{-16.841782\du}{20.119501\du}}
\pgfpathlineto{\pgfpoint{-16.847454\du}{20.116814\du}}
\pgfpathlineto{\pgfpoint{-16.853127\du}{20.113530\du}}
\pgfpathlineto{\pgfpoint{-16.859396\du}{20.111738\du}}
\pgfpathlineto{\pgfpoint{-16.865964\du}{20.109051\du}}
\pgfpathlineto{\pgfpoint{-16.873130\du}{20.106663\du}}
\pgfpathlineto{\pgfpoint{-16.880594\du}{20.104573\du}}
\pgfpathlineto{\pgfpoint{-16.888057\du}{20.102483\du}}
\pgfpathlineto{\pgfpoint{-16.896417\du}{20.100393\du}}
\pgfpathlineto{\pgfpoint{-16.904776\du}{20.098304\du}}
\pgfpathlineto{\pgfpoint{-16.912837\du}{20.096811\du}}
\pgfpathlineto{\pgfpoint{-16.922092\du}{20.095318\du}}
\pgfpathlineto{\pgfpoint{-16.930750\du}{20.094124\du}}
\pgfpathlineto{\pgfpoint{-16.940304\du}{20.092631\du}}
\pgfpathlineto{\pgfpoint{-16.949559\du}{20.091735\du}}
\pgfpathlineto{\pgfpoint{-16.958814\du}{20.090541\du}}
\pgfpathlineto{\pgfpoint{-16.968666\du}{20.089944\du}}
\pgfpathlineto{\pgfpoint{-16.978519\du}{20.089347\du}}
\pgfpathlineto{\pgfpoint{-16.988371\du}{20.089048\du}}
\pgfpathlineto{\pgfpoint{-16.998522\du}{20.088750\du}}
\pgfpathlineto{\pgfpoint{-17.008374\du}{20.088750\du}}
\pgfpathlineto{\pgfpoint{-17.008374\du}{20.088750\du}}
\pgfpathlineto{\pgfpoint{-17.018226\du}{20.088750\du}}
\pgfpathlineto{\pgfpoint{-17.028078\du}{20.089048\du}}
\pgfpathlineto{\pgfpoint{-17.038229\du}{20.089347\du}}
\pgfpathlineto{\pgfpoint{-17.048081\du}{20.089944\du}}
\pgfpathlineto{\pgfpoint{-17.057635\du}{20.090541\du}}
\pgfpathlineto{\pgfpoint{-17.067189\du}{20.091735\du}}
\pgfpathlineto{\pgfpoint{-17.076444\du}{20.092631\du}}
\pgfpathlineto{\pgfpoint{-17.085699\du}{20.094124\du}}
\pgfpathlineto{\pgfpoint{-17.094655\du}{20.095318\du}}
\pgfpathlineto{\pgfpoint{-17.103612\du}{20.096811\du}}
\pgfpathlineto{\pgfpoint{-17.111971\du}{20.098304\du}}
\pgfpathlineto{\pgfpoint{-17.120331\du}{20.100393\du}}
\pgfpathlineto{\pgfpoint{-17.128690\du}{20.102483\du}}
\pgfpathlineto{\pgfpoint{-17.135856\du}{20.104573\du}}
\pgfpathlineto{\pgfpoint{-17.143319\du}{20.106663\du}}
\pgfpathlineto{\pgfpoint{-17.150485\du}{20.109051\du}}
\pgfpathlineto{\pgfpoint{-17.157351\du}{20.111738\du}}
\pgfpathlineto{\pgfpoint{-17.163322\du}{20.113530\du}}
\pgfpathlineto{\pgfpoint{-17.169293\du}{20.116814\du}}
\pgfpathlineto{\pgfpoint{-17.174966\du}{20.119501\du}}
\pgfpathlineto{\pgfpoint{-17.180041\du}{20.122486\du}}
\pgfpathlineto{\pgfpoint{-17.184818\du}{20.125472\du}}
\pgfpathlineto{\pgfpoint{-17.188998\du}{20.128159\du}}
\pgfpathlineto{\pgfpoint{-17.192282\du}{20.131443\du}}
\pgfpathlineto{\pgfpoint{-17.195566\du}{20.134428\du}}
\pgfpathlineto{\pgfpoint{-17.198551\du}{20.137712\du}}
\pgfpathlineto{\pgfpoint{-17.200641\du}{20.140996\du}}
\pgfpathlineto{\pgfpoint{-17.202134\du}{20.144281\du}}
\pgfpathlineto{\pgfpoint{-17.203925\du}{20.147266\du}}
\pgfpathlineto{\pgfpoint{-17.204224\du}{20.150849\du}}
\pgfpathlineto{\pgfpoint{-17.204522\du}{20.154133\du}}
\pgfpathlineto{\pgfpoint{-17.204522\du}{20.154133\du}}
\pgfpathlineto{\pgfpoint{-17.204224\du}{20.157417\du}}
\pgfpathlineto{\pgfpoint{-17.203925\du}{20.160999\du}}
\pgfpathlineto{\pgfpoint{-17.202134\du}{20.163985\du}}
\pgfpathlineto{\pgfpoint{-17.200641\du}{20.167269\du}}
\pgfpathlineto{\pgfpoint{-17.198551\du}{20.170852\du}}
\pgfpathlineto{\pgfpoint{-17.195566\du}{20.173837\du}}
\pgfpathlineto{\pgfpoint{-17.192282\du}{20.176823\du}}
\pgfpathlineto{\pgfpoint{-17.188998\du}{20.180107\du}}
\pgfpathlineto{\pgfpoint{-17.184818\du}{20.182495\du}}
\pgfpathlineto{\pgfpoint{-17.180041\du}{20.185779\du}}
\pgfpathlineto{\pgfpoint{-17.174966\du}{20.189063\du}}
\pgfpathlineto{\pgfpoint{-17.169293\du}{20.191750\du}}
\pgfpathlineto{\pgfpoint{-17.163322\du}{20.194437\du}}
\pgfpathlineto{\pgfpoint{-17.157351\du}{20.196527\du}}
\pgfpathlineto{\pgfpoint{-17.150485\du}{20.199513\du}}
\pgfpathlineto{\pgfpoint{-17.143319\du}{20.201901\du}}
\pgfpathlineto{\pgfpoint{-17.135856\du}{20.203692\du}}
\pgfpathlineto{\pgfpoint{-17.128690\du}{20.205782\du}}
\pgfpathlineto{\pgfpoint{-17.120331\du}{20.207872\du}}
\pgfpathlineto{\pgfpoint{-17.111971\du}{20.209962\du}}
\pgfpathlineto{\pgfpoint{-17.103612\du}{20.211455\du}}
\pgfpathlineto{\pgfpoint{-17.094655\du}{20.212947\du}}
\pgfpathlineto{\pgfpoint{-17.085699\du}{20.214142\du}}
\pgfpathlineto{\pgfpoint{-17.076444\du}{20.215634\du}}
\pgfpathlineto{\pgfpoint{-17.067189\du}{20.216530\du}}
\pgfpathlineto{\pgfpoint{-17.057635\du}{20.217724\du}}
\pgfpathlineto{\pgfpoint{-17.048081\du}{20.218023\du}}
\pgfpathlineto{\pgfpoint{-17.038229\du}{20.218919\du}}
\pgfpathlineto{\pgfpoint{-17.028078\du}{20.218919\du}}
\pgfpathlineto{\pgfpoint{-17.018226\du}{20.219814\du}}
\pgfpathlineto{\pgfpoint{-17.008374\du}{20.219814\du}}
\pgfpathlineto{\pgfpoint{-17.008374\du}{20.219814\du}}
\pgfpathlineto{\pgfpoint{-16.998522\du}{20.219814\du}}
\pgfpathlineto{\pgfpoint{-16.988371\du}{20.218919\du}}
\pgfpathlineto{\pgfpoint{-16.978519\du}{20.218919\du}}
\pgfpathlineto{\pgfpoint{-16.968666\du}{20.218023\du}}
\pgfpathlineto{\pgfpoint{-16.958814\du}{20.217724\du}}
\pgfpathlineto{\pgfpoint{-16.949559\du}{20.216530\du}}
\pgfpathlineto{\pgfpoint{-16.940304\du}{20.215634\du}}
\pgfpathlineto{\pgfpoint{-16.930750\du}{20.214142\du}}
\pgfpathlineto{\pgfpoint{-16.922092\du}{20.212947\du}}
\pgfpathlineto{\pgfpoint{-16.912837\du}{20.211455\du}}
\pgfpathlineto{\pgfpoint{-16.904776\du}{20.209962\du}}
\pgfpathlineto{\pgfpoint{-16.896417\du}{20.207872\du}}
\pgfpathlineto{\pgfpoint{-16.888057\du}{20.205782\du}}
\pgfpathlineto{\pgfpoint{-16.880594\du}{20.203692\du}}
\pgfpathlineto{\pgfpoint{-16.873130\du}{20.201901\du}}
\pgfpathlineto{\pgfpoint{-16.865964\du}{20.199513\du}}
\pgfpathlineto{\pgfpoint{-16.859396\du}{20.196527\du}}
\pgfpathlineto{\pgfpoint{-16.853127\du}{20.194437\du}}
\pgfpathlineto{\pgfpoint{-16.847454\du}{20.191750\du}}
\pgfpathlineto{\pgfpoint{-16.841782\du}{20.189063\du}}
\pgfpathlineto{\pgfpoint{-16.836706\du}{20.185779\du}}
\pgfpathlineto{\pgfpoint{-16.831930\du}{20.182495\du}}
\pgfpathlineto{\pgfpoint{-16.827750\du}{20.180107\du}}
\pgfpathlineto{\pgfpoint{-16.824167\du}{20.176823\du}}
\pgfpathlineto{\pgfpoint{-16.821182\du}{20.173837\du}}
\pgfpathlineto{\pgfpoint{-16.818495\du}{20.170852\du}}
\pgfpathlineto{\pgfpoint{-16.816106\du}{20.167269\du}}
\pgfpathlineto{\pgfpoint{-16.814315\du}{20.163985\du}}
\pgfpathlineto{\pgfpoint{-16.813419\du}{20.160999\du}}
\pgfpathlineto{\pgfpoint{-16.812524\du}{20.157417\du}}
\pgfpathlineto{\pgfpoint{-16.811927\du}{20.154133\du}}
\pgfusepath{fill}
\pgfsetbuttcap
\pgfsetmiterjoin
\pgfsetdash{}{0pt}
\definecolor{dialinecolor}{rgb}{0.666667, 0.901961, 1.000000}
\pgfsetstrokecolor{dialinecolor}
\pgfpathmoveto{\pgfpoint{-16.824167\du}{20.148460\du}}
\pgfpathlineto{\pgfpoint{-16.824167\du}{20.145176\du}}
\pgfpathlineto{\pgfpoint{-16.825063\du}{20.142788\du}}
\pgfpathlineto{\pgfpoint{-16.826257\du}{20.139504\du}}
\pgfpathlineto{\pgfpoint{-16.827750\du}{20.135921\du}}
\pgfpathlineto{\pgfpoint{-16.830138\du}{20.133533\du}}
\pgfpathlineto{\pgfpoint{-16.832527\du}{20.130249\du}}
\pgfpathlineto{\pgfpoint{-16.835811\du}{20.127860\du}}
\pgfpathlineto{\pgfpoint{-16.839095\du}{20.125173\du}}
\pgfpathlineto{\pgfpoint{-16.843275\du}{20.122486\du}}
\pgfpathlineto{\pgfpoint{-16.847753\du}{20.119799\du}}
\pgfpathlineto{\pgfpoint{-16.852530\du}{20.117112\du}}
\pgfpathlineto{\pgfpoint{-16.858202\du}{20.114425\du}}
\pgfpathlineto{\pgfpoint{-16.863576\du}{20.112037\du}}
\pgfpathlineto{\pgfpoint{-16.869846\du}{20.109947\du}}
\pgfpathlineto{\pgfpoint{-16.875817\du}{20.107559\du}}
\pgfpathlineto{\pgfpoint{-16.883280\du}{20.105170\du}}
\pgfpathlineto{\pgfpoint{-16.890147\du}{20.103379\du}}
\pgfpathlineto{\pgfpoint{-16.897312\du}{20.101886\du}}
\pgfpathlineto{\pgfpoint{-16.905075\du}{20.099498\du}}
\pgfpathlineto{\pgfpoint{-16.913733\du}{20.098005\du}}
\pgfpathlineto{\pgfpoint{-16.922092\du}{20.096811\du}}
\pgfpathlineto{\pgfpoint{-16.930452\du}{20.095019\du}}
\pgfpathlineto{\pgfpoint{-16.939408\du}{20.094124\du}}
\pgfpathlineto{\pgfpoint{-16.948066\du}{20.092631\du}}
\pgfpathlineto{\pgfpoint{-16.957023\du}{20.091735\du}}
\pgfpathlineto{\pgfpoint{-16.966278\du}{20.091138\du}}
\pgfpathlineto{\pgfpoint{-16.975832\du}{20.090243\du}}
\pgfpathlineto{\pgfpoint{-16.985385\du}{20.089944\du}}
\pgfpathlineto{\pgfpoint{-16.994342\du}{20.089645\du}}
\pgfpathlineto{\pgfpoint{-17.004493\du}{20.089048\du}}
\pgfpathlineto{\pgfpoint{-17.014046\du}{20.089048\du}}
\pgfpathlineto{\pgfpoint{-17.014046\du}{20.089048\du}}
\pgfpathlineto{\pgfpoint{-17.023301\du}{20.089048\du}}
\pgfpathlineto{\pgfpoint{-17.033154\du}{20.089645\du}}
\pgfpathlineto{\pgfpoint{-17.042707\du}{20.089944\du}}
\pgfpathlineto{\pgfpoint{-17.052261\du}{20.090243\du}}
\pgfpathlineto{\pgfpoint{-17.061516\du}{20.091138\du}}
\pgfpathlineto{\pgfpoint{-17.071368\du}{20.091735\du}}
\pgfpathlineto{\pgfpoint{-17.080026\du}{20.092631\du}}
\pgfpathlineto{\pgfpoint{-17.088983\du}{20.094124\du}}
\pgfpathlineto{\pgfpoint{-17.097342\du}{20.095019\du}}
\pgfpathlineto{\pgfpoint{-17.106299\du}{20.096811\du}}
\pgfpathlineto{\pgfpoint{-17.114658\du}{20.098005\du}}
\pgfpathlineto{\pgfpoint{-17.122719\du}{20.099498\du}}
\pgfpathlineto{\pgfpoint{-17.130482\du}{20.101886\du}}
\pgfpathlineto{\pgfpoint{-17.137945\du}{20.103379\du}}
\pgfpathlineto{\pgfpoint{-17.144812\du}{20.105170\du}}
\pgfpathlineto{\pgfpoint{-17.152276\du}{20.107559\du}}
\pgfpathlineto{\pgfpoint{-17.158247\du}{20.109947\du}}
\pgfpathlineto{\pgfpoint{-17.164516\du}{20.112037\du}}
\pgfpathlineto{\pgfpoint{-17.170189\du}{20.114425\du}}
\pgfpathlineto{\pgfpoint{-17.175563\du}{20.117112\du}}
\pgfpathlineto{\pgfpoint{-17.180340\du}{20.119799\du}}
\pgfpathlineto{\pgfpoint{-17.184818\du}{20.122486\du}}
\pgfpathlineto{\pgfpoint{-17.188998\du}{20.125173\du}}
\pgfpathlineto{\pgfpoint{-17.192282\du}{20.127860\du}}
\pgfpathlineto{\pgfpoint{-17.195864\du}{20.130249\du}}
\pgfpathlineto{\pgfpoint{-17.197656\du}{20.133533\du}}
\pgfpathlineto{\pgfpoint{-17.200044\du}{20.135921\du}}
\pgfpathlineto{\pgfpoint{-17.201835\du}{20.139504\du}}
\pgfpathlineto{\pgfpoint{-17.203328\du}{20.142788\du}}
\pgfpathlineto{\pgfpoint{-17.203925\du}{20.145176\du}}
\pgfpathlineto{\pgfpoint{-17.204224\du}{20.148460\du}}
\pgfpathlineto{\pgfpoint{-17.204224\du}{20.148460\du}}
\pgfpathlineto{\pgfpoint{-17.203925\du}{20.151744\du}}
\pgfpathlineto{\pgfpoint{-17.203328\du}{20.154431\du}}
\pgfpathlineto{\pgfpoint{-17.201835\du}{20.157417\du}}
\pgfpathlineto{\pgfpoint{-17.200044\du}{20.160701\du}}
\pgfpathlineto{\pgfpoint{-17.197656\du}{20.163388\du}}
\pgfpathlineto{\pgfpoint{-17.195864\du}{20.166672\du}}
\pgfpathlineto{\pgfpoint{-17.192282\du}{20.169359\du}}
\pgfpathlineto{\pgfpoint{-17.188998\du}{20.172046\du}}
\pgfpathlineto{\pgfpoint{-17.184818\du}{20.174733\du}}
\pgfpathlineto{\pgfpoint{-17.180340\du}{20.177420\du}}
\pgfpathlineto{\pgfpoint{-17.175563\du}{20.180107\du}}
\pgfpathlineto{\pgfpoint{-17.170189\du}{20.182197\du}}
\pgfpathlineto{\pgfpoint{-17.164516\du}{20.184884\du}}
\pgfpathlineto{\pgfpoint{-17.158247\du}{20.187272\du}}
\pgfpathlineto{\pgfpoint{-17.152276\du}{20.189660\du}}
\pgfpathlineto{\pgfpoint{-17.144812\du}{20.191750\du}}
\pgfpathlineto{\pgfpoint{-17.137945\du}{20.193542\du}}
\pgfpathlineto{\pgfpoint{-17.130482\du}{20.195333\du}}
\pgfpathlineto{\pgfpoint{-17.122719\du}{20.197124\du}}
\pgfpathlineto{\pgfpoint{-17.114658\du}{20.199214\du}}
\pgfpathlineto{\pgfpoint{-17.106299\du}{20.200408\du}}
\pgfpathlineto{\pgfpoint{-17.097342\du}{20.202200\du}}
\pgfpathlineto{\pgfpoint{-17.088983\du}{20.203095\du}}
\pgfpathlineto{\pgfpoint{-17.080026\du}{20.204588\du}}
\pgfpathlineto{\pgfpoint{-17.071368\du}{20.204887\du}}
\pgfpathlineto{\pgfpoint{-17.061516\du}{20.206081\du}}
\pgfpathlineto{\pgfpoint{-17.052261\du}{20.206379\du}}
\pgfpathlineto{\pgfpoint{-17.042707\du}{20.207275\du}}
\pgfpathlineto{\pgfpoint{-17.033154\du}{20.207574\du}}
\pgfpathlineto{\pgfpoint{-17.023301\du}{20.207872\du}}
\pgfpathlineto{\pgfpoint{-17.014046\du}{20.207872\du}}
\pgfpathlineto{\pgfpoint{-17.014046\du}{20.207872\du}}
\pgfpathlineto{\pgfpoint{-17.004493\du}{20.207872\du}}
\pgfpathlineto{\pgfpoint{-16.994342\du}{20.207574\du}}
\pgfpathlineto{\pgfpoint{-16.985385\du}{20.207275\du}}
\pgfpathlineto{\pgfpoint{-16.975832\du}{20.206379\du}}
\pgfpathlineto{\pgfpoint{-16.966278\du}{20.206081\du}}
\pgfpathlineto{\pgfpoint{-16.957023\du}{20.204887\du}}
\pgfpathlineto{\pgfpoint{-16.948066\du}{20.204588\du}}
\pgfpathlineto{\pgfpoint{-16.939408\du}{20.203095\du}}
\pgfpathlineto{\pgfpoint{-16.930452\du}{20.202200\du}}
\pgfpathlineto{\pgfpoint{-16.922092\du}{20.200408\du}}
\pgfpathlineto{\pgfpoint{-16.913733\du}{20.199214\du}}
\pgfpathlineto{\pgfpoint{-16.905075\du}{20.197124\du}}
\pgfpathlineto{\pgfpoint{-16.897312\du}{20.195333\du}}
\pgfpathlineto{\pgfpoint{-16.890147\du}{20.193542\du}}
\pgfpathlineto{\pgfpoint{-16.883280\du}{20.191750\du}}
\pgfpathlineto{\pgfpoint{-16.875817\du}{20.189660\du}}
\pgfpathlineto{\pgfpoint{-16.869846\du}{20.187272\du}}
\pgfpathlineto{\pgfpoint{-16.863576\du}{20.184884\du}}
\pgfpathlineto{\pgfpoint{-16.858202\du}{20.182197\du}}
\pgfpathlineto{\pgfpoint{-16.852530\du}{20.180107\du}}
\pgfpathlineto{\pgfpoint{-16.847753\du}{20.177420\du}}
\pgfpathlineto{\pgfpoint{-16.843275\du}{20.174733\du}}
\pgfpathlineto{\pgfpoint{-16.839095\du}{20.172046\du}}
\pgfpathlineto{\pgfpoint{-16.835811\du}{20.169359\du}}
\pgfpathlineto{\pgfpoint{-16.832527\du}{20.166672\du}}
\pgfpathlineto{\pgfpoint{-16.830138\du}{20.163388\du}}
\pgfpathlineto{\pgfpoint{-16.827750\du}{20.160701\du}}
\pgfpathlineto{\pgfpoint{-16.826257\du}{20.157417\du}}
\pgfpathlineto{\pgfpoint{-16.825063\du}{20.154431\du}}
\pgfpathlineto{\pgfpoint{-16.824167\du}{20.151744\du}}
\pgfpathlineto{\pgfpoint{-16.824167\du}{20.148460\du}}
\pgfusepath{stroke}
\pgfsetbuttcap
\pgfsetmiterjoin
\pgfsetdash{}{0pt}
\definecolor{dialinecolor}{rgb}{0.000000, 0.000000, 0.000000}
\pgfsetfillcolor{dialinecolor}
\pgfpathmoveto{\pgfpoint{-17.001806\du}{20.136817\du}}
\pgfpathlineto{\pgfpoint{-16.978519\du}{20.148460\du}}
\pgfpathlineto{\pgfpoint{-16.906568\du}{20.124875\du}}
\pgfpathlineto{\pgfpoint{-16.871040\du}{20.136817\du}}
\pgfpathlineto{\pgfpoint{-16.894625\du}{20.112933\du}}
\pgfpathlineto{\pgfpoint{-16.978519\du}{20.112933\du}}
\pgfpathlineto{\pgfpoint{-16.942692\du}{20.112933\du}}
\pgfpathlineto{\pgfpoint{-17.001806\du}{20.136817\du}}
\pgfusepath{fill}
\pgfsetbuttcap
\pgfsetmiterjoin
\pgfsetdash{}{0pt}
\definecolor{dialinecolor}{rgb}{0.000000, 0.000000, 0.000000}
\pgfsetfillcolor{dialinecolor}
\pgfpathmoveto{\pgfpoint{-17.014046\du}{20.172344\du}}
\pgfpathlineto{\pgfpoint{-17.049873\du}{20.160701\du}}
\pgfpathlineto{\pgfpoint{-17.108986\du}{20.184585\du}}
\pgfpathlineto{\pgfpoint{-17.144812\du}{20.172344\du}}
\pgfpathlineto{\pgfpoint{-17.132571\du}{20.207872\du}}
\pgfpathlineto{\pgfpoint{-17.049873\du}{20.207872\du}}
\pgfpathlineto{\pgfpoint{-17.085400\du}{20.196527\du}}
\pgfpathlineto{\pgfpoint{-17.014046\du}{20.172344\du}}
\pgfusepath{fill}
\pgfsetbuttcap
\pgfsetmiterjoin
\pgfsetdash{}{0pt}
\definecolor{dialinecolor}{rgb}{0.000000, 0.000000, 0.000000}
\pgfsetfillcolor{dialinecolor}
\pgfpathmoveto{\pgfpoint{-17.132571\du}{20.112933\du}}
\pgfpathlineto{\pgfpoint{-17.108986\du}{20.100393\du}}
\pgfpathlineto{\pgfpoint{-17.037333\du}{20.136817\du}}
\pgfpathlineto{\pgfpoint{-17.014046\du}{20.124875\du}}
\pgfpathlineto{\pgfpoint{-17.025690\du}{20.148460\du}}
\pgfpathlineto{\pgfpoint{-17.108986\du}{20.148460\du}}
\pgfpathlineto{\pgfpoint{-17.073160\du}{20.136817\du}}
\pgfpathlineto{\pgfpoint{-17.132571\du}{20.112933\du}}
\pgfusepath{fill}
\pgfsetbuttcap
\pgfsetmiterjoin
\pgfsetdash{}{0pt}
\definecolor{dialinecolor}{rgb}{0.000000, 0.000000, 0.000000}
\pgfsetfillcolor{dialinecolor}
\pgfpathmoveto{\pgfpoint{-16.883280\du}{20.196527\du}}
\pgfpathlineto{\pgfpoint{-16.906568\du}{20.207872\du}}
\pgfpathlineto{\pgfpoint{-16.978519\du}{20.172344\du}}
\pgfpathlineto{\pgfpoint{-17.014046\du}{20.184585\du}}
\pgfpathlineto{\pgfpoint{-16.990162\du}{20.160701\du}}
\pgfpathlineto{\pgfpoint{-16.906568\du}{20.160701\du}}
\pgfpathlineto{\pgfpoint{-16.942692\du}{20.172344\du}}
\pgfpathlineto{\pgfpoint{-16.883280\du}{20.196527\du}}
\pgfusepath{fill}
\pgfsetbuttcap
\pgfsetmiterjoin
\pgfsetdash{}{0pt}
\definecolor{dialinecolor}{rgb}{1.000000, 1.000000, 1.000000}
\pgfsetfillcolor{dialinecolor}
\pgfpathmoveto{\pgfpoint{-17.001806\du}{20.148460\du}}
\pgfpathlineto{\pgfpoint{-16.978519\du}{20.148460\du}}
\pgfpathlineto{\pgfpoint{-16.906568\du}{20.124875\du}}
\pgfpathlineto{\pgfpoint{-16.871040\du}{20.136817\du}}
\pgfpathlineto{\pgfpoint{-16.894625\du}{20.112933\du}}
\pgfpathlineto{\pgfpoint{-16.966278\du}{20.112933\du}}
\pgfpathlineto{\pgfpoint{-16.942692\du}{20.112933\du}}
\pgfpathlineto{\pgfpoint{-17.001806\du}{20.148460\du}}
\pgfusepath{fill}
\pgfsetbuttcap
\pgfsetmiterjoin
\pgfsetdash{}{0pt}
\definecolor{dialinecolor}{rgb}{1.000000, 1.000000, 1.000000}
\pgfsetfillcolor{dialinecolor}
\pgfpathmoveto{\pgfpoint{-17.014046\du}{20.172344\du}}
\pgfpathlineto{\pgfpoint{-17.037333\du}{20.160701\du}}
\pgfpathlineto{\pgfpoint{-17.108986\du}{20.184585\du}}
\pgfpathlineto{\pgfpoint{-17.144812\du}{20.184585\du}}
\pgfpathlineto{\pgfpoint{-17.121226\du}{20.207872\du}}
\pgfpathlineto{\pgfpoint{-17.037333\du}{20.207872\du}}
\pgfpathlineto{\pgfpoint{-17.073160\du}{20.196527\du}}
\pgfpathlineto{\pgfpoint{-17.014046\du}{20.172344\du}}
\pgfusepath{fill}
\pgfsetbuttcap
\pgfsetmiterjoin
\pgfsetdash{}{0pt}
\definecolor{dialinecolor}{rgb}{1.000000, 1.000000, 1.000000}
\pgfsetfillcolor{dialinecolor}
\pgfpathmoveto{\pgfpoint{-17.132571\du}{20.112933\du}}
\pgfpathlineto{\pgfpoint{-17.108986\du}{20.112933\du}}
\pgfpathlineto{\pgfpoint{-17.037333\du}{20.136817\du}}
\pgfpathlineto{\pgfpoint{-17.001806\du}{20.124875\du}}
\pgfpathlineto{\pgfpoint{-17.025690\du}{20.148460\du}}
\pgfpathlineto{\pgfpoint{-17.108986\du}{20.148460\du}}
\pgfpathlineto{\pgfpoint{-17.073160\du}{20.148460\du}}
\pgfpathlineto{\pgfpoint{-17.132571\du}{20.112933\du}}
\pgfusepath{fill}
\pgfsetbuttcap
\pgfsetmiterjoin
\pgfsetdash{}{0pt}
\definecolor{dialinecolor}{rgb}{1.000000, 1.000000, 1.000000}
\pgfsetfillcolor{dialinecolor}
\pgfpathmoveto{\pgfpoint{-16.883280\du}{20.196527\du}}
\pgfpathlineto{\pgfpoint{-16.906568\du}{20.207872\du}}
\pgfpathlineto{\pgfpoint{-16.978519\du}{20.184585\du}}
\pgfpathlineto{\pgfpoint{-17.001806\du}{20.184585\du}}
\pgfpathlineto{\pgfpoint{-16.990162\du}{20.160701\du}}
\pgfpathlineto{\pgfpoint{-16.906568\du}{20.160701\du}}
\pgfpathlineto{\pgfpoint{-16.942692\du}{20.172344\du}}
\pgfpathlineto{\pgfpoint{-16.883280\du}{20.196527\du}}
\pgfusepath{fill}
\pgfsetbuttcap
\pgfsetmiterjoin
\pgfsetdash{}{0pt}
\definecolor{dialinecolor}{rgb}{0.666667, 0.901961, 1.000000}
\pgfsetstrokecolor{dialinecolor}
\pgfpathmoveto{\pgfpoint{-17.204224\du}{20.148460\du}}
\pgfpathlineto{\pgfpoint{-17.204224\du}{20.243400\du}}
\pgfusepath{stroke}
\pgfsetbuttcap
\pgfsetmiterjoin
\pgfsetdash{}{0pt}
\definecolor{dialinecolor}{rgb}{0.666667, 0.901961, 1.000000}
\pgfsetstrokecolor{dialinecolor}
\pgfpathmoveto{\pgfpoint{-16.823570\du}{20.148460\du}}
\pgfpathlineto{\pgfpoint{-16.823570\du}{20.243400\du}}
\pgfusepath{stroke}
\pgfsetbuttcap
\pgfsetmiterjoin
\pgfsetdash{}{0pt}
\definecolor{dialinecolor}{rgb}{0.000000, 0.470588, 0.666667}
\pgfsetfillcolor{dialinecolor}
\pgfpathmoveto{\pgfpoint{-16.097790\du}{20.249669\du}}
\pgfpathlineto{\pgfpoint{-16.097790\du}{20.246385\du}}
\pgfpathlineto{\pgfpoint{-16.098686\du}{20.243101\du}}
\pgfpathlineto{\pgfpoint{-16.099880\du}{20.239817\du}}
\pgfpathlineto{\pgfpoint{-16.101373\du}{20.236533\du}}
\pgfpathlineto{\pgfpoint{-16.103761\du}{20.233548\du}}
\pgfpathlineto{\pgfpoint{-16.106448\du}{20.230264\du}}
\pgfpathlineto{\pgfpoint{-16.109434\du}{20.226979\du}}
\pgfpathlineto{\pgfpoint{-16.113613\du}{20.223994\du}}
\pgfpathlineto{\pgfpoint{-16.117196\du}{20.221008\du}}
\pgfpathlineto{\pgfpoint{-16.121973\du}{20.218023\du}}
\pgfpathlineto{\pgfpoint{-16.127048\du}{20.215336\du}}
\pgfpathlineto{\pgfpoint{-16.132422\du}{20.212350\du}}
\pgfpathlineto{\pgfpoint{-16.138692\du}{20.209962\du}}
\pgfpathlineto{\pgfpoint{-16.144663\du}{20.207275\du}}
\pgfpathlineto{\pgfpoint{-16.151529\du}{20.204887\du}}
\pgfpathlineto{\pgfpoint{-16.158396\du}{20.202498\du}}
\pgfpathlineto{\pgfpoint{-16.165860\du}{20.200408\du}}
\pgfpathlineto{\pgfpoint{-16.173324\du}{20.198020\du}}
\pgfpathlineto{\pgfpoint{-16.181385\du}{20.196527\du}}
\pgfpathlineto{\pgfpoint{-16.190043\du}{20.194437\du}}
\pgfpathlineto{\pgfpoint{-16.198402\du}{20.192646\du}}
\pgfpathlineto{\pgfpoint{-16.207060\du}{20.191452\du}}
\pgfpathlineto{\pgfpoint{-16.216315\du}{20.189959\du}}
\pgfpathlineto{\pgfpoint{-16.225570\du}{20.188765\du}}
\pgfpathlineto{\pgfpoint{-16.234825\du}{20.187571\du}}
\pgfpathlineto{\pgfpoint{-16.244379\du}{20.186675\du}}
\pgfpathlineto{\pgfpoint{-16.254530\du}{20.185779\du}}
\pgfpathlineto{\pgfpoint{-16.264084\du}{20.185182\du}}
\pgfpathlineto{\pgfpoint{-16.273637\du}{20.184884\du}}
\pgfpathlineto{\pgfpoint{-16.283788\du}{20.184585\du}}
\pgfpathlineto{\pgfpoint{-16.293640\du}{20.184585\du}}
\pgfpathlineto{\pgfpoint{-16.293640\du}{20.184585\du}}
\pgfpathlineto{\pgfpoint{-16.303492\du}{20.184585\du}}
\pgfpathlineto{\pgfpoint{-16.313643\du}{20.184884\du}}
\pgfpathlineto{\pgfpoint{-16.323495\du}{20.185182\du}}
\pgfpathlineto{\pgfpoint{-16.333348\du}{20.185779\du}}
\pgfpathlineto{\pgfpoint{-16.342901\du}{20.186675\du}}
\pgfpathlineto{\pgfpoint{-16.352455\du}{20.187571\du}}
\pgfpathlineto{\pgfpoint{-16.362009\du}{20.188765\du}}
\pgfpathlineto{\pgfpoint{-16.370965\du}{20.189959\du}}
\pgfpathlineto{\pgfpoint{-16.380220\du}{20.191452\du}}
\pgfpathlineto{\pgfpoint{-16.388878\du}{20.192646\du}}
\pgfpathlineto{\pgfpoint{-16.397536\du}{20.194437\du}}
\pgfpathlineto{\pgfpoint{-16.405896\du}{20.196527\du}}
\pgfpathlineto{\pgfpoint{-16.413957\du}{20.198020\du}}
\pgfpathlineto{\pgfpoint{-16.421421\du}{20.200408\du}}
\pgfpathlineto{\pgfpoint{-16.428884\du}{20.202498\du}}
\pgfpathlineto{\pgfpoint{-16.436050\du}{20.204887\du}}
\pgfpathlineto{\pgfpoint{-16.442618\du}{20.207275\du}}
\pgfpathlineto{\pgfpoint{-16.448887\du}{20.209962\du}}
\pgfpathlineto{\pgfpoint{-16.454858\du}{20.212350\du}}
\pgfpathlineto{\pgfpoint{-16.460232\du}{20.215336\du}}
\pgfpathlineto{\pgfpoint{-16.465308\du}{20.218023\du}}
\pgfpathlineto{\pgfpoint{-16.470084\du}{20.221008\du}}
\pgfpathlineto{\pgfpoint{-16.473966\du}{20.223994\du}}
\pgfpathlineto{\pgfpoint{-16.477847\du}{20.226979\du}}
\pgfpathlineto{\pgfpoint{-16.480832\du}{20.230264\du}}
\pgfpathlineto{\pgfpoint{-16.483519\du}{20.233548\du}}
\pgfpathlineto{\pgfpoint{-16.485908\du}{20.236533\du}}
\pgfpathlineto{\pgfpoint{-16.487699\du}{20.239817\du}}
\pgfpathlineto{\pgfpoint{-16.489192\du}{20.243101\du}}
\pgfpathlineto{\pgfpoint{-16.490087\du}{20.246385\du}}
\pgfpathlineto{\pgfpoint{-16.490087\du}{20.249669\du}}
\pgfpathlineto{\pgfpoint{-16.490087\du}{20.249669\du}}
\pgfpathlineto{\pgfpoint{-16.490087\du}{20.253252\du}}
\pgfpathlineto{\pgfpoint{-16.489192\du}{20.256238\du}}
\pgfpathlineto{\pgfpoint{-16.487699\du}{20.259522\du}}
\pgfpathlineto{\pgfpoint{-16.485908\du}{20.262806\du}}
\pgfpathlineto{\pgfpoint{-16.483519\du}{20.266090\du}}
\pgfpathlineto{\pgfpoint{-16.480832\du}{20.269374\du}}
\pgfpathlineto{\pgfpoint{-16.477847\du}{20.272359\du}}
\pgfpathlineto{\pgfpoint{-16.473966\du}{20.275345\du}}
\pgfpathlineto{\pgfpoint{-16.470084\du}{20.278330\du}}
\pgfpathlineto{\pgfpoint{-16.465308\du}{20.281614\du}}
\pgfpathlineto{\pgfpoint{-16.460232\du}{20.284301\du}}
\pgfpathlineto{\pgfpoint{-16.454858\du}{20.286988\du}}
\pgfpathlineto{\pgfpoint{-16.448887\du}{20.289675\du}}
\pgfpathlineto{\pgfpoint{-16.442618\du}{20.292362\du}}
\pgfpathlineto{\pgfpoint{-16.436050\du}{20.294751\du}}
\pgfpathlineto{\pgfpoint{-16.428884\du}{20.297139\du}}
\pgfpathlineto{\pgfpoint{-16.421421\du}{20.298930\du}}
\pgfpathlineto{\pgfpoint{-16.413957\du}{20.301020\du}}
\pgfpathlineto{\pgfpoint{-16.405896\du}{20.303110\du}}
\pgfpathlineto{\pgfpoint{-16.397536\du}{20.305200\du}}
\pgfpathlineto{\pgfpoint{-16.388878\du}{20.306693\du}}
\pgfpathlineto{\pgfpoint{-16.380220\du}{20.308186\du}}
\pgfpathlineto{\pgfpoint{-16.370965\du}{20.309678\du}}
\pgfpathlineto{\pgfpoint{-16.362009\du}{20.310873\du}}
\pgfpathlineto{\pgfpoint{-16.352455\du}{20.311768\du}}
\pgfpathlineto{\pgfpoint{-16.342901\du}{20.312962\du}}
\pgfpathlineto{\pgfpoint{-16.333348\du}{20.313560\du}}
\pgfpathlineto{\pgfpoint{-16.323495\du}{20.314157\du}}
\pgfpathlineto{\pgfpoint{-16.313643\du}{20.314754\du}}
\pgfpathlineto{\pgfpoint{-16.303492\du}{20.315052\du}}
\pgfpathlineto{\pgfpoint{-16.293640\du}{20.315052\du}}
\pgfpathlineto{\pgfpoint{-16.293640\du}{20.315052\du}}
\pgfpathlineto{\pgfpoint{-16.283788\du}{20.315052\du}}
\pgfpathlineto{\pgfpoint{-16.273637\du}{20.314754\du}}
\pgfpathlineto{\pgfpoint{-16.264084\du}{20.314157\du}}
\pgfpathlineto{\pgfpoint{-16.254530\du}{20.313560\du}}
\pgfpathlineto{\pgfpoint{-16.244379\du}{20.312962\du}}
\pgfpathlineto{\pgfpoint{-16.234825\du}{20.311768\du}}
\pgfpathlineto{\pgfpoint{-16.225570\du}{20.310873\du}}
\pgfpathlineto{\pgfpoint{-16.216315\du}{20.309678\du}}
\pgfpathlineto{\pgfpoint{-16.207060\du}{20.308186\du}}
\pgfpathlineto{\pgfpoint{-16.198402\du}{20.306693\du}}
\pgfpathlineto{\pgfpoint{-16.190043\du}{20.305200\du}}
\pgfpathlineto{\pgfpoint{-16.181385\du}{20.303110\du}}
\pgfpathlineto{\pgfpoint{-16.173324\du}{20.301020\du}}
\pgfpathlineto{\pgfpoint{-16.165860\du}{20.298930\du}}
\pgfpathlineto{\pgfpoint{-16.158396\du}{20.297139\du}}
\pgfpathlineto{\pgfpoint{-16.151529\du}{20.294751\du}}
\pgfpathlineto{\pgfpoint{-16.144663\du}{20.292362\du}}
\pgfpathlineto{\pgfpoint{-16.138692\du}{20.289675\du}}
\pgfpathlineto{\pgfpoint{-16.132422\du}{20.286988\du}}
\pgfpathlineto{\pgfpoint{-16.127048\du}{20.284301\du}}
\pgfpathlineto{\pgfpoint{-16.121973\du}{20.281614\du}}
\pgfpathlineto{\pgfpoint{-16.117196\du}{20.278330\du}}
\pgfpathlineto{\pgfpoint{-16.113613\du}{20.275345\du}}
\pgfpathlineto{\pgfpoint{-16.109434\du}{20.272359\du}}
\pgfpathlineto{\pgfpoint{-16.106448\du}{20.269374\du}}
\pgfpathlineto{\pgfpoint{-16.103761\du}{20.266090\du}}
\pgfpathlineto{\pgfpoint{-16.101373\du}{20.262806\du}}
\pgfpathlineto{\pgfpoint{-16.099880\du}{20.259522\du}}
\pgfpathlineto{\pgfpoint{-16.098686\du}{20.256238\du}}
\pgfpathlineto{\pgfpoint{-16.097790\du}{20.253252\du}}
\pgfpathlineto{\pgfpoint{-16.097790\du}{20.249669\du}}
\pgfusepath{fill}
\pgfsetbuttcap
\pgfsetmiterjoin
\pgfsetdash{}{0pt}
\definecolor{dialinecolor}{rgb}{0.666667, 0.901961, 1.000000}
\pgfsetstrokecolor{dialinecolor}
\pgfpathmoveto{\pgfpoint{-16.109434\du}{20.243997\du}}
\pgfpathlineto{\pgfpoint{-16.109434\du}{20.240414\du}}
\pgfpathlineto{\pgfpoint{-16.110329\du}{20.238026\du}}
\pgfpathlineto{\pgfpoint{-16.111524\du}{20.234742\du}}
\pgfpathlineto{\pgfpoint{-16.113613\du}{20.231756\du}}
\pgfpathlineto{\pgfpoint{-16.115703\du}{20.229069\du}}
\pgfpathlineto{\pgfpoint{-16.117793\du}{20.225487\du}}
\pgfpathlineto{\pgfpoint{-16.121077\du}{20.223098\du}}
\pgfpathlineto{\pgfpoint{-16.124361\du}{20.220411\du}}
\pgfpathlineto{\pgfpoint{-16.128541\du}{20.217724\du}}
\pgfpathlineto{\pgfpoint{-16.133318\du}{20.215336\du}}
\pgfpathlineto{\pgfpoint{-16.137796\du}{20.212350\du}}
\pgfpathlineto{\pgfpoint{-16.143469\du}{20.210261\du}}
\pgfpathlineto{\pgfpoint{-16.149440\du}{20.207574\du}}
\pgfpathlineto{\pgfpoint{-16.155112\du}{20.204887\du}}
\pgfpathlineto{\pgfpoint{-16.161680\du}{20.203095\du}}
\pgfpathlineto{\pgfpoint{-16.168547\du}{20.200707\du}}
\pgfpathlineto{\pgfpoint{-16.175414\du}{20.199214\du}}
\pgfpathlineto{\pgfpoint{-16.183176\du}{20.197124\du}}
\pgfpathlineto{\pgfpoint{-16.190640\du}{20.195333\du}}
\pgfpathlineto{\pgfpoint{-16.198999\du}{20.193542\du}}
\pgfpathlineto{\pgfpoint{-16.207657\du}{20.192347\du}}
\pgfpathlineto{\pgfpoint{-16.216017\du}{20.189959\du}}
\pgfpathlineto{\pgfpoint{-16.224675\du}{20.189660\du}}
\pgfpathlineto{\pgfpoint{-16.233631\du}{20.188168\du}}
\pgfpathlineto{\pgfpoint{-16.242289\du}{20.187272\du}}
\pgfpathlineto{\pgfpoint{-16.251843\du}{20.186376\du}}
\pgfpathlineto{\pgfpoint{-16.261098\du}{20.185779\du}}
\pgfpathlineto{\pgfpoint{-16.270652\du}{20.185481\du}}
\pgfpathlineto{\pgfpoint{-16.280205\du}{20.184884\du}}
\pgfpathlineto{\pgfpoint{-16.290356\du}{20.184585\du}}
\pgfpathlineto{\pgfpoint{-16.300208\du}{20.184585\du}}
\pgfpathlineto{\pgfpoint{-16.300208\du}{20.184585\du}}
\pgfpathlineto{\pgfpoint{-16.309165\du}{20.184585\du}}
\pgfpathlineto{\pgfpoint{-16.319017\du}{20.184884\du}}
\pgfpathlineto{\pgfpoint{-16.328571\du}{20.185481\du}}
\pgfpathlineto{\pgfpoint{-16.337826\du}{20.185779\du}}
\pgfpathlineto{\pgfpoint{-16.347380\du}{20.186376\du}}
\pgfpathlineto{\pgfpoint{-16.356933\du}{20.187272\du}}
\pgfpathlineto{\pgfpoint{-16.365890\du}{20.188168\du}}
\pgfpathlineto{\pgfpoint{-16.374846\du}{20.189660\du}}
\pgfpathlineto{\pgfpoint{-16.383504\du}{20.189959\du}}
\pgfpathlineto{\pgfpoint{-16.391864\du}{20.192347\du}}
\pgfpathlineto{\pgfpoint{-16.400522\du}{20.193542\du}}
\pgfpathlineto{\pgfpoint{-16.408583\du}{20.195333\du}}
\pgfpathlineto{\pgfpoint{-16.416345\du}{20.197124\du}}
\pgfpathlineto{\pgfpoint{-16.423809\du}{20.199214\du}}
\pgfpathlineto{\pgfpoint{-16.430377\du}{20.200707\du}}
\pgfpathlineto{\pgfpoint{-16.437841\du}{20.203095\du}}
\pgfpathlineto{\pgfpoint{-16.444110\du}{20.204887\du}}
\pgfpathlineto{\pgfpoint{-16.450082\du}{20.207574\du}}
\pgfpathlineto{\pgfpoint{-16.456053\du}{20.210261\du}}
\pgfpathlineto{\pgfpoint{-16.461128\du}{20.212350\du}}
\pgfpathlineto{\pgfpoint{-16.466203\du}{20.215336\du}}
\pgfpathlineto{\pgfpoint{-16.470980\du}{20.217724\du}}
\pgfpathlineto{\pgfpoint{-16.475160\du}{20.220411\du}}
\pgfpathlineto{\pgfpoint{-16.478444\du}{20.223098\du}}
\pgfpathlineto{\pgfpoint{-16.481728\du}{20.225487\du}}
\pgfpathlineto{\pgfpoint{-16.484116\du}{20.229069\du}}
\pgfpathlineto{\pgfpoint{-16.485908\du}{20.231756\du}}
\pgfpathlineto{\pgfpoint{-16.488296\du}{20.234742\du}}
\pgfpathlineto{\pgfpoint{-16.489192\du}{20.238026\du}}
\pgfpathlineto{\pgfpoint{-16.490087\du}{20.240414\du}}
\pgfpathlineto{\pgfpoint{-16.490087\du}{20.243997\du}}
\pgfpathlineto{\pgfpoint{-16.490087\du}{20.243997\du}}
\pgfpathlineto{\pgfpoint{-16.490087\du}{20.246982\du}}
\pgfpathlineto{\pgfpoint{-16.489192\du}{20.249669\du}}
\pgfpathlineto{\pgfpoint{-16.488296\du}{20.252655\du}}
\pgfpathlineto{\pgfpoint{-16.485908\du}{20.256238\du}}
\pgfpathlineto{\pgfpoint{-16.484116\du}{20.258924\du}}
\pgfpathlineto{\pgfpoint{-16.481728\du}{20.261910\du}}
\pgfpathlineto{\pgfpoint{-16.478444\du}{20.264597\du}}
\pgfpathlineto{\pgfpoint{-16.475160\du}{20.267284\du}}
\pgfpathlineto{\pgfpoint{-16.470980\du}{20.269971\du}}
\pgfpathlineto{\pgfpoint{-16.466203\du}{20.272658\du}}
\pgfpathlineto{\pgfpoint{-16.461128\du}{20.275345\du}}
\pgfpathlineto{\pgfpoint{-16.456053\du}{20.277733\du}}
\pgfpathlineto{\pgfpoint{-16.450082\du}{20.280420\du}}
\pgfpathlineto{\pgfpoint{-16.444110\du}{20.282510\du}}
\pgfpathlineto{\pgfpoint{-16.437841\du}{20.284899\du}}
\pgfpathlineto{\pgfpoint{-16.430377\du}{20.287287\du}}
\pgfpathlineto{\pgfpoint{-16.423809\du}{20.289078\du}}
\pgfpathlineto{\pgfpoint{-16.416345\du}{20.290571\du}}
\pgfpathlineto{\pgfpoint{-16.408583\du}{20.292661\du}}
\pgfpathlineto{\pgfpoint{-16.400522\du}{20.294452\du}}
\pgfpathlineto{\pgfpoint{-16.391864\du}{20.295646\du}}
\pgfpathlineto{\pgfpoint{-16.383504\du}{20.297438\du}}
\pgfpathlineto{\pgfpoint{-16.374846\du}{20.298333\du}}
\pgfpathlineto{\pgfpoint{-16.365890\du}{20.299826\du}}
\pgfpathlineto{\pgfpoint{-16.356933\du}{20.300722\du}}
\pgfpathlineto{\pgfpoint{-16.347380\du}{20.301020\du}}
\pgfpathlineto{\pgfpoint{-16.337826\du}{20.302215\du}}
\pgfpathlineto{\pgfpoint{-16.328571\du}{20.302513\du}}
\pgfpathlineto{\pgfpoint{-16.319017\du}{20.302812\du}}
\pgfpathlineto{\pgfpoint{-16.309165\du}{20.303409\du}}
\pgfpathlineto{\pgfpoint{-16.300208\du}{20.303409\du}}
\pgfpathlineto{\pgfpoint{-16.300208\du}{20.303409\du}}
\pgfpathlineto{\pgfpoint{-16.290356\du}{20.303409\du}}
\pgfpathlineto{\pgfpoint{-16.280205\du}{20.302812\du}}
\pgfpathlineto{\pgfpoint{-16.270652\du}{20.302513\du}}
\pgfpathlineto{\pgfpoint{-16.261098\du}{20.302215\du}}
\pgfpathlineto{\pgfpoint{-16.251843\du}{20.301020\du}}
\pgfpathlineto{\pgfpoint{-16.242289\du}{20.300722\du}}
\pgfpathlineto{\pgfpoint{-16.233631\du}{20.299826\du}}
\pgfpathlineto{\pgfpoint{-16.224675\du}{20.298333\du}}
\pgfpathlineto{\pgfpoint{-16.216017\du}{20.297438\du}}
\pgfpathlineto{\pgfpoint{-16.207657\du}{20.295646\du}}
\pgfpathlineto{\pgfpoint{-16.198999\du}{20.294452\du}}
\pgfpathlineto{\pgfpoint{-16.190640\du}{20.292661\du}}
\pgfpathlineto{\pgfpoint{-16.183176\du}{20.290571\du}}
\pgfpathlineto{\pgfpoint{-16.175414\du}{20.289078\du}}
\pgfpathlineto{\pgfpoint{-16.168547\du}{20.287287\du}}
\pgfpathlineto{\pgfpoint{-16.161680\du}{20.284899\du}}
\pgfpathlineto{\pgfpoint{-16.155112\du}{20.282510\du}}
\pgfpathlineto{\pgfpoint{-16.149440\du}{20.280420\du}}
\pgfpathlineto{\pgfpoint{-16.143469\du}{20.277733\du}}
\pgfpathlineto{\pgfpoint{-16.137796\du}{20.275345\du}}
\pgfpathlineto{\pgfpoint{-16.133318\du}{20.272658\du}}
\pgfpathlineto{\pgfpoint{-16.128541\du}{20.269971\du}}
\pgfpathlineto{\pgfpoint{-16.124361\du}{20.267284\du}}
\pgfpathlineto{\pgfpoint{-16.121077\du}{20.264597\du}}
\pgfpathlineto{\pgfpoint{-16.117793\du}{20.261910\du}}
\pgfpathlineto{\pgfpoint{-16.115703\du}{20.258924\du}}
\pgfpathlineto{\pgfpoint{-16.113613\du}{20.256238\du}}
\pgfpathlineto{\pgfpoint{-16.111524\du}{20.252655\du}}
\pgfpathlineto{\pgfpoint{-16.110329\du}{20.249669\du}}
\pgfpathlineto{\pgfpoint{-16.109434\du}{20.246982\du}}
\pgfpathlineto{\pgfpoint{-16.109434\du}{20.243997\du}}
\pgfusepath{stroke}
\pgfsetbuttcap
\pgfsetmiterjoin
\pgfsetdash{}{0pt}
\definecolor{dialinecolor}{rgb}{0.000000, 0.470588, 0.666667}
\pgfsetfillcolor{dialinecolor}
\pgfpathmoveto{\pgfpoint{-16.490087\du}{20.160701\du}}
\pgfpathlineto{\pgfpoint{-16.490087\du}{20.255342\du}}
\pgfpathlineto{\pgfpoint{-16.097193\du}{20.255342\du}}
\pgfpathlineto{\pgfpoint{-16.097193\du}{20.160701\du}}
\pgfpathlineto{\pgfpoint{-16.490087\du}{20.160701\du}}
\pgfusepath{fill}
\pgfsetbuttcap
\pgfsetmiterjoin
\pgfsetdash{}{0pt}
\definecolor{dialinecolor}{rgb}{0.000000, 0.705882, 1.000000}
\pgfsetfillcolor{dialinecolor}
\pgfpathmoveto{\pgfpoint{-16.097790\du}{20.154133\du}}
\pgfpathlineto{\pgfpoint{-16.097790\du}{20.150849\du}}
\pgfpathlineto{\pgfpoint{-16.098686\du}{20.147266\du}}
\pgfpathlineto{\pgfpoint{-16.099880\du}{20.144281\du}}
\pgfpathlineto{\pgfpoint{-16.101373\du}{20.140996\du}}
\pgfpathlineto{\pgfpoint{-16.103761\du}{20.137712\du}}
\pgfpathlineto{\pgfpoint{-16.106448\du}{20.134428\du}}
\pgfpathlineto{\pgfpoint{-16.109434\du}{20.131443\du}}
\pgfpathlineto{\pgfpoint{-16.113613\du}{20.128159\du}}
\pgfpathlineto{\pgfpoint{-16.117196\du}{20.125472\du}}
\pgfpathlineto{\pgfpoint{-16.121973\du}{20.122486\du}}
\pgfpathlineto{\pgfpoint{-16.127048\du}{20.119501\du}}
\pgfpathlineto{\pgfpoint{-16.132422\du}{20.116814\du}}
\pgfpathlineto{\pgfpoint{-16.138692\du}{20.113530\du}}
\pgfpathlineto{\pgfpoint{-16.144663\du}{20.111738\du}}
\pgfpathlineto{\pgfpoint{-16.151529\du}{20.109051\du}}
\pgfpathlineto{\pgfpoint{-16.158396\du}{20.106663\du}}
\pgfpathlineto{\pgfpoint{-16.165860\du}{20.104573\du}}
\pgfpathlineto{\pgfpoint{-16.173324\du}{20.102483\du}}
\pgfpathlineto{\pgfpoint{-16.181385\du}{20.100393\du}}
\pgfpathlineto{\pgfpoint{-16.190043\du}{20.098304\du}}
\pgfpathlineto{\pgfpoint{-16.198402\du}{20.096811\du}}
\pgfpathlineto{\pgfpoint{-16.207060\du}{20.095318\du}}
\pgfpathlineto{\pgfpoint{-16.216315\du}{20.094124\du}}
\pgfpathlineto{\pgfpoint{-16.225570\du}{20.092631\du}}
\pgfpathlineto{\pgfpoint{-16.234825\du}{20.091735\du}}
\pgfpathlineto{\pgfpoint{-16.244379\du}{20.090541\du}}
\pgfpathlineto{\pgfpoint{-16.254530\du}{20.089944\du}}
\pgfpathlineto{\pgfpoint{-16.264084\du}{20.089347\du}}
\pgfpathlineto{\pgfpoint{-16.273637\du}{20.089048\du}}
\pgfpathlineto{\pgfpoint{-16.283788\du}{20.088750\du}}
\pgfpathlineto{\pgfpoint{-16.293640\du}{20.088750\du}}
\pgfpathlineto{\pgfpoint{-16.293640\du}{20.088750\du}}
\pgfpathlineto{\pgfpoint{-16.303492\du}{20.088750\du}}
\pgfpathlineto{\pgfpoint{-16.313643\du}{20.089048\du}}
\pgfpathlineto{\pgfpoint{-16.323495\du}{20.089347\du}}
\pgfpathlineto{\pgfpoint{-16.333348\du}{20.089944\du}}
\pgfpathlineto{\pgfpoint{-16.342901\du}{20.090541\du}}
\pgfpathlineto{\pgfpoint{-16.352455\du}{20.091735\du}}
\pgfpathlineto{\pgfpoint{-16.362009\du}{20.092631\du}}
\pgfpathlineto{\pgfpoint{-16.370965\du}{20.094124\du}}
\pgfpathlineto{\pgfpoint{-16.380220\du}{20.095318\du}}
\pgfpathlineto{\pgfpoint{-16.388878\du}{20.096811\du}}
\pgfpathlineto{\pgfpoint{-16.397536\du}{20.098304\du}}
\pgfpathlineto{\pgfpoint{-16.405896\du}{20.100393\du}}
\pgfpathlineto{\pgfpoint{-16.413957\du}{20.102483\du}}
\pgfpathlineto{\pgfpoint{-16.421421\du}{20.104573\du}}
\pgfpathlineto{\pgfpoint{-16.428884\du}{20.106663\du}}
\pgfpathlineto{\pgfpoint{-16.436050\du}{20.109051\du}}
\pgfpathlineto{\pgfpoint{-16.442618\du}{20.111738\du}}
\pgfpathlineto{\pgfpoint{-16.448887\du}{20.113530\du}}
\pgfpathlineto{\pgfpoint{-16.454858\du}{20.116814\du}}
\pgfpathlineto{\pgfpoint{-16.460232\du}{20.119501\du}}
\pgfpathlineto{\pgfpoint{-16.465308\du}{20.122486\du}}
\pgfpathlineto{\pgfpoint{-16.470084\du}{20.125472\du}}
\pgfpathlineto{\pgfpoint{-16.473966\du}{20.128159\du}}
\pgfpathlineto{\pgfpoint{-16.477847\du}{20.131443\du}}
\pgfpathlineto{\pgfpoint{-16.480832\du}{20.134428\du}}
\pgfpathlineto{\pgfpoint{-16.483519\du}{20.137712\du}}
\pgfpathlineto{\pgfpoint{-16.485908\du}{20.140996\du}}
\pgfpathlineto{\pgfpoint{-16.487699\du}{20.144281\du}}
\pgfpathlineto{\pgfpoint{-16.489192\du}{20.147266\du}}
\pgfpathlineto{\pgfpoint{-16.490087\du}{20.150849\du}}
\pgfpathlineto{\pgfpoint{-16.490087\du}{20.154133\du}}
\pgfpathlineto{\pgfpoint{-16.490087\du}{20.154133\du}}
\pgfpathlineto{\pgfpoint{-16.490087\du}{20.157417\du}}
\pgfpathlineto{\pgfpoint{-16.489192\du}{20.160999\du}}
\pgfpathlineto{\pgfpoint{-16.487699\du}{20.163985\du}}
\pgfpathlineto{\pgfpoint{-16.485908\du}{20.167269\du}}
\pgfpathlineto{\pgfpoint{-16.483519\du}{20.170852\du}}
\pgfpathlineto{\pgfpoint{-16.480832\du}{20.173837\du}}
\pgfpathlineto{\pgfpoint{-16.477847\du}{20.176823\du}}
\pgfpathlineto{\pgfpoint{-16.473966\du}{20.180107\du}}
\pgfpathlineto{\pgfpoint{-16.470084\du}{20.182495\du}}
\pgfpathlineto{\pgfpoint{-16.465308\du}{20.185779\du}}
\pgfpathlineto{\pgfpoint{-16.460232\du}{20.189063\du}}
\pgfpathlineto{\pgfpoint{-16.454858\du}{20.191750\du}}
\pgfpathlineto{\pgfpoint{-16.448887\du}{20.194437\du}}
\pgfpathlineto{\pgfpoint{-16.442618\du}{20.196527\du}}
\pgfpathlineto{\pgfpoint{-16.436050\du}{20.199513\du}}
\pgfpathlineto{\pgfpoint{-16.428884\du}{20.201901\du}}
\pgfpathlineto{\pgfpoint{-16.421421\du}{20.203692\du}}
\pgfpathlineto{\pgfpoint{-16.413957\du}{20.205782\du}}
\pgfpathlineto{\pgfpoint{-16.405896\du}{20.207872\du}}
\pgfpathlineto{\pgfpoint{-16.397536\du}{20.209962\du}}
\pgfpathlineto{\pgfpoint{-16.388878\du}{20.211455\du}}
\pgfpathlineto{\pgfpoint{-16.380220\du}{20.212947\du}}
\pgfpathlineto{\pgfpoint{-16.370965\du}{20.214142\du}}
\pgfpathlineto{\pgfpoint{-16.362009\du}{20.215634\du}}
\pgfpathlineto{\pgfpoint{-16.352455\du}{20.216530\du}}
\pgfpathlineto{\pgfpoint{-16.342901\du}{20.217724\du}}
\pgfpathlineto{\pgfpoint{-16.333348\du}{20.218023\du}}
\pgfpathlineto{\pgfpoint{-16.323495\du}{20.218919\du}}
\pgfpathlineto{\pgfpoint{-16.313643\du}{20.218919\du}}
\pgfpathlineto{\pgfpoint{-16.303492\du}{20.219814\du}}
\pgfpathlineto{\pgfpoint{-16.293640\du}{20.219814\du}}
\pgfpathlineto{\pgfpoint{-16.293640\du}{20.219814\du}}
\pgfpathlineto{\pgfpoint{-16.283788\du}{20.219814\du}}
\pgfpathlineto{\pgfpoint{-16.273637\du}{20.218919\du}}
\pgfpathlineto{\pgfpoint{-16.264084\du}{20.218919\du}}
\pgfpathlineto{\pgfpoint{-16.254530\du}{20.218023\du}}
\pgfpathlineto{\pgfpoint{-16.244379\du}{20.217724\du}}
\pgfpathlineto{\pgfpoint{-16.234825\du}{20.216530\du}}
\pgfpathlineto{\pgfpoint{-16.225570\du}{20.215634\du}}
\pgfpathlineto{\pgfpoint{-16.216315\du}{20.214142\du}}
\pgfpathlineto{\pgfpoint{-16.207060\du}{20.212947\du}}
\pgfpathlineto{\pgfpoint{-16.198402\du}{20.211455\du}}
\pgfpathlineto{\pgfpoint{-16.190043\du}{20.209962\du}}
\pgfpathlineto{\pgfpoint{-16.181385\du}{20.207872\du}}
\pgfpathlineto{\pgfpoint{-16.173324\du}{20.205782\du}}
\pgfpathlineto{\pgfpoint{-16.165860\du}{20.203692\du}}
\pgfpathlineto{\pgfpoint{-16.158396\du}{20.201901\du}}
\pgfpathlineto{\pgfpoint{-16.151529\du}{20.199513\du}}
\pgfpathlineto{\pgfpoint{-16.144663\du}{20.196527\du}}
\pgfpathlineto{\pgfpoint{-16.138692\du}{20.194437\du}}
\pgfpathlineto{\pgfpoint{-16.132422\du}{20.191750\du}}
\pgfpathlineto{\pgfpoint{-16.127048\du}{20.189063\du}}
\pgfpathlineto{\pgfpoint{-16.121973\du}{20.185779\du}}
\pgfpathlineto{\pgfpoint{-16.117196\du}{20.182495\du}}
\pgfpathlineto{\pgfpoint{-16.113613\du}{20.180107\du}}
\pgfpathlineto{\pgfpoint{-16.109434\du}{20.176823\du}}
\pgfpathlineto{\pgfpoint{-16.106448\du}{20.173837\du}}
\pgfpathlineto{\pgfpoint{-16.103761\du}{20.170852\du}}
\pgfpathlineto{\pgfpoint{-16.101373\du}{20.167269\du}}
\pgfpathlineto{\pgfpoint{-16.099880\du}{20.163985\du}}
\pgfpathlineto{\pgfpoint{-16.098686\du}{20.160999\du}}
\pgfpathlineto{\pgfpoint{-16.097790\du}{20.157417\du}}
\pgfpathlineto{\pgfpoint{-16.097790\du}{20.154133\du}}
\pgfusepath{fill}
\pgfsetbuttcap
\pgfsetmiterjoin
\pgfsetdash{}{0pt}
\definecolor{dialinecolor}{rgb}{0.666667, 0.901961, 1.000000}
\pgfsetstrokecolor{dialinecolor}
\pgfpathmoveto{\pgfpoint{-16.109434\du}{20.148460\du}}
\pgfpathlineto{\pgfpoint{-16.109434\du}{20.145176\du}}
\pgfpathlineto{\pgfpoint{-16.110329\du}{20.142788\du}}
\pgfpathlineto{\pgfpoint{-16.111524\du}{20.139504\du}}
\pgfpathlineto{\pgfpoint{-16.113613\du}{20.135921\du}}
\pgfpathlineto{\pgfpoint{-16.115703\du}{20.133533\du}}
\pgfpathlineto{\pgfpoint{-16.117793\du}{20.130249\du}}
\pgfpathlineto{\pgfpoint{-16.121077\du}{20.127860\du}}
\pgfpathlineto{\pgfpoint{-16.124361\du}{20.125173\du}}
\pgfpathlineto{\pgfpoint{-16.128541\du}{20.122486\du}}
\pgfpathlineto{\pgfpoint{-16.133318\du}{20.119799\du}}
\pgfpathlineto{\pgfpoint{-16.137796\du}{20.117112\du}}
\pgfpathlineto{\pgfpoint{-16.143469\du}{20.114425\du}}
\pgfpathlineto{\pgfpoint{-16.149440\du}{20.112037\du}}
\pgfpathlineto{\pgfpoint{-16.155112\du}{20.109947\du}}
\pgfpathlineto{\pgfpoint{-16.161680\du}{20.107559\du}}
\pgfpathlineto{\pgfpoint{-16.168547\du}{20.105170\du}}
\pgfpathlineto{\pgfpoint{-16.175414\du}{20.103379\du}}
\pgfpathlineto{\pgfpoint{-16.183176\du}{20.101886\du}}
\pgfpathlineto{\pgfpoint{-16.190640\du}{20.099498\du}}
\pgfpathlineto{\pgfpoint{-16.198999\du}{20.098005\du}}
\pgfpathlineto{\pgfpoint{-16.207657\du}{20.096811\du}}
\pgfpathlineto{\pgfpoint{-16.216017\du}{20.095019\du}}
\pgfpathlineto{\pgfpoint{-16.224675\du}{20.094124\du}}
\pgfpathlineto{\pgfpoint{-16.233631\du}{20.092631\du}}
\pgfpathlineto{\pgfpoint{-16.242289\du}{20.091735\du}}
\pgfpathlineto{\pgfpoint{-16.251843\du}{20.091138\du}}
\pgfpathlineto{\pgfpoint{-16.261098\du}{20.090243\du}}
\pgfpathlineto{\pgfpoint{-16.270652\du}{20.089944\du}}
\pgfpathlineto{\pgfpoint{-16.280205\du}{20.089645\du}}
\pgfpathlineto{\pgfpoint{-16.290356\du}{20.089048\du}}
\pgfpathlineto{\pgfpoint{-16.300208\du}{20.089048\du}}
\pgfpathlineto{\pgfpoint{-16.300208\du}{20.089048\du}}
\pgfpathlineto{\pgfpoint{-16.309165\du}{20.089048\du}}
\pgfpathlineto{\pgfpoint{-16.319017\du}{20.089645\du}}
\pgfpathlineto{\pgfpoint{-16.328571\du}{20.089944\du}}
\pgfpathlineto{\pgfpoint{-16.337826\du}{20.090243\du}}
\pgfpathlineto{\pgfpoint{-16.347380\du}{20.091138\du}}
\pgfpathlineto{\pgfpoint{-16.356933\du}{20.091735\du}}
\pgfpathlineto{\pgfpoint{-16.365890\du}{20.092631\du}}
\pgfpathlineto{\pgfpoint{-16.374846\du}{20.094124\du}}
\pgfpathlineto{\pgfpoint{-16.383504\du}{20.095019\du}}
\pgfpathlineto{\pgfpoint{-16.391864\du}{20.096811\du}}
\pgfpathlineto{\pgfpoint{-16.400522\du}{20.098005\du}}
\pgfpathlineto{\pgfpoint{-16.408583\du}{20.099498\du}}
\pgfpathlineto{\pgfpoint{-16.416345\du}{20.101886\du}}
\pgfpathlineto{\pgfpoint{-16.423809\du}{20.103379\du}}
\pgfpathlineto{\pgfpoint{-16.430377\du}{20.105170\du}}
\pgfpathlineto{\pgfpoint{-16.437841\du}{20.107559\du}}
\pgfpathlineto{\pgfpoint{-16.444110\du}{20.109947\du}}
\pgfpathlineto{\pgfpoint{-16.450082\du}{20.112037\du}}
\pgfpathlineto{\pgfpoint{-16.456053\du}{20.114425\du}}
\pgfpathlineto{\pgfpoint{-16.461128\du}{20.117112\du}}
\pgfpathlineto{\pgfpoint{-16.466203\du}{20.119799\du}}
\pgfpathlineto{\pgfpoint{-16.470980\du}{20.122486\du}}
\pgfpathlineto{\pgfpoint{-16.475160\du}{20.125173\du}}
\pgfpathlineto{\pgfpoint{-16.478444\du}{20.127860\du}}
\pgfpathlineto{\pgfpoint{-16.481728\du}{20.130249\du}}
\pgfpathlineto{\pgfpoint{-16.484116\du}{20.133533\du}}
\pgfpathlineto{\pgfpoint{-16.485908\du}{20.135921\du}}
\pgfpathlineto{\pgfpoint{-16.488296\du}{20.139504\du}}
\pgfpathlineto{\pgfpoint{-16.489192\du}{20.142788\du}}
\pgfpathlineto{\pgfpoint{-16.490087\du}{20.145176\du}}
\pgfpathlineto{\pgfpoint{-16.490087\du}{20.148460\du}}
\pgfpathlineto{\pgfpoint{-16.490087\du}{20.148460\du}}
\pgfpathlineto{\pgfpoint{-16.490087\du}{20.151744\du}}
\pgfpathlineto{\pgfpoint{-16.489192\du}{20.154431\du}}
\pgfpathlineto{\pgfpoint{-16.488296\du}{20.157417\du}}
\pgfpathlineto{\pgfpoint{-16.485908\du}{20.160701\du}}
\pgfpathlineto{\pgfpoint{-16.484116\du}{20.163388\du}}
\pgfpathlineto{\pgfpoint{-16.481728\du}{20.166672\du}}
\pgfpathlineto{\pgfpoint{-16.478444\du}{20.169359\du}}
\pgfpathlineto{\pgfpoint{-16.475160\du}{20.172046\du}}
\pgfpathlineto{\pgfpoint{-16.470980\du}{20.174733\du}}
\pgfpathlineto{\pgfpoint{-16.466203\du}{20.177420\du}}
\pgfpathlineto{\pgfpoint{-16.461128\du}{20.180107\du}}
\pgfpathlineto{\pgfpoint{-16.456053\du}{20.182197\du}}
\pgfpathlineto{\pgfpoint{-16.450082\du}{20.184884\du}}
\pgfpathlineto{\pgfpoint{-16.444110\du}{20.187272\du}}
\pgfpathlineto{\pgfpoint{-16.437841\du}{20.189660\du}}
\pgfpathlineto{\pgfpoint{-16.430377\du}{20.191750\du}}
\pgfpathlineto{\pgfpoint{-16.423809\du}{20.193542\du}}
\pgfpathlineto{\pgfpoint{-16.416345\du}{20.195333\du}}
\pgfpathlineto{\pgfpoint{-16.408583\du}{20.197124\du}}
\pgfpathlineto{\pgfpoint{-16.400522\du}{20.199214\du}}
\pgfpathlineto{\pgfpoint{-16.391864\du}{20.200408\du}}
\pgfpathlineto{\pgfpoint{-16.383504\du}{20.202200\du}}
\pgfpathlineto{\pgfpoint{-16.374846\du}{20.203095\du}}
\pgfpathlineto{\pgfpoint{-16.365890\du}{20.204588\du}}
\pgfpathlineto{\pgfpoint{-16.356933\du}{20.204887\du}}
\pgfpathlineto{\pgfpoint{-16.347380\du}{20.206081\du}}
\pgfpathlineto{\pgfpoint{-16.337826\du}{20.206379\du}}
\pgfpathlineto{\pgfpoint{-16.328571\du}{20.207275\du}}
\pgfpathlineto{\pgfpoint{-16.319017\du}{20.207574\du}}
\pgfpathlineto{\pgfpoint{-16.309165\du}{20.207872\du}}
\pgfpathlineto{\pgfpoint{-16.300208\du}{20.207872\du}}
\pgfpathlineto{\pgfpoint{-16.300208\du}{20.207872\du}}
\pgfpathlineto{\pgfpoint{-16.290356\du}{20.207872\du}}
\pgfpathlineto{\pgfpoint{-16.280205\du}{20.207574\du}}
\pgfpathlineto{\pgfpoint{-16.270652\du}{20.207275\du}}
\pgfpathlineto{\pgfpoint{-16.261098\du}{20.206379\du}}
\pgfpathlineto{\pgfpoint{-16.251843\du}{20.206081\du}}
\pgfpathlineto{\pgfpoint{-16.242289\du}{20.204887\du}}
\pgfpathlineto{\pgfpoint{-16.233631\du}{20.204588\du}}
\pgfpathlineto{\pgfpoint{-16.224675\du}{20.203095\du}}
\pgfpathlineto{\pgfpoint{-16.216017\du}{20.202200\du}}
\pgfpathlineto{\pgfpoint{-16.207657\du}{20.200408\du}}
\pgfpathlineto{\pgfpoint{-16.198999\du}{20.199214\du}}
\pgfpathlineto{\pgfpoint{-16.190640\du}{20.197124\du}}
\pgfpathlineto{\pgfpoint{-16.183176\du}{20.195333\du}}
\pgfpathlineto{\pgfpoint{-16.175414\du}{20.193542\du}}
\pgfpathlineto{\pgfpoint{-16.168547\du}{20.191750\du}}
\pgfpathlineto{\pgfpoint{-16.161680\du}{20.189660\du}}
\pgfpathlineto{\pgfpoint{-16.155112\du}{20.187272\du}}
\pgfpathlineto{\pgfpoint{-16.149440\du}{20.184884\du}}
\pgfpathlineto{\pgfpoint{-16.143469\du}{20.182197\du}}
\pgfpathlineto{\pgfpoint{-16.137796\du}{20.180107\du}}
\pgfpathlineto{\pgfpoint{-16.133318\du}{20.177420\du}}
\pgfpathlineto{\pgfpoint{-16.128541\du}{20.174733\du}}
\pgfpathlineto{\pgfpoint{-16.124361\du}{20.172046\du}}
\pgfpathlineto{\pgfpoint{-16.121077\du}{20.169359\du}}
\pgfpathlineto{\pgfpoint{-16.117793\du}{20.166672\du}}
\pgfpathlineto{\pgfpoint{-16.115703\du}{20.163388\du}}
\pgfpathlineto{\pgfpoint{-16.113613\du}{20.160701\du}}
\pgfpathlineto{\pgfpoint{-16.111524\du}{20.157417\du}}
\pgfpathlineto{\pgfpoint{-16.110329\du}{20.154431\du}}
\pgfpathlineto{\pgfpoint{-16.109434\du}{20.151744\du}}
\pgfpathlineto{\pgfpoint{-16.109434\du}{20.148460\du}}
\pgfusepath{stroke}
\pgfsetbuttcap
\pgfsetmiterjoin
\pgfsetdash{}{0pt}
\definecolor{dialinecolor}{rgb}{0.000000, 0.000000, 0.000000}
\pgfsetfillcolor{dialinecolor}
\pgfpathmoveto{\pgfpoint{-16.287968\du}{20.136817\du}}
\pgfpathlineto{\pgfpoint{-16.264084\du}{20.148460\du}}
\pgfpathlineto{\pgfpoint{-16.192730\du}{20.124875\du}}
\pgfpathlineto{\pgfpoint{-16.156903\du}{20.136817\du}}
\pgfpathlineto{\pgfpoint{-16.180788\du}{20.112933\du}}
\pgfpathlineto{\pgfpoint{-16.264084\du}{20.112933\du}}
\pgfpathlineto{\pgfpoint{-16.228257\du}{20.112933\du}}
\pgfpathlineto{\pgfpoint{-16.287968\du}{20.136817\du}}
\pgfusepath{fill}
\pgfsetbuttcap
\pgfsetmiterjoin
\pgfsetdash{}{0pt}
\definecolor{dialinecolor}{rgb}{0.000000, 0.000000, 0.000000}
\pgfsetfillcolor{dialinecolor}
\pgfpathmoveto{\pgfpoint{-16.300208\du}{20.172344\du}}
\pgfpathlineto{\pgfpoint{-16.335139\du}{20.160701\du}}
\pgfpathlineto{\pgfpoint{-16.395148\du}{20.184585\du}}
\pgfpathlineto{\pgfpoint{-16.430377\du}{20.172344\du}}
\pgfpathlineto{\pgfpoint{-16.419032\du}{20.207872\du}}
\pgfpathlineto{\pgfpoint{-16.335139\du}{20.207872\du}}
\pgfpathlineto{\pgfpoint{-16.371264\du}{20.196527\du}}
\pgfpathlineto{\pgfpoint{-16.300208\du}{20.172344\du}}
\pgfusepath{fill}
\pgfsetbuttcap
\pgfsetmiterjoin
\pgfsetdash{}{0pt}
\definecolor{dialinecolor}{rgb}{0.000000, 0.000000, 0.000000}
\pgfsetfillcolor{dialinecolor}
\pgfpathmoveto{\pgfpoint{-16.419032\du}{20.112933\du}}
\pgfpathlineto{\pgfpoint{-16.395148\du}{20.100393\du}}
\pgfpathlineto{\pgfpoint{-16.323495\du}{20.136817\du}}
\pgfpathlineto{\pgfpoint{-16.300208\du}{20.124875\du}}
\pgfpathlineto{\pgfpoint{-16.311852\du}{20.148460\du}}
\pgfpathlineto{\pgfpoint{-16.395148\du}{20.148460\du}}
\pgfpathlineto{\pgfpoint{-16.359322\du}{20.136817\du}}
\pgfpathlineto{\pgfpoint{-16.419032\du}{20.112933\du}}
\pgfusepath{fill}
\pgfsetbuttcap
\pgfsetmiterjoin
\pgfsetdash{}{0pt}
\definecolor{dialinecolor}{rgb}{0.000000, 0.000000, 0.000000}
\pgfsetfillcolor{dialinecolor}
\pgfpathmoveto{\pgfpoint{-16.168547\du}{20.196527\du}}
\pgfpathlineto{\pgfpoint{-16.192730\du}{20.207872\du}}
\pgfpathlineto{\pgfpoint{-16.264084\du}{20.172344\du}}
\pgfpathlineto{\pgfpoint{-16.300208\du}{20.184585\du}}
\pgfpathlineto{\pgfpoint{-16.275727\du}{20.160701\du}}
\pgfpathlineto{\pgfpoint{-16.192730\du}{20.160701\du}}
\pgfpathlineto{\pgfpoint{-16.228257\du}{20.172344\du}}
\pgfpathlineto{\pgfpoint{-16.168547\du}{20.196527\du}}
\pgfusepath{fill}
\pgfsetbuttcap
\pgfsetmiterjoin
\pgfsetdash{}{0pt}
\definecolor{dialinecolor}{rgb}{1.000000, 1.000000, 1.000000}
\pgfsetfillcolor{dialinecolor}
\pgfpathmoveto{\pgfpoint{-16.287968\du}{20.148460\du}}
\pgfpathlineto{\pgfpoint{-16.264084\du}{20.148460\du}}
\pgfpathlineto{\pgfpoint{-16.192730\du}{20.124875\du}}
\pgfpathlineto{\pgfpoint{-16.156903\du}{20.136817\du}}
\pgfpathlineto{\pgfpoint{-16.180788\du}{20.112933\du}}
\pgfpathlineto{\pgfpoint{-16.251843\du}{20.112933\du}}
\pgfpathlineto{\pgfpoint{-16.228257\du}{20.112933\du}}
\pgfpathlineto{\pgfpoint{-16.287968\du}{20.148460\du}}
\pgfusepath{fill}
\pgfsetbuttcap
\pgfsetmiterjoin
\pgfsetdash{}{0pt}
\definecolor{dialinecolor}{rgb}{1.000000, 1.000000, 1.000000}
\pgfsetfillcolor{dialinecolor}
\pgfpathmoveto{\pgfpoint{-16.300208\du}{20.172344\du}}
\pgfpathlineto{\pgfpoint{-16.323495\du}{20.160701\du}}
\pgfpathlineto{\pgfpoint{-16.395148\du}{20.184585\du}}
\pgfpathlineto{\pgfpoint{-16.430377\du}{20.184585\du}}
\pgfpathlineto{\pgfpoint{-16.406791\du}{20.207872\du}}
\pgfpathlineto{\pgfpoint{-16.323495\du}{20.207872\du}}
\pgfpathlineto{\pgfpoint{-16.359322\du}{20.196527\du}}
\pgfpathlineto{\pgfpoint{-16.300208\du}{20.172344\du}}
\pgfusepath{fill}
\pgfsetbuttcap
\pgfsetmiterjoin
\pgfsetdash{}{0pt}
\definecolor{dialinecolor}{rgb}{1.000000, 1.000000, 1.000000}
\pgfsetfillcolor{dialinecolor}
\pgfpathmoveto{\pgfpoint{-16.419032\du}{20.112933\du}}
\pgfpathlineto{\pgfpoint{-16.395148\du}{20.112933\du}}
\pgfpathlineto{\pgfpoint{-16.323495\du}{20.136817\du}}
\pgfpathlineto{\pgfpoint{-16.287968\du}{20.124875\du}}
\pgfpathlineto{\pgfpoint{-16.311852\du}{20.148460\du}}
\pgfpathlineto{\pgfpoint{-16.395148\du}{20.148460\du}}
\pgfpathlineto{\pgfpoint{-16.359322\du}{20.148460\du}}
\pgfpathlineto{\pgfpoint{-16.419032\du}{20.112933\du}}
\pgfusepath{fill}
\pgfsetbuttcap
\pgfsetmiterjoin
\pgfsetdash{}{0pt}
\definecolor{dialinecolor}{rgb}{1.000000, 1.000000, 1.000000}
\pgfsetfillcolor{dialinecolor}
\pgfpathmoveto{\pgfpoint{-16.168547\du}{20.196527\du}}
\pgfpathlineto{\pgfpoint{-16.192730\du}{20.207872\du}}
\pgfpathlineto{\pgfpoint{-16.264084\du}{20.184585\du}}
\pgfpathlineto{\pgfpoint{-16.287968\du}{20.184585\du}}
\pgfpathlineto{\pgfpoint{-16.275727\du}{20.160701\du}}
\pgfpathlineto{\pgfpoint{-16.192730\du}{20.160701\du}}
\pgfpathlineto{\pgfpoint{-16.228257\du}{20.172344\du}}
\pgfpathlineto{\pgfpoint{-16.168547\du}{20.196527\du}}
\pgfusepath{fill}
\pgfsetbuttcap
\pgfsetmiterjoin
\pgfsetdash{}{0pt}
\definecolor{dialinecolor}{rgb}{0.666667, 0.901961, 1.000000}
\pgfsetstrokecolor{dialinecolor}
\pgfpathmoveto{\pgfpoint{-16.490087\du}{20.148460\du}}
\pgfpathlineto{\pgfpoint{-16.490087\du}{20.243400\du}}
\pgfusepath{stroke}
\pgfsetbuttcap
\pgfsetmiterjoin
\pgfsetdash{}{0pt}
\definecolor{dialinecolor}{rgb}{0.666667, 0.901961, 1.000000}
\pgfsetstrokecolor{dialinecolor}
\pgfpathmoveto{\pgfpoint{-16.109434\du}{20.148460\du}}
\pgfpathlineto{\pgfpoint{-16.109434\du}{20.243400\du}}
\pgfusepath{stroke}
\pgfsetbuttcap
\pgfsetmiterjoin
\pgfsetdash{}{0pt}
\definecolor{dialinecolor}{rgb}{0.000000, 0.470588, 0.666667}
\pgfsetfillcolor{dialinecolor}
\pgfpathmoveto{\pgfpoint{-16.454858\du}{20.654506\du}}
\pgfpathlineto{\pgfpoint{-16.455157\du}{20.650923\du}}
\pgfpathlineto{\pgfpoint{-16.456053\du}{20.647639\du}}
\pgfpathlineto{\pgfpoint{-16.457247\du}{20.644355\du}}
\pgfpathlineto{\pgfpoint{-16.459038\du}{20.641370\du}}
\pgfpathlineto{\pgfpoint{-16.461128\du}{20.638086\du}}
\pgfpathlineto{\pgfpoint{-16.464412\du}{20.634801\du}}
\pgfpathlineto{\pgfpoint{-16.467099\du}{20.631816\du}}
\pgfpathlineto{\pgfpoint{-16.470980\du}{20.628532\du}}
\pgfpathlineto{\pgfpoint{-16.475160\du}{20.625845\du}}
\pgfpathlineto{\pgfpoint{-16.479638\du}{20.622561\du}}
\pgfpathlineto{\pgfpoint{-16.485012\du}{20.619874\du}}
\pgfpathlineto{\pgfpoint{-16.490087\du}{20.617187\du}}
\pgfpathlineto{\pgfpoint{-16.496059\du}{20.614500\du}}
\pgfpathlineto{\pgfpoint{-16.502030\du}{20.611813\du}}
\pgfpathlineto{\pgfpoint{-16.509195\du}{20.609425\du}}
\pgfpathlineto{\pgfpoint{-16.516062\du}{20.607036\du}}
\pgfpathlineto{\pgfpoint{-16.523525\du}{20.605245\du}}
\pgfpathlineto{\pgfpoint{-16.530989\du}{20.602856\du}}
\pgfpathlineto{\pgfpoint{-16.539050\du}{20.601065\du}}
\pgfpathlineto{\pgfpoint{-16.547409\du}{20.598677\du}}
\pgfpathlineto{\pgfpoint{-16.555769\du}{20.597184\du}}
\pgfpathlineto{\pgfpoint{-16.564725\du}{20.595990\du}}
\pgfpathlineto{\pgfpoint{-16.573682\du}{20.594497\du}}
\pgfpathlineto{\pgfpoint{-16.582937\du}{20.593303\du}}
\pgfpathlineto{\pgfpoint{-16.592789\du}{20.592109\du}}
\pgfpathlineto{\pgfpoint{-16.602044\du}{20.591213\du}}
\pgfpathlineto{\pgfpoint{-16.611897\du}{20.590616\du}}
\pgfpathlineto{\pgfpoint{-16.621450\du}{20.590019\du}}
\pgfpathlineto{\pgfpoint{-16.631303\du}{20.589422\du}}
\pgfpathlineto{\pgfpoint{-16.641752\du}{20.589123\du}}
\pgfpathlineto{\pgfpoint{-16.651007\du}{20.589123\du}}
\pgfpathlineto{\pgfpoint{-16.651007\du}{20.589123\du}}
\pgfpathlineto{\pgfpoint{-16.660859\du}{20.589123\du}}
\pgfpathlineto{\pgfpoint{-16.671010\du}{20.589422\du}}
\pgfpathlineto{\pgfpoint{-16.680862\du}{20.590019\du}}
\pgfpathlineto{\pgfpoint{-16.690714\du}{20.590616\du}}
\pgfpathlineto{\pgfpoint{-16.700567\du}{20.591213\du}}
\pgfpathlineto{\pgfpoint{-16.709822\du}{20.592109\du}}
\pgfpathlineto{\pgfpoint{-16.719674\du}{20.593303\du}}
\pgfpathlineto{\pgfpoint{-16.728929\du}{20.594497\du}}
\pgfpathlineto{\pgfpoint{-16.737886\du}{20.595990\du}}
\pgfpathlineto{\pgfpoint{-16.746544\du}{20.597184\du}}
\pgfpathlineto{\pgfpoint{-16.754903\du}{20.598677\du}}
\pgfpathlineto{\pgfpoint{-16.763263\du}{20.601065\du}}
\pgfpathlineto{\pgfpoint{-16.771323\du}{20.602856\du}}
\pgfpathlineto{\pgfpoint{-16.778787\du}{20.605245\du}}
\pgfpathlineto{\pgfpoint{-16.786251\du}{20.607036\du}}
\pgfpathlineto{\pgfpoint{-16.793416\du}{20.609425\du}}
\pgfpathlineto{\pgfpoint{-16.800283\du}{20.611813\du}}
\pgfpathlineto{\pgfpoint{-16.806254\du}{20.614500\du}}
\pgfpathlineto{\pgfpoint{-16.812524\du}{20.617187\du}}
\pgfpathlineto{\pgfpoint{-16.817599\du}{20.619874\du}}
\pgfpathlineto{\pgfpoint{-16.822674\du}{20.622561\du}}
\pgfpathlineto{\pgfpoint{-16.827451\du}{20.625845\du}}
\pgfpathlineto{\pgfpoint{-16.831631\du}{20.628532\du}}
\pgfpathlineto{\pgfpoint{-16.835214\du}{20.631816\du}}
\pgfpathlineto{\pgfpoint{-16.838199\du}{20.634801\du}}
\pgfpathlineto{\pgfpoint{-16.840886\du}{20.638086\du}}
\pgfpathlineto{\pgfpoint{-16.843275\du}{20.641370\du}}
\pgfpathlineto{\pgfpoint{-16.845066\du}{20.644355\du}}
\pgfpathlineto{\pgfpoint{-16.846559\du}{20.647639\du}}
\pgfpathlineto{\pgfpoint{-16.847454\du}{20.650923\du}}
\pgfpathlineto{\pgfpoint{-16.847454\du}{20.654506\du}}
\pgfpathlineto{\pgfpoint{-16.847454\du}{20.654506\du}}
\pgfpathlineto{\pgfpoint{-16.847454\du}{20.657790\du}}
\pgfpathlineto{\pgfpoint{-16.846559\du}{20.661074\du}}
\pgfpathlineto{\pgfpoint{-16.845066\du}{20.664358\du}}
\pgfpathlineto{\pgfpoint{-16.843275\du}{20.667642\du}}
\pgfpathlineto{\pgfpoint{-16.840886\du}{20.670628\du}}
\pgfpathlineto{\pgfpoint{-16.838199\du}{20.673912\du}}
\pgfpathlineto{\pgfpoint{-16.835214\du}{20.676897\du}}
\pgfpathlineto{\pgfpoint{-16.831631\du}{20.680181\du}}
\pgfpathlineto{\pgfpoint{-16.827451\du}{20.683167\du}}
\pgfpathlineto{\pgfpoint{-16.822674\du}{20.686152\du}}
\pgfpathlineto{\pgfpoint{-16.817599\du}{20.688839\du}}
\pgfpathlineto{\pgfpoint{-16.812524\du}{20.691526\du}}
\pgfpathlineto{\pgfpoint{-16.806254\du}{20.694213\du}}
\pgfpathlineto{\pgfpoint{-16.800283\du}{20.696602\du}}
\pgfpathlineto{\pgfpoint{-16.793416\du}{20.699289\du}}
\pgfpathlineto{\pgfpoint{-16.786251\du}{20.701677\du}}
\pgfpathlineto{\pgfpoint{-16.778787\du}{20.703767\du}}
\pgfpathlineto{\pgfpoint{-16.771323\du}{20.706155\du}}
\pgfpathlineto{\pgfpoint{-16.763263\du}{20.707648\du}}
\pgfpathlineto{\pgfpoint{-16.754903\du}{20.709738\du}}
\pgfpathlineto{\pgfpoint{-16.746544\du}{20.711529\du}}
\pgfpathlineto{\pgfpoint{-16.737886\du}{20.712724\du}}
\pgfpathlineto{\pgfpoint{-16.728929\du}{20.714216\du}}
\pgfpathlineto{\pgfpoint{-16.719674\du}{20.715411\du}}
\pgfpathlineto{\pgfpoint{-16.709822\du}{20.716605\du}}
\pgfpathlineto{\pgfpoint{-16.700567\du}{20.717500\du}}
\pgfpathlineto{\pgfpoint{-16.690714\du}{20.718098\du}}
\pgfpathlineto{\pgfpoint{-16.680862\du}{20.718993\du}}
\pgfpathlineto{\pgfpoint{-16.671010\du}{20.718993\du}}
\pgfpathlineto{\pgfpoint{-16.660859\du}{20.719590\du}}
\pgfpathlineto{\pgfpoint{-16.651007\du}{20.719590\du}}
\pgfpathlineto{\pgfpoint{-16.651007\du}{20.719590\du}}
\pgfpathlineto{\pgfpoint{-16.641752\du}{20.719590\du}}
\pgfpathlineto{\pgfpoint{-16.631303\du}{20.718993\du}}
\pgfpathlineto{\pgfpoint{-16.621450\du}{20.718993\du}}
\pgfpathlineto{\pgfpoint{-16.611897\du}{20.718098\du}}
\pgfpathlineto{\pgfpoint{-16.602044\du}{20.717500\du}}
\pgfpathlineto{\pgfpoint{-16.592789\du}{20.716605\du}}
\pgfpathlineto{\pgfpoint{-16.582937\du}{20.715411\du}}
\pgfpathlineto{\pgfpoint{-16.573682\du}{20.714216\du}}
\pgfpathlineto{\pgfpoint{-16.564725\du}{20.712724\du}}
\pgfpathlineto{\pgfpoint{-16.555769\du}{20.711529\du}}
\pgfpathlineto{\pgfpoint{-16.547409\du}{20.709738\du}}
\pgfpathlineto{\pgfpoint{-16.539050\du}{20.707648\du}}
\pgfpathlineto{\pgfpoint{-16.530989\du}{20.706155\du}}
\pgfpathlineto{\pgfpoint{-16.523525\du}{20.703767\du}}
\pgfpathlineto{\pgfpoint{-16.516062\du}{20.701677\du}}
\pgfpathlineto{\pgfpoint{-16.509195\du}{20.699289\du}}
\pgfpathlineto{\pgfpoint{-16.502030\du}{20.696602\du}}
\pgfpathlineto{\pgfpoint{-16.496059\du}{20.694213\du}}
\pgfpathlineto{\pgfpoint{-16.490087\du}{20.691526\du}}
\pgfpathlineto{\pgfpoint{-16.485012\du}{20.688839\du}}
\pgfpathlineto{\pgfpoint{-16.479638\du}{20.686152\du}}
\pgfpathlineto{\pgfpoint{-16.475160\du}{20.683167\du}}
\pgfpathlineto{\pgfpoint{-16.470980\du}{20.680181\du}}
\pgfpathlineto{\pgfpoint{-16.467099\du}{20.676897\du}}
\pgfpathlineto{\pgfpoint{-16.464412\du}{20.673912\du}}
\pgfpathlineto{\pgfpoint{-16.461128\du}{20.670628\du}}
\pgfpathlineto{\pgfpoint{-16.459038\du}{20.667642\du}}
\pgfpathlineto{\pgfpoint{-16.457247\du}{20.664358\du}}
\pgfpathlineto{\pgfpoint{-16.456053\du}{20.661074\du}}
\pgfpathlineto{\pgfpoint{-16.455157\du}{20.657790\du}}
\pgfpathlineto{\pgfpoint{-16.454858\du}{20.654506\du}}
\pgfusepath{fill}
\pgfsetbuttcap
\pgfsetmiterjoin
\pgfsetdash{}{0pt}
\definecolor{dialinecolor}{rgb}{0.666667, 0.901961, 1.000000}
\pgfsetstrokecolor{dialinecolor}
\pgfpathmoveto{\pgfpoint{-16.466800\du}{20.648236\du}}
\pgfpathlineto{\pgfpoint{-16.467099\du}{20.645251\du}}
\pgfpathlineto{\pgfpoint{-16.467696\du}{20.642564\du}}
\pgfpathlineto{\pgfpoint{-16.468890\du}{20.639578\du}}
\pgfpathlineto{\pgfpoint{-16.470980\du}{20.635996\du}}
\pgfpathlineto{\pgfpoint{-16.473070\du}{20.633607\du}}
\pgfpathlineto{\pgfpoint{-16.475160\du}{20.630323\du}}
\pgfpathlineto{\pgfpoint{-16.478444\du}{20.627636\du}}
\pgfpathlineto{\pgfpoint{-16.481728\du}{20.624949\du}}
\pgfpathlineto{\pgfpoint{-16.485908\du}{20.621964\du}}
\pgfpathlineto{\pgfpoint{-16.490685\du}{20.619575\du}}
\pgfpathlineto{\pgfpoint{-16.495760\du}{20.616888\du}}
\pgfpathlineto{\pgfpoint{-16.500835\du}{20.614500\du}}
\pgfpathlineto{\pgfpoint{-16.506508\du}{20.612112\du}}
\pgfpathlineto{\pgfpoint{-16.512479\du}{20.609723\du}}
\pgfpathlineto{\pgfpoint{-16.519047\du}{20.607036\du}}
\pgfpathlineto{\pgfpoint{-16.525914\du}{20.605543\du}}
\pgfpathlineto{\pgfpoint{-16.532482\du}{20.603454\du}}
\pgfpathlineto{\pgfpoint{-16.540543\du}{20.601662\du}}
\pgfpathlineto{\pgfpoint{-16.548007\du}{20.600169\du}}
\pgfpathlineto{\pgfpoint{-16.555769\du}{20.598080\du}}
\pgfpathlineto{\pgfpoint{-16.564725\du}{20.596587\du}}
\pgfpathlineto{\pgfpoint{-16.573085\du}{20.595094\du}}
\pgfpathlineto{\pgfpoint{-16.581444\du}{20.593900\du}}
\pgfpathlineto{\pgfpoint{-16.590700\du}{20.592706\du}}
\pgfpathlineto{\pgfpoint{-16.599358\du}{20.591810\du}}
\pgfpathlineto{\pgfpoint{-16.608911\du}{20.590914\du}}
\pgfpathlineto{\pgfpoint{-16.618465\du}{20.590616\du}}
\pgfpathlineto{\pgfpoint{-16.627720\du}{20.590019\du}}
\pgfpathlineto{\pgfpoint{-16.637572\du}{20.589422\du}}
\pgfpathlineto{\pgfpoint{-16.647424\du}{20.589123\du}}
\pgfpathlineto{\pgfpoint{-16.656978\du}{20.589123\du}}
\pgfpathlineto{\pgfpoint{-16.656978\du}{20.589123\du}}
\pgfpathlineto{\pgfpoint{-16.666233\du}{20.589123\du}}
\pgfpathlineto{\pgfpoint{-16.676085\du}{20.589422\du}}
\pgfpathlineto{\pgfpoint{-16.685639\du}{20.590019\du}}
\pgfpathlineto{\pgfpoint{-16.695193\du}{20.590616\du}}
\pgfpathlineto{\pgfpoint{-16.704746\du}{20.590914\du}}
\pgfpathlineto{\pgfpoint{-16.714001\du}{20.591810\du}}
\pgfpathlineto{\pgfpoint{-16.722958\du}{20.592706\du}}
\pgfpathlineto{\pgfpoint{-16.732213\du}{20.593900\du}}
\pgfpathlineto{\pgfpoint{-16.740573\du}{20.595094\du}}
\pgfpathlineto{\pgfpoint{-16.749529\du}{20.596587\du}}
\pgfpathlineto{\pgfpoint{-16.757889\du}{20.598080\du}}
\pgfpathlineto{\pgfpoint{-16.765651\du}{20.600169\du}}
\pgfpathlineto{\pgfpoint{-16.773712\du}{20.601662\du}}
\pgfpathlineto{\pgfpoint{-16.781176\du}{20.603454\du}}
\pgfpathlineto{\pgfpoint{-16.787744\du}{20.605543\du}}
\pgfpathlineto{\pgfpoint{-16.795208\du}{20.607036\du}}
\pgfpathlineto{\pgfpoint{-16.801179\du}{20.609723\du}}
\pgfpathlineto{\pgfpoint{-16.807150\du}{20.612112\du}}
\pgfpathlineto{\pgfpoint{-16.813419\du}{20.614500\du}}
\pgfpathlineto{\pgfpoint{-16.818495\du}{20.616888\du}}
\pgfpathlineto{\pgfpoint{-16.823570\du}{20.619575\du}}
\pgfpathlineto{\pgfpoint{-16.827750\du}{20.621964\du}}
\pgfpathlineto{\pgfpoint{-16.831930\du}{20.624949\du}}
\pgfpathlineto{\pgfpoint{-16.835811\du}{20.627636\du}}
\pgfpathlineto{\pgfpoint{-16.839095\du}{20.630323\du}}
\pgfpathlineto{\pgfpoint{-16.840886\du}{20.633607\du}}
\pgfpathlineto{\pgfpoint{-16.843275\du}{20.635996\du}}
\pgfpathlineto{\pgfpoint{-16.845066\du}{20.639578\du}}
\pgfpathlineto{\pgfpoint{-16.846559\du}{20.642564\du}}
\pgfpathlineto{\pgfpoint{-16.846857\du}{20.645251\du}}
\pgfpathlineto{\pgfpoint{-16.847454\du}{20.648236\du}}
\pgfpathlineto{\pgfpoint{-16.847454\du}{20.648236\du}}
\pgfpathlineto{\pgfpoint{-16.846857\du}{20.651819\du}}
\pgfpathlineto{\pgfpoint{-16.846559\du}{20.654506\du}}
\pgfpathlineto{\pgfpoint{-16.845066\du}{20.657491\du}}
\pgfpathlineto{\pgfpoint{-16.843275\du}{20.660477\du}}
\pgfpathlineto{\pgfpoint{-16.840886\du}{20.663164\du}}
\pgfpathlineto{\pgfpoint{-16.839095\du}{20.666448\du}}
\pgfpathlineto{\pgfpoint{-16.835811\du}{20.669135\du}}
\pgfpathlineto{\pgfpoint{-16.831930\du}{20.672120\du}}
\pgfpathlineto{\pgfpoint{-16.827750\du}{20.674807\du}}
\pgfpathlineto{\pgfpoint{-16.823570\du}{20.677494\du}}
\pgfpathlineto{\pgfpoint{-16.818495\du}{20.680181\du}}
\pgfpathlineto{\pgfpoint{-16.813419\du}{20.681674\du}}
\pgfpathlineto{\pgfpoint{-16.807150\du}{20.684958\du}}
\pgfpathlineto{\pgfpoint{-16.801179\du}{20.687048\du}}
\pgfpathlineto{\pgfpoint{-16.795208\du}{20.689138\du}}
\pgfpathlineto{\pgfpoint{-16.787744\du}{20.691526\du}}
\pgfpathlineto{\pgfpoint{-16.781176\du}{20.693616\du}}
\pgfpathlineto{\pgfpoint{-16.773712\du}{20.695408\du}}
\pgfpathlineto{\pgfpoint{-16.765651\du}{20.696602\du}}
\pgfpathlineto{\pgfpoint{-16.757889\du}{20.698990\du}}
\pgfpathlineto{\pgfpoint{-16.749529\du}{20.700483\du}}
\pgfpathlineto{\pgfpoint{-16.740573\du}{20.701976\du}}
\pgfpathlineto{\pgfpoint{-16.732213\du}{20.703170\du}}
\pgfpathlineto{\pgfpoint{-16.722958\du}{20.704066\du}}
\pgfpathlineto{\pgfpoint{-16.714001\du}{20.704961\du}}
\pgfpathlineto{\pgfpoint{-16.704746\du}{20.706155\du}}
\pgfpathlineto{\pgfpoint{-16.695193\du}{20.706454\du}}
\pgfpathlineto{\pgfpoint{-16.685639\du}{20.707051\du}}
\pgfpathlineto{\pgfpoint{-16.676085\du}{20.707350\du}}
\pgfpathlineto{\pgfpoint{-16.666233\du}{20.707648\du}}
\pgfpathlineto{\pgfpoint{-16.656978\du}{20.707648\du}}
\pgfpathlineto{\pgfpoint{-16.656978\du}{20.707648\du}}
\pgfpathlineto{\pgfpoint{-16.647424\du}{20.707648\du}}
\pgfpathlineto{\pgfpoint{-16.637572\du}{20.707350\du}}
\pgfpathlineto{\pgfpoint{-16.627720\du}{20.707051\du}}
\pgfpathlineto{\pgfpoint{-16.618465\du}{20.706454\du}}
\pgfpathlineto{\pgfpoint{-16.608911\du}{20.706155\du}}
\pgfpathlineto{\pgfpoint{-16.599358\du}{20.704961\du}}
\pgfpathlineto{\pgfpoint{-16.590700\du}{20.704066\du}}
\pgfpathlineto{\pgfpoint{-16.581444\du}{20.703170\du}}
\pgfpathlineto{\pgfpoint{-16.573085\du}{20.701976\du}}
\pgfpathlineto{\pgfpoint{-16.564725\du}{20.700483\du}}
\pgfpathlineto{\pgfpoint{-16.555769\du}{20.698990\du}}
\pgfpathlineto{\pgfpoint{-16.548007\du}{20.696602\du}}
\pgfpathlineto{\pgfpoint{-16.540543\du}{20.695408\du}}
\pgfpathlineto{\pgfpoint{-16.532482\du}{20.693616\du}}
\pgfpathlineto{\pgfpoint{-16.525914\du}{20.691526\du}}
\pgfpathlineto{\pgfpoint{-16.519047\du}{20.689138\du}}
\pgfpathlineto{\pgfpoint{-16.512479\du}{20.687048\du}}
\pgfpathlineto{\pgfpoint{-16.506508\du}{20.684958\du}}
\pgfpathlineto{\pgfpoint{-16.500835\du}{20.681674\du}}
\pgfpathlineto{\pgfpoint{-16.495760\du}{20.680181\du}}
\pgfpathlineto{\pgfpoint{-16.490685\du}{20.677494\du}}
\pgfpathlineto{\pgfpoint{-16.485908\du}{20.674807\du}}
\pgfpathlineto{\pgfpoint{-16.481728\du}{20.672120\du}}
\pgfpathlineto{\pgfpoint{-16.478444\du}{20.669135\du}}
\pgfpathlineto{\pgfpoint{-16.475160\du}{20.666448\du}}
\pgfpathlineto{\pgfpoint{-16.473070\du}{20.663164\du}}
\pgfpathlineto{\pgfpoint{-16.470980\du}{20.660477\du}}
\pgfpathlineto{\pgfpoint{-16.468890\du}{20.657491\du}}
\pgfpathlineto{\pgfpoint{-16.467696\du}{20.654506\du}}
\pgfpathlineto{\pgfpoint{-16.467099\du}{20.651819\du}}
\pgfpathlineto{\pgfpoint{-16.466800\du}{20.648236\du}}
\pgfusepath{stroke}
\pgfsetbuttcap
\pgfsetmiterjoin
\pgfsetdash{}{0pt}
\definecolor{dialinecolor}{rgb}{0.000000, 0.470588, 0.666667}
\pgfsetfillcolor{dialinecolor}
\pgfpathmoveto{\pgfpoint{-16.847454\du}{20.565239\du}}
\pgfpathlineto{\pgfpoint{-16.847454\du}{20.660477\du}}
\pgfpathlineto{\pgfpoint{-16.454560\du}{20.660477\du}}
\pgfpathlineto{\pgfpoint{-16.454560\du}{20.565239\du}}
\pgfpathlineto{\pgfpoint{-16.847454\du}{20.565239\du}}
\pgfusepath{fill}
\pgfsetbuttcap
\pgfsetmiterjoin
\pgfsetdash{}{0pt}
\definecolor{dialinecolor}{rgb}{0.000000, 0.705882, 1.000000}
\pgfsetfillcolor{dialinecolor}
\pgfpathmoveto{\pgfpoint{-16.454858\du}{20.558969\du}}
\pgfpathlineto{\pgfpoint{-16.455157\du}{20.555387\du}}
\pgfpathlineto{\pgfpoint{-16.456053\du}{20.552103\du}}
\pgfpathlineto{\pgfpoint{-16.457247\du}{20.549117\du}}
\pgfpathlineto{\pgfpoint{-16.459038\du}{20.545534\du}}
\pgfpathlineto{\pgfpoint{-16.461128\du}{20.542250\du}}
\pgfpathlineto{\pgfpoint{-16.464412\du}{20.538966\du}}
\pgfpathlineto{\pgfpoint{-16.467099\du}{20.536279\du}}
\pgfpathlineto{\pgfpoint{-16.470980\du}{20.532398\du}}
\pgfpathlineto{\pgfpoint{-16.475160\du}{20.530010\du}}
\pgfpathlineto{\pgfpoint{-16.479638\du}{20.527024\du}}
\pgfpathlineto{\pgfpoint{-16.485012\du}{20.524039\du}}
\pgfpathlineto{\pgfpoint{-16.490087\du}{20.521352\du}}
\pgfpathlineto{\pgfpoint{-16.496059\du}{20.518665\du}}
\pgfpathlineto{\pgfpoint{-16.502030\du}{20.516276\du}}
\pgfpathlineto{\pgfpoint{-16.509195\du}{20.513589\du}}
\pgfpathlineto{\pgfpoint{-16.516062\du}{20.511201\du}}
\pgfpathlineto{\pgfpoint{-16.523525\du}{20.509111\du}}
\pgfpathlineto{\pgfpoint{-16.530989\du}{20.507021\du}}
\pgfpathlineto{\pgfpoint{-16.539050\du}{20.505230\du}}
\pgfpathlineto{\pgfpoint{-16.547409\du}{20.503140\du}}
\pgfpathlineto{\pgfpoint{-16.555769\du}{20.501349\du}}
\pgfpathlineto{\pgfpoint{-16.564725\du}{20.500155\du}}
\pgfpathlineto{\pgfpoint{-16.573682\du}{20.498662\du}}
\pgfpathlineto{\pgfpoint{-16.582937\du}{20.497468\du}}
\pgfpathlineto{\pgfpoint{-16.592789\du}{20.496273\du}}
\pgfpathlineto{\pgfpoint{-16.602044\du}{20.495079\du}}
\pgfpathlineto{\pgfpoint{-16.611897\du}{20.494781\du}}
\pgfpathlineto{\pgfpoint{-16.621450\du}{20.493885\du}}
\pgfpathlineto{\pgfpoint{-16.631303\du}{20.493586\du}}
\pgfpathlineto{\pgfpoint{-16.641752\du}{20.493288\du}}
\pgfpathlineto{\pgfpoint{-16.651007\du}{20.493288\du}}
\pgfpathlineto{\pgfpoint{-16.651007\du}{20.493288\du}}
\pgfpathlineto{\pgfpoint{-16.660859\du}{20.493288\du}}
\pgfpathlineto{\pgfpoint{-16.671010\du}{20.493586\du}}
\pgfpathlineto{\pgfpoint{-16.680862\du}{20.493885\du}}
\pgfpathlineto{\pgfpoint{-16.690714\du}{20.494781\du}}
\pgfpathlineto{\pgfpoint{-16.700567\du}{20.495079\du}}
\pgfpathlineto{\pgfpoint{-16.709822\du}{20.496273\du}}
\pgfpathlineto{\pgfpoint{-16.719674\du}{20.497468\du}}
\pgfpathlineto{\pgfpoint{-16.728929\du}{20.498662\du}}
\pgfpathlineto{\pgfpoint{-16.737886\du}{20.500155\du}}
\pgfpathlineto{\pgfpoint{-16.746544\du}{20.501349\du}}
\pgfpathlineto{\pgfpoint{-16.754903\du}{20.503140\du}}
\pgfpathlineto{\pgfpoint{-16.763263\du}{20.505230\du}}
\pgfpathlineto{\pgfpoint{-16.771323\du}{20.507021\du}}
\pgfpathlineto{\pgfpoint{-16.778787\du}{20.509111\du}}
\pgfpathlineto{\pgfpoint{-16.786251\du}{20.511201\du}}
\pgfpathlineto{\pgfpoint{-16.793416\du}{20.513589\du}}
\pgfpathlineto{\pgfpoint{-16.800283\du}{20.516276\du}}
\pgfpathlineto{\pgfpoint{-16.806254\du}{20.518665\du}}
\pgfpathlineto{\pgfpoint{-16.812524\du}{20.521352\du}}
\pgfpathlineto{\pgfpoint{-16.817599\du}{20.524039\du}}
\pgfpathlineto{\pgfpoint{-16.822674\du}{20.527024\du}}
\pgfpathlineto{\pgfpoint{-16.827451\du}{20.530010\du}}
\pgfpathlineto{\pgfpoint{-16.831631\du}{20.532398\du}}
\pgfpathlineto{\pgfpoint{-16.835214\du}{20.536279\du}}
\pgfpathlineto{\pgfpoint{-16.838199\du}{20.538966\du}}
\pgfpathlineto{\pgfpoint{-16.840886\du}{20.542250\du}}
\pgfpathlineto{\pgfpoint{-16.843275\du}{20.545534\du}}
\pgfpathlineto{\pgfpoint{-16.845066\du}{20.549117\du}}
\pgfpathlineto{\pgfpoint{-16.846559\du}{20.552103\du}}
\pgfpathlineto{\pgfpoint{-16.847454\du}{20.555387\du}}
\pgfpathlineto{\pgfpoint{-16.847454\du}{20.558969\du}}
\pgfpathlineto{\pgfpoint{-16.847454\du}{20.558969\du}}
\pgfpathlineto{\pgfpoint{-16.847454\du}{20.562253\du}}
\pgfpathlineto{\pgfpoint{-16.846559\du}{20.565537\du}}
\pgfpathlineto{\pgfpoint{-16.845066\du}{20.568523\du}}
\pgfpathlineto{\pgfpoint{-16.843275\du}{20.572106\du}}
\pgfpathlineto{\pgfpoint{-16.840886\du}{20.575390\du}}
\pgfpathlineto{\pgfpoint{-16.838199\du}{20.578375\du}}
\pgfpathlineto{\pgfpoint{-16.835214\du}{20.581361\du}}
\pgfpathlineto{\pgfpoint{-16.831631\du}{20.584645\du}}
\pgfpathlineto{\pgfpoint{-16.827451\du}{20.587630\du}}
\pgfpathlineto{\pgfpoint{-16.822674\du}{20.590616\du}}
\pgfpathlineto{\pgfpoint{-16.817599\du}{20.593601\du}}
\pgfpathlineto{\pgfpoint{-16.812524\du}{20.596288\du}}
\pgfpathlineto{\pgfpoint{-16.806254\du}{20.598677\du}}
\pgfpathlineto{\pgfpoint{-16.800283\du}{20.601364\du}}
\pgfpathlineto{\pgfpoint{-16.793416\du}{20.604051\du}}
\pgfpathlineto{\pgfpoint{-16.786251\du}{20.606140\du}}
\pgfpathlineto{\pgfpoint{-16.778787\du}{20.608529\du}}
\pgfpathlineto{\pgfpoint{-16.771323\du}{20.610619\du}}
\pgfpathlineto{\pgfpoint{-16.763263\du}{20.612410\du}}
\pgfpathlineto{\pgfpoint{-16.754903\du}{20.614500\du}}
\pgfpathlineto{\pgfpoint{-16.746544\du}{20.616291\du}}
\pgfpathlineto{\pgfpoint{-16.737886\du}{20.617485\du}}
\pgfpathlineto{\pgfpoint{-16.728929\du}{20.618978\du}}
\pgfpathlineto{\pgfpoint{-16.719674\du}{20.620172\du}}
\pgfpathlineto{\pgfpoint{-16.709822\du}{20.621068\du}}
\pgfpathlineto{\pgfpoint{-16.700567\du}{20.621964\du}}
\pgfpathlineto{\pgfpoint{-16.690714\du}{20.622859\du}}
\pgfpathlineto{\pgfpoint{-16.680862\du}{20.623755\du}}
\pgfpathlineto{\pgfpoint{-16.671010\du}{20.624054\du}}
\pgfpathlineto{\pgfpoint{-16.660859\du}{20.624352\du}}
\pgfpathlineto{\pgfpoint{-16.651007\du}{20.624352\du}}
\pgfpathlineto{\pgfpoint{-16.651007\du}{20.624352\du}}
\pgfpathlineto{\pgfpoint{-16.641752\du}{20.624352\du}}
\pgfpathlineto{\pgfpoint{-16.631303\du}{20.624054\du}}
\pgfpathlineto{\pgfpoint{-16.621450\du}{20.623755\du}}
\pgfpathlineto{\pgfpoint{-16.611897\du}{20.622859\du}}
\pgfpathlineto{\pgfpoint{-16.602044\du}{20.621964\du}}
\pgfpathlineto{\pgfpoint{-16.592789\du}{20.621068\du}}
\pgfpathlineto{\pgfpoint{-16.582937\du}{20.620172\du}}
\pgfpathlineto{\pgfpoint{-16.573682\du}{20.618978\du}}
\pgfpathlineto{\pgfpoint{-16.564725\du}{20.617485\du}}
\pgfpathlineto{\pgfpoint{-16.555769\du}{20.616291\du}}
\pgfpathlineto{\pgfpoint{-16.547409\du}{20.614500\du}}
\pgfpathlineto{\pgfpoint{-16.539050\du}{20.612410\du}}
\pgfpathlineto{\pgfpoint{-16.530989\du}{20.610619\du}}
\pgfpathlineto{\pgfpoint{-16.523525\du}{20.608529\du}}
\pgfpathlineto{\pgfpoint{-16.516062\du}{20.606140\du}}
\pgfpathlineto{\pgfpoint{-16.509195\du}{20.604051\du}}
\pgfpathlineto{\pgfpoint{-16.502030\du}{20.601364\du}}
\pgfpathlineto{\pgfpoint{-16.496059\du}{20.598677\du}}
\pgfpathlineto{\pgfpoint{-16.490087\du}{20.596288\du}}
\pgfpathlineto{\pgfpoint{-16.485012\du}{20.593601\du}}
\pgfpathlineto{\pgfpoint{-16.479638\du}{20.590616\du}}
\pgfpathlineto{\pgfpoint{-16.475160\du}{20.587630\du}}
\pgfpathlineto{\pgfpoint{-16.470980\du}{20.584645\du}}
\pgfpathlineto{\pgfpoint{-16.467099\du}{20.581361\du}}
\pgfpathlineto{\pgfpoint{-16.464412\du}{20.578375\du}}
\pgfpathlineto{\pgfpoint{-16.461128\du}{20.575390\du}}
\pgfpathlineto{\pgfpoint{-16.459038\du}{20.572106\du}}
\pgfpathlineto{\pgfpoint{-16.457247\du}{20.568523\du}}
\pgfpathlineto{\pgfpoint{-16.456053\du}{20.565537\du}}
\pgfpathlineto{\pgfpoint{-16.455157\du}{20.562253\du}}
\pgfpathlineto{\pgfpoint{-16.454858\du}{20.558969\du}}
\pgfusepath{fill}
\pgfsetbuttcap
\pgfsetmiterjoin
\pgfsetdash{}{0pt}
\definecolor{dialinecolor}{rgb}{0.666667, 0.901961, 1.000000}
\pgfsetstrokecolor{dialinecolor}
\pgfpathmoveto{\pgfpoint{-16.466800\du}{20.553297\du}}
\pgfpathlineto{\pgfpoint{-16.467099\du}{20.550311\du}}
\pgfpathlineto{\pgfpoint{-16.467696\du}{20.547326\du}}
\pgfpathlineto{\pgfpoint{-16.468890\du}{20.544639\du}}
\pgfpathlineto{\pgfpoint{-16.470980\du}{20.541355\du}}
\pgfpathlineto{\pgfpoint{-16.473070\du}{20.538668\du}}
\pgfpathlineto{\pgfpoint{-16.475160\du}{20.535384\du}}
\pgfpathlineto{\pgfpoint{-16.478444\du}{20.532398\du}}
\pgfpathlineto{\pgfpoint{-16.481728\du}{20.530010\du}}
\pgfpathlineto{\pgfpoint{-16.485908\du}{20.527323\du}}
\pgfpathlineto{\pgfpoint{-16.490685\du}{20.524636\du}}
\pgfpathlineto{\pgfpoint{-16.495760\du}{20.521949\du}}
\pgfpathlineto{\pgfpoint{-16.500835\du}{20.519560\du}}
\pgfpathlineto{\pgfpoint{-16.506508\du}{20.516575\du}}
\pgfpathlineto{\pgfpoint{-16.512479\du}{20.514784\du}}
\pgfpathlineto{\pgfpoint{-16.519047\du}{20.512694\du}}
\pgfpathlineto{\pgfpoint{-16.525914\du}{20.510007\du}}
\pgfpathlineto{\pgfpoint{-16.532482\du}{20.508514\du}}
\pgfpathlineto{\pgfpoint{-16.540543\du}{20.506723\du}}
\pgfpathlineto{\pgfpoint{-16.548007\du}{20.504633\du}}
\pgfpathlineto{\pgfpoint{-16.555769\du}{20.503140\du}}
\pgfpathlineto{\pgfpoint{-16.564725\du}{20.501647\du}}
\pgfpathlineto{\pgfpoint{-16.573085\du}{20.500155\du}}
\pgfpathlineto{\pgfpoint{-16.581444\du}{20.498960\du}}
\pgfpathlineto{\pgfpoint{-16.590700\du}{20.497766\du}}
\pgfpathlineto{\pgfpoint{-16.599358\du}{20.496572\du}}
\pgfpathlineto{\pgfpoint{-16.608911\du}{20.495975\du}}
\pgfpathlineto{\pgfpoint{-16.618465\du}{20.495079\du}}
\pgfpathlineto{\pgfpoint{-16.627720\du}{20.495079\du}}
\pgfpathlineto{\pgfpoint{-16.637572\du}{20.494781\du}}
\pgfpathlineto{\pgfpoint{-16.647424\du}{20.493885\du}}
\pgfpathlineto{\pgfpoint{-16.656978\du}{20.493885\du}}
\pgfpathlineto{\pgfpoint{-16.656978\du}{20.493885\du}}
\pgfpathlineto{\pgfpoint{-16.666233\du}{20.493885\du}}
\pgfpathlineto{\pgfpoint{-16.676085\du}{20.494781\du}}
\pgfpathlineto{\pgfpoint{-16.685639\du}{20.495079\du}}
\pgfpathlineto{\pgfpoint{-16.695193\du}{20.495079\du}}
\pgfpathlineto{\pgfpoint{-16.704746\du}{20.495975\du}}
\pgfpathlineto{\pgfpoint{-16.714001\du}{20.496572\du}}
\pgfpathlineto{\pgfpoint{-16.722958\du}{20.497766\du}}
\pgfpathlineto{\pgfpoint{-16.732213\du}{20.498960\du}}
\pgfpathlineto{\pgfpoint{-16.740573\du}{20.500155\du}}
\pgfpathlineto{\pgfpoint{-16.749529\du}{20.501647\du}}
\pgfpathlineto{\pgfpoint{-16.757889\du}{20.503140\du}}
\pgfpathlineto{\pgfpoint{-16.765651\du}{20.504633\du}}
\pgfpathlineto{\pgfpoint{-16.773712\du}{20.506723\du}}
\pgfpathlineto{\pgfpoint{-16.781176\du}{20.508514\du}}
\pgfpathlineto{\pgfpoint{-16.787744\du}{20.510007\du}}
\pgfpathlineto{\pgfpoint{-16.795208\du}{20.512694\du}}
\pgfpathlineto{\pgfpoint{-16.801179\du}{20.514784\du}}
\pgfpathlineto{\pgfpoint{-16.807150\du}{20.516575\du}}
\pgfpathlineto{\pgfpoint{-16.813419\du}{20.519560\du}}
\pgfpathlineto{\pgfpoint{-16.818495\du}{20.521949\du}}
\pgfpathlineto{\pgfpoint{-16.823570\du}{20.524636\du}}
\pgfpathlineto{\pgfpoint{-16.827750\du}{20.527323\du}}
\pgfpathlineto{\pgfpoint{-16.831930\du}{20.530010\du}}
\pgfpathlineto{\pgfpoint{-16.835811\du}{20.532398\du}}
\pgfpathlineto{\pgfpoint{-16.839095\du}{20.535384\du}}
\pgfpathlineto{\pgfpoint{-16.840886\du}{20.538668\du}}
\pgfpathlineto{\pgfpoint{-16.843275\du}{20.541355\du}}
\pgfpathlineto{\pgfpoint{-16.845066\du}{20.544639\du}}
\pgfpathlineto{\pgfpoint{-16.846559\du}{20.547326\du}}
\pgfpathlineto{\pgfpoint{-16.846857\du}{20.550311\du}}
\pgfpathlineto{\pgfpoint{-16.847454\du}{20.553297\du}}
\pgfpathlineto{\pgfpoint{-16.847454\du}{20.553297\du}}
\pgfpathlineto{\pgfpoint{-16.846857\du}{20.556879\du}}
\pgfpathlineto{\pgfpoint{-16.846559\du}{20.559566\du}}
\pgfpathlineto{\pgfpoint{-16.845066\du}{20.562253\du}}
\pgfpathlineto{\pgfpoint{-16.843275\du}{20.565537\du}}
\pgfpathlineto{\pgfpoint{-16.840886\du}{20.568224\du}}
\pgfpathlineto{\pgfpoint{-16.839095\du}{20.571807\du}}
\pgfpathlineto{\pgfpoint{-16.835811\du}{20.574494\du}}
\pgfpathlineto{\pgfpoint{-16.831930\du}{20.577181\du}}
\pgfpathlineto{\pgfpoint{-16.827750\du}{20.579868\du}}
\pgfpathlineto{\pgfpoint{-16.823570\du}{20.582555\du}}
\pgfpathlineto{\pgfpoint{-16.818495\du}{20.585242\du}}
\pgfpathlineto{\pgfpoint{-16.813419\du}{20.587332\du}}
\pgfpathlineto{\pgfpoint{-16.807150\du}{20.590019\du}}
\pgfpathlineto{\pgfpoint{-16.801179\du}{20.592109\du}}
\pgfpathlineto{\pgfpoint{-16.795208\du}{20.594497\du}}
\pgfpathlineto{\pgfpoint{-16.787744\du}{20.596587\du}}
\pgfpathlineto{\pgfpoint{-16.781176\du}{20.598677\du}}
\pgfpathlineto{\pgfpoint{-16.773712\du}{20.600468\du}}
\pgfpathlineto{\pgfpoint{-16.765651\du}{20.601961\du}}
\pgfpathlineto{\pgfpoint{-16.757889\du}{20.604051\du}}
\pgfpathlineto{\pgfpoint{-16.749529\du}{20.605543\du}}
\pgfpathlineto{\pgfpoint{-16.740573\du}{20.607036\du}}
\pgfpathlineto{\pgfpoint{-16.732213\du}{20.608230\du}}
\pgfpathlineto{\pgfpoint{-16.722958\du}{20.609425\du}}
\pgfpathlineto{\pgfpoint{-16.714001\du}{20.610022\du}}
\pgfpathlineto{\pgfpoint{-16.704746\du}{20.611216\du}}
\pgfpathlineto{\pgfpoint{-16.695193\du}{20.611514\du}}
\pgfpathlineto{\pgfpoint{-16.685639\du}{20.612112\du}}
\pgfpathlineto{\pgfpoint{-16.676085\du}{20.612410\du}}
\pgfpathlineto{\pgfpoint{-16.666233\du}{20.613007\du}}
\pgfpathlineto{\pgfpoint{-16.656978\du}{20.613007\du}}
\pgfpathlineto{\pgfpoint{-16.656978\du}{20.613007\du}}
\pgfpathlineto{\pgfpoint{-16.647424\du}{20.613007\du}}
\pgfpathlineto{\pgfpoint{-16.637572\du}{20.612410\du}}
\pgfpathlineto{\pgfpoint{-16.627720\du}{20.612112\du}}
\pgfpathlineto{\pgfpoint{-16.618465\du}{20.611514\du}}
\pgfpathlineto{\pgfpoint{-16.608911\du}{20.611216\du}}
\pgfpathlineto{\pgfpoint{-16.599358\du}{20.610022\du}}
\pgfpathlineto{\pgfpoint{-16.590700\du}{20.609425\du}}
\pgfpathlineto{\pgfpoint{-16.581444\du}{20.608230\du}}
\pgfpathlineto{\pgfpoint{-16.573085\du}{20.607036\du}}
\pgfpathlineto{\pgfpoint{-16.564725\du}{20.605543\du}}
\pgfpathlineto{\pgfpoint{-16.555769\du}{20.604051\du}}
\pgfpathlineto{\pgfpoint{-16.548007\du}{20.601961\du}}
\pgfpathlineto{\pgfpoint{-16.540543\du}{20.600468\du}}
\pgfpathlineto{\pgfpoint{-16.532482\du}{20.598677\du}}
\pgfpathlineto{\pgfpoint{-16.525914\du}{20.596587\du}}
\pgfpathlineto{\pgfpoint{-16.519047\du}{20.594497\du}}
\pgfpathlineto{\pgfpoint{-16.512479\du}{20.592109\du}}
\pgfpathlineto{\pgfpoint{-16.506508\du}{20.590019\du}}
\pgfpathlineto{\pgfpoint{-16.500835\du}{20.587332\du}}
\pgfpathlineto{\pgfpoint{-16.495760\du}{20.585242\du}}
\pgfpathlineto{\pgfpoint{-16.490685\du}{20.582555\du}}
\pgfpathlineto{\pgfpoint{-16.485908\du}{20.579868\du}}
\pgfpathlineto{\pgfpoint{-16.481728\du}{20.577181\du}}
\pgfpathlineto{\pgfpoint{-16.478444\du}{20.574494\du}}
\pgfpathlineto{\pgfpoint{-16.475160\du}{20.571807\du}}
\pgfpathlineto{\pgfpoint{-16.473070\du}{20.568224\du}}
\pgfpathlineto{\pgfpoint{-16.470980\du}{20.565537\du}}
\pgfpathlineto{\pgfpoint{-16.468890\du}{20.562253\du}}
\pgfpathlineto{\pgfpoint{-16.467696\du}{20.559566\du}}
\pgfpathlineto{\pgfpoint{-16.467099\du}{20.556879\du}}
\pgfpathlineto{\pgfpoint{-16.466800\du}{20.553297\du}}
\pgfusepath{stroke}
\pgfsetbuttcap
\pgfsetmiterjoin
\pgfsetdash{}{0pt}
\definecolor{dialinecolor}{rgb}{0.000000, 0.000000, 0.000000}
\pgfsetfillcolor{dialinecolor}
\pgfpathmoveto{\pgfpoint{-16.645036\du}{20.541355\du}}
\pgfpathlineto{\pgfpoint{-16.621152\du}{20.553297\du}}
\pgfpathlineto{\pgfpoint{-16.549798\du}{20.529711\du}}
\pgfpathlineto{\pgfpoint{-16.513972\du}{20.541355\du}}
\pgfpathlineto{\pgfpoint{-16.537557\du}{20.517471\du}}
\pgfpathlineto{\pgfpoint{-16.621152\du}{20.517471\du}}
\pgfpathlineto{\pgfpoint{-16.585326\du}{20.517471\du}}
\pgfpathlineto{\pgfpoint{-16.645036\du}{20.541355\du}}
\pgfusepath{fill}
\pgfsetbuttcap
\pgfsetmiterjoin
\pgfsetdash{}{0pt}
\definecolor{dialinecolor}{rgb}{0.000000, 0.000000, 0.000000}
\pgfsetfillcolor{dialinecolor}
\pgfpathmoveto{\pgfpoint{-16.656978\du}{20.577181\du}}
\pgfpathlineto{\pgfpoint{-16.692506\du}{20.565239\du}}
\pgfpathlineto{\pgfpoint{-16.752216\du}{20.589123\du}}
\pgfpathlineto{\pgfpoint{-16.787744\du}{20.577181\du}}
\pgfpathlineto{\pgfpoint{-16.775503\du}{20.613007\du}}
\pgfpathlineto{\pgfpoint{-16.692506\du}{20.613007\du}}
\pgfpathlineto{\pgfpoint{-16.728033\du}{20.601364\du}}
\pgfpathlineto{\pgfpoint{-16.656978\du}{20.577181\du}}
\pgfusepath{fill}
\pgfsetbuttcap
\pgfsetmiterjoin
\pgfsetdash{}{0pt}
\definecolor{dialinecolor}{rgb}{0.000000, 0.000000, 0.000000}
\pgfsetfillcolor{dialinecolor}
\pgfpathmoveto{\pgfpoint{-16.775503\du}{20.517471\du}}
\pgfpathlineto{\pgfpoint{-16.752216\du}{20.505827\du}}
\pgfpathlineto{\pgfpoint{-16.680862\du}{20.541355\du}}
\pgfpathlineto{\pgfpoint{-16.656978\du}{20.529711\du}}
\pgfpathlineto{\pgfpoint{-16.668323\du}{20.553297\du}}
\pgfpathlineto{\pgfpoint{-16.752216\du}{20.553297\du}}
\pgfpathlineto{\pgfpoint{-16.716390\du}{20.541355\du}}
\pgfpathlineto{\pgfpoint{-16.775503\du}{20.517471\du}}
\pgfusepath{fill}
\pgfsetbuttcap
\pgfsetmiterjoin
\pgfsetdash{}{0pt}
\definecolor{dialinecolor}{rgb}{0.000000, 0.000000, 0.000000}
\pgfsetfillcolor{dialinecolor}
\pgfpathmoveto{\pgfpoint{-16.525914\du}{20.601364\du}}
\pgfpathlineto{\pgfpoint{-16.549798\du}{20.613007\du}}
\pgfpathlineto{\pgfpoint{-16.621152\du}{20.577181\du}}
\pgfpathlineto{\pgfpoint{-16.656978\du}{20.589123\du}}
\pgfpathlineto{\pgfpoint{-16.632795\du}{20.565239\du}}
\pgfpathlineto{\pgfpoint{-16.549798\du}{20.565239\du}}
\pgfpathlineto{\pgfpoint{-16.585326\du}{20.577181\du}}
\pgfpathlineto{\pgfpoint{-16.525914\du}{20.601364\du}}
\pgfusepath{fill}
\pgfsetbuttcap
\pgfsetmiterjoin
\pgfsetdash{}{0pt}
\definecolor{dialinecolor}{rgb}{1.000000, 1.000000, 1.000000}
\pgfsetfillcolor{dialinecolor}
\pgfpathmoveto{\pgfpoint{-16.645036\du}{20.553297\du}}
\pgfpathlineto{\pgfpoint{-16.621152\du}{20.553297\du}}
\pgfpathlineto{\pgfpoint{-16.549798\du}{20.529711\du}}
\pgfpathlineto{\pgfpoint{-16.513972\du}{20.541355\du}}
\pgfpathlineto{\pgfpoint{-16.537557\du}{20.517471\du}}
\pgfpathlineto{\pgfpoint{-16.608911\du}{20.517471\du}}
\pgfpathlineto{\pgfpoint{-16.585326\du}{20.517471\du}}
\pgfpathlineto{\pgfpoint{-16.645036\du}{20.553297\du}}
\pgfusepath{fill}
\pgfsetbuttcap
\pgfsetmiterjoin
\pgfsetdash{}{0pt}
\definecolor{dialinecolor}{rgb}{1.000000, 1.000000, 1.000000}
\pgfsetfillcolor{dialinecolor}
\pgfpathmoveto{\pgfpoint{-16.656978\du}{20.577181\du}}
\pgfpathlineto{\pgfpoint{-16.680862\du}{20.565239\du}}
\pgfpathlineto{\pgfpoint{-16.752216\du}{20.589123\du}}
\pgfpathlineto{\pgfpoint{-16.787744\du}{20.589123\du}}
\pgfpathlineto{\pgfpoint{-16.763860\du}{20.613007\du}}
\pgfpathlineto{\pgfpoint{-16.680862\du}{20.613007\du}}
\pgfpathlineto{\pgfpoint{-16.716390\du}{20.601364\du}}
\pgfpathlineto{\pgfpoint{-16.656978\du}{20.577181\du}}
\pgfusepath{fill}
\pgfsetbuttcap
\pgfsetmiterjoin
\pgfsetdash{}{0pt}
\definecolor{dialinecolor}{rgb}{1.000000, 1.000000, 1.000000}
\pgfsetfillcolor{dialinecolor}
\pgfpathmoveto{\pgfpoint{-16.775503\du}{20.517471\du}}
\pgfpathlineto{\pgfpoint{-16.752216\du}{20.517471\du}}
\pgfpathlineto{\pgfpoint{-16.680862\du}{20.541355\du}}
\pgfpathlineto{\pgfpoint{-16.645036\du}{20.529711\du}}
\pgfpathlineto{\pgfpoint{-16.668323\du}{20.553297\du}}
\pgfpathlineto{\pgfpoint{-16.752216\du}{20.553297\du}}
\pgfpathlineto{\pgfpoint{-16.716390\du}{20.553297\du}}
\pgfpathlineto{\pgfpoint{-16.775503\du}{20.517471\du}}
\pgfusepath{fill}
\pgfsetbuttcap
\pgfsetmiterjoin
\pgfsetdash{}{0pt}
\definecolor{dialinecolor}{rgb}{1.000000, 1.000000, 1.000000}
\pgfsetfillcolor{dialinecolor}
\pgfpathmoveto{\pgfpoint{-16.525914\du}{20.601364\du}}
\pgfpathlineto{\pgfpoint{-16.549798\du}{20.613007\du}}
\pgfpathlineto{\pgfpoint{-16.621152\du}{20.589123\du}}
\pgfpathlineto{\pgfpoint{-16.645036\du}{20.589123\du}}
\pgfpathlineto{\pgfpoint{-16.632795\du}{20.565239\du}}
\pgfpathlineto{\pgfpoint{-16.549798\du}{20.565239\du}}
\pgfpathlineto{\pgfpoint{-16.585326\du}{20.577181\du}}
\pgfpathlineto{\pgfpoint{-16.525914\du}{20.601364\du}}
\pgfusepath{fill}
\pgfsetbuttcap
\pgfsetmiterjoin
\pgfsetdash{}{0pt}
\definecolor{dialinecolor}{rgb}{0.666667, 0.901961, 1.000000}
\pgfsetstrokecolor{dialinecolor}
\pgfpathmoveto{\pgfpoint{-16.847454\du}{20.553297\du}}
\pgfpathlineto{\pgfpoint{-16.847454\du}{20.648236\du}}
\pgfusepath{stroke}
\pgfsetbuttcap
\pgfsetmiterjoin
\pgfsetdash{}{0pt}
\definecolor{dialinecolor}{rgb}{0.666667, 0.901961, 1.000000}
\pgfsetstrokecolor{dialinecolor}
\pgfpathmoveto{\pgfpoint{-16.466800\du}{20.553297\du}}
\pgfpathlineto{\pgfpoint{-16.466800\du}{20.648236\du}}
\pgfusepath{stroke}
\pgfsetlinewidth{0.100000\du}
\pgfsetdash{}{0pt}
\pgfsetdash{}{0pt}
\pgfsetbuttcap
\pgfsetmiterjoin
\pgfsetlinewidth{0.001000\du}
\pgfsetbuttcap
\pgfsetmiterjoin
\pgfsetdash{}{0pt}
\definecolor{dialinecolor}{rgb}{0.000000, 0.470588, 0.666667}
\pgfsetstrokecolor{dialinecolor}
\pgfpathmoveto{\pgfpoint{-9.385683\du}{21.076254\du}}
\pgfpathlineto{\pgfpoint{-9.379115\du}{21.067297\du}}
\pgfpathlineto{\pgfpoint{-9.372845\du}{21.057744\du}}
\pgfpathlineto{\pgfpoint{-9.367471\du}{21.047593\du}}
\pgfpathlineto{\pgfpoint{-9.362396\du}{21.038338\du}}
\pgfpathlineto{\pgfpoint{-9.358515\du}{21.028187\du}}
\pgfpathlineto{\pgfpoint{-9.355231\du}{21.018036\du}}
\pgfpathlineto{\pgfpoint{-9.352245\du}{21.007886\du}}
\pgfpathlineto{\pgfpoint{-9.350454\du}{20.997436\du}}
\pgfpathlineto{\pgfpoint{-9.349260\du}{20.987285\du}}
\pgfpathlineto{\pgfpoint{-9.348364\du}{20.976836\du}}
\pgfpathlineto{\pgfpoint{-9.348364\du}{20.966387\du}}
\pgfpathlineto{\pgfpoint{-9.349260\du}{20.955639\du}}
\pgfpathlineto{\pgfpoint{-9.350454\du}{20.945488\du}}
\pgfpathlineto{\pgfpoint{-9.352245\du}{20.935039\du}}
\pgfpathlineto{\pgfpoint{-9.354932\du}{20.925187\du}}
\pgfpathlineto{\pgfpoint{-9.358515\du}{20.914737\du}}
\pgfpathlineto{\pgfpoint{-9.362396\du}{20.904885\du}}
\pgfpathlineto{\pgfpoint{-9.367471\du}{20.894734\du}}
\pgfpathlineto{\pgfpoint{-9.372845\du}{20.885181\du}}
\pgfpathlineto{\pgfpoint{-9.379115\du}{20.876224\du}}
\pgfpathlineto{\pgfpoint{-9.385683\du}{20.866670\du}}
\pgfpathlineto{\pgfpoint{-9.393147\du}{20.858012\du}}
\pgfpathlineto{\pgfpoint{-9.400611\du}{20.849354\du}}
\pgfpathlineto{\pgfpoint{-9.408970\du}{20.840995\du}}
\pgfpathlineto{\pgfpoint{-9.417628\du}{20.832934\du}}
\pgfpathlineto{\pgfpoint{-9.426883\du}{20.825172\du}}
\pgfpathlineto{\pgfpoint{-9.436437\du}{20.818006\du}}
\pgfpathlineto{\pgfpoint{-9.447185\du}{20.811438\du}}
\pgfpathlineto{\pgfpoint{-9.457634\du}{20.804870\du}}
\pgfpathlineto{\pgfpoint{-9.468979\du}{20.798601\du}}
\pgfpathlineto{\pgfpoint{-9.480921\du}{20.792928\du}}
\pgfpathlineto{\pgfpoint{-9.492863\du}{20.787554\du}}
\pgfpathlineto{\pgfpoint{-9.505104\du}{20.783374\du}}
\pgfpathlineto{\pgfpoint{-9.517643\du}{20.779195\du}}
\pgfpathlineto{\pgfpoint{-9.530779\du}{20.775015\du}}
\pgfpathlineto{\pgfpoint{-9.543617\du}{20.772029\du}}
\pgfpathlineto{\pgfpoint{-9.556753\du}{20.769342\du}}
\pgfpathlineto{\pgfpoint{-9.570487\du}{20.766954\du}}
\pgfpathlineto{\pgfpoint{-9.583623\du}{20.765461\du}}
\pgfpathlineto{\pgfpoint{-9.597655\du}{20.764267\du}}
\pgfpathlineto{\pgfpoint{-9.611388\du}{20.763371\du}}
\pgfpathlineto{\pgfpoint{-9.625420\du}{20.763371\du}}
\pgfpathlineto{\pgfpoint{-9.638855\du}{20.763969\du}}
\pgfpathlineto{\pgfpoint{-9.652589\du}{20.764566\du}}
\pgfpathlineto{\pgfpoint{-9.666621\du}{20.766058\du}}
\pgfpathlineto{\pgfpoint{-9.679757\du}{20.768447\du}}
\pgfpathlineto{\pgfpoint{-9.693490\du}{20.770835\du}}
\pgfpathlineto{\pgfpoint{-9.706029\du}{20.773821\du}}
\pgfpathlineto{\pgfpoint{-9.719166\du}{20.776806\du}}
\pgfpathlineto{\pgfpoint{-9.731705\du}{20.780986\du}}
\pgfpathlineto{\pgfpoint{-9.744543\du}{20.785763\du}}
\pgfpathlineto{\pgfpoint{-9.756783\du}{20.790241\du}}
\pgfpathlineto{\pgfpoint{-9.768128\du}{20.795615\du}}
\pgfpathlineto{\pgfpoint{-9.779772\du}{20.801885\du}}
\pgfpathlineto{\pgfpoint{-9.790818\du}{20.807856\du}}
\pgfpathlineto{\pgfpoint{-9.801268\du}{20.814722\du}}
\pgfpathlineto{\pgfpoint{-9.811418\du}{20.821291\du}}
\pgfpathlineto{\pgfpoint{-9.820972\du}{20.829053\du}}
\pgfpathlineto{\pgfpoint{-9.830526\du}{20.837114\du}}
\pgfpathlineto{\pgfpoint{-9.838885\du}{20.845175\du}}
\pgfpathlineto{\pgfpoint{-9.846946\du}{20.853534\du}}
\pgfpathlineto{\pgfpoint{-9.854708\du}{20.862491\du}}
\pgfusepath{stroke}
\pgfsetlinewidth{0.080000\du}
\pgfsetbuttcap
\pgfsetmiterjoin
\pgfsetdash{}{0pt}
\definecolor{dialinecolor}{rgb}{0.666667, 0.901961, 1.000000}
\pgfsetstrokecolor{dialinecolor}
\pgfpathmoveto{\pgfpoint{-9.394938\du}{21.068790\du}}
\pgfpathlineto{\pgfpoint{-9.389564\du}{21.059834\du}}
\pgfpathlineto{\pgfpoint{-9.383295\du}{21.051176\du}}
\pgfpathlineto{\pgfpoint{-9.378219\du}{21.041920\du}}
\pgfpathlineto{\pgfpoint{-9.374040\du}{21.032665\du}}
\pgfpathlineto{\pgfpoint{-9.370158\du}{21.023112\du}}
\pgfpathlineto{\pgfpoint{-9.367173\du}{21.013558\du}}
\pgfpathlineto{\pgfpoint{-9.364486\du}{21.003706\du}}
\pgfpathlineto{\pgfpoint{-9.362396\du}{20.994451\du}}
\pgfpathlineto{\pgfpoint{-9.361500\du}{20.984598\du}}
\pgfpathlineto{\pgfpoint{-9.360903\du}{20.974448\du}}
\pgfpathlineto{\pgfpoint{-9.360903\du}{20.964595\du}}
\pgfpathlineto{\pgfpoint{-9.361799\du}{20.954743\du}}
\pgfpathlineto{\pgfpoint{-9.363292\du}{20.944593\du}}
\pgfpathlineto{\pgfpoint{-9.365083\du}{20.935337\du}}
\pgfpathlineto{\pgfpoint{-9.368069\du}{20.925485\du}}
\pgfpathlineto{\pgfpoint{-9.371651\du}{20.915932\du}}
\pgfpathlineto{\pgfpoint{-9.375532\du}{20.906676\du}}
\pgfpathlineto{\pgfpoint{-9.380011\du}{20.897421\du}}
\pgfpathlineto{\pgfpoint{-9.385086\du}{20.888763\du}}
\pgfpathlineto{\pgfpoint{-9.390758\du}{20.879508\du}}
\pgfpathlineto{\pgfpoint{-9.397924\du}{20.871149\du}}
\pgfpathlineto{\pgfpoint{-9.404492\du}{20.862491\du}}
\pgfpathlineto{\pgfpoint{-9.412254\du}{20.854430\du}}
\pgfpathlineto{\pgfpoint{-9.420315\du}{20.846667\du}}
\pgfpathlineto{\pgfpoint{-9.428675\du}{20.839801\du}}
\pgfpathlineto{\pgfpoint{-9.437333\du}{20.832337\du}}
\pgfpathlineto{\pgfpoint{-9.447185\du}{20.825470\du}}
\pgfpathlineto{\pgfpoint{-9.457634\du}{20.818604\du}}
\pgfpathlineto{\pgfpoint{-9.467785\du}{20.812931\du}}
\pgfpathlineto{\pgfpoint{-9.478234\du}{20.807557\du}}
\pgfpathlineto{\pgfpoint{-9.490176\du}{20.802183\du}}
\pgfpathlineto{\pgfpoint{-9.500924\du}{20.797108\du}}
\pgfpathlineto{\pgfpoint{-9.513463\du}{20.792928\du}}
\pgfpathlineto{\pgfpoint{-9.524808\du}{20.789047\du}}
\pgfpathlineto{\pgfpoint{-9.537348\du}{20.785763\du}}
\pgfpathlineto{\pgfpoint{-9.549887\du}{20.782777\du}}
\pgfpathlineto{\pgfpoint{-9.562426\du}{20.780090\du}}
\pgfpathlineto{\pgfpoint{-9.575562\du}{20.778299\du}}
\pgfpathlineto{\pgfpoint{-9.588101\du}{20.776806\du}}
\pgfpathlineto{\pgfpoint{-9.601238\du}{20.775612\du}}
\pgfpathlineto{\pgfpoint{-9.615270\du}{20.775015\du}}
\pgfpathlineto{\pgfpoint{-9.628406\du}{20.775015\du}}
\pgfpathlineto{\pgfpoint{-9.641244\du}{20.775612\du}}
\pgfpathlineto{\pgfpoint{-9.654380\du}{20.776508\du}}
\pgfpathlineto{\pgfpoint{-9.666919\du}{20.778299\du}}
\pgfpathlineto{\pgfpoint{-9.680055\du}{20.780090\du}}
\pgfpathlineto{\pgfpoint{-9.692595\du}{20.782180\du}}
\pgfpathlineto{\pgfpoint{-9.705432\du}{20.785763\du}}
\pgfpathlineto{\pgfpoint{-9.717673\du}{20.788748\du}}
\pgfpathlineto{\pgfpoint{-9.729615\du}{20.792630\du}}
\pgfpathlineto{\pgfpoint{-9.741856\du}{20.796809\du}}
\pgfpathlineto{\pgfpoint{-9.753201\du}{20.801885\du}}
\pgfpathlineto{\pgfpoint{-9.764247\du}{20.806661\du}}
\pgfpathlineto{\pgfpoint{-9.774995\du}{20.812633\du}}
\pgfpathlineto{\pgfpoint{-9.785444\du}{20.818305\du}}
\pgfpathlineto{\pgfpoint{-9.795894\du}{20.824873\du}}
\pgfpathlineto{\pgfpoint{-9.805447\du}{20.831441\du}}
\pgfpathlineto{\pgfpoint{-9.814105\du}{20.838607\du}}
\pgfpathlineto{\pgfpoint{-9.823659\du}{20.845772\du}}
\pgfpathlineto{\pgfpoint{-9.831421\du}{20.853534\du}}
\pgfpathlineto{\pgfpoint{-9.838885\du}{20.861595\du}}
\pgfpathlineto{\pgfpoint{-9.846050\du}{20.869955\du}}
\pgfusepath{stroke}
\pgfsetbuttcap
\pgfsetmiterjoin
\pgfsetdash{}{0pt}
\definecolor{dialinecolor}{rgb}{0.000000, 0.705882, 1.000000}
\pgfsetfillcolor{dialinecolor}
\pgfpathmoveto{\pgfpoint{-11.311941\du}{20.751728\du}}
\pgfpathlineto{\pgfpoint{-11.109224\du}{20.560953\du}}
\pgfpathlineto{\pgfpoint{-9.597655\du}{20.560953\du}}
\pgfpathlineto{\pgfpoint{-9.800372\du}{20.751728\du}}
\pgfpathlineto{\pgfpoint{-11.311941\du}{20.751728\du}}
\pgfusepath{fill}
\pgfsetlinewidth{0.001000\du}
\pgfsetbuttcap
\pgfsetmiterjoin
\pgfsetdash{}{0pt}
\definecolor{dialinecolor}{rgb}{0.666667, 0.901961, 1.000000}
\pgfsetstrokecolor{dialinecolor}
\pgfpathmoveto{\pgfpoint{-11.311941\du}{20.751728\du}}
\pgfpathlineto{\pgfpoint{-11.109224\du}{20.560953\du}}
\pgfpathlineto{\pgfpoint{-9.597655\du}{20.560953\du}}
\pgfpathlineto{\pgfpoint{-9.800372\du}{20.751728\du}}
\pgfpathlineto{\pgfpoint{-11.311941\du}{20.751728\du}}
\pgfusepath{stroke}
\pgfsetbuttcap
\pgfsetmiterjoin
\pgfsetdash{}{0pt}
\definecolor{dialinecolor}{rgb}{0.000000, 0.588235, 0.831373}
\pgfsetfillcolor{dialinecolor}
\pgfpathmoveto{\pgfpoint{-11.311941\du}{20.751728\du}}
\pgfpathlineto{\pgfpoint{-9.800372\du}{20.751728\du}}
\pgfpathlineto{\pgfpoint{-9.800372\du}{21.037741\du}}
\pgfpathlineto{\pgfpoint{-11.311941\du}{21.037741\du}}
\pgfpathlineto{\pgfpoint{-11.311941\du}{20.751728\du}}
\pgfusepath{fill}
\pgfsetbuttcap
\pgfsetmiterjoin
\pgfsetdash{}{0pt}
\definecolor{dialinecolor}{rgb}{0.666667, 0.901961, 1.000000}
\pgfsetstrokecolor{dialinecolor}
\pgfpathmoveto{\pgfpoint{-11.311941\du}{20.751728\du}}
\pgfpathlineto{\pgfpoint{-9.800372\du}{20.751728\du}}
\pgfpathlineto{\pgfpoint{-9.800372\du}{21.037741\du}}
\pgfpathlineto{\pgfpoint{-11.311941\du}{21.037741\du}}
\pgfpathlineto{\pgfpoint{-11.311941\du}{20.751728\du}}
\pgfusepath{stroke}
\pgfsetbuttcap
\pgfsetmiterjoin
\pgfsetdash{}{0pt}
\definecolor{dialinecolor}{rgb}{0.000000, 0.352941, 0.501961}
\pgfsetfillcolor{dialinecolor}
\pgfpathmoveto{\pgfpoint{-9.800372\du}{21.037741\du}}
\pgfpathlineto{\pgfpoint{-9.597655\du}{20.835024\du}}
\pgfpathlineto{\pgfpoint{-9.597655\du}{20.560953\du}}
\pgfpathlineto{\pgfpoint{-9.800372\du}{20.751728\du}}
\pgfpathlineto{\pgfpoint{-9.800372\du}{21.037741\du}}
\pgfusepath{fill}
\pgfsetbuttcap
\pgfsetmiterjoin
\pgfsetdash{}{0pt}
\definecolor{dialinecolor}{rgb}{0.666667, 0.901961, 1.000000}
\pgfsetstrokecolor{dialinecolor}
\pgfpathmoveto{\pgfpoint{-9.800372\du}{21.037741\du}}
\pgfpathlineto{\pgfpoint{-9.597655\du}{20.835024\du}}
\pgfpathlineto{\pgfpoint{-9.597655\du}{20.560953\du}}
\pgfpathlineto{\pgfpoint{-9.800372\du}{20.751728\du}}
\pgfpathlineto{\pgfpoint{-9.800372\du}{21.037741\du}}
\pgfusepath{stroke}
\pgfsetbuttcap
\pgfsetmiterjoin
\pgfsetdash{}{0pt}
\definecolor{dialinecolor}{rgb}{0.000000, 0.000000, 0.000000}
\pgfsetfillcolor{dialinecolor}
\pgfpathmoveto{\pgfpoint{-11.263874\du}{20.715603\du}}
\pgfpathlineto{\pgfpoint{-11.109224\du}{20.560953\du}}
\pgfpathlineto{\pgfpoint{-9.633780\du}{20.560953\du}}
\pgfpathlineto{\pgfpoint{-9.788131\du}{20.715603\du}}
\pgfpathlineto{\pgfpoint{-11.263874\du}{20.715603\du}}
\pgfusepath{fill}
\pgfsetbuttcap
\pgfsetmiterjoin
\pgfsetdash{}{0pt}
\definecolor{dialinecolor}{rgb}{0.000000, 0.000000, 0.000000}
\pgfsetstrokecolor{dialinecolor}
\pgfpathmoveto{\pgfpoint{-11.263874\du}{20.715603\du}}
\pgfpathlineto{\pgfpoint{-11.109224\du}{20.560953\du}}
\pgfpathlineto{\pgfpoint{-9.633780\du}{20.560953\du}}
\pgfpathlineto{\pgfpoint{-9.788131\du}{20.715603\du}}
\pgfpathlineto{\pgfpoint{-11.263874\du}{20.715603\du}}
\pgfusepath{stroke}
\pgfsetbuttcap
\pgfsetmiterjoin
\pgfsetdash{}{0pt}
\definecolor{dialinecolor}{rgb}{0.000000, 0.705882, 1.000000}
\pgfsetfillcolor{dialinecolor}
\pgfpathmoveto{\pgfpoint{-11.311941\du}{19.489450\du}}
\pgfpathlineto{\pgfpoint{-11.144752\du}{19.334800\du}}
\pgfpathlineto{\pgfpoint{-9.633780\du}{19.334800\du}}
\pgfpathlineto{\pgfpoint{-9.800372\du}{19.489450\du}}
\pgfpathlineto{\pgfpoint{-11.311941\du}{19.489450\du}}
\pgfusepath{fill}
\pgfsetbuttcap
\pgfsetmiterjoin
\pgfsetdash{}{0pt}
\definecolor{dialinecolor}{rgb}{0.666667, 0.901961, 1.000000}
\pgfsetstrokecolor{dialinecolor}
\pgfpathmoveto{\pgfpoint{-11.311941\du}{19.489450\du}}
\pgfpathlineto{\pgfpoint{-11.144752\du}{19.334800\du}}
\pgfpathlineto{\pgfpoint{-9.633780\du}{19.334800\du}}
\pgfpathlineto{\pgfpoint{-9.800372\du}{19.489450\du}}
\pgfpathlineto{\pgfpoint{-11.311941\du}{19.489450\du}}
\pgfusepath{stroke}
\pgfsetbuttcap
\pgfsetmiterjoin
\pgfsetdash{}{0pt}
\definecolor{dialinecolor}{rgb}{0.000000, 0.588235, 0.831373}
\pgfsetfillcolor{dialinecolor}
\pgfpathmoveto{\pgfpoint{-11.311941\du}{19.489450\du}}
\pgfpathlineto{\pgfpoint{-9.788131\du}{19.489450\du}}
\pgfpathlineto{\pgfpoint{-9.788131\du}{20.691719\du}}
\pgfpathlineto{\pgfpoint{-11.311941\du}{20.691719\du}}
\pgfpathlineto{\pgfpoint{-11.311941\du}{19.489450\du}}
\pgfusepath{fill}
\pgfsetbuttcap
\pgfsetmiterjoin
\pgfsetdash{}{0pt}
\definecolor{dialinecolor}{rgb}{0.666667, 0.901961, 1.000000}
\pgfsetstrokecolor{dialinecolor}
\pgfpathmoveto{\pgfpoint{-11.311941\du}{19.489450\du}}
\pgfpathlineto{\pgfpoint{-9.788430\du}{19.489450\du}}
\pgfpathlineto{\pgfpoint{-9.788430\du}{20.691420\du}}
\pgfpathlineto{\pgfpoint{-11.311941\du}{20.691420\du}}
\pgfpathlineto{\pgfpoint{-11.311941\du}{19.489450\du}}
\pgfusepath{stroke}
\pgfsetbuttcap
\pgfsetmiterjoin
\pgfsetdash{}{0pt}
\definecolor{dialinecolor}{rgb}{1.000000, 1.000000, 1.000000}
\pgfsetfillcolor{dialinecolor}
\pgfpathmoveto{\pgfpoint{-11.180876\du}{19.644398\du}}
\pgfpathlineto{\pgfpoint{-9.919196\du}{19.644398\du}}
\pgfpathlineto{\pgfpoint{-9.919196\du}{20.573194\du}}
\pgfpathlineto{\pgfpoint{-11.180876\du}{20.573194\du}}
\pgfpathlineto{\pgfpoint{-11.180876\du}{19.644398\du}}
\pgfusepath{fill}
\pgfsetbuttcap
\pgfsetmiterjoin
\pgfsetdash{}{0pt}
\definecolor{dialinecolor}{rgb}{0.000000, 0.235294, 0.333333}
\pgfsetstrokecolor{dialinecolor}
\pgfpathmoveto{\pgfpoint{-11.180876\du}{19.644398\du}}
\pgfpathlineto{\pgfpoint{-9.919494\du}{19.644398\du}}
\pgfpathlineto{\pgfpoint{-9.919494\du}{20.572895\du}}
\pgfpathlineto{\pgfpoint{-11.180876\du}{20.572895\du}}
\pgfpathlineto{\pgfpoint{-11.180876\du}{19.644398\du}}
\pgfusepath{stroke}
\pgfsetbuttcap
\pgfsetmiterjoin
\pgfsetdash{}{0pt}
\definecolor{dialinecolor}{rgb}{0.000000, 0.352941, 0.501961}
\pgfsetfillcolor{dialinecolor}
\pgfpathmoveto{\pgfpoint{-9.800372\du}{20.680374\du}}
\pgfpathlineto{\pgfpoint{-9.633780\du}{20.513483\du}}
\pgfpathlineto{\pgfpoint{-9.633780\du}{19.334800\du}}
\pgfpathlineto{\pgfpoint{-9.800372\du}{19.489450\du}}
\pgfpathlineto{\pgfpoint{-9.800372\du}{20.680374\du}}
\pgfusepath{fill}
\pgfsetbuttcap
\pgfsetmiterjoin
\pgfsetdash{}{0pt}
\definecolor{dialinecolor}{rgb}{0.666667, 0.901961, 1.000000}
\pgfsetstrokecolor{dialinecolor}
\pgfpathmoveto{\pgfpoint{-9.800372\du}{20.680374\du}}
\pgfpathlineto{\pgfpoint{-9.633780\du}{20.513483\du}}
\pgfpathlineto{\pgfpoint{-9.633780\du}{19.334800\du}}
\pgfpathlineto{\pgfpoint{-9.800372\du}{19.489450\du}}
\pgfpathlineto{\pgfpoint{-9.800372\du}{20.680374\du}}
\pgfusepath{stroke}
\pgfsetbuttcap
\pgfsetmiterjoin
\pgfsetdash{}{0pt}
\definecolor{dialinecolor}{rgb}{0.000000, 0.705882, 1.000000}
\pgfsetfillcolor{dialinecolor}
\pgfpathmoveto{\pgfpoint{-11.609000\du}{21.275687\du}}
\pgfpathlineto{\pgfpoint{-11.371353\du}{20.977732\du}}
\pgfpathlineto{\pgfpoint{-9.728719\du}{20.977732\du}}
\pgfpathlineto{\pgfpoint{-9.966665\du}{21.275687\du}}
\pgfpathlineto{\pgfpoint{-11.609000\du}{21.275687\du}}
\pgfusepath{fill}
\pgfsetbuttcap
\pgfsetmiterjoin
\pgfsetdash{}{0pt}
\definecolor{dialinecolor}{rgb}{0.666667, 0.901961, 1.000000}
\pgfsetstrokecolor{dialinecolor}
\pgfpathmoveto{\pgfpoint{-11.609000\du}{21.275687\du}}
\pgfpathlineto{\pgfpoint{-11.371353\du}{20.977732\du}}
\pgfpathlineto{\pgfpoint{-9.728719\du}{20.977732\du}}
\pgfpathlineto{\pgfpoint{-9.966665\du}{21.275687\du}}
\pgfpathlineto{\pgfpoint{-11.609000\du}{21.275687\du}}
\pgfusepath{stroke}
\pgfsetbuttcap
\pgfsetmiterjoin
\pgfsetdash{}{0pt}
\definecolor{dialinecolor}{rgb}{0.000000, 0.352941, 0.501961}
\pgfsetfillcolor{dialinecolor}
\pgfpathmoveto{\pgfpoint{-9.966665\du}{21.334800\du}}
\pgfpathlineto{\pgfpoint{-9.728719\du}{21.084912\du}}
\pgfpathlineto{\pgfpoint{-9.728719\du}{20.977732\du}}
\pgfpathlineto{\pgfpoint{-9.966665\du}{21.275687\du}}
\pgfpathlineto{\pgfpoint{-9.966665\du}{21.334800\du}}
\pgfusepath{fill}
\pgfsetbuttcap
\pgfsetmiterjoin
\pgfsetdash{}{0pt}
\definecolor{dialinecolor}{rgb}{0.666667, 0.901961, 1.000000}
\pgfsetstrokecolor{dialinecolor}
\pgfpathmoveto{\pgfpoint{-9.966665\du}{21.334800\du}}
\pgfpathlineto{\pgfpoint{-9.728719\du}{21.084912\du}}
\pgfpathlineto{\pgfpoint{-9.728719\du}{20.977732\du}}
\pgfpathlineto{\pgfpoint{-9.966665\du}{21.275687\du}}
\pgfpathlineto{\pgfpoint{-9.966665\du}{21.334800\du}}
\pgfusepath{stroke}
\pgfsetbuttcap
\pgfsetmiterjoin
\pgfsetdash{}{0pt}
\definecolor{dialinecolor}{rgb}{0.000000, 0.588235, 0.831373}
\pgfsetfillcolor{dialinecolor}
\pgfpathmoveto{\pgfpoint{-11.609000\du}{21.275687\du}}
\pgfpathlineto{\pgfpoint{-9.966665\du}{21.275687\du}}
\pgfpathlineto{\pgfpoint{-9.966665\du}{21.334800\du}}
\pgfpathlineto{\pgfpoint{-11.609000\du}{21.334800\du}}
\pgfpathlineto{\pgfpoint{-11.609000\du}{21.275687\du}}
\pgfusepath{fill}
\pgfsetbuttcap
\pgfsetmiterjoin
\pgfsetdash{}{0pt}
\definecolor{dialinecolor}{rgb}{0.666667, 0.901961, 1.000000}
\pgfsetstrokecolor{dialinecolor}
\pgfpathmoveto{\pgfpoint{-11.609000\du}{21.275687\du}}
\pgfpathlineto{\pgfpoint{-9.966665\du}{21.275687\du}}
\pgfpathlineto{\pgfpoint{-9.966665\du}{21.334800\du}}
\pgfpathlineto{\pgfpoint{-11.609000\du}{21.334800\du}}
\pgfpathlineto{\pgfpoint{-11.609000\du}{21.275687\du}}
\pgfusepath{stroke}
\pgfsetbuttcap
\pgfsetmiterjoin
\pgfsetdash{}{0pt}
\definecolor{dialinecolor}{rgb}{0.000000, 0.705882, 1.000000}
\pgfsetfillcolor{dialinecolor}
\pgfpathmoveto{\pgfpoint{-9.633780\du}{21.239562\du}}
\pgfpathlineto{\pgfpoint{-9.514658\du}{21.061326\du}}
\pgfpathlineto{\pgfpoint{-9.299999\du}{21.061326\du}}
\pgfpathlineto{\pgfpoint{-9.418822\du}{21.239562\du}}
\pgfpathlineto{\pgfpoint{-9.633780\du}{21.239562\du}}
\pgfusepath{fill}
\pgfsetbuttcap
\pgfsetmiterjoin
\pgfsetdash{}{0pt}
\definecolor{dialinecolor}{rgb}{0.666667, 0.901961, 1.000000}
\pgfsetstrokecolor{dialinecolor}
\pgfpathmoveto{\pgfpoint{-9.633780\du}{21.239562\du}}
\pgfpathlineto{\pgfpoint{-9.514658\du}{21.061326\du}}
\pgfpathlineto{\pgfpoint{-9.299999\du}{21.061326\du}}
\pgfpathlineto{\pgfpoint{-9.418822\du}{21.239562\du}}
\pgfpathlineto{\pgfpoint{-9.633780\du}{21.239562\du}}
\pgfusepath{stroke}
\pgfsetbuttcap
\pgfsetmiterjoin
\pgfsetdash{}{0pt}
\definecolor{dialinecolor}{rgb}{0.000000, 0.352941, 0.501961}
\pgfsetfillcolor{dialinecolor}
\pgfpathmoveto{\pgfpoint{-9.418822\du}{21.299272\du}}
\pgfpathlineto{\pgfpoint{-9.299999\du}{21.168208\du}}
\pgfpathlineto{\pgfpoint{-9.299999\du}{21.061326\du}}
\pgfpathlineto{\pgfpoint{-9.418822\du}{21.239562\du}}
\pgfpathlineto{\pgfpoint{-9.418822\du}{21.299272\du}}
\pgfusepath{fill}
\pgfsetbuttcap
\pgfsetmiterjoin
\pgfsetdash{}{0pt}
\definecolor{dialinecolor}{rgb}{0.666667, 0.901961, 1.000000}
\pgfsetstrokecolor{dialinecolor}
\pgfpathmoveto{\pgfpoint{-9.418822\du}{21.299272\du}}
\pgfpathlineto{\pgfpoint{-9.299999\du}{21.168208\du}}
\pgfpathlineto{\pgfpoint{-9.299999\du}{21.061326\du}}
\pgfpathlineto{\pgfpoint{-9.418822\du}{21.239562\du}}
\pgfpathlineto{\pgfpoint{-9.418822\du}{21.299272\du}}
\pgfusepath{stroke}
\pgfsetbuttcap
\pgfsetmiterjoin
\pgfsetdash{}{0pt}
\definecolor{dialinecolor}{rgb}{0.000000, 0.588235, 0.831373}
\pgfsetfillcolor{dialinecolor}
\pgfpathmoveto{\pgfpoint{-9.633780\du}{21.239562\du}}
\pgfpathlineto{\pgfpoint{-9.418822\du}{21.239562\du}}
\pgfpathlineto{\pgfpoint{-9.418822\du}{21.299272\du}}
\pgfpathlineto{\pgfpoint{-9.633780\du}{21.299272\du}}
\pgfpathlineto{\pgfpoint{-9.633780\du}{21.239562\du}}
\pgfusepath{fill}
\pgfsetbuttcap
\pgfsetmiterjoin
\pgfsetdash{}{0pt}
\definecolor{dialinecolor}{rgb}{0.666667, 0.901961, 1.000000}
\pgfsetstrokecolor{dialinecolor}
\pgfpathmoveto{\pgfpoint{-9.633780\du}{21.239562\du}}
\pgfpathlineto{\pgfpoint{-9.418822\du}{21.239562\du}}
\pgfpathlineto{\pgfpoint{-9.418822\du}{21.299272\du}}
\pgfpathlineto{\pgfpoint{-9.633780\du}{21.299272\du}}
\pgfpathlineto{\pgfpoint{-9.633780\du}{21.239562\du}}
\pgfusepath{stroke}
\pgfsetbuttcap
\pgfsetmiterjoin
\pgfsetdash{}{0pt}
\definecolor{dialinecolor}{rgb}{0.811765, 0.054902, 0.188235}
\pgfsetstrokecolor{dialinecolor}
\pgfpathmoveto{\pgfpoint{-10.847692\du}{19.870701\du}}
\pgfpathlineto{\pgfpoint{-10.513911\du}{19.870701\du}}
\pgfpathlineto{\pgfpoint{-10.597506\du}{19.953997\du}}
\pgfpathlineto{\pgfpoint{-10.311791\du}{19.953997\du}}
\pgfusepath{stroke}
\pgfsetbuttcap
\pgfsetmiterjoin
\pgfsetdash{}{0pt}
\definecolor{dialinecolor}{rgb}{0.811765, 0.054902, 0.188235}
\pgfsetstrokecolor{dialinecolor}
\pgfpathmoveto{\pgfpoint{-10.883519\du}{19.870701\du}}
\pgfpathlineto{\pgfpoint{-10.644976\du}{20.108647\du}}
\pgfpathlineto{\pgfpoint{-10.764396\du}{20.108647\du}}
\pgfpathlineto{\pgfpoint{-10.561978\du}{20.311364\du}}
\pgfusepath{stroke}
\pgfsetbuttcap
\pgfsetmiterjoin
\pgfsetdash{}{0pt}
\definecolor{dialinecolor}{rgb}{0.811765, 0.054902, 0.188235}
\pgfsetstrokecolor{dialinecolor}
\pgfpathmoveto{\pgfpoint{-10.205208\du}{19.870701\du}}
\pgfpathlineto{\pgfpoint{-10.442557\du}{20.108647\du}}
\pgfpathlineto{\pgfpoint{-10.324032\du}{20.108647\du}}
\pgfpathlineto{\pgfpoint{-10.526152\du}{20.311364\du}}
\pgfusepath{stroke}
\pgfsetbuttcap
\pgfsetmiterjoin
\pgfsetdash{}{0pt}
\definecolor{dialinecolor}{rgb}{0.000000, 0.470588, 0.666667}
\pgfsetfillcolor{dialinecolor}
\pgfpathmoveto{\pgfpoint{-10.716927\du}{19.959669\du}}
\pgfpathlineto{\pgfpoint{-10.717524\du}{19.956385\du}}
\pgfpathlineto{\pgfpoint{-10.718419\du}{19.953101\du}}
\pgfpathlineto{\pgfpoint{-10.719315\du}{19.949817\du}}
\pgfpathlineto{\pgfpoint{-10.721106\du}{19.946533\du}}
\pgfpathlineto{\pgfpoint{-10.723495\du}{19.943548\du}}
\pgfpathlineto{\pgfpoint{-10.726182\du}{19.940264\du}}
\pgfpathlineto{\pgfpoint{-10.729167\du}{19.936979\du}}
\pgfpathlineto{\pgfpoint{-10.732750\du}{19.933994\du}}
\pgfpathlineto{\pgfpoint{-10.736930\du}{19.931008\du}}
\pgfpathlineto{\pgfpoint{-10.741706\du}{19.928023\du}}
\pgfpathlineto{\pgfpoint{-10.746782\du}{19.925336\du}}
\pgfpathlineto{\pgfpoint{-10.752454\du}{19.922350\du}}
\pgfpathlineto{\pgfpoint{-10.758127\du}{19.919962\du}}
\pgfpathlineto{\pgfpoint{-10.764396\du}{19.917275\du}}
\pgfpathlineto{\pgfpoint{-10.770964\du}{19.914887\du}}
\pgfpathlineto{\pgfpoint{-10.778130\du}{19.912498\du}}
\pgfpathlineto{\pgfpoint{-10.785594\du}{19.910408\du}}
\pgfpathlineto{\pgfpoint{-10.793057\du}{19.908020\du}}
\pgfpathlineto{\pgfpoint{-10.801417\du}{19.906527\du}}
\pgfpathlineto{\pgfpoint{-10.809776\du}{19.904437\du}}
\pgfpathlineto{\pgfpoint{-10.817837\du}{19.902646\du}}
\pgfpathlineto{\pgfpoint{-10.827092\du}{19.901452\du}}
\pgfpathlineto{\pgfpoint{-10.835750\du}{19.899959\du}}
\pgfpathlineto{\pgfpoint{-10.845304\du}{19.898765\du}}
\pgfpathlineto{\pgfpoint{-10.854559\du}{19.897571\du}}
\pgfpathlineto{\pgfpoint{-10.863814\du}{19.896675\du}}
\pgfpathlineto{\pgfpoint{-10.873666\du}{19.895779\du}}
\pgfpathlineto{\pgfpoint{-10.883519\du}{19.895182\du}}
\pgfpathlineto{\pgfpoint{-10.893371\du}{19.894884\du}}
\pgfpathlineto{\pgfpoint{-10.903522\du}{19.894585\du}}
\pgfpathlineto{\pgfpoint{-10.913374\du}{19.894585\du}}
\pgfpathlineto{\pgfpoint{-10.913374\du}{19.894585\du}}
\pgfpathlineto{\pgfpoint{-10.923226\du}{19.894585\du}}
\pgfpathlineto{\pgfpoint{-10.933078\du}{19.894884\du}}
\pgfpathlineto{\pgfpoint{-10.943229\du}{19.895182\du}}
\pgfpathlineto{\pgfpoint{-10.953081\du}{19.895779\du}}
\pgfpathlineto{\pgfpoint{-10.962635\du}{19.896675\du}}
\pgfpathlineto{\pgfpoint{-10.972189\du}{19.897571\du}}
\pgfpathlineto{\pgfpoint{-10.981444\du}{19.898765\du}}
\pgfpathlineto{\pgfpoint{-10.990699\du}{19.899959\du}}
\pgfpathlineto{\pgfpoint{-10.999655\du}{19.901452\du}}
\pgfpathlineto{\pgfpoint{-11.008612\du}{19.902646\du}}
\pgfpathlineto{\pgfpoint{-11.016971\du}{19.904437\du}}
\pgfpathlineto{\pgfpoint{-11.025331\du}{19.906527\du}}
\pgfpathlineto{\pgfpoint{-11.033690\du}{19.908020\du}}
\pgfpathlineto{\pgfpoint{-11.040856\du}{19.910408\du}}
\pgfpathlineto{\pgfpoint{-11.048319\du}{19.912498\du}}
\pgfpathlineto{\pgfpoint{-11.055485\du}{19.914887\du}}
\pgfpathlineto{\pgfpoint{-11.062351\du}{19.917275\du}}
\pgfpathlineto{\pgfpoint{-11.068322\du}{19.919962\du}}
\pgfpathlineto{\pgfpoint{-11.074293\du}{19.922350\du}}
\pgfpathlineto{\pgfpoint{-11.079966\du}{19.925336\du}}
\pgfpathlineto{\pgfpoint{-11.085041\du}{19.928023\du}}
\pgfpathlineto{\pgfpoint{-11.089818\du}{19.931008\du}}
\pgfpathlineto{\pgfpoint{-11.093998\du}{19.933994\du}}
\pgfpathlineto{\pgfpoint{-11.097282\du}{19.936979\du}}
\pgfpathlineto{\pgfpoint{-11.100566\du}{19.940264\du}}
\pgfpathlineto{\pgfpoint{-11.103551\du}{19.943548\du}}
\pgfpathlineto{\pgfpoint{-11.105641\du}{19.946533\du}}
\pgfpathlineto{\pgfpoint{-11.107134\du}{19.949817\du}}
\pgfpathlineto{\pgfpoint{-11.108925\du}{19.953101\du}}
\pgfpathlineto{\pgfpoint{-11.109224\du}{19.956385\du}}
\pgfpathlineto{\pgfpoint{-11.109522\du}{19.959669\du}}
\pgfpathlineto{\pgfpoint{-11.109522\du}{19.959669\du}}
\pgfpathlineto{\pgfpoint{-11.109224\du}{19.963252\du}}
\pgfpathlineto{\pgfpoint{-11.108925\du}{19.966238\du}}
\pgfpathlineto{\pgfpoint{-11.107134\du}{19.969522\du}}
\pgfpathlineto{\pgfpoint{-11.105641\du}{19.972806\du}}
\pgfpathlineto{\pgfpoint{-11.103551\du}{19.976090\du}}
\pgfpathlineto{\pgfpoint{-11.100566\du}{19.979374\du}}
\pgfpathlineto{\pgfpoint{-11.097282\du}{19.982359\du}}
\pgfpathlineto{\pgfpoint{-11.093998\du}{19.985345\du}}
\pgfpathlineto{\pgfpoint{-11.089818\du}{19.988330\du}}
\pgfpathlineto{\pgfpoint{-11.085041\du}{19.991614\du}}
\pgfpathlineto{\pgfpoint{-11.079966\du}{19.994301\du}}
\pgfpathlineto{\pgfpoint{-11.074293\du}{19.996988\du}}
\pgfpathlineto{\pgfpoint{-11.068322\du}{19.999675\du}}
\pgfpathlineto{\pgfpoint{-11.062351\du}{20.002362\du}}
\pgfpathlineto{\pgfpoint{-11.055485\du}{20.004751\du}}
\pgfpathlineto{\pgfpoint{-11.048319\du}{20.007139\du}}
\pgfpathlineto{\pgfpoint{-11.040856\du}{20.008930\du}}
\pgfpathlineto{\pgfpoint{-11.033690\du}{20.011020\du}}
\pgfpathlineto{\pgfpoint{-11.025331\du}{20.013110\du}}
\pgfpathlineto{\pgfpoint{-11.016971\du}{20.015200\du}}
\pgfpathlineto{\pgfpoint{-11.008612\du}{20.016693\du}}
\pgfpathlineto{\pgfpoint{-10.999655\du}{20.018186\du}}
\pgfpathlineto{\pgfpoint{-10.990699\du}{20.019678\du}}
\pgfpathlineto{\pgfpoint{-10.981444\du}{20.020873\du}}
\pgfpathlineto{\pgfpoint{-10.972189\du}{20.021768\du}}
\pgfpathlineto{\pgfpoint{-10.962635\du}{20.022962\du}}
\pgfpathlineto{\pgfpoint{-10.953081\du}{20.023560\du}}
\pgfpathlineto{\pgfpoint{-10.943229\du}{20.024157\du}}
\pgfpathlineto{\pgfpoint{-10.933078\du}{20.024754\du}}
\pgfpathlineto{\pgfpoint{-10.923226\du}{20.025052\du}}
\pgfpathlineto{\pgfpoint{-10.913374\du}{20.025052\du}}
\pgfpathlineto{\pgfpoint{-10.913374\du}{20.025052\du}}
\pgfpathlineto{\pgfpoint{-10.903522\du}{20.025052\du}}
\pgfpathlineto{\pgfpoint{-10.893371\du}{20.024754\du}}
\pgfpathlineto{\pgfpoint{-10.883519\du}{20.024157\du}}
\pgfpathlineto{\pgfpoint{-10.873666\du}{20.023560\du}}
\pgfpathlineto{\pgfpoint{-10.863814\du}{20.022962\du}}
\pgfpathlineto{\pgfpoint{-10.854559\du}{20.021768\du}}
\pgfpathlineto{\pgfpoint{-10.845304\du}{20.020873\du}}
\pgfpathlineto{\pgfpoint{-10.835750\du}{20.019678\du}}
\pgfpathlineto{\pgfpoint{-10.827092\du}{20.018186\du}}
\pgfpathlineto{\pgfpoint{-10.817837\du}{20.016693\du}}
\pgfpathlineto{\pgfpoint{-10.809776\du}{20.015200\du}}
\pgfpathlineto{\pgfpoint{-10.801417\du}{20.013110\du}}
\pgfpathlineto{\pgfpoint{-10.793057\du}{20.011020\du}}
\pgfpathlineto{\pgfpoint{-10.785594\du}{20.008930\du}}
\pgfpathlineto{\pgfpoint{-10.778130\du}{20.007139\du}}
\pgfpathlineto{\pgfpoint{-10.770964\du}{20.004751\du}}
\pgfpathlineto{\pgfpoint{-10.764396\du}{20.002362\du}}
\pgfpathlineto{\pgfpoint{-10.758127\du}{19.999675\du}}
\pgfpathlineto{\pgfpoint{-10.752454\du}{19.996988\du}}
\pgfpathlineto{\pgfpoint{-10.746782\du}{19.994301\du}}
\pgfpathlineto{\pgfpoint{-10.741706\du}{19.991614\du}}
\pgfpathlineto{\pgfpoint{-10.736930\du}{19.988330\du}}
\pgfpathlineto{\pgfpoint{-10.732750\du}{19.985345\du}}
\pgfpathlineto{\pgfpoint{-10.729167\du}{19.982359\du}}
\pgfpathlineto{\pgfpoint{-10.726182\du}{19.979374\du}}
\pgfpathlineto{\pgfpoint{-10.723495\du}{19.976090\du}}
\pgfpathlineto{\pgfpoint{-10.721106\du}{19.972806\du}}
\pgfpathlineto{\pgfpoint{-10.719315\du}{19.969522\du}}
\pgfpathlineto{\pgfpoint{-10.718419\du}{19.966238\du}}
\pgfpathlineto{\pgfpoint{-10.717524\du}{19.963252\du}}
\pgfpathlineto{\pgfpoint{-10.716927\du}{19.959669\du}}
\pgfusepath{fill}
\pgfsetbuttcap
\pgfsetmiterjoin
\pgfsetdash{}{0pt}
\definecolor{dialinecolor}{rgb}{0.666667, 0.901961, 1.000000}
\pgfsetstrokecolor{dialinecolor}
\pgfpathmoveto{\pgfpoint{-10.729167\du}{19.953997\du}}
\pgfpathlineto{\pgfpoint{-10.729167\du}{19.950414\du}}
\pgfpathlineto{\pgfpoint{-10.730063\du}{19.948026\du}}
\pgfpathlineto{\pgfpoint{-10.731257\du}{19.944742\du}}
\pgfpathlineto{\pgfpoint{-10.732750\du}{19.941756\du}}
\pgfpathlineto{\pgfpoint{-10.735138\du}{19.939069\du}}
\pgfpathlineto{\pgfpoint{-10.737527\du}{19.935487\du}}
\pgfpathlineto{\pgfpoint{-10.740811\du}{19.933098\du}}
\pgfpathlineto{\pgfpoint{-10.744095\du}{19.930411\du}}
\pgfpathlineto{\pgfpoint{-10.748275\du}{19.927724\du}}
\pgfpathlineto{\pgfpoint{-10.752753\du}{19.925336\du}}
\pgfpathlineto{\pgfpoint{-10.757530\du}{19.922350\du}}
\pgfpathlineto{\pgfpoint{-10.763202\du}{19.920261\du}}
\pgfpathlineto{\pgfpoint{-10.768576\du}{19.917574\du}}
\pgfpathlineto{\pgfpoint{-10.774846\du}{19.914887\du}}
\pgfpathlineto{\pgfpoint{-10.780817\du}{19.913095\du}}
\pgfpathlineto{\pgfpoint{-10.788280\du}{19.910707\du}}
\pgfpathlineto{\pgfpoint{-10.795147\du}{19.909214\du}}
\pgfpathlineto{\pgfpoint{-10.802312\du}{19.907124\du}}
\pgfpathlineto{\pgfpoint{-10.810075\du}{19.905333\du}}
\pgfpathlineto{\pgfpoint{-10.818733\du}{19.903542\du}}
\pgfpathlineto{\pgfpoint{-10.827092\du}{19.902347\du}}
\pgfpathlineto{\pgfpoint{-10.835452\du}{19.899959\du}}
\pgfpathlineto{\pgfpoint{-10.844408\du}{19.899660\du}}
\pgfpathlineto{\pgfpoint{-10.853066\du}{19.898168\du}}
\pgfpathlineto{\pgfpoint{-10.862023\du}{19.897272\du}}
\pgfpathlineto{\pgfpoint{-10.871278\du}{19.896376\du}}
\pgfpathlineto{\pgfpoint{-10.880832\du}{19.895779\du}}
\pgfpathlineto{\pgfpoint{-10.890385\du}{19.895481\du}}
\pgfpathlineto{\pgfpoint{-10.899342\du}{19.894884\du}}
\pgfpathlineto{\pgfpoint{-10.909493\du}{19.894585\du}}
\pgfpathlineto{\pgfpoint{-10.919046\du}{19.894585\du}}
\pgfpathlineto{\pgfpoint{-10.919046\du}{19.894585\du}}
\pgfpathlineto{\pgfpoint{-10.928301\du}{19.894585\du}}
\pgfpathlineto{\pgfpoint{-10.938154\du}{19.894884\du}}
\pgfpathlineto{\pgfpoint{-10.947707\du}{19.895481\du}}
\pgfpathlineto{\pgfpoint{-10.957261\du}{19.895779\du}}
\pgfpathlineto{\pgfpoint{-10.966516\du}{19.896376\du}}
\pgfpathlineto{\pgfpoint{-10.976368\du}{19.897272\du}}
\pgfpathlineto{\pgfpoint{-10.985026\du}{19.898168\du}}
\pgfpathlineto{\pgfpoint{-10.993983\du}{19.899660\du}}
\pgfpathlineto{\pgfpoint{-11.002342\du}{19.899959\du}}
\pgfpathlineto{\pgfpoint{-11.011299\du}{19.902347\du}}
\pgfpathlineto{\pgfpoint{-11.019658\du}{19.903542\du}}
\pgfpathlineto{\pgfpoint{-11.027719\du}{19.905333\du}}
\pgfpathlineto{\pgfpoint{-11.035482\du}{19.907124\du}}
\pgfpathlineto{\pgfpoint{-11.042945\du}{19.909214\du}}
\pgfpathlineto{\pgfpoint{-11.049812\du}{19.910707\du}}
\pgfpathlineto{\pgfpoint{-11.057276\du}{19.913095\du}}
\pgfpathlineto{\pgfpoint{-11.063247\du}{19.914887\du}}
\pgfpathlineto{\pgfpoint{-11.069516\du}{19.917574\du}}
\pgfpathlineto{\pgfpoint{-11.075189\du}{19.920261\du}}
\pgfpathlineto{\pgfpoint{-11.080563\du}{19.922350\du}}
\pgfpathlineto{\pgfpoint{-11.085340\du}{19.925336\du}}
\pgfpathlineto{\pgfpoint{-11.089818\du}{19.927724\du}}
\pgfpathlineto{\pgfpoint{-11.093998\du}{19.930411\du}}
\pgfpathlineto{\pgfpoint{-11.097282\du}{19.933098\du}}
\pgfpathlineto{\pgfpoint{-11.100864\du}{19.935487\du}}
\pgfpathlineto{\pgfpoint{-11.102656\du}{19.939069\du}}
\pgfpathlineto{\pgfpoint{-11.105044\du}{19.941756\du}}
\pgfpathlineto{\pgfpoint{-11.106835\du}{19.944742\du}}
\pgfpathlineto{\pgfpoint{-11.108328\du}{19.948026\du}}
\pgfpathlineto{\pgfpoint{-11.108925\du}{19.950414\du}}
\pgfpathlineto{\pgfpoint{-11.109224\du}{19.953997\du}}
\pgfpathlineto{\pgfpoint{-11.109224\du}{19.953997\du}}
\pgfpathlineto{\pgfpoint{-11.108925\du}{19.956982\du}}
\pgfpathlineto{\pgfpoint{-11.108328\du}{19.959669\du}}
\pgfpathlineto{\pgfpoint{-11.106835\du}{19.962655\du}}
\pgfpathlineto{\pgfpoint{-11.105044\du}{19.966238\du}}
\pgfpathlineto{\pgfpoint{-11.102656\du}{19.968924\du}}
\pgfpathlineto{\pgfpoint{-11.100864\du}{19.971910\du}}
\pgfpathlineto{\pgfpoint{-11.097282\du}{19.974597\du}}
\pgfpathlineto{\pgfpoint{-11.093998\du}{19.977284\du}}
\pgfpathlineto{\pgfpoint{-11.089818\du}{19.979971\du}}
\pgfpathlineto{\pgfpoint{-11.085340\du}{19.982658\du}}
\pgfpathlineto{\pgfpoint{-11.080563\du}{19.985345\du}}
\pgfpathlineto{\pgfpoint{-11.075189\du}{19.987733\du}}
\pgfpathlineto{\pgfpoint{-11.069516\du}{19.990420\du}}
\pgfpathlineto{\pgfpoint{-11.063247\du}{19.992510\du}}
\pgfpathlineto{\pgfpoint{-11.057276\du}{19.994899\du}}
\pgfpathlineto{\pgfpoint{-11.049812\du}{19.997287\du}}
\pgfpathlineto{\pgfpoint{-11.042945\du}{19.999078\du}}
\pgfpathlineto{\pgfpoint{-11.035482\du}{20.000571\du}}
\pgfpathlineto{\pgfpoint{-11.027719\du}{20.002661\du}}
\pgfpathlineto{\pgfpoint{-11.019658\du}{20.004452\du}}
\pgfpathlineto{\pgfpoint{-11.011299\du}{20.005646\du}}
\pgfpathlineto{\pgfpoint{-11.002342\du}{20.007438\du}}
\pgfpathlineto{\pgfpoint{-10.993983\du}{20.008333\du}}
\pgfpathlineto{\pgfpoint{-10.985026\du}{20.009826\du}}
\pgfpathlineto{\pgfpoint{-10.976368\du}{20.010722\du}}
\pgfpathlineto{\pgfpoint{-10.966516\du}{20.011020\du}}
\pgfpathlineto{\pgfpoint{-10.957261\du}{20.012215\du}}
\pgfpathlineto{\pgfpoint{-10.947707\du}{20.012513\du}}
\pgfpathlineto{\pgfpoint{-10.938154\du}{20.012812\du}}
\pgfpathlineto{\pgfpoint{-10.928301\du}{20.013409\du}}
\pgfpathlineto{\pgfpoint{-10.919046\du}{20.013409\du}}
\pgfpathlineto{\pgfpoint{-10.919046\du}{20.013409\du}}
\pgfpathlineto{\pgfpoint{-10.909493\du}{20.013409\du}}
\pgfpathlineto{\pgfpoint{-10.899342\du}{20.012812\du}}
\pgfpathlineto{\pgfpoint{-10.890385\du}{20.012513\du}}
\pgfpathlineto{\pgfpoint{-10.880832\du}{20.012215\du}}
\pgfpathlineto{\pgfpoint{-10.871278\du}{20.011020\du}}
\pgfpathlineto{\pgfpoint{-10.862023\du}{20.010722\du}}
\pgfpathlineto{\pgfpoint{-10.853066\du}{20.009826\du}}
\pgfpathlineto{\pgfpoint{-10.844408\du}{20.008333\du}}
\pgfpathlineto{\pgfpoint{-10.835452\du}{20.007438\du}}
\pgfpathlineto{\pgfpoint{-10.827092\du}{20.005646\du}}
\pgfpathlineto{\pgfpoint{-10.818733\du}{20.004452\du}}
\pgfpathlineto{\pgfpoint{-10.810075\du}{20.002661\du}}
\pgfpathlineto{\pgfpoint{-10.802312\du}{20.000571\du}}
\pgfpathlineto{\pgfpoint{-10.795147\du}{19.999078\du}}
\pgfpathlineto{\pgfpoint{-10.788280\du}{19.997287\du}}
\pgfpathlineto{\pgfpoint{-10.780817\du}{19.994899\du}}
\pgfpathlineto{\pgfpoint{-10.774846\du}{19.992510\du}}
\pgfpathlineto{\pgfpoint{-10.768576\du}{19.990420\du}}
\pgfpathlineto{\pgfpoint{-10.763202\du}{19.987733\du}}
\pgfpathlineto{\pgfpoint{-10.757530\du}{19.985345\du}}
\pgfpathlineto{\pgfpoint{-10.752753\du}{19.982658\du}}
\pgfpathlineto{\pgfpoint{-10.748275\du}{19.979971\du}}
\pgfpathlineto{\pgfpoint{-10.744095\du}{19.977284\du}}
\pgfpathlineto{\pgfpoint{-10.740811\du}{19.974597\du}}
\pgfpathlineto{\pgfpoint{-10.737527\du}{19.971910\du}}
\pgfpathlineto{\pgfpoint{-10.735138\du}{19.968924\du}}
\pgfpathlineto{\pgfpoint{-10.732750\du}{19.966238\du}}
\pgfpathlineto{\pgfpoint{-10.731257\du}{19.962655\du}}
\pgfpathlineto{\pgfpoint{-10.730063\du}{19.959669\du}}
\pgfpathlineto{\pgfpoint{-10.729167\du}{19.956982\du}}
\pgfpathlineto{\pgfpoint{-10.729167\du}{19.953997\du}}
\pgfusepath{stroke}
\pgfsetbuttcap
\pgfsetmiterjoin
\pgfsetdash{}{0pt}
\definecolor{dialinecolor}{rgb}{0.000000, 0.470588, 0.666667}
\pgfsetfillcolor{dialinecolor}
\pgfpathmoveto{\pgfpoint{-11.109224\du}{19.870701\du}}
\pgfpathlineto{\pgfpoint{-11.109224\du}{19.965342\du}}
\pgfpathlineto{\pgfpoint{-10.716628\du}{19.965342\du}}
\pgfpathlineto{\pgfpoint{-10.716628\du}{19.870701\du}}
\pgfpathlineto{\pgfpoint{-11.109224\du}{19.870701\du}}
\pgfusepath{fill}
\pgfsetbuttcap
\pgfsetmiterjoin
\pgfsetdash{}{0pt}
\definecolor{dialinecolor}{rgb}{0.000000, 0.705882, 1.000000}
\pgfsetfillcolor{dialinecolor}
\pgfpathmoveto{\pgfpoint{-10.716927\du}{19.864133\du}}
\pgfpathlineto{\pgfpoint{-10.717524\du}{19.860849\du}}
\pgfpathlineto{\pgfpoint{-10.718419\du}{19.857266\du}}
\pgfpathlineto{\pgfpoint{-10.719315\du}{19.854281\du}}
\pgfpathlineto{\pgfpoint{-10.721106\du}{19.850996\du}}
\pgfpathlineto{\pgfpoint{-10.723495\du}{19.847712\du}}
\pgfpathlineto{\pgfpoint{-10.726182\du}{19.844428\du}}
\pgfpathlineto{\pgfpoint{-10.729167\du}{19.841443\du}}
\pgfpathlineto{\pgfpoint{-10.732750\du}{19.838159\du}}
\pgfpathlineto{\pgfpoint{-10.736930\du}{19.835472\du}}
\pgfpathlineto{\pgfpoint{-10.741706\du}{19.832486\du}}
\pgfpathlineto{\pgfpoint{-10.746782\du}{19.829501\du}}
\pgfpathlineto{\pgfpoint{-10.752454\du}{19.826814\du}}
\pgfpathlineto{\pgfpoint{-10.758127\du}{19.823530\du}}
\pgfpathlineto{\pgfpoint{-10.764396\du}{19.821738\du}}
\pgfpathlineto{\pgfpoint{-10.770964\du}{19.819051\du}}
\pgfpathlineto{\pgfpoint{-10.778130\du}{19.816663\du}}
\pgfpathlineto{\pgfpoint{-10.785594\du}{19.814573\du}}
\pgfpathlineto{\pgfpoint{-10.793057\du}{19.812483\du}}
\pgfpathlineto{\pgfpoint{-10.801417\du}{19.810393\du}}
\pgfpathlineto{\pgfpoint{-10.809776\du}{19.808304\du}}
\pgfpathlineto{\pgfpoint{-10.817837\du}{19.806811\du}}
\pgfpathlineto{\pgfpoint{-10.827092\du}{19.805318\du}}
\pgfpathlineto{\pgfpoint{-10.835750\du}{19.804124\du}}
\pgfpathlineto{\pgfpoint{-10.845304\du}{19.802631\du}}
\pgfpathlineto{\pgfpoint{-10.854559\du}{19.801735\du}}
\pgfpathlineto{\pgfpoint{-10.863814\du}{19.800541\du}}
\pgfpathlineto{\pgfpoint{-10.873666\du}{19.799944\du}}
\pgfpathlineto{\pgfpoint{-10.883519\du}{19.799347\du}}
\pgfpathlineto{\pgfpoint{-10.893371\du}{19.799048\du}}
\pgfpathlineto{\pgfpoint{-10.903522\du}{19.798750\du}}
\pgfpathlineto{\pgfpoint{-10.913374\du}{19.798750\du}}
\pgfpathlineto{\pgfpoint{-10.913374\du}{19.798750\du}}
\pgfpathlineto{\pgfpoint{-10.923226\du}{19.798750\du}}
\pgfpathlineto{\pgfpoint{-10.933078\du}{19.799048\du}}
\pgfpathlineto{\pgfpoint{-10.943229\du}{19.799347\du}}
\pgfpathlineto{\pgfpoint{-10.953081\du}{19.799944\du}}
\pgfpathlineto{\pgfpoint{-10.962635\du}{19.800541\du}}
\pgfpathlineto{\pgfpoint{-10.972189\du}{19.801735\du}}
\pgfpathlineto{\pgfpoint{-10.981444\du}{19.802631\du}}
\pgfpathlineto{\pgfpoint{-10.990699\du}{19.804124\du}}
\pgfpathlineto{\pgfpoint{-10.999655\du}{19.805318\du}}
\pgfpathlineto{\pgfpoint{-11.008612\du}{19.806811\du}}
\pgfpathlineto{\pgfpoint{-11.016971\du}{19.808304\du}}
\pgfpathlineto{\pgfpoint{-11.025331\du}{19.810393\du}}
\pgfpathlineto{\pgfpoint{-11.033690\du}{19.812483\du}}
\pgfpathlineto{\pgfpoint{-11.040856\du}{19.814573\du}}
\pgfpathlineto{\pgfpoint{-11.048319\du}{19.816663\du}}
\pgfpathlineto{\pgfpoint{-11.055485\du}{19.819051\du}}
\pgfpathlineto{\pgfpoint{-11.062351\du}{19.821738\du}}
\pgfpathlineto{\pgfpoint{-11.068322\du}{19.823530\du}}
\pgfpathlineto{\pgfpoint{-11.074293\du}{19.826814\du}}
\pgfpathlineto{\pgfpoint{-11.079966\du}{19.829501\du}}
\pgfpathlineto{\pgfpoint{-11.085041\du}{19.832486\du}}
\pgfpathlineto{\pgfpoint{-11.089818\du}{19.835472\du}}
\pgfpathlineto{\pgfpoint{-11.093998\du}{19.838159\du}}
\pgfpathlineto{\pgfpoint{-11.097282\du}{19.841443\du}}
\pgfpathlineto{\pgfpoint{-11.100566\du}{19.844428\du}}
\pgfpathlineto{\pgfpoint{-11.103551\du}{19.847712\du}}
\pgfpathlineto{\pgfpoint{-11.105641\du}{19.850996\du}}
\pgfpathlineto{\pgfpoint{-11.107134\du}{19.854281\du}}
\pgfpathlineto{\pgfpoint{-11.108925\du}{19.857266\du}}
\pgfpathlineto{\pgfpoint{-11.109224\du}{19.860849\du}}
\pgfpathlineto{\pgfpoint{-11.109522\du}{19.864133\du}}
\pgfpathlineto{\pgfpoint{-11.109522\du}{19.864133\du}}
\pgfpathlineto{\pgfpoint{-11.109224\du}{19.867417\du}}
\pgfpathlineto{\pgfpoint{-11.108925\du}{19.870999\du}}
\pgfpathlineto{\pgfpoint{-11.107134\du}{19.873985\du}}
\pgfpathlineto{\pgfpoint{-11.105641\du}{19.877269\du}}
\pgfpathlineto{\pgfpoint{-11.103551\du}{19.880852\du}}
\pgfpathlineto{\pgfpoint{-11.100566\du}{19.883837\du}}
\pgfpathlineto{\pgfpoint{-11.097282\du}{19.886823\du}}
\pgfpathlineto{\pgfpoint{-11.093998\du}{19.890107\du}}
\pgfpathlineto{\pgfpoint{-11.089818\du}{19.892495\du}}
\pgfpathlineto{\pgfpoint{-11.085041\du}{19.895779\du}}
\pgfpathlineto{\pgfpoint{-11.079966\du}{19.899063\du}}
\pgfpathlineto{\pgfpoint{-11.074293\du}{19.901750\du}}
\pgfpathlineto{\pgfpoint{-11.068322\du}{19.904437\du}}
\pgfpathlineto{\pgfpoint{-11.062351\du}{19.906527\du}}
\pgfpathlineto{\pgfpoint{-11.055485\du}{19.909513\du}}
\pgfpathlineto{\pgfpoint{-11.048319\du}{19.911901\du}}
\pgfpathlineto{\pgfpoint{-11.040856\du}{19.913692\du}}
\pgfpathlineto{\pgfpoint{-11.033690\du}{19.915782\du}}
\pgfpathlineto{\pgfpoint{-11.025331\du}{19.917872\du}}
\pgfpathlineto{\pgfpoint{-11.016971\du}{19.919962\du}}
\pgfpathlineto{\pgfpoint{-11.008612\du}{19.921455\du}}
\pgfpathlineto{\pgfpoint{-10.999655\du}{19.922947\du}}
\pgfpathlineto{\pgfpoint{-10.990699\du}{19.924142\du}}
\pgfpathlineto{\pgfpoint{-10.981444\du}{19.925634\du}}
\pgfpathlineto{\pgfpoint{-10.972189\du}{19.926530\du}}
\pgfpathlineto{\pgfpoint{-10.962635\du}{19.927724\du}}
\pgfpathlineto{\pgfpoint{-10.953081\du}{19.928023\du}}
\pgfpathlineto{\pgfpoint{-10.943229\du}{19.928919\du}}
\pgfpathlineto{\pgfpoint{-10.933078\du}{19.928919\du}}
\pgfpathlineto{\pgfpoint{-10.923226\du}{19.929814\du}}
\pgfpathlineto{\pgfpoint{-10.913374\du}{19.929814\du}}
\pgfpathlineto{\pgfpoint{-10.913374\du}{19.929814\du}}
\pgfpathlineto{\pgfpoint{-10.903522\du}{19.929814\du}}
\pgfpathlineto{\pgfpoint{-10.893371\du}{19.928919\du}}
\pgfpathlineto{\pgfpoint{-10.883519\du}{19.928919\du}}
\pgfpathlineto{\pgfpoint{-10.873666\du}{19.928023\du}}
\pgfpathlineto{\pgfpoint{-10.863814\du}{19.927724\du}}
\pgfpathlineto{\pgfpoint{-10.854559\du}{19.926530\du}}
\pgfpathlineto{\pgfpoint{-10.845304\du}{19.925634\du}}
\pgfpathlineto{\pgfpoint{-10.835750\du}{19.924142\du}}
\pgfpathlineto{\pgfpoint{-10.827092\du}{19.922947\du}}
\pgfpathlineto{\pgfpoint{-10.817837\du}{19.921455\du}}
\pgfpathlineto{\pgfpoint{-10.809776\du}{19.919962\du}}
\pgfpathlineto{\pgfpoint{-10.801417\du}{19.917872\du}}
\pgfpathlineto{\pgfpoint{-10.793057\du}{19.915782\du}}
\pgfpathlineto{\pgfpoint{-10.785594\du}{19.913692\du}}
\pgfpathlineto{\pgfpoint{-10.778130\du}{19.911901\du}}
\pgfpathlineto{\pgfpoint{-10.770964\du}{19.909513\du}}
\pgfpathlineto{\pgfpoint{-10.764396\du}{19.906527\du}}
\pgfpathlineto{\pgfpoint{-10.758127\du}{19.904437\du}}
\pgfpathlineto{\pgfpoint{-10.752454\du}{19.901750\du}}
\pgfpathlineto{\pgfpoint{-10.746782\du}{19.899063\du}}
\pgfpathlineto{\pgfpoint{-10.741706\du}{19.895779\du}}
\pgfpathlineto{\pgfpoint{-10.736930\du}{19.892495\du}}
\pgfpathlineto{\pgfpoint{-10.732750\du}{19.890107\du}}
\pgfpathlineto{\pgfpoint{-10.729167\du}{19.886823\du}}
\pgfpathlineto{\pgfpoint{-10.726182\du}{19.883837\du}}
\pgfpathlineto{\pgfpoint{-10.723495\du}{19.880852\du}}
\pgfpathlineto{\pgfpoint{-10.721106\du}{19.877269\du}}
\pgfpathlineto{\pgfpoint{-10.719315\du}{19.873985\du}}
\pgfpathlineto{\pgfpoint{-10.718419\du}{19.870999\du}}
\pgfpathlineto{\pgfpoint{-10.717524\du}{19.867417\du}}
\pgfpathlineto{\pgfpoint{-10.716927\du}{19.864133\du}}
\pgfusepath{fill}
\pgfsetbuttcap
\pgfsetmiterjoin
\pgfsetdash{}{0pt}
\definecolor{dialinecolor}{rgb}{0.666667, 0.901961, 1.000000}
\pgfsetstrokecolor{dialinecolor}
\pgfpathmoveto{\pgfpoint{-10.729167\du}{19.858460\du}}
\pgfpathlineto{\pgfpoint{-10.729167\du}{19.855176\du}}
\pgfpathlineto{\pgfpoint{-10.730063\du}{19.852788\du}}
\pgfpathlineto{\pgfpoint{-10.731257\du}{19.849504\du}}
\pgfpathlineto{\pgfpoint{-10.732750\du}{19.845921\du}}
\pgfpathlineto{\pgfpoint{-10.735138\du}{19.843533\du}}
\pgfpathlineto{\pgfpoint{-10.737527\du}{19.840249\du}}
\pgfpathlineto{\pgfpoint{-10.740811\du}{19.837860\du}}
\pgfpathlineto{\pgfpoint{-10.744095\du}{19.835173\du}}
\pgfpathlineto{\pgfpoint{-10.748275\du}{19.832486\du}}
\pgfpathlineto{\pgfpoint{-10.752753\du}{19.829799\du}}
\pgfpathlineto{\pgfpoint{-10.757530\du}{19.827112\du}}
\pgfpathlineto{\pgfpoint{-10.763202\du}{19.824425\du}}
\pgfpathlineto{\pgfpoint{-10.768576\du}{19.822037\du}}
\pgfpathlineto{\pgfpoint{-10.774846\du}{19.819947\du}}
\pgfpathlineto{\pgfpoint{-10.780817\du}{19.817559\du}}
\pgfpathlineto{\pgfpoint{-10.788280\du}{19.815170\du}}
\pgfpathlineto{\pgfpoint{-10.795147\du}{19.813379\du}}
\pgfpathlineto{\pgfpoint{-10.802312\du}{19.811886\du}}
\pgfpathlineto{\pgfpoint{-10.810075\du}{19.809498\du}}
\pgfpathlineto{\pgfpoint{-10.818733\du}{19.808005\du}}
\pgfpathlineto{\pgfpoint{-10.827092\du}{19.806811\du}}
\pgfpathlineto{\pgfpoint{-10.835452\du}{19.805019\du}}
\pgfpathlineto{\pgfpoint{-10.844408\du}{19.804124\du}}
\pgfpathlineto{\pgfpoint{-10.853066\du}{19.802631\du}}
\pgfpathlineto{\pgfpoint{-10.862023\du}{19.801735\du}}
\pgfpathlineto{\pgfpoint{-10.871278\du}{19.801138\du}}
\pgfpathlineto{\pgfpoint{-10.880832\du}{19.800243\du}}
\pgfpathlineto{\pgfpoint{-10.890385\du}{19.799944\du}}
\pgfpathlineto{\pgfpoint{-10.899342\du}{19.799645\du}}
\pgfpathlineto{\pgfpoint{-10.909493\du}{19.799048\du}}
\pgfpathlineto{\pgfpoint{-10.919046\du}{19.799048\du}}
\pgfpathlineto{\pgfpoint{-10.919046\du}{19.799048\du}}
\pgfpathlineto{\pgfpoint{-10.928301\du}{19.799048\du}}
\pgfpathlineto{\pgfpoint{-10.938154\du}{19.799645\du}}
\pgfpathlineto{\pgfpoint{-10.947707\du}{19.799944\du}}
\pgfpathlineto{\pgfpoint{-10.957261\du}{19.800243\du}}
\pgfpathlineto{\pgfpoint{-10.966516\du}{19.801138\du}}
\pgfpathlineto{\pgfpoint{-10.976368\du}{19.801735\du}}
\pgfpathlineto{\pgfpoint{-10.985026\du}{19.802631\du}}
\pgfpathlineto{\pgfpoint{-10.993983\du}{19.804124\du}}
\pgfpathlineto{\pgfpoint{-11.002342\du}{19.805019\du}}
\pgfpathlineto{\pgfpoint{-11.011299\du}{19.806811\du}}
\pgfpathlineto{\pgfpoint{-11.019658\du}{19.808005\du}}
\pgfpathlineto{\pgfpoint{-11.027719\du}{19.809498\du}}
\pgfpathlineto{\pgfpoint{-11.035482\du}{19.811886\du}}
\pgfpathlineto{\pgfpoint{-11.042945\du}{19.813379\du}}
\pgfpathlineto{\pgfpoint{-11.049812\du}{19.815170\du}}
\pgfpathlineto{\pgfpoint{-11.057276\du}{19.817559\du}}
\pgfpathlineto{\pgfpoint{-11.063247\du}{19.819947\du}}
\pgfpathlineto{\pgfpoint{-11.069516\du}{19.822037\du}}
\pgfpathlineto{\pgfpoint{-11.075189\du}{19.824425\du}}
\pgfpathlineto{\pgfpoint{-11.080563\du}{19.827112\du}}
\pgfpathlineto{\pgfpoint{-11.085340\du}{19.829799\du}}
\pgfpathlineto{\pgfpoint{-11.089818\du}{19.832486\du}}
\pgfpathlineto{\pgfpoint{-11.093998\du}{19.835173\du}}
\pgfpathlineto{\pgfpoint{-11.097282\du}{19.837860\du}}
\pgfpathlineto{\pgfpoint{-11.100864\du}{19.840249\du}}
\pgfpathlineto{\pgfpoint{-11.102656\du}{19.843533\du}}
\pgfpathlineto{\pgfpoint{-11.105044\du}{19.845921\du}}
\pgfpathlineto{\pgfpoint{-11.106835\du}{19.849504\du}}
\pgfpathlineto{\pgfpoint{-11.108328\du}{19.852788\du}}
\pgfpathlineto{\pgfpoint{-11.108925\du}{19.855176\du}}
\pgfpathlineto{\pgfpoint{-11.109224\du}{19.858460\du}}
\pgfpathlineto{\pgfpoint{-11.109224\du}{19.858460\du}}
\pgfpathlineto{\pgfpoint{-11.108925\du}{19.861744\du}}
\pgfpathlineto{\pgfpoint{-11.108328\du}{19.864431\du}}
\pgfpathlineto{\pgfpoint{-11.106835\du}{19.867417\du}}
\pgfpathlineto{\pgfpoint{-11.105044\du}{19.870701\du}}
\pgfpathlineto{\pgfpoint{-11.102656\du}{19.873388\du}}
\pgfpathlineto{\pgfpoint{-11.100864\du}{19.876672\du}}
\pgfpathlineto{\pgfpoint{-11.097282\du}{19.879359\du}}
\pgfpathlineto{\pgfpoint{-11.093998\du}{19.882046\du}}
\pgfpathlineto{\pgfpoint{-11.089818\du}{19.884733\du}}
\pgfpathlineto{\pgfpoint{-11.085340\du}{19.887420\du}}
\pgfpathlineto{\pgfpoint{-11.080563\du}{19.890107\du}}
\pgfpathlineto{\pgfpoint{-11.075189\du}{19.892197\du}}
\pgfpathlineto{\pgfpoint{-11.069516\du}{19.894884\du}}
\pgfpathlineto{\pgfpoint{-11.063247\du}{19.897272\du}}
\pgfpathlineto{\pgfpoint{-11.057276\du}{19.899660\du}}
\pgfpathlineto{\pgfpoint{-11.049812\du}{19.901750\du}}
\pgfpathlineto{\pgfpoint{-11.042945\du}{19.903542\du}}
\pgfpathlineto{\pgfpoint{-11.035482\du}{19.905333\du}}
\pgfpathlineto{\pgfpoint{-11.027719\du}{19.907124\du}}
\pgfpathlineto{\pgfpoint{-11.019658\du}{19.909214\du}}
\pgfpathlineto{\pgfpoint{-11.011299\du}{19.910408\du}}
\pgfpathlineto{\pgfpoint{-11.002342\du}{19.912200\du}}
\pgfpathlineto{\pgfpoint{-10.993983\du}{19.913095\du}}
\pgfpathlineto{\pgfpoint{-10.985026\du}{19.914588\du}}
\pgfpathlineto{\pgfpoint{-10.976368\du}{19.914887\du}}
\pgfpathlineto{\pgfpoint{-10.966516\du}{19.916081\du}}
\pgfpathlineto{\pgfpoint{-10.957261\du}{19.916379\du}}
\pgfpathlineto{\pgfpoint{-10.947707\du}{19.917275\du}}
\pgfpathlineto{\pgfpoint{-10.938154\du}{19.917574\du}}
\pgfpathlineto{\pgfpoint{-10.928301\du}{19.917872\du}}
\pgfpathlineto{\pgfpoint{-10.919046\du}{19.917872\du}}
\pgfpathlineto{\pgfpoint{-10.919046\du}{19.917872\du}}
\pgfpathlineto{\pgfpoint{-10.909493\du}{19.917872\du}}
\pgfpathlineto{\pgfpoint{-10.899342\du}{19.917574\du}}
\pgfpathlineto{\pgfpoint{-10.890385\du}{19.917275\du}}
\pgfpathlineto{\pgfpoint{-10.880832\du}{19.916379\du}}
\pgfpathlineto{\pgfpoint{-10.871278\du}{19.916081\du}}
\pgfpathlineto{\pgfpoint{-10.862023\du}{19.914887\du}}
\pgfpathlineto{\pgfpoint{-10.853066\du}{19.914588\du}}
\pgfpathlineto{\pgfpoint{-10.844408\du}{19.913095\du}}
\pgfpathlineto{\pgfpoint{-10.835452\du}{19.912200\du}}
\pgfpathlineto{\pgfpoint{-10.827092\du}{19.910408\du}}
\pgfpathlineto{\pgfpoint{-10.818733\du}{19.909214\du}}
\pgfpathlineto{\pgfpoint{-10.810075\du}{19.907124\du}}
\pgfpathlineto{\pgfpoint{-10.802312\du}{19.905333\du}}
\pgfpathlineto{\pgfpoint{-10.795147\du}{19.903542\du}}
\pgfpathlineto{\pgfpoint{-10.788280\du}{19.901750\du}}
\pgfpathlineto{\pgfpoint{-10.780817\du}{19.899660\du}}
\pgfpathlineto{\pgfpoint{-10.774846\du}{19.897272\du}}
\pgfpathlineto{\pgfpoint{-10.768576\du}{19.894884\du}}
\pgfpathlineto{\pgfpoint{-10.763202\du}{19.892197\du}}
\pgfpathlineto{\pgfpoint{-10.757530\du}{19.890107\du}}
\pgfpathlineto{\pgfpoint{-10.752753\du}{19.887420\du}}
\pgfpathlineto{\pgfpoint{-10.748275\du}{19.884733\du}}
\pgfpathlineto{\pgfpoint{-10.744095\du}{19.882046\du}}
\pgfpathlineto{\pgfpoint{-10.740811\du}{19.879359\du}}
\pgfpathlineto{\pgfpoint{-10.737527\du}{19.876672\du}}
\pgfpathlineto{\pgfpoint{-10.735138\du}{19.873388\du}}
\pgfpathlineto{\pgfpoint{-10.732750\du}{19.870701\du}}
\pgfpathlineto{\pgfpoint{-10.731257\du}{19.867417\du}}
\pgfpathlineto{\pgfpoint{-10.730063\du}{19.864431\du}}
\pgfpathlineto{\pgfpoint{-10.729167\du}{19.861744\du}}
\pgfpathlineto{\pgfpoint{-10.729167\du}{19.858460\du}}
\pgfusepath{stroke}
\pgfsetbuttcap
\pgfsetmiterjoin
\pgfsetdash{}{0pt}
\definecolor{dialinecolor}{rgb}{0.000000, 0.000000, 0.000000}
\pgfsetfillcolor{dialinecolor}
\pgfpathmoveto{\pgfpoint{-10.906806\du}{19.846817\du}}
\pgfpathlineto{\pgfpoint{-10.883519\du}{19.858460\du}}
\pgfpathlineto{\pgfpoint{-10.811568\du}{19.834875\du}}
\pgfpathlineto{\pgfpoint{-10.776040\du}{19.846817\du}}
\pgfpathlineto{\pgfpoint{-10.799625\du}{19.822933\du}}
\pgfpathlineto{\pgfpoint{-10.883519\du}{19.822933\du}}
\pgfpathlineto{\pgfpoint{-10.847692\du}{19.822933\du}}
\pgfpathlineto{\pgfpoint{-10.906806\du}{19.846817\du}}
\pgfusepath{fill}
\pgfsetbuttcap
\pgfsetmiterjoin
\pgfsetdash{}{0pt}
\definecolor{dialinecolor}{rgb}{0.000000, 0.000000, 0.000000}
\pgfsetfillcolor{dialinecolor}
\pgfpathmoveto{\pgfpoint{-10.919046\du}{19.882344\du}}
\pgfpathlineto{\pgfpoint{-10.954873\du}{19.870701\du}}
\pgfpathlineto{\pgfpoint{-11.013986\du}{19.894585\du}}
\pgfpathlineto{\pgfpoint{-11.049812\du}{19.882344\du}}
\pgfpathlineto{\pgfpoint{-11.037571\du}{19.917872\du}}
\pgfpathlineto{\pgfpoint{-10.954873\du}{19.917872\du}}
\pgfpathlineto{\pgfpoint{-10.990400\du}{19.906527\du}}
\pgfpathlineto{\pgfpoint{-10.919046\du}{19.882344\du}}
\pgfusepath{fill}
\pgfsetbuttcap
\pgfsetmiterjoin
\pgfsetdash{}{0pt}
\definecolor{dialinecolor}{rgb}{0.000000, 0.000000, 0.000000}
\pgfsetfillcolor{dialinecolor}
\pgfpathmoveto{\pgfpoint{-11.037571\du}{19.822933\du}}
\pgfpathlineto{\pgfpoint{-11.013986\du}{19.810393\du}}
\pgfpathlineto{\pgfpoint{-10.942333\du}{19.846817\du}}
\pgfpathlineto{\pgfpoint{-10.919046\du}{19.834875\du}}
\pgfpathlineto{\pgfpoint{-10.930690\du}{19.858460\du}}
\pgfpathlineto{\pgfpoint{-11.013986\du}{19.858460\du}}
\pgfpathlineto{\pgfpoint{-10.978160\du}{19.846817\du}}
\pgfpathlineto{\pgfpoint{-11.037571\du}{19.822933\du}}
\pgfusepath{fill}
\pgfsetbuttcap
\pgfsetmiterjoin
\pgfsetdash{}{0pt}
\definecolor{dialinecolor}{rgb}{0.000000, 0.000000, 0.000000}
\pgfsetfillcolor{dialinecolor}
\pgfpathmoveto{\pgfpoint{-10.788280\du}{19.906527\du}}
\pgfpathlineto{\pgfpoint{-10.811568\du}{19.917872\du}}
\pgfpathlineto{\pgfpoint{-10.883519\du}{19.882344\du}}
\pgfpathlineto{\pgfpoint{-10.919046\du}{19.894585\du}}
\pgfpathlineto{\pgfpoint{-10.895162\du}{19.870701\du}}
\pgfpathlineto{\pgfpoint{-10.811568\du}{19.870701\du}}
\pgfpathlineto{\pgfpoint{-10.847692\du}{19.882344\du}}
\pgfpathlineto{\pgfpoint{-10.788280\du}{19.906527\du}}
\pgfusepath{fill}
\pgfsetbuttcap
\pgfsetmiterjoin
\pgfsetdash{}{0pt}
\definecolor{dialinecolor}{rgb}{1.000000, 1.000000, 1.000000}
\pgfsetfillcolor{dialinecolor}
\pgfpathmoveto{\pgfpoint{-10.906806\du}{19.858460\du}}
\pgfpathlineto{\pgfpoint{-10.883519\du}{19.858460\du}}
\pgfpathlineto{\pgfpoint{-10.811568\du}{19.834875\du}}
\pgfpathlineto{\pgfpoint{-10.776040\du}{19.846817\du}}
\pgfpathlineto{\pgfpoint{-10.799625\du}{19.822933\du}}
\pgfpathlineto{\pgfpoint{-10.871278\du}{19.822933\du}}
\pgfpathlineto{\pgfpoint{-10.847692\du}{19.822933\du}}
\pgfpathlineto{\pgfpoint{-10.906806\du}{19.858460\du}}
\pgfusepath{fill}
\pgfsetbuttcap
\pgfsetmiterjoin
\pgfsetdash{}{0pt}
\definecolor{dialinecolor}{rgb}{1.000000, 1.000000, 1.000000}
\pgfsetfillcolor{dialinecolor}
\pgfpathmoveto{\pgfpoint{-10.919046\du}{19.882344\du}}
\pgfpathlineto{\pgfpoint{-10.942333\du}{19.870701\du}}
\pgfpathlineto{\pgfpoint{-11.013986\du}{19.894585\du}}
\pgfpathlineto{\pgfpoint{-11.049812\du}{19.894585\du}}
\pgfpathlineto{\pgfpoint{-11.026226\du}{19.917872\du}}
\pgfpathlineto{\pgfpoint{-10.942333\du}{19.917872\du}}
\pgfpathlineto{\pgfpoint{-10.978160\du}{19.906527\du}}
\pgfpathlineto{\pgfpoint{-10.919046\du}{19.882344\du}}
\pgfusepath{fill}
\pgfsetbuttcap
\pgfsetmiterjoin
\pgfsetdash{}{0pt}
\definecolor{dialinecolor}{rgb}{1.000000, 1.000000, 1.000000}
\pgfsetfillcolor{dialinecolor}
\pgfpathmoveto{\pgfpoint{-11.037571\du}{19.822933\du}}
\pgfpathlineto{\pgfpoint{-11.013986\du}{19.822933\du}}
\pgfpathlineto{\pgfpoint{-10.942333\du}{19.846817\du}}
\pgfpathlineto{\pgfpoint{-10.906806\du}{19.834875\du}}
\pgfpathlineto{\pgfpoint{-10.930690\du}{19.858460\du}}
\pgfpathlineto{\pgfpoint{-11.013986\du}{19.858460\du}}
\pgfpathlineto{\pgfpoint{-10.978160\du}{19.858460\du}}
\pgfpathlineto{\pgfpoint{-11.037571\du}{19.822933\du}}
\pgfusepath{fill}
\pgfsetbuttcap
\pgfsetmiterjoin
\pgfsetdash{}{0pt}
\definecolor{dialinecolor}{rgb}{1.000000, 1.000000, 1.000000}
\pgfsetfillcolor{dialinecolor}
\pgfpathmoveto{\pgfpoint{-10.788280\du}{19.906527\du}}
\pgfpathlineto{\pgfpoint{-10.811568\du}{19.917872\du}}
\pgfpathlineto{\pgfpoint{-10.883519\du}{19.894585\du}}
\pgfpathlineto{\pgfpoint{-10.906806\du}{19.894585\du}}
\pgfpathlineto{\pgfpoint{-10.895162\du}{19.870701\du}}
\pgfpathlineto{\pgfpoint{-10.811568\du}{19.870701\du}}
\pgfpathlineto{\pgfpoint{-10.847692\du}{19.882344\du}}
\pgfpathlineto{\pgfpoint{-10.788280\du}{19.906527\du}}
\pgfusepath{fill}
\pgfsetbuttcap
\pgfsetmiterjoin
\pgfsetdash{}{0pt}
\definecolor{dialinecolor}{rgb}{0.666667, 0.901961, 1.000000}
\pgfsetstrokecolor{dialinecolor}
\pgfpathmoveto{\pgfpoint{-11.109224\du}{19.858460\du}}
\pgfpathlineto{\pgfpoint{-11.109224\du}{19.953400\du}}
\pgfusepath{stroke}
\pgfsetbuttcap
\pgfsetmiterjoin
\pgfsetdash{}{0pt}
\definecolor{dialinecolor}{rgb}{0.666667, 0.901961, 1.000000}
\pgfsetstrokecolor{dialinecolor}
\pgfpathmoveto{\pgfpoint{-10.728570\du}{19.858460\du}}
\pgfpathlineto{\pgfpoint{-10.728570\du}{19.953400\du}}
\pgfusepath{stroke}
\pgfsetbuttcap
\pgfsetmiterjoin
\pgfsetdash{}{0pt}
\definecolor{dialinecolor}{rgb}{0.000000, 0.470588, 0.666667}
\pgfsetfillcolor{dialinecolor}
\pgfpathmoveto{\pgfpoint{-10.002790\du}{19.959669\du}}
\pgfpathlineto{\pgfpoint{-10.002790\du}{19.956385\du}}
\pgfpathlineto{\pgfpoint{-10.003686\du}{19.953101\du}}
\pgfpathlineto{\pgfpoint{-10.004880\du}{19.949817\du}}
\pgfpathlineto{\pgfpoint{-10.006373\du}{19.946533\du}}
\pgfpathlineto{\pgfpoint{-10.008761\du}{19.943548\du}}
\pgfpathlineto{\pgfpoint{-10.011448\du}{19.940264\du}}
\pgfpathlineto{\pgfpoint{-10.014434\du}{19.936979\du}}
\pgfpathlineto{\pgfpoint{-10.018613\du}{19.933994\du}}
\pgfpathlineto{\pgfpoint{-10.022196\du}{19.931008\du}}
\pgfpathlineto{\pgfpoint{-10.026973\du}{19.928023\du}}
\pgfpathlineto{\pgfpoint{-10.032048\du}{19.925336\du}}
\pgfpathlineto{\pgfpoint{-10.037422\du}{19.922350\du}}
\pgfpathlineto{\pgfpoint{-10.043692\du}{19.919962\du}}
\pgfpathlineto{\pgfpoint{-10.049663\du}{19.917275\du}}
\pgfpathlineto{\pgfpoint{-10.056529\du}{19.914887\du}}
\pgfpathlineto{\pgfpoint{-10.063396\du}{19.912498\du}}
\pgfpathlineto{\pgfpoint{-10.070860\du}{19.910408\du}}
\pgfpathlineto{\pgfpoint{-10.078324\du}{19.908020\du}}
\pgfpathlineto{\pgfpoint{-10.086385\du}{19.906527\du}}
\pgfpathlineto{\pgfpoint{-10.095043\du}{19.904437\du}}
\pgfpathlineto{\pgfpoint{-10.103402\du}{19.902646\du}}
\pgfpathlineto{\pgfpoint{-10.112060\du}{19.901452\du}}
\pgfpathlineto{\pgfpoint{-10.121315\du}{19.899959\du}}
\pgfpathlineto{\pgfpoint{-10.130570\du}{19.898765\du}}
\pgfpathlineto{\pgfpoint{-10.139825\du}{19.897571\du}}
\pgfpathlineto{\pgfpoint{-10.149379\du}{19.896675\du}}
\pgfpathlineto{\pgfpoint{-10.159530\du}{19.895779\du}}
\pgfpathlineto{\pgfpoint{-10.169084\du}{19.895182\du}}
\pgfpathlineto{\pgfpoint{-10.178637\du}{19.894884\du}}
\pgfpathlineto{\pgfpoint{-10.188788\du}{19.894585\du}}
\pgfpathlineto{\pgfpoint{-10.198640\du}{19.894585\du}}
\pgfpathlineto{\pgfpoint{-10.198640\du}{19.894585\du}}
\pgfpathlineto{\pgfpoint{-10.208492\du}{19.894585\du}}
\pgfpathlineto{\pgfpoint{-10.218643\du}{19.894884\du}}
\pgfpathlineto{\pgfpoint{-10.228495\du}{19.895182\du}}
\pgfpathlineto{\pgfpoint{-10.238348\du}{19.895779\du}}
\pgfpathlineto{\pgfpoint{-10.247901\du}{19.896675\du}}
\pgfpathlineto{\pgfpoint{-10.257455\du}{19.897571\du}}
\pgfpathlineto{\pgfpoint{-10.267009\du}{19.898765\du}}
\pgfpathlineto{\pgfpoint{-10.275965\du}{19.899959\du}}
\pgfpathlineto{\pgfpoint{-10.285220\du}{19.901452\du}}
\pgfpathlineto{\pgfpoint{-10.293878\du}{19.902646\du}}
\pgfpathlineto{\pgfpoint{-10.302536\du}{19.904437\du}}
\pgfpathlineto{\pgfpoint{-10.310896\du}{19.906527\du}}
\pgfpathlineto{\pgfpoint{-10.318957\du}{19.908020\du}}
\pgfpathlineto{\pgfpoint{-10.326421\du}{19.910408\du}}
\pgfpathlineto{\pgfpoint{-10.333884\du}{19.912498\du}}
\pgfpathlineto{\pgfpoint{-10.341050\du}{19.914887\du}}
\pgfpathlineto{\pgfpoint{-10.347618\du}{19.917275\du}}
\pgfpathlineto{\pgfpoint{-10.353887\du}{19.919962\du}}
\pgfpathlineto{\pgfpoint{-10.359858\du}{19.922350\du}}
\pgfpathlineto{\pgfpoint{-10.365232\du}{19.925336\du}}
\pgfpathlineto{\pgfpoint{-10.370308\du}{19.928023\du}}
\pgfpathlineto{\pgfpoint{-10.375084\du}{19.931008\du}}
\pgfpathlineto{\pgfpoint{-10.378966\du}{19.933994\du}}
\pgfpathlineto{\pgfpoint{-10.382847\du}{19.936979\du}}
\pgfpathlineto{\pgfpoint{-10.385832\du}{19.940264\du}}
\pgfpathlineto{\pgfpoint{-10.388519\du}{19.943548\du}}
\pgfpathlineto{\pgfpoint{-10.390908\du}{19.946533\du}}
\pgfpathlineto{\pgfpoint{-10.392699\du}{19.949817\du}}
\pgfpathlineto{\pgfpoint{-10.394192\du}{19.953101\du}}
\pgfpathlineto{\pgfpoint{-10.395087\du}{19.956385\du}}
\pgfpathlineto{\pgfpoint{-10.395087\du}{19.959669\du}}
\pgfpathlineto{\pgfpoint{-10.395087\du}{19.959669\du}}
\pgfpathlineto{\pgfpoint{-10.395087\du}{19.963252\du}}
\pgfpathlineto{\pgfpoint{-10.394192\du}{19.966238\du}}
\pgfpathlineto{\pgfpoint{-10.392699\du}{19.969522\du}}
\pgfpathlineto{\pgfpoint{-10.390908\du}{19.972806\du}}
\pgfpathlineto{\pgfpoint{-10.388519\du}{19.976090\du}}
\pgfpathlineto{\pgfpoint{-10.385832\du}{19.979374\du}}
\pgfpathlineto{\pgfpoint{-10.382847\du}{19.982359\du}}
\pgfpathlineto{\pgfpoint{-10.378966\du}{19.985345\du}}
\pgfpathlineto{\pgfpoint{-10.375084\du}{19.988330\du}}
\pgfpathlineto{\pgfpoint{-10.370308\du}{19.991614\du}}
\pgfpathlineto{\pgfpoint{-10.365232\du}{19.994301\du}}
\pgfpathlineto{\pgfpoint{-10.359858\du}{19.996988\du}}
\pgfpathlineto{\pgfpoint{-10.353887\du}{19.999675\du}}
\pgfpathlineto{\pgfpoint{-10.347618\du}{20.002362\du}}
\pgfpathlineto{\pgfpoint{-10.341050\du}{20.004751\du}}
\pgfpathlineto{\pgfpoint{-10.333884\du}{20.007139\du}}
\pgfpathlineto{\pgfpoint{-10.326421\du}{20.008930\du}}
\pgfpathlineto{\pgfpoint{-10.318957\du}{20.011020\du}}
\pgfpathlineto{\pgfpoint{-10.310896\du}{20.013110\du}}
\pgfpathlineto{\pgfpoint{-10.302536\du}{20.015200\du}}
\pgfpathlineto{\pgfpoint{-10.293878\du}{20.016693\du}}
\pgfpathlineto{\pgfpoint{-10.285220\du}{20.018186\du}}
\pgfpathlineto{\pgfpoint{-10.275965\du}{20.019678\du}}
\pgfpathlineto{\pgfpoint{-10.267009\du}{20.020873\du}}
\pgfpathlineto{\pgfpoint{-10.257455\du}{20.021768\du}}
\pgfpathlineto{\pgfpoint{-10.247901\du}{20.022962\du}}
\pgfpathlineto{\pgfpoint{-10.238348\du}{20.023560\du}}
\pgfpathlineto{\pgfpoint{-10.228495\du}{20.024157\du}}
\pgfpathlineto{\pgfpoint{-10.218643\du}{20.024754\du}}
\pgfpathlineto{\pgfpoint{-10.208492\du}{20.025052\du}}
\pgfpathlineto{\pgfpoint{-10.198640\du}{20.025052\du}}
\pgfpathlineto{\pgfpoint{-10.198640\du}{20.025052\du}}
\pgfpathlineto{\pgfpoint{-10.188788\du}{20.025052\du}}
\pgfpathlineto{\pgfpoint{-10.178637\du}{20.024754\du}}
\pgfpathlineto{\pgfpoint{-10.169084\du}{20.024157\du}}
\pgfpathlineto{\pgfpoint{-10.159530\du}{20.023560\du}}
\pgfpathlineto{\pgfpoint{-10.149379\du}{20.022962\du}}
\pgfpathlineto{\pgfpoint{-10.139825\du}{20.021768\du}}
\pgfpathlineto{\pgfpoint{-10.130570\du}{20.020873\du}}
\pgfpathlineto{\pgfpoint{-10.121315\du}{20.019678\du}}
\pgfpathlineto{\pgfpoint{-10.112060\du}{20.018186\du}}
\pgfpathlineto{\pgfpoint{-10.103402\du}{20.016693\du}}
\pgfpathlineto{\pgfpoint{-10.095043\du}{20.015200\du}}
\pgfpathlineto{\pgfpoint{-10.086385\du}{20.013110\du}}
\pgfpathlineto{\pgfpoint{-10.078324\du}{20.011020\du}}
\pgfpathlineto{\pgfpoint{-10.070860\du}{20.008930\du}}
\pgfpathlineto{\pgfpoint{-10.063396\du}{20.007139\du}}
\pgfpathlineto{\pgfpoint{-10.056529\du}{20.004751\du}}
\pgfpathlineto{\pgfpoint{-10.049663\du}{20.002362\du}}
\pgfpathlineto{\pgfpoint{-10.043692\du}{19.999675\du}}
\pgfpathlineto{\pgfpoint{-10.037422\du}{19.996988\du}}
\pgfpathlineto{\pgfpoint{-10.032048\du}{19.994301\du}}
\pgfpathlineto{\pgfpoint{-10.026973\du}{19.991614\du}}
\pgfpathlineto{\pgfpoint{-10.022196\du}{19.988330\du}}
\pgfpathlineto{\pgfpoint{-10.018613\du}{19.985345\du}}
\pgfpathlineto{\pgfpoint{-10.014434\du}{19.982359\du}}
\pgfpathlineto{\pgfpoint{-10.011448\du}{19.979374\du}}
\pgfpathlineto{\pgfpoint{-10.008761\du}{19.976090\du}}
\pgfpathlineto{\pgfpoint{-10.006373\du}{19.972806\du}}
\pgfpathlineto{\pgfpoint{-10.004880\du}{19.969522\du}}
\pgfpathlineto{\pgfpoint{-10.003686\du}{19.966238\du}}
\pgfpathlineto{\pgfpoint{-10.002790\du}{19.963252\du}}
\pgfpathlineto{\pgfpoint{-10.002790\du}{19.959669\du}}
\pgfusepath{fill}
\pgfsetbuttcap
\pgfsetmiterjoin
\pgfsetdash{}{0pt}
\definecolor{dialinecolor}{rgb}{0.666667, 0.901961, 1.000000}
\pgfsetstrokecolor{dialinecolor}
\pgfpathmoveto{\pgfpoint{-10.014434\du}{19.953997\du}}
\pgfpathlineto{\pgfpoint{-10.014434\du}{19.950414\du}}
\pgfpathlineto{\pgfpoint{-10.015329\du}{19.948026\du}}
\pgfpathlineto{\pgfpoint{-10.016524\du}{19.944742\du}}
\pgfpathlineto{\pgfpoint{-10.018613\du}{19.941756\du}}
\pgfpathlineto{\pgfpoint{-10.020703\du}{19.939069\du}}
\pgfpathlineto{\pgfpoint{-10.022793\du}{19.935487\du}}
\pgfpathlineto{\pgfpoint{-10.026077\du}{19.933098\du}}
\pgfpathlineto{\pgfpoint{-10.029361\du}{19.930411\du}}
\pgfpathlineto{\pgfpoint{-10.033541\du}{19.927724\du}}
\pgfpathlineto{\pgfpoint{-10.038318\du}{19.925336\du}}
\pgfpathlineto{\pgfpoint{-10.042796\du}{19.922350\du}}
\pgfpathlineto{\pgfpoint{-10.048469\du}{19.920261\du}}
\pgfpathlineto{\pgfpoint{-10.054440\du}{19.917574\du}}
\pgfpathlineto{\pgfpoint{-10.060112\du}{19.914887\du}}
\pgfpathlineto{\pgfpoint{-10.066680\du}{19.913095\du}}
\pgfpathlineto{\pgfpoint{-10.073547\du}{19.910707\du}}
\pgfpathlineto{\pgfpoint{-10.080414\du}{19.909214\du}}
\pgfpathlineto{\pgfpoint{-10.088176\du}{19.907124\du}}
\pgfpathlineto{\pgfpoint{-10.095640\du}{19.905333\du}}
\pgfpathlineto{\pgfpoint{-10.103999\du}{19.903542\du}}
\pgfpathlineto{\pgfpoint{-10.112657\du}{19.902347\du}}
\pgfpathlineto{\pgfpoint{-10.121017\du}{19.899959\du}}
\pgfpathlineto{\pgfpoint{-10.129675\du}{19.899660\du}}
\pgfpathlineto{\pgfpoint{-10.138631\du}{19.898168\du}}
\pgfpathlineto{\pgfpoint{-10.147289\du}{19.897272\du}}
\pgfpathlineto{\pgfpoint{-10.156843\du}{19.896376\du}}
\pgfpathlineto{\pgfpoint{-10.166098\du}{19.895779\du}}
\pgfpathlineto{\pgfpoint{-10.175652\du}{19.895481\du}}
\pgfpathlineto{\pgfpoint{-10.185205\du}{19.894884\du}}
\pgfpathlineto{\pgfpoint{-10.195356\du}{19.894585\du}}
\pgfpathlineto{\pgfpoint{-10.205208\du}{19.894585\du}}
\pgfpathlineto{\pgfpoint{-10.205208\du}{19.894585\du}}
\pgfpathlineto{\pgfpoint{-10.214165\du}{19.894585\du}}
\pgfpathlineto{\pgfpoint{-10.224017\du}{19.894884\du}}
\pgfpathlineto{\pgfpoint{-10.233571\du}{19.895481\du}}
\pgfpathlineto{\pgfpoint{-10.242826\du}{19.895779\du}}
\pgfpathlineto{\pgfpoint{-10.252380\du}{19.896376\du}}
\pgfpathlineto{\pgfpoint{-10.261933\du}{19.897272\du}}
\pgfpathlineto{\pgfpoint{-10.270890\du}{19.898168\du}}
\pgfpathlineto{\pgfpoint{-10.279846\du}{19.899660\du}}
\pgfpathlineto{\pgfpoint{-10.288504\du}{19.899959\du}}
\pgfpathlineto{\pgfpoint{-10.296864\du}{19.902347\du}}
\pgfpathlineto{\pgfpoint{-10.305522\du}{19.903542\du}}
\pgfpathlineto{\pgfpoint{-10.313583\du}{19.905333\du}}
\pgfpathlineto{\pgfpoint{-10.321345\du}{19.907124\du}}
\pgfpathlineto{\pgfpoint{-10.328809\du}{19.909214\du}}
\pgfpathlineto{\pgfpoint{-10.335377\du}{19.910707\du}}
\pgfpathlineto{\pgfpoint{-10.342841\du}{19.913095\du}}
\pgfpathlineto{\pgfpoint{-10.349110\du}{19.914887\du}}
\pgfpathlineto{\pgfpoint{-10.355082\du}{19.917574\du}}
\pgfpathlineto{\pgfpoint{-10.361053\du}{19.920261\du}}
\pgfpathlineto{\pgfpoint{-10.366128\du}{19.922350\du}}
\pgfpathlineto{\pgfpoint{-10.371203\du}{19.925336\du}}
\pgfpathlineto{\pgfpoint{-10.375980\du}{19.927724\du}}
\pgfpathlineto{\pgfpoint{-10.380160\du}{19.930411\du}}
\pgfpathlineto{\pgfpoint{-10.383444\du}{19.933098\du}}
\pgfpathlineto{\pgfpoint{-10.386728\du}{19.935487\du}}
\pgfpathlineto{\pgfpoint{-10.389116\du}{19.939069\du}}
\pgfpathlineto{\pgfpoint{-10.390908\du}{19.941756\du}}
\pgfpathlineto{\pgfpoint{-10.393296\du}{19.944742\du}}
\pgfpathlineto{\pgfpoint{-10.394192\du}{19.948026\du}}
\pgfpathlineto{\pgfpoint{-10.395087\du}{19.950414\du}}
\pgfpathlineto{\pgfpoint{-10.395087\du}{19.953997\du}}
\pgfpathlineto{\pgfpoint{-10.395087\du}{19.953997\du}}
\pgfpathlineto{\pgfpoint{-10.395087\du}{19.956982\du}}
\pgfpathlineto{\pgfpoint{-10.394192\du}{19.959669\du}}
\pgfpathlineto{\pgfpoint{-10.393296\du}{19.962655\du}}
\pgfpathlineto{\pgfpoint{-10.390908\du}{19.966238\du}}
\pgfpathlineto{\pgfpoint{-10.389116\du}{19.968924\du}}
\pgfpathlineto{\pgfpoint{-10.386728\du}{19.971910\du}}
\pgfpathlineto{\pgfpoint{-10.383444\du}{19.974597\du}}
\pgfpathlineto{\pgfpoint{-10.380160\du}{19.977284\du}}
\pgfpathlineto{\pgfpoint{-10.375980\du}{19.979971\du}}
\pgfpathlineto{\pgfpoint{-10.371203\du}{19.982658\du}}
\pgfpathlineto{\pgfpoint{-10.366128\du}{19.985345\du}}
\pgfpathlineto{\pgfpoint{-10.361053\du}{19.987733\du}}
\pgfpathlineto{\pgfpoint{-10.355082\du}{19.990420\du}}
\pgfpathlineto{\pgfpoint{-10.349110\du}{19.992510\du}}
\pgfpathlineto{\pgfpoint{-10.342841\du}{19.994899\du}}
\pgfpathlineto{\pgfpoint{-10.335377\du}{19.997287\du}}
\pgfpathlineto{\pgfpoint{-10.328809\du}{19.999078\du}}
\pgfpathlineto{\pgfpoint{-10.321345\du}{20.000571\du}}
\pgfpathlineto{\pgfpoint{-10.313583\du}{20.002661\du}}
\pgfpathlineto{\pgfpoint{-10.305522\du}{20.004452\du}}
\pgfpathlineto{\pgfpoint{-10.296864\du}{20.005646\du}}
\pgfpathlineto{\pgfpoint{-10.288504\du}{20.007438\du}}
\pgfpathlineto{\pgfpoint{-10.279846\du}{20.008333\du}}
\pgfpathlineto{\pgfpoint{-10.270890\du}{20.009826\du}}
\pgfpathlineto{\pgfpoint{-10.261933\du}{20.010722\du}}
\pgfpathlineto{\pgfpoint{-10.252380\du}{20.011020\du}}
\pgfpathlineto{\pgfpoint{-10.242826\du}{20.012215\du}}
\pgfpathlineto{\pgfpoint{-10.233571\du}{20.012513\du}}
\pgfpathlineto{\pgfpoint{-10.224017\du}{20.012812\du}}
\pgfpathlineto{\pgfpoint{-10.214165\du}{20.013409\du}}
\pgfpathlineto{\pgfpoint{-10.205208\du}{20.013409\du}}
\pgfpathlineto{\pgfpoint{-10.205208\du}{20.013409\du}}
\pgfpathlineto{\pgfpoint{-10.195356\du}{20.013409\du}}
\pgfpathlineto{\pgfpoint{-10.185205\du}{20.012812\du}}
\pgfpathlineto{\pgfpoint{-10.175652\du}{20.012513\du}}
\pgfpathlineto{\pgfpoint{-10.166098\du}{20.012215\du}}
\pgfpathlineto{\pgfpoint{-10.156843\du}{20.011020\du}}
\pgfpathlineto{\pgfpoint{-10.147289\du}{20.010722\du}}
\pgfpathlineto{\pgfpoint{-10.138631\du}{20.009826\du}}
\pgfpathlineto{\pgfpoint{-10.129675\du}{20.008333\du}}
\pgfpathlineto{\pgfpoint{-10.121017\du}{20.007438\du}}
\pgfpathlineto{\pgfpoint{-10.112657\du}{20.005646\du}}
\pgfpathlineto{\pgfpoint{-10.103999\du}{20.004452\du}}
\pgfpathlineto{\pgfpoint{-10.095640\du}{20.002661\du}}
\pgfpathlineto{\pgfpoint{-10.088176\du}{20.000571\du}}
\pgfpathlineto{\pgfpoint{-10.080414\du}{19.999078\du}}
\pgfpathlineto{\pgfpoint{-10.073547\du}{19.997287\du}}
\pgfpathlineto{\pgfpoint{-10.066680\du}{19.994899\du}}
\pgfpathlineto{\pgfpoint{-10.060112\du}{19.992510\du}}
\pgfpathlineto{\pgfpoint{-10.054440\du}{19.990420\du}}
\pgfpathlineto{\pgfpoint{-10.048469\du}{19.987733\du}}
\pgfpathlineto{\pgfpoint{-10.042796\du}{19.985345\du}}
\pgfpathlineto{\pgfpoint{-10.038318\du}{19.982658\du}}
\pgfpathlineto{\pgfpoint{-10.033541\du}{19.979971\du}}
\pgfpathlineto{\pgfpoint{-10.029361\du}{19.977284\du}}
\pgfpathlineto{\pgfpoint{-10.026077\du}{19.974597\du}}
\pgfpathlineto{\pgfpoint{-10.022793\du}{19.971910\du}}
\pgfpathlineto{\pgfpoint{-10.020703\du}{19.968924\du}}
\pgfpathlineto{\pgfpoint{-10.018613\du}{19.966238\du}}
\pgfpathlineto{\pgfpoint{-10.016524\du}{19.962655\du}}
\pgfpathlineto{\pgfpoint{-10.015329\du}{19.959669\du}}
\pgfpathlineto{\pgfpoint{-10.014434\du}{19.956982\du}}
\pgfpathlineto{\pgfpoint{-10.014434\du}{19.953997\du}}
\pgfusepath{stroke}
\pgfsetbuttcap
\pgfsetmiterjoin
\pgfsetdash{}{0pt}
\definecolor{dialinecolor}{rgb}{0.000000, 0.470588, 0.666667}
\pgfsetfillcolor{dialinecolor}
\pgfpathmoveto{\pgfpoint{-10.395087\du}{19.870701\du}}
\pgfpathlineto{\pgfpoint{-10.395087\du}{19.965342\du}}
\pgfpathlineto{\pgfpoint{-10.002193\du}{19.965342\du}}
\pgfpathlineto{\pgfpoint{-10.002193\du}{19.870701\du}}
\pgfpathlineto{\pgfpoint{-10.395087\du}{19.870701\du}}
\pgfusepath{fill}
\pgfsetbuttcap
\pgfsetmiterjoin
\pgfsetdash{}{0pt}
\definecolor{dialinecolor}{rgb}{0.000000, 0.705882, 1.000000}
\pgfsetfillcolor{dialinecolor}
\pgfpathmoveto{\pgfpoint{-10.002790\du}{19.864133\du}}
\pgfpathlineto{\pgfpoint{-10.002790\du}{19.860849\du}}
\pgfpathlineto{\pgfpoint{-10.003686\du}{19.857266\du}}
\pgfpathlineto{\pgfpoint{-10.004880\du}{19.854281\du}}
\pgfpathlineto{\pgfpoint{-10.006373\du}{19.850996\du}}
\pgfpathlineto{\pgfpoint{-10.008761\du}{19.847712\du}}
\pgfpathlineto{\pgfpoint{-10.011448\du}{19.844428\du}}
\pgfpathlineto{\pgfpoint{-10.014434\du}{19.841443\du}}
\pgfpathlineto{\pgfpoint{-10.018613\du}{19.838159\du}}
\pgfpathlineto{\pgfpoint{-10.022196\du}{19.835472\du}}
\pgfpathlineto{\pgfpoint{-10.026973\du}{19.832486\du}}
\pgfpathlineto{\pgfpoint{-10.032048\du}{19.829501\du}}
\pgfpathlineto{\pgfpoint{-10.037422\du}{19.826814\du}}
\pgfpathlineto{\pgfpoint{-10.043692\du}{19.823530\du}}
\pgfpathlineto{\pgfpoint{-10.049663\du}{19.821738\du}}
\pgfpathlineto{\pgfpoint{-10.056529\du}{19.819051\du}}
\pgfpathlineto{\pgfpoint{-10.063396\du}{19.816663\du}}
\pgfpathlineto{\pgfpoint{-10.070860\du}{19.814573\du}}
\pgfpathlineto{\pgfpoint{-10.078324\du}{19.812483\du}}
\pgfpathlineto{\pgfpoint{-10.086385\du}{19.810393\du}}
\pgfpathlineto{\pgfpoint{-10.095043\du}{19.808304\du}}
\pgfpathlineto{\pgfpoint{-10.103402\du}{19.806811\du}}
\pgfpathlineto{\pgfpoint{-10.112060\du}{19.805318\du}}
\pgfpathlineto{\pgfpoint{-10.121315\du}{19.804124\du}}
\pgfpathlineto{\pgfpoint{-10.130570\du}{19.802631\du}}
\pgfpathlineto{\pgfpoint{-10.139825\du}{19.801735\du}}
\pgfpathlineto{\pgfpoint{-10.149379\du}{19.800541\du}}
\pgfpathlineto{\pgfpoint{-10.159530\du}{19.799944\du}}
\pgfpathlineto{\pgfpoint{-10.169084\du}{19.799347\du}}
\pgfpathlineto{\pgfpoint{-10.178637\du}{19.799048\du}}
\pgfpathlineto{\pgfpoint{-10.188788\du}{19.798750\du}}
\pgfpathlineto{\pgfpoint{-10.198640\du}{19.798750\du}}
\pgfpathlineto{\pgfpoint{-10.198640\du}{19.798750\du}}
\pgfpathlineto{\pgfpoint{-10.208492\du}{19.798750\du}}
\pgfpathlineto{\pgfpoint{-10.218643\du}{19.799048\du}}
\pgfpathlineto{\pgfpoint{-10.228495\du}{19.799347\du}}
\pgfpathlineto{\pgfpoint{-10.238348\du}{19.799944\du}}
\pgfpathlineto{\pgfpoint{-10.247901\du}{19.800541\du}}
\pgfpathlineto{\pgfpoint{-10.257455\du}{19.801735\du}}
\pgfpathlineto{\pgfpoint{-10.267009\du}{19.802631\du}}
\pgfpathlineto{\pgfpoint{-10.275965\du}{19.804124\du}}
\pgfpathlineto{\pgfpoint{-10.285220\du}{19.805318\du}}
\pgfpathlineto{\pgfpoint{-10.293878\du}{19.806811\du}}
\pgfpathlineto{\pgfpoint{-10.302536\du}{19.808304\du}}
\pgfpathlineto{\pgfpoint{-10.310896\du}{19.810393\du}}
\pgfpathlineto{\pgfpoint{-10.318957\du}{19.812483\du}}
\pgfpathlineto{\pgfpoint{-10.326421\du}{19.814573\du}}
\pgfpathlineto{\pgfpoint{-10.333884\du}{19.816663\du}}
\pgfpathlineto{\pgfpoint{-10.341050\du}{19.819051\du}}
\pgfpathlineto{\pgfpoint{-10.347618\du}{19.821738\du}}
\pgfpathlineto{\pgfpoint{-10.353887\du}{19.823530\du}}
\pgfpathlineto{\pgfpoint{-10.359858\du}{19.826814\du}}
\pgfpathlineto{\pgfpoint{-10.365232\du}{19.829501\du}}
\pgfpathlineto{\pgfpoint{-10.370308\du}{19.832486\du}}
\pgfpathlineto{\pgfpoint{-10.375084\du}{19.835472\du}}
\pgfpathlineto{\pgfpoint{-10.378966\du}{19.838159\du}}
\pgfpathlineto{\pgfpoint{-10.382847\du}{19.841443\du}}
\pgfpathlineto{\pgfpoint{-10.385832\du}{19.844428\du}}
\pgfpathlineto{\pgfpoint{-10.388519\du}{19.847712\du}}
\pgfpathlineto{\pgfpoint{-10.390908\du}{19.850996\du}}
\pgfpathlineto{\pgfpoint{-10.392699\du}{19.854281\du}}
\pgfpathlineto{\pgfpoint{-10.394192\du}{19.857266\du}}
\pgfpathlineto{\pgfpoint{-10.395087\du}{19.860849\du}}
\pgfpathlineto{\pgfpoint{-10.395087\du}{19.864133\du}}
\pgfpathlineto{\pgfpoint{-10.395087\du}{19.864133\du}}
\pgfpathlineto{\pgfpoint{-10.395087\du}{19.867417\du}}
\pgfpathlineto{\pgfpoint{-10.394192\du}{19.870999\du}}
\pgfpathlineto{\pgfpoint{-10.392699\du}{19.873985\du}}
\pgfpathlineto{\pgfpoint{-10.390908\du}{19.877269\du}}
\pgfpathlineto{\pgfpoint{-10.388519\du}{19.880852\du}}
\pgfpathlineto{\pgfpoint{-10.385832\du}{19.883837\du}}
\pgfpathlineto{\pgfpoint{-10.382847\du}{19.886823\du}}
\pgfpathlineto{\pgfpoint{-10.378966\du}{19.890107\du}}
\pgfpathlineto{\pgfpoint{-10.375084\du}{19.892495\du}}
\pgfpathlineto{\pgfpoint{-10.370308\du}{19.895779\du}}
\pgfpathlineto{\pgfpoint{-10.365232\du}{19.899063\du}}
\pgfpathlineto{\pgfpoint{-10.359858\du}{19.901750\du}}
\pgfpathlineto{\pgfpoint{-10.353887\du}{19.904437\du}}
\pgfpathlineto{\pgfpoint{-10.347618\du}{19.906527\du}}
\pgfpathlineto{\pgfpoint{-10.341050\du}{19.909513\du}}
\pgfpathlineto{\pgfpoint{-10.333884\du}{19.911901\du}}
\pgfpathlineto{\pgfpoint{-10.326421\du}{19.913692\du}}
\pgfpathlineto{\pgfpoint{-10.318957\du}{19.915782\du}}
\pgfpathlineto{\pgfpoint{-10.310896\du}{19.917872\du}}
\pgfpathlineto{\pgfpoint{-10.302536\du}{19.919962\du}}
\pgfpathlineto{\pgfpoint{-10.293878\du}{19.921455\du}}
\pgfpathlineto{\pgfpoint{-10.285220\du}{19.922947\du}}
\pgfpathlineto{\pgfpoint{-10.275965\du}{19.924142\du}}
\pgfpathlineto{\pgfpoint{-10.267009\du}{19.925634\du}}
\pgfpathlineto{\pgfpoint{-10.257455\du}{19.926530\du}}
\pgfpathlineto{\pgfpoint{-10.247901\du}{19.927724\du}}
\pgfpathlineto{\pgfpoint{-10.238348\du}{19.928023\du}}
\pgfpathlineto{\pgfpoint{-10.228495\du}{19.928919\du}}
\pgfpathlineto{\pgfpoint{-10.218643\du}{19.928919\du}}
\pgfpathlineto{\pgfpoint{-10.208492\du}{19.929814\du}}
\pgfpathlineto{\pgfpoint{-10.198640\du}{19.929814\du}}
\pgfpathlineto{\pgfpoint{-10.198640\du}{19.929814\du}}
\pgfpathlineto{\pgfpoint{-10.188788\du}{19.929814\du}}
\pgfpathlineto{\pgfpoint{-10.178637\du}{19.928919\du}}
\pgfpathlineto{\pgfpoint{-10.169084\du}{19.928919\du}}
\pgfpathlineto{\pgfpoint{-10.159530\du}{19.928023\du}}
\pgfpathlineto{\pgfpoint{-10.149379\du}{19.927724\du}}
\pgfpathlineto{\pgfpoint{-10.139825\du}{19.926530\du}}
\pgfpathlineto{\pgfpoint{-10.130570\du}{19.925634\du}}
\pgfpathlineto{\pgfpoint{-10.121315\du}{19.924142\du}}
\pgfpathlineto{\pgfpoint{-10.112060\du}{19.922947\du}}
\pgfpathlineto{\pgfpoint{-10.103402\du}{19.921455\du}}
\pgfpathlineto{\pgfpoint{-10.095043\du}{19.919962\du}}
\pgfpathlineto{\pgfpoint{-10.086385\du}{19.917872\du}}
\pgfpathlineto{\pgfpoint{-10.078324\du}{19.915782\du}}
\pgfpathlineto{\pgfpoint{-10.070860\du}{19.913692\du}}
\pgfpathlineto{\pgfpoint{-10.063396\du}{19.911901\du}}
\pgfpathlineto{\pgfpoint{-10.056529\du}{19.909513\du}}
\pgfpathlineto{\pgfpoint{-10.049663\du}{19.906527\du}}
\pgfpathlineto{\pgfpoint{-10.043692\du}{19.904437\du}}
\pgfpathlineto{\pgfpoint{-10.037422\du}{19.901750\du}}
\pgfpathlineto{\pgfpoint{-10.032048\du}{19.899063\du}}
\pgfpathlineto{\pgfpoint{-10.026973\du}{19.895779\du}}
\pgfpathlineto{\pgfpoint{-10.022196\du}{19.892495\du}}
\pgfpathlineto{\pgfpoint{-10.018613\du}{19.890107\du}}
\pgfpathlineto{\pgfpoint{-10.014434\du}{19.886823\du}}
\pgfpathlineto{\pgfpoint{-10.011448\du}{19.883837\du}}
\pgfpathlineto{\pgfpoint{-10.008761\du}{19.880852\du}}
\pgfpathlineto{\pgfpoint{-10.006373\du}{19.877269\du}}
\pgfpathlineto{\pgfpoint{-10.004880\du}{19.873985\du}}
\pgfpathlineto{\pgfpoint{-10.003686\du}{19.870999\du}}
\pgfpathlineto{\pgfpoint{-10.002790\du}{19.867417\du}}
\pgfpathlineto{\pgfpoint{-10.002790\du}{19.864133\du}}
\pgfusepath{fill}
\pgfsetbuttcap
\pgfsetmiterjoin
\pgfsetdash{}{0pt}
\definecolor{dialinecolor}{rgb}{0.666667, 0.901961, 1.000000}
\pgfsetstrokecolor{dialinecolor}
\pgfpathmoveto{\pgfpoint{-10.014434\du}{19.858460\du}}
\pgfpathlineto{\pgfpoint{-10.014434\du}{19.855176\du}}
\pgfpathlineto{\pgfpoint{-10.015329\du}{19.852788\du}}
\pgfpathlineto{\pgfpoint{-10.016524\du}{19.849504\du}}
\pgfpathlineto{\pgfpoint{-10.018613\du}{19.845921\du}}
\pgfpathlineto{\pgfpoint{-10.020703\du}{19.843533\du}}
\pgfpathlineto{\pgfpoint{-10.022793\du}{19.840249\du}}
\pgfpathlineto{\pgfpoint{-10.026077\du}{19.837860\du}}
\pgfpathlineto{\pgfpoint{-10.029361\du}{19.835173\du}}
\pgfpathlineto{\pgfpoint{-10.033541\du}{19.832486\du}}
\pgfpathlineto{\pgfpoint{-10.038318\du}{19.829799\du}}
\pgfpathlineto{\pgfpoint{-10.042796\du}{19.827112\du}}
\pgfpathlineto{\pgfpoint{-10.048469\du}{19.824425\du}}
\pgfpathlineto{\pgfpoint{-10.054440\du}{19.822037\du}}
\pgfpathlineto{\pgfpoint{-10.060112\du}{19.819947\du}}
\pgfpathlineto{\pgfpoint{-10.066680\du}{19.817559\du}}
\pgfpathlineto{\pgfpoint{-10.073547\du}{19.815170\du}}
\pgfpathlineto{\pgfpoint{-10.080414\du}{19.813379\du}}
\pgfpathlineto{\pgfpoint{-10.088176\du}{19.811886\du}}
\pgfpathlineto{\pgfpoint{-10.095640\du}{19.809498\du}}
\pgfpathlineto{\pgfpoint{-10.103999\du}{19.808005\du}}
\pgfpathlineto{\pgfpoint{-10.112657\du}{19.806811\du}}
\pgfpathlineto{\pgfpoint{-10.121017\du}{19.805019\du}}
\pgfpathlineto{\pgfpoint{-10.129675\du}{19.804124\du}}
\pgfpathlineto{\pgfpoint{-10.138631\du}{19.802631\du}}
\pgfpathlineto{\pgfpoint{-10.147289\du}{19.801735\du}}
\pgfpathlineto{\pgfpoint{-10.156843\du}{19.801138\du}}
\pgfpathlineto{\pgfpoint{-10.166098\du}{19.800243\du}}
\pgfpathlineto{\pgfpoint{-10.175652\du}{19.799944\du}}
\pgfpathlineto{\pgfpoint{-10.185205\du}{19.799645\du}}
\pgfpathlineto{\pgfpoint{-10.195356\du}{19.799048\du}}
\pgfpathlineto{\pgfpoint{-10.205208\du}{19.799048\du}}
\pgfpathlineto{\pgfpoint{-10.205208\du}{19.799048\du}}
\pgfpathlineto{\pgfpoint{-10.214165\du}{19.799048\du}}
\pgfpathlineto{\pgfpoint{-10.224017\du}{19.799645\du}}
\pgfpathlineto{\pgfpoint{-10.233571\du}{19.799944\du}}
\pgfpathlineto{\pgfpoint{-10.242826\du}{19.800243\du}}
\pgfpathlineto{\pgfpoint{-10.252380\du}{19.801138\du}}
\pgfpathlineto{\pgfpoint{-10.261933\du}{19.801735\du}}
\pgfpathlineto{\pgfpoint{-10.270890\du}{19.802631\du}}
\pgfpathlineto{\pgfpoint{-10.279846\du}{19.804124\du}}
\pgfpathlineto{\pgfpoint{-10.288504\du}{19.805019\du}}
\pgfpathlineto{\pgfpoint{-10.296864\du}{19.806811\du}}
\pgfpathlineto{\pgfpoint{-10.305522\du}{19.808005\du}}
\pgfpathlineto{\pgfpoint{-10.313583\du}{19.809498\du}}
\pgfpathlineto{\pgfpoint{-10.321345\du}{19.811886\du}}
\pgfpathlineto{\pgfpoint{-10.328809\du}{19.813379\du}}
\pgfpathlineto{\pgfpoint{-10.335377\du}{19.815170\du}}
\pgfpathlineto{\pgfpoint{-10.342841\du}{19.817559\du}}
\pgfpathlineto{\pgfpoint{-10.349110\du}{19.819947\du}}
\pgfpathlineto{\pgfpoint{-10.355082\du}{19.822037\du}}
\pgfpathlineto{\pgfpoint{-10.361053\du}{19.824425\du}}
\pgfpathlineto{\pgfpoint{-10.366128\du}{19.827112\du}}
\pgfpathlineto{\pgfpoint{-10.371203\du}{19.829799\du}}
\pgfpathlineto{\pgfpoint{-10.375980\du}{19.832486\du}}
\pgfpathlineto{\pgfpoint{-10.380160\du}{19.835173\du}}
\pgfpathlineto{\pgfpoint{-10.383444\du}{19.837860\du}}
\pgfpathlineto{\pgfpoint{-10.386728\du}{19.840249\du}}
\pgfpathlineto{\pgfpoint{-10.389116\du}{19.843533\du}}
\pgfpathlineto{\pgfpoint{-10.390908\du}{19.845921\du}}
\pgfpathlineto{\pgfpoint{-10.393296\du}{19.849504\du}}
\pgfpathlineto{\pgfpoint{-10.394192\du}{19.852788\du}}
\pgfpathlineto{\pgfpoint{-10.395087\du}{19.855176\du}}
\pgfpathlineto{\pgfpoint{-10.395087\du}{19.858460\du}}
\pgfpathlineto{\pgfpoint{-10.395087\du}{19.858460\du}}
\pgfpathlineto{\pgfpoint{-10.395087\du}{19.861744\du}}
\pgfpathlineto{\pgfpoint{-10.394192\du}{19.864431\du}}
\pgfpathlineto{\pgfpoint{-10.393296\du}{19.867417\du}}
\pgfpathlineto{\pgfpoint{-10.390908\du}{19.870701\du}}
\pgfpathlineto{\pgfpoint{-10.389116\du}{19.873388\du}}
\pgfpathlineto{\pgfpoint{-10.386728\du}{19.876672\du}}
\pgfpathlineto{\pgfpoint{-10.383444\du}{19.879359\du}}
\pgfpathlineto{\pgfpoint{-10.380160\du}{19.882046\du}}
\pgfpathlineto{\pgfpoint{-10.375980\du}{19.884733\du}}
\pgfpathlineto{\pgfpoint{-10.371203\du}{19.887420\du}}
\pgfpathlineto{\pgfpoint{-10.366128\du}{19.890107\du}}
\pgfpathlineto{\pgfpoint{-10.361053\du}{19.892197\du}}
\pgfpathlineto{\pgfpoint{-10.355082\du}{19.894884\du}}
\pgfpathlineto{\pgfpoint{-10.349110\du}{19.897272\du}}
\pgfpathlineto{\pgfpoint{-10.342841\du}{19.899660\du}}
\pgfpathlineto{\pgfpoint{-10.335377\du}{19.901750\du}}
\pgfpathlineto{\pgfpoint{-10.328809\du}{19.903542\du}}
\pgfpathlineto{\pgfpoint{-10.321345\du}{19.905333\du}}
\pgfpathlineto{\pgfpoint{-10.313583\du}{19.907124\du}}
\pgfpathlineto{\pgfpoint{-10.305522\du}{19.909214\du}}
\pgfpathlineto{\pgfpoint{-10.296864\du}{19.910408\du}}
\pgfpathlineto{\pgfpoint{-10.288504\du}{19.912200\du}}
\pgfpathlineto{\pgfpoint{-10.279846\du}{19.913095\du}}
\pgfpathlineto{\pgfpoint{-10.270890\du}{19.914588\du}}
\pgfpathlineto{\pgfpoint{-10.261933\du}{19.914887\du}}
\pgfpathlineto{\pgfpoint{-10.252380\du}{19.916081\du}}
\pgfpathlineto{\pgfpoint{-10.242826\du}{19.916379\du}}
\pgfpathlineto{\pgfpoint{-10.233571\du}{19.917275\du}}
\pgfpathlineto{\pgfpoint{-10.224017\du}{19.917574\du}}
\pgfpathlineto{\pgfpoint{-10.214165\du}{19.917872\du}}
\pgfpathlineto{\pgfpoint{-10.205208\du}{19.917872\du}}
\pgfpathlineto{\pgfpoint{-10.205208\du}{19.917872\du}}
\pgfpathlineto{\pgfpoint{-10.195356\du}{19.917872\du}}
\pgfpathlineto{\pgfpoint{-10.185205\du}{19.917574\du}}
\pgfpathlineto{\pgfpoint{-10.175652\du}{19.917275\du}}
\pgfpathlineto{\pgfpoint{-10.166098\du}{19.916379\du}}
\pgfpathlineto{\pgfpoint{-10.156843\du}{19.916081\du}}
\pgfpathlineto{\pgfpoint{-10.147289\du}{19.914887\du}}
\pgfpathlineto{\pgfpoint{-10.138631\du}{19.914588\du}}
\pgfpathlineto{\pgfpoint{-10.129675\du}{19.913095\du}}
\pgfpathlineto{\pgfpoint{-10.121017\du}{19.912200\du}}
\pgfpathlineto{\pgfpoint{-10.112657\du}{19.910408\du}}
\pgfpathlineto{\pgfpoint{-10.103999\du}{19.909214\du}}
\pgfpathlineto{\pgfpoint{-10.095640\du}{19.907124\du}}
\pgfpathlineto{\pgfpoint{-10.088176\du}{19.905333\du}}
\pgfpathlineto{\pgfpoint{-10.080414\du}{19.903542\du}}
\pgfpathlineto{\pgfpoint{-10.073547\du}{19.901750\du}}
\pgfpathlineto{\pgfpoint{-10.066680\du}{19.899660\du}}
\pgfpathlineto{\pgfpoint{-10.060112\du}{19.897272\du}}
\pgfpathlineto{\pgfpoint{-10.054440\du}{19.894884\du}}
\pgfpathlineto{\pgfpoint{-10.048469\du}{19.892197\du}}
\pgfpathlineto{\pgfpoint{-10.042796\du}{19.890107\du}}
\pgfpathlineto{\pgfpoint{-10.038318\du}{19.887420\du}}
\pgfpathlineto{\pgfpoint{-10.033541\du}{19.884733\du}}
\pgfpathlineto{\pgfpoint{-10.029361\du}{19.882046\du}}
\pgfpathlineto{\pgfpoint{-10.026077\du}{19.879359\du}}
\pgfpathlineto{\pgfpoint{-10.022793\du}{19.876672\du}}
\pgfpathlineto{\pgfpoint{-10.020703\du}{19.873388\du}}
\pgfpathlineto{\pgfpoint{-10.018613\du}{19.870701\du}}
\pgfpathlineto{\pgfpoint{-10.016524\du}{19.867417\du}}
\pgfpathlineto{\pgfpoint{-10.015329\du}{19.864431\du}}
\pgfpathlineto{\pgfpoint{-10.014434\du}{19.861744\du}}
\pgfpathlineto{\pgfpoint{-10.014434\du}{19.858460\du}}
\pgfusepath{stroke}
\pgfsetbuttcap
\pgfsetmiterjoin
\pgfsetdash{}{0pt}
\definecolor{dialinecolor}{rgb}{0.000000, 0.000000, 0.000000}
\pgfsetfillcolor{dialinecolor}
\pgfpathmoveto{\pgfpoint{-10.192968\du}{19.846817\du}}
\pgfpathlineto{\pgfpoint{-10.169084\du}{19.858460\du}}
\pgfpathlineto{\pgfpoint{-10.097730\du}{19.834875\du}}
\pgfpathlineto{\pgfpoint{-10.061903\du}{19.846817\du}}
\pgfpathlineto{\pgfpoint{-10.085788\du}{19.822933\du}}
\pgfpathlineto{\pgfpoint{-10.169084\du}{19.822933\du}}
\pgfpathlineto{\pgfpoint{-10.133257\du}{19.822933\du}}
\pgfpathlineto{\pgfpoint{-10.192968\du}{19.846817\du}}
\pgfusepath{fill}
\pgfsetbuttcap
\pgfsetmiterjoin
\pgfsetdash{}{0pt}
\definecolor{dialinecolor}{rgb}{0.000000, 0.000000, 0.000000}
\pgfsetfillcolor{dialinecolor}
\pgfpathmoveto{\pgfpoint{-10.205208\du}{19.882344\du}}
\pgfpathlineto{\pgfpoint{-10.240139\du}{19.870701\du}}
\pgfpathlineto{\pgfpoint{-10.300148\du}{19.894585\du}}
\pgfpathlineto{\pgfpoint{-10.335377\du}{19.882344\du}}
\pgfpathlineto{\pgfpoint{-10.324032\du}{19.917872\du}}
\pgfpathlineto{\pgfpoint{-10.240139\du}{19.917872\du}}
\pgfpathlineto{\pgfpoint{-10.276264\du}{19.906527\du}}
\pgfpathlineto{\pgfpoint{-10.205208\du}{19.882344\du}}
\pgfusepath{fill}
\pgfsetbuttcap
\pgfsetmiterjoin
\pgfsetdash{}{0pt}
\definecolor{dialinecolor}{rgb}{0.000000, 0.000000, 0.000000}
\pgfsetfillcolor{dialinecolor}
\pgfpathmoveto{\pgfpoint{-10.324032\du}{19.822933\du}}
\pgfpathlineto{\pgfpoint{-10.300148\du}{19.810393\du}}
\pgfpathlineto{\pgfpoint{-10.228495\du}{19.846817\du}}
\pgfpathlineto{\pgfpoint{-10.205208\du}{19.834875\du}}
\pgfpathlineto{\pgfpoint{-10.216852\du}{19.858460\du}}
\pgfpathlineto{\pgfpoint{-10.300148\du}{19.858460\du}}
\pgfpathlineto{\pgfpoint{-10.264322\du}{19.846817\du}}
\pgfpathlineto{\pgfpoint{-10.324032\du}{19.822933\du}}
\pgfusepath{fill}
\pgfsetbuttcap
\pgfsetmiterjoin
\pgfsetdash{}{0pt}
\definecolor{dialinecolor}{rgb}{0.000000, 0.000000, 0.000000}
\pgfsetfillcolor{dialinecolor}
\pgfpathmoveto{\pgfpoint{-10.073547\du}{19.906527\du}}
\pgfpathlineto{\pgfpoint{-10.097730\du}{19.917872\du}}
\pgfpathlineto{\pgfpoint{-10.169084\du}{19.882344\du}}
\pgfpathlineto{\pgfpoint{-10.205208\du}{19.894585\du}}
\pgfpathlineto{\pgfpoint{-10.180727\du}{19.870701\du}}
\pgfpathlineto{\pgfpoint{-10.097730\du}{19.870701\du}}
\pgfpathlineto{\pgfpoint{-10.133257\du}{19.882344\du}}
\pgfpathlineto{\pgfpoint{-10.073547\du}{19.906527\du}}
\pgfusepath{fill}
\pgfsetbuttcap
\pgfsetmiterjoin
\pgfsetdash{}{0pt}
\definecolor{dialinecolor}{rgb}{1.000000, 1.000000, 1.000000}
\pgfsetfillcolor{dialinecolor}
\pgfpathmoveto{\pgfpoint{-10.192968\du}{19.858460\du}}
\pgfpathlineto{\pgfpoint{-10.169084\du}{19.858460\du}}
\pgfpathlineto{\pgfpoint{-10.097730\du}{19.834875\du}}
\pgfpathlineto{\pgfpoint{-10.061903\du}{19.846817\du}}
\pgfpathlineto{\pgfpoint{-10.085788\du}{19.822933\du}}
\pgfpathlineto{\pgfpoint{-10.156843\du}{19.822933\du}}
\pgfpathlineto{\pgfpoint{-10.133257\du}{19.822933\du}}
\pgfpathlineto{\pgfpoint{-10.192968\du}{19.858460\du}}
\pgfusepath{fill}
\pgfsetbuttcap
\pgfsetmiterjoin
\pgfsetdash{}{0pt}
\definecolor{dialinecolor}{rgb}{1.000000, 1.000000, 1.000000}
\pgfsetfillcolor{dialinecolor}
\pgfpathmoveto{\pgfpoint{-10.205208\du}{19.882344\du}}
\pgfpathlineto{\pgfpoint{-10.228495\du}{19.870701\du}}
\pgfpathlineto{\pgfpoint{-10.300148\du}{19.894585\du}}
\pgfpathlineto{\pgfpoint{-10.335377\du}{19.894585\du}}
\pgfpathlineto{\pgfpoint{-10.311791\du}{19.917872\du}}
\pgfpathlineto{\pgfpoint{-10.228495\du}{19.917872\du}}
\pgfpathlineto{\pgfpoint{-10.264322\du}{19.906527\du}}
\pgfpathlineto{\pgfpoint{-10.205208\du}{19.882344\du}}
\pgfusepath{fill}
\pgfsetbuttcap
\pgfsetmiterjoin
\pgfsetdash{}{0pt}
\definecolor{dialinecolor}{rgb}{1.000000, 1.000000, 1.000000}
\pgfsetfillcolor{dialinecolor}
\pgfpathmoveto{\pgfpoint{-10.324032\du}{19.822933\du}}
\pgfpathlineto{\pgfpoint{-10.300148\du}{19.822933\du}}
\pgfpathlineto{\pgfpoint{-10.228495\du}{19.846817\du}}
\pgfpathlineto{\pgfpoint{-10.192968\du}{19.834875\du}}
\pgfpathlineto{\pgfpoint{-10.216852\du}{19.858460\du}}
\pgfpathlineto{\pgfpoint{-10.300148\du}{19.858460\du}}
\pgfpathlineto{\pgfpoint{-10.264322\du}{19.858460\du}}
\pgfpathlineto{\pgfpoint{-10.324032\du}{19.822933\du}}
\pgfusepath{fill}
\pgfsetbuttcap
\pgfsetmiterjoin
\pgfsetdash{}{0pt}
\definecolor{dialinecolor}{rgb}{1.000000, 1.000000, 1.000000}
\pgfsetfillcolor{dialinecolor}
\pgfpathmoveto{\pgfpoint{-10.073547\du}{19.906527\du}}
\pgfpathlineto{\pgfpoint{-10.097730\du}{19.917872\du}}
\pgfpathlineto{\pgfpoint{-10.169084\du}{19.894585\du}}
\pgfpathlineto{\pgfpoint{-10.192968\du}{19.894585\du}}
\pgfpathlineto{\pgfpoint{-10.180727\du}{19.870701\du}}
\pgfpathlineto{\pgfpoint{-10.097730\du}{19.870701\du}}
\pgfpathlineto{\pgfpoint{-10.133257\du}{19.882344\du}}
\pgfpathlineto{\pgfpoint{-10.073547\du}{19.906527\du}}
\pgfusepath{fill}
\pgfsetbuttcap
\pgfsetmiterjoin
\pgfsetdash{}{0pt}
\definecolor{dialinecolor}{rgb}{0.666667, 0.901961, 1.000000}
\pgfsetstrokecolor{dialinecolor}
\pgfpathmoveto{\pgfpoint{-10.395087\du}{19.858460\du}}
\pgfpathlineto{\pgfpoint{-10.395087\du}{19.953400\du}}
\pgfusepath{stroke}
\pgfsetbuttcap
\pgfsetmiterjoin
\pgfsetdash{}{0pt}
\definecolor{dialinecolor}{rgb}{0.666667, 0.901961, 1.000000}
\pgfsetstrokecolor{dialinecolor}
\pgfpathmoveto{\pgfpoint{-10.014434\du}{19.858460\du}}
\pgfpathlineto{\pgfpoint{-10.014434\du}{19.953400\du}}
\pgfusepath{stroke}
\pgfsetbuttcap
\pgfsetmiterjoin
\pgfsetdash{}{0pt}
\definecolor{dialinecolor}{rgb}{0.000000, 0.470588, 0.666667}
\pgfsetfillcolor{dialinecolor}
\pgfpathmoveto{\pgfpoint{-10.359858\du}{20.364506\du}}
\pgfpathlineto{\pgfpoint{-10.360157\du}{20.360923\du}}
\pgfpathlineto{\pgfpoint{-10.361053\du}{20.357639\du}}
\pgfpathlineto{\pgfpoint{-10.362247\du}{20.354355\du}}
\pgfpathlineto{\pgfpoint{-10.364038\du}{20.351370\du}}
\pgfpathlineto{\pgfpoint{-10.366128\du}{20.348086\du}}
\pgfpathlineto{\pgfpoint{-10.369412\du}{20.344801\du}}
\pgfpathlineto{\pgfpoint{-10.372099\du}{20.341816\du}}
\pgfpathlineto{\pgfpoint{-10.375980\du}{20.338532\du}}
\pgfpathlineto{\pgfpoint{-10.380160\du}{20.335845\du}}
\pgfpathlineto{\pgfpoint{-10.384638\du}{20.332561\du}}
\pgfpathlineto{\pgfpoint{-10.390012\du}{20.329874\du}}
\pgfpathlineto{\pgfpoint{-10.395087\du}{20.327187\du}}
\pgfpathlineto{\pgfpoint{-10.401059\du}{20.324500\du}}
\pgfpathlineto{\pgfpoint{-10.407030\du}{20.321813\du}}
\pgfpathlineto{\pgfpoint{-10.414195\du}{20.319425\du}}
\pgfpathlineto{\pgfpoint{-10.421062\du}{20.317036\du}}
\pgfpathlineto{\pgfpoint{-10.428525\du}{20.315245\du}}
\pgfpathlineto{\pgfpoint{-10.435989\du}{20.312856\du}}
\pgfpathlineto{\pgfpoint{-10.444050\du}{20.311065\du}}
\pgfpathlineto{\pgfpoint{-10.452409\du}{20.308677\du}}
\pgfpathlineto{\pgfpoint{-10.460769\du}{20.307184\du}}
\pgfpathlineto{\pgfpoint{-10.469725\du}{20.305990\du}}
\pgfpathlineto{\pgfpoint{-10.478682\du}{20.304497\du}}
\pgfpathlineto{\pgfpoint{-10.487937\du}{20.303303\du}}
\pgfpathlineto{\pgfpoint{-10.497789\du}{20.302109\du}}
\pgfpathlineto{\pgfpoint{-10.507044\du}{20.301213\du}}
\pgfpathlineto{\pgfpoint{-10.516897\du}{20.300616\du}}
\pgfpathlineto{\pgfpoint{-10.526450\du}{20.300019\du}}
\pgfpathlineto{\pgfpoint{-10.536303\du}{20.299422\du}}
\pgfpathlineto{\pgfpoint{-10.546752\du}{20.299123\du}}
\pgfpathlineto{\pgfpoint{-10.556007\du}{20.299123\du}}
\pgfpathlineto{\pgfpoint{-10.556007\du}{20.299123\du}}
\pgfpathlineto{\pgfpoint{-10.565859\du}{20.299123\du}}
\pgfpathlineto{\pgfpoint{-10.576010\du}{20.299422\du}}
\pgfpathlineto{\pgfpoint{-10.585862\du}{20.300019\du}}
\pgfpathlineto{\pgfpoint{-10.595714\du}{20.300616\du}}
\pgfpathlineto{\pgfpoint{-10.605567\du}{20.301213\du}}
\pgfpathlineto{\pgfpoint{-10.614822\du}{20.302109\du}}
\pgfpathlineto{\pgfpoint{-10.624674\du}{20.303303\du}}
\pgfpathlineto{\pgfpoint{-10.633929\du}{20.304497\du}}
\pgfpathlineto{\pgfpoint{-10.642886\du}{20.305990\du}}
\pgfpathlineto{\pgfpoint{-10.651544\du}{20.307184\du}}
\pgfpathlineto{\pgfpoint{-10.659903\du}{20.308677\du}}
\pgfpathlineto{\pgfpoint{-10.668263\du}{20.311065\du}}
\pgfpathlineto{\pgfpoint{-10.676323\du}{20.312856\du}}
\pgfpathlineto{\pgfpoint{-10.683787\du}{20.315245\du}}
\pgfpathlineto{\pgfpoint{-10.691251\du}{20.317036\du}}
\pgfpathlineto{\pgfpoint{-10.698416\du}{20.319425\du}}
\pgfpathlineto{\pgfpoint{-10.705283\du}{20.321813\du}}
\pgfpathlineto{\pgfpoint{-10.711254\du}{20.324500\du}}
\pgfpathlineto{\pgfpoint{-10.717524\du}{20.327187\du}}
\pgfpathlineto{\pgfpoint{-10.722599\du}{20.329874\du}}
\pgfpathlineto{\pgfpoint{-10.727674\du}{20.332561\du}}
\pgfpathlineto{\pgfpoint{-10.732451\du}{20.335845\du}}
\pgfpathlineto{\pgfpoint{-10.736631\du}{20.338532\du}}
\pgfpathlineto{\pgfpoint{-10.740214\du}{20.341816\du}}
\pgfpathlineto{\pgfpoint{-10.743199\du}{20.344801\du}}
\pgfpathlineto{\pgfpoint{-10.745886\du}{20.348086\du}}
\pgfpathlineto{\pgfpoint{-10.748275\du}{20.351370\du}}
\pgfpathlineto{\pgfpoint{-10.750066\du}{20.354355\du}}
\pgfpathlineto{\pgfpoint{-10.751559\du}{20.357639\du}}
\pgfpathlineto{\pgfpoint{-10.752454\du}{20.360923\du}}
\pgfpathlineto{\pgfpoint{-10.752454\du}{20.364506\du}}
\pgfpathlineto{\pgfpoint{-10.752454\du}{20.364506\du}}
\pgfpathlineto{\pgfpoint{-10.752454\du}{20.367790\du}}
\pgfpathlineto{\pgfpoint{-10.751559\du}{20.371074\du}}
\pgfpathlineto{\pgfpoint{-10.750066\du}{20.374358\du}}
\pgfpathlineto{\pgfpoint{-10.748275\du}{20.377642\du}}
\pgfpathlineto{\pgfpoint{-10.745886\du}{20.380628\du}}
\pgfpathlineto{\pgfpoint{-10.743199\du}{20.383912\du}}
\pgfpathlineto{\pgfpoint{-10.740214\du}{20.386897\du}}
\pgfpathlineto{\pgfpoint{-10.736631\du}{20.390181\du}}
\pgfpathlineto{\pgfpoint{-10.732451\du}{20.393167\du}}
\pgfpathlineto{\pgfpoint{-10.727674\du}{20.396152\du}}
\pgfpathlineto{\pgfpoint{-10.722599\du}{20.398839\du}}
\pgfpathlineto{\pgfpoint{-10.717524\du}{20.401526\du}}
\pgfpathlineto{\pgfpoint{-10.711254\du}{20.404213\du}}
\pgfpathlineto{\pgfpoint{-10.705283\du}{20.406602\du}}
\pgfpathlineto{\pgfpoint{-10.698416\du}{20.409289\du}}
\pgfpathlineto{\pgfpoint{-10.691251\du}{20.411677\du}}
\pgfpathlineto{\pgfpoint{-10.683787\du}{20.413767\du}}
\pgfpathlineto{\pgfpoint{-10.676323\du}{20.416155\du}}
\pgfpathlineto{\pgfpoint{-10.668263\du}{20.417648\du}}
\pgfpathlineto{\pgfpoint{-10.659903\du}{20.419738\du}}
\pgfpathlineto{\pgfpoint{-10.651544\du}{20.421529\du}}
\pgfpathlineto{\pgfpoint{-10.642886\du}{20.422724\du}}
\pgfpathlineto{\pgfpoint{-10.633929\du}{20.424216\du}}
\pgfpathlineto{\pgfpoint{-10.624674\du}{20.425411\du}}
\pgfpathlineto{\pgfpoint{-10.614822\du}{20.426605\du}}
\pgfpathlineto{\pgfpoint{-10.605567\du}{20.427500\du}}
\pgfpathlineto{\pgfpoint{-10.595714\du}{20.428098\du}}
\pgfpathlineto{\pgfpoint{-10.585862\du}{20.428993\du}}
\pgfpathlineto{\pgfpoint{-10.576010\du}{20.428993\du}}
\pgfpathlineto{\pgfpoint{-10.565859\du}{20.429590\du}}
\pgfpathlineto{\pgfpoint{-10.556007\du}{20.429590\du}}
\pgfpathlineto{\pgfpoint{-10.556007\du}{20.429590\du}}
\pgfpathlineto{\pgfpoint{-10.546752\du}{20.429590\du}}
\pgfpathlineto{\pgfpoint{-10.536303\du}{20.428993\du}}
\pgfpathlineto{\pgfpoint{-10.526450\du}{20.428993\du}}
\pgfpathlineto{\pgfpoint{-10.516897\du}{20.428098\du}}
\pgfpathlineto{\pgfpoint{-10.507044\du}{20.427500\du}}
\pgfpathlineto{\pgfpoint{-10.497789\du}{20.426605\du}}
\pgfpathlineto{\pgfpoint{-10.487937\du}{20.425411\du}}
\pgfpathlineto{\pgfpoint{-10.478682\du}{20.424216\du}}
\pgfpathlineto{\pgfpoint{-10.469725\du}{20.422724\du}}
\pgfpathlineto{\pgfpoint{-10.460769\du}{20.421529\du}}
\pgfpathlineto{\pgfpoint{-10.452409\du}{20.419738\du}}
\pgfpathlineto{\pgfpoint{-10.444050\du}{20.417648\du}}
\pgfpathlineto{\pgfpoint{-10.435989\du}{20.416155\du}}
\pgfpathlineto{\pgfpoint{-10.428525\du}{20.413767\du}}
\pgfpathlineto{\pgfpoint{-10.421062\du}{20.411677\du}}
\pgfpathlineto{\pgfpoint{-10.414195\du}{20.409289\du}}
\pgfpathlineto{\pgfpoint{-10.407030\du}{20.406602\du}}
\pgfpathlineto{\pgfpoint{-10.401059\du}{20.404213\du}}
\pgfpathlineto{\pgfpoint{-10.395087\du}{20.401526\du}}
\pgfpathlineto{\pgfpoint{-10.390012\du}{20.398839\du}}
\pgfpathlineto{\pgfpoint{-10.384638\du}{20.396152\du}}
\pgfpathlineto{\pgfpoint{-10.380160\du}{20.393167\du}}
\pgfpathlineto{\pgfpoint{-10.375980\du}{20.390181\du}}
\pgfpathlineto{\pgfpoint{-10.372099\du}{20.386897\du}}
\pgfpathlineto{\pgfpoint{-10.369412\du}{20.383912\du}}
\pgfpathlineto{\pgfpoint{-10.366128\du}{20.380628\du}}
\pgfpathlineto{\pgfpoint{-10.364038\du}{20.377642\du}}
\pgfpathlineto{\pgfpoint{-10.362247\du}{20.374358\du}}
\pgfpathlineto{\pgfpoint{-10.361053\du}{20.371074\du}}
\pgfpathlineto{\pgfpoint{-10.360157\du}{20.367790\du}}
\pgfpathlineto{\pgfpoint{-10.359858\du}{20.364506\du}}
\pgfusepath{fill}
\pgfsetbuttcap
\pgfsetmiterjoin
\pgfsetdash{}{0pt}
\definecolor{dialinecolor}{rgb}{0.666667, 0.901961, 1.000000}
\pgfsetstrokecolor{dialinecolor}
\pgfpathmoveto{\pgfpoint{-10.371800\du}{20.358236\du}}
\pgfpathlineto{\pgfpoint{-10.372099\du}{20.355251\du}}
\pgfpathlineto{\pgfpoint{-10.372696\du}{20.352564\du}}
\pgfpathlineto{\pgfpoint{-10.373890\du}{20.349578\du}}
\pgfpathlineto{\pgfpoint{-10.375980\du}{20.345996\du}}
\pgfpathlineto{\pgfpoint{-10.378070\du}{20.343607\du}}
\pgfpathlineto{\pgfpoint{-10.380160\du}{20.340323\du}}
\pgfpathlineto{\pgfpoint{-10.383444\du}{20.337636\du}}
\pgfpathlineto{\pgfpoint{-10.386728\du}{20.334949\du}}
\pgfpathlineto{\pgfpoint{-10.390908\du}{20.331964\du}}
\pgfpathlineto{\pgfpoint{-10.395685\du}{20.329575\du}}
\pgfpathlineto{\pgfpoint{-10.400760\du}{20.326888\du}}
\pgfpathlineto{\pgfpoint{-10.405835\du}{20.324500\du}}
\pgfpathlineto{\pgfpoint{-10.411508\du}{20.322112\du}}
\pgfpathlineto{\pgfpoint{-10.417479\du}{20.319723\du}}
\pgfpathlineto{\pgfpoint{-10.424047\du}{20.317036\du}}
\pgfpathlineto{\pgfpoint{-10.430914\du}{20.315543\du}}
\pgfpathlineto{\pgfpoint{-10.437482\du}{20.313454\du}}
\pgfpathlineto{\pgfpoint{-10.445543\du}{20.311662\du}}
\pgfpathlineto{\pgfpoint{-10.453007\du}{20.310169\du}}
\pgfpathlineto{\pgfpoint{-10.460769\du}{20.308080\du}}
\pgfpathlineto{\pgfpoint{-10.469725\du}{20.306587\du}}
\pgfpathlineto{\pgfpoint{-10.478085\du}{20.305094\du}}
\pgfpathlineto{\pgfpoint{-10.486444\du}{20.303900\du}}
\pgfpathlineto{\pgfpoint{-10.495700\du}{20.302706\du}}
\pgfpathlineto{\pgfpoint{-10.504358\du}{20.301810\du}}
\pgfpathlineto{\pgfpoint{-10.513911\du}{20.300914\du}}
\pgfpathlineto{\pgfpoint{-10.523465\du}{20.300616\du}}
\pgfpathlineto{\pgfpoint{-10.532720\du}{20.300019\du}}
\pgfpathlineto{\pgfpoint{-10.542572\du}{20.299422\du}}
\pgfpathlineto{\pgfpoint{-10.552424\du}{20.299123\du}}
\pgfpathlineto{\pgfpoint{-10.561978\du}{20.299123\du}}
\pgfpathlineto{\pgfpoint{-10.561978\du}{20.299123\du}}
\pgfpathlineto{\pgfpoint{-10.571233\du}{20.299123\du}}
\pgfpathlineto{\pgfpoint{-10.581085\du}{20.299422\du}}
\pgfpathlineto{\pgfpoint{-10.590639\du}{20.300019\du}}
\pgfpathlineto{\pgfpoint{-10.600193\du}{20.300616\du}}
\pgfpathlineto{\pgfpoint{-10.609746\du}{20.300914\du}}
\pgfpathlineto{\pgfpoint{-10.619001\du}{20.301810\du}}
\pgfpathlineto{\pgfpoint{-10.627958\du}{20.302706\du}}
\pgfpathlineto{\pgfpoint{-10.637213\du}{20.303900\du}}
\pgfpathlineto{\pgfpoint{-10.645573\du}{20.305094\du}}
\pgfpathlineto{\pgfpoint{-10.654529\du}{20.306587\du}}
\pgfpathlineto{\pgfpoint{-10.662889\du}{20.308080\du}}
\pgfpathlineto{\pgfpoint{-10.670651\du}{20.310169\du}}
\pgfpathlineto{\pgfpoint{-10.678712\du}{20.311662\du}}
\pgfpathlineto{\pgfpoint{-10.686176\du}{20.313454\du}}
\pgfpathlineto{\pgfpoint{-10.692744\du}{20.315543\du}}
\pgfpathlineto{\pgfpoint{-10.700208\du}{20.317036\du}}
\pgfpathlineto{\pgfpoint{-10.706179\du}{20.319723\du}}
\pgfpathlineto{\pgfpoint{-10.712150\du}{20.322112\du}}
\pgfpathlineto{\pgfpoint{-10.718419\du}{20.324500\du}}
\pgfpathlineto{\pgfpoint{-10.723495\du}{20.326888\du}}
\pgfpathlineto{\pgfpoint{-10.728570\du}{20.329575\du}}
\pgfpathlineto{\pgfpoint{-10.732750\du}{20.331964\du}}
\pgfpathlineto{\pgfpoint{-10.736930\du}{20.334949\du}}
\pgfpathlineto{\pgfpoint{-10.740811\du}{20.337636\du}}
\pgfpathlineto{\pgfpoint{-10.744095\du}{20.340323\du}}
\pgfpathlineto{\pgfpoint{-10.745886\du}{20.343607\du}}
\pgfpathlineto{\pgfpoint{-10.748275\du}{20.345996\du}}
\pgfpathlineto{\pgfpoint{-10.750066\du}{20.349578\du}}
\pgfpathlineto{\pgfpoint{-10.751559\du}{20.352564\du}}
\pgfpathlineto{\pgfpoint{-10.751857\du}{20.355251\du}}
\pgfpathlineto{\pgfpoint{-10.752454\du}{20.358236\du}}
\pgfpathlineto{\pgfpoint{-10.752454\du}{20.358236\du}}
\pgfpathlineto{\pgfpoint{-10.751857\du}{20.361819\du}}
\pgfpathlineto{\pgfpoint{-10.751559\du}{20.364506\du}}
\pgfpathlineto{\pgfpoint{-10.750066\du}{20.367491\du}}
\pgfpathlineto{\pgfpoint{-10.748275\du}{20.370477\du}}
\pgfpathlineto{\pgfpoint{-10.745886\du}{20.373164\du}}
\pgfpathlineto{\pgfpoint{-10.744095\du}{20.376448\du}}
\pgfpathlineto{\pgfpoint{-10.740811\du}{20.379135\du}}
\pgfpathlineto{\pgfpoint{-10.736930\du}{20.382120\du}}
\pgfpathlineto{\pgfpoint{-10.732750\du}{20.384807\du}}
\pgfpathlineto{\pgfpoint{-10.728570\du}{20.387494\du}}
\pgfpathlineto{\pgfpoint{-10.723495\du}{20.390181\du}}
\pgfpathlineto{\pgfpoint{-10.718419\du}{20.391674\du}}
\pgfpathlineto{\pgfpoint{-10.712150\du}{20.394958\du}}
\pgfpathlineto{\pgfpoint{-10.706179\du}{20.397048\du}}
\pgfpathlineto{\pgfpoint{-10.700208\du}{20.399138\du}}
\pgfpathlineto{\pgfpoint{-10.692744\du}{20.401526\du}}
\pgfpathlineto{\pgfpoint{-10.686176\du}{20.403616\du}}
\pgfpathlineto{\pgfpoint{-10.678712\du}{20.405408\du}}
\pgfpathlineto{\pgfpoint{-10.670651\du}{20.406602\du}}
\pgfpathlineto{\pgfpoint{-10.662889\du}{20.408990\du}}
\pgfpathlineto{\pgfpoint{-10.654529\du}{20.410483\du}}
\pgfpathlineto{\pgfpoint{-10.645573\du}{20.411976\du}}
\pgfpathlineto{\pgfpoint{-10.637213\du}{20.413170\du}}
\pgfpathlineto{\pgfpoint{-10.627958\du}{20.414066\du}}
\pgfpathlineto{\pgfpoint{-10.619001\du}{20.414961\du}}
\pgfpathlineto{\pgfpoint{-10.609746\du}{20.416155\du}}
\pgfpathlineto{\pgfpoint{-10.600193\du}{20.416454\du}}
\pgfpathlineto{\pgfpoint{-10.590639\du}{20.417051\du}}
\pgfpathlineto{\pgfpoint{-10.581085\du}{20.417350\du}}
\pgfpathlineto{\pgfpoint{-10.571233\du}{20.417648\du}}
\pgfpathlineto{\pgfpoint{-10.561978\du}{20.417648\du}}
\pgfpathlineto{\pgfpoint{-10.561978\du}{20.417648\du}}
\pgfpathlineto{\pgfpoint{-10.552424\du}{20.417648\du}}
\pgfpathlineto{\pgfpoint{-10.542572\du}{20.417350\du}}
\pgfpathlineto{\pgfpoint{-10.532720\du}{20.417051\du}}
\pgfpathlineto{\pgfpoint{-10.523465\du}{20.416454\du}}
\pgfpathlineto{\pgfpoint{-10.513911\du}{20.416155\du}}
\pgfpathlineto{\pgfpoint{-10.504358\du}{20.414961\du}}
\pgfpathlineto{\pgfpoint{-10.495700\du}{20.414066\du}}
\pgfpathlineto{\pgfpoint{-10.486444\du}{20.413170\du}}
\pgfpathlineto{\pgfpoint{-10.478085\du}{20.411976\du}}
\pgfpathlineto{\pgfpoint{-10.469725\du}{20.410483\du}}
\pgfpathlineto{\pgfpoint{-10.460769\du}{20.408990\du}}
\pgfpathlineto{\pgfpoint{-10.453007\du}{20.406602\du}}
\pgfpathlineto{\pgfpoint{-10.445543\du}{20.405408\du}}
\pgfpathlineto{\pgfpoint{-10.437482\du}{20.403616\du}}
\pgfpathlineto{\pgfpoint{-10.430914\du}{20.401526\du}}
\pgfpathlineto{\pgfpoint{-10.424047\du}{20.399138\du}}
\pgfpathlineto{\pgfpoint{-10.417479\du}{20.397048\du}}
\pgfpathlineto{\pgfpoint{-10.411508\du}{20.394958\du}}
\pgfpathlineto{\pgfpoint{-10.405835\du}{20.391674\du}}
\pgfpathlineto{\pgfpoint{-10.400760\du}{20.390181\du}}
\pgfpathlineto{\pgfpoint{-10.395685\du}{20.387494\du}}
\pgfpathlineto{\pgfpoint{-10.390908\du}{20.384807\du}}
\pgfpathlineto{\pgfpoint{-10.386728\du}{20.382120\du}}
\pgfpathlineto{\pgfpoint{-10.383444\du}{20.379135\du}}
\pgfpathlineto{\pgfpoint{-10.380160\du}{20.376448\du}}
\pgfpathlineto{\pgfpoint{-10.378070\du}{20.373164\du}}
\pgfpathlineto{\pgfpoint{-10.375980\du}{20.370477\du}}
\pgfpathlineto{\pgfpoint{-10.373890\du}{20.367491\du}}
\pgfpathlineto{\pgfpoint{-10.372696\du}{20.364506\du}}
\pgfpathlineto{\pgfpoint{-10.372099\du}{20.361819\du}}
\pgfpathlineto{\pgfpoint{-10.371800\du}{20.358236\du}}
\pgfusepath{stroke}
\pgfsetbuttcap
\pgfsetmiterjoin
\pgfsetdash{}{0pt}
\definecolor{dialinecolor}{rgb}{0.000000, 0.470588, 0.666667}
\pgfsetfillcolor{dialinecolor}
\pgfpathmoveto{\pgfpoint{-10.752454\du}{20.275239\du}}
\pgfpathlineto{\pgfpoint{-10.752454\du}{20.370477\du}}
\pgfpathlineto{\pgfpoint{-10.359560\du}{20.370477\du}}
\pgfpathlineto{\pgfpoint{-10.359560\du}{20.275239\du}}
\pgfpathlineto{\pgfpoint{-10.752454\du}{20.275239\du}}
\pgfusepath{fill}
\pgfsetbuttcap
\pgfsetmiterjoin
\pgfsetdash{}{0pt}
\definecolor{dialinecolor}{rgb}{0.000000, 0.705882, 1.000000}
\pgfsetfillcolor{dialinecolor}
\pgfpathmoveto{\pgfpoint{-10.359858\du}{20.268969\du}}
\pgfpathlineto{\pgfpoint{-10.360157\du}{20.265387\du}}
\pgfpathlineto{\pgfpoint{-10.361053\du}{20.262103\du}}
\pgfpathlineto{\pgfpoint{-10.362247\du}{20.259117\du}}
\pgfpathlineto{\pgfpoint{-10.364038\du}{20.255534\du}}
\pgfpathlineto{\pgfpoint{-10.366128\du}{20.252250\du}}
\pgfpathlineto{\pgfpoint{-10.369412\du}{20.248966\du}}
\pgfpathlineto{\pgfpoint{-10.372099\du}{20.246279\du}}
\pgfpathlineto{\pgfpoint{-10.375980\du}{20.242398\du}}
\pgfpathlineto{\pgfpoint{-10.380160\du}{20.240010\du}}
\pgfpathlineto{\pgfpoint{-10.384638\du}{20.237024\du}}
\pgfpathlineto{\pgfpoint{-10.390012\du}{20.234039\du}}
\pgfpathlineto{\pgfpoint{-10.395087\du}{20.231352\du}}
\pgfpathlineto{\pgfpoint{-10.401059\du}{20.228665\du}}
\pgfpathlineto{\pgfpoint{-10.407030\du}{20.226276\du}}
\pgfpathlineto{\pgfpoint{-10.414195\du}{20.223589\du}}
\pgfpathlineto{\pgfpoint{-10.421062\du}{20.221201\du}}
\pgfpathlineto{\pgfpoint{-10.428525\du}{20.219111\du}}
\pgfpathlineto{\pgfpoint{-10.435989\du}{20.217021\du}}
\pgfpathlineto{\pgfpoint{-10.444050\du}{20.215230\du}}
\pgfpathlineto{\pgfpoint{-10.452409\du}{20.213140\du}}
\pgfpathlineto{\pgfpoint{-10.460769\du}{20.211349\du}}
\pgfpathlineto{\pgfpoint{-10.469725\du}{20.210155\du}}
\pgfpathlineto{\pgfpoint{-10.478682\du}{20.208662\du}}
\pgfpathlineto{\pgfpoint{-10.487937\du}{20.207468\du}}
\pgfpathlineto{\pgfpoint{-10.497789\du}{20.206273\du}}
\pgfpathlineto{\pgfpoint{-10.507044\du}{20.205079\du}}
\pgfpathlineto{\pgfpoint{-10.516897\du}{20.204781\du}}
\pgfpathlineto{\pgfpoint{-10.526450\du}{20.203885\du}}
\pgfpathlineto{\pgfpoint{-10.536303\du}{20.203586\du}}
\pgfpathlineto{\pgfpoint{-10.546752\du}{20.203288\du}}
\pgfpathlineto{\pgfpoint{-10.556007\du}{20.203288\du}}
\pgfpathlineto{\pgfpoint{-10.556007\du}{20.203288\du}}
\pgfpathlineto{\pgfpoint{-10.565859\du}{20.203288\du}}
\pgfpathlineto{\pgfpoint{-10.576010\du}{20.203586\du}}
\pgfpathlineto{\pgfpoint{-10.585862\du}{20.203885\du}}
\pgfpathlineto{\pgfpoint{-10.595714\du}{20.204781\du}}
\pgfpathlineto{\pgfpoint{-10.605567\du}{20.205079\du}}
\pgfpathlineto{\pgfpoint{-10.614822\du}{20.206273\du}}
\pgfpathlineto{\pgfpoint{-10.624674\du}{20.207468\du}}
\pgfpathlineto{\pgfpoint{-10.633929\du}{20.208662\du}}
\pgfpathlineto{\pgfpoint{-10.642886\du}{20.210155\du}}
\pgfpathlineto{\pgfpoint{-10.651544\du}{20.211349\du}}
\pgfpathlineto{\pgfpoint{-10.659903\du}{20.213140\du}}
\pgfpathlineto{\pgfpoint{-10.668263\du}{20.215230\du}}
\pgfpathlineto{\pgfpoint{-10.676323\du}{20.217021\du}}
\pgfpathlineto{\pgfpoint{-10.683787\du}{20.219111\du}}
\pgfpathlineto{\pgfpoint{-10.691251\du}{20.221201\du}}
\pgfpathlineto{\pgfpoint{-10.698416\du}{20.223589\du}}
\pgfpathlineto{\pgfpoint{-10.705283\du}{20.226276\du}}
\pgfpathlineto{\pgfpoint{-10.711254\du}{20.228665\du}}
\pgfpathlineto{\pgfpoint{-10.717524\du}{20.231352\du}}
\pgfpathlineto{\pgfpoint{-10.722599\du}{20.234039\du}}
\pgfpathlineto{\pgfpoint{-10.727674\du}{20.237024\du}}
\pgfpathlineto{\pgfpoint{-10.732451\du}{20.240010\du}}
\pgfpathlineto{\pgfpoint{-10.736631\du}{20.242398\du}}
\pgfpathlineto{\pgfpoint{-10.740214\du}{20.246279\du}}
\pgfpathlineto{\pgfpoint{-10.743199\du}{20.248966\du}}
\pgfpathlineto{\pgfpoint{-10.745886\du}{20.252250\du}}
\pgfpathlineto{\pgfpoint{-10.748275\du}{20.255534\du}}
\pgfpathlineto{\pgfpoint{-10.750066\du}{20.259117\du}}
\pgfpathlineto{\pgfpoint{-10.751559\du}{20.262103\du}}
\pgfpathlineto{\pgfpoint{-10.752454\du}{20.265387\du}}
\pgfpathlineto{\pgfpoint{-10.752454\du}{20.268969\du}}
\pgfpathlineto{\pgfpoint{-10.752454\du}{20.268969\du}}
\pgfpathlineto{\pgfpoint{-10.752454\du}{20.272253\du}}
\pgfpathlineto{\pgfpoint{-10.751559\du}{20.275537\du}}
\pgfpathlineto{\pgfpoint{-10.750066\du}{20.278523\du}}
\pgfpathlineto{\pgfpoint{-10.748275\du}{20.282106\du}}
\pgfpathlineto{\pgfpoint{-10.745886\du}{20.285390\du}}
\pgfpathlineto{\pgfpoint{-10.743199\du}{20.288375\du}}
\pgfpathlineto{\pgfpoint{-10.740214\du}{20.291361\du}}
\pgfpathlineto{\pgfpoint{-10.736631\du}{20.294645\du}}
\pgfpathlineto{\pgfpoint{-10.732451\du}{20.297630\du}}
\pgfpathlineto{\pgfpoint{-10.727674\du}{20.300616\du}}
\pgfpathlineto{\pgfpoint{-10.722599\du}{20.303601\du}}
\pgfpathlineto{\pgfpoint{-10.717524\du}{20.306288\du}}
\pgfpathlineto{\pgfpoint{-10.711254\du}{20.308677\du}}
\pgfpathlineto{\pgfpoint{-10.705283\du}{20.311364\du}}
\pgfpathlineto{\pgfpoint{-10.698416\du}{20.314051\du}}
\pgfpathlineto{\pgfpoint{-10.691251\du}{20.316140\du}}
\pgfpathlineto{\pgfpoint{-10.683787\du}{20.318529\du}}
\pgfpathlineto{\pgfpoint{-10.676323\du}{20.320619\du}}
\pgfpathlineto{\pgfpoint{-10.668263\du}{20.322410\du}}
\pgfpathlineto{\pgfpoint{-10.659903\du}{20.324500\du}}
\pgfpathlineto{\pgfpoint{-10.651544\du}{20.326291\du}}
\pgfpathlineto{\pgfpoint{-10.642886\du}{20.327485\du}}
\pgfpathlineto{\pgfpoint{-10.633929\du}{20.328978\du}}
\pgfpathlineto{\pgfpoint{-10.624674\du}{20.330172\du}}
\pgfpathlineto{\pgfpoint{-10.614822\du}{20.331068\du}}
\pgfpathlineto{\pgfpoint{-10.605567\du}{20.331964\du}}
\pgfpathlineto{\pgfpoint{-10.595714\du}{20.332859\du}}
\pgfpathlineto{\pgfpoint{-10.585862\du}{20.333755\du}}
\pgfpathlineto{\pgfpoint{-10.576010\du}{20.334054\du}}
\pgfpathlineto{\pgfpoint{-10.565859\du}{20.334352\du}}
\pgfpathlineto{\pgfpoint{-10.556007\du}{20.334352\du}}
\pgfpathlineto{\pgfpoint{-10.556007\du}{20.334352\du}}
\pgfpathlineto{\pgfpoint{-10.546752\du}{20.334352\du}}
\pgfpathlineto{\pgfpoint{-10.536303\du}{20.334054\du}}
\pgfpathlineto{\pgfpoint{-10.526450\du}{20.333755\du}}
\pgfpathlineto{\pgfpoint{-10.516897\du}{20.332859\du}}
\pgfpathlineto{\pgfpoint{-10.507044\du}{20.331964\du}}
\pgfpathlineto{\pgfpoint{-10.497789\du}{20.331068\du}}
\pgfpathlineto{\pgfpoint{-10.487937\du}{20.330172\du}}
\pgfpathlineto{\pgfpoint{-10.478682\du}{20.328978\du}}
\pgfpathlineto{\pgfpoint{-10.469725\du}{20.327485\du}}
\pgfpathlineto{\pgfpoint{-10.460769\du}{20.326291\du}}
\pgfpathlineto{\pgfpoint{-10.452409\du}{20.324500\du}}
\pgfpathlineto{\pgfpoint{-10.444050\du}{20.322410\du}}
\pgfpathlineto{\pgfpoint{-10.435989\du}{20.320619\du}}
\pgfpathlineto{\pgfpoint{-10.428525\du}{20.318529\du}}
\pgfpathlineto{\pgfpoint{-10.421062\du}{20.316140\du}}
\pgfpathlineto{\pgfpoint{-10.414195\du}{20.314051\du}}
\pgfpathlineto{\pgfpoint{-10.407030\du}{20.311364\du}}
\pgfpathlineto{\pgfpoint{-10.401059\du}{20.308677\du}}
\pgfpathlineto{\pgfpoint{-10.395087\du}{20.306288\du}}
\pgfpathlineto{\pgfpoint{-10.390012\du}{20.303601\du}}
\pgfpathlineto{\pgfpoint{-10.384638\du}{20.300616\du}}
\pgfpathlineto{\pgfpoint{-10.380160\du}{20.297630\du}}
\pgfpathlineto{\pgfpoint{-10.375980\du}{20.294645\du}}
\pgfpathlineto{\pgfpoint{-10.372099\du}{20.291361\du}}
\pgfpathlineto{\pgfpoint{-10.369412\du}{20.288375\du}}
\pgfpathlineto{\pgfpoint{-10.366128\du}{20.285390\du}}
\pgfpathlineto{\pgfpoint{-10.364038\du}{20.282106\du}}
\pgfpathlineto{\pgfpoint{-10.362247\du}{20.278523\du}}
\pgfpathlineto{\pgfpoint{-10.361053\du}{20.275537\du}}
\pgfpathlineto{\pgfpoint{-10.360157\du}{20.272253\du}}
\pgfpathlineto{\pgfpoint{-10.359858\du}{20.268969\du}}
\pgfusepath{fill}
\pgfsetbuttcap
\pgfsetmiterjoin
\pgfsetdash{}{0pt}
\definecolor{dialinecolor}{rgb}{0.666667, 0.901961, 1.000000}
\pgfsetstrokecolor{dialinecolor}
\pgfpathmoveto{\pgfpoint{-10.371800\du}{20.263297\du}}
\pgfpathlineto{\pgfpoint{-10.372099\du}{20.260311\du}}
\pgfpathlineto{\pgfpoint{-10.372696\du}{20.257326\du}}
\pgfpathlineto{\pgfpoint{-10.373890\du}{20.254639\du}}
\pgfpathlineto{\pgfpoint{-10.375980\du}{20.251355\du}}
\pgfpathlineto{\pgfpoint{-10.378070\du}{20.248668\du}}
\pgfpathlineto{\pgfpoint{-10.380160\du}{20.245384\du}}
\pgfpathlineto{\pgfpoint{-10.383444\du}{20.242398\du}}
\pgfpathlineto{\pgfpoint{-10.386728\du}{20.240010\du}}
\pgfpathlineto{\pgfpoint{-10.390908\du}{20.237323\du}}
\pgfpathlineto{\pgfpoint{-10.395685\du}{20.234636\du}}
\pgfpathlineto{\pgfpoint{-10.400760\du}{20.231949\du}}
\pgfpathlineto{\pgfpoint{-10.405835\du}{20.229560\du}}
\pgfpathlineto{\pgfpoint{-10.411508\du}{20.226575\du}}
\pgfpathlineto{\pgfpoint{-10.417479\du}{20.224784\du}}
\pgfpathlineto{\pgfpoint{-10.424047\du}{20.222694\du}}
\pgfpathlineto{\pgfpoint{-10.430914\du}{20.220007\du}}
\pgfpathlineto{\pgfpoint{-10.437482\du}{20.218514\du}}
\pgfpathlineto{\pgfpoint{-10.445543\du}{20.216723\du}}
\pgfpathlineto{\pgfpoint{-10.453007\du}{20.214633\du}}
\pgfpathlineto{\pgfpoint{-10.460769\du}{20.213140\du}}
\pgfpathlineto{\pgfpoint{-10.469725\du}{20.211647\du}}
\pgfpathlineto{\pgfpoint{-10.478085\du}{20.210155\du}}
\pgfpathlineto{\pgfpoint{-10.486444\du}{20.208960\du}}
\pgfpathlineto{\pgfpoint{-10.495700\du}{20.207766\du}}
\pgfpathlineto{\pgfpoint{-10.504358\du}{20.206572\du}}
\pgfpathlineto{\pgfpoint{-10.513911\du}{20.205975\du}}
\pgfpathlineto{\pgfpoint{-10.523465\du}{20.205079\du}}
\pgfpathlineto{\pgfpoint{-10.532720\du}{20.205079\du}}
\pgfpathlineto{\pgfpoint{-10.542572\du}{20.204781\du}}
\pgfpathlineto{\pgfpoint{-10.552424\du}{20.203885\du}}
\pgfpathlineto{\pgfpoint{-10.561978\du}{20.203885\du}}
\pgfpathlineto{\pgfpoint{-10.561978\du}{20.203885\du}}
\pgfpathlineto{\pgfpoint{-10.571233\du}{20.203885\du}}
\pgfpathlineto{\pgfpoint{-10.581085\du}{20.204781\du}}
\pgfpathlineto{\pgfpoint{-10.590639\du}{20.205079\du}}
\pgfpathlineto{\pgfpoint{-10.600193\du}{20.205079\du}}
\pgfpathlineto{\pgfpoint{-10.609746\du}{20.205975\du}}
\pgfpathlineto{\pgfpoint{-10.619001\du}{20.206572\du}}
\pgfpathlineto{\pgfpoint{-10.627958\du}{20.207766\du}}
\pgfpathlineto{\pgfpoint{-10.637213\du}{20.208960\du}}
\pgfpathlineto{\pgfpoint{-10.645573\du}{20.210155\du}}
\pgfpathlineto{\pgfpoint{-10.654529\du}{20.211647\du}}
\pgfpathlineto{\pgfpoint{-10.662889\du}{20.213140\du}}
\pgfpathlineto{\pgfpoint{-10.670651\du}{20.214633\du}}
\pgfpathlineto{\pgfpoint{-10.678712\du}{20.216723\du}}
\pgfpathlineto{\pgfpoint{-10.686176\du}{20.218514\du}}
\pgfpathlineto{\pgfpoint{-10.692744\du}{20.220007\du}}
\pgfpathlineto{\pgfpoint{-10.700208\du}{20.222694\du}}
\pgfpathlineto{\pgfpoint{-10.706179\du}{20.224784\du}}
\pgfpathlineto{\pgfpoint{-10.712150\du}{20.226575\du}}
\pgfpathlineto{\pgfpoint{-10.718419\du}{20.229560\du}}
\pgfpathlineto{\pgfpoint{-10.723495\du}{20.231949\du}}
\pgfpathlineto{\pgfpoint{-10.728570\du}{20.234636\du}}
\pgfpathlineto{\pgfpoint{-10.732750\du}{20.237323\du}}
\pgfpathlineto{\pgfpoint{-10.736930\du}{20.240010\du}}
\pgfpathlineto{\pgfpoint{-10.740811\du}{20.242398\du}}
\pgfpathlineto{\pgfpoint{-10.744095\du}{20.245384\du}}
\pgfpathlineto{\pgfpoint{-10.745886\du}{20.248668\du}}
\pgfpathlineto{\pgfpoint{-10.748275\du}{20.251355\du}}
\pgfpathlineto{\pgfpoint{-10.750066\du}{20.254639\du}}
\pgfpathlineto{\pgfpoint{-10.751559\du}{20.257326\du}}
\pgfpathlineto{\pgfpoint{-10.751857\du}{20.260311\du}}
\pgfpathlineto{\pgfpoint{-10.752454\du}{20.263297\du}}
\pgfpathlineto{\pgfpoint{-10.752454\du}{20.263297\du}}
\pgfpathlineto{\pgfpoint{-10.751857\du}{20.266879\du}}
\pgfpathlineto{\pgfpoint{-10.751559\du}{20.269566\du}}
\pgfpathlineto{\pgfpoint{-10.750066\du}{20.272253\du}}
\pgfpathlineto{\pgfpoint{-10.748275\du}{20.275537\du}}
\pgfpathlineto{\pgfpoint{-10.745886\du}{20.278224\du}}
\pgfpathlineto{\pgfpoint{-10.744095\du}{20.281807\du}}
\pgfpathlineto{\pgfpoint{-10.740811\du}{20.284494\du}}
\pgfpathlineto{\pgfpoint{-10.736930\du}{20.287181\du}}
\pgfpathlineto{\pgfpoint{-10.732750\du}{20.289868\du}}
\pgfpathlineto{\pgfpoint{-10.728570\du}{20.292555\du}}
\pgfpathlineto{\pgfpoint{-10.723495\du}{20.295242\du}}
\pgfpathlineto{\pgfpoint{-10.718419\du}{20.297332\du}}
\pgfpathlineto{\pgfpoint{-10.712150\du}{20.300019\du}}
\pgfpathlineto{\pgfpoint{-10.706179\du}{20.302109\du}}
\pgfpathlineto{\pgfpoint{-10.700208\du}{20.304497\du}}
\pgfpathlineto{\pgfpoint{-10.692744\du}{20.306587\du}}
\pgfpathlineto{\pgfpoint{-10.686176\du}{20.308677\du}}
\pgfpathlineto{\pgfpoint{-10.678712\du}{20.310468\du}}
\pgfpathlineto{\pgfpoint{-10.670651\du}{20.311961\du}}
\pgfpathlineto{\pgfpoint{-10.662889\du}{20.314051\du}}
\pgfpathlineto{\pgfpoint{-10.654529\du}{20.315543\du}}
\pgfpathlineto{\pgfpoint{-10.645573\du}{20.317036\du}}
\pgfpathlineto{\pgfpoint{-10.637213\du}{20.318230\du}}
\pgfpathlineto{\pgfpoint{-10.627958\du}{20.319425\du}}
\pgfpathlineto{\pgfpoint{-10.619001\du}{20.320022\du}}
\pgfpathlineto{\pgfpoint{-10.609746\du}{20.321216\du}}
\pgfpathlineto{\pgfpoint{-10.600193\du}{20.321514\du}}
\pgfpathlineto{\pgfpoint{-10.590639\du}{20.322112\du}}
\pgfpathlineto{\pgfpoint{-10.581085\du}{20.322410\du}}
\pgfpathlineto{\pgfpoint{-10.571233\du}{20.323007\du}}
\pgfpathlineto{\pgfpoint{-10.561978\du}{20.323007\du}}
\pgfpathlineto{\pgfpoint{-10.561978\du}{20.323007\du}}
\pgfpathlineto{\pgfpoint{-10.552424\du}{20.323007\du}}
\pgfpathlineto{\pgfpoint{-10.542572\du}{20.322410\du}}
\pgfpathlineto{\pgfpoint{-10.532720\du}{20.322112\du}}
\pgfpathlineto{\pgfpoint{-10.523465\du}{20.321514\du}}
\pgfpathlineto{\pgfpoint{-10.513911\du}{20.321216\du}}
\pgfpathlineto{\pgfpoint{-10.504358\du}{20.320022\du}}
\pgfpathlineto{\pgfpoint{-10.495700\du}{20.319425\du}}
\pgfpathlineto{\pgfpoint{-10.486444\du}{20.318230\du}}
\pgfpathlineto{\pgfpoint{-10.478085\du}{20.317036\du}}
\pgfpathlineto{\pgfpoint{-10.469725\du}{20.315543\du}}
\pgfpathlineto{\pgfpoint{-10.460769\du}{20.314051\du}}
\pgfpathlineto{\pgfpoint{-10.453007\du}{20.311961\du}}
\pgfpathlineto{\pgfpoint{-10.445543\du}{20.310468\du}}
\pgfpathlineto{\pgfpoint{-10.437482\du}{20.308677\du}}
\pgfpathlineto{\pgfpoint{-10.430914\du}{20.306587\du}}
\pgfpathlineto{\pgfpoint{-10.424047\du}{20.304497\du}}
\pgfpathlineto{\pgfpoint{-10.417479\du}{20.302109\du}}
\pgfpathlineto{\pgfpoint{-10.411508\du}{20.300019\du}}
\pgfpathlineto{\pgfpoint{-10.405835\du}{20.297332\du}}
\pgfpathlineto{\pgfpoint{-10.400760\du}{20.295242\du}}
\pgfpathlineto{\pgfpoint{-10.395685\du}{20.292555\du}}
\pgfpathlineto{\pgfpoint{-10.390908\du}{20.289868\du}}
\pgfpathlineto{\pgfpoint{-10.386728\du}{20.287181\du}}
\pgfpathlineto{\pgfpoint{-10.383444\du}{20.284494\du}}
\pgfpathlineto{\pgfpoint{-10.380160\du}{20.281807\du}}
\pgfpathlineto{\pgfpoint{-10.378070\du}{20.278224\du}}
\pgfpathlineto{\pgfpoint{-10.375980\du}{20.275537\du}}
\pgfpathlineto{\pgfpoint{-10.373890\du}{20.272253\du}}
\pgfpathlineto{\pgfpoint{-10.372696\du}{20.269566\du}}
\pgfpathlineto{\pgfpoint{-10.372099\du}{20.266879\du}}
\pgfpathlineto{\pgfpoint{-10.371800\du}{20.263297\du}}
\pgfusepath{stroke}
\pgfsetbuttcap
\pgfsetmiterjoin
\pgfsetdash{}{0pt}
\definecolor{dialinecolor}{rgb}{0.000000, 0.000000, 0.000000}
\pgfsetfillcolor{dialinecolor}
\pgfpathmoveto{\pgfpoint{-10.550036\du}{20.251355\du}}
\pgfpathlineto{\pgfpoint{-10.526152\du}{20.263297\du}}
\pgfpathlineto{\pgfpoint{-10.454798\du}{20.239711\du}}
\pgfpathlineto{\pgfpoint{-10.418972\du}{20.251355\du}}
\pgfpathlineto{\pgfpoint{-10.442557\du}{20.227471\du}}
\pgfpathlineto{\pgfpoint{-10.526152\du}{20.227471\du}}
\pgfpathlineto{\pgfpoint{-10.490326\du}{20.227471\du}}
\pgfpathlineto{\pgfpoint{-10.550036\du}{20.251355\du}}
\pgfusepath{fill}
\pgfsetbuttcap
\pgfsetmiterjoin
\pgfsetdash{}{0pt}
\definecolor{dialinecolor}{rgb}{0.000000, 0.000000, 0.000000}
\pgfsetfillcolor{dialinecolor}
\pgfpathmoveto{\pgfpoint{-10.561978\du}{20.287181\du}}
\pgfpathlineto{\pgfpoint{-10.597506\du}{20.275239\du}}
\pgfpathlineto{\pgfpoint{-10.657216\du}{20.299123\du}}
\pgfpathlineto{\pgfpoint{-10.692744\du}{20.287181\du}}
\pgfpathlineto{\pgfpoint{-10.680503\du}{20.323007\du}}
\pgfpathlineto{\pgfpoint{-10.597506\du}{20.323007\du}}
\pgfpathlineto{\pgfpoint{-10.633033\du}{20.311364\du}}
\pgfpathlineto{\pgfpoint{-10.561978\du}{20.287181\du}}
\pgfusepath{fill}
\pgfsetbuttcap
\pgfsetmiterjoin
\pgfsetdash{}{0pt}
\definecolor{dialinecolor}{rgb}{0.000000, 0.000000, 0.000000}
\pgfsetfillcolor{dialinecolor}
\pgfpathmoveto{\pgfpoint{-10.680503\du}{20.227471\du}}
\pgfpathlineto{\pgfpoint{-10.657216\du}{20.215827\du}}
\pgfpathlineto{\pgfpoint{-10.585862\du}{20.251355\du}}
\pgfpathlineto{\pgfpoint{-10.561978\du}{20.239711\du}}
\pgfpathlineto{\pgfpoint{-10.573323\du}{20.263297\du}}
\pgfpathlineto{\pgfpoint{-10.657216\du}{20.263297\du}}
\pgfpathlineto{\pgfpoint{-10.621390\du}{20.251355\du}}
\pgfpathlineto{\pgfpoint{-10.680503\du}{20.227471\du}}
\pgfusepath{fill}
\pgfsetbuttcap
\pgfsetmiterjoin
\pgfsetdash{}{0pt}
\definecolor{dialinecolor}{rgb}{0.000000, 0.000000, 0.000000}
\pgfsetfillcolor{dialinecolor}
\pgfpathmoveto{\pgfpoint{-10.430914\du}{20.311364\du}}
\pgfpathlineto{\pgfpoint{-10.454798\du}{20.323007\du}}
\pgfpathlineto{\pgfpoint{-10.526152\du}{20.287181\du}}
\pgfpathlineto{\pgfpoint{-10.561978\du}{20.299123\du}}
\pgfpathlineto{\pgfpoint{-10.537795\du}{20.275239\du}}
\pgfpathlineto{\pgfpoint{-10.454798\du}{20.275239\du}}
\pgfpathlineto{\pgfpoint{-10.490326\du}{20.287181\du}}
\pgfpathlineto{\pgfpoint{-10.430914\du}{20.311364\du}}
\pgfusepath{fill}
\pgfsetbuttcap
\pgfsetmiterjoin
\pgfsetdash{}{0pt}
\definecolor{dialinecolor}{rgb}{1.000000, 1.000000, 1.000000}
\pgfsetfillcolor{dialinecolor}
\pgfpathmoveto{\pgfpoint{-10.550036\du}{20.263297\du}}
\pgfpathlineto{\pgfpoint{-10.526152\du}{20.263297\du}}
\pgfpathlineto{\pgfpoint{-10.454798\du}{20.239711\du}}
\pgfpathlineto{\pgfpoint{-10.418972\du}{20.251355\du}}
\pgfpathlineto{\pgfpoint{-10.442557\du}{20.227471\du}}
\pgfpathlineto{\pgfpoint{-10.513911\du}{20.227471\du}}
\pgfpathlineto{\pgfpoint{-10.490326\du}{20.227471\du}}
\pgfpathlineto{\pgfpoint{-10.550036\du}{20.263297\du}}
\pgfusepath{fill}
\pgfsetbuttcap
\pgfsetmiterjoin
\pgfsetdash{}{0pt}
\definecolor{dialinecolor}{rgb}{1.000000, 1.000000, 1.000000}
\pgfsetfillcolor{dialinecolor}
\pgfpathmoveto{\pgfpoint{-10.561978\du}{20.287181\du}}
\pgfpathlineto{\pgfpoint{-10.585862\du}{20.275239\du}}
\pgfpathlineto{\pgfpoint{-10.657216\du}{20.299123\du}}
\pgfpathlineto{\pgfpoint{-10.692744\du}{20.299123\du}}
\pgfpathlineto{\pgfpoint{-10.668860\du}{20.323007\du}}
\pgfpathlineto{\pgfpoint{-10.585862\du}{20.323007\du}}
\pgfpathlineto{\pgfpoint{-10.621390\du}{20.311364\du}}
\pgfpathlineto{\pgfpoint{-10.561978\du}{20.287181\du}}
\pgfusepath{fill}
\pgfsetbuttcap
\pgfsetmiterjoin
\pgfsetdash{}{0pt}
\definecolor{dialinecolor}{rgb}{1.000000, 1.000000, 1.000000}
\pgfsetfillcolor{dialinecolor}
\pgfpathmoveto{\pgfpoint{-10.680503\du}{20.227471\du}}
\pgfpathlineto{\pgfpoint{-10.657216\du}{20.227471\du}}
\pgfpathlineto{\pgfpoint{-10.585862\du}{20.251355\du}}
\pgfpathlineto{\pgfpoint{-10.550036\du}{20.239711\du}}
\pgfpathlineto{\pgfpoint{-10.573323\du}{20.263297\du}}
\pgfpathlineto{\pgfpoint{-10.657216\du}{20.263297\du}}
\pgfpathlineto{\pgfpoint{-10.621390\du}{20.263297\du}}
\pgfpathlineto{\pgfpoint{-10.680503\du}{20.227471\du}}
\pgfusepath{fill}
\pgfsetbuttcap
\pgfsetmiterjoin
\pgfsetdash{}{0pt}
\definecolor{dialinecolor}{rgb}{1.000000, 1.000000, 1.000000}
\pgfsetfillcolor{dialinecolor}
\pgfpathmoveto{\pgfpoint{-10.430914\du}{20.311364\du}}
\pgfpathlineto{\pgfpoint{-10.454798\du}{20.323007\du}}
\pgfpathlineto{\pgfpoint{-10.526152\du}{20.299123\du}}
\pgfpathlineto{\pgfpoint{-10.550036\du}{20.299123\du}}
\pgfpathlineto{\pgfpoint{-10.537795\du}{20.275239\du}}
\pgfpathlineto{\pgfpoint{-10.454798\du}{20.275239\du}}
\pgfpathlineto{\pgfpoint{-10.490326\du}{20.287181\du}}
\pgfpathlineto{\pgfpoint{-10.430914\du}{20.311364\du}}
\pgfusepath{fill}
\pgfsetbuttcap
\pgfsetmiterjoin
\pgfsetdash{}{0pt}
\definecolor{dialinecolor}{rgb}{0.666667, 0.901961, 1.000000}
\pgfsetstrokecolor{dialinecolor}
\pgfpathmoveto{\pgfpoint{-10.752454\du}{20.263297\du}}
\pgfpathlineto{\pgfpoint{-10.752454\du}{20.358236\du}}
\pgfusepath{stroke}
\pgfsetbuttcap
\pgfsetmiterjoin
\pgfsetdash{}{0pt}
\definecolor{dialinecolor}{rgb}{0.666667, 0.901961, 1.000000}
\pgfsetstrokecolor{dialinecolor}
\pgfpathmoveto{\pgfpoint{-10.371800\du}{20.263297\du}}
\pgfpathlineto{\pgfpoint{-10.371800\du}{20.358236\du}}
\pgfusepath{stroke}
\pgfsetlinewidth{0.100000\du}
\pgfsetdash{}{0pt}
\pgfsetdash{}{0pt}
\pgfsetbuttcap
\pgfsetmiterjoin
\pgfsetlinewidth{0.001000\du}
\pgfsetbuttcap
\pgfsetmiterjoin
\pgfsetdash{}{0pt}
\definecolor{dialinecolor}{rgb}{0.000000, 0.470588, 0.666667}
\pgfsetfillcolor{dialinecolor}
\pgfpathmoveto{\pgfpoint{-12.370418\du}{17.216422\du}}
\pgfpathlineto{\pgfpoint{-12.371706\du}{17.201227\du}}
\pgfpathlineto{\pgfpoint{-12.375054\du}{17.185773\du}}
\pgfpathlineto{\pgfpoint{-12.380720\du}{17.170578\du}}
\pgfpathlineto{\pgfpoint{-12.388705\du}{17.155639\du}}
\pgfpathlineto{\pgfpoint{-12.399007\du}{17.140444\du}}
\pgfpathlineto{\pgfpoint{-12.411369\du}{17.125506\du}}
\pgfpathlineto{\pgfpoint{-12.426050\du}{17.111083\du}}
\pgfpathlineto{\pgfpoint{-12.442534\du}{17.096917\du}}
\pgfpathlineto{\pgfpoint{-12.461593\du}{17.083009\du}}
\pgfpathlineto{\pgfpoint{-12.482455\du}{17.069616\du}}
\pgfpathlineto{\pgfpoint{-12.505377\du}{17.056223\du}}
\pgfpathlineto{\pgfpoint{-12.530360\du}{17.043346\du}}
\pgfpathlineto{\pgfpoint{-12.557145\du}{17.030983\du}}
\pgfpathlineto{\pgfpoint{-12.585734\du}{17.018878\du}}
\pgfpathlineto{\pgfpoint{-12.616125\du}{17.007546\du}}
\pgfpathlineto{\pgfpoint{-12.647805\du}{16.996728\du}}
\pgfpathlineto{\pgfpoint{-12.681287\du}{16.986426\du}}
\pgfpathlineto{\pgfpoint{-12.716572\du}{16.976896\du}}
\pgfpathlineto{\pgfpoint{-12.752887\du}{16.967367\du}}
\pgfpathlineto{\pgfpoint{-12.790490\du}{16.959383\du}}
\pgfpathlineto{\pgfpoint{-12.829381\du}{16.951399\du}}
\pgfpathlineto{\pgfpoint{-12.869559\du}{16.944445\du}}
\pgfpathlineto{\pgfpoint{-12.910511\du}{16.937748\du}}
\pgfpathlineto{\pgfpoint{-12.952749\du}{16.932082\du}}
\pgfpathlineto{\pgfpoint{-12.995246\du}{16.927188\du}}
\pgfpathlineto{\pgfpoint{-13.039030\du}{16.923068\du}}
\pgfpathlineto{\pgfpoint{-13.083072\du}{16.919462\du}}
\pgfpathlineto{\pgfpoint{-13.127372\du}{16.916629\du}}
\pgfpathlineto{\pgfpoint{-13.172186\du}{16.914568\du}}
\pgfpathlineto{\pgfpoint{-13.217258\du}{16.913796\du}}
\pgfpathlineto{\pgfpoint{-13.262845\du}{16.913023\du}}
\pgfpathlineto{\pgfpoint{-13.262845\du}{16.913023\du}}
\pgfpathlineto{\pgfpoint{-13.307918\du}{16.913796\du}}
\pgfpathlineto{\pgfpoint{-13.352732\du}{16.914568\du}}
\pgfpathlineto{\pgfpoint{-13.397547\du}{16.916629\du}}
\pgfpathlineto{\pgfpoint{-13.441846\du}{16.919462\du}}
\pgfpathlineto{\pgfpoint{-13.485888\du}{16.923068\du}}
\pgfpathlineto{\pgfpoint{-13.529672\du}{16.927188\du}}
\pgfpathlineto{\pgfpoint{-13.572169\du}{16.932082\du}}
\pgfpathlineto{\pgfpoint{-13.614408\du}{16.937748\du}}
\pgfpathlineto{\pgfpoint{-13.655359\du}{16.944445\du}}
\pgfpathlineto{\pgfpoint{-13.695537\du}{16.951399\du}}
\pgfpathlineto{\pgfpoint{-13.734686\du}{16.959383\du}}
\pgfpathlineto{\pgfpoint{-13.772031\du}{16.967367\du}}
\pgfpathlineto{\pgfpoint{-13.808346\du}{16.976896\du}}
\pgfpathlineto{\pgfpoint{-13.843631\du}{16.986426\du}}
\pgfpathlineto{\pgfpoint{-13.877114\du}{16.996728\du}}
\pgfpathlineto{\pgfpoint{-13.908793\du}{17.007546\du}}
\pgfpathlineto{\pgfpoint{-13.939184\du}{17.018878\du}}
\pgfpathlineto{\pgfpoint{-13.967773\du}{17.030983\du}}
\pgfpathlineto{\pgfpoint{-13.994558\du}{17.043346\du}}
\pgfpathlineto{\pgfpoint{-14.019799\du}{17.056223\du}}
\pgfpathlineto{\pgfpoint{-14.042464\du}{17.069616\du}}
\pgfpathlineto{\pgfpoint{-14.063583\du}{17.083009\du}}
\pgfpathlineto{\pgfpoint{-14.082385\du}{17.096917\du}}
\pgfpathlineto{\pgfpoint{-14.098868\du}{17.111083\du}}
\pgfpathlineto{\pgfpoint{-14.113549\du}{17.125506\du}}
\pgfpathlineto{\pgfpoint{-14.125911\du}{17.140444\du}}
\pgfpathlineto{\pgfpoint{-14.136214\du}{17.155639\du}}
\pgfpathlineto{\pgfpoint{-14.144198\du}{17.170578\du}}
\pgfpathlineto{\pgfpoint{-14.150122\du}{17.185773\du}}
\pgfpathlineto{\pgfpoint{-14.153212\du}{17.201227\du}}
\pgfpathlineto{\pgfpoint{-14.154500\du}{17.216422\du}}
\pgfpathlineto{\pgfpoint{-14.154500\du}{17.216422\du}}
\pgfpathlineto{\pgfpoint{-14.153212\du}{17.232133\du}}
\pgfpathlineto{\pgfpoint{-14.150122\du}{17.247329\du}}
\pgfpathlineto{\pgfpoint{-14.144198\du}{17.262525\du}}
\pgfpathlineto{\pgfpoint{-14.136214\du}{17.277720\du}}
\pgfpathlineto{\pgfpoint{-14.125911\du}{17.292659\du}}
\pgfpathlineto{\pgfpoint{-14.113549\du}{17.307597\du}}
\pgfpathlineto{\pgfpoint{-14.098868\du}{17.322020\du}}
\pgfpathlineto{\pgfpoint{-14.082385\du}{17.336185\du}}
\pgfpathlineto{\pgfpoint{-14.063583\du}{17.350093\du}}
\pgfpathlineto{\pgfpoint{-14.042464\du}{17.363744\du}}
\pgfpathlineto{\pgfpoint{-14.019799\du}{17.376621\du}}
\pgfpathlineto{\pgfpoint{-13.994558\du}{17.390014\du}}
\pgfpathlineto{\pgfpoint{-13.967773\du}{17.402377\du}}
\pgfpathlineto{\pgfpoint{-13.939184\du}{17.414224\du}}
\pgfpathlineto{\pgfpoint{-13.908793\du}{17.425299\du}}
\pgfpathlineto{\pgfpoint{-13.877114\du}{17.436374\du}}
\pgfpathlineto{\pgfpoint{-13.843631\du}{17.446676\du}}
\pgfpathlineto{\pgfpoint{-13.808346\du}{17.456463\du}}
\pgfpathlineto{\pgfpoint{-13.772031\du}{17.465478\du}}
\pgfpathlineto{\pgfpoint{-13.734686\du}{17.473977\du}}
\pgfpathlineto{\pgfpoint{-13.695537\du}{17.481704\du}}
\pgfpathlineto{\pgfpoint{-13.655359\du}{17.488658\du}}
\pgfpathlineto{\pgfpoint{-13.614408\du}{17.495354\du}}
\pgfpathlineto{\pgfpoint{-13.572169\du}{17.501020\du}}
\pgfpathlineto{\pgfpoint{-13.529672\du}{17.505914\du}}
\pgfpathlineto{\pgfpoint{-13.485888\du}{17.510292\du}}
\pgfpathlineto{\pgfpoint{-13.441846\du}{17.513383\du}}
\pgfpathlineto{\pgfpoint{-13.397547\du}{17.516216\du}}
\pgfpathlineto{\pgfpoint{-13.352732\du}{17.518277\du}}
\pgfpathlineto{\pgfpoint{-13.307918\du}{17.519564\du}}
\pgfpathlineto{\pgfpoint{-13.262845\du}{17.519822\du}}
\pgfpathlineto{\pgfpoint{-13.262845\du}{17.519822\du}}
\pgfpathlineto{\pgfpoint{-13.217258\du}{17.519564\du}}
\pgfpathlineto{\pgfpoint{-13.172186\du}{17.518277\du}}
\pgfpathlineto{\pgfpoint{-13.127372\du}{17.516216\du}}
\pgfpathlineto{\pgfpoint{-13.083072\du}{17.513383\du}}
\pgfpathlineto{\pgfpoint{-13.039030\du}{17.510292\du}}
\pgfpathlineto{\pgfpoint{-12.995246\du}{17.505914\du}}
\pgfpathlineto{\pgfpoint{-12.952749\du}{17.501020\du}}
\pgfpathlineto{\pgfpoint{-12.910511\du}{17.495354\du}}
\pgfpathlineto{\pgfpoint{-12.869559\du}{17.488658\du}}
\pgfpathlineto{\pgfpoint{-12.829381\du}{17.481704\du}}
\pgfpathlineto{\pgfpoint{-12.790490\du}{17.473977\du}}
\pgfpathlineto{\pgfpoint{-12.752887\du}{17.465478\du}}
\pgfpathlineto{\pgfpoint{-12.716572\du}{17.456463\du}}
\pgfpathlineto{\pgfpoint{-12.681287\du}{17.446676\du}}
\pgfpathlineto{\pgfpoint{-12.647805\du}{17.436374\du}}
\pgfpathlineto{\pgfpoint{-12.616125\du}{17.425299\du}}
\pgfpathlineto{\pgfpoint{-12.585734\du}{17.414224\du}}
\pgfpathlineto{\pgfpoint{-12.557145\du}{17.402377\du}}
\pgfpathlineto{\pgfpoint{-12.530360\du}{17.390014\du}}
\pgfpathlineto{\pgfpoint{-12.505377\du}{17.376621\du}}
\pgfpathlineto{\pgfpoint{-12.482455\du}{17.363744\du}}
\pgfpathlineto{\pgfpoint{-12.461593\du}{17.350093\du}}
\pgfpathlineto{\pgfpoint{-12.442534\du}{17.336185\du}}
\pgfpathlineto{\pgfpoint{-12.426050\du}{17.322020\du}}
\pgfpathlineto{\pgfpoint{-12.411369\du}{17.307597\du}}
\pgfpathlineto{\pgfpoint{-12.399007\du}{17.292659\du}}
\pgfpathlineto{\pgfpoint{-12.388705\du}{17.277720\du}}
\pgfpathlineto{\pgfpoint{-12.380720\du}{17.262525\du}}
\pgfpathlineto{\pgfpoint{-12.375054\du}{17.247329\du}}
\pgfpathlineto{\pgfpoint{-12.371706\du}{17.232133\du}}
\pgfpathlineto{\pgfpoint{-12.370418\du}{17.216422\du}}
\pgfusepath{fill}
\pgfsetbuttcap
\pgfsetmiterjoin
\pgfsetdash{}{0pt}
\definecolor{dialinecolor}{rgb}{0.666667, 0.901961, 1.000000}
\pgfsetstrokecolor{dialinecolor}
\pgfpathmoveto{\pgfpoint{-12.380720\du}{17.211786\du}}
\pgfpathlineto{\pgfpoint{-12.382008\du}{17.196333\du}}
\pgfpathlineto{\pgfpoint{-12.385356\du}{17.181653\du}}
\pgfpathlineto{\pgfpoint{-12.391023\du}{17.166457\du}}
\pgfpathlineto{\pgfpoint{-12.398749\du}{17.151776\du}}
\pgfpathlineto{\pgfpoint{-12.409051\du}{17.136838\du}}
\pgfpathlineto{\pgfpoint{-12.421672\du}{17.122415\du}}
\pgfpathlineto{\pgfpoint{-12.435580\du}{17.108249\du}}
\pgfpathlineto{\pgfpoint{-12.452578\du}{17.094084\du}}
\pgfpathlineto{\pgfpoint{-12.471380\du}{17.080434\du}}
\pgfpathlineto{\pgfpoint{-12.492242\du}{17.066783\du}}
\pgfpathlineto{\pgfpoint{-12.514906\du}{17.054163\du}}
\pgfpathlineto{\pgfpoint{-12.539889\du}{17.041543\du}}
\pgfpathlineto{\pgfpoint{-12.566160\du}{17.028923\du}}
\pgfpathlineto{\pgfpoint{-12.594491\du}{17.017333\du}}
\pgfpathlineto{\pgfpoint{-12.624882\du}{17.006000\du}}
\pgfpathlineto{\pgfpoint{-12.656819\du}{16.995698\du}}
\pgfpathlineto{\pgfpoint{-12.690044\du}{16.985396\du}}
\pgfpathlineto{\pgfpoint{-12.724814\du}{16.975866\du}}
\pgfpathlineto{\pgfpoint{-12.760614\du}{16.967109\du}}
\pgfpathlineto{\pgfpoint{-12.798474\du}{16.958610\du}}
\pgfpathlineto{\pgfpoint{-12.837365\du}{16.951141\du}}
\pgfpathlineto{\pgfpoint{-12.876771\du}{16.944187\du}}
\pgfpathlineto{\pgfpoint{-12.917722\du}{16.937748\du}}
\pgfpathlineto{\pgfpoint{-12.959446\du}{16.932082\du}}
\pgfpathlineto{\pgfpoint{-13.001942\du}{16.927188\du}}
\pgfpathlineto{\pgfpoint{-13.045469\du}{16.923068\du}}
\pgfpathlineto{\pgfpoint{-13.088996\du}{16.919462\du}}
\pgfpathlineto{\pgfpoint{-13.133038\du}{16.917401\du}}
\pgfpathlineto{\pgfpoint{-13.177852\du}{16.915083\du}}
\pgfpathlineto{\pgfpoint{-13.222667\du}{16.913796\du}}
\pgfpathlineto{\pgfpoint{-13.267739\du}{16.913796\du}}
\pgfpathlineto{\pgfpoint{-13.267739\du}{16.913796\du}}
\pgfpathlineto{\pgfpoint{-13.312553\du}{16.913796\du}}
\pgfpathlineto{\pgfpoint{-13.357626\du}{16.915083\du}}
\pgfpathlineto{\pgfpoint{-13.402183\du}{16.917401\du}}
\pgfpathlineto{\pgfpoint{-13.446224\du}{16.919462\du}}
\pgfpathlineto{\pgfpoint{-13.490009\du}{16.923068\du}}
\pgfpathlineto{\pgfpoint{-13.533278\du}{16.927188\du}}
\pgfpathlineto{\pgfpoint{-13.575774\du}{16.932082\du}}
\pgfpathlineto{\pgfpoint{-13.617756\du}{16.937748\du}}
\pgfpathlineto{\pgfpoint{-13.658450\du}{16.944187\du}}
\pgfpathlineto{\pgfpoint{-13.697855\du}{16.951141\du}}
\pgfpathlineto{\pgfpoint{-13.737004\du}{16.958610\du}}
\pgfpathlineto{\pgfpoint{-13.774349\du}{16.967109\du}}
\pgfpathlineto{\pgfpoint{-13.810149\du}{16.975866\du}}
\pgfpathlineto{\pgfpoint{-13.845177\du}{16.985396\du}}
\pgfpathlineto{\pgfpoint{-13.878916\du}{16.995698\du}}
\pgfpathlineto{\pgfpoint{-13.910853\du}{17.006000\du}}
\pgfpathlineto{\pgfpoint{-13.940730\du}{17.017333\du}}
\pgfpathlineto{\pgfpoint{-13.969318\du}{17.028923\du}}
\pgfpathlineto{\pgfpoint{-13.995846\du}{17.041543\du}}
\pgfpathlineto{\pgfpoint{-14.020314\du}{17.054163\du}}
\pgfpathlineto{\pgfpoint{-14.042721\du}{17.066783\du}}
\pgfpathlineto{\pgfpoint{-14.063841\du}{17.080434\du}}
\pgfpathlineto{\pgfpoint{-14.082900\du}{17.094084\du}}
\pgfpathlineto{\pgfpoint{-14.099383\du}{17.108249\du}}
\pgfpathlineto{\pgfpoint{-14.113806\du}{17.122415\du}}
\pgfpathlineto{\pgfpoint{-14.126684\du}{17.136838\du}}
\pgfpathlineto{\pgfpoint{-14.136214\du}{17.151776\du}}
\pgfpathlineto{\pgfpoint{-14.144455\du}{17.166457\du}}
\pgfpathlineto{\pgfpoint{-14.150122\du}{17.181653\du}}
\pgfpathlineto{\pgfpoint{-14.153212\du}{17.196333\du}}
\pgfpathlineto{\pgfpoint{-14.154500\du}{17.211786\du}}
\pgfpathlineto{\pgfpoint{-14.154500\du}{17.211786\du}}
\pgfpathlineto{\pgfpoint{-14.153212\du}{17.226982\du}}
\pgfpathlineto{\pgfpoint{-14.150122\du}{17.241920\du}}
\pgfpathlineto{\pgfpoint{-14.144455\du}{17.256859\du}}
\pgfpathlineto{\pgfpoint{-14.136214\du}{17.271539\du}}
\pgfpathlineto{\pgfpoint{-14.126684\du}{17.286477\du}}
\pgfpathlineto{\pgfpoint{-14.113806\du}{17.300900\du}}
\pgfpathlineto{\pgfpoint{-14.099383\du}{17.315066\du}}
\pgfpathlineto{\pgfpoint{-14.082900\du}{17.328716\du}}
\pgfpathlineto{\pgfpoint{-14.063841\du}{17.342882\du}}
\pgfpathlineto{\pgfpoint{-14.042721\du}{17.356532\du}}
\pgfpathlineto{\pgfpoint{-14.020314\du}{17.369410\du}}
\pgfpathlineto{\pgfpoint{-13.995846\du}{17.381773\du}}
\pgfpathlineto{\pgfpoint{-13.969318\du}{17.394135\du}}
\pgfpathlineto{\pgfpoint{-13.940730\du}{17.405983\du}}
\pgfpathlineto{\pgfpoint{-13.910853\du}{17.417058\du}}
\pgfpathlineto{\pgfpoint{-13.878916\du}{17.427617\du}}
\pgfpathlineto{\pgfpoint{-13.845177\du}{17.437404\du}}
\pgfpathlineto{\pgfpoint{-13.810149\du}{17.447449\du}}
\pgfpathlineto{\pgfpoint{-13.774349\du}{17.456206\du}}
\pgfpathlineto{\pgfpoint{-13.737004\du}{17.464448\du}}
\pgfpathlineto{\pgfpoint{-13.697855\du}{17.472432\du}}
\pgfpathlineto{\pgfpoint{-13.658450\du}{17.479128\du}}
\pgfpathlineto{\pgfpoint{-13.617756\du}{17.485310\du}}
\pgfpathlineto{\pgfpoint{-13.575774\du}{17.491233\du}}
\pgfpathlineto{\pgfpoint{-13.533278\du}{17.495869\du}}
\pgfpathlineto{\pgfpoint{-13.490009\du}{17.499733\du}}
\pgfpathlineto{\pgfpoint{-13.446224\du}{17.503596\du}}
\pgfpathlineto{\pgfpoint{-13.402183\du}{17.505914\du}}
\pgfpathlineto{\pgfpoint{-13.357626\du}{17.507974\du}}
\pgfpathlineto{\pgfpoint{-13.312553\du}{17.509005\du}}
\pgfpathlineto{\pgfpoint{-13.267739\du}{17.509520\du}}
\pgfpathlineto{\pgfpoint{-13.267739\du}{17.509520\du}}
\pgfpathlineto{\pgfpoint{-13.222667\du}{17.509005\du}}
\pgfpathlineto{\pgfpoint{-13.177852\du}{17.507974\du}}
\pgfpathlineto{\pgfpoint{-13.133038\du}{17.505914\du}}
\pgfpathlineto{\pgfpoint{-13.088996\du}{17.503596\du}}
\pgfpathlineto{\pgfpoint{-13.045469\du}{17.499733\du}}
\pgfpathlineto{\pgfpoint{-13.001942\du}{17.495869\du}}
\pgfpathlineto{\pgfpoint{-12.959446\du}{17.491233\du}}
\pgfpathlineto{\pgfpoint{-12.917722\du}{17.485310\du}}
\pgfpathlineto{\pgfpoint{-12.876771\du}{17.479128\du}}
\pgfpathlineto{\pgfpoint{-12.837365\du}{17.472432\du}}
\pgfpathlineto{\pgfpoint{-12.798474\du}{17.464448\du}}
\pgfpathlineto{\pgfpoint{-12.760614\du}{17.456206\du}}
\pgfpathlineto{\pgfpoint{-12.724814\du}{17.447449\du}}
\pgfpathlineto{\pgfpoint{-12.690044\du}{17.437404\du}}
\pgfpathlineto{\pgfpoint{-12.656819\du}{17.427617\du}}
\pgfpathlineto{\pgfpoint{-12.624882\du}{17.417058\du}}
\pgfpathlineto{\pgfpoint{-12.594491\du}{17.405983\du}}
\pgfpathlineto{\pgfpoint{-12.566160\du}{17.394135\du}}
\pgfpathlineto{\pgfpoint{-12.539889\du}{17.381773\du}}
\pgfpathlineto{\pgfpoint{-12.514906\du}{17.369410\du}}
\pgfpathlineto{\pgfpoint{-12.492242\du}{17.356532\du}}
\pgfpathlineto{\pgfpoint{-12.471380\du}{17.342882\du}}
\pgfpathlineto{\pgfpoint{-12.452578\du}{17.328716\du}}
\pgfpathlineto{\pgfpoint{-12.435580\du}{17.315066\du}}
\pgfpathlineto{\pgfpoint{-12.421672\du}{17.300900\du}}
\pgfpathlineto{\pgfpoint{-12.409051\du}{17.286477\du}}
\pgfpathlineto{\pgfpoint{-12.398749\du}{17.271539\du}}
\pgfpathlineto{\pgfpoint{-12.391023\du}{17.256859\du}}
\pgfpathlineto{\pgfpoint{-12.385356\du}{17.241920\du}}
\pgfpathlineto{\pgfpoint{-12.382008\du}{17.226982\du}}
\pgfpathlineto{\pgfpoint{-12.380720\du}{17.211786\du}}
\pgfusepath{stroke}
\pgfsetbuttcap
\pgfsetmiterjoin
\pgfsetdash{}{0pt}
\definecolor{dialinecolor}{rgb}{0.000000, 0.470588, 0.666667}
\pgfsetfillcolor{dialinecolor}
\pgfpathmoveto{\pgfpoint{-14.154500\du}{16.796608\du}}
\pgfpathlineto{\pgfpoint{-14.154500\du}{17.222089\du}}
\pgfpathlineto{\pgfpoint{-12.380720\du}{17.222089\du}}
\pgfpathlineto{\pgfpoint{-12.380720\du}{16.796608\du}}
\pgfpathlineto{\pgfpoint{-14.154500\du}{16.796608\du}}
\pgfusepath{fill}
\pgfsetbuttcap
\pgfsetmiterjoin
\pgfsetdash{}{0pt}
\definecolor{dialinecolor}{rgb}{0.000000, 0.705882, 1.000000}
\pgfsetfillcolor{dialinecolor}
\pgfpathmoveto{\pgfpoint{-12.370418\du}{16.790684\du}}
\pgfpathlineto{\pgfpoint{-12.371706\du}{16.775746\du}}
\pgfpathlineto{\pgfpoint{-12.375054\du}{16.760293\du}}
\pgfpathlineto{\pgfpoint{-12.380720\du}{16.745097\du}}
\pgfpathlineto{\pgfpoint{-12.388705\du}{16.729901\du}}
\pgfpathlineto{\pgfpoint{-12.399007\du}{16.714706\du}}
\pgfpathlineto{\pgfpoint{-12.411369\du}{16.700025\du}}
\pgfpathlineto{\pgfpoint{-12.426050\du}{16.685602\du}}
\pgfpathlineto{\pgfpoint{-12.442534\du}{16.671179\du}}
\pgfpathlineto{\pgfpoint{-12.461593\du}{16.657529\du}}
\pgfpathlineto{\pgfpoint{-12.482455\du}{16.643878\du}}
\pgfpathlineto{\pgfpoint{-12.505377\du}{16.630743\du}}
\pgfpathlineto{\pgfpoint{-12.530360\du}{16.617350\du}}
\pgfpathlineto{\pgfpoint{-12.557145\du}{16.605245\du}}
\pgfpathlineto{\pgfpoint{-12.585734\du}{16.593140\du}}
\pgfpathlineto{\pgfpoint{-12.616125\du}{16.581807\du}}
\pgfpathlineto{\pgfpoint{-12.647805\du}{16.571248\du}}
\pgfpathlineto{\pgfpoint{-12.681287\du}{16.560946\du}}
\pgfpathlineto{\pgfpoint{-12.716572\du}{16.551158\du}}
\pgfpathlineto{\pgfpoint{-12.752887\du}{16.542144\du}}
\pgfpathlineto{\pgfpoint{-12.790490\du}{16.533645\du}}
\pgfpathlineto{\pgfpoint{-12.829381\du}{16.525661\du}}
\pgfpathlineto{\pgfpoint{-12.869559\du}{16.518964\du}}
\pgfpathlineto{\pgfpoint{-12.910511\du}{16.512010\du}}
\pgfpathlineto{\pgfpoint{-12.952749\du}{16.506344\du}}
\pgfpathlineto{\pgfpoint{-12.995246\du}{16.501450\du}}
\pgfpathlineto{\pgfpoint{-13.039030\du}{16.497330\du}}
\pgfpathlineto{\pgfpoint{-13.083072\du}{16.494239\du}}
\pgfpathlineto{\pgfpoint{-13.127372\du}{16.491406\du}}
\pgfpathlineto{\pgfpoint{-13.172186\du}{16.489345\du}}
\pgfpathlineto{\pgfpoint{-13.217258\du}{16.487800\du}}
\pgfpathlineto{\pgfpoint{-13.262845\du}{16.487800\du}}
\pgfpathlineto{\pgfpoint{-13.262845\du}{16.487800\du}}
\pgfpathlineto{\pgfpoint{-13.307918\du}{16.487800\du}}
\pgfpathlineto{\pgfpoint{-13.352732\du}{16.489345\du}}
\pgfpathlineto{\pgfpoint{-13.397547\du}{16.491406\du}}
\pgfpathlineto{\pgfpoint{-13.441846\du}{16.494239\du}}
\pgfpathlineto{\pgfpoint{-13.485888\du}{16.497330\du}}
\pgfpathlineto{\pgfpoint{-13.529672\du}{16.501450\du}}
\pgfpathlineto{\pgfpoint{-13.572169\du}{16.506344\du}}
\pgfpathlineto{\pgfpoint{-13.614408\du}{16.512010\du}}
\pgfpathlineto{\pgfpoint{-13.655359\du}{16.518964\du}}
\pgfpathlineto{\pgfpoint{-13.695537\du}{16.525661\du}}
\pgfpathlineto{\pgfpoint{-13.734686\du}{16.533645\du}}
\pgfpathlineto{\pgfpoint{-13.772031\du}{16.542144\du}}
\pgfpathlineto{\pgfpoint{-13.808346\du}{16.551158\du}}
\pgfpathlineto{\pgfpoint{-13.843631\du}{16.560946\du}}
\pgfpathlineto{\pgfpoint{-13.877114\du}{16.571248\du}}
\pgfpathlineto{\pgfpoint{-13.908793\du}{16.581807\du}}
\pgfpathlineto{\pgfpoint{-13.939184\du}{16.593140\du}}
\pgfpathlineto{\pgfpoint{-13.967773\du}{16.605245\du}}
\pgfpathlineto{\pgfpoint{-13.994558\du}{16.617350\du}}
\pgfpathlineto{\pgfpoint{-14.019799\du}{16.630743\du}}
\pgfpathlineto{\pgfpoint{-14.042464\du}{16.643878\du}}
\pgfpathlineto{\pgfpoint{-14.063583\du}{16.657529\du}}
\pgfpathlineto{\pgfpoint{-14.082385\du}{16.671179\du}}
\pgfpathlineto{\pgfpoint{-14.098868\du}{16.685602\du}}
\pgfpathlineto{\pgfpoint{-14.113549\du}{16.700025\du}}
\pgfpathlineto{\pgfpoint{-14.125911\du}{16.714706\du}}
\pgfpathlineto{\pgfpoint{-14.136214\du}{16.729901\du}}
\pgfpathlineto{\pgfpoint{-14.144198\du}{16.745097\du}}
\pgfpathlineto{\pgfpoint{-14.150122\du}{16.760293\du}}
\pgfpathlineto{\pgfpoint{-14.153212\du}{16.775746\du}}
\pgfpathlineto{\pgfpoint{-14.154500\du}{16.790684\du}}
\pgfpathlineto{\pgfpoint{-14.154500\du}{16.790684\du}}
\pgfpathlineto{\pgfpoint{-14.153212\du}{16.806395\du}}
\pgfpathlineto{\pgfpoint{-14.150122\du}{16.821849\du}}
\pgfpathlineto{\pgfpoint{-14.144198\du}{16.837044\du}}
\pgfpathlineto{\pgfpoint{-14.136214\du}{16.851982\du}}
\pgfpathlineto{\pgfpoint{-14.125911\du}{16.867178\du}}
\pgfpathlineto{\pgfpoint{-14.113549\du}{16.881601\du}}
\pgfpathlineto{\pgfpoint{-14.098868\du}{16.896539\du}}
\pgfpathlineto{\pgfpoint{-14.082385\du}{16.910705\du}}
\pgfpathlineto{\pgfpoint{-14.063583\du}{16.924613\du}}
\pgfpathlineto{\pgfpoint{-14.042464\du}{16.938263\du}}
\pgfpathlineto{\pgfpoint{-14.019799\du}{16.951399\du}}
\pgfpathlineto{\pgfpoint{-13.994558\du}{16.964534\du}}
\pgfpathlineto{\pgfpoint{-13.967773\du}{16.976896\du}}
\pgfpathlineto{\pgfpoint{-13.939184\du}{16.988229\du}}
\pgfpathlineto{\pgfpoint{-13.908793\du}{17.000076\du}}
\pgfpathlineto{\pgfpoint{-13.877114\du}{17.010894\du}}
\pgfpathlineto{\pgfpoint{-13.843631\du}{17.021196\du}}
\pgfpathlineto{\pgfpoint{-13.808346\du}{17.030983\du}}
\pgfpathlineto{\pgfpoint{-13.772031\du}{17.039997\du}}
\pgfpathlineto{\pgfpoint{-13.734686\du}{17.048239\du}}
\pgfpathlineto{\pgfpoint{-13.695537\du}{17.056223\du}}
\pgfpathlineto{\pgfpoint{-13.655359\du}{17.063177\du}}
\pgfpathlineto{\pgfpoint{-13.614408\du}{17.069874\du}}
\pgfpathlineto{\pgfpoint{-13.572169\du}{17.075540\du}}
\pgfpathlineto{\pgfpoint{-13.529672\du}{17.080434\du}}
\pgfpathlineto{\pgfpoint{-13.485888\du}{17.084554\du}}
\pgfpathlineto{\pgfpoint{-13.441846\du}{17.087645\du}}
\pgfpathlineto{\pgfpoint{-13.397547\du}{17.090478\du}}
\pgfpathlineto{\pgfpoint{-13.352732\du}{17.092796\du}}
\pgfpathlineto{\pgfpoint{-13.307918\du}{17.094084\du}}
\pgfpathlineto{\pgfpoint{-13.262845\du}{17.094084\du}}
\pgfpathlineto{\pgfpoint{-13.262845\du}{17.094084\du}}
\pgfpathlineto{\pgfpoint{-13.217258\du}{17.094084\du}}
\pgfpathlineto{\pgfpoint{-13.172186\du}{17.092796\du}}
\pgfpathlineto{\pgfpoint{-13.127372\du}{17.090478\du}}
\pgfpathlineto{\pgfpoint{-13.083072\du}{17.087645\du}}
\pgfpathlineto{\pgfpoint{-13.039030\du}{17.084554\du}}
\pgfpathlineto{\pgfpoint{-12.995246\du}{17.080434\du}}
\pgfpathlineto{\pgfpoint{-12.952749\du}{17.075540\du}}
\pgfpathlineto{\pgfpoint{-12.910511\du}{17.069874\du}}
\pgfpathlineto{\pgfpoint{-12.869559\du}{17.063177\du}}
\pgfpathlineto{\pgfpoint{-12.829381\du}{17.056223\du}}
\pgfpathlineto{\pgfpoint{-12.790490\du}{17.048239\du}}
\pgfpathlineto{\pgfpoint{-12.752887\du}{17.039997\du}}
\pgfpathlineto{\pgfpoint{-12.716572\du}{17.030983\du}}
\pgfpathlineto{\pgfpoint{-12.681287\du}{17.021196\du}}
\pgfpathlineto{\pgfpoint{-12.647805\du}{17.010894\du}}
\pgfpathlineto{\pgfpoint{-12.616125\du}{17.000076\du}}
\pgfpathlineto{\pgfpoint{-12.585734\du}{16.988229\du}}
\pgfpathlineto{\pgfpoint{-12.557145\du}{16.976896\du}}
\pgfpathlineto{\pgfpoint{-12.530360\du}{16.964534\du}}
\pgfpathlineto{\pgfpoint{-12.505377\du}{16.951399\du}}
\pgfpathlineto{\pgfpoint{-12.482455\du}{16.938263\du}}
\pgfpathlineto{\pgfpoint{-12.461593\du}{16.924613\du}}
\pgfpathlineto{\pgfpoint{-12.442534\du}{16.910705\du}}
\pgfpathlineto{\pgfpoint{-12.426050\du}{16.896539\du}}
\pgfpathlineto{\pgfpoint{-12.411369\du}{16.881601\du}}
\pgfpathlineto{\pgfpoint{-12.399007\du}{16.867178\du}}
\pgfpathlineto{\pgfpoint{-12.388705\du}{16.851982\du}}
\pgfpathlineto{\pgfpoint{-12.380720\du}{16.837044\du}}
\pgfpathlineto{\pgfpoint{-12.375054\du}{16.821849\du}}
\pgfpathlineto{\pgfpoint{-12.371706\du}{16.806395\du}}
\pgfpathlineto{\pgfpoint{-12.370418\du}{16.790684\du}}
\pgfusepath{fill}
\pgfsetbuttcap
\pgfsetmiterjoin
\pgfsetdash{}{0pt}
\definecolor{dialinecolor}{rgb}{0.666667, 0.901961, 1.000000}
\pgfsetstrokecolor{dialinecolor}
\pgfpathmoveto{\pgfpoint{-12.380720\du}{16.786306\du}}
\pgfpathlineto{\pgfpoint{-12.382008\du}{16.770853\du}}
\pgfpathlineto{\pgfpoint{-12.385356\du}{16.755657\du}}
\pgfpathlineto{\pgfpoint{-12.391023\du}{16.741234\du}}
\pgfpathlineto{\pgfpoint{-12.398749\du}{16.726038\du}}
\pgfpathlineto{\pgfpoint{-12.409051\du}{16.711358\du}}
\pgfpathlineto{\pgfpoint{-12.421672\du}{16.696934\du}}
\pgfpathlineto{\pgfpoint{-12.435580\du}{16.683027\du}}
\pgfpathlineto{\pgfpoint{-12.452578\du}{16.668603\du}}
\pgfpathlineto{\pgfpoint{-12.471380\du}{16.655211\du}}
\pgfpathlineto{\pgfpoint{-12.492242\du}{16.641303\du}}
\pgfpathlineto{\pgfpoint{-12.514906\du}{16.628425\du}}
\pgfpathlineto{\pgfpoint{-12.539889\du}{16.615547\du}}
\pgfpathlineto{\pgfpoint{-12.566160\du}{16.603442\du}}
\pgfpathlineto{\pgfpoint{-12.594491\du}{16.592110\du}}
\pgfpathlineto{\pgfpoint{-12.624882\du}{16.580262\du}}
\pgfpathlineto{\pgfpoint{-12.656819\du}{16.569960\du}}
\pgfpathlineto{\pgfpoint{-12.690044\du}{16.559915\du}}
\pgfpathlineto{\pgfpoint{-12.724814\du}{16.550386\du}}
\pgfpathlineto{\pgfpoint{-12.760614\du}{16.541629\du}}
\pgfpathlineto{\pgfpoint{-12.798474\du}{16.532872\du}}
\pgfpathlineto{\pgfpoint{-12.837365\du}{16.525403\du}}
\pgfpathlineto{\pgfpoint{-12.876771\du}{16.518449\du}}
\pgfpathlineto{\pgfpoint{-12.917722\du}{16.512010\du}}
\pgfpathlineto{\pgfpoint{-12.959446\du}{16.506344\du}}
\pgfpathlineto{\pgfpoint{-13.001942\du}{16.501450\du}}
\pgfpathlineto{\pgfpoint{-13.045469\du}{16.497845\du}}
\pgfpathlineto{\pgfpoint{-13.088996\du}{16.494239\du}}
\pgfpathlineto{\pgfpoint{-13.133038\du}{16.491663\du}}
\pgfpathlineto{\pgfpoint{-13.177852\du}{16.489345\du}}
\pgfpathlineto{\pgfpoint{-13.222667\du}{16.488315\du}}
\pgfpathlineto{\pgfpoint{-13.267739\du}{16.487800\du}}
\pgfpathlineto{\pgfpoint{-13.267739\du}{16.487800\du}}
\pgfpathlineto{\pgfpoint{-13.312553\du}{16.488315\du}}
\pgfpathlineto{\pgfpoint{-13.357626\du}{16.489345\du}}
\pgfpathlineto{\pgfpoint{-13.402183\du}{16.491663\du}}
\pgfpathlineto{\pgfpoint{-13.446224\du}{16.494239\du}}
\pgfpathlineto{\pgfpoint{-13.490009\du}{16.497845\du}}
\pgfpathlineto{\pgfpoint{-13.533278\du}{16.501450\du}}
\pgfpathlineto{\pgfpoint{-13.575774\du}{16.506344\du}}
\pgfpathlineto{\pgfpoint{-13.617756\du}{16.512010\du}}
\pgfpathlineto{\pgfpoint{-13.658450\du}{16.518449\du}}
\pgfpathlineto{\pgfpoint{-13.697855\du}{16.525403\du}}
\pgfpathlineto{\pgfpoint{-13.737004\du}{16.532872\du}}
\pgfpathlineto{\pgfpoint{-13.774349\du}{16.541629\du}}
\pgfpathlineto{\pgfpoint{-13.810149\du}{16.550386\du}}
\pgfpathlineto{\pgfpoint{-13.845177\du}{16.559915\du}}
\pgfpathlineto{\pgfpoint{-13.878916\du}{16.569960\du}}
\pgfpathlineto{\pgfpoint{-13.910853\du}{16.580262\du}}
\pgfpathlineto{\pgfpoint{-13.940730\du}{16.592110\du}}
\pgfpathlineto{\pgfpoint{-13.969318\du}{16.603442\du}}
\pgfpathlineto{\pgfpoint{-13.995846\du}{16.615547\du}}
\pgfpathlineto{\pgfpoint{-14.020314\du}{16.628425\du}}
\pgfpathlineto{\pgfpoint{-14.042721\du}{16.641303\du}}
\pgfpathlineto{\pgfpoint{-14.063841\du}{16.655211\du}}
\pgfpathlineto{\pgfpoint{-14.082900\du}{16.668603\du}}
\pgfpathlineto{\pgfpoint{-14.099383\du}{16.683027\du}}
\pgfpathlineto{\pgfpoint{-14.113806\du}{16.696934\du}}
\pgfpathlineto{\pgfpoint{-14.126684\du}{16.711358\du}}
\pgfpathlineto{\pgfpoint{-14.136214\du}{16.726038\du}}
\pgfpathlineto{\pgfpoint{-14.144455\du}{16.741234\du}}
\pgfpathlineto{\pgfpoint{-14.150122\du}{16.755657\du}}
\pgfpathlineto{\pgfpoint{-14.153212\du}{16.770853\du}}
\pgfpathlineto{\pgfpoint{-14.154500\du}{16.786306\du}}
\pgfpathlineto{\pgfpoint{-14.154500\du}{16.786306\du}}
\pgfpathlineto{\pgfpoint{-14.153212\du}{16.801244\du}}
\pgfpathlineto{\pgfpoint{-14.150122\du}{16.816440\du}}
\pgfpathlineto{\pgfpoint{-14.144455\du}{16.830863\du}}
\pgfpathlineto{\pgfpoint{-14.136214\du}{16.845801\du}}
\pgfpathlineto{\pgfpoint{-14.126684\du}{16.860739\du}}
\pgfpathlineto{\pgfpoint{-14.113806\du}{16.875162\du}}
\pgfpathlineto{\pgfpoint{-14.099383\du}{16.889585\du}}
\pgfpathlineto{\pgfpoint{-14.082900\du}{16.903751\du}}
\pgfpathlineto{\pgfpoint{-14.063841\du}{16.917401\du}}
\pgfpathlineto{\pgfpoint{-14.042721\du}{16.930794\du}}
\pgfpathlineto{\pgfpoint{-14.020314\du}{16.943929\du}}
\pgfpathlineto{\pgfpoint{-13.995846\du}{16.956035\du}}
\pgfpathlineto{\pgfpoint{-13.969318\du}{16.968655\du}}
\pgfpathlineto{\pgfpoint{-13.940730\du}{16.980245\du}}
\pgfpathlineto{\pgfpoint{-13.910853\du}{16.991577\du}}
\pgfpathlineto{\pgfpoint{-13.878916\du}{17.002137\du}}
\pgfpathlineto{\pgfpoint{-13.845177\du}{17.012181\du}}
\pgfpathlineto{\pgfpoint{-13.810149\du}{17.021711\du}}
\pgfpathlineto{\pgfpoint{-13.774349\du}{17.030468\du}}
\pgfpathlineto{\pgfpoint{-13.737004\du}{17.038967\du}}
\pgfpathlineto{\pgfpoint{-13.697855\du}{17.046951\du}}
\pgfpathlineto{\pgfpoint{-13.658450\du}{17.053390\du}}
\pgfpathlineto{\pgfpoint{-13.617756\du}{17.060087\du}}
\pgfpathlineto{\pgfpoint{-13.575774\du}{17.065495\du}}
\pgfpathlineto{\pgfpoint{-13.533278\du}{17.070389\du}}
\pgfpathlineto{\pgfpoint{-13.490009\du}{17.074767\du}}
\pgfpathlineto{\pgfpoint{-13.446224\du}{17.078373\du}}
\pgfpathlineto{\pgfpoint{-13.402183\du}{17.080434\du}}
\pgfpathlineto{\pgfpoint{-13.357626\du}{17.082494\du}}
\pgfpathlineto{\pgfpoint{-13.312553\du}{17.083782\du}}
\pgfpathlineto{\pgfpoint{-13.267739\du}{17.084039\du}}
\pgfpathlineto{\pgfpoint{-13.267739\du}{17.084039\du}}
\pgfpathlineto{\pgfpoint{-13.222667\du}{17.083782\du}}
\pgfpathlineto{\pgfpoint{-13.177852\du}{17.082494\du}}
\pgfpathlineto{\pgfpoint{-13.133038\du}{17.080434\du}}
\pgfpathlineto{\pgfpoint{-13.088996\du}{17.078373\du}}
\pgfpathlineto{\pgfpoint{-13.045469\du}{17.074767\du}}
\pgfpathlineto{\pgfpoint{-13.001942\du}{17.070389\du}}
\pgfpathlineto{\pgfpoint{-12.959446\du}{17.065495\du}}
\pgfpathlineto{\pgfpoint{-12.917722\du}{17.060087\du}}
\pgfpathlineto{\pgfpoint{-12.876771\du}{17.053390\du}}
\pgfpathlineto{\pgfpoint{-12.837365\du}{17.046951\du}}
\pgfpathlineto{\pgfpoint{-12.798474\du}{17.038967\du}}
\pgfpathlineto{\pgfpoint{-12.760614\du}{17.030468\du}}
\pgfpathlineto{\pgfpoint{-12.724814\du}{17.021711\du}}
\pgfpathlineto{\pgfpoint{-12.690044\du}{17.012181\du}}
\pgfpathlineto{\pgfpoint{-12.656819\du}{17.002137\du}}
\pgfpathlineto{\pgfpoint{-12.624882\du}{16.991577\du}}
\pgfpathlineto{\pgfpoint{-12.594491\du}{16.980245\du}}
\pgfpathlineto{\pgfpoint{-12.566160\du}{16.968655\du}}
\pgfpathlineto{\pgfpoint{-12.539889\du}{16.956035\du}}
\pgfpathlineto{\pgfpoint{-12.514906\du}{16.943929\du}}
\pgfpathlineto{\pgfpoint{-12.492242\du}{16.930794\du}}
\pgfpathlineto{\pgfpoint{-12.471380\du}{16.917401\du}}
\pgfpathlineto{\pgfpoint{-12.452578\du}{16.903751\du}}
\pgfpathlineto{\pgfpoint{-12.435580\du}{16.889585\du}}
\pgfpathlineto{\pgfpoint{-12.421672\du}{16.875162\du}}
\pgfpathlineto{\pgfpoint{-12.409051\du}{16.860739\du}}
\pgfpathlineto{\pgfpoint{-12.398749\du}{16.845801\du}}
\pgfpathlineto{\pgfpoint{-12.391023\du}{16.830863\du}}
\pgfpathlineto{\pgfpoint{-12.385356\du}{16.816440\du}}
\pgfpathlineto{\pgfpoint{-12.382008\du}{16.801244\du}}
\pgfpathlineto{\pgfpoint{-12.380720\du}{16.786306\du}}
\pgfusepath{stroke}
\pgfsetbuttcap
\pgfsetmiterjoin
\pgfsetdash{}{0pt}
\definecolor{dialinecolor}{rgb}{0.000000, 0.000000, 0.000000}
\pgfsetfillcolor{dialinecolor}
\pgfpathmoveto{\pgfpoint{-13.246362\du}{16.721917\du}}
\pgfpathlineto{\pgfpoint{-13.116297\du}{16.764671\du}}
\pgfpathlineto{\pgfpoint{-12.802338\du}{16.636667\du}}
\pgfpathlineto{\pgfpoint{-12.661970\du}{16.679421\du}}
\pgfpathlineto{\pgfpoint{-12.738206\du}{16.573051\du}}
\pgfpathlineto{\pgfpoint{-13.105222\du}{16.573051\du}}
\pgfpathlineto{\pgfpoint{-12.954295\du}{16.605245\du}}
\pgfpathlineto{\pgfpoint{-13.246362\du}{16.721917\du}}
\pgfusepath{fill}
\pgfsetbuttcap
\pgfsetmiterjoin
\pgfsetdash{}{0pt}
\definecolor{dialinecolor}{rgb}{0.000000, 0.000000, 0.000000}
\pgfsetfillcolor{dialinecolor}
\pgfpathmoveto{\pgfpoint{-13.300706\du}{16.839362\du}}
\pgfpathlineto{\pgfpoint{-13.429741\du}{16.796608\du}}
\pgfpathlineto{\pgfpoint{-13.732883\du}{16.924613\du}}
\pgfpathlineto{\pgfpoint{-13.884068\du}{16.881601\du}}
\pgfpathlineto{\pgfpoint{-13.808346\du}{16.998789\du}}
\pgfpathlineto{\pgfpoint{-13.429741\du}{16.998789\du}}
\pgfpathlineto{\pgfpoint{-13.592258\du}{16.956035\du}}
\pgfpathlineto{\pgfpoint{-13.300706\du}{16.839362\du}}
\pgfusepath{fill}
\pgfsetbuttcap
\pgfsetmiterjoin
\pgfsetdash{}{0pt}
\definecolor{dialinecolor}{rgb}{0.000000, 0.000000, 0.000000}
\pgfsetfillcolor{dialinecolor}
\pgfpathmoveto{\pgfpoint{-13.851873\du}{16.605245\du}}
\pgfpathlineto{\pgfpoint{-13.722066\du}{16.562491\du}}
\pgfpathlineto{\pgfpoint{-13.408621\du}{16.679421\du}}
\pgfpathlineto{\pgfpoint{-13.267739\du}{16.647742\du}}
\pgfpathlineto{\pgfpoint{-13.343203\du}{16.753854\du}}
\pgfpathlineto{\pgfpoint{-13.711248\du}{16.753854\du}}
\pgfpathlineto{\pgfpoint{-13.560064\du}{16.721917\du}}
\pgfpathlineto{\pgfpoint{-13.851873\du}{16.605245\du}}
\pgfusepath{fill}
\pgfsetbuttcap
\pgfsetmiterjoin
\pgfsetdash{}{0pt}
\definecolor{dialinecolor}{rgb}{0.000000, 0.000000, 0.000000}
\pgfsetfillcolor{dialinecolor}
\pgfpathmoveto{\pgfpoint{-12.683862\du}{16.966852\du}}
\pgfpathlineto{\pgfpoint{-12.813412\du}{17.009348\du}}
\pgfpathlineto{\pgfpoint{-13.116297\du}{16.881601\du}}
\pgfpathlineto{\pgfpoint{-13.267739\du}{16.924613\du}}
\pgfpathlineto{\pgfpoint{-13.192018\du}{16.817985\du}}
\pgfpathlineto{\pgfpoint{-12.813412\du}{16.817985\du}}
\pgfpathlineto{\pgfpoint{-12.975929\du}{16.849407\du}}
\pgfpathlineto{\pgfpoint{-12.683862\du}{16.966852\du}}
\pgfusepath{fill}
\pgfsetbuttcap
\pgfsetmiterjoin
\pgfsetdash{}{0pt}
\definecolor{dialinecolor}{rgb}{1.000000, 1.000000, 1.000000}
\pgfsetfillcolor{dialinecolor}
\pgfpathmoveto{\pgfpoint{-13.235545\du}{16.732477\du}}
\pgfpathlineto{\pgfpoint{-13.105222\du}{16.774974\du}}
\pgfpathlineto{\pgfpoint{-12.791778\du}{16.647742\du}}
\pgfpathlineto{\pgfpoint{-12.651153\du}{16.690496\du}}
\pgfpathlineto{\pgfpoint{-12.726874\du}{16.583610\du}}
\pgfpathlineto{\pgfpoint{-13.094920\du}{16.583610\du}}
\pgfpathlineto{\pgfpoint{-12.943220\du}{16.615547\du}}
\pgfpathlineto{\pgfpoint{-13.235545\du}{16.732477\du}}
\pgfusepath{fill}
\pgfsetbuttcap
\pgfsetmiterjoin
\pgfsetdash{}{0pt}
\definecolor{dialinecolor}{rgb}{1.000000, 1.000000, 1.000000}
\pgfsetfillcolor{dialinecolor}
\pgfpathmoveto{\pgfpoint{-13.289116\du}{16.849407\du}}
\pgfpathlineto{\pgfpoint{-13.419439\du}{16.807168\du}}
\pgfpathlineto{\pgfpoint{-13.722066\du}{16.934915\du}}
\pgfpathlineto{\pgfpoint{-13.873508\du}{16.892161\du}}
\pgfpathlineto{\pgfpoint{-13.797787\du}{17.009348\du}}
\pgfpathlineto{\pgfpoint{-13.419439\du}{17.009348\du}}
\pgfpathlineto{\pgfpoint{-13.581183\du}{16.966852\du}}
\pgfpathlineto{\pgfpoint{-13.289116\du}{16.849407\du}}
\pgfusepath{fill}
\pgfsetbuttcap
\pgfsetmiterjoin
\pgfsetdash{}{0pt}
\definecolor{dialinecolor}{rgb}{1.000000, 1.000000, 1.000000}
\pgfsetfillcolor{dialinecolor}
\pgfpathmoveto{\pgfpoint{-13.840798\du}{16.615547\du}}
\pgfpathlineto{\pgfpoint{-13.711248\du}{16.573051\du}}
\pgfpathlineto{\pgfpoint{-13.397547\du}{16.690496\du}}
\pgfpathlineto{\pgfpoint{-13.256922\du}{16.658044\du}}
\pgfpathlineto{\pgfpoint{-13.332643\du}{16.764671\du}}
\pgfpathlineto{\pgfpoint{-13.700431\du}{16.764671\du}}
\pgfpathlineto{\pgfpoint{-13.548731\du}{16.732477\du}}
\pgfpathlineto{\pgfpoint{-13.840798\du}{16.615547\du}}
\pgfusepath{fill}
\pgfsetbuttcap
\pgfsetmiterjoin
\pgfsetdash{}{0pt}
\definecolor{dialinecolor}{rgb}{1.000000, 1.000000, 1.000000}
\pgfsetfillcolor{dialinecolor}
\pgfpathmoveto{\pgfpoint{-12.672787\du}{16.977412\du}}
\pgfpathlineto{\pgfpoint{-12.802338\du}{17.020166\du}}
\pgfpathlineto{\pgfpoint{-13.105222\du}{16.892161\du}}
\pgfpathlineto{\pgfpoint{-13.256922\du}{16.934915\du}}
\pgfpathlineto{\pgfpoint{-13.181458\du}{16.828803\du}}
\pgfpathlineto{\pgfpoint{-12.802338\du}{16.828803\du}}
\pgfpathlineto{\pgfpoint{-12.964597\du}{16.860482\du}}
\pgfpathlineto{\pgfpoint{-12.672787\du}{16.977412\du}}
\pgfusepath{fill}
\pgfsetbuttcap
\pgfsetmiterjoin
\pgfsetdash{}{0pt}
\definecolor{dialinecolor}{rgb}{0.666667, 0.901961, 1.000000}
\pgfsetstrokecolor{dialinecolor}
\pgfpathmoveto{\pgfpoint{-14.154500\du}{16.786306\du}}
\pgfpathlineto{\pgfpoint{-14.154500\du}{17.211271\du}}
\pgfusepath{stroke}
\pgfsetbuttcap
\pgfsetmiterjoin
\pgfsetdash{}{0pt}
\definecolor{dialinecolor}{rgb}{0.666667, 0.901961, 1.000000}
\pgfsetstrokecolor{dialinecolor}
\pgfpathmoveto{\pgfpoint{-12.380720\du}{16.786306\du}}
\pgfpathlineto{\pgfpoint{-12.380720\du}{17.211271\du}}
\pgfusepath{stroke}
\pgfsetbuttcap
\pgfsetmiterjoin
\pgfsetdash{}{0pt}
\definecolor{dialinecolor}{rgb}{0.000000, 0.000000, 0.000000}
\pgfsetfillcolor{dialinecolor}
\pgfpathmoveto{\pgfpoint{-13.592258\du}{17.158473\du}}
\pgfpathlineto{\pgfpoint{-13.397547\du}{17.158473\du}}
\pgfpathlineto{\pgfpoint{-13.278556\du}{17.254283\du}}
\pgfpathlineto{\pgfpoint{-13.159824\du}{17.158473\du}}
\pgfpathlineto{\pgfpoint{-12.975929\du}{17.158473\du}}
\pgfpathlineto{\pgfpoint{-12.975929\du}{17.115976\du}}
\pgfpathlineto{\pgfpoint{-12.845864\du}{17.179335\du}}
\pgfpathlineto{\pgfpoint{-12.975929\du}{17.243466\du}}
\pgfpathlineto{\pgfpoint{-12.975929\du}{17.200969\du}}
\pgfpathlineto{\pgfpoint{-13.137931\du}{17.200969\du}}
\pgfpathlineto{\pgfpoint{-13.246362\du}{17.275402\du}}
\pgfpathlineto{\pgfpoint{-13.137931\du}{17.360138\du}}
\pgfpathlineto{\pgfpoint{-12.975929\du}{17.360138\du}}
\pgfpathlineto{\pgfpoint{-12.975929\du}{17.307597\du}}
\pgfpathlineto{\pgfpoint{-12.845864\du}{17.381773\du}}
\pgfpathlineto{\pgfpoint{-12.975929\du}{17.445904\du}}
\pgfpathlineto{\pgfpoint{-12.975929\du}{17.402892\du}}
\pgfpathlineto{\pgfpoint{-13.159824\du}{17.402892\du}}
\pgfpathlineto{\pgfpoint{-13.278556\du}{17.307597\du}}
\pgfpathlineto{\pgfpoint{-13.397547\du}{17.402892\du}}
\pgfpathlineto{\pgfpoint{-13.592258\du}{17.402892\du}}
\pgfpathlineto{\pgfpoint{-13.592258\du}{17.445904\du}}
\pgfpathlineto{\pgfpoint{-13.722066\du}{17.381773\du}}
\pgfpathlineto{\pgfpoint{-13.592258\du}{17.307597\du}}
\pgfpathlineto{\pgfpoint{-13.592258\du}{17.360138\du}}
\pgfpathlineto{\pgfpoint{-13.419439\du}{17.360138\du}}
\pgfpathlineto{\pgfpoint{-13.311008\du}{17.285962\du}}
\pgfpathlineto{\pgfpoint{-13.419439\du}{17.200969\du}}
\pgfpathlineto{\pgfpoint{-13.592258\du}{17.200969\du}}
\pgfpathlineto{\pgfpoint{-13.592258\du}{17.243466\du}}
\pgfpathlineto{\pgfpoint{-13.722066\du}{17.179335\du}}
\pgfpathlineto{\pgfpoint{-13.592258\du}{17.115976\du}}
\pgfpathlineto{\pgfpoint{-13.592258\du}{17.158473\du}}
\pgfusepath{fill}
\pgfsetbuttcap
\pgfsetmiterjoin
\pgfsetdash{}{0pt}
\definecolor{dialinecolor}{rgb}{1.000000, 1.000000, 1.000000}
\pgfsetfillcolor{dialinecolor}
\pgfpathmoveto{\pgfpoint{-13.581183\du}{17.168775\du}}
\pgfpathlineto{\pgfpoint{-13.386729\du}{17.168775\du}}
\pgfpathlineto{\pgfpoint{-13.267739\du}{17.264328\du}}
\pgfpathlineto{\pgfpoint{-13.149006\du}{17.168775\du}}
\pgfpathlineto{\pgfpoint{-12.964597\du}{17.168775\du}}
\pgfpathlineto{\pgfpoint{-12.964597\du}{17.126278\du}}
\pgfpathlineto{\pgfpoint{-12.835305\du}{17.190409\du}}
\pgfpathlineto{\pgfpoint{-12.964597\du}{17.254283\du}}
\pgfpathlineto{\pgfpoint{-12.964597\du}{17.211786\du}}
\pgfpathlineto{\pgfpoint{-13.127114\du}{17.211786\du}}
\pgfpathlineto{\pgfpoint{-13.235545\du}{17.285962\du}}
\pgfpathlineto{\pgfpoint{-13.127114\du}{17.370955\du}}
\pgfpathlineto{\pgfpoint{-12.964597\du}{17.370955\du}}
\pgfpathlineto{\pgfpoint{-12.964597\du}{17.317899\du}}
\pgfpathlineto{\pgfpoint{-12.835305\du}{17.392332\du}}
\pgfpathlineto{\pgfpoint{-12.964597\du}{17.456206\du}}
\pgfpathlineto{\pgfpoint{-12.964597\du}{17.413452\du}}
\pgfpathlineto{\pgfpoint{-13.149006\du}{17.413452\du}}
\pgfpathlineto{\pgfpoint{-13.267739\du}{17.317899\du}}
\pgfpathlineto{\pgfpoint{-13.386729\du}{17.413452\du}}
\pgfpathlineto{\pgfpoint{-13.581183\du}{17.413452\du}}
\pgfpathlineto{\pgfpoint{-13.581183\du}{17.456206\du}}
\pgfpathlineto{\pgfpoint{-13.711248\du}{17.392332\du}}
\pgfpathlineto{\pgfpoint{-13.581183\du}{17.317899\du}}
\pgfpathlineto{\pgfpoint{-13.581183\du}{17.370955\du}}
\pgfpathlineto{\pgfpoint{-13.408621\du}{17.370955\du}}
\pgfpathlineto{\pgfpoint{-13.300706\du}{17.296522\du}}
\pgfpathlineto{\pgfpoint{-13.408621\du}{17.211786\du}}
\pgfpathlineto{\pgfpoint{-13.581183\du}{17.211786\du}}
\pgfpathlineto{\pgfpoint{-13.581183\du}{17.254283\du}}
\pgfpathlineto{\pgfpoint{-13.711248\du}{17.190409\du}}
\pgfpathlineto{\pgfpoint{-13.581183\du}{17.126278\du}}
\pgfpathlineto{\pgfpoint{-13.581183\du}{17.168775\du}}
\pgfusepath{fill}
\pgfsetlinewidth{0.100000\du}
\pgfsetdash{}{0pt}
\pgfsetdash{}{0pt}
\pgfsetbuttcap
{
\definecolor{dialinecolor}{rgb}{0.000000, 0.000000, 0.000000}
\pgfsetfillcolor{dialinecolor}
% was here!!!
\definecolor{dialinecolor}{rgb}{0.000000, 0.000000, 0.000000}
\pgfsetstrokecolor{dialinecolor}
\draw (-16.929428\du,15.060522\du)--(-13.995735\du,16.615373\du);
}
\pgfsetlinewidth{0.100000\du}
\pgfsetdash{}{0pt}
\pgfsetdash{}{0pt}
\pgfsetbuttcap
{
\definecolor{dialinecolor}{rgb}{0.000000, 0.000000, 0.000000}
\pgfsetfillcolor{dialinecolor}
% was here!!!
\definecolor{dialinecolor}{rgb}{0.000000, 0.000000, 0.000000}
\pgfsetstrokecolor{dialinecolor}
\draw (-16.384046\du,3.334167\du)--(-16.274763\du,5.796111\du);
}
\pgfsetlinewidth{0.100000\du}
\pgfsetdash{}{0pt}
\pgfsetdash{}{0pt}
\pgfsetbuttcap
{
\definecolor{dialinecolor}{rgb}{0.000000, 0.000000, 0.000000}
\pgfsetfillcolor{dialinecolor}
% was here!!!
\definecolor{dialinecolor}{rgb}{0.000000, 0.000000, 0.000000}
\pgfsetstrokecolor{dialinecolor}
\draw (-15.511355\du,3.090648\du)--(-14.183845\du,5.621522\du);
}
\pgfsetlinewidth{0.100000\du}
\pgfsetdash{}{0pt}
\pgfsetdash{}{0pt}
\pgfsetbuttcap
{
\definecolor{dialinecolor}{rgb}{0.000000, 0.000000, 0.000000}
\pgfsetfillcolor{dialinecolor}
% was here!!!
\definecolor{dialinecolor}{rgb}{0.000000, 0.000000, 0.000000}
\pgfsetstrokecolor{dialinecolor}
\draw (-17.302500\du,3.090648\du)--(-18.984116\du,5.812300\du);
}
\pgfsetlinewidth{0.100000\du}
\pgfsetdash{}{0pt}
\pgfsetdash{}{0pt}
\pgfsetbuttcap
{
\definecolor{dialinecolor}{rgb}{0.000000, 0.000000, 0.000000}
\pgfsetfillcolor{dialinecolor}
% was here!!!
\definecolor{dialinecolor}{rgb}{0.000000, 0.000000, 0.000000}
\pgfsetstrokecolor{dialinecolor}
\draw (-12.370418\du,17.277720\du)--(-11.089864\du,19.336607\du);
}
\pgfsetlinewidth{0.100000\du}
\pgfsetdash{}{0pt}
\pgfsetdash{}{0pt}
\pgfsetbuttcap
{
\definecolor{dialinecolor}{rgb}{0.000000, 0.000000, 0.000000}
\pgfsetfillcolor{dialinecolor}
% was here!!!
\definecolor{dialinecolor}{rgb}{0.000000, 0.000000, 0.000000}
\pgfsetstrokecolor{dialinecolor}
\draw (-22.661828\du,18.077177\du)--(-23.889041\du,20.050920\du);
}
\pgfsetlinewidth{0.100000\du}
\pgfsetdash{}{0pt}
\pgfsetdash{}{0pt}
\pgfsetbuttcap
{
\definecolor{dialinecolor}{rgb}{0.000000, 0.000000, 0.000000}
\pgfsetfillcolor{dialinecolor}
% was here!!!
\definecolor{dialinecolor}{rgb}{0.000000, 0.000000, 0.000000}
\pgfsetstrokecolor{dialinecolor}
\draw (-13.258487\du,17.520212\du)--(-13.241800\du,19.620986\du);
}
\pgfsetlinewidth{0.100000\du}
\pgfsetdash{}{0pt}
\pgfsetdash{}{0pt}
\pgfsetbuttcap
{
\definecolor{dialinecolor}{rgb}{0.000000, 0.000000, 0.000000}
\pgfsetfillcolor{dialinecolor}
% was here!!!
\definecolor{dialinecolor}{rgb}{0.000000, 0.000000, 0.000000}
\pgfsetstrokecolor{dialinecolor}
\draw (-13.680237\du,17.484525\du)--(-15.734504\du,19.848350\du);
}
\pgfsetlinewidth{0.100000\du}
\pgfsetdash{}{0pt}
\pgfsetdash{}{0pt}
\pgfsetbuttcap
{
\definecolor{dialinecolor}{rgb}{0.000000, 0.000000, 0.000000}
\pgfsetfillcolor{dialinecolor}
% was here!!!
\definecolor{dialinecolor}{rgb}{0.000000, 0.000000, 0.000000}
\pgfsetstrokecolor{dialinecolor}
\draw (-21.460118\du,17.853120\du)--(-20.745356\du,19.841620\du);
}
\pgfsetlinewidth{0.100000\du}
\pgfsetdash{}{0pt}
\pgfsetdash{}{0pt}
\pgfsetbuttcap
{
\definecolor{dialinecolor}{rgb}{0.000000, 0.000000, 0.000000}
\pgfsetfillcolor{dialinecolor}
% was here!!!
\definecolor{dialinecolor}{rgb}{0.000000, 0.000000, 0.000000}
\pgfsetstrokecolor{dialinecolor}
\draw (-7.602427\du,3.668293\du)--(-6.854309\du,5.732470\du);
}
\pgfsetlinewidth{0.100000\du}
\pgfsetdash{}{0pt}
\pgfsetdash{}{0pt}
\pgfsetbuttcap
{
\definecolor{dialinecolor}{rgb}{0.000000, 0.000000, 0.000000}
\pgfsetfillcolor{dialinecolor}
% was here!!!
\definecolor{dialinecolor}{rgb}{0.000000, 0.000000, 0.000000}
\pgfsetstrokecolor{dialinecolor}
\draw (-8.184939\du,3.654432\du)--(-9.443825\du,5.632398\du);
}
\pgfsetlinewidth{0.100000\du}
\pgfsetdash{}{0pt}
\pgfsetdash{}{0pt}
\pgfsetbuttcap
{
\definecolor{dialinecolor}{rgb}{0.000000, 0.000000, 0.000000}
\pgfsetfillcolor{dialinecolor}
% was here!!!
\definecolor{dialinecolor}{rgb}{0.000000, 0.000000, 0.000000}
\pgfsetstrokecolor{dialinecolor}
\draw (-1.792482\du,3.286735\du)--(-2.748907\du,5.783785\du);
}
\pgfsetlinewidth{0.100000\du}
\pgfsetdash{}{0pt}
\pgfsetdash{}{0pt}
\pgfsetbuttcap
{
\definecolor{dialinecolor}{rgb}{0.000000, 0.000000, 0.000000}
\pgfsetfillcolor{dialinecolor}
% was here!!!
\definecolor{dialinecolor}{rgb}{0.000000, 0.000000, 0.000000}
\pgfsetstrokecolor{dialinecolor}
\draw (-1.566983\du,3.298100\du)--(-0.197944\du,5.821100\du);
}
\pgfsetlinewidth{0.100000\du}
\pgfsetdash{}{0pt}
\pgfsetdash{}{0pt}
\pgfsetbuttcap
{
\definecolor{dialinecolor}{rgb}{0.000000, 0.000000, 0.000000}
\pgfsetfillcolor{dialinecolor}
% was here!!!
\definecolor{dialinecolor}{rgb}{0.000000, 0.000000, 0.000000}
\pgfsetstrokecolor{dialinecolor}
\draw (-0.930494\du,3.224493\du)--(1.895248\du,5.561900\du);
}
\pgfsetlinewidth{0.100000\du}
\pgfsetdash{}{0pt}
\pgfsetdash{}{0pt}
\pgfsetbuttcap
{
\definecolor{dialinecolor}{rgb}{0.000000, 0.000000, 0.000000}
\pgfsetfillcolor{dialinecolor}
% was here!!!
\pgfsetarrowsend{stealth}
\definecolor{dialinecolor}{rgb}{0.000000, 0.000000, 0.000000}
\pgfsetstrokecolor{dialinecolor}
\draw (-18.850000\du,-2.590000\du)--(-18.850000\du,-0.940000\du);
}
\pgfsetlinewidth{0.100000\du}
\pgfsetdash{}{0pt}
\pgfsetdash{}{0pt}
\pgfsetbuttcap
{
\definecolor{dialinecolor}{rgb}{0.000000, 0.000000, 0.000000}
\pgfsetfillcolor{dialinecolor}
% was here!!!
\pgfsetarrowsend{stealth}
\definecolor{dialinecolor}{rgb}{0.000000, 0.000000, 0.000000}
\pgfsetstrokecolor{dialinecolor}
\draw (-18.293200\du,-2.630000\du)--(-18.293200\du,-0.980000\du);
}
\pgfsetlinewidth{0.100000\du}
\pgfsetdash{}{0pt}
\pgfsetdash{}{0pt}
\pgfsetbuttcap
{
\definecolor{dialinecolor}{rgb}{0.000000, 0.000000, 0.000000}
\pgfsetfillcolor{dialinecolor}
% was here!!!
\pgfsetarrowsend{stealth}
\definecolor{dialinecolor}{rgb}{0.000000, 0.000000, 0.000000}
\pgfsetstrokecolor{dialinecolor}
\draw (-17.798200\du,-2.620000\du)--(-17.798200\du,-0.970000\du);
}
\pgfsetlinewidth{0.100000\du}
\pgfsetdash{}{0pt}
\pgfsetdash{}{0pt}
\pgfsetbuttcap
{
\definecolor{dialinecolor}{rgb}{0.000000, 0.000000, 0.000000}
\pgfsetfillcolor{dialinecolor}
% was here!!!
\pgfsetarrowsend{stealth}
\definecolor{dialinecolor}{rgb}{0.000000, 0.000000, 0.000000}
\pgfsetstrokecolor{dialinecolor}
\draw (-3.903200\du,-2.860000\du)--(-3.903200\du,-1.210000\du);
}
\pgfsetlinewidth{0.100000\du}
\pgfsetdash{}{0pt}
\pgfsetdash{}{0pt}
\pgfsetbuttcap
{
\definecolor{dialinecolor}{rgb}{0.000000, 0.000000, 0.000000}
\pgfsetfillcolor{dialinecolor}
% was here!!!
\pgfsetarrowsend{stealth}
\definecolor{dialinecolor}{rgb}{0.000000, 0.000000, 0.000000}
\pgfsetstrokecolor{dialinecolor}
\draw (-17.258200\du,11.750000\du)--(-17.258200\du,13.400000\du);
}
\pgfsetlinewidth{0.100000\du}
\pgfsetdash{}{0pt}
\pgfsetdash{}{0pt}
\pgfsetbuttcap
{
\definecolor{dialinecolor}{rgb}{0.000000, 0.000000, 0.000000}
\pgfsetfillcolor{dialinecolor}
% was here!!!
\pgfsetarrowsend{stealth}
\definecolor{dialinecolor}{rgb}{0.000000, 0.000000, 0.000000}
\pgfsetstrokecolor{dialinecolor}
\draw (-17.800000\du,0.302900\du)--(-16.700000\du,1.660000\du);
}
\pgfsetlinewidth{0.100000\du}
\pgfsetdash{}{0pt}
\pgfsetdash{}{0pt}
\pgfsetbuttcap
{
\definecolor{dialinecolor}{rgb}{0.000000, 0.000000, 0.000000}
\pgfsetfillcolor{dialinecolor}
% was here!!!
\pgfsetarrowsend{stealth}
\definecolor{dialinecolor}{rgb}{0.000000, 0.000000, 0.000000}
\pgfsetstrokecolor{dialinecolor}
\draw (-14.900000\du,3.410000\du)--(-14.050000\du,4.910000\du);
}
\pgfsetlinewidth{0.100000\du}
\pgfsetdash{}{0pt}
\pgfsetdash{}{0pt}
\pgfsetbuttcap
{
\definecolor{dialinecolor}{rgb}{0.000000, 0.000000, 0.000000}
\pgfsetfillcolor{dialinecolor}
% was here!!!
\pgfsetarrowsend{stealth}
\definecolor{dialinecolor}{rgb}{0.000000, 0.000000, 0.000000}
\pgfsetstrokecolor{dialinecolor}
\draw (-18.000000\du,3.460000\du)--(-18.723200\du,4.730000\du);
}
\pgfsetlinewidth{0.100000\du}
\pgfsetdash{}{0pt}
\pgfsetdash{}{0pt}
\pgfsetbuttcap
{
\definecolor{dialinecolor}{rgb}{0.000000, 0.000000, 0.000000}
\pgfsetfillcolor{dialinecolor}
% was here!!!
\pgfsetarrowsend{stealth}
\definecolor{dialinecolor}{rgb}{0.000000, 0.000000, 0.000000}
\pgfsetstrokecolor{dialinecolor}
\draw (-20.800000\du,0.860000\du)--(-22.278200\du,1.890000\du);
}
\pgfsetlinewidth{0.100000\du}
\pgfsetdash{}{0pt}
\pgfsetdash{}{0pt}
\pgfsetbuttcap
{
\definecolor{dialinecolor}{rgb}{0.000000, 0.000000, 0.000000}
\pgfsetfillcolor{dialinecolor}
% was here!!!
\pgfsetarrowsend{stealth}
\definecolor{dialinecolor}{rgb}{0.000000, 0.000000, 0.000000}
\pgfsetstrokecolor{dialinecolor}
\draw (-21.550000\du,3.660000\du)--(-21.133200\du,5.150000\du);
}
\pgfsetlinewidth{0.100000\du}
\pgfsetdash{}{0pt}
\pgfsetdash{}{0pt}
\pgfsetbuttcap
{
\definecolor{dialinecolor}{rgb}{0.000000, 0.000000, 0.000000}
\pgfsetfillcolor{dialinecolor}
% was here!!!
\pgfsetarrowsend{stealth}
\definecolor{dialinecolor}{rgb}{0.000000, 0.000000, 0.000000}
\pgfsetstrokecolor{dialinecolor}
\draw (-3.039700\du,0.350300\du)--(-1.939700\du,1.707400\du);
}
\pgfsetlinewidth{0.100000\du}
\pgfsetdash{}{0pt}
\pgfsetdash{}{0pt}
\pgfsetbuttcap
{
\definecolor{dialinecolor}{rgb}{0.000000, 0.000000, 0.000000}
\pgfsetfillcolor{dialinecolor}
% was here!!!
\pgfsetarrowsend{stealth}
\definecolor{dialinecolor}{rgb}{0.000000, 0.000000, 0.000000}
\pgfsetstrokecolor{dialinecolor}
\draw (0.405300\du,3.310300\du)--(1.505300\du,4.667400\du);
}
\pgfsetlinewidth{0.100000\du}
\pgfsetdash{}{0pt}
\pgfsetdash{}{0pt}
\pgfsetbuttcap
{
\definecolor{dialinecolor}{rgb}{0.000000, 0.000000, 0.000000}
\pgfsetfillcolor{dialinecolor}
% was here!!!
\pgfsetarrowsend{stealth}
\definecolor{dialinecolor}{rgb}{0.000000, 0.000000, 0.000000}
\pgfsetstrokecolor{dialinecolor}
\draw (-5.885100\du,0.639600\du)--(-7.150000\du,1.710000\du);
}
\pgfsetlinewidth{0.100000\du}
\pgfsetdash{}{0pt}
\pgfsetdash{}{0pt}
\pgfsetbuttcap
{
\definecolor{dialinecolor}{rgb}{0.000000, 0.000000, 0.000000}
\pgfsetfillcolor{dialinecolor}
% was here!!!
\pgfsetarrowsend{stealth}
\definecolor{dialinecolor}{rgb}{0.000000, 0.000000, 0.000000}
\pgfsetstrokecolor{dialinecolor}
\draw (-2.500000\du,3.760000\du)--(-3.050000\du,5.160000\du);
}
\pgfsetlinewidth{0.100000\du}
\pgfsetdash{}{0pt}
\pgfsetdash{}{0pt}
\pgfsetbuttcap
{
\definecolor{dialinecolor}{rgb}{0.000000, 0.000000, 0.000000}
\pgfsetfillcolor{dialinecolor}
% was here!!!
\pgfsetarrowsend{stealth}
\definecolor{dialinecolor}{rgb}{0.000000, 0.000000, 0.000000}
\pgfsetstrokecolor{dialinecolor}
\draw (-7.100000\du,3.810000\du)--(-6.600000\du,5.010000\du);
}
% setfont left to latex
\definecolor{dialinecolor}{rgb}{0.000000, 0.000000, 0.000000}
\pgfsetstrokecolor{dialinecolor}
\node[anchor=west] at (-18.850000\du,8.727000\du){};
% setfont left to latex
\definecolor{dialinecolor}{rgb}{0.000000, 0.000000, 0.000000}
\pgfsetstrokecolor{dialinecolor}
\node[anchor=west] at (-23.350000\du,8.727000\du){Unicasting- per client};
% setfont left to latex
\definecolor{dialinecolor}{rgb}{0.000000, 0.000000, 0.000000}
\pgfsetstrokecolor{dialinecolor}
\node[anchor=west] at (-4.300000\du,8.727000\du){};
% setfont left to latex
\definecolor{dialinecolor}{rgb}{0.000000, 0.000000, 0.000000}
\pgfsetstrokecolor{dialinecolor}
\node[anchor=west] at (-8.100000\du,9.027000\du){Multicasting- to members};
\pgfsetlinewidth{0.100000\du}
\pgfsetdash{}{0pt}
\pgfsetdash{}{0pt}
\pgfsetbuttcap
{
\definecolor{dialinecolor}{rgb}{0.000000, 0.000000, 0.000000}
\pgfsetfillcolor{dialinecolor}
% was here!!!
\pgfsetarrowsend{stealth}
\definecolor{dialinecolor}{rgb}{0.000000, 0.000000, 0.000000}
\pgfsetstrokecolor{dialinecolor}
\draw (-16.000000\du,14.952900\du)--(-14.650000\du,15.852900\du);
}
\pgfsetlinewidth{0.100000\du}
\pgfsetdash{}{0pt}
\pgfsetdash{}{0pt}
\pgfsetbuttcap
{
\definecolor{dialinecolor}{rgb}{0.000000, 0.000000, 0.000000}
\pgfsetfillcolor{dialinecolor}
% was here!!!
\pgfsetarrowsend{stealth}
\definecolor{dialinecolor}{rgb}{0.000000, 0.000000, 0.000000}
\pgfsetstrokecolor{dialinecolor}
\draw (-11.800000\du,17.402900\du)--(-10.900000\du,18.652900\du);
}
\pgfsetlinewidth{0.100000\du}
\pgfsetdash{}{0pt}
\pgfsetdash{}{0pt}
\pgfsetbuttcap
{
\definecolor{dialinecolor}{rgb}{0.000000, 0.000000, 0.000000}
\pgfsetfillcolor{dialinecolor}
% was here!!!
\pgfsetarrowsend{stealth}
\definecolor{dialinecolor}{rgb}{0.000000, 0.000000, 0.000000}
\pgfsetstrokecolor{dialinecolor}
\draw (-12.800000\du,17.802900\du)--(-12.800000\du,19.152900\du);
}
\pgfsetlinewidth{0.100000\du}
\pgfsetdash{}{0pt}
\pgfsetdash{}{0pt}
\pgfsetbuttcap
{
\definecolor{dialinecolor}{rgb}{0.000000, 0.000000, 0.000000}
\pgfsetfillcolor{dialinecolor}
% was here!!!
\pgfsetarrowsend{stealth}
\definecolor{dialinecolor}{rgb}{0.000000, 0.000000, 0.000000}
\pgfsetstrokecolor{dialinecolor}
\draw (-14.450000\du,17.852900\du)--(-15.645700\du,19.102200\du);
}
\pgfsetlinewidth{0.100000\du}
\pgfsetdash{}{0pt}
\pgfsetdash{}{0pt}
\pgfsetbuttcap
{
\definecolor{dialinecolor}{rgb}{0.000000, 0.000000, 0.000000}
\pgfsetfillcolor{dialinecolor}
% was here!!!
\pgfsetarrowsend{stealth}
\definecolor{dialinecolor}{rgb}{0.000000, 0.000000, 0.000000}
\pgfsetstrokecolor{dialinecolor}
\draw (-19.550000\du,15.452900\du)--(-20.900000\du,16.352900\du);
}
\pgfsetlinewidth{0.100000\du}
\pgfsetdash{}{0pt}
\pgfsetdash{}{0pt}
\pgfsetbuttcap
{
\definecolor{dialinecolor}{rgb}{0.000000, 0.000000, 0.000000}
\pgfsetfillcolor{dialinecolor}
% was here!!!
\pgfsetarrowsend{stealth}
\definecolor{dialinecolor}{rgb}{0.000000, 0.000000, 0.000000}
\pgfsetstrokecolor{dialinecolor}
\draw (-21.000000\du,17.902900\du)--(-20.450000\du,19.302900\du);
}
\pgfsetlinewidth{0.100000\du}
\pgfsetdash{}{0pt}
\pgfsetdash{}{0pt}
\pgfsetbuttcap
{
\definecolor{dialinecolor}{rgb}{0.000000, 0.000000, 0.000000}
\pgfsetfillcolor{dialinecolor}
% was here!!!
\pgfsetarrowsend{stealth}
\definecolor{dialinecolor}{rgb}{0.000000, 0.000000, 0.000000}
\pgfsetstrokecolor{dialinecolor}
\draw (-23.250000\du,18.302900\du)--(-24.100000\du,19.502900\du);
}
% setfont left to latex
\definecolor{dialinecolor}{rgb}{0.000000, 0.000000, 0.000000}
\pgfsetstrokecolor{dialinecolor}
\node[anchor=west] at (-19.244999\du,21.565726\du){};
% setfont left to latex
\definecolor{dialinecolor}{rgb}{0.000000, 0.000000, 0.000000}
\pgfsetstrokecolor{dialinecolor}
\node[anchor=west] at (-19.244999\du,22.365726\du){};
% setfont left to latex
\definecolor{dialinecolor}{rgb}{0.000000, 0.000000, 0.000000}
\pgfsetstrokecolor{dialinecolor}
\node[anchor=west] at (-21.200000\du,22.352900\du){};
% setfont left to latex
\definecolor{dialinecolor}{rgb}{0.000000, 0.000000, 0.000000}
\pgfsetstrokecolor{dialinecolor}
\node[anchor=west] at (-20.950000\du,22.502900\du){};
% setfont left to latex
\definecolor{dialinecolor}{rgb}{0.000000, 0.000000, 0.000000}
\pgfsetstrokecolor{dialinecolor}
\node[anchor=west] at (-21.100000\du,23.352900\du){Broadcasting-to all};
\pgfsetlinewidth{0.100000\du}
\pgfsetdash{}{0pt}
\pgfsetdash{}{0pt}
\pgfsetbuttcap
\pgfsetmiterjoin
\pgfsetlinewidth{0.001000\du}
\pgfsetbuttcap
\pgfsetmiterjoin
\pgfsetdash{}{0pt}
\definecolor{dialinecolor}{rgb}{0.717647, 0.717647, 0.615686}
\pgfsetfillcolor{dialinecolor}
\pgfpathmoveto{\pgfpoint{-5.753800\du}{15.289588\du}}
\pgfpathlineto{\pgfpoint{-4.346388\du}{15.289588\du}}
\pgfpathlineto{\pgfpoint{-4.346388\du}{15.549694\du}}
\pgfpathlineto{\pgfpoint{-5.753800\du}{15.549694\du}}
\pgfpathlineto{\pgfpoint{-5.753800\du}{15.289588\du}}
\pgfusepath{fill}
\pgfsetbuttcap
\pgfsetmiterjoin
\pgfsetdash{}{0pt}
\definecolor{dialinecolor}{rgb}{0.286275, 0.286275, 0.211765}
\pgfsetstrokecolor{dialinecolor}
\pgfpathmoveto{\pgfpoint{-5.753800\du}{15.289588\du}}
\pgfpathlineto{\pgfpoint{-4.346388\du}{15.289588\du}}
\pgfpathlineto{\pgfpoint{-4.346388\du}{15.549694\du}}
\pgfpathlineto{\pgfpoint{-5.753800\du}{15.549694\du}}
\pgfpathlineto{\pgfpoint{-5.753800\du}{15.289588\du}}
\pgfusepath{stroke}
\pgfsetbuttcap
\pgfsetmiterjoin
\pgfsetdash{}{0pt}
\definecolor{dialinecolor}{rgb}{0.788235, 0.788235, 0.713726}
\pgfsetfillcolor{dialinecolor}
\pgfpathmoveto{\pgfpoint{-5.753800\du}{15.289588\du}}
\pgfpathlineto{\pgfpoint{-5.604566\du}{15.148115\du}}
\pgfpathlineto{\pgfpoint{-4.197154\du}{15.148115\du}}
\pgfpathlineto{\pgfpoint{-4.346388\du}{15.289588\du}}
\pgfpathlineto{\pgfpoint{-5.753800\du}{15.289588\du}}
\pgfusepath{fill}
\pgfsetbuttcap
\pgfsetmiterjoin
\pgfsetdash{}{0pt}
\definecolor{dialinecolor}{rgb}{0.286275, 0.286275, 0.211765}
\pgfsetstrokecolor{dialinecolor}
\pgfpathmoveto{\pgfpoint{-5.753800\du}{15.289588\du}}
\pgfpathlineto{\pgfpoint{-5.604566\du}{15.148115\du}}
\pgfpathlineto{\pgfpoint{-4.197154\du}{15.148115\du}}
\pgfpathlineto{\pgfpoint{-4.346388\du}{15.289588\du}}
\pgfpathlineto{\pgfpoint{-5.753800\du}{15.289588\du}}
\pgfusepath{stroke}
\pgfsetlinewidth{0.106000\du}
\pgfsetbuttcap
\pgfsetmiterjoin
\pgfsetdash{}{0pt}
\definecolor{dialinecolor}{rgb}{0.000000, 0.000000, 0.000000}
\pgfsetstrokecolor{dialinecolor}
\pgfpathmoveto{\pgfpoint{-4.425108\du}{15.407999\du}}
\pgfpathlineto{\pgfpoint{-4.762824\du}{15.407999\du}}
\pgfusepath{stroke}
\pgfsetlinewidth{0.001000\du}
\pgfsetbuttcap
\pgfsetmiterjoin
\pgfsetdash{}{0pt}
\definecolor{dialinecolor}{rgb}{0.478431, 0.478431, 0.352941}
\pgfsetfillcolor{dialinecolor}
\pgfpathmoveto{\pgfpoint{-4.346388\du}{15.549694\du}}
\pgfpathlineto{\pgfpoint{-4.197154\du}{15.400016\du}}
\pgfpathlineto{\pgfpoint{-4.197154\du}{15.148115\du}}
\pgfpathlineto{\pgfpoint{-4.346388\du}{15.289588\du}}
\pgfpathlineto{\pgfpoint{-4.346388\du}{15.549694\du}}
\pgfusepath{fill}
\pgfsetbuttcap
\pgfsetmiterjoin
\pgfsetdash{}{0pt}
\definecolor{dialinecolor}{rgb}{0.286275, 0.286275, 0.211765}
\pgfsetstrokecolor{dialinecolor}
\pgfpathmoveto{\pgfpoint{-4.346388\du}{15.549694\du}}
\pgfpathlineto{\pgfpoint{-4.197154\du}{15.400016\du}}
\pgfpathlineto{\pgfpoint{-4.197154\du}{15.148115\du}}
\pgfpathlineto{\pgfpoint{-4.346388\du}{15.289588\du}}
\pgfpathlineto{\pgfpoint{-4.346388\du}{15.549694\du}}
\pgfusepath{stroke}
\pgfsetbuttcap
\pgfsetmiterjoin
\pgfsetdash{}{0pt}
\definecolor{dialinecolor}{rgb}{0.788235, 0.788235, 0.713726}
\pgfsetfillcolor{dialinecolor}
\pgfpathmoveto{\pgfpoint{-5.745817\du}{15.714672\du}}
\pgfpathlineto{\pgfpoint{-5.588822\du}{15.517984\du}}
\pgfpathlineto{\pgfpoint{-4.503605\du}{15.517984\du}}
\pgfpathlineto{\pgfpoint{-4.660600\du}{15.714672\du}}
\pgfpathlineto{\pgfpoint{-5.745817\du}{15.714672\du}}
\pgfusepath{fill}
\pgfsetbuttcap
\pgfsetmiterjoin
\pgfsetdash{}{0pt}
\definecolor{dialinecolor}{rgb}{0.286275, 0.286275, 0.211765}
\pgfsetstrokecolor{dialinecolor}
\pgfpathmoveto{\pgfpoint{-5.745817\du}{15.714672\du}}
\pgfpathlineto{\pgfpoint{-5.588822\du}{15.517984\du}}
\pgfpathlineto{\pgfpoint{-4.503605\du}{15.517984\du}}
\pgfpathlineto{\pgfpoint{-4.660600\du}{15.714672\du}}
\pgfpathlineto{\pgfpoint{-5.745817\du}{15.714672\du}}
\pgfusepath{stroke}
\pgfsetbuttcap
\pgfsetmiterjoin
\pgfsetdash{}{0pt}
\definecolor{dialinecolor}{rgb}{0.478431, 0.478431, 0.352941}
\pgfsetfillcolor{dialinecolor}
\pgfpathmoveto{\pgfpoint{-4.660600\du}{15.753920\du}}
\pgfpathlineto{\pgfpoint{-4.503605\du}{15.588943\du}}
\pgfpathlineto{\pgfpoint{-4.503605\du}{15.517984\du}}
\pgfpathlineto{\pgfpoint{-4.660600\du}{15.714672\du}}
\pgfpathlineto{\pgfpoint{-4.660600\du}{15.753920\du}}
\pgfusepath{fill}
\pgfsetbuttcap
\pgfsetmiterjoin
\pgfsetdash{}{0pt}
\definecolor{dialinecolor}{rgb}{0.286275, 0.286275, 0.211765}
\pgfsetstrokecolor{dialinecolor}
\pgfpathmoveto{\pgfpoint{-4.660600\du}{15.753920\du}}
\pgfpathlineto{\pgfpoint{-4.503605\du}{15.588943\du}}
\pgfpathlineto{\pgfpoint{-4.503605\du}{15.517984\du}}
\pgfpathlineto{\pgfpoint{-4.660600\du}{15.714672\du}}
\pgfpathlineto{\pgfpoint{-4.660600\du}{15.753920\du}}
\pgfusepath{stroke}
\pgfsetbuttcap
\pgfsetmiterjoin
\pgfsetdash{}{0pt}
\definecolor{dialinecolor}{rgb}{0.717647, 0.717647, 0.615686}
\pgfsetfillcolor{dialinecolor}
\pgfpathmoveto{\pgfpoint{-5.745817\du}{15.714672\du}}
\pgfpathlineto{\pgfpoint{-4.660600\du}{15.714672\du}}
\pgfpathlineto{\pgfpoint{-4.660600\du}{15.753920\du}}
\pgfpathlineto{\pgfpoint{-5.745817\du}{15.753920\du}}
\pgfpathlineto{\pgfpoint{-5.745817\du}{15.714672\du}}
\pgfusepath{fill}
\pgfsetbuttcap
\pgfsetmiterjoin
\pgfsetdash{}{0pt}
\definecolor{dialinecolor}{rgb}{0.286275, 0.286275, 0.211765}
\pgfsetstrokecolor{dialinecolor}
\pgfpathmoveto{\pgfpoint{-5.745817\du}{15.714672\du}}
\pgfpathlineto{\pgfpoint{-4.660600\du}{15.714672\du}}
\pgfpathlineto{\pgfpoint{-4.660600\du}{15.753920\du}}
\pgfpathlineto{\pgfpoint{-5.745817\du}{15.753920\du}}
\pgfpathlineto{\pgfpoint{-5.745817\du}{15.714672\du}}
\pgfusepath{stroke}
\pgfsetbuttcap
\pgfsetmiterjoin
\pgfsetdash{}{0pt}
\definecolor{dialinecolor}{rgb}{0.000000, 0.000000, 0.000000}
\pgfsetfillcolor{dialinecolor}
\pgfpathmoveto{\pgfpoint{-5.541591\du}{15.258322\du}}
\pgfpathlineto{\pgfpoint{-5.423401\du}{15.148115\du}}
\pgfpathlineto{\pgfpoint{-4.425108\du}{15.148115\du}}
\pgfpathlineto{\pgfpoint{-4.534871\du}{15.258322\du}}
\pgfpathlineto{\pgfpoint{-5.541591\du}{15.258322\du}}
\pgfusepath{fill}
\pgfsetbuttcap
\pgfsetmiterjoin
\pgfsetdash{}{0pt}
\definecolor{dialinecolor}{rgb}{0.000000, 0.000000, 0.000000}
\pgfsetstrokecolor{dialinecolor}
\pgfpathmoveto{\pgfpoint{-5.541591\du}{15.258322\du}}
\pgfpathlineto{\pgfpoint{-5.423401\du}{15.148115\du}}
\pgfpathlineto{\pgfpoint{-4.425108\du}{15.148115\du}}
\pgfpathlineto{\pgfpoint{-4.534871\du}{15.258322\du}}
\pgfpathlineto{\pgfpoint{-5.541591\du}{15.258322\du}}
\pgfusepath{stroke}
\pgfsetbuttcap
\pgfsetmiterjoin
\pgfsetdash{}{0pt}
\definecolor{dialinecolor}{rgb}{0.788235, 0.788235, 0.713726}
\pgfsetfillcolor{dialinecolor}
\pgfpathmoveto{\pgfpoint{-5.549573\du}{14.447624\du}}
\pgfpathlineto{\pgfpoint{-5.439145\du}{14.345400\du}}
\pgfpathlineto{\pgfpoint{-4.440408\du}{14.345400\du}}
\pgfpathlineto{\pgfpoint{-4.550837\du}{14.447624\du}}
\pgfpathlineto{\pgfpoint{-5.549573\du}{14.447624\du}}
\pgfusepath{fill}
\pgfsetbuttcap
\pgfsetmiterjoin
\pgfsetdash{}{0pt}
\definecolor{dialinecolor}{rgb}{0.286275, 0.286275, 0.211765}
\pgfsetstrokecolor{dialinecolor}
\pgfpathmoveto{\pgfpoint{-5.549573\du}{14.447624\du}}
\pgfpathlineto{\pgfpoint{-5.439145\du}{14.345400\du}}
\pgfpathlineto{\pgfpoint{-4.440408\du}{14.345400\du}}
\pgfpathlineto{\pgfpoint{-4.550837\du}{14.447624\du}}
\pgfpathlineto{\pgfpoint{-5.549573\du}{14.447624\du}}
\pgfusepath{stroke}
\pgfsetbuttcap
\pgfsetmiterjoin
\pgfsetdash{}{0pt}
\definecolor{dialinecolor}{rgb}{0.717647, 0.717647, 0.615686}
\pgfsetfillcolor{dialinecolor}
\pgfpathmoveto{\pgfpoint{-5.549573\du}{14.447624\du}}
\pgfpathlineto{\pgfpoint{-4.542854\du}{14.447624\du}}
\pgfpathlineto{\pgfpoint{-4.542854\du}{15.242356\du}}
\pgfpathlineto{\pgfpoint{-5.549573\du}{15.242356\du}}
\pgfpathlineto{\pgfpoint{-5.549573\du}{14.447624\du}}
\pgfusepath{fill}
\pgfsetbuttcap
\pgfsetmiterjoin
\pgfsetdash{}{0pt}
\definecolor{dialinecolor}{rgb}{0.286275, 0.286275, 0.211765}
\pgfsetstrokecolor{dialinecolor}
\pgfpathmoveto{\pgfpoint{-5.549573\du}{14.447624\du}}
\pgfpathlineto{\pgfpoint{-4.543297\du}{14.447624\du}}
\pgfpathlineto{\pgfpoint{-4.543297\du}{15.242134\du}}
\pgfpathlineto{\pgfpoint{-5.549573\du}{15.242134\du}}
\pgfpathlineto{\pgfpoint{-5.549573\du}{14.447624\du}}
\pgfusepath{stroke}
\pgfsetbuttcap
\pgfsetmiterjoin
\pgfsetdash{}{0pt}
\definecolor{dialinecolor}{rgb}{1.000000, 1.000000, 1.000000}
\pgfsetfillcolor{dialinecolor}
\pgfpathmoveto{\pgfpoint{-5.463093\du}{14.549627\du}}
\pgfpathlineto{\pgfpoint{-4.629556\du}{14.549627\du}}
\pgfpathlineto{\pgfpoint{-4.629556\du}{15.163637\du}}
\pgfpathlineto{\pgfpoint{-5.463093\du}{15.163637\du}}
\pgfpathlineto{\pgfpoint{-5.463093\du}{14.549627\du}}
\pgfusepath{fill}
\pgfsetbuttcap
\pgfsetmiterjoin
\pgfsetdash{}{0pt}
\definecolor{dialinecolor}{rgb}{0.286275, 0.286275, 0.211765}
\pgfsetstrokecolor{dialinecolor}
\pgfpathmoveto{\pgfpoint{-5.463093\du}{14.549627\du}}
\pgfpathlineto{\pgfpoint{-4.629556\du}{14.549627\du}}
\pgfpathlineto{\pgfpoint{-4.629556\du}{15.163415\du}}
\pgfpathlineto{\pgfpoint{-5.463093\du}{15.163415\du}}
\pgfpathlineto{\pgfpoint{-5.463093\du}{14.549627\du}}
\pgfusepath{stroke}
\pgfsetbuttcap
\pgfsetmiterjoin
\pgfsetdash{}{0pt}
\definecolor{dialinecolor}{rgb}{0.478431, 0.478431, 0.352941}
\pgfsetfillcolor{dialinecolor}
\pgfpathmoveto{\pgfpoint{-4.550837\du}{15.234817\du}}
\pgfpathlineto{\pgfpoint{-4.440408\du}{15.124610\du}}
\pgfpathlineto{\pgfpoint{-4.440408\du}{14.345400\du}}
\pgfpathlineto{\pgfpoint{-4.550837\du}{14.447624\du}}
\pgfpathlineto{\pgfpoint{-4.550837\du}{15.234817\du}}
\pgfusepath{fill}
\pgfsetbuttcap
\pgfsetmiterjoin
\pgfsetdash{}{0pt}
\definecolor{dialinecolor}{rgb}{0.286275, 0.286275, 0.211765}
\pgfsetstrokecolor{dialinecolor}
\pgfpathmoveto{\pgfpoint{-4.550837\du}{15.234817\du}}
\pgfpathlineto{\pgfpoint{-4.440408\du}{15.124610\du}}
\pgfpathlineto{\pgfpoint{-4.440408\du}{14.345400\du}}
\pgfpathlineto{\pgfpoint{-4.550837\du}{14.447624\du}}
\pgfpathlineto{\pgfpoint{-4.550837\du}{15.234817\du}}
\pgfusepath{stroke}
\pgfsetlinewidth{0.100000\du}
\pgfsetdash{}{0pt}
\pgfsetdash{}{0pt}
\pgfsetbuttcap
\pgfsetmiterjoin
\pgfsetlinewidth{0.001000\du}
\pgfsetbuttcap
\pgfsetmiterjoin
\pgfsetdash{}{0pt}
\definecolor{dialinecolor}{rgb}{0.000000, 0.470588, 0.666667}
\pgfsetstrokecolor{dialinecolor}
\pgfpathmoveto{\pgfpoint{-4.219239\du}{13.412883\du}}
\pgfpathlineto{\pgfpoint{-4.214114\du}{13.405895\du}}
\pgfpathlineto{\pgfpoint{-4.209222\du}{13.398441\du}}
\pgfpathlineto{\pgfpoint{-4.205029\du}{13.390520\du}}
\pgfpathlineto{\pgfpoint{-4.201069\du}{13.383299\du}}
\pgfpathlineto{\pgfpoint{-4.198041\du}{13.375379\du}}
\pgfpathlineto{\pgfpoint{-4.195479\du}{13.367459\du}}
\pgfpathlineto{\pgfpoint{-4.193149\du}{13.359538\du}}
\pgfpathlineto{\pgfpoint{-4.191751\du}{13.351385\du}}
\pgfpathlineto{\pgfpoint{-4.190820\du}{13.343465\du}}
\pgfpathlineto{\pgfpoint{-4.190121\du}{13.335312\du}}
\pgfpathlineto{\pgfpoint{-4.190121\du}{13.327159\du}}
\pgfpathlineto{\pgfpoint{-4.190820\du}{13.318773\du}}
\pgfpathlineto{\pgfpoint{-4.191751\du}{13.310852\du}}
\pgfpathlineto{\pgfpoint{-4.193149\du}{13.302699\du}}
\pgfpathlineto{\pgfpoint{-4.195246\du}{13.295012\du}}
\pgfpathlineto{\pgfpoint{-4.198041\du}{13.286859\du}}
\pgfpathlineto{\pgfpoint{-4.201069\du}{13.279171\du}}
\pgfpathlineto{\pgfpoint{-4.205029\du}{13.271251\du}}
\pgfpathlineto{\pgfpoint{-4.209222\du}{13.263797\du}}
\pgfpathlineto{\pgfpoint{-4.214114\du}{13.256808\du}}
\pgfpathlineto{\pgfpoint{-4.219239\du}{13.249354\du}}
\pgfpathlineto{\pgfpoint{-4.225063\du}{13.242599\du}}
\pgfpathlineto{\pgfpoint{-4.230887\du}{13.235843\du}}
\pgfpathlineto{\pgfpoint{-4.237409\du}{13.229321\du}}
\pgfpathlineto{\pgfpoint{-4.244165\du}{13.223031\du}}
\pgfpathlineto{\pgfpoint{-4.251386\du}{13.216974\du}}
\pgfpathlineto{\pgfpoint{-4.258840\du}{13.211384\du}}
\pgfpathlineto{\pgfpoint{-4.267226\du}{13.206259\du}}
\pgfpathlineto{\pgfpoint{-4.275380\du}{13.201134\du}}
\pgfpathlineto{\pgfpoint{-4.284232\du}{13.196242\du}}
\pgfpathlineto{\pgfpoint{-4.293549\du}{13.191816\du}}
\pgfpathlineto{\pgfpoint{-4.302867\du}{13.187623\du}}
\pgfpathlineto{\pgfpoint{-4.312418\du}{13.184362\du}}
\pgfpathlineto{\pgfpoint{-4.322202\du}{13.181101\du}}
\pgfpathlineto{\pgfpoint{-4.332452\du}{13.177839\du}}
\pgfpathlineto{\pgfpoint{-4.342468\du}{13.175510\du}}
\pgfpathlineto{\pgfpoint{-4.352718\du}{13.173413\du}}
\pgfpathlineto{\pgfpoint{-4.363434\du}{13.171550\du}}
\pgfpathlineto{\pgfpoint{-4.373683\du}{13.170385\du}}
\pgfpathlineto{\pgfpoint{-4.384632\du}{13.169453\du}}
\pgfpathlineto{\pgfpoint{-4.395348\du}{13.168754\du}}
\pgfpathlineto{\pgfpoint{-4.406296\du}{13.168754\du}}
\pgfpathlineto{\pgfpoint{-4.416779\du}{13.169220\du}}
\pgfpathlineto{\pgfpoint{-4.427494\du}{13.169686\du}}
\pgfpathlineto{\pgfpoint{-4.438443\du}{13.170851\du}}
\pgfpathlineto{\pgfpoint{-4.448693\du}{13.172714\du}}
\pgfpathlineto{\pgfpoint{-4.459408\du}{13.174578\du}}
\pgfpathlineto{\pgfpoint{-4.469192\du}{13.176907\du}}
\pgfpathlineto{\pgfpoint{-4.479442\du}{13.179237\du}}
\pgfpathlineto{\pgfpoint{-4.489225\du}{13.182498\du}}
\pgfpathlineto{\pgfpoint{-4.499242\du}{13.186225\du}}
\pgfpathlineto{\pgfpoint{-4.508793\du}{13.189720\du}}
\pgfpathlineto{\pgfpoint{-4.517645\du}{13.193913\du}}
\pgfpathlineto{\pgfpoint{-4.526730\du}{13.198805\du}}
\pgfpathlineto{\pgfpoint{-4.535349\du}{13.203463\du}}
\pgfpathlineto{\pgfpoint{-4.543502\du}{13.208821\du}}
\pgfpathlineto{\pgfpoint{-4.551422\du}{13.213946\du}}
\pgfpathlineto{\pgfpoint{-4.558877\du}{13.220003\du}}
\pgfpathlineto{\pgfpoint{-4.566331\du}{13.226292\du}}
\pgfpathlineto{\pgfpoint{-4.572854\du}{13.232582\du}}
\pgfpathlineto{\pgfpoint{-4.579143\du}{13.239104\du}}
\pgfpathlineto{\pgfpoint{-4.585200\du}{13.246093\du}}
\pgfusepath{stroke}
\pgfsetlinewidth{0.080000\du}
\pgfsetbuttcap
\pgfsetmiterjoin
\pgfsetdash{}{0pt}
\definecolor{dialinecolor}{rgb}{0.666667, 0.901961, 1.000000}
\pgfsetstrokecolor{dialinecolor}
\pgfpathmoveto{\pgfpoint{-4.226461\du}{13.407060\du}}
\pgfpathlineto{\pgfpoint{-4.222267\du}{13.400071\du}}
\pgfpathlineto{\pgfpoint{-4.217376\du}{13.393316\du}}
\pgfpathlineto{\pgfpoint{-4.213415\du}{13.386094\du}}
\pgfpathlineto{\pgfpoint{-4.210154\du}{13.378873\du}}
\pgfpathlineto{\pgfpoint{-4.207126\du}{13.371419\du}}
\pgfpathlineto{\pgfpoint{-4.204796\du}{13.363964\du}}
\pgfpathlineto{\pgfpoint{-4.202700\du}{13.356277\du}}
\pgfpathlineto{\pgfpoint{-4.201069\du}{13.349056\du}}
\pgfpathlineto{\pgfpoint{-4.200370\du}{13.341368\du}}
\pgfpathlineto{\pgfpoint{-4.199905\du}{13.333448\du}}
\pgfpathlineto{\pgfpoint{-4.199905\du}{13.325761\du}}
\pgfpathlineto{\pgfpoint{-4.200603\du}{13.318074\du}}
\pgfpathlineto{\pgfpoint{-4.201768\du}{13.310153\du}}
\pgfpathlineto{\pgfpoint{-4.203166\du}{13.302932\du}}
\pgfpathlineto{\pgfpoint{-4.205495\du}{13.295245\du}}
\pgfpathlineto{\pgfpoint{-4.208291\du}{13.287791\du}}
\pgfpathlineto{\pgfpoint{-4.211319\du}{13.280569\du}}
\pgfpathlineto{\pgfpoint{-4.214813\du}{13.273348\du}}
\pgfpathlineto{\pgfpoint{-4.218773\du}{13.266592\du}}
\pgfpathlineto{\pgfpoint{-4.223199\du}{13.259371\du}}
\pgfpathlineto{\pgfpoint{-4.228790\du}{13.252848\du}}
\pgfpathlineto{\pgfpoint{-4.233915\du}{13.246093\du}}
\pgfpathlineto{\pgfpoint{-4.239972\du}{13.239803\du}}
\pgfpathlineto{\pgfpoint{-4.246261\du}{13.233747\du}}
\pgfpathlineto{\pgfpoint{-4.252784\du}{13.228389\du}}
\pgfpathlineto{\pgfpoint{-4.259539\du}{13.222565\du}}
\pgfpathlineto{\pgfpoint{-4.267226\du}{13.217207\du}}
\pgfpathlineto{\pgfpoint{-4.275380\du}{13.211850\du}}
\pgfpathlineto{\pgfpoint{-4.283300\du}{13.207424\du}}
\pgfpathlineto{\pgfpoint{-4.291453\du}{13.203231\du}}
\pgfpathlineto{\pgfpoint{-4.300771\du}{13.199037\du}}
\pgfpathlineto{\pgfpoint{-4.309157\du}{13.195077\du}}
\pgfpathlineto{\pgfpoint{-4.318941\du}{13.191816\du}}
\pgfpathlineto{\pgfpoint{-4.327793\du}{13.188788\du}}
\pgfpathlineto{\pgfpoint{-4.337577\du}{13.186225\du}}
\pgfpathlineto{\pgfpoint{-4.347360\du}{13.183896\du}}
\pgfpathlineto{\pgfpoint{-4.357144\du}{13.181799\du}}
\pgfpathlineto{\pgfpoint{-4.367394\du}{13.180402\du}}
\pgfpathlineto{\pgfpoint{-4.377178\du}{13.179237\du}}
\pgfpathlineto{\pgfpoint{-4.387427\du}{13.178305\du}}
\pgfpathlineto{\pgfpoint{-4.398376\du}{13.177839\du}}
\pgfpathlineto{\pgfpoint{-4.408626\du}{13.177839\du}}
\pgfpathlineto{\pgfpoint{-4.418642\du}{13.178305\du}}
\pgfpathlineto{\pgfpoint{-4.428892\du}{13.179004\du}}
\pgfpathlineto{\pgfpoint{-4.438676\du}{13.180402\du}}
\pgfpathlineto{\pgfpoint{-4.448926\du}{13.181799\du}}
\pgfpathlineto{\pgfpoint{-4.458709\du}{13.183430\du}}
\pgfpathlineto{\pgfpoint{-4.468726\du}{13.186225\du}}
\pgfpathlineto{\pgfpoint{-4.478277\du}{13.188555\du}}
\pgfpathlineto{\pgfpoint{-4.487595\du}{13.191583\du}}
\pgfpathlineto{\pgfpoint{-4.497146\du}{13.194844\du}}
\pgfpathlineto{\pgfpoint{-4.505998\du}{13.198805\du}}
\pgfpathlineto{\pgfpoint{-4.514617\du}{13.202532\du}}
\pgfpathlineto{\pgfpoint{-4.523003\du}{13.207191\du}}
\pgfpathlineto{\pgfpoint{-4.531156\du}{13.211617\du}}
\pgfpathlineto{\pgfpoint{-4.539309\du}{13.216742\du}}
\pgfpathlineto{\pgfpoint{-4.546764\du}{13.221866\du}}
\pgfpathlineto{\pgfpoint{-4.553519\du}{13.227457\du}}
\pgfpathlineto{\pgfpoint{-4.560973\du}{13.233048\du}}
\pgfpathlineto{\pgfpoint{-4.567030\du}{13.239104\du}}
\pgfpathlineto{\pgfpoint{-4.572854\du}{13.245394\du}}
\pgfpathlineto{\pgfpoint{-4.578444\du}{13.251917\du}}
\pgfusepath{stroke}
\pgfsetbuttcap
\pgfsetmiterjoin
\pgfsetdash{}{0pt}
\definecolor{dialinecolor}{rgb}{0.000000, 0.705882, 1.000000}
\pgfsetfillcolor{dialinecolor}
\pgfpathmoveto{\pgfpoint{-5.722217\du}{13.159669\du}}
\pgfpathlineto{\pgfpoint{-5.564046\du}{13.010816\du}}
\pgfpathlineto{\pgfpoint{-4.384632\du}{13.010816\du}}
\pgfpathlineto{\pgfpoint{-4.542803\du}{13.159669\du}}
\pgfpathlineto{\pgfpoint{-5.722217\du}{13.159669\du}}
\pgfusepath{fill}
\pgfsetlinewidth{0.001000\du}
\pgfsetbuttcap
\pgfsetmiterjoin
\pgfsetdash{}{0pt}
\definecolor{dialinecolor}{rgb}{0.666667, 0.901961, 1.000000}
\pgfsetstrokecolor{dialinecolor}
\pgfpathmoveto{\pgfpoint{-5.722217\du}{13.159669\du}}
\pgfpathlineto{\pgfpoint{-5.564046\du}{13.010816\du}}
\pgfpathlineto{\pgfpoint{-4.384632\du}{13.010816\du}}
\pgfpathlineto{\pgfpoint{-4.542803\du}{13.159669\du}}
\pgfpathlineto{\pgfpoint{-5.722217\du}{13.159669\du}}
\pgfusepath{stroke}
\pgfsetbuttcap
\pgfsetmiterjoin
\pgfsetdash{}{0pt}
\definecolor{dialinecolor}{rgb}{0.000000, 0.588235, 0.831373}
\pgfsetfillcolor{dialinecolor}
\pgfpathmoveto{\pgfpoint{-5.722217\du}{13.159669\du}}
\pgfpathlineto{\pgfpoint{-4.542803\du}{13.159669\du}}
\pgfpathlineto{\pgfpoint{-4.542803\du}{13.382833\du}}
\pgfpathlineto{\pgfpoint{-5.722217\du}{13.382833\du}}
\pgfpathlineto{\pgfpoint{-5.722217\du}{13.159669\du}}
\pgfusepath{fill}
\pgfsetbuttcap
\pgfsetmiterjoin
\pgfsetdash{}{0pt}
\definecolor{dialinecolor}{rgb}{0.666667, 0.901961, 1.000000}
\pgfsetstrokecolor{dialinecolor}
\pgfpathmoveto{\pgfpoint{-5.722217\du}{13.159669\du}}
\pgfpathlineto{\pgfpoint{-4.542803\du}{13.159669\du}}
\pgfpathlineto{\pgfpoint{-4.542803\du}{13.382833\du}}
\pgfpathlineto{\pgfpoint{-5.722217\du}{13.382833\du}}
\pgfpathlineto{\pgfpoint{-5.722217\du}{13.159669\du}}
\pgfusepath{stroke}
\pgfsetbuttcap
\pgfsetmiterjoin
\pgfsetdash{}{0pt}
\definecolor{dialinecolor}{rgb}{0.000000, 0.352941, 0.501961}
\pgfsetfillcolor{dialinecolor}
\pgfpathmoveto{\pgfpoint{-4.542803\du}{13.382833\du}}
\pgfpathlineto{\pgfpoint{-4.384632\du}{13.224662\du}}
\pgfpathlineto{\pgfpoint{-4.384632\du}{13.010816\du}}
\pgfpathlineto{\pgfpoint{-4.542803\du}{13.159669\du}}
\pgfpathlineto{\pgfpoint{-4.542803\du}{13.382833\du}}
\pgfusepath{fill}
\pgfsetbuttcap
\pgfsetmiterjoin
\pgfsetdash{}{0pt}
\definecolor{dialinecolor}{rgb}{0.666667, 0.901961, 1.000000}
\pgfsetstrokecolor{dialinecolor}
\pgfpathmoveto{\pgfpoint{-4.542803\du}{13.382833\du}}
\pgfpathlineto{\pgfpoint{-4.384632\du}{13.224662\du}}
\pgfpathlineto{\pgfpoint{-4.384632\du}{13.010816\du}}
\pgfpathlineto{\pgfpoint{-4.542803\du}{13.159669\du}}
\pgfpathlineto{\pgfpoint{-4.542803\du}{13.382833\du}}
\pgfusepath{stroke}
\pgfsetbuttcap
\pgfsetmiterjoin
\pgfsetdash{}{0pt}
\definecolor{dialinecolor}{rgb}{0.000000, 0.000000, 0.000000}
\pgfsetfillcolor{dialinecolor}
\pgfpathmoveto{\pgfpoint{-5.684713\du}{13.131483\du}}
\pgfpathlineto{\pgfpoint{-5.564046\du}{13.010816\du}}
\pgfpathlineto{\pgfpoint{-4.412819\du}{13.010816\du}}
\pgfpathlineto{\pgfpoint{-4.533253\du}{13.131483\du}}
\pgfpathlineto{\pgfpoint{-5.684713\du}{13.131483\du}}
\pgfusepath{fill}
\pgfsetbuttcap
\pgfsetmiterjoin
\pgfsetdash{}{0pt}
\definecolor{dialinecolor}{rgb}{0.000000, 0.000000, 0.000000}
\pgfsetstrokecolor{dialinecolor}
\pgfpathmoveto{\pgfpoint{-5.684713\du}{13.131483\du}}
\pgfpathlineto{\pgfpoint{-5.564046\du}{13.010816\du}}
\pgfpathlineto{\pgfpoint{-4.412819\du}{13.010816\du}}
\pgfpathlineto{\pgfpoint{-4.533253\du}{13.131483\du}}
\pgfpathlineto{\pgfpoint{-5.684713\du}{13.131483\du}}
\pgfusepath{stroke}
\pgfsetbuttcap
\pgfsetmiterjoin
\pgfsetdash{}{0pt}
\definecolor{dialinecolor}{rgb}{0.000000, 0.705882, 1.000000}
\pgfsetfillcolor{dialinecolor}
\pgfpathmoveto{\pgfpoint{-5.722217\du}{12.174767\du}}
\pgfpathlineto{\pgfpoint{-5.591766\du}{12.054100\du}}
\pgfpathlineto{\pgfpoint{-4.412819\du}{12.054100\du}}
\pgfpathlineto{\pgfpoint{-4.542803\du}{12.174767\du}}
\pgfpathlineto{\pgfpoint{-5.722217\du}{12.174767\du}}
\pgfusepath{fill}
\pgfsetbuttcap
\pgfsetmiterjoin
\pgfsetdash{}{0pt}
\definecolor{dialinecolor}{rgb}{0.666667, 0.901961, 1.000000}
\pgfsetstrokecolor{dialinecolor}
\pgfpathmoveto{\pgfpoint{-5.722217\du}{12.174767\du}}
\pgfpathlineto{\pgfpoint{-5.591766\du}{12.054100\du}}
\pgfpathlineto{\pgfpoint{-4.412819\du}{12.054100\du}}
\pgfpathlineto{\pgfpoint{-4.542803\du}{12.174767\du}}
\pgfpathlineto{\pgfpoint{-5.722217\du}{12.174767\du}}
\pgfusepath{stroke}
\pgfsetbuttcap
\pgfsetmiterjoin
\pgfsetdash{}{0pt}
\definecolor{dialinecolor}{rgb}{0.000000, 0.588235, 0.831373}
\pgfsetfillcolor{dialinecolor}
\pgfpathmoveto{\pgfpoint{-5.722217\du}{12.174767\du}}
\pgfpathlineto{\pgfpoint{-4.533253\du}{12.174767\du}}
\pgfpathlineto{\pgfpoint{-4.533253\du}{13.112847\du}}
\pgfpathlineto{\pgfpoint{-5.722217\du}{13.112847\du}}
\pgfpathlineto{\pgfpoint{-5.722217\du}{12.174767\du}}
\pgfusepath{fill}
\pgfsetbuttcap
\pgfsetmiterjoin
\pgfsetdash{}{0pt}
\definecolor{dialinecolor}{rgb}{0.666667, 0.901961, 1.000000}
\pgfsetstrokecolor{dialinecolor}
\pgfpathmoveto{\pgfpoint{-5.722217\du}{12.174767\du}}
\pgfpathlineto{\pgfpoint{-4.533485\du}{12.174767\du}}
\pgfpathlineto{\pgfpoint{-4.533485\du}{13.112614\du}}
\pgfpathlineto{\pgfpoint{-5.722217\du}{13.112614\du}}
\pgfpathlineto{\pgfpoint{-5.722217\du}{12.174767\du}}
\pgfusepath{stroke}
\pgfsetbuttcap
\pgfsetmiterjoin
\pgfsetdash{}{0pt}
\definecolor{dialinecolor}{rgb}{1.000000, 1.000000, 1.000000}
\pgfsetfillcolor{dialinecolor}
\pgfpathmoveto{\pgfpoint{-5.619953\du}{12.295667\du}}
\pgfpathlineto{\pgfpoint{-4.635517\du}{12.295667\du}}
\pgfpathlineto{\pgfpoint{-4.635517\du}{13.020367\du}}
\pgfpathlineto{\pgfpoint{-5.619953\du}{13.020367\du}}
\pgfpathlineto{\pgfpoint{-5.619953\du}{12.295667\du}}
\pgfusepath{fill}
\pgfsetbuttcap
\pgfsetmiterjoin
\pgfsetdash{}{0pt}
\definecolor{dialinecolor}{rgb}{0.000000, 0.235294, 0.333333}
\pgfsetstrokecolor{dialinecolor}
\pgfpathmoveto{\pgfpoint{-5.619953\du}{12.295667\du}}
\pgfpathlineto{\pgfpoint{-4.635749\du}{12.295667\du}}
\pgfpathlineto{\pgfpoint{-4.635749\du}{13.020134\du}}
\pgfpathlineto{\pgfpoint{-5.619953\du}{13.020134\du}}
\pgfpathlineto{\pgfpoint{-5.619953\du}{12.295667\du}}
\pgfusepath{stroke}
\pgfsetbuttcap
\pgfsetmiterjoin
\pgfsetdash{}{0pt}
\definecolor{dialinecolor}{rgb}{0.000000, 0.352941, 0.501961}
\pgfsetfillcolor{dialinecolor}
\pgfpathmoveto{\pgfpoint{-4.542803\du}{13.103995\du}}
\pgfpathlineto{\pgfpoint{-4.412819\du}{12.973777\du}}
\pgfpathlineto{\pgfpoint{-4.412819\du}{12.054100\du}}
\pgfpathlineto{\pgfpoint{-4.542803\du}{12.174767\du}}
\pgfpathlineto{\pgfpoint{-4.542803\du}{13.103995\du}}
\pgfusepath{fill}
\pgfsetbuttcap
\pgfsetmiterjoin
\pgfsetdash{}{0pt}
\definecolor{dialinecolor}{rgb}{0.666667, 0.901961, 1.000000}
\pgfsetstrokecolor{dialinecolor}
\pgfpathmoveto{\pgfpoint{-4.542803\du}{13.103995\du}}
\pgfpathlineto{\pgfpoint{-4.412819\du}{12.973777\du}}
\pgfpathlineto{\pgfpoint{-4.412819\du}{12.054100\du}}
\pgfpathlineto{\pgfpoint{-4.542803\du}{12.174767\du}}
\pgfpathlineto{\pgfpoint{-4.542803\du}{13.103995\du}}
\pgfusepath{stroke}
\pgfsetbuttcap
\pgfsetmiterjoin
\pgfsetdash{}{0pt}
\definecolor{dialinecolor}{rgb}{0.000000, 0.705882, 1.000000}
\pgfsetfillcolor{dialinecolor}
\pgfpathmoveto{\pgfpoint{-5.954000\du}{13.568492\du}}
\pgfpathlineto{\pgfpoint{-5.768574\du}{13.336011\du}}
\pgfpathlineto{\pgfpoint{-4.486896\du}{13.336011\du}}
\pgfpathlineto{\pgfpoint{-4.672555\du}{13.568492\du}}
\pgfpathlineto{\pgfpoint{-5.954000\du}{13.568492\du}}
\pgfusepath{fill}
\pgfsetbuttcap
\pgfsetmiterjoin
\pgfsetdash{}{0pt}
\definecolor{dialinecolor}{rgb}{0.666667, 0.901961, 1.000000}
\pgfsetstrokecolor{dialinecolor}
\pgfpathmoveto{\pgfpoint{-5.954000\du}{13.568492\du}}
\pgfpathlineto{\pgfpoint{-5.768574\du}{13.336011\du}}
\pgfpathlineto{\pgfpoint{-4.486896\du}{13.336011\du}}
\pgfpathlineto{\pgfpoint{-4.672555\du}{13.568492\du}}
\pgfpathlineto{\pgfpoint{-5.954000\du}{13.568492\du}}
\pgfusepath{stroke}
\pgfsetbuttcap
\pgfsetmiterjoin
\pgfsetdash{}{0pt}
\definecolor{dialinecolor}{rgb}{0.000000, 0.352941, 0.501961}
\pgfsetfillcolor{dialinecolor}
\pgfpathmoveto{\pgfpoint{-4.672555\du}{13.614616\du}}
\pgfpathlineto{\pgfpoint{-4.486896\du}{13.419639\du}}
\pgfpathlineto{\pgfpoint{-4.486896\du}{13.336011\du}}
\pgfpathlineto{\pgfpoint{-4.672555\du}{13.568492\du}}
\pgfpathlineto{\pgfpoint{-4.672555\du}{13.614616\du}}
\pgfusepath{fill}
\pgfsetbuttcap
\pgfsetmiterjoin
\pgfsetdash{}{0pt}
\definecolor{dialinecolor}{rgb}{0.666667, 0.901961, 1.000000}
\pgfsetstrokecolor{dialinecolor}
\pgfpathmoveto{\pgfpoint{-4.672555\du}{13.614616\du}}
\pgfpathlineto{\pgfpoint{-4.486896\du}{13.419639\du}}
\pgfpathlineto{\pgfpoint{-4.486896\du}{13.336011\du}}
\pgfpathlineto{\pgfpoint{-4.672555\du}{13.568492\du}}
\pgfpathlineto{\pgfpoint{-4.672555\du}{13.614616\du}}
\pgfusepath{stroke}
\pgfsetbuttcap
\pgfsetmiterjoin
\pgfsetdash{}{0pt}
\definecolor{dialinecolor}{rgb}{0.000000, 0.588235, 0.831373}
\pgfsetfillcolor{dialinecolor}
\pgfpathmoveto{\pgfpoint{-5.954000\du}{13.568492\du}}
\pgfpathlineto{\pgfpoint{-4.672555\du}{13.568492\du}}
\pgfpathlineto{\pgfpoint{-4.672555\du}{13.614616\du}}
\pgfpathlineto{\pgfpoint{-5.954000\du}{13.614616\du}}
\pgfpathlineto{\pgfpoint{-5.954000\du}{13.568492\du}}
\pgfusepath{fill}
\pgfsetbuttcap
\pgfsetmiterjoin
\pgfsetdash{}{0pt}
\definecolor{dialinecolor}{rgb}{0.666667, 0.901961, 1.000000}
\pgfsetstrokecolor{dialinecolor}
\pgfpathmoveto{\pgfpoint{-5.954000\du}{13.568492\du}}
\pgfpathlineto{\pgfpoint{-4.672555\du}{13.568492\du}}
\pgfpathlineto{\pgfpoint{-4.672555\du}{13.614616\du}}
\pgfpathlineto{\pgfpoint{-5.954000\du}{13.614616\du}}
\pgfpathlineto{\pgfpoint{-5.954000\du}{13.568492\du}}
\pgfusepath{stroke}
\pgfsetbuttcap
\pgfsetmiterjoin
\pgfsetdash{}{0pt}
\definecolor{dialinecolor}{rgb}{0.000000, 0.705882, 1.000000}
\pgfsetfillcolor{dialinecolor}
\pgfpathmoveto{\pgfpoint{-4.412819\du}{13.540306\du}}
\pgfpathlineto{\pgfpoint{-4.319873\du}{13.401236\du}}
\pgfpathlineto{\pgfpoint{-4.152383\du}{13.401236\du}}
\pgfpathlineto{\pgfpoint{-4.245096\du}{13.540306\du}}
\pgfpathlineto{\pgfpoint{-4.412819\du}{13.540306\du}}
\pgfusepath{fill}
\pgfsetbuttcap
\pgfsetmiterjoin
\pgfsetdash{}{0pt}
\definecolor{dialinecolor}{rgb}{0.666667, 0.901961, 1.000000}
\pgfsetstrokecolor{dialinecolor}
\pgfpathmoveto{\pgfpoint{-4.412819\du}{13.540306\du}}
\pgfpathlineto{\pgfpoint{-4.319873\du}{13.401236\du}}
\pgfpathlineto{\pgfpoint{-4.152383\du}{13.401236\du}}
\pgfpathlineto{\pgfpoint{-4.245096\du}{13.540306\du}}
\pgfpathlineto{\pgfpoint{-4.412819\du}{13.540306\du}}
\pgfusepath{stroke}
\pgfsetbuttcap
\pgfsetmiterjoin
\pgfsetdash{}{0pt}
\definecolor{dialinecolor}{rgb}{0.000000, 0.352941, 0.501961}
\pgfsetfillcolor{dialinecolor}
\pgfpathmoveto{\pgfpoint{-4.245096\du}{13.586895\du}}
\pgfpathlineto{\pgfpoint{-4.152383\du}{13.484631\du}}
\pgfpathlineto{\pgfpoint{-4.152383\du}{13.401236\du}}
\pgfpathlineto{\pgfpoint{-4.245096\du}{13.540306\du}}
\pgfpathlineto{\pgfpoint{-4.245096\du}{13.586895\du}}
\pgfusepath{fill}
\pgfsetbuttcap
\pgfsetmiterjoin
\pgfsetdash{}{0pt}
\definecolor{dialinecolor}{rgb}{0.666667, 0.901961, 1.000000}
\pgfsetstrokecolor{dialinecolor}
\pgfpathmoveto{\pgfpoint{-4.245096\du}{13.586895\du}}
\pgfpathlineto{\pgfpoint{-4.152383\du}{13.484631\du}}
\pgfpathlineto{\pgfpoint{-4.152383\du}{13.401236\du}}
\pgfpathlineto{\pgfpoint{-4.245096\du}{13.540306\du}}
\pgfpathlineto{\pgfpoint{-4.245096\du}{13.586895\du}}
\pgfusepath{stroke}
\pgfsetbuttcap
\pgfsetmiterjoin
\pgfsetdash{}{0pt}
\definecolor{dialinecolor}{rgb}{0.000000, 0.588235, 0.831373}
\pgfsetfillcolor{dialinecolor}
\pgfpathmoveto{\pgfpoint{-4.412819\du}{13.540306\du}}
\pgfpathlineto{\pgfpoint{-4.245096\du}{13.540306\du}}
\pgfpathlineto{\pgfpoint{-4.245096\du}{13.586895\du}}
\pgfpathlineto{\pgfpoint{-4.412819\du}{13.586895\du}}
\pgfpathlineto{\pgfpoint{-4.412819\du}{13.540306\du}}
\pgfusepath{fill}
\pgfsetbuttcap
\pgfsetmiterjoin
\pgfsetdash{}{0pt}
\definecolor{dialinecolor}{rgb}{0.666667, 0.901961, 1.000000}
\pgfsetstrokecolor{dialinecolor}
\pgfpathmoveto{\pgfpoint{-4.412819\du}{13.540306\du}}
\pgfpathlineto{\pgfpoint{-4.245096\du}{13.540306\du}}
\pgfpathlineto{\pgfpoint{-4.245096\du}{13.586895\du}}
\pgfpathlineto{\pgfpoint{-4.412819\du}{13.586895\du}}
\pgfpathlineto{\pgfpoint{-4.412819\du}{13.540306\du}}
\pgfusepath{stroke}
\pgfsetbuttcap
\pgfsetmiterjoin
\pgfsetdash{}{0pt}
\definecolor{dialinecolor}{rgb}{0.811765, 0.054902, 0.188235}
\pgfsetstrokecolor{dialinecolor}
\pgfpathmoveto{\pgfpoint{-5.359984\du}{12.472241\du}}
\pgfpathlineto{\pgfpoint{-5.099548\du}{12.472241\du}}
\pgfpathlineto{\pgfpoint{-5.164774\du}{12.537233\du}}
\pgfpathlineto{\pgfpoint{-4.941843\du}{12.537233\du}}
\pgfusepath{stroke}
\pgfsetbuttcap
\pgfsetmiterjoin
\pgfsetdash{}{0pt}
\definecolor{dialinecolor}{rgb}{0.811765, 0.054902, 0.188235}
\pgfsetstrokecolor{dialinecolor}
\pgfpathmoveto{\pgfpoint{-5.387937\du}{12.472241\du}}
\pgfpathlineto{\pgfpoint{-5.201812\du}{12.657900\du}}
\pgfpathlineto{\pgfpoint{-5.294991\du}{12.657900\du}}
\pgfpathlineto{\pgfpoint{-5.137053\du}{12.816072\du}}
\pgfusepath{stroke}
\pgfsetbuttcap
\pgfsetmiterjoin
\pgfsetdash{}{0pt}
\definecolor{dialinecolor}{rgb}{0.811765, 0.054902, 0.188235}
\pgfsetstrokecolor{dialinecolor}
\pgfpathmoveto{\pgfpoint{-4.858680\du}{12.472241\du}}
\pgfpathlineto{\pgfpoint{-5.043874\du}{12.657900\du}}
\pgfpathlineto{\pgfpoint{-4.951393\du}{12.657900\du}}
\pgfpathlineto{\pgfpoint{-5.109099\du}{12.816072\du}}
\pgfusepath{stroke}
\pgfsetbuttcap
\pgfsetmiterjoin
\pgfsetdash{}{0pt}
\definecolor{dialinecolor}{rgb}{0.000000, 0.470588, 0.666667}
\pgfsetfillcolor{dialinecolor}
\pgfpathmoveto{\pgfpoint{-5.257953\du}{12.541659\du}}
\pgfpathlineto{\pgfpoint{-5.258418\du}{12.539097\du}}
\pgfpathlineto{\pgfpoint{-5.259117\du}{12.536534\du}}
\pgfpathlineto{\pgfpoint{-5.259816\du}{12.533972\du}}
\pgfpathlineto{\pgfpoint{-5.261214\du}{12.531410\du}}
\pgfpathlineto{\pgfpoint{-5.263077\du}{12.529080\du}}
\pgfpathlineto{\pgfpoint{-5.265174\du}{12.526518\du}}
\pgfpathlineto{\pgfpoint{-5.267503\du}{12.523955\du}}
\pgfpathlineto{\pgfpoint{-5.270299\du}{12.521626\du}}
\pgfpathlineto{\pgfpoint{-5.273560\du}{12.519296\du}}
\pgfpathlineto{\pgfpoint{-5.277287\du}{12.516967\du}}
\pgfpathlineto{\pgfpoint{-5.281247\du}{12.514870\du}}
\pgfpathlineto{\pgfpoint{-5.285673\du}{12.512541\du}}
\pgfpathlineto{\pgfpoint{-5.290099\du}{12.510677\du}}
\pgfpathlineto{\pgfpoint{-5.294991\du}{12.508581\du}}
\pgfpathlineto{\pgfpoint{-5.300116\du}{12.506717\du}}
\pgfpathlineto{\pgfpoint{-5.305707\du}{12.504854\du}}
\pgfpathlineto{\pgfpoint{-5.311530\du}{12.503223\du}}
\pgfpathlineto{\pgfpoint{-5.317354\du}{12.501359\du}}
\pgfpathlineto{\pgfpoint{-5.323877\du}{12.500195\du}}
\pgfpathlineto{\pgfpoint{-5.330399\du}{12.498564\du}}
\pgfpathlineto{\pgfpoint{-5.336689\du}{12.497166\du}}
\pgfpathlineto{\pgfpoint{-5.343910\du}{12.496235\du}}
\pgfpathlineto{\pgfpoint{-5.350666\du}{12.495070\du}}
\pgfpathlineto{\pgfpoint{-5.358120\du}{12.494138\du}}
\pgfpathlineto{\pgfpoint{-5.365341\du}{12.493206\du}}
\pgfpathlineto{\pgfpoint{-5.372563\du}{12.492507\du}}
\pgfpathlineto{\pgfpoint{-5.380250\du}{12.491809\du}}
\pgfpathlineto{\pgfpoint{-5.387937\du}{12.491343\du}}
\pgfpathlineto{\pgfpoint{-5.395625\du}{12.491110\du}}
\pgfpathlineto{\pgfpoint{-5.403545\du}{12.490877\du}}
\pgfpathlineto{\pgfpoint{-5.411232\du}{12.490877\du}}
\pgfpathlineto{\pgfpoint{-5.411232\du}{12.490877\du}}
\pgfpathlineto{\pgfpoint{-5.418919\du}{12.490877\du}}
\pgfpathlineto{\pgfpoint{-5.426607\du}{12.491110\du}}
\pgfpathlineto{\pgfpoint{-5.434527\du}{12.491343\du}}
\pgfpathlineto{\pgfpoint{-5.442214\du}{12.491809\du}}
\pgfpathlineto{\pgfpoint{-5.449668\du}{12.492507\du}}
\pgfpathlineto{\pgfpoint{-5.457123\du}{12.493206\du}}
\pgfpathlineto{\pgfpoint{-5.464344\du}{12.494138\du}}
\pgfpathlineto{\pgfpoint{-5.471566\du}{12.495070\du}}
\pgfpathlineto{\pgfpoint{-5.478554\du}{12.496235\du}}
\pgfpathlineto{\pgfpoint{-5.485542\du}{12.497166\du}}
\pgfpathlineto{\pgfpoint{-5.492065\du}{12.498564\du}}
\pgfpathlineto{\pgfpoint{-5.498587\du}{12.500195\du}}
\pgfpathlineto{\pgfpoint{-5.505110\du}{12.501359\du}}
\pgfpathlineto{\pgfpoint{-5.510701\du}{12.503223\du}}
\pgfpathlineto{\pgfpoint{-5.516524\du}{12.504854\du}}
\pgfpathlineto{\pgfpoint{-5.522115\du}{12.506717\du}}
\pgfpathlineto{\pgfpoint{-5.527473\du}{12.508581\du}}
\pgfpathlineto{\pgfpoint{-5.532132\du}{12.510677\du}}
\pgfpathlineto{\pgfpoint{-5.536791\du}{12.512541\du}}
\pgfpathlineto{\pgfpoint{-5.541217\du}{12.514870\du}}
\pgfpathlineto{\pgfpoint{-5.545177\du}{12.516967\du}}
\pgfpathlineto{\pgfpoint{-5.548904\du}{12.519296\du}}
\pgfpathlineto{\pgfpoint{-5.552165\du}{12.521626\du}}
\pgfpathlineto{\pgfpoint{-5.554728\du}{12.523955\du}}
\pgfpathlineto{\pgfpoint{-5.557290\du}{12.526518\du}}
\pgfpathlineto{\pgfpoint{-5.559620\du}{12.529080\du}}
\pgfpathlineto{\pgfpoint{-5.561250\du}{12.531410\du}}
\pgfpathlineto{\pgfpoint{-5.562415\du}{12.533972\du}}
\pgfpathlineto{\pgfpoint{-5.563813\du}{12.536534\du}}
\pgfpathlineto{\pgfpoint{-5.564046\du}{12.539097\du}}
\pgfpathlineto{\pgfpoint{-5.564279\du}{12.541659\du}}
\pgfpathlineto{\pgfpoint{-5.564279\du}{12.541659\du}}
\pgfpathlineto{\pgfpoint{-5.564046\du}{12.544455\du}}
\pgfpathlineto{\pgfpoint{-5.563813\du}{12.546784\du}}
\pgfpathlineto{\pgfpoint{-5.562415\du}{12.549347\du}}
\pgfpathlineto{\pgfpoint{-5.561250\du}{12.551909\du}}
\pgfpathlineto{\pgfpoint{-5.559620\du}{12.554471\du}}
\pgfpathlineto{\pgfpoint{-5.557290\du}{12.557034\du}}
\pgfpathlineto{\pgfpoint{-5.554728\du}{12.559363\du}}
\pgfpathlineto{\pgfpoint{-5.552165\du}{12.561693\du}}
\pgfpathlineto{\pgfpoint{-5.548904\du}{12.564022\du}}
\pgfpathlineto{\pgfpoint{-5.545177\du}{12.566585\du}}
\pgfpathlineto{\pgfpoint{-5.541217\du}{12.568681\du}}
\pgfpathlineto{\pgfpoint{-5.536791\du}{12.570778\du}}
\pgfpathlineto{\pgfpoint{-5.532132\du}{12.572874\du}}
\pgfpathlineto{\pgfpoint{-5.527473\du}{12.574971\du}}
\pgfpathlineto{\pgfpoint{-5.522115\du}{12.576834\du}}
\pgfpathlineto{\pgfpoint{-5.516524\du}{12.578698\du}}
\pgfpathlineto{\pgfpoint{-5.510701\du}{12.580096\du}}
\pgfpathlineto{\pgfpoint{-5.505110\du}{12.581726\du}}
\pgfpathlineto{\pgfpoint{-5.498587\du}{12.583357\du}}
\pgfpathlineto{\pgfpoint{-5.492065\du}{12.584988\du}}
\pgfpathlineto{\pgfpoint{-5.485542\du}{12.586152\du}}
\pgfpathlineto{\pgfpoint{-5.478554\du}{12.587317\du}}
\pgfpathlineto{\pgfpoint{-5.471566\du}{12.588482\du}}
\pgfpathlineto{\pgfpoint{-5.464344\du}{12.589414\du}}
\pgfpathlineto{\pgfpoint{-5.457123\du}{12.590112\du}}
\pgfpathlineto{\pgfpoint{-5.449668\du}{12.591044\du}}
\pgfpathlineto{\pgfpoint{-5.442214\du}{12.591510\du}}
\pgfpathlineto{\pgfpoint{-5.434527\du}{12.591976\du}}
\pgfpathlineto{\pgfpoint{-5.426607\du}{12.592442\du}}
\pgfpathlineto{\pgfpoint{-5.418919\du}{12.592675\du}}
\pgfpathlineto{\pgfpoint{-5.411232\du}{12.592675\du}}
\pgfpathlineto{\pgfpoint{-5.411232\du}{12.592675\du}}
\pgfpathlineto{\pgfpoint{-5.403545\du}{12.592675\du}}
\pgfpathlineto{\pgfpoint{-5.395625\du}{12.592442\du}}
\pgfpathlineto{\pgfpoint{-5.387937\du}{12.591976\du}}
\pgfpathlineto{\pgfpoint{-5.380250\du}{12.591510\du}}
\pgfpathlineto{\pgfpoint{-5.372563\du}{12.591044\du}}
\pgfpathlineto{\pgfpoint{-5.365341\du}{12.590112\du}}
\pgfpathlineto{\pgfpoint{-5.358120\du}{12.589414\du}}
\pgfpathlineto{\pgfpoint{-5.350666\du}{12.588482\du}}
\pgfpathlineto{\pgfpoint{-5.343910\du}{12.587317\du}}
\pgfpathlineto{\pgfpoint{-5.336689\du}{12.586152\du}}
\pgfpathlineto{\pgfpoint{-5.330399\du}{12.584988\du}}
\pgfpathlineto{\pgfpoint{-5.323877\du}{12.583357\du}}
\pgfpathlineto{\pgfpoint{-5.317354\du}{12.581726\du}}
\pgfpathlineto{\pgfpoint{-5.311530\du}{12.580096\du}}
\pgfpathlineto{\pgfpoint{-5.305707\du}{12.578698\du}}
\pgfpathlineto{\pgfpoint{-5.300116\du}{12.576834\du}}
\pgfpathlineto{\pgfpoint{-5.294991\du}{12.574971\du}}
\pgfpathlineto{\pgfpoint{-5.290099\du}{12.572874\du}}
\pgfpathlineto{\pgfpoint{-5.285673\du}{12.570778\du}}
\pgfpathlineto{\pgfpoint{-5.281247\du}{12.568681\du}}
\pgfpathlineto{\pgfpoint{-5.277287\du}{12.566585\du}}
\pgfpathlineto{\pgfpoint{-5.273560\du}{12.564022\du}}
\pgfpathlineto{\pgfpoint{-5.270299\du}{12.561693\du}}
\pgfpathlineto{\pgfpoint{-5.267503\du}{12.559363\du}}
\pgfpathlineto{\pgfpoint{-5.265174\du}{12.557034\du}}
\pgfpathlineto{\pgfpoint{-5.263077\du}{12.554471\du}}
\pgfpathlineto{\pgfpoint{-5.261214\du}{12.551909\du}}
\pgfpathlineto{\pgfpoint{-5.259816\du}{12.549347\du}}
\pgfpathlineto{\pgfpoint{-5.259117\du}{12.546784\du}}
\pgfpathlineto{\pgfpoint{-5.258418\du}{12.544455\du}}
\pgfpathlineto{\pgfpoint{-5.257953\du}{12.541659\du}}
\pgfusepath{fill}
\pgfsetbuttcap
\pgfsetmiterjoin
\pgfsetdash{}{0pt}
\definecolor{dialinecolor}{rgb}{0.666667, 0.901961, 1.000000}
\pgfsetstrokecolor{dialinecolor}
\pgfpathmoveto{\pgfpoint{-5.267503\du}{12.537233\du}}
\pgfpathlineto{\pgfpoint{-5.267503\du}{12.534438\du}}
\pgfpathlineto{\pgfpoint{-5.268202\du}{12.532574\du}}
\pgfpathlineto{\pgfpoint{-5.269134\du}{12.530012\du}}
\pgfpathlineto{\pgfpoint{-5.270299\du}{12.527682\du}}
\pgfpathlineto{\pgfpoint{-5.272162\du}{12.525586\du}}
\pgfpathlineto{\pgfpoint{-5.274026\du}{12.522791\du}}
\pgfpathlineto{\pgfpoint{-5.276588\du}{12.520927\du}}
\pgfpathlineto{\pgfpoint{-5.279151\du}{12.518830\du}}
\pgfpathlineto{\pgfpoint{-5.282412\du}{12.516734\du}}
\pgfpathlineto{\pgfpoint{-5.285906\du}{12.514870\du}}
\pgfpathlineto{\pgfpoint{-5.289633\du}{12.512541\du}}
\pgfpathlineto{\pgfpoint{-5.294059\du}{12.510910\du}}
\pgfpathlineto{\pgfpoint{-5.298252\du}{12.508814\du}}
\pgfpathlineto{\pgfpoint{-5.303144\du}{12.506717\du}}
\pgfpathlineto{\pgfpoint{-5.307803\du}{12.505320\du}}
\pgfpathlineto{\pgfpoint{-5.313627\du}{12.503456\du}}
\pgfpathlineto{\pgfpoint{-5.318985\du}{12.502291\du}}
\pgfpathlineto{\pgfpoint{-5.324576\du}{12.500661\du}}
\pgfpathlineto{\pgfpoint{-5.330632\du}{12.499263\du}}
\pgfpathlineto{\pgfpoint{-5.337388\du}{12.497865\du}}
\pgfpathlineto{\pgfpoint{-5.343910\du}{12.496933\du}}
\pgfpathlineto{\pgfpoint{-5.350433\du}{12.495070\du}}
\pgfpathlineto{\pgfpoint{-5.357421\du}{12.494837\du}}
\pgfpathlineto{\pgfpoint{-5.364177\du}{12.493672\du}}
\pgfpathlineto{\pgfpoint{-5.371165\du}{12.492973\du}}
\pgfpathlineto{\pgfpoint{-5.378386\du}{12.492274\du}}
\pgfpathlineto{\pgfpoint{-5.385841\du}{12.491809\du}}
\pgfpathlineto{\pgfpoint{-5.393295\du}{12.491576\du}}
\pgfpathlineto{\pgfpoint{-5.400284\du}{12.491110\du}}
\pgfpathlineto{\pgfpoint{-5.408204\du}{12.490877\du}}
\pgfpathlineto{\pgfpoint{-5.415658\du}{12.490877\du}}
\pgfpathlineto{\pgfpoint{-5.415658\du}{12.490877\du}}
\pgfpathlineto{\pgfpoint{-5.422879\du}{12.490877\du}}
\pgfpathlineto{\pgfpoint{-5.430567\du}{12.491110\du}}
\pgfpathlineto{\pgfpoint{-5.438021\du}{12.491576\du}}
\pgfpathlineto{\pgfpoint{-5.445475\du}{12.491809\du}}
\pgfpathlineto{\pgfpoint{-5.452697\du}{12.492274\du}}
\pgfpathlineto{\pgfpoint{-5.460384\du}{12.492973\du}}
\pgfpathlineto{\pgfpoint{-5.467140\du}{12.493672\du}}
\pgfpathlineto{\pgfpoint{-5.474128\du}{12.494837\du}}
\pgfpathlineto{\pgfpoint{-5.480650\du}{12.495070\du}}
\pgfpathlineto{\pgfpoint{-5.487639\du}{12.496933\du}}
\pgfpathlineto{\pgfpoint{-5.494161\du}{12.497865\du}}
\pgfpathlineto{\pgfpoint{-5.500451\du}{12.499263\du}}
\pgfpathlineto{\pgfpoint{-5.506508\du}{12.500661\du}}
\pgfpathlineto{\pgfpoint{-5.512331\du}{12.502291\du}}
\pgfpathlineto{\pgfpoint{-5.517689\du}{12.503456\du}}
\pgfpathlineto{\pgfpoint{-5.523513\du}{12.505320\du}}
\pgfpathlineto{\pgfpoint{-5.528172\du}{12.506717\du}}
\pgfpathlineto{\pgfpoint{-5.533064\du}{12.508814\du}}
\pgfpathlineto{\pgfpoint{-5.537490\du}{12.510910\du}}
\pgfpathlineto{\pgfpoint{-5.541683\du}{12.512541\du}}
\pgfpathlineto{\pgfpoint{-5.545410\du}{12.514870\du}}
\pgfpathlineto{\pgfpoint{-5.548904\du}{12.516734\du}}
\pgfpathlineto{\pgfpoint{-5.552165\du}{12.518830\du}}
\pgfpathlineto{\pgfpoint{-5.554728\du}{12.520927\du}}
\pgfpathlineto{\pgfpoint{-5.557523\du}{12.522791\du}}
\pgfpathlineto{\pgfpoint{-5.558921\du}{12.525586\du}}
\pgfpathlineto{\pgfpoint{-5.560784\du}{12.527682\du}}
\pgfpathlineto{\pgfpoint{-5.562182\du}{12.530012\du}}
\pgfpathlineto{\pgfpoint{-5.563347\du}{12.532574\du}}
\pgfpathlineto{\pgfpoint{-5.563813\du}{12.534438\du}}
\pgfpathlineto{\pgfpoint{-5.564046\du}{12.537233\du}}
\pgfpathlineto{\pgfpoint{-5.564046\du}{12.537233\du}}
\pgfpathlineto{\pgfpoint{-5.563813\du}{12.539563\du}}
\pgfpathlineto{\pgfpoint{-5.563347\du}{12.541659\du}}
\pgfpathlineto{\pgfpoint{-5.562182\du}{12.543989\du}}
\pgfpathlineto{\pgfpoint{-5.560784\du}{12.546784\du}}
\pgfpathlineto{\pgfpoint{-5.558921\du}{12.548881\du}}
\pgfpathlineto{\pgfpoint{-5.557523\du}{12.551210\du}}
\pgfpathlineto{\pgfpoint{-5.554728\du}{12.553307\du}}
\pgfpathlineto{\pgfpoint{-5.552165\du}{12.555403\du}}
\pgfpathlineto{\pgfpoint{-5.548904\du}{12.557500\du}}
\pgfpathlineto{\pgfpoint{-5.545410\du}{12.559596\du}}
\pgfpathlineto{\pgfpoint{-5.541683\du}{12.561693\du}}
\pgfpathlineto{\pgfpoint{-5.537490\du}{12.563556\du}}
\pgfpathlineto{\pgfpoint{-5.533064\du}{12.565653\du}}
\pgfpathlineto{\pgfpoint{-5.528172\du}{12.567284\du}}
\pgfpathlineto{\pgfpoint{-5.523513\du}{12.569147\du}}
\pgfpathlineto{\pgfpoint{-5.517689\du}{12.571011\du}}
\pgfpathlineto{\pgfpoint{-5.512331\du}{12.572408\du}}
\pgfpathlineto{\pgfpoint{-5.506508\du}{12.573573\du}}
\pgfpathlineto{\pgfpoint{-5.500451\du}{12.575204\du}}
\pgfpathlineto{\pgfpoint{-5.494161\du}{12.576601\du}}
\pgfpathlineto{\pgfpoint{-5.487639\du}{12.577533\du}}
\pgfpathlineto{\pgfpoint{-5.480650\du}{12.578931\du}}
\pgfpathlineto{\pgfpoint{-5.474128\du}{12.579630\du}}
\pgfpathlineto{\pgfpoint{-5.467140\du}{12.580795\du}}
\pgfpathlineto{\pgfpoint{-5.460384\du}{12.581493\du}}
\pgfpathlineto{\pgfpoint{-5.452697\du}{12.581726\du}}
\pgfpathlineto{\pgfpoint{-5.445475\du}{12.582658\du}}
\pgfpathlineto{\pgfpoint{-5.438021\du}{12.582891\du}}
\pgfpathlineto{\pgfpoint{-5.430567\du}{12.583124\du}}
\pgfpathlineto{\pgfpoint{-5.422879\du}{12.583590\du}}
\pgfpathlineto{\pgfpoint{-5.415658\du}{12.583590\du}}
\pgfpathlineto{\pgfpoint{-5.415658\du}{12.583590\du}}
\pgfpathlineto{\pgfpoint{-5.408204\du}{12.583590\du}}
\pgfpathlineto{\pgfpoint{-5.400284\du}{12.583124\du}}
\pgfpathlineto{\pgfpoint{-5.393295\du}{12.582891\du}}
\pgfpathlineto{\pgfpoint{-5.385841\du}{12.582658\du}}
\pgfpathlineto{\pgfpoint{-5.378386\du}{12.581726\du}}
\pgfpathlineto{\pgfpoint{-5.371165\du}{12.581493\du}}
\pgfpathlineto{\pgfpoint{-5.364177\du}{12.580795\du}}
\pgfpathlineto{\pgfpoint{-5.357421\du}{12.579630\du}}
\pgfpathlineto{\pgfpoint{-5.350433\du}{12.578931\du}}
\pgfpathlineto{\pgfpoint{-5.343910\du}{12.577533\du}}
\pgfpathlineto{\pgfpoint{-5.337388\du}{12.576601\du}}
\pgfpathlineto{\pgfpoint{-5.330632\du}{12.575204\du}}
\pgfpathlineto{\pgfpoint{-5.324576\du}{12.573573\du}}
\pgfpathlineto{\pgfpoint{-5.318985\du}{12.572408\du}}
\pgfpathlineto{\pgfpoint{-5.313627\du}{12.571011\du}}
\pgfpathlineto{\pgfpoint{-5.307803\du}{12.569147\du}}
\pgfpathlineto{\pgfpoint{-5.303144\du}{12.567284\du}}
\pgfpathlineto{\pgfpoint{-5.298252\du}{12.565653\du}}
\pgfpathlineto{\pgfpoint{-5.294059\du}{12.563556\du}}
\pgfpathlineto{\pgfpoint{-5.289633\du}{12.561693\du}}
\pgfpathlineto{\pgfpoint{-5.285906\du}{12.559596\du}}
\pgfpathlineto{\pgfpoint{-5.282412\du}{12.557500\du}}
\pgfpathlineto{\pgfpoint{-5.279151\du}{12.555403\du}}
\pgfpathlineto{\pgfpoint{-5.276588\du}{12.553307\du}}
\pgfpathlineto{\pgfpoint{-5.274026\du}{12.551210\du}}
\pgfpathlineto{\pgfpoint{-5.272162\du}{12.548881\du}}
\pgfpathlineto{\pgfpoint{-5.270299\du}{12.546784\du}}
\pgfpathlineto{\pgfpoint{-5.269134\du}{12.543989\du}}
\pgfpathlineto{\pgfpoint{-5.268202\du}{12.541659\du}}
\pgfpathlineto{\pgfpoint{-5.267503\du}{12.539563\du}}
\pgfpathlineto{\pgfpoint{-5.267503\du}{12.537233\du}}
\pgfusepath{stroke}
\pgfsetbuttcap
\pgfsetmiterjoin
\pgfsetdash{}{0pt}
\definecolor{dialinecolor}{rgb}{0.000000, 0.470588, 0.666667}
\pgfsetfillcolor{dialinecolor}
\pgfpathmoveto{\pgfpoint{-5.564046\du}{12.472241\du}}
\pgfpathlineto{\pgfpoint{-5.564046\du}{12.546085\du}}
\pgfpathlineto{\pgfpoint{-5.257720\du}{12.546085\du}}
\pgfpathlineto{\pgfpoint{-5.257720\du}{12.472241\du}}
\pgfpathlineto{\pgfpoint{-5.564046\du}{12.472241\du}}
\pgfusepath{fill}
\pgfsetbuttcap
\pgfsetmiterjoin
\pgfsetdash{}{0pt}
\definecolor{dialinecolor}{rgb}{0.000000, 0.705882, 1.000000}
\pgfsetfillcolor{dialinecolor}
\pgfpathmoveto{\pgfpoint{-5.257953\du}{12.467116\du}}
\pgfpathlineto{\pgfpoint{-5.258418\du}{12.464554\du}}
\pgfpathlineto{\pgfpoint{-5.259117\du}{12.461758\du}}
\pgfpathlineto{\pgfpoint{-5.259816\du}{12.459429\du}}
\pgfpathlineto{\pgfpoint{-5.261214\du}{12.456866\du}}
\pgfpathlineto{\pgfpoint{-5.263077\du}{12.454304\du}}
\pgfpathlineto{\pgfpoint{-5.265174\du}{12.451742\du}}
\pgfpathlineto{\pgfpoint{-5.267503\du}{12.449412\du}}
\pgfpathlineto{\pgfpoint{-5.270299\du}{12.446850\du}}
\pgfpathlineto{\pgfpoint{-5.273560\du}{12.444753\du}}
\pgfpathlineto{\pgfpoint{-5.277287\du}{12.442424\du}}
\pgfpathlineto{\pgfpoint{-5.281247\du}{12.440094\du}}
\pgfpathlineto{\pgfpoint{-5.285673\du}{12.437998\du}}
\pgfpathlineto{\pgfpoint{-5.290099\du}{12.435435\du}}
\pgfpathlineto{\pgfpoint{-5.294991\du}{12.434038\du}}
\pgfpathlineto{\pgfpoint{-5.300116\du}{12.431941\du}}
\pgfpathlineto{\pgfpoint{-5.305707\du}{12.430077\du}}
\pgfpathlineto{\pgfpoint{-5.311530\du}{12.428447\du}}
\pgfpathlineto{\pgfpoint{-5.317354\du}{12.426816\du}}
\pgfpathlineto{\pgfpoint{-5.323877\du}{12.425186\du}}
\pgfpathlineto{\pgfpoint{-5.330399\du}{12.423555\du}}
\pgfpathlineto{\pgfpoint{-5.336689\du}{12.422390\du}}
\pgfpathlineto{\pgfpoint{-5.343910\du}{12.421225\du}}
\pgfpathlineto{\pgfpoint{-5.350666\du}{12.420294\du}}
\pgfpathlineto{\pgfpoint{-5.358120\du}{12.419129\du}}
\pgfpathlineto{\pgfpoint{-5.365341\du}{12.418430\du}}
\pgfpathlineto{\pgfpoint{-5.372563\du}{12.417498\du}}
\pgfpathlineto{\pgfpoint{-5.380250\du}{12.417032\du}}
\pgfpathlineto{\pgfpoint{-5.387937\du}{12.416566\du}}
\pgfpathlineto{\pgfpoint{-5.395625\du}{12.416334\du}}
\pgfpathlineto{\pgfpoint{-5.403545\du}{12.416101\du}}
\pgfpathlineto{\pgfpoint{-5.411232\du}{12.416101\du}}
\pgfpathlineto{\pgfpoint{-5.411232\du}{12.416101\du}}
\pgfpathlineto{\pgfpoint{-5.418919\du}{12.416101\du}}
\pgfpathlineto{\pgfpoint{-5.426607\du}{12.416334\du}}
\pgfpathlineto{\pgfpoint{-5.434527\du}{12.416566\du}}
\pgfpathlineto{\pgfpoint{-5.442214\du}{12.417032\du}}
\pgfpathlineto{\pgfpoint{-5.449668\du}{12.417498\du}}
\pgfpathlineto{\pgfpoint{-5.457123\du}{12.418430\du}}
\pgfpathlineto{\pgfpoint{-5.464344\du}{12.419129\du}}
\pgfpathlineto{\pgfpoint{-5.471566\du}{12.420294\du}}
\pgfpathlineto{\pgfpoint{-5.478554\du}{12.421225\du}}
\pgfpathlineto{\pgfpoint{-5.485542\du}{12.422390\du}}
\pgfpathlineto{\pgfpoint{-5.492065\du}{12.423555\du}}
\pgfpathlineto{\pgfpoint{-5.498587\du}{12.425186\du}}
\pgfpathlineto{\pgfpoint{-5.505110\du}{12.426816\du}}
\pgfpathlineto{\pgfpoint{-5.510701\du}{12.428447\du}}
\pgfpathlineto{\pgfpoint{-5.516524\du}{12.430077\du}}
\pgfpathlineto{\pgfpoint{-5.522115\du}{12.431941\du}}
\pgfpathlineto{\pgfpoint{-5.527473\du}{12.434038\du}}
\pgfpathlineto{\pgfpoint{-5.532132\du}{12.435435\du}}
\pgfpathlineto{\pgfpoint{-5.536791\du}{12.437998\du}}
\pgfpathlineto{\pgfpoint{-5.541217\du}{12.440094\du}}
\pgfpathlineto{\pgfpoint{-5.545177\du}{12.442424\du}}
\pgfpathlineto{\pgfpoint{-5.548904\du}{12.444753\du}}
\pgfpathlineto{\pgfpoint{-5.552165\du}{12.446850\du}}
\pgfpathlineto{\pgfpoint{-5.554728\du}{12.449412\du}}
\pgfpathlineto{\pgfpoint{-5.557290\du}{12.451742\du}}
\pgfpathlineto{\pgfpoint{-5.559620\du}{12.454304\du}}
\pgfpathlineto{\pgfpoint{-5.561250\du}{12.456866\du}}
\pgfpathlineto{\pgfpoint{-5.562415\du}{12.459429\du}}
\pgfpathlineto{\pgfpoint{-5.563813\du}{12.461758\du}}
\pgfpathlineto{\pgfpoint{-5.564046\du}{12.464554\du}}
\pgfpathlineto{\pgfpoint{-5.564279\du}{12.467116\du}}
\pgfpathlineto{\pgfpoint{-5.564279\du}{12.467116\du}}
\pgfpathlineto{\pgfpoint{-5.564046\du}{12.469679\du}}
\pgfpathlineto{\pgfpoint{-5.563813\du}{12.472474\du}}
\pgfpathlineto{\pgfpoint{-5.562415\du}{12.474803\du}}
\pgfpathlineto{\pgfpoint{-5.561250\du}{12.477366\du}}
\pgfpathlineto{\pgfpoint{-5.559620\du}{12.480161\du}}
\pgfpathlineto{\pgfpoint{-5.557290\du}{12.482491\du}}
\pgfpathlineto{\pgfpoint{-5.554728\du}{12.484820\du}}
\pgfpathlineto{\pgfpoint{-5.552165\du}{12.487383\du}}
\pgfpathlineto{\pgfpoint{-5.548904\du}{12.489246\du}}
\pgfpathlineto{\pgfpoint{-5.545177\du}{12.491809\du}}
\pgfpathlineto{\pgfpoint{-5.541217\du}{12.494371\du}}
\pgfpathlineto{\pgfpoint{-5.536791\du}{12.496468\du}}
\pgfpathlineto{\pgfpoint{-5.532132\du}{12.498564\du}}
\pgfpathlineto{\pgfpoint{-5.527473\du}{12.500195\du}}
\pgfpathlineto{\pgfpoint{-5.522115\du}{12.502524\du}}
\pgfpathlineto{\pgfpoint{-5.516524\du}{12.504388\du}}
\pgfpathlineto{\pgfpoint{-5.510701\du}{12.505785\du}}
\pgfpathlineto{\pgfpoint{-5.505110\du}{12.507416\du}}
\pgfpathlineto{\pgfpoint{-5.498587\du}{12.509047\du}}
\pgfpathlineto{\pgfpoint{-5.492065\du}{12.510677\du}}
\pgfpathlineto{\pgfpoint{-5.485542\du}{12.511842\du}}
\pgfpathlineto{\pgfpoint{-5.478554\du}{12.513007\du}}
\pgfpathlineto{\pgfpoint{-5.471566\du}{12.513939\du}}
\pgfpathlineto{\pgfpoint{-5.464344\du}{12.515103\du}}
\pgfpathlineto{\pgfpoint{-5.457123\du}{12.515802\du}}
\pgfpathlineto{\pgfpoint{-5.449668\du}{12.516734\du}}
\pgfpathlineto{\pgfpoint{-5.442214\du}{12.516967\du}}
\pgfpathlineto{\pgfpoint{-5.434527\du}{12.517666\du}}
\pgfpathlineto{\pgfpoint{-5.426607\du}{12.517666\du}}
\pgfpathlineto{\pgfpoint{-5.418919\du}{12.518365\du}}
\pgfpathlineto{\pgfpoint{-5.411232\du}{12.518365\du}}
\pgfpathlineto{\pgfpoint{-5.411232\du}{12.518365\du}}
\pgfpathlineto{\pgfpoint{-5.403545\du}{12.518365\du}}
\pgfpathlineto{\pgfpoint{-5.395625\du}{12.517666\du}}
\pgfpathlineto{\pgfpoint{-5.387937\du}{12.517666\du}}
\pgfpathlineto{\pgfpoint{-5.380250\du}{12.516967\du}}
\pgfpathlineto{\pgfpoint{-5.372563\du}{12.516734\du}}
\pgfpathlineto{\pgfpoint{-5.365341\du}{12.515802\du}}
\pgfpathlineto{\pgfpoint{-5.358120\du}{12.515103\du}}
\pgfpathlineto{\pgfpoint{-5.350666\du}{12.513939\du}}
\pgfpathlineto{\pgfpoint{-5.343910\du}{12.513007\du}}
\pgfpathlineto{\pgfpoint{-5.336689\du}{12.511842\du}}
\pgfpathlineto{\pgfpoint{-5.330399\du}{12.510677\du}}
\pgfpathlineto{\pgfpoint{-5.323877\du}{12.509047\du}}
\pgfpathlineto{\pgfpoint{-5.317354\du}{12.507416\du}}
\pgfpathlineto{\pgfpoint{-5.311530\du}{12.505785\du}}
\pgfpathlineto{\pgfpoint{-5.305707\du}{12.504388\du}}
\pgfpathlineto{\pgfpoint{-5.300116\du}{12.502524\du}}
\pgfpathlineto{\pgfpoint{-5.294991\du}{12.500195\du}}
\pgfpathlineto{\pgfpoint{-5.290099\du}{12.498564\du}}
\pgfpathlineto{\pgfpoint{-5.285673\du}{12.496468\du}}
\pgfpathlineto{\pgfpoint{-5.281247\du}{12.494371\du}}
\pgfpathlineto{\pgfpoint{-5.277287\du}{12.491809\du}}
\pgfpathlineto{\pgfpoint{-5.273560\du}{12.489246\du}}
\pgfpathlineto{\pgfpoint{-5.270299\du}{12.487383\du}}
\pgfpathlineto{\pgfpoint{-5.267503\du}{12.484820\du}}
\pgfpathlineto{\pgfpoint{-5.265174\du}{12.482491\du}}
\pgfpathlineto{\pgfpoint{-5.263077\du}{12.480161\du}}
\pgfpathlineto{\pgfpoint{-5.261214\du}{12.477366\du}}
\pgfpathlineto{\pgfpoint{-5.259816\du}{12.474803\du}}
\pgfpathlineto{\pgfpoint{-5.259117\du}{12.472474\du}}
\pgfpathlineto{\pgfpoint{-5.258418\du}{12.469679\du}}
\pgfpathlineto{\pgfpoint{-5.257953\du}{12.467116\du}}
\pgfusepath{fill}
\pgfsetbuttcap
\pgfsetmiterjoin
\pgfsetdash{}{0pt}
\definecolor{dialinecolor}{rgb}{0.666667, 0.901961, 1.000000}
\pgfsetstrokecolor{dialinecolor}
\pgfpathmoveto{\pgfpoint{-5.267503\du}{12.462690\du}}
\pgfpathlineto{\pgfpoint{-5.267503\du}{12.460128\du}}
\pgfpathlineto{\pgfpoint{-5.268202\du}{12.458264\du}}
\pgfpathlineto{\pgfpoint{-5.269134\du}{12.455702\du}}
\pgfpathlineto{\pgfpoint{-5.270299\du}{12.452906\du}}
\pgfpathlineto{\pgfpoint{-5.272162\du}{12.451043\du}}
\pgfpathlineto{\pgfpoint{-5.274026\du}{12.448480\du}}
\pgfpathlineto{\pgfpoint{-5.276588\du}{12.446617\du}}
\pgfpathlineto{\pgfpoint{-5.279151\du}{12.444520\du}}
\pgfpathlineto{\pgfpoint{-5.282412\du}{12.442424\du}}
\pgfpathlineto{\pgfpoint{-5.285906\du}{12.440327\du}}
\pgfpathlineto{\pgfpoint{-5.289633\du}{12.438231\du}}
\pgfpathlineto{\pgfpoint{-5.294059\du}{12.436134\du}}
\pgfpathlineto{\pgfpoint{-5.298252\du}{12.434270\du}}
\pgfpathlineto{\pgfpoint{-5.303144\du}{12.432640\du}}
\pgfpathlineto{\pgfpoint{-5.307803\du}{12.430776\du}}
\pgfpathlineto{\pgfpoint{-5.313627\du}{12.428913\du}}
\pgfpathlineto{\pgfpoint{-5.318985\du}{12.427515\du}}
\pgfpathlineto{\pgfpoint{-5.324576\du}{12.426350\du}}
\pgfpathlineto{\pgfpoint{-5.330632\du}{12.424487\du}}
\pgfpathlineto{\pgfpoint{-5.337388\du}{12.423322\du}}
\pgfpathlineto{\pgfpoint{-5.343910\du}{12.422390\du}}
\pgfpathlineto{\pgfpoint{-5.350433\du}{12.420992\du}}
\pgfpathlineto{\pgfpoint{-5.357421\du}{12.420294\du}}
\pgfpathlineto{\pgfpoint{-5.364177\du}{12.419129\du}}
\pgfpathlineto{\pgfpoint{-5.371165\du}{12.418430\du}}
\pgfpathlineto{\pgfpoint{-5.378386\du}{12.417964\du}}
\pgfpathlineto{\pgfpoint{-5.385841\du}{12.417265\du}}
\pgfpathlineto{\pgfpoint{-5.393295\du}{12.417032\du}}
\pgfpathlineto{\pgfpoint{-5.400284\du}{12.416799\du}}
\pgfpathlineto{\pgfpoint{-5.408204\du}{12.416334\du}}
\pgfpathlineto{\pgfpoint{-5.415658\du}{12.416334\du}}
\pgfpathlineto{\pgfpoint{-5.415658\du}{12.416334\du}}
\pgfpathlineto{\pgfpoint{-5.422879\du}{12.416334\du}}
\pgfpathlineto{\pgfpoint{-5.430567\du}{12.416799\du}}
\pgfpathlineto{\pgfpoint{-5.438021\du}{12.417032\du}}
\pgfpathlineto{\pgfpoint{-5.445475\du}{12.417265\du}}
\pgfpathlineto{\pgfpoint{-5.452697\du}{12.417964\du}}
\pgfpathlineto{\pgfpoint{-5.460384\du}{12.418430\du}}
\pgfpathlineto{\pgfpoint{-5.467140\du}{12.419129\du}}
\pgfpathlineto{\pgfpoint{-5.474128\du}{12.420294\du}}
\pgfpathlineto{\pgfpoint{-5.480650\du}{12.420992\du}}
\pgfpathlineto{\pgfpoint{-5.487639\du}{12.422390\du}}
\pgfpathlineto{\pgfpoint{-5.494161\du}{12.423322\du}}
\pgfpathlineto{\pgfpoint{-5.500451\du}{12.424487\du}}
\pgfpathlineto{\pgfpoint{-5.506508\du}{12.426350\du}}
\pgfpathlineto{\pgfpoint{-5.512331\du}{12.427515\du}}
\pgfpathlineto{\pgfpoint{-5.517689\du}{12.428913\du}}
\pgfpathlineto{\pgfpoint{-5.523513\du}{12.430776\du}}
\pgfpathlineto{\pgfpoint{-5.528172\du}{12.432640\du}}
\pgfpathlineto{\pgfpoint{-5.533064\du}{12.434270\du}}
\pgfpathlineto{\pgfpoint{-5.537490\du}{12.436134\du}}
\pgfpathlineto{\pgfpoint{-5.541683\du}{12.438231\du}}
\pgfpathlineto{\pgfpoint{-5.545410\du}{12.440327\du}}
\pgfpathlineto{\pgfpoint{-5.548904\du}{12.442424\du}}
\pgfpathlineto{\pgfpoint{-5.552165\du}{12.444520\du}}
\pgfpathlineto{\pgfpoint{-5.554728\du}{12.446617\du}}
\pgfpathlineto{\pgfpoint{-5.557523\du}{12.448480\du}}
\pgfpathlineto{\pgfpoint{-5.558921\du}{12.451043\du}}
\pgfpathlineto{\pgfpoint{-5.560784\du}{12.452906\du}}
\pgfpathlineto{\pgfpoint{-5.562182\du}{12.455702\du}}
\pgfpathlineto{\pgfpoint{-5.563347\du}{12.458264\du}}
\pgfpathlineto{\pgfpoint{-5.563813\du}{12.460128\du}}
\pgfpathlineto{\pgfpoint{-5.564046\du}{12.462690\du}}
\pgfpathlineto{\pgfpoint{-5.564046\du}{12.462690\du}}
\pgfpathlineto{\pgfpoint{-5.563813\du}{12.465253\du}}
\pgfpathlineto{\pgfpoint{-5.563347\du}{12.467349\du}}
\pgfpathlineto{\pgfpoint{-5.562182\du}{12.469679\du}}
\pgfpathlineto{\pgfpoint{-5.560784\du}{12.472241\du}}
\pgfpathlineto{\pgfpoint{-5.558921\du}{12.474337\du}}
\pgfpathlineto{\pgfpoint{-5.557523\du}{12.476900\du}}
\pgfpathlineto{\pgfpoint{-5.554728\du}{12.478996\du}}
\pgfpathlineto{\pgfpoint{-5.552165\du}{12.481093\du}}
\pgfpathlineto{\pgfpoint{-5.548904\du}{12.483189\du}}
\pgfpathlineto{\pgfpoint{-5.545410\du}{12.485286\du}}
\pgfpathlineto{\pgfpoint{-5.541683\du}{12.487383\du}}
\pgfpathlineto{\pgfpoint{-5.537490\du}{12.489013\du}}
\pgfpathlineto{\pgfpoint{-5.533064\du}{12.491110\du}}
\pgfpathlineto{\pgfpoint{-5.528172\du}{12.492973\du}}
\pgfpathlineto{\pgfpoint{-5.523513\du}{12.494837\du}}
\pgfpathlineto{\pgfpoint{-5.517689\du}{12.496468\du}}
\pgfpathlineto{\pgfpoint{-5.512331\du}{12.497865\du}}
\pgfpathlineto{\pgfpoint{-5.506508\du}{12.499263\du}}
\pgfpathlineto{\pgfpoint{-5.500451\du}{12.500661\du}}
\pgfpathlineto{\pgfpoint{-5.494161\du}{12.502291\du}}
\pgfpathlineto{\pgfpoint{-5.487639\du}{12.503223\du}}
\pgfpathlineto{\pgfpoint{-5.480650\du}{12.504621\du}}
\pgfpathlineto{\pgfpoint{-5.474128\du}{12.505320\du}}
\pgfpathlineto{\pgfpoint{-5.467140\du}{12.506484\du}}
\pgfpathlineto{\pgfpoint{-5.460384\du}{12.506717\du}}
\pgfpathlineto{\pgfpoint{-5.452697\du}{12.507649\du}}
\pgfpathlineto{\pgfpoint{-5.445475\du}{12.507882\du}}
\pgfpathlineto{\pgfpoint{-5.438021\du}{12.508581\du}}
\pgfpathlineto{\pgfpoint{-5.430567\du}{12.508814\du}}
\pgfpathlineto{\pgfpoint{-5.422879\du}{12.509047\du}}
\pgfpathlineto{\pgfpoint{-5.415658\du}{12.509047\du}}
\pgfpathlineto{\pgfpoint{-5.415658\du}{12.509047\du}}
\pgfpathlineto{\pgfpoint{-5.408204\du}{12.509047\du}}
\pgfpathlineto{\pgfpoint{-5.400284\du}{12.508814\du}}
\pgfpathlineto{\pgfpoint{-5.393295\du}{12.508581\du}}
\pgfpathlineto{\pgfpoint{-5.385841\du}{12.507882\du}}
\pgfpathlineto{\pgfpoint{-5.378386\du}{12.507649\du}}
\pgfpathlineto{\pgfpoint{-5.371165\du}{12.506717\du}}
\pgfpathlineto{\pgfpoint{-5.364177\du}{12.506484\du}}
\pgfpathlineto{\pgfpoint{-5.357421\du}{12.505320\du}}
\pgfpathlineto{\pgfpoint{-5.350433\du}{12.504621\du}}
\pgfpathlineto{\pgfpoint{-5.343910\du}{12.503223\du}}
\pgfpathlineto{\pgfpoint{-5.337388\du}{12.502291\du}}
\pgfpathlineto{\pgfpoint{-5.330632\du}{12.500661\du}}
\pgfpathlineto{\pgfpoint{-5.324576\du}{12.499263\du}}
\pgfpathlineto{\pgfpoint{-5.318985\du}{12.497865\du}}
\pgfpathlineto{\pgfpoint{-5.313627\du}{12.496468\du}}
\pgfpathlineto{\pgfpoint{-5.307803\du}{12.494837\du}}
\pgfpathlineto{\pgfpoint{-5.303144\du}{12.492973\du}}
\pgfpathlineto{\pgfpoint{-5.298252\du}{12.491110\du}}
\pgfpathlineto{\pgfpoint{-5.294059\du}{12.489013\du}}
\pgfpathlineto{\pgfpoint{-5.289633\du}{12.487383\du}}
\pgfpathlineto{\pgfpoint{-5.285906\du}{12.485286\du}}
\pgfpathlineto{\pgfpoint{-5.282412\du}{12.483189\du}}
\pgfpathlineto{\pgfpoint{-5.279151\du}{12.481093\du}}
\pgfpathlineto{\pgfpoint{-5.276588\du}{12.478996\du}}
\pgfpathlineto{\pgfpoint{-5.274026\du}{12.476900\du}}
\pgfpathlineto{\pgfpoint{-5.272162\du}{12.474337\du}}
\pgfpathlineto{\pgfpoint{-5.270299\du}{12.472241\du}}
\pgfpathlineto{\pgfpoint{-5.269134\du}{12.469679\du}}
\pgfpathlineto{\pgfpoint{-5.268202\du}{12.467349\du}}
\pgfpathlineto{\pgfpoint{-5.267503\du}{12.465253\du}}
\pgfpathlineto{\pgfpoint{-5.267503\du}{12.462690\du}}
\pgfusepath{stroke}
\pgfsetbuttcap
\pgfsetmiterjoin
\pgfsetdash{}{0pt}
\definecolor{dialinecolor}{rgb}{0.000000, 0.000000, 0.000000}
\pgfsetfillcolor{dialinecolor}
\pgfpathmoveto{\pgfpoint{-5.406107\du}{12.453605\du}}
\pgfpathlineto{\pgfpoint{-5.387937\du}{12.462690\du}}
\pgfpathlineto{\pgfpoint{-5.331797\du}{12.444287\du}}
\pgfpathlineto{\pgfpoint{-5.304076\du}{12.453605\du}}
\pgfpathlineto{\pgfpoint{-5.322479\du}{12.434969\du}}
\pgfpathlineto{\pgfpoint{-5.387937\du}{12.434969\du}}
\pgfpathlineto{\pgfpoint{-5.359984\du}{12.434969\du}}
\pgfpathlineto{\pgfpoint{-5.406107\du}{12.453605\du}}
\pgfusepath{fill}
\pgfsetbuttcap
\pgfsetmiterjoin
\pgfsetdash{}{0pt}
\definecolor{dialinecolor}{rgb}{0.000000, 0.000000, 0.000000}
\pgfsetfillcolor{dialinecolor}
\pgfpathmoveto{\pgfpoint{-5.415658\du}{12.481326\du}}
\pgfpathlineto{\pgfpoint{-5.443612\du}{12.472241\du}}
\pgfpathlineto{\pgfpoint{-5.489735\du}{12.490877\du}}
\pgfpathlineto{\pgfpoint{-5.517689\du}{12.481326\du}}
\pgfpathlineto{\pgfpoint{-5.508138\du}{12.509047\du}}
\pgfpathlineto{\pgfpoint{-5.443612\du}{12.509047\du}}
\pgfpathlineto{\pgfpoint{-5.471333\du}{12.500195\du}}
\pgfpathlineto{\pgfpoint{-5.415658\du}{12.481326\du}}
\pgfusepath{fill}
\pgfsetbuttcap
\pgfsetmiterjoin
\pgfsetdash{}{0pt}
\definecolor{dialinecolor}{rgb}{0.000000, 0.000000, 0.000000}
\pgfsetfillcolor{dialinecolor}
\pgfpathmoveto{\pgfpoint{-5.508138\du}{12.434969\du}}
\pgfpathlineto{\pgfpoint{-5.489735\du}{12.425186\du}}
\pgfpathlineto{\pgfpoint{-5.433828\du}{12.453605\du}}
\pgfpathlineto{\pgfpoint{-5.415658\du}{12.444287\du}}
\pgfpathlineto{\pgfpoint{-5.424743\du}{12.462690\du}}
\pgfpathlineto{\pgfpoint{-5.489735\du}{12.462690\du}}
\pgfpathlineto{\pgfpoint{-5.461782\du}{12.453605\du}}
\pgfpathlineto{\pgfpoint{-5.508138\du}{12.434969\du}}
\pgfusepath{fill}
\pgfsetbuttcap
\pgfsetmiterjoin
\pgfsetdash{}{0pt}
\definecolor{dialinecolor}{rgb}{0.000000, 0.000000, 0.000000}
\pgfsetfillcolor{dialinecolor}
\pgfpathmoveto{\pgfpoint{-5.313627\du}{12.500195\du}}
\pgfpathlineto{\pgfpoint{-5.331797\du}{12.509047\du}}
\pgfpathlineto{\pgfpoint{-5.387937\du}{12.481326\du}}
\pgfpathlineto{\pgfpoint{-5.415658\du}{12.490877\du}}
\pgfpathlineto{\pgfpoint{-5.397022\du}{12.472241\du}}
\pgfpathlineto{\pgfpoint{-5.331797\du}{12.472241\du}}
\pgfpathlineto{\pgfpoint{-5.359984\du}{12.481326\du}}
\pgfpathlineto{\pgfpoint{-5.313627\du}{12.500195\du}}
\pgfusepath{fill}
\pgfsetbuttcap
\pgfsetmiterjoin
\pgfsetdash{}{0pt}
\definecolor{dialinecolor}{rgb}{1.000000, 1.000000, 1.000000}
\pgfsetfillcolor{dialinecolor}
\pgfpathmoveto{\pgfpoint{-5.406107\du}{12.462690\du}}
\pgfpathlineto{\pgfpoint{-5.387937\du}{12.462690\du}}
\pgfpathlineto{\pgfpoint{-5.331797\du}{12.444287\du}}
\pgfpathlineto{\pgfpoint{-5.304076\du}{12.453605\du}}
\pgfpathlineto{\pgfpoint{-5.322479\du}{12.434969\du}}
\pgfpathlineto{\pgfpoint{-5.378386\du}{12.434969\du}}
\pgfpathlineto{\pgfpoint{-5.359984\du}{12.434969\du}}
\pgfpathlineto{\pgfpoint{-5.406107\du}{12.462690\du}}
\pgfusepath{fill}
\pgfsetbuttcap
\pgfsetmiterjoin
\pgfsetdash{}{0pt}
\definecolor{dialinecolor}{rgb}{1.000000, 1.000000, 1.000000}
\pgfsetfillcolor{dialinecolor}
\pgfpathmoveto{\pgfpoint{-5.415658\du}{12.481326\du}}
\pgfpathlineto{\pgfpoint{-5.433828\du}{12.472241\du}}
\pgfpathlineto{\pgfpoint{-5.489735\du}{12.490877\du}}
\pgfpathlineto{\pgfpoint{-5.517689\du}{12.490877\du}}
\pgfpathlineto{\pgfpoint{-5.499286\du}{12.509047\du}}
\pgfpathlineto{\pgfpoint{-5.433828\du}{12.509047\du}}
\pgfpathlineto{\pgfpoint{-5.461782\du}{12.500195\du}}
\pgfpathlineto{\pgfpoint{-5.415658\du}{12.481326\du}}
\pgfusepath{fill}
\pgfsetbuttcap
\pgfsetmiterjoin
\pgfsetdash{}{0pt}
\definecolor{dialinecolor}{rgb}{1.000000, 1.000000, 1.000000}
\pgfsetfillcolor{dialinecolor}
\pgfpathmoveto{\pgfpoint{-5.508138\du}{12.434969\du}}
\pgfpathlineto{\pgfpoint{-5.489735\du}{12.434969\du}}
\pgfpathlineto{\pgfpoint{-5.433828\du}{12.453605\du}}
\pgfpathlineto{\pgfpoint{-5.406107\du}{12.444287\du}}
\pgfpathlineto{\pgfpoint{-5.424743\du}{12.462690\du}}
\pgfpathlineto{\pgfpoint{-5.489735\du}{12.462690\du}}
\pgfpathlineto{\pgfpoint{-5.461782\du}{12.462690\du}}
\pgfpathlineto{\pgfpoint{-5.508138\du}{12.434969\du}}
\pgfusepath{fill}
\pgfsetbuttcap
\pgfsetmiterjoin
\pgfsetdash{}{0pt}
\definecolor{dialinecolor}{rgb}{1.000000, 1.000000, 1.000000}
\pgfsetfillcolor{dialinecolor}
\pgfpathmoveto{\pgfpoint{-5.313627\du}{12.500195\du}}
\pgfpathlineto{\pgfpoint{-5.331797\du}{12.509047\du}}
\pgfpathlineto{\pgfpoint{-5.387937\du}{12.490877\du}}
\pgfpathlineto{\pgfpoint{-5.406107\du}{12.490877\du}}
\pgfpathlineto{\pgfpoint{-5.397022\du}{12.472241\du}}
\pgfpathlineto{\pgfpoint{-5.331797\du}{12.472241\du}}
\pgfpathlineto{\pgfpoint{-5.359984\du}{12.481326\du}}
\pgfpathlineto{\pgfpoint{-5.313627\du}{12.500195\du}}
\pgfusepath{fill}
\pgfsetbuttcap
\pgfsetmiterjoin
\pgfsetdash{}{0pt}
\definecolor{dialinecolor}{rgb}{0.666667, 0.901961, 1.000000}
\pgfsetstrokecolor{dialinecolor}
\pgfpathmoveto{\pgfpoint{-5.564046\du}{12.462690\du}}
\pgfpathlineto{\pgfpoint{-5.564046\du}{12.536767\du}}
\pgfusepath{stroke}
\pgfsetbuttcap
\pgfsetmiterjoin
\pgfsetdash{}{0pt}
\definecolor{dialinecolor}{rgb}{0.666667, 0.901961, 1.000000}
\pgfsetstrokecolor{dialinecolor}
\pgfpathmoveto{\pgfpoint{-5.267038\du}{12.462690\du}}
\pgfpathlineto{\pgfpoint{-5.267038\du}{12.536767\du}}
\pgfusepath{stroke}
\pgfsetbuttcap
\pgfsetmiterjoin
\pgfsetdash{}{0pt}
\definecolor{dialinecolor}{rgb}{0.000000, 0.470588, 0.666667}
\pgfsetfillcolor{dialinecolor}
\pgfpathmoveto{\pgfpoint{-4.700742\du}{12.541659\du}}
\pgfpathlineto{\pgfpoint{-4.700742\du}{12.539097\du}}
\pgfpathlineto{\pgfpoint{-4.701441\du}{12.536534\du}}
\pgfpathlineto{\pgfpoint{-4.702373\du}{12.533972\du}}
\pgfpathlineto{\pgfpoint{-4.703537\du}{12.531410\du}}
\pgfpathlineto{\pgfpoint{-4.705401\du}{12.529080\du}}
\pgfpathlineto{\pgfpoint{-4.707497\du}{12.526518\du}}
\pgfpathlineto{\pgfpoint{-4.709827\du}{12.523955\du}}
\pgfpathlineto{\pgfpoint{-4.713088\du}{12.521626\du}}
\pgfpathlineto{\pgfpoint{-4.715883\du}{12.519296\du}}
\pgfpathlineto{\pgfpoint{-4.719611\du}{12.516967\du}}
\pgfpathlineto{\pgfpoint{-4.723571\du}{12.514870\du}}
\pgfpathlineto{\pgfpoint{-4.727764\du}{12.512541\du}}
\pgfpathlineto{\pgfpoint{-4.732656\du}{12.510677\du}}
\pgfpathlineto{\pgfpoint{-4.737315\du}{12.508581\du}}
\pgfpathlineto{\pgfpoint{-4.742672\du}{12.506717\du}}
\pgfpathlineto{\pgfpoint{-4.748030\du}{12.504854\du}}
\pgfpathlineto{\pgfpoint{-4.753854\du}{12.503223\du}}
\pgfpathlineto{\pgfpoint{-4.759678\du}{12.501359\du}}
\pgfpathlineto{\pgfpoint{-4.765967\du}{12.500195\du}}
\pgfpathlineto{\pgfpoint{-4.772723\du}{12.498564\du}}
\pgfpathlineto{\pgfpoint{-4.779245\du}{12.497166\du}}
\pgfpathlineto{\pgfpoint{-4.786001\du}{12.496235\du}}
\pgfpathlineto{\pgfpoint{-4.793222\du}{12.495070\du}}
\pgfpathlineto{\pgfpoint{-4.800443\du}{12.494138\du}}
\pgfpathlineto{\pgfpoint{-4.807665\du}{12.493206\du}}
\pgfpathlineto{\pgfpoint{-4.815119\du}{12.492507\du}}
\pgfpathlineto{\pgfpoint{-4.823039\du}{12.491809\du}}
\pgfpathlineto{\pgfpoint{-4.830494\du}{12.491343\du}}
\pgfpathlineto{\pgfpoint{-4.837948\du}{12.491110\du}}
\pgfpathlineto{\pgfpoint{-4.845868\du}{12.490877\du}}
\pgfpathlineto{\pgfpoint{-4.853556\du}{12.490877\du}}
\pgfpathlineto{\pgfpoint{-4.853556\du}{12.490877\du}}
\pgfpathlineto{\pgfpoint{-4.861243\du}{12.490877\du}}
\pgfpathlineto{\pgfpoint{-4.869163\du}{12.491110\du}}
\pgfpathlineto{\pgfpoint{-4.876850\du}{12.491343\du}}
\pgfpathlineto{\pgfpoint{-4.884538\du}{12.491809\du}}
\pgfpathlineto{\pgfpoint{-4.891992\du}{12.492507\du}}
\pgfpathlineto{\pgfpoint{-4.899446\du}{12.493206\du}}
\pgfpathlineto{\pgfpoint{-4.906901\du}{12.494138\du}}
\pgfpathlineto{\pgfpoint{-4.913889\du}{12.495070\du}}
\pgfpathlineto{\pgfpoint{-4.921110\du}{12.496235\du}}
\pgfpathlineto{\pgfpoint{-4.927866\du}{12.497166\du}}
\pgfpathlineto{\pgfpoint{-4.934621\du}{12.498564\du}}
\pgfpathlineto{\pgfpoint{-4.941144\du}{12.500195\du}}
\pgfpathlineto{\pgfpoint{-4.947433\du}{12.501359\du}}
\pgfpathlineto{\pgfpoint{-4.953257\du}{12.503223\du}}
\pgfpathlineto{\pgfpoint{-4.959081\du}{12.504854\du}}
\pgfpathlineto{\pgfpoint{-4.964672\du}{12.506717\du}}
\pgfpathlineto{\pgfpoint{-4.969796\du}{12.508581\du}}
\pgfpathlineto{\pgfpoint{-4.974688\du}{12.510677\du}}
\pgfpathlineto{\pgfpoint{-4.979347\du}{12.512541\du}}
\pgfpathlineto{\pgfpoint{-4.983540\du}{12.514870\du}}
\pgfpathlineto{\pgfpoint{-4.987500\du}{12.516967\du}}
\pgfpathlineto{\pgfpoint{-4.991228\du}{12.519296\du}}
\pgfpathlineto{\pgfpoint{-4.994256\du}{12.521626\du}}
\pgfpathlineto{\pgfpoint{-4.997284\du}{12.523955\du}}
\pgfpathlineto{\pgfpoint{-4.999614\du}{12.526518\du}}
\pgfpathlineto{\pgfpoint{-5.001710\du}{12.529080\du}}
\pgfpathlineto{\pgfpoint{-5.003574\du}{12.531410\du}}
\pgfpathlineto{\pgfpoint{-5.004971\du}{12.533972\du}}
\pgfpathlineto{\pgfpoint{-5.006136\du}{12.536534\du}}
\pgfpathlineto{\pgfpoint{-5.006835\du}{12.539097\du}}
\pgfpathlineto{\pgfpoint{-5.006835\du}{12.541659\du}}
\pgfpathlineto{\pgfpoint{-5.006835\du}{12.541659\du}}
\pgfpathlineto{\pgfpoint{-5.006835\du}{12.544455\du}}
\pgfpathlineto{\pgfpoint{-5.006136\du}{12.546784\du}}
\pgfpathlineto{\pgfpoint{-5.004971\du}{12.549347\du}}
\pgfpathlineto{\pgfpoint{-5.003574\du}{12.551909\du}}
\pgfpathlineto{\pgfpoint{-5.001710\du}{12.554471\du}}
\pgfpathlineto{\pgfpoint{-4.999614\du}{12.557034\du}}
\pgfpathlineto{\pgfpoint{-4.997284\du}{12.559363\du}}
\pgfpathlineto{\pgfpoint{-4.994256\du}{12.561693\du}}
\pgfpathlineto{\pgfpoint{-4.991228\du}{12.564022\du}}
\pgfpathlineto{\pgfpoint{-4.987500\du}{12.566585\du}}
\pgfpathlineto{\pgfpoint{-4.983540\du}{12.568681\du}}
\pgfpathlineto{\pgfpoint{-4.979347\du}{12.570778\du}}
\pgfpathlineto{\pgfpoint{-4.974688\du}{12.572874\du}}
\pgfpathlineto{\pgfpoint{-4.969796\du}{12.574971\du}}
\pgfpathlineto{\pgfpoint{-4.964672\du}{12.576834\du}}
\pgfpathlineto{\pgfpoint{-4.959081\du}{12.578698\du}}
\pgfpathlineto{\pgfpoint{-4.953257\du}{12.580096\du}}
\pgfpathlineto{\pgfpoint{-4.947433\du}{12.581726\du}}
\pgfpathlineto{\pgfpoint{-4.941144\du}{12.583357\du}}
\pgfpathlineto{\pgfpoint{-4.934621\du}{12.584988\du}}
\pgfpathlineto{\pgfpoint{-4.927866\du}{12.586152\du}}
\pgfpathlineto{\pgfpoint{-4.921110\du}{12.587317\du}}
\pgfpathlineto{\pgfpoint{-4.913889\du}{12.588482\du}}
\pgfpathlineto{\pgfpoint{-4.906901\du}{12.589414\du}}
\pgfpathlineto{\pgfpoint{-4.899446\du}{12.590112\du}}
\pgfpathlineto{\pgfpoint{-4.891992\du}{12.591044\du}}
\pgfpathlineto{\pgfpoint{-4.884538\du}{12.591510\du}}
\pgfpathlineto{\pgfpoint{-4.876850\du}{12.591976\du}}
\pgfpathlineto{\pgfpoint{-4.869163\du}{12.592442\du}}
\pgfpathlineto{\pgfpoint{-4.861243\du}{12.592675\du}}
\pgfpathlineto{\pgfpoint{-4.853556\du}{12.592675\du}}
\pgfpathlineto{\pgfpoint{-4.853556\du}{12.592675\du}}
\pgfpathlineto{\pgfpoint{-4.845868\du}{12.592675\du}}
\pgfpathlineto{\pgfpoint{-4.837948\du}{12.592442\du}}
\pgfpathlineto{\pgfpoint{-4.830494\du}{12.591976\du}}
\pgfpathlineto{\pgfpoint{-4.823039\du}{12.591510\du}}
\pgfpathlineto{\pgfpoint{-4.815119\du}{12.591044\du}}
\pgfpathlineto{\pgfpoint{-4.807665\du}{12.590112\du}}
\pgfpathlineto{\pgfpoint{-4.800443\du}{12.589414\du}}
\pgfpathlineto{\pgfpoint{-4.793222\du}{12.588482\du}}
\pgfpathlineto{\pgfpoint{-4.786001\du}{12.587317\du}}
\pgfpathlineto{\pgfpoint{-4.779245\du}{12.586152\du}}
\pgfpathlineto{\pgfpoint{-4.772723\du}{12.584988\du}}
\pgfpathlineto{\pgfpoint{-4.765967\du}{12.583357\du}}
\pgfpathlineto{\pgfpoint{-4.759678\du}{12.581726\du}}
\pgfpathlineto{\pgfpoint{-4.753854\du}{12.580096\du}}
\pgfpathlineto{\pgfpoint{-4.748030\du}{12.578698\du}}
\pgfpathlineto{\pgfpoint{-4.742672\du}{12.576834\du}}
\pgfpathlineto{\pgfpoint{-4.737315\du}{12.574971\du}}
\pgfpathlineto{\pgfpoint{-4.732656\du}{12.572874\du}}
\pgfpathlineto{\pgfpoint{-4.727764\du}{12.570778\du}}
\pgfpathlineto{\pgfpoint{-4.723571\du}{12.568681\du}}
\pgfpathlineto{\pgfpoint{-4.719611\du}{12.566585\du}}
\pgfpathlineto{\pgfpoint{-4.715883\du}{12.564022\du}}
\pgfpathlineto{\pgfpoint{-4.713088\du}{12.561693\du}}
\pgfpathlineto{\pgfpoint{-4.709827\du}{12.559363\du}}
\pgfpathlineto{\pgfpoint{-4.707497\du}{12.557034\du}}
\pgfpathlineto{\pgfpoint{-4.705401\du}{12.554471\du}}
\pgfpathlineto{\pgfpoint{-4.703537\du}{12.551909\du}}
\pgfpathlineto{\pgfpoint{-4.702373\du}{12.549347\du}}
\pgfpathlineto{\pgfpoint{-4.701441\du}{12.546784\du}}
\pgfpathlineto{\pgfpoint{-4.700742\du}{12.544455\du}}
\pgfpathlineto{\pgfpoint{-4.700742\du}{12.541659\du}}
\pgfusepath{fill}
\pgfsetbuttcap
\pgfsetmiterjoin
\pgfsetdash{}{0pt}
\definecolor{dialinecolor}{rgb}{0.666667, 0.901961, 1.000000}
\pgfsetstrokecolor{dialinecolor}
\pgfpathmoveto{\pgfpoint{-4.709827\du}{12.537233\du}}
\pgfpathlineto{\pgfpoint{-4.709827\du}{12.534438\du}}
\pgfpathlineto{\pgfpoint{-4.710526\du}{12.532574\du}}
\pgfpathlineto{\pgfpoint{-4.711457\du}{12.530012\du}}
\pgfpathlineto{\pgfpoint{-4.713088\du}{12.527682\du}}
\pgfpathlineto{\pgfpoint{-4.714719\du}{12.525586\du}}
\pgfpathlineto{\pgfpoint{-4.716349\du}{12.522791\du}}
\pgfpathlineto{\pgfpoint{-4.718912\du}{12.520927\du}}
\pgfpathlineto{\pgfpoint{-4.721474\du}{12.518830\du}}
\pgfpathlineto{\pgfpoint{-4.724735\du}{12.516734\du}}
\pgfpathlineto{\pgfpoint{-4.728463\du}{12.514870\du}}
\pgfpathlineto{\pgfpoint{-4.731957\du}{12.512541\du}}
\pgfpathlineto{\pgfpoint{-4.736383\du}{12.510910\du}}
\pgfpathlineto{\pgfpoint{-4.741042\du}{12.508814\du}}
\pgfpathlineto{\pgfpoint{-4.745468\du}{12.506717\du}}
\pgfpathlineto{\pgfpoint{-4.750593\du}{12.505320\du}}
\pgfpathlineto{\pgfpoint{-4.755950\du}{12.503456\du}}
\pgfpathlineto{\pgfpoint{-4.761308\du}{12.502291\du}}
\pgfpathlineto{\pgfpoint{-4.767365\du}{12.500661\du}}
\pgfpathlineto{\pgfpoint{-4.773189\du}{12.499263\du}}
\pgfpathlineto{\pgfpoint{-4.779711\du}{12.497865\du}}
\pgfpathlineto{\pgfpoint{-4.786467\du}{12.496933\du}}
\pgfpathlineto{\pgfpoint{-4.792989\du}{12.495070\du}}
\pgfpathlineto{\pgfpoint{-4.799745\du}{12.494837\du}}
\pgfpathlineto{\pgfpoint{-4.806733\du}{12.493672\du}}
\pgfpathlineto{\pgfpoint{-4.813489\du}{12.492973\du}}
\pgfpathlineto{\pgfpoint{-4.820943\du}{12.492274\du}}
\pgfpathlineto{\pgfpoint{-4.828164\du}{12.491809\du}}
\pgfpathlineto{\pgfpoint{-4.835619\du}{12.491576\du}}
\pgfpathlineto{\pgfpoint{-4.843073\du}{12.491110\du}}
\pgfpathlineto{\pgfpoint{-4.850993\du}{12.490877\du}}
\pgfpathlineto{\pgfpoint{-4.858680\du}{12.490877\du}}
\pgfpathlineto{\pgfpoint{-4.858680\du}{12.490877\du}}
\pgfpathlineto{\pgfpoint{-4.865669\du}{12.490877\du}}
\pgfpathlineto{\pgfpoint{-4.873356\du}{12.491110\du}}
\pgfpathlineto{\pgfpoint{-4.880810\du}{12.491576\du}}
\pgfpathlineto{\pgfpoint{-4.888032\du}{12.491809\du}}
\pgfpathlineto{\pgfpoint{-4.895486\du}{12.492274\du}}
\pgfpathlineto{\pgfpoint{-4.902940\du}{12.492973\du}}
\pgfpathlineto{\pgfpoint{-4.909929\du}{12.493672\du}}
\pgfpathlineto{\pgfpoint{-4.916917\du}{12.494837\du}}
\pgfpathlineto{\pgfpoint{-4.923673\du}{12.495070\du}}
\pgfpathlineto{\pgfpoint{-4.930195\du}{12.496933\du}}
\pgfpathlineto{\pgfpoint{-4.936951\du}{12.497865\du}}
\pgfpathlineto{\pgfpoint{-4.943240\du}{12.499263\du}}
\pgfpathlineto{\pgfpoint{-4.949297\du}{12.500661\du}}
\pgfpathlineto{\pgfpoint{-4.955121\du}{12.502291\du}}
\pgfpathlineto{\pgfpoint{-4.960246\du}{12.503456\du}}
\pgfpathlineto{\pgfpoint{-4.966069\du}{12.505320\du}}
\pgfpathlineto{\pgfpoint{-4.970961\du}{12.506717\du}}
\pgfpathlineto{\pgfpoint{-4.975620\du}{12.508814\du}}
\pgfpathlineto{\pgfpoint{-4.980279\du}{12.510910\du}}
\pgfpathlineto{\pgfpoint{-4.984239\du}{12.512541\du}}
\pgfpathlineto{\pgfpoint{-4.988199\du}{12.514870\du}}
\pgfpathlineto{\pgfpoint{-4.991926\du}{12.516734\du}}
\pgfpathlineto{\pgfpoint{-4.995188\du}{12.518830\du}}
\pgfpathlineto{\pgfpoint{-4.997750\du}{12.520927\du}}
\pgfpathlineto{\pgfpoint{-5.000312\du}{12.522791\du}}
\pgfpathlineto{\pgfpoint{-5.002176\du}{12.525586\du}}
\pgfpathlineto{\pgfpoint{-5.003574\du}{12.527682\du}}
\pgfpathlineto{\pgfpoint{-5.005437\du}{12.530012\du}}
\pgfpathlineto{\pgfpoint{-5.006136\du}{12.532574\du}}
\pgfpathlineto{\pgfpoint{-5.006835\du}{12.534438\du}}
\pgfpathlineto{\pgfpoint{-5.006835\du}{12.537233\du}}
\pgfpathlineto{\pgfpoint{-5.006835\du}{12.537233\du}}
\pgfpathlineto{\pgfpoint{-5.006835\du}{12.539563\du}}
\pgfpathlineto{\pgfpoint{-5.006136\du}{12.541659\du}}
\pgfpathlineto{\pgfpoint{-5.005437\du}{12.543989\du}}
\pgfpathlineto{\pgfpoint{-5.003574\du}{12.546784\du}}
\pgfpathlineto{\pgfpoint{-5.002176\du}{12.548881\du}}
\pgfpathlineto{\pgfpoint{-5.000312\du}{12.551210\du}}
\pgfpathlineto{\pgfpoint{-4.997750\du}{12.553307\du}}
\pgfpathlineto{\pgfpoint{-4.995188\du}{12.555403\du}}
\pgfpathlineto{\pgfpoint{-4.991926\du}{12.557500\du}}
\pgfpathlineto{\pgfpoint{-4.988199\du}{12.559596\du}}
\pgfpathlineto{\pgfpoint{-4.984239\du}{12.561693\du}}
\pgfpathlineto{\pgfpoint{-4.980279\du}{12.563556\du}}
\pgfpathlineto{\pgfpoint{-4.975620\du}{12.565653\du}}
\pgfpathlineto{\pgfpoint{-4.970961\du}{12.567284\du}}
\pgfpathlineto{\pgfpoint{-4.966069\du}{12.569147\du}}
\pgfpathlineto{\pgfpoint{-4.960246\du}{12.571011\du}}
\pgfpathlineto{\pgfpoint{-4.955121\du}{12.572408\du}}
\pgfpathlineto{\pgfpoint{-4.949297\du}{12.573573\du}}
\pgfpathlineto{\pgfpoint{-4.943240\du}{12.575204\du}}
\pgfpathlineto{\pgfpoint{-4.936951\du}{12.576601\du}}
\pgfpathlineto{\pgfpoint{-4.930195\du}{12.577533\du}}
\pgfpathlineto{\pgfpoint{-4.923673\du}{12.578931\du}}
\pgfpathlineto{\pgfpoint{-4.916917\du}{12.579630\du}}
\pgfpathlineto{\pgfpoint{-4.909929\du}{12.580795\du}}
\pgfpathlineto{\pgfpoint{-4.902940\du}{12.581493\du}}
\pgfpathlineto{\pgfpoint{-4.895486\du}{12.581726\du}}
\pgfpathlineto{\pgfpoint{-4.888032\du}{12.582658\du}}
\pgfpathlineto{\pgfpoint{-4.880810\du}{12.582891\du}}
\pgfpathlineto{\pgfpoint{-4.873356\du}{12.583124\du}}
\pgfpathlineto{\pgfpoint{-4.865669\du}{12.583590\du}}
\pgfpathlineto{\pgfpoint{-4.858680\du}{12.583590\du}}
\pgfpathlineto{\pgfpoint{-4.858680\du}{12.583590\du}}
\pgfpathlineto{\pgfpoint{-4.850993\du}{12.583590\du}}
\pgfpathlineto{\pgfpoint{-4.843073\du}{12.583124\du}}
\pgfpathlineto{\pgfpoint{-4.835619\du}{12.582891\du}}
\pgfpathlineto{\pgfpoint{-4.828164\du}{12.582658\du}}
\pgfpathlineto{\pgfpoint{-4.820943\du}{12.581726\du}}
\pgfpathlineto{\pgfpoint{-4.813489\du}{12.581493\du}}
\pgfpathlineto{\pgfpoint{-4.806733\du}{12.580795\du}}
\pgfpathlineto{\pgfpoint{-4.799745\du}{12.579630\du}}
\pgfpathlineto{\pgfpoint{-4.792989\du}{12.578931\du}}
\pgfpathlineto{\pgfpoint{-4.786467\du}{12.577533\du}}
\pgfpathlineto{\pgfpoint{-4.779711\du}{12.576601\du}}
\pgfpathlineto{\pgfpoint{-4.773189\du}{12.575204\du}}
\pgfpathlineto{\pgfpoint{-4.767365\du}{12.573573\du}}
\pgfpathlineto{\pgfpoint{-4.761308\du}{12.572408\du}}
\pgfpathlineto{\pgfpoint{-4.755950\du}{12.571011\du}}
\pgfpathlineto{\pgfpoint{-4.750593\du}{12.569147\du}}
\pgfpathlineto{\pgfpoint{-4.745468\du}{12.567284\du}}
\pgfpathlineto{\pgfpoint{-4.741042\du}{12.565653\du}}
\pgfpathlineto{\pgfpoint{-4.736383\du}{12.563556\du}}
\pgfpathlineto{\pgfpoint{-4.731957\du}{12.561693\du}}
\pgfpathlineto{\pgfpoint{-4.728463\du}{12.559596\du}}
\pgfpathlineto{\pgfpoint{-4.724735\du}{12.557500\du}}
\pgfpathlineto{\pgfpoint{-4.721474\du}{12.555403\du}}
\pgfpathlineto{\pgfpoint{-4.718912\du}{12.553307\du}}
\pgfpathlineto{\pgfpoint{-4.716349\du}{12.551210\du}}
\pgfpathlineto{\pgfpoint{-4.714719\du}{12.548881\du}}
\pgfpathlineto{\pgfpoint{-4.713088\du}{12.546784\du}}
\pgfpathlineto{\pgfpoint{-4.711457\du}{12.543989\du}}
\pgfpathlineto{\pgfpoint{-4.710526\du}{12.541659\du}}
\pgfpathlineto{\pgfpoint{-4.709827\du}{12.539563\du}}
\pgfpathlineto{\pgfpoint{-4.709827\du}{12.537233\du}}
\pgfusepath{stroke}
\pgfsetbuttcap
\pgfsetmiterjoin
\pgfsetdash{}{0pt}
\definecolor{dialinecolor}{rgb}{0.000000, 0.470588, 0.666667}
\pgfsetfillcolor{dialinecolor}
\pgfpathmoveto{\pgfpoint{-5.006835\du}{12.472241\du}}
\pgfpathlineto{\pgfpoint{-5.006835\du}{12.546085\du}}
\pgfpathlineto{\pgfpoint{-4.700276\du}{12.546085\du}}
\pgfpathlineto{\pgfpoint{-4.700276\du}{12.472241\du}}
\pgfpathlineto{\pgfpoint{-5.006835\du}{12.472241\du}}
\pgfusepath{fill}
\pgfsetbuttcap
\pgfsetmiterjoin
\pgfsetdash{}{0pt}
\definecolor{dialinecolor}{rgb}{0.000000, 0.705882, 1.000000}
\pgfsetfillcolor{dialinecolor}
\pgfpathmoveto{\pgfpoint{-4.700742\du}{12.467116\du}}
\pgfpathlineto{\pgfpoint{-4.700742\du}{12.464554\du}}
\pgfpathlineto{\pgfpoint{-4.701441\du}{12.461758\du}}
\pgfpathlineto{\pgfpoint{-4.702373\du}{12.459429\du}}
\pgfpathlineto{\pgfpoint{-4.703537\du}{12.456866\du}}
\pgfpathlineto{\pgfpoint{-4.705401\du}{12.454304\du}}
\pgfpathlineto{\pgfpoint{-4.707497\du}{12.451742\du}}
\pgfpathlineto{\pgfpoint{-4.709827\du}{12.449412\du}}
\pgfpathlineto{\pgfpoint{-4.713088\du}{12.446850\du}}
\pgfpathlineto{\pgfpoint{-4.715883\du}{12.444753\du}}
\pgfpathlineto{\pgfpoint{-4.719611\du}{12.442424\du}}
\pgfpathlineto{\pgfpoint{-4.723571\du}{12.440094\du}}
\pgfpathlineto{\pgfpoint{-4.727764\du}{12.437998\du}}
\pgfpathlineto{\pgfpoint{-4.732656\du}{12.435435\du}}
\pgfpathlineto{\pgfpoint{-4.737315\du}{12.434038\du}}
\pgfpathlineto{\pgfpoint{-4.742672\du}{12.431941\du}}
\pgfpathlineto{\pgfpoint{-4.748030\du}{12.430077\du}}
\pgfpathlineto{\pgfpoint{-4.753854\du}{12.428447\du}}
\pgfpathlineto{\pgfpoint{-4.759678\du}{12.426816\du}}
\pgfpathlineto{\pgfpoint{-4.765967\du}{12.425186\du}}
\pgfpathlineto{\pgfpoint{-4.772723\du}{12.423555\du}}
\pgfpathlineto{\pgfpoint{-4.779245\du}{12.422390\du}}
\pgfpathlineto{\pgfpoint{-4.786001\du}{12.421225\du}}
\pgfpathlineto{\pgfpoint{-4.793222\du}{12.420294\du}}
\pgfpathlineto{\pgfpoint{-4.800443\du}{12.419129\du}}
\pgfpathlineto{\pgfpoint{-4.807665\du}{12.418430\du}}
\pgfpathlineto{\pgfpoint{-4.815119\du}{12.417498\du}}
\pgfpathlineto{\pgfpoint{-4.823039\du}{12.417032\du}}
\pgfpathlineto{\pgfpoint{-4.830494\du}{12.416566\du}}
\pgfpathlineto{\pgfpoint{-4.837948\du}{12.416334\du}}
\pgfpathlineto{\pgfpoint{-4.845868\du}{12.416101\du}}
\pgfpathlineto{\pgfpoint{-4.853556\du}{12.416101\du}}
\pgfpathlineto{\pgfpoint{-4.853556\du}{12.416101\du}}
\pgfpathlineto{\pgfpoint{-4.861243\du}{12.416101\du}}
\pgfpathlineto{\pgfpoint{-4.869163\du}{12.416334\du}}
\pgfpathlineto{\pgfpoint{-4.876850\du}{12.416566\du}}
\pgfpathlineto{\pgfpoint{-4.884538\du}{12.417032\du}}
\pgfpathlineto{\pgfpoint{-4.891992\du}{12.417498\du}}
\pgfpathlineto{\pgfpoint{-4.899446\du}{12.418430\du}}
\pgfpathlineto{\pgfpoint{-4.906901\du}{12.419129\du}}
\pgfpathlineto{\pgfpoint{-4.913889\du}{12.420294\du}}
\pgfpathlineto{\pgfpoint{-4.921110\du}{12.421225\du}}
\pgfpathlineto{\pgfpoint{-4.927866\du}{12.422390\du}}
\pgfpathlineto{\pgfpoint{-4.934621\du}{12.423555\du}}
\pgfpathlineto{\pgfpoint{-4.941144\du}{12.425186\du}}
\pgfpathlineto{\pgfpoint{-4.947433\du}{12.426816\du}}
\pgfpathlineto{\pgfpoint{-4.953257\du}{12.428447\du}}
\pgfpathlineto{\pgfpoint{-4.959081\du}{12.430077\du}}
\pgfpathlineto{\pgfpoint{-4.964672\du}{12.431941\du}}
\pgfpathlineto{\pgfpoint{-4.969796\du}{12.434038\du}}
\pgfpathlineto{\pgfpoint{-4.974688\du}{12.435435\du}}
\pgfpathlineto{\pgfpoint{-4.979347\du}{12.437998\du}}
\pgfpathlineto{\pgfpoint{-4.983540\du}{12.440094\du}}
\pgfpathlineto{\pgfpoint{-4.987500\du}{12.442424\du}}
\pgfpathlineto{\pgfpoint{-4.991228\du}{12.444753\du}}
\pgfpathlineto{\pgfpoint{-4.994256\du}{12.446850\du}}
\pgfpathlineto{\pgfpoint{-4.997284\du}{12.449412\du}}
\pgfpathlineto{\pgfpoint{-4.999614\du}{12.451742\du}}
\pgfpathlineto{\pgfpoint{-5.001710\du}{12.454304\du}}
\pgfpathlineto{\pgfpoint{-5.003574\du}{12.456866\du}}
\pgfpathlineto{\pgfpoint{-5.004971\du}{12.459429\du}}
\pgfpathlineto{\pgfpoint{-5.006136\du}{12.461758\du}}
\pgfpathlineto{\pgfpoint{-5.006835\du}{12.464554\du}}
\pgfpathlineto{\pgfpoint{-5.006835\du}{12.467116\du}}
\pgfpathlineto{\pgfpoint{-5.006835\du}{12.467116\du}}
\pgfpathlineto{\pgfpoint{-5.006835\du}{12.469679\du}}
\pgfpathlineto{\pgfpoint{-5.006136\du}{12.472474\du}}
\pgfpathlineto{\pgfpoint{-5.004971\du}{12.474803\du}}
\pgfpathlineto{\pgfpoint{-5.003574\du}{12.477366\du}}
\pgfpathlineto{\pgfpoint{-5.001710\du}{12.480161\du}}
\pgfpathlineto{\pgfpoint{-4.999614\du}{12.482491\du}}
\pgfpathlineto{\pgfpoint{-4.997284\du}{12.484820\du}}
\pgfpathlineto{\pgfpoint{-4.994256\du}{12.487383\du}}
\pgfpathlineto{\pgfpoint{-4.991228\du}{12.489246\du}}
\pgfpathlineto{\pgfpoint{-4.987500\du}{12.491809\du}}
\pgfpathlineto{\pgfpoint{-4.983540\du}{12.494371\du}}
\pgfpathlineto{\pgfpoint{-4.979347\du}{12.496468\du}}
\pgfpathlineto{\pgfpoint{-4.974688\du}{12.498564\du}}
\pgfpathlineto{\pgfpoint{-4.969796\du}{12.500195\du}}
\pgfpathlineto{\pgfpoint{-4.964672\du}{12.502524\du}}
\pgfpathlineto{\pgfpoint{-4.959081\du}{12.504388\du}}
\pgfpathlineto{\pgfpoint{-4.953257\du}{12.505785\du}}
\pgfpathlineto{\pgfpoint{-4.947433\du}{12.507416\du}}
\pgfpathlineto{\pgfpoint{-4.941144\du}{12.509047\du}}
\pgfpathlineto{\pgfpoint{-4.934621\du}{12.510677\du}}
\pgfpathlineto{\pgfpoint{-4.927866\du}{12.511842\du}}
\pgfpathlineto{\pgfpoint{-4.921110\du}{12.513007\du}}
\pgfpathlineto{\pgfpoint{-4.913889\du}{12.513939\du}}
\pgfpathlineto{\pgfpoint{-4.906901\du}{12.515103\du}}
\pgfpathlineto{\pgfpoint{-4.899446\du}{12.515802\du}}
\pgfpathlineto{\pgfpoint{-4.891992\du}{12.516734\du}}
\pgfpathlineto{\pgfpoint{-4.884538\du}{12.516967\du}}
\pgfpathlineto{\pgfpoint{-4.876850\du}{12.517666\du}}
\pgfpathlineto{\pgfpoint{-4.869163\du}{12.517666\du}}
\pgfpathlineto{\pgfpoint{-4.861243\du}{12.518365\du}}
\pgfpathlineto{\pgfpoint{-4.853556\du}{12.518365\du}}
\pgfpathlineto{\pgfpoint{-4.853556\du}{12.518365\du}}
\pgfpathlineto{\pgfpoint{-4.845868\du}{12.518365\du}}
\pgfpathlineto{\pgfpoint{-4.837948\du}{12.517666\du}}
\pgfpathlineto{\pgfpoint{-4.830494\du}{12.517666\du}}
\pgfpathlineto{\pgfpoint{-4.823039\du}{12.516967\du}}
\pgfpathlineto{\pgfpoint{-4.815119\du}{12.516734\du}}
\pgfpathlineto{\pgfpoint{-4.807665\du}{12.515802\du}}
\pgfpathlineto{\pgfpoint{-4.800443\du}{12.515103\du}}
\pgfpathlineto{\pgfpoint{-4.793222\du}{12.513939\du}}
\pgfpathlineto{\pgfpoint{-4.786001\du}{12.513007\du}}
\pgfpathlineto{\pgfpoint{-4.779245\du}{12.511842\du}}
\pgfpathlineto{\pgfpoint{-4.772723\du}{12.510677\du}}
\pgfpathlineto{\pgfpoint{-4.765967\du}{12.509047\du}}
\pgfpathlineto{\pgfpoint{-4.759678\du}{12.507416\du}}
\pgfpathlineto{\pgfpoint{-4.753854\du}{12.505785\du}}
\pgfpathlineto{\pgfpoint{-4.748030\du}{12.504388\du}}
\pgfpathlineto{\pgfpoint{-4.742672\du}{12.502524\du}}
\pgfpathlineto{\pgfpoint{-4.737315\du}{12.500195\du}}
\pgfpathlineto{\pgfpoint{-4.732656\du}{12.498564\du}}
\pgfpathlineto{\pgfpoint{-4.727764\du}{12.496468\du}}
\pgfpathlineto{\pgfpoint{-4.723571\du}{12.494371\du}}
\pgfpathlineto{\pgfpoint{-4.719611\du}{12.491809\du}}
\pgfpathlineto{\pgfpoint{-4.715883\du}{12.489246\du}}
\pgfpathlineto{\pgfpoint{-4.713088\du}{12.487383\du}}
\pgfpathlineto{\pgfpoint{-4.709827\du}{12.484820\du}}
\pgfpathlineto{\pgfpoint{-4.707497\du}{12.482491\du}}
\pgfpathlineto{\pgfpoint{-4.705401\du}{12.480161\du}}
\pgfpathlineto{\pgfpoint{-4.703537\du}{12.477366\du}}
\pgfpathlineto{\pgfpoint{-4.702373\du}{12.474803\du}}
\pgfpathlineto{\pgfpoint{-4.701441\du}{12.472474\du}}
\pgfpathlineto{\pgfpoint{-4.700742\du}{12.469679\du}}
\pgfpathlineto{\pgfpoint{-4.700742\du}{12.467116\du}}
\pgfusepath{fill}
\pgfsetbuttcap
\pgfsetmiterjoin
\pgfsetdash{}{0pt}
\definecolor{dialinecolor}{rgb}{0.666667, 0.901961, 1.000000}
\pgfsetstrokecolor{dialinecolor}
\pgfpathmoveto{\pgfpoint{-4.709827\du}{12.462690\du}}
\pgfpathlineto{\pgfpoint{-4.709827\du}{12.460128\du}}
\pgfpathlineto{\pgfpoint{-4.710526\du}{12.458264\du}}
\pgfpathlineto{\pgfpoint{-4.711457\du}{12.455702\du}}
\pgfpathlineto{\pgfpoint{-4.713088\du}{12.452906\du}}
\pgfpathlineto{\pgfpoint{-4.714719\du}{12.451043\du}}
\pgfpathlineto{\pgfpoint{-4.716349\du}{12.448480\du}}
\pgfpathlineto{\pgfpoint{-4.718912\du}{12.446617\du}}
\pgfpathlineto{\pgfpoint{-4.721474\du}{12.444520\du}}
\pgfpathlineto{\pgfpoint{-4.724735\du}{12.442424\du}}
\pgfpathlineto{\pgfpoint{-4.728463\du}{12.440327\du}}
\pgfpathlineto{\pgfpoint{-4.731957\du}{12.438231\du}}
\pgfpathlineto{\pgfpoint{-4.736383\du}{12.436134\du}}
\pgfpathlineto{\pgfpoint{-4.741042\du}{12.434270\du}}
\pgfpathlineto{\pgfpoint{-4.745468\du}{12.432640\du}}
\pgfpathlineto{\pgfpoint{-4.750593\du}{12.430776\du}}
\pgfpathlineto{\pgfpoint{-4.755950\du}{12.428913\du}}
\pgfpathlineto{\pgfpoint{-4.761308\du}{12.427515\du}}
\pgfpathlineto{\pgfpoint{-4.767365\du}{12.426350\du}}
\pgfpathlineto{\pgfpoint{-4.773189\du}{12.424487\du}}
\pgfpathlineto{\pgfpoint{-4.779711\du}{12.423322\du}}
\pgfpathlineto{\pgfpoint{-4.786467\du}{12.422390\du}}
\pgfpathlineto{\pgfpoint{-4.792989\du}{12.420992\du}}
\pgfpathlineto{\pgfpoint{-4.799745\du}{12.420294\du}}
\pgfpathlineto{\pgfpoint{-4.806733\du}{12.419129\du}}
\pgfpathlineto{\pgfpoint{-4.813489\du}{12.418430\du}}
\pgfpathlineto{\pgfpoint{-4.820943\du}{12.417964\du}}
\pgfpathlineto{\pgfpoint{-4.828164\du}{12.417265\du}}
\pgfpathlineto{\pgfpoint{-4.835619\du}{12.417032\du}}
\pgfpathlineto{\pgfpoint{-4.843073\du}{12.416799\du}}
\pgfpathlineto{\pgfpoint{-4.850993\du}{12.416334\du}}
\pgfpathlineto{\pgfpoint{-4.858680\du}{12.416334\du}}
\pgfpathlineto{\pgfpoint{-4.858680\du}{12.416334\du}}
\pgfpathlineto{\pgfpoint{-4.865669\du}{12.416334\du}}
\pgfpathlineto{\pgfpoint{-4.873356\du}{12.416799\du}}
\pgfpathlineto{\pgfpoint{-4.880810\du}{12.417032\du}}
\pgfpathlineto{\pgfpoint{-4.888032\du}{12.417265\du}}
\pgfpathlineto{\pgfpoint{-4.895486\du}{12.417964\du}}
\pgfpathlineto{\pgfpoint{-4.902940\du}{12.418430\du}}
\pgfpathlineto{\pgfpoint{-4.909929\du}{12.419129\du}}
\pgfpathlineto{\pgfpoint{-4.916917\du}{12.420294\du}}
\pgfpathlineto{\pgfpoint{-4.923673\du}{12.420992\du}}
\pgfpathlineto{\pgfpoint{-4.930195\du}{12.422390\du}}
\pgfpathlineto{\pgfpoint{-4.936951\du}{12.423322\du}}
\pgfpathlineto{\pgfpoint{-4.943240\du}{12.424487\du}}
\pgfpathlineto{\pgfpoint{-4.949297\du}{12.426350\du}}
\pgfpathlineto{\pgfpoint{-4.955121\du}{12.427515\du}}
\pgfpathlineto{\pgfpoint{-4.960246\du}{12.428913\du}}
\pgfpathlineto{\pgfpoint{-4.966069\du}{12.430776\du}}
\pgfpathlineto{\pgfpoint{-4.970961\du}{12.432640\du}}
\pgfpathlineto{\pgfpoint{-4.975620\du}{12.434270\du}}
\pgfpathlineto{\pgfpoint{-4.980279\du}{12.436134\du}}
\pgfpathlineto{\pgfpoint{-4.984239\du}{12.438231\du}}
\pgfpathlineto{\pgfpoint{-4.988199\du}{12.440327\du}}
\pgfpathlineto{\pgfpoint{-4.991926\du}{12.442424\du}}
\pgfpathlineto{\pgfpoint{-4.995188\du}{12.444520\du}}
\pgfpathlineto{\pgfpoint{-4.997750\du}{12.446617\du}}
\pgfpathlineto{\pgfpoint{-5.000312\du}{12.448480\du}}
\pgfpathlineto{\pgfpoint{-5.002176\du}{12.451043\du}}
\pgfpathlineto{\pgfpoint{-5.003574\du}{12.452906\du}}
\pgfpathlineto{\pgfpoint{-5.005437\du}{12.455702\du}}
\pgfpathlineto{\pgfpoint{-5.006136\du}{12.458264\du}}
\pgfpathlineto{\pgfpoint{-5.006835\du}{12.460128\du}}
\pgfpathlineto{\pgfpoint{-5.006835\du}{12.462690\du}}
\pgfpathlineto{\pgfpoint{-5.006835\du}{12.462690\du}}
\pgfpathlineto{\pgfpoint{-5.006835\du}{12.465253\du}}
\pgfpathlineto{\pgfpoint{-5.006136\du}{12.467349\du}}
\pgfpathlineto{\pgfpoint{-5.005437\du}{12.469679\du}}
\pgfpathlineto{\pgfpoint{-5.003574\du}{12.472241\du}}
\pgfpathlineto{\pgfpoint{-5.002176\du}{12.474337\du}}
\pgfpathlineto{\pgfpoint{-5.000312\du}{12.476900\du}}
\pgfpathlineto{\pgfpoint{-4.997750\du}{12.478996\du}}
\pgfpathlineto{\pgfpoint{-4.995188\du}{12.481093\du}}
\pgfpathlineto{\pgfpoint{-4.991926\du}{12.483189\du}}
\pgfpathlineto{\pgfpoint{-4.988199\du}{12.485286\du}}
\pgfpathlineto{\pgfpoint{-4.984239\du}{12.487383\du}}
\pgfpathlineto{\pgfpoint{-4.980279\du}{12.489013\du}}
\pgfpathlineto{\pgfpoint{-4.975620\du}{12.491110\du}}
\pgfpathlineto{\pgfpoint{-4.970961\du}{12.492973\du}}
\pgfpathlineto{\pgfpoint{-4.966069\du}{12.494837\du}}
\pgfpathlineto{\pgfpoint{-4.960246\du}{12.496468\du}}
\pgfpathlineto{\pgfpoint{-4.955121\du}{12.497865\du}}
\pgfpathlineto{\pgfpoint{-4.949297\du}{12.499263\du}}
\pgfpathlineto{\pgfpoint{-4.943240\du}{12.500661\du}}
\pgfpathlineto{\pgfpoint{-4.936951\du}{12.502291\du}}
\pgfpathlineto{\pgfpoint{-4.930195\du}{12.503223\du}}
\pgfpathlineto{\pgfpoint{-4.923673\du}{12.504621\du}}
\pgfpathlineto{\pgfpoint{-4.916917\du}{12.505320\du}}
\pgfpathlineto{\pgfpoint{-4.909929\du}{12.506484\du}}
\pgfpathlineto{\pgfpoint{-4.902940\du}{12.506717\du}}
\pgfpathlineto{\pgfpoint{-4.895486\du}{12.507649\du}}
\pgfpathlineto{\pgfpoint{-4.888032\du}{12.507882\du}}
\pgfpathlineto{\pgfpoint{-4.880810\du}{12.508581\du}}
\pgfpathlineto{\pgfpoint{-4.873356\du}{12.508814\du}}
\pgfpathlineto{\pgfpoint{-4.865669\du}{12.509047\du}}
\pgfpathlineto{\pgfpoint{-4.858680\du}{12.509047\du}}
\pgfpathlineto{\pgfpoint{-4.858680\du}{12.509047\du}}
\pgfpathlineto{\pgfpoint{-4.850993\du}{12.509047\du}}
\pgfpathlineto{\pgfpoint{-4.843073\du}{12.508814\du}}
\pgfpathlineto{\pgfpoint{-4.835619\du}{12.508581\du}}
\pgfpathlineto{\pgfpoint{-4.828164\du}{12.507882\du}}
\pgfpathlineto{\pgfpoint{-4.820943\du}{12.507649\du}}
\pgfpathlineto{\pgfpoint{-4.813489\du}{12.506717\du}}
\pgfpathlineto{\pgfpoint{-4.806733\du}{12.506484\du}}
\pgfpathlineto{\pgfpoint{-4.799745\du}{12.505320\du}}
\pgfpathlineto{\pgfpoint{-4.792989\du}{12.504621\du}}
\pgfpathlineto{\pgfpoint{-4.786467\du}{12.503223\du}}
\pgfpathlineto{\pgfpoint{-4.779711\du}{12.502291\du}}
\pgfpathlineto{\pgfpoint{-4.773189\du}{12.500661\du}}
\pgfpathlineto{\pgfpoint{-4.767365\du}{12.499263\du}}
\pgfpathlineto{\pgfpoint{-4.761308\du}{12.497865\du}}
\pgfpathlineto{\pgfpoint{-4.755950\du}{12.496468\du}}
\pgfpathlineto{\pgfpoint{-4.750593\du}{12.494837\du}}
\pgfpathlineto{\pgfpoint{-4.745468\du}{12.492973\du}}
\pgfpathlineto{\pgfpoint{-4.741042\du}{12.491110\du}}
\pgfpathlineto{\pgfpoint{-4.736383\du}{12.489013\du}}
\pgfpathlineto{\pgfpoint{-4.731957\du}{12.487383\du}}
\pgfpathlineto{\pgfpoint{-4.728463\du}{12.485286\du}}
\pgfpathlineto{\pgfpoint{-4.724735\du}{12.483189\du}}
\pgfpathlineto{\pgfpoint{-4.721474\du}{12.481093\du}}
\pgfpathlineto{\pgfpoint{-4.718912\du}{12.478996\du}}
\pgfpathlineto{\pgfpoint{-4.716349\du}{12.476900\du}}
\pgfpathlineto{\pgfpoint{-4.714719\du}{12.474337\du}}
\pgfpathlineto{\pgfpoint{-4.713088\du}{12.472241\du}}
\pgfpathlineto{\pgfpoint{-4.711457\du}{12.469679\du}}
\pgfpathlineto{\pgfpoint{-4.710526\du}{12.467349\du}}
\pgfpathlineto{\pgfpoint{-4.709827\du}{12.465253\du}}
\pgfpathlineto{\pgfpoint{-4.709827\du}{12.462690\du}}
\pgfusepath{stroke}
\pgfsetbuttcap
\pgfsetmiterjoin
\pgfsetdash{}{0pt}
\definecolor{dialinecolor}{rgb}{0.000000, 0.000000, 0.000000}
\pgfsetfillcolor{dialinecolor}
\pgfpathmoveto{\pgfpoint{-4.849129\du}{12.453605\du}}
\pgfpathlineto{\pgfpoint{-4.830494\du}{12.462690\du}}
\pgfpathlineto{\pgfpoint{-4.774819\du}{12.444287\du}}
\pgfpathlineto{\pgfpoint{-4.746865\du}{12.453605\du}}
\pgfpathlineto{\pgfpoint{-4.765501\du}{12.434969\du}}
\pgfpathlineto{\pgfpoint{-4.830494\du}{12.434969\du}}
\pgfpathlineto{\pgfpoint{-4.802540\du}{12.434969\du}}
\pgfpathlineto{\pgfpoint{-4.849129\du}{12.453605\du}}
\pgfusepath{fill}
\pgfsetbuttcap
\pgfsetmiterjoin
\pgfsetdash{}{0pt}
\definecolor{dialinecolor}{rgb}{0.000000, 0.000000, 0.000000}
\pgfsetfillcolor{dialinecolor}
\pgfpathmoveto{\pgfpoint{-4.858680\du}{12.481326\du}}
\pgfpathlineto{\pgfpoint{-4.885935\du}{12.472241\du}}
\pgfpathlineto{\pgfpoint{-4.932758\du}{12.490877\du}}
\pgfpathlineto{\pgfpoint{-4.960246\du}{12.481326\du}}
\pgfpathlineto{\pgfpoint{-4.951393\du}{12.509047\du}}
\pgfpathlineto{\pgfpoint{-4.885935\du}{12.509047\du}}
\pgfpathlineto{\pgfpoint{-4.914122\du}{12.500195\du}}
\pgfpathlineto{\pgfpoint{-4.858680\du}{12.481326\du}}
\pgfusepath{fill}
\pgfsetbuttcap
\pgfsetmiterjoin
\pgfsetdash{}{0pt}
\definecolor{dialinecolor}{rgb}{0.000000, 0.000000, 0.000000}
\pgfsetfillcolor{dialinecolor}
\pgfpathmoveto{\pgfpoint{-4.951393\du}{12.434969\du}}
\pgfpathlineto{\pgfpoint{-4.932758\du}{12.425186\du}}
\pgfpathlineto{\pgfpoint{-4.876850\du}{12.453605\du}}
\pgfpathlineto{\pgfpoint{-4.858680\du}{12.444287\du}}
\pgfpathlineto{\pgfpoint{-4.867765\du}{12.462690\du}}
\pgfpathlineto{\pgfpoint{-4.932758\du}{12.462690\du}}
\pgfpathlineto{\pgfpoint{-4.904804\du}{12.453605\du}}
\pgfpathlineto{\pgfpoint{-4.951393\du}{12.434969\du}}
\pgfusepath{fill}
\pgfsetbuttcap
\pgfsetmiterjoin
\pgfsetdash{}{0pt}
\definecolor{dialinecolor}{rgb}{0.000000, 0.000000, 0.000000}
\pgfsetfillcolor{dialinecolor}
\pgfpathmoveto{\pgfpoint{-4.755950\du}{12.500195\du}}
\pgfpathlineto{\pgfpoint{-4.774819\du}{12.509047\du}}
\pgfpathlineto{\pgfpoint{-4.830494\du}{12.481326\du}}
\pgfpathlineto{\pgfpoint{-4.858680\du}{12.490877\du}}
\pgfpathlineto{\pgfpoint{-4.839579\du}{12.472241\du}}
\pgfpathlineto{\pgfpoint{-4.774819\du}{12.472241\du}}
\pgfpathlineto{\pgfpoint{-4.802540\du}{12.481326\du}}
\pgfpathlineto{\pgfpoint{-4.755950\du}{12.500195\du}}
\pgfusepath{fill}
\pgfsetbuttcap
\pgfsetmiterjoin
\pgfsetdash{}{0pt}
\definecolor{dialinecolor}{rgb}{1.000000, 1.000000, 1.000000}
\pgfsetfillcolor{dialinecolor}
\pgfpathmoveto{\pgfpoint{-4.849129\du}{12.462690\du}}
\pgfpathlineto{\pgfpoint{-4.830494\du}{12.462690\du}}
\pgfpathlineto{\pgfpoint{-4.774819\du}{12.444287\du}}
\pgfpathlineto{\pgfpoint{-4.746865\du}{12.453605\du}}
\pgfpathlineto{\pgfpoint{-4.765501\du}{12.434969\du}}
\pgfpathlineto{\pgfpoint{-4.820943\du}{12.434969\du}}
\pgfpathlineto{\pgfpoint{-4.802540\du}{12.434969\du}}
\pgfpathlineto{\pgfpoint{-4.849129\du}{12.462690\du}}
\pgfusepath{fill}
\pgfsetbuttcap
\pgfsetmiterjoin
\pgfsetdash{}{0pt}
\definecolor{dialinecolor}{rgb}{1.000000, 1.000000, 1.000000}
\pgfsetfillcolor{dialinecolor}
\pgfpathmoveto{\pgfpoint{-4.858680\du}{12.481326\du}}
\pgfpathlineto{\pgfpoint{-4.876850\du}{12.472241\du}}
\pgfpathlineto{\pgfpoint{-4.932758\du}{12.490877\du}}
\pgfpathlineto{\pgfpoint{-4.960246\du}{12.490877\du}}
\pgfpathlineto{\pgfpoint{-4.941843\du}{12.509047\du}}
\pgfpathlineto{\pgfpoint{-4.876850\du}{12.509047\du}}
\pgfpathlineto{\pgfpoint{-4.904804\du}{12.500195\du}}
\pgfpathlineto{\pgfpoint{-4.858680\du}{12.481326\du}}
\pgfusepath{fill}
\pgfsetbuttcap
\pgfsetmiterjoin
\pgfsetdash{}{0pt}
\definecolor{dialinecolor}{rgb}{1.000000, 1.000000, 1.000000}
\pgfsetfillcolor{dialinecolor}
\pgfpathmoveto{\pgfpoint{-4.951393\du}{12.434969\du}}
\pgfpathlineto{\pgfpoint{-4.932758\du}{12.434969\du}}
\pgfpathlineto{\pgfpoint{-4.876850\du}{12.453605\du}}
\pgfpathlineto{\pgfpoint{-4.849129\du}{12.444287\du}}
\pgfpathlineto{\pgfpoint{-4.867765\du}{12.462690\du}}
\pgfpathlineto{\pgfpoint{-4.932758\du}{12.462690\du}}
\pgfpathlineto{\pgfpoint{-4.904804\du}{12.462690\du}}
\pgfpathlineto{\pgfpoint{-4.951393\du}{12.434969\du}}
\pgfusepath{fill}
\pgfsetbuttcap
\pgfsetmiterjoin
\pgfsetdash{}{0pt}
\definecolor{dialinecolor}{rgb}{1.000000, 1.000000, 1.000000}
\pgfsetfillcolor{dialinecolor}
\pgfpathmoveto{\pgfpoint{-4.755950\du}{12.500195\du}}
\pgfpathlineto{\pgfpoint{-4.774819\du}{12.509047\du}}
\pgfpathlineto{\pgfpoint{-4.830494\du}{12.490877\du}}
\pgfpathlineto{\pgfpoint{-4.849129\du}{12.490877\du}}
\pgfpathlineto{\pgfpoint{-4.839579\du}{12.472241\du}}
\pgfpathlineto{\pgfpoint{-4.774819\du}{12.472241\du}}
\pgfpathlineto{\pgfpoint{-4.802540\du}{12.481326\du}}
\pgfpathlineto{\pgfpoint{-4.755950\du}{12.500195\du}}
\pgfusepath{fill}
\pgfsetbuttcap
\pgfsetmiterjoin
\pgfsetdash{}{0pt}
\definecolor{dialinecolor}{rgb}{0.666667, 0.901961, 1.000000}
\pgfsetstrokecolor{dialinecolor}
\pgfpathmoveto{\pgfpoint{-5.006835\du}{12.462690\du}}
\pgfpathlineto{\pgfpoint{-5.006835\du}{12.536767\du}}
\pgfusepath{stroke}
\pgfsetbuttcap
\pgfsetmiterjoin
\pgfsetdash{}{0pt}
\definecolor{dialinecolor}{rgb}{0.666667, 0.901961, 1.000000}
\pgfsetstrokecolor{dialinecolor}
\pgfpathmoveto{\pgfpoint{-4.709827\du}{12.462690\du}}
\pgfpathlineto{\pgfpoint{-4.709827\du}{12.536767\du}}
\pgfusepath{stroke}
\pgfsetbuttcap
\pgfsetmiterjoin
\pgfsetdash{}{0pt}
\definecolor{dialinecolor}{rgb}{0.000000, 0.470588, 0.666667}
\pgfsetfillcolor{dialinecolor}
\pgfpathmoveto{\pgfpoint{-4.979347\du}{12.857536\du}}
\pgfpathlineto{\pgfpoint{-4.979580\du}{12.854741\du}}
\pgfpathlineto{\pgfpoint{-4.980279\du}{12.852179\du}}
\pgfpathlineto{\pgfpoint{-4.981211\du}{12.849616\du}}
\pgfpathlineto{\pgfpoint{-4.982608\du}{12.847287\du}}
\pgfpathlineto{\pgfpoint{-4.984239\du}{12.844724\du}}
\pgfpathlineto{\pgfpoint{-4.986802\du}{12.842162\du}}
\pgfpathlineto{\pgfpoint{-4.988898\du}{12.839832\du}}
\pgfpathlineto{\pgfpoint{-4.991926\du}{12.837270\du}}
\pgfpathlineto{\pgfpoint{-4.995188\du}{12.835173\du}}
\pgfpathlineto{\pgfpoint{-4.998682\du}{12.832611\du}}
\pgfpathlineto{\pgfpoint{-5.002875\du}{12.830514\du}}
\pgfpathlineto{\pgfpoint{-5.006835\du}{12.828418\du}}
\pgfpathlineto{\pgfpoint{-5.011494\du}{12.826321\du}}
\pgfpathlineto{\pgfpoint{-5.016153\du}{12.824225\du}}
\pgfpathlineto{\pgfpoint{-5.021744\du}{12.822361\du}}
\pgfpathlineto{\pgfpoint{-5.027101\du}{12.820498\du}}
\pgfpathlineto{\pgfpoint{-5.032925\du}{12.819100\du}}
\pgfpathlineto{\pgfpoint{-5.038749\du}{12.817236\du}}
\pgfpathlineto{\pgfpoint{-5.045038\du}{12.815839\du}}
\pgfpathlineto{\pgfpoint{-5.051561\du}{12.813975\du}}
\pgfpathlineto{\pgfpoint{-5.058084\du}{12.812810\du}}
\pgfpathlineto{\pgfpoint{-5.065072\du}{12.811879\du}}
\pgfpathlineto{\pgfpoint{-5.072060\du}{12.810714\du}}
\pgfpathlineto{\pgfpoint{-5.079282\du}{12.809782\du}}
\pgfpathlineto{\pgfpoint{-5.086969\du}{12.808850\du}}
\pgfpathlineto{\pgfpoint{-5.094190\du}{12.808151\du}}
\pgfpathlineto{\pgfpoint{-5.101878\du}{12.807686\du}}
\pgfpathlineto{\pgfpoint{-5.109332\du}{12.807220\du}}
\pgfpathlineto{\pgfpoint{-5.117019\du}{12.806754\du}}
\pgfpathlineto{\pgfpoint{-5.125172\du}{12.806521\du}}
\pgfpathlineto{\pgfpoint{-5.132394\du}{12.806521\du}}
\pgfpathlineto{\pgfpoint{-5.132394\du}{12.806521\du}}
\pgfpathlineto{\pgfpoint{-5.140081\du}{12.806521\du}}
\pgfpathlineto{\pgfpoint{-5.148001\du}{12.806754\du}}
\pgfpathlineto{\pgfpoint{-5.155689\du}{12.807220\du}}
\pgfpathlineto{\pgfpoint{-5.163376\du}{12.807686\du}}
\pgfpathlineto{\pgfpoint{-5.171063\du}{12.808151\du}}
\pgfpathlineto{\pgfpoint{-5.178284\du}{12.808850\du}}
\pgfpathlineto{\pgfpoint{-5.185972\du}{12.809782\du}}
\pgfpathlineto{\pgfpoint{-5.193193\du}{12.810714\du}}
\pgfpathlineto{\pgfpoint{-5.200182\du}{12.811879\du}}
\pgfpathlineto{\pgfpoint{-5.206937\du}{12.812810\du}}
\pgfpathlineto{\pgfpoint{-5.213460\du}{12.813975\du}}
\pgfpathlineto{\pgfpoint{-5.219982\du}{12.815839\du}}
\pgfpathlineto{\pgfpoint{-5.226272\du}{12.817236\du}}
\pgfpathlineto{\pgfpoint{-5.232095\du}{12.819100\du}}
\pgfpathlineto{\pgfpoint{-5.237919\du}{12.820498\du}}
\pgfpathlineto{\pgfpoint{-5.243510\du}{12.822361\du}}
\pgfpathlineto{\pgfpoint{-5.248868\du}{12.824225\du}}
\pgfpathlineto{\pgfpoint{-5.253527\du}{12.826321\du}}
\pgfpathlineto{\pgfpoint{-5.258418\du}{12.828418\du}}
\pgfpathlineto{\pgfpoint{-5.262379\du}{12.830514\du}}
\pgfpathlineto{\pgfpoint{-5.266339\du}{12.832611\du}}
\pgfpathlineto{\pgfpoint{-5.270066\du}{12.835173\du}}
\pgfpathlineto{\pgfpoint{-5.273327\du}{12.837270\du}}
\pgfpathlineto{\pgfpoint{-5.276122\du}{12.839832\du}}
\pgfpathlineto{\pgfpoint{-5.278452\du}{12.842162\du}}
\pgfpathlineto{\pgfpoint{-5.280548\du}{12.844724\du}}
\pgfpathlineto{\pgfpoint{-5.282412\du}{12.847287\du}}
\pgfpathlineto{\pgfpoint{-5.283810\du}{12.849616\du}}
\pgfpathlineto{\pgfpoint{-5.284974\du}{12.852179\du}}
\pgfpathlineto{\pgfpoint{-5.285673\du}{12.854741\du}}
\pgfpathlineto{\pgfpoint{-5.285673\du}{12.857536\du}}
\pgfpathlineto{\pgfpoint{-5.285673\du}{12.857536\du}}
\pgfpathlineto{\pgfpoint{-5.285673\du}{12.860099\du}}
\pgfpathlineto{\pgfpoint{-5.284974\du}{12.862661\du}}
\pgfpathlineto{\pgfpoint{-5.283810\du}{12.865224\du}}
\pgfpathlineto{\pgfpoint{-5.282412\du}{12.867786\du}}
\pgfpathlineto{\pgfpoint{-5.280548\du}{12.870115\du}}
\pgfpathlineto{\pgfpoint{-5.278452\du}{12.872678\du}}
\pgfpathlineto{\pgfpoint{-5.276122\du}{12.875007\du}}
\pgfpathlineto{\pgfpoint{-5.273327\du}{12.877570\du}}
\pgfpathlineto{\pgfpoint{-5.270066\du}{12.879899\du}}
\pgfpathlineto{\pgfpoint{-5.266339\du}{12.882229\du}}
\pgfpathlineto{\pgfpoint{-5.262379\du}{12.884325\du}}
\pgfpathlineto{\pgfpoint{-5.258418\du}{12.886422\du}}
\pgfpathlineto{\pgfpoint{-5.253527\du}{12.888518\du}}
\pgfpathlineto{\pgfpoint{-5.248868\du}{12.890382\du}}
\pgfpathlineto{\pgfpoint{-5.243510\du}{12.892478\du}}
\pgfpathlineto{\pgfpoint{-5.237919\du}{12.894342\du}}
\pgfpathlineto{\pgfpoint{-5.232095\du}{12.895973\du}}
\pgfpathlineto{\pgfpoint{-5.226272\du}{12.897836\du}}
\pgfpathlineto{\pgfpoint{-5.219982\du}{12.899001\du}}
\pgfpathlineto{\pgfpoint{-5.213460\du}{12.900632\du}}
\pgfpathlineto{\pgfpoint{-5.206937\du}{12.902029\du}}
\pgfpathlineto{\pgfpoint{-5.200182\du}{12.902961\du}}
\pgfpathlineto{\pgfpoint{-5.193193\du}{12.904126\du}}
\pgfpathlineto{\pgfpoint{-5.185972\du}{12.905058\du}}
\pgfpathlineto{\pgfpoint{-5.178284\du}{12.905989\du}}
\pgfpathlineto{\pgfpoint{-5.171063\du}{12.906688\du}}
\pgfpathlineto{\pgfpoint{-5.163376\du}{12.907154\du}}
\pgfpathlineto{\pgfpoint{-5.155689\du}{12.907853\du}}
\pgfpathlineto{\pgfpoint{-5.148001\du}{12.907853\du}}
\pgfpathlineto{\pgfpoint{-5.140081\du}{12.908319\du}}
\pgfpathlineto{\pgfpoint{-5.132394\du}{12.908319\du}}
\pgfpathlineto{\pgfpoint{-5.132394\du}{12.908319\du}}
\pgfpathlineto{\pgfpoint{-5.125172\du}{12.908319\du}}
\pgfpathlineto{\pgfpoint{-5.117019\du}{12.907853\du}}
\pgfpathlineto{\pgfpoint{-5.109332\du}{12.907853\du}}
\pgfpathlineto{\pgfpoint{-5.101878\du}{12.907154\du}}
\pgfpathlineto{\pgfpoint{-5.094190\du}{12.906688\du}}
\pgfpathlineto{\pgfpoint{-5.086969\du}{12.905989\du}}
\pgfpathlineto{\pgfpoint{-5.079282\du}{12.905058\du}}
\pgfpathlineto{\pgfpoint{-5.072060\du}{12.904126\du}}
\pgfpathlineto{\pgfpoint{-5.065072\du}{12.902961\du}}
\pgfpathlineto{\pgfpoint{-5.058084\du}{12.902029\du}}
\pgfpathlineto{\pgfpoint{-5.051561\du}{12.900632\du}}
\pgfpathlineto{\pgfpoint{-5.045038\du}{12.899001\du}}
\pgfpathlineto{\pgfpoint{-5.038749\du}{12.897836\du}}
\pgfpathlineto{\pgfpoint{-5.032925\du}{12.895973\du}}
\pgfpathlineto{\pgfpoint{-5.027101\du}{12.894342\du}}
\pgfpathlineto{\pgfpoint{-5.021744\du}{12.892478\du}}
\pgfpathlineto{\pgfpoint{-5.016153\du}{12.890382\du}}
\pgfpathlineto{\pgfpoint{-5.011494\du}{12.888518\du}}
\pgfpathlineto{\pgfpoint{-5.006835\du}{12.886422\du}}
\pgfpathlineto{\pgfpoint{-5.002875\du}{12.884325\du}}
\pgfpathlineto{\pgfpoint{-4.998682\du}{12.882229\du}}
\pgfpathlineto{\pgfpoint{-4.995188\du}{12.879899\du}}
\pgfpathlineto{\pgfpoint{-4.991926\du}{12.877570\du}}
\pgfpathlineto{\pgfpoint{-4.988898\du}{12.875007\du}}
\pgfpathlineto{\pgfpoint{-4.986802\du}{12.872678\du}}
\pgfpathlineto{\pgfpoint{-4.984239\du}{12.870115\du}}
\pgfpathlineto{\pgfpoint{-4.982608\du}{12.867786\du}}
\pgfpathlineto{\pgfpoint{-4.981211\du}{12.865224\du}}
\pgfpathlineto{\pgfpoint{-4.980279\du}{12.862661\du}}
\pgfpathlineto{\pgfpoint{-4.979580\du}{12.860099\du}}
\pgfpathlineto{\pgfpoint{-4.979347\du}{12.857536\du}}
\pgfusepath{fill}
\pgfsetbuttcap
\pgfsetmiterjoin
\pgfsetdash{}{0pt}
\definecolor{dialinecolor}{rgb}{0.666667, 0.901961, 1.000000}
\pgfsetstrokecolor{dialinecolor}
\pgfpathmoveto{\pgfpoint{-4.988665\du}{12.852644\du}}
\pgfpathlineto{\pgfpoint{-4.988898\du}{12.850315\du}}
\pgfpathlineto{\pgfpoint{-4.989364\du}{12.848218\du}}
\pgfpathlineto{\pgfpoint{-4.990296\du}{12.845889\du}}
\pgfpathlineto{\pgfpoint{-4.991926\du}{12.843094\du}}
\pgfpathlineto{\pgfpoint{-4.993557\du}{12.841230\du}}
\pgfpathlineto{\pgfpoint{-4.995188\du}{12.838668\du}}
\pgfpathlineto{\pgfpoint{-4.997750\du}{12.836571\du}}
\pgfpathlineto{\pgfpoint{-5.000312\du}{12.834474\du}}
\pgfpathlineto{\pgfpoint{-5.003574\du}{12.832145\du}}
\pgfpathlineto{\pgfpoint{-5.007301\du}{12.830281\du}}
\pgfpathlineto{\pgfpoint{-5.011261\du}{12.828185\du}}
\pgfpathlineto{\pgfpoint{-5.015221\du}{12.826321\du}}
\pgfpathlineto{\pgfpoint{-5.019647\du}{12.824458\du}}
\pgfpathlineto{\pgfpoint{-5.024306\du}{12.822594\du}}
\pgfpathlineto{\pgfpoint{-5.029431\du}{12.820498\du}}
\pgfpathlineto{\pgfpoint{-5.034789\du}{12.819333\du}}
\pgfpathlineto{\pgfpoint{-5.039914\du}{12.817702\du}}
\pgfpathlineto{\pgfpoint{-5.046203\du}{12.816305\du}}
\pgfpathlineto{\pgfpoint{-5.052027\du}{12.815140\du}}
\pgfpathlineto{\pgfpoint{-5.058084\du}{12.813509\du}}
\pgfpathlineto{\pgfpoint{-5.065072\du}{12.812344\du}}
\pgfpathlineto{\pgfpoint{-5.071594\du}{12.811180\du}}
\pgfpathlineto{\pgfpoint{-5.078117\du}{12.810248\du}}
\pgfpathlineto{\pgfpoint{-5.085338\du}{12.809316\du}}
\pgfpathlineto{\pgfpoint{-5.092094\du}{12.808617\du}}
\pgfpathlineto{\pgfpoint{-5.099548\du}{12.807918\du}}
\pgfpathlineto{\pgfpoint{-5.107002\du}{12.807686\du}}
\pgfpathlineto{\pgfpoint{-5.114224\du}{12.807220\du}}
\pgfpathlineto{\pgfpoint{-5.121911\du}{12.806754\du}}
\pgfpathlineto{\pgfpoint{-5.129598\du}{12.806521\du}}
\pgfpathlineto{\pgfpoint{-5.137053\du}{12.806521\du}}
\pgfpathlineto{\pgfpoint{-5.137053\du}{12.806521\du}}
\pgfpathlineto{\pgfpoint{-5.144274\du}{12.806521\du}}
\pgfpathlineto{\pgfpoint{-5.151961\du}{12.806754\du}}
\pgfpathlineto{\pgfpoint{-5.159416\du}{12.807220\du}}
\pgfpathlineto{\pgfpoint{-5.166870\du}{12.807686\du}}
\pgfpathlineto{\pgfpoint{-5.174324\du}{12.807918\du}}
\pgfpathlineto{\pgfpoint{-5.181546\du}{12.808617\du}}
\pgfpathlineto{\pgfpoint{-5.188534\du}{12.809316\du}}
\pgfpathlineto{\pgfpoint{-5.195756\du}{12.810248\du}}
\pgfpathlineto{\pgfpoint{-5.202278\du}{12.811180\du}}
\pgfpathlineto{\pgfpoint{-5.209266\du}{12.812344\du}}
\pgfpathlineto{\pgfpoint{-5.215789\du}{12.813509\du}}
\pgfpathlineto{\pgfpoint{-5.221846\du}{12.815140\du}}
\pgfpathlineto{\pgfpoint{-5.228135\du}{12.816305\du}}
\pgfpathlineto{\pgfpoint{-5.233959\du}{12.817702\du}}
\pgfpathlineto{\pgfpoint{-5.239084\du}{12.819333\du}}
\pgfpathlineto{\pgfpoint{-5.244907\du}{12.820498\du}}
\pgfpathlineto{\pgfpoint{-5.249566\du}{12.822594\du}}
\pgfpathlineto{\pgfpoint{-5.254225\du}{12.824458\du}}
\pgfpathlineto{\pgfpoint{-5.259117\du}{12.826321\du}}
\pgfpathlineto{\pgfpoint{-5.263077\du}{12.828185\du}}
\pgfpathlineto{\pgfpoint{-5.267038\du}{12.830281\du}}
\pgfpathlineto{\pgfpoint{-5.270299\du}{12.832145\du}}
\pgfpathlineto{\pgfpoint{-5.273560\du}{12.834474\du}}
\pgfpathlineto{\pgfpoint{-5.276588\du}{12.836571\du}}
\pgfpathlineto{\pgfpoint{-5.279151\du}{12.838668\du}}
\pgfpathlineto{\pgfpoint{-5.280548\du}{12.841230\du}}
\pgfpathlineto{\pgfpoint{-5.282412\du}{12.843094\du}}
\pgfpathlineto{\pgfpoint{-5.283810\du}{12.845889\du}}
\pgfpathlineto{\pgfpoint{-5.284974\du}{12.848218\du}}
\pgfpathlineto{\pgfpoint{-5.285207\du}{12.850315\du}}
\pgfpathlineto{\pgfpoint{-5.285673\du}{12.852644\du}}
\pgfpathlineto{\pgfpoint{-5.285673\du}{12.852644\du}}
\pgfpathlineto{\pgfpoint{-5.285207\du}{12.855440\du}}
\pgfpathlineto{\pgfpoint{-5.284974\du}{12.857536\du}}
\pgfpathlineto{\pgfpoint{-5.283810\du}{12.859866\du}}
\pgfpathlineto{\pgfpoint{-5.282412\du}{12.862195\du}}
\pgfpathlineto{\pgfpoint{-5.280548\du}{12.864292\du}}
\pgfpathlineto{\pgfpoint{-5.279151\du}{12.866854\du}}
\pgfpathlineto{\pgfpoint{-5.276588\du}{12.868951\du}}
\pgfpathlineto{\pgfpoint{-5.273560\du}{12.871280\du}}
\pgfpathlineto{\pgfpoint{-5.270299\du}{12.873377\du}}
\pgfpathlineto{\pgfpoint{-5.267038\du}{12.875473\du}}
\pgfpathlineto{\pgfpoint{-5.263077\du}{12.877570\du}}
\pgfpathlineto{\pgfpoint{-5.259117\du}{12.878735\du}}
\pgfpathlineto{\pgfpoint{-5.254225\du}{12.881297\du}}
\pgfpathlineto{\pgfpoint{-5.249566\du}{12.882928\du}}
\pgfpathlineto{\pgfpoint{-5.244907\du}{12.884558\du}}
\pgfpathlineto{\pgfpoint{-5.239084\du}{12.886422\du}}
\pgfpathlineto{\pgfpoint{-5.233959\du}{12.888052\du}}
\pgfpathlineto{\pgfpoint{-5.228135\du}{12.889450\du}}
\pgfpathlineto{\pgfpoint{-5.221846\du}{12.890382\du}}
\pgfpathlineto{\pgfpoint{-5.215789\du}{12.892245\du}}
\pgfpathlineto{\pgfpoint{-5.209266\du}{12.893410\du}}
\pgfpathlineto{\pgfpoint{-5.202278\du}{12.894575\du}}
\pgfpathlineto{\pgfpoint{-5.195756\du}{12.895507\du}}
\pgfpathlineto{\pgfpoint{-5.188534\du}{12.896206\du}}
\pgfpathlineto{\pgfpoint{-5.181546\du}{12.896904\du}}
\pgfpathlineto{\pgfpoint{-5.174324\du}{12.897836\du}}
\pgfpathlineto{\pgfpoint{-5.166870\du}{12.898069\du}}
\pgfpathlineto{\pgfpoint{-5.159416\du}{12.898535\du}}
\pgfpathlineto{\pgfpoint{-5.151961\du}{12.898768\du}}
\pgfpathlineto{\pgfpoint{-5.144274\du}{12.899001\du}}
\pgfpathlineto{\pgfpoint{-5.137053\du}{12.899001\du}}
\pgfpathlineto{\pgfpoint{-5.137053\du}{12.899001\du}}
\pgfpathlineto{\pgfpoint{-5.129598\du}{12.899001\du}}
\pgfpathlineto{\pgfpoint{-5.121911\du}{12.898768\du}}
\pgfpathlineto{\pgfpoint{-5.114224\du}{12.898535\du}}
\pgfpathlineto{\pgfpoint{-5.107002\du}{12.898069\du}}
\pgfpathlineto{\pgfpoint{-5.099548\du}{12.897836\du}}
\pgfpathlineto{\pgfpoint{-5.092094\du}{12.896904\du}}
\pgfpathlineto{\pgfpoint{-5.085338\du}{12.896206\du}}
\pgfpathlineto{\pgfpoint{-5.078117\du}{12.895507\du}}
\pgfpathlineto{\pgfpoint{-5.071594\du}{12.894575\du}}
\pgfpathlineto{\pgfpoint{-5.065072\du}{12.893410\du}}
\pgfpathlineto{\pgfpoint{-5.058084\du}{12.892245\du}}
\pgfpathlineto{\pgfpoint{-5.052027\du}{12.890382\du}}
\pgfpathlineto{\pgfpoint{-5.046203\du}{12.889450\du}}
\pgfpathlineto{\pgfpoint{-5.039914\du}{12.888052\du}}
\pgfpathlineto{\pgfpoint{-5.034789\du}{12.886422\du}}
\pgfpathlineto{\pgfpoint{-5.029431\du}{12.884558\du}}
\pgfpathlineto{\pgfpoint{-5.024306\du}{12.882928\du}}
\pgfpathlineto{\pgfpoint{-5.019647\du}{12.881297\du}}
\pgfpathlineto{\pgfpoint{-5.015221\du}{12.878735\du}}
\pgfpathlineto{\pgfpoint{-5.011261\du}{12.877570\du}}
\pgfpathlineto{\pgfpoint{-5.007301\du}{12.875473\du}}
\pgfpathlineto{\pgfpoint{-5.003574\du}{12.873377\du}}
\pgfpathlineto{\pgfpoint{-5.000312\du}{12.871280\du}}
\pgfpathlineto{\pgfpoint{-4.997750\du}{12.868951\du}}
\pgfpathlineto{\pgfpoint{-4.995188\du}{12.866854\du}}
\pgfpathlineto{\pgfpoint{-4.993557\du}{12.864292\du}}
\pgfpathlineto{\pgfpoint{-4.991926\du}{12.862195\du}}
\pgfpathlineto{\pgfpoint{-4.990296\du}{12.859866\du}}
\pgfpathlineto{\pgfpoint{-4.989364\du}{12.857536\du}}
\pgfpathlineto{\pgfpoint{-4.988898\du}{12.855440\du}}
\pgfpathlineto{\pgfpoint{-4.988665\du}{12.852644\du}}
\pgfusepath{stroke}
\pgfsetbuttcap
\pgfsetmiterjoin
\pgfsetdash{}{0pt}
\definecolor{dialinecolor}{rgb}{0.000000, 0.470588, 0.666667}
\pgfsetfillcolor{dialinecolor}
\pgfpathmoveto{\pgfpoint{-5.285673\du}{12.787885\du}}
\pgfpathlineto{\pgfpoint{-5.285673\du}{12.862195\du}}
\pgfpathlineto{\pgfpoint{-4.979114\du}{12.862195\du}}
\pgfpathlineto{\pgfpoint{-4.979114\du}{12.787885\du}}
\pgfpathlineto{\pgfpoint{-5.285673\du}{12.787885\du}}
\pgfusepath{fill}
\pgfsetbuttcap
\pgfsetmiterjoin
\pgfsetdash{}{0pt}
\definecolor{dialinecolor}{rgb}{0.000000, 0.705882, 1.000000}
\pgfsetfillcolor{dialinecolor}
\pgfpathmoveto{\pgfpoint{-4.979347\du}{12.782993\du}}
\pgfpathlineto{\pgfpoint{-4.979580\du}{12.780198\du}}
\pgfpathlineto{\pgfpoint{-4.980279\du}{12.777635\du}}
\pgfpathlineto{\pgfpoint{-4.981211\du}{12.775306\du}}
\pgfpathlineto{\pgfpoint{-4.982608\du}{12.772510\du}}
\pgfpathlineto{\pgfpoint{-4.984239\du}{12.769948\du}}
\pgfpathlineto{\pgfpoint{-4.986802\du}{12.767386\du}}
\pgfpathlineto{\pgfpoint{-4.988898\du}{12.765289\du}}
\pgfpathlineto{\pgfpoint{-4.991926\du}{12.762261\du}}
\pgfpathlineto{\pgfpoint{-4.995188\du}{12.760397\du}}
\pgfpathlineto{\pgfpoint{-4.998682\du}{12.758068\du}}
\pgfpathlineto{\pgfpoint{-5.002875\du}{12.755738\du}}
\pgfpathlineto{\pgfpoint{-5.006835\du}{12.753642\du}}
\pgfpathlineto{\pgfpoint{-5.011494\du}{12.751545\du}}
\pgfpathlineto{\pgfpoint{-5.016153\du}{12.749682\du}}
\pgfpathlineto{\pgfpoint{-5.021744\du}{12.747585\du}}
\pgfpathlineto{\pgfpoint{-5.027101\du}{12.745721\du}}
\pgfpathlineto{\pgfpoint{-5.032925\du}{12.744091\du}}
\pgfpathlineto{\pgfpoint{-5.038749\du}{12.742460\du}}
\pgfpathlineto{\pgfpoint{-5.045038\du}{12.741062\du}}
\pgfpathlineto{\pgfpoint{-5.051561\du}{12.739432\du}}
\pgfpathlineto{\pgfpoint{-5.058084\du}{12.738034\du}}
\pgfpathlineto{\pgfpoint{-5.065072\du}{12.737102\du}}
\pgfpathlineto{\pgfpoint{-5.072060\du}{12.735938\du}}
\pgfpathlineto{\pgfpoint{-5.079282\du}{12.735006\du}}
\pgfpathlineto{\pgfpoint{-5.086969\du}{12.734074\du}}
\pgfpathlineto{\pgfpoint{-5.094190\du}{12.733142\du}}
\pgfpathlineto{\pgfpoint{-5.101878\du}{12.732909\du}}
\pgfpathlineto{\pgfpoint{-5.109332\du}{12.732210\du}}
\pgfpathlineto{\pgfpoint{-5.117019\du}{12.731978\du}}
\pgfpathlineto{\pgfpoint{-5.125172\du}{12.731745\du}}
\pgfpathlineto{\pgfpoint{-5.132394\du}{12.731745\du}}
\pgfpathlineto{\pgfpoint{-5.132394\du}{12.731745\du}}
\pgfpathlineto{\pgfpoint{-5.140081\du}{12.731745\du}}
\pgfpathlineto{\pgfpoint{-5.148001\du}{12.731978\du}}
\pgfpathlineto{\pgfpoint{-5.155689\du}{12.732210\du}}
\pgfpathlineto{\pgfpoint{-5.163376\du}{12.732909\du}}
\pgfpathlineto{\pgfpoint{-5.171063\du}{12.733142\du}}
\pgfpathlineto{\pgfpoint{-5.178284\du}{12.734074\du}}
\pgfpathlineto{\pgfpoint{-5.185972\du}{12.735006\du}}
\pgfpathlineto{\pgfpoint{-5.193193\du}{12.735938\du}}
\pgfpathlineto{\pgfpoint{-5.200182\du}{12.737102\du}}
\pgfpathlineto{\pgfpoint{-5.206937\du}{12.738034\du}}
\pgfpathlineto{\pgfpoint{-5.213460\du}{12.739432\du}}
\pgfpathlineto{\pgfpoint{-5.219982\du}{12.741062\du}}
\pgfpathlineto{\pgfpoint{-5.226272\du}{12.742460\du}}
\pgfpathlineto{\pgfpoint{-5.232095\du}{12.744091\du}}
\pgfpathlineto{\pgfpoint{-5.237919\du}{12.745721\du}}
\pgfpathlineto{\pgfpoint{-5.243510\du}{12.747585\du}}
\pgfpathlineto{\pgfpoint{-5.248868\du}{12.749682\du}}
\pgfpathlineto{\pgfpoint{-5.253527\du}{12.751545\du}}
\pgfpathlineto{\pgfpoint{-5.258418\du}{12.753642\du}}
\pgfpathlineto{\pgfpoint{-5.262379\du}{12.755738\du}}
\pgfpathlineto{\pgfpoint{-5.266339\du}{12.758068\du}}
\pgfpathlineto{\pgfpoint{-5.270066\du}{12.760397\du}}
\pgfpathlineto{\pgfpoint{-5.273327\du}{12.762261\du}}
\pgfpathlineto{\pgfpoint{-5.276122\du}{12.765289\du}}
\pgfpathlineto{\pgfpoint{-5.278452\du}{12.767386\du}}
\pgfpathlineto{\pgfpoint{-5.280548\du}{12.769948\du}}
\pgfpathlineto{\pgfpoint{-5.282412\du}{12.772510\du}}
\pgfpathlineto{\pgfpoint{-5.283810\du}{12.775306\du}}
\pgfpathlineto{\pgfpoint{-5.284974\du}{12.777635\du}}
\pgfpathlineto{\pgfpoint{-5.285673\du}{12.780198\du}}
\pgfpathlineto{\pgfpoint{-5.285673\du}{12.782993\du}}
\pgfpathlineto{\pgfpoint{-5.285673\du}{12.782993\du}}
\pgfpathlineto{\pgfpoint{-5.285673\du}{12.785555\du}}
\pgfpathlineto{\pgfpoint{-5.284974\du}{12.788118\du}}
\pgfpathlineto{\pgfpoint{-5.283810\du}{12.790447\du}}
\pgfpathlineto{\pgfpoint{-5.282412\du}{12.793243\du}}
\pgfpathlineto{\pgfpoint{-5.280548\du}{12.795805\du}}
\pgfpathlineto{\pgfpoint{-5.278452\du}{12.798135\du}}
\pgfpathlineto{\pgfpoint{-5.276122\du}{12.800464\du}}
\pgfpathlineto{\pgfpoint{-5.273327\du}{12.803027\du}}
\pgfpathlineto{\pgfpoint{-5.270066\du}{12.805356\du}}
\pgfpathlineto{\pgfpoint{-5.266339\du}{12.807686\du}}
\pgfpathlineto{\pgfpoint{-5.262379\du}{12.810015\du}}
\pgfpathlineto{\pgfpoint{-5.258418\du}{12.812112\du}}
\pgfpathlineto{\pgfpoint{-5.253527\du}{12.813975\du}}
\pgfpathlineto{\pgfpoint{-5.248868\du}{12.816072\du}}
\pgfpathlineto{\pgfpoint{-5.243510\du}{12.818168\du}}
\pgfpathlineto{\pgfpoint{-5.237919\du}{12.819799\du}}
\pgfpathlineto{\pgfpoint{-5.232095\du}{12.821662\du}}
\pgfpathlineto{\pgfpoint{-5.226272\du}{12.823293\du}}
\pgfpathlineto{\pgfpoint{-5.219982\du}{12.824691\du}}
\pgfpathlineto{\pgfpoint{-5.213460\du}{12.826321\du}}
\pgfpathlineto{\pgfpoint{-5.206937\du}{12.827719\du}}
\pgfpathlineto{\pgfpoint{-5.200182\du}{12.828651\du}}
\pgfpathlineto{\pgfpoint{-5.193193\du}{12.829816\du}}
\pgfpathlineto{\pgfpoint{-5.185972\du}{12.830747\du}}
\pgfpathlineto{\pgfpoint{-5.178284\du}{12.831446\du}}
\pgfpathlineto{\pgfpoint{-5.171063\du}{12.832145\du}}
\pgfpathlineto{\pgfpoint{-5.163376\du}{12.832844\du}}
\pgfpathlineto{\pgfpoint{-5.155689\du}{12.833543\du}}
\pgfpathlineto{\pgfpoint{-5.148001\du}{12.833776\du}}
\pgfpathlineto{\pgfpoint{-5.140081\du}{12.834009\du}}
\pgfpathlineto{\pgfpoint{-5.132394\du}{12.834009\du}}
\pgfpathlineto{\pgfpoint{-5.132394\du}{12.834009\du}}
\pgfpathlineto{\pgfpoint{-5.125172\du}{12.834009\du}}
\pgfpathlineto{\pgfpoint{-5.117019\du}{12.833776\du}}
\pgfpathlineto{\pgfpoint{-5.109332\du}{12.833543\du}}
\pgfpathlineto{\pgfpoint{-5.101878\du}{12.832844\du}}
\pgfpathlineto{\pgfpoint{-5.094190\du}{12.832145\du}}
\pgfpathlineto{\pgfpoint{-5.086969\du}{12.831446\du}}
\pgfpathlineto{\pgfpoint{-5.079282\du}{12.830747\du}}
\pgfpathlineto{\pgfpoint{-5.072060\du}{12.829816\du}}
\pgfpathlineto{\pgfpoint{-5.065072\du}{12.828651\du}}
\pgfpathlineto{\pgfpoint{-5.058084\du}{12.827719\du}}
\pgfpathlineto{\pgfpoint{-5.051561\du}{12.826321\du}}
\pgfpathlineto{\pgfpoint{-5.045038\du}{12.824691\du}}
\pgfpathlineto{\pgfpoint{-5.038749\du}{12.823293\du}}
\pgfpathlineto{\pgfpoint{-5.032925\du}{12.821662\du}}
\pgfpathlineto{\pgfpoint{-5.027101\du}{12.819799\du}}
\pgfpathlineto{\pgfpoint{-5.021744\du}{12.818168\du}}
\pgfpathlineto{\pgfpoint{-5.016153\du}{12.816072\du}}
\pgfpathlineto{\pgfpoint{-5.011494\du}{12.813975\du}}
\pgfpathlineto{\pgfpoint{-5.006835\du}{12.812112\du}}
\pgfpathlineto{\pgfpoint{-5.002875\du}{12.810015\du}}
\pgfpathlineto{\pgfpoint{-4.998682\du}{12.807686\du}}
\pgfpathlineto{\pgfpoint{-4.995188\du}{12.805356\du}}
\pgfpathlineto{\pgfpoint{-4.991926\du}{12.803027\du}}
\pgfpathlineto{\pgfpoint{-4.988898\du}{12.800464\du}}
\pgfpathlineto{\pgfpoint{-4.986802\du}{12.798135\du}}
\pgfpathlineto{\pgfpoint{-4.984239\du}{12.795805\du}}
\pgfpathlineto{\pgfpoint{-4.982608\du}{12.793243\du}}
\pgfpathlineto{\pgfpoint{-4.981211\du}{12.790447\du}}
\pgfpathlineto{\pgfpoint{-4.980279\du}{12.788118\du}}
\pgfpathlineto{\pgfpoint{-4.979580\du}{12.785555\du}}
\pgfpathlineto{\pgfpoint{-4.979347\du}{12.782993\du}}
\pgfusepath{fill}
\pgfsetbuttcap
\pgfsetmiterjoin
\pgfsetdash{}{0pt}
\definecolor{dialinecolor}{rgb}{0.666667, 0.901961, 1.000000}
\pgfsetstrokecolor{dialinecolor}
\pgfpathmoveto{\pgfpoint{-4.988665\du}{12.778567\du}}
\pgfpathlineto{\pgfpoint{-4.988898\du}{12.776238\du}}
\pgfpathlineto{\pgfpoint{-4.989364\du}{12.773908\du}}
\pgfpathlineto{\pgfpoint{-4.990296\du}{12.771812\du}}
\pgfpathlineto{\pgfpoint{-4.991926\du}{12.769249\du}}
\pgfpathlineto{\pgfpoint{-4.993557\du}{12.767153\du}}
\pgfpathlineto{\pgfpoint{-4.995188\du}{12.764590\du}}
\pgfpathlineto{\pgfpoint{-4.997750\du}{12.762261\du}}
\pgfpathlineto{\pgfpoint{-5.000312\du}{12.760397\du}}
\pgfpathlineto{\pgfpoint{-5.003574\du}{12.758301\du}}
\pgfpathlineto{\pgfpoint{-5.007301\du}{12.756204\du}}
\pgfpathlineto{\pgfpoint{-5.011261\du}{12.754108\du}}
\pgfpathlineto{\pgfpoint{-5.015221\du}{12.752244\du}}
\pgfpathlineto{\pgfpoint{-5.019647\du}{12.749915\du}}
\pgfpathlineto{\pgfpoint{-5.024306\du}{12.748517\du}}
\pgfpathlineto{\pgfpoint{-5.029431\du}{12.746886\du}}
\pgfpathlineto{\pgfpoint{-5.034789\du}{12.744790\du}}
\pgfpathlineto{\pgfpoint{-5.039914\du}{12.743625\du}}
\pgfpathlineto{\pgfpoint{-5.046203\du}{12.742227\du}}
\pgfpathlineto{\pgfpoint{-5.052027\du}{12.740597\du}}
\pgfpathlineto{\pgfpoint{-5.058084\du}{12.739432\du}}
\pgfpathlineto{\pgfpoint{-5.065072\du}{12.738267\du}}
\pgfpathlineto{\pgfpoint{-5.071594\du}{12.737102\du}}
\pgfpathlineto{\pgfpoint{-5.078117\du}{12.736171\du}}
\pgfpathlineto{\pgfpoint{-5.085338\du}{12.735239\du}}
\pgfpathlineto{\pgfpoint{-5.092094\du}{12.734307\du}}
\pgfpathlineto{\pgfpoint{-5.099548\du}{12.733841\du}}
\pgfpathlineto{\pgfpoint{-5.107002\du}{12.733142\du}}
\pgfpathlineto{\pgfpoint{-5.114224\du}{12.733142\du}}
\pgfpathlineto{\pgfpoint{-5.121911\du}{12.732909\du}}
\pgfpathlineto{\pgfpoint{-5.129598\du}{12.732210\du}}
\pgfpathlineto{\pgfpoint{-5.137053\du}{12.732210\du}}
\pgfpathlineto{\pgfpoint{-5.137053\du}{12.732210\du}}
\pgfpathlineto{\pgfpoint{-5.144274\du}{12.732210\du}}
\pgfpathlineto{\pgfpoint{-5.151961\du}{12.732909\du}}
\pgfpathlineto{\pgfpoint{-5.159416\du}{12.733142\du}}
\pgfpathlineto{\pgfpoint{-5.166870\du}{12.733142\du}}
\pgfpathlineto{\pgfpoint{-5.174324\du}{12.733841\du}}
\pgfpathlineto{\pgfpoint{-5.181546\du}{12.734307\du}}
\pgfpathlineto{\pgfpoint{-5.188534\du}{12.735239\du}}
\pgfpathlineto{\pgfpoint{-5.195756\du}{12.736171\du}}
\pgfpathlineto{\pgfpoint{-5.202278\du}{12.737102\du}}
\pgfpathlineto{\pgfpoint{-5.209266\du}{12.738267\du}}
\pgfpathlineto{\pgfpoint{-5.215789\du}{12.739432\du}}
\pgfpathlineto{\pgfpoint{-5.221846\du}{12.740597\du}}
\pgfpathlineto{\pgfpoint{-5.228135\du}{12.742227\du}}
\pgfpathlineto{\pgfpoint{-5.233959\du}{12.743625\du}}
\pgfpathlineto{\pgfpoint{-5.239084\du}{12.744790\du}}
\pgfpathlineto{\pgfpoint{-5.244907\du}{12.746886\du}}
\pgfpathlineto{\pgfpoint{-5.249566\du}{12.748517\du}}
\pgfpathlineto{\pgfpoint{-5.254225\du}{12.749915\du}}
\pgfpathlineto{\pgfpoint{-5.259117\du}{12.752244\du}}
\pgfpathlineto{\pgfpoint{-5.263077\du}{12.754108\du}}
\pgfpathlineto{\pgfpoint{-5.267038\du}{12.756204\du}}
\pgfpathlineto{\pgfpoint{-5.270299\du}{12.758301\du}}
\pgfpathlineto{\pgfpoint{-5.273560\du}{12.760397\du}}
\pgfpathlineto{\pgfpoint{-5.276588\du}{12.762261\du}}
\pgfpathlineto{\pgfpoint{-5.279151\du}{12.764590\du}}
\pgfpathlineto{\pgfpoint{-5.280548\du}{12.767153\du}}
\pgfpathlineto{\pgfpoint{-5.282412\du}{12.769249\du}}
\pgfpathlineto{\pgfpoint{-5.283810\du}{12.771812\du}}
\pgfpathlineto{\pgfpoint{-5.284974\du}{12.773908\du}}
\pgfpathlineto{\pgfpoint{-5.285207\du}{12.776238\du}}
\pgfpathlineto{\pgfpoint{-5.285673\du}{12.778567\du}}
\pgfpathlineto{\pgfpoint{-5.285673\du}{12.778567\du}}
\pgfpathlineto{\pgfpoint{-5.285207\du}{12.781362\du}}
\pgfpathlineto{\pgfpoint{-5.284974\du}{12.783459\du}}
\pgfpathlineto{\pgfpoint{-5.283810\du}{12.785555\du}}
\pgfpathlineto{\pgfpoint{-5.282412\du}{12.788118\du}}
\pgfpathlineto{\pgfpoint{-5.280548\du}{12.790214\du}}
\pgfpathlineto{\pgfpoint{-5.279151\du}{12.793010\du}}
\pgfpathlineto{\pgfpoint{-5.276588\du}{12.795106\du}}
\pgfpathlineto{\pgfpoint{-5.273560\du}{12.797203\du}}
\pgfpathlineto{\pgfpoint{-5.270299\du}{12.799299\du}}
\pgfpathlineto{\pgfpoint{-5.267038\du}{12.801396\du}}
\pgfpathlineto{\pgfpoint{-5.263077\du}{12.803492\du}}
\pgfpathlineto{\pgfpoint{-5.259117\du}{12.805123\du}}
\pgfpathlineto{\pgfpoint{-5.254225\du}{12.807220\du}}
\pgfpathlineto{\pgfpoint{-5.249566\du}{12.808850\du}}
\pgfpathlineto{\pgfpoint{-5.244907\du}{12.810714\du}}
\pgfpathlineto{\pgfpoint{-5.239084\du}{12.812344\du}}
\pgfpathlineto{\pgfpoint{-5.233959\du}{12.813975\du}}
\pgfpathlineto{\pgfpoint{-5.228135\du}{12.815373\du}}
\pgfpathlineto{\pgfpoint{-5.221846\du}{12.816538\du}}
\pgfpathlineto{\pgfpoint{-5.215789\du}{12.818168\du}}
\pgfpathlineto{\pgfpoint{-5.209266\du}{12.819333\du}}
\pgfpathlineto{\pgfpoint{-5.202278\du}{12.820498\du}}
\pgfpathlineto{\pgfpoint{-5.195756\du}{12.821429\du}}
\pgfpathlineto{\pgfpoint{-5.188534\du}{12.822361\du}}
\pgfpathlineto{\pgfpoint{-5.181546\du}{12.822827\du}}
\pgfpathlineto{\pgfpoint{-5.174324\du}{12.823759\du}}
\pgfpathlineto{\pgfpoint{-5.166870\du}{12.823992\du}}
\pgfpathlineto{\pgfpoint{-5.159416\du}{12.824458\du}}
\pgfpathlineto{\pgfpoint{-5.151961\du}{12.824691\du}}
\pgfpathlineto{\pgfpoint{-5.144274\du}{12.825157\du}}
\pgfpathlineto{\pgfpoint{-5.137053\du}{12.825157\du}}
\pgfpathlineto{\pgfpoint{-5.137053\du}{12.825157\du}}
\pgfpathlineto{\pgfpoint{-5.129598\du}{12.825157\du}}
\pgfpathlineto{\pgfpoint{-5.121911\du}{12.824691\du}}
\pgfpathlineto{\pgfpoint{-5.114224\du}{12.824458\du}}
\pgfpathlineto{\pgfpoint{-5.107002\du}{12.823992\du}}
\pgfpathlineto{\pgfpoint{-5.099548\du}{12.823759\du}}
\pgfpathlineto{\pgfpoint{-5.092094\du}{12.822827\du}}
\pgfpathlineto{\pgfpoint{-5.085338\du}{12.822361\du}}
\pgfpathlineto{\pgfpoint{-5.078117\du}{12.821429\du}}
\pgfpathlineto{\pgfpoint{-5.071594\du}{12.820498\du}}
\pgfpathlineto{\pgfpoint{-5.065072\du}{12.819333\du}}
\pgfpathlineto{\pgfpoint{-5.058084\du}{12.818168\du}}
\pgfpathlineto{\pgfpoint{-5.052027\du}{12.816538\du}}
\pgfpathlineto{\pgfpoint{-5.046203\du}{12.815373\du}}
\pgfpathlineto{\pgfpoint{-5.039914\du}{12.813975\du}}
\pgfpathlineto{\pgfpoint{-5.034789\du}{12.812344\du}}
\pgfpathlineto{\pgfpoint{-5.029431\du}{12.810714\du}}
\pgfpathlineto{\pgfpoint{-5.024306\du}{12.808850\du}}
\pgfpathlineto{\pgfpoint{-5.019647\du}{12.807220\du}}
\pgfpathlineto{\pgfpoint{-5.015221\du}{12.805123\du}}
\pgfpathlineto{\pgfpoint{-5.011261\du}{12.803492\du}}
\pgfpathlineto{\pgfpoint{-5.007301\du}{12.801396\du}}
\pgfpathlineto{\pgfpoint{-5.003574\du}{12.799299\du}}
\pgfpathlineto{\pgfpoint{-5.000312\du}{12.797203\du}}
\pgfpathlineto{\pgfpoint{-4.997750\du}{12.795106\du}}
\pgfpathlineto{\pgfpoint{-4.995188\du}{12.793010\du}}
\pgfpathlineto{\pgfpoint{-4.993557\du}{12.790214\du}}
\pgfpathlineto{\pgfpoint{-4.991926\du}{12.788118\du}}
\pgfpathlineto{\pgfpoint{-4.990296\du}{12.785555\du}}
\pgfpathlineto{\pgfpoint{-4.989364\du}{12.783459\du}}
\pgfpathlineto{\pgfpoint{-4.988898\du}{12.781362\du}}
\pgfpathlineto{\pgfpoint{-4.988665\du}{12.778567\du}}
\pgfusepath{stroke}
\pgfsetbuttcap
\pgfsetmiterjoin
\pgfsetdash{}{0pt}
\definecolor{dialinecolor}{rgb}{0.000000, 0.000000, 0.000000}
\pgfsetfillcolor{dialinecolor}
\pgfpathmoveto{\pgfpoint{-5.127735\du}{12.769249\du}}
\pgfpathlineto{\pgfpoint{-5.109099\du}{12.778567\du}}
\pgfpathlineto{\pgfpoint{-5.053425\du}{12.760164\du}}
\pgfpathlineto{\pgfpoint{-5.025471\du}{12.769249\du}}
\pgfpathlineto{\pgfpoint{-5.043874\du}{12.750613\du}}
\pgfpathlineto{\pgfpoint{-5.109099\du}{12.750613\du}}
\pgfpathlineto{\pgfpoint{-5.081145\du}{12.750613\du}}
\pgfpathlineto{\pgfpoint{-5.127735\du}{12.769249\du}}
\pgfusepath{fill}
\pgfsetbuttcap
\pgfsetmiterjoin
\pgfsetdash{}{0pt}
\definecolor{dialinecolor}{rgb}{0.000000, 0.000000, 0.000000}
\pgfsetfillcolor{dialinecolor}
\pgfpathmoveto{\pgfpoint{-5.137053\du}{12.797203\du}}
\pgfpathlineto{\pgfpoint{-5.164774\du}{12.787885\du}}
\pgfpathlineto{\pgfpoint{-5.211363\du}{12.806521\du}}
\pgfpathlineto{\pgfpoint{-5.239084\du}{12.797203\du}}
\pgfpathlineto{\pgfpoint{-5.229533\du}{12.825157\du}}
\pgfpathlineto{\pgfpoint{-5.164774\du}{12.825157\du}}
\pgfpathlineto{\pgfpoint{-5.192494\du}{12.816072\du}}
\pgfpathlineto{\pgfpoint{-5.137053\du}{12.797203\du}}
\pgfusepath{fill}
\pgfsetbuttcap
\pgfsetmiterjoin
\pgfsetdash{}{0pt}
\definecolor{dialinecolor}{rgb}{0.000000, 0.000000, 0.000000}
\pgfsetfillcolor{dialinecolor}
\pgfpathmoveto{\pgfpoint{-5.229533\du}{12.750613\du}}
\pgfpathlineto{\pgfpoint{-5.211363\du}{12.741528\du}}
\pgfpathlineto{\pgfpoint{-5.155689\du}{12.769249\du}}
\pgfpathlineto{\pgfpoint{-5.137053\du}{12.760164\du}}
\pgfpathlineto{\pgfpoint{-5.145905\du}{12.778567\du}}
\pgfpathlineto{\pgfpoint{-5.211363\du}{12.778567\du}}
\pgfpathlineto{\pgfpoint{-5.183409\du}{12.769249\du}}
\pgfpathlineto{\pgfpoint{-5.229533\du}{12.750613\du}}
\pgfusepath{fill}
\pgfsetbuttcap
\pgfsetmiterjoin
\pgfsetdash{}{0pt}
\definecolor{dialinecolor}{rgb}{0.000000, 0.000000, 0.000000}
\pgfsetfillcolor{dialinecolor}
\pgfpathmoveto{\pgfpoint{-5.034789\du}{12.816072\du}}
\pgfpathlineto{\pgfpoint{-5.053425\du}{12.825157\du}}
\pgfpathlineto{\pgfpoint{-5.109099\du}{12.797203\du}}
\pgfpathlineto{\pgfpoint{-5.137053\du}{12.806521\du}}
\pgfpathlineto{\pgfpoint{-5.118184\du}{12.787885\du}}
\pgfpathlineto{\pgfpoint{-5.053425\du}{12.787885\du}}
\pgfpathlineto{\pgfpoint{-5.081145\du}{12.797203\du}}
\pgfpathlineto{\pgfpoint{-5.034789\du}{12.816072\du}}
\pgfusepath{fill}
\pgfsetbuttcap
\pgfsetmiterjoin
\pgfsetdash{}{0pt}
\definecolor{dialinecolor}{rgb}{1.000000, 1.000000, 1.000000}
\pgfsetfillcolor{dialinecolor}
\pgfpathmoveto{\pgfpoint{-5.127735\du}{12.778567\du}}
\pgfpathlineto{\pgfpoint{-5.109099\du}{12.778567\du}}
\pgfpathlineto{\pgfpoint{-5.053425\du}{12.760164\du}}
\pgfpathlineto{\pgfpoint{-5.025471\du}{12.769249\du}}
\pgfpathlineto{\pgfpoint{-5.043874\du}{12.750613\du}}
\pgfpathlineto{\pgfpoint{-5.099548\du}{12.750613\du}}
\pgfpathlineto{\pgfpoint{-5.081145\du}{12.750613\du}}
\pgfpathlineto{\pgfpoint{-5.127735\du}{12.778567\du}}
\pgfusepath{fill}
\pgfsetbuttcap
\pgfsetmiterjoin
\pgfsetdash{}{0pt}
\definecolor{dialinecolor}{rgb}{1.000000, 1.000000, 1.000000}
\pgfsetfillcolor{dialinecolor}
\pgfpathmoveto{\pgfpoint{-5.137053\du}{12.797203\du}}
\pgfpathlineto{\pgfpoint{-5.155689\du}{12.787885\du}}
\pgfpathlineto{\pgfpoint{-5.211363\du}{12.806521\du}}
\pgfpathlineto{\pgfpoint{-5.239084\du}{12.806521\du}}
\pgfpathlineto{\pgfpoint{-5.220448\du}{12.825157\du}}
\pgfpathlineto{\pgfpoint{-5.155689\du}{12.825157\du}}
\pgfpathlineto{\pgfpoint{-5.183409\du}{12.816072\du}}
\pgfpathlineto{\pgfpoint{-5.137053\du}{12.797203\du}}
\pgfusepath{fill}
\pgfsetbuttcap
\pgfsetmiterjoin
\pgfsetdash{}{0pt}
\definecolor{dialinecolor}{rgb}{1.000000, 1.000000, 1.000000}
\pgfsetfillcolor{dialinecolor}
\pgfpathmoveto{\pgfpoint{-5.229533\du}{12.750613\du}}
\pgfpathlineto{\pgfpoint{-5.211363\du}{12.750613\du}}
\pgfpathlineto{\pgfpoint{-5.155689\du}{12.769249\du}}
\pgfpathlineto{\pgfpoint{-5.127735\du}{12.760164\du}}
\pgfpathlineto{\pgfpoint{-5.145905\du}{12.778567\du}}
\pgfpathlineto{\pgfpoint{-5.211363\du}{12.778567\du}}
\pgfpathlineto{\pgfpoint{-5.183409\du}{12.778567\du}}
\pgfpathlineto{\pgfpoint{-5.229533\du}{12.750613\du}}
\pgfusepath{fill}
\pgfsetbuttcap
\pgfsetmiterjoin
\pgfsetdash{}{0pt}
\definecolor{dialinecolor}{rgb}{1.000000, 1.000000, 1.000000}
\pgfsetfillcolor{dialinecolor}
\pgfpathmoveto{\pgfpoint{-5.034789\du}{12.816072\du}}
\pgfpathlineto{\pgfpoint{-5.053425\du}{12.825157\du}}
\pgfpathlineto{\pgfpoint{-5.109099\du}{12.806521\du}}
\pgfpathlineto{\pgfpoint{-5.127735\du}{12.806521\du}}
\pgfpathlineto{\pgfpoint{-5.118184\du}{12.787885\du}}
\pgfpathlineto{\pgfpoint{-5.053425\du}{12.787885\du}}
\pgfpathlineto{\pgfpoint{-5.081145\du}{12.797203\du}}
\pgfpathlineto{\pgfpoint{-5.034789\du}{12.816072\du}}
\pgfusepath{fill}
\pgfsetbuttcap
\pgfsetmiterjoin
\pgfsetdash{}{0pt}
\definecolor{dialinecolor}{rgb}{0.666667, 0.901961, 1.000000}
\pgfsetstrokecolor{dialinecolor}
\pgfpathmoveto{\pgfpoint{-5.285673\du}{12.778567\du}}
\pgfpathlineto{\pgfpoint{-5.285673\du}{12.852644\du}}
\pgfusepath{stroke}
\pgfsetbuttcap
\pgfsetmiterjoin
\pgfsetdash{}{0pt}
\definecolor{dialinecolor}{rgb}{0.666667, 0.901961, 1.000000}
\pgfsetstrokecolor{dialinecolor}
\pgfpathmoveto{\pgfpoint{-4.988665\du}{12.778567\du}}
\pgfpathlineto{\pgfpoint{-4.988665\du}{12.852644\du}}
\pgfusepath{stroke}
% setfont left to latex
\definecolor{dialinecolor}{rgb}{0.000000, 0.000000, 0.000000}
\pgfsetstrokecolor{dialinecolor}
\node[anchor=west] at (-3.200000\du,12.792800\du){Interested Client};
% setfont left to latex
\definecolor{dialinecolor}{rgb}{0.000000, 0.000000, 0.000000}
\pgfsetstrokecolor{dialinecolor}
\node[anchor=west] at (-3.200000\du,15.292800\du){Uninterested Client};
\end{tikzpicture}
\end{flushleft}
We refer to Figure 1.1 . The source generates a data steam with throughput
of 1Mbps, and the data stream is to be received by three recipients. On the upper
left-hand side we see a unicast transmission where the source sends three instances
of the same data to the receivers. The network link has to carry the same data
stream multiplied by three, and so 3Mbps of bandwidth easily consumed. Whereas,
the one on the bottom, which is the essence of multicasting, the data is replicated
only when it needs to at the junction points. That, in fact, is one of the performance
measures of multicasting routing protocols. We can tell from figure 1.1, it saves
2Mbps for the up until the fan-out point and saves another 1Mbps for the network
on the rightmost side. It decouples the link bandwidth dependence on the number
of hosts wanting that traffic. The disadvantage of broadcasting1 in relation can
1This has caused the creation of violent terms like packet flooding and broadcast storm.
inferred from Figure 1.1- which is uninterested hosts and routers get copy of the
data. Though we do not want to proceed with the point- in multicast routing, it is
desirable effect that the data travels long distance with the least “fan-out” possible
thus saving a lot of bandwidth along its way to its final destination.
\section{WHAT IS MULTICAST FIREWALL?}
Well it is all good, we can save all the bandwidth with scale of 1/n where n
is the number of unicast computers in the network that have to receive the
traffic.
The notion of unicast firewalling scheme breaks down when we are dealing with
multicast traffic. In this thesis, when we say firewall2 what we are referring
is the
filtering type, which drops or accepts traffic, based on packet header
information3.
A general golden rule for firewall that maintain state, known as stateful
firewall, is
to trust an outsider that is contacted by an insider. Keeping that in mind, we
now
deal with firewalling multicast.
A multicast source is a host that sends packets with the destination IP
address
set to a multicast group. Anyone can send traffic to that group! A source need
not
to be a member of the group; sourcing and listening are mutually exclusive.
The
reason why the firewall golden rule no longer works is the destination address
is a
group address which is in some sense virtual. How can we establish which
actual
outsider our client is trying to contact?
To compound the problem these addresses usually are volatile. There have
been/are efforts by The Internet Assigned Numbers Authority (IANA) and various
standard bodies to get some order on the address allocations and how ISPs
should
manage these addresses, should be used and various recommendations. For this
thesis, we are better off to assume that these addresses as mentioned above
are not
static. We are also going to be dealing with access list control exclusively.
The
authentication mechanism of sources is mainly the firewall golden rule.
So, let’s outline why a firewall maybe needed in controlling access list
(whose
traffic is authorized to come in) in multicasting:
• To protect the whole internal network from unwanted traffic
• To protect client (insider) group members from unwanted traffic (individ-
ual tailored) from outside or within the network and
• Possibly, preventing an unauthorized user within or outside the network
from accessing our traffic.
The list is not exhaustive but helps to flesh out what we want achieve. The
first two
are for network mantainablity while the last one is question of privacy.
Access list
is non-trivial especially if dealing with insiders. Some insecurities in
multicasting
are unavoidable from its designs that ironically make it attractive too. We
discuss
the last two points more in Chapter 6.
\section{Problem Statement}
From a Network Manager’s perspective the biggest worry is Denial of Ser-
vice/Distrubuted denial-of-service (DoS/DDoS)attack on the network as
multicas-
ting amplifies it because multiple resources can be members- (e.g an outsider
dumps
traffic on the stock exchange network). All innocent group members become
victim
as opposed to unicast where only one user4 is affected. To defend that, the
usual
modus operandi in SMEs is for the Network Manager to simply block all incom-
ing multicast traffic. Then, if someone must join a multicast group, the
Network
Manager allows the unique source and multicast address manually and traffic
starts
to come in from that particular source. Now one can imagine what could be the
bureaucratic process to achieve that.
The problem statement is : To design a multicast firewall that relieves the
Network Manager from entering static rules per session.
Being rather prescriptive, our design solution should satisfy the following
re-
quirements in order to relieve the Network Manager from doing that:
(1) Block all incoming multicast traffic
(2) Inspect unicast traffic that come out from clients inside our network-
(these usually have distinct ports)- to track which unicast source addresses
our clients are contacting
(3) Inspect all multicast group subscription of our clients
(4) Allow incoming multicast group traffic from sources
(a) That have associations with clients from step 2.
(b) Those clients are also subscribed that group in step 3.
2-4 are restatements of the paper written on this same topic by Vijay et
al[8]. To
elaborate on step 2 a bit, it can, for example, be watching over all TCP
connections
on Port 80. That is when a user is simply browsing the web. We say more on
this
on later parts.
\section{Contribution}
Contribution of this thesis involves:
• An algorithm that avoids as many memory allocations as possible.
• A data structure that help models all entities involved and algorithms
along with it
• The overall design
The first contribution needs acknowledgement. My supervisor, Vijay Sivaraman,
indicated this is area of bottleneck from previous student implementation.
This is
to early to discuss more contribution but hopefully we are able to show in
Chapter
4 during discussion of detailed design.
\section{Document Scope}
Questions we address.
• What kind of security model multicast entails and thus the paradigm we
follow?
• What is the design solution and its implementation?
• What is the prototype design environment?
• The comparison of the design solution with other implementations.
• How is the performance with the tests conducted so far?
We are quite careful to avoid discussion of multicast routing in most cases.
It is
quite involved on its own and surely this thesis will only address such issues
which
we can exploit in the design of our firewall. Often, what we want is to
extract the
unicast address of the multicast source efficiently- an address one of our
clients is
trying to contact.
Another key fact that one has to know when working on the network stack is
it is can never be decoupled how the underlying Operating System handles IP.
It
requires an intimate knowledge of operating systems networking stack.
Lastly on light tone, I will jestly tell owner of a computer with 32-bit x86
single
processor as very lucky person unlike me who owns lowest end of ’fancy’
multicore.
A newer issue has propped with the advancement of multi-core operating systems
too. We have to model our data and our programs to achieve performance taking
the fact that there is more than one processor on the computer. Oh before that
though, avoiding deadlocks, bus errors and racing are our challenges. This
entails
questions like do processors share the same memory address space, the same bus
and so on? In fact, the author has spent more time dealing with the two above
issues
than the logical design of this firewall. He had two manually reboot his
computer
one two many times on segmentation faults that result in computer hangs . We
will not discuss those topics in this thesis as they can be digressions.
\chapter{Background: Protocols, Firewalls and Related
works}
The purpose of this chapter is to review some prerequisite material in the
design
of our firewall. We try to filter in only those that are required for
dispensing the idea
our firewall. These are mainly IP multicast networking concepts.
Implementation
wise though, we also need to have a firm grasp of the prototyping environment,
the encapsulating framework ( for our case, these are Linux kernel and
Netfilter
respectively) and how the life of a packet evolves in those environments. We
try
to minimize, at least defer, discussion of these latter issues as they are
tied to
implementation. We focus mainly on IPv4 as our implemented module works on it
and the key ideas can be fully elaborated using that version of the IP. We try
to shed
some light at multicast enabler protocols with special emphasis given to the
ones we
need for this thesis most. It also introduces us to ’related’ work. The
discussion of
IPv6 diversion and design challenged is a great omission. But we give some
pointers
to equivalent protocols-like IGMP for Multicast Listener Discovery (MLD) - as
we
are about to see.
\section{Protocols in Multicasting}
When we talk about a working multicasting system, a failure to mention the
combination of the enabling protocols involved leaves a big gap of
information.
Usually, Network Engineers put these separated in hyphen or slash in triplet
or
so (e.g MSDP/MBGP/PIM in IPv4). This is only to stress multicasting is only a
model- not actual technology. So is unicasting but it perfectly maps to IP
without
any other protocol (needless to say, IP was designed with unicasting model in
mind).
\section{Group Management Protocols}
One of the concepts we referred in chapter 1, is this notion of predefined set
of
hosts. Internet Group Management Protocol (IGMP), as the name suggests, is the
MLD is IPv6 coun-
protocol that helps the group manager,which is usually the nearest router,
manipu- terpart of IGMP.
late the group members. A prospect member host must be able to construct IGMP
messages if its wants to participate in a multicast group. A Multicast
manager, on
the other hand, must also be able to manage the list. IGMP is encapsulated in
IP
datagram by setting the protocol field to IPPROTO\_IGMP1 - which is a constant
2. The following steps are involved while managing a group.
• The router asks every period (of about 125 seconds) if anyone wants any
multicast message. By default, every multicast capable device must always
Router IGMP Query
be subscribed to one address to receive this. This is, for example, 224.0.0.1
in IPv4.
• Hosts that want to join a group or keep their membership alive reply
(report) with the group address they are still participating or want to
Host IGMP Report
start participating.
• If the router receives no interest from at least a single host to be in a
group, that group is deleted (assuming it exists) and the router notifies
his upper peer (if inter-domain multicasting) that it is no longer interested
in that group.
The list outlines the baseline procedure that makes this group management
work.
This was exactly what was implemented in version 1 of IGMP. Further on the
next
iterations various optimization were added which are discussed below. The
latest
IGMP version number is 3. Version 0 is mere history.
We should take note from first point that a group’s age is about 2 minutes
which we exploit on later in our Implementation.
\subsection{The law of backward compatibility.} The IETF standards explic-
itly require the group management ’language’ be the lowest IGMP version among
members if at least one member has that version as its upper limit support.
This
entails for the greater versions of IGMP to be backward compatible with their
predecessors.
\subsection{IGMP (v1 and v2)}. We discuss both IGMP[3] versions for compar-
isons. Version 2 introduces ’Explicit Leave group’ message by the host when it
no longer wants to receive traffic. This is to decrease link load by issuing a
leave
message than waiting for the group membership timer to expire; which is
followed
by the router asking if any one wants to receive traffic. The ’leave latency’,
the
time between when a user stops wanting to receive traffic and the router knows
that, will decrease. It is an optimization on the version 1 which simply
’shuts up’
when it no longer wants to receive traffic. Thus, the discussion IGMPv2
includes
both. We are mainly concerned with report messages of clients as that is where
we extract the group a host wants to join. The following is header information
of
IGMPv22:
%packet
0
7
type
15
max resp time
31
checksum
group address
IGMP payload/ future expansion
Figure 2.1. IGMPv1/v2 Message Format
%figure
The fact is we don’t even need to dissect the IGMP header to extract the group
it is destined for. It can be found in the IP datagram header. That changes
for
leave messages; where the destination address of the IP header is set to
224.0.0.2
(the all-routers group). The following is some information about the fields:
• The type describes what kind of IGMP message it is. From it we know
if it router query, host leave message or host membership report. We are
often concerned with IGMPV{x}\_HOST\_MEMBERSHIP\_REPORT.
• The max resp time is heuristic based constant, which tells the host about
when to respond with a report after a router sent query to all hosts. It is
about 1 second and can be set to other values in IGMPv2 or later.
2.2.3. IGMPv3. As can be seen on Figure 2.2 [1]this version have a radi-
cal change from the previous in that the header did change dramatically. It
also
changes the service model of traditional multicasting by coupling unicasting
as the
underlying framework. A new idea of ’the user should know where the multicast
source is’ introduces that one can add:
MLDv2 is the coun-
• Only those unicast source addresses that one want to hear from. (Include
mode)
• All except those source addresses that one want hear from. (Exclude
mode)
The sudden complexity and bloatness of this design is worth as it is proven by
experience it works for one type of model (one to many multicasting). We
discuss
in detail multicasting models in Section 2.4.
ail multicasting models in Section 2.4.
0
7
type=0x22
15
31
Reserved
Reserved
checksum
Number of Group Records (M)
Group Record [1]
Group Record [2]
.
.
.
Group Record [M]
Figure 2.2. IGMP Version 3 Membership Report Message
The group record further decomposes to the following. We find the list of
sources we want to include/exclude per multicast address.
One of the great things about this protocol is the fact it has breathed fresh
air
to otherwise sleeping multicasting usage in the industry. It especially works
great
if we just have a few multicast sources per group. We want to see it from the
angle
of our problem though. Now, it has the capacity to block static addresses when
used in exclude \& include mode. Yes sure, one can use it to block an annoying
user
’on the fly’ participating in the multicast group using exclude mode (for
example,
quite handy when a user is playing Multiplayer online games and wants to block
0
7
Record Type
15
Aux Data Len
9
31
Number of Sources (N)
Multicast group address
Source Address [1]
Source Address [2]
.
.
.
Source Address [N]
Auxiliary Data
.
.
.
Figure 2.3. IGMP Version 3 Group Record Internal Format
someone). In the end, we can only list so much unicast addresses in static
form
that it cannot be of practical usage for protection against DoS/DDoS attacks
when
all members are both sources and receivers- the traditional model.
Essentially, the
verdict is that it is quite a capable protocol if we want a one to many
dissemination
type communication.3 But, multicasting has two other extra paradigms-many to
one and many to many- and as such it solves 1/3rd of the problem.
There is also argument against it in routing inefficiency when the number of
sources increase. The router could easily run out of memory if it is going to
keep
lots of state information per client. The data path distribution tree is also
sub-
optimal. In this thesis though, we do not pursue these as our arguments for
our
design. IGMPv3 is still considered fairly new (in 2009) and we expect to have
hosts and routers that only work with the earlier versions. As an example,
Linux
kernel networking stack integrated it in 2002 and FreeBSD4 still has not
integrated
it. Even if it is integrated, more often than not it is working with
compatibility
mode we discussed above. This adds some weight as why we should pursue such
endeavor.
A salient point in relation to backward compatibility is an insider can expose
the network for DoS/DDoS attack by reporting using lower IGMP versions. This
is because the language of that group is turned into one of the earlier
versions and
there is no longer source filtering in place.
\section{A Word About Multicast Routing Protocols}
It seems instructional to say a bit about multicast routing as some of the
protocols and service models that we are about to discuss have direct relation
to
them and lest the document may become incoherent. But, the filter module can
be
done away without much information on the routing protocols.
\footnote{’radio’ style communication where a subscriber/tuner can only receive.}
\footnote{Hard to claim but a little research on the kernel mailing list will show.
http://lists.freebsd.org/pipermail/freebsd-current/2009-March/004128.html}
Figure 2.4. Multicast Distribution Tree
Because there can be multiple receivers, the traffic path may have several
branches as we see in Figure 2.4. The whole data path from start the source to
receivers is known as a distribution tree. Data flow through the distribution
trees
is referred as either upstream and downstream. Downstream is in the direction
to-
ward the receivers. Upstream is in the direction toward the source. A
downstream
interface is an outgoing or outbound interface; and an upstream interface is
as an
incoming or inbound interface. This is all relative to source.
But let’s see how this distribution tree is formed following steps in Figure
2.5.
The following steps will take place.
(1) The source sends data packet to specific multicast group. Nobody is yet
interested; router A simply discards it.
(2) Listener 2 sends IGMP multicast report to router C.
(3) Router C sends a join (this is done with the routing protocol) on behalf
of Listener 2 to router B. But how did it know to ask B, not D? The
concept of Rendez-vous point-root router- of the distribution tree comes
here. Each multicast capable router (either dynamically or statically)
must know the root of the tree. We assumed the source is the root here-
which SSM also assumes as we will see soon.
(4) Router B to A as in (3) and then the route ABC is established.
As further elaboration of (3) if D was the Rendez-vous point configured in all
routers. Source contacts A, A contacts B and, B contacts D and notifies it
source is
multicasting. Then, listener contacts C, C contacts D. A path ABDC would have
been established. This path could be optimal in cases where there maybe
another
source but we leave that for the moment.
Figure 2.5. Multicast Pre-distribution Tree
Each sub network that contains at least one interested listener is a leaf on
the
tree. When a new listener tunes in, a new branch is built, joining the leaf to
the
tree. When a listener tunes out, its branch is pruned off the tree. Where the
tree
branches, routers replicate the data and send a single flow down each branch.
Thus
no link ever carries a duplicate flow of packets.
The above steps are roughly taken by Protocol Independent Multicasting-
Sparse Mode (PIM-SM) and its variants. The SM in PIM-SM refers to the fact the
multicast sources are sparse ( taken from graph-theory). That is as far we can
go
with multicast routing as Stevens[7] note it can easily consume a book! We
refer
to [6] if this section disappoints.
2.4. Any source vs. Source Specific Multicast
As we stressed earlier, multicasting is a model realized through various
under-
lying protocols. Multicasting itself is traditionally further divided into any
source
and source specific models5.SSM is due to Holbrook [9]. The any-source part is
of
\footnote{This can be misleading as SSM is subset of ASM. It would have been better to
say ASM is
divided into...But we stick to the literature.}
primary concern to this thesis.
If the user is active participant and does not know all the sources of the
groups
s/he wants to join, any-source multicasting is appropriate.
In our discussion of IGMPv3, we noted how the user can opt to specify sources
s/he wants to hear from. Then, it is no wonder that IGMPv3 is subset in real-
ization of SSM model. SSM uses PIM-SM for routing with the simplification that
we already know the root of the distribution tree leading a much more strict
ver-
sion called PIM-SSM. That simple assumption eliminates the need for RPTs, RPs,
and Multicast Source Discovery Protocol (a lot of the ’hyphens and slashes’ we
mentioned in Section 2.1), radically simplifying the mechanisms needed to
deliver
multicast. It somehow becomes the shortest path problem between the source and
the user and reduces to unicasting routing model. This no longer considers
’group
optimization. As mentioned Section2.2.3, support of IGMPv3 is a function of
the
operating system as of this writing. It may take some time before the majority
of hosts on the Internet are enabled for IGMPv3. But also, the restriction of
this
solution to one model makes it only part of the solution to multicast security
we
noted earlier. We emphasize the reliance of ASM on the ideas of the earlier
versions
of IGMP (v1 \& v2)- the less strict ones.
By now, we have established our playground, how it differs with its closest
ally
solution- SSM. We should take SSM in the form of ’related work’ to our design
solution as we target SSM’s super set, ASM.
\section{Related Works}
We have implicitly stated some existing mechanisms like using manual entry of
the static addresses in a firewall, and using SSM. We now add some more.
The solutions are roughly divided into two: those that use cryptography mech-
anisms and those that use various other mechanisms/heuristics to identify the
gen-
uine sources-like ours.
\subsection{IETF MSEC.} The IETF Multicast Security Working Group (MSEC
WG)MSEC WG[13] is the standard body that is working on securing multicasting.
The MSEC architecture focuses mainly to using cryptographic mechanisms. It
deals
with such things as privacy, key management, encryption and is biased towards
securing the data than the network. One of its attractive application geared
towards
our problem is the shared key mechanism to authenticate group members. To
elaborate it a little more, a member can only join a group if it has the
shared key
which can be established by the standard challenge-response mechanism. RFC3740
goes into the detail of how it should be done. It is dogged with questions
like who
is going to be the group manager? How is the membership dynamics? But most
importantly do we want to receive traffic from all those who joined the group?
The use of encryption and decryption on every packet has heavy performance
hit on real-time data- which multicasting is used for often. Further on this
can be
found in [11].
\subsection{Other mechanisms.} These mechanisms can be seen under one um-
brella including ours. The ultimate objective is to let the firewall know
’these are
the sources we trust’. For example a solution proposed by [12] seeks to
implement
a private protocol to notify the firewall dynamically the genuine sources that
are
being contacted. This has very strong assumption on the firewall and the
multicast
application developer.
Another one is by converging Any Source Multicasting (ASM) to Source Specific
Multicasting (SSM). We have thoroughly discussed the models in Section 2.4. It
naturally follows if we are using full potential of SSM and are only
interested in the
service model, we can be secure from DoS/DoS attacks. The following is done to
make SSM-incapable hosts (due to IGMPv3) who rely on traditional multicast to
join SSM-enabled multicasting
Figure 2.6. DNS-based Solution
\subsubsection{SSM-Mapping (from Cisco).} The following should not be confused with
multicast Domain Name Server (DNS) (mDNS) which is used for LAN resource dis-
covery in line with ousting broadcast resource discovery. SSM-Mapping[10]
changes
the earlier version of IGMP to IGMPv3 ’on the fly’ thus changing the model.
Figure
2.6 tries to show the process.
• Router maps IGMPv2 Joins (in SSM address range) to well- known sources
via DNS
• It usually allows only for one, or more, sources per group
• Router maps group to source (sources) – Uses either DNS or static internal
database
\chapter{Design Concepts}
\mquot{If I were stranded on a desert island and could only take one data
structure with me, it would be the hash table.\\
-Peter van der Linden\\
\footnote{It is actually parody on ’C-compiler’, Expert C programming: Deep C secrets}
}
The purpose this chapter is to introduce to some key concepts that propped in
the
design of our program. We also walk on the program on a system level what a
packet that enters our module will happen to it. Some of the discussion may
seem
appropriate for the previous chapter but its pertinence is immediate matter to
the
design we put it alongside here.
3.1. Desired Properties for the Multicast Module
We want to see this from the perspective of two types of people, namely, the
Network Engineer and the user- security \& usablity. Let’s further assume as an
example: our user is watching Multicast TV programs- on a multi-channel IPTV
capable box. This is to illustrate what users want and observe their behaviour
and
design our solution accordingly. We use the term channel interchangeably with
group in the example. So, from the perspective of the user:
(1) We want no delay to traffic, latency, and even worse varying delay on
traffic, jitter. The first one might be tolerable so long as it constant and
negligibly small for some applications such as Teleconferencing. This is re-
stating the classical satellite delay problem- what people encounter as con-
versational gaps and synchronisation problems. After certain threshold-
which depends on the application, delay creates confusion among the par-
ticipants. In short, making sure zero jitter, we still can allow ’tolerable’
non-zero latency.
(2) Tolerable delay (for the user) when the user wants to joins a channel.
(3) Tolerable delay (for the user) when a user wants to leave a channel.
The first point is the reason why multicast uses UDP (user datagram protocol),
which has the principle of ’an early packet with error is better than a late
correct
packet’ (best-effort delivery). We want to maintain this principle in our
multicast
firewall design. This entails one specific requirement when inspecting the
packet
for validity to either drop or accept it at our filter. We need to have a fast
look
up if the source is valid. This justifies the choice of our data structure
which we
discuss in Section 3.2.
The three list above takes into account the behaviour of the user in our
example.
One key observation is the following: when a user leaves a channel, it is
either to
join another one or is leaving for good (finished watching for the day). In
the first
part there is genuine leave, in which case the user actually has left the
channel for
considerable amount of time. The other leave is comeback leave. In this case,
the
user has left with the intention to come back soon. For example, the user
could
be avoiding watching commercials on that Channel. Users can be annoyed if they
cannot flip back to the previous channel with the back button on the remote.
This
same behaviour can be easily mapped to perspective of different types of
multicast
applications. This idea is captured in our algorithm design in Chapter 4. It
also
plays a role in what should be included our data structures and memory
allocation
block.
While most of the requirements from the perspective of the Network Engineer
has been outlined in Chapter 1, we still want to list some more desired
properties.
We stress again our multicast firewall is about resource maintainability in
the event
of attacks (mainly DoS/DDoS). Thus, so much as we want to avoid exposing our
network to attacks, we also want to avoid undesired heavy-handedness on
internal
users.
(1) Our filter should minimize blocking genuine traffic. One special reason
that haunts our design to cause false positives2 is unicast snooping short-
fall. In Chapter 6, we outline ways to mitigate this problem of ’network
auto-immune disease’.
(2) Our filter should minimize false-negatives. This is of lesser priority
than
the above. It can be simple ’annoyance’ unless orchestrated intentionally
attack is taking place. Remember we begin with a completely blocked
wall. We are trying to balance usability and maintainability. In this case,
the CEO will not ask “why is the teleconferencing program not working?”
as the case maybe for the first point.
This leads us to discussion of the data structure of our choice.
3.2. Hashtables
Hashing implements the mathematical concept of a function. We choose how
our domain changes to the set of values we want it to be transformed to via
our
hashing function. Depending on the input range and output range, hashing has
different applications.
Hashtables, which are one part of the application of hashing and the relevant
part for this thesis, are about storing values. If we know beforehand we have
n
distinct inputs (also known as key values) and n associated data values, we
store
these values in an array of n size by transforming the inputs into range 0-n
via our
hash function. This is known as perfect hashing. In most cases though, we do
not
know the value and number of the inputs beforehand. So, we simply allocate n
units
2\footnote{from firewall point of view, the positives are the ones that are tested to be
malicious}
of memory of the output types and put a mechanism for collision handling when
two inputs map to the same address location-bucket. These collisions are often are
handled by linked list data structures. Of course, we can use any data structure
including a chained hash table. In this thesis, we use linked lists. This is a design
choice of simplicity before being fancy without knowing if that is the bottleneck-
Rob Pike advices[18]. When we have no prior knowledge our hash key and values,
we aspire to have hashing function that tries to spread new values uniformly across
the array. Some of the desirable features (among others) of a hash function for the
spreading to happen include:
• Good mixing and combining
• Good speed (input to output)
• Good avalanche effect
The first and third points are to make sure every bit, as in 1/0, of the data has
effect on producing the hash value.
Before we digress too much, our main concerns is storing IP addresses. Imagine
an internal client with IP address, c, contacts a destination source address, s. In
our design, we want to store this association in buckets by using the destination
address (which is random) as our key value. For the points we mentioned above,
Jenkins Hash[16] due to Bob Jenkins is used after doing research. We refer for the
explanation of the above points to any Data Structures book .
\section{A word about Netfilter/Linux}
The Netfilter Framework comprises a set of hooks on the Linux networking
stack. A hook essentially is a function pointer and gives others the freedom to im-
plement the function themselves- it provides the prototype of the function though3.
Netfilter is Linux’s extensible firewall. Now, Netfilter has access to the packet as
argument in hook form at different stages once the packet enters the Linux kernel.
It then has the freedom to do stuff with packet including dropping it. There are
five different stages where the control of packet is temporally passed to Netfilter:
• PREROUTING: All the packets, with no exceptions, hit this hook, which
is reached before the routing decision and after all the IP header sanity
checks are fulfilled. Port Address Translation (NAPT) and Redirection
that is, Destination Network Translation (DNAT), are implemented in
this hook. We can manipulate the data too.
• LOCAL INPUT: All the packets going to the local machine reach this
hook4. This is the last hook in the incoming path for the local machine
traffic.
• FORWARD: Packets not going to the local machine (e.g., packets going
through the firewall) reach this hook. In a real system this will be how
3\footnote{This is C’s powerful way providing what one calls ’interface in object-oriented paradigm’}
4\footnote{As opposed to forwarding traffic.}
ur module functions. All the packets are checked before forwarding them
to one of our clients in the network.
• LOCAL OUTPUT: This is the first hook in the outgoing packet path.
Packets leaving the local machine always hit this hook.
• POSTROUTING: This hook is implemented after the routing decision.
Source Network Address Translation (SNAT) is registered to this hook.
All the packets that leave the local machine reach this hook.
Figure 3.1. Netfilter Components (From Jan Engelhardt)
This is a very powerful framework that is capable of doing much more than
we can possibly explain here. Then we have Xtables, which is shown on Figure5
3.1,on top of it as brick layer. It provides a simpler API so that we can extend
Netfilter functionalists by getting access to the packet we want in less complicated
way than Netfilter. So, our module is a plugin to this top layer. There is also a code
I wrote for userland interaction, this is at Iptables layer, which is mainly based on
boilerplate. This Section is very application specific and the appendix might have
been a better place.
\section{System Level Overview of the Module}
\subsection{Our Multicast Firewall.} The general concept of the filter was dis-
cussed before as ’prescriptive’ problem statement. We have also said the way we
discover who the sources are is by snooping outgoing unicast traffic. Figure 3.2 tries
to show what we are trying to achieve. Darth is the bad guy and Alice is the good
5\footnote{Jan have also graciously provided extensive tutorial on ’Writing Netfilter Modules’
and the glue API following changes in kernel. The former has Creative Commons and the latter
GPL licenses. All can be found including the picture: http://jengelh.medozas.de/}
source. The way we discovered Alice was good source is because someone inside
our network has been browsing alice.com. With this, we also expose our strongest
assumption in this thesis. That is, somehow that person/application is likely to
have the unicast contact of some form with the source. The Figure also shows how
Darth’s traffic is dropped by our firewall and Alice’s accepted.
Figure 3.2. The dynamic multicast firewall
The formal statement of the above is :
Notation 1. let C = {ci } be the set of internal clients.
Notation 2. let S = {si } be the set of external sources.
Notation 3. let G = {gi } be the set of multicast groups.
Then the dynamic multicast firewall claims:
Claim 4. If ci contacts sj and joined gi a traf f ic f rom sj to gi is trustworthy.
We can prove the claim by assuming ’trust chaining6’ and using transitivity
property.
The way we implemented the above claim is:
6\footnote{This is how digital certificates work too.}
For outgoing traffic. We keep two arrays of buckets:
(1) One for groups addresses we snoop when clients join groups. Either the
destination address is set to group address or the IGMP header has the
group address depending on the IGMP version as in 2.2.3.
(2) One for source addresses when clients contact contact sources. As a nor-
mal unicast, the destination address (of the IP datagram) is set to the
source address.
In both cases, we hash the destination address and put the client’s IP address that
contacted them.
For incoming traffic. The incoming traffic has both source and group addresses(as
the source is the one sending to the group). The IP datagram contains source ad-
dress as source address and group address as destination address. That means we
have both our hash keys. We follow the following steps:
(1) Hash the group address, check if there exists any client at that bucket.
(2) Hash the source address, check if there exists any client at that bucket.
(3) If either one of (1) \& (2) is false we drop the packet and finish here.
(4) We cross check those lists until at least one client is common in both.
Mathematically, the cardinality of the union of the two sets will be less
than sum of the cardinality of each.
The external IP addresses serve the dual purpose of authentication and being hash
keys for storing value of the internal clients.
Figure 3.3 shows code flow chart and the ideas mentioned here.
\section{Data Structures and algorithms}
We follow a key concept in Computer Science that ’if we keep certain property
when storing a data, we can reap some kind of benefit from it later on. The ’later
operations’ could be simple data retrieval, search, delete or modify. This is the case
with AVL trees, dictionaries, RB-trees and various property-based data structures.
The algorithm that is at the core of our module relies on a property we keep in our
insertion as we will discuss in Section 4.4. This key property is that we store the
last seen client on top. Our data structure is as follows:
struct xt\_mcast {
struct hlist\_node node;
\_\_be32 ip;
//big-endian ip
unsigned long timeout;
};
struct mtable {
struct hlist\_head members[0];
};
The timeout fields plays that role. The client with largest timeout (the freshest
client we saw) would be the head node of the list. Somehow our ways resemble that
of operation on a stack.
3.5.1. Code flow Chart and Algorithm. The explanation of the algorithm
proceeds in the next Chapter.
packet
enters
Unicast
Multicast
is the
packet a?
IGMP report
lazyadd:
s, c, sexpire;
f or xi :
hash(s), hash(g)
is bucketxi =
∅?
lazyadd:
g, c, gexpire;
No
Drop
Yes
Accept
Yes
∃cs ∈
bucketxs
∃cg ∈
bucketxg
cs = cg ?
Figure 3.3. Code Flow Chart
Input: kip, vip, expire, array
Result: inserts vip \& deletes expired entries at bucket hash[kip]
1 index = hash(kip);
2 entry= null;
3 head= ref = array[index];
4 while ref is not null do
5 time = ref − > timeout;
6 if now ≤ time then
/* The node is expired.
7 entry = ref
8 ref = ref − > next
9 
*/
while ref is not null and ref is not head do
/* hold on to the place-memory!
*/;
/* simply advance it */;
10 temp= ref ;
11 free(ref )
12 ref = temp− > next;
end
13
14
else
ref = ref − > next;
15
16
/* all are ripe to be deleted-expired */;
end
17 end
18 if entry is null then
19
/* if we are not lucky from line 7 */
allocate(memory,entry);
20 else
21 ;
22 end
23 entry− > timout= now + expire
24 entry− > ip = vip;
25 if entry is not head then /* freshest entry must always be head */
26
27
/* should have memory by now */;
hlist\_add\_head(entry,array[index]);
else
/* memory was from head- do nothing */;
28
end
Algorithm 1: the lazyadd algorithm

\chapter{Detailed Design}
\begin{quote}
Make everything as simple as possible, but not simpler.
-Albert Einstein\end{quote}
\section{Design Dilemmas}
We start with something we have implicitly deferred for a while. This is the
case of IGMP leave messages. They have been left out from the code flow chart of
the previous Chapter intentionally. Leave message can be deceptive in the spirit
that we discussed as comeback leave & genuine leave earlier. The design dilemma
posed by that is: do we go to that particular group hash and remove the client
entry from there or simply ignore them?. If we remove the client to just discover
it comes back would be inefficient. IGMP leave messages are, of course, good for
routing scalability. In our case though, they can destroy the cache like behaviour
we worked hard for by removing the client prematurely. Even though we included
that implementation in our module, simplicity and less work bears against it.
Another dilemma is ’how long should a source be valid for?’. That is once it had
unicast interaction with one of our clients. There is no good answer. In the module
we passed that question to the Network Manager to decide (using the userland
Iptables program s/he can set the value on command line). Imagine our IPTV user
again and a hypothetical scenario. S/he switches on the box and during booting
the box contacts the content providers IP address (maybe during that action the
filter notices the contacted IP address). The user is passively watching a movie
from then on and no unicast interaction occurs. It would be undesirable to drop
the packets while the user is watching.
\section{Where is the Module and Why?}
As can be seen on Figure 4.1, our Filter is located right after the switch. This
mainly due to physical reason of getting the whole traffic in one interface. That is,
the switch is where the traffic is multiplexed and demultiplexed and as such it is
the converging point of outgoing traffic. This is called Bridge firewall- Bridgewall-
which is also powerful as it can also do layer 2 filtering (we can, for instance, deny
access to a user by MAC address-IP spoofing impossible- who is contacting bad IP
addresses).
Figure 4.1. Module location
\section{The Memory Management Paradigm}
One performance issue that dogged a previous student design solution was a
periodic ’call back’ function that goes and cleans entries when a timer expires. And
later on, goes to ask to allocate new memory when new entries appear. This can
keep the processor quite busy with loads and stores to and from main memory
while also destroying possible localities maintained by cache1. Memory allocation
is also an expensive process. And when we request for more memory every now
and then, it leads to O(n) linked list complexity (one fragmented memory points to
next as standard linked list structure) to access what we want. This is restatement
of memory ’fragmentation’ problem. There are smart ways of allocating- what the
’slab allocater’ does in the Linux kernel. It is not in the document scope to discuss
further on this matter.We urge the reader for some faith that less call to any of
_malloc family functions is better. If not convinced we refer for discussion memory
fragmentation and all other memory related stuff to [19].
The fundamental basis behind the design the algorithm about to be discussed
is:
• Asking kernel’s memory management system (the slab allocator in Linux)
for a new block of memory right after freeing one is to be avoided. Instead
1\footnote{This can cause the familiar computer ’crunching’ sound- when programs that (ab)use the memory
too much are running.}
we want to overwrite that with a new entry. This is in the same sense of
’one will not break a plate after eating on it and go out to buy new one
when s/he want to eat again’.
On further resolution the above statement, we can ask when is the right time to
free it and how are we going to decide that memory is not needed anymore? As we
will see in a bit the right time is when we go to that bucket location and we free it
if the timer on that node is expired.
The Lazy Cleaner’s Dilemma* (Optional). We introduce a motivation
for our design that we call the lazy cleaner’s dilemma. The cleaner wants to maxi-
mize cleanness and orderliness of stuff s/he uses but minimize cleaning work. The
cleaning precedes the use; the stuff is considered dirty after some time of its use.
We can raise a questions like ’When is it optimal to wash a cup that he used now to
be considered dirty’. It is not unusual to see people have another tea on the same
cup. Examples of non-optimal solutions: periodic cleaning, cleaning stuff s/he is
not likely to use again. I urge you to map this analogies to memory allocation call
optimisation.
\section{The Lazyadd Algorithm (Insert/Delete)}
Let’s begin with the several assumptions
• We have allocated array of n size (take array size of 8 for illustration),
• Initialised (for each bucket-cell) a head pointer to a a linked list node
structure. That is what we call the initialisation step if it referred to
below. The full glory of the data structure can be seen in Appendix
A.2.22.
• Our value IP address vip to be stored and hash key IP address,kip, just
as below:
hash(kip)
kip
vip
One term that we use is dirty node which we define as a node that a timer
associated with it has expired. We can proceed to the steps.
Our algorithm follows roughly these steps to add vip in a bucket located by
hash[kip]:
\subsection{No collision.}
2The way a linked list is implemented (and how one can utilize that) in the Linux kernel is slightly
different from standard Computer Science courses and for us, it is another matter. We refer to the
source code /linux/list.h if one is interested. It safe to assume standard linked list data structure
with next, prev pointers for each structure to understand these algorithms.
(1) Check the bucket pointed by the hash.
kip
hash(kip)
vip
(2) If the head pointer (which is valid during the initialisation step) still points
to NO linked list, create a entry node- head node3.
kip
hash(kip)
NULL
head
vip
(3) This is a case when the head node is still there but it is dirty (i.e the timer
on it has expired).
kip
hash(ip)
head
vip
(4) From one of the two above steps, we have memory. In the first case there
was explicit allocation; as there was no entry before, and in the second
case there was entry but was dirty. In this step, the new entry is added.
kip
hash(ip)
head
vip
Now let’s assume we have had collisions (we have more entries in that bucket),
the list is not dirty up until some middle node- the entries are fresh and have not
expired. We put this general case so the other cases can thought out or inferred.
This leads us to the next section:
\subsection{Collision.}
3\footnote{This will be the “head node”. The head pointer is in the one in the array- small dots. Remember
that the picture covers the head pointers when we show nodes- the nodes are sort of “pinned” on
it}
(1) We have a structure that looks like the following
hash(ip)
kip
head
vip
next
nnex
nnnx
(2) This is a very important step and key part in our algorithm.
(a) Traverse the list until we hit the first expired- we grab its space- and
store the new value.
(b) Delete/clean the list downwards.
kip
hash(ip)
head
vip
next
nnex
nnnx
(3) We restore the property which allows us to do step (2) in two steps
(a) Activate the timer (which no longer makes that node dirty)
kip
hash(ip)
head
vip
next
nnex
(b) This is the last to expire and so make it the head node (the property).
This allows to delete the rest of the list once we find latest expired
node as we did in (2b).
kip
hash(ip)
nnex
vip
head
next
This completes the process of the lazy cleaner insert algorithm.
We have to be careful not to delete the head node at step 2.b if the linked list
structure is circular. This is the case with Linux kernel implementation of the data
structure. Therefore, we need to make sure to hold head node and compare it is
not the node we are at, while dropping4 the expired structures.
\section{Design Tradeoffs}
\subsection{Performance, Memory, Correctness, Design Complexity.} Should
we store the source values that our client is trying to contact? What could it pos-
sibly mean we are not storing those values? Imagine a scenario where a malicious
user’s address maps to an address location which has client (that is due to possible
collision with genuine source). The probability of that happening can be quite sig-
nificant depending the property of the hash function. The module does not check
this but it is slight modification to the data structure including one more ip ad-
dress something I missed. We lose extra storage of memory and one more compare
instruction for the processor in return for correctness.
The performance (which is packet latency at our module) has direct relation
with how many memory cells we allocate to group and client arrays. This is a
simple hashtable tradeoff. We see more of it in Chapter 5.
We can also ask if it is worth the design complexity to, for example, abandon
linked list and use other faster access complexity structures. The answer at this
design stage is no as we have not identified that as bottleneck.
\footnote{As not to cause another ’gravity theory moment’.}


\chapter{Experiments and Results}
\begin{quote}
pigram 95. Don't have good ideas if you aren't willing to be re-
sponsible for them.
-Alan J. Perlis
\end{quote}
\section{Basics Functionality Setup}
\subsection{Tools needed.}
\begin{itemize}
\item Linux based system as outlined in README in Appedix A.1.
\item Iptables Netfilter Modules as outlined in the README in AppendixA.1.
\end{itemize}
The easiest way to check the functionality is to start a streaming video on VLC1
server. This is simply done using GUI of the player. Then, start that session in
VLC client.\footnote{This a very good media player and can be found on www.videolan.org} Again this is done using GUI feature of VLC. One can use the shell
scripts given in Appendix to test the blocking and unblocking functionalities.
\section{Resilence Test (Simulation)}
\subsection{Test platform.}
Operating System: OpenSuSE11.0 with Linux kernel version 2.6.25
Architecture: x86\_64 Pentium D
Memory: 2GB (RAM 2DDR)
\subsection{What we are measuring.} The main performance concern as previ-
ously mentioned is the latency. Therefore, what we need to test is how long will a
packet stay in our module before a decision is made to either drop or accept that
packet. For that, a userland simulator is rewritten which is based on the main
kernel module. There was a time and resource constraint to test the module on
real systems.
For the sake of getting pattern on how increasing number of clients, groups and
sources contacted influence, the numbers are made large. This is in contrast with
real networks where hundreds of clients might only be. The source code is provided
in Appendix B.
As can be infered from the source code, we have made the following assumptions:
\begin{itemize}
\item The space is kept at only 20 buckets (array)- again this is for the sake
getting a pattern not because of lack of memory.
\item Every client is participating in at least one random group.
\item Every client contacts some random external host as its source.
\end{itemize}

\begin{tabular}[|l| l| l| l| l| l| l|l|l|]
Clients& 200 & 500 & 600 & 700 & 800 & 900 & 1800 & 3000 & 4000 & 10,000\\
Time (μs)& 2 & 3 & 4 & 5 & 8 & 9 & 24 & 45 & 60 & 150
\end{tabular}
Table 5.1. Number of clients vs raw processing time

\begin{tabular}[|l| l| l| l| l| l| l|l|l|]
Groups & 10 & 20 & 50 & 100 & 200 & 500 & 1000 & 2000\\
Time (μs)& 5 & 8 & 10 & 10 & 11 & 11 & 10 & 7
\end{tabular}
Table 5.2. Number of groups vs raw processing time
Then what we try to simulate is basically the searching (at peak time when none
of the groups or sources have expired). It would also be interesting to see what the
effect of expiring time has. We may have to defer that to our recommendation. We
have not made any assumption whether the source is malcious or not. We average
the time it takes for a packet to pass through or be dropped over three hundred
source requests (a source delivering a multicast packet) .We run that again on the
terminal 10 times and average the result.
\subsection{Number of clients as variable.} We keep number of groups at 100
and number of sources 10000
25
Clients effect
20
15
10
5
0
500
1,000
1,500
clients
Figure 5.1. Effect of Number of Clients on latency
\subsection{Number of groups as variable.} We keep the number of clients at
1000 and the sources at 10000. The assumption being one client contacts to ten
possible unicast addresses at a time (e.g web browsing different sites.)
\begin{tabular}[|l| l| l| l| l| l| l|l|l|]
Sources &1 & 10 & 20 & 50 & 100 & 200 & 500 & 1000 & 10000\\
Time (μs)& 5 & 6 & 10 & 8 & 11 & 10 & 10 & 7 & 6
\end{tabular}
Table 5.3. Number of groups vs raw processing time
Groups effect
10
8
6
0
500
1,000
groups
1,500
2,000
Figure 5.2. Effect of Number of Groups on latency
\section{Number of sources as variable.} We keep the number of groups 100
and clients at 1000.
sources effect
10
8
6
0
0.2
0.4
0.6
sources
0.8
1
·104
Figure 5.3. Effect of Number of Sources on latency
\section{Results Interpretation}
The results shows the module's dependence on internal clients number- as would
be expected. The more clients we have the more states we have to keep. In other
words, it is our clients that cause the state information creation by initiating those
contacts. If we have more clients, they are likely to contact more groups and group
sources.
More interestingly though, the more groups we have outside, the more likely
our clients will contact different groups. This causes a faster hash access to the
clients. Remember we use the group and source IP addresses as hash keys to store
the clients and then do a cross check. If these values become diverse, which is
what happens when we increase their domain, their hash value is likely to differ
too and thus causing less collisions. This is what is apparent on the the latter
two graphs and on my tests. This aligns perfectly with the design challenge of
the dynamic multicast firewall, which does not know what the source IP or group
IP are before and considers them random. Be warned though, this assumption of
randomness may not be true in practice. And also, the actual network load (not
our management) maybe unbearable when so many sources and groups exist.
Lastly, even if the results put are in the order of μs, we need to bear in mind
multimedia applications (e.g IPTV) that mulitcasting is used for are in the order of
M bps. The result can be that we are not processsing as fast or faster than the bitrate
which is undesired for the application. We would be introducing artifact bandwidth
throttling by our module. For instance, if we are watching real-time video, the
implication is we no longer can watch it properly not because of lack of bandwidth
but the fact that the module has clogged it. Therefore, the machine's raw processing
power and other features are also important in these type of situations and need
considerations.


\chapter{Discussion}
\section{Mitigating Unicast Starvation}
The strongest assumption we made as discussed was that of unicast interaction
of sources and clients. Here are some ways that may be able to help mitigate those
problems:
\begin{itemize}
\item If a user/multicast application is not getting its/his traffic, see what ports,
protocols and unicast control message the application use and add those
to the monitoring list for unicast snooping.
\item It is often helpful to monitor web browsing (but it may increase the state)-
it is again up to the Network Manager to decide by adjusting the memory
required and the expire time of sources.
\item If other means do not work , explicit pings to the unicast sources are
always easy fixes.
\end{itemize}
\section{Another Case Against NATs}
NATs are haunted by controversies but we add one more of our own. Imagine
the following scenario: a source from NATted network is a valid and one of clients
tries to contact it. Due to the nature of the network on the side, our module will
assume all sources from the NATed network are valid. The problem is artefactual
and we do not consider it in our design as based on Engineering principles.
\section{Why Access List Security can be Hard Multicast?}
Multicasting has a seemingly very strange semantics. If two processes open the
same port as a multicast socket and then join two different multicast groups then
traffic for both multicast groups is forwarded to either process. This means that
application will get surprising data that they did not ask for. Applications will
have to filter these out in order to work correctly if multiple applications run on
the same system.
The intent was that multicast address socket behavior would be just like unicast
address behavior, with the exception that multicast addresses are interface-specific.
If we add a new unicast address to the machine, and the binding on the socket
doesn't restrict it, you will receive packets from any of the addresses on the machine.
That's what INADDR\_ANY means. Your multicast application may also receive
unicast traffic on that port too.
Some operating systems supporting Multicast have since changed to only sup-
plying multicast traffic to a socket that was selected through multicast join op-
erations. (application will no longer receive traffic from multicast groups that it
did not subscribe to.) That does not totally avoid the problem except violating
the RFC which is explained in the second paragraph. The problem is anyone on
the network, or within multicast routing domain, may reuse both-multicast group
and port- for a different machine and the application will receive them. The other
machine will also receive the traffic destined for the application.
Why did we go to 'unnecessary' detail above? It should be clear the virtualness
of multicast address makes it easy for anyone to get others people data and fill
garbage on the group address. In security point of view, there is a breach if someone
has access to the data even if they cannot decrypt it.


\chapter{Recommendation, Future Work}
There can be ambitious work that can be done to better this work which include:
\begin{itemize}
\item Testing and using another data structure as collision handlers (static ar-
rays for example). This is in line with claims[14] that linked list spoils
the CPU cache behaviour.
\item To migrate the code to using RCU (Read Copy Update)[15] mechanism
on the data structure instead of traditional locking we used. This is new
elegant paradigm of synchronization of multiprocessors though patented
by IBM.
\end{itemize}
And yes, the Author has done some work on the inclusion of IPv6. Part of the code
can be found in the link given the Appendix . In fact, my wish was I complete the
work in protocol independent way. It is fairly easy to infer what protocol has been
used from either Netfilter or raw packet buffer. I discourage the use of protocol-
dependent 'tricks' unless it is explicitly used for that particular protocol. This is a
simple restatement of the principle of portability


\chapter{Conclusion}
In Chapter 1 we looked at the benefits of multicasting and what kind security
threat the use of it entails in SMEs. The main framework of thinking we emphasized
is ultimately the traffic is coming from source with real IP address. The whole
design challenge in this security issue is how we can do the mapping of the group
virtual address to the real unicast source address. The dynamic stateful multicast
paradigm works on the premise ’if one contacts the group they are also likely to
contact the source they are seeking’. This, we said, is the strongest assumption of
the thesis. But, it is also a realistic assumption.
In Chapter 3 and4, we looked at the algorithms and data structures we modelled
and implemented this idea with. We model it as to provide with a good analytical
view the problems that are we are faced. The way it can be conceived is as a
security glue between SSM and ASM. It tries to feel the crack that can be opened
when one steps to or remains in ASM service model.
We emphasized the design challenges, observations and algorithms than our
particular implementation on the Linux Kernel. They are of greater importance if
anyone wants construct similar endeavour and as life cyle of kernel, as one facet of
technology, changes rapidly.
\appendix
\chapter{Module Code}
The compatibility source code for different versions of Kernel from the link:
http://github.com/bewket/adhoc-multicast/tree/master. It is taken from Jan En-
glehardt’s (http://jengelh.medozas.de/) xtables-addons which seeks to provide that.
\section{README}
\begin{verbatim}
This module provides a filter for incoming multicast traffic by doing the f/g:
1. Snooping Subscibe (or join) messages and storing the client by its group hash.
2. We also snoop mulitcast source address (unicast type). 
This is the part where one sort of NEEDs
to pierce a hole for the multicast (which can be seen in the second last line)
 source-
its unicast address.
We store the client by the multicast source hash.
3. When multicast trafffic come we hash both group and source address and
Checking if the client exists in both hashes.
You need iptables >= 1.4.3 (x_tables module)
and kernel-source >= 2.6.17
Tools: socat-or netcat-, vlc
For testing purpose do the following (need to be root) :
# sh test.sh
Start vlc (www.videolan.org) stream video
on 238.234.234.1
Start extra vlc. Try to open the video on open network and entering the streaming
address. It should NOT open.
You will have to choose UDP in both cases.
Then do the following (after changing the 192.x.x.x to your ip)
echo "hi" | socat - udp-sendto:192.168.1.4:4444
and it should now open
You can download the code by using git. At the terminal type (this
$git clone git://github.com/bewket/adhoc-multicast.git
\end{verbatim}
\section{The filter Module}
\subsection{xt\_MCAST.h.}
\begin{verbatim}
#ifndef _XT_MCAST_TARGET_H
#define _XT_MCAST_TARGET_H
#include <linux/types.h>
struct xt_mcast_target_info {
uint32_t msrc_size;
uint32_t mgrp_size;
};
#endif /* XT_MCAST_TARGET_H */
\end{verbatim}
\subsection{xt\_MCAST.c.}
\begin{verbatim}
#define MSRC_SIZE 20 //src->cli_head_nodes
#define MGRP_SIZE 20 //grp->cli_head_nodes
#define MGRP_EXPIRE
120 //seconds
#include "compat_xtables.h"
#include <linux/skbuff.h>
#include <linux/module.h>
#include <linux/jiffies.h>
#include <net/ipv6.h>
#include <net/ip.h>
#include <linux/ip.h>
#include <linux/igmp.h>
#include <linux/list.h>
#include <linux/jhash.h>
#include <linux/netfilter.h>
#include <linux/netfilter/x_tables.h>
#include "xt_MCAST.h"
//you can compare addresses directly without htonl
//if (ipv6_addr_cmp(&iph->saddr, &new_addr) == 0)
static u_int32_t jhash_rnd;
struct xt_mcast {
struct hlist_node node;
__be32 ip;
unsigned long timeout;
//unsigned long
expire; //this is in jiffies (allowed per group)
};
struct mtable {
struct hlist_head members[0];
};
static struct mtable *msrc_hash;
//client list indexed by hash(src)->cli_ip
static struct mtable *mgrp_hash;
//client list indexed by hash(gid)->cli_ip
static unsigned int source_size = MSRC_SIZE;
static unsigned int group_size = MGRP_SIZE;
static unsigned int source_expire = 2 * 60 * 60;
//2-hours (has to
static DEFINE_MUTEX(mcast_mutex);
static DEFINE_SPINLOCK(mcast_lock);
static struct hlist_node *get_headnode(struct mtable *t, uint32_t h)
{
return t->members[h].first;
}
static unsigned int
init_new_entry(struct xt_mcast *entry, __be32 ip, uint32_t h,
struct mtable *t, uint16_t expire)
{
//we are not going to simply add it.
//go to the list if there is someone who has been expired
//start from the head downwards
//if the head is expired simple update the head with this
//we want to avoid as many memory alloc calls as possible
struct hlist_node *pos, *temp;
struct xt_mcast *temp_entry;
struct hlist_node *ref = get_headnode(t, h);
entry = NULL; //just in case
spin_lock_bh(&mcast_lock);
hlist_for_each(pos, &t->members[h]) {
temp_entry = hlist_entry(pos, struct xt_mcast, node);
if (temp_entry != NULL && time_after(jiffies, temp_entry->timeout)) {
//timeout
entry = temp_entry;
//the first time this we want to delete all but the first expired
//go until you find the first timeout and delete downwards.
pos = pos->next;
//don’t delete pos- just advance it
while (pos != NULL && pos != ref) {
temp_entry = hlist_entry(pos, struct xt_mcast, node);
temp = pos;
pos = NULL; //No pter reshuffle here
pos = temp->next;
kfree(temp_entry);
}
break;
}
}
if (entry == NULL)
entry = kmalloc(sizeof(struct xt_mcast), GFP_ATOMIC);
else
printk("Updating entry\n");
entry->ip = ip;
entry->timeout = jiffies + expire * HZ;
hlist_add_head(&entry->node, &t->members[h]);
spin_unlock_bh(&mcast_lock);
return NF_ACCEPT;
}
static inline uint32_t mcast_hash(__be32 name, unsigned int size)
{
return jhash_1word(name, jhash_rnd) & (size - 1);
}
static struct hlist_node *lookup(__be32 hkey, __be32 ip, struct mtable *t)
{
struct xt_mcast *entry;
struct hlist_node *n = NULL;
if (t == NULL) {
printk("Table is NULL\n");
return NULL;
}
//use of the mutex should be appropriate here-Noisy Kernel unless new!!
spin_lock_bh(&mcast_lock);
hlist_for_each_entry(entry, n, &t->members[hkey], node) {
if (entry->ip == ip) {
spin_unlock_bh(&mcast_lock);
return n;
}
}
spin_unlock_bh(&mcast_lock);
return NULL;
}
static unsigned int igmp_report(__be32 cli, __be32 grp)
{
struct xt_mcast *entry = NULL, *head_entry;
//aka head
uint32_t h = mcast_hash(grp, group_size);
struct hlist_node *head_node = get_headnode(mgrp_hash, h);
struct hlist_node *entry_node;
printk("IGMP REPORT\n");
printk(NIPQUAD_FMT "->" NIPQUAD_FMT "\n", NIPQUAD(grp), NIPQUAD(cli));
if (head_node == NULL)
return init_new_entry(entry, cli, h, mgrp_hash, MGRP_EXPIRE);
//go fish
entry_node = lookup(h, cli, mgrp_hash);
//here there is grp alrdy but is the client in it.
if (entry_node == NULL) {
return init_new_entry(entry, cli, h, mgrp_hash, MGRP_EXPIRE);
}
spin_lock_bh(&mcast_lock);
//good the client is in the grp.
head_entry = hlist_entry(head_node, struct xt_mcast, node);
entry = hlist_entry(entry_node, struct xt_mcast, node);
entry->timeout = jiffies + MGRP_EXPIRE * HZ;
//we need to also update the same entry in the source list
if (entry->ip == head_entry->ip)
goto out;
//it the head. don’t do anything
hlist_del(entry_node); //shuffle and fix
hlist_add_head(&entry->node, &mgrp_hash->members[h]);
//we want this node to be head so we need to del:
//NOTE: when deleting a grp if it the head pointer
out:spin_unlock_bh(&mcast_lock);
return NF_ACCEPT;
}
/* this probably will increase complexity and superflous somehow
* we can just ignore this and use
* our clean first and insert mechanism in fact this
removes
* temporal locality of our entries
*/
static unsigned igmp_leave(__be32 cli, __be32 grp)
{
uint32_t h = mcast_hash(grp, group_size);
struct hlist_node *entry_node;
struct xt_mcast *entry, *head_entry; //aka head
struct hlist_node *head_node = get_headnode(mgrp_hash, h);
/* entry is NULL => "I am leaving this grp"-
*yet there is no such grp or doesn’t have me
* shouldn’t be possible */
if (head_node == NULL)
return NF_ACCEPT;
entry_node = lookup(h, cli, mgrp_hash);
if (entry_node == NULL) //weird
return NF_ACCEPT;
entry = hlist_entry(entry_node, struct xt_mcast, node);
//some clients may just join and leave straight away after becoming last_seen
//Note: Also IGMPv1 don’t issue leave
//The simple solution here is to make the head of the group last_seen.
//which often is true
head_entry = hlist_entry(head_node, struct xt_mcast, node);
entry = hlist_entry(entry_node, struct xt_mcast, node);
spin_lock_bh(&mcast_lock);
if (head_entry->ip == entry->ip && entry_node->next != NULL)
//promote the successor to be head
hlist_add_head(entry_node->next, &mgrp_hash->members[h]);
//else it is the last entry
hlist_del(entry_node);
spin_unlock_bh(&mcast_lock);
kfree(entry);
return NF_ACCEPT;
}
//Incoming traffic MCAST traffic
static unsigned int mcast_handler4(const struct iphdr *iph)
{
__be32 msrc_ip = iph->saddr;
__be32 mgrp_ip = iph->daddr;
uint32_t hs = mcast_hash(msrc_ip, source_size);
//h(msrc_ip)->cli_ip
uint32_t hg = mcast_hash(mgrp_ip, group_size);
//h(mgrp_ip)->cli_ip
struct xt_mcast *msrc2cli, *mgrp2cli;
struct hlist_node *n, *n2;
//check if there exists any client(s) that are mapped to msrc_ip
//while looking hash them and check if they have
spin_lock_bh(&mcast_lock);
hlist_for_each_entry(msrc2cli, n, &msrc_hash->members[hs], node) {
hlist_for_each_entry(mgrp2cli, n2, &mgrp_hash->members[hg], node) {
if (msrc2cli->ip == mgrp2cli->ip) { //success
spin_unlock_bh(&mcast_lock);
return NF_ACCEPT;
//remember we won’t update anything here (
}
}
}
spin_unlock_bh(&mcast_lock);
return NF_DROP;
}
static unsigned
int igmp_handler(const struct iphdr *iph, const struct sk_buff *skb)
{
const struct igmphdr *igmph = igmp_hdr(skb);
const struct igmpv3_grec *rec3;
switch (igmph->type) {
case IGMP_HOST_MEMBERSHIP_REPORT:
case IGMPV2_HOST_MEMBERSHIP_REPORT:
return igmp_report(iph->saddr, iph->daddr);
case IGMPV3_HOST_MEMBERSHIP_REPORT:
//this special V2 equivalent case we can handle per rfc3376
if (ntohs(igmpv3_report_hdr(skb)->ngrec) == 1) {
//ONLY one grp
rec3 = igmpv3_report_hdr(skb)->grec;
if (ntohs(rec3->grec_nsrcs) == 0) //No src list
switch (rec3->grec_type) {
case IGMPV3_CHANGE_TO_EXCLUDE:
return igmp_report(iph->saddr, rec3->grec_mca);
case IGMPV3_CHANGE_TO_INCLUDE:
return igmp_leave(iph->saddr, rec3->grec_mca);
}
}
return NF_ACCEPT;
case IGMP_HOST_LEAVE_MESSAGE:
return igmp_leave(iph->saddr, igmph->group);
//from IGMP_V2 to delete a group (ehnancement on Time interval timer)
default:
return NF_ACCEPT;
}
return NF_ACCEPT;
}
static inline struct mtable *init_table(uint32_t nmembers)
{
uint32_t i;
struct mtable *t;
t = kzalloc(sizeof(*t) + nmembers * sizeof(t->members[0]), GFP_KERNEL);
for (i = 0; i < nmembers; i++)
INIT_HLIST_HEAD(&t->members[i]);
return t;
}
static void delete_list(struct hlist_head *head)
{
struct xt_mcast *entry;
struct hlist_node *pos, *temp;
pos = head->first;
while (pos != NULL) {
entry = hlist_entry(pos, struct xt_mcast, node);
temp = pos;
pos = NULL; //No pter reshuffle here
pos = temp->next;
kfree(entry);
}
}
static inline void delete_table(struct mtable *t, uint32_t nsize)
{
uint32_t i;
if (t == NULL)
return;
spin_lock_bh(&mcast_lock);
for (i = 0; i < nsize; i++)
delete_list(&t->members[i]);
spin_unlock_bh(&mcast_lock);
kfree(t);
}
/**OUTGOING traffic only thus client = ip_hrd(skb)->saddr;
* mcast_src_addr= ip_hdr(skb)->daddr;
*/
static unsigned int unicast_handler(const struct iphdr *iph)
{
__be32 mcli_ip = iph->saddr;
uint32_t hs = mcast_hash(iph->daddr, source_size);
struct hlist_node *mcli_node;
struct xt_mcast *mcli = NULL;
mcli_node = lookup(hs, mcli_ip, msrc_hash);
if (mcli_node == NULL)
return init_new_entry(mcli, mcli_ip, hs, msrc_hash, source_expire);
spin_lock_bh(&mcast_lock);
mcli = hlist_entry(mcli_node, struct xt_mcast, node);
mcli->timeout = jiffies + source_expire * HZ;
spin_unlock_bh(&mcast_lock);
return NF_ACCEPT;
}
static unsigned int mcast_tg4(const struct sk_buff **pskb,
const struct xt_target_param *par)
{
const struct sk_buff *skb = *pskb;
const struct xt_mcast_target_info *info = par->targinfo;
const struct iphdr *iph = ip_hdr(skb);
if (info) {
if (info->msrc_size > MGRP_SIZE)
source_size = info->msrc_size;
if (info->mgrp_size > MGRP_SIZE)
group_size = info->mgrp_size;
}
if (iph == NULL)
return NF_ACCEPT;
if (ipv4_is_multicast(iph->daddr))
//broadcast also mcast so: skb->pkt_type==PACKET_MULTICAST check??
{
if (ip_hdr(skb)->protocol == IPPROTO_UDP)
return mcast_handler4(iph);
//there maybe NF_DROPs
if (ip_hdr(skb)->protocol == IPPROTO_IGMP)
return igmp_handler(iph, skb);
//always NF_ACCEPT just for snooping
} else if (ip_hdr(skb)->protocol == IPPROTO_UDP) {
//the heartbeat protocol can be ICMP, or any other so long as we
// get its ip header to extract msrc.
return unicast_handler(iph);
//this needs to make fast construction of <s,c> pair
}
return NF_ACCEPT;
}
static struct xt_target mcast_tg_reg __read_mostly = {
.name = "MCAST",
.revision = 0,
.family = NFPROTO_IPV4,
.target = mcast_tg4,
.targetsize = XT_ALIGN(sizeof(struct xt_mcast_target_info)),
.me = THIS_MODULE,
};
static int __init mcast_tg_init(void)
{
mgrp_hash = init_table(group_size);
msrc_hash = init_table(source_size);
get_random_bytes(&jhash_rnd, sizeof(jhash_rnd));
return xt_register_target(&mcast_tg_reg);
}
static void __exit mcast_tg_exit(void)
{
delete_table(mgrp_hash, group_size);
delete_table(msrc_hash, source_size);
xt_unregister_target(&mcast_tg_reg);
}
module_init(mcast_tg_init);
module\_exit(mcast\_tg\_exit);
MODULE\_LICENSE("GPL");
MODULE\_ALIAS("ipt\_MCAST");
MODULE\_DESCRIPTION("Xtables: Mcast packet filter");
MODULE\_AUTHOR("Bewketu Tadilo<xxxx@xxxx.edu.au>");
\end{verbatim}
\subsection{libxt\_MCAST.c (for userland interaction to set values).}
\begin{verbatim}
#include <getopt.h>
#include <stdio.h>
#include <string.h>
#include <xtables.h>
#include <linux/netfilter/x_tables.h>
#include "xt_MCAST.h"
static void mcast_tg_help(void)
{
printf("MCAST target options:\n"
" --multicast-source-size value expected number of multicast sources\n"
" --multicast-group-size value expected number of multicast groups \n");
}
enum {
MCAST_OPT_MSRC_SIZE,
MCAST_OPT_MGRP_SIZE,
};
static const struct option mcast_opts[] = {
{"multicast-source-size", 1, NULL, MCAST_OPT_MSRC_SIZE},
{"multicast-group-size", 1, NULL, MCAST_OPT_MGRP_SIZE},
{.name = NULL},
};
static int mcast_tg_parse(int c, char **argv, int invert,
unsigned int *flags, const void *entry,
struct xt_entry_target **target)
{
struct xt_mcast_target_info *info = (void *) (*target)->data;
unsigned int num;
switch (c) {
case MCAST_OPT_MSRC_SIZE:
if (*flags & (1 << c)) {
xtables_error(PARAMETER_PROBLEM,
"MCAST: can’t specify ‘--multicast-source-size’ twice");
}
if (!xtables_strtoui(optarg, NULL, &num, 1, UINT32_MAX)) {
xtables_error(PARAMETER_PROBLEM,
"Unable to parse ‘%s’ in "
"‘--mulitcast-souce-size’", optarg);
}
info->msrc_size = num;
*flags |= 1 << c;
break;
case MCAST_OPT_MGRP_SIZE:
if (*flags & (1 << c))
xtables_error(PARAMETER_PROBLEM,
"MCAST: can only specify ‘--multicast-group-size’ once");
if (!xtables_strtoui(optarg, NULL, &num, 1, UINT32_MAX)) {
xtables_error(PARAMETER_PROBLEM,
"Unable to parse ‘%s’ in "
"‘--mulitcast-souce-size’", optarg);
}
info->mgrp_size = num;
*flags |= 1 << c;
break;
default:
return 0;
}
return 1;
}
static void mcast_tg_check(unsigned int flags)
{
}
static struct xtables_target mcast_tg_reg = {
.version = XTABLES_VERSION,
.name = "MCAST",
.revision = 0,
.family = PF_INET,
.size = XT_ALIGN(sizeof(struct xt_mcast_target_info)),
.userspacesize = XT_ALIGN(sizeof(struct xt_mcast_target_info)),
.help = mcast_tg_help,
.parse = mcast_tg_parse,
.final_check = mcast_tg_check,
};
static __attribute__ ((constructor))
void mcast_tg_ldr(void)
{
xtables_register_target(&mcast_tg_reg);
}
\end{verbatim}
\subsection{libxt\_MCAST.c (for userland interaction to set values).}
\subsection{Scripts.}
\subsubsection{Makefile.}
\begin{verbatim}
top_srcdir := ..
srcdir := .
prefix := /usr/
exec_prefix := ${prefix}
libdir := /usr/lib
libexecdir := ${exec_prefix}/libexec
xtlibdir := ${libexecdir}/xtables
xtables_CFLAGS :=
CC := gcc
CCLD := ${CC}
LCFLAGS
:= -g -O2
LDFLAGS :=
regular_CFLAGS := -D_LARGEFILE_SOURCE=1 -D_LARGE_FILES -D_FILE_OFFSET_BITS=64
-D_REENTRANT -Wall -Waggregate-return -Wmissing-declarations
-Wmissing-prototypes -Wredundant-decls -Wshadow -Wstrict-prototypes
-Winline -pipe -DXTABLES_LIBDIR=\"${xtlibdir}\"
kinclude_CFLAGS :=
AM_CFLAGS
-I /lib/modules/2.6.25.20-0.1-default/build/include
:= ${regular_CFLAGS} -I${top_srcdir}/include ${xtables_CFLAGS}
${kinclude_CFLAGS}
AM_DEPFLAGS
= -Wp,-MMD,$(@D)/.$(@F).d,-MT,$@
obj-m += compat_xtables.o xt_MCAST.o
all: lib
make -C /lib/modules/‘uname -r‘/build M=‘pwd‘
clean:
make -C /lib/modules/‘uname -r‘/build M=‘pwd‘ clean
rm -rf libxt_MCAST.so
install:
cp *.so
$(xtlibdir)
cp *.ko /lib/modules/‘uname -r‘/extra
lib:
libxt_MCAST.so
lib%.so: lib%.oo
${CCLD} ${AM_LDFLAGS} -shared ${LDFLAGS} -o $@ $<;
lib%.oo: ${srcdir}/lib%.c
${CC} ${AM_DEPFLAGS} ${AM_CFLAGS} -D_INIT=lib$*_init -DPIC -fPIC
${LCFLAGS} -o $@ -c $<;
\end{verbatim}
\subsubsection{iptables rule (test.sh).}
\begin{verbatim}
#!/bin/sh
IPT=/usr/sbin/iptables
modprobe x_tables
$IPT -F
rmmod
xt_MCAST
rmmod compat_xtables.ko
insmod compat_xtables.ko
insmod xt_MCAST.ko
$IPT -I OUTPUT -p udp --dport 4444 -j MCAST
#snoop ONLY igmp we don’t care about outgoing Multicast traffic
$IPT -I OUTPUT -p igmp -j MCAST
$IPT -I INPUT -d 224.0.0.0/4 ! -p igmp -j MCAST
#apply policy on incoming mcast traffic which is IGMP (it could be proxied)
\end{verbatim} 
\chapter{Userland Simulator}
To get the following work one needs /linux/list.h which can be dowloaded
(by easy google search) or in /usr/src/linux/include/linux/list.h if one has source
installed.
• Remove all the prefetch statements
• Remove all the kernel-specific constants
• put containerof implementation in it as follows and can include it in user
(as opposed to kenel) application by #include :
#define container_of(ptr, type, member) ({ \
const typeof( ((type *)0)->member ) *__mptr = (ptr);
\
(type *)( (char *)__mptr - offsetof(type,member) );})
#define offsetof(TYPE, MEMBER) ((size_t) &((TYPE *)0)->MEMBER)
\section{Header File}
\begin{verbatim}
#ifndef _MTEST_H
#define _MTEST_H
#include <stdint.h>
struct xt_mcast {
struct hlist_node node;
__be32 ip;
unsigned long timeout;
//unsigned long
expire; //this is in jiffies (allowed per group)
};
struct mtable {
struct hlist_head members[0];
};
#define time_after(a,b) \
((long)(b) - (long)(a) < 0)
#define time_before(a,b) time_after(b,a)
static unsigned int
init_new_entry(struct xt_mcast *entry, __be32 ip, uint32_t h,
struct mtable *t, uint16_t expire);
static inline uint32_t mcast_hash(__be32 name, unsigned int size);
static struct hlist_node *lookup(__be32 hkey, __be32 ip, struct mtable *t);
static struct hlist_node *get_headnode(struct mtable *t, uint32_t h);
static unsigned int igmp_report(__be32 cli, __be32 grp);
static unsigned int mcast_search(__be32 msrs_ip, __be32 mgrp_ip);
static unsigned igmp_leave(__be32 cli, __be32 grp);
static inline struct mtable *init_table(uint32_t nmembers);
static void delete_list(struct hlist_head *head);
static inline void delete_table(struct mtable *t, uint32_t nsize);
static unsigned int unicast_handler(__be32, __be32);
#endif /*_MTEST_H*/
\end{verbatim}
\section{Main.c}
\begin{verbatim}
#include <stdio.h>
#include <stdbool.h>
#include <time.h>
#include <linux/types.h>
#include <stdlib.h>
#include <netinet/ip.h>
#include <netinet/in.h>
#include <netinet/ip6.h>
#include <getopt.h>
#include "list.h"
#include "mtest.h"
#include <linux/netfilter.h>
#include "jhash.h"
#include <sys/time.h>
#define MSRC_SIZE 20 //src->cli_head_nodes
#define MGRP_SIZE 20 //grp->cli_head_nodes
#define MGRP_EXPIRE
120
//2-minutes about the same timetime router will issue Group Query.
#define MCAST_GRP_ID
10
#define ARRAR_SIZE(array) (sizeof(array)/sizeof(array[0]))
static void print_hash_list(struct hlist_head head);
void print_usage(FILE * stream, char **argv);
static inline void simulate(int, char **);
static struct mtable *msrc_hash;
//client list indexed by hash(src)->cli_ip
static struct mtable *mgrp_hash;
//client list indexed by hash(gid)->cli_ip
static unsigned int source_size = MSRC_SIZE;
static unsigned int group_size = MGRP_SIZE;
static unsigned int source_expire = 2 * 60 * 60;
//2-hours
static unsigned int *rand_src, *rand_grp;
static unsigned int *rand_cli;
static struct timeval now;
static u_int32_t jhash_rnd;
// all insert functions are f(keyip,cli)
format
int main(int argc, char **argv)
{
int i;
if (argc < 2)
print_usage(stderr, argv);
jhash_rnd = time(NULL);
mgrp_hash = init_table(group_size);
msrc_hash = init_table(source_size);
simulate(argc, argv);
// for(i = 0; i < 10; i++)
//
igmp_report(i,MCAST_GRP_ID);
//print_hash_list(mgrp_hash->members[mcast_hash(MCAST_GRP_ID,group_size)]);
//unicast_handler(10,12);
//mcast_search(MCAST_GRP_ID,9);
delete_table(mgrp_hash, group_size);
delete_table(msrc_hash, source_size);
return 0;
}
static struct option long_opt[] = {
{"ngrp", 1, 0, ’g’},
{"ncli", 1, 0, ’c’},
{"nsrc", 1, 0, ’s’},
{"percentag", 1, 0, ’r’},
{NULL, 0, NULL, 0}
};
static char *short_opt = "g:c:s:r:";
static inline void simulate(int argc, char **argv)
{
//15-groups
//15-clients
int i;
srand(time(NULL));
int c, sizec, sizes, sizeg, per;
long mtime, seconds, useconds;
struct timeval start;
struct timeval end;
long average = 0;
while (1) {
c = getopt_long(argc, argv, short_opt, long_opt, NULL);
if (c == -1)
break;
switch (c) {
case ’c’:{
sizec = atoi(optarg);
rand_cli = malloc(sizec * sizeof(int *));
for (i = 0; i < sizec; i++)
rand_cli[i] = rand();
}
break;
case ’s’:{
sizes = atoi(optarg);
rand_src = malloc(sizes * sizeof(int *));
for (i = 0; i < sizes; i++)
rand_src[i] = rand();
}
break;
case ’g’:{
sizeg = atoi(optarg);
rand_grp = malloc(sizes * sizeof(int *));
for (i = 0; i < sizes; i++)
rand_grp[i] = rand();
}
break;
case ’r’:
per = atoi(optarg);
break;
default:
print_usage(stderr, argv);
}
}
for (i = 0; i < sizec; i++) {
igmp_report(rand_cli[i], rand_grp[rand() % sizeg]);
unicast_handler(rand_cli[i], rand_src[rand() % sizes]);
}
//again here now we simulate a random source contacts our network
//here we measure how long does it take
//to block it to let it go with a relation to the our memory
//do 300 random searches printing the raw time it took.
for (i = 0; i < 300; i++) {
gettimeofday(&start, NULL);
mcast_search(rand_src[rand() % sizes], rand_grp[rand() % sizeg]);
gettimeofday(&end, NULL);
//seconds = end.tv_sec - start.tv_sec;
useconds = end.tv_usec - start.tv_usec;
average += useconds;
//
mtime=
(seconds *1000 + useconds/1000.0)+ 0.5;
//printf("Elapsed: %ld us\n",useconds);
}
printf("Averag elapsed %ld useconds\n", average / 300);
}
static void print_hash_list(struct hlist_head head)
// this is array of hashtables
{
struct xt_mcast *entry;
struct hlist_node *pos;
hlist_for_each_entry(entry, pos, &head, node)
printf("ip: %u\n", entry->ip);
}
static unsigned int
init_new_entry(struct xt_mcast *entry, __be32 ip, uint32_t h,
struct mtable *t, uint16_t expire)
{
struct hlist_node *pos, *temp;
struct xt_mcast *temp_entry;
struct xt_mcast *place;
struct hlist_node *ref = get_headnode(t, h);
hlist_for_each_entry_safe(place, pos, temp, &t->members[h], node) {
gettimeofday(&now, NULL);
if (place != NULL && time_after(now.tv_sec, place->timeout)) {
entry = place;
pos = pos->next;
//don’t delete pos- just advance it
while (pos != NULL && pos != ref) {
temp_entry = hlist_entry(pos, struct xt_mcast, node);
free(temp_entry);
temp = pos;
pos = NULL; //No pter reshuffle here
pos = temp->next;
}
break;
}
}
if (entry == NULL)
entry = (struct xt_mcast *) malloc(sizeof(struct xt_mcast));
entry->ip = ip;
gettimeofday(&now, NULL);
entry->timeout = now.tv_sec + expire;
hlist_add_head(&entry->node, &t->members[h]);
return NF_ACCEPT;
}
static inline uint32_t mcast_hash(__be32 name, unsigned int size)
{
return jhash_1word(name, jhash_rnd) & (size - 1);
}
static struct hlist_node *lookup(__be32 hkey, __be32 ip, struct mtable *t)
{
struct xt_mcast *entry;
struct hlist_node *n;
struct hlist_node *temp = n;
//use of the mutex should be appropriate here-Noisy Kernel unless new!!
hlist_for_each_entry_safe(entry, n, temp, &t->members[hkey], node)
if (ip == entry->ip)
return n;
return NULL;
}
static struct hlist_node *get_headnode(struct mtable *t, uint32_t h)
{
return t->members[h].first;
}
static unsigned int igmp_report(__be32 cli, __be32 grp)
{
struct xt_mcast *entry = NULL, *head_entry;
//aka head
uint32_t h = mcast_hash(grp, group_size);
struct hlist_node *head_node = get_headnode(mgrp_hash, h);
struct hlist_node *entry_node;
if (head_node == NULL)
return init_new_entry(entry, cli, h, mgrp_hash, MGRP_EXPIRE);
//go fish
entry_node = lookup(h, cli, mgrp_hash);
//here there is grp alrdy but is the client in it.
if (entry_node == NULL) {
return init_new_entry(entry, cli, h, mgrp_hash, MGRP_EXPIRE);
}
//good the client is in the grp.
head_entry = hlist_entry(head_node, struct xt_mcast, node);
entry = hlist_entry(entry_node, struct xt_mcast, node);
gettimeofday(&now, NULL);
entry->timeout = now.tv_sec + MGRP_EXPIRE;
//we need to also update the same entry in the source list
if (entry->ip == head_entry->ip)
return NF_ACCEPT;
//it the head. don’t do anything
hlist_del(entry_node); //shuffle and fix
hlist_add_head(&entry->node, &mgrp_hash->members[h]);
//we want this node to be head so we need to del:
//NOTE: when deleting a grp if it the head pointer
printf("IGMP REPORT END\n");
return NF_ACCEPT;
}
static unsigned int mcast_search(__be32 msrc_ip, __be32 mgrp_ip)
{
uint32_t hs = mcast_hash(msrc_ip, source_size);
//h(msrc_ip)->cli_ip
uint32_t hg = mcast_hash(mgrp_ip, group_size);
//h(mgrp_ip)->cli_ip
struct xt_mcast *msrc2cli, *mgrp2cli;
struct hlist_node *n, *n2;
//check if there exists any client(s) that are mapped to msrc_ip
//while looking hash them and check if they have
hlist_for_each_entry(msrc2cli, n, &msrc_hash->members[hs], node) {
hlist_for_each_entry(mgrp2cli, n2, &mgrp_hash->members[hg], node) {
if (msrc2cli->ip == mgrp2cli->ip) { //success
return NF_ACCEPT;
//remember we won’t update anything here (
}
}
}
return NF_DROP;
}
static unsigned igmp_leave(__be32 cli, __be32 grp)
{
uint32_t h = mcast_hash(grp, group_size);
struct hlist_node *entry_node;
struct xt_mcast *entry, *head_entry; //aka head
struct hlist_node *head_node = get_headnode(mgrp_hash, h);
/* entry is NULL => "I am leaving this grp"-
yet there is no such grp or doesn’t have me
* shouldn’t be possible */
if (head_node == NULL)
return NF_ACCEPT;
entry_node = lookup(h, cli, mgrp_hash);
if (entry_node == NULL) //weird
return NF_ACCEPT;
entry = hlist_entry(entry_node, struct xt_mcast, node);
head_entry = hlist_entry(head_node, struct xt_mcast, node);
entry = hlist_entry(entry_node, struct xt_mcast, node);
if (head_entry->ip == entry->ip && entry_node->next != NULL)
//promote the successor to be head
hlist_add_head(entry_node->next, &mgrp_hash->members[h]);
//else it is the last entry
hlist_del(entry_node);
free(entry);
return NF_ACCEPT;
}
static inline struct mtable *init_table(uint32_t nmembers)
{
uint32_t i;
struct mtable *t;
t = malloc(sizeof(*t) + nmembers * sizeof(t->members[0]));
for (i = 0; i < nmembers; i++)
INIT_HLIST_HEAD(&t->members[i]);
return t;
}
static void delete_list(struct hlist_head *head)
{
struct xt_mcast *entry;
struct hlist_node *pos, *temp;
pos = head->first;
while (pos != NULL) {
entry = hlist_entry(pos, struct xt_mcast, node);
temp = pos;
pos = NULL; //No pter reshuffle here
pos = temp->next;
free(entry);
}
}
static inline void delete_table(struct mtable *t, uint32_t nsize)
{
uint32_t i;
if (t == NULL)
return;
for (i = 0; i < nsize; i++)
delete_list(&t->members[i]);
free(t);
}
static unsigned int unicast_handler(__be32 mcli_ip, __be32 msrc_ip)
{
uint32_t hs = mcast_hash(msrc_ip, source_size);
struct hlist_node *mcli_node;
struct xt_mcast *mcli = NULL;
mcli_node = lookup(hs, mcli_ip, msrc_hash);
if (mcli_node == NULL)
return init_new_entry(mcli, mcli_ip, hs, msrc_hash, source_expire);
mcli = hlist_entry(mcli_node, struct xt_mcast, node);
gettimeofday(&now, NULL);
mcli->timeout = now.tv_sec + source_expire;
return NF_ACCEPT;
}
void print_usage(FILE * stream, char **argv)
{
fprintf(stream, "\n" "Usage: %s -c <ncli> " "-s <nsrc> \n", argv[0]);
fprintf(stream,
"\n"
"Opitions:\n"
" -c, --ncli number of clients\n"
" -s, --nsrc expected number of sources\n"
" -r, --percentage multicast connection percentage (from the above)\n"
"\n");
exit(0);
}
\end{verbatim}
\backmatter
\newcommand{\refname}{References}
\renewcommand{\bibname}{References}
%\newcommand{\myreftoc}{\verbatim{\contentsline{section}{\numberline{\secnumdepth}section}}
%\addtocontents{toc}{\reftoc {\thepage}{\refname}}
\refstepcounter{chapter}
\addcontentsline{toc}{chapter}{\bibname}
%\addtocounter{\secnumdepth}{1}
%\chapter
\begin{thebibliography}{10}
\bibitem{Deering} http://rfc.org/dee
\end{thebibliography}


\end{document}
